\documentclass[twoside]{article}
\usepackage[preprint]{twocolumnconf}

% Personalized ------------------------------------------
\usepackage[round]{natbib}
\renewcommand{\bibname}{References}
\renewcommand{\bibsection}{\subsubsection*{\bibname}}
\usepackage{amsmath,amssymb,amsthm,mathtools,bm}
\usepackage{xspace}
\usepackage{booktabs}
\usepackage{graphicx}
\usepackage{pgf}
\providecommand{\mathdefault}[1]{#1}
\providecommand{\mathregular}[1]{#1}
\usepackage{enumitem}
\usepackage{hyperref}
\usepackage{xcolor}
\usepackage{float}

\usepackage{algorithm}
\usepackage{algorithmic}
\renewcommand{\algorithmicrequire}{\textbf{Input:}}
\renewcommand{\algorithmicensure}{\textbf{Output:}}

% Section headings are written in all caps to match AISTATS convention.
\newtheorem{theorem}{Theorem}
\newtheorem{proposition}[theorem]{Proposition}
\newtheorem{lemma}[theorem]{Lemma}
\newtheorem{corollary}[theorem]{Corollary}
\theoremstyle{definition}
\newtheorem{definition}[theorem]{Definition}
\newtheorem{assumption}[theorem]{Assumption}
\theoremstyle{remark}
\newtheorem{remark}[theorem]{Remark}

\DeclareMathOperator{\KL}{KL}
\DeclareMathOperator{\vol}{vol}
\DeclareMathOperator{\tr}{tr}
\DeclareMathOperator{\Scal}{Scal}
\DeclareMathOperator{\diam}{diam}
\newcommand{\R}{\mathbb{R}}
\newcommand{\N}{\mathbb{N}}
\newcommand{\E}{\mathbb{E}}
\newcommand{\Prob}{\mathbb{P}}
\newcommand{\M}{\mathcal{M}}
\newcommand{\dd}{\,\mathrm{d}}
\newcommand{\what}{\widehat}
\newcommand{\ess}{\mathrm{ESS}}
\newcommand{\gess}{\mathrm{gESS}}
\newcommand{\ip}[2]{\left\langle #1,#2\right\rangle}
\newcommand{\calH}{\mathcal H}
\newcommand{\mmd}{\mathrm{MMD}}
\newcommand{\op}{O_p}

% Personalized ------------------------------------------



\begin{document}

% If your paper is accepted and the title of your paper is very long,
% the style will print as headings an error message. Use the following
% command to supply a shorter title of your paper so that it can be
% used as headings.
%
\runningtitle{ }

% If your paper is accepted and the number of authors is large, the
% style will print as headings an error message. Use the following
% command to supply a shorter version of the author names so that
% they can be used as headings (for example, use only the surnames)
%
\runningauthor{ }

\twocolumn[

\aistatstitle{Heat-Kernel Entropy Profiles and Geometric Effective Sample Size for Weighted Measures on Manifolds}

\aistatsauthor{ Kisung You \And Boram Cho}
\aistatsaddress{ Baruch College \And Yale University} ]
%\aistatsaddress{ Baruch College \And The Graduate Center, City University of New York } ]

\begin{abstract}
Weighted empirical measures on compact manifolds appear in importance sampling, particle approximations, posterior summaries, quadrature, and representation learning. Ordinary effective sample size and related weight summaries ignore the geometry of the support. We introduce heat-kernel entropy profiles to measure nonuniformity after intrinsic diffusion at a range of scales. For order-two R\'enyi entropy, pairwise heat-kernel overlaps give an exact profile and a geometric effective sample size. This effective sample size discounts nearby or duplicate particles. It approaches ordinary effective sample size as overlaps between distinct particles vanish. On compact boundaryless manifolds, we establish profile monotonicity, gESS scale limits, deterministic-weight consistency, and a bounded-ratio result for self-normalized importance sampling. On spheres, the unlogged profile decomposes into spherical-harmonic energies. The first terms are squared mean-resultant and traceless-second-moment energies, which give vMF- and Bingham-type scalar summaries. Experiments identify antipodal, girdle, multimodal, and duplicate-particle structures that weight-only and first-moment summaries miss.
\end{abstract}

%----%----%----%----%----%----%----%----%----%----%----%----%----%
\section{INTRODUCTION}
Let $\M$ be a compact Riemannian manifold and suppose a probability
measure of interest is represented by weighted atoms,
\begin{equation}
  \widehat P_w=\sum_{i=1}^n w_i\delta_{x_i},
  \quad x_i\in\M,\; w_i\ge0,\; \sum_iw_i=1.
  \label{eq:weighted_measure}
\end{equation}
Such measures arise in quadrature, importance sampling, particle methods,
posterior summaries, ensembles, and normalized representations. Our goal is to summarize the diffuseness of \eqref{eq:weighted_measure} using both its weights and support geometry.

Many standard diagnostics are label-based. The
ordinary effective sample size (ESS), used in importance sampling and
sequential Monte Carlo, monitors weight degeneracy
\citep{kong_1994_SequentialImputationsBayesian, liu_1998_SequentialMonteCarlo, delmoral_2006_SequentialMonteCarlo}. It is defined as
\begin{equation*}
  \ess(w)=\frac{1}{\sum_iw_i^2}.  
\end{equation*}
This quantity is also the inverse-Simpson, or order-two Hill, effective number of the
weights \citep{hill_1973_DiversityEvennessUnifying}. It diagnoses label-level
weight concentration but is not a geometric property of the represented
measure. Splitting an atom into identical labels should not increase its
diffuseness. Likewise, equally weighted particles should contribute differently when their locations are close rather than far apart.

Geometric uncertainty also depends on scale. A weighted measure may resemble
many particles at fine resolution, a few clusters at medium resolution, and
a nearly uniform measure after sufficient smoothing. On spheres, common first-moment summaries include mean resultant length and fitted von Mises-Fisher (vMF) concentration.
These summaries can miss antipodal, girdle-like, and symmetric multimodal structure
\citep{fisher_1953_DispersionSphere,mardia_2000_DirectionalStatistics,
banerjee_2005_ClusteringUnitHypersphere}. Our examples focus on these cases, but the method applies to general manifolds.


We apply intrinsic heat flow to the weighted atoms and measure nonuniformity over diffusion time. For order-two R\'enyi entropy, a pairwise heat-kernel identity avoids manifold integration. Normalizing these overlaps gives a geometric effective sample size (gESS) for each scale. The unlogged profile is also a heat-kernel maximum mean discrepancy (MMD) from uniformity. We interpret the resulting quantities as effective occupied volume and effective particle number. Small- and large-time results justify these interpretations.

The theory covers compact,
connected, boundaryless manifolds with normalized volume. Exact
implementation requires heat-kernel evaluations or controlled
approximations. Our experiments use $S^2$, where the heat kernel is explicit.
The vMF and Bingham interpretations are specific to spheres and do not enter the general definition.

\paragraph{Contributions.}
The paper makes five contributions. First, we define heat-kernel R\'enyi-2 profiles for weighted empirical measures on compact manifolds. Pairwise heat-kernel overlaps give an exact computation. Second, gESS discounts nearby and duplicate atoms, while approaching ordinary ESS when distinct-atom overlaps vanish. Third, small-time analysis describes local heat-kernel volume, curvature, duplicate coalescence, and two-atom merging. Fourth, a spectral expansion yields a spherical specialization. Its first two nonconstant energies are squared mean-resultant and traceless-second-moment energies. Finally, we prove stability and weighted-consistency results, including a bounded-ratio case for self-normalized importance sampling. The appendix contains proofs and further experiments. Replication code is available at \url{https://github.com/kisungyou/HeatKernelEntropyProfiles}.

%----%----%----%----%----%----%----%----%----%----%----%----%----%
\section{RELATED WORK}\label{sec:related-work}

\paragraph{Weight-only and similarity-aware effective numbers.}
Classical ESS diagnostics measure weight degeneracy in importance sampling and sequential Monte Carlo. Discrepancy-based alternatives address the same problem \citep{martino_2017_EffectiveSampleSize}. Hill numbers and inverse-Simpson diversity have an order-two effective-number interpretation \citep{hill_1973_DiversityEvennessUnifying}. Leinster-Cobbold diversity adds a user-specified similarity matrix \citep{leinster_2012_MeasuringDiversityImportance}. Our gESS is an order-two similarity diversity based on normalized heat-kernel overlap. This choice gives merge invariance, scale limits, and a heat asymptotic interpretation.

\paragraph{Heat smoothing, entropy, and kernel discrepancies.}
Kernel density estimation on spheres and compact manifolds has a substantial literature \citep{hall_1987_KernelDensityEstimation, pelletier_2005_KernelDensityEstimation, lebrigant_2019_ApproximationDensitiesRiemannian}. Nearest-neighbor entropy estimators are also available for hyperspherical data \citep{li_2011_KNearestNeighborBased}. These methods usually choose one smoothing scale to estimate a density or entropy. Here, diffusion scale determines the resolution. For \(A_P(t)=\iint k_{2t}(x,y)\dd P(x)\dd P(y)\), the unlogged profile satisfies \(A_P(t)-1=\mmd_{k_{2t}}^2(P,\nu)\). This identity links the profile to kernel mean embeddings \citep{sriperumbudur_2010_HilbertSpaceEmbeddings,gretton_2012_KernelTwoSampleTest, muandet_2017_KernelMeanEmbedding}. It also links the profile to Sobolev tests of uniformity on compact manifolds \citep{gine_1975_InvariantTestsUniformity, garcia-portugues_2018_OverviewUniformityTests}. Our transformations give the MMD curve two scale-dependent interpretations. They are an effective occupied volume fraction and an effective number of distinguishable atoms.

\paragraph{Spherical directional summaries.}
On spheres, vMF and Bingham/Fisher-Bingham models provide
classical first- and second-order summaries of directional structure \citep{kent_1982_FisherBinghamDistributionSphere,mardia_2000_DirectionalStatistics}. The proposed profile does not replace these parametric models in their unimodal or axial settings. Their moment energies appear as the first two nonconstant spectral components. Higher harmonic degrees distinguish structures that first- and second-order summaries cannot separate.


%----%----%----%----%----%----%----%----%----%----%----%----%----%
\section{HEAT ENTROPY PROFILES}

Throughout, $\M$ is a compact, connected, $m$-dimensional Riemannian manifold
without boundary. Let $\nu=\vol/\vol(\M)$ be normalized Riemannian volume, so
$\nu(\M)=1$ and the uniform density with respect to $\nu$ is identically one.
For vectors, $\|\cdot\|$ is the Euclidean norm. We use $\|\cdot\|_F$ for matrices and $\|\cdot\|_q$ for the $L^q(\nu)$ norm.
Let $\Delta$ be the Laplacian, with $\partial_t u=\Delta u$ and nonnegative eigenvalues of $-\Delta$. The heat kernel with respect to $\nu$ is $k_t(x,y)$
\citep{rosenberg_1997_LaplacianRiemannianManifold, grigoryan_2009_HeatKernelAnalysis}:
\begin{align*}
  \int_\M k_t(x,y)\dd\nu(y)&=1,\\
  k_{t+s}(x,y)&=\int_\M k_t(x,z)k_s(z,y)\dd\nu(z).
\end{align*}

For a probability measure $P$ on $\M$, write $H_tP$ for the heat-smoothed law.
For every $t>0$, $H_tP$ has a smooth density $p_t$ with respect to $\nu$.
Define
\begin{align*}
  D_{2,P}(t)&=\log\int_\M p_t(x)^2\dd\nu(x),\\
  U_{2,P}(t)&=\exp\{-D_{2,P}(t)\}.
\end{align*}
The quantity $D_{2,P}(t)$ is the order-two R\'enyi divergence of $H_tP$ from $\nu$. Its inverse exponential, $U_{2,P}(t)$, is the effective occupied volume fraction.

For the weighted empirical measure \eqref{eq:weighted_measure}, the
heat-smoothed density is
\begin{equation}
  \widehat p_t(y)=(H_t\widehat P_w)(y)
  =\sum_{i=1}^n w_i k_t(y,x_i).
  \label{eq:pt_def}
\end{equation}
Evaluating the same functional at $\widehat P_w$ gives the empirical heat entropy profile:
\[
  \widehat D_2(t)=\log\int_\M \widehat p_t(y)^2\dd\nu(y),
  \quad
  \widehat U_2(t)=\exp\{-\widehat D_2(t)\}.
\]
Write
\[
  \widehat A_w(t)=\int_\M \widehat p_t(y)^2\dd\nu(y),
\]
so that $\widehat D_2(t)=\log \widehat A_w(t)$ and
$\widehat U_2(t)=\widehat A_w(t)^{-1}$.

The quantity $U_{2,P}(t)$ measures effective occupied volume. For atomic input, $U_{2,\widehat P_w}(t)\to0$ as $t\downarrow0$ because the smoothed density becomes singular. The gESS defined below removes the self-overlap of each heat kernel. It therefore measures distinguishable atoms rather than volume.

\begin{theorem}[Pairwise identity and monotonicity]
\label{thm:pairwise_monotone}
For every $t>0$,
\begin{equation}
  \widehat A_w(t)
  =
  \int_\M \widehat p_t(y)^2\dd\nu(y)
  =
  \sum_{i,j=1}^n w_iw_jk_{2t}(x_i,x_j).
  \label{eq:pairwise_identity}
\end{equation}
Consequently,
\[
  \widehat U_2(t)
  =
  \left[\sum_{i,j}w_iw_jk_{2t}(x_i,x_j)\right]^{-1}.
\]
Moreover, $t\mapsto \widehat D_2(t)$ is nonincreasing and
$t\mapsto \widehat U_2(t)$ is nondecreasing. Finally,
\[
  \lim_{t\to\infty}\widehat D_2(t)=0,
  \qquad
  \lim_{t\to\infty}\widehat U_2(t)=1.
\]
\end{theorem}

Thus, intrinsic diffusion monotonically removes nonuniformity. The rate at which $\widehat U_2(t)$ approaches one defines the uncertainty profile.


%----%----%----%----%----%----%----%----%----%----%----%----%----%
\section{GEOMETRIC EFFECTIVE SAMPLE SIZE}\label{sec:gess}

The entropy profile measures effective volume. We obtain an effective number of atoms by normalizing pairwise heat overlaps by each kernel's self-overlap.

\begin{definition}[Geometric ESS]
\label{def:gess}
For $t>0$, define
\[
  s_t(x,y)=\frac{k_{2t}(x,y)}{\sqrt{k_{2t}(x,x)k_{2t}(y,y)}}
\]
and
\[
  \gess_w(t)=\left[\sum_{i,j}w_iw_js_t(x_i,x_j)\right]^{-1}.
\]
\end{definition}

On homogeneous manifolds, $k_{2t}(x,x)$ is constant. Examples include spheres and compact Lie groups with invariant metrics. Then
\[
  \gess_w(t)=k_{2t}(x,x)\widehat U_2(t).
\]
Thus, gESS measures effective volume in units of one heat kernel's self-volume. It is the order-two Hill/Leinster-Cobbold diversity for
$S_t=(s_t(x_i,x_j))_{ij}$ \citep{hill_1973_DiversityEvennessUnifying, leinster_2012_MeasuringDiversityImportance}.

\begin{theorem}[Properties and limits of gESS]
\label{thm:gess}
For every $t>0$,
\[
  0<s_t(x,y)\le1,
  \qquad
  s_t(x,x)=1.
\]
Consequently,
\[
  1\le \gess_w(t)\le \frac{1}{\sum_iw_i^2}=\ess(w).
\]
If exact duplicate atoms are merged by summing their weights, $\gess_w(t)$ is
unchanged. More generally, if the distinct support locations are
$z_1,\ldots,z_r$ with total masses
\[
  \alpha_a=\sum_{i:x_i=z_a}w_i,
\]
then
\[
  \lim_{t\downarrow0}\gess_w(t)=\frac{1}{\sum_{a=1}^r\alpha_a^2}.
\]
In particular, if all $x_i$ are distinct, the limit equals $\ess(w)$. Finally,
\[
  \lim_{t\to\infty}\gess_w(t)=1.
\]
\end{theorem}

Theorem~\ref{thm:gess} adds support geometry to ordinary ESS. Separated particles count individually at fine scales and become indistinguishable at coarse scales. Computation requires pairwise heat-kernel Gram matrices on a fixed grid of times. The cost is comparable to a kernel discrepancy or MMD diagnostic. Unlike one discrepancy value, gESS gives an effective-number curve. Algorithm~\ref{alg:heat-gess} combines Theorem~\ref{thm:pairwise_monotone} with the normalization in Definition~\ref{def:gess}.

\begin{algorithm}[t]
\caption{Heat entropy profile and geometric effective sample size}
\label{alg:heat-gess}
\small
\begin{algorithmic}[1]
\REQUIRE Weighted atoms $(x_i,w_i)_{i=1}^n$ with $\sum_iw_i=1$; heat scales
$0<t_1<\cdots<t_K$; heat kernel $k_t$ on $\M$.
\ENSURE Profiles $\{(t_k,\widehat U_2(t_k),\gess_w(t_k))\}_{k=1}^K$.
\FOR{$k=1,\ldots,K$}
  \STATE Form $G^{(k)}\in\mathbb R^{n\times n}$ with
  $G^{(k)}_{ij}\gets k_{2t_k}(x_i,x_j)$.
  \STATE Set $\widehat U_2(t_k)\gets \{w^\top G^{(k)}w\}^{-1}$.
  \IF{$\M$ is homogeneous}
    \STATE Set $S^{(k)}\gets G^{(k)}/k_{2t_k}(x_0,x_0)$ for any $x_0\in\M$.
  \ELSE
    \STATE Set $S^{(k)}_{ij}\gets
    G^{(k)}_{ij}/\sqrt{G^{(k)}_{ii}G^{(k)}_{jj}}$ for all $i,j$.
  \ENDIF
  \STATE Set $\gess_w(t_k)\gets\{w^\top S^{(k)}w\}^{-1}$.
\ENDFOR
\end{algorithmic}
\end{algorithm}

The two outputs have different interpretations. The $\widehat U_2$ profile is a normalized volume and measures effective occupancy after diffusion. The $\gess_w$ profile counts the heat-kernel footprints needed at a given scale. The same overlaps therefore give a density summary and a particle degeneracy diagnostic. Figure~\ref{fig:profiles} shows both quantities for the spherical examples.

\begin{figure}[ht]
\centering
%\includegraphics[width=0.98\linewidth]{hkep_fig_profiles_column.pdf}
\providecommand{\mathdefault}[1]{#1}
\providecommand{\mathregular}[1]{#1}
%% Creator: Matplotlib, PGF backend
%%
%% To include the figure in your LaTeX document, write
%%   \input{<filename>.pgf}
%%
%% Make sure the required packages are loaded in your preamble
%%   \usepackage{pgf}
%%
%% Also ensure that all the required font packages are loaded; for instance,
%% the lmodern package is sometimes necessary when using math font.
%%   \usepackage{lmodern}
%%
%% Figures using additional raster images can only be included by \input if
%% they are in the same directory as the main LaTeX file. For loading figures
%% from other directories you can use the `import` package
%%   \usepackage{import}
%%
%% and then include the figures with
%%   \import{<path to file>}{<filename>.pgf}
%%
%% Matplotlib used the following preamble
%%   \def\mathdefault#1{#1}
%%   \everymath=\expandafter{\the\everymath\displaystyle}
%%   \IfFileExists{scrextend.sty}{
%%     \usepackage[fontsize=8.000000pt]{scrextend}
%%   }{
%%     \renewcommand{\normalsize}{\fontsize{8.000000}{9.600000}\selectfont}
%%     \normalsize
%%   }
%%   \usepackage{amsmath,amssymb}
%%   \makeatletter\@ifpackageloaded{underscore}{}{\usepackage[strings]{underscore}}\makeatother
%%
\begingroup%
\makeatletter%
\begin{pgfpicture}%
\pgfpathrectangle{\pgfpointorigin}{\pgfqpoint{3.484853in}{2.568639in}}%
\pgfusepath{use as bounding box, clip}%
\begin{pgfscope}%
\pgfsetbuttcap%
\pgfsetmiterjoin%
\definecolor{currentfill}{rgb}{1.000000,1.000000,1.000000}%
\pgfsetfillcolor{currentfill}%
\pgfsetlinewidth{0.000000pt}%
\definecolor{currentstroke}{rgb}{1.000000,1.000000,1.000000}%
\pgfsetstrokecolor{currentstroke}%
\pgfsetdash{}{0pt}%
\pgfpathmoveto{\pgfqpoint{0.000000in}{0.000000in}}%
\pgfpathlineto{\pgfqpoint{3.484853in}{0.000000in}}%
\pgfpathlineto{\pgfqpoint{3.484853in}{2.568639in}}%
\pgfpathlineto{\pgfqpoint{0.000000in}{2.568639in}}%
\pgfpathlineto{\pgfqpoint{0.000000in}{0.000000in}}%
\pgfpathclose%
\pgfusepath{fill}%
\end{pgfscope}%
\begin{pgfscope}%
\pgfsetbuttcap%
\pgfsetmiterjoin%
\definecolor{currentfill}{rgb}{1.000000,1.000000,1.000000}%
\pgfsetfillcolor{currentfill}%
\pgfsetlinewidth{0.000000pt}%
\definecolor{currentstroke}{rgb}{0.000000,0.000000,0.000000}%
\pgfsetstrokecolor{currentstroke}%
\pgfsetstrokeopacity{0.000000}%
\pgfsetdash{}{0pt}%
\pgfpathmoveto{\pgfqpoint{0.420602in}{0.992139in}}%
\pgfpathlineto{\pgfqpoint{1.596892in}{0.992139in}}%
\pgfpathlineto{\pgfqpoint{1.596892in}{2.343639in}}%
\pgfpathlineto{\pgfqpoint{0.420602in}{2.343639in}}%
\pgfpathlineto{\pgfqpoint{0.420602in}{0.992139in}}%
\pgfpathclose%
\pgfusepath{fill}%
\end{pgfscope}%
\begin{pgfscope}%
\pgfpathrectangle{\pgfqpoint{0.420602in}{0.992139in}}{\pgfqpoint{1.176290in}{1.351500in}}%
\pgfusepath{clip}%
\pgfsetrectcap%
\pgfsetroundjoin%
\pgfsetlinewidth{0.451687pt}%
\definecolor{currentstroke}{rgb}{0.690196,0.690196,0.690196}%
\pgfsetstrokecolor{currentstroke}%
\pgfsetstrokeopacity{0.280000}%
\pgfsetdash{}{0pt}%
\pgfpathmoveto{\pgfqpoint{0.696512in}{0.992139in}}%
\pgfpathlineto{\pgfqpoint{0.696512in}{2.343639in}}%
\pgfusepath{stroke}%
\end{pgfscope}%
\begin{pgfscope}%
\pgfsetbuttcap%
\pgfsetroundjoin%
\definecolor{currentfill}{rgb}{0.000000,0.000000,0.000000}%
\pgfsetfillcolor{currentfill}%
\pgfsetlinewidth{0.803000pt}%
\definecolor{currentstroke}{rgb}{0.000000,0.000000,0.000000}%
\pgfsetstrokecolor{currentstroke}%
\pgfsetdash{}{0pt}%
\pgfsys@defobject{currentmarker}{\pgfqpoint{0.000000in}{-0.048611in}}{\pgfqpoint{0.000000in}{0.000000in}}{%
\pgfpathmoveto{\pgfqpoint{0.000000in}{0.000000in}}%
\pgfpathlineto{\pgfqpoint{0.000000in}{-0.048611in}}%
\pgfusepath{stroke,fill}%
}%
\begin{pgfscope}%
\pgfsys@transformshift{0.696512in}{0.992139in}%
\pgfsys@useobject{currentmarker}{}%
\end{pgfscope}%
\end{pgfscope}%
\begin{pgfscope}%
\definecolor{textcolor}{rgb}{0.000000,0.000000,0.000000}%
\pgfsetstrokecolor{textcolor}%
\pgfsetfillcolor{textcolor}%
\pgftext[x=0.696512in,y=0.894917in,,top]{\color{textcolor}{\rmfamily\fontsize{7.000000}{8.400000}\selectfont\catcode`\^=\active\def^{\ifmmode\sp\else\^{}\fi}\catcode`\%=\active\def%{\%}$\mathdefault{10^{1}}$}}%
\end{pgfscope}%
\begin{pgfscope}%
\pgfpathrectangle{\pgfqpoint{0.420602in}{0.992139in}}{\pgfqpoint{1.176290in}{1.351500in}}%
\pgfusepath{clip}%
\pgfsetrectcap%
\pgfsetroundjoin%
\pgfsetlinewidth{0.451687pt}%
\definecolor{currentstroke}{rgb}{0.690196,0.690196,0.690196}%
\pgfsetstrokecolor{currentstroke}%
\pgfsetstrokeopacity{0.280000}%
\pgfsetdash{}{0pt}%
\pgfpathmoveto{\pgfqpoint{1.435445in}{0.992139in}}%
\pgfpathlineto{\pgfqpoint{1.435445in}{2.343639in}}%
\pgfusepath{stroke}%
\end{pgfscope}%
\begin{pgfscope}%
\pgfsetbuttcap%
\pgfsetroundjoin%
\definecolor{currentfill}{rgb}{0.000000,0.000000,0.000000}%
\pgfsetfillcolor{currentfill}%
\pgfsetlinewidth{0.803000pt}%
\definecolor{currentstroke}{rgb}{0.000000,0.000000,0.000000}%
\pgfsetstrokecolor{currentstroke}%
\pgfsetdash{}{0pt}%
\pgfsys@defobject{currentmarker}{\pgfqpoint{0.000000in}{-0.048611in}}{\pgfqpoint{0.000000in}{0.000000in}}{%
\pgfpathmoveto{\pgfqpoint{0.000000in}{0.000000in}}%
\pgfpathlineto{\pgfqpoint{0.000000in}{-0.048611in}}%
\pgfusepath{stroke,fill}%
}%
\begin{pgfscope}%
\pgfsys@transformshift{1.435445in}{0.992139in}%
\pgfsys@useobject{currentmarker}{}%
\end{pgfscope}%
\end{pgfscope}%
\begin{pgfscope}%
\definecolor{textcolor}{rgb}{0.000000,0.000000,0.000000}%
\pgfsetstrokecolor{textcolor}%
\pgfsetfillcolor{textcolor}%
\pgftext[x=1.435445in,y=0.894917in,,top]{\color{textcolor}{\rmfamily\fontsize{7.000000}{8.400000}\selectfont\catcode`\^=\active\def^{\ifmmode\sp\else\^{}\fi}\catcode`\%=\active\def%{\%}$\mathdefault{10^{2}}$}}%
\end{pgfscope}%
\begin{pgfscope}%
\pgfsetbuttcap%
\pgfsetroundjoin%
\definecolor{currentfill}{rgb}{0.000000,0.000000,0.000000}%
\pgfsetfillcolor{currentfill}%
\pgfsetlinewidth{0.602250pt}%
\definecolor{currentstroke}{rgb}{0.000000,0.000000,0.000000}%
\pgfsetstrokecolor{currentstroke}%
\pgfsetdash{}{0pt}%
\pgfsys@defobject{currentmarker}{\pgfqpoint{0.000000in}{-0.027778in}}{\pgfqpoint{0.000000in}{0.000000in}}{%
\pgfpathmoveto{\pgfqpoint{0.000000in}{0.000000in}}%
\pgfpathlineto{\pgfqpoint{0.000000in}{-0.027778in}}%
\pgfusepath{stroke,fill}%
}%
\begin{pgfscope}%
\pgfsys@transformshift{0.474070in}{0.992139in}%
\pgfsys@useobject{currentmarker}{}%
\end{pgfscope}%
\end{pgfscope}%
\begin{pgfscope}%
\pgfsetbuttcap%
\pgfsetroundjoin%
\definecolor{currentfill}{rgb}{0.000000,0.000000,0.000000}%
\pgfsetfillcolor{currentfill}%
\pgfsetlinewidth{0.602250pt}%
\definecolor{currentstroke}{rgb}{0.000000,0.000000,0.000000}%
\pgfsetstrokecolor{currentstroke}%
\pgfsetdash{}{0pt}%
\pgfsys@defobject{currentmarker}{\pgfqpoint{0.000000in}{-0.027778in}}{\pgfqpoint{0.000000in}{0.000000in}}{%
\pgfpathmoveto{\pgfqpoint{0.000000in}{0.000000in}}%
\pgfpathlineto{\pgfqpoint{0.000000in}{-0.027778in}}%
\pgfusepath{stroke,fill}%
}%
\begin{pgfscope}%
\pgfsys@transformshift{0.532580in}{0.992139in}%
\pgfsys@useobject{currentmarker}{}%
\end{pgfscope}%
\end{pgfscope}%
\begin{pgfscope}%
\pgfsetbuttcap%
\pgfsetroundjoin%
\definecolor{currentfill}{rgb}{0.000000,0.000000,0.000000}%
\pgfsetfillcolor{currentfill}%
\pgfsetlinewidth{0.602250pt}%
\definecolor{currentstroke}{rgb}{0.000000,0.000000,0.000000}%
\pgfsetstrokecolor{currentstroke}%
\pgfsetdash{}{0pt}%
\pgfsys@defobject{currentmarker}{\pgfqpoint{0.000000in}{-0.027778in}}{\pgfqpoint{0.000000in}{0.000000in}}{%
\pgfpathmoveto{\pgfqpoint{0.000000in}{0.000000in}}%
\pgfpathlineto{\pgfqpoint{0.000000in}{-0.027778in}}%
\pgfusepath{stroke,fill}%
}%
\begin{pgfscope}%
\pgfsys@transformshift{0.582049in}{0.992139in}%
\pgfsys@useobject{currentmarker}{}%
\end{pgfscope}%
\end{pgfscope}%
\begin{pgfscope}%
\pgfsetbuttcap%
\pgfsetroundjoin%
\definecolor{currentfill}{rgb}{0.000000,0.000000,0.000000}%
\pgfsetfillcolor{currentfill}%
\pgfsetlinewidth{0.602250pt}%
\definecolor{currentstroke}{rgb}{0.000000,0.000000,0.000000}%
\pgfsetstrokecolor{currentstroke}%
\pgfsetdash{}{0pt}%
\pgfsys@defobject{currentmarker}{\pgfqpoint{0.000000in}{-0.027778in}}{\pgfqpoint{0.000000in}{0.000000in}}{%
\pgfpathmoveto{\pgfqpoint{0.000000in}{0.000000in}}%
\pgfpathlineto{\pgfqpoint{0.000000in}{-0.027778in}}%
\pgfusepath{stroke,fill}%
}%
\begin{pgfscope}%
\pgfsys@transformshift{0.624901in}{0.992139in}%
\pgfsys@useobject{currentmarker}{}%
\end{pgfscope}%
\end{pgfscope}%
\begin{pgfscope}%
\pgfsetbuttcap%
\pgfsetroundjoin%
\definecolor{currentfill}{rgb}{0.000000,0.000000,0.000000}%
\pgfsetfillcolor{currentfill}%
\pgfsetlinewidth{0.602250pt}%
\definecolor{currentstroke}{rgb}{0.000000,0.000000,0.000000}%
\pgfsetstrokecolor{currentstroke}%
\pgfsetdash{}{0pt}%
\pgfsys@defobject{currentmarker}{\pgfqpoint{0.000000in}{-0.027778in}}{\pgfqpoint{0.000000in}{0.000000in}}{%
\pgfpathmoveto{\pgfqpoint{0.000000in}{0.000000in}}%
\pgfpathlineto{\pgfqpoint{0.000000in}{-0.027778in}}%
\pgfusepath{stroke,fill}%
}%
\begin{pgfscope}%
\pgfsys@transformshift{0.662700in}{0.992139in}%
\pgfsys@useobject{currentmarker}{}%
\end{pgfscope}%
\end{pgfscope}%
\begin{pgfscope}%
\pgfsetbuttcap%
\pgfsetroundjoin%
\definecolor{currentfill}{rgb}{0.000000,0.000000,0.000000}%
\pgfsetfillcolor{currentfill}%
\pgfsetlinewidth{0.602250pt}%
\definecolor{currentstroke}{rgb}{0.000000,0.000000,0.000000}%
\pgfsetstrokecolor{currentstroke}%
\pgfsetdash{}{0pt}%
\pgfsys@defobject{currentmarker}{\pgfqpoint{0.000000in}{-0.027778in}}{\pgfqpoint{0.000000in}{0.000000in}}{%
\pgfpathmoveto{\pgfqpoint{0.000000in}{0.000000in}}%
\pgfpathlineto{\pgfqpoint{0.000000in}{-0.027778in}}%
\pgfusepath{stroke,fill}%
}%
\begin{pgfscope}%
\pgfsys@transformshift{0.918953in}{0.992139in}%
\pgfsys@useobject{currentmarker}{}%
\end{pgfscope}%
\end{pgfscope}%
\begin{pgfscope}%
\pgfsetbuttcap%
\pgfsetroundjoin%
\definecolor{currentfill}{rgb}{0.000000,0.000000,0.000000}%
\pgfsetfillcolor{currentfill}%
\pgfsetlinewidth{0.602250pt}%
\definecolor{currentstroke}{rgb}{0.000000,0.000000,0.000000}%
\pgfsetstrokecolor{currentstroke}%
\pgfsetdash{}{0pt}%
\pgfsys@defobject{currentmarker}{\pgfqpoint{0.000000in}{-0.027778in}}{\pgfqpoint{0.000000in}{0.000000in}}{%
\pgfpathmoveto{\pgfqpoint{0.000000in}{0.000000in}}%
\pgfpathlineto{\pgfqpoint{0.000000in}{-0.027778in}}%
\pgfusepath{stroke,fill}%
}%
\begin{pgfscope}%
\pgfsys@transformshift{1.049073in}{0.992139in}%
\pgfsys@useobject{currentmarker}{}%
\end{pgfscope}%
\end{pgfscope}%
\begin{pgfscope}%
\pgfsetbuttcap%
\pgfsetroundjoin%
\definecolor{currentfill}{rgb}{0.000000,0.000000,0.000000}%
\pgfsetfillcolor{currentfill}%
\pgfsetlinewidth{0.602250pt}%
\definecolor{currentstroke}{rgb}{0.000000,0.000000,0.000000}%
\pgfsetstrokecolor{currentstroke}%
\pgfsetdash{}{0pt}%
\pgfsys@defobject{currentmarker}{\pgfqpoint{0.000000in}{-0.027778in}}{\pgfqpoint{0.000000in}{0.000000in}}{%
\pgfpathmoveto{\pgfqpoint{0.000000in}{0.000000in}}%
\pgfpathlineto{\pgfqpoint{0.000000in}{-0.027778in}}%
\pgfusepath{stroke,fill}%
}%
\begin{pgfscope}%
\pgfsys@transformshift{1.141394in}{0.992139in}%
\pgfsys@useobject{currentmarker}{}%
\end{pgfscope}%
\end{pgfscope}%
\begin{pgfscope}%
\pgfsetbuttcap%
\pgfsetroundjoin%
\definecolor{currentfill}{rgb}{0.000000,0.000000,0.000000}%
\pgfsetfillcolor{currentfill}%
\pgfsetlinewidth{0.602250pt}%
\definecolor{currentstroke}{rgb}{0.000000,0.000000,0.000000}%
\pgfsetstrokecolor{currentstroke}%
\pgfsetdash{}{0pt}%
\pgfsys@defobject{currentmarker}{\pgfqpoint{0.000000in}{-0.027778in}}{\pgfqpoint{0.000000in}{0.000000in}}{%
\pgfpathmoveto{\pgfqpoint{0.000000in}{0.000000in}}%
\pgfpathlineto{\pgfqpoint{0.000000in}{-0.027778in}}%
\pgfusepath{stroke,fill}%
}%
\begin{pgfscope}%
\pgfsys@transformshift{1.213004in}{0.992139in}%
\pgfsys@useobject{currentmarker}{}%
\end{pgfscope}%
\end{pgfscope}%
\begin{pgfscope}%
\pgfsetbuttcap%
\pgfsetroundjoin%
\definecolor{currentfill}{rgb}{0.000000,0.000000,0.000000}%
\pgfsetfillcolor{currentfill}%
\pgfsetlinewidth{0.602250pt}%
\definecolor{currentstroke}{rgb}{0.000000,0.000000,0.000000}%
\pgfsetstrokecolor{currentstroke}%
\pgfsetdash{}{0pt}%
\pgfsys@defobject{currentmarker}{\pgfqpoint{0.000000in}{-0.027778in}}{\pgfqpoint{0.000000in}{0.000000in}}{%
\pgfpathmoveto{\pgfqpoint{0.000000in}{0.000000in}}%
\pgfpathlineto{\pgfqpoint{0.000000in}{-0.027778in}}%
\pgfusepath{stroke,fill}%
}%
\begin{pgfscope}%
\pgfsys@transformshift{1.271514in}{0.992139in}%
\pgfsys@useobject{currentmarker}{}%
\end{pgfscope}%
\end{pgfscope}%
\begin{pgfscope}%
\pgfsetbuttcap%
\pgfsetroundjoin%
\definecolor{currentfill}{rgb}{0.000000,0.000000,0.000000}%
\pgfsetfillcolor{currentfill}%
\pgfsetlinewidth{0.602250pt}%
\definecolor{currentstroke}{rgb}{0.000000,0.000000,0.000000}%
\pgfsetstrokecolor{currentstroke}%
\pgfsetdash{}{0pt}%
\pgfsys@defobject{currentmarker}{\pgfqpoint{0.000000in}{-0.027778in}}{\pgfqpoint{0.000000in}{0.000000in}}{%
\pgfpathmoveto{\pgfqpoint{0.000000in}{0.000000in}}%
\pgfpathlineto{\pgfqpoint{0.000000in}{-0.027778in}}%
\pgfusepath{stroke,fill}%
}%
\begin{pgfscope}%
\pgfsys@transformshift{1.320983in}{0.992139in}%
\pgfsys@useobject{currentmarker}{}%
\end{pgfscope}%
\end{pgfscope}%
\begin{pgfscope}%
\pgfsetbuttcap%
\pgfsetroundjoin%
\definecolor{currentfill}{rgb}{0.000000,0.000000,0.000000}%
\pgfsetfillcolor{currentfill}%
\pgfsetlinewidth{0.602250pt}%
\definecolor{currentstroke}{rgb}{0.000000,0.000000,0.000000}%
\pgfsetstrokecolor{currentstroke}%
\pgfsetdash{}{0pt}%
\pgfsys@defobject{currentmarker}{\pgfqpoint{0.000000in}{-0.027778in}}{\pgfqpoint{0.000000in}{0.000000in}}{%
\pgfpathmoveto{\pgfqpoint{0.000000in}{0.000000in}}%
\pgfpathlineto{\pgfqpoint{0.000000in}{-0.027778in}}%
\pgfusepath{stroke,fill}%
}%
\begin{pgfscope}%
\pgfsys@transformshift{1.363835in}{0.992139in}%
\pgfsys@useobject{currentmarker}{}%
\end{pgfscope}%
\end{pgfscope}%
\begin{pgfscope}%
\pgfsetbuttcap%
\pgfsetroundjoin%
\definecolor{currentfill}{rgb}{0.000000,0.000000,0.000000}%
\pgfsetfillcolor{currentfill}%
\pgfsetlinewidth{0.602250pt}%
\definecolor{currentstroke}{rgb}{0.000000,0.000000,0.000000}%
\pgfsetstrokecolor{currentstroke}%
\pgfsetdash{}{0pt}%
\pgfsys@defobject{currentmarker}{\pgfqpoint{0.000000in}{-0.027778in}}{\pgfqpoint{0.000000in}{0.000000in}}{%
\pgfpathmoveto{\pgfqpoint{0.000000in}{0.000000in}}%
\pgfpathlineto{\pgfqpoint{0.000000in}{-0.027778in}}%
\pgfusepath{stroke,fill}%
}%
\begin{pgfscope}%
\pgfsys@transformshift{1.401634in}{0.992139in}%
\pgfsys@useobject{currentmarker}{}%
\end{pgfscope}%
\end{pgfscope}%
\begin{pgfscope}%
\definecolor{textcolor}{rgb}{0.000000,0.000000,0.000000}%
\pgfsetstrokecolor{textcolor}%
\pgfsetfillcolor{textcolor}%
\pgftext[x=1.008747in,y=0.787417in,,top]{\color{textcolor}{\rmfamily\fontsize{8.000000}{9.600000}\selectfont\catcode`\^=\active\def^{\ifmmode\sp\else\^{}\fi}\catcode`\%=\active\def%{\%}Scale $\alpha$ (Degrees)}}%
\end{pgfscope}%
\begin{pgfscope}%
\pgfpathrectangle{\pgfqpoint{0.420602in}{0.992139in}}{\pgfqpoint{1.176290in}{1.351500in}}%
\pgfusepath{clip}%
\pgfsetrectcap%
\pgfsetroundjoin%
\pgfsetlinewidth{0.451687pt}%
\definecolor{currentstroke}{rgb}{0.690196,0.690196,0.690196}%
\pgfsetstrokecolor{currentstroke}%
\pgfsetstrokeopacity{0.280000}%
\pgfsetdash{}{0pt}%
\pgfpathmoveto{\pgfqpoint{0.420602in}{1.025928in}}%
\pgfpathlineto{\pgfqpoint{1.596892in}{1.025928in}}%
\pgfusepath{stroke}%
\end{pgfscope}%
\begin{pgfscope}%
\pgfsetbuttcap%
\pgfsetroundjoin%
\definecolor{currentfill}{rgb}{0.000000,0.000000,0.000000}%
\pgfsetfillcolor{currentfill}%
\pgfsetlinewidth{0.803000pt}%
\definecolor{currentstroke}{rgb}{0.000000,0.000000,0.000000}%
\pgfsetstrokecolor{currentstroke}%
\pgfsetdash{}{0pt}%
\pgfsys@defobject{currentmarker}{\pgfqpoint{-0.048611in}{0.000000in}}{\pgfqpoint{-0.000000in}{0.000000in}}{%
\pgfpathmoveto{\pgfqpoint{-0.000000in}{0.000000in}}%
\pgfpathlineto{\pgfqpoint{-0.048611in}{0.000000in}}%
\pgfusepath{stroke,fill}%
}%
\begin{pgfscope}%
\pgfsys@transformshift{0.420602in}{1.025928in}%
\pgfsys@useobject{currentmarker}{}%
\end{pgfscope}%
\end{pgfscope}%
\begin{pgfscope}%
\definecolor{textcolor}{rgb}{0.000000,0.000000,0.000000}%
\pgfsetstrokecolor{textcolor}%
\pgfsetfillcolor{textcolor}%
\pgftext[x=0.268017in, y=0.992171in, left, base]{\color{textcolor}{\rmfamily\fontsize{7.000000}{8.400000}\selectfont\catcode`\^=\active\def^{\ifmmode\sp\else\^{}\fi}\catcode`\%=\active\def%{\%}0}}%
\end{pgfscope}%
\begin{pgfscope}%
\pgfpathrectangle{\pgfqpoint{0.420602in}{0.992139in}}{\pgfqpoint{1.176290in}{1.351500in}}%
\pgfusepath{clip}%
\pgfsetrectcap%
\pgfsetroundjoin%
\pgfsetlinewidth{0.451687pt}%
\definecolor{currentstroke}{rgb}{0.690196,0.690196,0.690196}%
\pgfsetstrokecolor{currentstroke}%
\pgfsetstrokeopacity{0.280000}%
\pgfsetdash{}{0pt}%
\pgfpathmoveto{\pgfqpoint{0.420602in}{1.301833in}}%
\pgfpathlineto{\pgfqpoint{1.596892in}{1.301833in}}%
\pgfusepath{stroke}%
\end{pgfscope}%
\begin{pgfscope}%
\pgfsetbuttcap%
\pgfsetroundjoin%
\definecolor{currentfill}{rgb}{0.000000,0.000000,0.000000}%
\pgfsetfillcolor{currentfill}%
\pgfsetlinewidth{0.803000pt}%
\definecolor{currentstroke}{rgb}{0.000000,0.000000,0.000000}%
\pgfsetstrokecolor{currentstroke}%
\pgfsetdash{}{0pt}%
\pgfsys@defobject{currentmarker}{\pgfqpoint{-0.048611in}{0.000000in}}{\pgfqpoint{-0.000000in}{0.000000in}}{%
\pgfpathmoveto{\pgfqpoint{-0.000000in}{0.000000in}}%
\pgfpathlineto{\pgfqpoint{-0.048611in}{0.000000in}}%
\pgfusepath{stroke,fill}%
}%
\begin{pgfscope}%
\pgfsys@transformshift{0.420602in}{1.301833in}%
\pgfsys@useobject{currentmarker}{}%
\end{pgfscope}%
\end{pgfscope}%
\begin{pgfscope}%
\definecolor{textcolor}{rgb}{0.000000,0.000000,0.000000}%
\pgfsetstrokecolor{textcolor}%
\pgfsetfillcolor{textcolor}%
\pgftext[x=0.212654in, y=1.268075in, left, base]{\color{textcolor}{\rmfamily\fontsize{7.000000}{8.400000}\selectfont\catcode`\^=\active\def^{\ifmmode\sp\else\^{}\fi}\catcode`\%=\active\def%{\%}10}}%
\end{pgfscope}%
\begin{pgfscope}%
\pgfpathrectangle{\pgfqpoint{0.420602in}{0.992139in}}{\pgfqpoint{1.176290in}{1.351500in}}%
\pgfusepath{clip}%
\pgfsetrectcap%
\pgfsetroundjoin%
\pgfsetlinewidth{0.451687pt}%
\definecolor{currentstroke}{rgb}{0.690196,0.690196,0.690196}%
\pgfsetstrokecolor{currentstroke}%
\pgfsetstrokeopacity{0.280000}%
\pgfsetdash{}{0pt}%
\pgfpathmoveto{\pgfqpoint{0.420602in}{1.577738in}}%
\pgfpathlineto{\pgfqpoint{1.596892in}{1.577738in}}%
\pgfusepath{stroke}%
\end{pgfscope}%
\begin{pgfscope}%
\pgfsetbuttcap%
\pgfsetroundjoin%
\definecolor{currentfill}{rgb}{0.000000,0.000000,0.000000}%
\pgfsetfillcolor{currentfill}%
\pgfsetlinewidth{0.803000pt}%
\definecolor{currentstroke}{rgb}{0.000000,0.000000,0.000000}%
\pgfsetstrokecolor{currentstroke}%
\pgfsetdash{}{0pt}%
\pgfsys@defobject{currentmarker}{\pgfqpoint{-0.048611in}{0.000000in}}{\pgfqpoint{-0.000000in}{0.000000in}}{%
\pgfpathmoveto{\pgfqpoint{-0.000000in}{0.000000in}}%
\pgfpathlineto{\pgfqpoint{-0.048611in}{0.000000in}}%
\pgfusepath{stroke,fill}%
}%
\begin{pgfscope}%
\pgfsys@transformshift{0.420602in}{1.577738in}%
\pgfsys@useobject{currentmarker}{}%
\end{pgfscope}%
\end{pgfscope}%
\begin{pgfscope}%
\definecolor{textcolor}{rgb}{0.000000,0.000000,0.000000}%
\pgfsetstrokecolor{textcolor}%
\pgfsetfillcolor{textcolor}%
\pgftext[x=0.212654in, y=1.543980in, left, base]{\color{textcolor}{\rmfamily\fontsize{7.000000}{8.400000}\selectfont\catcode`\^=\active\def^{\ifmmode\sp\else\^{}\fi}\catcode`\%=\active\def%{\%}20}}%
\end{pgfscope}%
\begin{pgfscope}%
\pgfpathrectangle{\pgfqpoint{0.420602in}{0.992139in}}{\pgfqpoint{1.176290in}{1.351500in}}%
\pgfusepath{clip}%
\pgfsetrectcap%
\pgfsetroundjoin%
\pgfsetlinewidth{0.451687pt}%
\definecolor{currentstroke}{rgb}{0.690196,0.690196,0.690196}%
\pgfsetstrokecolor{currentstroke}%
\pgfsetstrokeopacity{0.280000}%
\pgfsetdash{}{0pt}%
\pgfpathmoveto{\pgfqpoint{0.420602in}{1.853642in}}%
\pgfpathlineto{\pgfqpoint{1.596892in}{1.853642in}}%
\pgfusepath{stroke}%
\end{pgfscope}%
\begin{pgfscope}%
\pgfsetbuttcap%
\pgfsetroundjoin%
\definecolor{currentfill}{rgb}{0.000000,0.000000,0.000000}%
\pgfsetfillcolor{currentfill}%
\pgfsetlinewidth{0.803000pt}%
\definecolor{currentstroke}{rgb}{0.000000,0.000000,0.000000}%
\pgfsetstrokecolor{currentstroke}%
\pgfsetdash{}{0pt}%
\pgfsys@defobject{currentmarker}{\pgfqpoint{-0.048611in}{0.000000in}}{\pgfqpoint{-0.000000in}{0.000000in}}{%
\pgfpathmoveto{\pgfqpoint{-0.000000in}{0.000000in}}%
\pgfpathlineto{\pgfqpoint{-0.048611in}{0.000000in}}%
\pgfusepath{stroke,fill}%
}%
\begin{pgfscope}%
\pgfsys@transformshift{0.420602in}{1.853642in}%
\pgfsys@useobject{currentmarker}{}%
\end{pgfscope}%
\end{pgfscope}%
\begin{pgfscope}%
\definecolor{textcolor}{rgb}{0.000000,0.000000,0.000000}%
\pgfsetstrokecolor{textcolor}%
\pgfsetfillcolor{textcolor}%
\pgftext[x=0.212654in, y=1.819885in, left, base]{\color{textcolor}{\rmfamily\fontsize{7.000000}{8.400000}\selectfont\catcode`\^=\active\def^{\ifmmode\sp\else\^{}\fi}\catcode`\%=\active\def%{\%}30}}%
\end{pgfscope}%
\begin{pgfscope}%
\pgfpathrectangle{\pgfqpoint{0.420602in}{0.992139in}}{\pgfqpoint{1.176290in}{1.351500in}}%
\pgfusepath{clip}%
\pgfsetrectcap%
\pgfsetroundjoin%
\pgfsetlinewidth{0.451687pt}%
\definecolor{currentstroke}{rgb}{0.690196,0.690196,0.690196}%
\pgfsetstrokecolor{currentstroke}%
\pgfsetstrokeopacity{0.280000}%
\pgfsetdash{}{0pt}%
\pgfpathmoveto{\pgfqpoint{0.420602in}{2.129547in}}%
\pgfpathlineto{\pgfqpoint{1.596892in}{2.129547in}}%
\pgfusepath{stroke}%
\end{pgfscope}%
\begin{pgfscope}%
\pgfsetbuttcap%
\pgfsetroundjoin%
\definecolor{currentfill}{rgb}{0.000000,0.000000,0.000000}%
\pgfsetfillcolor{currentfill}%
\pgfsetlinewidth{0.803000pt}%
\definecolor{currentstroke}{rgb}{0.000000,0.000000,0.000000}%
\pgfsetstrokecolor{currentstroke}%
\pgfsetdash{}{0pt}%
\pgfsys@defobject{currentmarker}{\pgfqpoint{-0.048611in}{0.000000in}}{\pgfqpoint{-0.000000in}{0.000000in}}{%
\pgfpathmoveto{\pgfqpoint{-0.000000in}{0.000000in}}%
\pgfpathlineto{\pgfqpoint{-0.048611in}{0.000000in}}%
\pgfusepath{stroke,fill}%
}%
\begin{pgfscope}%
\pgfsys@transformshift{0.420602in}{2.129547in}%
\pgfsys@useobject{currentmarker}{}%
\end{pgfscope}%
\end{pgfscope}%
\begin{pgfscope}%
\definecolor{textcolor}{rgb}{0.000000,0.000000,0.000000}%
\pgfsetstrokecolor{textcolor}%
\pgfsetfillcolor{textcolor}%
\pgftext[x=0.212654in, y=2.095789in, left, base]{\color{textcolor}{\rmfamily\fontsize{7.000000}{8.400000}\selectfont\catcode`\^=\active\def^{\ifmmode\sp\else\^{}\fi}\catcode`\%=\active\def%{\%}40}}%
\end{pgfscope}%
\begin{pgfscope}%
\definecolor{textcolor}{rgb}{0.000000,0.000000,0.000000}%
\pgfsetstrokecolor{textcolor}%
\pgfsetfillcolor{textcolor}%
\pgftext[x=0.198766in,y=1.667889in,,bottom,rotate=90.000000]{\color{textcolor}{\rmfamily\fontsize{8.000000}{9.600000}\selectfont\catcode`\^=\active\def^{\ifmmode\sp\else\^{}\fi}\catcode`\%=\active\def%{\%}Geometric ESS}}%
\end{pgfscope}%
\begin{pgfscope}%
\pgfpathrectangle{\pgfqpoint{0.420602in}{0.992139in}}{\pgfqpoint{1.176290in}{1.351500in}}%
\pgfusepath{clip}%
\pgfsetrectcap%
\pgfsetroundjoin%
\pgfsetlinewidth{1.355062pt}%
\definecolor{currentstroke}{rgb}{0.121569,0.466667,0.705882}%
\pgfsetstrokecolor{currentstroke}%
\pgfsetdash{}{0pt}%
\pgfpathmoveto{\pgfqpoint{0.474070in}{1.154484in}}%
\pgfpathlineto{\pgfqpoint{0.487606in}{1.146399in}}%
\pgfpathlineto{\pgfqpoint{0.501142in}{1.138959in}}%
\pgfpathlineto{\pgfqpoint{0.514679in}{1.132111in}}%
\pgfpathlineto{\pgfqpoint{0.528215in}{1.125810in}}%
\pgfpathlineto{\pgfqpoint{0.541751in}{1.120011in}}%
\pgfpathlineto{\pgfqpoint{0.555287in}{1.114675in}}%
\pgfpathlineto{\pgfqpoint{0.568823in}{1.109766in}}%
\pgfpathlineto{\pgfqpoint{0.582359in}{1.105248in}}%
\pgfpathlineto{\pgfqpoint{0.595895in}{1.101092in}}%
\pgfpathlineto{\pgfqpoint{0.609432in}{1.097268in}}%
\pgfpathlineto{\pgfqpoint{0.622968in}{1.093749in}}%
\pgfpathlineto{\pgfqpoint{0.636504in}{1.090512in}}%
\pgfpathlineto{\pgfqpoint{0.650040in}{1.087534in}}%
\pgfpathlineto{\pgfqpoint{0.663576in}{1.084795in}}%
\pgfpathlineto{\pgfqpoint{0.677112in}{1.082274in}}%
\pgfpathlineto{\pgfqpoint{0.690648in}{1.079956in}}%
\pgfpathlineto{\pgfqpoint{0.704184in}{1.077823in}}%
\pgfpathlineto{\pgfqpoint{0.717721in}{1.075862in}}%
\pgfpathlineto{\pgfqpoint{0.731257in}{1.074057in}}%
\pgfpathlineto{\pgfqpoint{0.744793in}{1.072398in}}%
\pgfpathlineto{\pgfqpoint{0.758329in}{1.070871in}}%
\pgfpathlineto{\pgfqpoint{0.771865in}{1.069467in}}%
\pgfpathlineto{\pgfqpoint{0.785401in}{1.068176in}}%
\pgfpathlineto{\pgfqpoint{0.798937in}{1.066989in}}%
\pgfpathlineto{\pgfqpoint{0.812473in}{1.065897in}}%
\pgfpathlineto{\pgfqpoint{0.826010in}{1.064893in}}%
\pgfpathlineto{\pgfqpoint{0.839546in}{1.063970in}}%
\pgfpathlineto{\pgfqpoint{0.853082in}{1.063121in}}%
\pgfpathlineto{\pgfqpoint{0.866618in}{1.062341in}}%
\pgfpathlineto{\pgfqpoint{0.880154in}{1.061623in}}%
\pgfpathlineto{\pgfqpoint{0.893690in}{1.060963in}}%
\pgfpathlineto{\pgfqpoint{0.907226in}{1.060357in}}%
\pgfpathlineto{\pgfqpoint{0.920762in}{1.059799in}}%
\pgfpathlineto{\pgfqpoint{0.934299in}{1.059286in}}%
\pgfpathlineto{\pgfqpoint{0.947835in}{1.058815in}}%
\pgfpathlineto{\pgfqpoint{0.961371in}{1.058382in}}%
\pgfpathlineto{\pgfqpoint{0.974907in}{1.057983in}}%
\pgfpathlineto{\pgfqpoint{0.988443in}{1.057617in}}%
\pgfpathlineto{\pgfqpoint{1.001979in}{1.057281in}}%
\pgfpathlineto{\pgfqpoint{1.015515in}{1.056971in}}%
\pgfpathlineto{\pgfqpoint{1.029052in}{1.056687in}}%
\pgfpathlineto{\pgfqpoint{1.042588in}{1.056425in}}%
\pgfpathlineto{\pgfqpoint{1.056124in}{1.056185in}}%
\pgfpathlineto{\pgfqpoint{1.069660in}{1.055964in}}%
\pgfpathlineto{\pgfqpoint{1.083196in}{1.055761in}}%
\pgfpathlineto{\pgfqpoint{1.096732in}{1.055574in}}%
\pgfpathlineto{\pgfqpoint{1.110268in}{1.055403in}}%
\pgfpathlineto{\pgfqpoint{1.123804in}{1.055245in}}%
\pgfpathlineto{\pgfqpoint{1.137341in}{1.055100in}}%
\pgfpathlineto{\pgfqpoint{1.150877in}{1.054967in}}%
\pgfpathlineto{\pgfqpoint{1.164413in}{1.054844in}}%
\pgfpathlineto{\pgfqpoint{1.177949in}{1.054732in}}%
\pgfpathlineto{\pgfqpoint{1.191485in}{1.054628in}}%
\pgfpathlineto{\pgfqpoint{1.205021in}{1.054533in}}%
\pgfpathlineto{\pgfqpoint{1.218557in}{1.054445in}}%
\pgfpathlineto{\pgfqpoint{1.232093in}{1.054365in}}%
\pgfpathlineto{\pgfqpoint{1.245630in}{1.054291in}}%
\pgfpathlineto{\pgfqpoint{1.259166in}{1.054223in}}%
\pgfpathlineto{\pgfqpoint{1.272702in}{1.054161in}}%
\pgfpathlineto{\pgfqpoint{1.286238in}{1.054103in}}%
\pgfpathlineto{\pgfqpoint{1.299774in}{1.054050in}}%
\pgfpathlineto{\pgfqpoint{1.313310in}{1.054002in}}%
\pgfpathlineto{\pgfqpoint{1.326846in}{1.053957in}}%
\pgfpathlineto{\pgfqpoint{1.340383in}{1.053916in}}%
\pgfpathlineto{\pgfqpoint{1.353919in}{1.053878in}}%
\pgfpathlineto{\pgfqpoint{1.367455in}{1.053843in}}%
\pgfpathlineto{\pgfqpoint{1.380991in}{1.053811in}}%
\pgfpathlineto{\pgfqpoint{1.394527in}{1.053782in}}%
\pgfpathlineto{\pgfqpoint{1.408063in}{1.053754in}}%
\pgfpathlineto{\pgfqpoint{1.421599in}{1.053729in}}%
\pgfpathlineto{\pgfqpoint{1.435135in}{1.053706in}}%
\pgfpathlineto{\pgfqpoint{1.448672in}{1.053684in}}%
\pgfpathlineto{\pgfqpoint{1.462208in}{1.053664in}}%
\pgfpathlineto{\pgfqpoint{1.475744in}{1.053645in}}%
\pgfpathlineto{\pgfqpoint{1.489280in}{1.053628in}}%
\pgfpathlineto{\pgfqpoint{1.502816in}{1.053612in}}%
\pgfpathlineto{\pgfqpoint{1.516352in}{1.053597in}}%
\pgfpathlineto{\pgfqpoint{1.529888in}{1.053583in}}%
\pgfpathlineto{\pgfqpoint{1.543424in}{1.053571in}}%
\pgfusepath{stroke}%
\end{pgfscope}%
\begin{pgfscope}%
\pgfpathrectangle{\pgfqpoint{0.420602in}{0.992139in}}{\pgfqpoint{1.176290in}{1.351500in}}%
\pgfusepath{clip}%
\pgfsetrectcap%
\pgfsetroundjoin%
\pgfsetlinewidth{1.355062pt}%
\definecolor{currentstroke}{rgb}{1.000000,0.498039,0.054902}%
\pgfsetstrokecolor{currentstroke}%
\pgfsetdash{}{0pt}%
\pgfpathmoveto{\pgfqpoint{0.474070in}{1.297043in}}%
\pgfpathlineto{\pgfqpoint{0.487606in}{1.279991in}}%
\pgfpathlineto{\pgfqpoint{0.501142in}{1.264241in}}%
\pgfpathlineto{\pgfqpoint{0.514679in}{1.249701in}}%
\pgfpathlineto{\pgfqpoint{0.528215in}{1.236282in}}%
\pgfpathlineto{\pgfqpoint{0.541751in}{1.223901in}}%
\pgfpathlineto{\pgfqpoint{0.555287in}{1.212483in}}%
\pgfpathlineto{\pgfqpoint{0.568823in}{1.201956in}}%
\pgfpathlineto{\pgfqpoint{0.582359in}{1.192255in}}%
\pgfpathlineto{\pgfqpoint{0.595895in}{1.183316in}}%
\pgfpathlineto{\pgfqpoint{0.609432in}{1.175083in}}%
\pgfpathlineto{\pgfqpoint{0.622968in}{1.167502in}}%
\pgfpathlineto{\pgfqpoint{0.636504in}{1.160523in}}%
\pgfpathlineto{\pgfqpoint{0.650040in}{1.154101in}}%
\pgfpathlineto{\pgfqpoint{0.663576in}{1.148192in}}%
\pgfpathlineto{\pgfqpoint{0.677112in}{1.142756in}}%
\pgfpathlineto{\pgfqpoint{0.690648in}{1.137757in}}%
\pgfpathlineto{\pgfqpoint{0.704184in}{1.133159in}}%
\pgfpathlineto{\pgfqpoint{0.717721in}{1.128932in}}%
\pgfpathlineto{\pgfqpoint{0.731257in}{1.125046in}}%
\pgfpathlineto{\pgfqpoint{0.744793in}{1.121474in}}%
\pgfpathlineto{\pgfqpoint{0.758329in}{1.118191in}}%
\pgfpathlineto{\pgfqpoint{0.771865in}{1.115173in}}%
\pgfpathlineto{\pgfqpoint{0.785401in}{1.112400in}}%
\pgfpathlineto{\pgfqpoint{0.798937in}{1.109851in}}%
\pgfpathlineto{\pgfqpoint{0.812473in}{1.107509in}}%
\pgfpathlineto{\pgfqpoint{0.826010in}{1.105357in}}%
\pgfpathlineto{\pgfqpoint{0.839546in}{1.103380in}}%
\pgfpathlineto{\pgfqpoint{0.853082in}{1.101563in}}%
\pgfpathlineto{\pgfqpoint{0.866618in}{1.099893in}}%
\pgfpathlineto{\pgfqpoint{0.880154in}{1.098359in}}%
\pgfpathlineto{\pgfqpoint{0.893690in}{1.096949in}}%
\pgfpathlineto{\pgfqpoint{0.907226in}{1.095653in}}%
\pgfpathlineto{\pgfqpoint{0.920762in}{1.094463in}}%
\pgfpathlineto{\pgfqpoint{0.934299in}{1.093370in}}%
\pgfpathlineto{\pgfqpoint{0.947835in}{1.092365in}}%
\pgfpathlineto{\pgfqpoint{0.961371in}{1.091441in}}%
\pgfpathlineto{\pgfqpoint{0.974907in}{1.090593in}}%
\pgfpathlineto{\pgfqpoint{0.988443in}{1.089813in}}%
\pgfpathlineto{\pgfqpoint{1.001979in}{1.089096in}}%
\pgfpathlineto{\pgfqpoint{1.015515in}{1.088438in}}%
\pgfpathlineto{\pgfqpoint{1.029052in}{1.087833in}}%
\pgfpathlineto{\pgfqpoint{1.042588in}{1.087277in}}%
\pgfpathlineto{\pgfqpoint{1.056124in}{1.086766in}}%
\pgfpathlineto{\pgfqpoint{1.069660in}{1.086297in}}%
\pgfpathlineto{\pgfqpoint{1.083196in}{1.085865in}}%
\pgfpathlineto{\pgfqpoint{1.096732in}{1.085469in}}%
\pgfpathlineto{\pgfqpoint{1.110268in}{1.085104in}}%
\pgfpathlineto{\pgfqpoint{1.123804in}{1.084770in}}%
\pgfpathlineto{\pgfqpoint{1.137341in}{1.084462in}}%
\pgfpathlineto{\pgfqpoint{1.150877in}{1.084179in}}%
\pgfpathlineto{\pgfqpoint{1.164413in}{1.083919in}}%
\pgfpathlineto{\pgfqpoint{1.177949in}{1.083680in}}%
\pgfpathlineto{\pgfqpoint{1.191485in}{1.083460in}}%
\pgfpathlineto{\pgfqpoint{1.205021in}{1.083257in}}%
\pgfpathlineto{\pgfqpoint{1.218557in}{1.083069in}}%
\pgfpathlineto{\pgfqpoint{1.232093in}{1.082893in}}%
\pgfpathlineto{\pgfqpoint{1.245630in}{1.082724in}}%
\pgfpathlineto{\pgfqpoint{1.259166in}{1.082554in}}%
\pgfpathlineto{\pgfqpoint{1.272702in}{1.082373in}}%
\pgfpathlineto{\pgfqpoint{1.286238in}{1.082165in}}%
\pgfpathlineto{\pgfqpoint{1.299774in}{1.081908in}}%
\pgfpathlineto{\pgfqpoint{1.313310in}{1.081573in}}%
\pgfpathlineto{\pgfqpoint{1.326846in}{1.081128in}}%
\pgfpathlineto{\pgfqpoint{1.340383in}{1.080536in}}%
\pgfpathlineto{\pgfqpoint{1.353919in}{1.079763in}}%
\pgfpathlineto{\pgfqpoint{1.367455in}{1.078780in}}%
\pgfpathlineto{\pgfqpoint{1.380991in}{1.077575in}}%
\pgfpathlineto{\pgfqpoint{1.394527in}{1.076148in}}%
\pgfpathlineto{\pgfqpoint{1.408063in}{1.074522in}}%
\pgfpathlineto{\pgfqpoint{1.421599in}{1.072734in}}%
\pgfpathlineto{\pgfqpoint{1.435135in}{1.070838in}}%
\pgfpathlineto{\pgfqpoint{1.448672in}{1.068889in}}%
\pgfpathlineto{\pgfqpoint{1.462208in}{1.066948in}}%
\pgfpathlineto{\pgfqpoint{1.475744in}{1.065065in}}%
\pgfpathlineto{\pgfqpoint{1.489280in}{1.063284in}}%
\pgfpathlineto{\pgfqpoint{1.502816in}{1.061640in}}%
\pgfpathlineto{\pgfqpoint{1.516352in}{1.060153in}}%
\pgfpathlineto{\pgfqpoint{1.529888in}{1.058838in}}%
\pgfpathlineto{\pgfqpoint{1.543424in}{1.057699in}}%
\pgfusepath{stroke}%
\end{pgfscope}%
\begin{pgfscope}%
\pgfpathrectangle{\pgfqpoint{0.420602in}{0.992139in}}{\pgfqpoint{1.176290in}{1.351500in}}%
\pgfusepath{clip}%
\pgfsetrectcap%
\pgfsetroundjoin%
\pgfsetlinewidth{1.355062pt}%
\definecolor{currentstroke}{rgb}{0.172549,0.627451,0.172549}%
\pgfsetstrokecolor{currentstroke}%
\pgfsetdash{}{0pt}%
\pgfpathmoveto{\pgfqpoint{0.474070in}{2.282207in}}%
\pgfpathlineto{\pgfqpoint{0.487606in}{2.220480in}}%
\pgfpathlineto{\pgfqpoint{0.501142in}{2.161958in}}%
\pgfpathlineto{\pgfqpoint{0.514679in}{2.106505in}}%
\pgfpathlineto{\pgfqpoint{0.528215in}{2.053981in}}%
\pgfpathlineto{\pgfqpoint{0.541751in}{2.004247in}}%
\pgfpathlineto{\pgfqpoint{0.555287in}{1.957166in}}%
\pgfpathlineto{\pgfqpoint{0.568823in}{1.912602in}}%
\pgfpathlineto{\pgfqpoint{0.582359in}{1.870420in}}%
\pgfpathlineto{\pgfqpoint{0.595895in}{1.830490in}}%
\pgfpathlineto{\pgfqpoint{0.609432in}{1.792686in}}%
\pgfpathlineto{\pgfqpoint{0.622968in}{1.756887in}}%
\pgfpathlineto{\pgfqpoint{0.636504in}{1.722976in}}%
\pgfpathlineto{\pgfqpoint{0.650040in}{1.690843in}}%
\pgfpathlineto{\pgfqpoint{0.663576in}{1.660383in}}%
\pgfpathlineto{\pgfqpoint{0.677112in}{1.631498in}}%
\pgfpathlineto{\pgfqpoint{0.690648in}{1.604094in}}%
\pgfpathlineto{\pgfqpoint{0.704184in}{1.578086in}}%
\pgfpathlineto{\pgfqpoint{0.717721in}{1.553391in}}%
\pgfpathlineto{\pgfqpoint{0.731257in}{1.529934in}}%
\pgfpathlineto{\pgfqpoint{0.744793in}{1.507645in}}%
\pgfpathlineto{\pgfqpoint{0.758329in}{1.486458in}}%
\pgfpathlineto{\pgfqpoint{0.771865in}{1.466310in}}%
\pgfpathlineto{\pgfqpoint{0.785401in}{1.447143in}}%
\pgfpathlineto{\pgfqpoint{0.798937in}{1.428903in}}%
\pgfpathlineto{\pgfqpoint{0.812473in}{1.411539in}}%
\pgfpathlineto{\pgfqpoint{0.826010in}{1.395001in}}%
\pgfpathlineto{\pgfqpoint{0.839546in}{1.379243in}}%
\pgfpathlineto{\pgfqpoint{0.853082in}{1.364224in}}%
\pgfpathlineto{\pgfqpoint{0.866618in}{1.349900in}}%
\pgfpathlineto{\pgfqpoint{0.880154in}{1.336234in}}%
\pgfpathlineto{\pgfqpoint{0.893690in}{1.323189in}}%
\pgfpathlineto{\pgfqpoint{0.907226in}{1.310731in}}%
\pgfpathlineto{\pgfqpoint{0.920762in}{1.298827in}}%
\pgfpathlineto{\pgfqpoint{0.934299in}{1.287447in}}%
\pgfpathlineto{\pgfqpoint{0.947835in}{1.276562in}}%
\pgfpathlineto{\pgfqpoint{0.961371in}{1.266147in}}%
\pgfpathlineto{\pgfqpoint{0.974907in}{1.256175in}}%
\pgfpathlineto{\pgfqpoint{0.988443in}{1.246624in}}%
\pgfpathlineto{\pgfqpoint{1.001979in}{1.237472in}}%
\pgfpathlineto{\pgfqpoint{1.015515in}{1.228698in}}%
\pgfpathlineto{\pgfqpoint{1.029052in}{1.220283in}}%
\pgfpathlineto{\pgfqpoint{1.042588in}{1.212210in}}%
\pgfpathlineto{\pgfqpoint{1.056124in}{1.204460in}}%
\pgfpathlineto{\pgfqpoint{1.069660in}{1.197020in}}%
\pgfpathlineto{\pgfqpoint{1.083196in}{1.189873in}}%
\pgfpathlineto{\pgfqpoint{1.096732in}{1.183006in}}%
\pgfpathlineto{\pgfqpoint{1.110268in}{1.176407in}}%
\pgfpathlineto{\pgfqpoint{1.123804in}{1.170061in}}%
\pgfpathlineto{\pgfqpoint{1.137341in}{1.163958in}}%
\pgfpathlineto{\pgfqpoint{1.150877in}{1.158087in}}%
\pgfpathlineto{\pgfqpoint{1.164413in}{1.152436in}}%
\pgfpathlineto{\pgfqpoint{1.177949in}{1.146995in}}%
\pgfpathlineto{\pgfqpoint{1.191485in}{1.141754in}}%
\pgfpathlineto{\pgfqpoint{1.205021in}{1.136703in}}%
\pgfpathlineto{\pgfqpoint{1.218557in}{1.131833in}}%
\pgfpathlineto{\pgfqpoint{1.232093in}{1.127133in}}%
\pgfpathlineto{\pgfqpoint{1.245630in}{1.122594in}}%
\pgfpathlineto{\pgfqpoint{1.259166in}{1.118206in}}%
\pgfpathlineto{\pgfqpoint{1.272702in}{1.113960in}}%
\pgfpathlineto{\pgfqpoint{1.286238in}{1.109845in}}%
\pgfpathlineto{\pgfqpoint{1.299774in}{1.105853in}}%
\pgfpathlineto{\pgfqpoint{1.313310in}{1.101977in}}%
\pgfpathlineto{\pgfqpoint{1.326846in}{1.098211in}}%
\pgfpathlineto{\pgfqpoint{1.340383in}{1.094554in}}%
\pgfpathlineto{\pgfqpoint{1.353919in}{1.091005in}}%
\pgfpathlineto{\pgfqpoint{1.367455in}{1.087570in}}%
\pgfpathlineto{\pgfqpoint{1.380991in}{1.084257in}}%
\pgfpathlineto{\pgfqpoint{1.394527in}{1.081077in}}%
\pgfpathlineto{\pgfqpoint{1.408063in}{1.078041in}}%
\pgfpathlineto{\pgfqpoint{1.421599in}{1.075162in}}%
\pgfpathlineto{\pgfqpoint{1.435135in}{1.072452in}}%
\pgfpathlineto{\pgfqpoint{1.448672in}{1.069923in}}%
\pgfpathlineto{\pgfqpoint{1.462208in}{1.067584in}}%
\pgfpathlineto{\pgfqpoint{1.475744in}{1.065439in}}%
\pgfpathlineto{\pgfqpoint{1.489280in}{1.063495in}}%
\pgfpathlineto{\pgfqpoint{1.502816in}{1.061752in}}%
\pgfpathlineto{\pgfqpoint{1.516352in}{1.060209in}}%
\pgfpathlineto{\pgfqpoint{1.529888in}{1.058864in}}%
\pgfpathlineto{\pgfqpoint{1.543424in}{1.057710in}}%
\pgfusepath{stroke}%
\end{pgfscope}%
\begin{pgfscope}%
\pgfpathrectangle{\pgfqpoint{0.420602in}{0.992139in}}{\pgfqpoint{1.176290in}{1.351500in}}%
\pgfusepath{clip}%
\pgfsetrectcap%
\pgfsetroundjoin%
\pgfsetlinewidth{1.355062pt}%
\definecolor{currentstroke}{rgb}{0.839216,0.152941,0.156863}%
\pgfsetstrokecolor{currentstroke}%
\pgfsetdash{}{0pt}%
\pgfpathmoveto{\pgfqpoint{0.474070in}{1.504770in}}%
\pgfpathlineto{\pgfqpoint{0.487606in}{1.475937in}}%
\pgfpathlineto{\pgfqpoint{0.501142in}{1.449271in}}%
\pgfpathlineto{\pgfqpoint{0.514679in}{1.424623in}}%
\pgfpathlineto{\pgfqpoint{0.528215in}{1.401854in}}%
\pgfpathlineto{\pgfqpoint{0.541751in}{1.380828in}}%
\pgfpathlineto{\pgfqpoint{0.555287in}{1.361422in}}%
\pgfpathlineto{\pgfqpoint{0.568823in}{1.343517in}}%
\pgfpathlineto{\pgfqpoint{0.582359in}{1.327003in}}%
\pgfpathlineto{\pgfqpoint{0.595895in}{1.311777in}}%
\pgfpathlineto{\pgfqpoint{0.609432in}{1.297742in}}%
\pgfpathlineto{\pgfqpoint{0.622968in}{1.284810in}}%
\pgfpathlineto{\pgfqpoint{0.636504in}{1.272895in}}%
\pgfpathlineto{\pgfqpoint{0.650040in}{1.261921in}}%
\pgfpathlineto{\pgfqpoint{0.663576in}{1.251815in}}%
\pgfpathlineto{\pgfqpoint{0.677112in}{1.242511in}}%
\pgfpathlineto{\pgfqpoint{0.690648in}{1.233947in}}%
\pgfpathlineto{\pgfqpoint{0.704184in}{1.226065in}}%
\pgfpathlineto{\pgfqpoint{0.717721in}{1.218811in}}%
\pgfpathlineto{\pgfqpoint{0.731257in}{1.212138in}}%
\pgfpathlineto{\pgfqpoint{0.744793in}{1.205999in}}%
\pgfpathlineto{\pgfqpoint{0.758329in}{1.200352in}}%
\pgfpathlineto{\pgfqpoint{0.771865in}{1.195158in}}%
\pgfpathlineto{\pgfqpoint{0.785401in}{1.190381in}}%
\pgfpathlineto{\pgfqpoint{0.798937in}{1.185989in}}%
\pgfpathlineto{\pgfqpoint{0.812473in}{1.181951in}}%
\pgfpathlineto{\pgfqpoint{0.826010in}{1.178238in}}%
\pgfpathlineto{\pgfqpoint{0.839546in}{1.174825in}}%
\pgfpathlineto{\pgfqpoint{0.853082in}{1.171687in}}%
\pgfpathlineto{\pgfqpoint{0.866618in}{1.168803in}}%
\pgfpathlineto{\pgfqpoint{0.880154in}{1.166152in}}%
\pgfpathlineto{\pgfqpoint{0.893690in}{1.163715in}}%
\pgfpathlineto{\pgfqpoint{0.907226in}{1.161475in}}%
\pgfpathlineto{\pgfqpoint{0.920762in}{1.159416in}}%
\pgfpathlineto{\pgfqpoint{0.934299in}{1.157524in}}%
\pgfpathlineto{\pgfqpoint{0.947835in}{1.155785in}}%
\pgfpathlineto{\pgfqpoint{0.961371in}{1.154187in}}%
\pgfpathlineto{\pgfqpoint{0.974907in}{1.152719in}}%
\pgfpathlineto{\pgfqpoint{0.988443in}{1.151369in}}%
\pgfpathlineto{\pgfqpoint{1.001979in}{1.150128in}}%
\pgfpathlineto{\pgfqpoint{1.015515in}{1.148988in}}%
\pgfpathlineto{\pgfqpoint{1.029052in}{1.147940in}}%
\pgfpathlineto{\pgfqpoint{1.042588in}{1.146976in}}%
\pgfpathlineto{\pgfqpoint{1.056124in}{1.146088in}}%
\pgfpathlineto{\pgfqpoint{1.069660in}{1.145269in}}%
\pgfpathlineto{\pgfqpoint{1.083196in}{1.144509in}}%
\pgfpathlineto{\pgfqpoint{1.096732in}{1.143796in}}%
\pgfpathlineto{\pgfqpoint{1.110268in}{1.143116in}}%
\pgfpathlineto{\pgfqpoint{1.123804in}{1.142446in}}%
\pgfpathlineto{\pgfqpoint{1.137341in}{1.141758in}}%
\pgfpathlineto{\pgfqpoint{1.150877in}{1.141014in}}%
\pgfpathlineto{\pgfqpoint{1.164413in}{1.140164in}}%
\pgfpathlineto{\pgfqpoint{1.177949in}{1.139149in}}%
\pgfpathlineto{\pgfqpoint{1.191485in}{1.137899in}}%
\pgfpathlineto{\pgfqpoint{1.205021in}{1.136342in}}%
\pgfpathlineto{\pgfqpoint{1.218557in}{1.134410in}}%
\pgfpathlineto{\pgfqpoint{1.232093in}{1.132046in}}%
\pgfpathlineto{\pgfqpoint{1.245630in}{1.129215in}}%
\pgfpathlineto{\pgfqpoint{1.259166in}{1.125914in}}%
\pgfpathlineto{\pgfqpoint{1.272702in}{1.122173in}}%
\pgfpathlineto{\pgfqpoint{1.286238in}{1.118054in}}%
\pgfpathlineto{\pgfqpoint{1.299774in}{1.113646in}}%
\pgfpathlineto{\pgfqpoint{1.313310in}{1.109055in}}%
\pgfpathlineto{\pgfqpoint{1.326846in}{1.104390in}}%
\pgfpathlineto{\pgfqpoint{1.340383in}{1.099751in}}%
\pgfpathlineto{\pgfqpoint{1.353919in}{1.095228in}}%
\pgfpathlineto{\pgfqpoint{1.367455in}{1.090888in}}%
\pgfpathlineto{\pgfqpoint{1.380991in}{1.086779in}}%
\pgfpathlineto{\pgfqpoint{1.394527in}{1.082933in}}%
\pgfpathlineto{\pgfqpoint{1.408063in}{1.079364in}}%
\pgfpathlineto{\pgfqpoint{1.421599in}{1.076075in}}%
\pgfpathlineto{\pgfqpoint{1.435135in}{1.073063in}}%
\pgfpathlineto{\pgfqpoint{1.448672in}{1.070318in}}%
\pgfpathlineto{\pgfqpoint{1.462208in}{1.067831in}}%
\pgfpathlineto{\pgfqpoint{1.475744in}{1.065589in}}%
\pgfpathlineto{\pgfqpoint{1.489280in}{1.063583in}}%
\pgfpathlineto{\pgfqpoint{1.502816in}{1.061802in}}%
\pgfpathlineto{\pgfqpoint{1.516352in}{1.060238in}}%
\pgfpathlineto{\pgfqpoint{1.529888in}{1.058880in}}%
\pgfpathlineto{\pgfqpoint{1.543424in}{1.057718in}}%
\pgfusepath{stroke}%
\end{pgfscope}%
\begin{pgfscope}%
\pgfpathrectangle{\pgfqpoint{0.420602in}{0.992139in}}{\pgfqpoint{1.176290in}{1.351500in}}%
\pgfusepath{clip}%
\pgfsetrectcap%
\pgfsetroundjoin%
\pgfsetlinewidth{1.355062pt}%
\definecolor{currentstroke}{rgb}{0.580392,0.403922,0.741176}%
\pgfsetstrokecolor{currentstroke}%
\pgfsetdash{}{0pt}%
\pgfpathmoveto{\pgfqpoint{0.474070in}{1.853642in}}%
\pgfpathlineto{\pgfqpoint{0.487606in}{1.853642in}}%
\pgfpathlineto{\pgfqpoint{0.501142in}{1.853642in}}%
\pgfpathlineto{\pgfqpoint{0.514679in}{1.853642in}}%
\pgfpathlineto{\pgfqpoint{0.528215in}{1.853642in}}%
\pgfpathlineto{\pgfqpoint{0.541751in}{1.853642in}}%
\pgfpathlineto{\pgfqpoint{0.555287in}{1.853642in}}%
\pgfpathlineto{\pgfqpoint{0.568823in}{1.853642in}}%
\pgfpathlineto{\pgfqpoint{0.582359in}{1.853642in}}%
\pgfpathlineto{\pgfqpoint{0.595895in}{1.853642in}}%
\pgfpathlineto{\pgfqpoint{0.609432in}{1.853642in}}%
\pgfpathlineto{\pgfqpoint{0.622968in}{1.853642in}}%
\pgfpathlineto{\pgfqpoint{0.636504in}{1.853642in}}%
\pgfpathlineto{\pgfqpoint{0.650040in}{1.853642in}}%
\pgfpathlineto{\pgfqpoint{0.663576in}{1.853642in}}%
\pgfpathlineto{\pgfqpoint{0.677112in}{1.853641in}}%
\pgfpathlineto{\pgfqpoint{0.690648in}{1.853637in}}%
\pgfpathlineto{\pgfqpoint{0.704184in}{1.853629in}}%
\pgfpathlineto{\pgfqpoint{0.717721in}{1.853607in}}%
\pgfpathlineto{\pgfqpoint{0.731257in}{1.853557in}}%
\pgfpathlineto{\pgfqpoint{0.744793in}{1.853451in}}%
\pgfpathlineto{\pgfqpoint{0.758329in}{1.853236in}}%
\pgfpathlineto{\pgfqpoint{0.771865in}{1.852827in}}%
\pgfpathlineto{\pgfqpoint{0.785401in}{1.852088in}}%
\pgfpathlineto{\pgfqpoint{0.798937in}{1.850823in}}%
\pgfpathlineto{\pgfqpoint{0.812473in}{1.848757in}}%
\pgfpathlineto{\pgfqpoint{0.826010in}{1.845532in}}%
\pgfpathlineto{\pgfqpoint{0.839546in}{1.840708in}}%
\pgfpathlineto{\pgfqpoint{0.853082in}{1.833783in}}%
\pgfpathlineto{\pgfqpoint{0.866618in}{1.824228in}}%
\pgfpathlineto{\pgfqpoint{0.880154in}{1.811538in}}%
\pgfpathlineto{\pgfqpoint{0.893690in}{1.795298in}}%
\pgfpathlineto{\pgfqpoint{0.907226in}{1.775251in}}%
\pgfpathlineto{\pgfqpoint{0.920762in}{1.751353in}}%
\pgfpathlineto{\pgfqpoint{0.934299in}{1.723803in}}%
\pgfpathlineto{\pgfqpoint{0.947835in}{1.693042in}}%
\pgfpathlineto{\pgfqpoint{0.961371in}{1.659705in}}%
\pgfpathlineto{\pgfqpoint{0.974907in}{1.624560in}}%
\pgfpathlineto{\pgfqpoint{0.988443in}{1.588422in}}%
\pgfpathlineto{\pgfqpoint{1.001979in}{1.552074in}}%
\pgfpathlineto{\pgfqpoint{1.015515in}{1.516209in}}%
\pgfpathlineto{\pgfqpoint{1.029052in}{1.481393in}}%
\pgfpathlineto{\pgfqpoint{1.042588in}{1.448047in}}%
\pgfpathlineto{\pgfqpoint{1.056124in}{1.416452in}}%
\pgfpathlineto{\pgfqpoint{1.069660in}{1.386771in}}%
\pgfpathlineto{\pgfqpoint{1.083196in}{1.359069in}}%
\pgfpathlineto{\pgfqpoint{1.096732in}{1.333338in}}%
\pgfpathlineto{\pgfqpoint{1.110268in}{1.309521in}}%
\pgfpathlineto{\pgfqpoint{1.123804in}{1.287528in}}%
\pgfpathlineto{\pgfqpoint{1.137341in}{1.267254in}}%
\pgfpathlineto{\pgfqpoint{1.150877in}{1.248584in}}%
\pgfpathlineto{\pgfqpoint{1.164413in}{1.231404in}}%
\pgfpathlineto{\pgfqpoint{1.177949in}{1.215603in}}%
\pgfpathlineto{\pgfqpoint{1.191485in}{1.201074in}}%
\pgfpathlineto{\pgfqpoint{1.205021in}{1.187719in}}%
\pgfpathlineto{\pgfqpoint{1.218557in}{1.175445in}}%
\pgfpathlineto{\pgfqpoint{1.232093in}{1.164166in}}%
\pgfpathlineto{\pgfqpoint{1.245630in}{1.153804in}}%
\pgfpathlineto{\pgfqpoint{1.259166in}{1.144285in}}%
\pgfpathlineto{\pgfqpoint{1.272702in}{1.135542in}}%
\pgfpathlineto{\pgfqpoint{1.286238in}{1.127515in}}%
\pgfpathlineto{\pgfqpoint{1.299774in}{1.120145in}}%
\pgfpathlineto{\pgfqpoint{1.313310in}{1.113381in}}%
\pgfpathlineto{\pgfqpoint{1.326846in}{1.107175in}}%
\pgfpathlineto{\pgfqpoint{1.340383in}{1.101484in}}%
\pgfpathlineto{\pgfqpoint{1.353919in}{1.096266in}}%
\pgfpathlineto{\pgfqpoint{1.367455in}{1.091486in}}%
\pgfpathlineto{\pgfqpoint{1.380991in}{1.087110in}}%
\pgfpathlineto{\pgfqpoint{1.394527in}{1.083108in}}%
\pgfpathlineto{\pgfqpoint{1.408063in}{1.079452in}}%
\pgfpathlineto{\pgfqpoint{1.421599in}{1.076117in}}%
\pgfpathlineto{\pgfqpoint{1.435135in}{1.073082in}}%
\pgfpathlineto{\pgfqpoint{1.448672in}{1.070327in}}%
\pgfpathlineto{\pgfqpoint{1.462208in}{1.067834in}}%
\pgfpathlineto{\pgfqpoint{1.475744in}{1.065591in}}%
\pgfpathlineto{\pgfqpoint{1.489280in}{1.063584in}}%
\pgfpathlineto{\pgfqpoint{1.502816in}{1.061803in}}%
\pgfpathlineto{\pgfqpoint{1.516352in}{1.060238in}}%
\pgfpathlineto{\pgfqpoint{1.529888in}{1.058880in}}%
\pgfpathlineto{\pgfqpoint{1.543424in}{1.057719in}}%
\pgfusepath{stroke}%
\end{pgfscope}%
\begin{pgfscope}%
\pgfsetrectcap%
\pgfsetmiterjoin%
\pgfsetlinewidth{0.803000pt}%
\definecolor{currentstroke}{rgb}{0.000000,0.000000,0.000000}%
\pgfsetstrokecolor{currentstroke}%
\pgfsetdash{}{0pt}%
\pgfpathmoveto{\pgfqpoint{0.420602in}{0.992139in}}%
\pgfpathlineto{\pgfqpoint{0.420602in}{2.343639in}}%
\pgfusepath{stroke}%
\end{pgfscope}%
\begin{pgfscope}%
\pgfsetrectcap%
\pgfsetmiterjoin%
\pgfsetlinewidth{0.803000pt}%
\definecolor{currentstroke}{rgb}{0.000000,0.000000,0.000000}%
\pgfsetstrokecolor{currentstroke}%
\pgfsetdash{}{0pt}%
\pgfpathmoveto{\pgfqpoint{1.596892in}{0.992139in}}%
\pgfpathlineto{\pgfqpoint{1.596892in}{2.343639in}}%
\pgfusepath{stroke}%
\end{pgfscope}%
\begin{pgfscope}%
\pgfsetrectcap%
\pgfsetmiterjoin%
\pgfsetlinewidth{0.803000pt}%
\definecolor{currentstroke}{rgb}{0.000000,0.000000,0.000000}%
\pgfsetstrokecolor{currentstroke}%
\pgfsetdash{}{0pt}%
\pgfpathmoveto{\pgfqpoint{0.420602in}{0.992139in}}%
\pgfpathlineto{\pgfqpoint{1.596892in}{0.992139in}}%
\pgfusepath{stroke}%
\end{pgfscope}%
\begin{pgfscope}%
\pgfsetrectcap%
\pgfsetmiterjoin%
\pgfsetlinewidth{0.803000pt}%
\definecolor{currentstroke}{rgb}{0.000000,0.000000,0.000000}%
\pgfsetstrokecolor{currentstroke}%
\pgfsetdash{}{0pt}%
\pgfpathmoveto{\pgfqpoint{0.420602in}{2.343639in}}%
\pgfpathlineto{\pgfqpoint{1.596892in}{2.343639in}}%
\pgfusepath{stroke}%
\end{pgfscope}%
\begin{pgfscope}%
\definecolor{textcolor}{rgb}{0.000000,0.000000,0.000000}%
\pgfsetstrokecolor{textcolor}%
\pgfsetfillcolor{textcolor}%
\pgftext[x=1.008747in,y=2.385306in,,base]{\color{textcolor}{\rmfamily\fontsize{8.000000}{9.600000}\selectfont\catcode`\^=\active\def^{\ifmmode\sp\else\^{}\fi}\catcode`\%=\active\def%{\%}(a) gESS}}%
\end{pgfscope}%
\begin{pgfscope}%
\pgfsetbuttcap%
\pgfsetmiterjoin%
\definecolor{currentfill}{rgb}{1.000000,1.000000,1.000000}%
\pgfsetfillcolor{currentfill}%
\pgfsetlinewidth{0.000000pt}%
\definecolor{currentstroke}{rgb}{0.000000,0.000000,0.000000}%
\pgfsetstrokecolor{currentstroke}%
\pgfsetstrokeopacity{0.000000}%
\pgfsetdash{}{0pt}%
\pgfpathmoveto{\pgfqpoint{2.208563in}{0.992139in}}%
\pgfpathlineto{\pgfqpoint{3.384853in}{0.992139in}}%
\pgfpathlineto{\pgfqpoint{3.384853in}{2.343639in}}%
\pgfpathlineto{\pgfqpoint{2.208563in}{2.343639in}}%
\pgfpathlineto{\pgfqpoint{2.208563in}{0.992139in}}%
\pgfpathclose%
\pgfusepath{fill}%
\end{pgfscope}%
\begin{pgfscope}%
\pgfpathrectangle{\pgfqpoint{2.208563in}{0.992139in}}{\pgfqpoint{1.176290in}{1.351500in}}%
\pgfusepath{clip}%
\pgfsetrectcap%
\pgfsetroundjoin%
\pgfsetlinewidth{0.451687pt}%
\definecolor{currentstroke}{rgb}{0.690196,0.690196,0.690196}%
\pgfsetstrokecolor{currentstroke}%
\pgfsetstrokeopacity{0.280000}%
\pgfsetdash{}{0pt}%
\pgfpathmoveto{\pgfqpoint{2.484472in}{0.992139in}}%
\pgfpathlineto{\pgfqpoint{2.484472in}{2.343639in}}%
\pgfusepath{stroke}%
\end{pgfscope}%
\begin{pgfscope}%
\pgfsetbuttcap%
\pgfsetroundjoin%
\definecolor{currentfill}{rgb}{0.000000,0.000000,0.000000}%
\pgfsetfillcolor{currentfill}%
\pgfsetlinewidth{0.803000pt}%
\definecolor{currentstroke}{rgb}{0.000000,0.000000,0.000000}%
\pgfsetstrokecolor{currentstroke}%
\pgfsetdash{}{0pt}%
\pgfsys@defobject{currentmarker}{\pgfqpoint{0.000000in}{-0.048611in}}{\pgfqpoint{0.000000in}{0.000000in}}{%
\pgfpathmoveto{\pgfqpoint{0.000000in}{0.000000in}}%
\pgfpathlineto{\pgfqpoint{0.000000in}{-0.048611in}}%
\pgfusepath{stroke,fill}%
}%
\begin{pgfscope}%
\pgfsys@transformshift{2.484472in}{0.992139in}%
\pgfsys@useobject{currentmarker}{}%
\end{pgfscope}%
\end{pgfscope}%
\begin{pgfscope}%
\definecolor{textcolor}{rgb}{0.000000,0.000000,0.000000}%
\pgfsetstrokecolor{textcolor}%
\pgfsetfillcolor{textcolor}%
\pgftext[x=2.484472in,y=0.894917in,,top]{\color{textcolor}{\rmfamily\fontsize{7.000000}{8.400000}\selectfont\catcode`\^=\active\def^{\ifmmode\sp\else\^{}\fi}\catcode`\%=\active\def%{\%}$\mathdefault{10^{1}}$}}%
\end{pgfscope}%
\begin{pgfscope}%
\pgfpathrectangle{\pgfqpoint{2.208563in}{0.992139in}}{\pgfqpoint{1.176290in}{1.351500in}}%
\pgfusepath{clip}%
\pgfsetrectcap%
\pgfsetroundjoin%
\pgfsetlinewidth{0.451687pt}%
\definecolor{currentstroke}{rgb}{0.690196,0.690196,0.690196}%
\pgfsetstrokecolor{currentstroke}%
\pgfsetstrokeopacity{0.280000}%
\pgfsetdash{}{0pt}%
\pgfpathmoveto{\pgfqpoint{3.223406in}{0.992139in}}%
\pgfpathlineto{\pgfqpoint{3.223406in}{2.343639in}}%
\pgfusepath{stroke}%
\end{pgfscope}%
\begin{pgfscope}%
\pgfsetbuttcap%
\pgfsetroundjoin%
\definecolor{currentfill}{rgb}{0.000000,0.000000,0.000000}%
\pgfsetfillcolor{currentfill}%
\pgfsetlinewidth{0.803000pt}%
\definecolor{currentstroke}{rgb}{0.000000,0.000000,0.000000}%
\pgfsetstrokecolor{currentstroke}%
\pgfsetdash{}{0pt}%
\pgfsys@defobject{currentmarker}{\pgfqpoint{0.000000in}{-0.048611in}}{\pgfqpoint{0.000000in}{0.000000in}}{%
\pgfpathmoveto{\pgfqpoint{0.000000in}{0.000000in}}%
\pgfpathlineto{\pgfqpoint{0.000000in}{-0.048611in}}%
\pgfusepath{stroke,fill}%
}%
\begin{pgfscope}%
\pgfsys@transformshift{3.223406in}{0.992139in}%
\pgfsys@useobject{currentmarker}{}%
\end{pgfscope}%
\end{pgfscope}%
\begin{pgfscope}%
\definecolor{textcolor}{rgb}{0.000000,0.000000,0.000000}%
\pgfsetstrokecolor{textcolor}%
\pgfsetfillcolor{textcolor}%
\pgftext[x=3.223406in,y=0.894917in,,top]{\color{textcolor}{\rmfamily\fontsize{7.000000}{8.400000}\selectfont\catcode`\^=\active\def^{\ifmmode\sp\else\^{}\fi}\catcode`\%=\active\def%{\%}$\mathdefault{10^{2}}$}}%
\end{pgfscope}%
\begin{pgfscope}%
\pgfsetbuttcap%
\pgfsetroundjoin%
\definecolor{currentfill}{rgb}{0.000000,0.000000,0.000000}%
\pgfsetfillcolor{currentfill}%
\pgfsetlinewidth{0.602250pt}%
\definecolor{currentstroke}{rgb}{0.000000,0.000000,0.000000}%
\pgfsetstrokecolor{currentstroke}%
\pgfsetdash{}{0pt}%
\pgfsys@defobject{currentmarker}{\pgfqpoint{0.000000in}{-0.027778in}}{\pgfqpoint{0.000000in}{0.000000in}}{%
\pgfpathmoveto{\pgfqpoint{0.000000in}{0.000000in}}%
\pgfpathlineto{\pgfqpoint{0.000000in}{-0.027778in}}%
\pgfusepath{stroke,fill}%
}%
\begin{pgfscope}%
\pgfsys@transformshift{2.262031in}{0.992139in}%
\pgfsys@useobject{currentmarker}{}%
\end{pgfscope}%
\end{pgfscope}%
\begin{pgfscope}%
\pgfsetbuttcap%
\pgfsetroundjoin%
\definecolor{currentfill}{rgb}{0.000000,0.000000,0.000000}%
\pgfsetfillcolor{currentfill}%
\pgfsetlinewidth{0.602250pt}%
\definecolor{currentstroke}{rgb}{0.000000,0.000000,0.000000}%
\pgfsetstrokecolor{currentstroke}%
\pgfsetdash{}{0pt}%
\pgfsys@defobject{currentmarker}{\pgfqpoint{0.000000in}{-0.027778in}}{\pgfqpoint{0.000000in}{0.000000in}}{%
\pgfpathmoveto{\pgfqpoint{0.000000in}{0.000000in}}%
\pgfpathlineto{\pgfqpoint{0.000000in}{-0.027778in}}%
\pgfusepath{stroke,fill}%
}%
\begin{pgfscope}%
\pgfsys@transformshift{2.320540in}{0.992139in}%
\pgfsys@useobject{currentmarker}{}%
\end{pgfscope}%
\end{pgfscope}%
\begin{pgfscope}%
\pgfsetbuttcap%
\pgfsetroundjoin%
\definecolor{currentfill}{rgb}{0.000000,0.000000,0.000000}%
\pgfsetfillcolor{currentfill}%
\pgfsetlinewidth{0.602250pt}%
\definecolor{currentstroke}{rgb}{0.000000,0.000000,0.000000}%
\pgfsetstrokecolor{currentstroke}%
\pgfsetdash{}{0pt}%
\pgfsys@defobject{currentmarker}{\pgfqpoint{0.000000in}{-0.027778in}}{\pgfqpoint{0.000000in}{0.000000in}}{%
\pgfpathmoveto{\pgfqpoint{0.000000in}{0.000000in}}%
\pgfpathlineto{\pgfqpoint{0.000000in}{-0.027778in}}%
\pgfusepath{stroke,fill}%
}%
\begin{pgfscope}%
\pgfsys@transformshift{2.370010in}{0.992139in}%
\pgfsys@useobject{currentmarker}{}%
\end{pgfscope}%
\end{pgfscope}%
\begin{pgfscope}%
\pgfsetbuttcap%
\pgfsetroundjoin%
\definecolor{currentfill}{rgb}{0.000000,0.000000,0.000000}%
\pgfsetfillcolor{currentfill}%
\pgfsetlinewidth{0.602250pt}%
\definecolor{currentstroke}{rgb}{0.000000,0.000000,0.000000}%
\pgfsetstrokecolor{currentstroke}%
\pgfsetdash{}{0pt}%
\pgfsys@defobject{currentmarker}{\pgfqpoint{0.000000in}{-0.027778in}}{\pgfqpoint{0.000000in}{0.000000in}}{%
\pgfpathmoveto{\pgfqpoint{0.000000in}{0.000000in}}%
\pgfpathlineto{\pgfqpoint{0.000000in}{-0.027778in}}%
\pgfusepath{stroke,fill}%
}%
\begin{pgfscope}%
\pgfsys@transformshift{2.412862in}{0.992139in}%
\pgfsys@useobject{currentmarker}{}%
\end{pgfscope}%
\end{pgfscope}%
\begin{pgfscope}%
\pgfsetbuttcap%
\pgfsetroundjoin%
\definecolor{currentfill}{rgb}{0.000000,0.000000,0.000000}%
\pgfsetfillcolor{currentfill}%
\pgfsetlinewidth{0.602250pt}%
\definecolor{currentstroke}{rgb}{0.000000,0.000000,0.000000}%
\pgfsetstrokecolor{currentstroke}%
\pgfsetdash{}{0pt}%
\pgfsys@defobject{currentmarker}{\pgfqpoint{0.000000in}{-0.027778in}}{\pgfqpoint{0.000000in}{0.000000in}}{%
\pgfpathmoveto{\pgfqpoint{0.000000in}{0.000000in}}%
\pgfpathlineto{\pgfqpoint{0.000000in}{-0.027778in}}%
\pgfusepath{stroke,fill}%
}%
\begin{pgfscope}%
\pgfsys@transformshift{2.450660in}{0.992139in}%
\pgfsys@useobject{currentmarker}{}%
\end{pgfscope}%
\end{pgfscope}%
\begin{pgfscope}%
\pgfsetbuttcap%
\pgfsetroundjoin%
\definecolor{currentfill}{rgb}{0.000000,0.000000,0.000000}%
\pgfsetfillcolor{currentfill}%
\pgfsetlinewidth{0.602250pt}%
\definecolor{currentstroke}{rgb}{0.000000,0.000000,0.000000}%
\pgfsetstrokecolor{currentstroke}%
\pgfsetdash{}{0pt}%
\pgfsys@defobject{currentmarker}{\pgfqpoint{0.000000in}{-0.027778in}}{\pgfqpoint{0.000000in}{0.000000in}}{%
\pgfpathmoveto{\pgfqpoint{0.000000in}{0.000000in}}%
\pgfpathlineto{\pgfqpoint{0.000000in}{-0.027778in}}%
\pgfusepath{stroke,fill}%
}%
\begin{pgfscope}%
\pgfsys@transformshift{2.706913in}{0.992139in}%
\pgfsys@useobject{currentmarker}{}%
\end{pgfscope}%
\end{pgfscope}%
\begin{pgfscope}%
\pgfsetbuttcap%
\pgfsetroundjoin%
\definecolor{currentfill}{rgb}{0.000000,0.000000,0.000000}%
\pgfsetfillcolor{currentfill}%
\pgfsetlinewidth{0.602250pt}%
\definecolor{currentstroke}{rgb}{0.000000,0.000000,0.000000}%
\pgfsetstrokecolor{currentstroke}%
\pgfsetdash{}{0pt}%
\pgfsys@defobject{currentmarker}{\pgfqpoint{0.000000in}{-0.027778in}}{\pgfqpoint{0.000000in}{0.000000in}}{%
\pgfpathmoveto{\pgfqpoint{0.000000in}{0.000000in}}%
\pgfpathlineto{\pgfqpoint{0.000000in}{-0.027778in}}%
\pgfusepath{stroke,fill}%
}%
\begin{pgfscope}%
\pgfsys@transformshift{2.837033in}{0.992139in}%
\pgfsys@useobject{currentmarker}{}%
\end{pgfscope}%
\end{pgfscope}%
\begin{pgfscope}%
\pgfsetbuttcap%
\pgfsetroundjoin%
\definecolor{currentfill}{rgb}{0.000000,0.000000,0.000000}%
\pgfsetfillcolor{currentfill}%
\pgfsetlinewidth{0.602250pt}%
\definecolor{currentstroke}{rgb}{0.000000,0.000000,0.000000}%
\pgfsetstrokecolor{currentstroke}%
\pgfsetdash{}{0pt}%
\pgfsys@defobject{currentmarker}{\pgfqpoint{0.000000in}{-0.027778in}}{\pgfqpoint{0.000000in}{0.000000in}}{%
\pgfpathmoveto{\pgfqpoint{0.000000in}{0.000000in}}%
\pgfpathlineto{\pgfqpoint{0.000000in}{-0.027778in}}%
\pgfusepath{stroke,fill}%
}%
\begin{pgfscope}%
\pgfsys@transformshift{2.929354in}{0.992139in}%
\pgfsys@useobject{currentmarker}{}%
\end{pgfscope}%
\end{pgfscope}%
\begin{pgfscope}%
\pgfsetbuttcap%
\pgfsetroundjoin%
\definecolor{currentfill}{rgb}{0.000000,0.000000,0.000000}%
\pgfsetfillcolor{currentfill}%
\pgfsetlinewidth{0.602250pt}%
\definecolor{currentstroke}{rgb}{0.000000,0.000000,0.000000}%
\pgfsetstrokecolor{currentstroke}%
\pgfsetdash{}{0pt}%
\pgfsys@defobject{currentmarker}{\pgfqpoint{0.000000in}{-0.027778in}}{\pgfqpoint{0.000000in}{0.000000in}}{%
\pgfpathmoveto{\pgfqpoint{0.000000in}{0.000000in}}%
\pgfpathlineto{\pgfqpoint{0.000000in}{-0.027778in}}%
\pgfusepath{stroke,fill}%
}%
\begin{pgfscope}%
\pgfsys@transformshift{3.000965in}{0.992139in}%
\pgfsys@useobject{currentmarker}{}%
\end{pgfscope}%
\end{pgfscope}%
\begin{pgfscope}%
\pgfsetbuttcap%
\pgfsetroundjoin%
\definecolor{currentfill}{rgb}{0.000000,0.000000,0.000000}%
\pgfsetfillcolor{currentfill}%
\pgfsetlinewidth{0.602250pt}%
\definecolor{currentstroke}{rgb}{0.000000,0.000000,0.000000}%
\pgfsetstrokecolor{currentstroke}%
\pgfsetdash{}{0pt}%
\pgfsys@defobject{currentmarker}{\pgfqpoint{0.000000in}{-0.027778in}}{\pgfqpoint{0.000000in}{0.000000in}}{%
\pgfpathmoveto{\pgfqpoint{0.000000in}{0.000000in}}%
\pgfpathlineto{\pgfqpoint{0.000000in}{-0.027778in}}%
\pgfusepath{stroke,fill}%
}%
\begin{pgfscope}%
\pgfsys@transformshift{3.059474in}{0.992139in}%
\pgfsys@useobject{currentmarker}{}%
\end{pgfscope}%
\end{pgfscope}%
\begin{pgfscope}%
\pgfsetbuttcap%
\pgfsetroundjoin%
\definecolor{currentfill}{rgb}{0.000000,0.000000,0.000000}%
\pgfsetfillcolor{currentfill}%
\pgfsetlinewidth{0.602250pt}%
\definecolor{currentstroke}{rgb}{0.000000,0.000000,0.000000}%
\pgfsetstrokecolor{currentstroke}%
\pgfsetdash{}{0pt}%
\pgfsys@defobject{currentmarker}{\pgfqpoint{0.000000in}{-0.027778in}}{\pgfqpoint{0.000000in}{0.000000in}}{%
\pgfpathmoveto{\pgfqpoint{0.000000in}{0.000000in}}%
\pgfpathlineto{\pgfqpoint{0.000000in}{-0.027778in}}%
\pgfusepath{stroke,fill}%
}%
\begin{pgfscope}%
\pgfsys@transformshift{3.108943in}{0.992139in}%
\pgfsys@useobject{currentmarker}{}%
\end{pgfscope}%
\end{pgfscope}%
\begin{pgfscope}%
\pgfsetbuttcap%
\pgfsetroundjoin%
\definecolor{currentfill}{rgb}{0.000000,0.000000,0.000000}%
\pgfsetfillcolor{currentfill}%
\pgfsetlinewidth{0.602250pt}%
\definecolor{currentstroke}{rgb}{0.000000,0.000000,0.000000}%
\pgfsetstrokecolor{currentstroke}%
\pgfsetdash{}{0pt}%
\pgfsys@defobject{currentmarker}{\pgfqpoint{0.000000in}{-0.027778in}}{\pgfqpoint{0.000000in}{0.000000in}}{%
\pgfpathmoveto{\pgfqpoint{0.000000in}{0.000000in}}%
\pgfpathlineto{\pgfqpoint{0.000000in}{-0.027778in}}%
\pgfusepath{stroke,fill}%
}%
\begin{pgfscope}%
\pgfsys@transformshift{3.151796in}{0.992139in}%
\pgfsys@useobject{currentmarker}{}%
\end{pgfscope}%
\end{pgfscope}%
\begin{pgfscope}%
\pgfsetbuttcap%
\pgfsetroundjoin%
\definecolor{currentfill}{rgb}{0.000000,0.000000,0.000000}%
\pgfsetfillcolor{currentfill}%
\pgfsetlinewidth{0.602250pt}%
\definecolor{currentstroke}{rgb}{0.000000,0.000000,0.000000}%
\pgfsetstrokecolor{currentstroke}%
\pgfsetdash{}{0pt}%
\pgfsys@defobject{currentmarker}{\pgfqpoint{0.000000in}{-0.027778in}}{\pgfqpoint{0.000000in}{0.000000in}}{%
\pgfpathmoveto{\pgfqpoint{0.000000in}{0.000000in}}%
\pgfpathlineto{\pgfqpoint{0.000000in}{-0.027778in}}%
\pgfusepath{stroke,fill}%
}%
\begin{pgfscope}%
\pgfsys@transformshift{3.189594in}{0.992139in}%
\pgfsys@useobject{currentmarker}{}%
\end{pgfscope}%
\end{pgfscope}%
\begin{pgfscope}%
\definecolor{textcolor}{rgb}{0.000000,0.000000,0.000000}%
\pgfsetstrokecolor{textcolor}%
\pgfsetfillcolor{textcolor}%
\pgftext[x=2.796708in,y=0.787417in,,top]{\color{textcolor}{\rmfamily\fontsize{8.000000}{9.600000}\selectfont\catcode`\^=\active\def^{\ifmmode\sp\else\^{}\fi}\catcode`\%=\active\def%{\%}Scale $\alpha$ (Degrees)}}%
\end{pgfscope}%
\begin{pgfscope}%
\pgfpathrectangle{\pgfqpoint{2.208563in}{0.992139in}}{\pgfqpoint{1.176290in}{1.351500in}}%
\pgfusepath{clip}%
\pgfsetrectcap%
\pgfsetroundjoin%
\pgfsetlinewidth{0.451687pt}%
\definecolor{currentstroke}{rgb}{0.690196,0.690196,0.690196}%
\pgfsetstrokecolor{currentstroke}%
\pgfsetstrokeopacity{0.280000}%
\pgfsetdash{}{0pt}%
\pgfpathmoveto{\pgfqpoint{2.208563in}{1.042581in}}%
\pgfpathlineto{\pgfqpoint{3.384853in}{1.042581in}}%
\pgfusepath{stroke}%
\end{pgfscope}%
\begin{pgfscope}%
\pgfsetbuttcap%
\pgfsetroundjoin%
\definecolor{currentfill}{rgb}{0.000000,0.000000,0.000000}%
\pgfsetfillcolor{currentfill}%
\pgfsetlinewidth{0.803000pt}%
\definecolor{currentstroke}{rgb}{0.000000,0.000000,0.000000}%
\pgfsetstrokecolor{currentstroke}%
\pgfsetdash{}{0pt}%
\pgfsys@defobject{currentmarker}{\pgfqpoint{-0.048611in}{0.000000in}}{\pgfqpoint{-0.000000in}{0.000000in}}{%
\pgfpathmoveto{\pgfqpoint{-0.000000in}{0.000000in}}%
\pgfpathlineto{\pgfqpoint{-0.048611in}{0.000000in}}%
\pgfusepath{stroke,fill}%
}%
\begin{pgfscope}%
\pgfsys@transformshift{2.208563in}{1.042581in}%
\pgfsys@useobject{currentmarker}{}%
\end{pgfscope}%
\end{pgfscope}%
\begin{pgfscope}%
\definecolor{textcolor}{rgb}{0.000000,0.000000,0.000000}%
\pgfsetstrokecolor{textcolor}%
\pgfsetfillcolor{textcolor}%
\pgftext[x=1.969172in, y=1.008823in, left, base]{\color{textcolor}{\rmfamily\fontsize{7.000000}{8.400000}\selectfont\catcode`\^=\active\def^{\ifmmode\sp\else\^{}\fi}\catcode`\%=\active\def%{\%}0.0}}%
\end{pgfscope}%
\begin{pgfscope}%
\pgfpathrectangle{\pgfqpoint{2.208563in}{0.992139in}}{\pgfqpoint{1.176290in}{1.351500in}}%
\pgfusepath{clip}%
\pgfsetrectcap%
\pgfsetroundjoin%
\pgfsetlinewidth{0.451687pt}%
\definecolor{currentstroke}{rgb}{0.690196,0.690196,0.690196}%
\pgfsetstrokecolor{currentstroke}%
\pgfsetstrokeopacity{0.280000}%
\pgfsetdash{}{0pt}%
\pgfpathmoveto{\pgfqpoint{2.208563in}{1.290507in}}%
\pgfpathlineto{\pgfqpoint{3.384853in}{1.290507in}}%
\pgfusepath{stroke}%
\end{pgfscope}%
\begin{pgfscope}%
\pgfsetbuttcap%
\pgfsetroundjoin%
\definecolor{currentfill}{rgb}{0.000000,0.000000,0.000000}%
\pgfsetfillcolor{currentfill}%
\pgfsetlinewidth{0.803000pt}%
\definecolor{currentstroke}{rgb}{0.000000,0.000000,0.000000}%
\pgfsetstrokecolor{currentstroke}%
\pgfsetdash{}{0pt}%
\pgfsys@defobject{currentmarker}{\pgfqpoint{-0.048611in}{0.000000in}}{\pgfqpoint{-0.000000in}{0.000000in}}{%
\pgfpathmoveto{\pgfqpoint{-0.000000in}{0.000000in}}%
\pgfpathlineto{\pgfqpoint{-0.048611in}{0.000000in}}%
\pgfusepath{stroke,fill}%
}%
\begin{pgfscope}%
\pgfsys@transformshift{2.208563in}{1.290507in}%
\pgfsys@useobject{currentmarker}{}%
\end{pgfscope}%
\end{pgfscope}%
\begin{pgfscope}%
\definecolor{textcolor}{rgb}{0.000000,0.000000,0.000000}%
\pgfsetstrokecolor{textcolor}%
\pgfsetfillcolor{textcolor}%
\pgftext[x=1.969172in, y=1.256750in, left, base]{\color{textcolor}{\rmfamily\fontsize{7.000000}{8.400000}\selectfont\catcode`\^=\active\def^{\ifmmode\sp\else\^{}\fi}\catcode`\%=\active\def%{\%}0.2}}%
\end{pgfscope}%
\begin{pgfscope}%
\pgfpathrectangle{\pgfqpoint{2.208563in}{0.992139in}}{\pgfqpoint{1.176290in}{1.351500in}}%
\pgfusepath{clip}%
\pgfsetrectcap%
\pgfsetroundjoin%
\pgfsetlinewidth{0.451687pt}%
\definecolor{currentstroke}{rgb}{0.690196,0.690196,0.690196}%
\pgfsetstrokecolor{currentstroke}%
\pgfsetstrokeopacity{0.280000}%
\pgfsetdash{}{0pt}%
\pgfpathmoveto{\pgfqpoint{2.208563in}{1.538433in}}%
\pgfpathlineto{\pgfqpoint{3.384853in}{1.538433in}}%
\pgfusepath{stroke}%
\end{pgfscope}%
\begin{pgfscope}%
\pgfsetbuttcap%
\pgfsetroundjoin%
\definecolor{currentfill}{rgb}{0.000000,0.000000,0.000000}%
\pgfsetfillcolor{currentfill}%
\pgfsetlinewidth{0.803000pt}%
\definecolor{currentstroke}{rgb}{0.000000,0.000000,0.000000}%
\pgfsetstrokecolor{currentstroke}%
\pgfsetdash{}{0pt}%
\pgfsys@defobject{currentmarker}{\pgfqpoint{-0.048611in}{0.000000in}}{\pgfqpoint{-0.000000in}{0.000000in}}{%
\pgfpathmoveto{\pgfqpoint{-0.000000in}{0.000000in}}%
\pgfpathlineto{\pgfqpoint{-0.048611in}{0.000000in}}%
\pgfusepath{stroke,fill}%
}%
\begin{pgfscope}%
\pgfsys@transformshift{2.208563in}{1.538433in}%
\pgfsys@useobject{currentmarker}{}%
\end{pgfscope}%
\end{pgfscope}%
\begin{pgfscope}%
\definecolor{textcolor}{rgb}{0.000000,0.000000,0.000000}%
\pgfsetstrokecolor{textcolor}%
\pgfsetfillcolor{textcolor}%
\pgftext[x=1.969172in, y=1.504676in, left, base]{\color{textcolor}{\rmfamily\fontsize{7.000000}{8.400000}\selectfont\catcode`\^=\active\def^{\ifmmode\sp\else\^{}\fi}\catcode`\%=\active\def%{\%}0.4}}%
\end{pgfscope}%
\begin{pgfscope}%
\pgfpathrectangle{\pgfqpoint{2.208563in}{0.992139in}}{\pgfqpoint{1.176290in}{1.351500in}}%
\pgfusepath{clip}%
\pgfsetrectcap%
\pgfsetroundjoin%
\pgfsetlinewidth{0.451687pt}%
\definecolor{currentstroke}{rgb}{0.690196,0.690196,0.690196}%
\pgfsetstrokecolor{currentstroke}%
\pgfsetstrokeopacity{0.280000}%
\pgfsetdash{}{0pt}%
\pgfpathmoveto{\pgfqpoint{2.208563in}{1.786360in}}%
\pgfpathlineto{\pgfqpoint{3.384853in}{1.786360in}}%
\pgfusepath{stroke}%
\end{pgfscope}%
\begin{pgfscope}%
\pgfsetbuttcap%
\pgfsetroundjoin%
\definecolor{currentfill}{rgb}{0.000000,0.000000,0.000000}%
\pgfsetfillcolor{currentfill}%
\pgfsetlinewidth{0.803000pt}%
\definecolor{currentstroke}{rgb}{0.000000,0.000000,0.000000}%
\pgfsetstrokecolor{currentstroke}%
\pgfsetdash{}{0pt}%
\pgfsys@defobject{currentmarker}{\pgfqpoint{-0.048611in}{0.000000in}}{\pgfqpoint{-0.000000in}{0.000000in}}{%
\pgfpathmoveto{\pgfqpoint{-0.000000in}{0.000000in}}%
\pgfpathlineto{\pgfqpoint{-0.048611in}{0.000000in}}%
\pgfusepath{stroke,fill}%
}%
\begin{pgfscope}%
\pgfsys@transformshift{2.208563in}{1.786360in}%
\pgfsys@useobject{currentmarker}{}%
\end{pgfscope}%
\end{pgfscope}%
\begin{pgfscope}%
\definecolor{textcolor}{rgb}{0.000000,0.000000,0.000000}%
\pgfsetstrokecolor{textcolor}%
\pgfsetfillcolor{textcolor}%
\pgftext[x=1.969172in, y=1.752602in, left, base]{\color{textcolor}{\rmfamily\fontsize{7.000000}{8.400000}\selectfont\catcode`\^=\active\def^{\ifmmode\sp\else\^{}\fi}\catcode`\%=\active\def%{\%}0.6}}%
\end{pgfscope}%
\begin{pgfscope}%
\pgfpathrectangle{\pgfqpoint{2.208563in}{0.992139in}}{\pgfqpoint{1.176290in}{1.351500in}}%
\pgfusepath{clip}%
\pgfsetrectcap%
\pgfsetroundjoin%
\pgfsetlinewidth{0.451687pt}%
\definecolor{currentstroke}{rgb}{0.690196,0.690196,0.690196}%
\pgfsetstrokecolor{currentstroke}%
\pgfsetstrokeopacity{0.280000}%
\pgfsetdash{}{0pt}%
\pgfpathmoveto{\pgfqpoint{2.208563in}{2.034286in}}%
\pgfpathlineto{\pgfqpoint{3.384853in}{2.034286in}}%
\pgfusepath{stroke}%
\end{pgfscope}%
\begin{pgfscope}%
\pgfsetbuttcap%
\pgfsetroundjoin%
\definecolor{currentfill}{rgb}{0.000000,0.000000,0.000000}%
\pgfsetfillcolor{currentfill}%
\pgfsetlinewidth{0.803000pt}%
\definecolor{currentstroke}{rgb}{0.000000,0.000000,0.000000}%
\pgfsetstrokecolor{currentstroke}%
\pgfsetdash{}{0pt}%
\pgfsys@defobject{currentmarker}{\pgfqpoint{-0.048611in}{0.000000in}}{\pgfqpoint{-0.000000in}{0.000000in}}{%
\pgfpathmoveto{\pgfqpoint{-0.000000in}{0.000000in}}%
\pgfpathlineto{\pgfqpoint{-0.048611in}{0.000000in}}%
\pgfusepath{stroke,fill}%
}%
\begin{pgfscope}%
\pgfsys@transformshift{2.208563in}{2.034286in}%
\pgfsys@useobject{currentmarker}{}%
\end{pgfscope}%
\end{pgfscope}%
\begin{pgfscope}%
\definecolor{textcolor}{rgb}{0.000000,0.000000,0.000000}%
\pgfsetstrokecolor{textcolor}%
\pgfsetfillcolor{textcolor}%
\pgftext[x=1.969172in, y=2.000528in, left, base]{\color{textcolor}{\rmfamily\fontsize{7.000000}{8.400000}\selectfont\catcode`\^=\active\def^{\ifmmode\sp\else\^{}\fi}\catcode`\%=\active\def%{\%}0.8}}%
\end{pgfscope}%
\begin{pgfscope}%
\pgfpathrectangle{\pgfqpoint{2.208563in}{0.992139in}}{\pgfqpoint{1.176290in}{1.351500in}}%
\pgfusepath{clip}%
\pgfsetrectcap%
\pgfsetroundjoin%
\pgfsetlinewidth{0.451687pt}%
\definecolor{currentstroke}{rgb}{0.690196,0.690196,0.690196}%
\pgfsetstrokecolor{currentstroke}%
\pgfsetstrokeopacity{0.280000}%
\pgfsetdash{}{0pt}%
\pgfpathmoveto{\pgfqpoint{2.208563in}{2.282212in}}%
\pgfpathlineto{\pgfqpoint{3.384853in}{2.282212in}}%
\pgfusepath{stroke}%
\end{pgfscope}%
\begin{pgfscope}%
\pgfsetbuttcap%
\pgfsetroundjoin%
\definecolor{currentfill}{rgb}{0.000000,0.000000,0.000000}%
\pgfsetfillcolor{currentfill}%
\pgfsetlinewidth{0.803000pt}%
\definecolor{currentstroke}{rgb}{0.000000,0.000000,0.000000}%
\pgfsetstrokecolor{currentstroke}%
\pgfsetdash{}{0pt}%
\pgfsys@defobject{currentmarker}{\pgfqpoint{-0.048611in}{0.000000in}}{\pgfqpoint{-0.000000in}{0.000000in}}{%
\pgfpathmoveto{\pgfqpoint{-0.000000in}{0.000000in}}%
\pgfpathlineto{\pgfqpoint{-0.048611in}{0.000000in}}%
\pgfusepath{stroke,fill}%
}%
\begin{pgfscope}%
\pgfsys@transformshift{2.208563in}{2.282212in}%
\pgfsys@useobject{currentmarker}{}%
\end{pgfscope}%
\end{pgfscope}%
\begin{pgfscope}%
\definecolor{textcolor}{rgb}{0.000000,0.000000,0.000000}%
\pgfsetstrokecolor{textcolor}%
\pgfsetfillcolor{textcolor}%
\pgftext[x=1.969172in, y=2.248454in, left, base]{\color{textcolor}{\rmfamily\fontsize{7.000000}{8.400000}\selectfont\catcode`\^=\active\def^{\ifmmode\sp\else\^{}\fi}\catcode`\%=\active\def%{\%}1.0}}%
\end{pgfscope}%
\begin{pgfscope}%
\definecolor{textcolor}{rgb}{0.000000,0.000000,0.000000}%
\pgfsetstrokecolor{textcolor}%
\pgfsetfillcolor{textcolor}%
\pgftext[x=1.955283in,y=1.667889in,,bottom,rotate=90.000000]{\color{textcolor}{\rmfamily\fontsize{8.000000}{9.600000}\selectfont\catcode`\^=\active\def^{\ifmmode\sp\else\^{}\fi}\catcode`\%=\active\def%{\%}Eff. Volume $\widehat U_2(t)$}}%
\end{pgfscope}%
\begin{pgfscope}%
\pgfpathrectangle{\pgfqpoint{2.208563in}{0.992139in}}{\pgfqpoint{1.176290in}{1.351500in}}%
\pgfusepath{clip}%
\pgfsetrectcap%
\pgfsetroundjoin%
\pgfsetlinewidth{1.355062pt}%
\definecolor{currentstroke}{rgb}{0.121569,0.466667,0.705882}%
\pgfsetstrokecolor{currentstroke}%
\pgfsetdash{}{0pt}%
\pgfpathmoveto{\pgfqpoint{2.262031in}{1.053571in}}%
\pgfpathlineto{\pgfqpoint{2.275567in}{1.053785in}}%
\pgfpathlineto{\pgfqpoint{2.289103in}{1.054018in}}%
\pgfpathlineto{\pgfqpoint{2.302639in}{1.054270in}}%
\pgfpathlineto{\pgfqpoint{2.316175in}{1.054543in}}%
\pgfpathlineto{\pgfqpoint{2.329711in}{1.054840in}}%
\pgfpathlineto{\pgfqpoint{2.343247in}{1.055161in}}%
\pgfpathlineto{\pgfqpoint{2.356783in}{1.055510in}}%
\pgfpathlineto{\pgfqpoint{2.370320in}{1.055889in}}%
\pgfpathlineto{\pgfqpoint{2.383856in}{1.056300in}}%
\pgfpathlineto{\pgfqpoint{2.397392in}{1.056746in}}%
\pgfpathlineto{\pgfqpoint{2.410928in}{1.057231in}}%
\pgfpathlineto{\pgfqpoint{2.424464in}{1.057758in}}%
\pgfpathlineto{\pgfqpoint{2.438000in}{1.058330in}}%
\pgfpathlineto{\pgfqpoint{2.451536in}{1.058951in}}%
\pgfpathlineto{\pgfqpoint{2.465072in}{1.059627in}}%
\pgfpathlineto{\pgfqpoint{2.478609in}{1.060360in}}%
\pgfpathlineto{\pgfqpoint{2.492145in}{1.061158in}}%
\pgfpathlineto{\pgfqpoint{2.505681in}{1.062024in}}%
\pgfpathlineto{\pgfqpoint{2.519217in}{1.062966in}}%
\pgfpathlineto{\pgfqpoint{2.532753in}{1.063989in}}%
\pgfpathlineto{\pgfqpoint{2.546289in}{1.065102in}}%
\pgfpathlineto{\pgfqpoint{2.559825in}{1.066311in}}%
\pgfpathlineto{\pgfqpoint{2.573362in}{1.067625in}}%
\pgfpathlineto{\pgfqpoint{2.586898in}{1.069054in}}%
\pgfpathlineto{\pgfqpoint{2.600434in}{1.070606in}}%
\pgfpathlineto{\pgfqpoint{2.613970in}{1.072293in}}%
\pgfpathlineto{\pgfqpoint{2.627506in}{1.074127in}}%
\pgfpathlineto{\pgfqpoint{2.641042in}{1.076119in}}%
\pgfpathlineto{\pgfqpoint{2.654578in}{1.078284in}}%
\pgfpathlineto{\pgfqpoint{2.668114in}{1.080636in}}%
\pgfpathlineto{\pgfqpoint{2.681651in}{1.083192in}}%
\pgfpathlineto{\pgfqpoint{2.695187in}{1.085969in}}%
\pgfpathlineto{\pgfqpoint{2.708723in}{1.088984in}}%
\pgfpathlineto{\pgfqpoint{2.722259in}{1.092259in}}%
\pgfpathlineto{\pgfqpoint{2.735795in}{1.095816in}}%
\pgfpathlineto{\pgfqpoint{2.749331in}{1.099677in}}%
\pgfpathlineto{\pgfqpoint{2.762867in}{1.103869in}}%
\pgfpathlineto{\pgfqpoint{2.776403in}{1.108418in}}%
\pgfpathlineto{\pgfqpoint{2.789940in}{1.113354in}}%
\pgfpathlineto{\pgfqpoint{2.803476in}{1.118709in}}%
\pgfpathlineto{\pgfqpoint{2.817012in}{1.124516in}}%
\pgfpathlineto{\pgfqpoint{2.830548in}{1.130813in}}%
\pgfpathlineto{\pgfqpoint{2.844084in}{1.137638in}}%
\pgfpathlineto{\pgfqpoint{2.857620in}{1.145033in}}%
\pgfpathlineto{\pgfqpoint{2.871156in}{1.153043in}}%
\pgfpathlineto{\pgfqpoint{2.884693in}{1.161715in}}%
\pgfpathlineto{\pgfqpoint{2.898229in}{1.171100in}}%
\pgfpathlineto{\pgfqpoint{2.911765in}{1.181253in}}%
\pgfpathlineto{\pgfqpoint{2.925301in}{1.192228in}}%
\pgfpathlineto{\pgfqpoint{2.938837in}{1.204088in}}%
\pgfpathlineto{\pgfqpoint{2.952373in}{1.216893in}}%
\pgfpathlineto{\pgfqpoint{2.965909in}{1.230712in}}%
\pgfpathlineto{\pgfqpoint{2.979445in}{1.245612in}}%
\pgfpathlineto{\pgfqpoint{2.992982in}{1.261664in}}%
\pgfpathlineto{\pgfqpoint{3.006518in}{1.278942in}}%
\pgfpathlineto{\pgfqpoint{3.020054in}{1.297521in}}%
\pgfpathlineto{\pgfqpoint{3.033590in}{1.317476in}}%
\pgfpathlineto{\pgfqpoint{3.047126in}{1.338883in}}%
\pgfpathlineto{\pgfqpoint{3.060662in}{1.361818in}}%
\pgfpathlineto{\pgfqpoint{3.074198in}{1.386353in}}%
\pgfpathlineto{\pgfqpoint{3.087734in}{1.412555in}}%
\pgfpathlineto{\pgfqpoint{3.101271in}{1.440487in}}%
\pgfpathlineto{\pgfqpoint{3.114807in}{1.470204in}}%
\pgfpathlineto{\pgfqpoint{3.128343in}{1.501747in}}%
\pgfpathlineto{\pgfqpoint{3.141879in}{1.535144in}}%
\pgfpathlineto{\pgfqpoint{3.155415in}{1.570402in}}%
\pgfpathlineto{\pgfqpoint{3.168951in}{1.607506in}}%
\pgfpathlineto{\pgfqpoint{3.182487in}{1.646409in}}%
\pgfpathlineto{\pgfqpoint{3.196023in}{1.687026in}}%
\pgfpathlineto{\pgfqpoint{3.209560in}{1.729225in}}%
\pgfpathlineto{\pgfqpoint{3.223096in}{1.772814in}}%
\pgfpathlineto{\pgfqpoint{3.236632in}{1.817531in}}%
\pgfpathlineto{\pgfqpoint{3.250168in}{1.863026in}}%
\pgfpathlineto{\pgfqpoint{3.263704in}{1.908852in}}%
\pgfpathlineto{\pgfqpoint{3.277240in}{1.954451in}}%
\pgfpathlineto{\pgfqpoint{3.290776in}{1.999167in}}%
\pgfpathlineto{\pgfqpoint{3.304313in}{2.042258in}}%
\pgfpathlineto{\pgfqpoint{3.317849in}{2.082944in}}%
\pgfpathlineto{\pgfqpoint{3.331385in}{2.120467in}}%
\pgfusepath{stroke}%
\end{pgfscope}%
\begin{pgfscope}%
\pgfpathrectangle{\pgfqpoint{2.208563in}{0.992139in}}{\pgfqpoint{1.176290in}{1.351500in}}%
\pgfusepath{clip}%
\pgfsetrectcap%
\pgfsetroundjoin%
\pgfsetlinewidth{1.355062pt}%
\definecolor{currentstroke}{rgb}{1.000000,0.498039,0.054902}%
\pgfsetstrokecolor{currentstroke}%
\pgfsetdash{}{0pt}%
\pgfpathmoveto{\pgfqpoint{2.262031in}{1.065757in}}%
\pgfpathlineto{\pgfqpoint{2.275567in}{1.066210in}}%
\pgfpathlineto{\pgfqpoint{2.289103in}{1.066695in}}%
\pgfpathlineto{\pgfqpoint{2.302639in}{1.067215in}}%
\pgfpathlineto{\pgfqpoint{2.316175in}{1.067774in}}%
\pgfpathlineto{\pgfqpoint{2.329711in}{1.068376in}}%
\pgfpathlineto{\pgfqpoint{2.343247in}{1.069026in}}%
\pgfpathlineto{\pgfqpoint{2.356783in}{1.069727in}}%
\pgfpathlineto{\pgfqpoint{2.370320in}{1.070486in}}%
\pgfpathlineto{\pgfqpoint{2.383856in}{1.071307in}}%
\pgfpathlineto{\pgfqpoint{2.397392in}{1.072197in}}%
\pgfpathlineto{\pgfqpoint{2.410928in}{1.073163in}}%
\pgfpathlineto{\pgfqpoint{2.424464in}{1.074210in}}%
\pgfpathlineto{\pgfqpoint{2.438000in}{1.075347in}}%
\pgfpathlineto{\pgfqpoint{2.451536in}{1.076581in}}%
\pgfpathlineto{\pgfqpoint{2.465072in}{1.077923in}}%
\pgfpathlineto{\pgfqpoint{2.478609in}{1.079381in}}%
\pgfpathlineto{\pgfqpoint{2.492145in}{1.080966in}}%
\pgfpathlineto{\pgfqpoint{2.505681in}{1.082689in}}%
\pgfpathlineto{\pgfqpoint{2.519217in}{1.084562in}}%
\pgfpathlineto{\pgfqpoint{2.532753in}{1.086599in}}%
\pgfpathlineto{\pgfqpoint{2.546289in}{1.088814in}}%
\pgfpathlineto{\pgfqpoint{2.559825in}{1.091222in}}%
\pgfpathlineto{\pgfqpoint{2.573362in}{1.093840in}}%
\pgfpathlineto{\pgfqpoint{2.586898in}{1.096687in}}%
\pgfpathlineto{\pgfqpoint{2.600434in}{1.099783in}}%
\pgfpathlineto{\pgfqpoint{2.613970in}{1.103148in}}%
\pgfpathlineto{\pgfqpoint{2.627506in}{1.106806in}}%
\pgfpathlineto{\pgfqpoint{2.641042in}{1.110782in}}%
\pgfpathlineto{\pgfqpoint{2.654578in}{1.115104in}}%
\pgfpathlineto{\pgfqpoint{2.668114in}{1.119801in}}%
\pgfpathlineto{\pgfqpoint{2.681651in}{1.124905in}}%
\pgfpathlineto{\pgfqpoint{2.695187in}{1.130450in}}%
\pgfpathlineto{\pgfqpoint{2.708723in}{1.136474in}}%
\pgfpathlineto{\pgfqpoint{2.722259in}{1.143017in}}%
\pgfpathlineto{\pgfqpoint{2.735795in}{1.150123in}}%
\pgfpathlineto{\pgfqpoint{2.749331in}{1.157839in}}%
\pgfpathlineto{\pgfqpoint{2.762867in}{1.166216in}}%
\pgfpathlineto{\pgfqpoint{2.776403in}{1.175307in}}%
\pgfpathlineto{\pgfqpoint{2.789940in}{1.185173in}}%
\pgfpathlineto{\pgfqpoint{2.803476in}{1.195876in}}%
\pgfpathlineto{\pgfqpoint{2.817012in}{1.207485in}}%
\pgfpathlineto{\pgfqpoint{2.830548in}{1.220071in}}%
\pgfpathlineto{\pgfqpoint{2.844084in}{1.233715in}}%
\pgfpathlineto{\pgfqpoint{2.857620in}{1.248498in}}%
\pgfpathlineto{\pgfqpoint{2.871156in}{1.264511in}}%
\pgfpathlineto{\pgfqpoint{2.884693in}{1.281849in}}%
\pgfpathlineto{\pgfqpoint{2.898229in}{1.300612in}}%
\pgfpathlineto{\pgfqpoint{2.911765in}{1.320908in}}%
\pgfpathlineto{\pgfqpoint{2.925301in}{1.342852in}}%
\pgfpathlineto{\pgfqpoint{2.938837in}{1.366562in}}%
\pgfpathlineto{\pgfqpoint{2.952373in}{1.392164in}}%
\pgfpathlineto{\pgfqpoint{2.965909in}{1.419791in}}%
\pgfpathlineto{\pgfqpoint{2.979445in}{1.449578in}}%
\pgfpathlineto{\pgfqpoint{2.992982in}{1.481664in}}%
\pgfpathlineto{\pgfqpoint{3.006518in}{1.516187in}}%
\pgfpathlineto{\pgfqpoint{3.020054in}{1.553277in}}%
\pgfpathlineto{\pgfqpoint{3.033590in}{1.593045in}}%
\pgfpathlineto{\pgfqpoint{3.047126in}{1.635562in}}%
\pgfpathlineto{\pgfqpoint{3.060662in}{1.680830in}}%
\pgfpathlineto{\pgfqpoint{3.074198in}{1.728745in}}%
\pgfpathlineto{\pgfqpoint{3.087734in}{1.779049in}}%
\pgfpathlineto{\pgfqpoint{3.101271in}{1.831281in}}%
\pgfpathlineto{\pgfqpoint{3.114807in}{1.884742in}}%
\pgfpathlineto{\pgfqpoint{3.128343in}{1.938477in}}%
\pgfpathlineto{\pgfqpoint{3.141879in}{1.991309in}}%
\pgfpathlineto{\pgfqpoint{3.155415in}{2.041916in}}%
\pgfpathlineto{\pgfqpoint{3.168951in}{2.088969in}}%
\pgfpathlineto{\pgfqpoint{3.182487in}{2.131285in}}%
\pgfpathlineto{\pgfqpoint{3.196023in}{2.167990in}}%
\pgfpathlineto{\pgfqpoint{3.209560in}{2.198625in}}%
\pgfpathlineto{\pgfqpoint{3.223096in}{2.223183in}}%
\pgfpathlineto{\pgfqpoint{3.236632in}{2.242062in}}%
\pgfpathlineto{\pgfqpoint{3.250168in}{2.255962in}}%
\pgfpathlineto{\pgfqpoint{3.263704in}{2.265753in}}%
\pgfpathlineto{\pgfqpoint{3.277240in}{2.272339in}}%
\pgfpathlineto{\pgfqpoint{3.290776in}{2.276563in}}%
\pgfpathlineto{\pgfqpoint{3.304313in}{2.279138in}}%
\pgfpathlineto{\pgfqpoint{3.317849in}{2.280627in}}%
\pgfpathlineto{\pgfqpoint{3.331385in}{2.281440in}}%
\pgfusepath{stroke}%
\end{pgfscope}%
\begin{pgfscope}%
\pgfpathrectangle{\pgfqpoint{2.208563in}{0.992139in}}{\pgfqpoint{1.176290in}{1.351500in}}%
\pgfusepath{clip}%
\pgfsetrectcap%
\pgfsetroundjoin%
\pgfsetlinewidth{1.355062pt}%
\definecolor{currentstroke}{rgb}{0.172549,0.627451,0.172549}%
\pgfsetstrokecolor{currentstroke}%
\pgfsetdash{}{0pt}%
\pgfpathmoveto{\pgfqpoint{2.262031in}{1.149975in}}%
\pgfpathlineto{\pgfqpoint{2.275567in}{1.153680in}}%
\pgfpathlineto{\pgfqpoint{2.289103in}{1.157530in}}%
\pgfpathlineto{\pgfqpoint{2.302639in}{1.161535in}}%
\pgfpathlineto{\pgfqpoint{2.316175in}{1.165706in}}%
\pgfpathlineto{\pgfqpoint{2.329711in}{1.170053in}}%
\pgfpathlineto{\pgfqpoint{2.343247in}{1.174587in}}%
\pgfpathlineto{\pgfqpoint{2.356783in}{1.179320in}}%
\pgfpathlineto{\pgfqpoint{2.370320in}{1.184264in}}%
\pgfpathlineto{\pgfqpoint{2.383856in}{1.189430in}}%
\pgfpathlineto{\pgfqpoint{2.397392in}{1.194830in}}%
\pgfpathlineto{\pgfqpoint{2.410928in}{1.200477in}}%
\pgfpathlineto{\pgfqpoint{2.424464in}{1.206381in}}%
\pgfpathlineto{\pgfqpoint{2.438000in}{1.212557in}}%
\pgfpathlineto{\pgfqpoint{2.451536in}{1.219017in}}%
\pgfpathlineto{\pgfqpoint{2.465072in}{1.225774in}}%
\pgfpathlineto{\pgfqpoint{2.478609in}{1.232841in}}%
\pgfpathlineto{\pgfqpoint{2.492145in}{1.240233in}}%
\pgfpathlineto{\pgfqpoint{2.505681in}{1.247965in}}%
\pgfpathlineto{\pgfqpoint{2.519217in}{1.256051in}}%
\pgfpathlineto{\pgfqpoint{2.532753in}{1.264507in}}%
\pgfpathlineto{\pgfqpoint{2.546289in}{1.273351in}}%
\pgfpathlineto{\pgfqpoint{2.559825in}{1.282600in}}%
\pgfpathlineto{\pgfqpoint{2.573362in}{1.292271in}}%
\pgfpathlineto{\pgfqpoint{2.586898in}{1.302384in}}%
\pgfpathlineto{\pgfqpoint{2.600434in}{1.312957in}}%
\pgfpathlineto{\pgfqpoint{2.613970in}{1.324009in}}%
\pgfpathlineto{\pgfqpoint{2.627506in}{1.335560in}}%
\pgfpathlineto{\pgfqpoint{2.641042in}{1.347631in}}%
\pgfpathlineto{\pgfqpoint{2.654578in}{1.360240in}}%
\pgfpathlineto{\pgfqpoint{2.668114in}{1.373407in}}%
\pgfpathlineto{\pgfqpoint{2.681651in}{1.387152in}}%
\pgfpathlineto{\pgfqpoint{2.695187in}{1.401495in}}%
\pgfpathlineto{\pgfqpoint{2.708723in}{1.416453in}}%
\pgfpathlineto{\pgfqpoint{2.722259in}{1.432045in}}%
\pgfpathlineto{\pgfqpoint{2.735795in}{1.448290in}}%
\pgfpathlineto{\pgfqpoint{2.749331in}{1.465203in}}%
\pgfpathlineto{\pgfqpoint{2.762867in}{1.482800in}}%
\pgfpathlineto{\pgfqpoint{2.776403in}{1.501097in}}%
\pgfpathlineto{\pgfqpoint{2.789940in}{1.520107in}}%
\pgfpathlineto{\pgfqpoint{2.803476in}{1.539841in}}%
\pgfpathlineto{\pgfqpoint{2.817012in}{1.560309in}}%
\pgfpathlineto{\pgfqpoint{2.830548in}{1.581517in}}%
\pgfpathlineto{\pgfqpoint{2.844084in}{1.603472in}}%
\pgfpathlineto{\pgfqpoint{2.857620in}{1.626175in}}%
\pgfpathlineto{\pgfqpoint{2.871156in}{1.649623in}}%
\pgfpathlineto{\pgfqpoint{2.884693in}{1.673809in}}%
\pgfpathlineto{\pgfqpoint{2.898229in}{1.698724in}}%
\pgfpathlineto{\pgfqpoint{2.911765in}{1.724349in}}%
\pgfpathlineto{\pgfqpoint{2.925301in}{1.750660in}}%
\pgfpathlineto{\pgfqpoint{2.938837in}{1.777624in}}%
\pgfpathlineto{\pgfqpoint{2.952373in}{1.805201in}}%
\pgfpathlineto{\pgfqpoint{2.965909in}{1.833336in}}%
\pgfpathlineto{\pgfqpoint{2.979445in}{1.861965in}}%
\pgfpathlineto{\pgfqpoint{2.992982in}{1.891005in}}%
\pgfpathlineto{\pgfqpoint{3.006518in}{1.920356in}}%
\pgfpathlineto{\pgfqpoint{3.020054in}{1.949895in}}%
\pgfpathlineto{\pgfqpoint{3.033590in}{1.979474in}}%
\pgfpathlineto{\pgfqpoint{3.047126in}{2.008912in}}%
\pgfpathlineto{\pgfqpoint{3.060662in}{2.037997in}}%
\pgfpathlineto{\pgfqpoint{3.074198in}{2.066478in}}%
\pgfpathlineto{\pgfqpoint{3.087734in}{2.094073in}}%
\pgfpathlineto{\pgfqpoint{3.101271in}{2.120477in}}%
\pgfpathlineto{\pgfqpoint{3.114807in}{2.145372in}}%
\pgfpathlineto{\pgfqpoint{3.128343in}{2.168450in}}%
\pgfpathlineto{\pgfqpoint{3.141879in}{2.189439in}}%
\pgfpathlineto{\pgfqpoint{3.155415in}{2.208122in}}%
\pgfpathlineto{\pgfqpoint{3.168951in}{2.224368in}}%
\pgfpathlineto{\pgfqpoint{3.182487in}{2.238136in}}%
\pgfpathlineto{\pgfqpoint{3.196023in}{2.249489in}}%
\pgfpathlineto{\pgfqpoint{3.209560in}{2.258579in}}%
\pgfpathlineto{\pgfqpoint{3.223096in}{2.265634in}}%
\pgfpathlineto{\pgfqpoint{3.236632in}{2.270931in}}%
\pgfpathlineto{\pgfqpoint{3.250168in}{2.274772in}}%
\pgfpathlineto{\pgfqpoint{3.263704in}{2.277458in}}%
\pgfpathlineto{\pgfqpoint{3.277240in}{2.279266in}}%
\pgfpathlineto{\pgfqpoint{3.290776in}{2.280437in}}%
\pgfpathlineto{\pgfqpoint{3.304313in}{2.281166in}}%
\pgfpathlineto{\pgfqpoint{3.317849in}{2.281603in}}%
\pgfpathlineto{\pgfqpoint{3.331385in}{2.281857in}}%
\pgfusepath{stroke}%
\end{pgfscope}%
\begin{pgfscope}%
\pgfpathrectangle{\pgfqpoint{2.208563in}{0.992139in}}{\pgfqpoint{1.176290in}{1.351500in}}%
\pgfusepath{clip}%
\pgfsetrectcap%
\pgfsetroundjoin%
\pgfsetlinewidth{1.355062pt}%
\definecolor{currentstroke}{rgb}{0.839216,0.152941,0.156863}%
\pgfsetstrokecolor{currentstroke}%
\pgfsetdash{}{0pt}%
\pgfpathmoveto{\pgfqpoint{2.262031in}{1.083515in}}%
\pgfpathlineto{\pgfqpoint{2.275567in}{1.084434in}}%
\pgfpathlineto{\pgfqpoint{2.289103in}{1.085417in}}%
\pgfpathlineto{\pgfqpoint{2.302639in}{1.086471in}}%
\pgfpathlineto{\pgfqpoint{2.316175in}{1.087604in}}%
\pgfpathlineto{\pgfqpoint{2.329711in}{1.088823in}}%
\pgfpathlineto{\pgfqpoint{2.343247in}{1.090138in}}%
\pgfpathlineto{\pgfqpoint{2.356783in}{1.091558in}}%
\pgfpathlineto{\pgfqpoint{2.370320in}{1.093093in}}%
\pgfpathlineto{\pgfqpoint{2.383856in}{1.094754in}}%
\pgfpathlineto{\pgfqpoint{2.397392in}{1.096553in}}%
\pgfpathlineto{\pgfqpoint{2.410928in}{1.098502in}}%
\pgfpathlineto{\pgfqpoint{2.424464in}{1.100616in}}%
\pgfpathlineto{\pgfqpoint{2.438000in}{1.102909in}}%
\pgfpathlineto{\pgfqpoint{2.451536in}{1.105398in}}%
\pgfpathlineto{\pgfqpoint{2.465072in}{1.108100in}}%
\pgfpathlineto{\pgfqpoint{2.478609in}{1.111035in}}%
\pgfpathlineto{\pgfqpoint{2.492145in}{1.114223in}}%
\pgfpathlineto{\pgfqpoint{2.505681in}{1.117686in}}%
\pgfpathlineto{\pgfqpoint{2.519217in}{1.121449in}}%
\pgfpathlineto{\pgfqpoint{2.532753in}{1.125539in}}%
\pgfpathlineto{\pgfqpoint{2.546289in}{1.129984in}}%
\pgfpathlineto{\pgfqpoint{2.559825in}{1.134815in}}%
\pgfpathlineto{\pgfqpoint{2.573362in}{1.140067in}}%
\pgfpathlineto{\pgfqpoint{2.586898in}{1.145774in}}%
\pgfpathlineto{\pgfqpoint{2.600434in}{1.151978in}}%
\pgfpathlineto{\pgfqpoint{2.613970in}{1.158722in}}%
\pgfpathlineto{\pgfqpoint{2.627506in}{1.166051in}}%
\pgfpathlineto{\pgfqpoint{2.641042in}{1.174016in}}%
\pgfpathlineto{\pgfqpoint{2.654578in}{1.182672in}}%
\pgfpathlineto{\pgfqpoint{2.668114in}{1.192077in}}%
\pgfpathlineto{\pgfqpoint{2.681651in}{1.202297in}}%
\pgfpathlineto{\pgfqpoint{2.695187in}{1.213400in}}%
\pgfpathlineto{\pgfqpoint{2.708723in}{1.225460in}}%
\pgfpathlineto{\pgfqpoint{2.722259in}{1.238559in}}%
\pgfpathlineto{\pgfqpoint{2.735795in}{1.252784in}}%
\pgfpathlineto{\pgfqpoint{2.749331in}{1.268230in}}%
\pgfpathlineto{\pgfqpoint{2.762867in}{1.284997in}}%
\pgfpathlineto{\pgfqpoint{2.776403in}{1.303195in}}%
\pgfpathlineto{\pgfqpoint{2.789940in}{1.322942in}}%
\pgfpathlineto{\pgfqpoint{2.803476in}{1.344365in}}%
\pgfpathlineto{\pgfqpoint{2.817012in}{1.367599in}}%
\pgfpathlineto{\pgfqpoint{2.830548in}{1.392788in}}%
\pgfpathlineto{\pgfqpoint{2.844084in}{1.420086in}}%
\pgfpathlineto{\pgfqpoint{2.857620in}{1.449653in}}%
\pgfpathlineto{\pgfqpoint{2.871156in}{1.481652in}}%
\pgfpathlineto{\pgfqpoint{2.884693in}{1.516241in}}%
\pgfpathlineto{\pgfqpoint{2.898229in}{1.553562in}}%
\pgfpathlineto{\pgfqpoint{2.911765in}{1.593724in}}%
\pgfpathlineto{\pgfqpoint{2.925301in}{1.636775in}}%
\pgfpathlineto{\pgfqpoint{2.938837in}{1.682669in}}%
\pgfpathlineto{\pgfqpoint{2.952373in}{1.731225in}}%
\pgfpathlineto{\pgfqpoint{2.965909in}{1.782090in}}%
\pgfpathlineto{\pgfqpoint{2.979445in}{1.834694in}}%
\pgfpathlineto{\pgfqpoint{2.992982in}{1.888242in}}%
\pgfpathlineto{\pgfqpoint{3.006518in}{1.941719in}}%
\pgfpathlineto{\pgfqpoint{3.020054in}{1.993939in}}%
\pgfpathlineto{\pgfqpoint{3.033590in}{2.043646in}}%
\pgfpathlineto{\pgfqpoint{3.047126in}{2.089629in}}%
\pgfpathlineto{\pgfqpoint{3.060662in}{2.130865in}}%
\pgfpathlineto{\pgfqpoint{3.074198in}{2.166636in}}%
\pgfpathlineto{\pgfqpoint{3.087734in}{2.196599in}}%
\pgfpathlineto{\pgfqpoint{3.101271in}{2.220802in}}%
\pgfpathlineto{\pgfqpoint{3.114807in}{2.239630in}}%
\pgfpathlineto{\pgfqpoint{3.128343in}{2.253722in}}%
\pgfpathlineto{\pgfqpoint{3.141879in}{2.263855in}}%
\pgfpathlineto{\pgfqpoint{3.155415in}{2.270846in}}%
\pgfpathlineto{\pgfqpoint{3.168951in}{2.275464in}}%
\pgfpathlineto{\pgfqpoint{3.182487in}{2.278379in}}%
\pgfpathlineto{\pgfqpoint{3.196023in}{2.280132in}}%
\pgfpathlineto{\pgfqpoint{3.209560in}{2.281133in}}%
\pgfpathlineto{\pgfqpoint{3.223096in}{2.281675in}}%
\pgfpathlineto{\pgfqpoint{3.236632in}{2.281952in}}%
\pgfpathlineto{\pgfqpoint{3.250168in}{2.282087in}}%
\pgfpathlineto{\pgfqpoint{3.263704in}{2.282148in}}%
\pgfpathlineto{\pgfqpoint{3.277240in}{2.282176in}}%
\pgfpathlineto{\pgfqpoint{3.290776in}{2.282189in}}%
\pgfpathlineto{\pgfqpoint{3.304313in}{2.282196in}}%
\pgfpathlineto{\pgfqpoint{3.317849in}{2.282200in}}%
\pgfpathlineto{\pgfqpoint{3.331385in}{2.282203in}}%
\pgfusepath{stroke}%
\end{pgfscope}%
\begin{pgfscope}%
\pgfpathrectangle{\pgfqpoint{2.208563in}{0.992139in}}{\pgfqpoint{1.176290in}{1.351500in}}%
\pgfusepath{clip}%
\pgfsetrectcap%
\pgfsetroundjoin%
\pgfsetlinewidth{1.355062pt}%
\definecolor{currentstroke}{rgb}{0.580392,0.403922,0.741176}%
\pgfsetstrokecolor{currentstroke}%
\pgfsetdash{}{0pt}%
\pgfpathmoveto{\pgfqpoint{2.262031in}{1.113339in}}%
\pgfpathlineto{\pgfqpoint{2.275567in}{1.119562in}}%
\pgfpathlineto{\pgfqpoint{2.289103in}{1.126333in}}%
\pgfpathlineto{\pgfqpoint{2.302639in}{1.133699in}}%
\pgfpathlineto{\pgfqpoint{2.316175in}{1.141712in}}%
\pgfpathlineto{\pgfqpoint{2.329711in}{1.150429in}}%
\pgfpathlineto{\pgfqpoint{2.343247in}{1.159912in}}%
\pgfpathlineto{\pgfqpoint{2.356783in}{1.170228in}}%
\pgfpathlineto{\pgfqpoint{2.370320in}{1.181449in}}%
\pgfpathlineto{\pgfqpoint{2.383856in}{1.193656in}}%
\pgfpathlineto{\pgfqpoint{2.397392in}{1.206934in}}%
\pgfpathlineto{\pgfqpoint{2.410928in}{1.221377in}}%
\pgfpathlineto{\pgfqpoint{2.424464in}{1.237087in}}%
\pgfpathlineto{\pgfqpoint{2.438000in}{1.254175in}}%
\pgfpathlineto{\pgfqpoint{2.451536in}{1.272761in}}%
\pgfpathlineto{\pgfqpoint{2.465072in}{1.292975in}}%
\pgfpathlineto{\pgfqpoint{2.478609in}{1.314960in}}%
\pgfpathlineto{\pgfqpoint{2.492145in}{1.338868in}}%
\pgfpathlineto{\pgfqpoint{2.505681in}{1.364863in}}%
\pgfpathlineto{\pgfqpoint{2.519217in}{1.393120in}}%
\pgfpathlineto{\pgfqpoint{2.532753in}{1.423819in}}%
\pgfpathlineto{\pgfqpoint{2.546289in}{1.457143in}}%
\pgfpathlineto{\pgfqpoint{2.559825in}{1.493261in}}%
\pgfpathlineto{\pgfqpoint{2.573362in}{1.532317in}}%
\pgfpathlineto{\pgfqpoint{2.586898in}{1.574401in}}%
\pgfpathlineto{\pgfqpoint{2.600434in}{1.619518in}}%
\pgfpathlineto{\pgfqpoint{2.613970in}{1.667551in}}%
\pgfpathlineto{\pgfqpoint{2.627506in}{1.718220in}}%
\pgfpathlineto{\pgfqpoint{2.641042in}{1.771045in}}%
\pgfpathlineto{\pgfqpoint{2.654578in}{1.825325in}}%
\pgfpathlineto{\pgfqpoint{2.668114in}{1.880143in}}%
\pgfpathlineto{\pgfqpoint{2.681651in}{1.934400in}}%
\pgfpathlineto{\pgfqpoint{2.695187in}{1.986892in}}%
\pgfpathlineto{\pgfqpoint{2.708723in}{2.036415in}}%
\pgfpathlineto{\pgfqpoint{2.722259in}{2.081886in}}%
\pgfpathlineto{\pgfqpoint{2.735795in}{2.122457in}}%
\pgfpathlineto{\pgfqpoint{2.749331in}{2.157598in}}%
\pgfpathlineto{\pgfqpoint{2.762867in}{2.187132in}}%
\pgfpathlineto{\pgfqpoint{2.776403in}{2.211212in}}%
\pgfpathlineto{\pgfqpoint{2.789940in}{2.230269in}}%
\pgfpathlineto{\pgfqpoint{2.803476in}{2.244915in}}%
\pgfpathlineto{\pgfqpoint{2.817012in}{2.255860in}}%
\pgfpathlineto{\pgfqpoint{2.830548in}{2.263825in}}%
\pgfpathlineto{\pgfqpoint{2.844084in}{2.269482in}}%
\pgfpathlineto{\pgfqpoint{2.857620in}{2.273416in}}%
\pgfpathlineto{\pgfqpoint{2.871156in}{2.276106in}}%
\pgfpathlineto{\pgfqpoint{2.884693in}{2.277926in}}%
\pgfpathlineto{\pgfqpoint{2.898229in}{2.279153in}}%
\pgfpathlineto{\pgfqpoint{2.911765in}{2.279984in}}%
\pgfpathlineto{\pgfqpoint{2.925301in}{2.280554in}}%
\pgfpathlineto{\pgfqpoint{2.938837in}{2.280954in}}%
\pgfpathlineto{\pgfqpoint{2.952373in}{2.281241in}}%
\pgfpathlineto{\pgfqpoint{2.965909in}{2.281452in}}%
\pgfpathlineto{\pgfqpoint{2.979445in}{2.281611in}}%
\pgfpathlineto{\pgfqpoint{2.992982in}{2.281732in}}%
\pgfpathlineto{\pgfqpoint{3.006518in}{2.281827in}}%
\pgfpathlineto{\pgfqpoint{3.020054in}{2.281900in}}%
\pgfpathlineto{\pgfqpoint{3.033590in}{2.281959in}}%
\pgfpathlineto{\pgfqpoint{3.047126in}{2.282005in}}%
\pgfpathlineto{\pgfqpoint{3.060662in}{2.282042in}}%
\pgfpathlineto{\pgfqpoint{3.074198in}{2.282072in}}%
\pgfpathlineto{\pgfqpoint{3.087734in}{2.282096in}}%
\pgfpathlineto{\pgfqpoint{3.101271in}{2.282115in}}%
\pgfpathlineto{\pgfqpoint{3.114807in}{2.282131in}}%
\pgfpathlineto{\pgfqpoint{3.128343in}{2.282144in}}%
\pgfpathlineto{\pgfqpoint{3.141879in}{2.282154in}}%
\pgfpathlineto{\pgfqpoint{3.155415in}{2.282163in}}%
\pgfpathlineto{\pgfqpoint{3.168951in}{2.282170in}}%
\pgfpathlineto{\pgfqpoint{3.182487in}{2.282176in}}%
\pgfpathlineto{\pgfqpoint{3.196023in}{2.282181in}}%
\pgfpathlineto{\pgfqpoint{3.209560in}{2.282185in}}%
\pgfpathlineto{\pgfqpoint{3.223096in}{2.282189in}}%
\pgfpathlineto{\pgfqpoint{3.236632in}{2.282192in}}%
\pgfpathlineto{\pgfqpoint{3.250168in}{2.282195in}}%
\pgfpathlineto{\pgfqpoint{3.263704in}{2.282198in}}%
\pgfpathlineto{\pgfqpoint{3.277240in}{2.282200in}}%
\pgfpathlineto{\pgfqpoint{3.290776in}{2.282202in}}%
\pgfpathlineto{\pgfqpoint{3.304313in}{2.282204in}}%
\pgfpathlineto{\pgfqpoint{3.317849in}{2.282206in}}%
\pgfpathlineto{\pgfqpoint{3.331385in}{2.282207in}}%
\pgfusepath{stroke}%
\end{pgfscope}%
\begin{pgfscope}%
\pgfsetrectcap%
\pgfsetmiterjoin%
\pgfsetlinewidth{0.803000pt}%
\definecolor{currentstroke}{rgb}{0.000000,0.000000,0.000000}%
\pgfsetstrokecolor{currentstroke}%
\pgfsetdash{}{0pt}%
\pgfpathmoveto{\pgfqpoint{2.208563in}{0.992139in}}%
\pgfpathlineto{\pgfqpoint{2.208563in}{2.343639in}}%
\pgfusepath{stroke}%
\end{pgfscope}%
\begin{pgfscope}%
\pgfsetrectcap%
\pgfsetmiterjoin%
\pgfsetlinewidth{0.803000pt}%
\definecolor{currentstroke}{rgb}{0.000000,0.000000,0.000000}%
\pgfsetstrokecolor{currentstroke}%
\pgfsetdash{}{0pt}%
\pgfpathmoveto{\pgfqpoint{3.384853in}{0.992139in}}%
\pgfpathlineto{\pgfqpoint{3.384853in}{2.343639in}}%
\pgfusepath{stroke}%
\end{pgfscope}%
\begin{pgfscope}%
\pgfsetrectcap%
\pgfsetmiterjoin%
\pgfsetlinewidth{0.803000pt}%
\definecolor{currentstroke}{rgb}{0.000000,0.000000,0.000000}%
\pgfsetstrokecolor{currentstroke}%
\pgfsetdash{}{0pt}%
\pgfpathmoveto{\pgfqpoint{2.208563in}{0.992139in}}%
\pgfpathlineto{\pgfqpoint{3.384853in}{0.992139in}}%
\pgfusepath{stroke}%
\end{pgfscope}%
\begin{pgfscope}%
\pgfsetrectcap%
\pgfsetmiterjoin%
\pgfsetlinewidth{0.803000pt}%
\definecolor{currentstroke}{rgb}{0.000000,0.000000,0.000000}%
\pgfsetstrokecolor{currentstroke}%
\pgfsetdash{}{0pt}%
\pgfpathmoveto{\pgfqpoint{2.208563in}{2.343639in}}%
\pgfpathlineto{\pgfqpoint{3.384853in}{2.343639in}}%
\pgfusepath{stroke}%
\end{pgfscope}%
\begin{pgfscope}%
\definecolor{textcolor}{rgb}{0.000000,0.000000,0.000000}%
\pgfsetstrokecolor{textcolor}%
\pgfsetfillcolor{textcolor}%
\pgftext[x=2.796708in,y=2.385306in,,base]{\color{textcolor}{\rmfamily\fontsize{8.000000}{9.600000}\selectfont\catcode`\^=\active\def^{\ifmmode\sp\else\^{}\fi}\catcode`\%=\active\def%{\%}(b) Effective Volume}}%
\end{pgfscope}%
\begin{pgfscope}%
\pgfsetrectcap%
\pgfsetroundjoin%
\pgfsetlinewidth{1.355062pt}%
\definecolor{currentstroke}{rgb}{0.121569,0.466667,0.705882}%
\pgfsetstrokecolor{currentstroke}%
\pgfsetdash{}{0pt}%
\pgfpathmoveto{\pgfqpoint{0.440616in}{0.327392in}}%
\pgfpathlineto{\pgfqpoint{0.518394in}{0.327392in}}%
\pgfpathlineto{\pgfqpoint{0.596172in}{0.327392in}}%
\pgfusepath{stroke}%
\end{pgfscope}%
\begin{pgfscope}%
\definecolor{textcolor}{rgb}{0.000000,0.000000,0.000000}%
\pgfsetstrokecolor{textcolor}%
\pgfsetfillcolor{textcolor}%
\pgftext[x=0.673950in,y=0.293364in,left,base]{\color{textcolor}{\rmfamily\fontsize{7.000000}{8.400000}\selectfont\catcode`\^=\active\def^{\ifmmode\sp\else\^{}\fi}\catcode`\%=\active\def%{\%}Unimodal}}%
\end{pgfscope}%
\begin{pgfscope}%
\pgfsetrectcap%
\pgfsetroundjoin%
\pgfsetlinewidth{1.355062pt}%
\definecolor{currentstroke}{rgb}{1.000000,0.498039,0.054902}%
\pgfsetstrokecolor{currentstroke}%
\pgfsetdash{}{0pt}%
\pgfpathmoveto{\pgfqpoint{0.440616in}{0.191821in}}%
\pgfpathlineto{\pgfqpoint{0.518394in}{0.191821in}}%
\pgfpathlineto{\pgfqpoint{0.596172in}{0.191821in}}%
\pgfusepath{stroke}%
\end{pgfscope}%
\begin{pgfscope}%
\definecolor{textcolor}{rgb}{0.000000,0.000000,0.000000}%
\pgfsetstrokecolor{textcolor}%
\pgfsetfillcolor{textcolor}%
\pgftext[x=0.673950in,y=0.157793in,left,base]{\color{textcolor}{\rmfamily\fontsize{7.000000}{8.400000}\selectfont\catcode`\^=\active\def^{\ifmmode\sp\else\^{}\fi}\catcode`\%=\active\def%{\%}Antipodal}}%
\end{pgfscope}%
\begin{pgfscope}%
\pgfsetrectcap%
\pgfsetroundjoin%
\pgfsetlinewidth{1.355062pt}%
\definecolor{currentstroke}{rgb}{0.172549,0.627451,0.172549}%
\pgfsetstrokecolor{currentstroke}%
\pgfsetdash{}{0pt}%
\pgfpathmoveto{\pgfqpoint{1.247466in}{0.327392in}}%
\pgfpathlineto{\pgfqpoint{1.325244in}{0.327392in}}%
\pgfpathlineto{\pgfqpoint{1.403022in}{0.327392in}}%
\pgfusepath{stroke}%
\end{pgfscope}%
\begin{pgfscope}%
\definecolor{textcolor}{rgb}{0.000000,0.000000,0.000000}%
\pgfsetstrokecolor{textcolor}%
\pgfsetfillcolor{textcolor}%
\pgftext[x=1.480800in,y=0.293364in,left,base]{\color{textcolor}{\rmfamily\fontsize{7.000000}{8.400000}\selectfont\catcode`\^=\active\def^{\ifmmode\sp\else\^{}\fi}\catcode`\%=\active\def%{\%}Girdle}}%
\end{pgfscope}%
\begin{pgfscope}%
\pgfsetrectcap%
\pgfsetroundjoin%
\pgfsetlinewidth{1.355062pt}%
\definecolor{currentstroke}{rgb}{0.839216,0.152941,0.156863}%
\pgfsetstrokecolor{currentstroke}%
\pgfsetdash{}{0pt}%
\pgfpathmoveto{\pgfqpoint{1.247466in}{0.191821in}}%
\pgfpathlineto{\pgfqpoint{1.325244in}{0.191821in}}%
\pgfpathlineto{\pgfqpoint{1.403022in}{0.191821in}}%
\pgfusepath{stroke}%
\end{pgfscope}%
\begin{pgfscope}%
\definecolor{textcolor}{rgb}{0.000000,0.000000,0.000000}%
\pgfsetstrokecolor{textcolor}%
\pgfsetfillcolor{textcolor}%
\pgftext[x=1.480800in,y=0.157793in,left,base]{\color{textcolor}{\rmfamily\fontsize{7.000000}{8.400000}\selectfont\catcode`\^=\active\def^{\ifmmode\sp\else\^{}\fi}\catcode`\%=\active\def%{\%}Four Modes}}%
\end{pgfscope}%
\begin{pgfscope}%
\pgfsetrectcap%
\pgfsetroundjoin%
\pgfsetlinewidth{1.355062pt}%
\definecolor{currentstroke}{rgb}{0.580392,0.403922,0.741176}%
\pgfsetstrokecolor{currentstroke}%
\pgfsetdash{}{0pt}%
\pgfpathmoveto{\pgfqpoint{2.141528in}{0.327392in}}%
\pgfpathlineto{\pgfqpoint{2.219305in}{0.327392in}}%
\pgfpathlineto{\pgfqpoint{2.297083in}{0.327392in}}%
\pgfusepath{stroke}%
\end{pgfscope}%
\begin{pgfscope}%
\definecolor{textcolor}{rgb}{0.000000,0.000000,0.000000}%
\pgfsetstrokecolor{textcolor}%
\pgfsetfillcolor{textcolor}%
\pgftext[x=2.374861in,y=0.293364in,left,base]{\color{textcolor}{\rmfamily\fontsize{7.000000}{8.400000}\selectfont\catcode`\^=\active\def^{\ifmmode\sp\else\^{}\fi}\catcode`\%=\active\def%{\%}Duplicates}}%
\end{pgfscope}%
\end{pgfpicture}%
\makeatother%
\endgroup%

\caption{Profiles for representative weighted measures on $S^2$. The horizontal axis is angular resolution $\alpha=\sqrt{8t}$ in degrees.
Panel (a) shows gESS. Panel (b) shows effective occupied volume.}
\label{fig:profiles}
\end{figure}


%----%----%----%----%----%----%----%----%----%----%----%----%----%
\section{ASYMPTOTIC INTERPRETATION}

The next results describe the profiles at small and large diffusion scales. At small times, each atom resembles a heat ball. Nearby atoms then begin to overlap. At large times, heat flow leaves only low-frequency geometry. These limits provide the profile's scale calibration.

\subsection{Small-time behavior}

Let $V=\vol(\M)$. For simplicity, collapse exact duplicates in $\widehat P_w$ and write
\begin{equation*}
  \widehat P_w=\sum_{a=1}^r\alpha_a\delta_{z_a},
  \quad z_a\ne z_b\ (a\ne b),\; \sum_a\alpha_a=1,
\end{equation*}
and define the local mass quantities
\begin{equation*}
  S_\alpha=\sum_a\alpha_a^2,
  \quad
  C_\alpha=\sum_a\alpha_a^2 \Scal(z_a).
\end{equation*}
The normalized-volume heat kernel is $V$ times the Riemannian-volume kernel. For fixed distinct support points, off-diagonal terms decay faster than every power of $t$. The following fixed-configuration law is therefore determined by the diagonal heat expansion.

\begin{theorem}[Fixed-configuration small-time law]
\label{thm:small_time}
As $t\downarrow0$,
\begin{align}
  \widehat A_w(t)
  &:=\sum_{a,b}\alpha_a\alpha_b k_{2t}(z_a,z_b) \notag\\
  &=V(8\pi t)^{-m/2}
  \left[S_\alpha+\frac{t}{3}C_\alpha+O(t^2)+O(e^{-c/t})\right].
  \label{eq:small_time_expansion}
\end{align}
for some $c>0$ depending on the finite support. Consequently,
\begin{align*}
  \widehat D_2(t)
  &=\frac{m}{2}\log\frac{1}{8\pi t}+\log V+\log S_\alpha
    +\frac{t}{3}\frac{C_\alpha}{S_\alpha} \notag\\
  &\quad+O(t^2)+O(e^{-c/t}),\\
  \widehat U_2(t)
  &=\frac{(8\pi t)^{m/2}}{V}\frac{1}{S_\alpha}
    \left[1-\frac{t}{3}\frac{C_\alpha}{S_\alpha}+O(t^2)+O(e^{-c/t})\right].
\end{align*}
If all labels are distinct, then $\alpha_a=w_a$.
\end{theorem}

Theorem~\ref{thm:small_time} gives the asymptotic effective volume. The next proposition quantifies merging when pairwise distances are comparable to the diffusion radius.

\begin{proposition}[Local two-atom merging law]
\label{prop:merging}
Let $P=a\delta_x+(1-a)\delta_y$ with $0<a<1$ and $r=d_g(x,y)$. Fix $r_0<\operatorname{inj}(\M)$ and $C<\infty$. As $t\downarrow0$, the following expansion is uniform over $x,y$ satisfying $r\le r_0$ and $r^2\le Ct$:
\begin{equation}
  s_t(x,y)=\exp\left\{-\frac{r^2}{8t}\right\}\{1+O(t+r^2)\}.
  \label{eq:local_overlap}
\end{equation}
Hence
\begin{align}
  \gess_P(t)^{-1}
  &=a^2+(1-a)^2 \notag\\
  &\quad+2a(1-a)\exp\{-r^2/(8t)\}\notag\\
  &\qquad\times\{1+O(t+r^2)\}.
  \label{eq:two_atom_gess}
\end{align}
\end{proposition}

Atoms at distance $r$ merge when $r^2$ is of order $8t$. We therefore report the spherical scale as $\alpha=\sqrt{8t}$ radians. At $r=\alpha$, the leading overlap is $e^{-1}$.

\subsection{Large-time spectral behavior}

Let $0=\lambda_0<\lambda_1\le\lambda_2\le\cdots$ be the Laplace-Beltrami eigenvalues. Let $\{\varphi_r\}_{r\ge0}$ be an orthonormal eigenbasis in $L^2(\nu)$, with $\varphi_0\equiv1$. Define
\begin{equation*}
  \widehat a_r=\sum_iw_i\varphi_r(x_i),\qquad r\ge1.
\end{equation*}
The heat-kernel spectral representation gives the next result
\citep{rosenberg_1997_LaplacianRiemannianManifold, grigoryan_2009_HeatKernelAnalysis}.

\begin{theorem}[Spectral expansion]
\label{thm:spectral}
For every $t>0$,
\begin{equation}
  \widehat p_t-1=\sum_{r\ge1}e^{-\lambda_rt}\widehat a_r\varphi_r,
  \label{eq:pt_spectral}
\end{equation}
and
\begin{equation}
  \widehat D_2(t)=\log\left(1+\sum_{r\ge1}e^{-2\lambda_rt}\widehat a_r^2\right).
  \label{eq:D2_spectral}
\end{equation}
Equivalently,
\begin{align}
  \widehat A_w(t)-1
  &=\frac{1}{\widehat U_2(t)}-1 \notag\\
  &=\sum_{r\ge1}e^{-2\lambda_rt}\widehat a_r^2.
  \label{eq:Aminus_spectral}
\end{align}
Consequently, as $t\to\infty$, if $\lambda_+$ is the next distinct eigenvalue after $\lambda_1$, then
\begin{align*}
  \widehat D_2(t)
  &=e^{-2\lambda_1t}\sum_{\lambda_r=\lambda_1}\widehat a_r^2 \notag\\
  &\quad+O\left(e^{-2\lambda_+t}+e^{-4\lambda_1t}\right).
\end{align*}
\end{theorem}

The exact expansion \eqref{eq:Aminus_spectral} has no logarithmic cross terms. Its lowest nonzero eigenmodes control large-scale uncertainty and persist longest under diffusion.

%----%----%----%----%----%----%----%----%----%----%----%----%----%
\section{SPHERICAL SPECIALIZATION: vMF AND BINGHAM STRUCTURE}\label{sec:spheres}


We next specialize the construction to the sphere $S^{p-1}$, where $m=p-1$. Its first Laplace-Beltrami eigenspaces have standard interpretations in directional statistics. Degree one corresponds to a preferred direction. Degree two describes axial or elliptical departures from uniformity. Higher degrees describe multimodality and finer angular structure. Let $\nu$ be normalized surface measure. The eigenvalues are
\begin{equation*}
  \lambda_\ell=\ell(\ell+p-2),\qquad \ell=0,1,2,\ldots.
\end{equation*}
The eigenspaces consist of degree-$\ell$ spherical harmonics. Let $\widehat a_{\ell j}=\sum_iw_iY_{\ell j}(x_i)$ in an orthonormal basis of the degree-$\ell$ eigenspace. This is the double-index form of Theorem~\ref{thm:spectral}. A sum over $r$ with $\lambda_r=\lambda_\ell$ becomes $\sum_j\widehat a_{\ell j}^2$.

Let
\begin{equation*}
  \bar x_w=\sum_iw_ix_i,
  \quad
  Q_w=\sum_iw_ix_ix_i^\top.
\end{equation*}
The degree-1 spectral energy is
\begin{equation}
  B_1:=\sum_j\widehat a_{1j}^2=p\|\bar x_w\|^2,
  \label{eq:B1}
\end{equation}
and the degree-2 spectral energy is
\begin{equation}
  B_2:=\sum_j\widehat a_{2j}^2=\frac{p(p+2)}{2}\left\|Q_w-\frac{I_p}{p}\right\|_F^2.
  \label{eq:B2}
\end{equation}

\begin{corollary}[Spherical power expansion]
\label{cor:sphere}
On $S^{p-1}$,
\begin{align*}
  \widehat A_w(t)-1
  &=\frac{1}{\widehat U_2(t)}-1 \notag\\
  &=\sum_{\ell\ge1}e^{-2\ell(\ell+p-2)t}\sum_j\widehat a_{\ell j}^2 \notag\\
  &=e^{-2(p-1)t}B_1+e^{-4pt}B_2+
  \sum_{\ell\ge3}e^{-2\ell(\ell+p-2)t}B_\ell,
\end{align*}
where $B_\ell=\sum_j\widehat a_{\ell j}^2$.
Moreover,
\begin{equation*}
  \widehat D_2(t)=\log\{1+\widehat A_w(t)-1\}.
\end{equation*}
Thus, when $B_1\ne0$, the logged entropy contains the cross term
\begin{equation*}
  -\frac12e^{-4(p-1)t}B_1^2
\end{equation*}
before the displayed degree-2 term. If $B_1=0$, the $B_2$ term is the leading logged contribution whenever degree two is nonzero.
\end{corollary}

Up to $p$, $B_1$ is the squared weighted mean resultant length used by Rayleigh tests and vMF concentration fitting. The term $B_2$ measures traceless second-moment anisotropy associated with Bingham or Fisher-Bingham structure \citep{kent_1982_FisherBinghamDistributionSphere}. These energies retain neither mean direction nor anisotropy axes. Higher degrees measure multimodality and finer angular structure. Corollary~\ref{cor:sphere} identifies the first two coefficients of $\widehat A_w(t)-1$ as vMF- and Bingham-type energies. The logarithm introduces interactions among lower-degree terms.


%----%----%----%----%----%----%----%----%----%----%----%----%----%
\section{WEIGHTED CONSISTENCY}


So far, the profile has been a deterministic functional of the weighted measure. In applications, weighted atoms may approximate a target distribution. We establish uniform convergence on scale intervals bounded away from zero. This lower cutoff is needed because heat kernels become singular as $t$ approaches zero. Let
\begin{align*}
  A_P(t)&=\iint k_{2t}(x,y)\dd P(x)\dd P(y),  \\
  \widehat A_w(t)&=\sum_{i,j}w_iw_jk_{2t}(x_i,x_j),
\end{align*}
and
\begin{equation*}
    D_{2,P}(t)=\log A_P(t).
\end{equation*}

\begin{theorem}[Stability and weighted convergence]
\label{thm:consistency}
Fix $0<t_0<T<\infty$. There is a constant $C=C(t_0,T,\M)$ for which every pair of probability measures $P,Q$ satisfies
\begin{equation}
  \sup_{t\in[t_0,T]} |A_P(t)-A_Q(t)|
  \le C W_1(P,Q),
  \label{eq:W1_stability}
\end{equation}
where $W_1$ is the geodesic Wasserstein-1 distance. Consequently, if $\widehat P_w\to P$ in $W_1$, then
\begin{equation*}
  \sup_{t\in[t_0,T]} |\widehat D_2(t)-D_{2,P}(t)|\to0.
\end{equation*}
If $X_1,\ldots,X_n\overset{iid}{\sim}P$ and deterministic weights satisfy $S_n:=\sum_iw_i^2\to0$, then
\begin{equation}
  \sup_{t\in[t_0,T]} |\widehat A_w(t)-A_P(t)|
  =O_p\left(\sqrt{S_n\log(1/S_n)}\right),
  \label{eq:weighted_rate}
\end{equation}
and the same rate holds for $\widehat D_2(t)$.

Finally, let $X_i\overset{iid}{\sim}Q$. Suppose $P$ has density ratio $\rho=\dd P/\dd Q$, where $0\le\rho\le M$ and $\E_Q\rho=1$. When the denominator is positive, set $w_i=\rho(X_i)/\sum_j\rho(X_j)$. Use any fixed convention otherwise. Then
\begin{equation*}
  \sup_{t\in[t_0,T]} |\widehat A_w(t)-A_P(t)|=O_p\left(\sqrt{\frac{\log n}{n}}\right),
\end{equation*}
with the same rate for $\widehat D_2(t)$.
\end{theorem}

The deterministic-weight rate is a worst-case bound. If $P=\nu$, the linear Hoeffding projection vanishes. At each fixed $t$, the error then improves to $O_p(S_n)$. The constants deteriorate as $t_0\downarrow0$ because heat-kernel suprema and derivatives scale as powers of $t_0^{-1}$. Fine-scale resolution therefore has weaker statistical stability.

%----%----%----%----%----%----%----%----%----%----%----%----%----%
\section{EXPERIMENTS}\label{sec:experiments}

The theory applies to general manifolds, while the experiments use $S^2$ because its heat kernel is explicit. Spherical data also make first-moment failures easy to visualize. The synthetic configurations are chosen to expose limitations of weight-only ESS and first-moment summaries. Figure~\ref{fig:clouds} shows these configurations.
\begin{figure*}[ht]
\centering
%\includegraphics[width=0.96\textwidth]{hkep_fig_clouds.pdf}
\providecommand{\mathdefault}[1]{#1}
\providecommand{\mathregular}[1]{#1}
%% Creator: Matplotlib, PGF backend
%%
%% To include the figure in your LaTeX document, write
%%   \input{<filename>.pgf}
%%
%% Make sure the required packages are loaded in your preamble
%%   \usepackage{pgf}
%%
%% Also ensure that all the required font packages are loaded; for instance,
%% the lmodern package is sometimes necessary when using math font.
%%   \usepackage{lmodern}
%%
%% Figures using additional raster images can only be included by \input if
%% they are in the same directory as the main LaTeX file. For loading figures
%% from other directories you can use the `import` package
%%   \usepackage{import}
%%
%% and then include the figures with
%%   \import{<path to file>}{<filename>.pgf}
%%
%% Matplotlib used the following preamble
%%   \def\mathdefault#1{#1}
%%   \everymath=\expandafter{\the\everymath\displaystyle}
%%   \IfFileExists{scrextend.sty}{
%%     \usepackage[fontsize=8.000000pt]{scrextend}
%%   }{
%%     \renewcommand{\normalsize}{\fontsize{8.000000}{9.600000}\selectfont}
%%     \normalsize
%%   }
%%   \usepackage{amsmath,amssymb}
%%   \makeatletter\@ifpackageloaded{underscore}{}{\usepackage[strings]{underscore}}\makeatother
%%
\begingroup%
\makeatletter%
\begin{pgfpicture}%
\pgfpathrectangle{\pgfpointorigin}{\pgfqpoint{6.920000in}{3.363234in}}%
\pgfusepath{use as bounding box, clip}%
\begin{pgfscope}%
\pgfsetbuttcap%
\pgfsetmiterjoin%
\definecolor{currentfill}{rgb}{1.000000,1.000000,1.000000}%
\pgfsetfillcolor{currentfill}%
\pgfsetlinewidth{0.000000pt}%
\definecolor{currentstroke}{rgb}{1.000000,1.000000,1.000000}%
\pgfsetstrokecolor{currentstroke}%
\pgfsetdash{}{0pt}%
\pgfpathmoveto{\pgfqpoint{0.000000in}{0.000000in}}%
\pgfpathlineto{\pgfqpoint{6.920000in}{0.000000in}}%
\pgfpathlineto{\pgfqpoint{6.920000in}{3.363234in}}%
\pgfpathlineto{\pgfqpoint{0.000000in}{3.363234in}}%
\pgfpathlineto{\pgfqpoint{0.000000in}{0.000000in}}%
\pgfpathclose%
\pgfusepath{fill}%
\end{pgfscope}%
\begin{pgfscope}%
\pgfsetbuttcap%
\pgfsetmiterjoin%
\definecolor{currentfill}{rgb}{1.000000,1.000000,1.000000}%
\pgfsetfillcolor{currentfill}%
\pgfsetlinewidth{0.000000pt}%
\definecolor{currentstroke}{rgb}{0.000000,0.000000,0.000000}%
\pgfsetstrokecolor{currentstroke}%
\pgfsetstrokeopacity{0.000000}%
\pgfsetdash{}{0pt}%
\pgfpathmoveto{\pgfqpoint{1.192000in}{2.052407in}}%
\pgfpathcurveto{\pgfqpoint{1.481602in}{2.052407in}}{\pgfqpoint{1.759381in}{2.109937in}}{\pgfqpoint{1.964161in}{2.212326in}}%
\pgfpathcurveto{\pgfqpoint{2.168940in}{2.314716in}}{\pgfqpoint{2.284000in}{2.453606in}}{\pgfqpoint{2.284000in}{2.598407in}}%
\pgfpathcurveto{\pgfqpoint{2.284000in}{2.743208in}}{\pgfqpoint{2.168940in}{2.882097in}}{\pgfqpoint{1.964161in}{2.984487in}}%
\pgfpathcurveto{\pgfqpoint{1.759381in}{3.086877in}}{\pgfqpoint{1.481602in}{3.144407in}}{\pgfqpoint{1.192000in}{3.144407in}}%
\pgfpathcurveto{\pgfqpoint{0.902398in}{3.144407in}}{\pgfqpoint{0.624619in}{3.086877in}}{\pgfqpoint{0.419839in}{2.984487in}}%
\pgfpathcurveto{\pgfqpoint{0.215060in}{2.882097in}}{\pgfqpoint{0.100000in}{2.743208in}}{\pgfqpoint{0.100000in}{2.598407in}}%
\pgfpathcurveto{\pgfqpoint{0.100000in}{2.453606in}}{\pgfqpoint{0.215060in}{2.314716in}}{\pgfqpoint{0.419839in}{2.212326in}}%
\pgfpathcurveto{\pgfqpoint{0.624619in}{2.109937in}}{\pgfqpoint{0.902398in}{2.052407in}}{\pgfqpoint{1.192000in}{2.052407in}}%
\pgfpathlineto{\pgfqpoint{1.192000in}{2.052407in}}%
\pgfpathclose%
\pgfusepath{fill}%
\end{pgfscope}%
\begin{pgfscope}%
\pgfpathmoveto{\pgfqpoint{1.192000in}{2.052407in}}%
\pgfpathcurveto{\pgfqpoint{1.481602in}{2.052407in}}{\pgfqpoint{1.759381in}{2.109937in}}{\pgfqpoint{1.964161in}{2.212326in}}%
\pgfpathcurveto{\pgfqpoint{2.168940in}{2.314716in}}{\pgfqpoint{2.284000in}{2.453606in}}{\pgfqpoint{2.284000in}{2.598407in}}%
\pgfpathcurveto{\pgfqpoint{2.284000in}{2.743208in}}{\pgfqpoint{2.168940in}{2.882097in}}{\pgfqpoint{1.964161in}{2.984487in}}%
\pgfpathcurveto{\pgfqpoint{1.759381in}{3.086877in}}{\pgfqpoint{1.481602in}{3.144407in}}{\pgfqpoint{1.192000in}{3.144407in}}%
\pgfpathcurveto{\pgfqpoint{0.902398in}{3.144407in}}{\pgfqpoint{0.624619in}{3.086877in}}{\pgfqpoint{0.419839in}{2.984487in}}%
\pgfpathcurveto{\pgfqpoint{0.215060in}{2.882097in}}{\pgfqpoint{0.100000in}{2.743208in}}{\pgfqpoint{0.100000in}{2.598407in}}%
\pgfpathcurveto{\pgfqpoint{0.100000in}{2.453606in}}{\pgfqpoint{0.215060in}{2.314716in}}{\pgfqpoint{0.419839in}{2.212326in}}%
\pgfpathcurveto{\pgfqpoint{0.624619in}{2.109937in}}{\pgfqpoint{0.902398in}{2.052407in}}{\pgfqpoint{1.192000in}{2.052407in}}%
\pgfpathlineto{\pgfqpoint{1.192000in}{2.052407in}}%
\pgfpathclose%
\pgfusepath{clip}%
\pgfsetbuttcap%
\pgfsetroundjoin%
\definecolor{currentfill}{rgb}{0.121569,0.466667,0.705882}%
\pgfsetfillcolor{currentfill}%
\pgfsetfillopacity{0.580000}%
\pgfsetlinewidth{0.000000pt}%
\definecolor{currentstroke}{rgb}{0.121569,0.466667,0.705882}%
\pgfsetstrokecolor{currentstroke}%
\pgfsetstrokeopacity{0.580000}%
\pgfsetdash{}{0pt}%
\pgfsys@defobject{currentmarker}{\pgfqpoint{-0.014232in}{-0.014232in}}{\pgfqpoint{0.014232in}{0.014232in}}{%
\pgfpathmoveto{\pgfqpoint{0.000000in}{-0.014232in}}%
\pgfpathcurveto{\pgfqpoint{0.003774in}{-0.014232in}}{\pgfqpoint{0.007395in}{-0.012732in}}{\pgfqpoint{0.010063in}{-0.010063in}}%
\pgfpathcurveto{\pgfqpoint{0.012732in}{-0.007395in}}{\pgfqpoint{0.014232in}{-0.003774in}}{\pgfqpoint{0.014232in}{0.000000in}}%
\pgfpathcurveto{\pgfqpoint{0.014232in}{0.003774in}}{\pgfqpoint{0.012732in}{0.007395in}}{\pgfqpoint{0.010063in}{0.010063in}}%
\pgfpathcurveto{\pgfqpoint{0.007395in}{0.012732in}}{\pgfqpoint{0.003774in}{0.014232in}}{\pgfqpoint{0.000000in}{0.014232in}}%
\pgfpathcurveto{\pgfqpoint{-0.003774in}{0.014232in}}{\pgfqpoint{-0.007395in}{0.012732in}}{\pgfqpoint{-0.010063in}{0.010063in}}%
\pgfpathcurveto{\pgfqpoint{-0.012732in}{0.007395in}}{\pgfqpoint{-0.014232in}{0.003774in}}{\pgfqpoint{-0.014232in}{0.000000in}}%
\pgfpathcurveto{\pgfqpoint{-0.014232in}{-0.003774in}}{\pgfqpoint{-0.012732in}{-0.007395in}}{\pgfqpoint{-0.010063in}{-0.010063in}}%
\pgfpathcurveto{\pgfqpoint{-0.007395in}{-0.012732in}}{\pgfqpoint{-0.003774in}{-0.014232in}}{\pgfqpoint{0.000000in}{-0.014232in}}%
\pgfpathlineto{\pgfqpoint{0.000000in}{-0.014232in}}%
\pgfpathclose%
\pgfusepath{fill}%
}%
\begin{pgfscope}%
\pgfsys@transformshift{1.036897in}{3.133603in}%
\pgfsys@useobject{currentmarker}{}%
\end{pgfscope}%
\begin{pgfscope}%
\pgfsys@transformshift{1.241111in}{3.140334in}%
\pgfsys@useobject{currentmarker}{}%
\end{pgfscope}%
\begin{pgfscope}%
\pgfsys@transformshift{1.107117in}{3.137230in}%
\pgfsys@useobject{currentmarker}{}%
\end{pgfscope}%
\begin{pgfscope}%
\pgfsys@transformshift{1.038233in}{3.120850in}%
\pgfsys@useobject{currentmarker}{}%
\end{pgfscope}%
\begin{pgfscope}%
\pgfsys@transformshift{1.027325in}{3.124046in}%
\pgfsys@useobject{currentmarker}{}%
\end{pgfscope}%
\begin{pgfscope}%
\pgfsys@transformshift{1.319407in}{3.139691in}%
\pgfsys@useobject{currentmarker}{}%
\end{pgfscope}%
\begin{pgfscope}%
\pgfsys@transformshift{1.552299in}{3.089718in}%
\pgfsys@useobject{currentmarker}{}%
\end{pgfscope}%
\begin{pgfscope}%
\pgfsys@transformshift{1.328198in}{3.138348in}%
\pgfsys@useobject{currentmarker}{}%
\end{pgfscope}%
\begin{pgfscope}%
\pgfsys@transformshift{1.203334in}{3.137751in}%
\pgfsys@useobject{currentmarker}{}%
\end{pgfscope}%
\begin{pgfscope}%
\pgfsys@transformshift{1.112938in}{3.129440in}%
\pgfsys@useobject{currentmarker}{}%
\end{pgfscope}%
\begin{pgfscope}%
\pgfsys@transformshift{0.922309in}{3.124154in}%
\pgfsys@useobject{currentmarker}{}%
\end{pgfscope}%
\begin{pgfscope}%
\pgfsys@transformshift{1.425567in}{3.123203in}%
\pgfsys@useobject{currentmarker}{}%
\end{pgfscope}%
\begin{pgfscope}%
\pgfsys@transformshift{0.913106in}{3.122161in}%
\pgfsys@useobject{currentmarker}{}%
\end{pgfscope}%
\begin{pgfscope}%
\pgfsys@transformshift{1.047630in}{3.128788in}%
\pgfsys@useobject{currentmarker}{}%
\end{pgfscope}%
\begin{pgfscope}%
\pgfsys@transformshift{1.379173in}{3.130449in}%
\pgfsys@useobject{currentmarker}{}%
\end{pgfscope}%
\begin{pgfscope}%
\pgfsys@transformshift{1.112787in}{3.131739in}%
\pgfsys@useobject{currentmarker}{}%
\end{pgfscope}%
\begin{pgfscope}%
\pgfsys@transformshift{1.184555in}{3.143917in}%
\pgfsys@useobject{currentmarker}{}%
\end{pgfscope}%
\begin{pgfscope}%
\pgfsys@transformshift{1.324771in}{3.137643in}%
\pgfsys@useobject{currentmarker}{}%
\end{pgfscope}%
\begin{pgfscope}%
\pgfsys@transformshift{1.251775in}{3.133531in}%
\pgfsys@useobject{currentmarker}{}%
\end{pgfscope}%
\begin{pgfscope}%
\pgfsys@transformshift{1.162098in}{3.143515in}%
\pgfsys@useobject{currentmarker}{}%
\end{pgfscope}%
\begin{pgfscope}%
\pgfsys@transformshift{1.187482in}{3.120382in}%
\pgfsys@useobject{currentmarker}{}%
\end{pgfscope}%
\begin{pgfscope}%
\pgfsys@transformshift{1.292386in}{3.117403in}%
\pgfsys@useobject{currentmarker}{}%
\end{pgfscope}%
\begin{pgfscope}%
\pgfsys@transformshift{1.066044in}{3.133295in}%
\pgfsys@useobject{currentmarker}{}%
\end{pgfscope}%
\begin{pgfscope}%
\pgfsys@transformshift{1.051104in}{3.105804in}%
\pgfsys@useobject{currentmarker}{}%
\end{pgfscope}%
\begin{pgfscope}%
\pgfsys@transformshift{1.126584in}{3.104096in}%
\pgfsys@useobject{currentmarker}{}%
\end{pgfscope}%
\begin{pgfscope}%
\pgfsys@transformshift{1.251942in}{3.130726in}%
\pgfsys@useobject{currentmarker}{}%
\end{pgfscope}%
\begin{pgfscope}%
\pgfsys@transformshift{1.238701in}{3.129391in}%
\pgfsys@useobject{currentmarker}{}%
\end{pgfscope}%
\begin{pgfscope}%
\pgfsys@transformshift{1.189005in}{3.140905in}%
\pgfsys@useobject{currentmarker}{}%
\end{pgfscope}%
\begin{pgfscope}%
\pgfsys@transformshift{1.242242in}{3.133704in}%
\pgfsys@useobject{currentmarker}{}%
\end{pgfscope}%
\begin{pgfscope}%
\pgfsys@transformshift{1.277962in}{3.130706in}%
\pgfsys@useobject{currentmarker}{}%
\end{pgfscope}%
\begin{pgfscope}%
\pgfsys@transformshift{0.981913in}{3.130241in}%
\pgfsys@useobject{currentmarker}{}%
\end{pgfscope}%
\begin{pgfscope}%
\pgfsys@transformshift{1.469269in}{3.121848in}%
\pgfsys@useobject{currentmarker}{}%
\end{pgfscope}%
\begin{pgfscope}%
\pgfsys@transformshift{0.966555in}{3.095529in}%
\pgfsys@useobject{currentmarker}{}%
\end{pgfscope}%
\begin{pgfscope}%
\pgfsys@transformshift{1.316922in}{3.119229in}%
\pgfsys@useobject{currentmarker}{}%
\end{pgfscope}%
\begin{pgfscope}%
\pgfsys@transformshift{1.157863in}{3.135221in}%
\pgfsys@useobject{currentmarker}{}%
\end{pgfscope}%
\begin{pgfscope}%
\pgfsys@transformshift{1.333205in}{3.119654in}%
\pgfsys@useobject{currentmarker}{}%
\end{pgfscope}%
\begin{pgfscope}%
\pgfsys@transformshift{1.063818in}{3.126472in}%
\pgfsys@useobject{currentmarker}{}%
\end{pgfscope}%
\begin{pgfscope}%
\pgfsys@transformshift{1.257722in}{3.087330in}%
\pgfsys@useobject{currentmarker}{}%
\end{pgfscope}%
\begin{pgfscope}%
\pgfsys@transformshift{1.342790in}{3.138569in}%
\pgfsys@useobject{currentmarker}{}%
\end{pgfscope}%
\begin{pgfscope}%
\pgfsys@transformshift{1.214228in}{3.117045in}%
\pgfsys@useobject{currentmarker}{}%
\end{pgfscope}%
\begin{pgfscope}%
\pgfsys@transformshift{1.383045in}{3.122765in}%
\pgfsys@useobject{currentmarker}{}%
\end{pgfscope}%
\begin{pgfscope}%
\pgfsys@transformshift{1.279705in}{3.139872in}%
\pgfsys@useobject{currentmarker}{}%
\end{pgfscope}%
\begin{pgfscope}%
\pgfsys@transformshift{1.055176in}{3.130590in}%
\pgfsys@useobject{currentmarker}{}%
\end{pgfscope}%
\begin{pgfscope}%
\pgfsys@transformshift{1.331475in}{3.139002in}%
\pgfsys@useobject{currentmarker}{}%
\end{pgfscope}%
\begin{pgfscope}%
\pgfsys@transformshift{1.216140in}{3.133959in}%
\pgfsys@useobject{currentmarker}{}%
\end{pgfscope}%
\begin{pgfscope}%
\pgfsys@transformshift{1.190273in}{3.136413in}%
\pgfsys@useobject{currentmarker}{}%
\end{pgfscope}%
\begin{pgfscope}%
\pgfsys@transformshift{1.109671in}{3.112137in}%
\pgfsys@useobject{currentmarker}{}%
\end{pgfscope}%
\begin{pgfscope}%
\pgfsys@transformshift{1.152377in}{3.131419in}%
\pgfsys@useobject{currentmarker}{}%
\end{pgfscope}%
\begin{pgfscope}%
\pgfsys@transformshift{1.312616in}{3.130535in}%
\pgfsys@useobject{currentmarker}{}%
\end{pgfscope}%
\begin{pgfscope}%
\pgfsys@transformshift{1.305585in}{3.139632in}%
\pgfsys@useobject{currentmarker}{}%
\end{pgfscope}%
\begin{pgfscope}%
\pgfsys@transformshift{1.319058in}{3.126200in}%
\pgfsys@useobject{currentmarker}{}%
\end{pgfscope}%
\begin{pgfscope}%
\pgfsys@transformshift{1.247613in}{3.132940in}%
\pgfsys@useobject{currentmarker}{}%
\end{pgfscope}%
\begin{pgfscope}%
\pgfsys@transformshift{1.217785in}{3.108321in}%
\pgfsys@useobject{currentmarker}{}%
\end{pgfscope}%
\begin{pgfscope}%
\pgfsys@transformshift{0.920696in}{3.127053in}%
\pgfsys@useobject{currentmarker}{}%
\end{pgfscope}%
\begin{pgfscope}%
\pgfsys@transformshift{1.263266in}{3.124886in}%
\pgfsys@useobject{currentmarker}{}%
\end{pgfscope}%
\begin{pgfscope}%
\pgfsys@transformshift{1.312946in}{3.116773in}%
\pgfsys@useobject{currentmarker}{}%
\end{pgfscope}%
\begin{pgfscope}%
\pgfsys@transformshift{1.146280in}{3.138226in}%
\pgfsys@useobject{currentmarker}{}%
\end{pgfscope}%
\begin{pgfscope}%
\pgfsys@transformshift{1.458464in}{3.126792in}%
\pgfsys@useobject{currentmarker}{}%
\end{pgfscope}%
\begin{pgfscope}%
\pgfsys@transformshift{1.220558in}{3.143022in}%
\pgfsys@useobject{currentmarker}{}%
\end{pgfscope}%
\begin{pgfscope}%
\pgfsys@transformshift{1.043725in}{3.132737in}%
\pgfsys@useobject{currentmarker}{}%
\end{pgfscope}%
\begin{pgfscope}%
\pgfsys@transformshift{1.309502in}{3.133110in}%
\pgfsys@useobject{currentmarker}{}%
\end{pgfscope}%
\begin{pgfscope}%
\pgfsys@transformshift{0.990457in}{3.133917in}%
\pgfsys@useobject{currentmarker}{}%
\end{pgfscope}%
\begin{pgfscope}%
\pgfsys@transformshift{1.307134in}{3.134847in}%
\pgfsys@useobject{currentmarker}{}%
\end{pgfscope}%
\begin{pgfscope}%
\pgfsys@transformshift{0.860306in}{3.116811in}%
\pgfsys@useobject{currentmarker}{}%
\end{pgfscope}%
\begin{pgfscope}%
\pgfsys@transformshift{0.998502in}{3.128649in}%
\pgfsys@useobject{currentmarker}{}%
\end{pgfscope}%
\begin{pgfscope}%
\pgfsys@transformshift{0.901686in}{3.121500in}%
\pgfsys@useobject{currentmarker}{}%
\end{pgfscope}%
\begin{pgfscope}%
\pgfsys@transformshift{1.405239in}{3.127516in}%
\pgfsys@useobject{currentmarker}{}%
\end{pgfscope}%
\begin{pgfscope}%
\pgfsys@transformshift{1.038437in}{3.112639in}%
\pgfsys@useobject{currentmarker}{}%
\end{pgfscope}%
\begin{pgfscope}%
\pgfsys@transformshift{1.169373in}{3.142574in}%
\pgfsys@useobject{currentmarker}{}%
\end{pgfscope}%
\begin{pgfscope}%
\pgfsys@transformshift{1.306961in}{3.120368in}%
\pgfsys@useobject{currentmarker}{}%
\end{pgfscope}%
\begin{pgfscope}%
\pgfsys@transformshift{1.247787in}{3.134734in}%
\pgfsys@useobject{currentmarker}{}%
\end{pgfscope}%
\begin{pgfscope}%
\pgfsys@transformshift{1.258865in}{3.124055in}%
\pgfsys@useobject{currentmarker}{}%
\end{pgfscope}%
\begin{pgfscope}%
\pgfsys@transformshift{1.186164in}{3.139705in}%
\pgfsys@useobject{currentmarker}{}%
\end{pgfscope}%
\begin{pgfscope}%
\pgfsys@transformshift{1.167815in}{3.134504in}%
\pgfsys@useobject{currentmarker}{}%
\end{pgfscope}%
\begin{pgfscope}%
\pgfsys@transformshift{1.353735in}{3.115504in}%
\pgfsys@useobject{currentmarker}{}%
\end{pgfscope}%
\begin{pgfscope}%
\pgfsys@transformshift{1.182731in}{3.138949in}%
\pgfsys@useobject{currentmarker}{}%
\end{pgfscope}%
\begin{pgfscope}%
\pgfsys@transformshift{1.236594in}{3.141765in}%
\pgfsys@useobject{currentmarker}{}%
\end{pgfscope}%
\begin{pgfscope}%
\pgfsys@transformshift{1.063656in}{3.140522in}%
\pgfsys@useobject{currentmarker}{}%
\end{pgfscope}%
\begin{pgfscope}%
\pgfsys@transformshift{1.420831in}{3.132155in}%
\pgfsys@useobject{currentmarker}{}%
\end{pgfscope}%
\begin{pgfscope}%
\pgfsys@transformshift{1.133341in}{3.116462in}%
\pgfsys@useobject{currentmarker}{}%
\end{pgfscope}%
\begin{pgfscope}%
\pgfsys@transformshift{1.218567in}{3.119232in}%
\pgfsys@useobject{currentmarker}{}%
\end{pgfscope}%
\begin{pgfscope}%
\pgfsys@transformshift{1.266232in}{3.141226in}%
\pgfsys@useobject{currentmarker}{}%
\end{pgfscope}%
\begin{pgfscope}%
\pgfsys@transformshift{1.004680in}{3.131709in}%
\pgfsys@useobject{currentmarker}{}%
\end{pgfscope}%
\begin{pgfscope}%
\pgfsys@transformshift{1.236065in}{3.118588in}%
\pgfsys@useobject{currentmarker}{}%
\end{pgfscope}%
\begin{pgfscope}%
\pgfsys@transformshift{1.257709in}{3.139972in}%
\pgfsys@useobject{currentmarker}{}%
\end{pgfscope}%
\begin{pgfscope}%
\pgfsys@transformshift{1.389321in}{3.134004in}%
\pgfsys@useobject{currentmarker}{}%
\end{pgfscope}%
\begin{pgfscope}%
\pgfsys@transformshift{1.137185in}{3.132154in}%
\pgfsys@useobject{currentmarker}{}%
\end{pgfscope}%
\begin{pgfscope}%
\pgfsys@transformshift{1.060091in}{3.126691in}%
\pgfsys@useobject{currentmarker}{}%
\end{pgfscope}%
\begin{pgfscope}%
\pgfsys@transformshift{1.100515in}{3.127775in}%
\pgfsys@useobject{currentmarker}{}%
\end{pgfscope}%
\begin{pgfscope}%
\pgfsys@transformshift{1.042429in}{3.121497in}%
\pgfsys@useobject{currentmarker}{}%
\end{pgfscope}%
\begin{pgfscope}%
\pgfsys@transformshift{0.952794in}{3.104621in}%
\pgfsys@useobject{currentmarker}{}%
\end{pgfscope}%
\begin{pgfscope}%
\pgfsys@transformshift{1.149271in}{3.139762in}%
\pgfsys@useobject{currentmarker}{}%
\end{pgfscope}%
\begin{pgfscope}%
\pgfsys@transformshift{1.030207in}{3.129434in}%
\pgfsys@useobject{currentmarker}{}%
\end{pgfscope}%
\begin{pgfscope}%
\pgfsys@transformshift{1.067645in}{3.131588in}%
\pgfsys@useobject{currentmarker}{}%
\end{pgfscope}%
\begin{pgfscope}%
\pgfsys@transformshift{1.366720in}{3.124854in}%
\pgfsys@useobject{currentmarker}{}%
\end{pgfscope}%
\begin{pgfscope}%
\pgfsys@transformshift{1.253526in}{3.136642in}%
\pgfsys@useobject{currentmarker}{}%
\end{pgfscope}%
\begin{pgfscope}%
\pgfsys@transformshift{0.984953in}{3.101317in}%
\pgfsys@useobject{currentmarker}{}%
\end{pgfscope}%
\begin{pgfscope}%
\pgfsys@transformshift{1.283850in}{3.126558in}%
\pgfsys@useobject{currentmarker}{}%
\end{pgfscope}%
\begin{pgfscope}%
\pgfsys@transformshift{1.124642in}{3.102793in}%
\pgfsys@useobject{currentmarker}{}%
\end{pgfscope}%
\begin{pgfscope}%
\pgfsys@transformshift{0.906113in}{3.114855in}%
\pgfsys@useobject{currentmarker}{}%
\end{pgfscope}%
\begin{pgfscope}%
\pgfsys@transformshift{1.251314in}{3.142548in}%
\pgfsys@useobject{currentmarker}{}%
\end{pgfscope}%
\begin{pgfscope}%
\pgfsys@transformshift{1.165625in}{3.134396in}%
\pgfsys@useobject{currentmarker}{}%
\end{pgfscope}%
\begin{pgfscope}%
\pgfsys@transformshift{1.273701in}{3.128294in}%
\pgfsys@useobject{currentmarker}{}%
\end{pgfscope}%
\begin{pgfscope}%
\pgfsys@transformshift{0.958126in}{3.130962in}%
\pgfsys@useobject{currentmarker}{}%
\end{pgfscope}%
\begin{pgfscope}%
\pgfsys@transformshift{1.218279in}{3.139671in}%
\pgfsys@useobject{currentmarker}{}%
\end{pgfscope}%
\begin{pgfscope}%
\pgfsys@transformshift{1.172367in}{3.125626in}%
\pgfsys@useobject{currentmarker}{}%
\end{pgfscope}%
\begin{pgfscope}%
\pgfsys@transformshift{1.001347in}{3.132744in}%
\pgfsys@useobject{currentmarker}{}%
\end{pgfscope}%
\begin{pgfscope}%
\pgfsys@transformshift{1.112835in}{3.135021in}%
\pgfsys@useobject{currentmarker}{}%
\end{pgfscope}%
\begin{pgfscope}%
\pgfsys@transformshift{1.350502in}{3.126006in}%
\pgfsys@useobject{currentmarker}{}%
\end{pgfscope}%
\begin{pgfscope}%
\pgfsys@transformshift{1.096097in}{3.130838in}%
\pgfsys@useobject{currentmarker}{}%
\end{pgfscope}%
\begin{pgfscope}%
\pgfsys@transformshift{1.085305in}{3.136978in}%
\pgfsys@useobject{currentmarker}{}%
\end{pgfscope}%
\begin{pgfscope}%
\pgfsys@transformshift{1.101139in}{3.140672in}%
\pgfsys@useobject{currentmarker}{}%
\end{pgfscope}%
\begin{pgfscope}%
\pgfsys@transformshift{1.158119in}{3.116828in}%
\pgfsys@useobject{currentmarker}{}%
\end{pgfscope}%
\begin{pgfscope}%
\pgfsys@transformshift{1.291284in}{3.141396in}%
\pgfsys@useobject{currentmarker}{}%
\end{pgfscope}%
\begin{pgfscope}%
\pgfsys@transformshift{1.587860in}{3.089602in}%
\pgfsys@useobject{currentmarker}{}%
\end{pgfscope}%
\begin{pgfscope}%
\pgfsys@transformshift{1.079293in}{3.136682in}%
\pgfsys@useobject{currentmarker}{}%
\end{pgfscope}%
\begin{pgfscope}%
\pgfsys@transformshift{1.137328in}{3.138082in}%
\pgfsys@useobject{currentmarker}{}%
\end{pgfscope}%
\begin{pgfscope}%
\pgfsys@transformshift{1.425491in}{3.115870in}%
\pgfsys@useobject{currentmarker}{}%
\end{pgfscope}%
\begin{pgfscope}%
\pgfsys@transformshift{1.016549in}{3.128016in}%
\pgfsys@useobject{currentmarker}{}%
\end{pgfscope}%
\begin{pgfscope}%
\pgfsys@transformshift{1.260645in}{3.138199in}%
\pgfsys@useobject{currentmarker}{}%
\end{pgfscope}%
\begin{pgfscope}%
\pgfsys@transformshift{0.792757in}{3.096951in}%
\pgfsys@useobject{currentmarker}{}%
\end{pgfscope}%
\begin{pgfscope}%
\pgfsys@transformshift{1.092586in}{3.133685in}%
\pgfsys@useobject{currentmarker}{}%
\end{pgfscope}%
\begin{pgfscope}%
\pgfsys@transformshift{1.214369in}{3.137652in}%
\pgfsys@useobject{currentmarker}{}%
\end{pgfscope}%
\begin{pgfscope}%
\pgfsys@transformshift{1.059537in}{3.130676in}%
\pgfsys@useobject{currentmarker}{}%
\end{pgfscope}%
\begin{pgfscope}%
\pgfsys@transformshift{1.197881in}{3.136031in}%
\pgfsys@useobject{currentmarker}{}%
\end{pgfscope}%
\begin{pgfscope}%
\pgfsys@transformshift{1.165600in}{3.120907in}%
\pgfsys@useobject{currentmarker}{}%
\end{pgfscope}%
\begin{pgfscope}%
\pgfsys@transformshift{1.497338in}{3.119548in}%
\pgfsys@useobject{currentmarker}{}%
\end{pgfscope}%
\begin{pgfscope}%
\pgfsys@transformshift{1.420407in}{3.126262in}%
\pgfsys@useobject{currentmarker}{}%
\end{pgfscope}%
\begin{pgfscope}%
\pgfsys@transformshift{1.006249in}{3.118525in}%
\pgfsys@useobject{currentmarker}{}%
\end{pgfscope}%
\begin{pgfscope}%
\pgfsys@transformshift{0.921739in}{3.125690in}%
\pgfsys@useobject{currentmarker}{}%
\end{pgfscope}%
\begin{pgfscope}%
\pgfsys@transformshift{1.135406in}{3.141870in}%
\pgfsys@useobject{currentmarker}{}%
\end{pgfscope}%
\begin{pgfscope}%
\pgfsys@transformshift{1.230850in}{3.132247in}%
\pgfsys@useobject{currentmarker}{}%
\end{pgfscope}%
\begin{pgfscope}%
\pgfsys@transformshift{0.942796in}{3.125484in}%
\pgfsys@useobject{currentmarker}{}%
\end{pgfscope}%
\begin{pgfscope}%
\pgfsys@transformshift{0.923958in}{3.122926in}%
\pgfsys@useobject{currentmarker}{}%
\end{pgfscope}%
\begin{pgfscope}%
\pgfsys@transformshift{1.178103in}{3.142003in}%
\pgfsys@useobject{currentmarker}{}%
\end{pgfscope}%
\begin{pgfscope}%
\pgfsys@transformshift{1.212552in}{3.128770in}%
\pgfsys@useobject{currentmarker}{}%
\end{pgfscope}%
\begin{pgfscope}%
\pgfsys@transformshift{1.260509in}{3.143095in}%
\pgfsys@useobject{currentmarker}{}%
\end{pgfscope}%
\begin{pgfscope}%
\pgfsys@transformshift{1.012873in}{3.130743in}%
\pgfsys@useobject{currentmarker}{}%
\end{pgfscope}%
\begin{pgfscope}%
\pgfsys@transformshift{0.975060in}{3.130893in}%
\pgfsys@useobject{currentmarker}{}%
\end{pgfscope}%
\begin{pgfscope}%
\pgfsys@transformshift{1.211791in}{3.140315in}%
\pgfsys@useobject{currentmarker}{}%
\end{pgfscope}%
\begin{pgfscope}%
\pgfsys@transformshift{1.040762in}{3.136437in}%
\pgfsys@useobject{currentmarker}{}%
\end{pgfscope}%
\begin{pgfscope}%
\pgfsys@transformshift{1.293075in}{3.132432in}%
\pgfsys@useobject{currentmarker}{}%
\end{pgfscope}%
\begin{pgfscope}%
\pgfsys@transformshift{1.295135in}{3.128247in}%
\pgfsys@useobject{currentmarker}{}%
\end{pgfscope}%
\begin{pgfscope}%
\pgfsys@transformshift{1.157710in}{3.139814in}%
\pgfsys@useobject{currentmarker}{}%
\end{pgfscope}%
\begin{pgfscope}%
\pgfsys@transformshift{1.297767in}{3.127796in}%
\pgfsys@useobject{currentmarker}{}%
\end{pgfscope}%
\begin{pgfscope}%
\pgfsys@transformshift{1.301514in}{3.141071in}%
\pgfsys@useobject{currentmarker}{}%
\end{pgfscope}%
\begin{pgfscope}%
\pgfsys@transformshift{1.000393in}{3.109600in}%
\pgfsys@useobject{currentmarker}{}%
\end{pgfscope}%
\begin{pgfscope}%
\pgfsys@transformshift{1.038103in}{3.128346in}%
\pgfsys@useobject{currentmarker}{}%
\end{pgfscope}%
\begin{pgfscope}%
\pgfsys@transformshift{0.976024in}{3.130850in}%
\pgfsys@useobject{currentmarker}{}%
\end{pgfscope}%
\begin{pgfscope}%
\pgfsys@transformshift{1.152192in}{3.141965in}%
\pgfsys@useobject{currentmarker}{}%
\end{pgfscope}%
\begin{pgfscope}%
\pgfsys@transformshift{1.480472in}{3.121997in}%
\pgfsys@useobject{currentmarker}{}%
\end{pgfscope}%
\begin{pgfscope}%
\pgfsys@transformshift{1.065085in}{3.137971in}%
\pgfsys@useobject{currentmarker}{}%
\end{pgfscope}%
\begin{pgfscope}%
\pgfsys@transformshift{1.118319in}{3.134847in}%
\pgfsys@useobject{currentmarker}{}%
\end{pgfscope}%
\begin{pgfscope}%
\pgfsys@transformshift{1.067813in}{3.135821in}%
\pgfsys@useobject{currentmarker}{}%
\end{pgfscope}%
\begin{pgfscope}%
\pgfsys@transformshift{1.037943in}{3.133713in}%
\pgfsys@useobject{currentmarker}{}%
\end{pgfscope}%
\begin{pgfscope}%
\pgfsys@transformshift{1.261645in}{3.142721in}%
\pgfsys@useobject{currentmarker}{}%
\end{pgfscope}%
\begin{pgfscope}%
\pgfsys@transformshift{0.938576in}{3.125227in}%
\pgfsys@useobject{currentmarker}{}%
\end{pgfscope}%
\begin{pgfscope}%
\pgfsys@transformshift{1.406308in}{3.127386in}%
\pgfsys@useobject{currentmarker}{}%
\end{pgfscope}%
\begin{pgfscope}%
\pgfsys@transformshift{1.026140in}{3.119771in}%
\pgfsys@useobject{currentmarker}{}%
\end{pgfscope}%
\begin{pgfscope}%
\pgfsys@transformshift{1.494243in}{3.106999in}%
\pgfsys@useobject{currentmarker}{}%
\end{pgfscope}%
\begin{pgfscope}%
\pgfsys@transformshift{1.491678in}{3.120266in}%
\pgfsys@useobject{currentmarker}{}%
\end{pgfscope}%
\begin{pgfscope}%
\pgfsys@transformshift{1.131177in}{3.140855in}%
\pgfsys@useobject{currentmarker}{}%
\end{pgfscope}%
\begin{pgfscope}%
\pgfsys@transformshift{1.037385in}{3.138824in}%
\pgfsys@useobject{currentmarker}{}%
\end{pgfscope}%
\begin{pgfscope}%
\pgfsys@transformshift{1.434848in}{3.114021in}%
\pgfsys@useobject{currentmarker}{}%
\end{pgfscope}%
\begin{pgfscope}%
\pgfsys@transformshift{1.115355in}{3.133079in}%
\pgfsys@useobject{currentmarker}{}%
\end{pgfscope}%
\begin{pgfscope}%
\pgfsys@transformshift{1.383456in}{3.129755in}%
\pgfsys@useobject{currentmarker}{}%
\end{pgfscope}%
\begin{pgfscope}%
\pgfsys@transformshift{1.122609in}{3.132853in}%
\pgfsys@useobject{currentmarker}{}%
\end{pgfscope}%
\begin{pgfscope}%
\pgfsys@transformshift{1.102760in}{3.134433in}%
\pgfsys@useobject{currentmarker}{}%
\end{pgfscope}%
\begin{pgfscope}%
\pgfsys@transformshift{1.454732in}{3.124289in}%
\pgfsys@useobject{currentmarker}{}%
\end{pgfscope}%
\begin{pgfscope}%
\pgfsys@transformshift{1.137381in}{3.142332in}%
\pgfsys@useobject{currentmarker}{}%
\end{pgfscope}%
\begin{pgfscope}%
\pgfsys@transformshift{1.209499in}{3.129381in}%
\pgfsys@useobject{currentmarker}{}%
\end{pgfscope}%
\begin{pgfscope}%
\pgfsys@transformshift{1.359949in}{3.133676in}%
\pgfsys@useobject{currentmarker}{}%
\end{pgfscope}%
\begin{pgfscope}%
\pgfsys@transformshift{1.132308in}{3.133850in}%
\pgfsys@useobject{currentmarker}{}%
\end{pgfscope}%
\begin{pgfscope}%
\pgfsys@transformshift{0.913932in}{3.118772in}%
\pgfsys@useobject{currentmarker}{}%
\end{pgfscope}%
\begin{pgfscope}%
\pgfsys@transformshift{0.848464in}{3.108691in}%
\pgfsys@useobject{currentmarker}{}%
\end{pgfscope}%
\begin{pgfscope}%
\pgfsys@transformshift{1.012747in}{3.127792in}%
\pgfsys@useobject{currentmarker}{}%
\end{pgfscope}%
\begin{pgfscope}%
\pgfsys@transformshift{1.076041in}{3.136948in}%
\pgfsys@useobject{currentmarker}{}%
\end{pgfscope}%
\begin{pgfscope}%
\pgfsys@transformshift{1.055249in}{3.138199in}%
\pgfsys@useobject{currentmarker}{}%
\end{pgfscope}%
\begin{pgfscope}%
\pgfsys@transformshift{1.351347in}{3.136130in}%
\pgfsys@useobject{currentmarker}{}%
\end{pgfscope}%
\begin{pgfscope}%
\pgfsys@transformshift{1.323679in}{3.114087in}%
\pgfsys@useobject{currentmarker}{}%
\end{pgfscope}%
\begin{pgfscope}%
\pgfsys@transformshift{1.265412in}{3.140793in}%
\pgfsys@useobject{currentmarker}{}%
\end{pgfscope}%
\begin{pgfscope}%
\pgfsys@transformshift{1.290674in}{3.137873in}%
\pgfsys@useobject{currentmarker}{}%
\end{pgfscope}%
\begin{pgfscope}%
\pgfsys@transformshift{1.313934in}{3.139801in}%
\pgfsys@useobject{currentmarker}{}%
\end{pgfscope}%
\begin{pgfscope}%
\pgfsys@transformshift{1.191546in}{3.130950in}%
\pgfsys@useobject{currentmarker}{}%
\end{pgfscope}%
\begin{pgfscope}%
\pgfsys@transformshift{1.279625in}{3.140860in}%
\pgfsys@useobject{currentmarker}{}%
\end{pgfscope}%
\begin{pgfscope}%
\pgfsys@transformshift{1.550613in}{3.106302in}%
\pgfsys@useobject{currentmarker}{}%
\end{pgfscope}%
\begin{pgfscope}%
\pgfsys@transformshift{1.502844in}{3.102777in}%
\pgfsys@useobject{currentmarker}{}%
\end{pgfscope}%
\begin{pgfscope}%
\pgfsys@transformshift{0.829333in}{3.099602in}%
\pgfsys@useobject{currentmarker}{}%
\end{pgfscope}%
\begin{pgfscope}%
\pgfsys@transformshift{1.159788in}{3.122072in}%
\pgfsys@useobject{currentmarker}{}%
\end{pgfscope}%
\begin{pgfscope}%
\pgfsys@transformshift{1.123149in}{3.121893in}%
\pgfsys@useobject{currentmarker}{}%
\end{pgfscope}%
\begin{pgfscope}%
\pgfsys@transformshift{1.440485in}{3.118968in}%
\pgfsys@useobject{currentmarker}{}%
\end{pgfscope}%
\begin{pgfscope}%
\pgfsys@transformshift{1.315129in}{3.132087in}%
\pgfsys@useobject{currentmarker}{}%
\end{pgfscope}%
\begin{pgfscope}%
\pgfsys@transformshift{1.097616in}{3.104819in}%
\pgfsys@useobject{currentmarker}{}%
\end{pgfscope}%
\begin{pgfscope}%
\pgfsys@transformshift{1.413705in}{3.132747in}%
\pgfsys@useobject{currentmarker}{}%
\end{pgfscope}%
\begin{pgfscope}%
\pgfsys@transformshift{0.980042in}{3.117753in}%
\pgfsys@useobject{currentmarker}{}%
\end{pgfscope}%
\begin{pgfscope}%
\pgfsys@transformshift{1.221769in}{3.134881in}%
\pgfsys@useobject{currentmarker}{}%
\end{pgfscope}%
\begin{pgfscope}%
\pgfsys@transformshift{1.076779in}{3.099924in}%
\pgfsys@useobject{currentmarker}{}%
\end{pgfscope}%
\begin{pgfscope}%
\pgfsys@transformshift{1.375183in}{3.124433in}%
\pgfsys@useobject{currentmarker}{}%
\end{pgfscope}%
\begin{pgfscope}%
\pgfsys@transformshift{1.302316in}{3.141551in}%
\pgfsys@useobject{currentmarker}{}%
\end{pgfscope}%
\begin{pgfscope}%
\pgfsys@transformshift{1.286023in}{3.131347in}%
\pgfsys@useobject{currentmarker}{}%
\end{pgfscope}%
\begin{pgfscope}%
\pgfsys@transformshift{1.331004in}{3.138108in}%
\pgfsys@useobject{currentmarker}{}%
\end{pgfscope}%
\begin{pgfscope}%
\pgfsys@transformshift{1.120441in}{3.134403in}%
\pgfsys@useobject{currentmarker}{}%
\end{pgfscope}%
\begin{pgfscope}%
\pgfsys@transformshift{1.116267in}{3.133251in}%
\pgfsys@useobject{currentmarker}{}%
\end{pgfscope}%
\begin{pgfscope}%
\pgfsys@transformshift{1.502886in}{3.119168in}%
\pgfsys@useobject{currentmarker}{}%
\end{pgfscope}%
\begin{pgfscope}%
\pgfsys@transformshift{1.085736in}{3.132274in}%
\pgfsys@useobject{currentmarker}{}%
\end{pgfscope}%
\begin{pgfscope}%
\pgfsys@transformshift{1.565330in}{3.104964in}%
\pgfsys@useobject{currentmarker}{}%
\end{pgfscope}%
\begin{pgfscope}%
\pgfsys@transformshift{1.317857in}{3.137864in}%
\pgfsys@useobject{currentmarker}{}%
\end{pgfscope}%
\begin{pgfscope}%
\pgfsys@transformshift{1.104651in}{3.142285in}%
\pgfsys@useobject{currentmarker}{}%
\end{pgfscope}%
\begin{pgfscope}%
\pgfsys@transformshift{1.252706in}{3.139179in}%
\pgfsys@useobject{currentmarker}{}%
\end{pgfscope}%
\begin{pgfscope}%
\pgfsys@transformshift{1.443498in}{3.107000in}%
\pgfsys@useobject{currentmarker}{}%
\end{pgfscope}%
\begin{pgfscope}%
\pgfsys@transformshift{1.024513in}{3.125435in}%
\pgfsys@useobject{currentmarker}{}%
\end{pgfscope}%
\begin{pgfscope}%
\pgfsys@transformshift{1.058577in}{3.124705in}%
\pgfsys@useobject{currentmarker}{}%
\end{pgfscope}%
\begin{pgfscope}%
\pgfsys@transformshift{1.232723in}{3.114047in}%
\pgfsys@useobject{currentmarker}{}%
\end{pgfscope}%
\begin{pgfscope}%
\pgfsys@transformshift{1.076692in}{3.133640in}%
\pgfsys@useobject{currentmarker}{}%
\end{pgfscope}%
\begin{pgfscope}%
\pgfsys@transformshift{1.129976in}{3.137761in}%
\pgfsys@useobject{currentmarker}{}%
\end{pgfscope}%
\begin{pgfscope}%
\pgfsys@transformshift{1.451837in}{3.124549in}%
\pgfsys@useobject{currentmarker}{}%
\end{pgfscope}%
\begin{pgfscope}%
\pgfsys@transformshift{1.210416in}{3.139407in}%
\pgfsys@useobject{currentmarker}{}%
\end{pgfscope}%
\begin{pgfscope}%
\pgfsys@transformshift{1.084212in}{3.137753in}%
\pgfsys@useobject{currentmarker}{}%
\end{pgfscope}%
\begin{pgfscope}%
\pgfsys@transformshift{1.365124in}{3.115323in}%
\pgfsys@useobject{currentmarker}{}%
\end{pgfscope}%
\begin{pgfscope}%
\pgfsys@transformshift{1.064281in}{3.137017in}%
\pgfsys@useobject{currentmarker}{}%
\end{pgfscope}%
\begin{pgfscope}%
\pgfsys@transformshift{1.297074in}{3.139934in}%
\pgfsys@useobject{currentmarker}{}%
\end{pgfscope}%
\begin{pgfscope}%
\pgfsys@transformshift{0.955316in}{3.119484in}%
\pgfsys@useobject{currentmarker}{}%
\end{pgfscope}%
\begin{pgfscope}%
\pgfsys@transformshift{1.310281in}{3.132255in}%
\pgfsys@useobject{currentmarker}{}%
\end{pgfscope}%
\begin{pgfscope}%
\pgfsys@transformshift{1.394848in}{3.133936in}%
\pgfsys@useobject{currentmarker}{}%
\end{pgfscope}%
\begin{pgfscope}%
\pgfsys@transformshift{0.982898in}{3.133216in}%
\pgfsys@useobject{currentmarker}{}%
\end{pgfscope}%
\begin{pgfscope}%
\pgfsys@transformshift{0.920949in}{3.112596in}%
\pgfsys@useobject{currentmarker}{}%
\end{pgfscope}%
\begin{pgfscope}%
\pgfsys@transformshift{1.276229in}{3.134479in}%
\pgfsys@useobject{currentmarker}{}%
\end{pgfscope}%
\begin{pgfscope}%
\pgfsys@transformshift{1.175720in}{3.133770in}%
\pgfsys@useobject{currentmarker}{}%
\end{pgfscope}%
\begin{pgfscope}%
\pgfsys@transformshift{1.285301in}{3.138414in}%
\pgfsys@useobject{currentmarker}{}%
\end{pgfscope}%
\begin{pgfscope}%
\pgfsys@transformshift{1.294801in}{3.137909in}%
\pgfsys@useobject{currentmarker}{}%
\end{pgfscope}%
\begin{pgfscope}%
\pgfsys@transformshift{1.339796in}{3.125033in}%
\pgfsys@useobject{currentmarker}{}%
\end{pgfscope}%
\begin{pgfscope}%
\pgfsys@transformshift{1.217852in}{3.135940in}%
\pgfsys@useobject{currentmarker}{}%
\end{pgfscope}%
\begin{pgfscope}%
\pgfsys@transformshift{1.315246in}{3.139524in}%
\pgfsys@useobject{currentmarker}{}%
\end{pgfscope}%
\begin{pgfscope}%
\pgfsys@transformshift{1.130376in}{3.140216in}%
\pgfsys@useobject{currentmarker}{}%
\end{pgfscope}%
\begin{pgfscope}%
\pgfsys@transformshift{0.995270in}{3.117482in}%
\pgfsys@useobject{currentmarker}{}%
\end{pgfscope}%
\begin{pgfscope}%
\pgfsys@transformshift{1.357462in}{3.101775in}%
\pgfsys@useobject{currentmarker}{}%
\end{pgfscope}%
\begin{pgfscope}%
\pgfsys@transformshift{1.249476in}{3.134228in}%
\pgfsys@useobject{currentmarker}{}%
\end{pgfscope}%
\begin{pgfscope}%
\pgfsys@transformshift{1.398073in}{3.120352in}%
\pgfsys@useobject{currentmarker}{}%
\end{pgfscope}%
\begin{pgfscope}%
\pgfsys@transformshift{1.084724in}{3.132002in}%
\pgfsys@useobject{currentmarker}{}%
\end{pgfscope}%
\begin{pgfscope}%
\pgfsys@transformshift{1.296826in}{3.141760in}%
\pgfsys@useobject{currentmarker}{}%
\end{pgfscope}%
\begin{pgfscope}%
\pgfsys@transformshift{1.149726in}{3.126788in}%
\pgfsys@useobject{currentmarker}{}%
\end{pgfscope}%
\begin{pgfscope}%
\pgfsys@transformshift{0.979555in}{3.122072in}%
\pgfsys@useobject{currentmarker}{}%
\end{pgfscope}%
\begin{pgfscope}%
\pgfsys@transformshift{0.997436in}{3.129116in}%
\pgfsys@useobject{currentmarker}{}%
\end{pgfscope}%
\begin{pgfscope}%
\pgfsys@transformshift{1.095233in}{3.134384in}%
\pgfsys@useobject{currentmarker}{}%
\end{pgfscope}%
\begin{pgfscope}%
\pgfsys@transformshift{1.009804in}{3.113044in}%
\pgfsys@useobject{currentmarker}{}%
\end{pgfscope}%
\begin{pgfscope}%
\pgfsys@transformshift{1.381011in}{3.126829in}%
\pgfsys@useobject{currentmarker}{}%
\end{pgfscope}%
\begin{pgfscope}%
\pgfsys@transformshift{1.371746in}{3.115655in}%
\pgfsys@useobject{currentmarker}{}%
\end{pgfscope}%
\begin{pgfscope}%
\pgfsys@transformshift{1.373154in}{3.134510in}%
\pgfsys@useobject{currentmarker}{}%
\end{pgfscope}%
\begin{pgfscope}%
\pgfsys@transformshift{1.347669in}{3.138581in}%
\pgfsys@useobject{currentmarker}{}%
\end{pgfscope}%
\begin{pgfscope}%
\pgfsys@transformshift{1.270167in}{3.126709in}%
\pgfsys@useobject{currentmarker}{}%
\end{pgfscope}%
\begin{pgfscope}%
\pgfsys@transformshift{1.436534in}{3.126543in}%
\pgfsys@useobject{currentmarker}{}%
\end{pgfscope}%
\begin{pgfscope}%
\pgfsys@transformshift{1.021173in}{3.131377in}%
\pgfsys@useobject{currentmarker}{}%
\end{pgfscope}%
\begin{pgfscope}%
\pgfsys@transformshift{1.278976in}{3.120370in}%
\pgfsys@useobject{currentmarker}{}%
\end{pgfscope}%
\begin{pgfscope}%
\pgfsys@transformshift{1.277771in}{3.121843in}%
\pgfsys@useobject{currentmarker}{}%
\end{pgfscope}%
\begin{pgfscope}%
\pgfsys@transformshift{1.291931in}{3.125128in}%
\pgfsys@useobject{currentmarker}{}%
\end{pgfscope}%
\begin{pgfscope}%
\pgfsys@transformshift{1.337294in}{3.129143in}%
\pgfsys@useobject{currentmarker}{}%
\end{pgfscope}%
\begin{pgfscope}%
\pgfsys@transformshift{1.163241in}{3.111105in}%
\pgfsys@useobject{currentmarker}{}%
\end{pgfscope}%
\begin{pgfscope}%
\pgfsys@transformshift{1.201445in}{3.136676in}%
\pgfsys@useobject{currentmarker}{}%
\end{pgfscope}%
\begin{pgfscope}%
\pgfsys@transformshift{1.275720in}{3.132392in}%
\pgfsys@useobject{currentmarker}{}%
\end{pgfscope}%
\begin{pgfscope}%
\pgfsys@transformshift{1.289562in}{3.124028in}%
\pgfsys@useobject{currentmarker}{}%
\end{pgfscope}%
\begin{pgfscope}%
\pgfsys@transformshift{1.550563in}{3.110660in}%
\pgfsys@useobject{currentmarker}{}%
\end{pgfscope}%
\begin{pgfscope}%
\pgfsys@transformshift{1.197430in}{3.136928in}%
\pgfsys@useobject{currentmarker}{}%
\end{pgfscope}%
\begin{pgfscope}%
\pgfsys@transformshift{1.513463in}{3.115465in}%
\pgfsys@useobject{currentmarker}{}%
\end{pgfscope}%
\begin{pgfscope}%
\pgfsys@transformshift{1.306015in}{3.115618in}%
\pgfsys@useobject{currentmarker}{}%
\end{pgfscope}%
\begin{pgfscope}%
\pgfsys@transformshift{1.381276in}{3.115436in}%
\pgfsys@useobject{currentmarker}{}%
\end{pgfscope}%
\begin{pgfscope}%
\pgfsys@transformshift{1.036041in}{3.111075in}%
\pgfsys@useobject{currentmarker}{}%
\end{pgfscope}%
\begin{pgfscope}%
\pgfsys@transformshift{1.239434in}{3.140593in}%
\pgfsys@useobject{currentmarker}{}%
\end{pgfscope}%
\begin{pgfscope}%
\pgfsys@transformshift{0.959267in}{3.122833in}%
\pgfsys@useobject{currentmarker}{}%
\end{pgfscope}%
\begin{pgfscope}%
\pgfsys@transformshift{1.430717in}{3.124280in}%
\pgfsys@useobject{currentmarker}{}%
\end{pgfscope}%
\begin{pgfscope}%
\pgfsys@transformshift{1.324229in}{3.138638in}%
\pgfsys@useobject{currentmarker}{}%
\end{pgfscope}%
\begin{pgfscope}%
\pgfsys@transformshift{1.112361in}{3.133476in}%
\pgfsys@useobject{currentmarker}{}%
\end{pgfscope}%
\begin{pgfscope}%
\pgfsys@transformshift{1.298061in}{3.118949in}%
\pgfsys@useobject{currentmarker}{}%
\end{pgfscope}%
\begin{pgfscope}%
\pgfsys@transformshift{1.152113in}{3.128488in}%
\pgfsys@useobject{currentmarker}{}%
\end{pgfscope}%
\begin{pgfscope}%
\pgfsys@transformshift{1.162402in}{3.139969in}%
\pgfsys@useobject{currentmarker}{}%
\end{pgfscope}%
\begin{pgfscope}%
\pgfsys@transformshift{1.008822in}{3.126661in}%
\pgfsys@useobject{currentmarker}{}%
\end{pgfscope}%
\begin{pgfscope}%
\pgfsys@transformshift{1.211239in}{3.135673in}%
\pgfsys@useobject{currentmarker}{}%
\end{pgfscope}%
\begin{pgfscope}%
\pgfsys@transformshift{0.944872in}{3.112649in}%
\pgfsys@useobject{currentmarker}{}%
\end{pgfscope}%
\begin{pgfscope}%
\pgfsys@transformshift{1.040022in}{3.136066in}%
\pgfsys@useobject{currentmarker}{}%
\end{pgfscope}%
\begin{pgfscope}%
\pgfsys@transformshift{1.111709in}{3.137249in}%
\pgfsys@useobject{currentmarker}{}%
\end{pgfscope}%
\begin{pgfscope}%
\pgfsys@transformshift{1.260731in}{3.138461in}%
\pgfsys@useobject{currentmarker}{}%
\end{pgfscope}%
\begin{pgfscope}%
\pgfsys@transformshift{1.172967in}{3.134793in}%
\pgfsys@useobject{currentmarker}{}%
\end{pgfscope}%
\begin{pgfscope}%
\pgfsys@transformshift{1.128888in}{3.126511in}%
\pgfsys@useobject{currentmarker}{}%
\end{pgfscope}%
\begin{pgfscope}%
\pgfsys@transformshift{1.280460in}{3.109012in}%
\pgfsys@useobject{currentmarker}{}%
\end{pgfscope}%
\begin{pgfscope}%
\pgfsys@transformshift{1.344561in}{3.130830in}%
\pgfsys@useobject{currentmarker}{}%
\end{pgfscope}%
\begin{pgfscope}%
\pgfsys@transformshift{1.203794in}{3.136790in}%
\pgfsys@useobject{currentmarker}{}%
\end{pgfscope}%
\begin{pgfscope}%
\pgfsys@transformshift{1.080998in}{3.119147in}%
\pgfsys@useobject{currentmarker}{}%
\end{pgfscope}%
\begin{pgfscope}%
\pgfsys@transformshift{1.140245in}{3.122709in}%
\pgfsys@useobject{currentmarker}{}%
\end{pgfscope}%
\begin{pgfscope}%
\pgfsys@transformshift{1.187150in}{3.131288in}%
\pgfsys@useobject{currentmarker}{}%
\end{pgfscope}%
\begin{pgfscope}%
\pgfsys@transformshift{1.048919in}{3.136570in}%
\pgfsys@useobject{currentmarker}{}%
\end{pgfscope}%
\begin{pgfscope}%
\pgfsys@transformshift{1.102875in}{3.140316in}%
\pgfsys@useobject{currentmarker}{}%
\end{pgfscope}%
\begin{pgfscope}%
\pgfsys@transformshift{0.964875in}{3.115071in}%
\pgfsys@useobject{currentmarker}{}%
\end{pgfscope}%
\begin{pgfscope}%
\pgfsys@transformshift{0.907834in}{3.118793in}%
\pgfsys@useobject{currentmarker}{}%
\end{pgfscope}%
\begin{pgfscope}%
\pgfsys@transformshift{1.175898in}{3.137812in}%
\pgfsys@useobject{currentmarker}{}%
\end{pgfscope}%
\begin{pgfscope}%
\pgfsys@transformshift{1.026189in}{3.134076in}%
\pgfsys@useobject{currentmarker}{}%
\end{pgfscope}%
\begin{pgfscope}%
\pgfsys@transformshift{1.320819in}{3.135914in}%
\pgfsys@useobject{currentmarker}{}%
\end{pgfscope}%
\begin{pgfscope}%
\pgfsys@transformshift{1.151691in}{3.144015in}%
\pgfsys@useobject{currentmarker}{}%
\end{pgfscope}%
\begin{pgfscope}%
\pgfsys@transformshift{1.196144in}{3.141578in}%
\pgfsys@useobject{currentmarker}{}%
\end{pgfscope}%
\begin{pgfscope}%
\pgfsys@transformshift{1.106621in}{3.138897in}%
\pgfsys@useobject{currentmarker}{}%
\end{pgfscope}%
\begin{pgfscope}%
\pgfsys@transformshift{1.084233in}{3.137265in}%
\pgfsys@useobject{currentmarker}{}%
\end{pgfscope}%
\begin{pgfscope}%
\pgfsys@transformshift{1.420484in}{3.127285in}%
\pgfsys@useobject{currentmarker}{}%
\end{pgfscope}%
\end{pgfscope}%
\begin{pgfscope}%
\pgfpathmoveto{\pgfqpoint{1.192000in}{2.052407in}}%
\pgfpathcurveto{\pgfqpoint{1.481602in}{2.052407in}}{\pgfqpoint{1.759381in}{2.109937in}}{\pgfqpoint{1.964161in}{2.212326in}}%
\pgfpathcurveto{\pgfqpoint{2.168940in}{2.314716in}}{\pgfqpoint{2.284000in}{2.453606in}}{\pgfqpoint{2.284000in}{2.598407in}}%
\pgfpathcurveto{\pgfqpoint{2.284000in}{2.743208in}}{\pgfqpoint{2.168940in}{2.882097in}}{\pgfqpoint{1.964161in}{2.984487in}}%
\pgfpathcurveto{\pgfqpoint{1.759381in}{3.086877in}}{\pgfqpoint{1.481602in}{3.144407in}}{\pgfqpoint{1.192000in}{3.144407in}}%
\pgfpathcurveto{\pgfqpoint{0.902398in}{3.144407in}}{\pgfqpoint{0.624619in}{3.086877in}}{\pgfqpoint{0.419839in}{2.984487in}}%
\pgfpathcurveto{\pgfqpoint{0.215060in}{2.882097in}}{\pgfqpoint{0.100000in}{2.743208in}}{\pgfqpoint{0.100000in}{2.598407in}}%
\pgfpathcurveto{\pgfqpoint{0.100000in}{2.453606in}}{\pgfqpoint{0.215060in}{2.314716in}}{\pgfqpoint{0.419839in}{2.212326in}}%
\pgfpathcurveto{\pgfqpoint{0.624619in}{2.109937in}}{\pgfqpoint{0.902398in}{2.052407in}}{\pgfqpoint{1.192000in}{2.052407in}}%
\pgfpathlineto{\pgfqpoint{1.192000in}{2.052407in}}%
\pgfpathclose%
\pgfusepath{clip}%
\pgfsetrectcap%
\pgfsetroundjoin%
\pgfsetlinewidth{0.451687pt}%
\definecolor{currentstroke}{rgb}{0.690196,0.690196,0.690196}%
\pgfsetstrokecolor{currentstroke}%
\pgfsetstrokeopacity{0.250000}%
\pgfsetdash{}{0pt}%
\pgfpathmoveto{\pgfqpoint{0.806921in}{2.103702in}}%
\pgfpathlineto{\pgfqpoint{0.775157in}{2.113059in}}%
\pgfpathlineto{\pgfqpoint{0.744700in}{2.122919in}}%
\pgfpathlineto{\pgfqpoint{0.716255in}{2.132965in}}%
\pgfpathlineto{\pgfqpoint{0.689090in}{2.143363in}}%
\pgfpathlineto{\pgfqpoint{0.663134in}{2.154082in}}%
\pgfpathlineto{\pgfqpoint{0.638319in}{2.165102in}}%
\pgfpathlineto{\pgfqpoint{0.614585in}{2.176401in}}%
\pgfpathlineto{\pgfqpoint{0.591881in}{2.187962in}}%
\pgfpathlineto{\pgfqpoint{0.570160in}{2.199772in}}%
\pgfpathlineto{\pgfqpoint{0.549384in}{2.211817in}}%
\pgfpathlineto{\pgfqpoint{0.529517in}{2.224083in}}%
\pgfpathlineto{\pgfqpoint{0.510532in}{2.236561in}}%
\pgfpathlineto{\pgfqpoint{0.492403in}{2.249238in}}%
\pgfpathlineto{\pgfqpoint{0.475106in}{2.262105in}}%
\pgfpathlineto{\pgfqpoint{0.458426in}{2.275314in}}%
\pgfpathlineto{\pgfqpoint{0.442806in}{2.288486in}}%
\pgfpathlineto{\pgfqpoint{0.427948in}{2.301831in}}%
\pgfpathlineto{\pgfqpoint{0.413847in}{2.315338in}}%
\pgfpathlineto{\pgfqpoint{0.400494in}{2.328996in}}%
\pgfpathlineto{\pgfqpoint{0.387882in}{2.342794in}}%
\pgfpathlineto{\pgfqpoint{0.376004in}{2.356724in}}%
\pgfpathlineto{\pgfqpoint{0.364855in}{2.370777in}}%
\pgfpathlineto{\pgfqpoint{0.354428in}{2.384944in}}%
\pgfpathlineto{\pgfqpoint{0.344718in}{2.399216in}}%
\pgfpathlineto{\pgfqpoint{0.335720in}{2.413586in}}%
\pgfpathlineto{\pgfqpoint{0.327430in}{2.428046in}}%
\pgfpathlineto{\pgfqpoint{0.319844in}{2.442588in}}%
\pgfpathlineto{\pgfqpoint{0.312957in}{2.457204in}}%
\pgfpathlineto{\pgfqpoint{0.306768in}{2.471888in}}%
\pgfpathlineto{\pgfqpoint{0.301272in}{2.486632in}}%
\pgfpathlineto{\pgfqpoint{0.296469in}{2.501429in}}%
\pgfpathlineto{\pgfqpoint{0.292314in}{2.516436in}}%
\pgfpathlineto{\pgfqpoint{0.288911in}{2.531243in}}%
\pgfpathlineto{\pgfqpoint{0.286184in}{2.546108in}}%
\pgfpathlineto{\pgfqpoint{0.284136in}{2.561018in}}%
\pgfpathlineto{\pgfqpoint{0.282769in}{2.575961in}}%
\pgfpathlineto{\pgfqpoint{0.282085in}{2.590923in}}%
\pgfpathlineto{\pgfqpoint{0.282085in}{2.605891in}}%
\pgfpathlineto{\pgfqpoint{0.282769in}{2.620853in}}%
\pgfpathlineto{\pgfqpoint{0.284136in}{2.635795in}}%
\pgfpathlineto{\pgfqpoint{0.286184in}{2.650705in}}%
\pgfpathlineto{\pgfqpoint{0.288911in}{2.665570in}}%
\pgfpathlineto{\pgfqpoint{0.292314in}{2.680378in}}%
\pgfpathlineto{\pgfqpoint{0.296469in}{2.695384in}}%
\pgfpathlineto{\pgfqpoint{0.301272in}{2.710181in}}%
\pgfpathlineto{\pgfqpoint{0.306768in}{2.724925in}}%
\pgfpathlineto{\pgfqpoint{0.312957in}{2.739609in}}%
\pgfpathlineto{\pgfqpoint{0.319844in}{2.754226in}}%
\pgfpathlineto{\pgfqpoint{0.327430in}{2.768768in}}%
\pgfpathlineto{\pgfqpoint{0.335720in}{2.783228in}}%
\pgfpathlineto{\pgfqpoint{0.344718in}{2.797598in}}%
\pgfpathlineto{\pgfqpoint{0.354428in}{2.811870in}}%
\pgfpathlineto{\pgfqpoint{0.364855in}{2.826037in}}%
\pgfpathlineto{\pgfqpoint{0.376004in}{2.840089in}}%
\pgfpathlineto{\pgfqpoint{0.387882in}{2.854019in}}%
\pgfpathlineto{\pgfqpoint{0.400494in}{2.867818in}}%
\pgfpathlineto{\pgfqpoint{0.413847in}{2.881476in}}%
\pgfpathlineto{\pgfqpoint{0.427948in}{2.894982in}}%
\pgfpathlineto{\pgfqpoint{0.442806in}{2.908327in}}%
\pgfpathlineto{\pgfqpoint{0.458426in}{2.921500in}}%
\pgfpathlineto{\pgfqpoint{0.475106in}{2.934708in}}%
\pgfpathlineto{\pgfqpoint{0.492403in}{2.947576in}}%
\pgfpathlineto{\pgfqpoint{0.510532in}{2.960253in}}%
\pgfpathlineto{\pgfqpoint{0.529517in}{2.972730in}}%
\pgfpathlineto{\pgfqpoint{0.549384in}{2.984997in}}%
\pgfpathlineto{\pgfqpoint{0.570160in}{2.997041in}}%
\pgfpathlineto{\pgfqpoint{0.591881in}{3.008851in}}%
\pgfpathlineto{\pgfqpoint{0.614585in}{3.020413in}}%
\pgfpathlineto{\pgfqpoint{0.638319in}{3.031712in}}%
\pgfpathlineto{\pgfqpoint{0.663134in}{3.042731in}}%
\pgfpathlineto{\pgfqpoint{0.689090in}{3.053451in}}%
\pgfpathlineto{\pgfqpoint{0.716255in}{3.063848in}}%
\pgfpathlineto{\pgfqpoint{0.744700in}{3.073894in}}%
\pgfpathlineto{\pgfqpoint{0.775157in}{3.083755in}}%
\pgfpathlineto{\pgfqpoint{0.806921in}{3.093112in}}%
\pgfusepath{stroke}%
\end{pgfscope}%
\begin{pgfscope}%
\pgfpathmoveto{\pgfqpoint{1.192000in}{2.052407in}}%
\pgfpathcurveto{\pgfqpoint{1.481602in}{2.052407in}}{\pgfqpoint{1.759381in}{2.109937in}}{\pgfqpoint{1.964161in}{2.212326in}}%
\pgfpathcurveto{\pgfqpoint{2.168940in}{2.314716in}}{\pgfqpoint{2.284000in}{2.453606in}}{\pgfqpoint{2.284000in}{2.598407in}}%
\pgfpathcurveto{\pgfqpoint{2.284000in}{2.743208in}}{\pgfqpoint{2.168940in}{2.882097in}}{\pgfqpoint{1.964161in}{2.984487in}}%
\pgfpathcurveto{\pgfqpoint{1.759381in}{3.086877in}}{\pgfqpoint{1.481602in}{3.144407in}}{\pgfqpoint{1.192000in}{3.144407in}}%
\pgfpathcurveto{\pgfqpoint{0.902398in}{3.144407in}}{\pgfqpoint{0.624619in}{3.086877in}}{\pgfqpoint{0.419839in}{2.984487in}}%
\pgfpathcurveto{\pgfqpoint{0.215060in}{2.882097in}}{\pgfqpoint{0.100000in}{2.743208in}}{\pgfqpoint{0.100000in}{2.598407in}}%
\pgfpathcurveto{\pgfqpoint{0.100000in}{2.453606in}}{\pgfqpoint{0.215060in}{2.314716in}}{\pgfqpoint{0.419839in}{2.212326in}}%
\pgfpathcurveto{\pgfqpoint{0.624619in}{2.109937in}}{\pgfqpoint{0.902398in}{2.052407in}}{\pgfqpoint{1.192000in}{2.052407in}}%
\pgfpathlineto{\pgfqpoint{1.192000in}{2.052407in}}%
\pgfpathclose%
\pgfusepath{clip}%
\pgfsetrectcap%
\pgfsetroundjoin%
\pgfsetlinewidth{0.451687pt}%
\definecolor{currentstroke}{rgb}{0.690196,0.690196,0.690196}%
\pgfsetstrokecolor{currentstroke}%
\pgfsetstrokeopacity{0.250000}%
\pgfsetdash{}{0pt}%
\pgfpathmoveto{\pgfqpoint{0.883937in}{2.103702in}}%
\pgfpathlineto{\pgfqpoint{0.858525in}{2.113059in}}%
\pgfpathlineto{\pgfqpoint{0.834160in}{2.122919in}}%
\pgfpathlineto{\pgfqpoint{0.811404in}{2.132965in}}%
\pgfpathlineto{\pgfqpoint{0.789672in}{2.143363in}}%
\pgfpathlineto{\pgfqpoint{0.768907in}{2.154082in}}%
\pgfpathlineto{\pgfqpoint{0.749055in}{2.165102in}}%
\pgfpathlineto{\pgfqpoint{0.730068in}{2.176401in}}%
\pgfpathlineto{\pgfqpoint{0.711905in}{2.187962in}}%
\pgfpathlineto{\pgfqpoint{0.694528in}{2.199772in}}%
\pgfpathlineto{\pgfqpoint{0.677907in}{2.211817in}}%
\pgfpathlineto{\pgfqpoint{0.662014in}{2.224083in}}%
\pgfpathlineto{\pgfqpoint{0.646826in}{2.236561in}}%
\pgfpathlineto{\pgfqpoint{0.632322in}{2.249238in}}%
\pgfpathlineto{\pgfqpoint{0.618485in}{2.262105in}}%
\pgfpathlineto{\pgfqpoint{0.605141in}{2.275314in}}%
\pgfpathlineto{\pgfqpoint{0.592644in}{2.288486in}}%
\pgfpathlineto{\pgfqpoint{0.580759in}{2.301831in}}%
\pgfpathlineto{\pgfqpoint{0.569477in}{2.315338in}}%
\pgfpathlineto{\pgfqpoint{0.558795in}{2.328996in}}%
\pgfpathlineto{\pgfqpoint{0.548705in}{2.342794in}}%
\pgfpathlineto{\pgfqpoint{0.539204in}{2.356724in}}%
\pgfpathlineto{\pgfqpoint{0.530284in}{2.370777in}}%
\pgfpathlineto{\pgfqpoint{0.521943in}{2.384944in}}%
\pgfpathlineto{\pgfqpoint{0.514175in}{2.399216in}}%
\pgfpathlineto{\pgfqpoint{0.506976in}{2.413586in}}%
\pgfpathlineto{\pgfqpoint{0.500344in}{2.428046in}}%
\pgfpathlineto{\pgfqpoint{0.494275in}{2.442588in}}%
\pgfpathlineto{\pgfqpoint{0.488766in}{2.457204in}}%
\pgfpathlineto{\pgfqpoint{0.483814in}{2.471888in}}%
\pgfpathlineto{\pgfqpoint{0.479418in}{2.486632in}}%
\pgfpathlineto{\pgfqpoint{0.475575in}{2.501429in}}%
\pgfpathlineto{\pgfqpoint{0.472251in}{2.516436in}}%
\pgfpathlineto{\pgfqpoint{0.469529in}{2.531243in}}%
\pgfpathlineto{\pgfqpoint{0.467347in}{2.546108in}}%
\pgfpathlineto{\pgfqpoint{0.465709in}{2.561018in}}%
\pgfpathlineto{\pgfqpoint{0.464615in}{2.575961in}}%
\pgfpathlineto{\pgfqpoint{0.464068in}{2.590923in}}%
\pgfpathlineto{\pgfqpoint{0.464068in}{2.605891in}}%
\pgfpathlineto{\pgfqpoint{0.464615in}{2.620853in}}%
\pgfpathlineto{\pgfqpoint{0.465709in}{2.635795in}}%
\pgfpathlineto{\pgfqpoint{0.467347in}{2.650705in}}%
\pgfpathlineto{\pgfqpoint{0.469529in}{2.665570in}}%
\pgfpathlineto{\pgfqpoint{0.472251in}{2.680378in}}%
\pgfpathlineto{\pgfqpoint{0.475575in}{2.695384in}}%
\pgfpathlineto{\pgfqpoint{0.479418in}{2.710181in}}%
\pgfpathlineto{\pgfqpoint{0.483814in}{2.724925in}}%
\pgfpathlineto{\pgfqpoint{0.488766in}{2.739609in}}%
\pgfpathlineto{\pgfqpoint{0.494275in}{2.754226in}}%
\pgfpathlineto{\pgfqpoint{0.500344in}{2.768768in}}%
\pgfpathlineto{\pgfqpoint{0.506976in}{2.783228in}}%
\pgfpathlineto{\pgfqpoint{0.514175in}{2.797598in}}%
\pgfpathlineto{\pgfqpoint{0.521943in}{2.811870in}}%
\pgfpathlineto{\pgfqpoint{0.530284in}{2.826037in}}%
\pgfpathlineto{\pgfqpoint{0.539204in}{2.840089in}}%
\pgfpathlineto{\pgfqpoint{0.548705in}{2.854019in}}%
\pgfpathlineto{\pgfqpoint{0.558795in}{2.867818in}}%
\pgfpathlineto{\pgfqpoint{0.569477in}{2.881476in}}%
\pgfpathlineto{\pgfqpoint{0.580759in}{2.894982in}}%
\pgfpathlineto{\pgfqpoint{0.592644in}{2.908327in}}%
\pgfpathlineto{\pgfqpoint{0.605141in}{2.921500in}}%
\pgfpathlineto{\pgfqpoint{0.618485in}{2.934708in}}%
\pgfpathlineto{\pgfqpoint{0.632322in}{2.947576in}}%
\pgfpathlineto{\pgfqpoint{0.646826in}{2.960253in}}%
\pgfpathlineto{\pgfqpoint{0.662014in}{2.972730in}}%
\pgfpathlineto{\pgfqpoint{0.677907in}{2.984997in}}%
\pgfpathlineto{\pgfqpoint{0.694528in}{2.997041in}}%
\pgfpathlineto{\pgfqpoint{0.711905in}{3.008851in}}%
\pgfpathlineto{\pgfqpoint{0.730068in}{3.020413in}}%
\pgfpathlineto{\pgfqpoint{0.749055in}{3.031712in}}%
\pgfpathlineto{\pgfqpoint{0.768907in}{3.042731in}}%
\pgfpathlineto{\pgfqpoint{0.789672in}{3.053451in}}%
\pgfpathlineto{\pgfqpoint{0.811404in}{3.063848in}}%
\pgfpathlineto{\pgfqpoint{0.834160in}{3.073894in}}%
\pgfpathlineto{\pgfqpoint{0.858525in}{3.083755in}}%
\pgfpathlineto{\pgfqpoint{0.883937in}{3.093112in}}%
\pgfusepath{stroke}%
\end{pgfscope}%
\begin{pgfscope}%
\pgfpathmoveto{\pgfqpoint{1.192000in}{2.052407in}}%
\pgfpathcurveto{\pgfqpoint{1.481602in}{2.052407in}}{\pgfqpoint{1.759381in}{2.109937in}}{\pgfqpoint{1.964161in}{2.212326in}}%
\pgfpathcurveto{\pgfqpoint{2.168940in}{2.314716in}}{\pgfqpoint{2.284000in}{2.453606in}}{\pgfqpoint{2.284000in}{2.598407in}}%
\pgfpathcurveto{\pgfqpoint{2.284000in}{2.743208in}}{\pgfqpoint{2.168940in}{2.882097in}}{\pgfqpoint{1.964161in}{2.984487in}}%
\pgfpathcurveto{\pgfqpoint{1.759381in}{3.086877in}}{\pgfqpoint{1.481602in}{3.144407in}}{\pgfqpoint{1.192000in}{3.144407in}}%
\pgfpathcurveto{\pgfqpoint{0.902398in}{3.144407in}}{\pgfqpoint{0.624619in}{3.086877in}}{\pgfqpoint{0.419839in}{2.984487in}}%
\pgfpathcurveto{\pgfqpoint{0.215060in}{2.882097in}}{\pgfqpoint{0.100000in}{2.743208in}}{\pgfqpoint{0.100000in}{2.598407in}}%
\pgfpathcurveto{\pgfqpoint{0.100000in}{2.453606in}}{\pgfqpoint{0.215060in}{2.314716in}}{\pgfqpoint{0.419839in}{2.212326in}}%
\pgfpathcurveto{\pgfqpoint{0.624619in}{2.109937in}}{\pgfqpoint{0.902398in}{2.052407in}}{\pgfqpoint{1.192000in}{2.052407in}}%
\pgfpathlineto{\pgfqpoint{1.192000in}{2.052407in}}%
\pgfpathclose%
\pgfusepath{clip}%
\pgfsetrectcap%
\pgfsetroundjoin%
\pgfsetlinewidth{0.451687pt}%
\definecolor{currentstroke}{rgb}{0.690196,0.690196,0.690196}%
\pgfsetstrokecolor{currentstroke}%
\pgfsetstrokeopacity{0.250000}%
\pgfsetdash{}{0pt}%
\pgfpathmoveto{\pgfqpoint{0.960953in}{2.103702in}}%
\pgfpathlineto{\pgfqpoint{0.941894in}{2.113059in}}%
\pgfpathlineto{\pgfqpoint{0.923620in}{2.122919in}}%
\pgfpathlineto{\pgfqpoint{0.906553in}{2.132965in}}%
\pgfpathlineto{\pgfqpoint{0.890254in}{2.143363in}}%
\pgfpathlineto{\pgfqpoint{0.874680in}{2.154082in}}%
\pgfpathlineto{\pgfqpoint{0.859791in}{2.165102in}}%
\pgfpathlineto{\pgfqpoint{0.845551in}{2.176401in}}%
\pgfpathlineto{\pgfqpoint{0.831929in}{2.187962in}}%
\pgfpathlineto{\pgfqpoint{0.818896in}{2.199772in}}%
\pgfpathlineto{\pgfqpoint{0.806430in}{2.211817in}}%
\pgfpathlineto{\pgfqpoint{0.794510in}{2.224083in}}%
\pgfpathlineto{\pgfqpoint{0.783119in}{2.236561in}}%
\pgfpathlineto{\pgfqpoint{0.772242in}{2.249238in}}%
\pgfpathlineto{\pgfqpoint{0.761864in}{2.262105in}}%
\pgfpathlineto{\pgfqpoint{0.751856in}{2.275314in}}%
\pgfpathlineto{\pgfqpoint{0.742483in}{2.288486in}}%
\pgfpathlineto{\pgfqpoint{0.733569in}{2.301831in}}%
\pgfpathlineto{\pgfqpoint{0.725108in}{2.315338in}}%
\pgfpathlineto{\pgfqpoint{0.717096in}{2.328996in}}%
\pgfpathlineto{\pgfqpoint{0.709529in}{2.342794in}}%
\pgfpathlineto{\pgfqpoint{0.702403in}{2.356724in}}%
\pgfpathlineto{\pgfqpoint{0.695713in}{2.370777in}}%
\pgfpathlineto{\pgfqpoint{0.689457in}{2.384944in}}%
\pgfpathlineto{\pgfqpoint{0.683631in}{2.399216in}}%
\pgfpathlineto{\pgfqpoint{0.678232in}{2.413586in}}%
\pgfpathlineto{\pgfqpoint{0.673258in}{2.428046in}}%
\pgfpathlineto{\pgfqpoint{0.668706in}{2.442588in}}%
\pgfpathlineto{\pgfqpoint{0.664574in}{2.457204in}}%
\pgfpathlineto{\pgfqpoint{0.660861in}{2.471888in}}%
\pgfpathlineto{\pgfqpoint{0.657563in}{2.486632in}}%
\pgfpathlineto{\pgfqpoint{0.654681in}{2.501429in}}%
\pgfpathlineto{\pgfqpoint{0.652188in}{2.516436in}}%
\pgfpathlineto{\pgfqpoint{0.650147in}{2.531243in}}%
\pgfpathlineto{\pgfqpoint{0.648510in}{2.546108in}}%
\pgfpathlineto{\pgfqpoint{0.647282in}{2.561018in}}%
\pgfpathlineto{\pgfqpoint{0.646462in}{2.575961in}}%
\pgfpathlineto{\pgfqpoint{0.646051in}{2.590923in}}%
\pgfpathlineto{\pgfqpoint{0.646051in}{2.605891in}}%
\pgfpathlineto{\pgfqpoint{0.646462in}{2.620853in}}%
\pgfpathlineto{\pgfqpoint{0.647282in}{2.635795in}}%
\pgfpathlineto{\pgfqpoint{0.648510in}{2.650705in}}%
\pgfpathlineto{\pgfqpoint{0.650147in}{2.665570in}}%
\pgfpathlineto{\pgfqpoint{0.652188in}{2.680378in}}%
\pgfpathlineto{\pgfqpoint{0.654681in}{2.695384in}}%
\pgfpathlineto{\pgfqpoint{0.657563in}{2.710181in}}%
\pgfpathlineto{\pgfqpoint{0.660861in}{2.724925in}}%
\pgfpathlineto{\pgfqpoint{0.664574in}{2.739609in}}%
\pgfpathlineto{\pgfqpoint{0.668706in}{2.754226in}}%
\pgfpathlineto{\pgfqpoint{0.673258in}{2.768768in}}%
\pgfpathlineto{\pgfqpoint{0.678232in}{2.783228in}}%
\pgfpathlineto{\pgfqpoint{0.683631in}{2.797598in}}%
\pgfpathlineto{\pgfqpoint{0.689457in}{2.811870in}}%
\pgfpathlineto{\pgfqpoint{0.695713in}{2.826037in}}%
\pgfpathlineto{\pgfqpoint{0.702403in}{2.840089in}}%
\pgfpathlineto{\pgfqpoint{0.709529in}{2.854019in}}%
\pgfpathlineto{\pgfqpoint{0.717096in}{2.867818in}}%
\pgfpathlineto{\pgfqpoint{0.725108in}{2.881476in}}%
\pgfpathlineto{\pgfqpoint{0.733569in}{2.894982in}}%
\pgfpathlineto{\pgfqpoint{0.742483in}{2.908327in}}%
\pgfpathlineto{\pgfqpoint{0.751856in}{2.921500in}}%
\pgfpathlineto{\pgfqpoint{0.761864in}{2.934708in}}%
\pgfpathlineto{\pgfqpoint{0.772242in}{2.947576in}}%
\pgfpathlineto{\pgfqpoint{0.783119in}{2.960253in}}%
\pgfpathlineto{\pgfqpoint{0.794510in}{2.972730in}}%
\pgfpathlineto{\pgfqpoint{0.806430in}{2.984997in}}%
\pgfpathlineto{\pgfqpoint{0.818896in}{2.997041in}}%
\pgfpathlineto{\pgfqpoint{0.831929in}{3.008851in}}%
\pgfpathlineto{\pgfqpoint{0.845551in}{3.020413in}}%
\pgfpathlineto{\pgfqpoint{0.859791in}{3.031712in}}%
\pgfpathlineto{\pgfqpoint{0.874680in}{3.042731in}}%
\pgfpathlineto{\pgfqpoint{0.890254in}{3.053451in}}%
\pgfpathlineto{\pgfqpoint{0.906553in}{3.063848in}}%
\pgfpathlineto{\pgfqpoint{0.923620in}{3.073894in}}%
\pgfpathlineto{\pgfqpoint{0.941894in}{3.083755in}}%
\pgfpathlineto{\pgfqpoint{0.960953in}{3.093112in}}%
\pgfusepath{stroke}%
\end{pgfscope}%
\begin{pgfscope}%
\pgfpathmoveto{\pgfqpoint{1.192000in}{2.052407in}}%
\pgfpathcurveto{\pgfqpoint{1.481602in}{2.052407in}}{\pgfqpoint{1.759381in}{2.109937in}}{\pgfqpoint{1.964161in}{2.212326in}}%
\pgfpathcurveto{\pgfqpoint{2.168940in}{2.314716in}}{\pgfqpoint{2.284000in}{2.453606in}}{\pgfqpoint{2.284000in}{2.598407in}}%
\pgfpathcurveto{\pgfqpoint{2.284000in}{2.743208in}}{\pgfqpoint{2.168940in}{2.882097in}}{\pgfqpoint{1.964161in}{2.984487in}}%
\pgfpathcurveto{\pgfqpoint{1.759381in}{3.086877in}}{\pgfqpoint{1.481602in}{3.144407in}}{\pgfqpoint{1.192000in}{3.144407in}}%
\pgfpathcurveto{\pgfqpoint{0.902398in}{3.144407in}}{\pgfqpoint{0.624619in}{3.086877in}}{\pgfqpoint{0.419839in}{2.984487in}}%
\pgfpathcurveto{\pgfqpoint{0.215060in}{2.882097in}}{\pgfqpoint{0.100000in}{2.743208in}}{\pgfqpoint{0.100000in}{2.598407in}}%
\pgfpathcurveto{\pgfqpoint{0.100000in}{2.453606in}}{\pgfqpoint{0.215060in}{2.314716in}}{\pgfqpoint{0.419839in}{2.212326in}}%
\pgfpathcurveto{\pgfqpoint{0.624619in}{2.109937in}}{\pgfqpoint{0.902398in}{2.052407in}}{\pgfqpoint{1.192000in}{2.052407in}}%
\pgfpathlineto{\pgfqpoint{1.192000in}{2.052407in}}%
\pgfpathclose%
\pgfusepath{clip}%
\pgfsetrectcap%
\pgfsetroundjoin%
\pgfsetlinewidth{0.451687pt}%
\definecolor{currentstroke}{rgb}{0.690196,0.690196,0.690196}%
\pgfsetstrokecolor{currentstroke}%
\pgfsetstrokeopacity{0.250000}%
\pgfsetdash{}{0pt}%
\pgfpathmoveto{\pgfqpoint{1.037969in}{2.103702in}}%
\pgfpathlineto{\pgfqpoint{1.025263in}{2.113059in}}%
\pgfpathlineto{\pgfqpoint{1.013080in}{2.122919in}}%
\pgfpathlineto{\pgfqpoint{1.001702in}{2.132965in}}%
\pgfpathlineto{\pgfqpoint{0.990836in}{2.143363in}}%
\pgfpathlineto{\pgfqpoint{0.980453in}{2.154082in}}%
\pgfpathlineto{\pgfqpoint{0.970527in}{2.165102in}}%
\pgfpathlineto{\pgfqpoint{0.961034in}{2.176401in}}%
\pgfpathlineto{\pgfqpoint{0.951952in}{2.187962in}}%
\pgfpathlineto{\pgfqpoint{0.943264in}{2.199772in}}%
\pgfpathlineto{\pgfqpoint{0.934953in}{2.211817in}}%
\pgfpathlineto{\pgfqpoint{0.927007in}{2.224083in}}%
\pgfpathlineto{\pgfqpoint{0.919413in}{2.236561in}}%
\pgfpathlineto{\pgfqpoint{0.912161in}{2.249238in}}%
\pgfpathlineto{\pgfqpoint{0.905243in}{2.262105in}}%
\pgfpathlineto{\pgfqpoint{0.898570in}{2.275314in}}%
\pgfpathlineto{\pgfqpoint{0.892322in}{2.288486in}}%
\pgfpathlineto{\pgfqpoint{0.886379in}{2.301831in}}%
\pgfpathlineto{\pgfqpoint{0.880739in}{2.315338in}}%
\pgfpathlineto{\pgfqpoint{0.875398in}{2.328996in}}%
\pgfpathlineto{\pgfqpoint{0.870353in}{2.342794in}}%
\pgfpathlineto{\pgfqpoint{0.865602in}{2.356724in}}%
\pgfpathlineto{\pgfqpoint{0.861142in}{2.370777in}}%
\pgfpathlineto{\pgfqpoint{0.856971in}{2.384944in}}%
\pgfpathlineto{\pgfqpoint{0.853087in}{2.399216in}}%
\pgfpathlineto{\pgfqpoint{0.849488in}{2.413586in}}%
\pgfpathlineto{\pgfqpoint{0.846172in}{2.428046in}}%
\pgfpathlineto{\pgfqpoint{0.843137in}{2.442588in}}%
\pgfpathlineto{\pgfqpoint{0.840383in}{2.457204in}}%
\pgfpathlineto{\pgfqpoint{0.837907in}{2.471888in}}%
\pgfpathlineto{\pgfqpoint{0.835709in}{2.486632in}}%
\pgfpathlineto{\pgfqpoint{0.833788in}{2.501429in}}%
\pgfpathlineto{\pgfqpoint{0.832125in}{2.516436in}}%
\pgfpathlineto{\pgfqpoint{0.830764in}{2.531243in}}%
\pgfpathlineto{\pgfqpoint{0.829674in}{2.546108in}}%
\pgfpathlineto{\pgfqpoint{0.828854in}{2.561018in}}%
\pgfpathlineto{\pgfqpoint{0.828308in}{2.575961in}}%
\pgfpathlineto{\pgfqpoint{0.828034in}{2.590923in}}%
\pgfpathlineto{\pgfqpoint{0.828034in}{2.605891in}}%
\pgfpathlineto{\pgfqpoint{0.828308in}{2.620853in}}%
\pgfpathlineto{\pgfqpoint{0.828854in}{2.635795in}}%
\pgfpathlineto{\pgfqpoint{0.829674in}{2.650705in}}%
\pgfpathlineto{\pgfqpoint{0.830764in}{2.665570in}}%
\pgfpathlineto{\pgfqpoint{0.832125in}{2.680378in}}%
\pgfpathlineto{\pgfqpoint{0.833788in}{2.695384in}}%
\pgfpathlineto{\pgfqpoint{0.835709in}{2.710181in}}%
\pgfpathlineto{\pgfqpoint{0.837907in}{2.724925in}}%
\pgfpathlineto{\pgfqpoint{0.840383in}{2.739609in}}%
\pgfpathlineto{\pgfqpoint{0.843137in}{2.754226in}}%
\pgfpathlineto{\pgfqpoint{0.846172in}{2.768768in}}%
\pgfpathlineto{\pgfqpoint{0.849488in}{2.783228in}}%
\pgfpathlineto{\pgfqpoint{0.853087in}{2.797598in}}%
\pgfpathlineto{\pgfqpoint{0.856971in}{2.811870in}}%
\pgfpathlineto{\pgfqpoint{0.861142in}{2.826037in}}%
\pgfpathlineto{\pgfqpoint{0.865602in}{2.840089in}}%
\pgfpathlineto{\pgfqpoint{0.870353in}{2.854019in}}%
\pgfpathlineto{\pgfqpoint{0.875398in}{2.867818in}}%
\pgfpathlineto{\pgfqpoint{0.880739in}{2.881476in}}%
\pgfpathlineto{\pgfqpoint{0.886379in}{2.894982in}}%
\pgfpathlineto{\pgfqpoint{0.892322in}{2.908327in}}%
\pgfpathlineto{\pgfqpoint{0.898570in}{2.921500in}}%
\pgfpathlineto{\pgfqpoint{0.905243in}{2.934708in}}%
\pgfpathlineto{\pgfqpoint{0.912161in}{2.947576in}}%
\pgfpathlineto{\pgfqpoint{0.919413in}{2.960253in}}%
\pgfpathlineto{\pgfqpoint{0.927007in}{2.972730in}}%
\pgfpathlineto{\pgfqpoint{0.934953in}{2.984997in}}%
\pgfpathlineto{\pgfqpoint{0.943264in}{2.997041in}}%
\pgfpathlineto{\pgfqpoint{0.951952in}{3.008851in}}%
\pgfpathlineto{\pgfqpoint{0.961034in}{3.020413in}}%
\pgfpathlineto{\pgfqpoint{0.970527in}{3.031712in}}%
\pgfpathlineto{\pgfqpoint{0.980453in}{3.042731in}}%
\pgfpathlineto{\pgfqpoint{0.990836in}{3.053451in}}%
\pgfpathlineto{\pgfqpoint{1.001702in}{3.063848in}}%
\pgfpathlineto{\pgfqpoint{1.013080in}{3.073894in}}%
\pgfpathlineto{\pgfqpoint{1.025263in}{3.083755in}}%
\pgfpathlineto{\pgfqpoint{1.037969in}{3.093112in}}%
\pgfusepath{stroke}%
\end{pgfscope}%
\begin{pgfscope}%
\pgfpathmoveto{\pgfqpoint{1.192000in}{2.052407in}}%
\pgfpathcurveto{\pgfqpoint{1.481602in}{2.052407in}}{\pgfqpoint{1.759381in}{2.109937in}}{\pgfqpoint{1.964161in}{2.212326in}}%
\pgfpathcurveto{\pgfqpoint{2.168940in}{2.314716in}}{\pgfqpoint{2.284000in}{2.453606in}}{\pgfqpoint{2.284000in}{2.598407in}}%
\pgfpathcurveto{\pgfqpoint{2.284000in}{2.743208in}}{\pgfqpoint{2.168940in}{2.882097in}}{\pgfqpoint{1.964161in}{2.984487in}}%
\pgfpathcurveto{\pgfqpoint{1.759381in}{3.086877in}}{\pgfqpoint{1.481602in}{3.144407in}}{\pgfqpoint{1.192000in}{3.144407in}}%
\pgfpathcurveto{\pgfqpoint{0.902398in}{3.144407in}}{\pgfqpoint{0.624619in}{3.086877in}}{\pgfqpoint{0.419839in}{2.984487in}}%
\pgfpathcurveto{\pgfqpoint{0.215060in}{2.882097in}}{\pgfqpoint{0.100000in}{2.743208in}}{\pgfqpoint{0.100000in}{2.598407in}}%
\pgfpathcurveto{\pgfqpoint{0.100000in}{2.453606in}}{\pgfqpoint{0.215060in}{2.314716in}}{\pgfqpoint{0.419839in}{2.212326in}}%
\pgfpathcurveto{\pgfqpoint{0.624619in}{2.109937in}}{\pgfqpoint{0.902398in}{2.052407in}}{\pgfqpoint{1.192000in}{2.052407in}}%
\pgfpathlineto{\pgfqpoint{1.192000in}{2.052407in}}%
\pgfpathclose%
\pgfusepath{clip}%
\pgfsetrectcap%
\pgfsetroundjoin%
\pgfsetlinewidth{0.451687pt}%
\definecolor{currentstroke}{rgb}{0.690196,0.690196,0.690196}%
\pgfsetstrokecolor{currentstroke}%
\pgfsetstrokeopacity{0.250000}%
\pgfsetdash{}{0pt}%
\pgfpathmoveto{\pgfqpoint{1.114984in}{2.103702in}}%
\pgfpathlineto{\pgfqpoint{1.108631in}{2.113059in}}%
\pgfpathlineto{\pgfqpoint{1.102540in}{2.122919in}}%
\pgfpathlineto{\pgfqpoint{1.096851in}{2.132965in}}%
\pgfpathlineto{\pgfqpoint{1.091418in}{2.143363in}}%
\pgfpathlineto{\pgfqpoint{1.086227in}{2.154082in}}%
\pgfpathlineto{\pgfqpoint{1.081264in}{2.165102in}}%
\pgfpathlineto{\pgfqpoint{1.076517in}{2.176401in}}%
\pgfpathlineto{\pgfqpoint{1.071976in}{2.187962in}}%
\pgfpathlineto{\pgfqpoint{1.067632in}{2.199772in}}%
\pgfpathlineto{\pgfqpoint{1.063477in}{2.211817in}}%
\pgfpathlineto{\pgfqpoint{1.059503in}{2.224083in}}%
\pgfpathlineto{\pgfqpoint{1.055706in}{2.236561in}}%
\pgfpathlineto{\pgfqpoint{1.052081in}{2.249238in}}%
\pgfpathlineto{\pgfqpoint{1.048621in}{2.262105in}}%
\pgfpathlineto{\pgfqpoint{1.045285in}{2.275314in}}%
\pgfpathlineto{\pgfqpoint{1.042161in}{2.288486in}}%
\pgfpathlineto{\pgfqpoint{1.039190in}{2.301831in}}%
\pgfpathlineto{\pgfqpoint{1.036369in}{2.315338in}}%
\pgfpathlineto{\pgfqpoint{1.033699in}{2.328996in}}%
\pgfpathlineto{\pgfqpoint{1.031176in}{2.342794in}}%
\pgfpathlineto{\pgfqpoint{1.028801in}{2.356724in}}%
\pgfpathlineto{\pgfqpoint{1.026571in}{2.370777in}}%
\pgfpathlineto{\pgfqpoint{1.024486in}{2.384944in}}%
\pgfpathlineto{\pgfqpoint{1.022544in}{2.399216in}}%
\pgfpathlineto{\pgfqpoint{1.020744in}{2.413586in}}%
\pgfpathlineto{\pgfqpoint{1.019086in}{2.428046in}}%
\pgfpathlineto{\pgfqpoint{1.017569in}{2.442588in}}%
\pgfpathlineto{\pgfqpoint{1.016191in}{2.457204in}}%
\pgfpathlineto{\pgfqpoint{1.014954in}{2.471888in}}%
\pgfpathlineto{\pgfqpoint{1.013854in}{2.486632in}}%
\pgfpathlineto{\pgfqpoint{1.012894in}{2.501429in}}%
\pgfpathlineto{\pgfqpoint{1.012063in}{2.516436in}}%
\pgfpathlineto{\pgfqpoint{1.011382in}{2.531243in}}%
\pgfpathlineto{\pgfqpoint{1.010837in}{2.546108in}}%
\pgfpathlineto{\pgfqpoint{1.010427in}{2.561018in}}%
\pgfpathlineto{\pgfqpoint{1.010154in}{2.575961in}}%
\pgfpathlineto{\pgfqpoint{1.010017in}{2.590923in}}%
\pgfpathlineto{\pgfqpoint{1.010017in}{2.605891in}}%
\pgfpathlineto{\pgfqpoint{1.010154in}{2.620853in}}%
\pgfpathlineto{\pgfqpoint{1.010427in}{2.635795in}}%
\pgfpathlineto{\pgfqpoint{1.010837in}{2.650705in}}%
\pgfpathlineto{\pgfqpoint{1.011382in}{2.665570in}}%
\pgfpathlineto{\pgfqpoint{1.012063in}{2.680378in}}%
\pgfpathlineto{\pgfqpoint{1.012894in}{2.695384in}}%
\pgfpathlineto{\pgfqpoint{1.013854in}{2.710181in}}%
\pgfpathlineto{\pgfqpoint{1.014954in}{2.724925in}}%
\pgfpathlineto{\pgfqpoint{1.016191in}{2.739609in}}%
\pgfpathlineto{\pgfqpoint{1.017569in}{2.754226in}}%
\pgfpathlineto{\pgfqpoint{1.019086in}{2.768768in}}%
\pgfpathlineto{\pgfqpoint{1.020744in}{2.783228in}}%
\pgfpathlineto{\pgfqpoint{1.022544in}{2.797598in}}%
\pgfpathlineto{\pgfqpoint{1.024486in}{2.811870in}}%
\pgfpathlineto{\pgfqpoint{1.026571in}{2.826037in}}%
\pgfpathlineto{\pgfqpoint{1.028801in}{2.840089in}}%
\pgfpathlineto{\pgfqpoint{1.031176in}{2.854019in}}%
\pgfpathlineto{\pgfqpoint{1.033699in}{2.867818in}}%
\pgfpathlineto{\pgfqpoint{1.036369in}{2.881476in}}%
\pgfpathlineto{\pgfqpoint{1.039190in}{2.894982in}}%
\pgfpathlineto{\pgfqpoint{1.042161in}{2.908327in}}%
\pgfpathlineto{\pgfqpoint{1.045285in}{2.921500in}}%
\pgfpathlineto{\pgfqpoint{1.048621in}{2.934708in}}%
\pgfpathlineto{\pgfqpoint{1.052081in}{2.947576in}}%
\pgfpathlineto{\pgfqpoint{1.055706in}{2.960253in}}%
\pgfpathlineto{\pgfqpoint{1.059503in}{2.972730in}}%
\pgfpathlineto{\pgfqpoint{1.063477in}{2.984997in}}%
\pgfpathlineto{\pgfqpoint{1.067632in}{2.997041in}}%
\pgfpathlineto{\pgfqpoint{1.071976in}{3.008851in}}%
\pgfpathlineto{\pgfqpoint{1.076517in}{3.020413in}}%
\pgfpathlineto{\pgfqpoint{1.081264in}{3.031712in}}%
\pgfpathlineto{\pgfqpoint{1.086227in}{3.042731in}}%
\pgfpathlineto{\pgfqpoint{1.091418in}{3.053451in}}%
\pgfpathlineto{\pgfqpoint{1.096851in}{3.063848in}}%
\pgfpathlineto{\pgfqpoint{1.102540in}{3.073894in}}%
\pgfpathlineto{\pgfqpoint{1.108631in}{3.083755in}}%
\pgfpathlineto{\pgfqpoint{1.114984in}{3.093112in}}%
\pgfusepath{stroke}%
\end{pgfscope}%
\begin{pgfscope}%
\pgfpathmoveto{\pgfqpoint{1.192000in}{2.052407in}}%
\pgfpathcurveto{\pgfqpoint{1.481602in}{2.052407in}}{\pgfqpoint{1.759381in}{2.109937in}}{\pgfqpoint{1.964161in}{2.212326in}}%
\pgfpathcurveto{\pgfqpoint{2.168940in}{2.314716in}}{\pgfqpoint{2.284000in}{2.453606in}}{\pgfqpoint{2.284000in}{2.598407in}}%
\pgfpathcurveto{\pgfqpoint{2.284000in}{2.743208in}}{\pgfqpoint{2.168940in}{2.882097in}}{\pgfqpoint{1.964161in}{2.984487in}}%
\pgfpathcurveto{\pgfqpoint{1.759381in}{3.086877in}}{\pgfqpoint{1.481602in}{3.144407in}}{\pgfqpoint{1.192000in}{3.144407in}}%
\pgfpathcurveto{\pgfqpoint{0.902398in}{3.144407in}}{\pgfqpoint{0.624619in}{3.086877in}}{\pgfqpoint{0.419839in}{2.984487in}}%
\pgfpathcurveto{\pgfqpoint{0.215060in}{2.882097in}}{\pgfqpoint{0.100000in}{2.743208in}}{\pgfqpoint{0.100000in}{2.598407in}}%
\pgfpathcurveto{\pgfqpoint{0.100000in}{2.453606in}}{\pgfqpoint{0.215060in}{2.314716in}}{\pgfqpoint{0.419839in}{2.212326in}}%
\pgfpathcurveto{\pgfqpoint{0.624619in}{2.109937in}}{\pgfqpoint{0.902398in}{2.052407in}}{\pgfqpoint{1.192000in}{2.052407in}}%
\pgfpathlineto{\pgfqpoint{1.192000in}{2.052407in}}%
\pgfpathclose%
\pgfusepath{clip}%
\pgfsetrectcap%
\pgfsetroundjoin%
\pgfsetlinewidth{0.451687pt}%
\definecolor{currentstroke}{rgb}{0.690196,0.690196,0.690196}%
\pgfsetstrokecolor{currentstroke}%
\pgfsetstrokeopacity{0.250000}%
\pgfsetdash{}{0pt}%
\pgfpathmoveto{\pgfqpoint{1.192000in}{2.103702in}}%
\pgfpathlineto{\pgfqpoint{1.192000in}{2.113059in}}%
\pgfpathlineto{\pgfqpoint{1.192000in}{2.122919in}}%
\pgfpathlineto{\pgfqpoint{1.192000in}{2.132965in}}%
\pgfpathlineto{\pgfqpoint{1.192000in}{2.143363in}}%
\pgfpathlineto{\pgfqpoint{1.192000in}{2.154082in}}%
\pgfpathlineto{\pgfqpoint{1.192000in}{2.165102in}}%
\pgfpathlineto{\pgfqpoint{1.192000in}{2.176401in}}%
\pgfpathlineto{\pgfqpoint{1.192000in}{2.187962in}}%
\pgfpathlineto{\pgfqpoint{1.192000in}{2.199772in}}%
\pgfpathlineto{\pgfqpoint{1.192000in}{2.211817in}}%
\pgfpathlineto{\pgfqpoint{1.192000in}{2.224083in}}%
\pgfpathlineto{\pgfqpoint{1.192000in}{2.236561in}}%
\pgfpathlineto{\pgfqpoint{1.192000in}{2.249238in}}%
\pgfpathlineto{\pgfqpoint{1.192000in}{2.262105in}}%
\pgfpathlineto{\pgfqpoint{1.192000in}{2.275314in}}%
\pgfpathlineto{\pgfqpoint{1.192000in}{2.288486in}}%
\pgfpathlineto{\pgfqpoint{1.192000in}{2.301831in}}%
\pgfpathlineto{\pgfqpoint{1.192000in}{2.315338in}}%
\pgfpathlineto{\pgfqpoint{1.192000in}{2.328996in}}%
\pgfpathlineto{\pgfqpoint{1.192000in}{2.342794in}}%
\pgfpathlineto{\pgfqpoint{1.192000in}{2.356724in}}%
\pgfpathlineto{\pgfqpoint{1.192000in}{2.370777in}}%
\pgfpathlineto{\pgfqpoint{1.192000in}{2.384944in}}%
\pgfpathlineto{\pgfqpoint{1.192000in}{2.399216in}}%
\pgfpathlineto{\pgfqpoint{1.192000in}{2.413586in}}%
\pgfpathlineto{\pgfqpoint{1.192000in}{2.428046in}}%
\pgfpathlineto{\pgfqpoint{1.192000in}{2.442588in}}%
\pgfpathlineto{\pgfqpoint{1.192000in}{2.457204in}}%
\pgfpathlineto{\pgfqpoint{1.192000in}{2.471888in}}%
\pgfpathlineto{\pgfqpoint{1.192000in}{2.486632in}}%
\pgfpathlineto{\pgfqpoint{1.192000in}{2.501429in}}%
\pgfpathlineto{\pgfqpoint{1.192000in}{2.516436in}}%
\pgfpathlineto{\pgfqpoint{1.192000in}{2.531243in}}%
\pgfpathlineto{\pgfqpoint{1.192000in}{2.546108in}}%
\pgfpathlineto{\pgfqpoint{1.192000in}{2.561018in}}%
\pgfpathlineto{\pgfqpoint{1.192000in}{2.575961in}}%
\pgfpathlineto{\pgfqpoint{1.192000in}{2.590923in}}%
\pgfpathlineto{\pgfqpoint{1.192000in}{2.605891in}}%
\pgfpathlineto{\pgfqpoint{1.192000in}{2.620853in}}%
\pgfpathlineto{\pgfqpoint{1.192000in}{2.635795in}}%
\pgfpathlineto{\pgfqpoint{1.192000in}{2.650705in}}%
\pgfpathlineto{\pgfqpoint{1.192000in}{2.665570in}}%
\pgfpathlineto{\pgfqpoint{1.192000in}{2.680378in}}%
\pgfpathlineto{\pgfqpoint{1.192000in}{2.695384in}}%
\pgfpathlineto{\pgfqpoint{1.192000in}{2.710181in}}%
\pgfpathlineto{\pgfqpoint{1.192000in}{2.724925in}}%
\pgfpathlineto{\pgfqpoint{1.192000in}{2.739609in}}%
\pgfpathlineto{\pgfqpoint{1.192000in}{2.754226in}}%
\pgfpathlineto{\pgfqpoint{1.192000in}{2.768768in}}%
\pgfpathlineto{\pgfqpoint{1.192000in}{2.783228in}}%
\pgfpathlineto{\pgfqpoint{1.192000in}{2.797598in}}%
\pgfpathlineto{\pgfqpoint{1.192000in}{2.811870in}}%
\pgfpathlineto{\pgfqpoint{1.192000in}{2.826037in}}%
\pgfpathlineto{\pgfqpoint{1.192000in}{2.840089in}}%
\pgfpathlineto{\pgfqpoint{1.192000in}{2.854019in}}%
\pgfpathlineto{\pgfqpoint{1.192000in}{2.867818in}}%
\pgfpathlineto{\pgfqpoint{1.192000in}{2.881476in}}%
\pgfpathlineto{\pgfqpoint{1.192000in}{2.894982in}}%
\pgfpathlineto{\pgfqpoint{1.192000in}{2.908327in}}%
\pgfpathlineto{\pgfqpoint{1.192000in}{2.921500in}}%
\pgfpathlineto{\pgfqpoint{1.192000in}{2.934708in}}%
\pgfpathlineto{\pgfqpoint{1.192000in}{2.947576in}}%
\pgfpathlineto{\pgfqpoint{1.192000in}{2.960253in}}%
\pgfpathlineto{\pgfqpoint{1.192000in}{2.972730in}}%
\pgfpathlineto{\pgfqpoint{1.192000in}{2.984997in}}%
\pgfpathlineto{\pgfqpoint{1.192000in}{2.997041in}}%
\pgfpathlineto{\pgfqpoint{1.192000in}{3.008851in}}%
\pgfpathlineto{\pgfqpoint{1.192000in}{3.020413in}}%
\pgfpathlineto{\pgfqpoint{1.192000in}{3.031712in}}%
\pgfpathlineto{\pgfqpoint{1.192000in}{3.042731in}}%
\pgfpathlineto{\pgfqpoint{1.192000in}{3.053451in}}%
\pgfpathlineto{\pgfqpoint{1.192000in}{3.063848in}}%
\pgfpathlineto{\pgfqpoint{1.192000in}{3.073894in}}%
\pgfpathlineto{\pgfqpoint{1.192000in}{3.083755in}}%
\pgfpathlineto{\pgfqpoint{1.192000in}{3.093112in}}%
\pgfusepath{stroke}%
\end{pgfscope}%
\begin{pgfscope}%
\pgfpathmoveto{\pgfqpoint{1.192000in}{2.052407in}}%
\pgfpathcurveto{\pgfqpoint{1.481602in}{2.052407in}}{\pgfqpoint{1.759381in}{2.109937in}}{\pgfqpoint{1.964161in}{2.212326in}}%
\pgfpathcurveto{\pgfqpoint{2.168940in}{2.314716in}}{\pgfqpoint{2.284000in}{2.453606in}}{\pgfqpoint{2.284000in}{2.598407in}}%
\pgfpathcurveto{\pgfqpoint{2.284000in}{2.743208in}}{\pgfqpoint{2.168940in}{2.882097in}}{\pgfqpoint{1.964161in}{2.984487in}}%
\pgfpathcurveto{\pgfqpoint{1.759381in}{3.086877in}}{\pgfqpoint{1.481602in}{3.144407in}}{\pgfqpoint{1.192000in}{3.144407in}}%
\pgfpathcurveto{\pgfqpoint{0.902398in}{3.144407in}}{\pgfqpoint{0.624619in}{3.086877in}}{\pgfqpoint{0.419839in}{2.984487in}}%
\pgfpathcurveto{\pgfqpoint{0.215060in}{2.882097in}}{\pgfqpoint{0.100000in}{2.743208in}}{\pgfqpoint{0.100000in}{2.598407in}}%
\pgfpathcurveto{\pgfqpoint{0.100000in}{2.453606in}}{\pgfqpoint{0.215060in}{2.314716in}}{\pgfqpoint{0.419839in}{2.212326in}}%
\pgfpathcurveto{\pgfqpoint{0.624619in}{2.109937in}}{\pgfqpoint{0.902398in}{2.052407in}}{\pgfqpoint{1.192000in}{2.052407in}}%
\pgfpathlineto{\pgfqpoint{1.192000in}{2.052407in}}%
\pgfpathclose%
\pgfusepath{clip}%
\pgfsetrectcap%
\pgfsetroundjoin%
\pgfsetlinewidth{0.451687pt}%
\definecolor{currentstroke}{rgb}{0.690196,0.690196,0.690196}%
\pgfsetstrokecolor{currentstroke}%
\pgfsetstrokeopacity{0.250000}%
\pgfsetdash{}{0pt}%
\pgfpathmoveto{\pgfqpoint{1.269016in}{2.103702in}}%
\pgfpathlineto{\pgfqpoint{1.275369in}{2.113059in}}%
\pgfpathlineto{\pgfqpoint{1.281460in}{2.122919in}}%
\pgfpathlineto{\pgfqpoint{1.287149in}{2.132965in}}%
\pgfpathlineto{\pgfqpoint{1.292582in}{2.143363in}}%
\pgfpathlineto{\pgfqpoint{1.297773in}{2.154082in}}%
\pgfpathlineto{\pgfqpoint{1.302736in}{2.165102in}}%
\pgfpathlineto{\pgfqpoint{1.307483in}{2.176401in}}%
\pgfpathlineto{\pgfqpoint{1.312024in}{2.187962in}}%
\pgfpathlineto{\pgfqpoint{1.316368in}{2.199772in}}%
\pgfpathlineto{\pgfqpoint{1.320523in}{2.211817in}}%
\pgfpathlineto{\pgfqpoint{1.324497in}{2.224083in}}%
\pgfpathlineto{\pgfqpoint{1.328294in}{2.236561in}}%
\pgfpathlineto{\pgfqpoint{1.331919in}{2.249238in}}%
\pgfpathlineto{\pgfqpoint{1.335379in}{2.262105in}}%
\pgfpathlineto{\pgfqpoint{1.338715in}{2.275314in}}%
\pgfpathlineto{\pgfqpoint{1.341839in}{2.288486in}}%
\pgfpathlineto{\pgfqpoint{1.344810in}{2.301831in}}%
\pgfpathlineto{\pgfqpoint{1.347631in}{2.315338in}}%
\pgfpathlineto{\pgfqpoint{1.350301in}{2.328996in}}%
\pgfpathlineto{\pgfqpoint{1.352824in}{2.342794in}}%
\pgfpathlineto{\pgfqpoint{1.355199in}{2.356724in}}%
\pgfpathlineto{\pgfqpoint{1.357429in}{2.370777in}}%
\pgfpathlineto{\pgfqpoint{1.359514in}{2.384944in}}%
\pgfpathlineto{\pgfqpoint{1.361456in}{2.399216in}}%
\pgfpathlineto{\pgfqpoint{1.363256in}{2.413586in}}%
\pgfpathlineto{\pgfqpoint{1.364914in}{2.428046in}}%
\pgfpathlineto{\pgfqpoint{1.366431in}{2.442588in}}%
\pgfpathlineto{\pgfqpoint{1.367809in}{2.457204in}}%
\pgfpathlineto{\pgfqpoint{1.369046in}{2.471888in}}%
\pgfpathlineto{\pgfqpoint{1.370146in}{2.486632in}}%
\pgfpathlineto{\pgfqpoint{1.371106in}{2.501429in}}%
\pgfpathlineto{\pgfqpoint{1.371937in}{2.516436in}}%
\pgfpathlineto{\pgfqpoint{1.372618in}{2.531243in}}%
\pgfpathlineto{\pgfqpoint{1.373163in}{2.546108in}}%
\pgfpathlineto{\pgfqpoint{1.373573in}{2.561018in}}%
\pgfpathlineto{\pgfqpoint{1.373846in}{2.575961in}}%
\pgfpathlineto{\pgfqpoint{1.373983in}{2.590923in}}%
\pgfpathlineto{\pgfqpoint{1.373983in}{2.605891in}}%
\pgfpathlineto{\pgfqpoint{1.373846in}{2.620853in}}%
\pgfpathlineto{\pgfqpoint{1.373573in}{2.635795in}}%
\pgfpathlineto{\pgfqpoint{1.373163in}{2.650705in}}%
\pgfpathlineto{\pgfqpoint{1.372618in}{2.665570in}}%
\pgfpathlineto{\pgfqpoint{1.371937in}{2.680378in}}%
\pgfpathlineto{\pgfqpoint{1.371106in}{2.695384in}}%
\pgfpathlineto{\pgfqpoint{1.370146in}{2.710181in}}%
\pgfpathlineto{\pgfqpoint{1.369046in}{2.724925in}}%
\pgfpathlineto{\pgfqpoint{1.367809in}{2.739609in}}%
\pgfpathlineto{\pgfqpoint{1.366431in}{2.754226in}}%
\pgfpathlineto{\pgfqpoint{1.364914in}{2.768768in}}%
\pgfpathlineto{\pgfqpoint{1.363256in}{2.783228in}}%
\pgfpathlineto{\pgfqpoint{1.361456in}{2.797598in}}%
\pgfpathlineto{\pgfqpoint{1.359514in}{2.811870in}}%
\pgfpathlineto{\pgfqpoint{1.357429in}{2.826037in}}%
\pgfpathlineto{\pgfqpoint{1.355199in}{2.840089in}}%
\pgfpathlineto{\pgfqpoint{1.352824in}{2.854019in}}%
\pgfpathlineto{\pgfqpoint{1.350301in}{2.867818in}}%
\pgfpathlineto{\pgfqpoint{1.347631in}{2.881476in}}%
\pgfpathlineto{\pgfqpoint{1.344810in}{2.894982in}}%
\pgfpathlineto{\pgfqpoint{1.341839in}{2.908327in}}%
\pgfpathlineto{\pgfqpoint{1.338715in}{2.921500in}}%
\pgfpathlineto{\pgfqpoint{1.335379in}{2.934708in}}%
\pgfpathlineto{\pgfqpoint{1.331919in}{2.947576in}}%
\pgfpathlineto{\pgfqpoint{1.328294in}{2.960253in}}%
\pgfpathlineto{\pgfqpoint{1.324497in}{2.972730in}}%
\pgfpathlineto{\pgfqpoint{1.320523in}{2.984997in}}%
\pgfpathlineto{\pgfqpoint{1.316368in}{2.997041in}}%
\pgfpathlineto{\pgfqpoint{1.312024in}{3.008851in}}%
\pgfpathlineto{\pgfqpoint{1.307483in}{3.020413in}}%
\pgfpathlineto{\pgfqpoint{1.302736in}{3.031712in}}%
\pgfpathlineto{\pgfqpoint{1.297773in}{3.042731in}}%
\pgfpathlineto{\pgfqpoint{1.292582in}{3.053451in}}%
\pgfpathlineto{\pgfqpoint{1.287149in}{3.063848in}}%
\pgfpathlineto{\pgfqpoint{1.281460in}{3.073894in}}%
\pgfpathlineto{\pgfqpoint{1.275369in}{3.083755in}}%
\pgfpathlineto{\pgfqpoint{1.269016in}{3.093112in}}%
\pgfusepath{stroke}%
\end{pgfscope}%
\begin{pgfscope}%
\pgfpathmoveto{\pgfqpoint{1.192000in}{2.052407in}}%
\pgfpathcurveto{\pgfqpoint{1.481602in}{2.052407in}}{\pgfqpoint{1.759381in}{2.109937in}}{\pgfqpoint{1.964161in}{2.212326in}}%
\pgfpathcurveto{\pgfqpoint{2.168940in}{2.314716in}}{\pgfqpoint{2.284000in}{2.453606in}}{\pgfqpoint{2.284000in}{2.598407in}}%
\pgfpathcurveto{\pgfqpoint{2.284000in}{2.743208in}}{\pgfqpoint{2.168940in}{2.882097in}}{\pgfqpoint{1.964161in}{2.984487in}}%
\pgfpathcurveto{\pgfqpoint{1.759381in}{3.086877in}}{\pgfqpoint{1.481602in}{3.144407in}}{\pgfqpoint{1.192000in}{3.144407in}}%
\pgfpathcurveto{\pgfqpoint{0.902398in}{3.144407in}}{\pgfqpoint{0.624619in}{3.086877in}}{\pgfqpoint{0.419839in}{2.984487in}}%
\pgfpathcurveto{\pgfqpoint{0.215060in}{2.882097in}}{\pgfqpoint{0.100000in}{2.743208in}}{\pgfqpoint{0.100000in}{2.598407in}}%
\pgfpathcurveto{\pgfqpoint{0.100000in}{2.453606in}}{\pgfqpoint{0.215060in}{2.314716in}}{\pgfqpoint{0.419839in}{2.212326in}}%
\pgfpathcurveto{\pgfqpoint{0.624619in}{2.109937in}}{\pgfqpoint{0.902398in}{2.052407in}}{\pgfqpoint{1.192000in}{2.052407in}}%
\pgfpathlineto{\pgfqpoint{1.192000in}{2.052407in}}%
\pgfpathclose%
\pgfusepath{clip}%
\pgfsetrectcap%
\pgfsetroundjoin%
\pgfsetlinewidth{0.451687pt}%
\definecolor{currentstroke}{rgb}{0.690196,0.690196,0.690196}%
\pgfsetstrokecolor{currentstroke}%
\pgfsetstrokeopacity{0.250000}%
\pgfsetdash{}{0pt}%
\pgfpathmoveto{\pgfqpoint{1.346031in}{2.103702in}}%
\pgfpathlineto{\pgfqpoint{1.358737in}{2.113059in}}%
\pgfpathlineto{\pgfqpoint{1.370920in}{2.122919in}}%
\pgfpathlineto{\pgfqpoint{1.382298in}{2.132965in}}%
\pgfpathlineto{\pgfqpoint{1.393164in}{2.143363in}}%
\pgfpathlineto{\pgfqpoint{1.403547in}{2.154082in}}%
\pgfpathlineto{\pgfqpoint{1.413473in}{2.165102in}}%
\pgfpathlineto{\pgfqpoint{1.422966in}{2.176401in}}%
\pgfpathlineto{\pgfqpoint{1.432048in}{2.187962in}}%
\pgfpathlineto{\pgfqpoint{1.440736in}{2.199772in}}%
\pgfpathlineto{\pgfqpoint{1.449047in}{2.211817in}}%
\pgfpathlineto{\pgfqpoint{1.456993in}{2.224083in}}%
\pgfpathlineto{\pgfqpoint{1.464587in}{2.236561in}}%
\pgfpathlineto{\pgfqpoint{1.471839in}{2.249238in}}%
\pgfpathlineto{\pgfqpoint{1.478757in}{2.262105in}}%
\pgfpathlineto{\pgfqpoint{1.485430in}{2.275314in}}%
\pgfpathlineto{\pgfqpoint{1.491678in}{2.288486in}}%
\pgfpathlineto{\pgfqpoint{1.497621in}{2.301831in}}%
\pgfpathlineto{\pgfqpoint{1.503261in}{2.315338in}}%
\pgfpathlineto{\pgfqpoint{1.508602in}{2.328996in}}%
\pgfpathlineto{\pgfqpoint{1.513647in}{2.342794in}}%
\pgfpathlineto{\pgfqpoint{1.518398in}{2.356724in}}%
\pgfpathlineto{\pgfqpoint{1.522858in}{2.370777in}}%
\pgfpathlineto{\pgfqpoint{1.527029in}{2.384944in}}%
\pgfpathlineto{\pgfqpoint{1.530913in}{2.399216in}}%
\pgfpathlineto{\pgfqpoint{1.534512in}{2.413586in}}%
\pgfpathlineto{\pgfqpoint{1.537828in}{2.428046in}}%
\pgfpathlineto{\pgfqpoint{1.540863in}{2.442588in}}%
\pgfpathlineto{\pgfqpoint{1.543617in}{2.457204in}}%
\pgfpathlineto{\pgfqpoint{1.546093in}{2.471888in}}%
\pgfpathlineto{\pgfqpoint{1.548291in}{2.486632in}}%
\pgfpathlineto{\pgfqpoint{1.550212in}{2.501429in}}%
\pgfpathlineto{\pgfqpoint{1.551875in}{2.516436in}}%
\pgfpathlineto{\pgfqpoint{1.553236in}{2.531243in}}%
\pgfpathlineto{\pgfqpoint{1.554326in}{2.546108in}}%
\pgfpathlineto{\pgfqpoint{1.555146in}{2.561018in}}%
\pgfpathlineto{\pgfqpoint{1.555692in}{2.575961in}}%
\pgfpathlineto{\pgfqpoint{1.555966in}{2.590923in}}%
\pgfpathlineto{\pgfqpoint{1.555966in}{2.605891in}}%
\pgfpathlineto{\pgfqpoint{1.555692in}{2.620853in}}%
\pgfpathlineto{\pgfqpoint{1.555146in}{2.635795in}}%
\pgfpathlineto{\pgfqpoint{1.554326in}{2.650705in}}%
\pgfpathlineto{\pgfqpoint{1.553236in}{2.665570in}}%
\pgfpathlineto{\pgfqpoint{1.551875in}{2.680378in}}%
\pgfpathlineto{\pgfqpoint{1.550212in}{2.695384in}}%
\pgfpathlineto{\pgfqpoint{1.548291in}{2.710181in}}%
\pgfpathlineto{\pgfqpoint{1.546093in}{2.724925in}}%
\pgfpathlineto{\pgfqpoint{1.543617in}{2.739609in}}%
\pgfpathlineto{\pgfqpoint{1.540863in}{2.754226in}}%
\pgfpathlineto{\pgfqpoint{1.537828in}{2.768768in}}%
\pgfpathlineto{\pgfqpoint{1.534512in}{2.783228in}}%
\pgfpathlineto{\pgfqpoint{1.530913in}{2.797598in}}%
\pgfpathlineto{\pgfqpoint{1.527029in}{2.811870in}}%
\pgfpathlineto{\pgfqpoint{1.522858in}{2.826037in}}%
\pgfpathlineto{\pgfqpoint{1.518398in}{2.840089in}}%
\pgfpathlineto{\pgfqpoint{1.513647in}{2.854019in}}%
\pgfpathlineto{\pgfqpoint{1.508602in}{2.867818in}}%
\pgfpathlineto{\pgfqpoint{1.503261in}{2.881476in}}%
\pgfpathlineto{\pgfqpoint{1.497621in}{2.894982in}}%
\pgfpathlineto{\pgfqpoint{1.491678in}{2.908327in}}%
\pgfpathlineto{\pgfqpoint{1.485430in}{2.921500in}}%
\pgfpathlineto{\pgfqpoint{1.478757in}{2.934708in}}%
\pgfpathlineto{\pgfqpoint{1.471839in}{2.947576in}}%
\pgfpathlineto{\pgfqpoint{1.464587in}{2.960253in}}%
\pgfpathlineto{\pgfqpoint{1.456993in}{2.972730in}}%
\pgfpathlineto{\pgfqpoint{1.449047in}{2.984997in}}%
\pgfpathlineto{\pgfqpoint{1.440736in}{2.997041in}}%
\pgfpathlineto{\pgfqpoint{1.432048in}{3.008851in}}%
\pgfpathlineto{\pgfqpoint{1.422966in}{3.020413in}}%
\pgfpathlineto{\pgfqpoint{1.413473in}{3.031712in}}%
\pgfpathlineto{\pgfqpoint{1.403547in}{3.042731in}}%
\pgfpathlineto{\pgfqpoint{1.393164in}{3.053451in}}%
\pgfpathlineto{\pgfqpoint{1.382298in}{3.063848in}}%
\pgfpathlineto{\pgfqpoint{1.370920in}{3.073894in}}%
\pgfpathlineto{\pgfqpoint{1.358737in}{3.083755in}}%
\pgfpathlineto{\pgfqpoint{1.346031in}{3.093112in}}%
\pgfusepath{stroke}%
\end{pgfscope}%
\begin{pgfscope}%
\pgfpathmoveto{\pgfqpoint{1.192000in}{2.052407in}}%
\pgfpathcurveto{\pgfqpoint{1.481602in}{2.052407in}}{\pgfqpoint{1.759381in}{2.109937in}}{\pgfqpoint{1.964161in}{2.212326in}}%
\pgfpathcurveto{\pgfqpoint{2.168940in}{2.314716in}}{\pgfqpoint{2.284000in}{2.453606in}}{\pgfqpoint{2.284000in}{2.598407in}}%
\pgfpathcurveto{\pgfqpoint{2.284000in}{2.743208in}}{\pgfqpoint{2.168940in}{2.882097in}}{\pgfqpoint{1.964161in}{2.984487in}}%
\pgfpathcurveto{\pgfqpoint{1.759381in}{3.086877in}}{\pgfqpoint{1.481602in}{3.144407in}}{\pgfqpoint{1.192000in}{3.144407in}}%
\pgfpathcurveto{\pgfqpoint{0.902398in}{3.144407in}}{\pgfqpoint{0.624619in}{3.086877in}}{\pgfqpoint{0.419839in}{2.984487in}}%
\pgfpathcurveto{\pgfqpoint{0.215060in}{2.882097in}}{\pgfqpoint{0.100000in}{2.743208in}}{\pgfqpoint{0.100000in}{2.598407in}}%
\pgfpathcurveto{\pgfqpoint{0.100000in}{2.453606in}}{\pgfqpoint{0.215060in}{2.314716in}}{\pgfqpoint{0.419839in}{2.212326in}}%
\pgfpathcurveto{\pgfqpoint{0.624619in}{2.109937in}}{\pgfqpoint{0.902398in}{2.052407in}}{\pgfqpoint{1.192000in}{2.052407in}}%
\pgfpathlineto{\pgfqpoint{1.192000in}{2.052407in}}%
\pgfpathclose%
\pgfusepath{clip}%
\pgfsetrectcap%
\pgfsetroundjoin%
\pgfsetlinewidth{0.451687pt}%
\definecolor{currentstroke}{rgb}{0.690196,0.690196,0.690196}%
\pgfsetstrokecolor{currentstroke}%
\pgfsetstrokeopacity{0.250000}%
\pgfsetdash{}{0pt}%
\pgfpathmoveto{\pgfqpoint{1.423047in}{2.103702in}}%
\pgfpathlineto{\pgfqpoint{1.442106in}{2.113059in}}%
\pgfpathlineto{\pgfqpoint{1.460380in}{2.122919in}}%
\pgfpathlineto{\pgfqpoint{1.477447in}{2.132965in}}%
\pgfpathlineto{\pgfqpoint{1.493746in}{2.143363in}}%
\pgfpathlineto{\pgfqpoint{1.509320in}{2.154082in}}%
\pgfpathlineto{\pgfqpoint{1.524209in}{2.165102in}}%
\pgfpathlineto{\pgfqpoint{1.538449in}{2.176401in}}%
\pgfpathlineto{\pgfqpoint{1.552071in}{2.187962in}}%
\pgfpathlineto{\pgfqpoint{1.565104in}{2.199772in}}%
\pgfpathlineto{\pgfqpoint{1.577570in}{2.211817in}}%
\pgfpathlineto{\pgfqpoint{1.589490in}{2.224083in}}%
\pgfpathlineto{\pgfqpoint{1.600881in}{2.236561in}}%
\pgfpathlineto{\pgfqpoint{1.611758in}{2.249238in}}%
\pgfpathlineto{\pgfqpoint{1.622136in}{2.262105in}}%
\pgfpathlineto{\pgfqpoint{1.632144in}{2.275314in}}%
\pgfpathlineto{\pgfqpoint{1.641517in}{2.288486in}}%
\pgfpathlineto{\pgfqpoint{1.650431in}{2.301831in}}%
\pgfpathlineto{\pgfqpoint{1.658892in}{2.315338in}}%
\pgfpathlineto{\pgfqpoint{1.666904in}{2.328996in}}%
\pgfpathlineto{\pgfqpoint{1.674471in}{2.342794in}}%
\pgfpathlineto{\pgfqpoint{1.681597in}{2.356724in}}%
\pgfpathlineto{\pgfqpoint{1.688287in}{2.370777in}}%
\pgfpathlineto{\pgfqpoint{1.694543in}{2.384944in}}%
\pgfpathlineto{\pgfqpoint{1.700369in}{2.399216in}}%
\pgfpathlineto{\pgfqpoint{1.705768in}{2.413586in}}%
\pgfpathlineto{\pgfqpoint{1.710742in}{2.428046in}}%
\pgfpathlineto{\pgfqpoint{1.715294in}{2.442588in}}%
\pgfpathlineto{\pgfqpoint{1.719426in}{2.457204in}}%
\pgfpathlineto{\pgfqpoint{1.723139in}{2.471888in}}%
\pgfpathlineto{\pgfqpoint{1.726437in}{2.486632in}}%
\pgfpathlineto{\pgfqpoint{1.729319in}{2.501429in}}%
\pgfpathlineto{\pgfqpoint{1.731812in}{2.516436in}}%
\pgfpathlineto{\pgfqpoint{1.733853in}{2.531243in}}%
\pgfpathlineto{\pgfqpoint{1.735490in}{2.546108in}}%
\pgfpathlineto{\pgfqpoint{1.736718in}{2.561018in}}%
\pgfpathlineto{\pgfqpoint{1.737538in}{2.575961in}}%
\pgfpathlineto{\pgfqpoint{1.737949in}{2.590923in}}%
\pgfpathlineto{\pgfqpoint{1.737949in}{2.605891in}}%
\pgfpathlineto{\pgfqpoint{1.737538in}{2.620853in}}%
\pgfpathlineto{\pgfqpoint{1.736718in}{2.635795in}}%
\pgfpathlineto{\pgfqpoint{1.735490in}{2.650705in}}%
\pgfpathlineto{\pgfqpoint{1.733853in}{2.665570in}}%
\pgfpathlineto{\pgfqpoint{1.731812in}{2.680378in}}%
\pgfpathlineto{\pgfqpoint{1.729319in}{2.695384in}}%
\pgfpathlineto{\pgfqpoint{1.726437in}{2.710181in}}%
\pgfpathlineto{\pgfqpoint{1.723139in}{2.724925in}}%
\pgfpathlineto{\pgfqpoint{1.719426in}{2.739609in}}%
\pgfpathlineto{\pgfqpoint{1.715294in}{2.754226in}}%
\pgfpathlineto{\pgfqpoint{1.710742in}{2.768768in}}%
\pgfpathlineto{\pgfqpoint{1.705768in}{2.783228in}}%
\pgfpathlineto{\pgfqpoint{1.700369in}{2.797598in}}%
\pgfpathlineto{\pgfqpoint{1.694543in}{2.811870in}}%
\pgfpathlineto{\pgfqpoint{1.688287in}{2.826037in}}%
\pgfpathlineto{\pgfqpoint{1.681597in}{2.840089in}}%
\pgfpathlineto{\pgfqpoint{1.674471in}{2.854019in}}%
\pgfpathlineto{\pgfqpoint{1.666904in}{2.867818in}}%
\pgfpathlineto{\pgfqpoint{1.658892in}{2.881476in}}%
\pgfpathlineto{\pgfqpoint{1.650431in}{2.894982in}}%
\pgfpathlineto{\pgfqpoint{1.641517in}{2.908327in}}%
\pgfpathlineto{\pgfqpoint{1.632144in}{2.921500in}}%
\pgfpathlineto{\pgfqpoint{1.622136in}{2.934708in}}%
\pgfpathlineto{\pgfqpoint{1.611758in}{2.947576in}}%
\pgfpathlineto{\pgfqpoint{1.600881in}{2.960253in}}%
\pgfpathlineto{\pgfqpoint{1.589490in}{2.972730in}}%
\pgfpathlineto{\pgfqpoint{1.577570in}{2.984997in}}%
\pgfpathlineto{\pgfqpoint{1.565104in}{2.997041in}}%
\pgfpathlineto{\pgfqpoint{1.552071in}{3.008851in}}%
\pgfpathlineto{\pgfqpoint{1.538449in}{3.020413in}}%
\pgfpathlineto{\pgfqpoint{1.524209in}{3.031712in}}%
\pgfpathlineto{\pgfqpoint{1.509320in}{3.042731in}}%
\pgfpathlineto{\pgfqpoint{1.493746in}{3.053451in}}%
\pgfpathlineto{\pgfqpoint{1.477447in}{3.063848in}}%
\pgfpathlineto{\pgfqpoint{1.460380in}{3.073894in}}%
\pgfpathlineto{\pgfqpoint{1.442106in}{3.083755in}}%
\pgfpathlineto{\pgfqpoint{1.423047in}{3.093112in}}%
\pgfusepath{stroke}%
\end{pgfscope}%
\begin{pgfscope}%
\pgfpathmoveto{\pgfqpoint{1.192000in}{2.052407in}}%
\pgfpathcurveto{\pgfqpoint{1.481602in}{2.052407in}}{\pgfqpoint{1.759381in}{2.109937in}}{\pgfqpoint{1.964161in}{2.212326in}}%
\pgfpathcurveto{\pgfqpoint{2.168940in}{2.314716in}}{\pgfqpoint{2.284000in}{2.453606in}}{\pgfqpoint{2.284000in}{2.598407in}}%
\pgfpathcurveto{\pgfqpoint{2.284000in}{2.743208in}}{\pgfqpoint{2.168940in}{2.882097in}}{\pgfqpoint{1.964161in}{2.984487in}}%
\pgfpathcurveto{\pgfqpoint{1.759381in}{3.086877in}}{\pgfqpoint{1.481602in}{3.144407in}}{\pgfqpoint{1.192000in}{3.144407in}}%
\pgfpathcurveto{\pgfqpoint{0.902398in}{3.144407in}}{\pgfqpoint{0.624619in}{3.086877in}}{\pgfqpoint{0.419839in}{2.984487in}}%
\pgfpathcurveto{\pgfqpoint{0.215060in}{2.882097in}}{\pgfqpoint{0.100000in}{2.743208in}}{\pgfqpoint{0.100000in}{2.598407in}}%
\pgfpathcurveto{\pgfqpoint{0.100000in}{2.453606in}}{\pgfqpoint{0.215060in}{2.314716in}}{\pgfqpoint{0.419839in}{2.212326in}}%
\pgfpathcurveto{\pgfqpoint{0.624619in}{2.109937in}}{\pgfqpoint{0.902398in}{2.052407in}}{\pgfqpoint{1.192000in}{2.052407in}}%
\pgfpathlineto{\pgfqpoint{1.192000in}{2.052407in}}%
\pgfpathclose%
\pgfusepath{clip}%
\pgfsetrectcap%
\pgfsetroundjoin%
\pgfsetlinewidth{0.451687pt}%
\definecolor{currentstroke}{rgb}{0.690196,0.690196,0.690196}%
\pgfsetstrokecolor{currentstroke}%
\pgfsetstrokeopacity{0.250000}%
\pgfsetdash{}{0pt}%
\pgfpathmoveto{\pgfqpoint{1.500063in}{2.103702in}}%
\pgfpathlineto{\pgfqpoint{1.525475in}{2.113059in}}%
\pgfpathlineto{\pgfqpoint{1.549840in}{2.122919in}}%
\pgfpathlineto{\pgfqpoint{1.572596in}{2.132965in}}%
\pgfpathlineto{\pgfqpoint{1.594328in}{2.143363in}}%
\pgfpathlineto{\pgfqpoint{1.615093in}{2.154082in}}%
\pgfpathlineto{\pgfqpoint{1.634945in}{2.165102in}}%
\pgfpathlineto{\pgfqpoint{1.653932in}{2.176401in}}%
\pgfpathlineto{\pgfqpoint{1.672095in}{2.187962in}}%
\pgfpathlineto{\pgfqpoint{1.689472in}{2.199772in}}%
\pgfpathlineto{\pgfqpoint{1.706093in}{2.211817in}}%
\pgfpathlineto{\pgfqpoint{1.721986in}{2.224083in}}%
\pgfpathlineto{\pgfqpoint{1.737174in}{2.236561in}}%
\pgfpathlineto{\pgfqpoint{1.751678in}{2.249238in}}%
\pgfpathlineto{\pgfqpoint{1.765515in}{2.262105in}}%
\pgfpathlineto{\pgfqpoint{1.778859in}{2.275314in}}%
\pgfpathlineto{\pgfqpoint{1.791356in}{2.288486in}}%
\pgfpathlineto{\pgfqpoint{1.803241in}{2.301831in}}%
\pgfpathlineto{\pgfqpoint{1.814523in}{2.315338in}}%
\pgfpathlineto{\pgfqpoint{1.825205in}{2.328996in}}%
\pgfpathlineto{\pgfqpoint{1.835295in}{2.342794in}}%
\pgfpathlineto{\pgfqpoint{1.844796in}{2.356724in}}%
\pgfpathlineto{\pgfqpoint{1.853716in}{2.370777in}}%
\pgfpathlineto{\pgfqpoint{1.862057in}{2.384944in}}%
\pgfpathlineto{\pgfqpoint{1.869825in}{2.399216in}}%
\pgfpathlineto{\pgfqpoint{1.877024in}{2.413586in}}%
\pgfpathlineto{\pgfqpoint{1.883656in}{2.428046in}}%
\pgfpathlineto{\pgfqpoint{1.889725in}{2.442588in}}%
\pgfpathlineto{\pgfqpoint{1.895234in}{2.457204in}}%
\pgfpathlineto{\pgfqpoint{1.900186in}{2.471888in}}%
\pgfpathlineto{\pgfqpoint{1.904582in}{2.486632in}}%
\pgfpathlineto{\pgfqpoint{1.908425in}{2.501429in}}%
\pgfpathlineto{\pgfqpoint{1.911749in}{2.516436in}}%
\pgfpathlineto{\pgfqpoint{1.914471in}{2.531243in}}%
\pgfpathlineto{\pgfqpoint{1.916653in}{2.546108in}}%
\pgfpathlineto{\pgfqpoint{1.918291in}{2.561018in}}%
\pgfpathlineto{\pgfqpoint{1.919385in}{2.575961in}}%
\pgfpathlineto{\pgfqpoint{1.919932in}{2.590923in}}%
\pgfpathlineto{\pgfqpoint{1.919932in}{2.605891in}}%
\pgfpathlineto{\pgfqpoint{1.919385in}{2.620853in}}%
\pgfpathlineto{\pgfqpoint{1.918291in}{2.635795in}}%
\pgfpathlineto{\pgfqpoint{1.916653in}{2.650705in}}%
\pgfpathlineto{\pgfqpoint{1.914471in}{2.665570in}}%
\pgfpathlineto{\pgfqpoint{1.911749in}{2.680378in}}%
\pgfpathlineto{\pgfqpoint{1.908425in}{2.695384in}}%
\pgfpathlineto{\pgfqpoint{1.904582in}{2.710181in}}%
\pgfpathlineto{\pgfqpoint{1.900186in}{2.724925in}}%
\pgfpathlineto{\pgfqpoint{1.895234in}{2.739609in}}%
\pgfpathlineto{\pgfqpoint{1.889725in}{2.754226in}}%
\pgfpathlineto{\pgfqpoint{1.883656in}{2.768768in}}%
\pgfpathlineto{\pgfqpoint{1.877024in}{2.783228in}}%
\pgfpathlineto{\pgfqpoint{1.869825in}{2.797598in}}%
\pgfpathlineto{\pgfqpoint{1.862057in}{2.811870in}}%
\pgfpathlineto{\pgfqpoint{1.853716in}{2.826037in}}%
\pgfpathlineto{\pgfqpoint{1.844796in}{2.840089in}}%
\pgfpathlineto{\pgfqpoint{1.835295in}{2.854019in}}%
\pgfpathlineto{\pgfqpoint{1.825205in}{2.867818in}}%
\pgfpathlineto{\pgfqpoint{1.814523in}{2.881476in}}%
\pgfpathlineto{\pgfqpoint{1.803241in}{2.894982in}}%
\pgfpathlineto{\pgfqpoint{1.791356in}{2.908327in}}%
\pgfpathlineto{\pgfqpoint{1.778859in}{2.921500in}}%
\pgfpathlineto{\pgfqpoint{1.765515in}{2.934708in}}%
\pgfpathlineto{\pgfqpoint{1.751678in}{2.947576in}}%
\pgfpathlineto{\pgfqpoint{1.737174in}{2.960253in}}%
\pgfpathlineto{\pgfqpoint{1.721986in}{2.972730in}}%
\pgfpathlineto{\pgfqpoint{1.706093in}{2.984997in}}%
\pgfpathlineto{\pgfqpoint{1.689472in}{2.997041in}}%
\pgfpathlineto{\pgfqpoint{1.672095in}{3.008851in}}%
\pgfpathlineto{\pgfqpoint{1.653932in}{3.020413in}}%
\pgfpathlineto{\pgfqpoint{1.634945in}{3.031712in}}%
\pgfpathlineto{\pgfqpoint{1.615093in}{3.042731in}}%
\pgfpathlineto{\pgfqpoint{1.594328in}{3.053451in}}%
\pgfpathlineto{\pgfqpoint{1.572596in}{3.063848in}}%
\pgfpathlineto{\pgfqpoint{1.549840in}{3.073894in}}%
\pgfpathlineto{\pgfqpoint{1.525475in}{3.083755in}}%
\pgfpathlineto{\pgfqpoint{1.500063in}{3.093112in}}%
\pgfusepath{stroke}%
\end{pgfscope}%
\begin{pgfscope}%
\pgfpathmoveto{\pgfqpoint{1.192000in}{2.052407in}}%
\pgfpathcurveto{\pgfqpoint{1.481602in}{2.052407in}}{\pgfqpoint{1.759381in}{2.109937in}}{\pgfqpoint{1.964161in}{2.212326in}}%
\pgfpathcurveto{\pgfqpoint{2.168940in}{2.314716in}}{\pgfqpoint{2.284000in}{2.453606in}}{\pgfqpoint{2.284000in}{2.598407in}}%
\pgfpathcurveto{\pgfqpoint{2.284000in}{2.743208in}}{\pgfqpoint{2.168940in}{2.882097in}}{\pgfqpoint{1.964161in}{2.984487in}}%
\pgfpathcurveto{\pgfqpoint{1.759381in}{3.086877in}}{\pgfqpoint{1.481602in}{3.144407in}}{\pgfqpoint{1.192000in}{3.144407in}}%
\pgfpathcurveto{\pgfqpoint{0.902398in}{3.144407in}}{\pgfqpoint{0.624619in}{3.086877in}}{\pgfqpoint{0.419839in}{2.984487in}}%
\pgfpathcurveto{\pgfqpoint{0.215060in}{2.882097in}}{\pgfqpoint{0.100000in}{2.743208in}}{\pgfqpoint{0.100000in}{2.598407in}}%
\pgfpathcurveto{\pgfqpoint{0.100000in}{2.453606in}}{\pgfqpoint{0.215060in}{2.314716in}}{\pgfqpoint{0.419839in}{2.212326in}}%
\pgfpathcurveto{\pgfqpoint{0.624619in}{2.109937in}}{\pgfqpoint{0.902398in}{2.052407in}}{\pgfqpoint{1.192000in}{2.052407in}}%
\pgfpathlineto{\pgfqpoint{1.192000in}{2.052407in}}%
\pgfpathclose%
\pgfusepath{clip}%
\pgfsetrectcap%
\pgfsetroundjoin%
\pgfsetlinewidth{0.451687pt}%
\definecolor{currentstroke}{rgb}{0.690196,0.690196,0.690196}%
\pgfsetstrokecolor{currentstroke}%
\pgfsetstrokeopacity{0.250000}%
\pgfsetdash{}{0pt}%
\pgfpathmoveto{\pgfqpoint{1.577079in}{2.103702in}}%
\pgfpathlineto{\pgfqpoint{1.608843in}{2.113059in}}%
\pgfpathlineto{\pgfqpoint{1.639300in}{2.122919in}}%
\pgfpathlineto{\pgfqpoint{1.667745in}{2.132965in}}%
\pgfpathlineto{\pgfqpoint{1.694910in}{2.143363in}}%
\pgfpathlineto{\pgfqpoint{1.720866in}{2.154082in}}%
\pgfpathlineto{\pgfqpoint{1.745681in}{2.165102in}}%
\pgfpathlineto{\pgfqpoint{1.769415in}{2.176401in}}%
\pgfpathlineto{\pgfqpoint{1.792119in}{2.187962in}}%
\pgfpathlineto{\pgfqpoint{1.813840in}{2.199772in}}%
\pgfpathlineto{\pgfqpoint{1.834616in}{2.211817in}}%
\pgfpathlineto{\pgfqpoint{1.854483in}{2.224083in}}%
\pgfpathlineto{\pgfqpoint{1.873468in}{2.236561in}}%
\pgfpathlineto{\pgfqpoint{1.891597in}{2.249238in}}%
\pgfpathlineto{\pgfqpoint{1.908894in}{2.262105in}}%
\pgfpathlineto{\pgfqpoint{1.925574in}{2.275314in}}%
\pgfpathlineto{\pgfqpoint{1.941194in}{2.288486in}}%
\pgfpathlineto{\pgfqpoint{1.956052in}{2.301831in}}%
\pgfpathlineto{\pgfqpoint{1.970153in}{2.315338in}}%
\pgfpathlineto{\pgfqpoint{1.983506in}{2.328996in}}%
\pgfpathlineto{\pgfqpoint{1.996118in}{2.342794in}}%
\pgfpathlineto{\pgfqpoint{2.007996in}{2.356724in}}%
\pgfpathlineto{\pgfqpoint{2.019145in}{2.370777in}}%
\pgfpathlineto{\pgfqpoint{2.029572in}{2.384944in}}%
\pgfpathlineto{\pgfqpoint{2.039282in}{2.399216in}}%
\pgfpathlineto{\pgfqpoint{2.048280in}{2.413586in}}%
\pgfpathlineto{\pgfqpoint{2.056570in}{2.428046in}}%
\pgfpathlineto{\pgfqpoint{2.064156in}{2.442588in}}%
\pgfpathlineto{\pgfqpoint{2.071043in}{2.457204in}}%
\pgfpathlineto{\pgfqpoint{2.077232in}{2.471888in}}%
\pgfpathlineto{\pgfqpoint{2.082728in}{2.486632in}}%
\pgfpathlineto{\pgfqpoint{2.087531in}{2.501429in}}%
\pgfpathlineto{\pgfqpoint{2.091686in}{2.516436in}}%
\pgfpathlineto{\pgfqpoint{2.095089in}{2.531243in}}%
\pgfpathlineto{\pgfqpoint{2.097816in}{2.546108in}}%
\pgfpathlineto{\pgfqpoint{2.099864in}{2.561018in}}%
\pgfpathlineto{\pgfqpoint{2.101231in}{2.575961in}}%
\pgfpathlineto{\pgfqpoint{2.101915in}{2.590923in}}%
\pgfpathlineto{\pgfqpoint{2.101915in}{2.605891in}}%
\pgfpathlineto{\pgfqpoint{2.101231in}{2.620853in}}%
\pgfpathlineto{\pgfqpoint{2.099864in}{2.635795in}}%
\pgfpathlineto{\pgfqpoint{2.097816in}{2.650705in}}%
\pgfpathlineto{\pgfqpoint{2.095089in}{2.665570in}}%
\pgfpathlineto{\pgfqpoint{2.091686in}{2.680378in}}%
\pgfpathlineto{\pgfqpoint{2.087531in}{2.695384in}}%
\pgfpathlineto{\pgfqpoint{2.082728in}{2.710181in}}%
\pgfpathlineto{\pgfqpoint{2.077232in}{2.724925in}}%
\pgfpathlineto{\pgfqpoint{2.071043in}{2.739609in}}%
\pgfpathlineto{\pgfqpoint{2.064156in}{2.754226in}}%
\pgfpathlineto{\pgfqpoint{2.056570in}{2.768768in}}%
\pgfpathlineto{\pgfqpoint{2.048280in}{2.783228in}}%
\pgfpathlineto{\pgfqpoint{2.039282in}{2.797598in}}%
\pgfpathlineto{\pgfqpoint{2.029572in}{2.811870in}}%
\pgfpathlineto{\pgfqpoint{2.019145in}{2.826037in}}%
\pgfpathlineto{\pgfqpoint{2.007996in}{2.840089in}}%
\pgfpathlineto{\pgfqpoint{1.996118in}{2.854019in}}%
\pgfpathlineto{\pgfqpoint{1.983506in}{2.867818in}}%
\pgfpathlineto{\pgfqpoint{1.970153in}{2.881476in}}%
\pgfpathlineto{\pgfqpoint{1.956052in}{2.894982in}}%
\pgfpathlineto{\pgfqpoint{1.941194in}{2.908327in}}%
\pgfpathlineto{\pgfqpoint{1.925574in}{2.921500in}}%
\pgfpathlineto{\pgfqpoint{1.908894in}{2.934708in}}%
\pgfpathlineto{\pgfqpoint{1.891597in}{2.947576in}}%
\pgfpathlineto{\pgfqpoint{1.873468in}{2.960253in}}%
\pgfpathlineto{\pgfqpoint{1.854483in}{2.972730in}}%
\pgfpathlineto{\pgfqpoint{1.834616in}{2.984997in}}%
\pgfpathlineto{\pgfqpoint{1.813840in}{2.997041in}}%
\pgfpathlineto{\pgfqpoint{1.792119in}{3.008851in}}%
\pgfpathlineto{\pgfqpoint{1.769415in}{3.020413in}}%
\pgfpathlineto{\pgfqpoint{1.745681in}{3.031712in}}%
\pgfpathlineto{\pgfqpoint{1.720866in}{3.042731in}}%
\pgfpathlineto{\pgfqpoint{1.694910in}{3.053451in}}%
\pgfpathlineto{\pgfqpoint{1.667745in}{3.063848in}}%
\pgfpathlineto{\pgfqpoint{1.639300in}{3.073894in}}%
\pgfpathlineto{\pgfqpoint{1.608843in}{3.083755in}}%
\pgfpathlineto{\pgfqpoint{1.577079in}{3.093112in}}%
\pgfusepath{stroke}%
\end{pgfscope}%
\begin{pgfscope}%
\pgfpathmoveto{\pgfqpoint{1.192000in}{2.052407in}}%
\pgfpathcurveto{\pgfqpoint{1.481602in}{2.052407in}}{\pgfqpoint{1.759381in}{2.109937in}}{\pgfqpoint{1.964161in}{2.212326in}}%
\pgfpathcurveto{\pgfqpoint{2.168940in}{2.314716in}}{\pgfqpoint{2.284000in}{2.453606in}}{\pgfqpoint{2.284000in}{2.598407in}}%
\pgfpathcurveto{\pgfqpoint{2.284000in}{2.743208in}}{\pgfqpoint{2.168940in}{2.882097in}}{\pgfqpoint{1.964161in}{2.984487in}}%
\pgfpathcurveto{\pgfqpoint{1.759381in}{3.086877in}}{\pgfqpoint{1.481602in}{3.144407in}}{\pgfqpoint{1.192000in}{3.144407in}}%
\pgfpathcurveto{\pgfqpoint{0.902398in}{3.144407in}}{\pgfqpoint{0.624619in}{3.086877in}}{\pgfqpoint{0.419839in}{2.984487in}}%
\pgfpathcurveto{\pgfqpoint{0.215060in}{2.882097in}}{\pgfqpoint{0.100000in}{2.743208in}}{\pgfqpoint{0.100000in}{2.598407in}}%
\pgfpathcurveto{\pgfqpoint{0.100000in}{2.453606in}}{\pgfqpoint{0.215060in}{2.314716in}}{\pgfqpoint{0.419839in}{2.212326in}}%
\pgfpathcurveto{\pgfqpoint{0.624619in}{2.109937in}}{\pgfqpoint{0.902398in}{2.052407in}}{\pgfqpoint{1.192000in}{2.052407in}}%
\pgfpathlineto{\pgfqpoint{1.192000in}{2.052407in}}%
\pgfpathclose%
\pgfusepath{clip}%
\pgfsetrectcap%
\pgfsetroundjoin%
\pgfsetlinewidth{0.451687pt}%
\definecolor{currentstroke}{rgb}{0.690196,0.690196,0.690196}%
\pgfsetstrokecolor{currentstroke}%
\pgfsetstrokeopacity{0.250000}%
\pgfsetdash{}{0pt}%
\pgfpathmoveto{\pgfqpoint{0.729906in}{2.103702in}}%
\pgfpathlineto{\pgfqpoint{0.742228in}{2.103702in}}%
\pgfpathlineto{\pgfqpoint{0.754551in}{2.103702in}}%
\pgfpathlineto{\pgfqpoint{0.766873in}{2.103702in}}%
\pgfpathlineto{\pgfqpoint{0.779196in}{2.103702in}}%
\pgfpathlineto{\pgfqpoint{0.791518in}{2.103702in}}%
\pgfpathlineto{\pgfqpoint{0.803841in}{2.103702in}}%
\pgfpathlineto{\pgfqpoint{0.816163in}{2.103702in}}%
\pgfpathlineto{\pgfqpoint{0.828486in}{2.103702in}}%
\pgfpathlineto{\pgfqpoint{0.840808in}{2.103702in}}%
\pgfpathlineto{\pgfqpoint{0.853131in}{2.103702in}}%
\pgfpathlineto{\pgfqpoint{0.865453in}{2.103702in}}%
\pgfpathlineto{\pgfqpoint{0.877776in}{2.103702in}}%
\pgfpathlineto{\pgfqpoint{0.890098in}{2.103702in}}%
\pgfpathlineto{\pgfqpoint{0.902421in}{2.103702in}}%
\pgfpathlineto{\pgfqpoint{0.914743in}{2.103702in}}%
\pgfpathlineto{\pgfqpoint{0.927066in}{2.103702in}}%
\pgfpathlineto{\pgfqpoint{0.939388in}{2.103702in}}%
\pgfpathlineto{\pgfqpoint{0.951711in}{2.103702in}}%
\pgfpathlineto{\pgfqpoint{0.964033in}{2.103702in}}%
\pgfpathlineto{\pgfqpoint{0.976356in}{2.103702in}}%
\pgfpathlineto{\pgfqpoint{0.988678in}{2.103702in}}%
\pgfpathlineto{\pgfqpoint{1.001001in}{2.103702in}}%
\pgfpathlineto{\pgfqpoint{1.013323in}{2.103702in}}%
\pgfpathlineto{\pgfqpoint{1.025646in}{2.103702in}}%
\pgfpathlineto{\pgfqpoint{1.037969in}{2.103702in}}%
\pgfpathlineto{\pgfqpoint{1.050291in}{2.103702in}}%
\pgfpathlineto{\pgfqpoint{1.062614in}{2.103702in}}%
\pgfpathlineto{\pgfqpoint{1.074936in}{2.103702in}}%
\pgfpathlineto{\pgfqpoint{1.087259in}{2.103702in}}%
\pgfpathlineto{\pgfqpoint{1.099581in}{2.103702in}}%
\pgfpathlineto{\pgfqpoint{1.111904in}{2.103702in}}%
\pgfpathlineto{\pgfqpoint{1.124226in}{2.103702in}}%
\pgfpathlineto{\pgfqpoint{1.136549in}{2.103702in}}%
\pgfpathlineto{\pgfqpoint{1.148871in}{2.103702in}}%
\pgfpathlineto{\pgfqpoint{1.161194in}{2.103702in}}%
\pgfpathlineto{\pgfqpoint{1.173516in}{2.103702in}}%
\pgfpathlineto{\pgfqpoint{1.185839in}{2.103702in}}%
\pgfpathlineto{\pgfqpoint{1.198161in}{2.103702in}}%
\pgfpathlineto{\pgfqpoint{1.210484in}{2.103702in}}%
\pgfpathlineto{\pgfqpoint{1.222806in}{2.103702in}}%
\pgfpathlineto{\pgfqpoint{1.235129in}{2.103702in}}%
\pgfpathlineto{\pgfqpoint{1.247451in}{2.103702in}}%
\pgfpathlineto{\pgfqpoint{1.259774in}{2.103702in}}%
\pgfpathlineto{\pgfqpoint{1.272096in}{2.103702in}}%
\pgfpathlineto{\pgfqpoint{1.284419in}{2.103702in}}%
\pgfpathlineto{\pgfqpoint{1.296741in}{2.103702in}}%
\pgfpathlineto{\pgfqpoint{1.309064in}{2.103702in}}%
\pgfpathlineto{\pgfqpoint{1.321386in}{2.103702in}}%
\pgfpathlineto{\pgfqpoint{1.333709in}{2.103702in}}%
\pgfpathlineto{\pgfqpoint{1.346031in}{2.103702in}}%
\pgfpathlineto{\pgfqpoint{1.358354in}{2.103702in}}%
\pgfpathlineto{\pgfqpoint{1.370677in}{2.103702in}}%
\pgfpathlineto{\pgfqpoint{1.382999in}{2.103702in}}%
\pgfpathlineto{\pgfqpoint{1.395322in}{2.103702in}}%
\pgfpathlineto{\pgfqpoint{1.407644in}{2.103702in}}%
\pgfpathlineto{\pgfqpoint{1.419967in}{2.103702in}}%
\pgfpathlineto{\pgfqpoint{1.432289in}{2.103702in}}%
\pgfpathlineto{\pgfqpoint{1.444612in}{2.103702in}}%
\pgfpathlineto{\pgfqpoint{1.456934in}{2.103702in}}%
\pgfpathlineto{\pgfqpoint{1.469257in}{2.103702in}}%
\pgfpathlineto{\pgfqpoint{1.481579in}{2.103702in}}%
\pgfpathlineto{\pgfqpoint{1.493902in}{2.103702in}}%
\pgfpathlineto{\pgfqpoint{1.506224in}{2.103702in}}%
\pgfpathlineto{\pgfqpoint{1.518547in}{2.103702in}}%
\pgfpathlineto{\pgfqpoint{1.530869in}{2.103702in}}%
\pgfpathlineto{\pgfqpoint{1.543192in}{2.103702in}}%
\pgfpathlineto{\pgfqpoint{1.555514in}{2.103702in}}%
\pgfpathlineto{\pgfqpoint{1.567837in}{2.103702in}}%
\pgfpathlineto{\pgfqpoint{1.580159in}{2.103702in}}%
\pgfpathlineto{\pgfqpoint{1.592482in}{2.103702in}}%
\pgfpathlineto{\pgfqpoint{1.604804in}{2.103702in}}%
\pgfpathlineto{\pgfqpoint{1.617127in}{2.103702in}}%
\pgfpathlineto{\pgfqpoint{1.629449in}{2.103702in}}%
\pgfpathlineto{\pgfqpoint{1.641772in}{2.103702in}}%
\pgfpathlineto{\pgfqpoint{1.654094in}{2.103702in}}%
\pgfusepath{stroke}%
\end{pgfscope}%
\begin{pgfscope}%
\pgfpathmoveto{\pgfqpoint{1.192000in}{2.052407in}}%
\pgfpathcurveto{\pgfqpoint{1.481602in}{2.052407in}}{\pgfqpoint{1.759381in}{2.109937in}}{\pgfqpoint{1.964161in}{2.212326in}}%
\pgfpathcurveto{\pgfqpoint{2.168940in}{2.314716in}}{\pgfqpoint{2.284000in}{2.453606in}}{\pgfqpoint{2.284000in}{2.598407in}}%
\pgfpathcurveto{\pgfqpoint{2.284000in}{2.743208in}}{\pgfqpoint{2.168940in}{2.882097in}}{\pgfqpoint{1.964161in}{2.984487in}}%
\pgfpathcurveto{\pgfqpoint{1.759381in}{3.086877in}}{\pgfqpoint{1.481602in}{3.144407in}}{\pgfqpoint{1.192000in}{3.144407in}}%
\pgfpathcurveto{\pgfqpoint{0.902398in}{3.144407in}}{\pgfqpoint{0.624619in}{3.086877in}}{\pgfqpoint{0.419839in}{2.984487in}}%
\pgfpathcurveto{\pgfqpoint{0.215060in}{2.882097in}}{\pgfqpoint{0.100000in}{2.743208in}}{\pgfqpoint{0.100000in}{2.598407in}}%
\pgfpathcurveto{\pgfqpoint{0.100000in}{2.453606in}}{\pgfqpoint{0.215060in}{2.314716in}}{\pgfqpoint{0.419839in}{2.212326in}}%
\pgfpathcurveto{\pgfqpoint{0.624619in}{2.109937in}}{\pgfqpoint{0.902398in}{2.052407in}}{\pgfqpoint{1.192000in}{2.052407in}}%
\pgfpathlineto{\pgfqpoint{1.192000in}{2.052407in}}%
\pgfpathclose%
\pgfusepath{clip}%
\pgfsetrectcap%
\pgfsetroundjoin%
\pgfsetlinewidth{0.451687pt}%
\definecolor{currentstroke}{rgb}{0.690196,0.690196,0.690196}%
\pgfsetstrokecolor{currentstroke}%
\pgfsetstrokeopacity{0.250000}%
\pgfsetdash{}{0pt}%
\pgfpathmoveto{\pgfqpoint{0.485329in}{2.182150in}}%
\pgfpathlineto{\pgfqpoint{0.504174in}{2.182150in}}%
\pgfpathlineto{\pgfqpoint{0.523018in}{2.182150in}}%
\pgfpathlineto{\pgfqpoint{0.541863in}{2.182150in}}%
\pgfpathlineto{\pgfqpoint{0.560707in}{2.182150in}}%
\pgfpathlineto{\pgfqpoint{0.579552in}{2.182150in}}%
\pgfpathlineto{\pgfqpoint{0.598396in}{2.182150in}}%
\pgfpathlineto{\pgfqpoint{0.617241in}{2.182150in}}%
\pgfpathlineto{\pgfqpoint{0.636086in}{2.182150in}}%
\pgfpathlineto{\pgfqpoint{0.654930in}{2.182150in}}%
\pgfpathlineto{\pgfqpoint{0.673775in}{2.182150in}}%
\pgfpathlineto{\pgfqpoint{0.692619in}{2.182150in}}%
\pgfpathlineto{\pgfqpoint{0.711464in}{2.182150in}}%
\pgfpathlineto{\pgfqpoint{0.730308in}{2.182150in}}%
\pgfpathlineto{\pgfqpoint{0.749153in}{2.182150in}}%
\pgfpathlineto{\pgfqpoint{0.767997in}{2.182150in}}%
\pgfpathlineto{\pgfqpoint{0.786842in}{2.182150in}}%
\pgfpathlineto{\pgfqpoint{0.805687in}{2.182150in}}%
\pgfpathlineto{\pgfqpoint{0.824531in}{2.182150in}}%
\pgfpathlineto{\pgfqpoint{0.843376in}{2.182150in}}%
\pgfpathlineto{\pgfqpoint{0.862220in}{2.182150in}}%
\pgfpathlineto{\pgfqpoint{0.881065in}{2.182150in}}%
\pgfpathlineto{\pgfqpoint{0.899909in}{2.182150in}}%
\pgfpathlineto{\pgfqpoint{0.918754in}{2.182150in}}%
\pgfpathlineto{\pgfqpoint{0.937598in}{2.182150in}}%
\pgfpathlineto{\pgfqpoint{0.956443in}{2.182150in}}%
\pgfpathlineto{\pgfqpoint{0.975288in}{2.182150in}}%
\pgfpathlineto{\pgfqpoint{0.994132in}{2.182150in}}%
\pgfpathlineto{\pgfqpoint{1.012977in}{2.182150in}}%
\pgfpathlineto{\pgfqpoint{1.031821in}{2.182150in}}%
\pgfpathlineto{\pgfqpoint{1.050666in}{2.182150in}}%
\pgfpathlineto{\pgfqpoint{1.069510in}{2.182150in}}%
\pgfpathlineto{\pgfqpoint{1.088355in}{2.182150in}}%
\pgfpathlineto{\pgfqpoint{1.107199in}{2.182150in}}%
\pgfpathlineto{\pgfqpoint{1.126044in}{2.182150in}}%
\pgfpathlineto{\pgfqpoint{1.144889in}{2.182150in}}%
\pgfpathlineto{\pgfqpoint{1.163733in}{2.182150in}}%
\pgfpathlineto{\pgfqpoint{1.182578in}{2.182150in}}%
\pgfpathlineto{\pgfqpoint{1.201422in}{2.182150in}}%
\pgfpathlineto{\pgfqpoint{1.220267in}{2.182150in}}%
\pgfpathlineto{\pgfqpoint{1.239111in}{2.182150in}}%
\pgfpathlineto{\pgfqpoint{1.257956in}{2.182150in}}%
\pgfpathlineto{\pgfqpoint{1.276801in}{2.182150in}}%
\pgfpathlineto{\pgfqpoint{1.295645in}{2.182150in}}%
\pgfpathlineto{\pgfqpoint{1.314490in}{2.182150in}}%
\pgfpathlineto{\pgfqpoint{1.333334in}{2.182150in}}%
\pgfpathlineto{\pgfqpoint{1.352179in}{2.182150in}}%
\pgfpathlineto{\pgfqpoint{1.371023in}{2.182150in}}%
\pgfpathlineto{\pgfqpoint{1.389868in}{2.182150in}}%
\pgfpathlineto{\pgfqpoint{1.408712in}{2.182150in}}%
\pgfpathlineto{\pgfqpoint{1.427557in}{2.182150in}}%
\pgfpathlineto{\pgfqpoint{1.446402in}{2.182150in}}%
\pgfpathlineto{\pgfqpoint{1.465246in}{2.182150in}}%
\pgfpathlineto{\pgfqpoint{1.484091in}{2.182150in}}%
\pgfpathlineto{\pgfqpoint{1.502935in}{2.182150in}}%
\pgfpathlineto{\pgfqpoint{1.521780in}{2.182150in}}%
\pgfpathlineto{\pgfqpoint{1.540624in}{2.182150in}}%
\pgfpathlineto{\pgfqpoint{1.559469in}{2.182150in}}%
\pgfpathlineto{\pgfqpoint{1.578313in}{2.182150in}}%
\pgfpathlineto{\pgfqpoint{1.597158in}{2.182150in}}%
\pgfpathlineto{\pgfqpoint{1.616003in}{2.182150in}}%
\pgfpathlineto{\pgfqpoint{1.634847in}{2.182150in}}%
\pgfpathlineto{\pgfqpoint{1.653692in}{2.182150in}}%
\pgfpathlineto{\pgfqpoint{1.672536in}{2.182150in}}%
\pgfpathlineto{\pgfqpoint{1.691381in}{2.182150in}}%
\pgfpathlineto{\pgfqpoint{1.710225in}{2.182150in}}%
\pgfpathlineto{\pgfqpoint{1.729070in}{2.182150in}}%
\pgfpathlineto{\pgfqpoint{1.747914in}{2.182150in}}%
\pgfpathlineto{\pgfqpoint{1.766759in}{2.182150in}}%
\pgfpathlineto{\pgfqpoint{1.785604in}{2.182150in}}%
\pgfpathlineto{\pgfqpoint{1.804448in}{2.182150in}}%
\pgfpathlineto{\pgfqpoint{1.823293in}{2.182150in}}%
\pgfpathlineto{\pgfqpoint{1.842137in}{2.182150in}}%
\pgfpathlineto{\pgfqpoint{1.860982in}{2.182150in}}%
\pgfpathlineto{\pgfqpoint{1.879826in}{2.182150in}}%
\pgfpathlineto{\pgfqpoint{1.898671in}{2.182150in}}%
\pgfusepath{stroke}%
\end{pgfscope}%
\begin{pgfscope}%
\pgfpathmoveto{\pgfqpoint{1.192000in}{2.052407in}}%
\pgfpathcurveto{\pgfqpoint{1.481602in}{2.052407in}}{\pgfqpoint{1.759381in}{2.109937in}}{\pgfqpoint{1.964161in}{2.212326in}}%
\pgfpathcurveto{\pgfqpoint{2.168940in}{2.314716in}}{\pgfqpoint{2.284000in}{2.453606in}}{\pgfqpoint{2.284000in}{2.598407in}}%
\pgfpathcurveto{\pgfqpoint{2.284000in}{2.743208in}}{\pgfqpoint{2.168940in}{2.882097in}}{\pgfqpoint{1.964161in}{2.984487in}}%
\pgfpathcurveto{\pgfqpoint{1.759381in}{3.086877in}}{\pgfqpoint{1.481602in}{3.144407in}}{\pgfqpoint{1.192000in}{3.144407in}}%
\pgfpathcurveto{\pgfqpoint{0.902398in}{3.144407in}}{\pgfqpoint{0.624619in}{3.086877in}}{\pgfqpoint{0.419839in}{2.984487in}}%
\pgfpathcurveto{\pgfqpoint{0.215060in}{2.882097in}}{\pgfqpoint{0.100000in}{2.743208in}}{\pgfqpoint{0.100000in}{2.598407in}}%
\pgfpathcurveto{\pgfqpoint{0.100000in}{2.453606in}}{\pgfqpoint{0.215060in}{2.314716in}}{\pgfqpoint{0.419839in}{2.212326in}}%
\pgfpathcurveto{\pgfqpoint{0.624619in}{2.109937in}}{\pgfqpoint{0.902398in}{2.052407in}}{\pgfqpoint{1.192000in}{2.052407in}}%
\pgfpathlineto{\pgfqpoint{1.192000in}{2.052407in}}%
\pgfpathclose%
\pgfusepath{clip}%
\pgfsetrectcap%
\pgfsetroundjoin%
\pgfsetlinewidth{0.451687pt}%
\definecolor{currentstroke}{rgb}{0.690196,0.690196,0.690196}%
\pgfsetstrokecolor{currentstroke}%
\pgfsetstrokeopacity{0.250000}%
\pgfsetdash{}{0pt}%
\pgfpathmoveto{\pgfqpoint{0.311711in}{2.275314in}}%
\pgfpathlineto{\pgfqpoint{0.335186in}{2.275314in}}%
\pgfpathlineto{\pgfqpoint{0.358660in}{2.275314in}}%
\pgfpathlineto{\pgfqpoint{0.382134in}{2.275314in}}%
\pgfpathlineto{\pgfqpoint{0.405609in}{2.275314in}}%
\pgfpathlineto{\pgfqpoint{0.429083in}{2.275314in}}%
\pgfpathlineto{\pgfqpoint{0.452558in}{2.275314in}}%
\pgfpathlineto{\pgfqpoint{0.476032in}{2.275314in}}%
\pgfpathlineto{\pgfqpoint{0.499506in}{2.275314in}}%
\pgfpathlineto{\pgfqpoint{0.522981in}{2.275314in}}%
\pgfpathlineto{\pgfqpoint{0.546455in}{2.275314in}}%
\pgfpathlineto{\pgfqpoint{0.569929in}{2.275314in}}%
\pgfpathlineto{\pgfqpoint{0.593404in}{2.275314in}}%
\pgfpathlineto{\pgfqpoint{0.616878in}{2.275314in}}%
\pgfpathlineto{\pgfqpoint{0.640352in}{2.275314in}}%
\pgfpathlineto{\pgfqpoint{0.663827in}{2.275314in}}%
\pgfpathlineto{\pgfqpoint{0.687301in}{2.275314in}}%
\pgfpathlineto{\pgfqpoint{0.710776in}{2.275314in}}%
\pgfpathlineto{\pgfqpoint{0.734250in}{2.275314in}}%
\pgfpathlineto{\pgfqpoint{0.757724in}{2.275314in}}%
\pgfpathlineto{\pgfqpoint{0.781199in}{2.275314in}}%
\pgfpathlineto{\pgfqpoint{0.804673in}{2.275314in}}%
\pgfpathlineto{\pgfqpoint{0.828147in}{2.275314in}}%
\pgfpathlineto{\pgfqpoint{0.851622in}{2.275314in}}%
\pgfpathlineto{\pgfqpoint{0.875096in}{2.275314in}}%
\pgfpathlineto{\pgfqpoint{0.898570in}{2.275314in}}%
\pgfpathlineto{\pgfqpoint{0.922045in}{2.275314in}}%
\pgfpathlineto{\pgfqpoint{0.945519in}{2.275314in}}%
\pgfpathlineto{\pgfqpoint{0.968994in}{2.275314in}}%
\pgfpathlineto{\pgfqpoint{0.992468in}{2.275314in}}%
\pgfpathlineto{\pgfqpoint{1.015942in}{2.275314in}}%
\pgfpathlineto{\pgfqpoint{1.039417in}{2.275314in}}%
\pgfpathlineto{\pgfqpoint{1.062891in}{2.275314in}}%
\pgfpathlineto{\pgfqpoint{1.086365in}{2.275314in}}%
\pgfpathlineto{\pgfqpoint{1.109840in}{2.275314in}}%
\pgfpathlineto{\pgfqpoint{1.133314in}{2.275314in}}%
\pgfpathlineto{\pgfqpoint{1.156788in}{2.275314in}}%
\pgfpathlineto{\pgfqpoint{1.180263in}{2.275314in}}%
\pgfpathlineto{\pgfqpoint{1.203737in}{2.275314in}}%
\pgfpathlineto{\pgfqpoint{1.227212in}{2.275314in}}%
\pgfpathlineto{\pgfqpoint{1.250686in}{2.275314in}}%
\pgfpathlineto{\pgfqpoint{1.274160in}{2.275314in}}%
\pgfpathlineto{\pgfqpoint{1.297635in}{2.275314in}}%
\pgfpathlineto{\pgfqpoint{1.321109in}{2.275314in}}%
\pgfpathlineto{\pgfqpoint{1.344583in}{2.275314in}}%
\pgfpathlineto{\pgfqpoint{1.368058in}{2.275314in}}%
\pgfpathlineto{\pgfqpoint{1.391532in}{2.275314in}}%
\pgfpathlineto{\pgfqpoint{1.415006in}{2.275314in}}%
\pgfpathlineto{\pgfqpoint{1.438481in}{2.275314in}}%
\pgfpathlineto{\pgfqpoint{1.461955in}{2.275314in}}%
\pgfpathlineto{\pgfqpoint{1.485430in}{2.275314in}}%
\pgfpathlineto{\pgfqpoint{1.508904in}{2.275314in}}%
\pgfpathlineto{\pgfqpoint{1.532378in}{2.275314in}}%
\pgfpathlineto{\pgfqpoint{1.555853in}{2.275314in}}%
\pgfpathlineto{\pgfqpoint{1.579327in}{2.275314in}}%
\pgfpathlineto{\pgfqpoint{1.602801in}{2.275314in}}%
\pgfpathlineto{\pgfqpoint{1.626276in}{2.275314in}}%
\pgfpathlineto{\pgfqpoint{1.649750in}{2.275314in}}%
\pgfpathlineto{\pgfqpoint{1.673224in}{2.275314in}}%
\pgfpathlineto{\pgfqpoint{1.696699in}{2.275314in}}%
\pgfpathlineto{\pgfqpoint{1.720173in}{2.275314in}}%
\pgfpathlineto{\pgfqpoint{1.743648in}{2.275314in}}%
\pgfpathlineto{\pgfqpoint{1.767122in}{2.275314in}}%
\pgfpathlineto{\pgfqpoint{1.790596in}{2.275314in}}%
\pgfpathlineto{\pgfqpoint{1.814071in}{2.275314in}}%
\pgfpathlineto{\pgfqpoint{1.837545in}{2.275314in}}%
\pgfpathlineto{\pgfqpoint{1.861019in}{2.275314in}}%
\pgfpathlineto{\pgfqpoint{1.884494in}{2.275314in}}%
\pgfpathlineto{\pgfqpoint{1.907968in}{2.275314in}}%
\pgfpathlineto{\pgfqpoint{1.931442in}{2.275314in}}%
\pgfpathlineto{\pgfqpoint{1.954917in}{2.275314in}}%
\pgfpathlineto{\pgfqpoint{1.978391in}{2.275314in}}%
\pgfpathlineto{\pgfqpoint{2.001866in}{2.275314in}}%
\pgfpathlineto{\pgfqpoint{2.025340in}{2.275314in}}%
\pgfpathlineto{\pgfqpoint{2.048814in}{2.275314in}}%
\pgfpathlineto{\pgfqpoint{2.072289in}{2.275314in}}%
\pgfusepath{stroke}%
\end{pgfscope}%
\begin{pgfscope}%
\pgfpathmoveto{\pgfqpoint{1.192000in}{2.052407in}}%
\pgfpathcurveto{\pgfqpoint{1.481602in}{2.052407in}}{\pgfqpoint{1.759381in}{2.109937in}}{\pgfqpoint{1.964161in}{2.212326in}}%
\pgfpathcurveto{\pgfqpoint{2.168940in}{2.314716in}}{\pgfqpoint{2.284000in}{2.453606in}}{\pgfqpoint{2.284000in}{2.598407in}}%
\pgfpathcurveto{\pgfqpoint{2.284000in}{2.743208in}}{\pgfqpoint{2.168940in}{2.882097in}}{\pgfqpoint{1.964161in}{2.984487in}}%
\pgfpathcurveto{\pgfqpoint{1.759381in}{3.086877in}}{\pgfqpoint{1.481602in}{3.144407in}}{\pgfqpoint{1.192000in}{3.144407in}}%
\pgfpathcurveto{\pgfqpoint{0.902398in}{3.144407in}}{\pgfqpoint{0.624619in}{3.086877in}}{\pgfqpoint{0.419839in}{2.984487in}}%
\pgfpathcurveto{\pgfqpoint{0.215060in}{2.882097in}}{\pgfqpoint{0.100000in}{2.743208in}}{\pgfqpoint{0.100000in}{2.598407in}}%
\pgfpathcurveto{\pgfqpoint{0.100000in}{2.453606in}}{\pgfqpoint{0.215060in}{2.314716in}}{\pgfqpoint{0.419839in}{2.212326in}}%
\pgfpathcurveto{\pgfqpoint{0.624619in}{2.109937in}}{\pgfqpoint{0.902398in}{2.052407in}}{\pgfqpoint{1.192000in}{2.052407in}}%
\pgfpathlineto{\pgfqpoint{1.192000in}{2.052407in}}%
\pgfpathclose%
\pgfusepath{clip}%
\pgfsetrectcap%
\pgfsetroundjoin%
\pgfsetlinewidth{0.451687pt}%
\definecolor{currentstroke}{rgb}{0.690196,0.690196,0.690196}%
\pgfsetstrokecolor{currentstroke}%
\pgfsetstrokeopacity{0.250000}%
\pgfsetdash{}{0pt}%
\pgfpathmoveto{\pgfqpoint{0.193062in}{2.377847in}}%
\pgfpathlineto{\pgfqpoint{0.219701in}{2.377847in}}%
\pgfpathlineto{\pgfqpoint{0.246339in}{2.377847in}}%
\pgfpathlineto{\pgfqpoint{0.272977in}{2.377847in}}%
\pgfpathlineto{\pgfqpoint{0.299616in}{2.377847in}}%
\pgfpathlineto{\pgfqpoint{0.326254in}{2.377847in}}%
\pgfpathlineto{\pgfqpoint{0.352892in}{2.377847in}}%
\pgfpathlineto{\pgfqpoint{0.379531in}{2.377847in}}%
\pgfpathlineto{\pgfqpoint{0.406169in}{2.377847in}}%
\pgfpathlineto{\pgfqpoint{0.432807in}{2.377847in}}%
\pgfpathlineto{\pgfqpoint{0.459446in}{2.377847in}}%
\pgfpathlineto{\pgfqpoint{0.486084in}{2.377847in}}%
\pgfpathlineto{\pgfqpoint{0.512722in}{2.377847in}}%
\pgfpathlineto{\pgfqpoint{0.539361in}{2.377847in}}%
\pgfpathlineto{\pgfqpoint{0.565999in}{2.377847in}}%
\pgfpathlineto{\pgfqpoint{0.592637in}{2.377847in}}%
\pgfpathlineto{\pgfqpoint{0.619276in}{2.377847in}}%
\pgfpathlineto{\pgfqpoint{0.645914in}{2.377847in}}%
\pgfpathlineto{\pgfqpoint{0.672552in}{2.377847in}}%
\pgfpathlineto{\pgfqpoint{0.699191in}{2.377847in}}%
\pgfpathlineto{\pgfqpoint{0.725829in}{2.377847in}}%
\pgfpathlineto{\pgfqpoint{0.752467in}{2.377847in}}%
\pgfpathlineto{\pgfqpoint{0.779106in}{2.377847in}}%
\pgfpathlineto{\pgfqpoint{0.805744in}{2.377847in}}%
\pgfpathlineto{\pgfqpoint{0.832382in}{2.377847in}}%
\pgfpathlineto{\pgfqpoint{0.859021in}{2.377847in}}%
\pgfpathlineto{\pgfqpoint{0.885659in}{2.377847in}}%
\pgfpathlineto{\pgfqpoint{0.912297in}{2.377847in}}%
\pgfpathlineto{\pgfqpoint{0.938936in}{2.377847in}}%
\pgfpathlineto{\pgfqpoint{0.965574in}{2.377847in}}%
\pgfpathlineto{\pgfqpoint{0.992212in}{2.377847in}}%
\pgfpathlineto{\pgfqpoint{1.018851in}{2.377847in}}%
\pgfpathlineto{\pgfqpoint{1.045489in}{2.377847in}}%
\pgfpathlineto{\pgfqpoint{1.072127in}{2.377847in}}%
\pgfpathlineto{\pgfqpoint{1.098766in}{2.377847in}}%
\pgfpathlineto{\pgfqpoint{1.125404in}{2.377847in}}%
\pgfpathlineto{\pgfqpoint{1.152042in}{2.377847in}}%
\pgfpathlineto{\pgfqpoint{1.178681in}{2.377847in}}%
\pgfpathlineto{\pgfqpoint{1.205319in}{2.377847in}}%
\pgfpathlineto{\pgfqpoint{1.231958in}{2.377847in}}%
\pgfpathlineto{\pgfqpoint{1.258596in}{2.377847in}}%
\pgfpathlineto{\pgfqpoint{1.285234in}{2.377847in}}%
\pgfpathlineto{\pgfqpoint{1.311873in}{2.377847in}}%
\pgfpathlineto{\pgfqpoint{1.338511in}{2.377847in}}%
\pgfpathlineto{\pgfqpoint{1.365149in}{2.377847in}}%
\pgfpathlineto{\pgfqpoint{1.391788in}{2.377847in}}%
\pgfpathlineto{\pgfqpoint{1.418426in}{2.377847in}}%
\pgfpathlineto{\pgfqpoint{1.445064in}{2.377847in}}%
\pgfpathlineto{\pgfqpoint{1.471703in}{2.377847in}}%
\pgfpathlineto{\pgfqpoint{1.498341in}{2.377847in}}%
\pgfpathlineto{\pgfqpoint{1.524979in}{2.377847in}}%
\pgfpathlineto{\pgfqpoint{1.551618in}{2.377847in}}%
\pgfpathlineto{\pgfqpoint{1.578256in}{2.377847in}}%
\pgfpathlineto{\pgfqpoint{1.604894in}{2.377847in}}%
\pgfpathlineto{\pgfqpoint{1.631533in}{2.377847in}}%
\pgfpathlineto{\pgfqpoint{1.658171in}{2.377847in}}%
\pgfpathlineto{\pgfqpoint{1.684809in}{2.377847in}}%
\pgfpathlineto{\pgfqpoint{1.711448in}{2.377847in}}%
\pgfpathlineto{\pgfqpoint{1.738086in}{2.377847in}}%
\pgfpathlineto{\pgfqpoint{1.764724in}{2.377847in}}%
\pgfpathlineto{\pgfqpoint{1.791363in}{2.377847in}}%
\pgfpathlineto{\pgfqpoint{1.818001in}{2.377847in}}%
\pgfpathlineto{\pgfqpoint{1.844639in}{2.377847in}}%
\pgfpathlineto{\pgfqpoint{1.871278in}{2.377847in}}%
\pgfpathlineto{\pgfqpoint{1.897916in}{2.377847in}}%
\pgfpathlineto{\pgfqpoint{1.924554in}{2.377847in}}%
\pgfpathlineto{\pgfqpoint{1.951193in}{2.377847in}}%
\pgfpathlineto{\pgfqpoint{1.977831in}{2.377847in}}%
\pgfpathlineto{\pgfqpoint{2.004469in}{2.377847in}}%
\pgfpathlineto{\pgfqpoint{2.031108in}{2.377847in}}%
\pgfpathlineto{\pgfqpoint{2.057746in}{2.377847in}}%
\pgfpathlineto{\pgfqpoint{2.084384in}{2.377847in}}%
\pgfpathlineto{\pgfqpoint{2.111023in}{2.377847in}}%
\pgfpathlineto{\pgfqpoint{2.137661in}{2.377847in}}%
\pgfpathlineto{\pgfqpoint{2.164299in}{2.377847in}}%
\pgfpathlineto{\pgfqpoint{2.190938in}{2.377847in}}%
\pgfusepath{stroke}%
\end{pgfscope}%
\begin{pgfscope}%
\pgfpathmoveto{\pgfqpoint{1.192000in}{2.052407in}}%
\pgfpathcurveto{\pgfqpoint{1.481602in}{2.052407in}}{\pgfqpoint{1.759381in}{2.109937in}}{\pgfqpoint{1.964161in}{2.212326in}}%
\pgfpathcurveto{\pgfqpoint{2.168940in}{2.314716in}}{\pgfqpoint{2.284000in}{2.453606in}}{\pgfqpoint{2.284000in}{2.598407in}}%
\pgfpathcurveto{\pgfqpoint{2.284000in}{2.743208in}}{\pgfqpoint{2.168940in}{2.882097in}}{\pgfqpoint{1.964161in}{2.984487in}}%
\pgfpathcurveto{\pgfqpoint{1.759381in}{3.086877in}}{\pgfqpoint{1.481602in}{3.144407in}}{\pgfqpoint{1.192000in}{3.144407in}}%
\pgfpathcurveto{\pgfqpoint{0.902398in}{3.144407in}}{\pgfqpoint{0.624619in}{3.086877in}}{\pgfqpoint{0.419839in}{2.984487in}}%
\pgfpathcurveto{\pgfqpoint{0.215060in}{2.882097in}}{\pgfqpoint{0.100000in}{2.743208in}}{\pgfqpoint{0.100000in}{2.598407in}}%
\pgfpathcurveto{\pgfqpoint{0.100000in}{2.453606in}}{\pgfqpoint{0.215060in}{2.314716in}}{\pgfqpoint{0.419839in}{2.212326in}}%
\pgfpathcurveto{\pgfqpoint{0.624619in}{2.109937in}}{\pgfqpoint{0.902398in}{2.052407in}}{\pgfqpoint{1.192000in}{2.052407in}}%
\pgfpathlineto{\pgfqpoint{1.192000in}{2.052407in}}%
\pgfpathclose%
\pgfusepath{clip}%
\pgfsetrectcap%
\pgfsetroundjoin%
\pgfsetlinewidth{0.451687pt}%
\definecolor{currentstroke}{rgb}{0.690196,0.690196,0.690196}%
\pgfsetstrokecolor{currentstroke}%
\pgfsetstrokeopacity{0.250000}%
\pgfsetdash{}{0pt}%
\pgfpathmoveto{\pgfqpoint{0.123127in}{2.486632in}}%
\pgfpathlineto{\pgfqpoint{0.151630in}{2.486632in}}%
\pgfpathlineto{\pgfqpoint{0.180133in}{2.486632in}}%
\pgfpathlineto{\pgfqpoint{0.208637in}{2.486632in}}%
\pgfpathlineto{\pgfqpoint{0.237140in}{2.486632in}}%
\pgfpathlineto{\pgfqpoint{0.265643in}{2.486632in}}%
\pgfpathlineto{\pgfqpoint{0.294146in}{2.486632in}}%
\pgfpathlineto{\pgfqpoint{0.322650in}{2.486632in}}%
\pgfpathlineto{\pgfqpoint{0.351153in}{2.486632in}}%
\pgfpathlineto{\pgfqpoint{0.379656in}{2.486632in}}%
\pgfpathlineto{\pgfqpoint{0.408160in}{2.486632in}}%
\pgfpathlineto{\pgfqpoint{0.436663in}{2.486632in}}%
\pgfpathlineto{\pgfqpoint{0.465166in}{2.486632in}}%
\pgfpathlineto{\pgfqpoint{0.493669in}{2.486632in}}%
\pgfpathlineto{\pgfqpoint{0.522173in}{2.486632in}}%
\pgfpathlineto{\pgfqpoint{0.550676in}{2.486632in}}%
\pgfpathlineto{\pgfqpoint{0.579179in}{2.486632in}}%
\pgfpathlineto{\pgfqpoint{0.607683in}{2.486632in}}%
\pgfpathlineto{\pgfqpoint{0.636186in}{2.486632in}}%
\pgfpathlineto{\pgfqpoint{0.664689in}{2.486632in}}%
\pgfpathlineto{\pgfqpoint{0.693192in}{2.486632in}}%
\pgfpathlineto{\pgfqpoint{0.721696in}{2.486632in}}%
\pgfpathlineto{\pgfqpoint{0.750199in}{2.486632in}}%
\pgfpathlineto{\pgfqpoint{0.778702in}{2.486632in}}%
\pgfpathlineto{\pgfqpoint{0.807206in}{2.486632in}}%
\pgfpathlineto{\pgfqpoint{0.835709in}{2.486632in}}%
\pgfpathlineto{\pgfqpoint{0.864212in}{2.486632in}}%
\pgfpathlineto{\pgfqpoint{0.892715in}{2.486632in}}%
\pgfpathlineto{\pgfqpoint{0.921219in}{2.486632in}}%
\pgfpathlineto{\pgfqpoint{0.949722in}{2.486632in}}%
\pgfpathlineto{\pgfqpoint{0.978225in}{2.486632in}}%
\pgfpathlineto{\pgfqpoint{1.006729in}{2.486632in}}%
\pgfpathlineto{\pgfqpoint{1.035232in}{2.486632in}}%
\pgfpathlineto{\pgfqpoint{1.063735in}{2.486632in}}%
\pgfpathlineto{\pgfqpoint{1.092238in}{2.486632in}}%
\pgfpathlineto{\pgfqpoint{1.120742in}{2.486632in}}%
\pgfpathlineto{\pgfqpoint{1.149245in}{2.486632in}}%
\pgfpathlineto{\pgfqpoint{1.177748in}{2.486632in}}%
\pgfpathlineto{\pgfqpoint{1.206252in}{2.486632in}}%
\pgfpathlineto{\pgfqpoint{1.234755in}{2.486632in}}%
\pgfpathlineto{\pgfqpoint{1.263258in}{2.486632in}}%
\pgfpathlineto{\pgfqpoint{1.291762in}{2.486632in}}%
\pgfpathlineto{\pgfqpoint{1.320265in}{2.486632in}}%
\pgfpathlineto{\pgfqpoint{1.348768in}{2.486632in}}%
\pgfpathlineto{\pgfqpoint{1.377271in}{2.486632in}}%
\pgfpathlineto{\pgfqpoint{1.405775in}{2.486632in}}%
\pgfpathlineto{\pgfqpoint{1.434278in}{2.486632in}}%
\pgfpathlineto{\pgfqpoint{1.462781in}{2.486632in}}%
\pgfpathlineto{\pgfqpoint{1.491285in}{2.486632in}}%
\pgfpathlineto{\pgfqpoint{1.519788in}{2.486632in}}%
\pgfpathlineto{\pgfqpoint{1.548291in}{2.486632in}}%
\pgfpathlineto{\pgfqpoint{1.576794in}{2.486632in}}%
\pgfpathlineto{\pgfqpoint{1.605298in}{2.486632in}}%
\pgfpathlineto{\pgfqpoint{1.633801in}{2.486632in}}%
\pgfpathlineto{\pgfqpoint{1.662304in}{2.486632in}}%
\pgfpathlineto{\pgfqpoint{1.690808in}{2.486632in}}%
\pgfpathlineto{\pgfqpoint{1.719311in}{2.486632in}}%
\pgfpathlineto{\pgfqpoint{1.747814in}{2.486632in}}%
\pgfpathlineto{\pgfqpoint{1.776317in}{2.486632in}}%
\pgfpathlineto{\pgfqpoint{1.804821in}{2.486632in}}%
\pgfpathlineto{\pgfqpoint{1.833324in}{2.486632in}}%
\pgfpathlineto{\pgfqpoint{1.861827in}{2.486632in}}%
\pgfpathlineto{\pgfqpoint{1.890331in}{2.486632in}}%
\pgfpathlineto{\pgfqpoint{1.918834in}{2.486632in}}%
\pgfpathlineto{\pgfqpoint{1.947337in}{2.486632in}}%
\pgfpathlineto{\pgfqpoint{1.975840in}{2.486632in}}%
\pgfpathlineto{\pgfqpoint{2.004344in}{2.486632in}}%
\pgfpathlineto{\pgfqpoint{2.032847in}{2.486632in}}%
\pgfpathlineto{\pgfqpoint{2.061350in}{2.486632in}}%
\pgfpathlineto{\pgfqpoint{2.089854in}{2.486632in}}%
\pgfpathlineto{\pgfqpoint{2.118357in}{2.486632in}}%
\pgfpathlineto{\pgfqpoint{2.146860in}{2.486632in}}%
\pgfpathlineto{\pgfqpoint{2.175363in}{2.486632in}}%
\pgfpathlineto{\pgfqpoint{2.203867in}{2.486632in}}%
\pgfpathlineto{\pgfqpoint{2.232370in}{2.486632in}}%
\pgfpathlineto{\pgfqpoint{2.260873in}{2.486632in}}%
\pgfusepath{stroke}%
\end{pgfscope}%
\begin{pgfscope}%
\pgfpathmoveto{\pgfqpoint{1.192000in}{2.052407in}}%
\pgfpathcurveto{\pgfqpoint{1.481602in}{2.052407in}}{\pgfqpoint{1.759381in}{2.109937in}}{\pgfqpoint{1.964161in}{2.212326in}}%
\pgfpathcurveto{\pgfqpoint{2.168940in}{2.314716in}}{\pgfqpoint{2.284000in}{2.453606in}}{\pgfqpoint{2.284000in}{2.598407in}}%
\pgfpathcurveto{\pgfqpoint{2.284000in}{2.743208in}}{\pgfqpoint{2.168940in}{2.882097in}}{\pgfqpoint{1.964161in}{2.984487in}}%
\pgfpathcurveto{\pgfqpoint{1.759381in}{3.086877in}}{\pgfqpoint{1.481602in}{3.144407in}}{\pgfqpoint{1.192000in}{3.144407in}}%
\pgfpathcurveto{\pgfqpoint{0.902398in}{3.144407in}}{\pgfqpoint{0.624619in}{3.086877in}}{\pgfqpoint{0.419839in}{2.984487in}}%
\pgfpathcurveto{\pgfqpoint{0.215060in}{2.882097in}}{\pgfqpoint{0.100000in}{2.743208in}}{\pgfqpoint{0.100000in}{2.598407in}}%
\pgfpathcurveto{\pgfqpoint{0.100000in}{2.453606in}}{\pgfqpoint{0.215060in}{2.314716in}}{\pgfqpoint{0.419839in}{2.212326in}}%
\pgfpathcurveto{\pgfqpoint{0.624619in}{2.109937in}}{\pgfqpoint{0.902398in}{2.052407in}}{\pgfqpoint{1.192000in}{2.052407in}}%
\pgfpathlineto{\pgfqpoint{1.192000in}{2.052407in}}%
\pgfpathclose%
\pgfusepath{clip}%
\pgfsetrectcap%
\pgfsetroundjoin%
\pgfsetlinewidth{0.451687pt}%
\definecolor{currentstroke}{rgb}{0.690196,0.690196,0.690196}%
\pgfsetstrokecolor{currentstroke}%
\pgfsetstrokeopacity{0.250000}%
\pgfsetdash{}{0pt}%
\pgfpathmoveto{\pgfqpoint{0.100000in}{2.598407in}}%
\pgfpathlineto{\pgfqpoint{0.129120in}{2.598407in}}%
\pgfpathlineto{\pgfqpoint{0.158240in}{2.598407in}}%
\pgfpathlineto{\pgfqpoint{0.187360in}{2.598407in}}%
\pgfpathlineto{\pgfqpoint{0.216480in}{2.598407in}}%
\pgfpathlineto{\pgfqpoint{0.245600in}{2.598407in}}%
\pgfpathlineto{\pgfqpoint{0.274720in}{2.598407in}}%
\pgfpathlineto{\pgfqpoint{0.303840in}{2.598407in}}%
\pgfpathlineto{\pgfqpoint{0.332960in}{2.598407in}}%
\pgfpathlineto{\pgfqpoint{0.362080in}{2.598407in}}%
\pgfpathlineto{\pgfqpoint{0.391200in}{2.598407in}}%
\pgfpathlineto{\pgfqpoint{0.420320in}{2.598407in}}%
\pgfpathlineto{\pgfqpoint{0.449440in}{2.598407in}}%
\pgfpathlineto{\pgfqpoint{0.478560in}{2.598407in}}%
\pgfpathlineto{\pgfqpoint{0.507680in}{2.598407in}}%
\pgfpathlineto{\pgfqpoint{0.536800in}{2.598407in}}%
\pgfpathlineto{\pgfqpoint{0.565920in}{2.598407in}}%
\pgfpathlineto{\pgfqpoint{0.595040in}{2.598407in}}%
\pgfpathlineto{\pgfqpoint{0.624160in}{2.598407in}}%
\pgfpathlineto{\pgfqpoint{0.653280in}{2.598407in}}%
\pgfpathlineto{\pgfqpoint{0.682400in}{2.598407in}}%
\pgfpathlineto{\pgfqpoint{0.711520in}{2.598407in}}%
\pgfpathlineto{\pgfqpoint{0.740640in}{2.598407in}}%
\pgfpathlineto{\pgfqpoint{0.769760in}{2.598407in}}%
\pgfpathlineto{\pgfqpoint{0.798880in}{2.598407in}}%
\pgfpathlineto{\pgfqpoint{0.828000in}{2.598407in}}%
\pgfpathlineto{\pgfqpoint{0.857120in}{2.598407in}}%
\pgfpathlineto{\pgfqpoint{0.886240in}{2.598407in}}%
\pgfpathlineto{\pgfqpoint{0.915360in}{2.598407in}}%
\pgfpathlineto{\pgfqpoint{0.944480in}{2.598407in}}%
\pgfpathlineto{\pgfqpoint{0.973600in}{2.598407in}}%
\pgfpathlineto{\pgfqpoint{1.002720in}{2.598407in}}%
\pgfpathlineto{\pgfqpoint{1.031840in}{2.598407in}}%
\pgfpathlineto{\pgfqpoint{1.060960in}{2.598407in}}%
\pgfpathlineto{\pgfqpoint{1.090080in}{2.598407in}}%
\pgfpathlineto{\pgfqpoint{1.119200in}{2.598407in}}%
\pgfpathlineto{\pgfqpoint{1.148320in}{2.598407in}}%
\pgfpathlineto{\pgfqpoint{1.177440in}{2.598407in}}%
\pgfpathlineto{\pgfqpoint{1.206560in}{2.598407in}}%
\pgfpathlineto{\pgfqpoint{1.235680in}{2.598407in}}%
\pgfpathlineto{\pgfqpoint{1.264800in}{2.598407in}}%
\pgfpathlineto{\pgfqpoint{1.293920in}{2.598407in}}%
\pgfpathlineto{\pgfqpoint{1.323040in}{2.598407in}}%
\pgfpathlineto{\pgfqpoint{1.352160in}{2.598407in}}%
\pgfpathlineto{\pgfqpoint{1.381280in}{2.598407in}}%
\pgfpathlineto{\pgfqpoint{1.410400in}{2.598407in}}%
\pgfpathlineto{\pgfqpoint{1.439520in}{2.598407in}}%
\pgfpathlineto{\pgfqpoint{1.468640in}{2.598407in}}%
\pgfpathlineto{\pgfqpoint{1.497760in}{2.598407in}}%
\pgfpathlineto{\pgfqpoint{1.526880in}{2.598407in}}%
\pgfpathlineto{\pgfqpoint{1.556000in}{2.598407in}}%
\pgfpathlineto{\pgfqpoint{1.585120in}{2.598407in}}%
\pgfpathlineto{\pgfqpoint{1.614240in}{2.598407in}}%
\pgfpathlineto{\pgfqpoint{1.643360in}{2.598407in}}%
\pgfpathlineto{\pgfqpoint{1.672480in}{2.598407in}}%
\pgfpathlineto{\pgfqpoint{1.701600in}{2.598407in}}%
\pgfpathlineto{\pgfqpoint{1.730720in}{2.598407in}}%
\pgfpathlineto{\pgfqpoint{1.759840in}{2.598407in}}%
\pgfpathlineto{\pgfqpoint{1.788960in}{2.598407in}}%
\pgfpathlineto{\pgfqpoint{1.818080in}{2.598407in}}%
\pgfpathlineto{\pgfqpoint{1.847200in}{2.598407in}}%
\pgfpathlineto{\pgfqpoint{1.876320in}{2.598407in}}%
\pgfpathlineto{\pgfqpoint{1.905440in}{2.598407in}}%
\pgfpathlineto{\pgfqpoint{1.934560in}{2.598407in}}%
\pgfpathlineto{\pgfqpoint{1.963680in}{2.598407in}}%
\pgfpathlineto{\pgfqpoint{1.992800in}{2.598407in}}%
\pgfpathlineto{\pgfqpoint{2.021920in}{2.598407in}}%
\pgfpathlineto{\pgfqpoint{2.051040in}{2.598407in}}%
\pgfpathlineto{\pgfqpoint{2.080160in}{2.598407in}}%
\pgfpathlineto{\pgfqpoint{2.109280in}{2.598407in}}%
\pgfpathlineto{\pgfqpoint{2.138400in}{2.598407in}}%
\pgfpathlineto{\pgfqpoint{2.167520in}{2.598407in}}%
\pgfpathlineto{\pgfqpoint{2.196640in}{2.598407in}}%
\pgfpathlineto{\pgfqpoint{2.225760in}{2.598407in}}%
\pgfpathlineto{\pgfqpoint{2.254880in}{2.598407in}}%
\pgfpathlineto{\pgfqpoint{2.284000in}{2.598407in}}%
\pgfusepath{stroke}%
\end{pgfscope}%
\begin{pgfscope}%
\pgfpathmoveto{\pgfqpoint{1.192000in}{2.052407in}}%
\pgfpathcurveto{\pgfqpoint{1.481602in}{2.052407in}}{\pgfqpoint{1.759381in}{2.109937in}}{\pgfqpoint{1.964161in}{2.212326in}}%
\pgfpathcurveto{\pgfqpoint{2.168940in}{2.314716in}}{\pgfqpoint{2.284000in}{2.453606in}}{\pgfqpoint{2.284000in}{2.598407in}}%
\pgfpathcurveto{\pgfqpoint{2.284000in}{2.743208in}}{\pgfqpoint{2.168940in}{2.882097in}}{\pgfqpoint{1.964161in}{2.984487in}}%
\pgfpathcurveto{\pgfqpoint{1.759381in}{3.086877in}}{\pgfqpoint{1.481602in}{3.144407in}}{\pgfqpoint{1.192000in}{3.144407in}}%
\pgfpathcurveto{\pgfqpoint{0.902398in}{3.144407in}}{\pgfqpoint{0.624619in}{3.086877in}}{\pgfqpoint{0.419839in}{2.984487in}}%
\pgfpathcurveto{\pgfqpoint{0.215060in}{2.882097in}}{\pgfqpoint{0.100000in}{2.743208in}}{\pgfqpoint{0.100000in}{2.598407in}}%
\pgfpathcurveto{\pgfqpoint{0.100000in}{2.453606in}}{\pgfqpoint{0.215060in}{2.314716in}}{\pgfqpoint{0.419839in}{2.212326in}}%
\pgfpathcurveto{\pgfqpoint{0.624619in}{2.109937in}}{\pgfqpoint{0.902398in}{2.052407in}}{\pgfqpoint{1.192000in}{2.052407in}}%
\pgfpathlineto{\pgfqpoint{1.192000in}{2.052407in}}%
\pgfpathclose%
\pgfusepath{clip}%
\pgfsetrectcap%
\pgfsetroundjoin%
\pgfsetlinewidth{0.451687pt}%
\definecolor{currentstroke}{rgb}{0.690196,0.690196,0.690196}%
\pgfsetstrokecolor{currentstroke}%
\pgfsetstrokeopacity{0.250000}%
\pgfsetdash{}{0pt}%
\pgfpathmoveto{\pgfqpoint{0.123127in}{2.710181in}}%
\pgfpathlineto{\pgfqpoint{0.151630in}{2.710181in}}%
\pgfpathlineto{\pgfqpoint{0.180133in}{2.710181in}}%
\pgfpathlineto{\pgfqpoint{0.208637in}{2.710181in}}%
\pgfpathlineto{\pgfqpoint{0.237140in}{2.710181in}}%
\pgfpathlineto{\pgfqpoint{0.265643in}{2.710181in}}%
\pgfpathlineto{\pgfqpoint{0.294146in}{2.710181in}}%
\pgfpathlineto{\pgfqpoint{0.322650in}{2.710181in}}%
\pgfpathlineto{\pgfqpoint{0.351153in}{2.710181in}}%
\pgfpathlineto{\pgfqpoint{0.379656in}{2.710181in}}%
\pgfpathlineto{\pgfqpoint{0.408160in}{2.710181in}}%
\pgfpathlineto{\pgfqpoint{0.436663in}{2.710181in}}%
\pgfpathlineto{\pgfqpoint{0.465166in}{2.710181in}}%
\pgfpathlineto{\pgfqpoint{0.493669in}{2.710181in}}%
\pgfpathlineto{\pgfqpoint{0.522173in}{2.710181in}}%
\pgfpathlineto{\pgfqpoint{0.550676in}{2.710181in}}%
\pgfpathlineto{\pgfqpoint{0.579179in}{2.710181in}}%
\pgfpathlineto{\pgfqpoint{0.607683in}{2.710181in}}%
\pgfpathlineto{\pgfqpoint{0.636186in}{2.710181in}}%
\pgfpathlineto{\pgfqpoint{0.664689in}{2.710181in}}%
\pgfpathlineto{\pgfqpoint{0.693192in}{2.710181in}}%
\pgfpathlineto{\pgfqpoint{0.721696in}{2.710181in}}%
\pgfpathlineto{\pgfqpoint{0.750199in}{2.710181in}}%
\pgfpathlineto{\pgfqpoint{0.778702in}{2.710181in}}%
\pgfpathlineto{\pgfqpoint{0.807206in}{2.710181in}}%
\pgfpathlineto{\pgfqpoint{0.835709in}{2.710181in}}%
\pgfpathlineto{\pgfqpoint{0.864212in}{2.710181in}}%
\pgfpathlineto{\pgfqpoint{0.892715in}{2.710181in}}%
\pgfpathlineto{\pgfqpoint{0.921219in}{2.710181in}}%
\pgfpathlineto{\pgfqpoint{0.949722in}{2.710181in}}%
\pgfpathlineto{\pgfqpoint{0.978225in}{2.710181in}}%
\pgfpathlineto{\pgfqpoint{1.006729in}{2.710181in}}%
\pgfpathlineto{\pgfqpoint{1.035232in}{2.710181in}}%
\pgfpathlineto{\pgfqpoint{1.063735in}{2.710181in}}%
\pgfpathlineto{\pgfqpoint{1.092238in}{2.710181in}}%
\pgfpathlineto{\pgfqpoint{1.120742in}{2.710181in}}%
\pgfpathlineto{\pgfqpoint{1.149245in}{2.710181in}}%
\pgfpathlineto{\pgfqpoint{1.177748in}{2.710181in}}%
\pgfpathlineto{\pgfqpoint{1.206252in}{2.710181in}}%
\pgfpathlineto{\pgfqpoint{1.234755in}{2.710181in}}%
\pgfpathlineto{\pgfqpoint{1.263258in}{2.710181in}}%
\pgfpathlineto{\pgfqpoint{1.291762in}{2.710181in}}%
\pgfpathlineto{\pgfqpoint{1.320265in}{2.710181in}}%
\pgfpathlineto{\pgfqpoint{1.348768in}{2.710181in}}%
\pgfpathlineto{\pgfqpoint{1.377271in}{2.710181in}}%
\pgfpathlineto{\pgfqpoint{1.405775in}{2.710181in}}%
\pgfpathlineto{\pgfqpoint{1.434278in}{2.710181in}}%
\pgfpathlineto{\pgfqpoint{1.462781in}{2.710181in}}%
\pgfpathlineto{\pgfqpoint{1.491285in}{2.710181in}}%
\pgfpathlineto{\pgfqpoint{1.519788in}{2.710181in}}%
\pgfpathlineto{\pgfqpoint{1.548291in}{2.710181in}}%
\pgfpathlineto{\pgfqpoint{1.576794in}{2.710181in}}%
\pgfpathlineto{\pgfqpoint{1.605298in}{2.710181in}}%
\pgfpathlineto{\pgfqpoint{1.633801in}{2.710181in}}%
\pgfpathlineto{\pgfqpoint{1.662304in}{2.710181in}}%
\pgfpathlineto{\pgfqpoint{1.690808in}{2.710181in}}%
\pgfpathlineto{\pgfqpoint{1.719311in}{2.710181in}}%
\pgfpathlineto{\pgfqpoint{1.747814in}{2.710181in}}%
\pgfpathlineto{\pgfqpoint{1.776317in}{2.710181in}}%
\pgfpathlineto{\pgfqpoint{1.804821in}{2.710181in}}%
\pgfpathlineto{\pgfqpoint{1.833324in}{2.710181in}}%
\pgfpathlineto{\pgfqpoint{1.861827in}{2.710181in}}%
\pgfpathlineto{\pgfqpoint{1.890331in}{2.710181in}}%
\pgfpathlineto{\pgfqpoint{1.918834in}{2.710181in}}%
\pgfpathlineto{\pgfqpoint{1.947337in}{2.710181in}}%
\pgfpathlineto{\pgfqpoint{1.975840in}{2.710181in}}%
\pgfpathlineto{\pgfqpoint{2.004344in}{2.710181in}}%
\pgfpathlineto{\pgfqpoint{2.032847in}{2.710181in}}%
\pgfpathlineto{\pgfqpoint{2.061350in}{2.710181in}}%
\pgfpathlineto{\pgfqpoint{2.089854in}{2.710181in}}%
\pgfpathlineto{\pgfqpoint{2.118357in}{2.710181in}}%
\pgfpathlineto{\pgfqpoint{2.146860in}{2.710181in}}%
\pgfpathlineto{\pgfqpoint{2.175363in}{2.710181in}}%
\pgfpathlineto{\pgfqpoint{2.203867in}{2.710181in}}%
\pgfpathlineto{\pgfqpoint{2.232370in}{2.710181in}}%
\pgfpathlineto{\pgfqpoint{2.260873in}{2.710181in}}%
\pgfusepath{stroke}%
\end{pgfscope}%
\begin{pgfscope}%
\pgfpathmoveto{\pgfqpoint{1.192000in}{2.052407in}}%
\pgfpathcurveto{\pgfqpoint{1.481602in}{2.052407in}}{\pgfqpoint{1.759381in}{2.109937in}}{\pgfqpoint{1.964161in}{2.212326in}}%
\pgfpathcurveto{\pgfqpoint{2.168940in}{2.314716in}}{\pgfqpoint{2.284000in}{2.453606in}}{\pgfqpoint{2.284000in}{2.598407in}}%
\pgfpathcurveto{\pgfqpoint{2.284000in}{2.743208in}}{\pgfqpoint{2.168940in}{2.882097in}}{\pgfqpoint{1.964161in}{2.984487in}}%
\pgfpathcurveto{\pgfqpoint{1.759381in}{3.086877in}}{\pgfqpoint{1.481602in}{3.144407in}}{\pgfqpoint{1.192000in}{3.144407in}}%
\pgfpathcurveto{\pgfqpoint{0.902398in}{3.144407in}}{\pgfqpoint{0.624619in}{3.086877in}}{\pgfqpoint{0.419839in}{2.984487in}}%
\pgfpathcurveto{\pgfqpoint{0.215060in}{2.882097in}}{\pgfqpoint{0.100000in}{2.743208in}}{\pgfqpoint{0.100000in}{2.598407in}}%
\pgfpathcurveto{\pgfqpoint{0.100000in}{2.453606in}}{\pgfqpoint{0.215060in}{2.314716in}}{\pgfqpoint{0.419839in}{2.212326in}}%
\pgfpathcurveto{\pgfqpoint{0.624619in}{2.109937in}}{\pgfqpoint{0.902398in}{2.052407in}}{\pgfqpoint{1.192000in}{2.052407in}}%
\pgfpathlineto{\pgfqpoint{1.192000in}{2.052407in}}%
\pgfpathclose%
\pgfusepath{clip}%
\pgfsetrectcap%
\pgfsetroundjoin%
\pgfsetlinewidth{0.451687pt}%
\definecolor{currentstroke}{rgb}{0.690196,0.690196,0.690196}%
\pgfsetstrokecolor{currentstroke}%
\pgfsetstrokeopacity{0.250000}%
\pgfsetdash{}{0pt}%
\pgfpathmoveto{\pgfqpoint{0.193062in}{2.818967in}}%
\pgfpathlineto{\pgfqpoint{0.219701in}{2.818967in}}%
\pgfpathlineto{\pgfqpoint{0.246339in}{2.818967in}}%
\pgfpathlineto{\pgfqpoint{0.272977in}{2.818967in}}%
\pgfpathlineto{\pgfqpoint{0.299616in}{2.818967in}}%
\pgfpathlineto{\pgfqpoint{0.326254in}{2.818967in}}%
\pgfpathlineto{\pgfqpoint{0.352892in}{2.818967in}}%
\pgfpathlineto{\pgfqpoint{0.379531in}{2.818967in}}%
\pgfpathlineto{\pgfqpoint{0.406169in}{2.818967in}}%
\pgfpathlineto{\pgfqpoint{0.432807in}{2.818967in}}%
\pgfpathlineto{\pgfqpoint{0.459446in}{2.818967in}}%
\pgfpathlineto{\pgfqpoint{0.486084in}{2.818967in}}%
\pgfpathlineto{\pgfqpoint{0.512722in}{2.818967in}}%
\pgfpathlineto{\pgfqpoint{0.539361in}{2.818967in}}%
\pgfpathlineto{\pgfqpoint{0.565999in}{2.818967in}}%
\pgfpathlineto{\pgfqpoint{0.592637in}{2.818967in}}%
\pgfpathlineto{\pgfqpoint{0.619276in}{2.818967in}}%
\pgfpathlineto{\pgfqpoint{0.645914in}{2.818967in}}%
\pgfpathlineto{\pgfqpoint{0.672552in}{2.818967in}}%
\pgfpathlineto{\pgfqpoint{0.699191in}{2.818967in}}%
\pgfpathlineto{\pgfqpoint{0.725829in}{2.818967in}}%
\pgfpathlineto{\pgfqpoint{0.752467in}{2.818967in}}%
\pgfpathlineto{\pgfqpoint{0.779106in}{2.818967in}}%
\pgfpathlineto{\pgfqpoint{0.805744in}{2.818967in}}%
\pgfpathlineto{\pgfqpoint{0.832382in}{2.818967in}}%
\pgfpathlineto{\pgfqpoint{0.859021in}{2.818967in}}%
\pgfpathlineto{\pgfqpoint{0.885659in}{2.818967in}}%
\pgfpathlineto{\pgfqpoint{0.912297in}{2.818967in}}%
\pgfpathlineto{\pgfqpoint{0.938936in}{2.818967in}}%
\pgfpathlineto{\pgfqpoint{0.965574in}{2.818967in}}%
\pgfpathlineto{\pgfqpoint{0.992212in}{2.818967in}}%
\pgfpathlineto{\pgfqpoint{1.018851in}{2.818967in}}%
\pgfpathlineto{\pgfqpoint{1.045489in}{2.818967in}}%
\pgfpathlineto{\pgfqpoint{1.072127in}{2.818967in}}%
\pgfpathlineto{\pgfqpoint{1.098766in}{2.818967in}}%
\pgfpathlineto{\pgfqpoint{1.125404in}{2.818967in}}%
\pgfpathlineto{\pgfqpoint{1.152042in}{2.818967in}}%
\pgfpathlineto{\pgfqpoint{1.178681in}{2.818967in}}%
\pgfpathlineto{\pgfqpoint{1.205319in}{2.818967in}}%
\pgfpathlineto{\pgfqpoint{1.231958in}{2.818967in}}%
\pgfpathlineto{\pgfqpoint{1.258596in}{2.818967in}}%
\pgfpathlineto{\pgfqpoint{1.285234in}{2.818967in}}%
\pgfpathlineto{\pgfqpoint{1.311873in}{2.818967in}}%
\pgfpathlineto{\pgfqpoint{1.338511in}{2.818967in}}%
\pgfpathlineto{\pgfqpoint{1.365149in}{2.818967in}}%
\pgfpathlineto{\pgfqpoint{1.391788in}{2.818967in}}%
\pgfpathlineto{\pgfqpoint{1.418426in}{2.818967in}}%
\pgfpathlineto{\pgfqpoint{1.445064in}{2.818967in}}%
\pgfpathlineto{\pgfqpoint{1.471703in}{2.818967in}}%
\pgfpathlineto{\pgfqpoint{1.498341in}{2.818967in}}%
\pgfpathlineto{\pgfqpoint{1.524979in}{2.818967in}}%
\pgfpathlineto{\pgfqpoint{1.551618in}{2.818967in}}%
\pgfpathlineto{\pgfqpoint{1.578256in}{2.818967in}}%
\pgfpathlineto{\pgfqpoint{1.604894in}{2.818967in}}%
\pgfpathlineto{\pgfqpoint{1.631533in}{2.818967in}}%
\pgfpathlineto{\pgfqpoint{1.658171in}{2.818967in}}%
\pgfpathlineto{\pgfqpoint{1.684809in}{2.818967in}}%
\pgfpathlineto{\pgfqpoint{1.711448in}{2.818967in}}%
\pgfpathlineto{\pgfqpoint{1.738086in}{2.818967in}}%
\pgfpathlineto{\pgfqpoint{1.764724in}{2.818967in}}%
\pgfpathlineto{\pgfqpoint{1.791363in}{2.818967in}}%
\pgfpathlineto{\pgfqpoint{1.818001in}{2.818967in}}%
\pgfpathlineto{\pgfqpoint{1.844639in}{2.818967in}}%
\pgfpathlineto{\pgfqpoint{1.871278in}{2.818967in}}%
\pgfpathlineto{\pgfqpoint{1.897916in}{2.818967in}}%
\pgfpathlineto{\pgfqpoint{1.924554in}{2.818967in}}%
\pgfpathlineto{\pgfqpoint{1.951193in}{2.818967in}}%
\pgfpathlineto{\pgfqpoint{1.977831in}{2.818967in}}%
\pgfpathlineto{\pgfqpoint{2.004469in}{2.818967in}}%
\pgfpathlineto{\pgfqpoint{2.031108in}{2.818967in}}%
\pgfpathlineto{\pgfqpoint{2.057746in}{2.818967in}}%
\pgfpathlineto{\pgfqpoint{2.084384in}{2.818967in}}%
\pgfpathlineto{\pgfqpoint{2.111023in}{2.818967in}}%
\pgfpathlineto{\pgfqpoint{2.137661in}{2.818967in}}%
\pgfpathlineto{\pgfqpoint{2.164299in}{2.818967in}}%
\pgfpathlineto{\pgfqpoint{2.190938in}{2.818967in}}%
\pgfusepath{stroke}%
\end{pgfscope}%
\begin{pgfscope}%
\pgfpathmoveto{\pgfqpoint{1.192000in}{2.052407in}}%
\pgfpathcurveto{\pgfqpoint{1.481602in}{2.052407in}}{\pgfqpoint{1.759381in}{2.109937in}}{\pgfqpoint{1.964161in}{2.212326in}}%
\pgfpathcurveto{\pgfqpoint{2.168940in}{2.314716in}}{\pgfqpoint{2.284000in}{2.453606in}}{\pgfqpoint{2.284000in}{2.598407in}}%
\pgfpathcurveto{\pgfqpoint{2.284000in}{2.743208in}}{\pgfqpoint{2.168940in}{2.882097in}}{\pgfqpoint{1.964161in}{2.984487in}}%
\pgfpathcurveto{\pgfqpoint{1.759381in}{3.086877in}}{\pgfqpoint{1.481602in}{3.144407in}}{\pgfqpoint{1.192000in}{3.144407in}}%
\pgfpathcurveto{\pgfqpoint{0.902398in}{3.144407in}}{\pgfqpoint{0.624619in}{3.086877in}}{\pgfqpoint{0.419839in}{2.984487in}}%
\pgfpathcurveto{\pgfqpoint{0.215060in}{2.882097in}}{\pgfqpoint{0.100000in}{2.743208in}}{\pgfqpoint{0.100000in}{2.598407in}}%
\pgfpathcurveto{\pgfqpoint{0.100000in}{2.453606in}}{\pgfqpoint{0.215060in}{2.314716in}}{\pgfqpoint{0.419839in}{2.212326in}}%
\pgfpathcurveto{\pgfqpoint{0.624619in}{2.109937in}}{\pgfqpoint{0.902398in}{2.052407in}}{\pgfqpoint{1.192000in}{2.052407in}}%
\pgfpathlineto{\pgfqpoint{1.192000in}{2.052407in}}%
\pgfpathclose%
\pgfusepath{clip}%
\pgfsetrectcap%
\pgfsetroundjoin%
\pgfsetlinewidth{0.451687pt}%
\definecolor{currentstroke}{rgb}{0.690196,0.690196,0.690196}%
\pgfsetstrokecolor{currentstroke}%
\pgfsetstrokeopacity{0.250000}%
\pgfsetdash{}{0pt}%
\pgfpathmoveto{\pgfqpoint{0.311711in}{2.921500in}}%
\pgfpathlineto{\pgfqpoint{0.335186in}{2.921500in}}%
\pgfpathlineto{\pgfqpoint{0.358660in}{2.921500in}}%
\pgfpathlineto{\pgfqpoint{0.382134in}{2.921500in}}%
\pgfpathlineto{\pgfqpoint{0.405609in}{2.921500in}}%
\pgfpathlineto{\pgfqpoint{0.429083in}{2.921500in}}%
\pgfpathlineto{\pgfqpoint{0.452558in}{2.921500in}}%
\pgfpathlineto{\pgfqpoint{0.476032in}{2.921500in}}%
\pgfpathlineto{\pgfqpoint{0.499506in}{2.921500in}}%
\pgfpathlineto{\pgfqpoint{0.522981in}{2.921500in}}%
\pgfpathlineto{\pgfqpoint{0.546455in}{2.921500in}}%
\pgfpathlineto{\pgfqpoint{0.569929in}{2.921500in}}%
\pgfpathlineto{\pgfqpoint{0.593404in}{2.921500in}}%
\pgfpathlineto{\pgfqpoint{0.616878in}{2.921500in}}%
\pgfpathlineto{\pgfqpoint{0.640352in}{2.921500in}}%
\pgfpathlineto{\pgfqpoint{0.663827in}{2.921500in}}%
\pgfpathlineto{\pgfqpoint{0.687301in}{2.921500in}}%
\pgfpathlineto{\pgfqpoint{0.710776in}{2.921500in}}%
\pgfpathlineto{\pgfqpoint{0.734250in}{2.921500in}}%
\pgfpathlineto{\pgfqpoint{0.757724in}{2.921500in}}%
\pgfpathlineto{\pgfqpoint{0.781199in}{2.921500in}}%
\pgfpathlineto{\pgfqpoint{0.804673in}{2.921500in}}%
\pgfpathlineto{\pgfqpoint{0.828147in}{2.921500in}}%
\pgfpathlineto{\pgfqpoint{0.851622in}{2.921500in}}%
\pgfpathlineto{\pgfqpoint{0.875096in}{2.921500in}}%
\pgfpathlineto{\pgfqpoint{0.898570in}{2.921500in}}%
\pgfpathlineto{\pgfqpoint{0.922045in}{2.921500in}}%
\pgfpathlineto{\pgfqpoint{0.945519in}{2.921500in}}%
\pgfpathlineto{\pgfqpoint{0.968994in}{2.921500in}}%
\pgfpathlineto{\pgfqpoint{0.992468in}{2.921500in}}%
\pgfpathlineto{\pgfqpoint{1.015942in}{2.921500in}}%
\pgfpathlineto{\pgfqpoint{1.039417in}{2.921500in}}%
\pgfpathlineto{\pgfqpoint{1.062891in}{2.921500in}}%
\pgfpathlineto{\pgfqpoint{1.086365in}{2.921500in}}%
\pgfpathlineto{\pgfqpoint{1.109840in}{2.921500in}}%
\pgfpathlineto{\pgfqpoint{1.133314in}{2.921500in}}%
\pgfpathlineto{\pgfqpoint{1.156788in}{2.921500in}}%
\pgfpathlineto{\pgfqpoint{1.180263in}{2.921500in}}%
\pgfpathlineto{\pgfqpoint{1.203737in}{2.921500in}}%
\pgfpathlineto{\pgfqpoint{1.227212in}{2.921500in}}%
\pgfpathlineto{\pgfqpoint{1.250686in}{2.921500in}}%
\pgfpathlineto{\pgfqpoint{1.274160in}{2.921500in}}%
\pgfpathlineto{\pgfqpoint{1.297635in}{2.921500in}}%
\pgfpathlineto{\pgfqpoint{1.321109in}{2.921500in}}%
\pgfpathlineto{\pgfqpoint{1.344583in}{2.921500in}}%
\pgfpathlineto{\pgfqpoint{1.368058in}{2.921500in}}%
\pgfpathlineto{\pgfqpoint{1.391532in}{2.921500in}}%
\pgfpathlineto{\pgfqpoint{1.415006in}{2.921500in}}%
\pgfpathlineto{\pgfqpoint{1.438481in}{2.921500in}}%
\pgfpathlineto{\pgfqpoint{1.461955in}{2.921500in}}%
\pgfpathlineto{\pgfqpoint{1.485430in}{2.921500in}}%
\pgfpathlineto{\pgfqpoint{1.508904in}{2.921500in}}%
\pgfpathlineto{\pgfqpoint{1.532378in}{2.921500in}}%
\pgfpathlineto{\pgfqpoint{1.555853in}{2.921500in}}%
\pgfpathlineto{\pgfqpoint{1.579327in}{2.921500in}}%
\pgfpathlineto{\pgfqpoint{1.602801in}{2.921500in}}%
\pgfpathlineto{\pgfqpoint{1.626276in}{2.921500in}}%
\pgfpathlineto{\pgfqpoint{1.649750in}{2.921500in}}%
\pgfpathlineto{\pgfqpoint{1.673224in}{2.921500in}}%
\pgfpathlineto{\pgfqpoint{1.696699in}{2.921500in}}%
\pgfpathlineto{\pgfqpoint{1.720173in}{2.921500in}}%
\pgfpathlineto{\pgfqpoint{1.743648in}{2.921500in}}%
\pgfpathlineto{\pgfqpoint{1.767122in}{2.921500in}}%
\pgfpathlineto{\pgfqpoint{1.790596in}{2.921500in}}%
\pgfpathlineto{\pgfqpoint{1.814071in}{2.921500in}}%
\pgfpathlineto{\pgfqpoint{1.837545in}{2.921500in}}%
\pgfpathlineto{\pgfqpoint{1.861019in}{2.921500in}}%
\pgfpathlineto{\pgfqpoint{1.884494in}{2.921500in}}%
\pgfpathlineto{\pgfqpoint{1.907968in}{2.921500in}}%
\pgfpathlineto{\pgfqpoint{1.931442in}{2.921500in}}%
\pgfpathlineto{\pgfqpoint{1.954917in}{2.921500in}}%
\pgfpathlineto{\pgfqpoint{1.978391in}{2.921500in}}%
\pgfpathlineto{\pgfqpoint{2.001866in}{2.921500in}}%
\pgfpathlineto{\pgfqpoint{2.025340in}{2.921500in}}%
\pgfpathlineto{\pgfqpoint{2.048814in}{2.921500in}}%
\pgfpathlineto{\pgfqpoint{2.072289in}{2.921500in}}%
\pgfusepath{stroke}%
\end{pgfscope}%
\begin{pgfscope}%
\pgfpathmoveto{\pgfqpoint{1.192000in}{2.052407in}}%
\pgfpathcurveto{\pgfqpoint{1.481602in}{2.052407in}}{\pgfqpoint{1.759381in}{2.109937in}}{\pgfqpoint{1.964161in}{2.212326in}}%
\pgfpathcurveto{\pgfqpoint{2.168940in}{2.314716in}}{\pgfqpoint{2.284000in}{2.453606in}}{\pgfqpoint{2.284000in}{2.598407in}}%
\pgfpathcurveto{\pgfqpoint{2.284000in}{2.743208in}}{\pgfqpoint{2.168940in}{2.882097in}}{\pgfqpoint{1.964161in}{2.984487in}}%
\pgfpathcurveto{\pgfqpoint{1.759381in}{3.086877in}}{\pgfqpoint{1.481602in}{3.144407in}}{\pgfqpoint{1.192000in}{3.144407in}}%
\pgfpathcurveto{\pgfqpoint{0.902398in}{3.144407in}}{\pgfqpoint{0.624619in}{3.086877in}}{\pgfqpoint{0.419839in}{2.984487in}}%
\pgfpathcurveto{\pgfqpoint{0.215060in}{2.882097in}}{\pgfqpoint{0.100000in}{2.743208in}}{\pgfqpoint{0.100000in}{2.598407in}}%
\pgfpathcurveto{\pgfqpoint{0.100000in}{2.453606in}}{\pgfqpoint{0.215060in}{2.314716in}}{\pgfqpoint{0.419839in}{2.212326in}}%
\pgfpathcurveto{\pgfqpoint{0.624619in}{2.109937in}}{\pgfqpoint{0.902398in}{2.052407in}}{\pgfqpoint{1.192000in}{2.052407in}}%
\pgfpathlineto{\pgfqpoint{1.192000in}{2.052407in}}%
\pgfpathclose%
\pgfusepath{clip}%
\pgfsetrectcap%
\pgfsetroundjoin%
\pgfsetlinewidth{0.451687pt}%
\definecolor{currentstroke}{rgb}{0.690196,0.690196,0.690196}%
\pgfsetstrokecolor{currentstroke}%
\pgfsetstrokeopacity{0.250000}%
\pgfsetdash{}{0pt}%
\pgfpathmoveto{\pgfqpoint{0.485329in}{3.014664in}}%
\pgfpathlineto{\pgfqpoint{0.504174in}{3.014664in}}%
\pgfpathlineto{\pgfqpoint{0.523018in}{3.014664in}}%
\pgfpathlineto{\pgfqpoint{0.541863in}{3.014664in}}%
\pgfpathlineto{\pgfqpoint{0.560707in}{3.014664in}}%
\pgfpathlineto{\pgfqpoint{0.579552in}{3.014664in}}%
\pgfpathlineto{\pgfqpoint{0.598396in}{3.014664in}}%
\pgfpathlineto{\pgfqpoint{0.617241in}{3.014664in}}%
\pgfpathlineto{\pgfqpoint{0.636086in}{3.014664in}}%
\pgfpathlineto{\pgfqpoint{0.654930in}{3.014664in}}%
\pgfpathlineto{\pgfqpoint{0.673775in}{3.014664in}}%
\pgfpathlineto{\pgfqpoint{0.692619in}{3.014664in}}%
\pgfpathlineto{\pgfqpoint{0.711464in}{3.014664in}}%
\pgfpathlineto{\pgfqpoint{0.730308in}{3.014664in}}%
\pgfpathlineto{\pgfqpoint{0.749153in}{3.014664in}}%
\pgfpathlineto{\pgfqpoint{0.767997in}{3.014664in}}%
\pgfpathlineto{\pgfqpoint{0.786842in}{3.014664in}}%
\pgfpathlineto{\pgfqpoint{0.805687in}{3.014664in}}%
\pgfpathlineto{\pgfqpoint{0.824531in}{3.014664in}}%
\pgfpathlineto{\pgfqpoint{0.843376in}{3.014664in}}%
\pgfpathlineto{\pgfqpoint{0.862220in}{3.014664in}}%
\pgfpathlineto{\pgfqpoint{0.881065in}{3.014664in}}%
\pgfpathlineto{\pgfqpoint{0.899909in}{3.014664in}}%
\pgfpathlineto{\pgfqpoint{0.918754in}{3.014664in}}%
\pgfpathlineto{\pgfqpoint{0.937598in}{3.014664in}}%
\pgfpathlineto{\pgfqpoint{0.956443in}{3.014664in}}%
\pgfpathlineto{\pgfqpoint{0.975288in}{3.014664in}}%
\pgfpathlineto{\pgfqpoint{0.994132in}{3.014664in}}%
\pgfpathlineto{\pgfqpoint{1.012977in}{3.014664in}}%
\pgfpathlineto{\pgfqpoint{1.031821in}{3.014664in}}%
\pgfpathlineto{\pgfqpoint{1.050666in}{3.014664in}}%
\pgfpathlineto{\pgfqpoint{1.069510in}{3.014664in}}%
\pgfpathlineto{\pgfqpoint{1.088355in}{3.014664in}}%
\pgfpathlineto{\pgfqpoint{1.107199in}{3.014664in}}%
\pgfpathlineto{\pgfqpoint{1.126044in}{3.014664in}}%
\pgfpathlineto{\pgfqpoint{1.144889in}{3.014664in}}%
\pgfpathlineto{\pgfqpoint{1.163733in}{3.014664in}}%
\pgfpathlineto{\pgfqpoint{1.182578in}{3.014664in}}%
\pgfpathlineto{\pgfqpoint{1.201422in}{3.014664in}}%
\pgfpathlineto{\pgfqpoint{1.220267in}{3.014664in}}%
\pgfpathlineto{\pgfqpoint{1.239111in}{3.014664in}}%
\pgfpathlineto{\pgfqpoint{1.257956in}{3.014664in}}%
\pgfpathlineto{\pgfqpoint{1.276801in}{3.014664in}}%
\pgfpathlineto{\pgfqpoint{1.295645in}{3.014664in}}%
\pgfpathlineto{\pgfqpoint{1.314490in}{3.014664in}}%
\pgfpathlineto{\pgfqpoint{1.333334in}{3.014664in}}%
\pgfpathlineto{\pgfqpoint{1.352179in}{3.014664in}}%
\pgfpathlineto{\pgfqpoint{1.371023in}{3.014664in}}%
\pgfpathlineto{\pgfqpoint{1.389868in}{3.014664in}}%
\pgfpathlineto{\pgfqpoint{1.408712in}{3.014664in}}%
\pgfpathlineto{\pgfqpoint{1.427557in}{3.014664in}}%
\pgfpathlineto{\pgfqpoint{1.446402in}{3.014664in}}%
\pgfpathlineto{\pgfqpoint{1.465246in}{3.014664in}}%
\pgfpathlineto{\pgfqpoint{1.484091in}{3.014664in}}%
\pgfpathlineto{\pgfqpoint{1.502935in}{3.014664in}}%
\pgfpathlineto{\pgfqpoint{1.521780in}{3.014664in}}%
\pgfpathlineto{\pgfqpoint{1.540624in}{3.014664in}}%
\pgfpathlineto{\pgfqpoint{1.559469in}{3.014664in}}%
\pgfpathlineto{\pgfqpoint{1.578313in}{3.014664in}}%
\pgfpathlineto{\pgfqpoint{1.597158in}{3.014664in}}%
\pgfpathlineto{\pgfqpoint{1.616003in}{3.014664in}}%
\pgfpathlineto{\pgfqpoint{1.634847in}{3.014664in}}%
\pgfpathlineto{\pgfqpoint{1.653692in}{3.014664in}}%
\pgfpathlineto{\pgfqpoint{1.672536in}{3.014664in}}%
\pgfpathlineto{\pgfqpoint{1.691381in}{3.014664in}}%
\pgfpathlineto{\pgfqpoint{1.710225in}{3.014664in}}%
\pgfpathlineto{\pgfqpoint{1.729070in}{3.014664in}}%
\pgfpathlineto{\pgfqpoint{1.747914in}{3.014664in}}%
\pgfpathlineto{\pgfqpoint{1.766759in}{3.014664in}}%
\pgfpathlineto{\pgfqpoint{1.785604in}{3.014664in}}%
\pgfpathlineto{\pgfqpoint{1.804448in}{3.014664in}}%
\pgfpathlineto{\pgfqpoint{1.823293in}{3.014664in}}%
\pgfpathlineto{\pgfqpoint{1.842137in}{3.014664in}}%
\pgfpathlineto{\pgfqpoint{1.860982in}{3.014664in}}%
\pgfpathlineto{\pgfqpoint{1.879826in}{3.014664in}}%
\pgfpathlineto{\pgfqpoint{1.898671in}{3.014664in}}%
\pgfusepath{stroke}%
\end{pgfscope}%
\begin{pgfscope}%
\pgfpathmoveto{\pgfqpoint{1.192000in}{2.052407in}}%
\pgfpathcurveto{\pgfqpoint{1.481602in}{2.052407in}}{\pgfqpoint{1.759381in}{2.109937in}}{\pgfqpoint{1.964161in}{2.212326in}}%
\pgfpathcurveto{\pgfqpoint{2.168940in}{2.314716in}}{\pgfqpoint{2.284000in}{2.453606in}}{\pgfqpoint{2.284000in}{2.598407in}}%
\pgfpathcurveto{\pgfqpoint{2.284000in}{2.743208in}}{\pgfqpoint{2.168940in}{2.882097in}}{\pgfqpoint{1.964161in}{2.984487in}}%
\pgfpathcurveto{\pgfqpoint{1.759381in}{3.086877in}}{\pgfqpoint{1.481602in}{3.144407in}}{\pgfqpoint{1.192000in}{3.144407in}}%
\pgfpathcurveto{\pgfqpoint{0.902398in}{3.144407in}}{\pgfqpoint{0.624619in}{3.086877in}}{\pgfqpoint{0.419839in}{2.984487in}}%
\pgfpathcurveto{\pgfqpoint{0.215060in}{2.882097in}}{\pgfqpoint{0.100000in}{2.743208in}}{\pgfqpoint{0.100000in}{2.598407in}}%
\pgfpathcurveto{\pgfqpoint{0.100000in}{2.453606in}}{\pgfqpoint{0.215060in}{2.314716in}}{\pgfqpoint{0.419839in}{2.212326in}}%
\pgfpathcurveto{\pgfqpoint{0.624619in}{2.109937in}}{\pgfqpoint{0.902398in}{2.052407in}}{\pgfqpoint{1.192000in}{2.052407in}}%
\pgfpathlineto{\pgfqpoint{1.192000in}{2.052407in}}%
\pgfpathclose%
\pgfusepath{clip}%
\pgfsetrectcap%
\pgfsetroundjoin%
\pgfsetlinewidth{0.451687pt}%
\definecolor{currentstroke}{rgb}{0.690196,0.690196,0.690196}%
\pgfsetstrokecolor{currentstroke}%
\pgfsetstrokeopacity{0.250000}%
\pgfsetdash{}{0pt}%
\pgfpathmoveto{\pgfqpoint{0.729906in}{3.093112in}}%
\pgfpathlineto{\pgfqpoint{0.742228in}{3.093112in}}%
\pgfpathlineto{\pgfqpoint{0.754551in}{3.093112in}}%
\pgfpathlineto{\pgfqpoint{0.766873in}{3.093112in}}%
\pgfpathlineto{\pgfqpoint{0.779196in}{3.093112in}}%
\pgfpathlineto{\pgfqpoint{0.791518in}{3.093112in}}%
\pgfpathlineto{\pgfqpoint{0.803841in}{3.093112in}}%
\pgfpathlineto{\pgfqpoint{0.816163in}{3.093112in}}%
\pgfpathlineto{\pgfqpoint{0.828486in}{3.093112in}}%
\pgfpathlineto{\pgfqpoint{0.840808in}{3.093112in}}%
\pgfpathlineto{\pgfqpoint{0.853131in}{3.093112in}}%
\pgfpathlineto{\pgfqpoint{0.865453in}{3.093112in}}%
\pgfpathlineto{\pgfqpoint{0.877776in}{3.093112in}}%
\pgfpathlineto{\pgfqpoint{0.890098in}{3.093112in}}%
\pgfpathlineto{\pgfqpoint{0.902421in}{3.093112in}}%
\pgfpathlineto{\pgfqpoint{0.914743in}{3.093112in}}%
\pgfpathlineto{\pgfqpoint{0.927066in}{3.093112in}}%
\pgfpathlineto{\pgfqpoint{0.939388in}{3.093112in}}%
\pgfpathlineto{\pgfqpoint{0.951711in}{3.093112in}}%
\pgfpathlineto{\pgfqpoint{0.964033in}{3.093112in}}%
\pgfpathlineto{\pgfqpoint{0.976356in}{3.093112in}}%
\pgfpathlineto{\pgfqpoint{0.988678in}{3.093112in}}%
\pgfpathlineto{\pgfqpoint{1.001001in}{3.093112in}}%
\pgfpathlineto{\pgfqpoint{1.013323in}{3.093112in}}%
\pgfpathlineto{\pgfqpoint{1.025646in}{3.093112in}}%
\pgfpathlineto{\pgfqpoint{1.037969in}{3.093112in}}%
\pgfpathlineto{\pgfqpoint{1.050291in}{3.093112in}}%
\pgfpathlineto{\pgfqpoint{1.062614in}{3.093112in}}%
\pgfpathlineto{\pgfqpoint{1.074936in}{3.093112in}}%
\pgfpathlineto{\pgfqpoint{1.087259in}{3.093112in}}%
\pgfpathlineto{\pgfqpoint{1.099581in}{3.093112in}}%
\pgfpathlineto{\pgfqpoint{1.111904in}{3.093112in}}%
\pgfpathlineto{\pgfqpoint{1.124226in}{3.093112in}}%
\pgfpathlineto{\pgfqpoint{1.136549in}{3.093112in}}%
\pgfpathlineto{\pgfqpoint{1.148871in}{3.093112in}}%
\pgfpathlineto{\pgfqpoint{1.161194in}{3.093112in}}%
\pgfpathlineto{\pgfqpoint{1.173516in}{3.093112in}}%
\pgfpathlineto{\pgfqpoint{1.185839in}{3.093112in}}%
\pgfpathlineto{\pgfqpoint{1.198161in}{3.093112in}}%
\pgfpathlineto{\pgfqpoint{1.210484in}{3.093112in}}%
\pgfpathlineto{\pgfqpoint{1.222806in}{3.093112in}}%
\pgfpathlineto{\pgfqpoint{1.235129in}{3.093112in}}%
\pgfpathlineto{\pgfqpoint{1.247451in}{3.093112in}}%
\pgfpathlineto{\pgfqpoint{1.259774in}{3.093112in}}%
\pgfpathlineto{\pgfqpoint{1.272096in}{3.093112in}}%
\pgfpathlineto{\pgfqpoint{1.284419in}{3.093112in}}%
\pgfpathlineto{\pgfqpoint{1.296741in}{3.093112in}}%
\pgfpathlineto{\pgfqpoint{1.309064in}{3.093112in}}%
\pgfpathlineto{\pgfqpoint{1.321386in}{3.093112in}}%
\pgfpathlineto{\pgfqpoint{1.333709in}{3.093112in}}%
\pgfpathlineto{\pgfqpoint{1.346031in}{3.093112in}}%
\pgfpathlineto{\pgfqpoint{1.358354in}{3.093112in}}%
\pgfpathlineto{\pgfqpoint{1.370677in}{3.093112in}}%
\pgfpathlineto{\pgfqpoint{1.382999in}{3.093112in}}%
\pgfpathlineto{\pgfqpoint{1.395322in}{3.093112in}}%
\pgfpathlineto{\pgfqpoint{1.407644in}{3.093112in}}%
\pgfpathlineto{\pgfqpoint{1.419967in}{3.093112in}}%
\pgfpathlineto{\pgfqpoint{1.432289in}{3.093112in}}%
\pgfpathlineto{\pgfqpoint{1.444612in}{3.093112in}}%
\pgfpathlineto{\pgfqpoint{1.456934in}{3.093112in}}%
\pgfpathlineto{\pgfqpoint{1.469257in}{3.093112in}}%
\pgfpathlineto{\pgfqpoint{1.481579in}{3.093112in}}%
\pgfpathlineto{\pgfqpoint{1.493902in}{3.093112in}}%
\pgfpathlineto{\pgfqpoint{1.506224in}{3.093112in}}%
\pgfpathlineto{\pgfqpoint{1.518547in}{3.093112in}}%
\pgfpathlineto{\pgfqpoint{1.530869in}{3.093112in}}%
\pgfpathlineto{\pgfqpoint{1.543192in}{3.093112in}}%
\pgfpathlineto{\pgfqpoint{1.555514in}{3.093112in}}%
\pgfpathlineto{\pgfqpoint{1.567837in}{3.093112in}}%
\pgfpathlineto{\pgfqpoint{1.580159in}{3.093112in}}%
\pgfpathlineto{\pgfqpoint{1.592482in}{3.093112in}}%
\pgfpathlineto{\pgfqpoint{1.604804in}{3.093112in}}%
\pgfpathlineto{\pgfqpoint{1.617127in}{3.093112in}}%
\pgfpathlineto{\pgfqpoint{1.629449in}{3.093112in}}%
\pgfpathlineto{\pgfqpoint{1.641772in}{3.093112in}}%
\pgfpathlineto{\pgfqpoint{1.654094in}{3.093112in}}%
\pgfusepath{stroke}%
\end{pgfscope}%
\begin{pgfscope}%
\pgfsetrectcap%
\pgfsetmiterjoin%
\pgfsetlinewidth{0.803000pt}%
\definecolor{currentstroke}{rgb}{0.000000,0.000000,0.000000}%
\pgfsetstrokecolor{currentstroke}%
\pgfsetdash{}{0pt}%
\pgfpathmoveto{\pgfqpoint{1.192000in}{2.052407in}}%
\pgfpathcurveto{\pgfqpoint{1.481602in}{2.052407in}}{\pgfqpoint{1.759381in}{2.109937in}}{\pgfqpoint{1.964161in}{2.212326in}}%
\pgfpathcurveto{\pgfqpoint{2.168940in}{2.314716in}}{\pgfqpoint{2.284000in}{2.453606in}}{\pgfqpoint{2.284000in}{2.598407in}}%
\pgfpathcurveto{\pgfqpoint{2.284000in}{2.743208in}}{\pgfqpoint{2.168940in}{2.882097in}}{\pgfqpoint{1.964161in}{2.984487in}}%
\pgfpathcurveto{\pgfqpoint{1.759381in}{3.086877in}}{\pgfqpoint{1.481602in}{3.144407in}}{\pgfqpoint{1.192000in}{3.144407in}}%
\pgfpathcurveto{\pgfqpoint{0.902398in}{3.144407in}}{\pgfqpoint{0.624619in}{3.086877in}}{\pgfqpoint{0.419839in}{2.984487in}}%
\pgfpathcurveto{\pgfqpoint{0.215060in}{2.882097in}}{\pgfqpoint{0.100000in}{2.743208in}}{\pgfqpoint{0.100000in}{2.598407in}}%
\pgfpathcurveto{\pgfqpoint{0.100000in}{2.453606in}}{\pgfqpoint{0.215060in}{2.314716in}}{\pgfqpoint{0.419839in}{2.212326in}}%
\pgfpathcurveto{\pgfqpoint{0.624619in}{2.109937in}}{\pgfqpoint{0.902398in}{2.052407in}}{\pgfqpoint{1.192000in}{2.052407in}}%
\pgfpathlineto{\pgfqpoint{1.192000in}{2.052407in}}%
\pgfpathclose%
\pgfusepath{stroke}%
\end{pgfscope}%
\begin{pgfscope}%
\definecolor{textcolor}{rgb}{0.000000,0.000000,0.000000}%
\pgfsetstrokecolor{textcolor}%
\pgfsetfillcolor{textcolor}%
\pgftext[x=1.192000in,y=3.186073in,,base]{\color{textcolor}{\rmfamily\fontsize{8.200000}{9.840000}\selectfont\catcode`\^=\active\def^{\ifmmode\sp\else\^{}\fi}\catcode`\%=\active\def%{\%}Unimodal}}%
\end{pgfscope}%
\begin{pgfscope}%
\pgfsetbuttcap%
\pgfsetmiterjoin%
\definecolor{currentfill}{rgb}{1.000000,1.000000,1.000000}%
\pgfsetfillcolor{currentfill}%
\pgfsetlinewidth{0.000000pt}%
\definecolor{currentstroke}{rgb}{0.000000,0.000000,0.000000}%
\pgfsetstrokecolor{currentstroke}%
\pgfsetstrokeopacity{0.000000}%
\pgfsetdash{}{0pt}%
\pgfpathmoveto{\pgfqpoint{3.460000in}{2.052407in}}%
\pgfpathcurveto{\pgfqpoint{3.749602in}{2.052407in}}{\pgfqpoint{4.027381in}{2.109937in}}{\pgfqpoint{4.232161in}{2.212326in}}%
\pgfpathcurveto{\pgfqpoint{4.436940in}{2.314716in}}{\pgfqpoint{4.552000in}{2.453606in}}{\pgfqpoint{4.552000in}{2.598407in}}%
\pgfpathcurveto{\pgfqpoint{4.552000in}{2.743208in}}{\pgfqpoint{4.436940in}{2.882097in}}{\pgfqpoint{4.232161in}{2.984487in}}%
\pgfpathcurveto{\pgfqpoint{4.027381in}{3.086877in}}{\pgfqpoint{3.749602in}{3.144407in}}{\pgfqpoint{3.460000in}{3.144407in}}%
\pgfpathcurveto{\pgfqpoint{3.170398in}{3.144407in}}{\pgfqpoint{2.892619in}{3.086877in}}{\pgfqpoint{2.687839in}{2.984487in}}%
\pgfpathcurveto{\pgfqpoint{2.483060in}{2.882097in}}{\pgfqpoint{2.368000in}{2.743208in}}{\pgfqpoint{2.368000in}{2.598407in}}%
\pgfpathcurveto{\pgfqpoint{2.368000in}{2.453606in}}{\pgfqpoint{2.483060in}{2.314716in}}{\pgfqpoint{2.687839in}{2.212326in}}%
\pgfpathcurveto{\pgfqpoint{2.892619in}{2.109937in}}{\pgfqpoint{3.170398in}{2.052407in}}{\pgfqpoint{3.460000in}{2.052407in}}%
\pgfpathlineto{\pgfqpoint{3.460000in}{2.052407in}}%
\pgfpathclose%
\pgfusepath{fill}%
\end{pgfscope}%
\begin{pgfscope}%
\pgfpathmoveto{\pgfqpoint{3.460000in}{2.052407in}}%
\pgfpathcurveto{\pgfqpoint{3.749602in}{2.052407in}}{\pgfqpoint{4.027381in}{2.109937in}}{\pgfqpoint{4.232161in}{2.212326in}}%
\pgfpathcurveto{\pgfqpoint{4.436940in}{2.314716in}}{\pgfqpoint{4.552000in}{2.453606in}}{\pgfqpoint{4.552000in}{2.598407in}}%
\pgfpathcurveto{\pgfqpoint{4.552000in}{2.743208in}}{\pgfqpoint{4.436940in}{2.882097in}}{\pgfqpoint{4.232161in}{2.984487in}}%
\pgfpathcurveto{\pgfqpoint{4.027381in}{3.086877in}}{\pgfqpoint{3.749602in}{3.144407in}}{\pgfqpoint{3.460000in}{3.144407in}}%
\pgfpathcurveto{\pgfqpoint{3.170398in}{3.144407in}}{\pgfqpoint{2.892619in}{3.086877in}}{\pgfqpoint{2.687839in}{2.984487in}}%
\pgfpathcurveto{\pgfqpoint{2.483060in}{2.882097in}}{\pgfqpoint{2.368000in}{2.743208in}}{\pgfqpoint{2.368000in}{2.598407in}}%
\pgfpathcurveto{\pgfqpoint{2.368000in}{2.453606in}}{\pgfqpoint{2.483060in}{2.314716in}}{\pgfqpoint{2.687839in}{2.212326in}}%
\pgfpathcurveto{\pgfqpoint{2.892619in}{2.109937in}}{\pgfqpoint{3.170398in}{2.052407in}}{\pgfqpoint{3.460000in}{2.052407in}}%
\pgfpathlineto{\pgfqpoint{3.460000in}{2.052407in}}%
\pgfpathclose%
\pgfusepath{clip}%
\pgfsetbuttcap%
\pgfsetroundjoin%
\definecolor{currentfill}{rgb}{0.121569,0.466667,0.705882}%
\pgfsetfillcolor{currentfill}%
\pgfsetfillopacity{0.580000}%
\pgfsetlinewidth{0.000000pt}%
\definecolor{currentstroke}{rgb}{0.121569,0.466667,0.705882}%
\pgfsetstrokecolor{currentstroke}%
\pgfsetstrokeopacity{0.580000}%
\pgfsetdash{}{0pt}%
\pgfsys@defobject{currentmarker}{\pgfqpoint{-0.014232in}{-0.014232in}}{\pgfqpoint{0.014232in}{0.014232in}}{%
\pgfpathmoveto{\pgfqpoint{0.000000in}{-0.014232in}}%
\pgfpathcurveto{\pgfqpoint{0.003774in}{-0.014232in}}{\pgfqpoint{0.007395in}{-0.012732in}}{\pgfqpoint{0.010063in}{-0.010063in}}%
\pgfpathcurveto{\pgfqpoint{0.012732in}{-0.007395in}}{\pgfqpoint{0.014232in}{-0.003774in}}{\pgfqpoint{0.014232in}{0.000000in}}%
\pgfpathcurveto{\pgfqpoint{0.014232in}{0.003774in}}{\pgfqpoint{0.012732in}{0.007395in}}{\pgfqpoint{0.010063in}{0.010063in}}%
\pgfpathcurveto{\pgfqpoint{0.007395in}{0.012732in}}{\pgfqpoint{0.003774in}{0.014232in}}{\pgfqpoint{0.000000in}{0.014232in}}%
\pgfpathcurveto{\pgfqpoint{-0.003774in}{0.014232in}}{\pgfqpoint{-0.007395in}{0.012732in}}{\pgfqpoint{-0.010063in}{0.010063in}}%
\pgfpathcurveto{\pgfqpoint{-0.012732in}{0.007395in}}{\pgfqpoint{-0.014232in}{0.003774in}}{\pgfqpoint{-0.014232in}{0.000000in}}%
\pgfpathcurveto{\pgfqpoint{-0.014232in}{-0.003774in}}{\pgfqpoint{-0.012732in}{-0.007395in}}{\pgfqpoint{-0.010063in}{-0.010063in}}%
\pgfpathcurveto{\pgfqpoint{-0.007395in}{-0.012732in}}{\pgfqpoint{-0.003774in}{-0.014232in}}{\pgfqpoint{0.000000in}{-0.014232in}}%
\pgfpathlineto{\pgfqpoint{0.000000in}{-0.014232in}}%
\pgfpathclose%
\pgfusepath{fill}%
}%
\begin{pgfscope}%
\pgfsys@transformshift{3.598262in}{3.139114in}%
\pgfsys@useobject{currentmarker}{}%
\end{pgfscope}%
\begin{pgfscope}%
\pgfsys@transformshift{3.336666in}{3.136743in}%
\pgfsys@useobject{currentmarker}{}%
\end{pgfscope}%
\begin{pgfscope}%
\pgfsys@transformshift{3.178639in}{3.125367in}%
\pgfsys@useobject{currentmarker}{}%
\end{pgfscope}%
\begin{pgfscope}%
\pgfsys@transformshift{3.585035in}{3.110655in}%
\pgfsys@useobject{currentmarker}{}%
\end{pgfscope}%
\begin{pgfscope}%
\pgfsys@transformshift{3.557473in}{3.115548in}%
\pgfsys@useobject{currentmarker}{}%
\end{pgfscope}%
\begin{pgfscope}%
\pgfsys@transformshift{3.648100in}{3.118793in}%
\pgfsys@useobject{currentmarker}{}%
\end{pgfscope}%
\begin{pgfscope}%
\pgfsys@transformshift{3.490851in}{3.124187in}%
\pgfsys@useobject{currentmarker}{}%
\end{pgfscope}%
\begin{pgfscope}%
\pgfsys@transformshift{3.572059in}{3.128858in}%
\pgfsys@useobject{currentmarker}{}%
\end{pgfscope}%
\begin{pgfscope}%
\pgfsys@transformshift{3.490613in}{3.141229in}%
\pgfsys@useobject{currentmarker}{}%
\end{pgfscope}%
\begin{pgfscope}%
\pgfsys@transformshift{3.740146in}{3.123203in}%
\pgfsys@useobject{currentmarker}{}%
\end{pgfscope}%
\begin{pgfscope}%
\pgfsys@transformshift{3.554479in}{3.140372in}%
\pgfsys@useobject{currentmarker}{}%
\end{pgfscope}%
\begin{pgfscope}%
\pgfsys@transformshift{3.292724in}{3.134943in}%
\pgfsys@useobject{currentmarker}{}%
\end{pgfscope}%
\begin{pgfscope}%
\pgfsys@transformshift{3.544795in}{3.139165in}%
\pgfsys@useobject{currentmarker}{}%
\end{pgfscope}%
\begin{pgfscope}%
\pgfsys@transformshift{3.478266in}{3.127540in}%
\pgfsys@useobject{currentmarker}{}%
\end{pgfscope}%
\begin{pgfscope}%
\pgfsys@transformshift{3.784273in}{3.086434in}%
\pgfsys@useobject{currentmarker}{}%
\end{pgfscope}%
\begin{pgfscope}%
\pgfsys@transformshift{3.375250in}{3.137504in}%
\pgfsys@useobject{currentmarker}{}%
\end{pgfscope}%
\begin{pgfscope}%
\pgfsys@transformshift{3.619256in}{3.121759in}%
\pgfsys@useobject{currentmarker}{}%
\end{pgfscope}%
\begin{pgfscope}%
\pgfsys@transformshift{3.740005in}{3.116263in}%
\pgfsys@useobject{currentmarker}{}%
\end{pgfscope}%
\begin{pgfscope}%
\pgfsys@transformshift{3.377069in}{3.137811in}%
\pgfsys@useobject{currentmarker}{}%
\end{pgfscope}%
\begin{pgfscope}%
\pgfsys@transformshift{3.623431in}{3.135084in}%
\pgfsys@useobject{currentmarker}{}%
\end{pgfscope}%
\begin{pgfscope}%
\pgfsys@transformshift{3.387473in}{3.120242in}%
\pgfsys@useobject{currentmarker}{}%
\end{pgfscope}%
\begin{pgfscope}%
\pgfsys@transformshift{3.410120in}{3.122765in}%
\pgfsys@useobject{currentmarker}{}%
\end{pgfscope}%
\begin{pgfscope}%
\pgfsys@transformshift{3.528931in}{3.141579in}%
\pgfsys@useobject{currentmarker}{}%
\end{pgfscope}%
\begin{pgfscope}%
\pgfsys@transformshift{3.299397in}{3.122804in}%
\pgfsys@useobject{currentmarker}{}%
\end{pgfscope}%
\begin{pgfscope}%
\pgfsys@transformshift{3.559611in}{3.135470in}%
\pgfsys@useobject{currentmarker}{}%
\end{pgfscope}%
\begin{pgfscope}%
\pgfsys@transformshift{3.558075in}{3.115354in}%
\pgfsys@useobject{currentmarker}{}%
\end{pgfscope}%
\begin{pgfscope}%
\pgfsys@transformshift{3.603858in}{3.138194in}%
\pgfsys@useobject{currentmarker}{}%
\end{pgfscope}%
\begin{pgfscope}%
\pgfsys@transformshift{3.176335in}{3.114441in}%
\pgfsys@useobject{currentmarker}{}%
\end{pgfscope}%
\begin{pgfscope}%
\pgfsys@transformshift{3.428743in}{3.130757in}%
\pgfsys@useobject{currentmarker}{}%
\end{pgfscope}%
\begin{pgfscope}%
\pgfsys@transformshift{3.580315in}{3.103990in}%
\pgfsys@useobject{currentmarker}{}%
\end{pgfscope}%
\begin{pgfscope}%
\pgfsys@transformshift{3.325004in}{3.131651in}%
\pgfsys@useobject{currentmarker}{}%
\end{pgfscope}%
\begin{pgfscope}%
\pgfsys@transformshift{3.605924in}{3.138989in}%
\pgfsys@useobject{currentmarker}{}%
\end{pgfscope}%
\begin{pgfscope}%
\pgfsys@transformshift{3.558479in}{3.111413in}%
\pgfsys@useobject{currentmarker}{}%
\end{pgfscope}%
\begin{pgfscope}%
\pgfsys@transformshift{3.443457in}{3.123036in}%
\pgfsys@useobject{currentmarker}{}%
\end{pgfscope}%
\begin{pgfscope}%
\pgfsys@transformshift{3.350668in}{3.136964in}%
\pgfsys@useobject{currentmarker}{}%
\end{pgfscope}%
\begin{pgfscope}%
\pgfsys@transformshift{3.516763in}{3.133185in}%
\pgfsys@useobject{currentmarker}{}%
\end{pgfscope}%
\begin{pgfscope}%
\pgfsys@transformshift{3.404948in}{3.124515in}%
\pgfsys@useobject{currentmarker}{}%
\end{pgfscope}%
\begin{pgfscope}%
\pgfsys@transformshift{3.504379in}{3.133696in}%
\pgfsys@useobject{currentmarker}{}%
\end{pgfscope}%
\begin{pgfscope}%
\pgfsys@transformshift{3.507635in}{3.129065in}%
\pgfsys@useobject{currentmarker}{}%
\end{pgfscope}%
\begin{pgfscope}%
\pgfsys@transformshift{3.557881in}{3.139396in}%
\pgfsys@useobject{currentmarker}{}%
\end{pgfscope}%
\begin{pgfscope}%
\pgfsys@transformshift{3.423101in}{3.123771in}%
\pgfsys@useobject{currentmarker}{}%
\end{pgfscope}%
\begin{pgfscope}%
\pgfsys@transformshift{3.472000in}{3.136796in}%
\pgfsys@useobject{currentmarker}{}%
\end{pgfscope}%
\begin{pgfscope}%
\pgfsys@transformshift{3.425060in}{3.133861in}%
\pgfsys@useobject{currentmarker}{}%
\end{pgfscope}%
\begin{pgfscope}%
\pgfsys@transformshift{3.315093in}{3.136381in}%
\pgfsys@useobject{currentmarker}{}%
\end{pgfscope}%
\begin{pgfscope}%
\pgfsys@transformshift{3.496936in}{3.126754in}%
\pgfsys@useobject{currentmarker}{}%
\end{pgfscope}%
\begin{pgfscope}%
\pgfsys@transformshift{3.295385in}{3.136009in}%
\pgfsys@useobject{currentmarker}{}%
\end{pgfscope}%
\begin{pgfscope}%
\pgfsys@transformshift{3.641505in}{3.135374in}%
\pgfsys@useobject{currentmarker}{}%
\end{pgfscope}%
\begin{pgfscope}%
\pgfsys@transformshift{3.221794in}{3.127048in}%
\pgfsys@useobject{currentmarker}{}%
\end{pgfscope}%
\begin{pgfscope}%
\pgfsys@transformshift{3.237478in}{3.101701in}%
\pgfsys@useobject{currentmarker}{}%
\end{pgfscope}%
\begin{pgfscope}%
\pgfsys@transformshift{3.399711in}{3.142389in}%
\pgfsys@useobject{currentmarker}{}%
\end{pgfscope}%
\begin{pgfscope}%
\pgfsys@transformshift{3.458925in}{3.135738in}%
\pgfsys@useobject{currentmarker}{}%
\end{pgfscope}%
\begin{pgfscope}%
\pgfsys@transformshift{3.561435in}{3.119863in}%
\pgfsys@useobject{currentmarker}{}%
\end{pgfscope}%
\begin{pgfscope}%
\pgfsys@transformshift{3.389188in}{3.136150in}%
\pgfsys@useobject{currentmarker}{}%
\end{pgfscope}%
\begin{pgfscope}%
\pgfsys@transformshift{3.707717in}{3.123213in}%
\pgfsys@useobject{currentmarker}{}%
\end{pgfscope}%
\begin{pgfscope}%
\pgfsys@transformshift{3.350090in}{3.129478in}%
\pgfsys@useobject{currentmarker}{}%
\end{pgfscope}%
\begin{pgfscope}%
\pgfsys@transformshift{3.427847in}{3.139331in}%
\pgfsys@useobject{currentmarker}{}%
\end{pgfscope}%
\begin{pgfscope}%
\pgfsys@transformshift{3.455861in}{3.128874in}%
\pgfsys@useobject{currentmarker}{}%
\end{pgfscope}%
\begin{pgfscope}%
\pgfsys@transformshift{3.618945in}{3.123580in}%
\pgfsys@useobject{currentmarker}{}%
\end{pgfscope}%
\begin{pgfscope}%
\pgfsys@transformshift{3.726332in}{3.124245in}%
\pgfsys@useobject{currentmarker}{}%
\end{pgfscope}%
\begin{pgfscope}%
\pgfsys@transformshift{3.416709in}{3.110600in}%
\pgfsys@useobject{currentmarker}{}%
\end{pgfscope}%
\begin{pgfscope}%
\pgfsys@transformshift{3.425376in}{3.120583in}%
\pgfsys@useobject{currentmarker}{}%
\end{pgfscope}%
\begin{pgfscope}%
\pgfsys@transformshift{3.686265in}{3.122909in}%
\pgfsys@useobject{currentmarker}{}%
\end{pgfscope}%
\begin{pgfscope}%
\pgfsys@transformshift{3.420000in}{3.118160in}%
\pgfsys@useobject{currentmarker}{}%
\end{pgfscope}%
\begin{pgfscope}%
\pgfsys@transformshift{3.774640in}{3.108857in}%
\pgfsys@useobject{currentmarker}{}%
\end{pgfscope}%
\begin{pgfscope}%
\pgfsys@transformshift{3.650704in}{3.129380in}%
\pgfsys@useobject{currentmarker}{}%
\end{pgfscope}%
\begin{pgfscope}%
\pgfsys@transformshift{3.303922in}{3.120338in}%
\pgfsys@useobject{currentmarker}{}%
\end{pgfscope}%
\begin{pgfscope}%
\pgfsys@transformshift{3.636241in}{3.136341in}%
\pgfsys@useobject{currentmarker}{}%
\end{pgfscope}%
\begin{pgfscope}%
\pgfsys@transformshift{3.568157in}{3.141064in}%
\pgfsys@useobject{currentmarker}{}%
\end{pgfscope}%
\begin{pgfscope}%
\pgfsys@transformshift{3.693895in}{3.122985in}%
\pgfsys@useobject{currentmarker}{}%
\end{pgfscope}%
\begin{pgfscope}%
\pgfsys@transformshift{3.658602in}{3.132219in}%
\pgfsys@useobject{currentmarker}{}%
\end{pgfscope}%
\begin{pgfscope}%
\pgfsys@transformshift{3.528462in}{3.130507in}%
\pgfsys@useobject{currentmarker}{}%
\end{pgfscope}%
\begin{pgfscope}%
\pgfsys@transformshift{3.506324in}{3.110379in}%
\pgfsys@useobject{currentmarker}{}%
\end{pgfscope}%
\begin{pgfscope}%
\pgfsys@transformshift{3.286071in}{3.115006in}%
\pgfsys@useobject{currentmarker}{}%
\end{pgfscope}%
\begin{pgfscope}%
\pgfsys@transformshift{3.339293in}{3.139992in}%
\pgfsys@useobject{currentmarker}{}%
\end{pgfscope}%
\begin{pgfscope}%
\pgfsys@transformshift{3.294329in}{3.129614in}%
\pgfsys@useobject{currentmarker}{}%
\end{pgfscope}%
\begin{pgfscope}%
\pgfsys@transformshift{3.393736in}{3.137936in}%
\pgfsys@useobject{currentmarker}{}%
\end{pgfscope}%
\begin{pgfscope}%
\pgfsys@transformshift{3.637233in}{3.109326in}%
\pgfsys@useobject{currentmarker}{}%
\end{pgfscope}%
\begin{pgfscope}%
\pgfsys@transformshift{3.564460in}{3.132857in}%
\pgfsys@useobject{currentmarker}{}%
\end{pgfscope}%
\begin{pgfscope}%
\pgfsys@transformshift{3.541411in}{3.139473in}%
\pgfsys@useobject{currentmarker}{}%
\end{pgfscope}%
\begin{pgfscope}%
\pgfsys@transformshift{3.548618in}{3.139519in}%
\pgfsys@useobject{currentmarker}{}%
\end{pgfscope}%
\begin{pgfscope}%
\pgfsys@transformshift{3.233167in}{3.090304in}%
\pgfsys@useobject{currentmarker}{}%
\end{pgfscope}%
\begin{pgfscope}%
\pgfsys@transformshift{3.640592in}{3.130908in}%
\pgfsys@useobject{currentmarker}{}%
\end{pgfscope}%
\begin{pgfscope}%
\pgfsys@transformshift{3.639509in}{3.126455in}%
\pgfsys@useobject{currentmarker}{}%
\end{pgfscope}%
\begin{pgfscope}%
\pgfsys@transformshift{3.328323in}{3.136852in}%
\pgfsys@useobject{currentmarker}{}%
\end{pgfscope}%
\begin{pgfscope}%
\pgfsys@transformshift{3.602049in}{3.110155in}%
\pgfsys@useobject{currentmarker}{}%
\end{pgfscope}%
\begin{pgfscope}%
\pgfsys@transformshift{3.524949in}{3.122655in}%
\pgfsys@useobject{currentmarker}{}%
\end{pgfscope}%
\begin{pgfscope}%
\pgfsys@transformshift{3.360027in}{3.138135in}%
\pgfsys@useobject{currentmarker}{}%
\end{pgfscope}%
\begin{pgfscope}%
\pgfsys@transformshift{3.336496in}{3.129486in}%
\pgfsys@useobject{currentmarker}{}%
\end{pgfscope}%
\begin{pgfscope}%
\pgfsys@transformshift{3.612575in}{3.129020in}%
\pgfsys@useobject{currentmarker}{}%
\end{pgfscope}%
\begin{pgfscope}%
\pgfsys@transformshift{3.395435in}{3.142603in}%
\pgfsys@useobject{currentmarker}{}%
\end{pgfscope}%
\begin{pgfscope}%
\pgfsys@transformshift{3.479778in}{3.124773in}%
\pgfsys@useobject{currentmarker}{}%
\end{pgfscope}%
\begin{pgfscope}%
\pgfsys@transformshift{3.787866in}{3.105376in}%
\pgfsys@useobject{currentmarker}{}%
\end{pgfscope}%
\begin{pgfscope}%
\pgfsys@transformshift{3.230400in}{3.127196in}%
\pgfsys@useobject{currentmarker}{}%
\end{pgfscope}%
\begin{pgfscope}%
\pgfsys@transformshift{3.633222in}{3.116976in}%
\pgfsys@useobject{currentmarker}{}%
\end{pgfscope}%
\begin{pgfscope}%
\pgfsys@transformshift{3.721386in}{3.111144in}%
\pgfsys@useobject{currentmarker}{}%
\end{pgfscope}%
\begin{pgfscope}%
\pgfsys@transformshift{3.280023in}{3.132108in}%
\pgfsys@useobject{currentmarker}{}%
\end{pgfscope}%
\begin{pgfscope}%
\pgfsys@transformshift{3.376970in}{3.142056in}%
\pgfsys@useobject{currentmarker}{}%
\end{pgfscope}%
\begin{pgfscope}%
\pgfsys@transformshift{3.530231in}{3.132401in}%
\pgfsys@useobject{currentmarker}{}%
\end{pgfscope}%
\begin{pgfscope}%
\pgfsys@transformshift{3.336306in}{3.134139in}%
\pgfsys@useobject{currentmarker}{}%
\end{pgfscope}%
\begin{pgfscope}%
\pgfsys@transformshift{3.506294in}{3.124907in}%
\pgfsys@useobject{currentmarker}{}%
\end{pgfscope}%
\begin{pgfscope}%
\pgfsys@transformshift{3.446091in}{3.132191in}%
\pgfsys@useobject{currentmarker}{}%
\end{pgfscope}%
\begin{pgfscope}%
\pgfsys@transformshift{3.534164in}{3.136664in}%
\pgfsys@useobject{currentmarker}{}%
\end{pgfscope}%
\begin{pgfscope}%
\pgfsys@transformshift{3.573962in}{3.132719in}%
\pgfsys@useobject{currentmarker}{}%
\end{pgfscope}%
\begin{pgfscope}%
\pgfsys@transformshift{3.527380in}{3.137891in}%
\pgfsys@useobject{currentmarker}{}%
\end{pgfscope}%
\begin{pgfscope}%
\pgfsys@transformshift{3.456012in}{3.131661in}%
\pgfsys@useobject{currentmarker}{}%
\end{pgfscope}%
\begin{pgfscope}%
\pgfsys@transformshift{3.646787in}{3.119528in}%
\pgfsys@useobject{currentmarker}{}%
\end{pgfscope}%
\begin{pgfscope}%
\pgfsys@transformshift{3.280840in}{3.132485in}%
\pgfsys@useobject{currentmarker}{}%
\end{pgfscope}%
\begin{pgfscope}%
\pgfsys@transformshift{3.455345in}{3.130078in}%
\pgfsys@useobject{currentmarker}{}%
\end{pgfscope}%
\begin{pgfscope}%
\pgfsys@transformshift{3.426341in}{3.117748in}%
\pgfsys@useobject{currentmarker}{}%
\end{pgfscope}%
\begin{pgfscope}%
\pgfsys@transformshift{3.478430in}{3.133458in}%
\pgfsys@useobject{currentmarker}{}%
\end{pgfscope}%
\begin{pgfscope}%
\pgfsys@transformshift{3.303296in}{3.138242in}%
\pgfsys@useobject{currentmarker}{}%
\end{pgfscope}%
\begin{pgfscope}%
\pgfsys@transformshift{3.181620in}{3.115752in}%
\pgfsys@useobject{currentmarker}{}%
\end{pgfscope}%
\begin{pgfscope}%
\pgfsys@transformshift{3.319608in}{3.131518in}%
\pgfsys@useobject{currentmarker}{}%
\end{pgfscope}%
\begin{pgfscope}%
\pgfsys@transformshift{3.586574in}{3.136219in}%
\pgfsys@useobject{currentmarker}{}%
\end{pgfscope}%
\begin{pgfscope}%
\pgfsys@transformshift{3.636667in}{3.127532in}%
\pgfsys@useobject{currentmarker}{}%
\end{pgfscope}%
\begin{pgfscope}%
\pgfsys@transformshift{3.462602in}{3.122865in}%
\pgfsys@useobject{currentmarker}{}%
\end{pgfscope}%
\begin{pgfscope}%
\pgfsys@transformshift{3.479514in}{3.126561in}%
\pgfsys@useobject{currentmarker}{}%
\end{pgfscope}%
\begin{pgfscope}%
\pgfsys@transformshift{3.247058in}{3.131433in}%
\pgfsys@useobject{currentmarker}{}%
\end{pgfscope}%
\begin{pgfscope}%
\pgfsys@transformshift{3.490483in}{3.135092in}%
\pgfsys@useobject{currentmarker}{}%
\end{pgfscope}%
\begin{pgfscope}%
\pgfsys@transformshift{3.554163in}{3.131885in}%
\pgfsys@useobject{currentmarker}{}%
\end{pgfscope}%
\begin{pgfscope}%
\pgfsys@transformshift{3.701158in}{3.113567in}%
\pgfsys@useobject{currentmarker}{}%
\end{pgfscope}%
\begin{pgfscope}%
\pgfsys@transformshift{3.483021in}{3.138265in}%
\pgfsys@useobject{currentmarker}{}%
\end{pgfscope}%
\begin{pgfscope}%
\pgfsys@transformshift{3.467402in}{3.141104in}%
\pgfsys@useobject{currentmarker}{}%
\end{pgfscope}%
\begin{pgfscope}%
\pgfsys@transformshift{3.406700in}{3.113062in}%
\pgfsys@useobject{currentmarker}{}%
\end{pgfscope}%
\begin{pgfscope}%
\pgfsys@transformshift{3.729889in}{3.121934in}%
\pgfsys@useobject{currentmarker}{}%
\end{pgfscope}%
\begin{pgfscope}%
\pgfsys@transformshift{3.596964in}{3.118201in}%
\pgfsys@useobject{currentmarker}{}%
\end{pgfscope}%
\begin{pgfscope}%
\pgfsys@transformshift{3.371057in}{3.139054in}%
\pgfsys@useobject{currentmarker}{}%
\end{pgfscope}%
\begin{pgfscope}%
\pgfsys@transformshift{3.462398in}{3.140704in}%
\pgfsys@useobject{currentmarker}{}%
\end{pgfscope}%
\begin{pgfscope}%
\pgfsys@transformshift{3.657480in}{3.105867in}%
\pgfsys@useobject{currentmarker}{}%
\end{pgfscope}%
\begin{pgfscope}%
\pgfsys@transformshift{3.437474in}{3.125218in}%
\pgfsys@useobject{currentmarker}{}%
\end{pgfscope}%
\begin{pgfscope}%
\pgfsys@transformshift{3.392936in}{3.119234in}%
\pgfsys@useobject{currentmarker}{}%
\end{pgfscope}%
\begin{pgfscope}%
\pgfsys@transformshift{3.275378in}{3.129485in}%
\pgfsys@useobject{currentmarker}{}%
\end{pgfscope}%
\begin{pgfscope}%
\pgfsys@transformshift{3.482585in}{3.141036in}%
\pgfsys@useobject{currentmarker}{}%
\end{pgfscope}%
\begin{pgfscope}%
\pgfsys@transformshift{3.703421in}{3.109783in}%
\pgfsys@useobject{currentmarker}{}%
\end{pgfscope}%
\begin{pgfscope}%
\pgfsys@transformshift{3.478721in}{3.118062in}%
\pgfsys@useobject{currentmarker}{}%
\end{pgfscope}%
\begin{pgfscope}%
\pgfsys@transformshift{3.569350in}{3.129971in}%
\pgfsys@useobject{currentmarker}{}%
\end{pgfscope}%
\begin{pgfscope}%
\pgfsys@transformshift{3.472248in}{3.130969in}%
\pgfsys@useobject{currentmarker}{}%
\end{pgfscope}%
\begin{pgfscope}%
\pgfsys@transformshift{3.506716in}{3.136502in}%
\pgfsys@useobject{currentmarker}{}%
\end{pgfscope}%
\begin{pgfscope}%
\pgfsys@transformshift{3.268482in}{3.128573in}%
\pgfsys@useobject{currentmarker}{}%
\end{pgfscope}%
\begin{pgfscope}%
\pgfsys@transformshift{3.529630in}{3.116171in}%
\pgfsys@useobject{currentmarker}{}%
\end{pgfscope}%
\begin{pgfscope}%
\pgfsys@transformshift{3.547714in}{3.125945in}%
\pgfsys@useobject{currentmarker}{}%
\end{pgfscope}%
\begin{pgfscope}%
\pgfsys@transformshift{3.445736in}{3.143173in}%
\pgfsys@useobject{currentmarker}{}%
\end{pgfscope}%
\begin{pgfscope}%
\pgfsys@transformshift{3.583228in}{3.136628in}%
\pgfsys@useobject{currentmarker}{}%
\end{pgfscope}%
\begin{pgfscope}%
\pgfsys@transformshift{3.698001in}{3.114342in}%
\pgfsys@useobject{currentmarker}{}%
\end{pgfscope}%
\begin{pgfscope}%
\pgfsys@transformshift{3.488813in}{3.138025in}%
\pgfsys@useobject{currentmarker}{}%
\end{pgfscope}%
\begin{pgfscope}%
\pgfsys@transformshift{3.340594in}{3.124447in}%
\pgfsys@useobject{currentmarker}{}%
\end{pgfscope}%
\begin{pgfscope}%
\pgfsys@transformshift{3.346882in}{3.137369in}%
\pgfsys@useobject{currentmarker}{}%
\end{pgfscope}%
\begin{pgfscope}%
\pgfsys@transformshift{3.053221in}{3.096091in}%
\pgfsys@useobject{currentmarker}{}%
\end{pgfscope}%
\begin{pgfscope}%
\pgfsys@transformshift{3.349859in}{3.139630in}%
\pgfsys@useobject{currentmarker}{}%
\end{pgfscope}%
\begin{pgfscope}%
\pgfsys@transformshift{3.527935in}{3.117531in}%
\pgfsys@useobject{currentmarker}{}%
\end{pgfscope}%
\begin{pgfscope}%
\pgfsys@transformshift{3.544569in}{2.074426in}%
\pgfsys@useobject{currentmarker}{}%
\end{pgfscope}%
\begin{pgfscope}%
\pgfsys@transformshift{3.252270in}{2.064306in}%
\pgfsys@useobject{currentmarker}{}%
\end{pgfscope}%
\begin{pgfscope}%
\pgfsys@transformshift{3.384261in}{2.054804in}%
\pgfsys@useobject{currentmarker}{}%
\end{pgfscope}%
\begin{pgfscope}%
\pgfsys@transformshift{3.600121in}{2.058568in}%
\pgfsys@useobject{currentmarker}{}%
\end{pgfscope}%
\begin{pgfscope}%
\pgfsys@transformshift{3.402183in}{2.062830in}%
\pgfsys@useobject{currentmarker}{}%
\end{pgfscope}%
\begin{pgfscope}%
\pgfsys@transformshift{3.648778in}{2.065983in}%
\pgfsys@useobject{currentmarker}{}%
\end{pgfscope}%
\begin{pgfscope}%
\pgfsys@transformshift{3.026347in}{2.099422in}%
\pgfsys@useobject{currentmarker}{}%
\end{pgfscope}%
\begin{pgfscope}%
\pgfsys@transformshift{3.491053in}{2.056659in}%
\pgfsys@useobject{currentmarker}{}%
\end{pgfscope}%
\begin{pgfscope}%
\pgfsys@transformshift{3.124001in}{2.087452in}%
\pgfsys@useobject{currentmarker}{}%
\end{pgfscope}%
\begin{pgfscope}%
\pgfsys@transformshift{3.295734in}{2.081494in}%
\pgfsys@useobject{currentmarker}{}%
\end{pgfscope}%
\begin{pgfscope}%
\pgfsys@transformshift{3.352935in}{2.061568in}%
\pgfsys@useobject{currentmarker}{}%
\end{pgfscope}%
\begin{pgfscope}%
\pgfsys@transformshift{3.544909in}{2.058691in}%
\pgfsys@useobject{currentmarker}{}%
\end{pgfscope}%
\begin{pgfscope}%
\pgfsys@transformshift{3.457182in}{2.059319in}%
\pgfsys@useobject{currentmarker}{}%
\end{pgfscope}%
\begin{pgfscope}%
\pgfsys@transformshift{3.315717in}{2.067815in}%
\pgfsys@useobject{currentmarker}{}%
\end{pgfscope}%
\begin{pgfscope}%
\pgfsys@transformshift{3.742262in}{2.078421in}%
\pgfsys@useobject{currentmarker}{}%
\end{pgfscope}%
\begin{pgfscope}%
\pgfsys@transformshift{3.802068in}{2.092804in}%
\pgfsys@useobject{currentmarker}{}%
\end{pgfscope}%
\begin{pgfscope}%
\pgfsys@transformshift{3.393037in}{2.057703in}%
\pgfsys@useobject{currentmarker}{}%
\end{pgfscope}%
\begin{pgfscope}%
\pgfsys@transformshift{3.487758in}{2.053837in}%
\pgfsys@useobject{currentmarker}{}%
\end{pgfscope}%
\begin{pgfscope}%
\pgfsys@transformshift{3.235391in}{2.075112in}%
\pgfsys@useobject{currentmarker}{}%
\end{pgfscope}%
\begin{pgfscope}%
\pgfsys@transformshift{3.377803in}{2.057321in}%
\pgfsys@useobject{currentmarker}{}%
\end{pgfscope}%
\begin{pgfscope}%
\pgfsys@transformshift{3.199305in}{2.069296in}%
\pgfsys@useobject{currentmarker}{}%
\end{pgfscope}%
\begin{pgfscope}%
\pgfsys@transformshift{3.385833in}{2.072352in}%
\pgfsys@useobject{currentmarker}{}%
\end{pgfscope}%
\begin{pgfscope}%
\pgfsys@transformshift{3.689814in}{2.066349in}%
\pgfsys@useobject{currentmarker}{}%
\end{pgfscope}%
\begin{pgfscope}%
\pgfsys@transformshift{3.890716in}{2.103970in}%
\pgfsys@useobject{currentmarker}{}%
\end{pgfscope}%
\begin{pgfscope}%
\pgfsys@transformshift{3.607966in}{2.068094in}%
\pgfsys@useobject{currentmarker}{}%
\end{pgfscope}%
\begin{pgfscope}%
\pgfsys@transformshift{3.407525in}{2.077213in}%
\pgfsys@useobject{currentmarker}{}%
\end{pgfscope}%
\begin{pgfscope}%
\pgfsys@transformshift{3.582790in}{2.061103in}%
\pgfsys@useobject{currentmarker}{}%
\end{pgfscope}%
\begin{pgfscope}%
\pgfsys@transformshift{3.310196in}{2.066418in}%
\pgfsys@useobject{currentmarker}{}%
\end{pgfscope}%
\begin{pgfscope}%
\pgfsys@transformshift{3.578784in}{2.057247in}%
\pgfsys@useobject{currentmarker}{}%
\end{pgfscope}%
\begin{pgfscope}%
\pgfsys@transformshift{3.442773in}{2.056495in}%
\pgfsys@useobject{currentmarker}{}%
\end{pgfscope}%
\begin{pgfscope}%
\pgfsys@transformshift{3.429710in}{2.072154in}%
\pgfsys@useobject{currentmarker}{}%
\end{pgfscope}%
\begin{pgfscope}%
\pgfsys@transformshift{3.478070in}{2.073533in}%
\pgfsys@useobject{currentmarker}{}%
\end{pgfscope}%
\begin{pgfscope}%
\pgfsys@transformshift{3.458786in}{2.055237in}%
\pgfsys@useobject{currentmarker}{}%
\end{pgfscope}%
\begin{pgfscope}%
\pgfsys@transformshift{3.266051in}{2.066380in}%
\pgfsys@useobject{currentmarker}{}%
\end{pgfscope}%
\begin{pgfscope}%
\pgfsys@transformshift{3.514230in}{2.084192in}%
\pgfsys@useobject{currentmarker}{}%
\end{pgfscope}%
\begin{pgfscope}%
\pgfsys@transformshift{3.859108in}{2.090650in}%
\pgfsys@useobject{currentmarker}{}%
\end{pgfscope}%
\begin{pgfscope}%
\pgfsys@transformshift{3.465886in}{2.053218in}%
\pgfsys@useobject{currentmarker}{}%
\end{pgfscope}%
\begin{pgfscope}%
\pgfsys@transformshift{3.254478in}{2.064419in}%
\pgfsys@useobject{currentmarker}{}%
\end{pgfscope}%
\begin{pgfscope}%
\pgfsys@transformshift{3.606178in}{2.089206in}%
\pgfsys@useobject{currentmarker}{}%
\end{pgfscope}%
\begin{pgfscope}%
\pgfsys@transformshift{3.639383in}{2.066185in}%
\pgfsys@useobject{currentmarker}{}%
\end{pgfscope}%
\begin{pgfscope}%
\pgfsys@transformshift{3.609571in}{2.066377in}%
\pgfsys@useobject{currentmarker}{}%
\end{pgfscope}%
\begin{pgfscope}%
\pgfsys@transformshift{3.251860in}{2.068308in}%
\pgfsys@useobject{currentmarker}{}%
\end{pgfscope}%
\begin{pgfscope}%
\pgfsys@transformshift{3.423439in}{2.054837in}%
\pgfsys@useobject{currentmarker}{}%
\end{pgfscope}%
\begin{pgfscope}%
\pgfsys@transformshift{3.446713in}{2.054921in}%
\pgfsys@useobject{currentmarker}{}%
\end{pgfscope}%
\begin{pgfscope}%
\pgfsys@transformshift{3.692479in}{2.081801in}%
\pgfsys@useobject{currentmarker}{}%
\end{pgfscope}%
\begin{pgfscope}%
\pgfsys@transformshift{3.165315in}{2.079943in}%
\pgfsys@useobject{currentmarker}{}%
\end{pgfscope}%
\begin{pgfscope}%
\pgfsys@transformshift{3.463166in}{2.069144in}%
\pgfsys@useobject{currentmarker}{}%
\end{pgfscope}%
\begin{pgfscope}%
\pgfsys@transformshift{3.577527in}{2.058145in}%
\pgfsys@useobject{currentmarker}{}%
\end{pgfscope}%
\begin{pgfscope}%
\pgfsys@transformshift{3.442986in}{2.063244in}%
\pgfsys@useobject{currentmarker}{}%
\end{pgfscope}%
\begin{pgfscope}%
\pgfsys@transformshift{3.505030in}{2.060687in}%
\pgfsys@useobject{currentmarker}{}%
\end{pgfscope}%
\begin{pgfscope}%
\pgfsys@transformshift{3.400502in}{2.064498in}%
\pgfsys@useobject{currentmarker}{}%
\end{pgfscope}%
\begin{pgfscope}%
\pgfsys@transformshift{3.236200in}{2.076687in}%
\pgfsys@useobject{currentmarker}{}%
\end{pgfscope}%
\begin{pgfscope}%
\pgfsys@transformshift{3.556028in}{2.058552in}%
\pgfsys@useobject{currentmarker}{}%
\end{pgfscope}%
\begin{pgfscope}%
\pgfsys@transformshift{3.262173in}{2.099995in}%
\pgfsys@useobject{currentmarker}{}%
\end{pgfscope}%
\begin{pgfscope}%
\pgfsys@transformshift{3.574435in}{2.061438in}%
\pgfsys@useobject{currentmarker}{}%
\end{pgfscope}%
\begin{pgfscope}%
\pgfsys@transformshift{3.391181in}{2.077483in}%
\pgfsys@useobject{currentmarker}{}%
\end{pgfscope}%
\begin{pgfscope}%
\pgfsys@transformshift{3.501906in}{2.067285in}%
\pgfsys@useobject{currentmarker}{}%
\end{pgfscope}%
\begin{pgfscope}%
\pgfsys@transformshift{3.479502in}{2.053946in}%
\pgfsys@useobject{currentmarker}{}%
\end{pgfscope}%
\begin{pgfscope}%
\pgfsys@transformshift{3.341377in}{2.065450in}%
\pgfsys@useobject{currentmarker}{}%
\end{pgfscope}%
\begin{pgfscope}%
\pgfsys@transformshift{3.812305in}{2.088831in}%
\pgfsys@useobject{currentmarker}{}%
\end{pgfscope}%
\begin{pgfscope}%
\pgfsys@transformshift{3.537624in}{2.068503in}%
\pgfsys@useobject{currentmarker}{}%
\end{pgfscope}%
\begin{pgfscope}%
\pgfsys@transformshift{3.258055in}{2.062357in}%
\pgfsys@useobject{currentmarker}{}%
\end{pgfscope}%
\begin{pgfscope}%
\pgfsys@transformshift{3.693034in}{2.088867in}%
\pgfsys@useobject{currentmarker}{}%
\end{pgfscope}%
\begin{pgfscope}%
\pgfsys@transformshift{3.181328in}{2.074746in}%
\pgfsys@useobject{currentmarker}{}%
\end{pgfscope}%
\begin{pgfscope}%
\pgfsys@transformshift{3.236173in}{2.097962in}%
\pgfsys@useobject{currentmarker}{}%
\end{pgfscope}%
\begin{pgfscope}%
\pgfsys@transformshift{3.392138in}{2.065423in}%
\pgfsys@useobject{currentmarker}{}%
\end{pgfscope}%
\begin{pgfscope}%
\pgfsys@transformshift{3.492407in}{2.069956in}%
\pgfsys@useobject{currentmarker}{}%
\end{pgfscope}%
\begin{pgfscope}%
\pgfsys@transformshift{3.259816in}{2.066388in}%
\pgfsys@useobject{currentmarker}{}%
\end{pgfscope}%
\begin{pgfscope}%
\pgfsys@transformshift{3.450311in}{2.072932in}%
\pgfsys@useobject{currentmarker}{}%
\end{pgfscope}%
\begin{pgfscope}%
\pgfsys@transformshift{3.651627in}{2.062440in}%
\pgfsys@useobject{currentmarker}{}%
\end{pgfscope}%
\begin{pgfscope}%
\pgfsys@transformshift{3.305859in}{2.078647in}%
\pgfsys@useobject{currentmarker}{}%
\end{pgfscope}%
\begin{pgfscope}%
\pgfsys@transformshift{3.485040in}{2.056214in}%
\pgfsys@useobject{currentmarker}{}%
\end{pgfscope}%
\begin{pgfscope}%
\pgfsys@transformshift{3.435671in}{2.054528in}%
\pgfsys@useobject{currentmarker}{}%
\end{pgfscope}%
\begin{pgfscope}%
\pgfsys@transformshift{3.252662in}{2.066679in}%
\pgfsys@useobject{currentmarker}{}%
\end{pgfscope}%
\begin{pgfscope}%
\pgfsys@transformshift{3.403120in}{2.067445in}%
\pgfsys@useobject{currentmarker}{}%
\end{pgfscope}%
\begin{pgfscope}%
\pgfsys@transformshift{3.396709in}{2.057489in}%
\pgfsys@useobject{currentmarker}{}%
\end{pgfscope}%
\begin{pgfscope}%
\pgfsys@transformshift{3.382009in}{2.073162in}%
\pgfsys@useobject{currentmarker}{}%
\end{pgfscope}%
\begin{pgfscope}%
\pgfsys@transformshift{3.479262in}{2.058914in}%
\pgfsys@useobject{currentmarker}{}%
\end{pgfscope}%
\begin{pgfscope}%
\pgfsys@transformshift{3.289933in}{2.073260in}%
\pgfsys@useobject{currentmarker}{}%
\end{pgfscope}%
\begin{pgfscope}%
\pgfsys@transformshift{3.535689in}{2.060923in}%
\pgfsys@useobject{currentmarker}{}%
\end{pgfscope}%
\begin{pgfscope}%
\pgfsys@transformshift{3.612334in}{2.071540in}%
\pgfsys@useobject{currentmarker}{}%
\end{pgfscope}%
\begin{pgfscope}%
\pgfsys@transformshift{3.603922in}{2.063764in}%
\pgfsys@useobject{currentmarker}{}%
\end{pgfscope}%
\begin{pgfscope}%
\pgfsys@transformshift{3.362868in}{2.055649in}%
\pgfsys@useobject{currentmarker}{}%
\end{pgfscope}%
\begin{pgfscope}%
\pgfsys@transformshift{3.166573in}{2.089232in}%
\pgfsys@useobject{currentmarker}{}%
\end{pgfscope}%
\begin{pgfscope}%
\pgfsys@transformshift{3.437132in}{2.058812in}%
\pgfsys@useobject{currentmarker}{}%
\end{pgfscope}%
\begin{pgfscope}%
\pgfsys@transformshift{3.466018in}{2.065376in}%
\pgfsys@useobject{currentmarker}{}%
\end{pgfscope}%
\begin{pgfscope}%
\pgfsys@transformshift{3.566479in}{2.070033in}%
\pgfsys@useobject{currentmarker}{}%
\end{pgfscope}%
\begin{pgfscope}%
\pgfsys@transformshift{3.664461in}{2.070981in}%
\pgfsys@useobject{currentmarker}{}%
\end{pgfscope}%
\begin{pgfscope}%
\pgfsys@transformshift{3.320526in}{2.087126in}%
\pgfsys@useobject{currentmarker}{}%
\end{pgfscope}%
\begin{pgfscope}%
\pgfsys@transformshift{3.183166in}{2.074506in}%
\pgfsys@useobject{currentmarker}{}%
\end{pgfscope}%
\begin{pgfscope}%
\pgfsys@transformshift{3.248562in}{2.063259in}%
\pgfsys@useobject{currentmarker}{}%
\end{pgfscope}%
\begin{pgfscope}%
\pgfsys@transformshift{3.284176in}{2.084092in}%
\pgfsys@useobject{currentmarker}{}%
\end{pgfscope}%
\begin{pgfscope}%
\pgfsys@transformshift{3.318415in}{2.062136in}%
\pgfsys@useobject{currentmarker}{}%
\end{pgfscope}%
\begin{pgfscope}%
\pgfsys@transformshift{3.514701in}{2.054449in}%
\pgfsys@useobject{currentmarker}{}%
\end{pgfscope}%
\begin{pgfscope}%
\pgfsys@transformshift{3.473094in}{2.053701in}%
\pgfsys@useobject{currentmarker}{}%
\end{pgfscope}%
\begin{pgfscope}%
\pgfsys@transformshift{3.586366in}{2.056700in}%
\pgfsys@useobject{currentmarker}{}%
\end{pgfscope}%
\begin{pgfscope}%
\pgfsys@transformshift{3.419980in}{2.058828in}%
\pgfsys@useobject{currentmarker}{}%
\end{pgfscope}%
\begin{pgfscope}%
\pgfsys@transformshift{3.192575in}{2.078370in}%
\pgfsys@useobject{currentmarker}{}%
\end{pgfscope}%
\begin{pgfscope}%
\pgfsys@transformshift{3.469288in}{2.054148in}%
\pgfsys@useobject{currentmarker}{}%
\end{pgfscope}%
\begin{pgfscope}%
\pgfsys@transformshift{3.745501in}{2.077052in}%
\pgfsys@useobject{currentmarker}{}%
\end{pgfscope}%
\begin{pgfscope}%
\pgfsys@transformshift{3.571545in}{2.057242in}%
\pgfsys@useobject{currentmarker}{}%
\end{pgfscope}%
\begin{pgfscope}%
\pgfsys@transformshift{3.341856in}{2.057008in}%
\pgfsys@useobject{currentmarker}{}%
\end{pgfscope}%
\begin{pgfscope}%
\pgfsys@transformshift{3.280871in}{2.063289in}%
\pgfsys@useobject{currentmarker}{}%
\end{pgfscope}%
\begin{pgfscope}%
\pgfsys@transformshift{3.436502in}{2.065916in}%
\pgfsys@useobject{currentmarker}{}%
\end{pgfscope}%
\begin{pgfscope}%
\pgfsys@transformshift{3.679846in}{2.077168in}%
\pgfsys@useobject{currentmarker}{}%
\end{pgfscope}%
\begin{pgfscope}%
\pgfsys@transformshift{3.657390in}{2.065965in}%
\pgfsys@useobject{currentmarker}{}%
\end{pgfscope}%
\begin{pgfscope}%
\pgfsys@transformshift{3.229647in}{2.073132in}%
\pgfsys@useobject{currentmarker}{}%
\end{pgfscope}%
\begin{pgfscope}%
\pgfsys@transformshift{3.379778in}{2.081732in}%
\pgfsys@useobject{currentmarker}{}%
\end{pgfscope}%
\begin{pgfscope}%
\pgfsys@transformshift{3.253094in}{2.081560in}%
\pgfsys@useobject{currentmarker}{}%
\end{pgfscope}%
\begin{pgfscope}%
\pgfsys@transformshift{3.618879in}{2.058962in}%
\pgfsys@useobject{currentmarker}{}%
\end{pgfscope}%
\begin{pgfscope}%
\pgfsys@transformshift{3.415227in}{2.055754in}%
\pgfsys@useobject{currentmarker}{}%
\end{pgfscope}%
\begin{pgfscope}%
\pgfsys@transformshift{3.796141in}{2.080592in}%
\pgfsys@useobject{currentmarker}{}%
\end{pgfscope}%
\begin{pgfscope}%
\pgfsys@transformshift{3.427336in}{2.100412in}%
\pgfsys@useobject{currentmarker}{}%
\end{pgfscope}%
\begin{pgfscope}%
\pgfsys@transformshift{3.632620in}{2.074212in}%
\pgfsys@useobject{currentmarker}{}%
\end{pgfscope}%
\begin{pgfscope}%
\pgfsys@transformshift{3.549018in}{2.081108in}%
\pgfsys@useobject{currentmarker}{}%
\end{pgfscope}%
\begin{pgfscope}%
\pgfsys@transformshift{3.348385in}{2.077785in}%
\pgfsys@useobject{currentmarker}{}%
\end{pgfscope}%
\begin{pgfscope}%
\pgfsys@transformshift{3.594861in}{2.064314in}%
\pgfsys@useobject{currentmarker}{}%
\end{pgfscope}%
\begin{pgfscope}%
\pgfsys@transformshift{3.379197in}{2.068696in}%
\pgfsys@useobject{currentmarker}{}%
\end{pgfscope}%
\begin{pgfscope}%
\pgfsys@transformshift{3.254477in}{2.068490in}%
\pgfsys@useobject{currentmarker}{}%
\end{pgfscope}%
\begin{pgfscope}%
\pgfsys@transformshift{3.142354in}{2.076907in}%
\pgfsys@useobject{currentmarker}{}%
\end{pgfscope}%
\begin{pgfscope}%
\pgfsys@transformshift{3.495934in}{2.062252in}%
\pgfsys@useobject{currentmarker}{}%
\end{pgfscope}%
\begin{pgfscope}%
\pgfsys@transformshift{3.332511in}{2.073175in}%
\pgfsys@useobject{currentmarker}{}%
\end{pgfscope}%
\begin{pgfscope}%
\pgfsys@transformshift{3.495642in}{2.069401in}%
\pgfsys@useobject{currentmarker}{}%
\end{pgfscope}%
\begin{pgfscope}%
\pgfsys@transformshift{3.442228in}{2.067501in}%
\pgfsys@useobject{currentmarker}{}%
\end{pgfscope}%
\begin{pgfscope}%
\pgfsys@transformshift{3.428066in}{2.056277in}%
\pgfsys@useobject{currentmarker}{}%
\end{pgfscope}%
\begin{pgfscope}%
\pgfsys@transformshift{3.648438in}{2.075016in}%
\pgfsys@useobject{currentmarker}{}%
\end{pgfscope}%
\begin{pgfscope}%
\pgfsys@transformshift{3.509221in}{2.060955in}%
\pgfsys@useobject{currentmarker}{}%
\end{pgfscope}%
\begin{pgfscope}%
\pgfsys@transformshift{3.393011in}{2.072424in}%
\pgfsys@useobject{currentmarker}{}%
\end{pgfscope}%
\begin{pgfscope}%
\pgfsys@transformshift{3.588970in}{2.056557in}%
\pgfsys@useobject{currentmarker}{}%
\end{pgfscope}%
\begin{pgfscope}%
\pgfsys@transformshift{3.527820in}{2.069558in}%
\pgfsys@useobject{currentmarker}{}%
\end{pgfscope}%
\begin{pgfscope}%
\pgfsys@transformshift{3.704468in}{2.066348in}%
\pgfsys@useobject{currentmarker}{}%
\end{pgfscope}%
\begin{pgfscope}%
\pgfsys@transformshift{3.529965in}{2.057025in}%
\pgfsys@useobject{currentmarker}{}%
\end{pgfscope}%
\begin{pgfscope}%
\pgfsys@transformshift{3.626711in}{2.062055in}%
\pgfsys@useobject{currentmarker}{}%
\end{pgfscope}%
\begin{pgfscope}%
\pgfsys@transformshift{3.195005in}{2.084397in}%
\pgfsys@useobject{currentmarker}{}%
\end{pgfscope}%
\begin{pgfscope}%
\pgfsys@transformshift{3.401029in}{2.080939in}%
\pgfsys@useobject{currentmarker}{}%
\end{pgfscope}%
\begin{pgfscope}%
\pgfsys@transformshift{3.436572in}{2.053200in}%
\pgfsys@useobject{currentmarker}{}%
\end{pgfscope}%
\begin{pgfscope}%
\pgfsys@transformshift{3.595167in}{2.094355in}%
\pgfsys@useobject{currentmarker}{}%
\end{pgfscope}%
\begin{pgfscope}%
\pgfsys@transformshift{3.741857in}{2.072741in}%
\pgfsys@useobject{currentmarker}{}%
\end{pgfscope}%
\begin{pgfscope}%
\pgfsys@transformshift{3.555560in}{2.054571in}%
\pgfsys@useobject{currentmarker}{}%
\end{pgfscope}%
\begin{pgfscope}%
\pgfsys@transformshift{3.450674in}{2.099435in}%
\pgfsys@useobject{currentmarker}{}%
\end{pgfscope}%
\begin{pgfscope}%
\pgfsys@transformshift{3.518869in}{2.079338in}%
\pgfsys@useobject{currentmarker}{}%
\end{pgfscope}%
\begin{pgfscope}%
\pgfsys@transformshift{3.145941in}{2.076525in}%
\pgfsys@useobject{currentmarker}{}%
\end{pgfscope}%
\begin{pgfscope}%
\pgfsys@transformshift{3.396166in}{2.085192in}%
\pgfsys@useobject{currentmarker}{}%
\end{pgfscope}%
\begin{pgfscope}%
\pgfsys@transformshift{3.661801in}{2.067260in}%
\pgfsys@useobject{currentmarker}{}%
\end{pgfscope}%
\begin{pgfscope}%
\pgfsys@transformshift{3.577243in}{2.064016in}%
\pgfsys@useobject{currentmarker}{}%
\end{pgfscope}%
\begin{pgfscope}%
\pgfsys@transformshift{3.201894in}{2.072097in}%
\pgfsys@useobject{currentmarker}{}%
\end{pgfscope}%
\begin{pgfscope}%
\pgfsys@transformshift{3.691111in}{2.083107in}%
\pgfsys@useobject{currentmarker}{}%
\end{pgfscope}%
\begin{pgfscope}%
\pgfsys@transformshift{3.435850in}{2.058721in}%
\pgfsys@useobject{currentmarker}{}%
\end{pgfscope}%
\begin{pgfscope}%
\pgfsys@transformshift{3.488435in}{2.058881in}%
\pgfsys@useobject{currentmarker}{}%
\end{pgfscope}%
\begin{pgfscope}%
\pgfsys@transformshift{3.555843in}{2.064853in}%
\pgfsys@useobject{currentmarker}{}%
\end{pgfscope}%
\end{pgfscope}%
\begin{pgfscope}%
\pgfpathmoveto{\pgfqpoint{3.460000in}{2.052407in}}%
\pgfpathcurveto{\pgfqpoint{3.749602in}{2.052407in}}{\pgfqpoint{4.027381in}{2.109937in}}{\pgfqpoint{4.232161in}{2.212326in}}%
\pgfpathcurveto{\pgfqpoint{4.436940in}{2.314716in}}{\pgfqpoint{4.552000in}{2.453606in}}{\pgfqpoint{4.552000in}{2.598407in}}%
\pgfpathcurveto{\pgfqpoint{4.552000in}{2.743208in}}{\pgfqpoint{4.436940in}{2.882097in}}{\pgfqpoint{4.232161in}{2.984487in}}%
\pgfpathcurveto{\pgfqpoint{4.027381in}{3.086877in}}{\pgfqpoint{3.749602in}{3.144407in}}{\pgfqpoint{3.460000in}{3.144407in}}%
\pgfpathcurveto{\pgfqpoint{3.170398in}{3.144407in}}{\pgfqpoint{2.892619in}{3.086877in}}{\pgfqpoint{2.687839in}{2.984487in}}%
\pgfpathcurveto{\pgfqpoint{2.483060in}{2.882097in}}{\pgfqpoint{2.368000in}{2.743208in}}{\pgfqpoint{2.368000in}{2.598407in}}%
\pgfpathcurveto{\pgfqpoint{2.368000in}{2.453606in}}{\pgfqpoint{2.483060in}{2.314716in}}{\pgfqpoint{2.687839in}{2.212326in}}%
\pgfpathcurveto{\pgfqpoint{2.892619in}{2.109937in}}{\pgfqpoint{3.170398in}{2.052407in}}{\pgfqpoint{3.460000in}{2.052407in}}%
\pgfpathlineto{\pgfqpoint{3.460000in}{2.052407in}}%
\pgfpathclose%
\pgfusepath{clip}%
\pgfsetrectcap%
\pgfsetroundjoin%
\pgfsetlinewidth{0.451687pt}%
\definecolor{currentstroke}{rgb}{0.690196,0.690196,0.690196}%
\pgfsetstrokecolor{currentstroke}%
\pgfsetstrokeopacity{0.250000}%
\pgfsetdash{}{0pt}%
\pgfpathmoveto{\pgfqpoint{3.074921in}{2.103702in}}%
\pgfpathlineto{\pgfqpoint{3.043157in}{2.113059in}}%
\pgfpathlineto{\pgfqpoint{3.012700in}{2.122919in}}%
\pgfpathlineto{\pgfqpoint{2.984255in}{2.132965in}}%
\pgfpathlineto{\pgfqpoint{2.957090in}{2.143363in}}%
\pgfpathlineto{\pgfqpoint{2.931134in}{2.154082in}}%
\pgfpathlineto{\pgfqpoint{2.906319in}{2.165102in}}%
\pgfpathlineto{\pgfqpoint{2.882585in}{2.176401in}}%
\pgfpathlineto{\pgfqpoint{2.859881in}{2.187962in}}%
\pgfpathlineto{\pgfqpoint{2.838160in}{2.199772in}}%
\pgfpathlineto{\pgfqpoint{2.817384in}{2.211817in}}%
\pgfpathlineto{\pgfqpoint{2.797517in}{2.224083in}}%
\pgfpathlineto{\pgfqpoint{2.778532in}{2.236561in}}%
\pgfpathlineto{\pgfqpoint{2.760403in}{2.249238in}}%
\pgfpathlineto{\pgfqpoint{2.743106in}{2.262105in}}%
\pgfpathlineto{\pgfqpoint{2.726426in}{2.275314in}}%
\pgfpathlineto{\pgfqpoint{2.710806in}{2.288486in}}%
\pgfpathlineto{\pgfqpoint{2.695948in}{2.301831in}}%
\pgfpathlineto{\pgfqpoint{2.681847in}{2.315338in}}%
\pgfpathlineto{\pgfqpoint{2.668494in}{2.328996in}}%
\pgfpathlineto{\pgfqpoint{2.655882in}{2.342794in}}%
\pgfpathlineto{\pgfqpoint{2.644004in}{2.356724in}}%
\pgfpathlineto{\pgfqpoint{2.632855in}{2.370777in}}%
\pgfpathlineto{\pgfqpoint{2.622428in}{2.384944in}}%
\pgfpathlineto{\pgfqpoint{2.612718in}{2.399216in}}%
\pgfpathlineto{\pgfqpoint{2.603720in}{2.413586in}}%
\pgfpathlineto{\pgfqpoint{2.595430in}{2.428046in}}%
\pgfpathlineto{\pgfqpoint{2.587844in}{2.442588in}}%
\pgfpathlineto{\pgfqpoint{2.580957in}{2.457204in}}%
\pgfpathlineto{\pgfqpoint{2.574768in}{2.471888in}}%
\pgfpathlineto{\pgfqpoint{2.569272in}{2.486632in}}%
\pgfpathlineto{\pgfqpoint{2.564469in}{2.501429in}}%
\pgfpathlineto{\pgfqpoint{2.560314in}{2.516436in}}%
\pgfpathlineto{\pgfqpoint{2.556911in}{2.531243in}}%
\pgfpathlineto{\pgfqpoint{2.554184in}{2.546108in}}%
\pgfpathlineto{\pgfqpoint{2.552136in}{2.561018in}}%
\pgfpathlineto{\pgfqpoint{2.550769in}{2.575961in}}%
\pgfpathlineto{\pgfqpoint{2.550085in}{2.590923in}}%
\pgfpathlineto{\pgfqpoint{2.550085in}{2.605891in}}%
\pgfpathlineto{\pgfqpoint{2.550769in}{2.620853in}}%
\pgfpathlineto{\pgfqpoint{2.552136in}{2.635795in}}%
\pgfpathlineto{\pgfqpoint{2.554184in}{2.650705in}}%
\pgfpathlineto{\pgfqpoint{2.556911in}{2.665570in}}%
\pgfpathlineto{\pgfqpoint{2.560314in}{2.680378in}}%
\pgfpathlineto{\pgfqpoint{2.564469in}{2.695384in}}%
\pgfpathlineto{\pgfqpoint{2.569272in}{2.710181in}}%
\pgfpathlineto{\pgfqpoint{2.574768in}{2.724925in}}%
\pgfpathlineto{\pgfqpoint{2.580957in}{2.739609in}}%
\pgfpathlineto{\pgfqpoint{2.587844in}{2.754226in}}%
\pgfpathlineto{\pgfqpoint{2.595430in}{2.768768in}}%
\pgfpathlineto{\pgfqpoint{2.603720in}{2.783228in}}%
\pgfpathlineto{\pgfqpoint{2.612718in}{2.797598in}}%
\pgfpathlineto{\pgfqpoint{2.622428in}{2.811870in}}%
\pgfpathlineto{\pgfqpoint{2.632855in}{2.826037in}}%
\pgfpathlineto{\pgfqpoint{2.644004in}{2.840089in}}%
\pgfpathlineto{\pgfqpoint{2.655882in}{2.854019in}}%
\pgfpathlineto{\pgfqpoint{2.668494in}{2.867818in}}%
\pgfpathlineto{\pgfqpoint{2.681847in}{2.881476in}}%
\pgfpathlineto{\pgfqpoint{2.695948in}{2.894982in}}%
\pgfpathlineto{\pgfqpoint{2.710806in}{2.908327in}}%
\pgfpathlineto{\pgfqpoint{2.726426in}{2.921500in}}%
\pgfpathlineto{\pgfqpoint{2.743106in}{2.934708in}}%
\pgfpathlineto{\pgfqpoint{2.760403in}{2.947576in}}%
\pgfpathlineto{\pgfqpoint{2.778532in}{2.960253in}}%
\pgfpathlineto{\pgfqpoint{2.797517in}{2.972730in}}%
\pgfpathlineto{\pgfqpoint{2.817384in}{2.984997in}}%
\pgfpathlineto{\pgfqpoint{2.838160in}{2.997041in}}%
\pgfpathlineto{\pgfqpoint{2.859881in}{3.008851in}}%
\pgfpathlineto{\pgfqpoint{2.882585in}{3.020413in}}%
\pgfpathlineto{\pgfqpoint{2.906319in}{3.031712in}}%
\pgfpathlineto{\pgfqpoint{2.931134in}{3.042731in}}%
\pgfpathlineto{\pgfqpoint{2.957090in}{3.053451in}}%
\pgfpathlineto{\pgfqpoint{2.984255in}{3.063848in}}%
\pgfpathlineto{\pgfqpoint{3.012700in}{3.073894in}}%
\pgfpathlineto{\pgfqpoint{3.043157in}{3.083755in}}%
\pgfpathlineto{\pgfqpoint{3.074921in}{3.093112in}}%
\pgfusepath{stroke}%
\end{pgfscope}%
\begin{pgfscope}%
\pgfpathmoveto{\pgfqpoint{3.460000in}{2.052407in}}%
\pgfpathcurveto{\pgfqpoint{3.749602in}{2.052407in}}{\pgfqpoint{4.027381in}{2.109937in}}{\pgfqpoint{4.232161in}{2.212326in}}%
\pgfpathcurveto{\pgfqpoint{4.436940in}{2.314716in}}{\pgfqpoint{4.552000in}{2.453606in}}{\pgfqpoint{4.552000in}{2.598407in}}%
\pgfpathcurveto{\pgfqpoint{4.552000in}{2.743208in}}{\pgfqpoint{4.436940in}{2.882097in}}{\pgfqpoint{4.232161in}{2.984487in}}%
\pgfpathcurveto{\pgfqpoint{4.027381in}{3.086877in}}{\pgfqpoint{3.749602in}{3.144407in}}{\pgfqpoint{3.460000in}{3.144407in}}%
\pgfpathcurveto{\pgfqpoint{3.170398in}{3.144407in}}{\pgfqpoint{2.892619in}{3.086877in}}{\pgfqpoint{2.687839in}{2.984487in}}%
\pgfpathcurveto{\pgfqpoint{2.483060in}{2.882097in}}{\pgfqpoint{2.368000in}{2.743208in}}{\pgfqpoint{2.368000in}{2.598407in}}%
\pgfpathcurveto{\pgfqpoint{2.368000in}{2.453606in}}{\pgfqpoint{2.483060in}{2.314716in}}{\pgfqpoint{2.687839in}{2.212326in}}%
\pgfpathcurveto{\pgfqpoint{2.892619in}{2.109937in}}{\pgfqpoint{3.170398in}{2.052407in}}{\pgfqpoint{3.460000in}{2.052407in}}%
\pgfpathlineto{\pgfqpoint{3.460000in}{2.052407in}}%
\pgfpathclose%
\pgfusepath{clip}%
\pgfsetrectcap%
\pgfsetroundjoin%
\pgfsetlinewidth{0.451687pt}%
\definecolor{currentstroke}{rgb}{0.690196,0.690196,0.690196}%
\pgfsetstrokecolor{currentstroke}%
\pgfsetstrokeopacity{0.250000}%
\pgfsetdash{}{0pt}%
\pgfpathmoveto{\pgfqpoint{3.151937in}{2.103702in}}%
\pgfpathlineto{\pgfqpoint{3.126525in}{2.113059in}}%
\pgfpathlineto{\pgfqpoint{3.102160in}{2.122919in}}%
\pgfpathlineto{\pgfqpoint{3.079404in}{2.132965in}}%
\pgfpathlineto{\pgfqpoint{3.057672in}{2.143363in}}%
\pgfpathlineto{\pgfqpoint{3.036907in}{2.154082in}}%
\pgfpathlineto{\pgfqpoint{3.017055in}{2.165102in}}%
\pgfpathlineto{\pgfqpoint{2.998068in}{2.176401in}}%
\pgfpathlineto{\pgfqpoint{2.979905in}{2.187962in}}%
\pgfpathlineto{\pgfqpoint{2.962528in}{2.199772in}}%
\pgfpathlineto{\pgfqpoint{2.945907in}{2.211817in}}%
\pgfpathlineto{\pgfqpoint{2.930014in}{2.224083in}}%
\pgfpathlineto{\pgfqpoint{2.914826in}{2.236561in}}%
\pgfpathlineto{\pgfqpoint{2.900322in}{2.249238in}}%
\pgfpathlineto{\pgfqpoint{2.886485in}{2.262105in}}%
\pgfpathlineto{\pgfqpoint{2.873141in}{2.275314in}}%
\pgfpathlineto{\pgfqpoint{2.860644in}{2.288486in}}%
\pgfpathlineto{\pgfqpoint{2.848759in}{2.301831in}}%
\pgfpathlineto{\pgfqpoint{2.837477in}{2.315338in}}%
\pgfpathlineto{\pgfqpoint{2.826795in}{2.328996in}}%
\pgfpathlineto{\pgfqpoint{2.816705in}{2.342794in}}%
\pgfpathlineto{\pgfqpoint{2.807204in}{2.356724in}}%
\pgfpathlineto{\pgfqpoint{2.798284in}{2.370777in}}%
\pgfpathlineto{\pgfqpoint{2.789943in}{2.384944in}}%
\pgfpathlineto{\pgfqpoint{2.782175in}{2.399216in}}%
\pgfpathlineto{\pgfqpoint{2.774976in}{2.413586in}}%
\pgfpathlineto{\pgfqpoint{2.768344in}{2.428046in}}%
\pgfpathlineto{\pgfqpoint{2.762275in}{2.442588in}}%
\pgfpathlineto{\pgfqpoint{2.756766in}{2.457204in}}%
\pgfpathlineto{\pgfqpoint{2.751814in}{2.471888in}}%
\pgfpathlineto{\pgfqpoint{2.747418in}{2.486632in}}%
\pgfpathlineto{\pgfqpoint{2.743575in}{2.501429in}}%
\pgfpathlineto{\pgfqpoint{2.740251in}{2.516436in}}%
\pgfpathlineto{\pgfqpoint{2.737529in}{2.531243in}}%
\pgfpathlineto{\pgfqpoint{2.735347in}{2.546108in}}%
\pgfpathlineto{\pgfqpoint{2.733709in}{2.561018in}}%
\pgfpathlineto{\pgfqpoint{2.732615in}{2.575961in}}%
\pgfpathlineto{\pgfqpoint{2.732068in}{2.590923in}}%
\pgfpathlineto{\pgfqpoint{2.732068in}{2.605891in}}%
\pgfpathlineto{\pgfqpoint{2.732615in}{2.620853in}}%
\pgfpathlineto{\pgfqpoint{2.733709in}{2.635795in}}%
\pgfpathlineto{\pgfqpoint{2.735347in}{2.650705in}}%
\pgfpathlineto{\pgfqpoint{2.737529in}{2.665570in}}%
\pgfpathlineto{\pgfqpoint{2.740251in}{2.680378in}}%
\pgfpathlineto{\pgfqpoint{2.743575in}{2.695384in}}%
\pgfpathlineto{\pgfqpoint{2.747418in}{2.710181in}}%
\pgfpathlineto{\pgfqpoint{2.751814in}{2.724925in}}%
\pgfpathlineto{\pgfqpoint{2.756766in}{2.739609in}}%
\pgfpathlineto{\pgfqpoint{2.762275in}{2.754226in}}%
\pgfpathlineto{\pgfqpoint{2.768344in}{2.768768in}}%
\pgfpathlineto{\pgfqpoint{2.774976in}{2.783228in}}%
\pgfpathlineto{\pgfqpoint{2.782175in}{2.797598in}}%
\pgfpathlineto{\pgfqpoint{2.789943in}{2.811870in}}%
\pgfpathlineto{\pgfqpoint{2.798284in}{2.826037in}}%
\pgfpathlineto{\pgfqpoint{2.807204in}{2.840089in}}%
\pgfpathlineto{\pgfqpoint{2.816705in}{2.854019in}}%
\pgfpathlineto{\pgfqpoint{2.826795in}{2.867818in}}%
\pgfpathlineto{\pgfqpoint{2.837477in}{2.881476in}}%
\pgfpathlineto{\pgfqpoint{2.848759in}{2.894982in}}%
\pgfpathlineto{\pgfqpoint{2.860644in}{2.908327in}}%
\pgfpathlineto{\pgfqpoint{2.873141in}{2.921500in}}%
\pgfpathlineto{\pgfqpoint{2.886485in}{2.934708in}}%
\pgfpathlineto{\pgfqpoint{2.900322in}{2.947576in}}%
\pgfpathlineto{\pgfqpoint{2.914826in}{2.960253in}}%
\pgfpathlineto{\pgfqpoint{2.930014in}{2.972730in}}%
\pgfpathlineto{\pgfqpoint{2.945907in}{2.984997in}}%
\pgfpathlineto{\pgfqpoint{2.962528in}{2.997041in}}%
\pgfpathlineto{\pgfqpoint{2.979905in}{3.008851in}}%
\pgfpathlineto{\pgfqpoint{2.998068in}{3.020413in}}%
\pgfpathlineto{\pgfqpoint{3.017055in}{3.031712in}}%
\pgfpathlineto{\pgfqpoint{3.036907in}{3.042731in}}%
\pgfpathlineto{\pgfqpoint{3.057672in}{3.053451in}}%
\pgfpathlineto{\pgfqpoint{3.079404in}{3.063848in}}%
\pgfpathlineto{\pgfqpoint{3.102160in}{3.073894in}}%
\pgfpathlineto{\pgfqpoint{3.126525in}{3.083755in}}%
\pgfpathlineto{\pgfqpoint{3.151937in}{3.093112in}}%
\pgfusepath{stroke}%
\end{pgfscope}%
\begin{pgfscope}%
\pgfpathmoveto{\pgfqpoint{3.460000in}{2.052407in}}%
\pgfpathcurveto{\pgfqpoint{3.749602in}{2.052407in}}{\pgfqpoint{4.027381in}{2.109937in}}{\pgfqpoint{4.232161in}{2.212326in}}%
\pgfpathcurveto{\pgfqpoint{4.436940in}{2.314716in}}{\pgfqpoint{4.552000in}{2.453606in}}{\pgfqpoint{4.552000in}{2.598407in}}%
\pgfpathcurveto{\pgfqpoint{4.552000in}{2.743208in}}{\pgfqpoint{4.436940in}{2.882097in}}{\pgfqpoint{4.232161in}{2.984487in}}%
\pgfpathcurveto{\pgfqpoint{4.027381in}{3.086877in}}{\pgfqpoint{3.749602in}{3.144407in}}{\pgfqpoint{3.460000in}{3.144407in}}%
\pgfpathcurveto{\pgfqpoint{3.170398in}{3.144407in}}{\pgfqpoint{2.892619in}{3.086877in}}{\pgfqpoint{2.687839in}{2.984487in}}%
\pgfpathcurveto{\pgfqpoint{2.483060in}{2.882097in}}{\pgfqpoint{2.368000in}{2.743208in}}{\pgfqpoint{2.368000in}{2.598407in}}%
\pgfpathcurveto{\pgfqpoint{2.368000in}{2.453606in}}{\pgfqpoint{2.483060in}{2.314716in}}{\pgfqpoint{2.687839in}{2.212326in}}%
\pgfpathcurveto{\pgfqpoint{2.892619in}{2.109937in}}{\pgfqpoint{3.170398in}{2.052407in}}{\pgfqpoint{3.460000in}{2.052407in}}%
\pgfpathlineto{\pgfqpoint{3.460000in}{2.052407in}}%
\pgfpathclose%
\pgfusepath{clip}%
\pgfsetrectcap%
\pgfsetroundjoin%
\pgfsetlinewidth{0.451687pt}%
\definecolor{currentstroke}{rgb}{0.690196,0.690196,0.690196}%
\pgfsetstrokecolor{currentstroke}%
\pgfsetstrokeopacity{0.250000}%
\pgfsetdash{}{0pt}%
\pgfpathmoveto{\pgfqpoint{3.228953in}{2.103702in}}%
\pgfpathlineto{\pgfqpoint{3.209894in}{2.113059in}}%
\pgfpathlineto{\pgfqpoint{3.191620in}{2.122919in}}%
\pgfpathlineto{\pgfqpoint{3.174553in}{2.132965in}}%
\pgfpathlineto{\pgfqpoint{3.158254in}{2.143363in}}%
\pgfpathlineto{\pgfqpoint{3.142680in}{2.154082in}}%
\pgfpathlineto{\pgfqpoint{3.127791in}{2.165102in}}%
\pgfpathlineto{\pgfqpoint{3.113551in}{2.176401in}}%
\pgfpathlineto{\pgfqpoint{3.099929in}{2.187962in}}%
\pgfpathlineto{\pgfqpoint{3.086896in}{2.199772in}}%
\pgfpathlineto{\pgfqpoint{3.074430in}{2.211817in}}%
\pgfpathlineto{\pgfqpoint{3.062510in}{2.224083in}}%
\pgfpathlineto{\pgfqpoint{3.051119in}{2.236561in}}%
\pgfpathlineto{\pgfqpoint{3.040242in}{2.249238in}}%
\pgfpathlineto{\pgfqpoint{3.029864in}{2.262105in}}%
\pgfpathlineto{\pgfqpoint{3.019856in}{2.275314in}}%
\pgfpathlineto{\pgfqpoint{3.010483in}{2.288486in}}%
\pgfpathlineto{\pgfqpoint{3.001569in}{2.301831in}}%
\pgfpathlineto{\pgfqpoint{2.993108in}{2.315338in}}%
\pgfpathlineto{\pgfqpoint{2.985096in}{2.328996in}}%
\pgfpathlineto{\pgfqpoint{2.977529in}{2.342794in}}%
\pgfpathlineto{\pgfqpoint{2.970403in}{2.356724in}}%
\pgfpathlineto{\pgfqpoint{2.963713in}{2.370777in}}%
\pgfpathlineto{\pgfqpoint{2.957457in}{2.384944in}}%
\pgfpathlineto{\pgfqpoint{2.951631in}{2.399216in}}%
\pgfpathlineto{\pgfqpoint{2.946232in}{2.413586in}}%
\pgfpathlineto{\pgfqpoint{2.941258in}{2.428046in}}%
\pgfpathlineto{\pgfqpoint{2.936706in}{2.442588in}}%
\pgfpathlineto{\pgfqpoint{2.932574in}{2.457204in}}%
\pgfpathlineto{\pgfqpoint{2.928861in}{2.471888in}}%
\pgfpathlineto{\pgfqpoint{2.925563in}{2.486632in}}%
\pgfpathlineto{\pgfqpoint{2.922681in}{2.501429in}}%
\pgfpathlineto{\pgfqpoint{2.920188in}{2.516436in}}%
\pgfpathlineto{\pgfqpoint{2.918147in}{2.531243in}}%
\pgfpathlineto{\pgfqpoint{2.916510in}{2.546108in}}%
\pgfpathlineto{\pgfqpoint{2.915282in}{2.561018in}}%
\pgfpathlineto{\pgfqpoint{2.914462in}{2.575961in}}%
\pgfpathlineto{\pgfqpoint{2.914051in}{2.590923in}}%
\pgfpathlineto{\pgfqpoint{2.914051in}{2.605891in}}%
\pgfpathlineto{\pgfqpoint{2.914462in}{2.620853in}}%
\pgfpathlineto{\pgfqpoint{2.915282in}{2.635795in}}%
\pgfpathlineto{\pgfqpoint{2.916510in}{2.650705in}}%
\pgfpathlineto{\pgfqpoint{2.918147in}{2.665570in}}%
\pgfpathlineto{\pgfqpoint{2.920188in}{2.680378in}}%
\pgfpathlineto{\pgfqpoint{2.922681in}{2.695384in}}%
\pgfpathlineto{\pgfqpoint{2.925563in}{2.710181in}}%
\pgfpathlineto{\pgfqpoint{2.928861in}{2.724925in}}%
\pgfpathlineto{\pgfqpoint{2.932574in}{2.739609in}}%
\pgfpathlineto{\pgfqpoint{2.936706in}{2.754226in}}%
\pgfpathlineto{\pgfqpoint{2.941258in}{2.768768in}}%
\pgfpathlineto{\pgfqpoint{2.946232in}{2.783228in}}%
\pgfpathlineto{\pgfqpoint{2.951631in}{2.797598in}}%
\pgfpathlineto{\pgfqpoint{2.957457in}{2.811870in}}%
\pgfpathlineto{\pgfqpoint{2.963713in}{2.826037in}}%
\pgfpathlineto{\pgfqpoint{2.970403in}{2.840089in}}%
\pgfpathlineto{\pgfqpoint{2.977529in}{2.854019in}}%
\pgfpathlineto{\pgfqpoint{2.985096in}{2.867818in}}%
\pgfpathlineto{\pgfqpoint{2.993108in}{2.881476in}}%
\pgfpathlineto{\pgfqpoint{3.001569in}{2.894982in}}%
\pgfpathlineto{\pgfqpoint{3.010483in}{2.908327in}}%
\pgfpathlineto{\pgfqpoint{3.019856in}{2.921500in}}%
\pgfpathlineto{\pgfqpoint{3.029864in}{2.934708in}}%
\pgfpathlineto{\pgfqpoint{3.040242in}{2.947576in}}%
\pgfpathlineto{\pgfqpoint{3.051119in}{2.960253in}}%
\pgfpathlineto{\pgfqpoint{3.062510in}{2.972730in}}%
\pgfpathlineto{\pgfqpoint{3.074430in}{2.984997in}}%
\pgfpathlineto{\pgfqpoint{3.086896in}{2.997041in}}%
\pgfpathlineto{\pgfqpoint{3.099929in}{3.008851in}}%
\pgfpathlineto{\pgfqpoint{3.113551in}{3.020413in}}%
\pgfpathlineto{\pgfqpoint{3.127791in}{3.031712in}}%
\pgfpathlineto{\pgfqpoint{3.142680in}{3.042731in}}%
\pgfpathlineto{\pgfqpoint{3.158254in}{3.053451in}}%
\pgfpathlineto{\pgfqpoint{3.174553in}{3.063848in}}%
\pgfpathlineto{\pgfqpoint{3.191620in}{3.073894in}}%
\pgfpathlineto{\pgfqpoint{3.209894in}{3.083755in}}%
\pgfpathlineto{\pgfqpoint{3.228953in}{3.093112in}}%
\pgfusepath{stroke}%
\end{pgfscope}%
\begin{pgfscope}%
\pgfpathmoveto{\pgfqpoint{3.460000in}{2.052407in}}%
\pgfpathcurveto{\pgfqpoint{3.749602in}{2.052407in}}{\pgfqpoint{4.027381in}{2.109937in}}{\pgfqpoint{4.232161in}{2.212326in}}%
\pgfpathcurveto{\pgfqpoint{4.436940in}{2.314716in}}{\pgfqpoint{4.552000in}{2.453606in}}{\pgfqpoint{4.552000in}{2.598407in}}%
\pgfpathcurveto{\pgfqpoint{4.552000in}{2.743208in}}{\pgfqpoint{4.436940in}{2.882097in}}{\pgfqpoint{4.232161in}{2.984487in}}%
\pgfpathcurveto{\pgfqpoint{4.027381in}{3.086877in}}{\pgfqpoint{3.749602in}{3.144407in}}{\pgfqpoint{3.460000in}{3.144407in}}%
\pgfpathcurveto{\pgfqpoint{3.170398in}{3.144407in}}{\pgfqpoint{2.892619in}{3.086877in}}{\pgfqpoint{2.687839in}{2.984487in}}%
\pgfpathcurveto{\pgfqpoint{2.483060in}{2.882097in}}{\pgfqpoint{2.368000in}{2.743208in}}{\pgfqpoint{2.368000in}{2.598407in}}%
\pgfpathcurveto{\pgfqpoint{2.368000in}{2.453606in}}{\pgfqpoint{2.483060in}{2.314716in}}{\pgfqpoint{2.687839in}{2.212326in}}%
\pgfpathcurveto{\pgfqpoint{2.892619in}{2.109937in}}{\pgfqpoint{3.170398in}{2.052407in}}{\pgfqpoint{3.460000in}{2.052407in}}%
\pgfpathlineto{\pgfqpoint{3.460000in}{2.052407in}}%
\pgfpathclose%
\pgfusepath{clip}%
\pgfsetrectcap%
\pgfsetroundjoin%
\pgfsetlinewidth{0.451687pt}%
\definecolor{currentstroke}{rgb}{0.690196,0.690196,0.690196}%
\pgfsetstrokecolor{currentstroke}%
\pgfsetstrokeopacity{0.250000}%
\pgfsetdash{}{0pt}%
\pgfpathmoveto{\pgfqpoint{3.305969in}{2.103702in}}%
\pgfpathlineto{\pgfqpoint{3.293263in}{2.113059in}}%
\pgfpathlineto{\pgfqpoint{3.281080in}{2.122919in}}%
\pgfpathlineto{\pgfqpoint{3.269702in}{2.132965in}}%
\pgfpathlineto{\pgfqpoint{3.258836in}{2.143363in}}%
\pgfpathlineto{\pgfqpoint{3.248453in}{2.154082in}}%
\pgfpathlineto{\pgfqpoint{3.238527in}{2.165102in}}%
\pgfpathlineto{\pgfqpoint{3.229034in}{2.176401in}}%
\pgfpathlineto{\pgfqpoint{3.219952in}{2.187962in}}%
\pgfpathlineto{\pgfqpoint{3.211264in}{2.199772in}}%
\pgfpathlineto{\pgfqpoint{3.202953in}{2.211817in}}%
\pgfpathlineto{\pgfqpoint{3.195007in}{2.224083in}}%
\pgfpathlineto{\pgfqpoint{3.187413in}{2.236561in}}%
\pgfpathlineto{\pgfqpoint{3.180161in}{2.249238in}}%
\pgfpathlineto{\pgfqpoint{3.173243in}{2.262105in}}%
\pgfpathlineto{\pgfqpoint{3.166570in}{2.275314in}}%
\pgfpathlineto{\pgfqpoint{3.160322in}{2.288486in}}%
\pgfpathlineto{\pgfqpoint{3.154379in}{2.301831in}}%
\pgfpathlineto{\pgfqpoint{3.148739in}{2.315338in}}%
\pgfpathlineto{\pgfqpoint{3.143398in}{2.328996in}}%
\pgfpathlineto{\pgfqpoint{3.138353in}{2.342794in}}%
\pgfpathlineto{\pgfqpoint{3.133602in}{2.356724in}}%
\pgfpathlineto{\pgfqpoint{3.129142in}{2.370777in}}%
\pgfpathlineto{\pgfqpoint{3.124971in}{2.384944in}}%
\pgfpathlineto{\pgfqpoint{3.121087in}{2.399216in}}%
\pgfpathlineto{\pgfqpoint{3.117488in}{2.413586in}}%
\pgfpathlineto{\pgfqpoint{3.114172in}{2.428046in}}%
\pgfpathlineto{\pgfqpoint{3.111137in}{2.442588in}}%
\pgfpathlineto{\pgfqpoint{3.108383in}{2.457204in}}%
\pgfpathlineto{\pgfqpoint{3.105907in}{2.471888in}}%
\pgfpathlineto{\pgfqpoint{3.103709in}{2.486632in}}%
\pgfpathlineto{\pgfqpoint{3.101788in}{2.501429in}}%
\pgfpathlineto{\pgfqpoint{3.100125in}{2.516436in}}%
\pgfpathlineto{\pgfqpoint{3.098764in}{2.531243in}}%
\pgfpathlineto{\pgfqpoint{3.097674in}{2.546108in}}%
\pgfpathlineto{\pgfqpoint{3.096854in}{2.561018in}}%
\pgfpathlineto{\pgfqpoint{3.096308in}{2.575961in}}%
\pgfpathlineto{\pgfqpoint{3.096034in}{2.590923in}}%
\pgfpathlineto{\pgfqpoint{3.096034in}{2.605891in}}%
\pgfpathlineto{\pgfqpoint{3.096308in}{2.620853in}}%
\pgfpathlineto{\pgfqpoint{3.096854in}{2.635795in}}%
\pgfpathlineto{\pgfqpoint{3.097674in}{2.650705in}}%
\pgfpathlineto{\pgfqpoint{3.098764in}{2.665570in}}%
\pgfpathlineto{\pgfqpoint{3.100125in}{2.680378in}}%
\pgfpathlineto{\pgfqpoint{3.101788in}{2.695384in}}%
\pgfpathlineto{\pgfqpoint{3.103709in}{2.710181in}}%
\pgfpathlineto{\pgfqpoint{3.105907in}{2.724925in}}%
\pgfpathlineto{\pgfqpoint{3.108383in}{2.739609in}}%
\pgfpathlineto{\pgfqpoint{3.111137in}{2.754226in}}%
\pgfpathlineto{\pgfqpoint{3.114172in}{2.768768in}}%
\pgfpathlineto{\pgfqpoint{3.117488in}{2.783228in}}%
\pgfpathlineto{\pgfqpoint{3.121087in}{2.797598in}}%
\pgfpathlineto{\pgfqpoint{3.124971in}{2.811870in}}%
\pgfpathlineto{\pgfqpoint{3.129142in}{2.826037in}}%
\pgfpathlineto{\pgfqpoint{3.133602in}{2.840089in}}%
\pgfpathlineto{\pgfqpoint{3.138353in}{2.854019in}}%
\pgfpathlineto{\pgfqpoint{3.143398in}{2.867818in}}%
\pgfpathlineto{\pgfqpoint{3.148739in}{2.881476in}}%
\pgfpathlineto{\pgfqpoint{3.154379in}{2.894982in}}%
\pgfpathlineto{\pgfqpoint{3.160322in}{2.908327in}}%
\pgfpathlineto{\pgfqpoint{3.166570in}{2.921500in}}%
\pgfpathlineto{\pgfqpoint{3.173243in}{2.934708in}}%
\pgfpathlineto{\pgfqpoint{3.180161in}{2.947576in}}%
\pgfpathlineto{\pgfqpoint{3.187413in}{2.960253in}}%
\pgfpathlineto{\pgfqpoint{3.195007in}{2.972730in}}%
\pgfpathlineto{\pgfqpoint{3.202953in}{2.984997in}}%
\pgfpathlineto{\pgfqpoint{3.211264in}{2.997041in}}%
\pgfpathlineto{\pgfqpoint{3.219952in}{3.008851in}}%
\pgfpathlineto{\pgfqpoint{3.229034in}{3.020413in}}%
\pgfpathlineto{\pgfqpoint{3.238527in}{3.031712in}}%
\pgfpathlineto{\pgfqpoint{3.248453in}{3.042731in}}%
\pgfpathlineto{\pgfqpoint{3.258836in}{3.053451in}}%
\pgfpathlineto{\pgfqpoint{3.269702in}{3.063848in}}%
\pgfpathlineto{\pgfqpoint{3.281080in}{3.073894in}}%
\pgfpathlineto{\pgfqpoint{3.293263in}{3.083755in}}%
\pgfpathlineto{\pgfqpoint{3.305969in}{3.093112in}}%
\pgfusepath{stroke}%
\end{pgfscope}%
\begin{pgfscope}%
\pgfpathmoveto{\pgfqpoint{3.460000in}{2.052407in}}%
\pgfpathcurveto{\pgfqpoint{3.749602in}{2.052407in}}{\pgfqpoint{4.027381in}{2.109937in}}{\pgfqpoint{4.232161in}{2.212326in}}%
\pgfpathcurveto{\pgfqpoint{4.436940in}{2.314716in}}{\pgfqpoint{4.552000in}{2.453606in}}{\pgfqpoint{4.552000in}{2.598407in}}%
\pgfpathcurveto{\pgfqpoint{4.552000in}{2.743208in}}{\pgfqpoint{4.436940in}{2.882097in}}{\pgfqpoint{4.232161in}{2.984487in}}%
\pgfpathcurveto{\pgfqpoint{4.027381in}{3.086877in}}{\pgfqpoint{3.749602in}{3.144407in}}{\pgfqpoint{3.460000in}{3.144407in}}%
\pgfpathcurveto{\pgfqpoint{3.170398in}{3.144407in}}{\pgfqpoint{2.892619in}{3.086877in}}{\pgfqpoint{2.687839in}{2.984487in}}%
\pgfpathcurveto{\pgfqpoint{2.483060in}{2.882097in}}{\pgfqpoint{2.368000in}{2.743208in}}{\pgfqpoint{2.368000in}{2.598407in}}%
\pgfpathcurveto{\pgfqpoint{2.368000in}{2.453606in}}{\pgfqpoint{2.483060in}{2.314716in}}{\pgfqpoint{2.687839in}{2.212326in}}%
\pgfpathcurveto{\pgfqpoint{2.892619in}{2.109937in}}{\pgfqpoint{3.170398in}{2.052407in}}{\pgfqpoint{3.460000in}{2.052407in}}%
\pgfpathlineto{\pgfqpoint{3.460000in}{2.052407in}}%
\pgfpathclose%
\pgfusepath{clip}%
\pgfsetrectcap%
\pgfsetroundjoin%
\pgfsetlinewidth{0.451687pt}%
\definecolor{currentstroke}{rgb}{0.690196,0.690196,0.690196}%
\pgfsetstrokecolor{currentstroke}%
\pgfsetstrokeopacity{0.250000}%
\pgfsetdash{}{0pt}%
\pgfpathmoveto{\pgfqpoint{3.382984in}{2.103702in}}%
\pgfpathlineto{\pgfqpoint{3.376631in}{2.113059in}}%
\pgfpathlineto{\pgfqpoint{3.370540in}{2.122919in}}%
\pgfpathlineto{\pgfqpoint{3.364851in}{2.132965in}}%
\pgfpathlineto{\pgfqpoint{3.359418in}{2.143363in}}%
\pgfpathlineto{\pgfqpoint{3.354227in}{2.154082in}}%
\pgfpathlineto{\pgfqpoint{3.349264in}{2.165102in}}%
\pgfpathlineto{\pgfqpoint{3.344517in}{2.176401in}}%
\pgfpathlineto{\pgfqpoint{3.339976in}{2.187962in}}%
\pgfpathlineto{\pgfqpoint{3.335632in}{2.199772in}}%
\pgfpathlineto{\pgfqpoint{3.331477in}{2.211817in}}%
\pgfpathlineto{\pgfqpoint{3.327503in}{2.224083in}}%
\pgfpathlineto{\pgfqpoint{3.323706in}{2.236561in}}%
\pgfpathlineto{\pgfqpoint{3.320081in}{2.249238in}}%
\pgfpathlineto{\pgfqpoint{3.316621in}{2.262105in}}%
\pgfpathlineto{\pgfqpoint{3.313285in}{2.275314in}}%
\pgfpathlineto{\pgfqpoint{3.310161in}{2.288486in}}%
\pgfpathlineto{\pgfqpoint{3.307190in}{2.301831in}}%
\pgfpathlineto{\pgfqpoint{3.304369in}{2.315338in}}%
\pgfpathlineto{\pgfqpoint{3.301699in}{2.328996in}}%
\pgfpathlineto{\pgfqpoint{3.299176in}{2.342794in}}%
\pgfpathlineto{\pgfqpoint{3.296801in}{2.356724in}}%
\pgfpathlineto{\pgfqpoint{3.294571in}{2.370777in}}%
\pgfpathlineto{\pgfqpoint{3.292486in}{2.384944in}}%
\pgfpathlineto{\pgfqpoint{3.290544in}{2.399216in}}%
\pgfpathlineto{\pgfqpoint{3.288744in}{2.413586in}}%
\pgfpathlineto{\pgfqpoint{3.287086in}{2.428046in}}%
\pgfpathlineto{\pgfqpoint{3.285569in}{2.442588in}}%
\pgfpathlineto{\pgfqpoint{3.284191in}{2.457204in}}%
\pgfpathlineto{\pgfqpoint{3.282954in}{2.471888in}}%
\pgfpathlineto{\pgfqpoint{3.281854in}{2.486632in}}%
\pgfpathlineto{\pgfqpoint{3.280894in}{2.501429in}}%
\pgfpathlineto{\pgfqpoint{3.280063in}{2.516436in}}%
\pgfpathlineto{\pgfqpoint{3.279382in}{2.531243in}}%
\pgfpathlineto{\pgfqpoint{3.278837in}{2.546108in}}%
\pgfpathlineto{\pgfqpoint{3.278427in}{2.561018in}}%
\pgfpathlineto{\pgfqpoint{3.278154in}{2.575961in}}%
\pgfpathlineto{\pgfqpoint{3.278017in}{2.590923in}}%
\pgfpathlineto{\pgfqpoint{3.278017in}{2.605891in}}%
\pgfpathlineto{\pgfqpoint{3.278154in}{2.620853in}}%
\pgfpathlineto{\pgfqpoint{3.278427in}{2.635795in}}%
\pgfpathlineto{\pgfqpoint{3.278837in}{2.650705in}}%
\pgfpathlineto{\pgfqpoint{3.279382in}{2.665570in}}%
\pgfpathlineto{\pgfqpoint{3.280063in}{2.680378in}}%
\pgfpathlineto{\pgfqpoint{3.280894in}{2.695384in}}%
\pgfpathlineto{\pgfqpoint{3.281854in}{2.710181in}}%
\pgfpathlineto{\pgfqpoint{3.282954in}{2.724925in}}%
\pgfpathlineto{\pgfqpoint{3.284191in}{2.739609in}}%
\pgfpathlineto{\pgfqpoint{3.285569in}{2.754226in}}%
\pgfpathlineto{\pgfqpoint{3.287086in}{2.768768in}}%
\pgfpathlineto{\pgfqpoint{3.288744in}{2.783228in}}%
\pgfpathlineto{\pgfqpoint{3.290544in}{2.797598in}}%
\pgfpathlineto{\pgfqpoint{3.292486in}{2.811870in}}%
\pgfpathlineto{\pgfqpoint{3.294571in}{2.826037in}}%
\pgfpathlineto{\pgfqpoint{3.296801in}{2.840089in}}%
\pgfpathlineto{\pgfqpoint{3.299176in}{2.854019in}}%
\pgfpathlineto{\pgfqpoint{3.301699in}{2.867818in}}%
\pgfpathlineto{\pgfqpoint{3.304369in}{2.881476in}}%
\pgfpathlineto{\pgfqpoint{3.307190in}{2.894982in}}%
\pgfpathlineto{\pgfqpoint{3.310161in}{2.908327in}}%
\pgfpathlineto{\pgfqpoint{3.313285in}{2.921500in}}%
\pgfpathlineto{\pgfqpoint{3.316621in}{2.934708in}}%
\pgfpathlineto{\pgfqpoint{3.320081in}{2.947576in}}%
\pgfpathlineto{\pgfqpoint{3.323706in}{2.960253in}}%
\pgfpathlineto{\pgfqpoint{3.327503in}{2.972730in}}%
\pgfpathlineto{\pgfqpoint{3.331477in}{2.984997in}}%
\pgfpathlineto{\pgfqpoint{3.335632in}{2.997041in}}%
\pgfpathlineto{\pgfqpoint{3.339976in}{3.008851in}}%
\pgfpathlineto{\pgfqpoint{3.344517in}{3.020413in}}%
\pgfpathlineto{\pgfqpoint{3.349264in}{3.031712in}}%
\pgfpathlineto{\pgfqpoint{3.354227in}{3.042731in}}%
\pgfpathlineto{\pgfqpoint{3.359418in}{3.053451in}}%
\pgfpathlineto{\pgfqpoint{3.364851in}{3.063848in}}%
\pgfpathlineto{\pgfqpoint{3.370540in}{3.073894in}}%
\pgfpathlineto{\pgfqpoint{3.376631in}{3.083755in}}%
\pgfpathlineto{\pgfqpoint{3.382984in}{3.093112in}}%
\pgfusepath{stroke}%
\end{pgfscope}%
\begin{pgfscope}%
\pgfpathmoveto{\pgfqpoint{3.460000in}{2.052407in}}%
\pgfpathcurveto{\pgfqpoint{3.749602in}{2.052407in}}{\pgfqpoint{4.027381in}{2.109937in}}{\pgfqpoint{4.232161in}{2.212326in}}%
\pgfpathcurveto{\pgfqpoint{4.436940in}{2.314716in}}{\pgfqpoint{4.552000in}{2.453606in}}{\pgfqpoint{4.552000in}{2.598407in}}%
\pgfpathcurveto{\pgfqpoint{4.552000in}{2.743208in}}{\pgfqpoint{4.436940in}{2.882097in}}{\pgfqpoint{4.232161in}{2.984487in}}%
\pgfpathcurveto{\pgfqpoint{4.027381in}{3.086877in}}{\pgfqpoint{3.749602in}{3.144407in}}{\pgfqpoint{3.460000in}{3.144407in}}%
\pgfpathcurveto{\pgfqpoint{3.170398in}{3.144407in}}{\pgfqpoint{2.892619in}{3.086877in}}{\pgfqpoint{2.687839in}{2.984487in}}%
\pgfpathcurveto{\pgfqpoint{2.483060in}{2.882097in}}{\pgfqpoint{2.368000in}{2.743208in}}{\pgfqpoint{2.368000in}{2.598407in}}%
\pgfpathcurveto{\pgfqpoint{2.368000in}{2.453606in}}{\pgfqpoint{2.483060in}{2.314716in}}{\pgfqpoint{2.687839in}{2.212326in}}%
\pgfpathcurveto{\pgfqpoint{2.892619in}{2.109937in}}{\pgfqpoint{3.170398in}{2.052407in}}{\pgfqpoint{3.460000in}{2.052407in}}%
\pgfpathlineto{\pgfqpoint{3.460000in}{2.052407in}}%
\pgfpathclose%
\pgfusepath{clip}%
\pgfsetrectcap%
\pgfsetroundjoin%
\pgfsetlinewidth{0.451687pt}%
\definecolor{currentstroke}{rgb}{0.690196,0.690196,0.690196}%
\pgfsetstrokecolor{currentstroke}%
\pgfsetstrokeopacity{0.250000}%
\pgfsetdash{}{0pt}%
\pgfpathmoveto{\pgfqpoint{3.460000in}{2.103702in}}%
\pgfpathlineto{\pgfqpoint{3.460000in}{2.113059in}}%
\pgfpathlineto{\pgfqpoint{3.460000in}{2.122919in}}%
\pgfpathlineto{\pgfqpoint{3.460000in}{2.132965in}}%
\pgfpathlineto{\pgfqpoint{3.460000in}{2.143363in}}%
\pgfpathlineto{\pgfqpoint{3.460000in}{2.154082in}}%
\pgfpathlineto{\pgfqpoint{3.460000in}{2.165102in}}%
\pgfpathlineto{\pgfqpoint{3.460000in}{2.176401in}}%
\pgfpathlineto{\pgfqpoint{3.460000in}{2.187962in}}%
\pgfpathlineto{\pgfqpoint{3.460000in}{2.199772in}}%
\pgfpathlineto{\pgfqpoint{3.460000in}{2.211817in}}%
\pgfpathlineto{\pgfqpoint{3.460000in}{2.224083in}}%
\pgfpathlineto{\pgfqpoint{3.460000in}{2.236561in}}%
\pgfpathlineto{\pgfqpoint{3.460000in}{2.249238in}}%
\pgfpathlineto{\pgfqpoint{3.460000in}{2.262105in}}%
\pgfpathlineto{\pgfqpoint{3.460000in}{2.275314in}}%
\pgfpathlineto{\pgfqpoint{3.460000in}{2.288486in}}%
\pgfpathlineto{\pgfqpoint{3.460000in}{2.301831in}}%
\pgfpathlineto{\pgfqpoint{3.460000in}{2.315338in}}%
\pgfpathlineto{\pgfqpoint{3.460000in}{2.328996in}}%
\pgfpathlineto{\pgfqpoint{3.460000in}{2.342794in}}%
\pgfpathlineto{\pgfqpoint{3.460000in}{2.356724in}}%
\pgfpathlineto{\pgfqpoint{3.460000in}{2.370777in}}%
\pgfpathlineto{\pgfqpoint{3.460000in}{2.384944in}}%
\pgfpathlineto{\pgfqpoint{3.460000in}{2.399216in}}%
\pgfpathlineto{\pgfqpoint{3.460000in}{2.413586in}}%
\pgfpathlineto{\pgfqpoint{3.460000in}{2.428046in}}%
\pgfpathlineto{\pgfqpoint{3.460000in}{2.442588in}}%
\pgfpathlineto{\pgfqpoint{3.460000in}{2.457204in}}%
\pgfpathlineto{\pgfqpoint{3.460000in}{2.471888in}}%
\pgfpathlineto{\pgfqpoint{3.460000in}{2.486632in}}%
\pgfpathlineto{\pgfqpoint{3.460000in}{2.501429in}}%
\pgfpathlineto{\pgfqpoint{3.460000in}{2.516436in}}%
\pgfpathlineto{\pgfqpoint{3.460000in}{2.531243in}}%
\pgfpathlineto{\pgfqpoint{3.460000in}{2.546108in}}%
\pgfpathlineto{\pgfqpoint{3.460000in}{2.561018in}}%
\pgfpathlineto{\pgfqpoint{3.460000in}{2.575961in}}%
\pgfpathlineto{\pgfqpoint{3.460000in}{2.590923in}}%
\pgfpathlineto{\pgfqpoint{3.460000in}{2.605891in}}%
\pgfpathlineto{\pgfqpoint{3.460000in}{2.620853in}}%
\pgfpathlineto{\pgfqpoint{3.460000in}{2.635795in}}%
\pgfpathlineto{\pgfqpoint{3.460000in}{2.650705in}}%
\pgfpathlineto{\pgfqpoint{3.460000in}{2.665570in}}%
\pgfpathlineto{\pgfqpoint{3.460000in}{2.680378in}}%
\pgfpathlineto{\pgfqpoint{3.460000in}{2.695384in}}%
\pgfpathlineto{\pgfqpoint{3.460000in}{2.710181in}}%
\pgfpathlineto{\pgfqpoint{3.460000in}{2.724925in}}%
\pgfpathlineto{\pgfqpoint{3.460000in}{2.739609in}}%
\pgfpathlineto{\pgfqpoint{3.460000in}{2.754226in}}%
\pgfpathlineto{\pgfqpoint{3.460000in}{2.768768in}}%
\pgfpathlineto{\pgfqpoint{3.460000in}{2.783228in}}%
\pgfpathlineto{\pgfqpoint{3.460000in}{2.797598in}}%
\pgfpathlineto{\pgfqpoint{3.460000in}{2.811870in}}%
\pgfpathlineto{\pgfqpoint{3.460000in}{2.826037in}}%
\pgfpathlineto{\pgfqpoint{3.460000in}{2.840089in}}%
\pgfpathlineto{\pgfqpoint{3.460000in}{2.854019in}}%
\pgfpathlineto{\pgfqpoint{3.460000in}{2.867818in}}%
\pgfpathlineto{\pgfqpoint{3.460000in}{2.881476in}}%
\pgfpathlineto{\pgfqpoint{3.460000in}{2.894982in}}%
\pgfpathlineto{\pgfqpoint{3.460000in}{2.908327in}}%
\pgfpathlineto{\pgfqpoint{3.460000in}{2.921500in}}%
\pgfpathlineto{\pgfqpoint{3.460000in}{2.934708in}}%
\pgfpathlineto{\pgfqpoint{3.460000in}{2.947576in}}%
\pgfpathlineto{\pgfqpoint{3.460000in}{2.960253in}}%
\pgfpathlineto{\pgfqpoint{3.460000in}{2.972730in}}%
\pgfpathlineto{\pgfqpoint{3.460000in}{2.984997in}}%
\pgfpathlineto{\pgfqpoint{3.460000in}{2.997041in}}%
\pgfpathlineto{\pgfqpoint{3.460000in}{3.008851in}}%
\pgfpathlineto{\pgfqpoint{3.460000in}{3.020413in}}%
\pgfpathlineto{\pgfqpoint{3.460000in}{3.031712in}}%
\pgfpathlineto{\pgfqpoint{3.460000in}{3.042731in}}%
\pgfpathlineto{\pgfqpoint{3.460000in}{3.053451in}}%
\pgfpathlineto{\pgfqpoint{3.460000in}{3.063848in}}%
\pgfpathlineto{\pgfqpoint{3.460000in}{3.073894in}}%
\pgfpathlineto{\pgfqpoint{3.460000in}{3.083755in}}%
\pgfpathlineto{\pgfqpoint{3.460000in}{3.093112in}}%
\pgfusepath{stroke}%
\end{pgfscope}%
\begin{pgfscope}%
\pgfpathmoveto{\pgfqpoint{3.460000in}{2.052407in}}%
\pgfpathcurveto{\pgfqpoint{3.749602in}{2.052407in}}{\pgfqpoint{4.027381in}{2.109937in}}{\pgfqpoint{4.232161in}{2.212326in}}%
\pgfpathcurveto{\pgfqpoint{4.436940in}{2.314716in}}{\pgfqpoint{4.552000in}{2.453606in}}{\pgfqpoint{4.552000in}{2.598407in}}%
\pgfpathcurveto{\pgfqpoint{4.552000in}{2.743208in}}{\pgfqpoint{4.436940in}{2.882097in}}{\pgfqpoint{4.232161in}{2.984487in}}%
\pgfpathcurveto{\pgfqpoint{4.027381in}{3.086877in}}{\pgfqpoint{3.749602in}{3.144407in}}{\pgfqpoint{3.460000in}{3.144407in}}%
\pgfpathcurveto{\pgfqpoint{3.170398in}{3.144407in}}{\pgfqpoint{2.892619in}{3.086877in}}{\pgfqpoint{2.687839in}{2.984487in}}%
\pgfpathcurveto{\pgfqpoint{2.483060in}{2.882097in}}{\pgfqpoint{2.368000in}{2.743208in}}{\pgfqpoint{2.368000in}{2.598407in}}%
\pgfpathcurveto{\pgfqpoint{2.368000in}{2.453606in}}{\pgfqpoint{2.483060in}{2.314716in}}{\pgfqpoint{2.687839in}{2.212326in}}%
\pgfpathcurveto{\pgfqpoint{2.892619in}{2.109937in}}{\pgfqpoint{3.170398in}{2.052407in}}{\pgfqpoint{3.460000in}{2.052407in}}%
\pgfpathlineto{\pgfqpoint{3.460000in}{2.052407in}}%
\pgfpathclose%
\pgfusepath{clip}%
\pgfsetrectcap%
\pgfsetroundjoin%
\pgfsetlinewidth{0.451687pt}%
\definecolor{currentstroke}{rgb}{0.690196,0.690196,0.690196}%
\pgfsetstrokecolor{currentstroke}%
\pgfsetstrokeopacity{0.250000}%
\pgfsetdash{}{0pt}%
\pgfpathmoveto{\pgfqpoint{3.537016in}{2.103702in}}%
\pgfpathlineto{\pgfqpoint{3.543369in}{2.113059in}}%
\pgfpathlineto{\pgfqpoint{3.549460in}{2.122919in}}%
\pgfpathlineto{\pgfqpoint{3.555149in}{2.132965in}}%
\pgfpathlineto{\pgfqpoint{3.560582in}{2.143363in}}%
\pgfpathlineto{\pgfqpoint{3.565773in}{2.154082in}}%
\pgfpathlineto{\pgfqpoint{3.570736in}{2.165102in}}%
\pgfpathlineto{\pgfqpoint{3.575483in}{2.176401in}}%
\pgfpathlineto{\pgfqpoint{3.580024in}{2.187962in}}%
\pgfpathlineto{\pgfqpoint{3.584368in}{2.199772in}}%
\pgfpathlineto{\pgfqpoint{3.588523in}{2.211817in}}%
\pgfpathlineto{\pgfqpoint{3.592497in}{2.224083in}}%
\pgfpathlineto{\pgfqpoint{3.596294in}{2.236561in}}%
\pgfpathlineto{\pgfqpoint{3.599919in}{2.249238in}}%
\pgfpathlineto{\pgfqpoint{3.603379in}{2.262105in}}%
\pgfpathlineto{\pgfqpoint{3.606715in}{2.275314in}}%
\pgfpathlineto{\pgfqpoint{3.609839in}{2.288486in}}%
\pgfpathlineto{\pgfqpoint{3.612810in}{2.301831in}}%
\pgfpathlineto{\pgfqpoint{3.615631in}{2.315338in}}%
\pgfpathlineto{\pgfqpoint{3.618301in}{2.328996in}}%
\pgfpathlineto{\pgfqpoint{3.620824in}{2.342794in}}%
\pgfpathlineto{\pgfqpoint{3.623199in}{2.356724in}}%
\pgfpathlineto{\pgfqpoint{3.625429in}{2.370777in}}%
\pgfpathlineto{\pgfqpoint{3.627514in}{2.384944in}}%
\pgfpathlineto{\pgfqpoint{3.629456in}{2.399216in}}%
\pgfpathlineto{\pgfqpoint{3.631256in}{2.413586in}}%
\pgfpathlineto{\pgfqpoint{3.632914in}{2.428046in}}%
\pgfpathlineto{\pgfqpoint{3.634431in}{2.442588in}}%
\pgfpathlineto{\pgfqpoint{3.635809in}{2.457204in}}%
\pgfpathlineto{\pgfqpoint{3.637046in}{2.471888in}}%
\pgfpathlineto{\pgfqpoint{3.638146in}{2.486632in}}%
\pgfpathlineto{\pgfqpoint{3.639106in}{2.501429in}}%
\pgfpathlineto{\pgfqpoint{3.639937in}{2.516436in}}%
\pgfpathlineto{\pgfqpoint{3.640618in}{2.531243in}}%
\pgfpathlineto{\pgfqpoint{3.641163in}{2.546108in}}%
\pgfpathlineto{\pgfqpoint{3.641573in}{2.561018in}}%
\pgfpathlineto{\pgfqpoint{3.641846in}{2.575961in}}%
\pgfpathlineto{\pgfqpoint{3.641983in}{2.590923in}}%
\pgfpathlineto{\pgfqpoint{3.641983in}{2.605891in}}%
\pgfpathlineto{\pgfqpoint{3.641846in}{2.620853in}}%
\pgfpathlineto{\pgfqpoint{3.641573in}{2.635795in}}%
\pgfpathlineto{\pgfqpoint{3.641163in}{2.650705in}}%
\pgfpathlineto{\pgfqpoint{3.640618in}{2.665570in}}%
\pgfpathlineto{\pgfqpoint{3.639937in}{2.680378in}}%
\pgfpathlineto{\pgfqpoint{3.639106in}{2.695384in}}%
\pgfpathlineto{\pgfqpoint{3.638146in}{2.710181in}}%
\pgfpathlineto{\pgfqpoint{3.637046in}{2.724925in}}%
\pgfpathlineto{\pgfqpoint{3.635809in}{2.739609in}}%
\pgfpathlineto{\pgfqpoint{3.634431in}{2.754226in}}%
\pgfpathlineto{\pgfqpoint{3.632914in}{2.768768in}}%
\pgfpathlineto{\pgfqpoint{3.631256in}{2.783228in}}%
\pgfpathlineto{\pgfqpoint{3.629456in}{2.797598in}}%
\pgfpathlineto{\pgfqpoint{3.627514in}{2.811870in}}%
\pgfpathlineto{\pgfqpoint{3.625429in}{2.826037in}}%
\pgfpathlineto{\pgfqpoint{3.623199in}{2.840089in}}%
\pgfpathlineto{\pgfqpoint{3.620824in}{2.854019in}}%
\pgfpathlineto{\pgfqpoint{3.618301in}{2.867818in}}%
\pgfpathlineto{\pgfqpoint{3.615631in}{2.881476in}}%
\pgfpathlineto{\pgfqpoint{3.612810in}{2.894982in}}%
\pgfpathlineto{\pgfqpoint{3.609839in}{2.908327in}}%
\pgfpathlineto{\pgfqpoint{3.606715in}{2.921500in}}%
\pgfpathlineto{\pgfqpoint{3.603379in}{2.934708in}}%
\pgfpathlineto{\pgfqpoint{3.599919in}{2.947576in}}%
\pgfpathlineto{\pgfqpoint{3.596294in}{2.960253in}}%
\pgfpathlineto{\pgfqpoint{3.592497in}{2.972730in}}%
\pgfpathlineto{\pgfqpoint{3.588523in}{2.984997in}}%
\pgfpathlineto{\pgfqpoint{3.584368in}{2.997041in}}%
\pgfpathlineto{\pgfqpoint{3.580024in}{3.008851in}}%
\pgfpathlineto{\pgfqpoint{3.575483in}{3.020413in}}%
\pgfpathlineto{\pgfqpoint{3.570736in}{3.031712in}}%
\pgfpathlineto{\pgfqpoint{3.565773in}{3.042731in}}%
\pgfpathlineto{\pgfqpoint{3.560582in}{3.053451in}}%
\pgfpathlineto{\pgfqpoint{3.555149in}{3.063848in}}%
\pgfpathlineto{\pgfqpoint{3.549460in}{3.073894in}}%
\pgfpathlineto{\pgfqpoint{3.543369in}{3.083755in}}%
\pgfpathlineto{\pgfqpoint{3.537016in}{3.093112in}}%
\pgfusepath{stroke}%
\end{pgfscope}%
\begin{pgfscope}%
\pgfpathmoveto{\pgfqpoint{3.460000in}{2.052407in}}%
\pgfpathcurveto{\pgfqpoint{3.749602in}{2.052407in}}{\pgfqpoint{4.027381in}{2.109937in}}{\pgfqpoint{4.232161in}{2.212326in}}%
\pgfpathcurveto{\pgfqpoint{4.436940in}{2.314716in}}{\pgfqpoint{4.552000in}{2.453606in}}{\pgfqpoint{4.552000in}{2.598407in}}%
\pgfpathcurveto{\pgfqpoint{4.552000in}{2.743208in}}{\pgfqpoint{4.436940in}{2.882097in}}{\pgfqpoint{4.232161in}{2.984487in}}%
\pgfpathcurveto{\pgfqpoint{4.027381in}{3.086877in}}{\pgfqpoint{3.749602in}{3.144407in}}{\pgfqpoint{3.460000in}{3.144407in}}%
\pgfpathcurveto{\pgfqpoint{3.170398in}{3.144407in}}{\pgfqpoint{2.892619in}{3.086877in}}{\pgfqpoint{2.687839in}{2.984487in}}%
\pgfpathcurveto{\pgfqpoint{2.483060in}{2.882097in}}{\pgfqpoint{2.368000in}{2.743208in}}{\pgfqpoint{2.368000in}{2.598407in}}%
\pgfpathcurveto{\pgfqpoint{2.368000in}{2.453606in}}{\pgfqpoint{2.483060in}{2.314716in}}{\pgfqpoint{2.687839in}{2.212326in}}%
\pgfpathcurveto{\pgfqpoint{2.892619in}{2.109937in}}{\pgfqpoint{3.170398in}{2.052407in}}{\pgfqpoint{3.460000in}{2.052407in}}%
\pgfpathlineto{\pgfqpoint{3.460000in}{2.052407in}}%
\pgfpathclose%
\pgfusepath{clip}%
\pgfsetrectcap%
\pgfsetroundjoin%
\pgfsetlinewidth{0.451687pt}%
\definecolor{currentstroke}{rgb}{0.690196,0.690196,0.690196}%
\pgfsetstrokecolor{currentstroke}%
\pgfsetstrokeopacity{0.250000}%
\pgfsetdash{}{0pt}%
\pgfpathmoveto{\pgfqpoint{3.614031in}{2.103702in}}%
\pgfpathlineto{\pgfqpoint{3.626737in}{2.113059in}}%
\pgfpathlineto{\pgfqpoint{3.638920in}{2.122919in}}%
\pgfpathlineto{\pgfqpoint{3.650298in}{2.132965in}}%
\pgfpathlineto{\pgfqpoint{3.661164in}{2.143363in}}%
\pgfpathlineto{\pgfqpoint{3.671547in}{2.154082in}}%
\pgfpathlineto{\pgfqpoint{3.681473in}{2.165102in}}%
\pgfpathlineto{\pgfqpoint{3.690966in}{2.176401in}}%
\pgfpathlineto{\pgfqpoint{3.700048in}{2.187962in}}%
\pgfpathlineto{\pgfqpoint{3.708736in}{2.199772in}}%
\pgfpathlineto{\pgfqpoint{3.717047in}{2.211817in}}%
\pgfpathlineto{\pgfqpoint{3.724993in}{2.224083in}}%
\pgfpathlineto{\pgfqpoint{3.732587in}{2.236561in}}%
\pgfpathlineto{\pgfqpoint{3.739839in}{2.249238in}}%
\pgfpathlineto{\pgfqpoint{3.746757in}{2.262105in}}%
\pgfpathlineto{\pgfqpoint{3.753430in}{2.275314in}}%
\pgfpathlineto{\pgfqpoint{3.759678in}{2.288486in}}%
\pgfpathlineto{\pgfqpoint{3.765621in}{2.301831in}}%
\pgfpathlineto{\pgfqpoint{3.771261in}{2.315338in}}%
\pgfpathlineto{\pgfqpoint{3.776602in}{2.328996in}}%
\pgfpathlineto{\pgfqpoint{3.781647in}{2.342794in}}%
\pgfpathlineto{\pgfqpoint{3.786398in}{2.356724in}}%
\pgfpathlineto{\pgfqpoint{3.790858in}{2.370777in}}%
\pgfpathlineto{\pgfqpoint{3.795029in}{2.384944in}}%
\pgfpathlineto{\pgfqpoint{3.798913in}{2.399216in}}%
\pgfpathlineto{\pgfqpoint{3.802512in}{2.413586in}}%
\pgfpathlineto{\pgfqpoint{3.805828in}{2.428046in}}%
\pgfpathlineto{\pgfqpoint{3.808863in}{2.442588in}}%
\pgfpathlineto{\pgfqpoint{3.811617in}{2.457204in}}%
\pgfpathlineto{\pgfqpoint{3.814093in}{2.471888in}}%
\pgfpathlineto{\pgfqpoint{3.816291in}{2.486632in}}%
\pgfpathlineto{\pgfqpoint{3.818212in}{2.501429in}}%
\pgfpathlineto{\pgfqpoint{3.819875in}{2.516436in}}%
\pgfpathlineto{\pgfqpoint{3.821236in}{2.531243in}}%
\pgfpathlineto{\pgfqpoint{3.822326in}{2.546108in}}%
\pgfpathlineto{\pgfqpoint{3.823146in}{2.561018in}}%
\pgfpathlineto{\pgfqpoint{3.823692in}{2.575961in}}%
\pgfpathlineto{\pgfqpoint{3.823966in}{2.590923in}}%
\pgfpathlineto{\pgfqpoint{3.823966in}{2.605891in}}%
\pgfpathlineto{\pgfqpoint{3.823692in}{2.620853in}}%
\pgfpathlineto{\pgfqpoint{3.823146in}{2.635795in}}%
\pgfpathlineto{\pgfqpoint{3.822326in}{2.650705in}}%
\pgfpathlineto{\pgfqpoint{3.821236in}{2.665570in}}%
\pgfpathlineto{\pgfqpoint{3.819875in}{2.680378in}}%
\pgfpathlineto{\pgfqpoint{3.818212in}{2.695384in}}%
\pgfpathlineto{\pgfqpoint{3.816291in}{2.710181in}}%
\pgfpathlineto{\pgfqpoint{3.814093in}{2.724925in}}%
\pgfpathlineto{\pgfqpoint{3.811617in}{2.739609in}}%
\pgfpathlineto{\pgfqpoint{3.808863in}{2.754226in}}%
\pgfpathlineto{\pgfqpoint{3.805828in}{2.768768in}}%
\pgfpathlineto{\pgfqpoint{3.802512in}{2.783228in}}%
\pgfpathlineto{\pgfqpoint{3.798913in}{2.797598in}}%
\pgfpathlineto{\pgfqpoint{3.795029in}{2.811870in}}%
\pgfpathlineto{\pgfqpoint{3.790858in}{2.826037in}}%
\pgfpathlineto{\pgfqpoint{3.786398in}{2.840089in}}%
\pgfpathlineto{\pgfqpoint{3.781647in}{2.854019in}}%
\pgfpathlineto{\pgfqpoint{3.776602in}{2.867818in}}%
\pgfpathlineto{\pgfqpoint{3.771261in}{2.881476in}}%
\pgfpathlineto{\pgfqpoint{3.765621in}{2.894982in}}%
\pgfpathlineto{\pgfqpoint{3.759678in}{2.908327in}}%
\pgfpathlineto{\pgfqpoint{3.753430in}{2.921500in}}%
\pgfpathlineto{\pgfqpoint{3.746757in}{2.934708in}}%
\pgfpathlineto{\pgfqpoint{3.739839in}{2.947576in}}%
\pgfpathlineto{\pgfqpoint{3.732587in}{2.960253in}}%
\pgfpathlineto{\pgfqpoint{3.724993in}{2.972730in}}%
\pgfpathlineto{\pgfqpoint{3.717047in}{2.984997in}}%
\pgfpathlineto{\pgfqpoint{3.708736in}{2.997041in}}%
\pgfpathlineto{\pgfqpoint{3.700048in}{3.008851in}}%
\pgfpathlineto{\pgfqpoint{3.690966in}{3.020413in}}%
\pgfpathlineto{\pgfqpoint{3.681473in}{3.031712in}}%
\pgfpathlineto{\pgfqpoint{3.671547in}{3.042731in}}%
\pgfpathlineto{\pgfqpoint{3.661164in}{3.053451in}}%
\pgfpathlineto{\pgfqpoint{3.650298in}{3.063848in}}%
\pgfpathlineto{\pgfqpoint{3.638920in}{3.073894in}}%
\pgfpathlineto{\pgfqpoint{3.626737in}{3.083755in}}%
\pgfpathlineto{\pgfqpoint{3.614031in}{3.093112in}}%
\pgfusepath{stroke}%
\end{pgfscope}%
\begin{pgfscope}%
\pgfpathmoveto{\pgfqpoint{3.460000in}{2.052407in}}%
\pgfpathcurveto{\pgfqpoint{3.749602in}{2.052407in}}{\pgfqpoint{4.027381in}{2.109937in}}{\pgfqpoint{4.232161in}{2.212326in}}%
\pgfpathcurveto{\pgfqpoint{4.436940in}{2.314716in}}{\pgfqpoint{4.552000in}{2.453606in}}{\pgfqpoint{4.552000in}{2.598407in}}%
\pgfpathcurveto{\pgfqpoint{4.552000in}{2.743208in}}{\pgfqpoint{4.436940in}{2.882097in}}{\pgfqpoint{4.232161in}{2.984487in}}%
\pgfpathcurveto{\pgfqpoint{4.027381in}{3.086877in}}{\pgfqpoint{3.749602in}{3.144407in}}{\pgfqpoint{3.460000in}{3.144407in}}%
\pgfpathcurveto{\pgfqpoint{3.170398in}{3.144407in}}{\pgfqpoint{2.892619in}{3.086877in}}{\pgfqpoint{2.687839in}{2.984487in}}%
\pgfpathcurveto{\pgfqpoint{2.483060in}{2.882097in}}{\pgfqpoint{2.368000in}{2.743208in}}{\pgfqpoint{2.368000in}{2.598407in}}%
\pgfpathcurveto{\pgfqpoint{2.368000in}{2.453606in}}{\pgfqpoint{2.483060in}{2.314716in}}{\pgfqpoint{2.687839in}{2.212326in}}%
\pgfpathcurveto{\pgfqpoint{2.892619in}{2.109937in}}{\pgfqpoint{3.170398in}{2.052407in}}{\pgfqpoint{3.460000in}{2.052407in}}%
\pgfpathlineto{\pgfqpoint{3.460000in}{2.052407in}}%
\pgfpathclose%
\pgfusepath{clip}%
\pgfsetrectcap%
\pgfsetroundjoin%
\pgfsetlinewidth{0.451687pt}%
\definecolor{currentstroke}{rgb}{0.690196,0.690196,0.690196}%
\pgfsetstrokecolor{currentstroke}%
\pgfsetstrokeopacity{0.250000}%
\pgfsetdash{}{0pt}%
\pgfpathmoveto{\pgfqpoint{3.691047in}{2.103702in}}%
\pgfpathlineto{\pgfqpoint{3.710106in}{2.113059in}}%
\pgfpathlineto{\pgfqpoint{3.728380in}{2.122919in}}%
\pgfpathlineto{\pgfqpoint{3.745447in}{2.132965in}}%
\pgfpathlineto{\pgfqpoint{3.761746in}{2.143363in}}%
\pgfpathlineto{\pgfqpoint{3.777320in}{2.154082in}}%
\pgfpathlineto{\pgfqpoint{3.792209in}{2.165102in}}%
\pgfpathlineto{\pgfqpoint{3.806449in}{2.176401in}}%
\pgfpathlineto{\pgfqpoint{3.820071in}{2.187962in}}%
\pgfpathlineto{\pgfqpoint{3.833104in}{2.199772in}}%
\pgfpathlineto{\pgfqpoint{3.845570in}{2.211817in}}%
\pgfpathlineto{\pgfqpoint{3.857490in}{2.224083in}}%
\pgfpathlineto{\pgfqpoint{3.868881in}{2.236561in}}%
\pgfpathlineto{\pgfqpoint{3.879758in}{2.249238in}}%
\pgfpathlineto{\pgfqpoint{3.890136in}{2.262105in}}%
\pgfpathlineto{\pgfqpoint{3.900144in}{2.275314in}}%
\pgfpathlineto{\pgfqpoint{3.909517in}{2.288486in}}%
\pgfpathlineto{\pgfqpoint{3.918431in}{2.301831in}}%
\pgfpathlineto{\pgfqpoint{3.926892in}{2.315338in}}%
\pgfpathlineto{\pgfqpoint{3.934904in}{2.328996in}}%
\pgfpathlineto{\pgfqpoint{3.942471in}{2.342794in}}%
\pgfpathlineto{\pgfqpoint{3.949597in}{2.356724in}}%
\pgfpathlineto{\pgfqpoint{3.956287in}{2.370777in}}%
\pgfpathlineto{\pgfqpoint{3.962543in}{2.384944in}}%
\pgfpathlineto{\pgfqpoint{3.968369in}{2.399216in}}%
\pgfpathlineto{\pgfqpoint{3.973768in}{2.413586in}}%
\pgfpathlineto{\pgfqpoint{3.978742in}{2.428046in}}%
\pgfpathlineto{\pgfqpoint{3.983294in}{2.442588in}}%
\pgfpathlineto{\pgfqpoint{3.987426in}{2.457204in}}%
\pgfpathlineto{\pgfqpoint{3.991139in}{2.471888in}}%
\pgfpathlineto{\pgfqpoint{3.994437in}{2.486632in}}%
\pgfpathlineto{\pgfqpoint{3.997319in}{2.501429in}}%
\pgfpathlineto{\pgfqpoint{3.999812in}{2.516436in}}%
\pgfpathlineto{\pgfqpoint{4.001853in}{2.531243in}}%
\pgfpathlineto{\pgfqpoint{4.003490in}{2.546108in}}%
\pgfpathlineto{\pgfqpoint{4.004718in}{2.561018in}}%
\pgfpathlineto{\pgfqpoint{4.005538in}{2.575961in}}%
\pgfpathlineto{\pgfqpoint{4.005949in}{2.590923in}}%
\pgfpathlineto{\pgfqpoint{4.005949in}{2.605891in}}%
\pgfpathlineto{\pgfqpoint{4.005538in}{2.620853in}}%
\pgfpathlineto{\pgfqpoint{4.004718in}{2.635795in}}%
\pgfpathlineto{\pgfqpoint{4.003490in}{2.650705in}}%
\pgfpathlineto{\pgfqpoint{4.001853in}{2.665570in}}%
\pgfpathlineto{\pgfqpoint{3.999812in}{2.680378in}}%
\pgfpathlineto{\pgfqpoint{3.997319in}{2.695384in}}%
\pgfpathlineto{\pgfqpoint{3.994437in}{2.710181in}}%
\pgfpathlineto{\pgfqpoint{3.991139in}{2.724925in}}%
\pgfpathlineto{\pgfqpoint{3.987426in}{2.739609in}}%
\pgfpathlineto{\pgfqpoint{3.983294in}{2.754226in}}%
\pgfpathlineto{\pgfqpoint{3.978742in}{2.768768in}}%
\pgfpathlineto{\pgfqpoint{3.973768in}{2.783228in}}%
\pgfpathlineto{\pgfqpoint{3.968369in}{2.797598in}}%
\pgfpathlineto{\pgfqpoint{3.962543in}{2.811870in}}%
\pgfpathlineto{\pgfqpoint{3.956287in}{2.826037in}}%
\pgfpathlineto{\pgfqpoint{3.949597in}{2.840089in}}%
\pgfpathlineto{\pgfqpoint{3.942471in}{2.854019in}}%
\pgfpathlineto{\pgfqpoint{3.934904in}{2.867818in}}%
\pgfpathlineto{\pgfqpoint{3.926892in}{2.881476in}}%
\pgfpathlineto{\pgfqpoint{3.918431in}{2.894982in}}%
\pgfpathlineto{\pgfqpoint{3.909517in}{2.908327in}}%
\pgfpathlineto{\pgfqpoint{3.900144in}{2.921500in}}%
\pgfpathlineto{\pgfqpoint{3.890136in}{2.934708in}}%
\pgfpathlineto{\pgfqpoint{3.879758in}{2.947576in}}%
\pgfpathlineto{\pgfqpoint{3.868881in}{2.960253in}}%
\pgfpathlineto{\pgfqpoint{3.857490in}{2.972730in}}%
\pgfpathlineto{\pgfqpoint{3.845570in}{2.984997in}}%
\pgfpathlineto{\pgfqpoint{3.833104in}{2.997041in}}%
\pgfpathlineto{\pgfqpoint{3.820071in}{3.008851in}}%
\pgfpathlineto{\pgfqpoint{3.806449in}{3.020413in}}%
\pgfpathlineto{\pgfqpoint{3.792209in}{3.031712in}}%
\pgfpathlineto{\pgfqpoint{3.777320in}{3.042731in}}%
\pgfpathlineto{\pgfqpoint{3.761746in}{3.053451in}}%
\pgfpathlineto{\pgfqpoint{3.745447in}{3.063848in}}%
\pgfpathlineto{\pgfqpoint{3.728380in}{3.073894in}}%
\pgfpathlineto{\pgfqpoint{3.710106in}{3.083755in}}%
\pgfpathlineto{\pgfqpoint{3.691047in}{3.093112in}}%
\pgfusepath{stroke}%
\end{pgfscope}%
\begin{pgfscope}%
\pgfpathmoveto{\pgfqpoint{3.460000in}{2.052407in}}%
\pgfpathcurveto{\pgfqpoint{3.749602in}{2.052407in}}{\pgfqpoint{4.027381in}{2.109937in}}{\pgfqpoint{4.232161in}{2.212326in}}%
\pgfpathcurveto{\pgfqpoint{4.436940in}{2.314716in}}{\pgfqpoint{4.552000in}{2.453606in}}{\pgfqpoint{4.552000in}{2.598407in}}%
\pgfpathcurveto{\pgfqpoint{4.552000in}{2.743208in}}{\pgfqpoint{4.436940in}{2.882097in}}{\pgfqpoint{4.232161in}{2.984487in}}%
\pgfpathcurveto{\pgfqpoint{4.027381in}{3.086877in}}{\pgfqpoint{3.749602in}{3.144407in}}{\pgfqpoint{3.460000in}{3.144407in}}%
\pgfpathcurveto{\pgfqpoint{3.170398in}{3.144407in}}{\pgfqpoint{2.892619in}{3.086877in}}{\pgfqpoint{2.687839in}{2.984487in}}%
\pgfpathcurveto{\pgfqpoint{2.483060in}{2.882097in}}{\pgfqpoint{2.368000in}{2.743208in}}{\pgfqpoint{2.368000in}{2.598407in}}%
\pgfpathcurveto{\pgfqpoint{2.368000in}{2.453606in}}{\pgfqpoint{2.483060in}{2.314716in}}{\pgfqpoint{2.687839in}{2.212326in}}%
\pgfpathcurveto{\pgfqpoint{2.892619in}{2.109937in}}{\pgfqpoint{3.170398in}{2.052407in}}{\pgfqpoint{3.460000in}{2.052407in}}%
\pgfpathlineto{\pgfqpoint{3.460000in}{2.052407in}}%
\pgfpathclose%
\pgfusepath{clip}%
\pgfsetrectcap%
\pgfsetroundjoin%
\pgfsetlinewidth{0.451687pt}%
\definecolor{currentstroke}{rgb}{0.690196,0.690196,0.690196}%
\pgfsetstrokecolor{currentstroke}%
\pgfsetstrokeopacity{0.250000}%
\pgfsetdash{}{0pt}%
\pgfpathmoveto{\pgfqpoint{3.768063in}{2.103702in}}%
\pgfpathlineto{\pgfqpoint{3.793475in}{2.113059in}}%
\pgfpathlineto{\pgfqpoint{3.817840in}{2.122919in}}%
\pgfpathlineto{\pgfqpoint{3.840596in}{2.132965in}}%
\pgfpathlineto{\pgfqpoint{3.862328in}{2.143363in}}%
\pgfpathlineto{\pgfqpoint{3.883093in}{2.154082in}}%
\pgfpathlineto{\pgfqpoint{3.902945in}{2.165102in}}%
\pgfpathlineto{\pgfqpoint{3.921932in}{2.176401in}}%
\pgfpathlineto{\pgfqpoint{3.940095in}{2.187962in}}%
\pgfpathlineto{\pgfqpoint{3.957472in}{2.199772in}}%
\pgfpathlineto{\pgfqpoint{3.974093in}{2.211817in}}%
\pgfpathlineto{\pgfqpoint{3.989986in}{2.224083in}}%
\pgfpathlineto{\pgfqpoint{4.005174in}{2.236561in}}%
\pgfpathlineto{\pgfqpoint{4.019678in}{2.249238in}}%
\pgfpathlineto{\pgfqpoint{4.033515in}{2.262105in}}%
\pgfpathlineto{\pgfqpoint{4.046859in}{2.275314in}}%
\pgfpathlineto{\pgfqpoint{4.059356in}{2.288486in}}%
\pgfpathlineto{\pgfqpoint{4.071241in}{2.301831in}}%
\pgfpathlineto{\pgfqpoint{4.082523in}{2.315338in}}%
\pgfpathlineto{\pgfqpoint{4.093205in}{2.328996in}}%
\pgfpathlineto{\pgfqpoint{4.103295in}{2.342794in}}%
\pgfpathlineto{\pgfqpoint{4.112796in}{2.356724in}}%
\pgfpathlineto{\pgfqpoint{4.121716in}{2.370777in}}%
\pgfpathlineto{\pgfqpoint{4.130057in}{2.384944in}}%
\pgfpathlineto{\pgfqpoint{4.137825in}{2.399216in}}%
\pgfpathlineto{\pgfqpoint{4.145024in}{2.413586in}}%
\pgfpathlineto{\pgfqpoint{4.151656in}{2.428046in}}%
\pgfpathlineto{\pgfqpoint{4.157725in}{2.442588in}}%
\pgfpathlineto{\pgfqpoint{4.163234in}{2.457204in}}%
\pgfpathlineto{\pgfqpoint{4.168186in}{2.471888in}}%
\pgfpathlineto{\pgfqpoint{4.172582in}{2.486632in}}%
\pgfpathlineto{\pgfqpoint{4.176425in}{2.501429in}}%
\pgfpathlineto{\pgfqpoint{4.179749in}{2.516436in}}%
\pgfpathlineto{\pgfqpoint{4.182471in}{2.531243in}}%
\pgfpathlineto{\pgfqpoint{4.184653in}{2.546108in}}%
\pgfpathlineto{\pgfqpoint{4.186291in}{2.561018in}}%
\pgfpathlineto{\pgfqpoint{4.187385in}{2.575961in}}%
\pgfpathlineto{\pgfqpoint{4.187932in}{2.590923in}}%
\pgfpathlineto{\pgfqpoint{4.187932in}{2.605891in}}%
\pgfpathlineto{\pgfqpoint{4.187385in}{2.620853in}}%
\pgfpathlineto{\pgfqpoint{4.186291in}{2.635795in}}%
\pgfpathlineto{\pgfqpoint{4.184653in}{2.650705in}}%
\pgfpathlineto{\pgfqpoint{4.182471in}{2.665570in}}%
\pgfpathlineto{\pgfqpoint{4.179749in}{2.680378in}}%
\pgfpathlineto{\pgfqpoint{4.176425in}{2.695384in}}%
\pgfpathlineto{\pgfqpoint{4.172582in}{2.710181in}}%
\pgfpathlineto{\pgfqpoint{4.168186in}{2.724925in}}%
\pgfpathlineto{\pgfqpoint{4.163234in}{2.739609in}}%
\pgfpathlineto{\pgfqpoint{4.157725in}{2.754226in}}%
\pgfpathlineto{\pgfqpoint{4.151656in}{2.768768in}}%
\pgfpathlineto{\pgfqpoint{4.145024in}{2.783228in}}%
\pgfpathlineto{\pgfqpoint{4.137825in}{2.797598in}}%
\pgfpathlineto{\pgfqpoint{4.130057in}{2.811870in}}%
\pgfpathlineto{\pgfqpoint{4.121716in}{2.826037in}}%
\pgfpathlineto{\pgfqpoint{4.112796in}{2.840089in}}%
\pgfpathlineto{\pgfqpoint{4.103295in}{2.854019in}}%
\pgfpathlineto{\pgfqpoint{4.093205in}{2.867818in}}%
\pgfpathlineto{\pgfqpoint{4.082523in}{2.881476in}}%
\pgfpathlineto{\pgfqpoint{4.071241in}{2.894982in}}%
\pgfpathlineto{\pgfqpoint{4.059356in}{2.908327in}}%
\pgfpathlineto{\pgfqpoint{4.046859in}{2.921500in}}%
\pgfpathlineto{\pgfqpoint{4.033515in}{2.934708in}}%
\pgfpathlineto{\pgfqpoint{4.019678in}{2.947576in}}%
\pgfpathlineto{\pgfqpoint{4.005174in}{2.960253in}}%
\pgfpathlineto{\pgfqpoint{3.989986in}{2.972730in}}%
\pgfpathlineto{\pgfqpoint{3.974093in}{2.984997in}}%
\pgfpathlineto{\pgfqpoint{3.957472in}{2.997041in}}%
\pgfpathlineto{\pgfqpoint{3.940095in}{3.008851in}}%
\pgfpathlineto{\pgfqpoint{3.921932in}{3.020413in}}%
\pgfpathlineto{\pgfqpoint{3.902945in}{3.031712in}}%
\pgfpathlineto{\pgfqpoint{3.883093in}{3.042731in}}%
\pgfpathlineto{\pgfqpoint{3.862328in}{3.053451in}}%
\pgfpathlineto{\pgfqpoint{3.840596in}{3.063848in}}%
\pgfpathlineto{\pgfqpoint{3.817840in}{3.073894in}}%
\pgfpathlineto{\pgfqpoint{3.793475in}{3.083755in}}%
\pgfpathlineto{\pgfqpoint{3.768063in}{3.093112in}}%
\pgfusepath{stroke}%
\end{pgfscope}%
\begin{pgfscope}%
\pgfpathmoveto{\pgfqpoint{3.460000in}{2.052407in}}%
\pgfpathcurveto{\pgfqpoint{3.749602in}{2.052407in}}{\pgfqpoint{4.027381in}{2.109937in}}{\pgfqpoint{4.232161in}{2.212326in}}%
\pgfpathcurveto{\pgfqpoint{4.436940in}{2.314716in}}{\pgfqpoint{4.552000in}{2.453606in}}{\pgfqpoint{4.552000in}{2.598407in}}%
\pgfpathcurveto{\pgfqpoint{4.552000in}{2.743208in}}{\pgfqpoint{4.436940in}{2.882097in}}{\pgfqpoint{4.232161in}{2.984487in}}%
\pgfpathcurveto{\pgfqpoint{4.027381in}{3.086877in}}{\pgfqpoint{3.749602in}{3.144407in}}{\pgfqpoint{3.460000in}{3.144407in}}%
\pgfpathcurveto{\pgfqpoint{3.170398in}{3.144407in}}{\pgfqpoint{2.892619in}{3.086877in}}{\pgfqpoint{2.687839in}{2.984487in}}%
\pgfpathcurveto{\pgfqpoint{2.483060in}{2.882097in}}{\pgfqpoint{2.368000in}{2.743208in}}{\pgfqpoint{2.368000in}{2.598407in}}%
\pgfpathcurveto{\pgfqpoint{2.368000in}{2.453606in}}{\pgfqpoint{2.483060in}{2.314716in}}{\pgfqpoint{2.687839in}{2.212326in}}%
\pgfpathcurveto{\pgfqpoint{2.892619in}{2.109937in}}{\pgfqpoint{3.170398in}{2.052407in}}{\pgfqpoint{3.460000in}{2.052407in}}%
\pgfpathlineto{\pgfqpoint{3.460000in}{2.052407in}}%
\pgfpathclose%
\pgfusepath{clip}%
\pgfsetrectcap%
\pgfsetroundjoin%
\pgfsetlinewidth{0.451687pt}%
\definecolor{currentstroke}{rgb}{0.690196,0.690196,0.690196}%
\pgfsetstrokecolor{currentstroke}%
\pgfsetstrokeopacity{0.250000}%
\pgfsetdash{}{0pt}%
\pgfpathmoveto{\pgfqpoint{3.845079in}{2.103702in}}%
\pgfpathlineto{\pgfqpoint{3.876843in}{2.113059in}}%
\pgfpathlineto{\pgfqpoint{3.907300in}{2.122919in}}%
\pgfpathlineto{\pgfqpoint{3.935745in}{2.132965in}}%
\pgfpathlineto{\pgfqpoint{3.962910in}{2.143363in}}%
\pgfpathlineto{\pgfqpoint{3.988866in}{2.154082in}}%
\pgfpathlineto{\pgfqpoint{4.013681in}{2.165102in}}%
\pgfpathlineto{\pgfqpoint{4.037415in}{2.176401in}}%
\pgfpathlineto{\pgfqpoint{4.060119in}{2.187962in}}%
\pgfpathlineto{\pgfqpoint{4.081840in}{2.199772in}}%
\pgfpathlineto{\pgfqpoint{4.102616in}{2.211817in}}%
\pgfpathlineto{\pgfqpoint{4.122483in}{2.224083in}}%
\pgfpathlineto{\pgfqpoint{4.141468in}{2.236561in}}%
\pgfpathlineto{\pgfqpoint{4.159597in}{2.249238in}}%
\pgfpathlineto{\pgfqpoint{4.176894in}{2.262105in}}%
\pgfpathlineto{\pgfqpoint{4.193574in}{2.275314in}}%
\pgfpathlineto{\pgfqpoint{4.209194in}{2.288486in}}%
\pgfpathlineto{\pgfqpoint{4.224052in}{2.301831in}}%
\pgfpathlineto{\pgfqpoint{4.238153in}{2.315338in}}%
\pgfpathlineto{\pgfqpoint{4.251506in}{2.328996in}}%
\pgfpathlineto{\pgfqpoint{4.264118in}{2.342794in}}%
\pgfpathlineto{\pgfqpoint{4.275996in}{2.356724in}}%
\pgfpathlineto{\pgfqpoint{4.287145in}{2.370777in}}%
\pgfpathlineto{\pgfqpoint{4.297572in}{2.384944in}}%
\pgfpathlineto{\pgfqpoint{4.307282in}{2.399216in}}%
\pgfpathlineto{\pgfqpoint{4.316280in}{2.413586in}}%
\pgfpathlineto{\pgfqpoint{4.324570in}{2.428046in}}%
\pgfpathlineto{\pgfqpoint{4.332156in}{2.442588in}}%
\pgfpathlineto{\pgfqpoint{4.339043in}{2.457204in}}%
\pgfpathlineto{\pgfqpoint{4.345232in}{2.471888in}}%
\pgfpathlineto{\pgfqpoint{4.350728in}{2.486632in}}%
\pgfpathlineto{\pgfqpoint{4.355531in}{2.501429in}}%
\pgfpathlineto{\pgfqpoint{4.359686in}{2.516436in}}%
\pgfpathlineto{\pgfqpoint{4.363089in}{2.531243in}}%
\pgfpathlineto{\pgfqpoint{4.365816in}{2.546108in}}%
\pgfpathlineto{\pgfqpoint{4.367864in}{2.561018in}}%
\pgfpathlineto{\pgfqpoint{4.369231in}{2.575961in}}%
\pgfpathlineto{\pgfqpoint{4.369915in}{2.590923in}}%
\pgfpathlineto{\pgfqpoint{4.369915in}{2.605891in}}%
\pgfpathlineto{\pgfqpoint{4.369231in}{2.620853in}}%
\pgfpathlineto{\pgfqpoint{4.367864in}{2.635795in}}%
\pgfpathlineto{\pgfqpoint{4.365816in}{2.650705in}}%
\pgfpathlineto{\pgfqpoint{4.363089in}{2.665570in}}%
\pgfpathlineto{\pgfqpoint{4.359686in}{2.680378in}}%
\pgfpathlineto{\pgfqpoint{4.355531in}{2.695384in}}%
\pgfpathlineto{\pgfqpoint{4.350728in}{2.710181in}}%
\pgfpathlineto{\pgfqpoint{4.345232in}{2.724925in}}%
\pgfpathlineto{\pgfqpoint{4.339043in}{2.739609in}}%
\pgfpathlineto{\pgfqpoint{4.332156in}{2.754226in}}%
\pgfpathlineto{\pgfqpoint{4.324570in}{2.768768in}}%
\pgfpathlineto{\pgfqpoint{4.316280in}{2.783228in}}%
\pgfpathlineto{\pgfqpoint{4.307282in}{2.797598in}}%
\pgfpathlineto{\pgfqpoint{4.297572in}{2.811870in}}%
\pgfpathlineto{\pgfqpoint{4.287145in}{2.826037in}}%
\pgfpathlineto{\pgfqpoint{4.275996in}{2.840089in}}%
\pgfpathlineto{\pgfqpoint{4.264118in}{2.854019in}}%
\pgfpathlineto{\pgfqpoint{4.251506in}{2.867818in}}%
\pgfpathlineto{\pgfqpoint{4.238153in}{2.881476in}}%
\pgfpathlineto{\pgfqpoint{4.224052in}{2.894982in}}%
\pgfpathlineto{\pgfqpoint{4.209194in}{2.908327in}}%
\pgfpathlineto{\pgfqpoint{4.193574in}{2.921500in}}%
\pgfpathlineto{\pgfqpoint{4.176894in}{2.934708in}}%
\pgfpathlineto{\pgfqpoint{4.159597in}{2.947576in}}%
\pgfpathlineto{\pgfqpoint{4.141468in}{2.960253in}}%
\pgfpathlineto{\pgfqpoint{4.122483in}{2.972730in}}%
\pgfpathlineto{\pgfqpoint{4.102616in}{2.984997in}}%
\pgfpathlineto{\pgfqpoint{4.081840in}{2.997041in}}%
\pgfpathlineto{\pgfqpoint{4.060119in}{3.008851in}}%
\pgfpathlineto{\pgfqpoint{4.037415in}{3.020413in}}%
\pgfpathlineto{\pgfqpoint{4.013681in}{3.031712in}}%
\pgfpathlineto{\pgfqpoint{3.988866in}{3.042731in}}%
\pgfpathlineto{\pgfqpoint{3.962910in}{3.053451in}}%
\pgfpathlineto{\pgfqpoint{3.935745in}{3.063848in}}%
\pgfpathlineto{\pgfqpoint{3.907300in}{3.073894in}}%
\pgfpathlineto{\pgfqpoint{3.876843in}{3.083755in}}%
\pgfpathlineto{\pgfqpoint{3.845079in}{3.093112in}}%
\pgfusepath{stroke}%
\end{pgfscope}%
\begin{pgfscope}%
\pgfpathmoveto{\pgfqpoint{3.460000in}{2.052407in}}%
\pgfpathcurveto{\pgfqpoint{3.749602in}{2.052407in}}{\pgfqpoint{4.027381in}{2.109937in}}{\pgfqpoint{4.232161in}{2.212326in}}%
\pgfpathcurveto{\pgfqpoint{4.436940in}{2.314716in}}{\pgfqpoint{4.552000in}{2.453606in}}{\pgfqpoint{4.552000in}{2.598407in}}%
\pgfpathcurveto{\pgfqpoint{4.552000in}{2.743208in}}{\pgfqpoint{4.436940in}{2.882097in}}{\pgfqpoint{4.232161in}{2.984487in}}%
\pgfpathcurveto{\pgfqpoint{4.027381in}{3.086877in}}{\pgfqpoint{3.749602in}{3.144407in}}{\pgfqpoint{3.460000in}{3.144407in}}%
\pgfpathcurveto{\pgfqpoint{3.170398in}{3.144407in}}{\pgfqpoint{2.892619in}{3.086877in}}{\pgfqpoint{2.687839in}{2.984487in}}%
\pgfpathcurveto{\pgfqpoint{2.483060in}{2.882097in}}{\pgfqpoint{2.368000in}{2.743208in}}{\pgfqpoint{2.368000in}{2.598407in}}%
\pgfpathcurveto{\pgfqpoint{2.368000in}{2.453606in}}{\pgfqpoint{2.483060in}{2.314716in}}{\pgfqpoint{2.687839in}{2.212326in}}%
\pgfpathcurveto{\pgfqpoint{2.892619in}{2.109937in}}{\pgfqpoint{3.170398in}{2.052407in}}{\pgfqpoint{3.460000in}{2.052407in}}%
\pgfpathlineto{\pgfqpoint{3.460000in}{2.052407in}}%
\pgfpathclose%
\pgfusepath{clip}%
\pgfsetrectcap%
\pgfsetroundjoin%
\pgfsetlinewidth{0.451687pt}%
\definecolor{currentstroke}{rgb}{0.690196,0.690196,0.690196}%
\pgfsetstrokecolor{currentstroke}%
\pgfsetstrokeopacity{0.250000}%
\pgfsetdash{}{0pt}%
\pgfpathmoveto{\pgfqpoint{2.997906in}{2.103702in}}%
\pgfpathlineto{\pgfqpoint{3.010228in}{2.103702in}}%
\pgfpathlineto{\pgfqpoint{3.022551in}{2.103702in}}%
\pgfpathlineto{\pgfqpoint{3.034873in}{2.103702in}}%
\pgfpathlineto{\pgfqpoint{3.047196in}{2.103702in}}%
\pgfpathlineto{\pgfqpoint{3.059518in}{2.103702in}}%
\pgfpathlineto{\pgfqpoint{3.071841in}{2.103702in}}%
\pgfpathlineto{\pgfqpoint{3.084163in}{2.103702in}}%
\pgfpathlineto{\pgfqpoint{3.096486in}{2.103702in}}%
\pgfpathlineto{\pgfqpoint{3.108808in}{2.103702in}}%
\pgfpathlineto{\pgfqpoint{3.121131in}{2.103702in}}%
\pgfpathlineto{\pgfqpoint{3.133453in}{2.103702in}}%
\pgfpathlineto{\pgfqpoint{3.145776in}{2.103702in}}%
\pgfpathlineto{\pgfqpoint{3.158098in}{2.103702in}}%
\pgfpathlineto{\pgfqpoint{3.170421in}{2.103702in}}%
\pgfpathlineto{\pgfqpoint{3.182743in}{2.103702in}}%
\pgfpathlineto{\pgfqpoint{3.195066in}{2.103702in}}%
\pgfpathlineto{\pgfqpoint{3.207388in}{2.103702in}}%
\pgfpathlineto{\pgfqpoint{3.219711in}{2.103702in}}%
\pgfpathlineto{\pgfqpoint{3.232033in}{2.103702in}}%
\pgfpathlineto{\pgfqpoint{3.244356in}{2.103702in}}%
\pgfpathlineto{\pgfqpoint{3.256678in}{2.103702in}}%
\pgfpathlineto{\pgfqpoint{3.269001in}{2.103702in}}%
\pgfpathlineto{\pgfqpoint{3.281323in}{2.103702in}}%
\pgfpathlineto{\pgfqpoint{3.293646in}{2.103702in}}%
\pgfpathlineto{\pgfqpoint{3.305969in}{2.103702in}}%
\pgfpathlineto{\pgfqpoint{3.318291in}{2.103702in}}%
\pgfpathlineto{\pgfqpoint{3.330614in}{2.103702in}}%
\pgfpathlineto{\pgfqpoint{3.342936in}{2.103702in}}%
\pgfpathlineto{\pgfqpoint{3.355259in}{2.103702in}}%
\pgfpathlineto{\pgfqpoint{3.367581in}{2.103702in}}%
\pgfpathlineto{\pgfqpoint{3.379904in}{2.103702in}}%
\pgfpathlineto{\pgfqpoint{3.392226in}{2.103702in}}%
\pgfpathlineto{\pgfqpoint{3.404549in}{2.103702in}}%
\pgfpathlineto{\pgfqpoint{3.416871in}{2.103702in}}%
\pgfpathlineto{\pgfqpoint{3.429194in}{2.103702in}}%
\pgfpathlineto{\pgfqpoint{3.441516in}{2.103702in}}%
\pgfpathlineto{\pgfqpoint{3.453839in}{2.103702in}}%
\pgfpathlineto{\pgfqpoint{3.466161in}{2.103702in}}%
\pgfpathlineto{\pgfqpoint{3.478484in}{2.103702in}}%
\pgfpathlineto{\pgfqpoint{3.490806in}{2.103702in}}%
\pgfpathlineto{\pgfqpoint{3.503129in}{2.103702in}}%
\pgfpathlineto{\pgfqpoint{3.515451in}{2.103702in}}%
\pgfpathlineto{\pgfqpoint{3.527774in}{2.103702in}}%
\pgfpathlineto{\pgfqpoint{3.540096in}{2.103702in}}%
\pgfpathlineto{\pgfqpoint{3.552419in}{2.103702in}}%
\pgfpathlineto{\pgfqpoint{3.564741in}{2.103702in}}%
\pgfpathlineto{\pgfqpoint{3.577064in}{2.103702in}}%
\pgfpathlineto{\pgfqpoint{3.589386in}{2.103702in}}%
\pgfpathlineto{\pgfqpoint{3.601709in}{2.103702in}}%
\pgfpathlineto{\pgfqpoint{3.614031in}{2.103702in}}%
\pgfpathlineto{\pgfqpoint{3.626354in}{2.103702in}}%
\pgfpathlineto{\pgfqpoint{3.638677in}{2.103702in}}%
\pgfpathlineto{\pgfqpoint{3.650999in}{2.103702in}}%
\pgfpathlineto{\pgfqpoint{3.663322in}{2.103702in}}%
\pgfpathlineto{\pgfqpoint{3.675644in}{2.103702in}}%
\pgfpathlineto{\pgfqpoint{3.687967in}{2.103702in}}%
\pgfpathlineto{\pgfqpoint{3.700289in}{2.103702in}}%
\pgfpathlineto{\pgfqpoint{3.712612in}{2.103702in}}%
\pgfpathlineto{\pgfqpoint{3.724934in}{2.103702in}}%
\pgfpathlineto{\pgfqpoint{3.737257in}{2.103702in}}%
\pgfpathlineto{\pgfqpoint{3.749579in}{2.103702in}}%
\pgfpathlineto{\pgfqpoint{3.761902in}{2.103702in}}%
\pgfpathlineto{\pgfqpoint{3.774224in}{2.103702in}}%
\pgfpathlineto{\pgfqpoint{3.786547in}{2.103702in}}%
\pgfpathlineto{\pgfqpoint{3.798869in}{2.103702in}}%
\pgfpathlineto{\pgfqpoint{3.811192in}{2.103702in}}%
\pgfpathlineto{\pgfqpoint{3.823514in}{2.103702in}}%
\pgfpathlineto{\pgfqpoint{3.835837in}{2.103702in}}%
\pgfpathlineto{\pgfqpoint{3.848159in}{2.103702in}}%
\pgfpathlineto{\pgfqpoint{3.860482in}{2.103702in}}%
\pgfpathlineto{\pgfqpoint{3.872804in}{2.103702in}}%
\pgfpathlineto{\pgfqpoint{3.885127in}{2.103702in}}%
\pgfpathlineto{\pgfqpoint{3.897449in}{2.103702in}}%
\pgfpathlineto{\pgfqpoint{3.909772in}{2.103702in}}%
\pgfpathlineto{\pgfqpoint{3.922094in}{2.103702in}}%
\pgfusepath{stroke}%
\end{pgfscope}%
\begin{pgfscope}%
\pgfpathmoveto{\pgfqpoint{3.460000in}{2.052407in}}%
\pgfpathcurveto{\pgfqpoint{3.749602in}{2.052407in}}{\pgfqpoint{4.027381in}{2.109937in}}{\pgfqpoint{4.232161in}{2.212326in}}%
\pgfpathcurveto{\pgfqpoint{4.436940in}{2.314716in}}{\pgfqpoint{4.552000in}{2.453606in}}{\pgfqpoint{4.552000in}{2.598407in}}%
\pgfpathcurveto{\pgfqpoint{4.552000in}{2.743208in}}{\pgfqpoint{4.436940in}{2.882097in}}{\pgfqpoint{4.232161in}{2.984487in}}%
\pgfpathcurveto{\pgfqpoint{4.027381in}{3.086877in}}{\pgfqpoint{3.749602in}{3.144407in}}{\pgfqpoint{3.460000in}{3.144407in}}%
\pgfpathcurveto{\pgfqpoint{3.170398in}{3.144407in}}{\pgfqpoint{2.892619in}{3.086877in}}{\pgfqpoint{2.687839in}{2.984487in}}%
\pgfpathcurveto{\pgfqpoint{2.483060in}{2.882097in}}{\pgfqpoint{2.368000in}{2.743208in}}{\pgfqpoint{2.368000in}{2.598407in}}%
\pgfpathcurveto{\pgfqpoint{2.368000in}{2.453606in}}{\pgfqpoint{2.483060in}{2.314716in}}{\pgfqpoint{2.687839in}{2.212326in}}%
\pgfpathcurveto{\pgfqpoint{2.892619in}{2.109937in}}{\pgfqpoint{3.170398in}{2.052407in}}{\pgfqpoint{3.460000in}{2.052407in}}%
\pgfpathlineto{\pgfqpoint{3.460000in}{2.052407in}}%
\pgfpathclose%
\pgfusepath{clip}%
\pgfsetrectcap%
\pgfsetroundjoin%
\pgfsetlinewidth{0.451687pt}%
\definecolor{currentstroke}{rgb}{0.690196,0.690196,0.690196}%
\pgfsetstrokecolor{currentstroke}%
\pgfsetstrokeopacity{0.250000}%
\pgfsetdash{}{0pt}%
\pgfpathmoveto{\pgfqpoint{2.753329in}{2.182150in}}%
\pgfpathlineto{\pgfqpoint{2.772174in}{2.182150in}}%
\pgfpathlineto{\pgfqpoint{2.791018in}{2.182150in}}%
\pgfpathlineto{\pgfqpoint{2.809863in}{2.182150in}}%
\pgfpathlineto{\pgfqpoint{2.828707in}{2.182150in}}%
\pgfpathlineto{\pgfqpoint{2.847552in}{2.182150in}}%
\pgfpathlineto{\pgfqpoint{2.866396in}{2.182150in}}%
\pgfpathlineto{\pgfqpoint{2.885241in}{2.182150in}}%
\pgfpathlineto{\pgfqpoint{2.904086in}{2.182150in}}%
\pgfpathlineto{\pgfqpoint{2.922930in}{2.182150in}}%
\pgfpathlineto{\pgfqpoint{2.941775in}{2.182150in}}%
\pgfpathlineto{\pgfqpoint{2.960619in}{2.182150in}}%
\pgfpathlineto{\pgfqpoint{2.979464in}{2.182150in}}%
\pgfpathlineto{\pgfqpoint{2.998308in}{2.182150in}}%
\pgfpathlineto{\pgfqpoint{3.017153in}{2.182150in}}%
\pgfpathlineto{\pgfqpoint{3.035997in}{2.182150in}}%
\pgfpathlineto{\pgfqpoint{3.054842in}{2.182150in}}%
\pgfpathlineto{\pgfqpoint{3.073687in}{2.182150in}}%
\pgfpathlineto{\pgfqpoint{3.092531in}{2.182150in}}%
\pgfpathlineto{\pgfqpoint{3.111376in}{2.182150in}}%
\pgfpathlineto{\pgfqpoint{3.130220in}{2.182150in}}%
\pgfpathlineto{\pgfqpoint{3.149065in}{2.182150in}}%
\pgfpathlineto{\pgfqpoint{3.167909in}{2.182150in}}%
\pgfpathlineto{\pgfqpoint{3.186754in}{2.182150in}}%
\pgfpathlineto{\pgfqpoint{3.205598in}{2.182150in}}%
\pgfpathlineto{\pgfqpoint{3.224443in}{2.182150in}}%
\pgfpathlineto{\pgfqpoint{3.243288in}{2.182150in}}%
\pgfpathlineto{\pgfqpoint{3.262132in}{2.182150in}}%
\pgfpathlineto{\pgfqpoint{3.280977in}{2.182150in}}%
\pgfpathlineto{\pgfqpoint{3.299821in}{2.182150in}}%
\pgfpathlineto{\pgfqpoint{3.318666in}{2.182150in}}%
\pgfpathlineto{\pgfqpoint{3.337510in}{2.182150in}}%
\pgfpathlineto{\pgfqpoint{3.356355in}{2.182150in}}%
\pgfpathlineto{\pgfqpoint{3.375199in}{2.182150in}}%
\pgfpathlineto{\pgfqpoint{3.394044in}{2.182150in}}%
\pgfpathlineto{\pgfqpoint{3.412889in}{2.182150in}}%
\pgfpathlineto{\pgfqpoint{3.431733in}{2.182150in}}%
\pgfpathlineto{\pgfqpoint{3.450578in}{2.182150in}}%
\pgfpathlineto{\pgfqpoint{3.469422in}{2.182150in}}%
\pgfpathlineto{\pgfqpoint{3.488267in}{2.182150in}}%
\pgfpathlineto{\pgfqpoint{3.507111in}{2.182150in}}%
\pgfpathlineto{\pgfqpoint{3.525956in}{2.182150in}}%
\pgfpathlineto{\pgfqpoint{3.544801in}{2.182150in}}%
\pgfpathlineto{\pgfqpoint{3.563645in}{2.182150in}}%
\pgfpathlineto{\pgfqpoint{3.582490in}{2.182150in}}%
\pgfpathlineto{\pgfqpoint{3.601334in}{2.182150in}}%
\pgfpathlineto{\pgfqpoint{3.620179in}{2.182150in}}%
\pgfpathlineto{\pgfqpoint{3.639023in}{2.182150in}}%
\pgfpathlineto{\pgfqpoint{3.657868in}{2.182150in}}%
\pgfpathlineto{\pgfqpoint{3.676712in}{2.182150in}}%
\pgfpathlineto{\pgfqpoint{3.695557in}{2.182150in}}%
\pgfpathlineto{\pgfqpoint{3.714402in}{2.182150in}}%
\pgfpathlineto{\pgfqpoint{3.733246in}{2.182150in}}%
\pgfpathlineto{\pgfqpoint{3.752091in}{2.182150in}}%
\pgfpathlineto{\pgfqpoint{3.770935in}{2.182150in}}%
\pgfpathlineto{\pgfqpoint{3.789780in}{2.182150in}}%
\pgfpathlineto{\pgfqpoint{3.808624in}{2.182150in}}%
\pgfpathlineto{\pgfqpoint{3.827469in}{2.182150in}}%
\pgfpathlineto{\pgfqpoint{3.846313in}{2.182150in}}%
\pgfpathlineto{\pgfqpoint{3.865158in}{2.182150in}}%
\pgfpathlineto{\pgfqpoint{3.884003in}{2.182150in}}%
\pgfpathlineto{\pgfqpoint{3.902847in}{2.182150in}}%
\pgfpathlineto{\pgfqpoint{3.921692in}{2.182150in}}%
\pgfpathlineto{\pgfqpoint{3.940536in}{2.182150in}}%
\pgfpathlineto{\pgfqpoint{3.959381in}{2.182150in}}%
\pgfpathlineto{\pgfqpoint{3.978225in}{2.182150in}}%
\pgfpathlineto{\pgfqpoint{3.997070in}{2.182150in}}%
\pgfpathlineto{\pgfqpoint{4.015914in}{2.182150in}}%
\pgfpathlineto{\pgfqpoint{4.034759in}{2.182150in}}%
\pgfpathlineto{\pgfqpoint{4.053604in}{2.182150in}}%
\pgfpathlineto{\pgfqpoint{4.072448in}{2.182150in}}%
\pgfpathlineto{\pgfqpoint{4.091293in}{2.182150in}}%
\pgfpathlineto{\pgfqpoint{4.110137in}{2.182150in}}%
\pgfpathlineto{\pgfqpoint{4.128982in}{2.182150in}}%
\pgfpathlineto{\pgfqpoint{4.147826in}{2.182150in}}%
\pgfpathlineto{\pgfqpoint{4.166671in}{2.182150in}}%
\pgfusepath{stroke}%
\end{pgfscope}%
\begin{pgfscope}%
\pgfpathmoveto{\pgfqpoint{3.460000in}{2.052407in}}%
\pgfpathcurveto{\pgfqpoint{3.749602in}{2.052407in}}{\pgfqpoint{4.027381in}{2.109937in}}{\pgfqpoint{4.232161in}{2.212326in}}%
\pgfpathcurveto{\pgfqpoint{4.436940in}{2.314716in}}{\pgfqpoint{4.552000in}{2.453606in}}{\pgfqpoint{4.552000in}{2.598407in}}%
\pgfpathcurveto{\pgfqpoint{4.552000in}{2.743208in}}{\pgfqpoint{4.436940in}{2.882097in}}{\pgfqpoint{4.232161in}{2.984487in}}%
\pgfpathcurveto{\pgfqpoint{4.027381in}{3.086877in}}{\pgfqpoint{3.749602in}{3.144407in}}{\pgfqpoint{3.460000in}{3.144407in}}%
\pgfpathcurveto{\pgfqpoint{3.170398in}{3.144407in}}{\pgfqpoint{2.892619in}{3.086877in}}{\pgfqpoint{2.687839in}{2.984487in}}%
\pgfpathcurveto{\pgfqpoint{2.483060in}{2.882097in}}{\pgfqpoint{2.368000in}{2.743208in}}{\pgfqpoint{2.368000in}{2.598407in}}%
\pgfpathcurveto{\pgfqpoint{2.368000in}{2.453606in}}{\pgfqpoint{2.483060in}{2.314716in}}{\pgfqpoint{2.687839in}{2.212326in}}%
\pgfpathcurveto{\pgfqpoint{2.892619in}{2.109937in}}{\pgfqpoint{3.170398in}{2.052407in}}{\pgfqpoint{3.460000in}{2.052407in}}%
\pgfpathlineto{\pgfqpoint{3.460000in}{2.052407in}}%
\pgfpathclose%
\pgfusepath{clip}%
\pgfsetrectcap%
\pgfsetroundjoin%
\pgfsetlinewidth{0.451687pt}%
\definecolor{currentstroke}{rgb}{0.690196,0.690196,0.690196}%
\pgfsetstrokecolor{currentstroke}%
\pgfsetstrokeopacity{0.250000}%
\pgfsetdash{}{0pt}%
\pgfpathmoveto{\pgfqpoint{2.579711in}{2.275314in}}%
\pgfpathlineto{\pgfqpoint{2.603186in}{2.275314in}}%
\pgfpathlineto{\pgfqpoint{2.626660in}{2.275314in}}%
\pgfpathlineto{\pgfqpoint{2.650134in}{2.275314in}}%
\pgfpathlineto{\pgfqpoint{2.673609in}{2.275314in}}%
\pgfpathlineto{\pgfqpoint{2.697083in}{2.275314in}}%
\pgfpathlineto{\pgfqpoint{2.720558in}{2.275314in}}%
\pgfpathlineto{\pgfqpoint{2.744032in}{2.275314in}}%
\pgfpathlineto{\pgfqpoint{2.767506in}{2.275314in}}%
\pgfpathlineto{\pgfqpoint{2.790981in}{2.275314in}}%
\pgfpathlineto{\pgfqpoint{2.814455in}{2.275314in}}%
\pgfpathlineto{\pgfqpoint{2.837929in}{2.275314in}}%
\pgfpathlineto{\pgfqpoint{2.861404in}{2.275314in}}%
\pgfpathlineto{\pgfqpoint{2.884878in}{2.275314in}}%
\pgfpathlineto{\pgfqpoint{2.908352in}{2.275314in}}%
\pgfpathlineto{\pgfqpoint{2.931827in}{2.275314in}}%
\pgfpathlineto{\pgfqpoint{2.955301in}{2.275314in}}%
\pgfpathlineto{\pgfqpoint{2.978776in}{2.275314in}}%
\pgfpathlineto{\pgfqpoint{3.002250in}{2.275314in}}%
\pgfpathlineto{\pgfqpoint{3.025724in}{2.275314in}}%
\pgfpathlineto{\pgfqpoint{3.049199in}{2.275314in}}%
\pgfpathlineto{\pgfqpoint{3.072673in}{2.275314in}}%
\pgfpathlineto{\pgfqpoint{3.096147in}{2.275314in}}%
\pgfpathlineto{\pgfqpoint{3.119622in}{2.275314in}}%
\pgfpathlineto{\pgfqpoint{3.143096in}{2.275314in}}%
\pgfpathlineto{\pgfqpoint{3.166570in}{2.275314in}}%
\pgfpathlineto{\pgfqpoint{3.190045in}{2.275314in}}%
\pgfpathlineto{\pgfqpoint{3.213519in}{2.275314in}}%
\pgfpathlineto{\pgfqpoint{3.236994in}{2.275314in}}%
\pgfpathlineto{\pgfqpoint{3.260468in}{2.275314in}}%
\pgfpathlineto{\pgfqpoint{3.283942in}{2.275314in}}%
\pgfpathlineto{\pgfqpoint{3.307417in}{2.275314in}}%
\pgfpathlineto{\pgfqpoint{3.330891in}{2.275314in}}%
\pgfpathlineto{\pgfqpoint{3.354365in}{2.275314in}}%
\pgfpathlineto{\pgfqpoint{3.377840in}{2.275314in}}%
\pgfpathlineto{\pgfqpoint{3.401314in}{2.275314in}}%
\pgfpathlineto{\pgfqpoint{3.424788in}{2.275314in}}%
\pgfpathlineto{\pgfqpoint{3.448263in}{2.275314in}}%
\pgfpathlineto{\pgfqpoint{3.471737in}{2.275314in}}%
\pgfpathlineto{\pgfqpoint{3.495212in}{2.275314in}}%
\pgfpathlineto{\pgfqpoint{3.518686in}{2.275314in}}%
\pgfpathlineto{\pgfqpoint{3.542160in}{2.275314in}}%
\pgfpathlineto{\pgfqpoint{3.565635in}{2.275314in}}%
\pgfpathlineto{\pgfqpoint{3.589109in}{2.275314in}}%
\pgfpathlineto{\pgfqpoint{3.612583in}{2.275314in}}%
\pgfpathlineto{\pgfqpoint{3.636058in}{2.275314in}}%
\pgfpathlineto{\pgfqpoint{3.659532in}{2.275314in}}%
\pgfpathlineto{\pgfqpoint{3.683006in}{2.275314in}}%
\pgfpathlineto{\pgfqpoint{3.706481in}{2.275314in}}%
\pgfpathlineto{\pgfqpoint{3.729955in}{2.275314in}}%
\pgfpathlineto{\pgfqpoint{3.753430in}{2.275314in}}%
\pgfpathlineto{\pgfqpoint{3.776904in}{2.275314in}}%
\pgfpathlineto{\pgfqpoint{3.800378in}{2.275314in}}%
\pgfpathlineto{\pgfqpoint{3.823853in}{2.275314in}}%
\pgfpathlineto{\pgfqpoint{3.847327in}{2.275314in}}%
\pgfpathlineto{\pgfqpoint{3.870801in}{2.275314in}}%
\pgfpathlineto{\pgfqpoint{3.894276in}{2.275314in}}%
\pgfpathlineto{\pgfqpoint{3.917750in}{2.275314in}}%
\pgfpathlineto{\pgfqpoint{3.941224in}{2.275314in}}%
\pgfpathlineto{\pgfqpoint{3.964699in}{2.275314in}}%
\pgfpathlineto{\pgfqpoint{3.988173in}{2.275314in}}%
\pgfpathlineto{\pgfqpoint{4.011648in}{2.275314in}}%
\pgfpathlineto{\pgfqpoint{4.035122in}{2.275314in}}%
\pgfpathlineto{\pgfqpoint{4.058596in}{2.275314in}}%
\pgfpathlineto{\pgfqpoint{4.082071in}{2.275314in}}%
\pgfpathlineto{\pgfqpoint{4.105545in}{2.275314in}}%
\pgfpathlineto{\pgfqpoint{4.129019in}{2.275314in}}%
\pgfpathlineto{\pgfqpoint{4.152494in}{2.275314in}}%
\pgfpathlineto{\pgfqpoint{4.175968in}{2.275314in}}%
\pgfpathlineto{\pgfqpoint{4.199442in}{2.275314in}}%
\pgfpathlineto{\pgfqpoint{4.222917in}{2.275314in}}%
\pgfpathlineto{\pgfqpoint{4.246391in}{2.275314in}}%
\pgfpathlineto{\pgfqpoint{4.269866in}{2.275314in}}%
\pgfpathlineto{\pgfqpoint{4.293340in}{2.275314in}}%
\pgfpathlineto{\pgfqpoint{4.316814in}{2.275314in}}%
\pgfpathlineto{\pgfqpoint{4.340289in}{2.275314in}}%
\pgfusepath{stroke}%
\end{pgfscope}%
\begin{pgfscope}%
\pgfpathmoveto{\pgfqpoint{3.460000in}{2.052407in}}%
\pgfpathcurveto{\pgfqpoint{3.749602in}{2.052407in}}{\pgfqpoint{4.027381in}{2.109937in}}{\pgfqpoint{4.232161in}{2.212326in}}%
\pgfpathcurveto{\pgfqpoint{4.436940in}{2.314716in}}{\pgfqpoint{4.552000in}{2.453606in}}{\pgfqpoint{4.552000in}{2.598407in}}%
\pgfpathcurveto{\pgfqpoint{4.552000in}{2.743208in}}{\pgfqpoint{4.436940in}{2.882097in}}{\pgfqpoint{4.232161in}{2.984487in}}%
\pgfpathcurveto{\pgfqpoint{4.027381in}{3.086877in}}{\pgfqpoint{3.749602in}{3.144407in}}{\pgfqpoint{3.460000in}{3.144407in}}%
\pgfpathcurveto{\pgfqpoint{3.170398in}{3.144407in}}{\pgfqpoint{2.892619in}{3.086877in}}{\pgfqpoint{2.687839in}{2.984487in}}%
\pgfpathcurveto{\pgfqpoint{2.483060in}{2.882097in}}{\pgfqpoint{2.368000in}{2.743208in}}{\pgfqpoint{2.368000in}{2.598407in}}%
\pgfpathcurveto{\pgfqpoint{2.368000in}{2.453606in}}{\pgfqpoint{2.483060in}{2.314716in}}{\pgfqpoint{2.687839in}{2.212326in}}%
\pgfpathcurveto{\pgfqpoint{2.892619in}{2.109937in}}{\pgfqpoint{3.170398in}{2.052407in}}{\pgfqpoint{3.460000in}{2.052407in}}%
\pgfpathlineto{\pgfqpoint{3.460000in}{2.052407in}}%
\pgfpathclose%
\pgfusepath{clip}%
\pgfsetrectcap%
\pgfsetroundjoin%
\pgfsetlinewidth{0.451687pt}%
\definecolor{currentstroke}{rgb}{0.690196,0.690196,0.690196}%
\pgfsetstrokecolor{currentstroke}%
\pgfsetstrokeopacity{0.250000}%
\pgfsetdash{}{0pt}%
\pgfpathmoveto{\pgfqpoint{2.461062in}{2.377847in}}%
\pgfpathlineto{\pgfqpoint{2.487701in}{2.377847in}}%
\pgfpathlineto{\pgfqpoint{2.514339in}{2.377847in}}%
\pgfpathlineto{\pgfqpoint{2.540977in}{2.377847in}}%
\pgfpathlineto{\pgfqpoint{2.567616in}{2.377847in}}%
\pgfpathlineto{\pgfqpoint{2.594254in}{2.377847in}}%
\pgfpathlineto{\pgfqpoint{2.620892in}{2.377847in}}%
\pgfpathlineto{\pgfqpoint{2.647531in}{2.377847in}}%
\pgfpathlineto{\pgfqpoint{2.674169in}{2.377847in}}%
\pgfpathlineto{\pgfqpoint{2.700807in}{2.377847in}}%
\pgfpathlineto{\pgfqpoint{2.727446in}{2.377847in}}%
\pgfpathlineto{\pgfqpoint{2.754084in}{2.377847in}}%
\pgfpathlineto{\pgfqpoint{2.780722in}{2.377847in}}%
\pgfpathlineto{\pgfqpoint{2.807361in}{2.377847in}}%
\pgfpathlineto{\pgfqpoint{2.833999in}{2.377847in}}%
\pgfpathlineto{\pgfqpoint{2.860637in}{2.377847in}}%
\pgfpathlineto{\pgfqpoint{2.887276in}{2.377847in}}%
\pgfpathlineto{\pgfqpoint{2.913914in}{2.377847in}}%
\pgfpathlineto{\pgfqpoint{2.940552in}{2.377847in}}%
\pgfpathlineto{\pgfqpoint{2.967191in}{2.377847in}}%
\pgfpathlineto{\pgfqpoint{2.993829in}{2.377847in}}%
\pgfpathlineto{\pgfqpoint{3.020467in}{2.377847in}}%
\pgfpathlineto{\pgfqpoint{3.047106in}{2.377847in}}%
\pgfpathlineto{\pgfqpoint{3.073744in}{2.377847in}}%
\pgfpathlineto{\pgfqpoint{3.100382in}{2.377847in}}%
\pgfpathlineto{\pgfqpoint{3.127021in}{2.377847in}}%
\pgfpathlineto{\pgfqpoint{3.153659in}{2.377847in}}%
\pgfpathlineto{\pgfqpoint{3.180297in}{2.377847in}}%
\pgfpathlineto{\pgfqpoint{3.206936in}{2.377847in}}%
\pgfpathlineto{\pgfqpoint{3.233574in}{2.377847in}}%
\pgfpathlineto{\pgfqpoint{3.260212in}{2.377847in}}%
\pgfpathlineto{\pgfqpoint{3.286851in}{2.377847in}}%
\pgfpathlineto{\pgfqpoint{3.313489in}{2.377847in}}%
\pgfpathlineto{\pgfqpoint{3.340127in}{2.377847in}}%
\pgfpathlineto{\pgfqpoint{3.366766in}{2.377847in}}%
\pgfpathlineto{\pgfqpoint{3.393404in}{2.377847in}}%
\pgfpathlineto{\pgfqpoint{3.420042in}{2.377847in}}%
\pgfpathlineto{\pgfqpoint{3.446681in}{2.377847in}}%
\pgfpathlineto{\pgfqpoint{3.473319in}{2.377847in}}%
\pgfpathlineto{\pgfqpoint{3.499958in}{2.377847in}}%
\pgfpathlineto{\pgfqpoint{3.526596in}{2.377847in}}%
\pgfpathlineto{\pgfqpoint{3.553234in}{2.377847in}}%
\pgfpathlineto{\pgfqpoint{3.579873in}{2.377847in}}%
\pgfpathlineto{\pgfqpoint{3.606511in}{2.377847in}}%
\pgfpathlineto{\pgfqpoint{3.633149in}{2.377847in}}%
\pgfpathlineto{\pgfqpoint{3.659788in}{2.377847in}}%
\pgfpathlineto{\pgfqpoint{3.686426in}{2.377847in}}%
\pgfpathlineto{\pgfqpoint{3.713064in}{2.377847in}}%
\pgfpathlineto{\pgfqpoint{3.739703in}{2.377847in}}%
\pgfpathlineto{\pgfqpoint{3.766341in}{2.377847in}}%
\pgfpathlineto{\pgfqpoint{3.792979in}{2.377847in}}%
\pgfpathlineto{\pgfqpoint{3.819618in}{2.377847in}}%
\pgfpathlineto{\pgfqpoint{3.846256in}{2.377847in}}%
\pgfpathlineto{\pgfqpoint{3.872894in}{2.377847in}}%
\pgfpathlineto{\pgfqpoint{3.899533in}{2.377847in}}%
\pgfpathlineto{\pgfqpoint{3.926171in}{2.377847in}}%
\pgfpathlineto{\pgfqpoint{3.952809in}{2.377847in}}%
\pgfpathlineto{\pgfqpoint{3.979448in}{2.377847in}}%
\pgfpathlineto{\pgfqpoint{4.006086in}{2.377847in}}%
\pgfpathlineto{\pgfqpoint{4.032724in}{2.377847in}}%
\pgfpathlineto{\pgfqpoint{4.059363in}{2.377847in}}%
\pgfpathlineto{\pgfqpoint{4.086001in}{2.377847in}}%
\pgfpathlineto{\pgfqpoint{4.112639in}{2.377847in}}%
\pgfpathlineto{\pgfqpoint{4.139278in}{2.377847in}}%
\pgfpathlineto{\pgfqpoint{4.165916in}{2.377847in}}%
\pgfpathlineto{\pgfqpoint{4.192554in}{2.377847in}}%
\pgfpathlineto{\pgfqpoint{4.219193in}{2.377847in}}%
\pgfpathlineto{\pgfqpoint{4.245831in}{2.377847in}}%
\pgfpathlineto{\pgfqpoint{4.272469in}{2.377847in}}%
\pgfpathlineto{\pgfqpoint{4.299108in}{2.377847in}}%
\pgfpathlineto{\pgfqpoint{4.325746in}{2.377847in}}%
\pgfpathlineto{\pgfqpoint{4.352384in}{2.377847in}}%
\pgfpathlineto{\pgfqpoint{4.379023in}{2.377847in}}%
\pgfpathlineto{\pgfqpoint{4.405661in}{2.377847in}}%
\pgfpathlineto{\pgfqpoint{4.432299in}{2.377847in}}%
\pgfpathlineto{\pgfqpoint{4.458938in}{2.377847in}}%
\pgfusepath{stroke}%
\end{pgfscope}%
\begin{pgfscope}%
\pgfpathmoveto{\pgfqpoint{3.460000in}{2.052407in}}%
\pgfpathcurveto{\pgfqpoint{3.749602in}{2.052407in}}{\pgfqpoint{4.027381in}{2.109937in}}{\pgfqpoint{4.232161in}{2.212326in}}%
\pgfpathcurveto{\pgfqpoint{4.436940in}{2.314716in}}{\pgfqpoint{4.552000in}{2.453606in}}{\pgfqpoint{4.552000in}{2.598407in}}%
\pgfpathcurveto{\pgfqpoint{4.552000in}{2.743208in}}{\pgfqpoint{4.436940in}{2.882097in}}{\pgfqpoint{4.232161in}{2.984487in}}%
\pgfpathcurveto{\pgfqpoint{4.027381in}{3.086877in}}{\pgfqpoint{3.749602in}{3.144407in}}{\pgfqpoint{3.460000in}{3.144407in}}%
\pgfpathcurveto{\pgfqpoint{3.170398in}{3.144407in}}{\pgfqpoint{2.892619in}{3.086877in}}{\pgfqpoint{2.687839in}{2.984487in}}%
\pgfpathcurveto{\pgfqpoint{2.483060in}{2.882097in}}{\pgfqpoint{2.368000in}{2.743208in}}{\pgfqpoint{2.368000in}{2.598407in}}%
\pgfpathcurveto{\pgfqpoint{2.368000in}{2.453606in}}{\pgfqpoint{2.483060in}{2.314716in}}{\pgfqpoint{2.687839in}{2.212326in}}%
\pgfpathcurveto{\pgfqpoint{2.892619in}{2.109937in}}{\pgfqpoint{3.170398in}{2.052407in}}{\pgfqpoint{3.460000in}{2.052407in}}%
\pgfpathlineto{\pgfqpoint{3.460000in}{2.052407in}}%
\pgfpathclose%
\pgfusepath{clip}%
\pgfsetrectcap%
\pgfsetroundjoin%
\pgfsetlinewidth{0.451687pt}%
\definecolor{currentstroke}{rgb}{0.690196,0.690196,0.690196}%
\pgfsetstrokecolor{currentstroke}%
\pgfsetstrokeopacity{0.250000}%
\pgfsetdash{}{0pt}%
\pgfpathmoveto{\pgfqpoint{2.391127in}{2.486632in}}%
\pgfpathlineto{\pgfqpoint{2.419630in}{2.486632in}}%
\pgfpathlineto{\pgfqpoint{2.448133in}{2.486632in}}%
\pgfpathlineto{\pgfqpoint{2.476637in}{2.486632in}}%
\pgfpathlineto{\pgfqpoint{2.505140in}{2.486632in}}%
\pgfpathlineto{\pgfqpoint{2.533643in}{2.486632in}}%
\pgfpathlineto{\pgfqpoint{2.562146in}{2.486632in}}%
\pgfpathlineto{\pgfqpoint{2.590650in}{2.486632in}}%
\pgfpathlineto{\pgfqpoint{2.619153in}{2.486632in}}%
\pgfpathlineto{\pgfqpoint{2.647656in}{2.486632in}}%
\pgfpathlineto{\pgfqpoint{2.676160in}{2.486632in}}%
\pgfpathlineto{\pgfqpoint{2.704663in}{2.486632in}}%
\pgfpathlineto{\pgfqpoint{2.733166in}{2.486632in}}%
\pgfpathlineto{\pgfqpoint{2.761669in}{2.486632in}}%
\pgfpathlineto{\pgfqpoint{2.790173in}{2.486632in}}%
\pgfpathlineto{\pgfqpoint{2.818676in}{2.486632in}}%
\pgfpathlineto{\pgfqpoint{2.847179in}{2.486632in}}%
\pgfpathlineto{\pgfqpoint{2.875683in}{2.486632in}}%
\pgfpathlineto{\pgfqpoint{2.904186in}{2.486632in}}%
\pgfpathlineto{\pgfqpoint{2.932689in}{2.486632in}}%
\pgfpathlineto{\pgfqpoint{2.961192in}{2.486632in}}%
\pgfpathlineto{\pgfqpoint{2.989696in}{2.486632in}}%
\pgfpathlineto{\pgfqpoint{3.018199in}{2.486632in}}%
\pgfpathlineto{\pgfqpoint{3.046702in}{2.486632in}}%
\pgfpathlineto{\pgfqpoint{3.075206in}{2.486632in}}%
\pgfpathlineto{\pgfqpoint{3.103709in}{2.486632in}}%
\pgfpathlineto{\pgfqpoint{3.132212in}{2.486632in}}%
\pgfpathlineto{\pgfqpoint{3.160715in}{2.486632in}}%
\pgfpathlineto{\pgfqpoint{3.189219in}{2.486632in}}%
\pgfpathlineto{\pgfqpoint{3.217722in}{2.486632in}}%
\pgfpathlineto{\pgfqpoint{3.246225in}{2.486632in}}%
\pgfpathlineto{\pgfqpoint{3.274729in}{2.486632in}}%
\pgfpathlineto{\pgfqpoint{3.303232in}{2.486632in}}%
\pgfpathlineto{\pgfqpoint{3.331735in}{2.486632in}}%
\pgfpathlineto{\pgfqpoint{3.360238in}{2.486632in}}%
\pgfpathlineto{\pgfqpoint{3.388742in}{2.486632in}}%
\pgfpathlineto{\pgfqpoint{3.417245in}{2.486632in}}%
\pgfpathlineto{\pgfqpoint{3.445748in}{2.486632in}}%
\pgfpathlineto{\pgfqpoint{3.474252in}{2.486632in}}%
\pgfpathlineto{\pgfqpoint{3.502755in}{2.486632in}}%
\pgfpathlineto{\pgfqpoint{3.531258in}{2.486632in}}%
\pgfpathlineto{\pgfqpoint{3.559762in}{2.486632in}}%
\pgfpathlineto{\pgfqpoint{3.588265in}{2.486632in}}%
\pgfpathlineto{\pgfqpoint{3.616768in}{2.486632in}}%
\pgfpathlineto{\pgfqpoint{3.645271in}{2.486632in}}%
\pgfpathlineto{\pgfqpoint{3.673775in}{2.486632in}}%
\pgfpathlineto{\pgfqpoint{3.702278in}{2.486632in}}%
\pgfpathlineto{\pgfqpoint{3.730781in}{2.486632in}}%
\pgfpathlineto{\pgfqpoint{3.759285in}{2.486632in}}%
\pgfpathlineto{\pgfqpoint{3.787788in}{2.486632in}}%
\pgfpathlineto{\pgfqpoint{3.816291in}{2.486632in}}%
\pgfpathlineto{\pgfqpoint{3.844794in}{2.486632in}}%
\pgfpathlineto{\pgfqpoint{3.873298in}{2.486632in}}%
\pgfpathlineto{\pgfqpoint{3.901801in}{2.486632in}}%
\pgfpathlineto{\pgfqpoint{3.930304in}{2.486632in}}%
\pgfpathlineto{\pgfqpoint{3.958808in}{2.486632in}}%
\pgfpathlineto{\pgfqpoint{3.987311in}{2.486632in}}%
\pgfpathlineto{\pgfqpoint{4.015814in}{2.486632in}}%
\pgfpathlineto{\pgfqpoint{4.044317in}{2.486632in}}%
\pgfpathlineto{\pgfqpoint{4.072821in}{2.486632in}}%
\pgfpathlineto{\pgfqpoint{4.101324in}{2.486632in}}%
\pgfpathlineto{\pgfqpoint{4.129827in}{2.486632in}}%
\pgfpathlineto{\pgfqpoint{4.158331in}{2.486632in}}%
\pgfpathlineto{\pgfqpoint{4.186834in}{2.486632in}}%
\pgfpathlineto{\pgfqpoint{4.215337in}{2.486632in}}%
\pgfpathlineto{\pgfqpoint{4.243840in}{2.486632in}}%
\pgfpathlineto{\pgfqpoint{4.272344in}{2.486632in}}%
\pgfpathlineto{\pgfqpoint{4.300847in}{2.486632in}}%
\pgfpathlineto{\pgfqpoint{4.329350in}{2.486632in}}%
\pgfpathlineto{\pgfqpoint{4.357854in}{2.486632in}}%
\pgfpathlineto{\pgfqpoint{4.386357in}{2.486632in}}%
\pgfpathlineto{\pgfqpoint{4.414860in}{2.486632in}}%
\pgfpathlineto{\pgfqpoint{4.443363in}{2.486632in}}%
\pgfpathlineto{\pgfqpoint{4.471867in}{2.486632in}}%
\pgfpathlineto{\pgfqpoint{4.500370in}{2.486632in}}%
\pgfpathlineto{\pgfqpoint{4.528873in}{2.486632in}}%
\pgfusepath{stroke}%
\end{pgfscope}%
\begin{pgfscope}%
\pgfpathmoveto{\pgfqpoint{3.460000in}{2.052407in}}%
\pgfpathcurveto{\pgfqpoint{3.749602in}{2.052407in}}{\pgfqpoint{4.027381in}{2.109937in}}{\pgfqpoint{4.232161in}{2.212326in}}%
\pgfpathcurveto{\pgfqpoint{4.436940in}{2.314716in}}{\pgfqpoint{4.552000in}{2.453606in}}{\pgfqpoint{4.552000in}{2.598407in}}%
\pgfpathcurveto{\pgfqpoint{4.552000in}{2.743208in}}{\pgfqpoint{4.436940in}{2.882097in}}{\pgfqpoint{4.232161in}{2.984487in}}%
\pgfpathcurveto{\pgfqpoint{4.027381in}{3.086877in}}{\pgfqpoint{3.749602in}{3.144407in}}{\pgfqpoint{3.460000in}{3.144407in}}%
\pgfpathcurveto{\pgfqpoint{3.170398in}{3.144407in}}{\pgfqpoint{2.892619in}{3.086877in}}{\pgfqpoint{2.687839in}{2.984487in}}%
\pgfpathcurveto{\pgfqpoint{2.483060in}{2.882097in}}{\pgfqpoint{2.368000in}{2.743208in}}{\pgfqpoint{2.368000in}{2.598407in}}%
\pgfpathcurveto{\pgfqpoint{2.368000in}{2.453606in}}{\pgfqpoint{2.483060in}{2.314716in}}{\pgfqpoint{2.687839in}{2.212326in}}%
\pgfpathcurveto{\pgfqpoint{2.892619in}{2.109937in}}{\pgfqpoint{3.170398in}{2.052407in}}{\pgfqpoint{3.460000in}{2.052407in}}%
\pgfpathlineto{\pgfqpoint{3.460000in}{2.052407in}}%
\pgfpathclose%
\pgfusepath{clip}%
\pgfsetrectcap%
\pgfsetroundjoin%
\pgfsetlinewidth{0.451687pt}%
\definecolor{currentstroke}{rgb}{0.690196,0.690196,0.690196}%
\pgfsetstrokecolor{currentstroke}%
\pgfsetstrokeopacity{0.250000}%
\pgfsetdash{}{0pt}%
\pgfpathmoveto{\pgfqpoint{2.368000in}{2.598407in}}%
\pgfpathlineto{\pgfqpoint{2.397120in}{2.598407in}}%
\pgfpathlineto{\pgfqpoint{2.426240in}{2.598407in}}%
\pgfpathlineto{\pgfqpoint{2.455360in}{2.598407in}}%
\pgfpathlineto{\pgfqpoint{2.484480in}{2.598407in}}%
\pgfpathlineto{\pgfqpoint{2.513600in}{2.598407in}}%
\pgfpathlineto{\pgfqpoint{2.542720in}{2.598407in}}%
\pgfpathlineto{\pgfqpoint{2.571840in}{2.598407in}}%
\pgfpathlineto{\pgfqpoint{2.600960in}{2.598407in}}%
\pgfpathlineto{\pgfqpoint{2.630080in}{2.598407in}}%
\pgfpathlineto{\pgfqpoint{2.659200in}{2.598407in}}%
\pgfpathlineto{\pgfqpoint{2.688320in}{2.598407in}}%
\pgfpathlineto{\pgfqpoint{2.717440in}{2.598407in}}%
\pgfpathlineto{\pgfqpoint{2.746560in}{2.598407in}}%
\pgfpathlineto{\pgfqpoint{2.775680in}{2.598407in}}%
\pgfpathlineto{\pgfqpoint{2.804800in}{2.598407in}}%
\pgfpathlineto{\pgfqpoint{2.833920in}{2.598407in}}%
\pgfpathlineto{\pgfqpoint{2.863040in}{2.598407in}}%
\pgfpathlineto{\pgfqpoint{2.892160in}{2.598407in}}%
\pgfpathlineto{\pgfqpoint{2.921280in}{2.598407in}}%
\pgfpathlineto{\pgfqpoint{2.950400in}{2.598407in}}%
\pgfpathlineto{\pgfqpoint{2.979520in}{2.598407in}}%
\pgfpathlineto{\pgfqpoint{3.008640in}{2.598407in}}%
\pgfpathlineto{\pgfqpoint{3.037760in}{2.598407in}}%
\pgfpathlineto{\pgfqpoint{3.066880in}{2.598407in}}%
\pgfpathlineto{\pgfqpoint{3.096000in}{2.598407in}}%
\pgfpathlineto{\pgfqpoint{3.125120in}{2.598407in}}%
\pgfpathlineto{\pgfqpoint{3.154240in}{2.598407in}}%
\pgfpathlineto{\pgfqpoint{3.183360in}{2.598407in}}%
\pgfpathlineto{\pgfqpoint{3.212480in}{2.598407in}}%
\pgfpathlineto{\pgfqpoint{3.241600in}{2.598407in}}%
\pgfpathlineto{\pgfqpoint{3.270720in}{2.598407in}}%
\pgfpathlineto{\pgfqpoint{3.299840in}{2.598407in}}%
\pgfpathlineto{\pgfqpoint{3.328960in}{2.598407in}}%
\pgfpathlineto{\pgfqpoint{3.358080in}{2.598407in}}%
\pgfpathlineto{\pgfqpoint{3.387200in}{2.598407in}}%
\pgfpathlineto{\pgfqpoint{3.416320in}{2.598407in}}%
\pgfpathlineto{\pgfqpoint{3.445440in}{2.598407in}}%
\pgfpathlineto{\pgfqpoint{3.474560in}{2.598407in}}%
\pgfpathlineto{\pgfqpoint{3.503680in}{2.598407in}}%
\pgfpathlineto{\pgfqpoint{3.532800in}{2.598407in}}%
\pgfpathlineto{\pgfqpoint{3.561920in}{2.598407in}}%
\pgfpathlineto{\pgfqpoint{3.591040in}{2.598407in}}%
\pgfpathlineto{\pgfqpoint{3.620160in}{2.598407in}}%
\pgfpathlineto{\pgfqpoint{3.649280in}{2.598407in}}%
\pgfpathlineto{\pgfqpoint{3.678400in}{2.598407in}}%
\pgfpathlineto{\pgfqpoint{3.707520in}{2.598407in}}%
\pgfpathlineto{\pgfqpoint{3.736640in}{2.598407in}}%
\pgfpathlineto{\pgfqpoint{3.765760in}{2.598407in}}%
\pgfpathlineto{\pgfqpoint{3.794880in}{2.598407in}}%
\pgfpathlineto{\pgfqpoint{3.824000in}{2.598407in}}%
\pgfpathlineto{\pgfqpoint{3.853120in}{2.598407in}}%
\pgfpathlineto{\pgfqpoint{3.882240in}{2.598407in}}%
\pgfpathlineto{\pgfqpoint{3.911360in}{2.598407in}}%
\pgfpathlineto{\pgfqpoint{3.940480in}{2.598407in}}%
\pgfpathlineto{\pgfqpoint{3.969600in}{2.598407in}}%
\pgfpathlineto{\pgfqpoint{3.998720in}{2.598407in}}%
\pgfpathlineto{\pgfqpoint{4.027840in}{2.598407in}}%
\pgfpathlineto{\pgfqpoint{4.056960in}{2.598407in}}%
\pgfpathlineto{\pgfqpoint{4.086080in}{2.598407in}}%
\pgfpathlineto{\pgfqpoint{4.115200in}{2.598407in}}%
\pgfpathlineto{\pgfqpoint{4.144320in}{2.598407in}}%
\pgfpathlineto{\pgfqpoint{4.173440in}{2.598407in}}%
\pgfpathlineto{\pgfqpoint{4.202560in}{2.598407in}}%
\pgfpathlineto{\pgfqpoint{4.231680in}{2.598407in}}%
\pgfpathlineto{\pgfqpoint{4.260800in}{2.598407in}}%
\pgfpathlineto{\pgfqpoint{4.289920in}{2.598407in}}%
\pgfpathlineto{\pgfqpoint{4.319040in}{2.598407in}}%
\pgfpathlineto{\pgfqpoint{4.348160in}{2.598407in}}%
\pgfpathlineto{\pgfqpoint{4.377280in}{2.598407in}}%
\pgfpathlineto{\pgfqpoint{4.406400in}{2.598407in}}%
\pgfpathlineto{\pgfqpoint{4.435520in}{2.598407in}}%
\pgfpathlineto{\pgfqpoint{4.464640in}{2.598407in}}%
\pgfpathlineto{\pgfqpoint{4.493760in}{2.598407in}}%
\pgfpathlineto{\pgfqpoint{4.522880in}{2.598407in}}%
\pgfpathlineto{\pgfqpoint{4.552000in}{2.598407in}}%
\pgfusepath{stroke}%
\end{pgfscope}%
\begin{pgfscope}%
\pgfpathmoveto{\pgfqpoint{3.460000in}{2.052407in}}%
\pgfpathcurveto{\pgfqpoint{3.749602in}{2.052407in}}{\pgfqpoint{4.027381in}{2.109937in}}{\pgfqpoint{4.232161in}{2.212326in}}%
\pgfpathcurveto{\pgfqpoint{4.436940in}{2.314716in}}{\pgfqpoint{4.552000in}{2.453606in}}{\pgfqpoint{4.552000in}{2.598407in}}%
\pgfpathcurveto{\pgfqpoint{4.552000in}{2.743208in}}{\pgfqpoint{4.436940in}{2.882097in}}{\pgfqpoint{4.232161in}{2.984487in}}%
\pgfpathcurveto{\pgfqpoint{4.027381in}{3.086877in}}{\pgfqpoint{3.749602in}{3.144407in}}{\pgfqpoint{3.460000in}{3.144407in}}%
\pgfpathcurveto{\pgfqpoint{3.170398in}{3.144407in}}{\pgfqpoint{2.892619in}{3.086877in}}{\pgfqpoint{2.687839in}{2.984487in}}%
\pgfpathcurveto{\pgfqpoint{2.483060in}{2.882097in}}{\pgfqpoint{2.368000in}{2.743208in}}{\pgfqpoint{2.368000in}{2.598407in}}%
\pgfpathcurveto{\pgfqpoint{2.368000in}{2.453606in}}{\pgfqpoint{2.483060in}{2.314716in}}{\pgfqpoint{2.687839in}{2.212326in}}%
\pgfpathcurveto{\pgfqpoint{2.892619in}{2.109937in}}{\pgfqpoint{3.170398in}{2.052407in}}{\pgfqpoint{3.460000in}{2.052407in}}%
\pgfpathlineto{\pgfqpoint{3.460000in}{2.052407in}}%
\pgfpathclose%
\pgfusepath{clip}%
\pgfsetrectcap%
\pgfsetroundjoin%
\pgfsetlinewidth{0.451687pt}%
\definecolor{currentstroke}{rgb}{0.690196,0.690196,0.690196}%
\pgfsetstrokecolor{currentstroke}%
\pgfsetstrokeopacity{0.250000}%
\pgfsetdash{}{0pt}%
\pgfpathmoveto{\pgfqpoint{2.391127in}{2.710181in}}%
\pgfpathlineto{\pgfqpoint{2.419630in}{2.710181in}}%
\pgfpathlineto{\pgfqpoint{2.448133in}{2.710181in}}%
\pgfpathlineto{\pgfqpoint{2.476637in}{2.710181in}}%
\pgfpathlineto{\pgfqpoint{2.505140in}{2.710181in}}%
\pgfpathlineto{\pgfqpoint{2.533643in}{2.710181in}}%
\pgfpathlineto{\pgfqpoint{2.562146in}{2.710181in}}%
\pgfpathlineto{\pgfqpoint{2.590650in}{2.710181in}}%
\pgfpathlineto{\pgfqpoint{2.619153in}{2.710181in}}%
\pgfpathlineto{\pgfqpoint{2.647656in}{2.710181in}}%
\pgfpathlineto{\pgfqpoint{2.676160in}{2.710181in}}%
\pgfpathlineto{\pgfqpoint{2.704663in}{2.710181in}}%
\pgfpathlineto{\pgfqpoint{2.733166in}{2.710181in}}%
\pgfpathlineto{\pgfqpoint{2.761669in}{2.710181in}}%
\pgfpathlineto{\pgfqpoint{2.790173in}{2.710181in}}%
\pgfpathlineto{\pgfqpoint{2.818676in}{2.710181in}}%
\pgfpathlineto{\pgfqpoint{2.847179in}{2.710181in}}%
\pgfpathlineto{\pgfqpoint{2.875683in}{2.710181in}}%
\pgfpathlineto{\pgfqpoint{2.904186in}{2.710181in}}%
\pgfpathlineto{\pgfqpoint{2.932689in}{2.710181in}}%
\pgfpathlineto{\pgfqpoint{2.961192in}{2.710181in}}%
\pgfpathlineto{\pgfqpoint{2.989696in}{2.710181in}}%
\pgfpathlineto{\pgfqpoint{3.018199in}{2.710181in}}%
\pgfpathlineto{\pgfqpoint{3.046702in}{2.710181in}}%
\pgfpathlineto{\pgfqpoint{3.075206in}{2.710181in}}%
\pgfpathlineto{\pgfqpoint{3.103709in}{2.710181in}}%
\pgfpathlineto{\pgfqpoint{3.132212in}{2.710181in}}%
\pgfpathlineto{\pgfqpoint{3.160715in}{2.710181in}}%
\pgfpathlineto{\pgfqpoint{3.189219in}{2.710181in}}%
\pgfpathlineto{\pgfqpoint{3.217722in}{2.710181in}}%
\pgfpathlineto{\pgfqpoint{3.246225in}{2.710181in}}%
\pgfpathlineto{\pgfqpoint{3.274729in}{2.710181in}}%
\pgfpathlineto{\pgfqpoint{3.303232in}{2.710181in}}%
\pgfpathlineto{\pgfqpoint{3.331735in}{2.710181in}}%
\pgfpathlineto{\pgfqpoint{3.360238in}{2.710181in}}%
\pgfpathlineto{\pgfqpoint{3.388742in}{2.710181in}}%
\pgfpathlineto{\pgfqpoint{3.417245in}{2.710181in}}%
\pgfpathlineto{\pgfqpoint{3.445748in}{2.710181in}}%
\pgfpathlineto{\pgfqpoint{3.474252in}{2.710181in}}%
\pgfpathlineto{\pgfqpoint{3.502755in}{2.710181in}}%
\pgfpathlineto{\pgfqpoint{3.531258in}{2.710181in}}%
\pgfpathlineto{\pgfqpoint{3.559762in}{2.710181in}}%
\pgfpathlineto{\pgfqpoint{3.588265in}{2.710181in}}%
\pgfpathlineto{\pgfqpoint{3.616768in}{2.710181in}}%
\pgfpathlineto{\pgfqpoint{3.645271in}{2.710181in}}%
\pgfpathlineto{\pgfqpoint{3.673775in}{2.710181in}}%
\pgfpathlineto{\pgfqpoint{3.702278in}{2.710181in}}%
\pgfpathlineto{\pgfqpoint{3.730781in}{2.710181in}}%
\pgfpathlineto{\pgfqpoint{3.759285in}{2.710181in}}%
\pgfpathlineto{\pgfqpoint{3.787788in}{2.710181in}}%
\pgfpathlineto{\pgfqpoint{3.816291in}{2.710181in}}%
\pgfpathlineto{\pgfqpoint{3.844794in}{2.710181in}}%
\pgfpathlineto{\pgfqpoint{3.873298in}{2.710181in}}%
\pgfpathlineto{\pgfqpoint{3.901801in}{2.710181in}}%
\pgfpathlineto{\pgfqpoint{3.930304in}{2.710181in}}%
\pgfpathlineto{\pgfqpoint{3.958808in}{2.710181in}}%
\pgfpathlineto{\pgfqpoint{3.987311in}{2.710181in}}%
\pgfpathlineto{\pgfqpoint{4.015814in}{2.710181in}}%
\pgfpathlineto{\pgfqpoint{4.044317in}{2.710181in}}%
\pgfpathlineto{\pgfqpoint{4.072821in}{2.710181in}}%
\pgfpathlineto{\pgfqpoint{4.101324in}{2.710181in}}%
\pgfpathlineto{\pgfqpoint{4.129827in}{2.710181in}}%
\pgfpathlineto{\pgfqpoint{4.158331in}{2.710181in}}%
\pgfpathlineto{\pgfqpoint{4.186834in}{2.710181in}}%
\pgfpathlineto{\pgfqpoint{4.215337in}{2.710181in}}%
\pgfpathlineto{\pgfqpoint{4.243840in}{2.710181in}}%
\pgfpathlineto{\pgfqpoint{4.272344in}{2.710181in}}%
\pgfpathlineto{\pgfqpoint{4.300847in}{2.710181in}}%
\pgfpathlineto{\pgfqpoint{4.329350in}{2.710181in}}%
\pgfpathlineto{\pgfqpoint{4.357854in}{2.710181in}}%
\pgfpathlineto{\pgfqpoint{4.386357in}{2.710181in}}%
\pgfpathlineto{\pgfqpoint{4.414860in}{2.710181in}}%
\pgfpathlineto{\pgfqpoint{4.443363in}{2.710181in}}%
\pgfpathlineto{\pgfqpoint{4.471867in}{2.710181in}}%
\pgfpathlineto{\pgfqpoint{4.500370in}{2.710181in}}%
\pgfpathlineto{\pgfqpoint{4.528873in}{2.710181in}}%
\pgfusepath{stroke}%
\end{pgfscope}%
\begin{pgfscope}%
\pgfpathmoveto{\pgfqpoint{3.460000in}{2.052407in}}%
\pgfpathcurveto{\pgfqpoint{3.749602in}{2.052407in}}{\pgfqpoint{4.027381in}{2.109937in}}{\pgfqpoint{4.232161in}{2.212326in}}%
\pgfpathcurveto{\pgfqpoint{4.436940in}{2.314716in}}{\pgfqpoint{4.552000in}{2.453606in}}{\pgfqpoint{4.552000in}{2.598407in}}%
\pgfpathcurveto{\pgfqpoint{4.552000in}{2.743208in}}{\pgfqpoint{4.436940in}{2.882097in}}{\pgfqpoint{4.232161in}{2.984487in}}%
\pgfpathcurveto{\pgfqpoint{4.027381in}{3.086877in}}{\pgfqpoint{3.749602in}{3.144407in}}{\pgfqpoint{3.460000in}{3.144407in}}%
\pgfpathcurveto{\pgfqpoint{3.170398in}{3.144407in}}{\pgfqpoint{2.892619in}{3.086877in}}{\pgfqpoint{2.687839in}{2.984487in}}%
\pgfpathcurveto{\pgfqpoint{2.483060in}{2.882097in}}{\pgfqpoint{2.368000in}{2.743208in}}{\pgfqpoint{2.368000in}{2.598407in}}%
\pgfpathcurveto{\pgfqpoint{2.368000in}{2.453606in}}{\pgfqpoint{2.483060in}{2.314716in}}{\pgfqpoint{2.687839in}{2.212326in}}%
\pgfpathcurveto{\pgfqpoint{2.892619in}{2.109937in}}{\pgfqpoint{3.170398in}{2.052407in}}{\pgfqpoint{3.460000in}{2.052407in}}%
\pgfpathlineto{\pgfqpoint{3.460000in}{2.052407in}}%
\pgfpathclose%
\pgfusepath{clip}%
\pgfsetrectcap%
\pgfsetroundjoin%
\pgfsetlinewidth{0.451687pt}%
\definecolor{currentstroke}{rgb}{0.690196,0.690196,0.690196}%
\pgfsetstrokecolor{currentstroke}%
\pgfsetstrokeopacity{0.250000}%
\pgfsetdash{}{0pt}%
\pgfpathmoveto{\pgfqpoint{2.461062in}{2.818967in}}%
\pgfpathlineto{\pgfqpoint{2.487701in}{2.818967in}}%
\pgfpathlineto{\pgfqpoint{2.514339in}{2.818967in}}%
\pgfpathlineto{\pgfqpoint{2.540977in}{2.818967in}}%
\pgfpathlineto{\pgfqpoint{2.567616in}{2.818967in}}%
\pgfpathlineto{\pgfqpoint{2.594254in}{2.818967in}}%
\pgfpathlineto{\pgfqpoint{2.620892in}{2.818967in}}%
\pgfpathlineto{\pgfqpoint{2.647531in}{2.818967in}}%
\pgfpathlineto{\pgfqpoint{2.674169in}{2.818967in}}%
\pgfpathlineto{\pgfqpoint{2.700807in}{2.818967in}}%
\pgfpathlineto{\pgfqpoint{2.727446in}{2.818967in}}%
\pgfpathlineto{\pgfqpoint{2.754084in}{2.818967in}}%
\pgfpathlineto{\pgfqpoint{2.780722in}{2.818967in}}%
\pgfpathlineto{\pgfqpoint{2.807361in}{2.818967in}}%
\pgfpathlineto{\pgfqpoint{2.833999in}{2.818967in}}%
\pgfpathlineto{\pgfqpoint{2.860637in}{2.818967in}}%
\pgfpathlineto{\pgfqpoint{2.887276in}{2.818967in}}%
\pgfpathlineto{\pgfqpoint{2.913914in}{2.818967in}}%
\pgfpathlineto{\pgfqpoint{2.940552in}{2.818967in}}%
\pgfpathlineto{\pgfqpoint{2.967191in}{2.818967in}}%
\pgfpathlineto{\pgfqpoint{2.993829in}{2.818967in}}%
\pgfpathlineto{\pgfqpoint{3.020467in}{2.818967in}}%
\pgfpathlineto{\pgfqpoint{3.047106in}{2.818967in}}%
\pgfpathlineto{\pgfqpoint{3.073744in}{2.818967in}}%
\pgfpathlineto{\pgfqpoint{3.100382in}{2.818967in}}%
\pgfpathlineto{\pgfqpoint{3.127021in}{2.818967in}}%
\pgfpathlineto{\pgfqpoint{3.153659in}{2.818967in}}%
\pgfpathlineto{\pgfqpoint{3.180297in}{2.818967in}}%
\pgfpathlineto{\pgfqpoint{3.206936in}{2.818967in}}%
\pgfpathlineto{\pgfqpoint{3.233574in}{2.818967in}}%
\pgfpathlineto{\pgfqpoint{3.260212in}{2.818967in}}%
\pgfpathlineto{\pgfqpoint{3.286851in}{2.818967in}}%
\pgfpathlineto{\pgfqpoint{3.313489in}{2.818967in}}%
\pgfpathlineto{\pgfqpoint{3.340127in}{2.818967in}}%
\pgfpathlineto{\pgfqpoint{3.366766in}{2.818967in}}%
\pgfpathlineto{\pgfqpoint{3.393404in}{2.818967in}}%
\pgfpathlineto{\pgfqpoint{3.420042in}{2.818967in}}%
\pgfpathlineto{\pgfqpoint{3.446681in}{2.818967in}}%
\pgfpathlineto{\pgfqpoint{3.473319in}{2.818967in}}%
\pgfpathlineto{\pgfqpoint{3.499958in}{2.818967in}}%
\pgfpathlineto{\pgfqpoint{3.526596in}{2.818967in}}%
\pgfpathlineto{\pgfqpoint{3.553234in}{2.818967in}}%
\pgfpathlineto{\pgfqpoint{3.579873in}{2.818967in}}%
\pgfpathlineto{\pgfqpoint{3.606511in}{2.818967in}}%
\pgfpathlineto{\pgfqpoint{3.633149in}{2.818967in}}%
\pgfpathlineto{\pgfqpoint{3.659788in}{2.818967in}}%
\pgfpathlineto{\pgfqpoint{3.686426in}{2.818967in}}%
\pgfpathlineto{\pgfqpoint{3.713064in}{2.818967in}}%
\pgfpathlineto{\pgfqpoint{3.739703in}{2.818967in}}%
\pgfpathlineto{\pgfqpoint{3.766341in}{2.818967in}}%
\pgfpathlineto{\pgfqpoint{3.792979in}{2.818967in}}%
\pgfpathlineto{\pgfqpoint{3.819618in}{2.818967in}}%
\pgfpathlineto{\pgfqpoint{3.846256in}{2.818967in}}%
\pgfpathlineto{\pgfqpoint{3.872894in}{2.818967in}}%
\pgfpathlineto{\pgfqpoint{3.899533in}{2.818967in}}%
\pgfpathlineto{\pgfqpoint{3.926171in}{2.818967in}}%
\pgfpathlineto{\pgfqpoint{3.952809in}{2.818967in}}%
\pgfpathlineto{\pgfqpoint{3.979448in}{2.818967in}}%
\pgfpathlineto{\pgfqpoint{4.006086in}{2.818967in}}%
\pgfpathlineto{\pgfqpoint{4.032724in}{2.818967in}}%
\pgfpathlineto{\pgfqpoint{4.059363in}{2.818967in}}%
\pgfpathlineto{\pgfqpoint{4.086001in}{2.818967in}}%
\pgfpathlineto{\pgfqpoint{4.112639in}{2.818967in}}%
\pgfpathlineto{\pgfqpoint{4.139278in}{2.818967in}}%
\pgfpathlineto{\pgfqpoint{4.165916in}{2.818967in}}%
\pgfpathlineto{\pgfqpoint{4.192554in}{2.818967in}}%
\pgfpathlineto{\pgfqpoint{4.219193in}{2.818967in}}%
\pgfpathlineto{\pgfqpoint{4.245831in}{2.818967in}}%
\pgfpathlineto{\pgfqpoint{4.272469in}{2.818967in}}%
\pgfpathlineto{\pgfqpoint{4.299108in}{2.818967in}}%
\pgfpathlineto{\pgfqpoint{4.325746in}{2.818967in}}%
\pgfpathlineto{\pgfqpoint{4.352384in}{2.818967in}}%
\pgfpathlineto{\pgfqpoint{4.379023in}{2.818967in}}%
\pgfpathlineto{\pgfqpoint{4.405661in}{2.818967in}}%
\pgfpathlineto{\pgfqpoint{4.432299in}{2.818967in}}%
\pgfpathlineto{\pgfqpoint{4.458938in}{2.818967in}}%
\pgfusepath{stroke}%
\end{pgfscope}%
\begin{pgfscope}%
\pgfpathmoveto{\pgfqpoint{3.460000in}{2.052407in}}%
\pgfpathcurveto{\pgfqpoint{3.749602in}{2.052407in}}{\pgfqpoint{4.027381in}{2.109937in}}{\pgfqpoint{4.232161in}{2.212326in}}%
\pgfpathcurveto{\pgfqpoint{4.436940in}{2.314716in}}{\pgfqpoint{4.552000in}{2.453606in}}{\pgfqpoint{4.552000in}{2.598407in}}%
\pgfpathcurveto{\pgfqpoint{4.552000in}{2.743208in}}{\pgfqpoint{4.436940in}{2.882097in}}{\pgfqpoint{4.232161in}{2.984487in}}%
\pgfpathcurveto{\pgfqpoint{4.027381in}{3.086877in}}{\pgfqpoint{3.749602in}{3.144407in}}{\pgfqpoint{3.460000in}{3.144407in}}%
\pgfpathcurveto{\pgfqpoint{3.170398in}{3.144407in}}{\pgfqpoint{2.892619in}{3.086877in}}{\pgfqpoint{2.687839in}{2.984487in}}%
\pgfpathcurveto{\pgfqpoint{2.483060in}{2.882097in}}{\pgfqpoint{2.368000in}{2.743208in}}{\pgfqpoint{2.368000in}{2.598407in}}%
\pgfpathcurveto{\pgfqpoint{2.368000in}{2.453606in}}{\pgfqpoint{2.483060in}{2.314716in}}{\pgfqpoint{2.687839in}{2.212326in}}%
\pgfpathcurveto{\pgfqpoint{2.892619in}{2.109937in}}{\pgfqpoint{3.170398in}{2.052407in}}{\pgfqpoint{3.460000in}{2.052407in}}%
\pgfpathlineto{\pgfqpoint{3.460000in}{2.052407in}}%
\pgfpathclose%
\pgfusepath{clip}%
\pgfsetrectcap%
\pgfsetroundjoin%
\pgfsetlinewidth{0.451687pt}%
\definecolor{currentstroke}{rgb}{0.690196,0.690196,0.690196}%
\pgfsetstrokecolor{currentstroke}%
\pgfsetstrokeopacity{0.250000}%
\pgfsetdash{}{0pt}%
\pgfpathmoveto{\pgfqpoint{2.579711in}{2.921500in}}%
\pgfpathlineto{\pgfqpoint{2.603186in}{2.921500in}}%
\pgfpathlineto{\pgfqpoint{2.626660in}{2.921500in}}%
\pgfpathlineto{\pgfqpoint{2.650134in}{2.921500in}}%
\pgfpathlineto{\pgfqpoint{2.673609in}{2.921500in}}%
\pgfpathlineto{\pgfqpoint{2.697083in}{2.921500in}}%
\pgfpathlineto{\pgfqpoint{2.720558in}{2.921500in}}%
\pgfpathlineto{\pgfqpoint{2.744032in}{2.921500in}}%
\pgfpathlineto{\pgfqpoint{2.767506in}{2.921500in}}%
\pgfpathlineto{\pgfqpoint{2.790981in}{2.921500in}}%
\pgfpathlineto{\pgfqpoint{2.814455in}{2.921500in}}%
\pgfpathlineto{\pgfqpoint{2.837929in}{2.921500in}}%
\pgfpathlineto{\pgfqpoint{2.861404in}{2.921500in}}%
\pgfpathlineto{\pgfqpoint{2.884878in}{2.921500in}}%
\pgfpathlineto{\pgfqpoint{2.908352in}{2.921500in}}%
\pgfpathlineto{\pgfqpoint{2.931827in}{2.921500in}}%
\pgfpathlineto{\pgfqpoint{2.955301in}{2.921500in}}%
\pgfpathlineto{\pgfqpoint{2.978776in}{2.921500in}}%
\pgfpathlineto{\pgfqpoint{3.002250in}{2.921500in}}%
\pgfpathlineto{\pgfqpoint{3.025724in}{2.921500in}}%
\pgfpathlineto{\pgfqpoint{3.049199in}{2.921500in}}%
\pgfpathlineto{\pgfqpoint{3.072673in}{2.921500in}}%
\pgfpathlineto{\pgfqpoint{3.096147in}{2.921500in}}%
\pgfpathlineto{\pgfqpoint{3.119622in}{2.921500in}}%
\pgfpathlineto{\pgfqpoint{3.143096in}{2.921500in}}%
\pgfpathlineto{\pgfqpoint{3.166570in}{2.921500in}}%
\pgfpathlineto{\pgfqpoint{3.190045in}{2.921500in}}%
\pgfpathlineto{\pgfqpoint{3.213519in}{2.921500in}}%
\pgfpathlineto{\pgfqpoint{3.236994in}{2.921500in}}%
\pgfpathlineto{\pgfqpoint{3.260468in}{2.921500in}}%
\pgfpathlineto{\pgfqpoint{3.283942in}{2.921500in}}%
\pgfpathlineto{\pgfqpoint{3.307417in}{2.921500in}}%
\pgfpathlineto{\pgfqpoint{3.330891in}{2.921500in}}%
\pgfpathlineto{\pgfqpoint{3.354365in}{2.921500in}}%
\pgfpathlineto{\pgfqpoint{3.377840in}{2.921500in}}%
\pgfpathlineto{\pgfqpoint{3.401314in}{2.921500in}}%
\pgfpathlineto{\pgfqpoint{3.424788in}{2.921500in}}%
\pgfpathlineto{\pgfqpoint{3.448263in}{2.921500in}}%
\pgfpathlineto{\pgfqpoint{3.471737in}{2.921500in}}%
\pgfpathlineto{\pgfqpoint{3.495212in}{2.921500in}}%
\pgfpathlineto{\pgfqpoint{3.518686in}{2.921500in}}%
\pgfpathlineto{\pgfqpoint{3.542160in}{2.921500in}}%
\pgfpathlineto{\pgfqpoint{3.565635in}{2.921500in}}%
\pgfpathlineto{\pgfqpoint{3.589109in}{2.921500in}}%
\pgfpathlineto{\pgfqpoint{3.612583in}{2.921500in}}%
\pgfpathlineto{\pgfqpoint{3.636058in}{2.921500in}}%
\pgfpathlineto{\pgfqpoint{3.659532in}{2.921500in}}%
\pgfpathlineto{\pgfqpoint{3.683006in}{2.921500in}}%
\pgfpathlineto{\pgfqpoint{3.706481in}{2.921500in}}%
\pgfpathlineto{\pgfqpoint{3.729955in}{2.921500in}}%
\pgfpathlineto{\pgfqpoint{3.753430in}{2.921500in}}%
\pgfpathlineto{\pgfqpoint{3.776904in}{2.921500in}}%
\pgfpathlineto{\pgfqpoint{3.800378in}{2.921500in}}%
\pgfpathlineto{\pgfqpoint{3.823853in}{2.921500in}}%
\pgfpathlineto{\pgfqpoint{3.847327in}{2.921500in}}%
\pgfpathlineto{\pgfqpoint{3.870801in}{2.921500in}}%
\pgfpathlineto{\pgfqpoint{3.894276in}{2.921500in}}%
\pgfpathlineto{\pgfqpoint{3.917750in}{2.921500in}}%
\pgfpathlineto{\pgfqpoint{3.941224in}{2.921500in}}%
\pgfpathlineto{\pgfqpoint{3.964699in}{2.921500in}}%
\pgfpathlineto{\pgfqpoint{3.988173in}{2.921500in}}%
\pgfpathlineto{\pgfqpoint{4.011648in}{2.921500in}}%
\pgfpathlineto{\pgfqpoint{4.035122in}{2.921500in}}%
\pgfpathlineto{\pgfqpoint{4.058596in}{2.921500in}}%
\pgfpathlineto{\pgfqpoint{4.082071in}{2.921500in}}%
\pgfpathlineto{\pgfqpoint{4.105545in}{2.921500in}}%
\pgfpathlineto{\pgfqpoint{4.129019in}{2.921500in}}%
\pgfpathlineto{\pgfqpoint{4.152494in}{2.921500in}}%
\pgfpathlineto{\pgfqpoint{4.175968in}{2.921500in}}%
\pgfpathlineto{\pgfqpoint{4.199442in}{2.921500in}}%
\pgfpathlineto{\pgfqpoint{4.222917in}{2.921500in}}%
\pgfpathlineto{\pgfqpoint{4.246391in}{2.921500in}}%
\pgfpathlineto{\pgfqpoint{4.269866in}{2.921500in}}%
\pgfpathlineto{\pgfqpoint{4.293340in}{2.921500in}}%
\pgfpathlineto{\pgfqpoint{4.316814in}{2.921500in}}%
\pgfpathlineto{\pgfqpoint{4.340289in}{2.921500in}}%
\pgfusepath{stroke}%
\end{pgfscope}%
\begin{pgfscope}%
\pgfpathmoveto{\pgfqpoint{3.460000in}{2.052407in}}%
\pgfpathcurveto{\pgfqpoint{3.749602in}{2.052407in}}{\pgfqpoint{4.027381in}{2.109937in}}{\pgfqpoint{4.232161in}{2.212326in}}%
\pgfpathcurveto{\pgfqpoint{4.436940in}{2.314716in}}{\pgfqpoint{4.552000in}{2.453606in}}{\pgfqpoint{4.552000in}{2.598407in}}%
\pgfpathcurveto{\pgfqpoint{4.552000in}{2.743208in}}{\pgfqpoint{4.436940in}{2.882097in}}{\pgfqpoint{4.232161in}{2.984487in}}%
\pgfpathcurveto{\pgfqpoint{4.027381in}{3.086877in}}{\pgfqpoint{3.749602in}{3.144407in}}{\pgfqpoint{3.460000in}{3.144407in}}%
\pgfpathcurveto{\pgfqpoint{3.170398in}{3.144407in}}{\pgfqpoint{2.892619in}{3.086877in}}{\pgfqpoint{2.687839in}{2.984487in}}%
\pgfpathcurveto{\pgfqpoint{2.483060in}{2.882097in}}{\pgfqpoint{2.368000in}{2.743208in}}{\pgfqpoint{2.368000in}{2.598407in}}%
\pgfpathcurveto{\pgfqpoint{2.368000in}{2.453606in}}{\pgfqpoint{2.483060in}{2.314716in}}{\pgfqpoint{2.687839in}{2.212326in}}%
\pgfpathcurveto{\pgfqpoint{2.892619in}{2.109937in}}{\pgfqpoint{3.170398in}{2.052407in}}{\pgfqpoint{3.460000in}{2.052407in}}%
\pgfpathlineto{\pgfqpoint{3.460000in}{2.052407in}}%
\pgfpathclose%
\pgfusepath{clip}%
\pgfsetrectcap%
\pgfsetroundjoin%
\pgfsetlinewidth{0.451687pt}%
\definecolor{currentstroke}{rgb}{0.690196,0.690196,0.690196}%
\pgfsetstrokecolor{currentstroke}%
\pgfsetstrokeopacity{0.250000}%
\pgfsetdash{}{0pt}%
\pgfpathmoveto{\pgfqpoint{2.753329in}{3.014664in}}%
\pgfpathlineto{\pgfqpoint{2.772174in}{3.014664in}}%
\pgfpathlineto{\pgfqpoint{2.791018in}{3.014664in}}%
\pgfpathlineto{\pgfqpoint{2.809863in}{3.014664in}}%
\pgfpathlineto{\pgfqpoint{2.828707in}{3.014664in}}%
\pgfpathlineto{\pgfqpoint{2.847552in}{3.014664in}}%
\pgfpathlineto{\pgfqpoint{2.866396in}{3.014664in}}%
\pgfpathlineto{\pgfqpoint{2.885241in}{3.014664in}}%
\pgfpathlineto{\pgfqpoint{2.904086in}{3.014664in}}%
\pgfpathlineto{\pgfqpoint{2.922930in}{3.014664in}}%
\pgfpathlineto{\pgfqpoint{2.941775in}{3.014664in}}%
\pgfpathlineto{\pgfqpoint{2.960619in}{3.014664in}}%
\pgfpathlineto{\pgfqpoint{2.979464in}{3.014664in}}%
\pgfpathlineto{\pgfqpoint{2.998308in}{3.014664in}}%
\pgfpathlineto{\pgfqpoint{3.017153in}{3.014664in}}%
\pgfpathlineto{\pgfqpoint{3.035997in}{3.014664in}}%
\pgfpathlineto{\pgfqpoint{3.054842in}{3.014664in}}%
\pgfpathlineto{\pgfqpoint{3.073687in}{3.014664in}}%
\pgfpathlineto{\pgfqpoint{3.092531in}{3.014664in}}%
\pgfpathlineto{\pgfqpoint{3.111376in}{3.014664in}}%
\pgfpathlineto{\pgfqpoint{3.130220in}{3.014664in}}%
\pgfpathlineto{\pgfqpoint{3.149065in}{3.014664in}}%
\pgfpathlineto{\pgfqpoint{3.167909in}{3.014664in}}%
\pgfpathlineto{\pgfqpoint{3.186754in}{3.014664in}}%
\pgfpathlineto{\pgfqpoint{3.205598in}{3.014664in}}%
\pgfpathlineto{\pgfqpoint{3.224443in}{3.014664in}}%
\pgfpathlineto{\pgfqpoint{3.243288in}{3.014664in}}%
\pgfpathlineto{\pgfqpoint{3.262132in}{3.014664in}}%
\pgfpathlineto{\pgfqpoint{3.280977in}{3.014664in}}%
\pgfpathlineto{\pgfqpoint{3.299821in}{3.014664in}}%
\pgfpathlineto{\pgfqpoint{3.318666in}{3.014664in}}%
\pgfpathlineto{\pgfqpoint{3.337510in}{3.014664in}}%
\pgfpathlineto{\pgfqpoint{3.356355in}{3.014664in}}%
\pgfpathlineto{\pgfqpoint{3.375199in}{3.014664in}}%
\pgfpathlineto{\pgfqpoint{3.394044in}{3.014664in}}%
\pgfpathlineto{\pgfqpoint{3.412889in}{3.014664in}}%
\pgfpathlineto{\pgfqpoint{3.431733in}{3.014664in}}%
\pgfpathlineto{\pgfqpoint{3.450578in}{3.014664in}}%
\pgfpathlineto{\pgfqpoint{3.469422in}{3.014664in}}%
\pgfpathlineto{\pgfqpoint{3.488267in}{3.014664in}}%
\pgfpathlineto{\pgfqpoint{3.507111in}{3.014664in}}%
\pgfpathlineto{\pgfqpoint{3.525956in}{3.014664in}}%
\pgfpathlineto{\pgfqpoint{3.544801in}{3.014664in}}%
\pgfpathlineto{\pgfqpoint{3.563645in}{3.014664in}}%
\pgfpathlineto{\pgfqpoint{3.582490in}{3.014664in}}%
\pgfpathlineto{\pgfqpoint{3.601334in}{3.014664in}}%
\pgfpathlineto{\pgfqpoint{3.620179in}{3.014664in}}%
\pgfpathlineto{\pgfqpoint{3.639023in}{3.014664in}}%
\pgfpathlineto{\pgfqpoint{3.657868in}{3.014664in}}%
\pgfpathlineto{\pgfqpoint{3.676712in}{3.014664in}}%
\pgfpathlineto{\pgfqpoint{3.695557in}{3.014664in}}%
\pgfpathlineto{\pgfqpoint{3.714402in}{3.014664in}}%
\pgfpathlineto{\pgfqpoint{3.733246in}{3.014664in}}%
\pgfpathlineto{\pgfqpoint{3.752091in}{3.014664in}}%
\pgfpathlineto{\pgfqpoint{3.770935in}{3.014664in}}%
\pgfpathlineto{\pgfqpoint{3.789780in}{3.014664in}}%
\pgfpathlineto{\pgfqpoint{3.808624in}{3.014664in}}%
\pgfpathlineto{\pgfqpoint{3.827469in}{3.014664in}}%
\pgfpathlineto{\pgfqpoint{3.846313in}{3.014664in}}%
\pgfpathlineto{\pgfqpoint{3.865158in}{3.014664in}}%
\pgfpathlineto{\pgfqpoint{3.884003in}{3.014664in}}%
\pgfpathlineto{\pgfqpoint{3.902847in}{3.014664in}}%
\pgfpathlineto{\pgfqpoint{3.921692in}{3.014664in}}%
\pgfpathlineto{\pgfqpoint{3.940536in}{3.014664in}}%
\pgfpathlineto{\pgfqpoint{3.959381in}{3.014664in}}%
\pgfpathlineto{\pgfqpoint{3.978225in}{3.014664in}}%
\pgfpathlineto{\pgfqpoint{3.997070in}{3.014664in}}%
\pgfpathlineto{\pgfqpoint{4.015914in}{3.014664in}}%
\pgfpathlineto{\pgfqpoint{4.034759in}{3.014664in}}%
\pgfpathlineto{\pgfqpoint{4.053604in}{3.014664in}}%
\pgfpathlineto{\pgfqpoint{4.072448in}{3.014664in}}%
\pgfpathlineto{\pgfqpoint{4.091293in}{3.014664in}}%
\pgfpathlineto{\pgfqpoint{4.110137in}{3.014664in}}%
\pgfpathlineto{\pgfqpoint{4.128982in}{3.014664in}}%
\pgfpathlineto{\pgfqpoint{4.147826in}{3.014664in}}%
\pgfpathlineto{\pgfqpoint{4.166671in}{3.014664in}}%
\pgfusepath{stroke}%
\end{pgfscope}%
\begin{pgfscope}%
\pgfpathmoveto{\pgfqpoint{3.460000in}{2.052407in}}%
\pgfpathcurveto{\pgfqpoint{3.749602in}{2.052407in}}{\pgfqpoint{4.027381in}{2.109937in}}{\pgfqpoint{4.232161in}{2.212326in}}%
\pgfpathcurveto{\pgfqpoint{4.436940in}{2.314716in}}{\pgfqpoint{4.552000in}{2.453606in}}{\pgfqpoint{4.552000in}{2.598407in}}%
\pgfpathcurveto{\pgfqpoint{4.552000in}{2.743208in}}{\pgfqpoint{4.436940in}{2.882097in}}{\pgfqpoint{4.232161in}{2.984487in}}%
\pgfpathcurveto{\pgfqpoint{4.027381in}{3.086877in}}{\pgfqpoint{3.749602in}{3.144407in}}{\pgfqpoint{3.460000in}{3.144407in}}%
\pgfpathcurveto{\pgfqpoint{3.170398in}{3.144407in}}{\pgfqpoint{2.892619in}{3.086877in}}{\pgfqpoint{2.687839in}{2.984487in}}%
\pgfpathcurveto{\pgfqpoint{2.483060in}{2.882097in}}{\pgfqpoint{2.368000in}{2.743208in}}{\pgfqpoint{2.368000in}{2.598407in}}%
\pgfpathcurveto{\pgfqpoint{2.368000in}{2.453606in}}{\pgfqpoint{2.483060in}{2.314716in}}{\pgfqpoint{2.687839in}{2.212326in}}%
\pgfpathcurveto{\pgfqpoint{2.892619in}{2.109937in}}{\pgfqpoint{3.170398in}{2.052407in}}{\pgfqpoint{3.460000in}{2.052407in}}%
\pgfpathlineto{\pgfqpoint{3.460000in}{2.052407in}}%
\pgfpathclose%
\pgfusepath{clip}%
\pgfsetrectcap%
\pgfsetroundjoin%
\pgfsetlinewidth{0.451687pt}%
\definecolor{currentstroke}{rgb}{0.690196,0.690196,0.690196}%
\pgfsetstrokecolor{currentstroke}%
\pgfsetstrokeopacity{0.250000}%
\pgfsetdash{}{0pt}%
\pgfpathmoveto{\pgfqpoint{2.997906in}{3.093112in}}%
\pgfpathlineto{\pgfqpoint{3.010228in}{3.093112in}}%
\pgfpathlineto{\pgfqpoint{3.022551in}{3.093112in}}%
\pgfpathlineto{\pgfqpoint{3.034873in}{3.093112in}}%
\pgfpathlineto{\pgfqpoint{3.047196in}{3.093112in}}%
\pgfpathlineto{\pgfqpoint{3.059518in}{3.093112in}}%
\pgfpathlineto{\pgfqpoint{3.071841in}{3.093112in}}%
\pgfpathlineto{\pgfqpoint{3.084163in}{3.093112in}}%
\pgfpathlineto{\pgfqpoint{3.096486in}{3.093112in}}%
\pgfpathlineto{\pgfqpoint{3.108808in}{3.093112in}}%
\pgfpathlineto{\pgfqpoint{3.121131in}{3.093112in}}%
\pgfpathlineto{\pgfqpoint{3.133453in}{3.093112in}}%
\pgfpathlineto{\pgfqpoint{3.145776in}{3.093112in}}%
\pgfpathlineto{\pgfqpoint{3.158098in}{3.093112in}}%
\pgfpathlineto{\pgfqpoint{3.170421in}{3.093112in}}%
\pgfpathlineto{\pgfqpoint{3.182743in}{3.093112in}}%
\pgfpathlineto{\pgfqpoint{3.195066in}{3.093112in}}%
\pgfpathlineto{\pgfqpoint{3.207388in}{3.093112in}}%
\pgfpathlineto{\pgfqpoint{3.219711in}{3.093112in}}%
\pgfpathlineto{\pgfqpoint{3.232033in}{3.093112in}}%
\pgfpathlineto{\pgfqpoint{3.244356in}{3.093112in}}%
\pgfpathlineto{\pgfqpoint{3.256678in}{3.093112in}}%
\pgfpathlineto{\pgfqpoint{3.269001in}{3.093112in}}%
\pgfpathlineto{\pgfqpoint{3.281323in}{3.093112in}}%
\pgfpathlineto{\pgfqpoint{3.293646in}{3.093112in}}%
\pgfpathlineto{\pgfqpoint{3.305969in}{3.093112in}}%
\pgfpathlineto{\pgfqpoint{3.318291in}{3.093112in}}%
\pgfpathlineto{\pgfqpoint{3.330614in}{3.093112in}}%
\pgfpathlineto{\pgfqpoint{3.342936in}{3.093112in}}%
\pgfpathlineto{\pgfqpoint{3.355259in}{3.093112in}}%
\pgfpathlineto{\pgfqpoint{3.367581in}{3.093112in}}%
\pgfpathlineto{\pgfqpoint{3.379904in}{3.093112in}}%
\pgfpathlineto{\pgfqpoint{3.392226in}{3.093112in}}%
\pgfpathlineto{\pgfqpoint{3.404549in}{3.093112in}}%
\pgfpathlineto{\pgfqpoint{3.416871in}{3.093112in}}%
\pgfpathlineto{\pgfqpoint{3.429194in}{3.093112in}}%
\pgfpathlineto{\pgfqpoint{3.441516in}{3.093112in}}%
\pgfpathlineto{\pgfqpoint{3.453839in}{3.093112in}}%
\pgfpathlineto{\pgfqpoint{3.466161in}{3.093112in}}%
\pgfpathlineto{\pgfqpoint{3.478484in}{3.093112in}}%
\pgfpathlineto{\pgfqpoint{3.490806in}{3.093112in}}%
\pgfpathlineto{\pgfqpoint{3.503129in}{3.093112in}}%
\pgfpathlineto{\pgfqpoint{3.515451in}{3.093112in}}%
\pgfpathlineto{\pgfqpoint{3.527774in}{3.093112in}}%
\pgfpathlineto{\pgfqpoint{3.540096in}{3.093112in}}%
\pgfpathlineto{\pgfqpoint{3.552419in}{3.093112in}}%
\pgfpathlineto{\pgfqpoint{3.564741in}{3.093112in}}%
\pgfpathlineto{\pgfqpoint{3.577064in}{3.093112in}}%
\pgfpathlineto{\pgfqpoint{3.589386in}{3.093112in}}%
\pgfpathlineto{\pgfqpoint{3.601709in}{3.093112in}}%
\pgfpathlineto{\pgfqpoint{3.614031in}{3.093112in}}%
\pgfpathlineto{\pgfqpoint{3.626354in}{3.093112in}}%
\pgfpathlineto{\pgfqpoint{3.638677in}{3.093112in}}%
\pgfpathlineto{\pgfqpoint{3.650999in}{3.093112in}}%
\pgfpathlineto{\pgfqpoint{3.663322in}{3.093112in}}%
\pgfpathlineto{\pgfqpoint{3.675644in}{3.093112in}}%
\pgfpathlineto{\pgfqpoint{3.687967in}{3.093112in}}%
\pgfpathlineto{\pgfqpoint{3.700289in}{3.093112in}}%
\pgfpathlineto{\pgfqpoint{3.712612in}{3.093112in}}%
\pgfpathlineto{\pgfqpoint{3.724934in}{3.093112in}}%
\pgfpathlineto{\pgfqpoint{3.737257in}{3.093112in}}%
\pgfpathlineto{\pgfqpoint{3.749579in}{3.093112in}}%
\pgfpathlineto{\pgfqpoint{3.761902in}{3.093112in}}%
\pgfpathlineto{\pgfqpoint{3.774224in}{3.093112in}}%
\pgfpathlineto{\pgfqpoint{3.786547in}{3.093112in}}%
\pgfpathlineto{\pgfqpoint{3.798869in}{3.093112in}}%
\pgfpathlineto{\pgfqpoint{3.811192in}{3.093112in}}%
\pgfpathlineto{\pgfqpoint{3.823514in}{3.093112in}}%
\pgfpathlineto{\pgfqpoint{3.835837in}{3.093112in}}%
\pgfpathlineto{\pgfqpoint{3.848159in}{3.093112in}}%
\pgfpathlineto{\pgfqpoint{3.860482in}{3.093112in}}%
\pgfpathlineto{\pgfqpoint{3.872804in}{3.093112in}}%
\pgfpathlineto{\pgfqpoint{3.885127in}{3.093112in}}%
\pgfpathlineto{\pgfqpoint{3.897449in}{3.093112in}}%
\pgfpathlineto{\pgfqpoint{3.909772in}{3.093112in}}%
\pgfpathlineto{\pgfqpoint{3.922094in}{3.093112in}}%
\pgfusepath{stroke}%
\end{pgfscope}%
\begin{pgfscope}%
\pgfsetrectcap%
\pgfsetmiterjoin%
\pgfsetlinewidth{0.803000pt}%
\definecolor{currentstroke}{rgb}{0.000000,0.000000,0.000000}%
\pgfsetstrokecolor{currentstroke}%
\pgfsetdash{}{0pt}%
\pgfpathmoveto{\pgfqpoint{3.460000in}{2.052407in}}%
\pgfpathcurveto{\pgfqpoint{3.749602in}{2.052407in}}{\pgfqpoint{4.027381in}{2.109937in}}{\pgfqpoint{4.232161in}{2.212326in}}%
\pgfpathcurveto{\pgfqpoint{4.436940in}{2.314716in}}{\pgfqpoint{4.552000in}{2.453606in}}{\pgfqpoint{4.552000in}{2.598407in}}%
\pgfpathcurveto{\pgfqpoint{4.552000in}{2.743208in}}{\pgfqpoint{4.436940in}{2.882097in}}{\pgfqpoint{4.232161in}{2.984487in}}%
\pgfpathcurveto{\pgfqpoint{4.027381in}{3.086877in}}{\pgfqpoint{3.749602in}{3.144407in}}{\pgfqpoint{3.460000in}{3.144407in}}%
\pgfpathcurveto{\pgfqpoint{3.170398in}{3.144407in}}{\pgfqpoint{2.892619in}{3.086877in}}{\pgfqpoint{2.687839in}{2.984487in}}%
\pgfpathcurveto{\pgfqpoint{2.483060in}{2.882097in}}{\pgfqpoint{2.368000in}{2.743208in}}{\pgfqpoint{2.368000in}{2.598407in}}%
\pgfpathcurveto{\pgfqpoint{2.368000in}{2.453606in}}{\pgfqpoint{2.483060in}{2.314716in}}{\pgfqpoint{2.687839in}{2.212326in}}%
\pgfpathcurveto{\pgfqpoint{2.892619in}{2.109937in}}{\pgfqpoint{3.170398in}{2.052407in}}{\pgfqpoint{3.460000in}{2.052407in}}%
\pgfpathlineto{\pgfqpoint{3.460000in}{2.052407in}}%
\pgfpathclose%
\pgfusepath{stroke}%
\end{pgfscope}%
\begin{pgfscope}%
\definecolor{textcolor}{rgb}{0.000000,0.000000,0.000000}%
\pgfsetstrokecolor{textcolor}%
\pgfsetfillcolor{textcolor}%
\pgftext[x=3.460000in,y=3.186073in,,base]{\color{textcolor}{\rmfamily\fontsize{8.200000}{9.840000}\selectfont\catcode`\^=\active\def^{\ifmmode\sp\else\^{}\fi}\catcode`\%=\active\def%{\%}Antipodal}}%
\end{pgfscope}%
\begin{pgfscope}%
\pgfsetbuttcap%
\pgfsetmiterjoin%
\definecolor{currentfill}{rgb}{1.000000,1.000000,1.000000}%
\pgfsetfillcolor{currentfill}%
\pgfsetlinewidth{0.000000pt}%
\definecolor{currentstroke}{rgb}{0.000000,0.000000,0.000000}%
\pgfsetstrokecolor{currentstroke}%
\pgfsetstrokeopacity{0.000000}%
\pgfsetdash{}{0pt}%
\pgfpathmoveto{\pgfqpoint{5.728000in}{2.052407in}}%
\pgfpathcurveto{\pgfqpoint{6.017602in}{2.052407in}}{\pgfqpoint{6.295381in}{2.109937in}}{\pgfqpoint{6.500161in}{2.212326in}}%
\pgfpathcurveto{\pgfqpoint{6.704940in}{2.314716in}}{\pgfqpoint{6.820000in}{2.453606in}}{\pgfqpoint{6.820000in}{2.598407in}}%
\pgfpathcurveto{\pgfqpoint{6.820000in}{2.743208in}}{\pgfqpoint{6.704940in}{2.882097in}}{\pgfqpoint{6.500161in}{2.984487in}}%
\pgfpathcurveto{\pgfqpoint{6.295381in}{3.086877in}}{\pgfqpoint{6.017602in}{3.144407in}}{\pgfqpoint{5.728000in}{3.144407in}}%
\pgfpathcurveto{\pgfqpoint{5.438398in}{3.144407in}}{\pgfqpoint{5.160619in}{3.086877in}}{\pgfqpoint{4.955839in}{2.984487in}}%
\pgfpathcurveto{\pgfqpoint{4.751060in}{2.882097in}}{\pgfqpoint{4.636000in}{2.743208in}}{\pgfqpoint{4.636000in}{2.598407in}}%
\pgfpathcurveto{\pgfqpoint{4.636000in}{2.453606in}}{\pgfqpoint{4.751060in}{2.314716in}}{\pgfqpoint{4.955839in}{2.212326in}}%
\pgfpathcurveto{\pgfqpoint{5.160619in}{2.109937in}}{\pgfqpoint{5.438398in}{2.052407in}}{\pgfqpoint{5.728000in}{2.052407in}}%
\pgfpathlineto{\pgfqpoint{5.728000in}{2.052407in}}%
\pgfpathclose%
\pgfusepath{fill}%
\end{pgfscope}%
\begin{pgfscope}%
\pgfpathmoveto{\pgfqpoint{5.728000in}{2.052407in}}%
\pgfpathcurveto{\pgfqpoint{6.017602in}{2.052407in}}{\pgfqpoint{6.295381in}{2.109937in}}{\pgfqpoint{6.500161in}{2.212326in}}%
\pgfpathcurveto{\pgfqpoint{6.704940in}{2.314716in}}{\pgfqpoint{6.820000in}{2.453606in}}{\pgfqpoint{6.820000in}{2.598407in}}%
\pgfpathcurveto{\pgfqpoint{6.820000in}{2.743208in}}{\pgfqpoint{6.704940in}{2.882097in}}{\pgfqpoint{6.500161in}{2.984487in}}%
\pgfpathcurveto{\pgfqpoint{6.295381in}{3.086877in}}{\pgfqpoint{6.017602in}{3.144407in}}{\pgfqpoint{5.728000in}{3.144407in}}%
\pgfpathcurveto{\pgfqpoint{5.438398in}{3.144407in}}{\pgfqpoint{5.160619in}{3.086877in}}{\pgfqpoint{4.955839in}{2.984487in}}%
\pgfpathcurveto{\pgfqpoint{4.751060in}{2.882097in}}{\pgfqpoint{4.636000in}{2.743208in}}{\pgfqpoint{4.636000in}{2.598407in}}%
\pgfpathcurveto{\pgfqpoint{4.636000in}{2.453606in}}{\pgfqpoint{4.751060in}{2.314716in}}{\pgfqpoint{4.955839in}{2.212326in}}%
\pgfpathcurveto{\pgfqpoint{5.160619in}{2.109937in}}{\pgfqpoint{5.438398in}{2.052407in}}{\pgfqpoint{5.728000in}{2.052407in}}%
\pgfpathlineto{\pgfqpoint{5.728000in}{2.052407in}}%
\pgfpathclose%
\pgfusepath{clip}%
\pgfsetbuttcap%
\pgfsetroundjoin%
\definecolor{currentfill}{rgb}{0.121569,0.466667,0.705882}%
\pgfsetfillcolor{currentfill}%
\pgfsetfillopacity{0.580000}%
\pgfsetlinewidth{0.000000pt}%
\definecolor{currentstroke}{rgb}{0.121569,0.466667,0.705882}%
\pgfsetstrokecolor{currentstroke}%
\pgfsetstrokeopacity{0.580000}%
\pgfsetdash{}{0pt}%
\pgfsys@defobject{currentmarker}{\pgfqpoint{-0.014232in}{-0.014232in}}{\pgfqpoint{0.014232in}{0.014232in}}{%
\pgfpathmoveto{\pgfqpoint{0.000000in}{-0.014232in}}%
\pgfpathcurveto{\pgfqpoint{0.003774in}{-0.014232in}}{\pgfqpoint{0.007395in}{-0.012732in}}{\pgfqpoint{0.010063in}{-0.010063in}}%
\pgfpathcurveto{\pgfqpoint{0.012732in}{-0.007395in}}{\pgfqpoint{0.014232in}{-0.003774in}}{\pgfqpoint{0.014232in}{0.000000in}}%
\pgfpathcurveto{\pgfqpoint{0.014232in}{0.003774in}}{\pgfqpoint{0.012732in}{0.007395in}}{\pgfqpoint{0.010063in}{0.010063in}}%
\pgfpathcurveto{\pgfqpoint{0.007395in}{0.012732in}}{\pgfqpoint{0.003774in}{0.014232in}}{\pgfqpoint{0.000000in}{0.014232in}}%
\pgfpathcurveto{\pgfqpoint{-0.003774in}{0.014232in}}{\pgfqpoint{-0.007395in}{0.012732in}}{\pgfqpoint{-0.010063in}{0.010063in}}%
\pgfpathcurveto{\pgfqpoint{-0.012732in}{0.007395in}}{\pgfqpoint{-0.014232in}{0.003774in}}{\pgfqpoint{-0.014232in}{0.000000in}}%
\pgfpathcurveto{\pgfqpoint{-0.014232in}{-0.003774in}}{\pgfqpoint{-0.012732in}{-0.007395in}}{\pgfqpoint{-0.010063in}{-0.010063in}}%
\pgfpathcurveto{\pgfqpoint{-0.007395in}{-0.012732in}}{\pgfqpoint{-0.003774in}{-0.014232in}}{\pgfqpoint{0.000000in}{-0.014232in}}%
\pgfpathlineto{\pgfqpoint{0.000000in}{-0.014232in}}%
\pgfpathclose%
\pgfusepath{fill}%
}%
\begin{pgfscope}%
\pgfsys@transformshift{5.911170in}{2.602299in}%
\pgfsys@useobject{currentmarker}{}%
\end{pgfscope}%
\begin{pgfscope}%
\pgfsys@transformshift{5.099640in}{2.569639in}%
\pgfsys@useobject{currentmarker}{}%
\end{pgfscope}%
\begin{pgfscope}%
\pgfsys@transformshift{6.367842in}{2.579597in}%
\pgfsys@useobject{currentmarker}{}%
\end{pgfscope}%
\begin{pgfscope}%
\pgfsys@transformshift{5.510115in}{2.618349in}%
\pgfsys@useobject{currentmarker}{}%
\end{pgfscope}%
\begin{pgfscope}%
\pgfsys@transformshift{5.495769in}{2.584256in}%
\pgfsys@useobject{currentmarker}{}%
\end{pgfscope}%
\begin{pgfscope}%
\pgfsys@transformshift{5.589083in}{2.606372in}%
\pgfsys@useobject{currentmarker}{}%
\end{pgfscope}%
\begin{pgfscope}%
\pgfsys@transformshift{6.287383in}{2.597226in}%
\pgfsys@useobject{currentmarker}{}%
\end{pgfscope}%
\begin{pgfscope}%
\pgfsys@transformshift{5.729434in}{2.582728in}%
\pgfsys@useobject{currentmarker}{}%
\end{pgfscope}%
\begin{pgfscope}%
\pgfsys@transformshift{6.574491in}{2.593180in}%
\pgfsys@useobject{currentmarker}{}%
\end{pgfscope}%
\begin{pgfscope}%
\pgfsys@transformshift{5.646246in}{2.590689in}%
\pgfsys@useobject{currentmarker}{}%
\end{pgfscope}%
\begin{pgfscope}%
\pgfsys@transformshift{5.212059in}{2.591092in}%
\pgfsys@useobject{currentmarker}{}%
\end{pgfscope}%
\begin{pgfscope}%
\pgfsys@transformshift{5.858269in}{2.609718in}%
\pgfsys@useobject{currentmarker}{}%
\end{pgfscope}%
\begin{pgfscope}%
\pgfsys@transformshift{5.078522in}{2.608198in}%
\pgfsys@useobject{currentmarker}{}%
\end{pgfscope}%
\begin{pgfscope}%
\pgfsys@transformshift{5.546890in}{2.589740in}%
\pgfsys@useobject{currentmarker}{}%
\end{pgfscope}%
\begin{pgfscope}%
\pgfsys@transformshift{5.657522in}{2.584791in}%
\pgfsys@useobject{currentmarker}{}%
\end{pgfscope}%
\begin{pgfscope}%
\pgfsys@transformshift{5.107738in}{2.602644in}%
\pgfsys@useobject{currentmarker}{}%
\end{pgfscope}%
\begin{pgfscope}%
\pgfsys@transformshift{5.428783in}{2.596906in}%
\pgfsys@useobject{currentmarker}{}%
\end{pgfscope}%
\begin{pgfscope}%
\pgfsys@transformshift{5.382416in}{2.613841in}%
\pgfsys@useobject{currentmarker}{}%
\end{pgfscope}%
\begin{pgfscope}%
\pgfsys@transformshift{6.508520in}{2.638277in}%
\pgfsys@useobject{currentmarker}{}%
\end{pgfscope}%
\begin{pgfscope}%
\pgfsys@transformshift{5.944254in}{2.611059in}%
\pgfsys@useobject{currentmarker}{}%
\end{pgfscope}%
\begin{pgfscope}%
\pgfsys@transformshift{5.132782in}{2.597743in}%
\pgfsys@useobject{currentmarker}{}%
\end{pgfscope}%
\begin{pgfscope}%
\pgfsys@transformshift{6.027015in}{2.595825in}%
\pgfsys@useobject{currentmarker}{}%
\end{pgfscope}%
\begin{pgfscope}%
\pgfsys@transformshift{6.566857in}{2.585548in}%
\pgfsys@useobject{currentmarker}{}%
\end{pgfscope}%
\begin{pgfscope}%
\pgfsys@transformshift{5.718690in}{2.584798in}%
\pgfsys@useobject{currentmarker}{}%
\end{pgfscope}%
\begin{pgfscope}%
\pgfsys@transformshift{4.766937in}{2.583346in}%
\pgfsys@useobject{currentmarker}{}%
\end{pgfscope}%
\begin{pgfscope}%
\pgfsys@transformshift{4.683542in}{2.592473in}%
\pgfsys@useobject{currentmarker}{}%
\end{pgfscope}%
\begin{pgfscope}%
\pgfsys@transformshift{5.614178in}{2.591996in}%
\pgfsys@useobject{currentmarker}{}%
\end{pgfscope}%
\begin{pgfscope}%
\pgfsys@transformshift{6.348214in}{2.609570in}%
\pgfsys@useobject{currentmarker}{}%
\end{pgfscope}%
\begin{pgfscope}%
\pgfsys@transformshift{5.726757in}{2.592418in}%
\pgfsys@useobject{currentmarker}{}%
\end{pgfscope}%
\begin{pgfscope}%
\pgfsys@transformshift{5.535767in}{2.595443in}%
\pgfsys@useobject{currentmarker}{}%
\end{pgfscope}%
\begin{pgfscope}%
\pgfsys@transformshift{6.522320in}{2.580120in}%
\pgfsys@useobject{currentmarker}{}%
\end{pgfscope}%
\begin{pgfscope}%
\pgfsys@transformshift{6.441961in}{2.623094in}%
\pgfsys@useobject{currentmarker}{}%
\end{pgfscope}%
\begin{pgfscope}%
\pgfsys@transformshift{6.004388in}{2.606016in}%
\pgfsys@useobject{currentmarker}{}%
\end{pgfscope}%
\begin{pgfscope}%
\pgfsys@transformshift{5.543981in}{2.625261in}%
\pgfsys@useobject{currentmarker}{}%
\end{pgfscope}%
\begin{pgfscope}%
\pgfsys@transformshift{6.133711in}{2.610368in}%
\pgfsys@useobject{currentmarker}{}%
\end{pgfscope}%
\begin{pgfscope}%
\pgfsys@transformshift{5.747593in}{2.620768in}%
\pgfsys@useobject{currentmarker}{}%
\end{pgfscope}%
\begin{pgfscope}%
\pgfsys@transformshift{5.487933in}{2.594677in}%
\pgfsys@useobject{currentmarker}{}%
\end{pgfscope}%
\begin{pgfscope}%
\pgfsys@transformshift{4.850454in}{2.609171in}%
\pgfsys@useobject{currentmarker}{}%
\end{pgfscope}%
\begin{pgfscope}%
\pgfsys@transformshift{5.013304in}{2.639031in}%
\pgfsys@useobject{currentmarker}{}%
\end{pgfscope}%
\begin{pgfscope}%
\pgfsys@transformshift{5.145510in}{2.580041in}%
\pgfsys@useobject{currentmarker}{}%
\end{pgfscope}%
\begin{pgfscope}%
\pgfsys@transformshift{5.821124in}{2.616063in}%
\pgfsys@useobject{currentmarker}{}%
\end{pgfscope}%
\begin{pgfscope}%
\pgfsys@transformshift{5.928393in}{2.623896in}%
\pgfsys@useobject{currentmarker}{}%
\end{pgfscope}%
\begin{pgfscope}%
\pgfsys@transformshift{6.819571in}{2.604728in}%
\pgfsys@useobject{currentmarker}{}%
\end{pgfscope}%
\begin{pgfscope}%
\pgfsys@transformshift{5.837746in}{2.609892in}%
\pgfsys@useobject{currentmarker}{}%
\end{pgfscope}%
\begin{pgfscope}%
\pgfsys@transformshift{4.646108in}{2.617615in}%
\pgfsys@useobject{currentmarker}{}%
\end{pgfscope}%
\begin{pgfscope}%
\pgfsys@transformshift{5.152261in}{2.588346in}%
\pgfsys@useobject{currentmarker}{}%
\end{pgfscope}%
\begin{pgfscope}%
\pgfsys@transformshift{4.979831in}{2.604672in}%
\pgfsys@useobject{currentmarker}{}%
\end{pgfscope}%
\begin{pgfscope}%
\pgfsys@transformshift{4.891008in}{2.585078in}%
\pgfsys@useobject{currentmarker}{}%
\end{pgfscope}%
\begin{pgfscope}%
\pgfsys@transformshift{5.829332in}{2.588032in}%
\pgfsys@useobject{currentmarker}{}%
\end{pgfscope}%
\begin{pgfscope}%
\pgfsys@transformshift{6.244776in}{2.588927in}%
\pgfsys@useobject{currentmarker}{}%
\end{pgfscope}%
\begin{pgfscope}%
\pgfsys@transformshift{5.504130in}{2.596211in}%
\pgfsys@useobject{currentmarker}{}%
\end{pgfscope}%
\begin{pgfscope}%
\pgfsys@transformshift{5.997111in}{2.600655in}%
\pgfsys@useobject{currentmarker}{}%
\end{pgfscope}%
\begin{pgfscope}%
\pgfsys@transformshift{4.708643in}{2.605357in}%
\pgfsys@useobject{currentmarker}{}%
\end{pgfscope}%
\begin{pgfscope}%
\pgfsys@transformshift{5.837219in}{2.576627in}%
\pgfsys@useobject{currentmarker}{}%
\end{pgfscope}%
\begin{pgfscope}%
\pgfsys@transformshift{5.595425in}{2.628834in}%
\pgfsys@useobject{currentmarker}{}%
\end{pgfscope}%
\begin{pgfscope}%
\pgfsys@transformshift{6.102496in}{2.616172in}%
\pgfsys@useobject{currentmarker}{}%
\end{pgfscope}%
\begin{pgfscope}%
\pgfsys@transformshift{6.623301in}{2.609567in}%
\pgfsys@useobject{currentmarker}{}%
\end{pgfscope}%
\begin{pgfscope}%
\pgfsys@transformshift{4.738220in}{2.592905in}%
\pgfsys@useobject{currentmarker}{}%
\end{pgfscope}%
\begin{pgfscope}%
\pgfsys@transformshift{5.113537in}{2.597036in}%
\pgfsys@useobject{currentmarker}{}%
\end{pgfscope}%
\begin{pgfscope}%
\pgfsys@transformshift{5.591412in}{2.606829in}%
\pgfsys@useobject{currentmarker}{}%
\end{pgfscope}%
\begin{pgfscope}%
\pgfsys@transformshift{5.152256in}{2.604732in}%
\pgfsys@useobject{currentmarker}{}%
\end{pgfscope}%
\begin{pgfscope}%
\pgfsys@transformshift{6.485803in}{2.572571in}%
\pgfsys@useobject{currentmarker}{}%
\end{pgfscope}%
\begin{pgfscope}%
\pgfsys@transformshift{5.892506in}{2.609803in}%
\pgfsys@useobject{currentmarker}{}%
\end{pgfscope}%
\begin{pgfscope}%
\pgfsys@transformshift{5.091135in}{2.643281in}%
\pgfsys@useobject{currentmarker}{}%
\end{pgfscope}%
\begin{pgfscope}%
\pgfsys@transformshift{5.668044in}{2.569984in}%
\pgfsys@useobject{currentmarker}{}%
\end{pgfscope}%
\begin{pgfscope}%
\pgfsys@transformshift{5.765148in}{2.613712in}%
\pgfsys@useobject{currentmarker}{}%
\end{pgfscope}%
\begin{pgfscope}%
\pgfsys@transformshift{5.235927in}{2.604088in}%
\pgfsys@useobject{currentmarker}{}%
\end{pgfscope}%
\begin{pgfscope}%
\pgfsys@transformshift{6.112692in}{2.596690in}%
\pgfsys@useobject{currentmarker}{}%
\end{pgfscope}%
\begin{pgfscope}%
\pgfsys@transformshift{6.731787in}{2.619527in}%
\pgfsys@useobject{currentmarker}{}%
\end{pgfscope}%
\begin{pgfscope}%
\pgfsys@transformshift{4.686019in}{2.580190in}%
\pgfsys@useobject{currentmarker}{}%
\end{pgfscope}%
\begin{pgfscope}%
\pgfsys@transformshift{4.885916in}{2.600915in}%
\pgfsys@useobject{currentmarker}{}%
\end{pgfscope}%
\begin{pgfscope}%
\pgfsys@transformshift{5.705168in}{2.620953in}%
\pgfsys@useobject{currentmarker}{}%
\end{pgfscope}%
\begin{pgfscope}%
\pgfsys@transformshift{5.045297in}{2.595373in}%
\pgfsys@useobject{currentmarker}{}%
\end{pgfscope}%
\begin{pgfscope}%
\pgfsys@transformshift{6.584914in}{2.591964in}%
\pgfsys@useobject{currentmarker}{}%
\end{pgfscope}%
\begin{pgfscope}%
\pgfsys@transformshift{5.105609in}{2.600177in}%
\pgfsys@useobject{currentmarker}{}%
\end{pgfscope}%
\begin{pgfscope}%
\pgfsys@transformshift{6.747261in}{2.591551in}%
\pgfsys@useobject{currentmarker}{}%
\end{pgfscope}%
\begin{pgfscope}%
\pgfsys@transformshift{6.629392in}{2.622262in}%
\pgfsys@useobject{currentmarker}{}%
\end{pgfscope}%
\begin{pgfscope}%
\pgfsys@transformshift{6.526325in}{2.584420in}%
\pgfsys@useobject{currentmarker}{}%
\end{pgfscope}%
\begin{pgfscope}%
\pgfsys@transformshift{6.805271in}{2.611869in}%
\pgfsys@useobject{currentmarker}{}%
\end{pgfscope}%
\begin{pgfscope}%
\pgfsys@transformshift{6.341718in}{2.579093in}%
\pgfsys@useobject{currentmarker}{}%
\end{pgfscope}%
\begin{pgfscope}%
\pgfsys@transformshift{6.256653in}{2.608727in}%
\pgfsys@useobject{currentmarker}{}%
\end{pgfscope}%
\begin{pgfscope}%
\pgfsys@transformshift{6.052250in}{2.596789in}%
\pgfsys@useobject{currentmarker}{}%
\end{pgfscope}%
\begin{pgfscope}%
\pgfsys@transformshift{6.686884in}{2.558539in}%
\pgfsys@useobject{currentmarker}{}%
\end{pgfscope}%
\begin{pgfscope}%
\pgfsys@transformshift{6.368483in}{2.598289in}%
\pgfsys@useobject{currentmarker}{}%
\end{pgfscope}%
\begin{pgfscope}%
\pgfsys@transformshift{5.372042in}{2.581944in}%
\pgfsys@useobject{currentmarker}{}%
\end{pgfscope}%
\begin{pgfscope}%
\pgfsys@transformshift{5.395827in}{2.591551in}%
\pgfsys@useobject{currentmarker}{}%
\end{pgfscope}%
\begin{pgfscope}%
\pgfsys@transformshift{6.478316in}{2.622148in}%
\pgfsys@useobject{currentmarker}{}%
\end{pgfscope}%
\begin{pgfscope}%
\pgfsys@transformshift{4.896621in}{2.581166in}%
\pgfsys@useobject{currentmarker}{}%
\end{pgfscope}%
\begin{pgfscope}%
\pgfsys@transformshift{5.514699in}{2.582047in}%
\pgfsys@useobject{currentmarker}{}%
\end{pgfscope}%
\begin{pgfscope}%
\pgfsys@transformshift{6.511759in}{2.620636in}%
\pgfsys@useobject{currentmarker}{}%
\end{pgfscope}%
\begin{pgfscope}%
\pgfsys@transformshift{4.755504in}{2.600196in}%
\pgfsys@useobject{currentmarker}{}%
\end{pgfscope}%
\begin{pgfscope}%
\pgfsys@transformshift{6.695812in}{2.622195in}%
\pgfsys@useobject{currentmarker}{}%
\end{pgfscope}%
\begin{pgfscope}%
\pgfsys@transformshift{4.964211in}{2.603399in}%
\pgfsys@useobject{currentmarker}{}%
\end{pgfscope}%
\begin{pgfscope}%
\pgfsys@transformshift{5.640167in}{2.584792in}%
\pgfsys@useobject{currentmarker}{}%
\end{pgfscope}%
\begin{pgfscope}%
\pgfsys@transformshift{6.030754in}{2.597384in}%
\pgfsys@useobject{currentmarker}{}%
\end{pgfscope}%
\begin{pgfscope}%
\pgfsys@transformshift{5.652447in}{2.617382in}%
\pgfsys@useobject{currentmarker}{}%
\end{pgfscope}%
\begin{pgfscope}%
\pgfsys@transformshift{4.998220in}{2.612929in}%
\pgfsys@useobject{currentmarker}{}%
\end{pgfscope}%
\begin{pgfscope}%
\pgfsys@transformshift{5.282181in}{2.598424in}%
\pgfsys@useobject{currentmarker}{}%
\end{pgfscope}%
\begin{pgfscope}%
\pgfsys@transformshift{4.931107in}{2.599365in}%
\pgfsys@useobject{currentmarker}{}%
\end{pgfscope}%
\begin{pgfscope}%
\pgfsys@transformshift{5.444774in}{2.596670in}%
\pgfsys@useobject{currentmarker}{}%
\end{pgfscope}%
\begin{pgfscope}%
\pgfsys@transformshift{6.440410in}{2.589481in}%
\pgfsys@useobject{currentmarker}{}%
\end{pgfscope}%
\begin{pgfscope}%
\pgfsys@transformshift{5.091327in}{2.628220in}%
\pgfsys@useobject{currentmarker}{}%
\end{pgfscope}%
\begin{pgfscope}%
\pgfsys@transformshift{6.304823in}{2.598548in}%
\pgfsys@useobject{currentmarker}{}%
\end{pgfscope}%
\begin{pgfscope}%
\pgfsys@transformshift{4.655456in}{2.580951in}%
\pgfsys@useobject{currentmarker}{}%
\end{pgfscope}%
\begin{pgfscope}%
\pgfsys@transformshift{5.032438in}{2.591046in}%
\pgfsys@useobject{currentmarker}{}%
\end{pgfscope}%
\begin{pgfscope}%
\pgfsys@transformshift{6.167163in}{2.609616in}%
\pgfsys@useobject{currentmarker}{}%
\end{pgfscope}%
\begin{pgfscope}%
\pgfsys@transformshift{5.861550in}{2.608306in}%
\pgfsys@useobject{currentmarker}{}%
\end{pgfscope}%
\begin{pgfscope}%
\pgfsys@transformshift{5.271434in}{2.596291in}%
\pgfsys@useobject{currentmarker}{}%
\end{pgfscope}%
\begin{pgfscope}%
\pgfsys@transformshift{5.310813in}{2.588091in}%
\pgfsys@useobject{currentmarker}{}%
\end{pgfscope}%
\begin{pgfscope}%
\pgfsys@transformshift{4.997812in}{2.598877in}%
\pgfsys@useobject{currentmarker}{}%
\end{pgfscope}%
\begin{pgfscope}%
\pgfsys@transformshift{6.429221in}{2.572577in}%
\pgfsys@useobject{currentmarker}{}%
\end{pgfscope}%
\begin{pgfscope}%
\pgfsys@transformshift{4.981955in}{2.578316in}%
\pgfsys@useobject{currentmarker}{}%
\end{pgfscope}%
\begin{pgfscope}%
\pgfsys@transformshift{5.911148in}{2.610736in}%
\pgfsys@useobject{currentmarker}{}%
\end{pgfscope}%
\begin{pgfscope}%
\pgfsys@transformshift{4.990989in}{2.591011in}%
\pgfsys@useobject{currentmarker}{}%
\end{pgfscope}%
\begin{pgfscope}%
\pgfsys@transformshift{4.802470in}{2.604193in}%
\pgfsys@useobject{currentmarker}{}%
\end{pgfscope}%
\begin{pgfscope}%
\pgfsys@transformshift{4.892470in}{2.617209in}%
\pgfsys@useobject{currentmarker}{}%
\end{pgfscope}%
\begin{pgfscope}%
\pgfsys@transformshift{6.704049in}{2.608432in}%
\pgfsys@useobject{currentmarker}{}%
\end{pgfscope}%
\begin{pgfscope}%
\pgfsys@transformshift{6.219746in}{2.598812in}%
\pgfsys@useobject{currentmarker}{}%
\end{pgfscope}%
\begin{pgfscope}%
\pgfsys@transformshift{5.130400in}{2.631255in}%
\pgfsys@useobject{currentmarker}{}%
\end{pgfscope}%
\begin{pgfscope}%
\pgfsys@transformshift{6.741480in}{2.608711in}%
\pgfsys@useobject{currentmarker}{}%
\end{pgfscope}%
\begin{pgfscope}%
\pgfsys@transformshift{5.171938in}{2.592937in}%
\pgfsys@useobject{currentmarker}{}%
\end{pgfscope}%
\begin{pgfscope}%
\pgfsys@transformshift{6.523199in}{2.609704in}%
\pgfsys@useobject{currentmarker}{}%
\end{pgfscope}%
\begin{pgfscope}%
\pgfsys@transformshift{5.919719in}{2.604701in}%
\pgfsys@useobject{currentmarker}{}%
\end{pgfscope}%
\begin{pgfscope}%
\pgfsys@transformshift{4.729020in}{2.585199in}%
\pgfsys@useobject{currentmarker}{}%
\end{pgfscope}%
\begin{pgfscope}%
\pgfsys@transformshift{4.904113in}{2.599533in}%
\pgfsys@useobject{currentmarker}{}%
\end{pgfscope}%
\begin{pgfscope}%
\pgfsys@transformshift{5.557163in}{2.595861in}%
\pgfsys@useobject{currentmarker}{}%
\end{pgfscope}%
\begin{pgfscope}%
\pgfsys@transformshift{5.483235in}{2.568058in}%
\pgfsys@useobject{currentmarker}{}%
\end{pgfscope}%
\begin{pgfscope}%
\pgfsys@transformshift{5.852812in}{2.595746in}%
\pgfsys@useobject{currentmarker}{}%
\end{pgfscope}%
\begin{pgfscope}%
\pgfsys@transformshift{5.229715in}{2.594038in}%
\pgfsys@useobject{currentmarker}{}%
\end{pgfscope}%
\begin{pgfscope}%
\pgfsys@transformshift{6.298773in}{2.591643in}%
\pgfsys@useobject{currentmarker}{}%
\end{pgfscope}%
\begin{pgfscope}%
\pgfsys@transformshift{5.067396in}{2.573021in}%
\pgfsys@useobject{currentmarker}{}%
\end{pgfscope}%
\begin{pgfscope}%
\pgfsys@transformshift{6.556747in}{2.592934in}%
\pgfsys@useobject{currentmarker}{}%
\end{pgfscope}%
\begin{pgfscope}%
\pgfsys@transformshift{6.431418in}{2.597627in}%
\pgfsys@useobject{currentmarker}{}%
\end{pgfscope}%
\begin{pgfscope}%
\pgfsys@transformshift{6.794407in}{2.599980in}%
\pgfsys@useobject{currentmarker}{}%
\end{pgfscope}%
\begin{pgfscope}%
\pgfsys@transformshift{4.703208in}{2.580885in}%
\pgfsys@useobject{currentmarker}{}%
\end{pgfscope}%
\begin{pgfscope}%
\pgfsys@transformshift{5.060858in}{2.608166in}%
\pgfsys@useobject{currentmarker}{}%
\end{pgfscope}%
\begin{pgfscope}%
\pgfsys@transformshift{4.937604in}{2.579649in}%
\pgfsys@useobject{currentmarker}{}%
\end{pgfscope}%
\begin{pgfscope}%
\pgfsys@transformshift{5.025324in}{2.594954in}%
\pgfsys@useobject{currentmarker}{}%
\end{pgfscope}%
\begin{pgfscope}%
\pgfsys@transformshift{6.566253in}{2.618546in}%
\pgfsys@useobject{currentmarker}{}%
\end{pgfscope}%
\begin{pgfscope}%
\pgfsys@transformshift{5.023446in}{2.587724in}%
\pgfsys@useobject{currentmarker}{}%
\end{pgfscope}%
\begin{pgfscope}%
\pgfsys@transformshift{6.482247in}{2.590527in}%
\pgfsys@useobject{currentmarker}{}%
\end{pgfscope}%
\begin{pgfscope}%
\pgfsys@transformshift{4.871247in}{2.614968in}%
\pgfsys@useobject{currentmarker}{}%
\end{pgfscope}%
\begin{pgfscope}%
\pgfsys@transformshift{6.731561in}{2.592482in}%
\pgfsys@useobject{currentmarker}{}%
\end{pgfscope}%
\begin{pgfscope}%
\pgfsys@transformshift{4.773186in}{2.593825in}%
\pgfsys@useobject{currentmarker}{}%
\end{pgfscope}%
\begin{pgfscope}%
\pgfsys@transformshift{5.766303in}{2.601819in}%
\pgfsys@useobject{currentmarker}{}%
\end{pgfscope}%
\begin{pgfscope}%
\pgfsys@transformshift{6.816871in}{2.614594in}%
\pgfsys@useobject{currentmarker}{}%
\end{pgfscope}%
\begin{pgfscope}%
\pgfsys@transformshift{6.207478in}{2.564407in}%
\pgfsys@useobject{currentmarker}{}%
\end{pgfscope}%
\begin{pgfscope}%
\pgfsys@transformshift{6.570384in}{2.585631in}%
\pgfsys@useobject{currentmarker}{}%
\end{pgfscope}%
\begin{pgfscope}%
\pgfsys@transformshift{5.889050in}{2.582873in}%
\pgfsys@useobject{currentmarker}{}%
\end{pgfscope}%
\begin{pgfscope}%
\pgfsys@transformshift{4.979024in}{2.581679in}%
\pgfsys@useobject{currentmarker}{}%
\end{pgfscope}%
\begin{pgfscope}%
\pgfsys@transformshift{5.495519in}{2.585232in}%
\pgfsys@useobject{currentmarker}{}%
\end{pgfscope}%
\begin{pgfscope}%
\pgfsys@transformshift{6.502223in}{2.585475in}%
\pgfsys@useobject{currentmarker}{}%
\end{pgfscope}%
\begin{pgfscope}%
\pgfsys@transformshift{5.957460in}{2.603048in}%
\pgfsys@useobject{currentmarker}{}%
\end{pgfscope}%
\begin{pgfscope}%
\pgfsys@transformshift{4.898926in}{2.584716in}%
\pgfsys@useobject{currentmarker}{}%
\end{pgfscope}%
\begin{pgfscope}%
\pgfsys@transformshift{6.005446in}{2.600292in}%
\pgfsys@useobject{currentmarker}{}%
\end{pgfscope}%
\begin{pgfscope}%
\pgfsys@transformshift{5.491316in}{2.604582in}%
\pgfsys@useobject{currentmarker}{}%
\end{pgfscope}%
\begin{pgfscope}%
\pgfsys@transformshift{5.995389in}{2.587179in}%
\pgfsys@useobject{currentmarker}{}%
\end{pgfscope}%
\begin{pgfscope}%
\pgfsys@transformshift{6.518320in}{2.600416in}%
\pgfsys@useobject{currentmarker}{}%
\end{pgfscope}%
\begin{pgfscope}%
\pgfsys@transformshift{6.807187in}{2.623589in}%
\pgfsys@useobject{currentmarker}{}%
\end{pgfscope}%
\begin{pgfscope}%
\pgfsys@transformshift{5.985424in}{2.603361in}%
\pgfsys@useobject{currentmarker}{}%
\end{pgfscope}%
\begin{pgfscope}%
\pgfsys@transformshift{6.424033in}{2.588606in}%
\pgfsys@useobject{currentmarker}{}%
\end{pgfscope}%
\begin{pgfscope}%
\pgfsys@transformshift{6.329208in}{2.597136in}%
\pgfsys@useobject{currentmarker}{}%
\end{pgfscope}%
\begin{pgfscope}%
\pgfsys@transformshift{5.102056in}{2.584329in}%
\pgfsys@useobject{currentmarker}{}%
\end{pgfscope}%
\begin{pgfscope}%
\pgfsys@transformshift{5.698022in}{2.602566in}%
\pgfsys@useobject{currentmarker}{}%
\end{pgfscope}%
\begin{pgfscope}%
\pgfsys@transformshift{6.520204in}{2.611652in}%
\pgfsys@useobject{currentmarker}{}%
\end{pgfscope}%
\begin{pgfscope}%
\pgfsys@transformshift{5.271053in}{2.585830in}%
\pgfsys@useobject{currentmarker}{}%
\end{pgfscope}%
\begin{pgfscope}%
\pgfsys@transformshift{6.034499in}{2.595196in}%
\pgfsys@useobject{currentmarker}{}%
\end{pgfscope}%
\begin{pgfscope}%
\pgfsys@transformshift{4.810980in}{2.615478in}%
\pgfsys@useobject{currentmarker}{}%
\end{pgfscope}%
\begin{pgfscope}%
\pgfsys@transformshift{6.672953in}{2.590775in}%
\pgfsys@useobject{currentmarker}{}%
\end{pgfscope}%
\begin{pgfscope}%
\pgfsys@transformshift{6.357385in}{2.601290in}%
\pgfsys@useobject{currentmarker}{}%
\end{pgfscope}%
\begin{pgfscope}%
\pgfsys@transformshift{5.338776in}{2.582704in}%
\pgfsys@useobject{currentmarker}{}%
\end{pgfscope}%
\begin{pgfscope}%
\pgfsys@transformshift{5.174332in}{2.585847in}%
\pgfsys@useobject{currentmarker}{}%
\end{pgfscope}%
\begin{pgfscope}%
\pgfsys@transformshift{6.020619in}{2.590550in}%
\pgfsys@useobject{currentmarker}{}%
\end{pgfscope}%
\begin{pgfscope}%
\pgfsys@transformshift{6.365473in}{2.631080in}%
\pgfsys@useobject{currentmarker}{}%
\end{pgfscope}%
\begin{pgfscope}%
\pgfsys@transformshift{5.333824in}{2.592267in}%
\pgfsys@useobject{currentmarker}{}%
\end{pgfscope}%
\begin{pgfscope}%
\pgfsys@transformshift{5.711472in}{2.590453in}%
\pgfsys@useobject{currentmarker}{}%
\end{pgfscope}%
\begin{pgfscope}%
\pgfsys@transformshift{4.879759in}{2.589954in}%
\pgfsys@useobject{currentmarker}{}%
\end{pgfscope}%
\begin{pgfscope}%
\pgfsys@transformshift{5.447793in}{2.595921in}%
\pgfsys@useobject{currentmarker}{}%
\end{pgfscope}%
\begin{pgfscope}%
\pgfsys@transformshift{5.693545in}{2.608252in}%
\pgfsys@useobject{currentmarker}{}%
\end{pgfscope}%
\begin{pgfscope}%
\pgfsys@transformshift{5.161359in}{2.601251in}%
\pgfsys@useobject{currentmarker}{}%
\end{pgfscope}%
\begin{pgfscope}%
\pgfsys@transformshift{4.953338in}{2.629300in}%
\pgfsys@useobject{currentmarker}{}%
\end{pgfscope}%
\begin{pgfscope}%
\pgfsys@transformshift{6.211178in}{2.601616in}%
\pgfsys@useobject{currentmarker}{}%
\end{pgfscope}%
\begin{pgfscope}%
\pgfsys@transformshift{6.536367in}{2.620134in}%
\pgfsys@useobject{currentmarker}{}%
\end{pgfscope}%
\begin{pgfscope}%
\pgfsys@transformshift{4.905867in}{2.602560in}%
\pgfsys@useobject{currentmarker}{}%
\end{pgfscope}%
\begin{pgfscope}%
\pgfsys@transformshift{5.567283in}{2.607496in}%
\pgfsys@useobject{currentmarker}{}%
\end{pgfscope}%
\begin{pgfscope}%
\pgfsys@transformshift{6.386715in}{2.576737in}%
\pgfsys@useobject{currentmarker}{}%
\end{pgfscope}%
\begin{pgfscope}%
\pgfsys@transformshift{6.591131in}{2.596500in}%
\pgfsys@useobject{currentmarker}{}%
\end{pgfscope}%
\begin{pgfscope}%
\pgfsys@transformshift{5.776565in}{2.597372in}%
\pgfsys@useobject{currentmarker}{}%
\end{pgfscope}%
\begin{pgfscope}%
\pgfsys@transformshift{5.065666in}{2.594760in}%
\pgfsys@useobject{currentmarker}{}%
\end{pgfscope}%
\begin{pgfscope}%
\pgfsys@transformshift{6.695059in}{2.566085in}%
\pgfsys@useobject{currentmarker}{}%
\end{pgfscope}%
\begin{pgfscope}%
\pgfsys@transformshift{5.286376in}{2.611625in}%
\pgfsys@useobject{currentmarker}{}%
\end{pgfscope}%
\begin{pgfscope}%
\pgfsys@transformshift{6.238274in}{2.592596in}%
\pgfsys@useobject{currentmarker}{}%
\end{pgfscope}%
\begin{pgfscope}%
\pgfsys@transformshift{4.984721in}{2.591832in}%
\pgfsys@useobject{currentmarker}{}%
\end{pgfscope}%
\begin{pgfscope}%
\pgfsys@transformshift{5.014052in}{2.594577in}%
\pgfsys@useobject{currentmarker}{}%
\end{pgfscope}%
\begin{pgfscope}%
\pgfsys@transformshift{6.742064in}{2.594984in}%
\pgfsys@useobject{currentmarker}{}%
\end{pgfscope}%
\begin{pgfscope}%
\pgfsys@transformshift{6.139889in}{2.604517in}%
\pgfsys@useobject{currentmarker}{}%
\end{pgfscope}%
\begin{pgfscope}%
\pgfsys@transformshift{6.607394in}{2.579328in}%
\pgfsys@useobject{currentmarker}{}%
\end{pgfscope}%
\begin{pgfscope}%
\pgfsys@transformshift{5.171356in}{2.584432in}%
\pgfsys@useobject{currentmarker}{}%
\end{pgfscope}%
\begin{pgfscope}%
\pgfsys@transformshift{6.141168in}{2.602769in}%
\pgfsys@useobject{currentmarker}{}%
\end{pgfscope}%
\begin{pgfscope}%
\pgfsys@transformshift{5.924648in}{2.600592in}%
\pgfsys@useobject{currentmarker}{}%
\end{pgfscope}%
\begin{pgfscope}%
\pgfsys@transformshift{6.738663in}{2.592754in}%
\pgfsys@useobject{currentmarker}{}%
\end{pgfscope}%
\begin{pgfscope}%
\pgfsys@transformshift{6.240131in}{2.589793in}%
\pgfsys@useobject{currentmarker}{}%
\end{pgfscope}%
\begin{pgfscope}%
\pgfsys@transformshift{5.292518in}{2.585882in}%
\pgfsys@useobject{currentmarker}{}%
\end{pgfscope}%
\begin{pgfscope}%
\pgfsys@transformshift{5.328866in}{2.597292in}%
\pgfsys@useobject{currentmarker}{}%
\end{pgfscope}%
\begin{pgfscope}%
\pgfsys@transformshift{5.608555in}{2.616130in}%
\pgfsys@useobject{currentmarker}{}%
\end{pgfscope}%
\begin{pgfscope}%
\pgfsys@transformshift{5.378189in}{2.582525in}%
\pgfsys@useobject{currentmarker}{}%
\end{pgfscope}%
\begin{pgfscope}%
\pgfsys@transformshift{4.958693in}{2.621273in}%
\pgfsys@useobject{currentmarker}{}%
\end{pgfscope}%
\begin{pgfscope}%
\pgfsys@transformshift{6.136411in}{2.611075in}%
\pgfsys@useobject{currentmarker}{}%
\end{pgfscope}%
\begin{pgfscope}%
\pgfsys@transformshift{5.341476in}{2.591608in}%
\pgfsys@useobject{currentmarker}{}%
\end{pgfscope}%
\begin{pgfscope}%
\pgfsys@transformshift{5.134296in}{2.610160in}%
\pgfsys@useobject{currentmarker}{}%
\end{pgfscope}%
\begin{pgfscope}%
\pgfsys@transformshift{5.175606in}{2.576402in}%
\pgfsys@useobject{currentmarker}{}%
\end{pgfscope}%
\begin{pgfscope}%
\pgfsys@transformshift{5.067945in}{2.572413in}%
\pgfsys@useobject{currentmarker}{}%
\end{pgfscope}%
\begin{pgfscope}%
\pgfsys@transformshift{5.357004in}{2.579047in}%
\pgfsys@useobject{currentmarker}{}%
\end{pgfscope}%
\begin{pgfscope}%
\pgfsys@transformshift{6.378449in}{2.592257in}%
\pgfsys@useobject{currentmarker}{}%
\end{pgfscope}%
\begin{pgfscope}%
\pgfsys@transformshift{5.012661in}{2.593523in}%
\pgfsys@useobject{currentmarker}{}%
\end{pgfscope}%
\begin{pgfscope}%
\pgfsys@transformshift{4.673887in}{2.590675in}%
\pgfsys@useobject{currentmarker}{}%
\end{pgfscope}%
\begin{pgfscope}%
\pgfsys@transformshift{6.465848in}{2.599899in}%
\pgfsys@useobject{currentmarker}{}%
\end{pgfscope}%
\begin{pgfscope}%
\pgfsys@transformshift{6.440586in}{2.566352in}%
\pgfsys@useobject{currentmarker}{}%
\end{pgfscope}%
\begin{pgfscope}%
\pgfsys@transformshift{6.067511in}{2.605842in}%
\pgfsys@useobject{currentmarker}{}%
\end{pgfscope}%
\begin{pgfscope}%
\pgfsys@transformshift{6.800696in}{2.577102in}%
\pgfsys@useobject{currentmarker}{}%
\end{pgfscope}%
\begin{pgfscope}%
\pgfsys@transformshift{6.533146in}{2.588488in}%
\pgfsys@useobject{currentmarker}{}%
\end{pgfscope}%
\begin{pgfscope}%
\pgfsys@transformshift{6.648532in}{2.592210in}%
\pgfsys@useobject{currentmarker}{}%
\end{pgfscope}%
\begin{pgfscope}%
\pgfsys@transformshift{5.404206in}{2.581052in}%
\pgfsys@useobject{currentmarker}{}%
\end{pgfscope}%
\begin{pgfscope}%
\pgfsys@transformshift{5.032917in}{2.577827in}%
\pgfsys@useobject{currentmarker}{}%
\end{pgfscope}%
\begin{pgfscope}%
\pgfsys@transformshift{5.029929in}{2.583306in}%
\pgfsys@useobject{currentmarker}{}%
\end{pgfscope}%
\begin{pgfscope}%
\pgfsys@transformshift{5.693157in}{2.598179in}%
\pgfsys@useobject{currentmarker}{}%
\end{pgfscope}%
\begin{pgfscope}%
\pgfsys@transformshift{6.203668in}{2.607262in}%
\pgfsys@useobject{currentmarker}{}%
\end{pgfscope}%
\begin{pgfscope}%
\pgfsys@transformshift{4.653373in}{2.606876in}%
\pgfsys@useobject{currentmarker}{}%
\end{pgfscope}%
\begin{pgfscope}%
\pgfsys@transformshift{5.929985in}{2.583480in}%
\pgfsys@useobject{currentmarker}{}%
\end{pgfscope}%
\begin{pgfscope}%
\pgfsys@transformshift{6.585653in}{2.586526in}%
\pgfsys@useobject{currentmarker}{}%
\end{pgfscope}%
\begin{pgfscope}%
\pgfsys@transformshift{5.619603in}{2.584589in}%
\pgfsys@useobject{currentmarker}{}%
\end{pgfscope}%
\begin{pgfscope}%
\pgfsys@transformshift{5.915636in}{2.605158in}%
\pgfsys@useobject{currentmarker}{}%
\end{pgfscope}%
\begin{pgfscope}%
\pgfsys@transformshift{5.199697in}{2.567820in}%
\pgfsys@useobject{currentmarker}{}%
\end{pgfscope}%
\begin{pgfscope}%
\pgfsys@transformshift{5.434279in}{2.616995in}%
\pgfsys@useobject{currentmarker}{}%
\end{pgfscope}%
\begin{pgfscope}%
\pgfsys@transformshift{4.655634in}{2.591130in}%
\pgfsys@useobject{currentmarker}{}%
\end{pgfscope}%
\begin{pgfscope}%
\pgfsys@transformshift{5.286084in}{2.586646in}%
\pgfsys@useobject{currentmarker}{}%
\end{pgfscope}%
\begin{pgfscope}%
\pgfsys@transformshift{6.691202in}{2.603502in}%
\pgfsys@useobject{currentmarker}{}%
\end{pgfscope}%
\begin{pgfscope}%
\pgfsys@transformshift{6.499415in}{2.613906in}%
\pgfsys@useobject{currentmarker}{}%
\end{pgfscope}%
\begin{pgfscope}%
\pgfsys@transformshift{4.648564in}{2.603549in}%
\pgfsys@useobject{currentmarker}{}%
\end{pgfscope}%
\begin{pgfscope}%
\pgfsys@transformshift{6.153141in}{2.614208in}%
\pgfsys@useobject{currentmarker}{}%
\end{pgfscope}%
\begin{pgfscope}%
\pgfsys@transformshift{6.399546in}{2.600423in}%
\pgfsys@useobject{currentmarker}{}%
\end{pgfscope}%
\begin{pgfscope}%
\pgfsys@transformshift{6.394560in}{2.616864in}%
\pgfsys@useobject{currentmarker}{}%
\end{pgfscope}%
\begin{pgfscope}%
\pgfsys@transformshift{5.702699in}{2.617097in}%
\pgfsys@useobject{currentmarker}{}%
\end{pgfscope}%
\begin{pgfscope}%
\pgfsys@transformshift{6.459993in}{2.596914in}%
\pgfsys@useobject{currentmarker}{}%
\end{pgfscope}%
\begin{pgfscope}%
\pgfsys@transformshift{6.227458in}{2.620010in}%
\pgfsys@useobject{currentmarker}{}%
\end{pgfscope}%
\begin{pgfscope}%
\pgfsys@transformshift{5.596453in}{2.600879in}%
\pgfsys@useobject{currentmarker}{}%
\end{pgfscope}%
\begin{pgfscope}%
\pgfsys@transformshift{6.030663in}{2.599711in}%
\pgfsys@useobject{currentmarker}{}%
\end{pgfscope}%
\begin{pgfscope}%
\pgfsys@transformshift{5.653334in}{2.584576in}%
\pgfsys@useobject{currentmarker}{}%
\end{pgfscope}%
\begin{pgfscope}%
\pgfsys@transformshift{5.407642in}{2.594072in}%
\pgfsys@useobject{currentmarker}{}%
\end{pgfscope}%
\begin{pgfscope}%
\pgfsys@transformshift{4.957352in}{2.610078in}%
\pgfsys@useobject{currentmarker}{}%
\end{pgfscope}%
\begin{pgfscope}%
\pgfsys@transformshift{6.661724in}{2.578703in}%
\pgfsys@useobject{currentmarker}{}%
\end{pgfscope}%
\begin{pgfscope}%
\pgfsys@transformshift{6.158704in}{2.608157in}%
\pgfsys@useobject{currentmarker}{}%
\end{pgfscope}%
\begin{pgfscope}%
\pgfsys@transformshift{5.959937in}{2.603639in}%
\pgfsys@useobject{currentmarker}{}%
\end{pgfscope}%
\begin{pgfscope}%
\pgfsys@transformshift{5.954420in}{2.602573in}%
\pgfsys@useobject{currentmarker}{}%
\end{pgfscope}%
\begin{pgfscope}%
\pgfsys@transformshift{6.489121in}{2.565280in}%
\pgfsys@useobject{currentmarker}{}%
\end{pgfscope}%
\begin{pgfscope}%
\pgfsys@transformshift{6.617208in}{2.596857in}%
\pgfsys@useobject{currentmarker}{}%
\end{pgfscope}%
\begin{pgfscope}%
\pgfsys@transformshift{5.561409in}{2.608137in}%
\pgfsys@useobject{currentmarker}{}%
\end{pgfscope}%
\begin{pgfscope}%
\pgfsys@transformshift{4.715343in}{2.586955in}%
\pgfsys@useobject{currentmarker}{}%
\end{pgfscope}%
\begin{pgfscope}%
\pgfsys@transformshift{5.707538in}{2.599762in}%
\pgfsys@useobject{currentmarker}{}%
\end{pgfscope}%
\begin{pgfscope}%
\pgfsys@transformshift{5.119724in}{2.562702in}%
\pgfsys@useobject{currentmarker}{}%
\end{pgfscope}%
\begin{pgfscope}%
\pgfsys@transformshift{6.047453in}{2.608484in}%
\pgfsys@useobject{currentmarker}{}%
\end{pgfscope}%
\begin{pgfscope}%
\pgfsys@transformshift{4.753153in}{2.588464in}%
\pgfsys@useobject{currentmarker}{}%
\end{pgfscope}%
\begin{pgfscope}%
\pgfsys@transformshift{6.075965in}{2.611486in}%
\pgfsys@useobject{currentmarker}{}%
\end{pgfscope}%
\begin{pgfscope}%
\pgfsys@transformshift{5.813508in}{2.570743in}%
\pgfsys@useobject{currentmarker}{}%
\end{pgfscope}%
\begin{pgfscope}%
\pgfsys@transformshift{5.605253in}{2.608797in}%
\pgfsys@useobject{currentmarker}{}%
\end{pgfscope}%
\begin{pgfscope}%
\pgfsys@transformshift{6.465858in}{2.594851in}%
\pgfsys@useobject{currentmarker}{}%
\end{pgfscope}%
\begin{pgfscope}%
\pgfsys@transformshift{6.141954in}{2.610529in}%
\pgfsys@useobject{currentmarker}{}%
\end{pgfscope}%
\begin{pgfscope}%
\pgfsys@transformshift{5.950241in}{2.584512in}%
\pgfsys@useobject{currentmarker}{}%
\end{pgfscope}%
\begin{pgfscope}%
\pgfsys@transformshift{4.968385in}{2.589124in}%
\pgfsys@useobject{currentmarker}{}%
\end{pgfscope}%
\begin{pgfscope}%
\pgfsys@transformshift{6.733233in}{2.632074in}%
\pgfsys@useobject{currentmarker}{}%
\end{pgfscope}%
\begin{pgfscope}%
\pgfsys@transformshift{5.753494in}{2.585950in}%
\pgfsys@useobject{currentmarker}{}%
\end{pgfscope}%
\begin{pgfscope}%
\pgfsys@transformshift{5.106073in}{2.587810in}%
\pgfsys@useobject{currentmarker}{}%
\end{pgfscope}%
\begin{pgfscope}%
\pgfsys@transformshift{6.288571in}{2.594256in}%
\pgfsys@useobject{currentmarker}{}%
\end{pgfscope}%
\begin{pgfscope}%
\pgfsys@transformshift{5.248928in}{2.583058in}%
\pgfsys@useobject{currentmarker}{}%
\end{pgfscope}%
\begin{pgfscope}%
\pgfsys@transformshift{6.794543in}{2.616276in}%
\pgfsys@useobject{currentmarker}{}%
\end{pgfscope}%
\begin{pgfscope}%
\pgfsys@transformshift{4.714850in}{2.624302in}%
\pgfsys@useobject{currentmarker}{}%
\end{pgfscope}%
\begin{pgfscope}%
\pgfsys@transformshift{6.726333in}{2.587839in}%
\pgfsys@useobject{currentmarker}{}%
\end{pgfscope}%
\begin{pgfscope}%
\pgfsys@transformshift{5.021142in}{2.572496in}%
\pgfsys@useobject{currentmarker}{}%
\end{pgfscope}%
\begin{pgfscope}%
\pgfsys@transformshift{5.706270in}{2.615998in}%
\pgfsys@useobject{currentmarker}{}%
\end{pgfscope}%
\begin{pgfscope}%
\pgfsys@transformshift{5.286759in}{2.615044in}%
\pgfsys@useobject{currentmarker}{}%
\end{pgfscope}%
\begin{pgfscope}%
\pgfsys@transformshift{5.878880in}{2.610740in}%
\pgfsys@useobject{currentmarker}{}%
\end{pgfscope}%
\begin{pgfscope}%
\pgfsys@transformshift{6.808102in}{2.615766in}%
\pgfsys@useobject{currentmarker}{}%
\end{pgfscope}%
\begin{pgfscope}%
\pgfsys@transformshift{6.239539in}{2.585386in}%
\pgfsys@useobject{currentmarker}{}%
\end{pgfscope}%
\begin{pgfscope}%
\pgfsys@transformshift{5.951637in}{2.595878in}%
\pgfsys@useobject{currentmarker}{}%
\end{pgfscope}%
\begin{pgfscope}%
\pgfsys@transformshift{6.153955in}{2.577630in}%
\pgfsys@useobject{currentmarker}{}%
\end{pgfscope}%
\begin{pgfscope}%
\pgfsys@transformshift{6.269470in}{2.577119in}%
\pgfsys@useobject{currentmarker}{}%
\end{pgfscope}%
\begin{pgfscope}%
\pgfsys@transformshift{5.640078in}{2.613133in}%
\pgfsys@useobject{currentmarker}{}%
\end{pgfscope}%
\begin{pgfscope}%
\pgfsys@transformshift{5.170880in}{2.588322in}%
\pgfsys@useobject{currentmarker}{}%
\end{pgfscope}%
\begin{pgfscope}%
\pgfsys@transformshift{5.407516in}{2.632631in}%
\pgfsys@useobject{currentmarker}{}%
\end{pgfscope}%
\begin{pgfscope}%
\pgfsys@transformshift{5.070229in}{2.599520in}%
\pgfsys@useobject{currentmarker}{}%
\end{pgfscope}%
\begin{pgfscope}%
\pgfsys@transformshift{6.417282in}{2.588794in}%
\pgfsys@useobject{currentmarker}{}%
\end{pgfscope}%
\begin{pgfscope}%
\pgfsys@transformshift{6.581969in}{2.610383in}%
\pgfsys@useobject{currentmarker}{}%
\end{pgfscope}%
\begin{pgfscope}%
\pgfsys@transformshift{6.689267in}{2.625447in}%
\pgfsys@useobject{currentmarker}{}%
\end{pgfscope}%
\begin{pgfscope}%
\pgfsys@transformshift{6.368745in}{2.605703in}%
\pgfsys@useobject{currentmarker}{}%
\end{pgfscope}%
\begin{pgfscope}%
\pgfsys@transformshift{6.647915in}{2.580137in}%
\pgfsys@useobject{currentmarker}{}%
\end{pgfscope}%
\begin{pgfscope}%
\pgfsys@transformshift{5.228298in}{2.603975in}%
\pgfsys@useobject{currentmarker}{}%
\end{pgfscope}%
\begin{pgfscope}%
\pgfsys@transformshift{5.893218in}{2.619055in}%
\pgfsys@useobject{currentmarker}{}%
\end{pgfscope}%
\begin{pgfscope}%
\pgfsys@transformshift{6.331984in}{2.591004in}%
\pgfsys@useobject{currentmarker}{}%
\end{pgfscope}%
\begin{pgfscope}%
\pgfsys@transformshift{6.773956in}{2.586795in}%
\pgfsys@useobject{currentmarker}{}%
\end{pgfscope}%
\begin{pgfscope}%
\pgfsys@transformshift{5.021151in}{2.613664in}%
\pgfsys@useobject{currentmarker}{}%
\end{pgfscope}%
\end{pgfscope}%
\begin{pgfscope}%
\pgfpathmoveto{\pgfqpoint{5.728000in}{2.052407in}}%
\pgfpathcurveto{\pgfqpoint{6.017602in}{2.052407in}}{\pgfqpoint{6.295381in}{2.109937in}}{\pgfqpoint{6.500161in}{2.212326in}}%
\pgfpathcurveto{\pgfqpoint{6.704940in}{2.314716in}}{\pgfqpoint{6.820000in}{2.453606in}}{\pgfqpoint{6.820000in}{2.598407in}}%
\pgfpathcurveto{\pgfqpoint{6.820000in}{2.743208in}}{\pgfqpoint{6.704940in}{2.882097in}}{\pgfqpoint{6.500161in}{2.984487in}}%
\pgfpathcurveto{\pgfqpoint{6.295381in}{3.086877in}}{\pgfqpoint{6.017602in}{3.144407in}}{\pgfqpoint{5.728000in}{3.144407in}}%
\pgfpathcurveto{\pgfqpoint{5.438398in}{3.144407in}}{\pgfqpoint{5.160619in}{3.086877in}}{\pgfqpoint{4.955839in}{2.984487in}}%
\pgfpathcurveto{\pgfqpoint{4.751060in}{2.882097in}}{\pgfqpoint{4.636000in}{2.743208in}}{\pgfqpoint{4.636000in}{2.598407in}}%
\pgfpathcurveto{\pgfqpoint{4.636000in}{2.453606in}}{\pgfqpoint{4.751060in}{2.314716in}}{\pgfqpoint{4.955839in}{2.212326in}}%
\pgfpathcurveto{\pgfqpoint{5.160619in}{2.109937in}}{\pgfqpoint{5.438398in}{2.052407in}}{\pgfqpoint{5.728000in}{2.052407in}}%
\pgfpathlineto{\pgfqpoint{5.728000in}{2.052407in}}%
\pgfpathclose%
\pgfusepath{clip}%
\pgfsetrectcap%
\pgfsetroundjoin%
\pgfsetlinewidth{0.451687pt}%
\definecolor{currentstroke}{rgb}{0.690196,0.690196,0.690196}%
\pgfsetstrokecolor{currentstroke}%
\pgfsetstrokeopacity{0.250000}%
\pgfsetdash{}{0pt}%
\pgfpathmoveto{\pgfqpoint{5.342921in}{2.103702in}}%
\pgfpathlineto{\pgfqpoint{5.311157in}{2.113059in}}%
\pgfpathlineto{\pgfqpoint{5.280700in}{2.122919in}}%
\pgfpathlineto{\pgfqpoint{5.252255in}{2.132965in}}%
\pgfpathlineto{\pgfqpoint{5.225090in}{2.143363in}}%
\pgfpathlineto{\pgfqpoint{5.199134in}{2.154082in}}%
\pgfpathlineto{\pgfqpoint{5.174319in}{2.165102in}}%
\pgfpathlineto{\pgfqpoint{5.150585in}{2.176401in}}%
\pgfpathlineto{\pgfqpoint{5.127881in}{2.187962in}}%
\pgfpathlineto{\pgfqpoint{5.106160in}{2.199772in}}%
\pgfpathlineto{\pgfqpoint{5.085384in}{2.211817in}}%
\pgfpathlineto{\pgfqpoint{5.065517in}{2.224083in}}%
\pgfpathlineto{\pgfqpoint{5.046532in}{2.236561in}}%
\pgfpathlineto{\pgfqpoint{5.028403in}{2.249238in}}%
\pgfpathlineto{\pgfqpoint{5.011106in}{2.262105in}}%
\pgfpathlineto{\pgfqpoint{4.994426in}{2.275314in}}%
\pgfpathlineto{\pgfqpoint{4.978806in}{2.288486in}}%
\pgfpathlineto{\pgfqpoint{4.963948in}{2.301831in}}%
\pgfpathlineto{\pgfqpoint{4.949847in}{2.315338in}}%
\pgfpathlineto{\pgfqpoint{4.936494in}{2.328996in}}%
\pgfpathlineto{\pgfqpoint{4.923882in}{2.342794in}}%
\pgfpathlineto{\pgfqpoint{4.912004in}{2.356724in}}%
\pgfpathlineto{\pgfqpoint{4.900855in}{2.370777in}}%
\pgfpathlineto{\pgfqpoint{4.890428in}{2.384944in}}%
\pgfpathlineto{\pgfqpoint{4.880718in}{2.399216in}}%
\pgfpathlineto{\pgfqpoint{4.871720in}{2.413586in}}%
\pgfpathlineto{\pgfqpoint{4.863430in}{2.428046in}}%
\pgfpathlineto{\pgfqpoint{4.855844in}{2.442588in}}%
\pgfpathlineto{\pgfqpoint{4.848957in}{2.457204in}}%
\pgfpathlineto{\pgfqpoint{4.842768in}{2.471888in}}%
\pgfpathlineto{\pgfqpoint{4.837272in}{2.486632in}}%
\pgfpathlineto{\pgfqpoint{4.832469in}{2.501429in}}%
\pgfpathlineto{\pgfqpoint{4.828314in}{2.516436in}}%
\pgfpathlineto{\pgfqpoint{4.824911in}{2.531243in}}%
\pgfpathlineto{\pgfqpoint{4.822184in}{2.546108in}}%
\pgfpathlineto{\pgfqpoint{4.820136in}{2.561018in}}%
\pgfpathlineto{\pgfqpoint{4.818769in}{2.575961in}}%
\pgfpathlineto{\pgfqpoint{4.818085in}{2.590923in}}%
\pgfpathlineto{\pgfqpoint{4.818085in}{2.605891in}}%
\pgfpathlineto{\pgfqpoint{4.818769in}{2.620853in}}%
\pgfpathlineto{\pgfqpoint{4.820136in}{2.635795in}}%
\pgfpathlineto{\pgfqpoint{4.822184in}{2.650705in}}%
\pgfpathlineto{\pgfqpoint{4.824911in}{2.665570in}}%
\pgfpathlineto{\pgfqpoint{4.828314in}{2.680378in}}%
\pgfpathlineto{\pgfqpoint{4.832469in}{2.695384in}}%
\pgfpathlineto{\pgfqpoint{4.837272in}{2.710181in}}%
\pgfpathlineto{\pgfqpoint{4.842768in}{2.724925in}}%
\pgfpathlineto{\pgfqpoint{4.848957in}{2.739609in}}%
\pgfpathlineto{\pgfqpoint{4.855844in}{2.754226in}}%
\pgfpathlineto{\pgfqpoint{4.863430in}{2.768768in}}%
\pgfpathlineto{\pgfqpoint{4.871720in}{2.783228in}}%
\pgfpathlineto{\pgfqpoint{4.880718in}{2.797598in}}%
\pgfpathlineto{\pgfqpoint{4.890428in}{2.811870in}}%
\pgfpathlineto{\pgfqpoint{4.900855in}{2.826037in}}%
\pgfpathlineto{\pgfqpoint{4.912004in}{2.840089in}}%
\pgfpathlineto{\pgfqpoint{4.923882in}{2.854019in}}%
\pgfpathlineto{\pgfqpoint{4.936494in}{2.867818in}}%
\pgfpathlineto{\pgfqpoint{4.949847in}{2.881476in}}%
\pgfpathlineto{\pgfqpoint{4.963948in}{2.894982in}}%
\pgfpathlineto{\pgfqpoint{4.978806in}{2.908327in}}%
\pgfpathlineto{\pgfqpoint{4.994426in}{2.921500in}}%
\pgfpathlineto{\pgfqpoint{5.011106in}{2.934708in}}%
\pgfpathlineto{\pgfqpoint{5.028403in}{2.947576in}}%
\pgfpathlineto{\pgfqpoint{5.046532in}{2.960253in}}%
\pgfpathlineto{\pgfqpoint{5.065517in}{2.972730in}}%
\pgfpathlineto{\pgfqpoint{5.085384in}{2.984997in}}%
\pgfpathlineto{\pgfqpoint{5.106160in}{2.997041in}}%
\pgfpathlineto{\pgfqpoint{5.127881in}{3.008851in}}%
\pgfpathlineto{\pgfqpoint{5.150585in}{3.020413in}}%
\pgfpathlineto{\pgfqpoint{5.174319in}{3.031712in}}%
\pgfpathlineto{\pgfqpoint{5.199134in}{3.042731in}}%
\pgfpathlineto{\pgfqpoint{5.225090in}{3.053451in}}%
\pgfpathlineto{\pgfqpoint{5.252255in}{3.063848in}}%
\pgfpathlineto{\pgfqpoint{5.280700in}{3.073894in}}%
\pgfpathlineto{\pgfqpoint{5.311157in}{3.083755in}}%
\pgfpathlineto{\pgfqpoint{5.342921in}{3.093112in}}%
\pgfusepath{stroke}%
\end{pgfscope}%
\begin{pgfscope}%
\pgfpathmoveto{\pgfqpoint{5.728000in}{2.052407in}}%
\pgfpathcurveto{\pgfqpoint{6.017602in}{2.052407in}}{\pgfqpoint{6.295381in}{2.109937in}}{\pgfqpoint{6.500161in}{2.212326in}}%
\pgfpathcurveto{\pgfqpoint{6.704940in}{2.314716in}}{\pgfqpoint{6.820000in}{2.453606in}}{\pgfqpoint{6.820000in}{2.598407in}}%
\pgfpathcurveto{\pgfqpoint{6.820000in}{2.743208in}}{\pgfqpoint{6.704940in}{2.882097in}}{\pgfqpoint{6.500161in}{2.984487in}}%
\pgfpathcurveto{\pgfqpoint{6.295381in}{3.086877in}}{\pgfqpoint{6.017602in}{3.144407in}}{\pgfqpoint{5.728000in}{3.144407in}}%
\pgfpathcurveto{\pgfqpoint{5.438398in}{3.144407in}}{\pgfqpoint{5.160619in}{3.086877in}}{\pgfqpoint{4.955839in}{2.984487in}}%
\pgfpathcurveto{\pgfqpoint{4.751060in}{2.882097in}}{\pgfqpoint{4.636000in}{2.743208in}}{\pgfqpoint{4.636000in}{2.598407in}}%
\pgfpathcurveto{\pgfqpoint{4.636000in}{2.453606in}}{\pgfqpoint{4.751060in}{2.314716in}}{\pgfqpoint{4.955839in}{2.212326in}}%
\pgfpathcurveto{\pgfqpoint{5.160619in}{2.109937in}}{\pgfqpoint{5.438398in}{2.052407in}}{\pgfqpoint{5.728000in}{2.052407in}}%
\pgfpathlineto{\pgfqpoint{5.728000in}{2.052407in}}%
\pgfpathclose%
\pgfusepath{clip}%
\pgfsetrectcap%
\pgfsetroundjoin%
\pgfsetlinewidth{0.451687pt}%
\definecolor{currentstroke}{rgb}{0.690196,0.690196,0.690196}%
\pgfsetstrokecolor{currentstroke}%
\pgfsetstrokeopacity{0.250000}%
\pgfsetdash{}{0pt}%
\pgfpathmoveto{\pgfqpoint{5.419937in}{2.103702in}}%
\pgfpathlineto{\pgfqpoint{5.394525in}{2.113059in}}%
\pgfpathlineto{\pgfqpoint{5.370160in}{2.122919in}}%
\pgfpathlineto{\pgfqpoint{5.347404in}{2.132965in}}%
\pgfpathlineto{\pgfqpoint{5.325672in}{2.143363in}}%
\pgfpathlineto{\pgfqpoint{5.304907in}{2.154082in}}%
\pgfpathlineto{\pgfqpoint{5.285055in}{2.165102in}}%
\pgfpathlineto{\pgfqpoint{5.266068in}{2.176401in}}%
\pgfpathlineto{\pgfqpoint{5.247905in}{2.187962in}}%
\pgfpathlineto{\pgfqpoint{5.230528in}{2.199772in}}%
\pgfpathlineto{\pgfqpoint{5.213907in}{2.211817in}}%
\pgfpathlineto{\pgfqpoint{5.198014in}{2.224083in}}%
\pgfpathlineto{\pgfqpoint{5.182826in}{2.236561in}}%
\pgfpathlineto{\pgfqpoint{5.168322in}{2.249238in}}%
\pgfpathlineto{\pgfqpoint{5.154485in}{2.262105in}}%
\pgfpathlineto{\pgfqpoint{5.141141in}{2.275314in}}%
\pgfpathlineto{\pgfqpoint{5.128644in}{2.288486in}}%
\pgfpathlineto{\pgfqpoint{5.116759in}{2.301831in}}%
\pgfpathlineto{\pgfqpoint{5.105477in}{2.315338in}}%
\pgfpathlineto{\pgfqpoint{5.094795in}{2.328996in}}%
\pgfpathlineto{\pgfqpoint{5.084705in}{2.342794in}}%
\pgfpathlineto{\pgfqpoint{5.075204in}{2.356724in}}%
\pgfpathlineto{\pgfqpoint{5.066284in}{2.370777in}}%
\pgfpathlineto{\pgfqpoint{5.057943in}{2.384944in}}%
\pgfpathlineto{\pgfqpoint{5.050175in}{2.399216in}}%
\pgfpathlineto{\pgfqpoint{5.042976in}{2.413586in}}%
\pgfpathlineto{\pgfqpoint{5.036344in}{2.428046in}}%
\pgfpathlineto{\pgfqpoint{5.030275in}{2.442588in}}%
\pgfpathlineto{\pgfqpoint{5.024766in}{2.457204in}}%
\pgfpathlineto{\pgfqpoint{5.019814in}{2.471888in}}%
\pgfpathlineto{\pgfqpoint{5.015418in}{2.486632in}}%
\pgfpathlineto{\pgfqpoint{5.011575in}{2.501429in}}%
\pgfpathlineto{\pgfqpoint{5.008251in}{2.516436in}}%
\pgfpathlineto{\pgfqpoint{5.005529in}{2.531243in}}%
\pgfpathlineto{\pgfqpoint{5.003347in}{2.546108in}}%
\pgfpathlineto{\pgfqpoint{5.001709in}{2.561018in}}%
\pgfpathlineto{\pgfqpoint{5.000615in}{2.575961in}}%
\pgfpathlineto{\pgfqpoint{5.000068in}{2.590923in}}%
\pgfpathlineto{\pgfqpoint{5.000068in}{2.605891in}}%
\pgfpathlineto{\pgfqpoint{5.000615in}{2.620853in}}%
\pgfpathlineto{\pgfqpoint{5.001709in}{2.635795in}}%
\pgfpathlineto{\pgfqpoint{5.003347in}{2.650705in}}%
\pgfpathlineto{\pgfqpoint{5.005529in}{2.665570in}}%
\pgfpathlineto{\pgfqpoint{5.008251in}{2.680378in}}%
\pgfpathlineto{\pgfqpoint{5.011575in}{2.695384in}}%
\pgfpathlineto{\pgfqpoint{5.015418in}{2.710181in}}%
\pgfpathlineto{\pgfqpoint{5.019814in}{2.724925in}}%
\pgfpathlineto{\pgfqpoint{5.024766in}{2.739609in}}%
\pgfpathlineto{\pgfqpoint{5.030275in}{2.754226in}}%
\pgfpathlineto{\pgfqpoint{5.036344in}{2.768768in}}%
\pgfpathlineto{\pgfqpoint{5.042976in}{2.783228in}}%
\pgfpathlineto{\pgfqpoint{5.050175in}{2.797598in}}%
\pgfpathlineto{\pgfqpoint{5.057943in}{2.811870in}}%
\pgfpathlineto{\pgfqpoint{5.066284in}{2.826037in}}%
\pgfpathlineto{\pgfqpoint{5.075204in}{2.840089in}}%
\pgfpathlineto{\pgfqpoint{5.084705in}{2.854019in}}%
\pgfpathlineto{\pgfqpoint{5.094795in}{2.867818in}}%
\pgfpathlineto{\pgfqpoint{5.105477in}{2.881476in}}%
\pgfpathlineto{\pgfqpoint{5.116759in}{2.894982in}}%
\pgfpathlineto{\pgfqpoint{5.128644in}{2.908327in}}%
\pgfpathlineto{\pgfqpoint{5.141141in}{2.921500in}}%
\pgfpathlineto{\pgfqpoint{5.154485in}{2.934708in}}%
\pgfpathlineto{\pgfqpoint{5.168322in}{2.947576in}}%
\pgfpathlineto{\pgfqpoint{5.182826in}{2.960253in}}%
\pgfpathlineto{\pgfqpoint{5.198014in}{2.972730in}}%
\pgfpathlineto{\pgfqpoint{5.213907in}{2.984997in}}%
\pgfpathlineto{\pgfqpoint{5.230528in}{2.997041in}}%
\pgfpathlineto{\pgfqpoint{5.247905in}{3.008851in}}%
\pgfpathlineto{\pgfqpoint{5.266068in}{3.020413in}}%
\pgfpathlineto{\pgfqpoint{5.285055in}{3.031712in}}%
\pgfpathlineto{\pgfqpoint{5.304907in}{3.042731in}}%
\pgfpathlineto{\pgfqpoint{5.325672in}{3.053451in}}%
\pgfpathlineto{\pgfqpoint{5.347404in}{3.063848in}}%
\pgfpathlineto{\pgfqpoint{5.370160in}{3.073894in}}%
\pgfpathlineto{\pgfqpoint{5.394525in}{3.083755in}}%
\pgfpathlineto{\pgfqpoint{5.419937in}{3.093112in}}%
\pgfusepath{stroke}%
\end{pgfscope}%
\begin{pgfscope}%
\pgfpathmoveto{\pgfqpoint{5.728000in}{2.052407in}}%
\pgfpathcurveto{\pgfqpoint{6.017602in}{2.052407in}}{\pgfqpoint{6.295381in}{2.109937in}}{\pgfqpoint{6.500161in}{2.212326in}}%
\pgfpathcurveto{\pgfqpoint{6.704940in}{2.314716in}}{\pgfqpoint{6.820000in}{2.453606in}}{\pgfqpoint{6.820000in}{2.598407in}}%
\pgfpathcurveto{\pgfqpoint{6.820000in}{2.743208in}}{\pgfqpoint{6.704940in}{2.882097in}}{\pgfqpoint{6.500161in}{2.984487in}}%
\pgfpathcurveto{\pgfqpoint{6.295381in}{3.086877in}}{\pgfqpoint{6.017602in}{3.144407in}}{\pgfqpoint{5.728000in}{3.144407in}}%
\pgfpathcurveto{\pgfqpoint{5.438398in}{3.144407in}}{\pgfqpoint{5.160619in}{3.086877in}}{\pgfqpoint{4.955839in}{2.984487in}}%
\pgfpathcurveto{\pgfqpoint{4.751060in}{2.882097in}}{\pgfqpoint{4.636000in}{2.743208in}}{\pgfqpoint{4.636000in}{2.598407in}}%
\pgfpathcurveto{\pgfqpoint{4.636000in}{2.453606in}}{\pgfqpoint{4.751060in}{2.314716in}}{\pgfqpoint{4.955839in}{2.212326in}}%
\pgfpathcurveto{\pgfqpoint{5.160619in}{2.109937in}}{\pgfqpoint{5.438398in}{2.052407in}}{\pgfqpoint{5.728000in}{2.052407in}}%
\pgfpathlineto{\pgfqpoint{5.728000in}{2.052407in}}%
\pgfpathclose%
\pgfusepath{clip}%
\pgfsetrectcap%
\pgfsetroundjoin%
\pgfsetlinewidth{0.451687pt}%
\definecolor{currentstroke}{rgb}{0.690196,0.690196,0.690196}%
\pgfsetstrokecolor{currentstroke}%
\pgfsetstrokeopacity{0.250000}%
\pgfsetdash{}{0pt}%
\pgfpathmoveto{\pgfqpoint{5.496953in}{2.103702in}}%
\pgfpathlineto{\pgfqpoint{5.477894in}{2.113059in}}%
\pgfpathlineto{\pgfqpoint{5.459620in}{2.122919in}}%
\pgfpathlineto{\pgfqpoint{5.442553in}{2.132965in}}%
\pgfpathlineto{\pgfqpoint{5.426254in}{2.143363in}}%
\pgfpathlineto{\pgfqpoint{5.410680in}{2.154082in}}%
\pgfpathlineto{\pgfqpoint{5.395791in}{2.165102in}}%
\pgfpathlineto{\pgfqpoint{5.381551in}{2.176401in}}%
\pgfpathlineto{\pgfqpoint{5.367929in}{2.187962in}}%
\pgfpathlineto{\pgfqpoint{5.354896in}{2.199772in}}%
\pgfpathlineto{\pgfqpoint{5.342430in}{2.211817in}}%
\pgfpathlineto{\pgfqpoint{5.330510in}{2.224083in}}%
\pgfpathlineto{\pgfqpoint{5.319119in}{2.236561in}}%
\pgfpathlineto{\pgfqpoint{5.308242in}{2.249238in}}%
\pgfpathlineto{\pgfqpoint{5.297864in}{2.262105in}}%
\pgfpathlineto{\pgfqpoint{5.287856in}{2.275314in}}%
\pgfpathlineto{\pgfqpoint{5.278483in}{2.288486in}}%
\pgfpathlineto{\pgfqpoint{5.269569in}{2.301831in}}%
\pgfpathlineto{\pgfqpoint{5.261108in}{2.315338in}}%
\pgfpathlineto{\pgfqpoint{5.253096in}{2.328996in}}%
\pgfpathlineto{\pgfqpoint{5.245529in}{2.342794in}}%
\pgfpathlineto{\pgfqpoint{5.238403in}{2.356724in}}%
\pgfpathlineto{\pgfqpoint{5.231713in}{2.370777in}}%
\pgfpathlineto{\pgfqpoint{5.225457in}{2.384944in}}%
\pgfpathlineto{\pgfqpoint{5.219631in}{2.399216in}}%
\pgfpathlineto{\pgfqpoint{5.214232in}{2.413586in}}%
\pgfpathlineto{\pgfqpoint{5.209258in}{2.428046in}}%
\pgfpathlineto{\pgfqpoint{5.204706in}{2.442588in}}%
\pgfpathlineto{\pgfqpoint{5.200574in}{2.457204in}}%
\pgfpathlineto{\pgfqpoint{5.196861in}{2.471888in}}%
\pgfpathlineto{\pgfqpoint{5.193563in}{2.486632in}}%
\pgfpathlineto{\pgfqpoint{5.190681in}{2.501429in}}%
\pgfpathlineto{\pgfqpoint{5.188188in}{2.516436in}}%
\pgfpathlineto{\pgfqpoint{5.186147in}{2.531243in}}%
\pgfpathlineto{\pgfqpoint{5.184510in}{2.546108in}}%
\pgfpathlineto{\pgfqpoint{5.183282in}{2.561018in}}%
\pgfpathlineto{\pgfqpoint{5.182462in}{2.575961in}}%
\pgfpathlineto{\pgfqpoint{5.182051in}{2.590923in}}%
\pgfpathlineto{\pgfqpoint{5.182051in}{2.605891in}}%
\pgfpathlineto{\pgfqpoint{5.182462in}{2.620853in}}%
\pgfpathlineto{\pgfqpoint{5.183282in}{2.635795in}}%
\pgfpathlineto{\pgfqpoint{5.184510in}{2.650705in}}%
\pgfpathlineto{\pgfqpoint{5.186147in}{2.665570in}}%
\pgfpathlineto{\pgfqpoint{5.188188in}{2.680378in}}%
\pgfpathlineto{\pgfqpoint{5.190681in}{2.695384in}}%
\pgfpathlineto{\pgfqpoint{5.193563in}{2.710181in}}%
\pgfpathlineto{\pgfqpoint{5.196861in}{2.724925in}}%
\pgfpathlineto{\pgfqpoint{5.200574in}{2.739609in}}%
\pgfpathlineto{\pgfqpoint{5.204706in}{2.754226in}}%
\pgfpathlineto{\pgfqpoint{5.209258in}{2.768768in}}%
\pgfpathlineto{\pgfqpoint{5.214232in}{2.783228in}}%
\pgfpathlineto{\pgfqpoint{5.219631in}{2.797598in}}%
\pgfpathlineto{\pgfqpoint{5.225457in}{2.811870in}}%
\pgfpathlineto{\pgfqpoint{5.231713in}{2.826037in}}%
\pgfpathlineto{\pgfqpoint{5.238403in}{2.840089in}}%
\pgfpathlineto{\pgfqpoint{5.245529in}{2.854019in}}%
\pgfpathlineto{\pgfqpoint{5.253096in}{2.867818in}}%
\pgfpathlineto{\pgfqpoint{5.261108in}{2.881476in}}%
\pgfpathlineto{\pgfqpoint{5.269569in}{2.894982in}}%
\pgfpathlineto{\pgfqpoint{5.278483in}{2.908327in}}%
\pgfpathlineto{\pgfqpoint{5.287856in}{2.921500in}}%
\pgfpathlineto{\pgfqpoint{5.297864in}{2.934708in}}%
\pgfpathlineto{\pgfqpoint{5.308242in}{2.947576in}}%
\pgfpathlineto{\pgfqpoint{5.319119in}{2.960253in}}%
\pgfpathlineto{\pgfqpoint{5.330510in}{2.972730in}}%
\pgfpathlineto{\pgfqpoint{5.342430in}{2.984997in}}%
\pgfpathlineto{\pgfqpoint{5.354896in}{2.997041in}}%
\pgfpathlineto{\pgfqpoint{5.367929in}{3.008851in}}%
\pgfpathlineto{\pgfqpoint{5.381551in}{3.020413in}}%
\pgfpathlineto{\pgfqpoint{5.395791in}{3.031712in}}%
\pgfpathlineto{\pgfqpoint{5.410680in}{3.042731in}}%
\pgfpathlineto{\pgfqpoint{5.426254in}{3.053451in}}%
\pgfpathlineto{\pgfqpoint{5.442553in}{3.063848in}}%
\pgfpathlineto{\pgfqpoint{5.459620in}{3.073894in}}%
\pgfpathlineto{\pgfqpoint{5.477894in}{3.083755in}}%
\pgfpathlineto{\pgfqpoint{5.496953in}{3.093112in}}%
\pgfusepath{stroke}%
\end{pgfscope}%
\begin{pgfscope}%
\pgfpathmoveto{\pgfqpoint{5.728000in}{2.052407in}}%
\pgfpathcurveto{\pgfqpoint{6.017602in}{2.052407in}}{\pgfqpoint{6.295381in}{2.109937in}}{\pgfqpoint{6.500161in}{2.212326in}}%
\pgfpathcurveto{\pgfqpoint{6.704940in}{2.314716in}}{\pgfqpoint{6.820000in}{2.453606in}}{\pgfqpoint{6.820000in}{2.598407in}}%
\pgfpathcurveto{\pgfqpoint{6.820000in}{2.743208in}}{\pgfqpoint{6.704940in}{2.882097in}}{\pgfqpoint{6.500161in}{2.984487in}}%
\pgfpathcurveto{\pgfqpoint{6.295381in}{3.086877in}}{\pgfqpoint{6.017602in}{3.144407in}}{\pgfqpoint{5.728000in}{3.144407in}}%
\pgfpathcurveto{\pgfqpoint{5.438398in}{3.144407in}}{\pgfqpoint{5.160619in}{3.086877in}}{\pgfqpoint{4.955839in}{2.984487in}}%
\pgfpathcurveto{\pgfqpoint{4.751060in}{2.882097in}}{\pgfqpoint{4.636000in}{2.743208in}}{\pgfqpoint{4.636000in}{2.598407in}}%
\pgfpathcurveto{\pgfqpoint{4.636000in}{2.453606in}}{\pgfqpoint{4.751060in}{2.314716in}}{\pgfqpoint{4.955839in}{2.212326in}}%
\pgfpathcurveto{\pgfqpoint{5.160619in}{2.109937in}}{\pgfqpoint{5.438398in}{2.052407in}}{\pgfqpoint{5.728000in}{2.052407in}}%
\pgfpathlineto{\pgfqpoint{5.728000in}{2.052407in}}%
\pgfpathclose%
\pgfusepath{clip}%
\pgfsetrectcap%
\pgfsetroundjoin%
\pgfsetlinewidth{0.451687pt}%
\definecolor{currentstroke}{rgb}{0.690196,0.690196,0.690196}%
\pgfsetstrokecolor{currentstroke}%
\pgfsetstrokeopacity{0.250000}%
\pgfsetdash{}{0pt}%
\pgfpathmoveto{\pgfqpoint{5.573969in}{2.103702in}}%
\pgfpathlineto{\pgfqpoint{5.561263in}{2.113059in}}%
\pgfpathlineto{\pgfqpoint{5.549080in}{2.122919in}}%
\pgfpathlineto{\pgfqpoint{5.537702in}{2.132965in}}%
\pgfpathlineto{\pgfqpoint{5.526836in}{2.143363in}}%
\pgfpathlineto{\pgfqpoint{5.516453in}{2.154082in}}%
\pgfpathlineto{\pgfqpoint{5.506527in}{2.165102in}}%
\pgfpathlineto{\pgfqpoint{5.497034in}{2.176401in}}%
\pgfpathlineto{\pgfqpoint{5.487952in}{2.187962in}}%
\pgfpathlineto{\pgfqpoint{5.479264in}{2.199772in}}%
\pgfpathlineto{\pgfqpoint{5.470953in}{2.211817in}}%
\pgfpathlineto{\pgfqpoint{5.463007in}{2.224083in}}%
\pgfpathlineto{\pgfqpoint{5.455413in}{2.236561in}}%
\pgfpathlineto{\pgfqpoint{5.448161in}{2.249238in}}%
\pgfpathlineto{\pgfqpoint{5.441243in}{2.262105in}}%
\pgfpathlineto{\pgfqpoint{5.434570in}{2.275314in}}%
\pgfpathlineto{\pgfqpoint{5.428322in}{2.288486in}}%
\pgfpathlineto{\pgfqpoint{5.422379in}{2.301831in}}%
\pgfpathlineto{\pgfqpoint{5.416739in}{2.315338in}}%
\pgfpathlineto{\pgfqpoint{5.411398in}{2.328996in}}%
\pgfpathlineto{\pgfqpoint{5.406353in}{2.342794in}}%
\pgfpathlineto{\pgfqpoint{5.401602in}{2.356724in}}%
\pgfpathlineto{\pgfqpoint{5.397142in}{2.370777in}}%
\pgfpathlineto{\pgfqpoint{5.392971in}{2.384944in}}%
\pgfpathlineto{\pgfqpoint{5.389087in}{2.399216in}}%
\pgfpathlineto{\pgfqpoint{5.385488in}{2.413586in}}%
\pgfpathlineto{\pgfqpoint{5.382172in}{2.428046in}}%
\pgfpathlineto{\pgfqpoint{5.379137in}{2.442588in}}%
\pgfpathlineto{\pgfqpoint{5.376383in}{2.457204in}}%
\pgfpathlineto{\pgfqpoint{5.373907in}{2.471888in}}%
\pgfpathlineto{\pgfqpoint{5.371709in}{2.486632in}}%
\pgfpathlineto{\pgfqpoint{5.369788in}{2.501429in}}%
\pgfpathlineto{\pgfqpoint{5.368125in}{2.516436in}}%
\pgfpathlineto{\pgfqpoint{5.366764in}{2.531243in}}%
\pgfpathlineto{\pgfqpoint{5.365674in}{2.546108in}}%
\pgfpathlineto{\pgfqpoint{5.364854in}{2.561018in}}%
\pgfpathlineto{\pgfqpoint{5.364308in}{2.575961in}}%
\pgfpathlineto{\pgfqpoint{5.364034in}{2.590923in}}%
\pgfpathlineto{\pgfqpoint{5.364034in}{2.605891in}}%
\pgfpathlineto{\pgfqpoint{5.364308in}{2.620853in}}%
\pgfpathlineto{\pgfqpoint{5.364854in}{2.635795in}}%
\pgfpathlineto{\pgfqpoint{5.365674in}{2.650705in}}%
\pgfpathlineto{\pgfqpoint{5.366764in}{2.665570in}}%
\pgfpathlineto{\pgfqpoint{5.368125in}{2.680378in}}%
\pgfpathlineto{\pgfqpoint{5.369788in}{2.695384in}}%
\pgfpathlineto{\pgfqpoint{5.371709in}{2.710181in}}%
\pgfpathlineto{\pgfqpoint{5.373907in}{2.724925in}}%
\pgfpathlineto{\pgfqpoint{5.376383in}{2.739609in}}%
\pgfpathlineto{\pgfqpoint{5.379137in}{2.754226in}}%
\pgfpathlineto{\pgfqpoint{5.382172in}{2.768768in}}%
\pgfpathlineto{\pgfqpoint{5.385488in}{2.783228in}}%
\pgfpathlineto{\pgfqpoint{5.389087in}{2.797598in}}%
\pgfpathlineto{\pgfqpoint{5.392971in}{2.811870in}}%
\pgfpathlineto{\pgfqpoint{5.397142in}{2.826037in}}%
\pgfpathlineto{\pgfqpoint{5.401602in}{2.840089in}}%
\pgfpathlineto{\pgfqpoint{5.406353in}{2.854019in}}%
\pgfpathlineto{\pgfqpoint{5.411398in}{2.867818in}}%
\pgfpathlineto{\pgfqpoint{5.416739in}{2.881476in}}%
\pgfpathlineto{\pgfqpoint{5.422379in}{2.894982in}}%
\pgfpathlineto{\pgfqpoint{5.428322in}{2.908327in}}%
\pgfpathlineto{\pgfqpoint{5.434570in}{2.921500in}}%
\pgfpathlineto{\pgfqpoint{5.441243in}{2.934708in}}%
\pgfpathlineto{\pgfqpoint{5.448161in}{2.947576in}}%
\pgfpathlineto{\pgfqpoint{5.455413in}{2.960253in}}%
\pgfpathlineto{\pgfqpoint{5.463007in}{2.972730in}}%
\pgfpathlineto{\pgfqpoint{5.470953in}{2.984997in}}%
\pgfpathlineto{\pgfqpoint{5.479264in}{2.997041in}}%
\pgfpathlineto{\pgfqpoint{5.487952in}{3.008851in}}%
\pgfpathlineto{\pgfqpoint{5.497034in}{3.020413in}}%
\pgfpathlineto{\pgfqpoint{5.506527in}{3.031712in}}%
\pgfpathlineto{\pgfqpoint{5.516453in}{3.042731in}}%
\pgfpathlineto{\pgfqpoint{5.526836in}{3.053451in}}%
\pgfpathlineto{\pgfqpoint{5.537702in}{3.063848in}}%
\pgfpathlineto{\pgfqpoint{5.549080in}{3.073894in}}%
\pgfpathlineto{\pgfqpoint{5.561263in}{3.083755in}}%
\pgfpathlineto{\pgfqpoint{5.573969in}{3.093112in}}%
\pgfusepath{stroke}%
\end{pgfscope}%
\begin{pgfscope}%
\pgfpathmoveto{\pgfqpoint{5.728000in}{2.052407in}}%
\pgfpathcurveto{\pgfqpoint{6.017602in}{2.052407in}}{\pgfqpoint{6.295381in}{2.109937in}}{\pgfqpoint{6.500161in}{2.212326in}}%
\pgfpathcurveto{\pgfqpoint{6.704940in}{2.314716in}}{\pgfqpoint{6.820000in}{2.453606in}}{\pgfqpoint{6.820000in}{2.598407in}}%
\pgfpathcurveto{\pgfqpoint{6.820000in}{2.743208in}}{\pgfqpoint{6.704940in}{2.882097in}}{\pgfqpoint{6.500161in}{2.984487in}}%
\pgfpathcurveto{\pgfqpoint{6.295381in}{3.086877in}}{\pgfqpoint{6.017602in}{3.144407in}}{\pgfqpoint{5.728000in}{3.144407in}}%
\pgfpathcurveto{\pgfqpoint{5.438398in}{3.144407in}}{\pgfqpoint{5.160619in}{3.086877in}}{\pgfqpoint{4.955839in}{2.984487in}}%
\pgfpathcurveto{\pgfqpoint{4.751060in}{2.882097in}}{\pgfqpoint{4.636000in}{2.743208in}}{\pgfqpoint{4.636000in}{2.598407in}}%
\pgfpathcurveto{\pgfqpoint{4.636000in}{2.453606in}}{\pgfqpoint{4.751060in}{2.314716in}}{\pgfqpoint{4.955839in}{2.212326in}}%
\pgfpathcurveto{\pgfqpoint{5.160619in}{2.109937in}}{\pgfqpoint{5.438398in}{2.052407in}}{\pgfqpoint{5.728000in}{2.052407in}}%
\pgfpathlineto{\pgfqpoint{5.728000in}{2.052407in}}%
\pgfpathclose%
\pgfusepath{clip}%
\pgfsetrectcap%
\pgfsetroundjoin%
\pgfsetlinewidth{0.451687pt}%
\definecolor{currentstroke}{rgb}{0.690196,0.690196,0.690196}%
\pgfsetstrokecolor{currentstroke}%
\pgfsetstrokeopacity{0.250000}%
\pgfsetdash{}{0pt}%
\pgfpathmoveto{\pgfqpoint{5.650984in}{2.103702in}}%
\pgfpathlineto{\pgfqpoint{5.644631in}{2.113059in}}%
\pgfpathlineto{\pgfqpoint{5.638540in}{2.122919in}}%
\pgfpathlineto{\pgfqpoint{5.632851in}{2.132965in}}%
\pgfpathlineto{\pgfqpoint{5.627418in}{2.143363in}}%
\pgfpathlineto{\pgfqpoint{5.622227in}{2.154082in}}%
\pgfpathlineto{\pgfqpoint{5.617264in}{2.165102in}}%
\pgfpathlineto{\pgfqpoint{5.612517in}{2.176401in}}%
\pgfpathlineto{\pgfqpoint{5.607976in}{2.187962in}}%
\pgfpathlineto{\pgfqpoint{5.603632in}{2.199772in}}%
\pgfpathlineto{\pgfqpoint{5.599477in}{2.211817in}}%
\pgfpathlineto{\pgfqpoint{5.595503in}{2.224083in}}%
\pgfpathlineto{\pgfqpoint{5.591706in}{2.236561in}}%
\pgfpathlineto{\pgfqpoint{5.588081in}{2.249238in}}%
\pgfpathlineto{\pgfqpoint{5.584621in}{2.262105in}}%
\pgfpathlineto{\pgfqpoint{5.581285in}{2.275314in}}%
\pgfpathlineto{\pgfqpoint{5.578161in}{2.288486in}}%
\pgfpathlineto{\pgfqpoint{5.575190in}{2.301831in}}%
\pgfpathlineto{\pgfqpoint{5.572369in}{2.315338in}}%
\pgfpathlineto{\pgfqpoint{5.569699in}{2.328996in}}%
\pgfpathlineto{\pgfqpoint{5.567176in}{2.342794in}}%
\pgfpathlineto{\pgfqpoint{5.564801in}{2.356724in}}%
\pgfpathlineto{\pgfqpoint{5.562571in}{2.370777in}}%
\pgfpathlineto{\pgfqpoint{5.560486in}{2.384944in}}%
\pgfpathlineto{\pgfqpoint{5.558544in}{2.399216in}}%
\pgfpathlineto{\pgfqpoint{5.556744in}{2.413586in}}%
\pgfpathlineto{\pgfqpoint{5.555086in}{2.428046in}}%
\pgfpathlineto{\pgfqpoint{5.553569in}{2.442588in}}%
\pgfpathlineto{\pgfqpoint{5.552191in}{2.457204in}}%
\pgfpathlineto{\pgfqpoint{5.550954in}{2.471888in}}%
\pgfpathlineto{\pgfqpoint{5.549854in}{2.486632in}}%
\pgfpathlineto{\pgfqpoint{5.548894in}{2.501429in}}%
\pgfpathlineto{\pgfqpoint{5.548063in}{2.516436in}}%
\pgfpathlineto{\pgfqpoint{5.547382in}{2.531243in}}%
\pgfpathlineto{\pgfqpoint{5.546837in}{2.546108in}}%
\pgfpathlineto{\pgfqpoint{5.546427in}{2.561018in}}%
\pgfpathlineto{\pgfqpoint{5.546154in}{2.575961in}}%
\pgfpathlineto{\pgfqpoint{5.546017in}{2.590923in}}%
\pgfpathlineto{\pgfqpoint{5.546017in}{2.605891in}}%
\pgfpathlineto{\pgfqpoint{5.546154in}{2.620853in}}%
\pgfpathlineto{\pgfqpoint{5.546427in}{2.635795in}}%
\pgfpathlineto{\pgfqpoint{5.546837in}{2.650705in}}%
\pgfpathlineto{\pgfqpoint{5.547382in}{2.665570in}}%
\pgfpathlineto{\pgfqpoint{5.548063in}{2.680378in}}%
\pgfpathlineto{\pgfqpoint{5.548894in}{2.695384in}}%
\pgfpathlineto{\pgfqpoint{5.549854in}{2.710181in}}%
\pgfpathlineto{\pgfqpoint{5.550954in}{2.724925in}}%
\pgfpathlineto{\pgfqpoint{5.552191in}{2.739609in}}%
\pgfpathlineto{\pgfqpoint{5.553569in}{2.754226in}}%
\pgfpathlineto{\pgfqpoint{5.555086in}{2.768768in}}%
\pgfpathlineto{\pgfqpoint{5.556744in}{2.783228in}}%
\pgfpathlineto{\pgfqpoint{5.558544in}{2.797598in}}%
\pgfpathlineto{\pgfqpoint{5.560486in}{2.811870in}}%
\pgfpathlineto{\pgfqpoint{5.562571in}{2.826037in}}%
\pgfpathlineto{\pgfqpoint{5.564801in}{2.840089in}}%
\pgfpathlineto{\pgfqpoint{5.567176in}{2.854019in}}%
\pgfpathlineto{\pgfqpoint{5.569699in}{2.867818in}}%
\pgfpathlineto{\pgfqpoint{5.572369in}{2.881476in}}%
\pgfpathlineto{\pgfqpoint{5.575190in}{2.894982in}}%
\pgfpathlineto{\pgfqpoint{5.578161in}{2.908327in}}%
\pgfpathlineto{\pgfqpoint{5.581285in}{2.921500in}}%
\pgfpathlineto{\pgfqpoint{5.584621in}{2.934708in}}%
\pgfpathlineto{\pgfqpoint{5.588081in}{2.947576in}}%
\pgfpathlineto{\pgfqpoint{5.591706in}{2.960253in}}%
\pgfpathlineto{\pgfqpoint{5.595503in}{2.972730in}}%
\pgfpathlineto{\pgfqpoint{5.599477in}{2.984997in}}%
\pgfpathlineto{\pgfqpoint{5.603632in}{2.997041in}}%
\pgfpathlineto{\pgfqpoint{5.607976in}{3.008851in}}%
\pgfpathlineto{\pgfqpoint{5.612517in}{3.020413in}}%
\pgfpathlineto{\pgfqpoint{5.617264in}{3.031712in}}%
\pgfpathlineto{\pgfqpoint{5.622227in}{3.042731in}}%
\pgfpathlineto{\pgfqpoint{5.627418in}{3.053451in}}%
\pgfpathlineto{\pgfqpoint{5.632851in}{3.063848in}}%
\pgfpathlineto{\pgfqpoint{5.638540in}{3.073894in}}%
\pgfpathlineto{\pgfqpoint{5.644631in}{3.083755in}}%
\pgfpathlineto{\pgfqpoint{5.650984in}{3.093112in}}%
\pgfusepath{stroke}%
\end{pgfscope}%
\begin{pgfscope}%
\pgfpathmoveto{\pgfqpoint{5.728000in}{2.052407in}}%
\pgfpathcurveto{\pgfqpoint{6.017602in}{2.052407in}}{\pgfqpoint{6.295381in}{2.109937in}}{\pgfqpoint{6.500161in}{2.212326in}}%
\pgfpathcurveto{\pgfqpoint{6.704940in}{2.314716in}}{\pgfqpoint{6.820000in}{2.453606in}}{\pgfqpoint{6.820000in}{2.598407in}}%
\pgfpathcurveto{\pgfqpoint{6.820000in}{2.743208in}}{\pgfqpoint{6.704940in}{2.882097in}}{\pgfqpoint{6.500161in}{2.984487in}}%
\pgfpathcurveto{\pgfqpoint{6.295381in}{3.086877in}}{\pgfqpoint{6.017602in}{3.144407in}}{\pgfqpoint{5.728000in}{3.144407in}}%
\pgfpathcurveto{\pgfqpoint{5.438398in}{3.144407in}}{\pgfqpoint{5.160619in}{3.086877in}}{\pgfqpoint{4.955839in}{2.984487in}}%
\pgfpathcurveto{\pgfqpoint{4.751060in}{2.882097in}}{\pgfqpoint{4.636000in}{2.743208in}}{\pgfqpoint{4.636000in}{2.598407in}}%
\pgfpathcurveto{\pgfqpoint{4.636000in}{2.453606in}}{\pgfqpoint{4.751060in}{2.314716in}}{\pgfqpoint{4.955839in}{2.212326in}}%
\pgfpathcurveto{\pgfqpoint{5.160619in}{2.109937in}}{\pgfqpoint{5.438398in}{2.052407in}}{\pgfqpoint{5.728000in}{2.052407in}}%
\pgfpathlineto{\pgfqpoint{5.728000in}{2.052407in}}%
\pgfpathclose%
\pgfusepath{clip}%
\pgfsetrectcap%
\pgfsetroundjoin%
\pgfsetlinewidth{0.451687pt}%
\definecolor{currentstroke}{rgb}{0.690196,0.690196,0.690196}%
\pgfsetstrokecolor{currentstroke}%
\pgfsetstrokeopacity{0.250000}%
\pgfsetdash{}{0pt}%
\pgfpathmoveto{\pgfqpoint{5.728000in}{2.103702in}}%
\pgfpathlineto{\pgfqpoint{5.728000in}{2.113059in}}%
\pgfpathlineto{\pgfqpoint{5.728000in}{2.122919in}}%
\pgfpathlineto{\pgfqpoint{5.728000in}{2.132965in}}%
\pgfpathlineto{\pgfqpoint{5.728000in}{2.143363in}}%
\pgfpathlineto{\pgfqpoint{5.728000in}{2.154082in}}%
\pgfpathlineto{\pgfqpoint{5.728000in}{2.165102in}}%
\pgfpathlineto{\pgfqpoint{5.728000in}{2.176401in}}%
\pgfpathlineto{\pgfqpoint{5.728000in}{2.187962in}}%
\pgfpathlineto{\pgfqpoint{5.728000in}{2.199772in}}%
\pgfpathlineto{\pgfqpoint{5.728000in}{2.211817in}}%
\pgfpathlineto{\pgfqpoint{5.728000in}{2.224083in}}%
\pgfpathlineto{\pgfqpoint{5.728000in}{2.236561in}}%
\pgfpathlineto{\pgfqpoint{5.728000in}{2.249238in}}%
\pgfpathlineto{\pgfqpoint{5.728000in}{2.262105in}}%
\pgfpathlineto{\pgfqpoint{5.728000in}{2.275314in}}%
\pgfpathlineto{\pgfqpoint{5.728000in}{2.288486in}}%
\pgfpathlineto{\pgfqpoint{5.728000in}{2.301831in}}%
\pgfpathlineto{\pgfqpoint{5.728000in}{2.315338in}}%
\pgfpathlineto{\pgfqpoint{5.728000in}{2.328996in}}%
\pgfpathlineto{\pgfqpoint{5.728000in}{2.342794in}}%
\pgfpathlineto{\pgfqpoint{5.728000in}{2.356724in}}%
\pgfpathlineto{\pgfqpoint{5.728000in}{2.370777in}}%
\pgfpathlineto{\pgfqpoint{5.728000in}{2.384944in}}%
\pgfpathlineto{\pgfqpoint{5.728000in}{2.399216in}}%
\pgfpathlineto{\pgfqpoint{5.728000in}{2.413586in}}%
\pgfpathlineto{\pgfqpoint{5.728000in}{2.428046in}}%
\pgfpathlineto{\pgfqpoint{5.728000in}{2.442588in}}%
\pgfpathlineto{\pgfqpoint{5.728000in}{2.457204in}}%
\pgfpathlineto{\pgfqpoint{5.728000in}{2.471888in}}%
\pgfpathlineto{\pgfqpoint{5.728000in}{2.486632in}}%
\pgfpathlineto{\pgfqpoint{5.728000in}{2.501429in}}%
\pgfpathlineto{\pgfqpoint{5.728000in}{2.516436in}}%
\pgfpathlineto{\pgfqpoint{5.728000in}{2.531243in}}%
\pgfpathlineto{\pgfqpoint{5.728000in}{2.546108in}}%
\pgfpathlineto{\pgfqpoint{5.728000in}{2.561018in}}%
\pgfpathlineto{\pgfqpoint{5.728000in}{2.575961in}}%
\pgfpathlineto{\pgfqpoint{5.728000in}{2.590923in}}%
\pgfpathlineto{\pgfqpoint{5.728000in}{2.605891in}}%
\pgfpathlineto{\pgfqpoint{5.728000in}{2.620853in}}%
\pgfpathlineto{\pgfqpoint{5.728000in}{2.635795in}}%
\pgfpathlineto{\pgfqpoint{5.728000in}{2.650705in}}%
\pgfpathlineto{\pgfqpoint{5.728000in}{2.665570in}}%
\pgfpathlineto{\pgfqpoint{5.728000in}{2.680378in}}%
\pgfpathlineto{\pgfqpoint{5.728000in}{2.695384in}}%
\pgfpathlineto{\pgfqpoint{5.728000in}{2.710181in}}%
\pgfpathlineto{\pgfqpoint{5.728000in}{2.724925in}}%
\pgfpathlineto{\pgfqpoint{5.728000in}{2.739609in}}%
\pgfpathlineto{\pgfqpoint{5.728000in}{2.754226in}}%
\pgfpathlineto{\pgfqpoint{5.728000in}{2.768768in}}%
\pgfpathlineto{\pgfqpoint{5.728000in}{2.783228in}}%
\pgfpathlineto{\pgfqpoint{5.728000in}{2.797598in}}%
\pgfpathlineto{\pgfqpoint{5.728000in}{2.811870in}}%
\pgfpathlineto{\pgfqpoint{5.728000in}{2.826037in}}%
\pgfpathlineto{\pgfqpoint{5.728000in}{2.840089in}}%
\pgfpathlineto{\pgfqpoint{5.728000in}{2.854019in}}%
\pgfpathlineto{\pgfqpoint{5.728000in}{2.867818in}}%
\pgfpathlineto{\pgfqpoint{5.728000in}{2.881476in}}%
\pgfpathlineto{\pgfqpoint{5.728000in}{2.894982in}}%
\pgfpathlineto{\pgfqpoint{5.728000in}{2.908327in}}%
\pgfpathlineto{\pgfqpoint{5.728000in}{2.921500in}}%
\pgfpathlineto{\pgfqpoint{5.728000in}{2.934708in}}%
\pgfpathlineto{\pgfqpoint{5.728000in}{2.947576in}}%
\pgfpathlineto{\pgfqpoint{5.728000in}{2.960253in}}%
\pgfpathlineto{\pgfqpoint{5.728000in}{2.972730in}}%
\pgfpathlineto{\pgfqpoint{5.728000in}{2.984997in}}%
\pgfpathlineto{\pgfqpoint{5.728000in}{2.997041in}}%
\pgfpathlineto{\pgfqpoint{5.728000in}{3.008851in}}%
\pgfpathlineto{\pgfqpoint{5.728000in}{3.020413in}}%
\pgfpathlineto{\pgfqpoint{5.728000in}{3.031712in}}%
\pgfpathlineto{\pgfqpoint{5.728000in}{3.042731in}}%
\pgfpathlineto{\pgfqpoint{5.728000in}{3.053451in}}%
\pgfpathlineto{\pgfqpoint{5.728000in}{3.063848in}}%
\pgfpathlineto{\pgfqpoint{5.728000in}{3.073894in}}%
\pgfpathlineto{\pgfqpoint{5.728000in}{3.083755in}}%
\pgfpathlineto{\pgfqpoint{5.728000in}{3.093112in}}%
\pgfusepath{stroke}%
\end{pgfscope}%
\begin{pgfscope}%
\pgfpathmoveto{\pgfqpoint{5.728000in}{2.052407in}}%
\pgfpathcurveto{\pgfqpoint{6.017602in}{2.052407in}}{\pgfqpoint{6.295381in}{2.109937in}}{\pgfqpoint{6.500161in}{2.212326in}}%
\pgfpathcurveto{\pgfqpoint{6.704940in}{2.314716in}}{\pgfqpoint{6.820000in}{2.453606in}}{\pgfqpoint{6.820000in}{2.598407in}}%
\pgfpathcurveto{\pgfqpoint{6.820000in}{2.743208in}}{\pgfqpoint{6.704940in}{2.882097in}}{\pgfqpoint{6.500161in}{2.984487in}}%
\pgfpathcurveto{\pgfqpoint{6.295381in}{3.086877in}}{\pgfqpoint{6.017602in}{3.144407in}}{\pgfqpoint{5.728000in}{3.144407in}}%
\pgfpathcurveto{\pgfqpoint{5.438398in}{3.144407in}}{\pgfqpoint{5.160619in}{3.086877in}}{\pgfqpoint{4.955839in}{2.984487in}}%
\pgfpathcurveto{\pgfqpoint{4.751060in}{2.882097in}}{\pgfqpoint{4.636000in}{2.743208in}}{\pgfqpoint{4.636000in}{2.598407in}}%
\pgfpathcurveto{\pgfqpoint{4.636000in}{2.453606in}}{\pgfqpoint{4.751060in}{2.314716in}}{\pgfqpoint{4.955839in}{2.212326in}}%
\pgfpathcurveto{\pgfqpoint{5.160619in}{2.109937in}}{\pgfqpoint{5.438398in}{2.052407in}}{\pgfqpoint{5.728000in}{2.052407in}}%
\pgfpathlineto{\pgfqpoint{5.728000in}{2.052407in}}%
\pgfpathclose%
\pgfusepath{clip}%
\pgfsetrectcap%
\pgfsetroundjoin%
\pgfsetlinewidth{0.451687pt}%
\definecolor{currentstroke}{rgb}{0.690196,0.690196,0.690196}%
\pgfsetstrokecolor{currentstroke}%
\pgfsetstrokeopacity{0.250000}%
\pgfsetdash{}{0pt}%
\pgfpathmoveto{\pgfqpoint{5.805016in}{2.103702in}}%
\pgfpathlineto{\pgfqpoint{5.811369in}{2.113059in}}%
\pgfpathlineto{\pgfqpoint{5.817460in}{2.122919in}}%
\pgfpathlineto{\pgfqpoint{5.823149in}{2.132965in}}%
\pgfpathlineto{\pgfqpoint{5.828582in}{2.143363in}}%
\pgfpathlineto{\pgfqpoint{5.833773in}{2.154082in}}%
\pgfpathlineto{\pgfqpoint{5.838736in}{2.165102in}}%
\pgfpathlineto{\pgfqpoint{5.843483in}{2.176401in}}%
\pgfpathlineto{\pgfqpoint{5.848024in}{2.187962in}}%
\pgfpathlineto{\pgfqpoint{5.852368in}{2.199772in}}%
\pgfpathlineto{\pgfqpoint{5.856523in}{2.211817in}}%
\pgfpathlineto{\pgfqpoint{5.860497in}{2.224083in}}%
\pgfpathlineto{\pgfqpoint{5.864294in}{2.236561in}}%
\pgfpathlineto{\pgfqpoint{5.867919in}{2.249238in}}%
\pgfpathlineto{\pgfqpoint{5.871379in}{2.262105in}}%
\pgfpathlineto{\pgfqpoint{5.874715in}{2.275314in}}%
\pgfpathlineto{\pgfqpoint{5.877839in}{2.288486in}}%
\pgfpathlineto{\pgfqpoint{5.880810in}{2.301831in}}%
\pgfpathlineto{\pgfqpoint{5.883631in}{2.315338in}}%
\pgfpathlineto{\pgfqpoint{5.886301in}{2.328996in}}%
\pgfpathlineto{\pgfqpoint{5.888824in}{2.342794in}}%
\pgfpathlineto{\pgfqpoint{5.891199in}{2.356724in}}%
\pgfpathlineto{\pgfqpoint{5.893429in}{2.370777in}}%
\pgfpathlineto{\pgfqpoint{5.895514in}{2.384944in}}%
\pgfpathlineto{\pgfqpoint{5.897456in}{2.399216in}}%
\pgfpathlineto{\pgfqpoint{5.899256in}{2.413586in}}%
\pgfpathlineto{\pgfqpoint{5.900914in}{2.428046in}}%
\pgfpathlineto{\pgfqpoint{5.902431in}{2.442588in}}%
\pgfpathlineto{\pgfqpoint{5.903809in}{2.457204in}}%
\pgfpathlineto{\pgfqpoint{5.905046in}{2.471888in}}%
\pgfpathlineto{\pgfqpoint{5.906146in}{2.486632in}}%
\pgfpathlineto{\pgfqpoint{5.907106in}{2.501429in}}%
\pgfpathlineto{\pgfqpoint{5.907937in}{2.516436in}}%
\pgfpathlineto{\pgfqpoint{5.908618in}{2.531243in}}%
\pgfpathlineto{\pgfqpoint{5.909163in}{2.546108in}}%
\pgfpathlineto{\pgfqpoint{5.909573in}{2.561018in}}%
\pgfpathlineto{\pgfqpoint{5.909846in}{2.575961in}}%
\pgfpathlineto{\pgfqpoint{5.909983in}{2.590923in}}%
\pgfpathlineto{\pgfqpoint{5.909983in}{2.605891in}}%
\pgfpathlineto{\pgfqpoint{5.909846in}{2.620853in}}%
\pgfpathlineto{\pgfqpoint{5.909573in}{2.635795in}}%
\pgfpathlineto{\pgfqpoint{5.909163in}{2.650705in}}%
\pgfpathlineto{\pgfqpoint{5.908618in}{2.665570in}}%
\pgfpathlineto{\pgfqpoint{5.907937in}{2.680378in}}%
\pgfpathlineto{\pgfqpoint{5.907106in}{2.695384in}}%
\pgfpathlineto{\pgfqpoint{5.906146in}{2.710181in}}%
\pgfpathlineto{\pgfqpoint{5.905046in}{2.724925in}}%
\pgfpathlineto{\pgfqpoint{5.903809in}{2.739609in}}%
\pgfpathlineto{\pgfqpoint{5.902431in}{2.754226in}}%
\pgfpathlineto{\pgfqpoint{5.900914in}{2.768768in}}%
\pgfpathlineto{\pgfqpoint{5.899256in}{2.783228in}}%
\pgfpathlineto{\pgfqpoint{5.897456in}{2.797598in}}%
\pgfpathlineto{\pgfqpoint{5.895514in}{2.811870in}}%
\pgfpathlineto{\pgfqpoint{5.893429in}{2.826037in}}%
\pgfpathlineto{\pgfqpoint{5.891199in}{2.840089in}}%
\pgfpathlineto{\pgfqpoint{5.888824in}{2.854019in}}%
\pgfpathlineto{\pgfqpoint{5.886301in}{2.867818in}}%
\pgfpathlineto{\pgfqpoint{5.883631in}{2.881476in}}%
\pgfpathlineto{\pgfqpoint{5.880810in}{2.894982in}}%
\pgfpathlineto{\pgfqpoint{5.877839in}{2.908327in}}%
\pgfpathlineto{\pgfqpoint{5.874715in}{2.921500in}}%
\pgfpathlineto{\pgfqpoint{5.871379in}{2.934708in}}%
\pgfpathlineto{\pgfqpoint{5.867919in}{2.947576in}}%
\pgfpathlineto{\pgfqpoint{5.864294in}{2.960253in}}%
\pgfpathlineto{\pgfqpoint{5.860497in}{2.972730in}}%
\pgfpathlineto{\pgfqpoint{5.856523in}{2.984997in}}%
\pgfpathlineto{\pgfqpoint{5.852368in}{2.997041in}}%
\pgfpathlineto{\pgfqpoint{5.848024in}{3.008851in}}%
\pgfpathlineto{\pgfqpoint{5.843483in}{3.020413in}}%
\pgfpathlineto{\pgfqpoint{5.838736in}{3.031712in}}%
\pgfpathlineto{\pgfqpoint{5.833773in}{3.042731in}}%
\pgfpathlineto{\pgfqpoint{5.828582in}{3.053451in}}%
\pgfpathlineto{\pgfqpoint{5.823149in}{3.063848in}}%
\pgfpathlineto{\pgfqpoint{5.817460in}{3.073894in}}%
\pgfpathlineto{\pgfqpoint{5.811369in}{3.083755in}}%
\pgfpathlineto{\pgfqpoint{5.805016in}{3.093112in}}%
\pgfusepath{stroke}%
\end{pgfscope}%
\begin{pgfscope}%
\pgfpathmoveto{\pgfqpoint{5.728000in}{2.052407in}}%
\pgfpathcurveto{\pgfqpoint{6.017602in}{2.052407in}}{\pgfqpoint{6.295381in}{2.109937in}}{\pgfqpoint{6.500161in}{2.212326in}}%
\pgfpathcurveto{\pgfqpoint{6.704940in}{2.314716in}}{\pgfqpoint{6.820000in}{2.453606in}}{\pgfqpoint{6.820000in}{2.598407in}}%
\pgfpathcurveto{\pgfqpoint{6.820000in}{2.743208in}}{\pgfqpoint{6.704940in}{2.882097in}}{\pgfqpoint{6.500161in}{2.984487in}}%
\pgfpathcurveto{\pgfqpoint{6.295381in}{3.086877in}}{\pgfqpoint{6.017602in}{3.144407in}}{\pgfqpoint{5.728000in}{3.144407in}}%
\pgfpathcurveto{\pgfqpoint{5.438398in}{3.144407in}}{\pgfqpoint{5.160619in}{3.086877in}}{\pgfqpoint{4.955839in}{2.984487in}}%
\pgfpathcurveto{\pgfqpoint{4.751060in}{2.882097in}}{\pgfqpoint{4.636000in}{2.743208in}}{\pgfqpoint{4.636000in}{2.598407in}}%
\pgfpathcurveto{\pgfqpoint{4.636000in}{2.453606in}}{\pgfqpoint{4.751060in}{2.314716in}}{\pgfqpoint{4.955839in}{2.212326in}}%
\pgfpathcurveto{\pgfqpoint{5.160619in}{2.109937in}}{\pgfqpoint{5.438398in}{2.052407in}}{\pgfqpoint{5.728000in}{2.052407in}}%
\pgfpathlineto{\pgfqpoint{5.728000in}{2.052407in}}%
\pgfpathclose%
\pgfusepath{clip}%
\pgfsetrectcap%
\pgfsetroundjoin%
\pgfsetlinewidth{0.451687pt}%
\definecolor{currentstroke}{rgb}{0.690196,0.690196,0.690196}%
\pgfsetstrokecolor{currentstroke}%
\pgfsetstrokeopacity{0.250000}%
\pgfsetdash{}{0pt}%
\pgfpathmoveto{\pgfqpoint{5.882031in}{2.103702in}}%
\pgfpathlineto{\pgfqpoint{5.894737in}{2.113059in}}%
\pgfpathlineto{\pgfqpoint{5.906920in}{2.122919in}}%
\pgfpathlineto{\pgfqpoint{5.918298in}{2.132965in}}%
\pgfpathlineto{\pgfqpoint{5.929164in}{2.143363in}}%
\pgfpathlineto{\pgfqpoint{5.939547in}{2.154082in}}%
\pgfpathlineto{\pgfqpoint{5.949473in}{2.165102in}}%
\pgfpathlineto{\pgfqpoint{5.958966in}{2.176401in}}%
\pgfpathlineto{\pgfqpoint{5.968048in}{2.187962in}}%
\pgfpathlineto{\pgfqpoint{5.976736in}{2.199772in}}%
\pgfpathlineto{\pgfqpoint{5.985047in}{2.211817in}}%
\pgfpathlineto{\pgfqpoint{5.992993in}{2.224083in}}%
\pgfpathlineto{\pgfqpoint{6.000587in}{2.236561in}}%
\pgfpathlineto{\pgfqpoint{6.007839in}{2.249238in}}%
\pgfpathlineto{\pgfqpoint{6.014757in}{2.262105in}}%
\pgfpathlineto{\pgfqpoint{6.021430in}{2.275314in}}%
\pgfpathlineto{\pgfqpoint{6.027678in}{2.288486in}}%
\pgfpathlineto{\pgfqpoint{6.033621in}{2.301831in}}%
\pgfpathlineto{\pgfqpoint{6.039261in}{2.315338in}}%
\pgfpathlineto{\pgfqpoint{6.044602in}{2.328996in}}%
\pgfpathlineto{\pgfqpoint{6.049647in}{2.342794in}}%
\pgfpathlineto{\pgfqpoint{6.054398in}{2.356724in}}%
\pgfpathlineto{\pgfqpoint{6.058858in}{2.370777in}}%
\pgfpathlineto{\pgfqpoint{6.063029in}{2.384944in}}%
\pgfpathlineto{\pgfqpoint{6.066913in}{2.399216in}}%
\pgfpathlineto{\pgfqpoint{6.070512in}{2.413586in}}%
\pgfpathlineto{\pgfqpoint{6.073828in}{2.428046in}}%
\pgfpathlineto{\pgfqpoint{6.076863in}{2.442588in}}%
\pgfpathlineto{\pgfqpoint{6.079617in}{2.457204in}}%
\pgfpathlineto{\pgfqpoint{6.082093in}{2.471888in}}%
\pgfpathlineto{\pgfqpoint{6.084291in}{2.486632in}}%
\pgfpathlineto{\pgfqpoint{6.086212in}{2.501429in}}%
\pgfpathlineto{\pgfqpoint{6.087875in}{2.516436in}}%
\pgfpathlineto{\pgfqpoint{6.089236in}{2.531243in}}%
\pgfpathlineto{\pgfqpoint{6.090326in}{2.546108in}}%
\pgfpathlineto{\pgfqpoint{6.091146in}{2.561018in}}%
\pgfpathlineto{\pgfqpoint{6.091692in}{2.575961in}}%
\pgfpathlineto{\pgfqpoint{6.091966in}{2.590923in}}%
\pgfpathlineto{\pgfqpoint{6.091966in}{2.605891in}}%
\pgfpathlineto{\pgfqpoint{6.091692in}{2.620853in}}%
\pgfpathlineto{\pgfqpoint{6.091146in}{2.635795in}}%
\pgfpathlineto{\pgfqpoint{6.090326in}{2.650705in}}%
\pgfpathlineto{\pgfqpoint{6.089236in}{2.665570in}}%
\pgfpathlineto{\pgfqpoint{6.087875in}{2.680378in}}%
\pgfpathlineto{\pgfqpoint{6.086212in}{2.695384in}}%
\pgfpathlineto{\pgfqpoint{6.084291in}{2.710181in}}%
\pgfpathlineto{\pgfqpoint{6.082093in}{2.724925in}}%
\pgfpathlineto{\pgfqpoint{6.079617in}{2.739609in}}%
\pgfpathlineto{\pgfqpoint{6.076863in}{2.754226in}}%
\pgfpathlineto{\pgfqpoint{6.073828in}{2.768768in}}%
\pgfpathlineto{\pgfqpoint{6.070512in}{2.783228in}}%
\pgfpathlineto{\pgfqpoint{6.066913in}{2.797598in}}%
\pgfpathlineto{\pgfqpoint{6.063029in}{2.811870in}}%
\pgfpathlineto{\pgfqpoint{6.058858in}{2.826037in}}%
\pgfpathlineto{\pgfqpoint{6.054398in}{2.840089in}}%
\pgfpathlineto{\pgfqpoint{6.049647in}{2.854019in}}%
\pgfpathlineto{\pgfqpoint{6.044602in}{2.867818in}}%
\pgfpathlineto{\pgfqpoint{6.039261in}{2.881476in}}%
\pgfpathlineto{\pgfqpoint{6.033621in}{2.894982in}}%
\pgfpathlineto{\pgfqpoint{6.027678in}{2.908327in}}%
\pgfpathlineto{\pgfqpoint{6.021430in}{2.921500in}}%
\pgfpathlineto{\pgfqpoint{6.014757in}{2.934708in}}%
\pgfpathlineto{\pgfqpoint{6.007839in}{2.947576in}}%
\pgfpathlineto{\pgfqpoint{6.000587in}{2.960253in}}%
\pgfpathlineto{\pgfqpoint{5.992993in}{2.972730in}}%
\pgfpathlineto{\pgfqpoint{5.985047in}{2.984997in}}%
\pgfpathlineto{\pgfqpoint{5.976736in}{2.997041in}}%
\pgfpathlineto{\pgfqpoint{5.968048in}{3.008851in}}%
\pgfpathlineto{\pgfqpoint{5.958966in}{3.020413in}}%
\pgfpathlineto{\pgfqpoint{5.949473in}{3.031712in}}%
\pgfpathlineto{\pgfqpoint{5.939547in}{3.042731in}}%
\pgfpathlineto{\pgfqpoint{5.929164in}{3.053451in}}%
\pgfpathlineto{\pgfqpoint{5.918298in}{3.063848in}}%
\pgfpathlineto{\pgfqpoint{5.906920in}{3.073894in}}%
\pgfpathlineto{\pgfqpoint{5.894737in}{3.083755in}}%
\pgfpathlineto{\pgfqpoint{5.882031in}{3.093112in}}%
\pgfusepath{stroke}%
\end{pgfscope}%
\begin{pgfscope}%
\pgfpathmoveto{\pgfqpoint{5.728000in}{2.052407in}}%
\pgfpathcurveto{\pgfqpoint{6.017602in}{2.052407in}}{\pgfqpoint{6.295381in}{2.109937in}}{\pgfqpoint{6.500161in}{2.212326in}}%
\pgfpathcurveto{\pgfqpoint{6.704940in}{2.314716in}}{\pgfqpoint{6.820000in}{2.453606in}}{\pgfqpoint{6.820000in}{2.598407in}}%
\pgfpathcurveto{\pgfqpoint{6.820000in}{2.743208in}}{\pgfqpoint{6.704940in}{2.882097in}}{\pgfqpoint{6.500161in}{2.984487in}}%
\pgfpathcurveto{\pgfqpoint{6.295381in}{3.086877in}}{\pgfqpoint{6.017602in}{3.144407in}}{\pgfqpoint{5.728000in}{3.144407in}}%
\pgfpathcurveto{\pgfqpoint{5.438398in}{3.144407in}}{\pgfqpoint{5.160619in}{3.086877in}}{\pgfqpoint{4.955839in}{2.984487in}}%
\pgfpathcurveto{\pgfqpoint{4.751060in}{2.882097in}}{\pgfqpoint{4.636000in}{2.743208in}}{\pgfqpoint{4.636000in}{2.598407in}}%
\pgfpathcurveto{\pgfqpoint{4.636000in}{2.453606in}}{\pgfqpoint{4.751060in}{2.314716in}}{\pgfqpoint{4.955839in}{2.212326in}}%
\pgfpathcurveto{\pgfqpoint{5.160619in}{2.109937in}}{\pgfqpoint{5.438398in}{2.052407in}}{\pgfqpoint{5.728000in}{2.052407in}}%
\pgfpathlineto{\pgfqpoint{5.728000in}{2.052407in}}%
\pgfpathclose%
\pgfusepath{clip}%
\pgfsetrectcap%
\pgfsetroundjoin%
\pgfsetlinewidth{0.451687pt}%
\definecolor{currentstroke}{rgb}{0.690196,0.690196,0.690196}%
\pgfsetstrokecolor{currentstroke}%
\pgfsetstrokeopacity{0.250000}%
\pgfsetdash{}{0pt}%
\pgfpathmoveto{\pgfqpoint{5.959047in}{2.103702in}}%
\pgfpathlineto{\pgfqpoint{5.978106in}{2.113059in}}%
\pgfpathlineto{\pgfqpoint{5.996380in}{2.122919in}}%
\pgfpathlineto{\pgfqpoint{6.013447in}{2.132965in}}%
\pgfpathlineto{\pgfqpoint{6.029746in}{2.143363in}}%
\pgfpathlineto{\pgfqpoint{6.045320in}{2.154082in}}%
\pgfpathlineto{\pgfqpoint{6.060209in}{2.165102in}}%
\pgfpathlineto{\pgfqpoint{6.074449in}{2.176401in}}%
\pgfpathlineto{\pgfqpoint{6.088071in}{2.187962in}}%
\pgfpathlineto{\pgfqpoint{6.101104in}{2.199772in}}%
\pgfpathlineto{\pgfqpoint{6.113570in}{2.211817in}}%
\pgfpathlineto{\pgfqpoint{6.125490in}{2.224083in}}%
\pgfpathlineto{\pgfqpoint{6.136881in}{2.236561in}}%
\pgfpathlineto{\pgfqpoint{6.147758in}{2.249238in}}%
\pgfpathlineto{\pgfqpoint{6.158136in}{2.262105in}}%
\pgfpathlineto{\pgfqpoint{6.168144in}{2.275314in}}%
\pgfpathlineto{\pgfqpoint{6.177517in}{2.288486in}}%
\pgfpathlineto{\pgfqpoint{6.186431in}{2.301831in}}%
\pgfpathlineto{\pgfqpoint{6.194892in}{2.315338in}}%
\pgfpathlineto{\pgfqpoint{6.202904in}{2.328996in}}%
\pgfpathlineto{\pgfqpoint{6.210471in}{2.342794in}}%
\pgfpathlineto{\pgfqpoint{6.217597in}{2.356724in}}%
\pgfpathlineto{\pgfqpoint{6.224287in}{2.370777in}}%
\pgfpathlineto{\pgfqpoint{6.230543in}{2.384944in}}%
\pgfpathlineto{\pgfqpoint{6.236369in}{2.399216in}}%
\pgfpathlineto{\pgfqpoint{6.241768in}{2.413586in}}%
\pgfpathlineto{\pgfqpoint{6.246742in}{2.428046in}}%
\pgfpathlineto{\pgfqpoint{6.251294in}{2.442588in}}%
\pgfpathlineto{\pgfqpoint{6.255426in}{2.457204in}}%
\pgfpathlineto{\pgfqpoint{6.259139in}{2.471888in}}%
\pgfpathlineto{\pgfqpoint{6.262437in}{2.486632in}}%
\pgfpathlineto{\pgfqpoint{6.265319in}{2.501429in}}%
\pgfpathlineto{\pgfqpoint{6.267812in}{2.516436in}}%
\pgfpathlineto{\pgfqpoint{6.269853in}{2.531243in}}%
\pgfpathlineto{\pgfqpoint{6.271490in}{2.546108in}}%
\pgfpathlineto{\pgfqpoint{6.272718in}{2.561018in}}%
\pgfpathlineto{\pgfqpoint{6.273538in}{2.575961in}}%
\pgfpathlineto{\pgfqpoint{6.273949in}{2.590923in}}%
\pgfpathlineto{\pgfqpoint{6.273949in}{2.605891in}}%
\pgfpathlineto{\pgfqpoint{6.273538in}{2.620853in}}%
\pgfpathlineto{\pgfqpoint{6.272718in}{2.635795in}}%
\pgfpathlineto{\pgfqpoint{6.271490in}{2.650705in}}%
\pgfpathlineto{\pgfqpoint{6.269853in}{2.665570in}}%
\pgfpathlineto{\pgfqpoint{6.267812in}{2.680378in}}%
\pgfpathlineto{\pgfqpoint{6.265319in}{2.695384in}}%
\pgfpathlineto{\pgfqpoint{6.262437in}{2.710181in}}%
\pgfpathlineto{\pgfqpoint{6.259139in}{2.724925in}}%
\pgfpathlineto{\pgfqpoint{6.255426in}{2.739609in}}%
\pgfpathlineto{\pgfqpoint{6.251294in}{2.754226in}}%
\pgfpathlineto{\pgfqpoint{6.246742in}{2.768768in}}%
\pgfpathlineto{\pgfqpoint{6.241768in}{2.783228in}}%
\pgfpathlineto{\pgfqpoint{6.236369in}{2.797598in}}%
\pgfpathlineto{\pgfqpoint{6.230543in}{2.811870in}}%
\pgfpathlineto{\pgfqpoint{6.224287in}{2.826037in}}%
\pgfpathlineto{\pgfqpoint{6.217597in}{2.840089in}}%
\pgfpathlineto{\pgfqpoint{6.210471in}{2.854019in}}%
\pgfpathlineto{\pgfqpoint{6.202904in}{2.867818in}}%
\pgfpathlineto{\pgfqpoint{6.194892in}{2.881476in}}%
\pgfpathlineto{\pgfqpoint{6.186431in}{2.894982in}}%
\pgfpathlineto{\pgfqpoint{6.177517in}{2.908327in}}%
\pgfpathlineto{\pgfqpoint{6.168144in}{2.921500in}}%
\pgfpathlineto{\pgfqpoint{6.158136in}{2.934708in}}%
\pgfpathlineto{\pgfqpoint{6.147758in}{2.947576in}}%
\pgfpathlineto{\pgfqpoint{6.136881in}{2.960253in}}%
\pgfpathlineto{\pgfqpoint{6.125490in}{2.972730in}}%
\pgfpathlineto{\pgfqpoint{6.113570in}{2.984997in}}%
\pgfpathlineto{\pgfqpoint{6.101104in}{2.997041in}}%
\pgfpathlineto{\pgfqpoint{6.088071in}{3.008851in}}%
\pgfpathlineto{\pgfqpoint{6.074449in}{3.020413in}}%
\pgfpathlineto{\pgfqpoint{6.060209in}{3.031712in}}%
\pgfpathlineto{\pgfqpoint{6.045320in}{3.042731in}}%
\pgfpathlineto{\pgfqpoint{6.029746in}{3.053451in}}%
\pgfpathlineto{\pgfqpoint{6.013447in}{3.063848in}}%
\pgfpathlineto{\pgfqpoint{5.996380in}{3.073894in}}%
\pgfpathlineto{\pgfqpoint{5.978106in}{3.083755in}}%
\pgfpathlineto{\pgfqpoint{5.959047in}{3.093112in}}%
\pgfusepath{stroke}%
\end{pgfscope}%
\begin{pgfscope}%
\pgfpathmoveto{\pgfqpoint{5.728000in}{2.052407in}}%
\pgfpathcurveto{\pgfqpoint{6.017602in}{2.052407in}}{\pgfqpoint{6.295381in}{2.109937in}}{\pgfqpoint{6.500161in}{2.212326in}}%
\pgfpathcurveto{\pgfqpoint{6.704940in}{2.314716in}}{\pgfqpoint{6.820000in}{2.453606in}}{\pgfqpoint{6.820000in}{2.598407in}}%
\pgfpathcurveto{\pgfqpoint{6.820000in}{2.743208in}}{\pgfqpoint{6.704940in}{2.882097in}}{\pgfqpoint{6.500161in}{2.984487in}}%
\pgfpathcurveto{\pgfqpoint{6.295381in}{3.086877in}}{\pgfqpoint{6.017602in}{3.144407in}}{\pgfqpoint{5.728000in}{3.144407in}}%
\pgfpathcurveto{\pgfqpoint{5.438398in}{3.144407in}}{\pgfqpoint{5.160619in}{3.086877in}}{\pgfqpoint{4.955839in}{2.984487in}}%
\pgfpathcurveto{\pgfqpoint{4.751060in}{2.882097in}}{\pgfqpoint{4.636000in}{2.743208in}}{\pgfqpoint{4.636000in}{2.598407in}}%
\pgfpathcurveto{\pgfqpoint{4.636000in}{2.453606in}}{\pgfqpoint{4.751060in}{2.314716in}}{\pgfqpoint{4.955839in}{2.212326in}}%
\pgfpathcurveto{\pgfqpoint{5.160619in}{2.109937in}}{\pgfqpoint{5.438398in}{2.052407in}}{\pgfqpoint{5.728000in}{2.052407in}}%
\pgfpathlineto{\pgfqpoint{5.728000in}{2.052407in}}%
\pgfpathclose%
\pgfusepath{clip}%
\pgfsetrectcap%
\pgfsetroundjoin%
\pgfsetlinewidth{0.451687pt}%
\definecolor{currentstroke}{rgb}{0.690196,0.690196,0.690196}%
\pgfsetstrokecolor{currentstroke}%
\pgfsetstrokeopacity{0.250000}%
\pgfsetdash{}{0pt}%
\pgfpathmoveto{\pgfqpoint{6.036063in}{2.103702in}}%
\pgfpathlineto{\pgfqpoint{6.061475in}{2.113059in}}%
\pgfpathlineto{\pgfqpoint{6.085840in}{2.122919in}}%
\pgfpathlineto{\pgfqpoint{6.108596in}{2.132965in}}%
\pgfpathlineto{\pgfqpoint{6.130328in}{2.143363in}}%
\pgfpathlineto{\pgfqpoint{6.151093in}{2.154082in}}%
\pgfpathlineto{\pgfqpoint{6.170945in}{2.165102in}}%
\pgfpathlineto{\pgfqpoint{6.189932in}{2.176401in}}%
\pgfpathlineto{\pgfqpoint{6.208095in}{2.187962in}}%
\pgfpathlineto{\pgfqpoint{6.225472in}{2.199772in}}%
\pgfpathlineto{\pgfqpoint{6.242093in}{2.211817in}}%
\pgfpathlineto{\pgfqpoint{6.257986in}{2.224083in}}%
\pgfpathlineto{\pgfqpoint{6.273174in}{2.236561in}}%
\pgfpathlineto{\pgfqpoint{6.287678in}{2.249238in}}%
\pgfpathlineto{\pgfqpoint{6.301515in}{2.262105in}}%
\pgfpathlineto{\pgfqpoint{6.314859in}{2.275314in}}%
\pgfpathlineto{\pgfqpoint{6.327356in}{2.288486in}}%
\pgfpathlineto{\pgfqpoint{6.339241in}{2.301831in}}%
\pgfpathlineto{\pgfqpoint{6.350523in}{2.315338in}}%
\pgfpathlineto{\pgfqpoint{6.361205in}{2.328996in}}%
\pgfpathlineto{\pgfqpoint{6.371295in}{2.342794in}}%
\pgfpathlineto{\pgfqpoint{6.380796in}{2.356724in}}%
\pgfpathlineto{\pgfqpoint{6.389716in}{2.370777in}}%
\pgfpathlineto{\pgfqpoint{6.398057in}{2.384944in}}%
\pgfpathlineto{\pgfqpoint{6.405825in}{2.399216in}}%
\pgfpathlineto{\pgfqpoint{6.413024in}{2.413586in}}%
\pgfpathlineto{\pgfqpoint{6.419656in}{2.428046in}}%
\pgfpathlineto{\pgfqpoint{6.425725in}{2.442588in}}%
\pgfpathlineto{\pgfqpoint{6.431234in}{2.457204in}}%
\pgfpathlineto{\pgfqpoint{6.436186in}{2.471888in}}%
\pgfpathlineto{\pgfqpoint{6.440582in}{2.486632in}}%
\pgfpathlineto{\pgfqpoint{6.444425in}{2.501429in}}%
\pgfpathlineto{\pgfqpoint{6.447749in}{2.516436in}}%
\pgfpathlineto{\pgfqpoint{6.450471in}{2.531243in}}%
\pgfpathlineto{\pgfqpoint{6.452653in}{2.546108in}}%
\pgfpathlineto{\pgfqpoint{6.454291in}{2.561018in}}%
\pgfpathlineto{\pgfqpoint{6.455385in}{2.575961in}}%
\pgfpathlineto{\pgfqpoint{6.455932in}{2.590923in}}%
\pgfpathlineto{\pgfqpoint{6.455932in}{2.605891in}}%
\pgfpathlineto{\pgfqpoint{6.455385in}{2.620853in}}%
\pgfpathlineto{\pgfqpoint{6.454291in}{2.635795in}}%
\pgfpathlineto{\pgfqpoint{6.452653in}{2.650705in}}%
\pgfpathlineto{\pgfqpoint{6.450471in}{2.665570in}}%
\pgfpathlineto{\pgfqpoint{6.447749in}{2.680378in}}%
\pgfpathlineto{\pgfqpoint{6.444425in}{2.695384in}}%
\pgfpathlineto{\pgfqpoint{6.440582in}{2.710181in}}%
\pgfpathlineto{\pgfqpoint{6.436186in}{2.724925in}}%
\pgfpathlineto{\pgfqpoint{6.431234in}{2.739609in}}%
\pgfpathlineto{\pgfqpoint{6.425725in}{2.754226in}}%
\pgfpathlineto{\pgfqpoint{6.419656in}{2.768768in}}%
\pgfpathlineto{\pgfqpoint{6.413024in}{2.783228in}}%
\pgfpathlineto{\pgfqpoint{6.405825in}{2.797598in}}%
\pgfpathlineto{\pgfqpoint{6.398057in}{2.811870in}}%
\pgfpathlineto{\pgfqpoint{6.389716in}{2.826037in}}%
\pgfpathlineto{\pgfqpoint{6.380796in}{2.840089in}}%
\pgfpathlineto{\pgfqpoint{6.371295in}{2.854019in}}%
\pgfpathlineto{\pgfqpoint{6.361205in}{2.867818in}}%
\pgfpathlineto{\pgfqpoint{6.350523in}{2.881476in}}%
\pgfpathlineto{\pgfqpoint{6.339241in}{2.894982in}}%
\pgfpathlineto{\pgfqpoint{6.327356in}{2.908327in}}%
\pgfpathlineto{\pgfqpoint{6.314859in}{2.921500in}}%
\pgfpathlineto{\pgfqpoint{6.301515in}{2.934708in}}%
\pgfpathlineto{\pgfqpoint{6.287678in}{2.947576in}}%
\pgfpathlineto{\pgfqpoint{6.273174in}{2.960253in}}%
\pgfpathlineto{\pgfqpoint{6.257986in}{2.972730in}}%
\pgfpathlineto{\pgfqpoint{6.242093in}{2.984997in}}%
\pgfpathlineto{\pgfqpoint{6.225472in}{2.997041in}}%
\pgfpathlineto{\pgfqpoint{6.208095in}{3.008851in}}%
\pgfpathlineto{\pgfqpoint{6.189932in}{3.020413in}}%
\pgfpathlineto{\pgfqpoint{6.170945in}{3.031712in}}%
\pgfpathlineto{\pgfqpoint{6.151093in}{3.042731in}}%
\pgfpathlineto{\pgfqpoint{6.130328in}{3.053451in}}%
\pgfpathlineto{\pgfqpoint{6.108596in}{3.063848in}}%
\pgfpathlineto{\pgfqpoint{6.085840in}{3.073894in}}%
\pgfpathlineto{\pgfqpoint{6.061475in}{3.083755in}}%
\pgfpathlineto{\pgfqpoint{6.036063in}{3.093112in}}%
\pgfusepath{stroke}%
\end{pgfscope}%
\begin{pgfscope}%
\pgfpathmoveto{\pgfqpoint{5.728000in}{2.052407in}}%
\pgfpathcurveto{\pgfqpoint{6.017602in}{2.052407in}}{\pgfqpoint{6.295381in}{2.109937in}}{\pgfqpoint{6.500161in}{2.212326in}}%
\pgfpathcurveto{\pgfqpoint{6.704940in}{2.314716in}}{\pgfqpoint{6.820000in}{2.453606in}}{\pgfqpoint{6.820000in}{2.598407in}}%
\pgfpathcurveto{\pgfqpoint{6.820000in}{2.743208in}}{\pgfqpoint{6.704940in}{2.882097in}}{\pgfqpoint{6.500161in}{2.984487in}}%
\pgfpathcurveto{\pgfqpoint{6.295381in}{3.086877in}}{\pgfqpoint{6.017602in}{3.144407in}}{\pgfqpoint{5.728000in}{3.144407in}}%
\pgfpathcurveto{\pgfqpoint{5.438398in}{3.144407in}}{\pgfqpoint{5.160619in}{3.086877in}}{\pgfqpoint{4.955839in}{2.984487in}}%
\pgfpathcurveto{\pgfqpoint{4.751060in}{2.882097in}}{\pgfqpoint{4.636000in}{2.743208in}}{\pgfqpoint{4.636000in}{2.598407in}}%
\pgfpathcurveto{\pgfqpoint{4.636000in}{2.453606in}}{\pgfqpoint{4.751060in}{2.314716in}}{\pgfqpoint{4.955839in}{2.212326in}}%
\pgfpathcurveto{\pgfqpoint{5.160619in}{2.109937in}}{\pgfqpoint{5.438398in}{2.052407in}}{\pgfqpoint{5.728000in}{2.052407in}}%
\pgfpathlineto{\pgfqpoint{5.728000in}{2.052407in}}%
\pgfpathclose%
\pgfusepath{clip}%
\pgfsetrectcap%
\pgfsetroundjoin%
\pgfsetlinewidth{0.451687pt}%
\definecolor{currentstroke}{rgb}{0.690196,0.690196,0.690196}%
\pgfsetstrokecolor{currentstroke}%
\pgfsetstrokeopacity{0.250000}%
\pgfsetdash{}{0pt}%
\pgfpathmoveto{\pgfqpoint{6.113079in}{2.103702in}}%
\pgfpathlineto{\pgfqpoint{6.144843in}{2.113059in}}%
\pgfpathlineto{\pgfqpoint{6.175300in}{2.122919in}}%
\pgfpathlineto{\pgfqpoint{6.203745in}{2.132965in}}%
\pgfpathlineto{\pgfqpoint{6.230910in}{2.143363in}}%
\pgfpathlineto{\pgfqpoint{6.256866in}{2.154082in}}%
\pgfpathlineto{\pgfqpoint{6.281681in}{2.165102in}}%
\pgfpathlineto{\pgfqpoint{6.305415in}{2.176401in}}%
\pgfpathlineto{\pgfqpoint{6.328119in}{2.187962in}}%
\pgfpathlineto{\pgfqpoint{6.349840in}{2.199772in}}%
\pgfpathlineto{\pgfqpoint{6.370616in}{2.211817in}}%
\pgfpathlineto{\pgfqpoint{6.390483in}{2.224083in}}%
\pgfpathlineto{\pgfqpoint{6.409468in}{2.236561in}}%
\pgfpathlineto{\pgfqpoint{6.427597in}{2.249238in}}%
\pgfpathlineto{\pgfqpoint{6.444894in}{2.262105in}}%
\pgfpathlineto{\pgfqpoint{6.461574in}{2.275314in}}%
\pgfpathlineto{\pgfqpoint{6.477194in}{2.288486in}}%
\pgfpathlineto{\pgfqpoint{6.492052in}{2.301831in}}%
\pgfpathlineto{\pgfqpoint{6.506153in}{2.315338in}}%
\pgfpathlineto{\pgfqpoint{6.519506in}{2.328996in}}%
\pgfpathlineto{\pgfqpoint{6.532118in}{2.342794in}}%
\pgfpathlineto{\pgfqpoint{6.543996in}{2.356724in}}%
\pgfpathlineto{\pgfqpoint{6.555145in}{2.370777in}}%
\pgfpathlineto{\pgfqpoint{6.565572in}{2.384944in}}%
\pgfpathlineto{\pgfqpoint{6.575282in}{2.399216in}}%
\pgfpathlineto{\pgfqpoint{6.584280in}{2.413586in}}%
\pgfpathlineto{\pgfqpoint{6.592570in}{2.428046in}}%
\pgfpathlineto{\pgfqpoint{6.600156in}{2.442588in}}%
\pgfpathlineto{\pgfqpoint{6.607043in}{2.457204in}}%
\pgfpathlineto{\pgfqpoint{6.613232in}{2.471888in}}%
\pgfpathlineto{\pgfqpoint{6.618728in}{2.486632in}}%
\pgfpathlineto{\pgfqpoint{6.623531in}{2.501429in}}%
\pgfpathlineto{\pgfqpoint{6.627686in}{2.516436in}}%
\pgfpathlineto{\pgfqpoint{6.631089in}{2.531243in}}%
\pgfpathlineto{\pgfqpoint{6.633816in}{2.546108in}}%
\pgfpathlineto{\pgfqpoint{6.635864in}{2.561018in}}%
\pgfpathlineto{\pgfqpoint{6.637231in}{2.575961in}}%
\pgfpathlineto{\pgfqpoint{6.637915in}{2.590923in}}%
\pgfpathlineto{\pgfqpoint{6.637915in}{2.605891in}}%
\pgfpathlineto{\pgfqpoint{6.637231in}{2.620853in}}%
\pgfpathlineto{\pgfqpoint{6.635864in}{2.635795in}}%
\pgfpathlineto{\pgfqpoint{6.633816in}{2.650705in}}%
\pgfpathlineto{\pgfqpoint{6.631089in}{2.665570in}}%
\pgfpathlineto{\pgfqpoint{6.627686in}{2.680378in}}%
\pgfpathlineto{\pgfqpoint{6.623531in}{2.695384in}}%
\pgfpathlineto{\pgfqpoint{6.618728in}{2.710181in}}%
\pgfpathlineto{\pgfqpoint{6.613232in}{2.724925in}}%
\pgfpathlineto{\pgfqpoint{6.607043in}{2.739609in}}%
\pgfpathlineto{\pgfqpoint{6.600156in}{2.754226in}}%
\pgfpathlineto{\pgfqpoint{6.592570in}{2.768768in}}%
\pgfpathlineto{\pgfqpoint{6.584280in}{2.783228in}}%
\pgfpathlineto{\pgfqpoint{6.575282in}{2.797598in}}%
\pgfpathlineto{\pgfqpoint{6.565572in}{2.811870in}}%
\pgfpathlineto{\pgfqpoint{6.555145in}{2.826037in}}%
\pgfpathlineto{\pgfqpoint{6.543996in}{2.840089in}}%
\pgfpathlineto{\pgfqpoint{6.532118in}{2.854019in}}%
\pgfpathlineto{\pgfqpoint{6.519506in}{2.867818in}}%
\pgfpathlineto{\pgfqpoint{6.506153in}{2.881476in}}%
\pgfpathlineto{\pgfqpoint{6.492052in}{2.894982in}}%
\pgfpathlineto{\pgfqpoint{6.477194in}{2.908327in}}%
\pgfpathlineto{\pgfqpoint{6.461574in}{2.921500in}}%
\pgfpathlineto{\pgfqpoint{6.444894in}{2.934708in}}%
\pgfpathlineto{\pgfqpoint{6.427597in}{2.947576in}}%
\pgfpathlineto{\pgfqpoint{6.409468in}{2.960253in}}%
\pgfpathlineto{\pgfqpoint{6.390483in}{2.972730in}}%
\pgfpathlineto{\pgfqpoint{6.370616in}{2.984997in}}%
\pgfpathlineto{\pgfqpoint{6.349840in}{2.997041in}}%
\pgfpathlineto{\pgfqpoint{6.328119in}{3.008851in}}%
\pgfpathlineto{\pgfqpoint{6.305415in}{3.020413in}}%
\pgfpathlineto{\pgfqpoint{6.281681in}{3.031712in}}%
\pgfpathlineto{\pgfqpoint{6.256866in}{3.042731in}}%
\pgfpathlineto{\pgfqpoint{6.230910in}{3.053451in}}%
\pgfpathlineto{\pgfqpoint{6.203745in}{3.063848in}}%
\pgfpathlineto{\pgfqpoint{6.175300in}{3.073894in}}%
\pgfpathlineto{\pgfqpoint{6.144843in}{3.083755in}}%
\pgfpathlineto{\pgfqpoint{6.113079in}{3.093112in}}%
\pgfusepath{stroke}%
\end{pgfscope}%
\begin{pgfscope}%
\pgfpathmoveto{\pgfqpoint{5.728000in}{2.052407in}}%
\pgfpathcurveto{\pgfqpoint{6.017602in}{2.052407in}}{\pgfqpoint{6.295381in}{2.109937in}}{\pgfqpoint{6.500161in}{2.212326in}}%
\pgfpathcurveto{\pgfqpoint{6.704940in}{2.314716in}}{\pgfqpoint{6.820000in}{2.453606in}}{\pgfqpoint{6.820000in}{2.598407in}}%
\pgfpathcurveto{\pgfqpoint{6.820000in}{2.743208in}}{\pgfqpoint{6.704940in}{2.882097in}}{\pgfqpoint{6.500161in}{2.984487in}}%
\pgfpathcurveto{\pgfqpoint{6.295381in}{3.086877in}}{\pgfqpoint{6.017602in}{3.144407in}}{\pgfqpoint{5.728000in}{3.144407in}}%
\pgfpathcurveto{\pgfqpoint{5.438398in}{3.144407in}}{\pgfqpoint{5.160619in}{3.086877in}}{\pgfqpoint{4.955839in}{2.984487in}}%
\pgfpathcurveto{\pgfqpoint{4.751060in}{2.882097in}}{\pgfqpoint{4.636000in}{2.743208in}}{\pgfqpoint{4.636000in}{2.598407in}}%
\pgfpathcurveto{\pgfqpoint{4.636000in}{2.453606in}}{\pgfqpoint{4.751060in}{2.314716in}}{\pgfqpoint{4.955839in}{2.212326in}}%
\pgfpathcurveto{\pgfqpoint{5.160619in}{2.109937in}}{\pgfqpoint{5.438398in}{2.052407in}}{\pgfqpoint{5.728000in}{2.052407in}}%
\pgfpathlineto{\pgfqpoint{5.728000in}{2.052407in}}%
\pgfpathclose%
\pgfusepath{clip}%
\pgfsetrectcap%
\pgfsetroundjoin%
\pgfsetlinewidth{0.451687pt}%
\definecolor{currentstroke}{rgb}{0.690196,0.690196,0.690196}%
\pgfsetstrokecolor{currentstroke}%
\pgfsetstrokeopacity{0.250000}%
\pgfsetdash{}{0pt}%
\pgfpathmoveto{\pgfqpoint{5.265906in}{2.103702in}}%
\pgfpathlineto{\pgfqpoint{5.278228in}{2.103702in}}%
\pgfpathlineto{\pgfqpoint{5.290551in}{2.103702in}}%
\pgfpathlineto{\pgfqpoint{5.302873in}{2.103702in}}%
\pgfpathlineto{\pgfqpoint{5.315196in}{2.103702in}}%
\pgfpathlineto{\pgfqpoint{5.327518in}{2.103702in}}%
\pgfpathlineto{\pgfqpoint{5.339841in}{2.103702in}}%
\pgfpathlineto{\pgfqpoint{5.352163in}{2.103702in}}%
\pgfpathlineto{\pgfqpoint{5.364486in}{2.103702in}}%
\pgfpathlineto{\pgfqpoint{5.376808in}{2.103702in}}%
\pgfpathlineto{\pgfqpoint{5.389131in}{2.103702in}}%
\pgfpathlineto{\pgfqpoint{5.401453in}{2.103702in}}%
\pgfpathlineto{\pgfqpoint{5.413776in}{2.103702in}}%
\pgfpathlineto{\pgfqpoint{5.426098in}{2.103702in}}%
\pgfpathlineto{\pgfqpoint{5.438421in}{2.103702in}}%
\pgfpathlineto{\pgfqpoint{5.450743in}{2.103702in}}%
\pgfpathlineto{\pgfqpoint{5.463066in}{2.103702in}}%
\pgfpathlineto{\pgfqpoint{5.475388in}{2.103702in}}%
\pgfpathlineto{\pgfqpoint{5.487711in}{2.103702in}}%
\pgfpathlineto{\pgfqpoint{5.500033in}{2.103702in}}%
\pgfpathlineto{\pgfqpoint{5.512356in}{2.103702in}}%
\pgfpathlineto{\pgfqpoint{5.524678in}{2.103702in}}%
\pgfpathlineto{\pgfqpoint{5.537001in}{2.103702in}}%
\pgfpathlineto{\pgfqpoint{5.549323in}{2.103702in}}%
\pgfpathlineto{\pgfqpoint{5.561646in}{2.103702in}}%
\pgfpathlineto{\pgfqpoint{5.573969in}{2.103702in}}%
\pgfpathlineto{\pgfqpoint{5.586291in}{2.103702in}}%
\pgfpathlineto{\pgfqpoint{5.598614in}{2.103702in}}%
\pgfpathlineto{\pgfqpoint{5.610936in}{2.103702in}}%
\pgfpathlineto{\pgfqpoint{5.623259in}{2.103702in}}%
\pgfpathlineto{\pgfqpoint{5.635581in}{2.103702in}}%
\pgfpathlineto{\pgfqpoint{5.647904in}{2.103702in}}%
\pgfpathlineto{\pgfqpoint{5.660226in}{2.103702in}}%
\pgfpathlineto{\pgfqpoint{5.672549in}{2.103702in}}%
\pgfpathlineto{\pgfqpoint{5.684871in}{2.103702in}}%
\pgfpathlineto{\pgfqpoint{5.697194in}{2.103702in}}%
\pgfpathlineto{\pgfqpoint{5.709516in}{2.103702in}}%
\pgfpathlineto{\pgfqpoint{5.721839in}{2.103702in}}%
\pgfpathlineto{\pgfqpoint{5.734161in}{2.103702in}}%
\pgfpathlineto{\pgfqpoint{5.746484in}{2.103702in}}%
\pgfpathlineto{\pgfqpoint{5.758806in}{2.103702in}}%
\pgfpathlineto{\pgfqpoint{5.771129in}{2.103702in}}%
\pgfpathlineto{\pgfqpoint{5.783451in}{2.103702in}}%
\pgfpathlineto{\pgfqpoint{5.795774in}{2.103702in}}%
\pgfpathlineto{\pgfqpoint{5.808096in}{2.103702in}}%
\pgfpathlineto{\pgfqpoint{5.820419in}{2.103702in}}%
\pgfpathlineto{\pgfqpoint{5.832741in}{2.103702in}}%
\pgfpathlineto{\pgfqpoint{5.845064in}{2.103702in}}%
\pgfpathlineto{\pgfqpoint{5.857386in}{2.103702in}}%
\pgfpathlineto{\pgfqpoint{5.869709in}{2.103702in}}%
\pgfpathlineto{\pgfqpoint{5.882031in}{2.103702in}}%
\pgfpathlineto{\pgfqpoint{5.894354in}{2.103702in}}%
\pgfpathlineto{\pgfqpoint{5.906677in}{2.103702in}}%
\pgfpathlineto{\pgfqpoint{5.918999in}{2.103702in}}%
\pgfpathlineto{\pgfqpoint{5.931322in}{2.103702in}}%
\pgfpathlineto{\pgfqpoint{5.943644in}{2.103702in}}%
\pgfpathlineto{\pgfqpoint{5.955967in}{2.103702in}}%
\pgfpathlineto{\pgfqpoint{5.968289in}{2.103702in}}%
\pgfpathlineto{\pgfqpoint{5.980612in}{2.103702in}}%
\pgfpathlineto{\pgfqpoint{5.992934in}{2.103702in}}%
\pgfpathlineto{\pgfqpoint{6.005257in}{2.103702in}}%
\pgfpathlineto{\pgfqpoint{6.017579in}{2.103702in}}%
\pgfpathlineto{\pgfqpoint{6.029902in}{2.103702in}}%
\pgfpathlineto{\pgfqpoint{6.042224in}{2.103702in}}%
\pgfpathlineto{\pgfqpoint{6.054547in}{2.103702in}}%
\pgfpathlineto{\pgfqpoint{6.066869in}{2.103702in}}%
\pgfpathlineto{\pgfqpoint{6.079192in}{2.103702in}}%
\pgfpathlineto{\pgfqpoint{6.091514in}{2.103702in}}%
\pgfpathlineto{\pgfqpoint{6.103837in}{2.103702in}}%
\pgfpathlineto{\pgfqpoint{6.116159in}{2.103702in}}%
\pgfpathlineto{\pgfqpoint{6.128482in}{2.103702in}}%
\pgfpathlineto{\pgfqpoint{6.140804in}{2.103702in}}%
\pgfpathlineto{\pgfqpoint{6.153127in}{2.103702in}}%
\pgfpathlineto{\pgfqpoint{6.165449in}{2.103702in}}%
\pgfpathlineto{\pgfqpoint{6.177772in}{2.103702in}}%
\pgfpathlineto{\pgfqpoint{6.190094in}{2.103702in}}%
\pgfusepath{stroke}%
\end{pgfscope}%
\begin{pgfscope}%
\pgfpathmoveto{\pgfqpoint{5.728000in}{2.052407in}}%
\pgfpathcurveto{\pgfqpoint{6.017602in}{2.052407in}}{\pgfqpoint{6.295381in}{2.109937in}}{\pgfqpoint{6.500161in}{2.212326in}}%
\pgfpathcurveto{\pgfqpoint{6.704940in}{2.314716in}}{\pgfqpoint{6.820000in}{2.453606in}}{\pgfqpoint{6.820000in}{2.598407in}}%
\pgfpathcurveto{\pgfqpoint{6.820000in}{2.743208in}}{\pgfqpoint{6.704940in}{2.882097in}}{\pgfqpoint{6.500161in}{2.984487in}}%
\pgfpathcurveto{\pgfqpoint{6.295381in}{3.086877in}}{\pgfqpoint{6.017602in}{3.144407in}}{\pgfqpoint{5.728000in}{3.144407in}}%
\pgfpathcurveto{\pgfqpoint{5.438398in}{3.144407in}}{\pgfqpoint{5.160619in}{3.086877in}}{\pgfqpoint{4.955839in}{2.984487in}}%
\pgfpathcurveto{\pgfqpoint{4.751060in}{2.882097in}}{\pgfqpoint{4.636000in}{2.743208in}}{\pgfqpoint{4.636000in}{2.598407in}}%
\pgfpathcurveto{\pgfqpoint{4.636000in}{2.453606in}}{\pgfqpoint{4.751060in}{2.314716in}}{\pgfqpoint{4.955839in}{2.212326in}}%
\pgfpathcurveto{\pgfqpoint{5.160619in}{2.109937in}}{\pgfqpoint{5.438398in}{2.052407in}}{\pgfqpoint{5.728000in}{2.052407in}}%
\pgfpathlineto{\pgfqpoint{5.728000in}{2.052407in}}%
\pgfpathclose%
\pgfusepath{clip}%
\pgfsetrectcap%
\pgfsetroundjoin%
\pgfsetlinewidth{0.451687pt}%
\definecolor{currentstroke}{rgb}{0.690196,0.690196,0.690196}%
\pgfsetstrokecolor{currentstroke}%
\pgfsetstrokeopacity{0.250000}%
\pgfsetdash{}{0pt}%
\pgfpathmoveto{\pgfqpoint{5.021329in}{2.182150in}}%
\pgfpathlineto{\pgfqpoint{5.040174in}{2.182150in}}%
\pgfpathlineto{\pgfqpoint{5.059018in}{2.182150in}}%
\pgfpathlineto{\pgfqpoint{5.077863in}{2.182150in}}%
\pgfpathlineto{\pgfqpoint{5.096707in}{2.182150in}}%
\pgfpathlineto{\pgfqpoint{5.115552in}{2.182150in}}%
\pgfpathlineto{\pgfqpoint{5.134396in}{2.182150in}}%
\pgfpathlineto{\pgfqpoint{5.153241in}{2.182150in}}%
\pgfpathlineto{\pgfqpoint{5.172086in}{2.182150in}}%
\pgfpathlineto{\pgfqpoint{5.190930in}{2.182150in}}%
\pgfpathlineto{\pgfqpoint{5.209775in}{2.182150in}}%
\pgfpathlineto{\pgfqpoint{5.228619in}{2.182150in}}%
\pgfpathlineto{\pgfqpoint{5.247464in}{2.182150in}}%
\pgfpathlineto{\pgfqpoint{5.266308in}{2.182150in}}%
\pgfpathlineto{\pgfqpoint{5.285153in}{2.182150in}}%
\pgfpathlineto{\pgfqpoint{5.303997in}{2.182150in}}%
\pgfpathlineto{\pgfqpoint{5.322842in}{2.182150in}}%
\pgfpathlineto{\pgfqpoint{5.341687in}{2.182150in}}%
\pgfpathlineto{\pgfqpoint{5.360531in}{2.182150in}}%
\pgfpathlineto{\pgfqpoint{5.379376in}{2.182150in}}%
\pgfpathlineto{\pgfqpoint{5.398220in}{2.182150in}}%
\pgfpathlineto{\pgfqpoint{5.417065in}{2.182150in}}%
\pgfpathlineto{\pgfqpoint{5.435909in}{2.182150in}}%
\pgfpathlineto{\pgfqpoint{5.454754in}{2.182150in}}%
\pgfpathlineto{\pgfqpoint{5.473598in}{2.182150in}}%
\pgfpathlineto{\pgfqpoint{5.492443in}{2.182150in}}%
\pgfpathlineto{\pgfqpoint{5.511288in}{2.182150in}}%
\pgfpathlineto{\pgfqpoint{5.530132in}{2.182150in}}%
\pgfpathlineto{\pgfqpoint{5.548977in}{2.182150in}}%
\pgfpathlineto{\pgfqpoint{5.567821in}{2.182150in}}%
\pgfpathlineto{\pgfqpoint{5.586666in}{2.182150in}}%
\pgfpathlineto{\pgfqpoint{5.605510in}{2.182150in}}%
\pgfpathlineto{\pgfqpoint{5.624355in}{2.182150in}}%
\pgfpathlineto{\pgfqpoint{5.643199in}{2.182150in}}%
\pgfpathlineto{\pgfqpoint{5.662044in}{2.182150in}}%
\pgfpathlineto{\pgfqpoint{5.680889in}{2.182150in}}%
\pgfpathlineto{\pgfqpoint{5.699733in}{2.182150in}}%
\pgfpathlineto{\pgfqpoint{5.718578in}{2.182150in}}%
\pgfpathlineto{\pgfqpoint{5.737422in}{2.182150in}}%
\pgfpathlineto{\pgfqpoint{5.756267in}{2.182150in}}%
\pgfpathlineto{\pgfqpoint{5.775111in}{2.182150in}}%
\pgfpathlineto{\pgfqpoint{5.793956in}{2.182150in}}%
\pgfpathlineto{\pgfqpoint{5.812801in}{2.182150in}}%
\pgfpathlineto{\pgfqpoint{5.831645in}{2.182150in}}%
\pgfpathlineto{\pgfqpoint{5.850490in}{2.182150in}}%
\pgfpathlineto{\pgfqpoint{5.869334in}{2.182150in}}%
\pgfpathlineto{\pgfqpoint{5.888179in}{2.182150in}}%
\pgfpathlineto{\pgfqpoint{5.907023in}{2.182150in}}%
\pgfpathlineto{\pgfqpoint{5.925868in}{2.182150in}}%
\pgfpathlineto{\pgfqpoint{5.944712in}{2.182150in}}%
\pgfpathlineto{\pgfqpoint{5.963557in}{2.182150in}}%
\pgfpathlineto{\pgfqpoint{5.982402in}{2.182150in}}%
\pgfpathlineto{\pgfqpoint{6.001246in}{2.182150in}}%
\pgfpathlineto{\pgfqpoint{6.020091in}{2.182150in}}%
\pgfpathlineto{\pgfqpoint{6.038935in}{2.182150in}}%
\pgfpathlineto{\pgfqpoint{6.057780in}{2.182150in}}%
\pgfpathlineto{\pgfqpoint{6.076624in}{2.182150in}}%
\pgfpathlineto{\pgfqpoint{6.095469in}{2.182150in}}%
\pgfpathlineto{\pgfqpoint{6.114313in}{2.182150in}}%
\pgfpathlineto{\pgfqpoint{6.133158in}{2.182150in}}%
\pgfpathlineto{\pgfqpoint{6.152003in}{2.182150in}}%
\pgfpathlineto{\pgfqpoint{6.170847in}{2.182150in}}%
\pgfpathlineto{\pgfqpoint{6.189692in}{2.182150in}}%
\pgfpathlineto{\pgfqpoint{6.208536in}{2.182150in}}%
\pgfpathlineto{\pgfqpoint{6.227381in}{2.182150in}}%
\pgfpathlineto{\pgfqpoint{6.246225in}{2.182150in}}%
\pgfpathlineto{\pgfqpoint{6.265070in}{2.182150in}}%
\pgfpathlineto{\pgfqpoint{6.283914in}{2.182150in}}%
\pgfpathlineto{\pgfqpoint{6.302759in}{2.182150in}}%
\pgfpathlineto{\pgfqpoint{6.321604in}{2.182150in}}%
\pgfpathlineto{\pgfqpoint{6.340448in}{2.182150in}}%
\pgfpathlineto{\pgfqpoint{6.359293in}{2.182150in}}%
\pgfpathlineto{\pgfqpoint{6.378137in}{2.182150in}}%
\pgfpathlineto{\pgfqpoint{6.396982in}{2.182150in}}%
\pgfpathlineto{\pgfqpoint{6.415826in}{2.182150in}}%
\pgfpathlineto{\pgfqpoint{6.434671in}{2.182150in}}%
\pgfusepath{stroke}%
\end{pgfscope}%
\begin{pgfscope}%
\pgfpathmoveto{\pgfqpoint{5.728000in}{2.052407in}}%
\pgfpathcurveto{\pgfqpoint{6.017602in}{2.052407in}}{\pgfqpoint{6.295381in}{2.109937in}}{\pgfqpoint{6.500161in}{2.212326in}}%
\pgfpathcurveto{\pgfqpoint{6.704940in}{2.314716in}}{\pgfqpoint{6.820000in}{2.453606in}}{\pgfqpoint{6.820000in}{2.598407in}}%
\pgfpathcurveto{\pgfqpoint{6.820000in}{2.743208in}}{\pgfqpoint{6.704940in}{2.882097in}}{\pgfqpoint{6.500161in}{2.984487in}}%
\pgfpathcurveto{\pgfqpoint{6.295381in}{3.086877in}}{\pgfqpoint{6.017602in}{3.144407in}}{\pgfqpoint{5.728000in}{3.144407in}}%
\pgfpathcurveto{\pgfqpoint{5.438398in}{3.144407in}}{\pgfqpoint{5.160619in}{3.086877in}}{\pgfqpoint{4.955839in}{2.984487in}}%
\pgfpathcurveto{\pgfqpoint{4.751060in}{2.882097in}}{\pgfqpoint{4.636000in}{2.743208in}}{\pgfqpoint{4.636000in}{2.598407in}}%
\pgfpathcurveto{\pgfqpoint{4.636000in}{2.453606in}}{\pgfqpoint{4.751060in}{2.314716in}}{\pgfqpoint{4.955839in}{2.212326in}}%
\pgfpathcurveto{\pgfqpoint{5.160619in}{2.109937in}}{\pgfqpoint{5.438398in}{2.052407in}}{\pgfqpoint{5.728000in}{2.052407in}}%
\pgfpathlineto{\pgfqpoint{5.728000in}{2.052407in}}%
\pgfpathclose%
\pgfusepath{clip}%
\pgfsetrectcap%
\pgfsetroundjoin%
\pgfsetlinewidth{0.451687pt}%
\definecolor{currentstroke}{rgb}{0.690196,0.690196,0.690196}%
\pgfsetstrokecolor{currentstroke}%
\pgfsetstrokeopacity{0.250000}%
\pgfsetdash{}{0pt}%
\pgfpathmoveto{\pgfqpoint{4.847711in}{2.275314in}}%
\pgfpathlineto{\pgfqpoint{4.871186in}{2.275314in}}%
\pgfpathlineto{\pgfqpoint{4.894660in}{2.275314in}}%
\pgfpathlineto{\pgfqpoint{4.918134in}{2.275314in}}%
\pgfpathlineto{\pgfqpoint{4.941609in}{2.275314in}}%
\pgfpathlineto{\pgfqpoint{4.965083in}{2.275314in}}%
\pgfpathlineto{\pgfqpoint{4.988558in}{2.275314in}}%
\pgfpathlineto{\pgfqpoint{5.012032in}{2.275314in}}%
\pgfpathlineto{\pgfqpoint{5.035506in}{2.275314in}}%
\pgfpathlineto{\pgfqpoint{5.058981in}{2.275314in}}%
\pgfpathlineto{\pgfqpoint{5.082455in}{2.275314in}}%
\pgfpathlineto{\pgfqpoint{5.105929in}{2.275314in}}%
\pgfpathlineto{\pgfqpoint{5.129404in}{2.275314in}}%
\pgfpathlineto{\pgfqpoint{5.152878in}{2.275314in}}%
\pgfpathlineto{\pgfqpoint{5.176352in}{2.275314in}}%
\pgfpathlineto{\pgfqpoint{5.199827in}{2.275314in}}%
\pgfpathlineto{\pgfqpoint{5.223301in}{2.275314in}}%
\pgfpathlineto{\pgfqpoint{5.246776in}{2.275314in}}%
\pgfpathlineto{\pgfqpoint{5.270250in}{2.275314in}}%
\pgfpathlineto{\pgfqpoint{5.293724in}{2.275314in}}%
\pgfpathlineto{\pgfqpoint{5.317199in}{2.275314in}}%
\pgfpathlineto{\pgfqpoint{5.340673in}{2.275314in}}%
\pgfpathlineto{\pgfqpoint{5.364147in}{2.275314in}}%
\pgfpathlineto{\pgfqpoint{5.387622in}{2.275314in}}%
\pgfpathlineto{\pgfqpoint{5.411096in}{2.275314in}}%
\pgfpathlineto{\pgfqpoint{5.434570in}{2.275314in}}%
\pgfpathlineto{\pgfqpoint{5.458045in}{2.275314in}}%
\pgfpathlineto{\pgfqpoint{5.481519in}{2.275314in}}%
\pgfpathlineto{\pgfqpoint{5.504994in}{2.275314in}}%
\pgfpathlineto{\pgfqpoint{5.528468in}{2.275314in}}%
\pgfpathlineto{\pgfqpoint{5.551942in}{2.275314in}}%
\pgfpathlineto{\pgfqpoint{5.575417in}{2.275314in}}%
\pgfpathlineto{\pgfqpoint{5.598891in}{2.275314in}}%
\pgfpathlineto{\pgfqpoint{5.622365in}{2.275314in}}%
\pgfpathlineto{\pgfqpoint{5.645840in}{2.275314in}}%
\pgfpathlineto{\pgfqpoint{5.669314in}{2.275314in}}%
\pgfpathlineto{\pgfqpoint{5.692788in}{2.275314in}}%
\pgfpathlineto{\pgfqpoint{5.716263in}{2.275314in}}%
\pgfpathlineto{\pgfqpoint{5.739737in}{2.275314in}}%
\pgfpathlineto{\pgfqpoint{5.763212in}{2.275314in}}%
\pgfpathlineto{\pgfqpoint{5.786686in}{2.275314in}}%
\pgfpathlineto{\pgfqpoint{5.810160in}{2.275314in}}%
\pgfpathlineto{\pgfqpoint{5.833635in}{2.275314in}}%
\pgfpathlineto{\pgfqpoint{5.857109in}{2.275314in}}%
\pgfpathlineto{\pgfqpoint{5.880583in}{2.275314in}}%
\pgfpathlineto{\pgfqpoint{5.904058in}{2.275314in}}%
\pgfpathlineto{\pgfqpoint{5.927532in}{2.275314in}}%
\pgfpathlineto{\pgfqpoint{5.951006in}{2.275314in}}%
\pgfpathlineto{\pgfqpoint{5.974481in}{2.275314in}}%
\pgfpathlineto{\pgfqpoint{5.997955in}{2.275314in}}%
\pgfpathlineto{\pgfqpoint{6.021430in}{2.275314in}}%
\pgfpathlineto{\pgfqpoint{6.044904in}{2.275314in}}%
\pgfpathlineto{\pgfqpoint{6.068378in}{2.275314in}}%
\pgfpathlineto{\pgfqpoint{6.091853in}{2.275314in}}%
\pgfpathlineto{\pgfqpoint{6.115327in}{2.275314in}}%
\pgfpathlineto{\pgfqpoint{6.138801in}{2.275314in}}%
\pgfpathlineto{\pgfqpoint{6.162276in}{2.275314in}}%
\pgfpathlineto{\pgfqpoint{6.185750in}{2.275314in}}%
\pgfpathlineto{\pgfqpoint{6.209224in}{2.275314in}}%
\pgfpathlineto{\pgfqpoint{6.232699in}{2.275314in}}%
\pgfpathlineto{\pgfqpoint{6.256173in}{2.275314in}}%
\pgfpathlineto{\pgfqpoint{6.279648in}{2.275314in}}%
\pgfpathlineto{\pgfqpoint{6.303122in}{2.275314in}}%
\pgfpathlineto{\pgfqpoint{6.326596in}{2.275314in}}%
\pgfpathlineto{\pgfqpoint{6.350071in}{2.275314in}}%
\pgfpathlineto{\pgfqpoint{6.373545in}{2.275314in}}%
\pgfpathlineto{\pgfqpoint{6.397019in}{2.275314in}}%
\pgfpathlineto{\pgfqpoint{6.420494in}{2.275314in}}%
\pgfpathlineto{\pgfqpoint{6.443968in}{2.275314in}}%
\pgfpathlineto{\pgfqpoint{6.467442in}{2.275314in}}%
\pgfpathlineto{\pgfqpoint{6.490917in}{2.275314in}}%
\pgfpathlineto{\pgfqpoint{6.514391in}{2.275314in}}%
\pgfpathlineto{\pgfqpoint{6.537866in}{2.275314in}}%
\pgfpathlineto{\pgfqpoint{6.561340in}{2.275314in}}%
\pgfpathlineto{\pgfqpoint{6.584814in}{2.275314in}}%
\pgfpathlineto{\pgfqpoint{6.608289in}{2.275314in}}%
\pgfusepath{stroke}%
\end{pgfscope}%
\begin{pgfscope}%
\pgfpathmoveto{\pgfqpoint{5.728000in}{2.052407in}}%
\pgfpathcurveto{\pgfqpoint{6.017602in}{2.052407in}}{\pgfqpoint{6.295381in}{2.109937in}}{\pgfqpoint{6.500161in}{2.212326in}}%
\pgfpathcurveto{\pgfqpoint{6.704940in}{2.314716in}}{\pgfqpoint{6.820000in}{2.453606in}}{\pgfqpoint{6.820000in}{2.598407in}}%
\pgfpathcurveto{\pgfqpoint{6.820000in}{2.743208in}}{\pgfqpoint{6.704940in}{2.882097in}}{\pgfqpoint{6.500161in}{2.984487in}}%
\pgfpathcurveto{\pgfqpoint{6.295381in}{3.086877in}}{\pgfqpoint{6.017602in}{3.144407in}}{\pgfqpoint{5.728000in}{3.144407in}}%
\pgfpathcurveto{\pgfqpoint{5.438398in}{3.144407in}}{\pgfqpoint{5.160619in}{3.086877in}}{\pgfqpoint{4.955839in}{2.984487in}}%
\pgfpathcurveto{\pgfqpoint{4.751060in}{2.882097in}}{\pgfqpoint{4.636000in}{2.743208in}}{\pgfqpoint{4.636000in}{2.598407in}}%
\pgfpathcurveto{\pgfqpoint{4.636000in}{2.453606in}}{\pgfqpoint{4.751060in}{2.314716in}}{\pgfqpoint{4.955839in}{2.212326in}}%
\pgfpathcurveto{\pgfqpoint{5.160619in}{2.109937in}}{\pgfqpoint{5.438398in}{2.052407in}}{\pgfqpoint{5.728000in}{2.052407in}}%
\pgfpathlineto{\pgfqpoint{5.728000in}{2.052407in}}%
\pgfpathclose%
\pgfusepath{clip}%
\pgfsetrectcap%
\pgfsetroundjoin%
\pgfsetlinewidth{0.451687pt}%
\definecolor{currentstroke}{rgb}{0.690196,0.690196,0.690196}%
\pgfsetstrokecolor{currentstroke}%
\pgfsetstrokeopacity{0.250000}%
\pgfsetdash{}{0pt}%
\pgfpathmoveto{\pgfqpoint{4.729062in}{2.377847in}}%
\pgfpathlineto{\pgfqpoint{4.755701in}{2.377847in}}%
\pgfpathlineto{\pgfqpoint{4.782339in}{2.377847in}}%
\pgfpathlineto{\pgfqpoint{4.808977in}{2.377847in}}%
\pgfpathlineto{\pgfqpoint{4.835616in}{2.377847in}}%
\pgfpathlineto{\pgfqpoint{4.862254in}{2.377847in}}%
\pgfpathlineto{\pgfqpoint{4.888892in}{2.377847in}}%
\pgfpathlineto{\pgfqpoint{4.915531in}{2.377847in}}%
\pgfpathlineto{\pgfqpoint{4.942169in}{2.377847in}}%
\pgfpathlineto{\pgfqpoint{4.968807in}{2.377847in}}%
\pgfpathlineto{\pgfqpoint{4.995446in}{2.377847in}}%
\pgfpathlineto{\pgfqpoint{5.022084in}{2.377847in}}%
\pgfpathlineto{\pgfqpoint{5.048722in}{2.377847in}}%
\pgfpathlineto{\pgfqpoint{5.075361in}{2.377847in}}%
\pgfpathlineto{\pgfqpoint{5.101999in}{2.377847in}}%
\pgfpathlineto{\pgfqpoint{5.128637in}{2.377847in}}%
\pgfpathlineto{\pgfqpoint{5.155276in}{2.377847in}}%
\pgfpathlineto{\pgfqpoint{5.181914in}{2.377847in}}%
\pgfpathlineto{\pgfqpoint{5.208552in}{2.377847in}}%
\pgfpathlineto{\pgfqpoint{5.235191in}{2.377847in}}%
\pgfpathlineto{\pgfqpoint{5.261829in}{2.377847in}}%
\pgfpathlineto{\pgfqpoint{5.288467in}{2.377847in}}%
\pgfpathlineto{\pgfqpoint{5.315106in}{2.377847in}}%
\pgfpathlineto{\pgfqpoint{5.341744in}{2.377847in}}%
\pgfpathlineto{\pgfqpoint{5.368382in}{2.377847in}}%
\pgfpathlineto{\pgfqpoint{5.395021in}{2.377847in}}%
\pgfpathlineto{\pgfqpoint{5.421659in}{2.377847in}}%
\pgfpathlineto{\pgfqpoint{5.448297in}{2.377847in}}%
\pgfpathlineto{\pgfqpoint{5.474936in}{2.377847in}}%
\pgfpathlineto{\pgfqpoint{5.501574in}{2.377847in}}%
\pgfpathlineto{\pgfqpoint{5.528212in}{2.377847in}}%
\pgfpathlineto{\pgfqpoint{5.554851in}{2.377847in}}%
\pgfpathlineto{\pgfqpoint{5.581489in}{2.377847in}}%
\pgfpathlineto{\pgfqpoint{5.608127in}{2.377847in}}%
\pgfpathlineto{\pgfqpoint{5.634766in}{2.377847in}}%
\pgfpathlineto{\pgfqpoint{5.661404in}{2.377847in}}%
\pgfpathlineto{\pgfqpoint{5.688042in}{2.377847in}}%
\pgfpathlineto{\pgfqpoint{5.714681in}{2.377847in}}%
\pgfpathlineto{\pgfqpoint{5.741319in}{2.377847in}}%
\pgfpathlineto{\pgfqpoint{5.767958in}{2.377847in}}%
\pgfpathlineto{\pgfqpoint{5.794596in}{2.377847in}}%
\pgfpathlineto{\pgfqpoint{5.821234in}{2.377847in}}%
\pgfpathlineto{\pgfqpoint{5.847873in}{2.377847in}}%
\pgfpathlineto{\pgfqpoint{5.874511in}{2.377847in}}%
\pgfpathlineto{\pgfqpoint{5.901149in}{2.377847in}}%
\pgfpathlineto{\pgfqpoint{5.927788in}{2.377847in}}%
\pgfpathlineto{\pgfqpoint{5.954426in}{2.377847in}}%
\pgfpathlineto{\pgfqpoint{5.981064in}{2.377847in}}%
\pgfpathlineto{\pgfqpoint{6.007703in}{2.377847in}}%
\pgfpathlineto{\pgfqpoint{6.034341in}{2.377847in}}%
\pgfpathlineto{\pgfqpoint{6.060979in}{2.377847in}}%
\pgfpathlineto{\pgfqpoint{6.087618in}{2.377847in}}%
\pgfpathlineto{\pgfqpoint{6.114256in}{2.377847in}}%
\pgfpathlineto{\pgfqpoint{6.140894in}{2.377847in}}%
\pgfpathlineto{\pgfqpoint{6.167533in}{2.377847in}}%
\pgfpathlineto{\pgfqpoint{6.194171in}{2.377847in}}%
\pgfpathlineto{\pgfqpoint{6.220809in}{2.377847in}}%
\pgfpathlineto{\pgfqpoint{6.247448in}{2.377847in}}%
\pgfpathlineto{\pgfqpoint{6.274086in}{2.377847in}}%
\pgfpathlineto{\pgfqpoint{6.300724in}{2.377847in}}%
\pgfpathlineto{\pgfqpoint{6.327363in}{2.377847in}}%
\pgfpathlineto{\pgfqpoint{6.354001in}{2.377847in}}%
\pgfpathlineto{\pgfqpoint{6.380639in}{2.377847in}}%
\pgfpathlineto{\pgfqpoint{6.407278in}{2.377847in}}%
\pgfpathlineto{\pgfqpoint{6.433916in}{2.377847in}}%
\pgfpathlineto{\pgfqpoint{6.460554in}{2.377847in}}%
\pgfpathlineto{\pgfqpoint{6.487193in}{2.377847in}}%
\pgfpathlineto{\pgfqpoint{6.513831in}{2.377847in}}%
\pgfpathlineto{\pgfqpoint{6.540469in}{2.377847in}}%
\pgfpathlineto{\pgfqpoint{6.567108in}{2.377847in}}%
\pgfpathlineto{\pgfqpoint{6.593746in}{2.377847in}}%
\pgfpathlineto{\pgfqpoint{6.620384in}{2.377847in}}%
\pgfpathlineto{\pgfqpoint{6.647023in}{2.377847in}}%
\pgfpathlineto{\pgfqpoint{6.673661in}{2.377847in}}%
\pgfpathlineto{\pgfqpoint{6.700299in}{2.377847in}}%
\pgfpathlineto{\pgfqpoint{6.726938in}{2.377847in}}%
\pgfusepath{stroke}%
\end{pgfscope}%
\begin{pgfscope}%
\pgfpathmoveto{\pgfqpoint{5.728000in}{2.052407in}}%
\pgfpathcurveto{\pgfqpoint{6.017602in}{2.052407in}}{\pgfqpoint{6.295381in}{2.109937in}}{\pgfqpoint{6.500161in}{2.212326in}}%
\pgfpathcurveto{\pgfqpoint{6.704940in}{2.314716in}}{\pgfqpoint{6.820000in}{2.453606in}}{\pgfqpoint{6.820000in}{2.598407in}}%
\pgfpathcurveto{\pgfqpoint{6.820000in}{2.743208in}}{\pgfqpoint{6.704940in}{2.882097in}}{\pgfqpoint{6.500161in}{2.984487in}}%
\pgfpathcurveto{\pgfqpoint{6.295381in}{3.086877in}}{\pgfqpoint{6.017602in}{3.144407in}}{\pgfqpoint{5.728000in}{3.144407in}}%
\pgfpathcurveto{\pgfqpoint{5.438398in}{3.144407in}}{\pgfqpoint{5.160619in}{3.086877in}}{\pgfqpoint{4.955839in}{2.984487in}}%
\pgfpathcurveto{\pgfqpoint{4.751060in}{2.882097in}}{\pgfqpoint{4.636000in}{2.743208in}}{\pgfqpoint{4.636000in}{2.598407in}}%
\pgfpathcurveto{\pgfqpoint{4.636000in}{2.453606in}}{\pgfqpoint{4.751060in}{2.314716in}}{\pgfqpoint{4.955839in}{2.212326in}}%
\pgfpathcurveto{\pgfqpoint{5.160619in}{2.109937in}}{\pgfqpoint{5.438398in}{2.052407in}}{\pgfqpoint{5.728000in}{2.052407in}}%
\pgfpathlineto{\pgfqpoint{5.728000in}{2.052407in}}%
\pgfpathclose%
\pgfusepath{clip}%
\pgfsetrectcap%
\pgfsetroundjoin%
\pgfsetlinewidth{0.451687pt}%
\definecolor{currentstroke}{rgb}{0.690196,0.690196,0.690196}%
\pgfsetstrokecolor{currentstroke}%
\pgfsetstrokeopacity{0.250000}%
\pgfsetdash{}{0pt}%
\pgfpathmoveto{\pgfqpoint{4.659127in}{2.486632in}}%
\pgfpathlineto{\pgfqpoint{4.687630in}{2.486632in}}%
\pgfpathlineto{\pgfqpoint{4.716133in}{2.486632in}}%
\pgfpathlineto{\pgfqpoint{4.744637in}{2.486632in}}%
\pgfpathlineto{\pgfqpoint{4.773140in}{2.486632in}}%
\pgfpathlineto{\pgfqpoint{4.801643in}{2.486632in}}%
\pgfpathlineto{\pgfqpoint{4.830146in}{2.486632in}}%
\pgfpathlineto{\pgfqpoint{4.858650in}{2.486632in}}%
\pgfpathlineto{\pgfqpoint{4.887153in}{2.486632in}}%
\pgfpathlineto{\pgfqpoint{4.915656in}{2.486632in}}%
\pgfpathlineto{\pgfqpoint{4.944160in}{2.486632in}}%
\pgfpathlineto{\pgfqpoint{4.972663in}{2.486632in}}%
\pgfpathlineto{\pgfqpoint{5.001166in}{2.486632in}}%
\pgfpathlineto{\pgfqpoint{5.029669in}{2.486632in}}%
\pgfpathlineto{\pgfqpoint{5.058173in}{2.486632in}}%
\pgfpathlineto{\pgfqpoint{5.086676in}{2.486632in}}%
\pgfpathlineto{\pgfqpoint{5.115179in}{2.486632in}}%
\pgfpathlineto{\pgfqpoint{5.143683in}{2.486632in}}%
\pgfpathlineto{\pgfqpoint{5.172186in}{2.486632in}}%
\pgfpathlineto{\pgfqpoint{5.200689in}{2.486632in}}%
\pgfpathlineto{\pgfqpoint{5.229192in}{2.486632in}}%
\pgfpathlineto{\pgfqpoint{5.257696in}{2.486632in}}%
\pgfpathlineto{\pgfqpoint{5.286199in}{2.486632in}}%
\pgfpathlineto{\pgfqpoint{5.314702in}{2.486632in}}%
\pgfpathlineto{\pgfqpoint{5.343206in}{2.486632in}}%
\pgfpathlineto{\pgfqpoint{5.371709in}{2.486632in}}%
\pgfpathlineto{\pgfqpoint{5.400212in}{2.486632in}}%
\pgfpathlineto{\pgfqpoint{5.428715in}{2.486632in}}%
\pgfpathlineto{\pgfqpoint{5.457219in}{2.486632in}}%
\pgfpathlineto{\pgfqpoint{5.485722in}{2.486632in}}%
\pgfpathlineto{\pgfqpoint{5.514225in}{2.486632in}}%
\pgfpathlineto{\pgfqpoint{5.542729in}{2.486632in}}%
\pgfpathlineto{\pgfqpoint{5.571232in}{2.486632in}}%
\pgfpathlineto{\pgfqpoint{5.599735in}{2.486632in}}%
\pgfpathlineto{\pgfqpoint{5.628238in}{2.486632in}}%
\pgfpathlineto{\pgfqpoint{5.656742in}{2.486632in}}%
\pgfpathlineto{\pgfqpoint{5.685245in}{2.486632in}}%
\pgfpathlineto{\pgfqpoint{5.713748in}{2.486632in}}%
\pgfpathlineto{\pgfqpoint{5.742252in}{2.486632in}}%
\pgfpathlineto{\pgfqpoint{5.770755in}{2.486632in}}%
\pgfpathlineto{\pgfqpoint{5.799258in}{2.486632in}}%
\pgfpathlineto{\pgfqpoint{5.827762in}{2.486632in}}%
\pgfpathlineto{\pgfqpoint{5.856265in}{2.486632in}}%
\pgfpathlineto{\pgfqpoint{5.884768in}{2.486632in}}%
\pgfpathlineto{\pgfqpoint{5.913271in}{2.486632in}}%
\pgfpathlineto{\pgfqpoint{5.941775in}{2.486632in}}%
\pgfpathlineto{\pgfqpoint{5.970278in}{2.486632in}}%
\pgfpathlineto{\pgfqpoint{5.998781in}{2.486632in}}%
\pgfpathlineto{\pgfqpoint{6.027285in}{2.486632in}}%
\pgfpathlineto{\pgfqpoint{6.055788in}{2.486632in}}%
\pgfpathlineto{\pgfqpoint{6.084291in}{2.486632in}}%
\pgfpathlineto{\pgfqpoint{6.112794in}{2.486632in}}%
\pgfpathlineto{\pgfqpoint{6.141298in}{2.486632in}}%
\pgfpathlineto{\pgfqpoint{6.169801in}{2.486632in}}%
\pgfpathlineto{\pgfqpoint{6.198304in}{2.486632in}}%
\pgfpathlineto{\pgfqpoint{6.226808in}{2.486632in}}%
\pgfpathlineto{\pgfqpoint{6.255311in}{2.486632in}}%
\pgfpathlineto{\pgfqpoint{6.283814in}{2.486632in}}%
\pgfpathlineto{\pgfqpoint{6.312317in}{2.486632in}}%
\pgfpathlineto{\pgfqpoint{6.340821in}{2.486632in}}%
\pgfpathlineto{\pgfqpoint{6.369324in}{2.486632in}}%
\pgfpathlineto{\pgfqpoint{6.397827in}{2.486632in}}%
\pgfpathlineto{\pgfqpoint{6.426331in}{2.486632in}}%
\pgfpathlineto{\pgfqpoint{6.454834in}{2.486632in}}%
\pgfpathlineto{\pgfqpoint{6.483337in}{2.486632in}}%
\pgfpathlineto{\pgfqpoint{6.511840in}{2.486632in}}%
\pgfpathlineto{\pgfqpoint{6.540344in}{2.486632in}}%
\pgfpathlineto{\pgfqpoint{6.568847in}{2.486632in}}%
\pgfpathlineto{\pgfqpoint{6.597350in}{2.486632in}}%
\pgfpathlineto{\pgfqpoint{6.625854in}{2.486632in}}%
\pgfpathlineto{\pgfqpoint{6.654357in}{2.486632in}}%
\pgfpathlineto{\pgfqpoint{6.682860in}{2.486632in}}%
\pgfpathlineto{\pgfqpoint{6.711363in}{2.486632in}}%
\pgfpathlineto{\pgfqpoint{6.739867in}{2.486632in}}%
\pgfpathlineto{\pgfqpoint{6.768370in}{2.486632in}}%
\pgfpathlineto{\pgfqpoint{6.796873in}{2.486632in}}%
\pgfusepath{stroke}%
\end{pgfscope}%
\begin{pgfscope}%
\pgfpathmoveto{\pgfqpoint{5.728000in}{2.052407in}}%
\pgfpathcurveto{\pgfqpoint{6.017602in}{2.052407in}}{\pgfqpoint{6.295381in}{2.109937in}}{\pgfqpoint{6.500161in}{2.212326in}}%
\pgfpathcurveto{\pgfqpoint{6.704940in}{2.314716in}}{\pgfqpoint{6.820000in}{2.453606in}}{\pgfqpoint{6.820000in}{2.598407in}}%
\pgfpathcurveto{\pgfqpoint{6.820000in}{2.743208in}}{\pgfqpoint{6.704940in}{2.882097in}}{\pgfqpoint{6.500161in}{2.984487in}}%
\pgfpathcurveto{\pgfqpoint{6.295381in}{3.086877in}}{\pgfqpoint{6.017602in}{3.144407in}}{\pgfqpoint{5.728000in}{3.144407in}}%
\pgfpathcurveto{\pgfqpoint{5.438398in}{3.144407in}}{\pgfqpoint{5.160619in}{3.086877in}}{\pgfqpoint{4.955839in}{2.984487in}}%
\pgfpathcurveto{\pgfqpoint{4.751060in}{2.882097in}}{\pgfqpoint{4.636000in}{2.743208in}}{\pgfqpoint{4.636000in}{2.598407in}}%
\pgfpathcurveto{\pgfqpoint{4.636000in}{2.453606in}}{\pgfqpoint{4.751060in}{2.314716in}}{\pgfqpoint{4.955839in}{2.212326in}}%
\pgfpathcurveto{\pgfqpoint{5.160619in}{2.109937in}}{\pgfqpoint{5.438398in}{2.052407in}}{\pgfqpoint{5.728000in}{2.052407in}}%
\pgfpathlineto{\pgfqpoint{5.728000in}{2.052407in}}%
\pgfpathclose%
\pgfusepath{clip}%
\pgfsetrectcap%
\pgfsetroundjoin%
\pgfsetlinewidth{0.451687pt}%
\definecolor{currentstroke}{rgb}{0.690196,0.690196,0.690196}%
\pgfsetstrokecolor{currentstroke}%
\pgfsetstrokeopacity{0.250000}%
\pgfsetdash{}{0pt}%
\pgfpathmoveto{\pgfqpoint{4.636000in}{2.598407in}}%
\pgfpathlineto{\pgfqpoint{4.665120in}{2.598407in}}%
\pgfpathlineto{\pgfqpoint{4.694240in}{2.598407in}}%
\pgfpathlineto{\pgfqpoint{4.723360in}{2.598407in}}%
\pgfpathlineto{\pgfqpoint{4.752480in}{2.598407in}}%
\pgfpathlineto{\pgfqpoint{4.781600in}{2.598407in}}%
\pgfpathlineto{\pgfqpoint{4.810720in}{2.598407in}}%
\pgfpathlineto{\pgfqpoint{4.839840in}{2.598407in}}%
\pgfpathlineto{\pgfqpoint{4.868960in}{2.598407in}}%
\pgfpathlineto{\pgfqpoint{4.898080in}{2.598407in}}%
\pgfpathlineto{\pgfqpoint{4.927200in}{2.598407in}}%
\pgfpathlineto{\pgfqpoint{4.956320in}{2.598407in}}%
\pgfpathlineto{\pgfqpoint{4.985440in}{2.598407in}}%
\pgfpathlineto{\pgfqpoint{5.014560in}{2.598407in}}%
\pgfpathlineto{\pgfqpoint{5.043680in}{2.598407in}}%
\pgfpathlineto{\pgfqpoint{5.072800in}{2.598407in}}%
\pgfpathlineto{\pgfqpoint{5.101920in}{2.598407in}}%
\pgfpathlineto{\pgfqpoint{5.131040in}{2.598407in}}%
\pgfpathlineto{\pgfqpoint{5.160160in}{2.598407in}}%
\pgfpathlineto{\pgfqpoint{5.189280in}{2.598407in}}%
\pgfpathlineto{\pgfqpoint{5.218400in}{2.598407in}}%
\pgfpathlineto{\pgfqpoint{5.247520in}{2.598407in}}%
\pgfpathlineto{\pgfqpoint{5.276640in}{2.598407in}}%
\pgfpathlineto{\pgfqpoint{5.305760in}{2.598407in}}%
\pgfpathlineto{\pgfqpoint{5.334880in}{2.598407in}}%
\pgfpathlineto{\pgfqpoint{5.364000in}{2.598407in}}%
\pgfpathlineto{\pgfqpoint{5.393120in}{2.598407in}}%
\pgfpathlineto{\pgfqpoint{5.422240in}{2.598407in}}%
\pgfpathlineto{\pgfqpoint{5.451360in}{2.598407in}}%
\pgfpathlineto{\pgfqpoint{5.480480in}{2.598407in}}%
\pgfpathlineto{\pgfqpoint{5.509600in}{2.598407in}}%
\pgfpathlineto{\pgfqpoint{5.538720in}{2.598407in}}%
\pgfpathlineto{\pgfqpoint{5.567840in}{2.598407in}}%
\pgfpathlineto{\pgfqpoint{5.596960in}{2.598407in}}%
\pgfpathlineto{\pgfqpoint{5.626080in}{2.598407in}}%
\pgfpathlineto{\pgfqpoint{5.655200in}{2.598407in}}%
\pgfpathlineto{\pgfqpoint{5.684320in}{2.598407in}}%
\pgfpathlineto{\pgfqpoint{5.713440in}{2.598407in}}%
\pgfpathlineto{\pgfqpoint{5.742560in}{2.598407in}}%
\pgfpathlineto{\pgfqpoint{5.771680in}{2.598407in}}%
\pgfpathlineto{\pgfqpoint{5.800800in}{2.598407in}}%
\pgfpathlineto{\pgfqpoint{5.829920in}{2.598407in}}%
\pgfpathlineto{\pgfqpoint{5.859040in}{2.598407in}}%
\pgfpathlineto{\pgfqpoint{5.888160in}{2.598407in}}%
\pgfpathlineto{\pgfqpoint{5.917280in}{2.598407in}}%
\pgfpathlineto{\pgfqpoint{5.946400in}{2.598407in}}%
\pgfpathlineto{\pgfqpoint{5.975520in}{2.598407in}}%
\pgfpathlineto{\pgfqpoint{6.004640in}{2.598407in}}%
\pgfpathlineto{\pgfqpoint{6.033760in}{2.598407in}}%
\pgfpathlineto{\pgfqpoint{6.062880in}{2.598407in}}%
\pgfpathlineto{\pgfqpoint{6.092000in}{2.598407in}}%
\pgfpathlineto{\pgfqpoint{6.121120in}{2.598407in}}%
\pgfpathlineto{\pgfqpoint{6.150240in}{2.598407in}}%
\pgfpathlineto{\pgfqpoint{6.179360in}{2.598407in}}%
\pgfpathlineto{\pgfqpoint{6.208480in}{2.598407in}}%
\pgfpathlineto{\pgfqpoint{6.237600in}{2.598407in}}%
\pgfpathlineto{\pgfqpoint{6.266720in}{2.598407in}}%
\pgfpathlineto{\pgfqpoint{6.295840in}{2.598407in}}%
\pgfpathlineto{\pgfqpoint{6.324960in}{2.598407in}}%
\pgfpathlineto{\pgfqpoint{6.354080in}{2.598407in}}%
\pgfpathlineto{\pgfqpoint{6.383200in}{2.598407in}}%
\pgfpathlineto{\pgfqpoint{6.412320in}{2.598407in}}%
\pgfpathlineto{\pgfqpoint{6.441440in}{2.598407in}}%
\pgfpathlineto{\pgfqpoint{6.470560in}{2.598407in}}%
\pgfpathlineto{\pgfqpoint{6.499680in}{2.598407in}}%
\pgfpathlineto{\pgfqpoint{6.528800in}{2.598407in}}%
\pgfpathlineto{\pgfqpoint{6.557920in}{2.598407in}}%
\pgfpathlineto{\pgfqpoint{6.587040in}{2.598407in}}%
\pgfpathlineto{\pgfqpoint{6.616160in}{2.598407in}}%
\pgfpathlineto{\pgfqpoint{6.645280in}{2.598407in}}%
\pgfpathlineto{\pgfqpoint{6.674400in}{2.598407in}}%
\pgfpathlineto{\pgfqpoint{6.703520in}{2.598407in}}%
\pgfpathlineto{\pgfqpoint{6.732640in}{2.598407in}}%
\pgfpathlineto{\pgfqpoint{6.761760in}{2.598407in}}%
\pgfpathlineto{\pgfqpoint{6.790880in}{2.598407in}}%
\pgfpathlineto{\pgfqpoint{6.820000in}{2.598407in}}%
\pgfusepath{stroke}%
\end{pgfscope}%
\begin{pgfscope}%
\pgfpathmoveto{\pgfqpoint{5.728000in}{2.052407in}}%
\pgfpathcurveto{\pgfqpoint{6.017602in}{2.052407in}}{\pgfqpoint{6.295381in}{2.109937in}}{\pgfqpoint{6.500161in}{2.212326in}}%
\pgfpathcurveto{\pgfqpoint{6.704940in}{2.314716in}}{\pgfqpoint{6.820000in}{2.453606in}}{\pgfqpoint{6.820000in}{2.598407in}}%
\pgfpathcurveto{\pgfqpoint{6.820000in}{2.743208in}}{\pgfqpoint{6.704940in}{2.882097in}}{\pgfqpoint{6.500161in}{2.984487in}}%
\pgfpathcurveto{\pgfqpoint{6.295381in}{3.086877in}}{\pgfqpoint{6.017602in}{3.144407in}}{\pgfqpoint{5.728000in}{3.144407in}}%
\pgfpathcurveto{\pgfqpoint{5.438398in}{3.144407in}}{\pgfqpoint{5.160619in}{3.086877in}}{\pgfqpoint{4.955839in}{2.984487in}}%
\pgfpathcurveto{\pgfqpoint{4.751060in}{2.882097in}}{\pgfqpoint{4.636000in}{2.743208in}}{\pgfqpoint{4.636000in}{2.598407in}}%
\pgfpathcurveto{\pgfqpoint{4.636000in}{2.453606in}}{\pgfqpoint{4.751060in}{2.314716in}}{\pgfqpoint{4.955839in}{2.212326in}}%
\pgfpathcurveto{\pgfqpoint{5.160619in}{2.109937in}}{\pgfqpoint{5.438398in}{2.052407in}}{\pgfqpoint{5.728000in}{2.052407in}}%
\pgfpathlineto{\pgfqpoint{5.728000in}{2.052407in}}%
\pgfpathclose%
\pgfusepath{clip}%
\pgfsetrectcap%
\pgfsetroundjoin%
\pgfsetlinewidth{0.451687pt}%
\definecolor{currentstroke}{rgb}{0.690196,0.690196,0.690196}%
\pgfsetstrokecolor{currentstroke}%
\pgfsetstrokeopacity{0.250000}%
\pgfsetdash{}{0pt}%
\pgfpathmoveto{\pgfqpoint{4.659127in}{2.710181in}}%
\pgfpathlineto{\pgfqpoint{4.687630in}{2.710181in}}%
\pgfpathlineto{\pgfqpoint{4.716133in}{2.710181in}}%
\pgfpathlineto{\pgfqpoint{4.744637in}{2.710181in}}%
\pgfpathlineto{\pgfqpoint{4.773140in}{2.710181in}}%
\pgfpathlineto{\pgfqpoint{4.801643in}{2.710181in}}%
\pgfpathlineto{\pgfqpoint{4.830146in}{2.710181in}}%
\pgfpathlineto{\pgfqpoint{4.858650in}{2.710181in}}%
\pgfpathlineto{\pgfqpoint{4.887153in}{2.710181in}}%
\pgfpathlineto{\pgfqpoint{4.915656in}{2.710181in}}%
\pgfpathlineto{\pgfqpoint{4.944160in}{2.710181in}}%
\pgfpathlineto{\pgfqpoint{4.972663in}{2.710181in}}%
\pgfpathlineto{\pgfqpoint{5.001166in}{2.710181in}}%
\pgfpathlineto{\pgfqpoint{5.029669in}{2.710181in}}%
\pgfpathlineto{\pgfqpoint{5.058173in}{2.710181in}}%
\pgfpathlineto{\pgfqpoint{5.086676in}{2.710181in}}%
\pgfpathlineto{\pgfqpoint{5.115179in}{2.710181in}}%
\pgfpathlineto{\pgfqpoint{5.143683in}{2.710181in}}%
\pgfpathlineto{\pgfqpoint{5.172186in}{2.710181in}}%
\pgfpathlineto{\pgfqpoint{5.200689in}{2.710181in}}%
\pgfpathlineto{\pgfqpoint{5.229192in}{2.710181in}}%
\pgfpathlineto{\pgfqpoint{5.257696in}{2.710181in}}%
\pgfpathlineto{\pgfqpoint{5.286199in}{2.710181in}}%
\pgfpathlineto{\pgfqpoint{5.314702in}{2.710181in}}%
\pgfpathlineto{\pgfqpoint{5.343206in}{2.710181in}}%
\pgfpathlineto{\pgfqpoint{5.371709in}{2.710181in}}%
\pgfpathlineto{\pgfqpoint{5.400212in}{2.710181in}}%
\pgfpathlineto{\pgfqpoint{5.428715in}{2.710181in}}%
\pgfpathlineto{\pgfqpoint{5.457219in}{2.710181in}}%
\pgfpathlineto{\pgfqpoint{5.485722in}{2.710181in}}%
\pgfpathlineto{\pgfqpoint{5.514225in}{2.710181in}}%
\pgfpathlineto{\pgfqpoint{5.542729in}{2.710181in}}%
\pgfpathlineto{\pgfqpoint{5.571232in}{2.710181in}}%
\pgfpathlineto{\pgfqpoint{5.599735in}{2.710181in}}%
\pgfpathlineto{\pgfqpoint{5.628238in}{2.710181in}}%
\pgfpathlineto{\pgfqpoint{5.656742in}{2.710181in}}%
\pgfpathlineto{\pgfqpoint{5.685245in}{2.710181in}}%
\pgfpathlineto{\pgfqpoint{5.713748in}{2.710181in}}%
\pgfpathlineto{\pgfqpoint{5.742252in}{2.710181in}}%
\pgfpathlineto{\pgfqpoint{5.770755in}{2.710181in}}%
\pgfpathlineto{\pgfqpoint{5.799258in}{2.710181in}}%
\pgfpathlineto{\pgfqpoint{5.827762in}{2.710181in}}%
\pgfpathlineto{\pgfqpoint{5.856265in}{2.710181in}}%
\pgfpathlineto{\pgfqpoint{5.884768in}{2.710181in}}%
\pgfpathlineto{\pgfqpoint{5.913271in}{2.710181in}}%
\pgfpathlineto{\pgfqpoint{5.941775in}{2.710181in}}%
\pgfpathlineto{\pgfqpoint{5.970278in}{2.710181in}}%
\pgfpathlineto{\pgfqpoint{5.998781in}{2.710181in}}%
\pgfpathlineto{\pgfqpoint{6.027285in}{2.710181in}}%
\pgfpathlineto{\pgfqpoint{6.055788in}{2.710181in}}%
\pgfpathlineto{\pgfqpoint{6.084291in}{2.710181in}}%
\pgfpathlineto{\pgfqpoint{6.112794in}{2.710181in}}%
\pgfpathlineto{\pgfqpoint{6.141298in}{2.710181in}}%
\pgfpathlineto{\pgfqpoint{6.169801in}{2.710181in}}%
\pgfpathlineto{\pgfqpoint{6.198304in}{2.710181in}}%
\pgfpathlineto{\pgfqpoint{6.226808in}{2.710181in}}%
\pgfpathlineto{\pgfqpoint{6.255311in}{2.710181in}}%
\pgfpathlineto{\pgfqpoint{6.283814in}{2.710181in}}%
\pgfpathlineto{\pgfqpoint{6.312317in}{2.710181in}}%
\pgfpathlineto{\pgfqpoint{6.340821in}{2.710181in}}%
\pgfpathlineto{\pgfqpoint{6.369324in}{2.710181in}}%
\pgfpathlineto{\pgfqpoint{6.397827in}{2.710181in}}%
\pgfpathlineto{\pgfqpoint{6.426331in}{2.710181in}}%
\pgfpathlineto{\pgfqpoint{6.454834in}{2.710181in}}%
\pgfpathlineto{\pgfqpoint{6.483337in}{2.710181in}}%
\pgfpathlineto{\pgfqpoint{6.511840in}{2.710181in}}%
\pgfpathlineto{\pgfqpoint{6.540344in}{2.710181in}}%
\pgfpathlineto{\pgfqpoint{6.568847in}{2.710181in}}%
\pgfpathlineto{\pgfqpoint{6.597350in}{2.710181in}}%
\pgfpathlineto{\pgfqpoint{6.625854in}{2.710181in}}%
\pgfpathlineto{\pgfqpoint{6.654357in}{2.710181in}}%
\pgfpathlineto{\pgfqpoint{6.682860in}{2.710181in}}%
\pgfpathlineto{\pgfqpoint{6.711363in}{2.710181in}}%
\pgfpathlineto{\pgfqpoint{6.739867in}{2.710181in}}%
\pgfpathlineto{\pgfqpoint{6.768370in}{2.710181in}}%
\pgfpathlineto{\pgfqpoint{6.796873in}{2.710181in}}%
\pgfusepath{stroke}%
\end{pgfscope}%
\begin{pgfscope}%
\pgfpathmoveto{\pgfqpoint{5.728000in}{2.052407in}}%
\pgfpathcurveto{\pgfqpoint{6.017602in}{2.052407in}}{\pgfqpoint{6.295381in}{2.109937in}}{\pgfqpoint{6.500161in}{2.212326in}}%
\pgfpathcurveto{\pgfqpoint{6.704940in}{2.314716in}}{\pgfqpoint{6.820000in}{2.453606in}}{\pgfqpoint{6.820000in}{2.598407in}}%
\pgfpathcurveto{\pgfqpoint{6.820000in}{2.743208in}}{\pgfqpoint{6.704940in}{2.882097in}}{\pgfqpoint{6.500161in}{2.984487in}}%
\pgfpathcurveto{\pgfqpoint{6.295381in}{3.086877in}}{\pgfqpoint{6.017602in}{3.144407in}}{\pgfqpoint{5.728000in}{3.144407in}}%
\pgfpathcurveto{\pgfqpoint{5.438398in}{3.144407in}}{\pgfqpoint{5.160619in}{3.086877in}}{\pgfqpoint{4.955839in}{2.984487in}}%
\pgfpathcurveto{\pgfqpoint{4.751060in}{2.882097in}}{\pgfqpoint{4.636000in}{2.743208in}}{\pgfqpoint{4.636000in}{2.598407in}}%
\pgfpathcurveto{\pgfqpoint{4.636000in}{2.453606in}}{\pgfqpoint{4.751060in}{2.314716in}}{\pgfqpoint{4.955839in}{2.212326in}}%
\pgfpathcurveto{\pgfqpoint{5.160619in}{2.109937in}}{\pgfqpoint{5.438398in}{2.052407in}}{\pgfqpoint{5.728000in}{2.052407in}}%
\pgfpathlineto{\pgfqpoint{5.728000in}{2.052407in}}%
\pgfpathclose%
\pgfusepath{clip}%
\pgfsetrectcap%
\pgfsetroundjoin%
\pgfsetlinewidth{0.451687pt}%
\definecolor{currentstroke}{rgb}{0.690196,0.690196,0.690196}%
\pgfsetstrokecolor{currentstroke}%
\pgfsetstrokeopacity{0.250000}%
\pgfsetdash{}{0pt}%
\pgfpathmoveto{\pgfqpoint{4.729062in}{2.818967in}}%
\pgfpathlineto{\pgfqpoint{4.755701in}{2.818967in}}%
\pgfpathlineto{\pgfqpoint{4.782339in}{2.818967in}}%
\pgfpathlineto{\pgfqpoint{4.808977in}{2.818967in}}%
\pgfpathlineto{\pgfqpoint{4.835616in}{2.818967in}}%
\pgfpathlineto{\pgfqpoint{4.862254in}{2.818967in}}%
\pgfpathlineto{\pgfqpoint{4.888892in}{2.818967in}}%
\pgfpathlineto{\pgfqpoint{4.915531in}{2.818967in}}%
\pgfpathlineto{\pgfqpoint{4.942169in}{2.818967in}}%
\pgfpathlineto{\pgfqpoint{4.968807in}{2.818967in}}%
\pgfpathlineto{\pgfqpoint{4.995446in}{2.818967in}}%
\pgfpathlineto{\pgfqpoint{5.022084in}{2.818967in}}%
\pgfpathlineto{\pgfqpoint{5.048722in}{2.818967in}}%
\pgfpathlineto{\pgfqpoint{5.075361in}{2.818967in}}%
\pgfpathlineto{\pgfqpoint{5.101999in}{2.818967in}}%
\pgfpathlineto{\pgfqpoint{5.128637in}{2.818967in}}%
\pgfpathlineto{\pgfqpoint{5.155276in}{2.818967in}}%
\pgfpathlineto{\pgfqpoint{5.181914in}{2.818967in}}%
\pgfpathlineto{\pgfqpoint{5.208552in}{2.818967in}}%
\pgfpathlineto{\pgfqpoint{5.235191in}{2.818967in}}%
\pgfpathlineto{\pgfqpoint{5.261829in}{2.818967in}}%
\pgfpathlineto{\pgfqpoint{5.288467in}{2.818967in}}%
\pgfpathlineto{\pgfqpoint{5.315106in}{2.818967in}}%
\pgfpathlineto{\pgfqpoint{5.341744in}{2.818967in}}%
\pgfpathlineto{\pgfqpoint{5.368382in}{2.818967in}}%
\pgfpathlineto{\pgfqpoint{5.395021in}{2.818967in}}%
\pgfpathlineto{\pgfqpoint{5.421659in}{2.818967in}}%
\pgfpathlineto{\pgfqpoint{5.448297in}{2.818967in}}%
\pgfpathlineto{\pgfqpoint{5.474936in}{2.818967in}}%
\pgfpathlineto{\pgfqpoint{5.501574in}{2.818967in}}%
\pgfpathlineto{\pgfqpoint{5.528212in}{2.818967in}}%
\pgfpathlineto{\pgfqpoint{5.554851in}{2.818967in}}%
\pgfpathlineto{\pgfqpoint{5.581489in}{2.818967in}}%
\pgfpathlineto{\pgfqpoint{5.608127in}{2.818967in}}%
\pgfpathlineto{\pgfqpoint{5.634766in}{2.818967in}}%
\pgfpathlineto{\pgfqpoint{5.661404in}{2.818967in}}%
\pgfpathlineto{\pgfqpoint{5.688042in}{2.818967in}}%
\pgfpathlineto{\pgfqpoint{5.714681in}{2.818967in}}%
\pgfpathlineto{\pgfqpoint{5.741319in}{2.818967in}}%
\pgfpathlineto{\pgfqpoint{5.767958in}{2.818967in}}%
\pgfpathlineto{\pgfqpoint{5.794596in}{2.818967in}}%
\pgfpathlineto{\pgfqpoint{5.821234in}{2.818967in}}%
\pgfpathlineto{\pgfqpoint{5.847873in}{2.818967in}}%
\pgfpathlineto{\pgfqpoint{5.874511in}{2.818967in}}%
\pgfpathlineto{\pgfqpoint{5.901149in}{2.818967in}}%
\pgfpathlineto{\pgfqpoint{5.927788in}{2.818967in}}%
\pgfpathlineto{\pgfqpoint{5.954426in}{2.818967in}}%
\pgfpathlineto{\pgfqpoint{5.981064in}{2.818967in}}%
\pgfpathlineto{\pgfqpoint{6.007703in}{2.818967in}}%
\pgfpathlineto{\pgfqpoint{6.034341in}{2.818967in}}%
\pgfpathlineto{\pgfqpoint{6.060979in}{2.818967in}}%
\pgfpathlineto{\pgfqpoint{6.087618in}{2.818967in}}%
\pgfpathlineto{\pgfqpoint{6.114256in}{2.818967in}}%
\pgfpathlineto{\pgfqpoint{6.140894in}{2.818967in}}%
\pgfpathlineto{\pgfqpoint{6.167533in}{2.818967in}}%
\pgfpathlineto{\pgfqpoint{6.194171in}{2.818967in}}%
\pgfpathlineto{\pgfqpoint{6.220809in}{2.818967in}}%
\pgfpathlineto{\pgfqpoint{6.247448in}{2.818967in}}%
\pgfpathlineto{\pgfqpoint{6.274086in}{2.818967in}}%
\pgfpathlineto{\pgfqpoint{6.300724in}{2.818967in}}%
\pgfpathlineto{\pgfqpoint{6.327363in}{2.818967in}}%
\pgfpathlineto{\pgfqpoint{6.354001in}{2.818967in}}%
\pgfpathlineto{\pgfqpoint{6.380639in}{2.818967in}}%
\pgfpathlineto{\pgfqpoint{6.407278in}{2.818967in}}%
\pgfpathlineto{\pgfqpoint{6.433916in}{2.818967in}}%
\pgfpathlineto{\pgfqpoint{6.460554in}{2.818967in}}%
\pgfpathlineto{\pgfqpoint{6.487193in}{2.818967in}}%
\pgfpathlineto{\pgfqpoint{6.513831in}{2.818967in}}%
\pgfpathlineto{\pgfqpoint{6.540469in}{2.818967in}}%
\pgfpathlineto{\pgfqpoint{6.567108in}{2.818967in}}%
\pgfpathlineto{\pgfqpoint{6.593746in}{2.818967in}}%
\pgfpathlineto{\pgfqpoint{6.620384in}{2.818967in}}%
\pgfpathlineto{\pgfqpoint{6.647023in}{2.818967in}}%
\pgfpathlineto{\pgfqpoint{6.673661in}{2.818967in}}%
\pgfpathlineto{\pgfqpoint{6.700299in}{2.818967in}}%
\pgfpathlineto{\pgfqpoint{6.726938in}{2.818967in}}%
\pgfusepath{stroke}%
\end{pgfscope}%
\begin{pgfscope}%
\pgfpathmoveto{\pgfqpoint{5.728000in}{2.052407in}}%
\pgfpathcurveto{\pgfqpoint{6.017602in}{2.052407in}}{\pgfqpoint{6.295381in}{2.109937in}}{\pgfqpoint{6.500161in}{2.212326in}}%
\pgfpathcurveto{\pgfqpoint{6.704940in}{2.314716in}}{\pgfqpoint{6.820000in}{2.453606in}}{\pgfqpoint{6.820000in}{2.598407in}}%
\pgfpathcurveto{\pgfqpoint{6.820000in}{2.743208in}}{\pgfqpoint{6.704940in}{2.882097in}}{\pgfqpoint{6.500161in}{2.984487in}}%
\pgfpathcurveto{\pgfqpoint{6.295381in}{3.086877in}}{\pgfqpoint{6.017602in}{3.144407in}}{\pgfqpoint{5.728000in}{3.144407in}}%
\pgfpathcurveto{\pgfqpoint{5.438398in}{3.144407in}}{\pgfqpoint{5.160619in}{3.086877in}}{\pgfqpoint{4.955839in}{2.984487in}}%
\pgfpathcurveto{\pgfqpoint{4.751060in}{2.882097in}}{\pgfqpoint{4.636000in}{2.743208in}}{\pgfqpoint{4.636000in}{2.598407in}}%
\pgfpathcurveto{\pgfqpoint{4.636000in}{2.453606in}}{\pgfqpoint{4.751060in}{2.314716in}}{\pgfqpoint{4.955839in}{2.212326in}}%
\pgfpathcurveto{\pgfqpoint{5.160619in}{2.109937in}}{\pgfqpoint{5.438398in}{2.052407in}}{\pgfqpoint{5.728000in}{2.052407in}}%
\pgfpathlineto{\pgfqpoint{5.728000in}{2.052407in}}%
\pgfpathclose%
\pgfusepath{clip}%
\pgfsetrectcap%
\pgfsetroundjoin%
\pgfsetlinewidth{0.451687pt}%
\definecolor{currentstroke}{rgb}{0.690196,0.690196,0.690196}%
\pgfsetstrokecolor{currentstroke}%
\pgfsetstrokeopacity{0.250000}%
\pgfsetdash{}{0pt}%
\pgfpathmoveto{\pgfqpoint{4.847711in}{2.921500in}}%
\pgfpathlineto{\pgfqpoint{4.871186in}{2.921500in}}%
\pgfpathlineto{\pgfqpoint{4.894660in}{2.921500in}}%
\pgfpathlineto{\pgfqpoint{4.918134in}{2.921500in}}%
\pgfpathlineto{\pgfqpoint{4.941609in}{2.921500in}}%
\pgfpathlineto{\pgfqpoint{4.965083in}{2.921500in}}%
\pgfpathlineto{\pgfqpoint{4.988558in}{2.921500in}}%
\pgfpathlineto{\pgfqpoint{5.012032in}{2.921500in}}%
\pgfpathlineto{\pgfqpoint{5.035506in}{2.921500in}}%
\pgfpathlineto{\pgfqpoint{5.058981in}{2.921500in}}%
\pgfpathlineto{\pgfqpoint{5.082455in}{2.921500in}}%
\pgfpathlineto{\pgfqpoint{5.105929in}{2.921500in}}%
\pgfpathlineto{\pgfqpoint{5.129404in}{2.921500in}}%
\pgfpathlineto{\pgfqpoint{5.152878in}{2.921500in}}%
\pgfpathlineto{\pgfqpoint{5.176352in}{2.921500in}}%
\pgfpathlineto{\pgfqpoint{5.199827in}{2.921500in}}%
\pgfpathlineto{\pgfqpoint{5.223301in}{2.921500in}}%
\pgfpathlineto{\pgfqpoint{5.246776in}{2.921500in}}%
\pgfpathlineto{\pgfqpoint{5.270250in}{2.921500in}}%
\pgfpathlineto{\pgfqpoint{5.293724in}{2.921500in}}%
\pgfpathlineto{\pgfqpoint{5.317199in}{2.921500in}}%
\pgfpathlineto{\pgfqpoint{5.340673in}{2.921500in}}%
\pgfpathlineto{\pgfqpoint{5.364147in}{2.921500in}}%
\pgfpathlineto{\pgfqpoint{5.387622in}{2.921500in}}%
\pgfpathlineto{\pgfqpoint{5.411096in}{2.921500in}}%
\pgfpathlineto{\pgfqpoint{5.434570in}{2.921500in}}%
\pgfpathlineto{\pgfqpoint{5.458045in}{2.921500in}}%
\pgfpathlineto{\pgfqpoint{5.481519in}{2.921500in}}%
\pgfpathlineto{\pgfqpoint{5.504994in}{2.921500in}}%
\pgfpathlineto{\pgfqpoint{5.528468in}{2.921500in}}%
\pgfpathlineto{\pgfqpoint{5.551942in}{2.921500in}}%
\pgfpathlineto{\pgfqpoint{5.575417in}{2.921500in}}%
\pgfpathlineto{\pgfqpoint{5.598891in}{2.921500in}}%
\pgfpathlineto{\pgfqpoint{5.622365in}{2.921500in}}%
\pgfpathlineto{\pgfqpoint{5.645840in}{2.921500in}}%
\pgfpathlineto{\pgfqpoint{5.669314in}{2.921500in}}%
\pgfpathlineto{\pgfqpoint{5.692788in}{2.921500in}}%
\pgfpathlineto{\pgfqpoint{5.716263in}{2.921500in}}%
\pgfpathlineto{\pgfqpoint{5.739737in}{2.921500in}}%
\pgfpathlineto{\pgfqpoint{5.763212in}{2.921500in}}%
\pgfpathlineto{\pgfqpoint{5.786686in}{2.921500in}}%
\pgfpathlineto{\pgfqpoint{5.810160in}{2.921500in}}%
\pgfpathlineto{\pgfqpoint{5.833635in}{2.921500in}}%
\pgfpathlineto{\pgfqpoint{5.857109in}{2.921500in}}%
\pgfpathlineto{\pgfqpoint{5.880583in}{2.921500in}}%
\pgfpathlineto{\pgfqpoint{5.904058in}{2.921500in}}%
\pgfpathlineto{\pgfqpoint{5.927532in}{2.921500in}}%
\pgfpathlineto{\pgfqpoint{5.951006in}{2.921500in}}%
\pgfpathlineto{\pgfqpoint{5.974481in}{2.921500in}}%
\pgfpathlineto{\pgfqpoint{5.997955in}{2.921500in}}%
\pgfpathlineto{\pgfqpoint{6.021430in}{2.921500in}}%
\pgfpathlineto{\pgfqpoint{6.044904in}{2.921500in}}%
\pgfpathlineto{\pgfqpoint{6.068378in}{2.921500in}}%
\pgfpathlineto{\pgfqpoint{6.091853in}{2.921500in}}%
\pgfpathlineto{\pgfqpoint{6.115327in}{2.921500in}}%
\pgfpathlineto{\pgfqpoint{6.138801in}{2.921500in}}%
\pgfpathlineto{\pgfqpoint{6.162276in}{2.921500in}}%
\pgfpathlineto{\pgfqpoint{6.185750in}{2.921500in}}%
\pgfpathlineto{\pgfqpoint{6.209224in}{2.921500in}}%
\pgfpathlineto{\pgfqpoint{6.232699in}{2.921500in}}%
\pgfpathlineto{\pgfqpoint{6.256173in}{2.921500in}}%
\pgfpathlineto{\pgfqpoint{6.279648in}{2.921500in}}%
\pgfpathlineto{\pgfqpoint{6.303122in}{2.921500in}}%
\pgfpathlineto{\pgfqpoint{6.326596in}{2.921500in}}%
\pgfpathlineto{\pgfqpoint{6.350071in}{2.921500in}}%
\pgfpathlineto{\pgfqpoint{6.373545in}{2.921500in}}%
\pgfpathlineto{\pgfqpoint{6.397019in}{2.921500in}}%
\pgfpathlineto{\pgfqpoint{6.420494in}{2.921500in}}%
\pgfpathlineto{\pgfqpoint{6.443968in}{2.921500in}}%
\pgfpathlineto{\pgfqpoint{6.467442in}{2.921500in}}%
\pgfpathlineto{\pgfqpoint{6.490917in}{2.921500in}}%
\pgfpathlineto{\pgfqpoint{6.514391in}{2.921500in}}%
\pgfpathlineto{\pgfqpoint{6.537866in}{2.921500in}}%
\pgfpathlineto{\pgfqpoint{6.561340in}{2.921500in}}%
\pgfpathlineto{\pgfqpoint{6.584814in}{2.921500in}}%
\pgfpathlineto{\pgfqpoint{6.608289in}{2.921500in}}%
\pgfusepath{stroke}%
\end{pgfscope}%
\begin{pgfscope}%
\pgfpathmoveto{\pgfqpoint{5.728000in}{2.052407in}}%
\pgfpathcurveto{\pgfqpoint{6.017602in}{2.052407in}}{\pgfqpoint{6.295381in}{2.109937in}}{\pgfqpoint{6.500161in}{2.212326in}}%
\pgfpathcurveto{\pgfqpoint{6.704940in}{2.314716in}}{\pgfqpoint{6.820000in}{2.453606in}}{\pgfqpoint{6.820000in}{2.598407in}}%
\pgfpathcurveto{\pgfqpoint{6.820000in}{2.743208in}}{\pgfqpoint{6.704940in}{2.882097in}}{\pgfqpoint{6.500161in}{2.984487in}}%
\pgfpathcurveto{\pgfqpoint{6.295381in}{3.086877in}}{\pgfqpoint{6.017602in}{3.144407in}}{\pgfqpoint{5.728000in}{3.144407in}}%
\pgfpathcurveto{\pgfqpoint{5.438398in}{3.144407in}}{\pgfqpoint{5.160619in}{3.086877in}}{\pgfqpoint{4.955839in}{2.984487in}}%
\pgfpathcurveto{\pgfqpoint{4.751060in}{2.882097in}}{\pgfqpoint{4.636000in}{2.743208in}}{\pgfqpoint{4.636000in}{2.598407in}}%
\pgfpathcurveto{\pgfqpoint{4.636000in}{2.453606in}}{\pgfqpoint{4.751060in}{2.314716in}}{\pgfqpoint{4.955839in}{2.212326in}}%
\pgfpathcurveto{\pgfqpoint{5.160619in}{2.109937in}}{\pgfqpoint{5.438398in}{2.052407in}}{\pgfqpoint{5.728000in}{2.052407in}}%
\pgfpathlineto{\pgfqpoint{5.728000in}{2.052407in}}%
\pgfpathclose%
\pgfusepath{clip}%
\pgfsetrectcap%
\pgfsetroundjoin%
\pgfsetlinewidth{0.451687pt}%
\definecolor{currentstroke}{rgb}{0.690196,0.690196,0.690196}%
\pgfsetstrokecolor{currentstroke}%
\pgfsetstrokeopacity{0.250000}%
\pgfsetdash{}{0pt}%
\pgfpathmoveto{\pgfqpoint{5.021329in}{3.014664in}}%
\pgfpathlineto{\pgfqpoint{5.040174in}{3.014664in}}%
\pgfpathlineto{\pgfqpoint{5.059018in}{3.014664in}}%
\pgfpathlineto{\pgfqpoint{5.077863in}{3.014664in}}%
\pgfpathlineto{\pgfqpoint{5.096707in}{3.014664in}}%
\pgfpathlineto{\pgfqpoint{5.115552in}{3.014664in}}%
\pgfpathlineto{\pgfqpoint{5.134396in}{3.014664in}}%
\pgfpathlineto{\pgfqpoint{5.153241in}{3.014664in}}%
\pgfpathlineto{\pgfqpoint{5.172086in}{3.014664in}}%
\pgfpathlineto{\pgfqpoint{5.190930in}{3.014664in}}%
\pgfpathlineto{\pgfqpoint{5.209775in}{3.014664in}}%
\pgfpathlineto{\pgfqpoint{5.228619in}{3.014664in}}%
\pgfpathlineto{\pgfqpoint{5.247464in}{3.014664in}}%
\pgfpathlineto{\pgfqpoint{5.266308in}{3.014664in}}%
\pgfpathlineto{\pgfqpoint{5.285153in}{3.014664in}}%
\pgfpathlineto{\pgfqpoint{5.303997in}{3.014664in}}%
\pgfpathlineto{\pgfqpoint{5.322842in}{3.014664in}}%
\pgfpathlineto{\pgfqpoint{5.341687in}{3.014664in}}%
\pgfpathlineto{\pgfqpoint{5.360531in}{3.014664in}}%
\pgfpathlineto{\pgfqpoint{5.379376in}{3.014664in}}%
\pgfpathlineto{\pgfqpoint{5.398220in}{3.014664in}}%
\pgfpathlineto{\pgfqpoint{5.417065in}{3.014664in}}%
\pgfpathlineto{\pgfqpoint{5.435909in}{3.014664in}}%
\pgfpathlineto{\pgfqpoint{5.454754in}{3.014664in}}%
\pgfpathlineto{\pgfqpoint{5.473598in}{3.014664in}}%
\pgfpathlineto{\pgfqpoint{5.492443in}{3.014664in}}%
\pgfpathlineto{\pgfqpoint{5.511288in}{3.014664in}}%
\pgfpathlineto{\pgfqpoint{5.530132in}{3.014664in}}%
\pgfpathlineto{\pgfqpoint{5.548977in}{3.014664in}}%
\pgfpathlineto{\pgfqpoint{5.567821in}{3.014664in}}%
\pgfpathlineto{\pgfqpoint{5.586666in}{3.014664in}}%
\pgfpathlineto{\pgfqpoint{5.605510in}{3.014664in}}%
\pgfpathlineto{\pgfqpoint{5.624355in}{3.014664in}}%
\pgfpathlineto{\pgfqpoint{5.643199in}{3.014664in}}%
\pgfpathlineto{\pgfqpoint{5.662044in}{3.014664in}}%
\pgfpathlineto{\pgfqpoint{5.680889in}{3.014664in}}%
\pgfpathlineto{\pgfqpoint{5.699733in}{3.014664in}}%
\pgfpathlineto{\pgfqpoint{5.718578in}{3.014664in}}%
\pgfpathlineto{\pgfqpoint{5.737422in}{3.014664in}}%
\pgfpathlineto{\pgfqpoint{5.756267in}{3.014664in}}%
\pgfpathlineto{\pgfqpoint{5.775111in}{3.014664in}}%
\pgfpathlineto{\pgfqpoint{5.793956in}{3.014664in}}%
\pgfpathlineto{\pgfqpoint{5.812801in}{3.014664in}}%
\pgfpathlineto{\pgfqpoint{5.831645in}{3.014664in}}%
\pgfpathlineto{\pgfqpoint{5.850490in}{3.014664in}}%
\pgfpathlineto{\pgfqpoint{5.869334in}{3.014664in}}%
\pgfpathlineto{\pgfqpoint{5.888179in}{3.014664in}}%
\pgfpathlineto{\pgfqpoint{5.907023in}{3.014664in}}%
\pgfpathlineto{\pgfqpoint{5.925868in}{3.014664in}}%
\pgfpathlineto{\pgfqpoint{5.944712in}{3.014664in}}%
\pgfpathlineto{\pgfqpoint{5.963557in}{3.014664in}}%
\pgfpathlineto{\pgfqpoint{5.982402in}{3.014664in}}%
\pgfpathlineto{\pgfqpoint{6.001246in}{3.014664in}}%
\pgfpathlineto{\pgfqpoint{6.020091in}{3.014664in}}%
\pgfpathlineto{\pgfqpoint{6.038935in}{3.014664in}}%
\pgfpathlineto{\pgfqpoint{6.057780in}{3.014664in}}%
\pgfpathlineto{\pgfqpoint{6.076624in}{3.014664in}}%
\pgfpathlineto{\pgfqpoint{6.095469in}{3.014664in}}%
\pgfpathlineto{\pgfqpoint{6.114313in}{3.014664in}}%
\pgfpathlineto{\pgfqpoint{6.133158in}{3.014664in}}%
\pgfpathlineto{\pgfqpoint{6.152003in}{3.014664in}}%
\pgfpathlineto{\pgfqpoint{6.170847in}{3.014664in}}%
\pgfpathlineto{\pgfqpoint{6.189692in}{3.014664in}}%
\pgfpathlineto{\pgfqpoint{6.208536in}{3.014664in}}%
\pgfpathlineto{\pgfqpoint{6.227381in}{3.014664in}}%
\pgfpathlineto{\pgfqpoint{6.246225in}{3.014664in}}%
\pgfpathlineto{\pgfqpoint{6.265070in}{3.014664in}}%
\pgfpathlineto{\pgfqpoint{6.283914in}{3.014664in}}%
\pgfpathlineto{\pgfqpoint{6.302759in}{3.014664in}}%
\pgfpathlineto{\pgfqpoint{6.321604in}{3.014664in}}%
\pgfpathlineto{\pgfqpoint{6.340448in}{3.014664in}}%
\pgfpathlineto{\pgfqpoint{6.359293in}{3.014664in}}%
\pgfpathlineto{\pgfqpoint{6.378137in}{3.014664in}}%
\pgfpathlineto{\pgfqpoint{6.396982in}{3.014664in}}%
\pgfpathlineto{\pgfqpoint{6.415826in}{3.014664in}}%
\pgfpathlineto{\pgfqpoint{6.434671in}{3.014664in}}%
\pgfusepath{stroke}%
\end{pgfscope}%
\begin{pgfscope}%
\pgfpathmoveto{\pgfqpoint{5.728000in}{2.052407in}}%
\pgfpathcurveto{\pgfqpoint{6.017602in}{2.052407in}}{\pgfqpoint{6.295381in}{2.109937in}}{\pgfqpoint{6.500161in}{2.212326in}}%
\pgfpathcurveto{\pgfqpoint{6.704940in}{2.314716in}}{\pgfqpoint{6.820000in}{2.453606in}}{\pgfqpoint{6.820000in}{2.598407in}}%
\pgfpathcurveto{\pgfqpoint{6.820000in}{2.743208in}}{\pgfqpoint{6.704940in}{2.882097in}}{\pgfqpoint{6.500161in}{2.984487in}}%
\pgfpathcurveto{\pgfqpoint{6.295381in}{3.086877in}}{\pgfqpoint{6.017602in}{3.144407in}}{\pgfqpoint{5.728000in}{3.144407in}}%
\pgfpathcurveto{\pgfqpoint{5.438398in}{3.144407in}}{\pgfqpoint{5.160619in}{3.086877in}}{\pgfqpoint{4.955839in}{2.984487in}}%
\pgfpathcurveto{\pgfqpoint{4.751060in}{2.882097in}}{\pgfqpoint{4.636000in}{2.743208in}}{\pgfqpoint{4.636000in}{2.598407in}}%
\pgfpathcurveto{\pgfqpoint{4.636000in}{2.453606in}}{\pgfqpoint{4.751060in}{2.314716in}}{\pgfqpoint{4.955839in}{2.212326in}}%
\pgfpathcurveto{\pgfqpoint{5.160619in}{2.109937in}}{\pgfqpoint{5.438398in}{2.052407in}}{\pgfqpoint{5.728000in}{2.052407in}}%
\pgfpathlineto{\pgfqpoint{5.728000in}{2.052407in}}%
\pgfpathclose%
\pgfusepath{clip}%
\pgfsetrectcap%
\pgfsetroundjoin%
\pgfsetlinewidth{0.451687pt}%
\definecolor{currentstroke}{rgb}{0.690196,0.690196,0.690196}%
\pgfsetstrokecolor{currentstroke}%
\pgfsetstrokeopacity{0.250000}%
\pgfsetdash{}{0pt}%
\pgfpathmoveto{\pgfqpoint{5.265906in}{3.093112in}}%
\pgfpathlineto{\pgfqpoint{5.278228in}{3.093112in}}%
\pgfpathlineto{\pgfqpoint{5.290551in}{3.093112in}}%
\pgfpathlineto{\pgfqpoint{5.302873in}{3.093112in}}%
\pgfpathlineto{\pgfqpoint{5.315196in}{3.093112in}}%
\pgfpathlineto{\pgfqpoint{5.327518in}{3.093112in}}%
\pgfpathlineto{\pgfqpoint{5.339841in}{3.093112in}}%
\pgfpathlineto{\pgfqpoint{5.352163in}{3.093112in}}%
\pgfpathlineto{\pgfqpoint{5.364486in}{3.093112in}}%
\pgfpathlineto{\pgfqpoint{5.376808in}{3.093112in}}%
\pgfpathlineto{\pgfqpoint{5.389131in}{3.093112in}}%
\pgfpathlineto{\pgfqpoint{5.401453in}{3.093112in}}%
\pgfpathlineto{\pgfqpoint{5.413776in}{3.093112in}}%
\pgfpathlineto{\pgfqpoint{5.426098in}{3.093112in}}%
\pgfpathlineto{\pgfqpoint{5.438421in}{3.093112in}}%
\pgfpathlineto{\pgfqpoint{5.450743in}{3.093112in}}%
\pgfpathlineto{\pgfqpoint{5.463066in}{3.093112in}}%
\pgfpathlineto{\pgfqpoint{5.475388in}{3.093112in}}%
\pgfpathlineto{\pgfqpoint{5.487711in}{3.093112in}}%
\pgfpathlineto{\pgfqpoint{5.500033in}{3.093112in}}%
\pgfpathlineto{\pgfqpoint{5.512356in}{3.093112in}}%
\pgfpathlineto{\pgfqpoint{5.524678in}{3.093112in}}%
\pgfpathlineto{\pgfqpoint{5.537001in}{3.093112in}}%
\pgfpathlineto{\pgfqpoint{5.549323in}{3.093112in}}%
\pgfpathlineto{\pgfqpoint{5.561646in}{3.093112in}}%
\pgfpathlineto{\pgfqpoint{5.573969in}{3.093112in}}%
\pgfpathlineto{\pgfqpoint{5.586291in}{3.093112in}}%
\pgfpathlineto{\pgfqpoint{5.598614in}{3.093112in}}%
\pgfpathlineto{\pgfqpoint{5.610936in}{3.093112in}}%
\pgfpathlineto{\pgfqpoint{5.623259in}{3.093112in}}%
\pgfpathlineto{\pgfqpoint{5.635581in}{3.093112in}}%
\pgfpathlineto{\pgfqpoint{5.647904in}{3.093112in}}%
\pgfpathlineto{\pgfqpoint{5.660226in}{3.093112in}}%
\pgfpathlineto{\pgfqpoint{5.672549in}{3.093112in}}%
\pgfpathlineto{\pgfqpoint{5.684871in}{3.093112in}}%
\pgfpathlineto{\pgfqpoint{5.697194in}{3.093112in}}%
\pgfpathlineto{\pgfqpoint{5.709516in}{3.093112in}}%
\pgfpathlineto{\pgfqpoint{5.721839in}{3.093112in}}%
\pgfpathlineto{\pgfqpoint{5.734161in}{3.093112in}}%
\pgfpathlineto{\pgfqpoint{5.746484in}{3.093112in}}%
\pgfpathlineto{\pgfqpoint{5.758806in}{3.093112in}}%
\pgfpathlineto{\pgfqpoint{5.771129in}{3.093112in}}%
\pgfpathlineto{\pgfqpoint{5.783451in}{3.093112in}}%
\pgfpathlineto{\pgfqpoint{5.795774in}{3.093112in}}%
\pgfpathlineto{\pgfqpoint{5.808096in}{3.093112in}}%
\pgfpathlineto{\pgfqpoint{5.820419in}{3.093112in}}%
\pgfpathlineto{\pgfqpoint{5.832741in}{3.093112in}}%
\pgfpathlineto{\pgfqpoint{5.845064in}{3.093112in}}%
\pgfpathlineto{\pgfqpoint{5.857386in}{3.093112in}}%
\pgfpathlineto{\pgfqpoint{5.869709in}{3.093112in}}%
\pgfpathlineto{\pgfqpoint{5.882031in}{3.093112in}}%
\pgfpathlineto{\pgfqpoint{5.894354in}{3.093112in}}%
\pgfpathlineto{\pgfqpoint{5.906677in}{3.093112in}}%
\pgfpathlineto{\pgfqpoint{5.918999in}{3.093112in}}%
\pgfpathlineto{\pgfqpoint{5.931322in}{3.093112in}}%
\pgfpathlineto{\pgfqpoint{5.943644in}{3.093112in}}%
\pgfpathlineto{\pgfqpoint{5.955967in}{3.093112in}}%
\pgfpathlineto{\pgfqpoint{5.968289in}{3.093112in}}%
\pgfpathlineto{\pgfqpoint{5.980612in}{3.093112in}}%
\pgfpathlineto{\pgfqpoint{5.992934in}{3.093112in}}%
\pgfpathlineto{\pgfqpoint{6.005257in}{3.093112in}}%
\pgfpathlineto{\pgfqpoint{6.017579in}{3.093112in}}%
\pgfpathlineto{\pgfqpoint{6.029902in}{3.093112in}}%
\pgfpathlineto{\pgfqpoint{6.042224in}{3.093112in}}%
\pgfpathlineto{\pgfqpoint{6.054547in}{3.093112in}}%
\pgfpathlineto{\pgfqpoint{6.066869in}{3.093112in}}%
\pgfpathlineto{\pgfqpoint{6.079192in}{3.093112in}}%
\pgfpathlineto{\pgfqpoint{6.091514in}{3.093112in}}%
\pgfpathlineto{\pgfqpoint{6.103837in}{3.093112in}}%
\pgfpathlineto{\pgfqpoint{6.116159in}{3.093112in}}%
\pgfpathlineto{\pgfqpoint{6.128482in}{3.093112in}}%
\pgfpathlineto{\pgfqpoint{6.140804in}{3.093112in}}%
\pgfpathlineto{\pgfqpoint{6.153127in}{3.093112in}}%
\pgfpathlineto{\pgfqpoint{6.165449in}{3.093112in}}%
\pgfpathlineto{\pgfqpoint{6.177772in}{3.093112in}}%
\pgfpathlineto{\pgfqpoint{6.190094in}{3.093112in}}%
\pgfusepath{stroke}%
\end{pgfscope}%
\begin{pgfscope}%
\pgfsetrectcap%
\pgfsetmiterjoin%
\pgfsetlinewidth{0.803000pt}%
\definecolor{currentstroke}{rgb}{0.000000,0.000000,0.000000}%
\pgfsetstrokecolor{currentstroke}%
\pgfsetdash{}{0pt}%
\pgfpathmoveto{\pgfqpoint{5.728000in}{2.052407in}}%
\pgfpathcurveto{\pgfqpoint{6.017602in}{2.052407in}}{\pgfqpoint{6.295381in}{2.109937in}}{\pgfqpoint{6.500161in}{2.212326in}}%
\pgfpathcurveto{\pgfqpoint{6.704940in}{2.314716in}}{\pgfqpoint{6.820000in}{2.453606in}}{\pgfqpoint{6.820000in}{2.598407in}}%
\pgfpathcurveto{\pgfqpoint{6.820000in}{2.743208in}}{\pgfqpoint{6.704940in}{2.882097in}}{\pgfqpoint{6.500161in}{2.984487in}}%
\pgfpathcurveto{\pgfqpoint{6.295381in}{3.086877in}}{\pgfqpoint{6.017602in}{3.144407in}}{\pgfqpoint{5.728000in}{3.144407in}}%
\pgfpathcurveto{\pgfqpoint{5.438398in}{3.144407in}}{\pgfqpoint{5.160619in}{3.086877in}}{\pgfqpoint{4.955839in}{2.984487in}}%
\pgfpathcurveto{\pgfqpoint{4.751060in}{2.882097in}}{\pgfqpoint{4.636000in}{2.743208in}}{\pgfqpoint{4.636000in}{2.598407in}}%
\pgfpathcurveto{\pgfqpoint{4.636000in}{2.453606in}}{\pgfqpoint{4.751060in}{2.314716in}}{\pgfqpoint{4.955839in}{2.212326in}}%
\pgfpathcurveto{\pgfqpoint{5.160619in}{2.109937in}}{\pgfqpoint{5.438398in}{2.052407in}}{\pgfqpoint{5.728000in}{2.052407in}}%
\pgfpathlineto{\pgfqpoint{5.728000in}{2.052407in}}%
\pgfpathclose%
\pgfusepath{stroke}%
\end{pgfscope}%
\begin{pgfscope}%
\definecolor{textcolor}{rgb}{0.000000,0.000000,0.000000}%
\pgfsetstrokecolor{textcolor}%
\pgfsetfillcolor{textcolor}%
\pgftext[x=5.728000in,y=3.186073in,,base]{\color{textcolor}{\rmfamily\fontsize{8.200000}{9.840000}\selectfont\catcode`\^=\active\def^{\ifmmode\sp\else\^{}\fi}\catcode`\%=\active\def%{\%}Girdle}}%
\end{pgfscope}%
\begin{pgfscope}%
\pgfsetbuttcap%
\pgfsetmiterjoin%
\definecolor{currentfill}{rgb}{1.000000,1.000000,1.000000}%
\pgfsetfillcolor{currentfill}%
\pgfsetlinewidth{0.000000pt}%
\definecolor{currentstroke}{rgb}{0.000000,0.000000,0.000000}%
\pgfsetstrokecolor{currentstroke}%
\pgfsetstrokeopacity{0.000000}%
\pgfsetdash{}{0pt}%
\pgfpathmoveto{\pgfqpoint{2.326000in}{0.100000in}}%
\pgfpathcurveto{\pgfqpoint{2.615602in}{0.100000in}}{\pgfqpoint{2.893381in}{0.157530in}}{\pgfqpoint{3.098161in}{0.259920in}}%
\pgfpathcurveto{\pgfqpoint{3.302940in}{0.362309in}}{\pgfqpoint{3.418000in}{0.501199in}}{\pgfqpoint{3.418000in}{0.646000in}}%
\pgfpathcurveto{\pgfqpoint{3.418000in}{0.790801in}}{\pgfqpoint{3.302940in}{0.929691in}}{\pgfqpoint{3.098161in}{1.032080in}}%
\pgfpathcurveto{\pgfqpoint{2.893381in}{1.134470in}}{\pgfqpoint{2.615602in}{1.192000in}}{\pgfqpoint{2.326000in}{1.192000in}}%
\pgfpathcurveto{\pgfqpoint{2.036398in}{1.192000in}}{\pgfqpoint{1.758619in}{1.134470in}}{\pgfqpoint{1.553839in}{1.032080in}}%
\pgfpathcurveto{\pgfqpoint{1.349060in}{0.929691in}}{\pgfqpoint{1.234000in}{0.790801in}}{\pgfqpoint{1.234000in}{0.646000in}}%
\pgfpathcurveto{\pgfqpoint{1.234000in}{0.501199in}}{\pgfqpoint{1.349060in}{0.362309in}}{\pgfqpoint{1.553839in}{0.259920in}}%
\pgfpathcurveto{\pgfqpoint{1.758619in}{0.157530in}}{\pgfqpoint{2.036398in}{0.100000in}}{\pgfqpoint{2.326000in}{0.100000in}}%
\pgfpathlineto{\pgfqpoint{2.326000in}{0.100000in}}%
\pgfpathclose%
\pgfusepath{fill}%
\end{pgfscope}%
\begin{pgfscope}%
\pgfpathmoveto{\pgfqpoint{2.326000in}{0.100000in}}%
\pgfpathcurveto{\pgfqpoint{2.615602in}{0.100000in}}{\pgfqpoint{2.893381in}{0.157530in}}{\pgfqpoint{3.098161in}{0.259920in}}%
\pgfpathcurveto{\pgfqpoint{3.302940in}{0.362309in}}{\pgfqpoint{3.418000in}{0.501199in}}{\pgfqpoint{3.418000in}{0.646000in}}%
\pgfpathcurveto{\pgfqpoint{3.418000in}{0.790801in}}{\pgfqpoint{3.302940in}{0.929691in}}{\pgfqpoint{3.098161in}{1.032080in}}%
\pgfpathcurveto{\pgfqpoint{2.893381in}{1.134470in}}{\pgfqpoint{2.615602in}{1.192000in}}{\pgfqpoint{2.326000in}{1.192000in}}%
\pgfpathcurveto{\pgfqpoint{2.036398in}{1.192000in}}{\pgfqpoint{1.758619in}{1.134470in}}{\pgfqpoint{1.553839in}{1.032080in}}%
\pgfpathcurveto{\pgfqpoint{1.349060in}{0.929691in}}{\pgfqpoint{1.234000in}{0.790801in}}{\pgfqpoint{1.234000in}{0.646000in}}%
\pgfpathcurveto{\pgfqpoint{1.234000in}{0.501199in}}{\pgfqpoint{1.349060in}{0.362309in}}{\pgfqpoint{1.553839in}{0.259920in}}%
\pgfpathcurveto{\pgfqpoint{1.758619in}{0.157530in}}{\pgfqpoint{2.036398in}{0.100000in}}{\pgfqpoint{2.326000in}{0.100000in}}%
\pgfpathlineto{\pgfqpoint{2.326000in}{0.100000in}}%
\pgfpathclose%
\pgfusepath{clip}%
\pgfsetbuttcap%
\pgfsetroundjoin%
\definecolor{currentfill}{rgb}{0.121569,0.466667,0.705882}%
\pgfsetfillcolor{currentfill}%
\pgfsetfillopacity{0.580000}%
\pgfsetlinewidth{0.000000pt}%
\definecolor{currentstroke}{rgb}{0.121569,0.466667,0.705882}%
\pgfsetstrokecolor{currentstroke}%
\pgfsetstrokeopacity{0.580000}%
\pgfsetdash{}{0pt}%
\pgfsys@defobject{currentmarker}{\pgfqpoint{-0.014232in}{-0.014232in}}{\pgfqpoint{0.014232in}{0.014232in}}{%
\pgfpathmoveto{\pgfqpoint{0.000000in}{-0.014232in}}%
\pgfpathcurveto{\pgfqpoint{0.003774in}{-0.014232in}}{\pgfqpoint{0.007395in}{-0.012732in}}{\pgfqpoint{0.010063in}{-0.010063in}}%
\pgfpathcurveto{\pgfqpoint{0.012732in}{-0.007395in}}{\pgfqpoint{0.014232in}{-0.003774in}}{\pgfqpoint{0.014232in}{0.000000in}}%
\pgfpathcurveto{\pgfqpoint{0.014232in}{0.003774in}}{\pgfqpoint{0.012732in}{0.007395in}}{\pgfqpoint{0.010063in}{0.010063in}}%
\pgfpathcurveto{\pgfqpoint{0.007395in}{0.012732in}}{\pgfqpoint{0.003774in}{0.014232in}}{\pgfqpoint{0.000000in}{0.014232in}}%
\pgfpathcurveto{\pgfqpoint{-0.003774in}{0.014232in}}{\pgfqpoint{-0.007395in}{0.012732in}}{\pgfqpoint{-0.010063in}{0.010063in}}%
\pgfpathcurveto{\pgfqpoint{-0.012732in}{0.007395in}}{\pgfqpoint{-0.014232in}{0.003774in}}{\pgfqpoint{-0.014232in}{0.000000in}}%
\pgfpathcurveto{\pgfqpoint{-0.014232in}{-0.003774in}}{\pgfqpoint{-0.012732in}{-0.007395in}}{\pgfqpoint{-0.010063in}{-0.010063in}}%
\pgfpathcurveto{\pgfqpoint{-0.007395in}{-0.012732in}}{\pgfqpoint{-0.003774in}{-0.014232in}}{\pgfqpoint{0.000000in}{-0.014232in}}%
\pgfpathlineto{\pgfqpoint{0.000000in}{-0.014232in}}%
\pgfpathclose%
\pgfusepath{fill}%
}%
\begin{pgfscope}%
\pgfsys@transformshift{2.546097in}{0.868269in}%
\pgfsys@useobject{currentmarker}{}%
\end{pgfscope}%
\begin{pgfscope}%
\pgfsys@transformshift{2.610019in}{0.857412in}%
\pgfsys@useobject{currentmarker}{}%
\end{pgfscope}%
\begin{pgfscope}%
\pgfsys@transformshift{2.618499in}{0.886961in}%
\pgfsys@useobject{currentmarker}{}%
\end{pgfscope}%
\begin{pgfscope}%
\pgfsys@transformshift{2.575269in}{0.839951in}%
\pgfsys@useobject{currentmarker}{}%
\end{pgfscope}%
\begin{pgfscope}%
\pgfsys@transformshift{2.553116in}{0.912831in}%
\pgfsys@useobject{currentmarker}{}%
\end{pgfscope}%
\begin{pgfscope}%
\pgfsys@transformshift{2.580128in}{0.908473in}%
\pgfsys@useobject{currentmarker}{}%
\end{pgfscope}%
\begin{pgfscope}%
\pgfsys@transformshift{2.570016in}{0.947350in}%
\pgfsys@useobject{currentmarker}{}%
\end{pgfscope}%
\begin{pgfscope}%
\pgfsys@transformshift{2.593532in}{0.887417in}%
\pgfsys@useobject{currentmarker}{}%
\end{pgfscope}%
\begin{pgfscope}%
\pgfsys@transformshift{2.605733in}{0.898355in}%
\pgfsys@useobject{currentmarker}{}%
\end{pgfscope}%
\begin{pgfscope}%
\pgfsys@transformshift{2.519024in}{0.938499in}%
\pgfsys@useobject{currentmarker}{}%
\end{pgfscope}%
\begin{pgfscope}%
\pgfsys@transformshift{2.598310in}{0.928933in}%
\pgfsys@useobject{currentmarker}{}%
\end{pgfscope}%
\begin{pgfscope}%
\pgfsys@transformshift{2.565868in}{0.912203in}%
\pgfsys@useobject{currentmarker}{}%
\end{pgfscope}%
\begin{pgfscope}%
\pgfsys@transformshift{2.608232in}{0.915855in}%
\pgfsys@useobject{currentmarker}{}%
\end{pgfscope}%
\begin{pgfscope}%
\pgfsys@transformshift{2.571743in}{0.934802in}%
\pgfsys@useobject{currentmarker}{}%
\end{pgfscope}%
\begin{pgfscope}%
\pgfsys@transformshift{2.556651in}{0.901755in}%
\pgfsys@useobject{currentmarker}{}%
\end{pgfscope}%
\begin{pgfscope}%
\pgfsys@transformshift{2.612929in}{0.869033in}%
\pgfsys@useobject{currentmarker}{}%
\end{pgfscope}%
\begin{pgfscope}%
\pgfsys@transformshift{2.574947in}{0.920347in}%
\pgfsys@useobject{currentmarker}{}%
\end{pgfscope}%
\begin{pgfscope}%
\pgfsys@transformshift{2.520840in}{0.965768in}%
\pgfsys@useobject{currentmarker}{}%
\end{pgfscope}%
\begin{pgfscope}%
\pgfsys@transformshift{2.567394in}{0.877538in}%
\pgfsys@useobject{currentmarker}{}%
\end{pgfscope}%
\begin{pgfscope}%
\pgfsys@transformshift{2.552171in}{0.891667in}%
\pgfsys@useobject{currentmarker}{}%
\end{pgfscope}%
\begin{pgfscope}%
\pgfsys@transformshift{2.576540in}{0.917202in}%
\pgfsys@useobject{currentmarker}{}%
\end{pgfscope}%
\begin{pgfscope}%
\pgfsys@transformshift{2.573078in}{0.897927in}%
\pgfsys@useobject{currentmarker}{}%
\end{pgfscope}%
\begin{pgfscope}%
\pgfsys@transformshift{2.541832in}{0.921186in}%
\pgfsys@useobject{currentmarker}{}%
\end{pgfscope}%
\begin{pgfscope}%
\pgfsys@transformshift{2.559698in}{0.906028in}%
\pgfsys@useobject{currentmarker}{}%
\end{pgfscope}%
\begin{pgfscope}%
\pgfsys@transformshift{2.548874in}{0.955982in}%
\pgfsys@useobject{currentmarker}{}%
\end{pgfscope}%
\begin{pgfscope}%
\pgfsys@transformshift{2.534831in}{0.892432in}%
\pgfsys@useobject{currentmarker}{}%
\end{pgfscope}%
\begin{pgfscope}%
\pgfsys@transformshift{2.568165in}{0.844854in}%
\pgfsys@useobject{currentmarker}{}%
\end{pgfscope}%
\begin{pgfscope}%
\pgfsys@transformshift{2.591382in}{0.961351in}%
\pgfsys@useobject{currentmarker}{}%
\end{pgfscope}%
\begin{pgfscope}%
\pgfsys@transformshift{2.529818in}{0.934037in}%
\pgfsys@useobject{currentmarker}{}%
\end{pgfscope}%
\begin{pgfscope}%
\pgfsys@transformshift{2.544025in}{0.895757in}%
\pgfsys@useobject{currentmarker}{}%
\end{pgfscope}%
\begin{pgfscope}%
\pgfsys@transformshift{2.557761in}{0.926846in}%
\pgfsys@useobject{currentmarker}{}%
\end{pgfscope}%
\begin{pgfscope}%
\pgfsys@transformshift{2.549876in}{0.994845in}%
\pgfsys@useobject{currentmarker}{}%
\end{pgfscope}%
\begin{pgfscope}%
\pgfsys@transformshift{2.577439in}{0.904286in}%
\pgfsys@useobject{currentmarker}{}%
\end{pgfscope}%
\begin{pgfscope}%
\pgfsys@transformshift{2.569223in}{0.955012in}%
\pgfsys@useobject{currentmarker}{}%
\end{pgfscope}%
\begin{pgfscope}%
\pgfsys@transformshift{2.574871in}{0.949689in}%
\pgfsys@useobject{currentmarker}{}%
\end{pgfscope}%
\begin{pgfscope}%
\pgfsys@transformshift{2.549020in}{0.899279in}%
\pgfsys@useobject{currentmarker}{}%
\end{pgfscope}%
\begin{pgfscope}%
\pgfsys@transformshift{2.551159in}{0.907931in}%
\pgfsys@useobject{currentmarker}{}%
\end{pgfscope}%
\begin{pgfscope}%
\pgfsys@transformshift{2.544948in}{0.923012in}%
\pgfsys@useobject{currentmarker}{}%
\end{pgfscope}%
\begin{pgfscope}%
\pgfsys@transformshift{2.526662in}{0.886416in}%
\pgfsys@useobject{currentmarker}{}%
\end{pgfscope}%
\begin{pgfscope}%
\pgfsys@transformshift{2.562425in}{0.863946in}%
\pgfsys@useobject{currentmarker}{}%
\end{pgfscope}%
\begin{pgfscope}%
\pgfsys@transformshift{2.548234in}{0.915556in}%
\pgfsys@useobject{currentmarker}{}%
\end{pgfscope}%
\begin{pgfscope}%
\pgfsys@transformshift{2.599399in}{0.927073in}%
\pgfsys@useobject{currentmarker}{}%
\end{pgfscope}%
\begin{pgfscope}%
\pgfsys@transformshift{2.547462in}{0.918557in}%
\pgfsys@useobject{currentmarker}{}%
\end{pgfscope}%
\begin{pgfscope}%
\pgfsys@transformshift{2.589192in}{0.891854in}%
\pgfsys@useobject{currentmarker}{}%
\end{pgfscope}%
\begin{pgfscope}%
\pgfsys@transformshift{2.593930in}{0.924780in}%
\pgfsys@useobject{currentmarker}{}%
\end{pgfscope}%
\begin{pgfscope}%
\pgfsys@transformshift{2.551347in}{0.908117in}%
\pgfsys@useobject{currentmarker}{}%
\end{pgfscope}%
\begin{pgfscope}%
\pgfsys@transformshift{2.553571in}{0.863254in}%
\pgfsys@useobject{currentmarker}{}%
\end{pgfscope}%
\begin{pgfscope}%
\pgfsys@transformshift{2.590812in}{0.890175in}%
\pgfsys@useobject{currentmarker}{}%
\end{pgfscope}%
\begin{pgfscope}%
\pgfsys@transformshift{2.584133in}{0.923486in}%
\pgfsys@useobject{currentmarker}{}%
\end{pgfscope}%
\begin{pgfscope}%
\pgfsys@transformshift{2.518454in}{0.908563in}%
\pgfsys@useobject{currentmarker}{}%
\end{pgfscope}%
\begin{pgfscope}%
\pgfsys@transformshift{2.600560in}{0.879419in}%
\pgfsys@useobject{currentmarker}{}%
\end{pgfscope}%
\begin{pgfscope}%
\pgfsys@transformshift{2.533757in}{0.913239in}%
\pgfsys@useobject{currentmarker}{}%
\end{pgfscope}%
\begin{pgfscope}%
\pgfsys@transformshift{2.597490in}{0.875305in}%
\pgfsys@useobject{currentmarker}{}%
\end{pgfscope}%
\begin{pgfscope}%
\pgfsys@transformshift{2.633940in}{0.869077in}%
\pgfsys@useobject{currentmarker}{}%
\end{pgfscope}%
\begin{pgfscope}%
\pgfsys@transformshift{2.567484in}{0.907397in}%
\pgfsys@useobject{currentmarker}{}%
\end{pgfscope}%
\begin{pgfscope}%
\pgfsys@transformshift{2.562906in}{0.912012in}%
\pgfsys@useobject{currentmarker}{}%
\end{pgfscope}%
\begin{pgfscope}%
\pgfsys@transformshift{2.487927in}{0.894189in}%
\pgfsys@useobject{currentmarker}{}%
\end{pgfscope}%
\begin{pgfscope}%
\pgfsys@transformshift{2.556126in}{0.903011in}%
\pgfsys@useobject{currentmarker}{}%
\end{pgfscope}%
\begin{pgfscope}%
\pgfsys@transformshift{2.598425in}{0.875972in}%
\pgfsys@useobject{currentmarker}{}%
\end{pgfscope}%
\begin{pgfscope}%
\pgfsys@transformshift{2.506714in}{0.949975in}%
\pgfsys@useobject{currentmarker}{}%
\end{pgfscope}%
\begin{pgfscope}%
\pgfsys@transformshift{2.559595in}{0.924332in}%
\pgfsys@useobject{currentmarker}{}%
\end{pgfscope}%
\begin{pgfscope}%
\pgfsys@transformshift{2.607437in}{0.904799in}%
\pgfsys@useobject{currentmarker}{}%
\end{pgfscope}%
\begin{pgfscope}%
\pgfsys@transformshift{2.540997in}{0.918950in}%
\pgfsys@useobject{currentmarker}{}%
\end{pgfscope}%
\begin{pgfscope}%
\pgfsys@transformshift{2.548308in}{0.900159in}%
\pgfsys@useobject{currentmarker}{}%
\end{pgfscope}%
\begin{pgfscope}%
\pgfsys@transformshift{2.589090in}{0.925175in}%
\pgfsys@useobject{currentmarker}{}%
\end{pgfscope}%
\begin{pgfscope}%
\pgfsys@transformshift{2.499150in}{0.876923in}%
\pgfsys@useobject{currentmarker}{}%
\end{pgfscope}%
\begin{pgfscope}%
\pgfsys@transformshift{2.572423in}{0.882335in}%
\pgfsys@useobject{currentmarker}{}%
\end{pgfscope}%
\begin{pgfscope}%
\pgfsys@transformshift{2.580544in}{0.875438in}%
\pgfsys@useobject{currentmarker}{}%
\end{pgfscope}%
\begin{pgfscope}%
\pgfsys@transformshift{2.553603in}{0.912632in}%
\pgfsys@useobject{currentmarker}{}%
\end{pgfscope}%
\begin{pgfscope}%
\pgfsys@transformshift{2.618918in}{0.878280in}%
\pgfsys@useobject{currentmarker}{}%
\end{pgfscope}%
\begin{pgfscope}%
\pgfsys@transformshift{2.565944in}{0.925833in}%
\pgfsys@useobject{currentmarker}{}%
\end{pgfscope}%
\begin{pgfscope}%
\pgfsys@transformshift{2.575464in}{0.915370in}%
\pgfsys@useobject{currentmarker}{}%
\end{pgfscope}%
\begin{pgfscope}%
\pgfsys@transformshift{2.584017in}{0.854810in}%
\pgfsys@useobject{currentmarker}{}%
\end{pgfscope}%
\begin{pgfscope}%
\pgfsys@transformshift{2.526853in}{0.932904in}%
\pgfsys@useobject{currentmarker}{}%
\end{pgfscope}%
\begin{pgfscope}%
\pgfsys@transformshift{2.524176in}{0.938848in}%
\pgfsys@useobject{currentmarker}{}%
\end{pgfscope}%
\begin{pgfscope}%
\pgfsys@transformshift{2.112217in}{0.440960in}%
\pgfsys@useobject{currentmarker}{}%
\end{pgfscope}%
\begin{pgfscope}%
\pgfsys@transformshift{2.063115in}{0.347909in}%
\pgfsys@useobject{currentmarker}{}%
\end{pgfscope}%
\begin{pgfscope}%
\pgfsys@transformshift{2.106397in}{0.388636in}%
\pgfsys@useobject{currentmarker}{}%
\end{pgfscope}%
\begin{pgfscope}%
\pgfsys@transformshift{2.098222in}{0.379452in}%
\pgfsys@useobject{currentmarker}{}%
\end{pgfscope}%
\begin{pgfscope}%
\pgfsys@transformshift{2.100425in}{0.376007in}%
\pgfsys@useobject{currentmarker}{}%
\end{pgfscope}%
\begin{pgfscope}%
\pgfsys@transformshift{2.090059in}{0.425606in}%
\pgfsys@useobject{currentmarker}{}%
\end{pgfscope}%
\begin{pgfscope}%
\pgfsys@transformshift{2.044198in}{0.421374in}%
\pgfsys@useobject{currentmarker}{}%
\end{pgfscope}%
\begin{pgfscope}%
\pgfsys@transformshift{2.113695in}{0.411589in}%
\pgfsys@useobject{currentmarker}{}%
\end{pgfscope}%
\begin{pgfscope}%
\pgfsys@transformshift{2.126973in}{0.422720in}%
\pgfsys@useobject{currentmarker}{}%
\end{pgfscope}%
\begin{pgfscope}%
\pgfsys@transformshift{2.069398in}{0.428230in}%
\pgfsys@useobject{currentmarker}{}%
\end{pgfscope}%
\begin{pgfscope}%
\pgfsys@transformshift{2.070864in}{0.340912in}%
\pgfsys@useobject{currentmarker}{}%
\end{pgfscope}%
\begin{pgfscope}%
\pgfsys@transformshift{2.146964in}{0.358896in}%
\pgfsys@useobject{currentmarker}{}%
\end{pgfscope}%
\begin{pgfscope}%
\pgfsys@transformshift{2.087833in}{0.344113in}%
\pgfsys@useobject{currentmarker}{}%
\end{pgfscope}%
\begin{pgfscope}%
\pgfsys@transformshift{2.153603in}{0.333672in}%
\pgfsys@useobject{currentmarker}{}%
\end{pgfscope}%
\begin{pgfscope}%
\pgfsys@transformshift{2.045376in}{0.406429in}%
\pgfsys@useobject{currentmarker}{}%
\end{pgfscope}%
\begin{pgfscope}%
\pgfsys@transformshift{2.080609in}{0.417246in}%
\pgfsys@useobject{currentmarker}{}%
\end{pgfscope}%
\begin{pgfscope}%
\pgfsys@transformshift{2.102577in}{0.408811in}%
\pgfsys@useobject{currentmarker}{}%
\end{pgfscope}%
\begin{pgfscope}%
\pgfsys@transformshift{2.091155in}{0.391274in}%
\pgfsys@useobject{currentmarker}{}%
\end{pgfscope}%
\begin{pgfscope}%
\pgfsys@transformshift{2.086135in}{0.404313in}%
\pgfsys@useobject{currentmarker}{}%
\end{pgfscope}%
\begin{pgfscope}%
\pgfsys@transformshift{2.089770in}{0.357254in}%
\pgfsys@useobject{currentmarker}{}%
\end{pgfscope}%
\begin{pgfscope}%
\pgfsys@transformshift{2.098889in}{0.377920in}%
\pgfsys@useobject{currentmarker}{}%
\end{pgfscope}%
\begin{pgfscope}%
\pgfsys@transformshift{2.073533in}{0.353325in}%
\pgfsys@useobject{currentmarker}{}%
\end{pgfscope}%
\begin{pgfscope}%
\pgfsys@transformshift{2.066142in}{0.383071in}%
\pgfsys@useobject{currentmarker}{}%
\end{pgfscope}%
\begin{pgfscope}%
\pgfsys@transformshift{2.134647in}{0.357761in}%
\pgfsys@useobject{currentmarker}{}%
\end{pgfscope}%
\begin{pgfscope}%
\pgfsys@transformshift{2.122563in}{0.376816in}%
\pgfsys@useobject{currentmarker}{}%
\end{pgfscope}%
\begin{pgfscope}%
\pgfsys@transformshift{2.081466in}{0.384274in}%
\pgfsys@useobject{currentmarker}{}%
\end{pgfscope}%
\begin{pgfscope}%
\pgfsys@transformshift{2.056858in}{0.379425in}%
\pgfsys@useobject{currentmarker}{}%
\end{pgfscope}%
\begin{pgfscope}%
\pgfsys@transformshift{2.089902in}{0.425385in}%
\pgfsys@useobject{currentmarker}{}%
\end{pgfscope}%
\begin{pgfscope}%
\pgfsys@transformshift{2.089818in}{0.365646in}%
\pgfsys@useobject{currentmarker}{}%
\end{pgfscope}%
\begin{pgfscope}%
\pgfsys@transformshift{2.119240in}{0.416506in}%
\pgfsys@useobject{currentmarker}{}%
\end{pgfscope}%
\begin{pgfscope}%
\pgfsys@transformshift{2.074183in}{0.408043in}%
\pgfsys@useobject{currentmarker}{}%
\end{pgfscope}%
\begin{pgfscope}%
\pgfsys@transformshift{2.174639in}{0.390643in}%
\pgfsys@useobject{currentmarker}{}%
\end{pgfscope}%
\begin{pgfscope}%
\pgfsys@transformshift{2.032875in}{0.384074in}%
\pgfsys@useobject{currentmarker}{}%
\end{pgfscope}%
\begin{pgfscope}%
\pgfsys@transformshift{2.146337in}{0.377255in}%
\pgfsys@useobject{currentmarker}{}%
\end{pgfscope}%
\begin{pgfscope}%
\pgfsys@transformshift{2.128602in}{0.409542in}%
\pgfsys@useobject{currentmarker}{}%
\end{pgfscope}%
\begin{pgfscope}%
\pgfsys@transformshift{2.080872in}{0.379920in}%
\pgfsys@useobject{currentmarker}{}%
\end{pgfscope}%
\begin{pgfscope}%
\pgfsys@transformshift{2.010667in}{0.457951in}%
\pgfsys@useobject{currentmarker}{}%
\end{pgfscope}%
\begin{pgfscope}%
\pgfsys@transformshift{2.039449in}{0.378986in}%
\pgfsys@useobject{currentmarker}{}%
\end{pgfscope}%
\begin{pgfscope}%
\pgfsys@transformshift{2.065010in}{0.367102in}%
\pgfsys@useobject{currentmarker}{}%
\end{pgfscope}%
\begin{pgfscope}%
\pgfsys@transformshift{2.057806in}{0.392506in}%
\pgfsys@useobject{currentmarker}{}%
\end{pgfscope}%
\begin{pgfscope}%
\pgfsys@transformshift{2.102630in}{0.399599in}%
\pgfsys@useobject{currentmarker}{}%
\end{pgfscope}%
\begin{pgfscope}%
\pgfsys@transformshift{2.050637in}{0.369622in}%
\pgfsys@useobject{currentmarker}{}%
\end{pgfscope}%
\begin{pgfscope}%
\pgfsys@transformshift{2.140597in}{0.388277in}%
\pgfsys@useobject{currentmarker}{}%
\end{pgfscope}%
\begin{pgfscope}%
\pgfsys@transformshift{2.092253in}{0.403711in}%
\pgfsys@useobject{currentmarker}{}%
\end{pgfscope}%
\begin{pgfscope}%
\pgfsys@transformshift{2.042997in}{0.452175in}%
\pgfsys@useobject{currentmarker}{}%
\end{pgfscope}%
\begin{pgfscope}%
\pgfsys@transformshift{2.159538in}{0.388860in}%
\pgfsys@useobject{currentmarker}{}%
\end{pgfscope}%
\begin{pgfscope}%
\pgfsys@transformshift{2.140473in}{0.384663in}%
\pgfsys@useobject{currentmarker}{}%
\end{pgfscope}%
\begin{pgfscope}%
\pgfsys@transformshift{2.049324in}{0.383164in}%
\pgfsys@useobject{currentmarker}{}%
\end{pgfscope}%
\begin{pgfscope}%
\pgfsys@transformshift{2.079174in}{0.380883in}%
\pgfsys@useobject{currentmarker}{}%
\end{pgfscope}%
\begin{pgfscope}%
\pgfsys@transformshift{2.079202in}{0.415454in}%
\pgfsys@useobject{currentmarker}{}%
\end{pgfscope}%
\begin{pgfscope}%
\pgfsys@transformshift{2.073331in}{0.424723in}%
\pgfsys@useobject{currentmarker}{}%
\end{pgfscope}%
\begin{pgfscope}%
\pgfsys@transformshift{2.111222in}{0.408612in}%
\pgfsys@useobject{currentmarker}{}%
\end{pgfscope}%
\begin{pgfscope}%
\pgfsys@transformshift{2.059296in}{0.406460in}%
\pgfsys@useobject{currentmarker}{}%
\end{pgfscope}%
\begin{pgfscope}%
\pgfsys@transformshift{2.053365in}{0.398408in}%
\pgfsys@useobject{currentmarker}{}%
\end{pgfscope}%
\begin{pgfscope}%
\pgfsys@transformshift{2.116176in}{0.343459in}%
\pgfsys@useobject{currentmarker}{}%
\end{pgfscope}%
\begin{pgfscope}%
\pgfsys@transformshift{2.156065in}{0.355872in}%
\pgfsys@useobject{currentmarker}{}%
\end{pgfscope}%
\begin{pgfscope}%
\pgfsys@transformshift{2.115408in}{0.331257in}%
\pgfsys@useobject{currentmarker}{}%
\end{pgfscope}%
\begin{pgfscope}%
\pgfsys@transformshift{2.104990in}{0.371346in}%
\pgfsys@useobject{currentmarker}{}%
\end{pgfscope}%
\begin{pgfscope}%
\pgfsys@transformshift{2.082720in}{0.420382in}%
\pgfsys@useobject{currentmarker}{}%
\end{pgfscope}%
\begin{pgfscope}%
\pgfsys@transformshift{2.071462in}{0.375559in}%
\pgfsys@useobject{currentmarker}{}%
\end{pgfscope}%
\begin{pgfscope}%
\pgfsys@transformshift{2.083910in}{0.312114in}%
\pgfsys@useobject{currentmarker}{}%
\end{pgfscope}%
\begin{pgfscope}%
\pgfsys@transformshift{2.071538in}{0.386258in}%
\pgfsys@useobject{currentmarker}{}%
\end{pgfscope}%
\begin{pgfscope}%
\pgfsys@transformshift{2.097490in}{0.429422in}%
\pgfsys@useobject{currentmarker}{}%
\end{pgfscope}%
\begin{pgfscope}%
\pgfsys@transformshift{2.072740in}{0.388882in}%
\pgfsys@useobject{currentmarker}{}%
\end{pgfscope}%
\begin{pgfscope}%
\pgfsys@transformshift{2.095471in}{0.354783in}%
\pgfsys@useobject{currentmarker}{}%
\end{pgfscope}%
\begin{pgfscope}%
\pgfsys@transformshift{2.120764in}{0.385619in}%
\pgfsys@useobject{currentmarker}{}%
\end{pgfscope}%
\begin{pgfscope}%
\pgfsys@transformshift{2.068702in}{0.405589in}%
\pgfsys@useobject{currentmarker}{}%
\end{pgfscope}%
\begin{pgfscope}%
\pgfsys@transformshift{2.115432in}{0.369835in}%
\pgfsys@useobject{currentmarker}{}%
\end{pgfscope}%
\begin{pgfscope}%
\pgfsys@transformshift{2.088160in}{0.359665in}%
\pgfsys@useobject{currentmarker}{}%
\end{pgfscope}%
\begin{pgfscope}%
\pgfsys@transformshift{2.064453in}{0.337960in}%
\pgfsys@useobject{currentmarker}{}%
\end{pgfscope}%
\begin{pgfscope}%
\pgfsys@transformshift{2.117763in}{0.373435in}%
\pgfsys@useobject{currentmarker}{}%
\end{pgfscope}%
\begin{pgfscope}%
\pgfsys@transformshift{2.060535in}{0.356967in}%
\pgfsys@useobject{currentmarker}{}%
\end{pgfscope}%
\begin{pgfscope}%
\pgfsys@transformshift{2.025253in}{0.362543in}%
\pgfsys@useobject{currentmarker}{}%
\end{pgfscope}%
\begin{pgfscope}%
\pgfsys@transformshift{2.147789in}{0.393111in}%
\pgfsys@useobject{currentmarker}{}%
\end{pgfscope}%
\begin{pgfscope}%
\pgfsys@transformshift{2.048323in}{0.413519in}%
\pgfsys@useobject{currentmarker}{}%
\end{pgfscope}%
\begin{pgfscope}%
\pgfsys@transformshift{3.064148in}{0.417174in}%
\pgfsys@useobject{currentmarker}{}%
\end{pgfscope}%
\begin{pgfscope}%
\pgfsys@transformshift{3.001234in}{0.339997in}%
\pgfsys@useobject{currentmarker}{}%
\end{pgfscope}%
\begin{pgfscope}%
\pgfsys@transformshift{3.061286in}{0.425441in}%
\pgfsys@useobject{currentmarker}{}%
\end{pgfscope}%
\begin{pgfscope}%
\pgfsys@transformshift{3.039628in}{0.382054in}%
\pgfsys@useobject{currentmarker}{}%
\end{pgfscope}%
\begin{pgfscope}%
\pgfsys@transformshift{3.058619in}{0.418281in}%
\pgfsys@useobject{currentmarker}{}%
\end{pgfscope}%
\begin{pgfscope}%
\pgfsys@transformshift{3.066352in}{0.427012in}%
\pgfsys@useobject{currentmarker}{}%
\end{pgfscope}%
\begin{pgfscope}%
\pgfsys@transformshift{2.967222in}{0.379378in}%
\pgfsys@useobject{currentmarker}{}%
\end{pgfscope}%
\begin{pgfscope}%
\pgfsys@transformshift{3.100896in}{0.458739in}%
\pgfsys@useobject{currentmarker}{}%
\end{pgfscope}%
\begin{pgfscope}%
\pgfsys@transformshift{3.071331in}{0.397876in}%
\pgfsys@useobject{currentmarker}{}%
\end{pgfscope}%
\begin{pgfscope}%
\pgfsys@transformshift{3.079504in}{0.401156in}%
\pgfsys@useobject{currentmarker}{}%
\end{pgfscope}%
\begin{pgfscope}%
\pgfsys@transformshift{3.067412in}{0.408359in}%
\pgfsys@useobject{currentmarker}{}%
\end{pgfscope}%
\begin{pgfscope}%
\pgfsys@transformshift{3.062135in}{0.371001in}%
\pgfsys@useobject{currentmarker}{}%
\end{pgfscope}%
\begin{pgfscope}%
\pgfsys@transformshift{3.053677in}{0.365564in}%
\pgfsys@useobject{currentmarker}{}%
\end{pgfscope}%
\begin{pgfscope}%
\pgfsys@transformshift{3.013686in}{0.351367in}%
\pgfsys@useobject{currentmarker}{}%
\end{pgfscope}%
\begin{pgfscope}%
\pgfsys@transformshift{3.069623in}{0.415048in}%
\pgfsys@useobject{currentmarker}{}%
\end{pgfscope}%
\begin{pgfscope}%
\pgfsys@transformshift{3.076174in}{0.438659in}%
\pgfsys@useobject{currentmarker}{}%
\end{pgfscope}%
\begin{pgfscope}%
\pgfsys@transformshift{3.005362in}{0.381153in}%
\pgfsys@useobject{currentmarker}{}%
\end{pgfscope}%
\begin{pgfscope}%
\pgfsys@transformshift{3.096129in}{0.422343in}%
\pgfsys@useobject{currentmarker}{}%
\end{pgfscope}%
\begin{pgfscope}%
\pgfsys@transformshift{3.010461in}{0.364358in}%
\pgfsys@useobject{currentmarker}{}%
\end{pgfscope}%
\begin{pgfscope}%
\pgfsys@transformshift{2.980331in}{0.370285in}%
\pgfsys@useobject{currentmarker}{}%
\end{pgfscope}%
\begin{pgfscope}%
\pgfsys@transformshift{2.939837in}{0.355684in}%
\pgfsys@useobject{currentmarker}{}%
\end{pgfscope}%
\begin{pgfscope}%
\pgfsys@transformshift{3.099818in}{0.369757in}%
\pgfsys@useobject{currentmarker}{}%
\end{pgfscope}%
\begin{pgfscope}%
\pgfsys@transformshift{2.969234in}{0.339365in}%
\pgfsys@useobject{currentmarker}{}%
\end{pgfscope}%
\begin{pgfscope}%
\pgfsys@transformshift{3.024937in}{0.376278in}%
\pgfsys@useobject{currentmarker}{}%
\end{pgfscope}%
\begin{pgfscope}%
\pgfsys@transformshift{3.026891in}{0.393740in}%
\pgfsys@useobject{currentmarker}{}%
\end{pgfscope}%
\begin{pgfscope}%
\pgfsys@transformshift{3.020284in}{0.364645in}%
\pgfsys@useobject{currentmarker}{}%
\end{pgfscope}%
\begin{pgfscope}%
\pgfsys@transformshift{3.097603in}{0.423900in}%
\pgfsys@useobject{currentmarker}{}%
\end{pgfscope}%
\begin{pgfscope}%
\pgfsys@transformshift{3.028979in}{0.435110in}%
\pgfsys@useobject{currentmarker}{}%
\end{pgfscope}%
\begin{pgfscope}%
\pgfsys@transformshift{3.046221in}{0.357227in}%
\pgfsys@useobject{currentmarker}{}%
\end{pgfscope}%
\begin{pgfscope}%
\pgfsys@transformshift{3.063528in}{0.386319in}%
\pgfsys@useobject{currentmarker}{}%
\end{pgfscope}%
\begin{pgfscope}%
\pgfsys@transformshift{3.030275in}{0.407467in}%
\pgfsys@useobject{currentmarker}{}%
\end{pgfscope}%
\begin{pgfscope}%
\pgfsys@transformshift{3.009470in}{0.361012in}%
\pgfsys@useobject{currentmarker}{}%
\end{pgfscope}%
\begin{pgfscope}%
\pgfsys@transformshift{3.073152in}{0.435429in}%
\pgfsys@useobject{currentmarker}{}%
\end{pgfscope}%
\begin{pgfscope}%
\pgfsys@transformshift{3.042808in}{0.438716in}%
\pgfsys@useobject{currentmarker}{}%
\end{pgfscope}%
\begin{pgfscope}%
\pgfsys@transformshift{3.076837in}{0.420951in}%
\pgfsys@useobject{currentmarker}{}%
\end{pgfscope}%
\begin{pgfscope}%
\pgfsys@transformshift{3.003248in}{0.378787in}%
\pgfsys@useobject{currentmarker}{}%
\end{pgfscope}%
\begin{pgfscope}%
\pgfsys@transformshift{3.105425in}{0.427336in}%
\pgfsys@useobject{currentmarker}{}%
\end{pgfscope}%
\begin{pgfscope}%
\pgfsys@transformshift{3.030179in}{0.341057in}%
\pgfsys@useobject{currentmarker}{}%
\end{pgfscope}%
\begin{pgfscope}%
\pgfsys@transformshift{3.016164in}{0.381182in}%
\pgfsys@useobject{currentmarker}{}%
\end{pgfscope}%
\begin{pgfscope}%
\pgfsys@transformshift{3.044860in}{0.350144in}%
\pgfsys@useobject{currentmarker}{}%
\end{pgfscope}%
\begin{pgfscope}%
\pgfsys@transformshift{3.063687in}{0.422666in}%
\pgfsys@useobject{currentmarker}{}%
\end{pgfscope}%
\begin{pgfscope}%
\pgfsys@transformshift{3.086539in}{0.444608in}%
\pgfsys@useobject{currentmarker}{}%
\end{pgfscope}%
\begin{pgfscope}%
\pgfsys@transformshift{3.029546in}{0.363689in}%
\pgfsys@useobject{currentmarker}{}%
\end{pgfscope}%
\begin{pgfscope}%
\pgfsys@transformshift{3.022431in}{0.408030in}%
\pgfsys@useobject{currentmarker}{}%
\end{pgfscope}%
\begin{pgfscope}%
\pgfsys@transformshift{3.071742in}{0.429318in}%
\pgfsys@useobject{currentmarker}{}%
\end{pgfscope}%
\begin{pgfscope}%
\pgfsys@transformshift{3.048234in}{0.398601in}%
\pgfsys@useobject{currentmarker}{}%
\end{pgfscope}%
\begin{pgfscope}%
\pgfsys@transformshift{3.053464in}{0.371206in}%
\pgfsys@useobject{currentmarker}{}%
\end{pgfscope}%
\begin{pgfscope}%
\pgfsys@transformshift{3.015516in}{0.373595in}%
\pgfsys@useobject{currentmarker}{}%
\end{pgfscope}%
\begin{pgfscope}%
\pgfsys@transformshift{3.070891in}{0.366416in}%
\pgfsys@useobject{currentmarker}{}%
\end{pgfscope}%
\begin{pgfscope}%
\pgfsys@transformshift{3.071974in}{0.397423in}%
\pgfsys@useobject{currentmarker}{}%
\end{pgfscope}%
\begin{pgfscope}%
\pgfsys@transformshift{3.045744in}{0.353654in}%
\pgfsys@useobject{currentmarker}{}%
\end{pgfscope}%
\begin{pgfscope}%
\pgfsys@transformshift{3.043429in}{0.383473in}%
\pgfsys@useobject{currentmarker}{}%
\end{pgfscope}%
\begin{pgfscope}%
\pgfsys@transformshift{3.028597in}{0.377752in}%
\pgfsys@useobject{currentmarker}{}%
\end{pgfscope}%
\begin{pgfscope}%
\pgfsys@transformshift{3.011580in}{0.362012in}%
\pgfsys@useobject{currentmarker}{}%
\end{pgfscope}%
\begin{pgfscope}%
\pgfsys@transformshift{3.038688in}{0.398142in}%
\pgfsys@useobject{currentmarker}{}%
\end{pgfscope}%
\begin{pgfscope}%
\pgfsys@transformshift{3.001474in}{0.362436in}%
\pgfsys@useobject{currentmarker}{}%
\end{pgfscope}%
\begin{pgfscope}%
\pgfsys@transformshift{3.054486in}{0.412174in}%
\pgfsys@useobject{currentmarker}{}%
\end{pgfscope}%
\begin{pgfscope}%
\pgfsys@transformshift{3.036064in}{0.437210in}%
\pgfsys@useobject{currentmarker}{}%
\end{pgfscope}%
\begin{pgfscope}%
\pgfsys@transformshift{3.054523in}{0.391623in}%
\pgfsys@useobject{currentmarker}{}%
\end{pgfscope}%
\begin{pgfscope}%
\pgfsys@transformshift{3.073858in}{0.407985in}%
\pgfsys@useobject{currentmarker}{}%
\end{pgfscope}%
\begin{pgfscope}%
\pgfsys@transformshift{3.000964in}{0.325230in}%
\pgfsys@useobject{currentmarker}{}%
\end{pgfscope}%
\begin{pgfscope}%
\pgfsys@transformshift{3.094863in}{0.380897in}%
\pgfsys@useobject{currentmarker}{}%
\end{pgfscope}%
\begin{pgfscope}%
\pgfsys@transformshift{3.004317in}{0.369491in}%
\pgfsys@useobject{currentmarker}{}%
\end{pgfscope}%
\begin{pgfscope}%
\pgfsys@transformshift{3.076248in}{0.395661in}%
\pgfsys@useobject{currentmarker}{}%
\end{pgfscope}%
\begin{pgfscope}%
\pgfsys@transformshift{3.126992in}{0.424150in}%
\pgfsys@useobject{currentmarker}{}%
\end{pgfscope}%
\begin{pgfscope}%
\pgfsys@transformshift{3.059976in}{0.394653in}%
\pgfsys@useobject{currentmarker}{}%
\end{pgfscope}%
\begin{pgfscope}%
\pgfsys@transformshift{3.049574in}{0.353761in}%
\pgfsys@useobject{currentmarker}{}%
\end{pgfscope}%
\begin{pgfscope}%
\pgfsys@transformshift{3.048668in}{0.391340in}%
\pgfsys@useobject{currentmarker}{}%
\end{pgfscope}%
\begin{pgfscope}%
\pgfsys@transformshift{3.068683in}{0.414016in}%
\pgfsys@useobject{currentmarker}{}%
\end{pgfscope}%
\begin{pgfscope}%
\pgfsys@transformshift{3.075019in}{0.376878in}%
\pgfsys@useobject{currentmarker}{}%
\end{pgfscope}%
\begin{pgfscope}%
\pgfsys@transformshift{3.104360in}{0.407214in}%
\pgfsys@useobject{currentmarker}{}%
\end{pgfscope}%
\begin{pgfscope}%
\pgfsys@transformshift{2.954911in}{0.333051in}%
\pgfsys@useobject{currentmarker}{}%
\end{pgfscope}%
\begin{pgfscope}%
\pgfsys@transformshift{3.066833in}{0.457959in}%
\pgfsys@useobject{currentmarker}{}%
\end{pgfscope}%
\begin{pgfscope}%
\pgfsys@transformshift{3.082399in}{0.393538in}%
\pgfsys@useobject{currentmarker}{}%
\end{pgfscope}%
\begin{pgfscope}%
\pgfsys@transformshift{3.063602in}{0.412597in}%
\pgfsys@useobject{currentmarker}{}%
\end{pgfscope}%
\begin{pgfscope}%
\pgfsys@transformshift{1.544576in}{0.857836in}%
\pgfsys@useobject{currentmarker}{}%
\end{pgfscope}%
\begin{pgfscope}%
\pgfsys@transformshift{1.606814in}{0.909450in}%
\pgfsys@useobject{currentmarker}{}%
\end{pgfscope}%
\begin{pgfscope}%
\pgfsys@transformshift{1.694175in}{0.944106in}%
\pgfsys@useobject{currentmarker}{}%
\end{pgfscope}%
\begin{pgfscope}%
\pgfsys@transformshift{1.649405in}{0.919616in}%
\pgfsys@useobject{currentmarker}{}%
\end{pgfscope}%
\begin{pgfscope}%
\pgfsys@transformshift{1.683265in}{0.921777in}%
\pgfsys@useobject{currentmarker}{}%
\end{pgfscope}%
\begin{pgfscope}%
\pgfsys@transformshift{1.642117in}{0.882122in}%
\pgfsys@useobject{currentmarker}{}%
\end{pgfscope}%
\begin{pgfscope}%
\pgfsys@transformshift{1.575219in}{0.882666in}%
\pgfsys@useobject{currentmarker}{}%
\end{pgfscope}%
\begin{pgfscope}%
\pgfsys@transformshift{1.562122in}{0.933635in}%
\pgfsys@useobject{currentmarker}{}%
\end{pgfscope}%
\begin{pgfscope}%
\pgfsys@transformshift{1.577685in}{0.856878in}%
\pgfsys@useobject{currentmarker}{}%
\end{pgfscope}%
\begin{pgfscope}%
\pgfsys@transformshift{1.648302in}{0.962127in}%
\pgfsys@useobject{currentmarker}{}%
\end{pgfscope}%
\begin{pgfscope}%
\pgfsys@transformshift{1.588166in}{0.884248in}%
\pgfsys@useobject{currentmarker}{}%
\end{pgfscope}%
\begin{pgfscope}%
\pgfsys@transformshift{1.694619in}{0.970711in}%
\pgfsys@useobject{currentmarker}{}%
\end{pgfscope}%
\begin{pgfscope}%
\pgfsys@transformshift{1.666671in}{0.939795in}%
\pgfsys@useobject{currentmarker}{}%
\end{pgfscope}%
\begin{pgfscope}%
\pgfsys@transformshift{1.600915in}{0.887533in}%
\pgfsys@useobject{currentmarker}{}%
\end{pgfscope}%
\begin{pgfscope}%
\pgfsys@transformshift{1.580864in}{0.908558in}%
\pgfsys@useobject{currentmarker}{}%
\end{pgfscope}%
\begin{pgfscope}%
\pgfsys@transformshift{1.595967in}{0.877818in}%
\pgfsys@useobject{currentmarker}{}%
\end{pgfscope}%
\begin{pgfscope}%
\pgfsys@transformshift{1.605230in}{0.900922in}%
\pgfsys@useobject{currentmarker}{}%
\end{pgfscope}%
\begin{pgfscope}%
\pgfsys@transformshift{1.664851in}{0.873852in}%
\pgfsys@useobject{currentmarker}{}%
\end{pgfscope}%
\begin{pgfscope}%
\pgfsys@transformshift{1.579440in}{0.872887in}%
\pgfsys@useobject{currentmarker}{}%
\end{pgfscope}%
\begin{pgfscope}%
\pgfsys@transformshift{1.623598in}{0.898326in}%
\pgfsys@useobject{currentmarker}{}%
\end{pgfscope}%
\begin{pgfscope}%
\pgfsys@transformshift{1.589806in}{0.878244in}%
\pgfsys@useobject{currentmarker}{}%
\end{pgfscope}%
\begin{pgfscope}%
\pgfsys@transformshift{1.635301in}{0.911680in}%
\pgfsys@useobject{currentmarker}{}%
\end{pgfscope}%
\begin{pgfscope}%
\pgfsys@transformshift{1.680660in}{0.908664in}%
\pgfsys@useobject{currentmarker}{}%
\end{pgfscope}%
\begin{pgfscope}%
\pgfsys@transformshift{1.698311in}{0.965519in}%
\pgfsys@useobject{currentmarker}{}%
\end{pgfscope}%
\begin{pgfscope}%
\pgfsys@transformshift{1.597948in}{0.918075in}%
\pgfsys@useobject{currentmarker}{}%
\end{pgfscope}%
\begin{pgfscope}%
\pgfsys@transformshift{1.569409in}{0.899586in}%
\pgfsys@useobject{currentmarker}{}%
\end{pgfscope}%
\begin{pgfscope}%
\pgfsys@transformshift{1.597314in}{0.939659in}%
\pgfsys@useobject{currentmarker}{}%
\end{pgfscope}%
\begin{pgfscope}%
\pgfsys@transformshift{1.663936in}{0.996084in}%
\pgfsys@useobject{currentmarker}{}%
\end{pgfscope}%
\begin{pgfscope}%
\pgfsys@transformshift{1.638033in}{0.901827in}%
\pgfsys@useobject{currentmarker}{}%
\end{pgfscope}%
\begin{pgfscope}%
\pgfsys@transformshift{1.698442in}{0.959728in}%
\pgfsys@useobject{currentmarker}{}%
\end{pgfscope}%
\begin{pgfscope}%
\pgfsys@transformshift{1.615276in}{0.927084in}%
\pgfsys@useobject{currentmarker}{}%
\end{pgfscope}%
\begin{pgfscope}%
\pgfsys@transformshift{1.612993in}{0.869464in}%
\pgfsys@useobject{currentmarker}{}%
\end{pgfscope}%
\begin{pgfscope}%
\pgfsys@transformshift{1.606350in}{0.921058in}%
\pgfsys@useobject{currentmarker}{}%
\end{pgfscope}%
\begin{pgfscope}%
\pgfsys@transformshift{1.591037in}{0.948325in}%
\pgfsys@useobject{currentmarker}{}%
\end{pgfscope}%
\begin{pgfscope}%
\pgfsys@transformshift{1.653187in}{0.924101in}%
\pgfsys@useobject{currentmarker}{}%
\end{pgfscope}%
\begin{pgfscope}%
\pgfsys@transformshift{1.618853in}{0.853569in}%
\pgfsys@useobject{currentmarker}{}%
\end{pgfscope}%
\begin{pgfscope}%
\pgfsys@transformshift{1.575338in}{0.892473in}%
\pgfsys@useobject{currentmarker}{}%
\end{pgfscope}%
\begin{pgfscope}%
\pgfsys@transformshift{1.632440in}{0.938992in}%
\pgfsys@useobject{currentmarker}{}%
\end{pgfscope}%
\begin{pgfscope}%
\pgfsys@transformshift{1.588799in}{0.936329in}%
\pgfsys@useobject{currentmarker}{}%
\end{pgfscope}%
\begin{pgfscope}%
\pgfsys@transformshift{1.615237in}{0.909094in}%
\pgfsys@useobject{currentmarker}{}%
\end{pgfscope}%
\begin{pgfscope}%
\pgfsys@transformshift{1.615638in}{0.888052in}%
\pgfsys@useobject{currentmarker}{}%
\end{pgfscope}%
\begin{pgfscope}%
\pgfsys@transformshift{1.661069in}{0.948940in}%
\pgfsys@useobject{currentmarker}{}%
\end{pgfscope}%
\begin{pgfscope}%
\pgfsys@transformshift{1.599702in}{0.886960in}%
\pgfsys@useobject{currentmarker}{}%
\end{pgfscope}%
\begin{pgfscope}%
\pgfsys@transformshift{1.557824in}{0.901207in}%
\pgfsys@useobject{currentmarker}{}%
\end{pgfscope}%
\begin{pgfscope}%
\pgfsys@transformshift{1.602321in}{0.893826in}%
\pgfsys@useobject{currentmarker}{}%
\end{pgfscope}%
\begin{pgfscope}%
\pgfsys@transformshift{1.578983in}{0.926291in}%
\pgfsys@useobject{currentmarker}{}%
\end{pgfscope}%
\begin{pgfscope}%
\pgfsys@transformshift{1.608406in}{0.909382in}%
\pgfsys@useobject{currentmarker}{}%
\end{pgfscope}%
\begin{pgfscope}%
\pgfsys@transformshift{1.706128in}{0.961919in}%
\pgfsys@useobject{currentmarker}{}%
\end{pgfscope}%
\begin{pgfscope}%
\pgfsys@transformshift{1.591627in}{0.915488in}%
\pgfsys@useobject{currentmarker}{}%
\end{pgfscope}%
\begin{pgfscope}%
\pgfsys@transformshift{1.647509in}{0.896171in}%
\pgfsys@useobject{currentmarker}{}%
\end{pgfscope}%
\begin{pgfscope}%
\pgfsys@transformshift{1.647471in}{0.906194in}%
\pgfsys@useobject{currentmarker}{}%
\end{pgfscope}%
\begin{pgfscope}%
\pgfsys@transformshift{1.591892in}{0.873254in}%
\pgfsys@useobject{currentmarker}{}%
\end{pgfscope}%
\begin{pgfscope}%
\pgfsys@transformshift{1.613396in}{0.839935in}%
\pgfsys@useobject{currentmarker}{}%
\end{pgfscope}%
\begin{pgfscope}%
\pgfsys@transformshift{1.625983in}{0.954407in}%
\pgfsys@useobject{currentmarker}{}%
\end{pgfscope}%
\begin{pgfscope}%
\pgfsys@transformshift{1.540154in}{0.908153in}%
\pgfsys@useobject{currentmarker}{}%
\end{pgfscope}%
\begin{pgfscope}%
\pgfsys@transformshift{1.630890in}{0.906794in}%
\pgfsys@useobject{currentmarker}{}%
\end{pgfscope}%
\begin{pgfscope}%
\pgfsys@transformshift{1.677788in}{0.938512in}%
\pgfsys@useobject{currentmarker}{}%
\end{pgfscope}%
\begin{pgfscope}%
\pgfsys@transformshift{1.556515in}{0.885823in}%
\pgfsys@useobject{currentmarker}{}%
\end{pgfscope}%
\begin{pgfscope}%
\pgfsys@transformshift{1.618632in}{0.872607in}%
\pgfsys@useobject{currentmarker}{}%
\end{pgfscope}%
\begin{pgfscope}%
\pgfsys@transformshift{1.669642in}{0.921129in}%
\pgfsys@useobject{currentmarker}{}%
\end{pgfscope}%
\begin{pgfscope}%
\pgfsys@transformshift{1.619068in}{0.911863in}%
\pgfsys@useobject{currentmarker}{}%
\end{pgfscope}%
\begin{pgfscope}%
\pgfsys@transformshift{1.585467in}{0.920884in}%
\pgfsys@useobject{currentmarker}{}%
\end{pgfscope}%
\begin{pgfscope}%
\pgfsys@transformshift{1.586052in}{0.872524in}%
\pgfsys@useobject{currentmarker}{}%
\end{pgfscope}%
\begin{pgfscope}%
\pgfsys@transformshift{1.574703in}{0.900112in}%
\pgfsys@useobject{currentmarker}{}%
\end{pgfscope}%
\begin{pgfscope}%
\pgfsys@transformshift{1.578026in}{0.891838in}%
\pgfsys@useobject{currentmarker}{}%
\end{pgfscope}%
\begin{pgfscope}%
\pgfsys@transformshift{1.537539in}{0.872870in}%
\pgfsys@useobject{currentmarker}{}%
\end{pgfscope}%
\begin{pgfscope}%
\pgfsys@transformshift{1.577185in}{0.942105in}%
\pgfsys@useobject{currentmarker}{}%
\end{pgfscope}%
\begin{pgfscope}%
\pgfsys@transformshift{1.655850in}{0.934978in}%
\pgfsys@useobject{currentmarker}{}%
\end{pgfscope}%
\begin{pgfscope}%
\pgfsys@transformshift{1.618521in}{0.861167in}%
\pgfsys@useobject{currentmarker}{}%
\end{pgfscope}%
\begin{pgfscope}%
\pgfsys@transformshift{1.583753in}{0.815510in}%
\pgfsys@useobject{currentmarker}{}%
\end{pgfscope}%
\begin{pgfscope}%
\pgfsys@transformshift{1.613942in}{0.900340in}%
\pgfsys@useobject{currentmarker}{}%
\end{pgfscope}%
\begin{pgfscope}%
\pgfsys@transformshift{1.625545in}{0.934959in}%
\pgfsys@useobject{currentmarker}{}%
\end{pgfscope}%
\begin{pgfscope}%
\pgfsys@transformshift{1.604303in}{0.936168in}%
\pgfsys@useobject{currentmarker}{}%
\end{pgfscope}%
\begin{pgfscope}%
\pgfsys@transformshift{1.547264in}{0.836084in}%
\pgfsys@useobject{currentmarker}{}%
\end{pgfscope}%
\begin{pgfscope}%
\pgfsys@transformshift{1.632850in}{0.912289in}%
\pgfsys@useobject{currentmarker}{}%
\end{pgfscope}%
\end{pgfscope}%
\begin{pgfscope}%
\pgfpathmoveto{\pgfqpoint{2.326000in}{0.100000in}}%
\pgfpathcurveto{\pgfqpoint{2.615602in}{0.100000in}}{\pgfqpoint{2.893381in}{0.157530in}}{\pgfqpoint{3.098161in}{0.259920in}}%
\pgfpathcurveto{\pgfqpoint{3.302940in}{0.362309in}}{\pgfqpoint{3.418000in}{0.501199in}}{\pgfqpoint{3.418000in}{0.646000in}}%
\pgfpathcurveto{\pgfqpoint{3.418000in}{0.790801in}}{\pgfqpoint{3.302940in}{0.929691in}}{\pgfqpoint{3.098161in}{1.032080in}}%
\pgfpathcurveto{\pgfqpoint{2.893381in}{1.134470in}}{\pgfqpoint{2.615602in}{1.192000in}}{\pgfqpoint{2.326000in}{1.192000in}}%
\pgfpathcurveto{\pgfqpoint{2.036398in}{1.192000in}}{\pgfqpoint{1.758619in}{1.134470in}}{\pgfqpoint{1.553839in}{1.032080in}}%
\pgfpathcurveto{\pgfqpoint{1.349060in}{0.929691in}}{\pgfqpoint{1.234000in}{0.790801in}}{\pgfqpoint{1.234000in}{0.646000in}}%
\pgfpathcurveto{\pgfqpoint{1.234000in}{0.501199in}}{\pgfqpoint{1.349060in}{0.362309in}}{\pgfqpoint{1.553839in}{0.259920in}}%
\pgfpathcurveto{\pgfqpoint{1.758619in}{0.157530in}}{\pgfqpoint{2.036398in}{0.100000in}}{\pgfqpoint{2.326000in}{0.100000in}}%
\pgfpathlineto{\pgfqpoint{2.326000in}{0.100000in}}%
\pgfpathclose%
\pgfusepath{clip}%
\pgfsetrectcap%
\pgfsetroundjoin%
\pgfsetlinewidth{0.451687pt}%
\definecolor{currentstroke}{rgb}{0.690196,0.690196,0.690196}%
\pgfsetstrokecolor{currentstroke}%
\pgfsetstrokeopacity{0.250000}%
\pgfsetdash{}{0pt}%
\pgfpathmoveto{\pgfqpoint{1.940921in}{0.151295in}}%
\pgfpathlineto{\pgfqpoint{1.909157in}{0.160652in}}%
\pgfpathlineto{\pgfqpoint{1.878700in}{0.170513in}}%
\pgfpathlineto{\pgfqpoint{1.850255in}{0.180558in}}%
\pgfpathlineto{\pgfqpoint{1.823090in}{0.190956in}}%
\pgfpathlineto{\pgfqpoint{1.797134in}{0.201676in}}%
\pgfpathlineto{\pgfqpoint{1.772319in}{0.212695in}}%
\pgfpathlineto{\pgfqpoint{1.748585in}{0.223994in}}%
\pgfpathlineto{\pgfqpoint{1.725881in}{0.235556in}}%
\pgfpathlineto{\pgfqpoint{1.704160in}{0.247365in}}%
\pgfpathlineto{\pgfqpoint{1.683384in}{0.259410in}}%
\pgfpathlineto{\pgfqpoint{1.663517in}{0.271677in}}%
\pgfpathlineto{\pgfqpoint{1.644532in}{0.284154in}}%
\pgfpathlineto{\pgfqpoint{1.626403in}{0.296831in}}%
\pgfpathlineto{\pgfqpoint{1.609106in}{0.309698in}}%
\pgfpathlineto{\pgfqpoint{1.592426in}{0.322907in}}%
\pgfpathlineto{\pgfqpoint{1.576806in}{0.336079in}}%
\pgfpathlineto{\pgfqpoint{1.561948in}{0.349425in}}%
\pgfpathlineto{\pgfqpoint{1.547847in}{0.362931in}}%
\pgfpathlineto{\pgfqpoint{1.534494in}{0.376589in}}%
\pgfpathlineto{\pgfqpoint{1.521882in}{0.390387in}}%
\pgfpathlineto{\pgfqpoint{1.510004in}{0.404317in}}%
\pgfpathlineto{\pgfqpoint{1.498855in}{0.418370in}}%
\pgfpathlineto{\pgfqpoint{1.488428in}{0.432537in}}%
\pgfpathlineto{\pgfqpoint{1.478718in}{0.446809in}}%
\pgfpathlineto{\pgfqpoint{1.469720in}{0.461179in}}%
\pgfpathlineto{\pgfqpoint{1.461430in}{0.475639in}}%
\pgfpathlineto{\pgfqpoint{1.453844in}{0.490181in}}%
\pgfpathlineto{\pgfqpoint{1.446957in}{0.504798in}}%
\pgfpathlineto{\pgfqpoint{1.440768in}{0.519482in}}%
\pgfpathlineto{\pgfqpoint{1.435272in}{0.534226in}}%
\pgfpathlineto{\pgfqpoint{1.430469in}{0.549023in}}%
\pgfpathlineto{\pgfqpoint{1.426314in}{0.564029in}}%
\pgfpathlineto{\pgfqpoint{1.422911in}{0.578836in}}%
\pgfpathlineto{\pgfqpoint{1.420184in}{0.593701in}}%
\pgfpathlineto{\pgfqpoint{1.418136in}{0.608611in}}%
\pgfpathlineto{\pgfqpoint{1.416769in}{0.623554in}}%
\pgfpathlineto{\pgfqpoint{1.416085in}{0.638516in}}%
\pgfpathlineto{\pgfqpoint{1.416085in}{0.653484in}}%
\pgfpathlineto{\pgfqpoint{1.416769in}{0.668446in}}%
\pgfpathlineto{\pgfqpoint{1.418136in}{0.683389in}}%
\pgfpathlineto{\pgfqpoint{1.420184in}{0.698299in}}%
\pgfpathlineto{\pgfqpoint{1.422911in}{0.713164in}}%
\pgfpathlineto{\pgfqpoint{1.426314in}{0.727971in}}%
\pgfpathlineto{\pgfqpoint{1.430469in}{0.742977in}}%
\pgfpathlineto{\pgfqpoint{1.435272in}{0.757774in}}%
\pgfpathlineto{\pgfqpoint{1.440768in}{0.772518in}}%
\pgfpathlineto{\pgfqpoint{1.446957in}{0.787202in}}%
\pgfpathlineto{\pgfqpoint{1.453844in}{0.801819in}}%
\pgfpathlineto{\pgfqpoint{1.461430in}{0.816361in}}%
\pgfpathlineto{\pgfqpoint{1.469720in}{0.830821in}}%
\pgfpathlineto{\pgfqpoint{1.478718in}{0.845191in}}%
\pgfpathlineto{\pgfqpoint{1.488428in}{0.859463in}}%
\pgfpathlineto{\pgfqpoint{1.498855in}{0.873630in}}%
\pgfpathlineto{\pgfqpoint{1.510004in}{0.887683in}}%
\pgfpathlineto{\pgfqpoint{1.521882in}{0.901613in}}%
\pgfpathlineto{\pgfqpoint{1.534494in}{0.915411in}}%
\pgfpathlineto{\pgfqpoint{1.547847in}{0.929069in}}%
\pgfpathlineto{\pgfqpoint{1.561948in}{0.942575in}}%
\pgfpathlineto{\pgfqpoint{1.576806in}{0.955921in}}%
\pgfpathlineto{\pgfqpoint{1.592426in}{0.969093in}}%
\pgfpathlineto{\pgfqpoint{1.609106in}{0.982302in}}%
\pgfpathlineto{\pgfqpoint{1.626403in}{0.995169in}}%
\pgfpathlineto{\pgfqpoint{1.644532in}{1.007846in}}%
\pgfpathlineto{\pgfqpoint{1.663517in}{1.020323in}}%
\pgfpathlineto{\pgfqpoint{1.683384in}{1.032590in}}%
\pgfpathlineto{\pgfqpoint{1.704160in}{1.044635in}}%
\pgfpathlineto{\pgfqpoint{1.725881in}{1.056444in}}%
\pgfpathlineto{\pgfqpoint{1.748585in}{1.068006in}}%
\pgfpathlineto{\pgfqpoint{1.772319in}{1.079305in}}%
\pgfpathlineto{\pgfqpoint{1.797134in}{1.090324in}}%
\pgfpathlineto{\pgfqpoint{1.823090in}{1.101044in}}%
\pgfpathlineto{\pgfqpoint{1.850255in}{1.111442in}}%
\pgfpathlineto{\pgfqpoint{1.878700in}{1.121487in}}%
\pgfpathlineto{\pgfqpoint{1.909157in}{1.131348in}}%
\pgfpathlineto{\pgfqpoint{1.940921in}{1.140705in}}%
\pgfusepath{stroke}%
\end{pgfscope}%
\begin{pgfscope}%
\pgfpathmoveto{\pgfqpoint{2.326000in}{0.100000in}}%
\pgfpathcurveto{\pgfqpoint{2.615602in}{0.100000in}}{\pgfqpoint{2.893381in}{0.157530in}}{\pgfqpoint{3.098161in}{0.259920in}}%
\pgfpathcurveto{\pgfqpoint{3.302940in}{0.362309in}}{\pgfqpoint{3.418000in}{0.501199in}}{\pgfqpoint{3.418000in}{0.646000in}}%
\pgfpathcurveto{\pgfqpoint{3.418000in}{0.790801in}}{\pgfqpoint{3.302940in}{0.929691in}}{\pgfqpoint{3.098161in}{1.032080in}}%
\pgfpathcurveto{\pgfqpoint{2.893381in}{1.134470in}}{\pgfqpoint{2.615602in}{1.192000in}}{\pgfqpoint{2.326000in}{1.192000in}}%
\pgfpathcurveto{\pgfqpoint{2.036398in}{1.192000in}}{\pgfqpoint{1.758619in}{1.134470in}}{\pgfqpoint{1.553839in}{1.032080in}}%
\pgfpathcurveto{\pgfqpoint{1.349060in}{0.929691in}}{\pgfqpoint{1.234000in}{0.790801in}}{\pgfqpoint{1.234000in}{0.646000in}}%
\pgfpathcurveto{\pgfqpoint{1.234000in}{0.501199in}}{\pgfqpoint{1.349060in}{0.362309in}}{\pgfqpoint{1.553839in}{0.259920in}}%
\pgfpathcurveto{\pgfqpoint{1.758619in}{0.157530in}}{\pgfqpoint{2.036398in}{0.100000in}}{\pgfqpoint{2.326000in}{0.100000in}}%
\pgfpathlineto{\pgfqpoint{2.326000in}{0.100000in}}%
\pgfpathclose%
\pgfusepath{clip}%
\pgfsetrectcap%
\pgfsetroundjoin%
\pgfsetlinewidth{0.451687pt}%
\definecolor{currentstroke}{rgb}{0.690196,0.690196,0.690196}%
\pgfsetstrokecolor{currentstroke}%
\pgfsetstrokeopacity{0.250000}%
\pgfsetdash{}{0pt}%
\pgfpathmoveto{\pgfqpoint{2.017937in}{0.151295in}}%
\pgfpathlineto{\pgfqpoint{1.992525in}{0.160652in}}%
\pgfpathlineto{\pgfqpoint{1.968160in}{0.170513in}}%
\pgfpathlineto{\pgfqpoint{1.945404in}{0.180558in}}%
\pgfpathlineto{\pgfqpoint{1.923672in}{0.190956in}}%
\pgfpathlineto{\pgfqpoint{1.902907in}{0.201676in}}%
\pgfpathlineto{\pgfqpoint{1.883055in}{0.212695in}}%
\pgfpathlineto{\pgfqpoint{1.864068in}{0.223994in}}%
\pgfpathlineto{\pgfqpoint{1.845905in}{0.235556in}}%
\pgfpathlineto{\pgfqpoint{1.828528in}{0.247365in}}%
\pgfpathlineto{\pgfqpoint{1.811907in}{0.259410in}}%
\pgfpathlineto{\pgfqpoint{1.796014in}{0.271677in}}%
\pgfpathlineto{\pgfqpoint{1.780826in}{0.284154in}}%
\pgfpathlineto{\pgfqpoint{1.766322in}{0.296831in}}%
\pgfpathlineto{\pgfqpoint{1.752485in}{0.309698in}}%
\pgfpathlineto{\pgfqpoint{1.739141in}{0.322907in}}%
\pgfpathlineto{\pgfqpoint{1.726644in}{0.336079in}}%
\pgfpathlineto{\pgfqpoint{1.714759in}{0.349425in}}%
\pgfpathlineto{\pgfqpoint{1.703477in}{0.362931in}}%
\pgfpathlineto{\pgfqpoint{1.692795in}{0.376589in}}%
\pgfpathlineto{\pgfqpoint{1.682705in}{0.390387in}}%
\pgfpathlineto{\pgfqpoint{1.673204in}{0.404317in}}%
\pgfpathlineto{\pgfqpoint{1.664284in}{0.418370in}}%
\pgfpathlineto{\pgfqpoint{1.655943in}{0.432537in}}%
\pgfpathlineto{\pgfqpoint{1.648175in}{0.446809in}}%
\pgfpathlineto{\pgfqpoint{1.640976in}{0.461179in}}%
\pgfpathlineto{\pgfqpoint{1.634344in}{0.475639in}}%
\pgfpathlineto{\pgfqpoint{1.628275in}{0.490181in}}%
\pgfpathlineto{\pgfqpoint{1.622766in}{0.504798in}}%
\pgfpathlineto{\pgfqpoint{1.617814in}{0.519482in}}%
\pgfpathlineto{\pgfqpoint{1.613418in}{0.534226in}}%
\pgfpathlineto{\pgfqpoint{1.609575in}{0.549023in}}%
\pgfpathlineto{\pgfqpoint{1.606251in}{0.564029in}}%
\pgfpathlineto{\pgfqpoint{1.603529in}{0.578836in}}%
\pgfpathlineto{\pgfqpoint{1.601347in}{0.593701in}}%
\pgfpathlineto{\pgfqpoint{1.599709in}{0.608611in}}%
\pgfpathlineto{\pgfqpoint{1.598615in}{0.623554in}}%
\pgfpathlineto{\pgfqpoint{1.598068in}{0.638516in}}%
\pgfpathlineto{\pgfqpoint{1.598068in}{0.653484in}}%
\pgfpathlineto{\pgfqpoint{1.598615in}{0.668446in}}%
\pgfpathlineto{\pgfqpoint{1.599709in}{0.683389in}}%
\pgfpathlineto{\pgfqpoint{1.601347in}{0.698299in}}%
\pgfpathlineto{\pgfqpoint{1.603529in}{0.713164in}}%
\pgfpathlineto{\pgfqpoint{1.606251in}{0.727971in}}%
\pgfpathlineto{\pgfqpoint{1.609575in}{0.742977in}}%
\pgfpathlineto{\pgfqpoint{1.613418in}{0.757774in}}%
\pgfpathlineto{\pgfqpoint{1.617814in}{0.772518in}}%
\pgfpathlineto{\pgfqpoint{1.622766in}{0.787202in}}%
\pgfpathlineto{\pgfqpoint{1.628275in}{0.801819in}}%
\pgfpathlineto{\pgfqpoint{1.634344in}{0.816361in}}%
\pgfpathlineto{\pgfqpoint{1.640976in}{0.830821in}}%
\pgfpathlineto{\pgfqpoint{1.648175in}{0.845191in}}%
\pgfpathlineto{\pgfqpoint{1.655943in}{0.859463in}}%
\pgfpathlineto{\pgfqpoint{1.664284in}{0.873630in}}%
\pgfpathlineto{\pgfqpoint{1.673204in}{0.887683in}}%
\pgfpathlineto{\pgfqpoint{1.682705in}{0.901613in}}%
\pgfpathlineto{\pgfqpoint{1.692795in}{0.915411in}}%
\pgfpathlineto{\pgfqpoint{1.703477in}{0.929069in}}%
\pgfpathlineto{\pgfqpoint{1.714759in}{0.942575in}}%
\pgfpathlineto{\pgfqpoint{1.726644in}{0.955921in}}%
\pgfpathlineto{\pgfqpoint{1.739141in}{0.969093in}}%
\pgfpathlineto{\pgfqpoint{1.752485in}{0.982302in}}%
\pgfpathlineto{\pgfqpoint{1.766322in}{0.995169in}}%
\pgfpathlineto{\pgfqpoint{1.780826in}{1.007846in}}%
\pgfpathlineto{\pgfqpoint{1.796014in}{1.020323in}}%
\pgfpathlineto{\pgfqpoint{1.811907in}{1.032590in}}%
\pgfpathlineto{\pgfqpoint{1.828528in}{1.044635in}}%
\pgfpathlineto{\pgfqpoint{1.845905in}{1.056444in}}%
\pgfpathlineto{\pgfqpoint{1.864068in}{1.068006in}}%
\pgfpathlineto{\pgfqpoint{1.883055in}{1.079305in}}%
\pgfpathlineto{\pgfqpoint{1.902907in}{1.090324in}}%
\pgfpathlineto{\pgfqpoint{1.923672in}{1.101044in}}%
\pgfpathlineto{\pgfqpoint{1.945404in}{1.111442in}}%
\pgfpathlineto{\pgfqpoint{1.968160in}{1.121487in}}%
\pgfpathlineto{\pgfqpoint{1.992525in}{1.131348in}}%
\pgfpathlineto{\pgfqpoint{2.017937in}{1.140705in}}%
\pgfusepath{stroke}%
\end{pgfscope}%
\begin{pgfscope}%
\pgfpathmoveto{\pgfqpoint{2.326000in}{0.100000in}}%
\pgfpathcurveto{\pgfqpoint{2.615602in}{0.100000in}}{\pgfqpoint{2.893381in}{0.157530in}}{\pgfqpoint{3.098161in}{0.259920in}}%
\pgfpathcurveto{\pgfqpoint{3.302940in}{0.362309in}}{\pgfqpoint{3.418000in}{0.501199in}}{\pgfqpoint{3.418000in}{0.646000in}}%
\pgfpathcurveto{\pgfqpoint{3.418000in}{0.790801in}}{\pgfqpoint{3.302940in}{0.929691in}}{\pgfqpoint{3.098161in}{1.032080in}}%
\pgfpathcurveto{\pgfqpoint{2.893381in}{1.134470in}}{\pgfqpoint{2.615602in}{1.192000in}}{\pgfqpoint{2.326000in}{1.192000in}}%
\pgfpathcurveto{\pgfqpoint{2.036398in}{1.192000in}}{\pgfqpoint{1.758619in}{1.134470in}}{\pgfqpoint{1.553839in}{1.032080in}}%
\pgfpathcurveto{\pgfqpoint{1.349060in}{0.929691in}}{\pgfqpoint{1.234000in}{0.790801in}}{\pgfqpoint{1.234000in}{0.646000in}}%
\pgfpathcurveto{\pgfqpoint{1.234000in}{0.501199in}}{\pgfqpoint{1.349060in}{0.362309in}}{\pgfqpoint{1.553839in}{0.259920in}}%
\pgfpathcurveto{\pgfqpoint{1.758619in}{0.157530in}}{\pgfqpoint{2.036398in}{0.100000in}}{\pgfqpoint{2.326000in}{0.100000in}}%
\pgfpathlineto{\pgfqpoint{2.326000in}{0.100000in}}%
\pgfpathclose%
\pgfusepath{clip}%
\pgfsetrectcap%
\pgfsetroundjoin%
\pgfsetlinewidth{0.451687pt}%
\definecolor{currentstroke}{rgb}{0.690196,0.690196,0.690196}%
\pgfsetstrokecolor{currentstroke}%
\pgfsetstrokeopacity{0.250000}%
\pgfsetdash{}{0pt}%
\pgfpathmoveto{\pgfqpoint{2.094953in}{0.151295in}}%
\pgfpathlineto{\pgfqpoint{2.075894in}{0.160652in}}%
\pgfpathlineto{\pgfqpoint{2.057620in}{0.170513in}}%
\pgfpathlineto{\pgfqpoint{2.040553in}{0.180558in}}%
\pgfpathlineto{\pgfqpoint{2.024254in}{0.190956in}}%
\pgfpathlineto{\pgfqpoint{2.008680in}{0.201676in}}%
\pgfpathlineto{\pgfqpoint{1.993791in}{0.212695in}}%
\pgfpathlineto{\pgfqpoint{1.979551in}{0.223994in}}%
\pgfpathlineto{\pgfqpoint{1.965929in}{0.235556in}}%
\pgfpathlineto{\pgfqpoint{1.952896in}{0.247365in}}%
\pgfpathlineto{\pgfqpoint{1.940430in}{0.259410in}}%
\pgfpathlineto{\pgfqpoint{1.928510in}{0.271677in}}%
\pgfpathlineto{\pgfqpoint{1.917119in}{0.284154in}}%
\pgfpathlineto{\pgfqpoint{1.906242in}{0.296831in}}%
\pgfpathlineto{\pgfqpoint{1.895864in}{0.309698in}}%
\pgfpathlineto{\pgfqpoint{1.885856in}{0.322907in}}%
\pgfpathlineto{\pgfqpoint{1.876483in}{0.336079in}}%
\pgfpathlineto{\pgfqpoint{1.867569in}{0.349425in}}%
\pgfpathlineto{\pgfqpoint{1.859108in}{0.362931in}}%
\pgfpathlineto{\pgfqpoint{1.851096in}{0.376589in}}%
\pgfpathlineto{\pgfqpoint{1.843529in}{0.390387in}}%
\pgfpathlineto{\pgfqpoint{1.836403in}{0.404317in}}%
\pgfpathlineto{\pgfqpoint{1.829713in}{0.418370in}}%
\pgfpathlineto{\pgfqpoint{1.823457in}{0.432537in}}%
\pgfpathlineto{\pgfqpoint{1.817631in}{0.446809in}}%
\pgfpathlineto{\pgfqpoint{1.812232in}{0.461179in}}%
\pgfpathlineto{\pgfqpoint{1.807258in}{0.475639in}}%
\pgfpathlineto{\pgfqpoint{1.802706in}{0.490181in}}%
\pgfpathlineto{\pgfqpoint{1.798574in}{0.504798in}}%
\pgfpathlineto{\pgfqpoint{1.794861in}{0.519482in}}%
\pgfpathlineto{\pgfqpoint{1.791563in}{0.534226in}}%
\pgfpathlineto{\pgfqpoint{1.788681in}{0.549023in}}%
\pgfpathlineto{\pgfqpoint{1.786188in}{0.564029in}}%
\pgfpathlineto{\pgfqpoint{1.784147in}{0.578836in}}%
\pgfpathlineto{\pgfqpoint{1.782510in}{0.593701in}}%
\pgfpathlineto{\pgfqpoint{1.781282in}{0.608611in}}%
\pgfpathlineto{\pgfqpoint{1.780462in}{0.623554in}}%
\pgfpathlineto{\pgfqpoint{1.780051in}{0.638516in}}%
\pgfpathlineto{\pgfqpoint{1.780051in}{0.653484in}}%
\pgfpathlineto{\pgfqpoint{1.780462in}{0.668446in}}%
\pgfpathlineto{\pgfqpoint{1.781282in}{0.683389in}}%
\pgfpathlineto{\pgfqpoint{1.782510in}{0.698299in}}%
\pgfpathlineto{\pgfqpoint{1.784147in}{0.713164in}}%
\pgfpathlineto{\pgfqpoint{1.786188in}{0.727971in}}%
\pgfpathlineto{\pgfqpoint{1.788681in}{0.742977in}}%
\pgfpathlineto{\pgfqpoint{1.791563in}{0.757774in}}%
\pgfpathlineto{\pgfqpoint{1.794861in}{0.772518in}}%
\pgfpathlineto{\pgfqpoint{1.798574in}{0.787202in}}%
\pgfpathlineto{\pgfqpoint{1.802706in}{0.801819in}}%
\pgfpathlineto{\pgfqpoint{1.807258in}{0.816361in}}%
\pgfpathlineto{\pgfqpoint{1.812232in}{0.830821in}}%
\pgfpathlineto{\pgfqpoint{1.817631in}{0.845191in}}%
\pgfpathlineto{\pgfqpoint{1.823457in}{0.859463in}}%
\pgfpathlineto{\pgfqpoint{1.829713in}{0.873630in}}%
\pgfpathlineto{\pgfqpoint{1.836403in}{0.887683in}}%
\pgfpathlineto{\pgfqpoint{1.843529in}{0.901613in}}%
\pgfpathlineto{\pgfqpoint{1.851096in}{0.915411in}}%
\pgfpathlineto{\pgfqpoint{1.859108in}{0.929069in}}%
\pgfpathlineto{\pgfqpoint{1.867569in}{0.942575in}}%
\pgfpathlineto{\pgfqpoint{1.876483in}{0.955921in}}%
\pgfpathlineto{\pgfqpoint{1.885856in}{0.969093in}}%
\pgfpathlineto{\pgfqpoint{1.895864in}{0.982302in}}%
\pgfpathlineto{\pgfqpoint{1.906242in}{0.995169in}}%
\pgfpathlineto{\pgfqpoint{1.917119in}{1.007846in}}%
\pgfpathlineto{\pgfqpoint{1.928510in}{1.020323in}}%
\pgfpathlineto{\pgfqpoint{1.940430in}{1.032590in}}%
\pgfpathlineto{\pgfqpoint{1.952896in}{1.044635in}}%
\pgfpathlineto{\pgfqpoint{1.965929in}{1.056444in}}%
\pgfpathlineto{\pgfqpoint{1.979551in}{1.068006in}}%
\pgfpathlineto{\pgfqpoint{1.993791in}{1.079305in}}%
\pgfpathlineto{\pgfqpoint{2.008680in}{1.090324in}}%
\pgfpathlineto{\pgfqpoint{2.024254in}{1.101044in}}%
\pgfpathlineto{\pgfqpoint{2.040553in}{1.111442in}}%
\pgfpathlineto{\pgfqpoint{2.057620in}{1.121487in}}%
\pgfpathlineto{\pgfqpoint{2.075894in}{1.131348in}}%
\pgfpathlineto{\pgfqpoint{2.094953in}{1.140705in}}%
\pgfusepath{stroke}%
\end{pgfscope}%
\begin{pgfscope}%
\pgfpathmoveto{\pgfqpoint{2.326000in}{0.100000in}}%
\pgfpathcurveto{\pgfqpoint{2.615602in}{0.100000in}}{\pgfqpoint{2.893381in}{0.157530in}}{\pgfqpoint{3.098161in}{0.259920in}}%
\pgfpathcurveto{\pgfqpoint{3.302940in}{0.362309in}}{\pgfqpoint{3.418000in}{0.501199in}}{\pgfqpoint{3.418000in}{0.646000in}}%
\pgfpathcurveto{\pgfqpoint{3.418000in}{0.790801in}}{\pgfqpoint{3.302940in}{0.929691in}}{\pgfqpoint{3.098161in}{1.032080in}}%
\pgfpathcurveto{\pgfqpoint{2.893381in}{1.134470in}}{\pgfqpoint{2.615602in}{1.192000in}}{\pgfqpoint{2.326000in}{1.192000in}}%
\pgfpathcurveto{\pgfqpoint{2.036398in}{1.192000in}}{\pgfqpoint{1.758619in}{1.134470in}}{\pgfqpoint{1.553839in}{1.032080in}}%
\pgfpathcurveto{\pgfqpoint{1.349060in}{0.929691in}}{\pgfqpoint{1.234000in}{0.790801in}}{\pgfqpoint{1.234000in}{0.646000in}}%
\pgfpathcurveto{\pgfqpoint{1.234000in}{0.501199in}}{\pgfqpoint{1.349060in}{0.362309in}}{\pgfqpoint{1.553839in}{0.259920in}}%
\pgfpathcurveto{\pgfqpoint{1.758619in}{0.157530in}}{\pgfqpoint{2.036398in}{0.100000in}}{\pgfqpoint{2.326000in}{0.100000in}}%
\pgfpathlineto{\pgfqpoint{2.326000in}{0.100000in}}%
\pgfpathclose%
\pgfusepath{clip}%
\pgfsetrectcap%
\pgfsetroundjoin%
\pgfsetlinewidth{0.451687pt}%
\definecolor{currentstroke}{rgb}{0.690196,0.690196,0.690196}%
\pgfsetstrokecolor{currentstroke}%
\pgfsetstrokeopacity{0.250000}%
\pgfsetdash{}{0pt}%
\pgfpathmoveto{\pgfqpoint{2.171969in}{0.151295in}}%
\pgfpathlineto{\pgfqpoint{2.159263in}{0.160652in}}%
\pgfpathlineto{\pgfqpoint{2.147080in}{0.170513in}}%
\pgfpathlineto{\pgfqpoint{2.135702in}{0.180558in}}%
\pgfpathlineto{\pgfqpoint{2.124836in}{0.190956in}}%
\pgfpathlineto{\pgfqpoint{2.114453in}{0.201676in}}%
\pgfpathlineto{\pgfqpoint{2.104527in}{0.212695in}}%
\pgfpathlineto{\pgfqpoint{2.095034in}{0.223994in}}%
\pgfpathlineto{\pgfqpoint{2.085952in}{0.235556in}}%
\pgfpathlineto{\pgfqpoint{2.077264in}{0.247365in}}%
\pgfpathlineto{\pgfqpoint{2.068953in}{0.259410in}}%
\pgfpathlineto{\pgfqpoint{2.061007in}{0.271677in}}%
\pgfpathlineto{\pgfqpoint{2.053413in}{0.284154in}}%
\pgfpathlineto{\pgfqpoint{2.046161in}{0.296831in}}%
\pgfpathlineto{\pgfqpoint{2.039243in}{0.309698in}}%
\pgfpathlineto{\pgfqpoint{2.032570in}{0.322907in}}%
\pgfpathlineto{\pgfqpoint{2.026322in}{0.336079in}}%
\pgfpathlineto{\pgfqpoint{2.020379in}{0.349425in}}%
\pgfpathlineto{\pgfqpoint{2.014739in}{0.362931in}}%
\pgfpathlineto{\pgfqpoint{2.009398in}{0.376589in}}%
\pgfpathlineto{\pgfqpoint{2.004353in}{0.390387in}}%
\pgfpathlineto{\pgfqpoint{1.999602in}{0.404317in}}%
\pgfpathlineto{\pgfqpoint{1.995142in}{0.418370in}}%
\pgfpathlineto{\pgfqpoint{1.990971in}{0.432537in}}%
\pgfpathlineto{\pgfqpoint{1.987087in}{0.446809in}}%
\pgfpathlineto{\pgfqpoint{1.983488in}{0.461179in}}%
\pgfpathlineto{\pgfqpoint{1.980172in}{0.475639in}}%
\pgfpathlineto{\pgfqpoint{1.977137in}{0.490181in}}%
\pgfpathlineto{\pgfqpoint{1.974383in}{0.504798in}}%
\pgfpathlineto{\pgfqpoint{1.971907in}{0.519482in}}%
\pgfpathlineto{\pgfqpoint{1.969709in}{0.534226in}}%
\pgfpathlineto{\pgfqpoint{1.967788in}{0.549023in}}%
\pgfpathlineto{\pgfqpoint{1.966125in}{0.564029in}}%
\pgfpathlineto{\pgfqpoint{1.964764in}{0.578836in}}%
\pgfpathlineto{\pgfqpoint{1.963674in}{0.593701in}}%
\pgfpathlineto{\pgfqpoint{1.962854in}{0.608611in}}%
\pgfpathlineto{\pgfqpoint{1.962308in}{0.623554in}}%
\pgfpathlineto{\pgfqpoint{1.962034in}{0.638516in}}%
\pgfpathlineto{\pgfqpoint{1.962034in}{0.653484in}}%
\pgfpathlineto{\pgfqpoint{1.962308in}{0.668446in}}%
\pgfpathlineto{\pgfqpoint{1.962854in}{0.683389in}}%
\pgfpathlineto{\pgfqpoint{1.963674in}{0.698299in}}%
\pgfpathlineto{\pgfqpoint{1.964764in}{0.713164in}}%
\pgfpathlineto{\pgfqpoint{1.966125in}{0.727971in}}%
\pgfpathlineto{\pgfqpoint{1.967788in}{0.742977in}}%
\pgfpathlineto{\pgfqpoint{1.969709in}{0.757774in}}%
\pgfpathlineto{\pgfqpoint{1.971907in}{0.772518in}}%
\pgfpathlineto{\pgfqpoint{1.974383in}{0.787202in}}%
\pgfpathlineto{\pgfqpoint{1.977137in}{0.801819in}}%
\pgfpathlineto{\pgfqpoint{1.980172in}{0.816361in}}%
\pgfpathlineto{\pgfqpoint{1.983488in}{0.830821in}}%
\pgfpathlineto{\pgfqpoint{1.987087in}{0.845191in}}%
\pgfpathlineto{\pgfqpoint{1.990971in}{0.859463in}}%
\pgfpathlineto{\pgfqpoint{1.995142in}{0.873630in}}%
\pgfpathlineto{\pgfqpoint{1.999602in}{0.887683in}}%
\pgfpathlineto{\pgfqpoint{2.004353in}{0.901613in}}%
\pgfpathlineto{\pgfqpoint{2.009398in}{0.915411in}}%
\pgfpathlineto{\pgfqpoint{2.014739in}{0.929069in}}%
\pgfpathlineto{\pgfqpoint{2.020379in}{0.942575in}}%
\pgfpathlineto{\pgfqpoint{2.026322in}{0.955921in}}%
\pgfpathlineto{\pgfqpoint{2.032570in}{0.969093in}}%
\pgfpathlineto{\pgfqpoint{2.039243in}{0.982302in}}%
\pgfpathlineto{\pgfqpoint{2.046161in}{0.995169in}}%
\pgfpathlineto{\pgfqpoint{2.053413in}{1.007846in}}%
\pgfpathlineto{\pgfqpoint{2.061007in}{1.020323in}}%
\pgfpathlineto{\pgfqpoint{2.068953in}{1.032590in}}%
\pgfpathlineto{\pgfqpoint{2.077264in}{1.044635in}}%
\pgfpathlineto{\pgfqpoint{2.085952in}{1.056444in}}%
\pgfpathlineto{\pgfqpoint{2.095034in}{1.068006in}}%
\pgfpathlineto{\pgfqpoint{2.104527in}{1.079305in}}%
\pgfpathlineto{\pgfqpoint{2.114453in}{1.090324in}}%
\pgfpathlineto{\pgfqpoint{2.124836in}{1.101044in}}%
\pgfpathlineto{\pgfqpoint{2.135702in}{1.111442in}}%
\pgfpathlineto{\pgfqpoint{2.147080in}{1.121487in}}%
\pgfpathlineto{\pgfqpoint{2.159263in}{1.131348in}}%
\pgfpathlineto{\pgfqpoint{2.171969in}{1.140705in}}%
\pgfusepath{stroke}%
\end{pgfscope}%
\begin{pgfscope}%
\pgfpathmoveto{\pgfqpoint{2.326000in}{0.100000in}}%
\pgfpathcurveto{\pgfqpoint{2.615602in}{0.100000in}}{\pgfqpoint{2.893381in}{0.157530in}}{\pgfqpoint{3.098161in}{0.259920in}}%
\pgfpathcurveto{\pgfqpoint{3.302940in}{0.362309in}}{\pgfqpoint{3.418000in}{0.501199in}}{\pgfqpoint{3.418000in}{0.646000in}}%
\pgfpathcurveto{\pgfqpoint{3.418000in}{0.790801in}}{\pgfqpoint{3.302940in}{0.929691in}}{\pgfqpoint{3.098161in}{1.032080in}}%
\pgfpathcurveto{\pgfqpoint{2.893381in}{1.134470in}}{\pgfqpoint{2.615602in}{1.192000in}}{\pgfqpoint{2.326000in}{1.192000in}}%
\pgfpathcurveto{\pgfqpoint{2.036398in}{1.192000in}}{\pgfqpoint{1.758619in}{1.134470in}}{\pgfqpoint{1.553839in}{1.032080in}}%
\pgfpathcurveto{\pgfqpoint{1.349060in}{0.929691in}}{\pgfqpoint{1.234000in}{0.790801in}}{\pgfqpoint{1.234000in}{0.646000in}}%
\pgfpathcurveto{\pgfqpoint{1.234000in}{0.501199in}}{\pgfqpoint{1.349060in}{0.362309in}}{\pgfqpoint{1.553839in}{0.259920in}}%
\pgfpathcurveto{\pgfqpoint{1.758619in}{0.157530in}}{\pgfqpoint{2.036398in}{0.100000in}}{\pgfqpoint{2.326000in}{0.100000in}}%
\pgfpathlineto{\pgfqpoint{2.326000in}{0.100000in}}%
\pgfpathclose%
\pgfusepath{clip}%
\pgfsetrectcap%
\pgfsetroundjoin%
\pgfsetlinewidth{0.451687pt}%
\definecolor{currentstroke}{rgb}{0.690196,0.690196,0.690196}%
\pgfsetstrokecolor{currentstroke}%
\pgfsetstrokeopacity{0.250000}%
\pgfsetdash{}{0pt}%
\pgfpathmoveto{\pgfqpoint{2.248984in}{0.151295in}}%
\pgfpathlineto{\pgfqpoint{2.242631in}{0.160652in}}%
\pgfpathlineto{\pgfqpoint{2.236540in}{0.170513in}}%
\pgfpathlineto{\pgfqpoint{2.230851in}{0.180558in}}%
\pgfpathlineto{\pgfqpoint{2.225418in}{0.190956in}}%
\pgfpathlineto{\pgfqpoint{2.220227in}{0.201676in}}%
\pgfpathlineto{\pgfqpoint{2.215264in}{0.212695in}}%
\pgfpathlineto{\pgfqpoint{2.210517in}{0.223994in}}%
\pgfpathlineto{\pgfqpoint{2.205976in}{0.235556in}}%
\pgfpathlineto{\pgfqpoint{2.201632in}{0.247365in}}%
\pgfpathlineto{\pgfqpoint{2.197477in}{0.259410in}}%
\pgfpathlineto{\pgfqpoint{2.193503in}{0.271677in}}%
\pgfpathlineto{\pgfqpoint{2.189706in}{0.284154in}}%
\pgfpathlineto{\pgfqpoint{2.186081in}{0.296831in}}%
\pgfpathlineto{\pgfqpoint{2.182621in}{0.309698in}}%
\pgfpathlineto{\pgfqpoint{2.179285in}{0.322907in}}%
\pgfpathlineto{\pgfqpoint{2.176161in}{0.336079in}}%
\pgfpathlineto{\pgfqpoint{2.173190in}{0.349425in}}%
\pgfpathlineto{\pgfqpoint{2.170369in}{0.362931in}}%
\pgfpathlineto{\pgfqpoint{2.167699in}{0.376589in}}%
\pgfpathlineto{\pgfqpoint{2.165176in}{0.390387in}}%
\pgfpathlineto{\pgfqpoint{2.162801in}{0.404317in}}%
\pgfpathlineto{\pgfqpoint{2.160571in}{0.418370in}}%
\pgfpathlineto{\pgfqpoint{2.158486in}{0.432537in}}%
\pgfpathlineto{\pgfqpoint{2.156544in}{0.446809in}}%
\pgfpathlineto{\pgfqpoint{2.154744in}{0.461179in}}%
\pgfpathlineto{\pgfqpoint{2.153086in}{0.475639in}}%
\pgfpathlineto{\pgfqpoint{2.151569in}{0.490181in}}%
\pgfpathlineto{\pgfqpoint{2.150191in}{0.504798in}}%
\pgfpathlineto{\pgfqpoint{2.148954in}{0.519482in}}%
\pgfpathlineto{\pgfqpoint{2.147854in}{0.534226in}}%
\pgfpathlineto{\pgfqpoint{2.146894in}{0.549023in}}%
\pgfpathlineto{\pgfqpoint{2.146063in}{0.564029in}}%
\pgfpathlineto{\pgfqpoint{2.145382in}{0.578836in}}%
\pgfpathlineto{\pgfqpoint{2.144837in}{0.593701in}}%
\pgfpathlineto{\pgfqpoint{2.144427in}{0.608611in}}%
\pgfpathlineto{\pgfqpoint{2.144154in}{0.623554in}}%
\pgfpathlineto{\pgfqpoint{2.144017in}{0.638516in}}%
\pgfpathlineto{\pgfqpoint{2.144017in}{0.653484in}}%
\pgfpathlineto{\pgfqpoint{2.144154in}{0.668446in}}%
\pgfpathlineto{\pgfqpoint{2.144427in}{0.683389in}}%
\pgfpathlineto{\pgfqpoint{2.144837in}{0.698299in}}%
\pgfpathlineto{\pgfqpoint{2.145382in}{0.713164in}}%
\pgfpathlineto{\pgfqpoint{2.146063in}{0.727971in}}%
\pgfpathlineto{\pgfqpoint{2.146894in}{0.742977in}}%
\pgfpathlineto{\pgfqpoint{2.147854in}{0.757774in}}%
\pgfpathlineto{\pgfqpoint{2.148954in}{0.772518in}}%
\pgfpathlineto{\pgfqpoint{2.150191in}{0.787202in}}%
\pgfpathlineto{\pgfqpoint{2.151569in}{0.801819in}}%
\pgfpathlineto{\pgfqpoint{2.153086in}{0.816361in}}%
\pgfpathlineto{\pgfqpoint{2.154744in}{0.830821in}}%
\pgfpathlineto{\pgfqpoint{2.156544in}{0.845191in}}%
\pgfpathlineto{\pgfqpoint{2.158486in}{0.859463in}}%
\pgfpathlineto{\pgfqpoint{2.160571in}{0.873630in}}%
\pgfpathlineto{\pgfqpoint{2.162801in}{0.887683in}}%
\pgfpathlineto{\pgfqpoint{2.165176in}{0.901613in}}%
\pgfpathlineto{\pgfqpoint{2.167699in}{0.915411in}}%
\pgfpathlineto{\pgfqpoint{2.170369in}{0.929069in}}%
\pgfpathlineto{\pgfqpoint{2.173190in}{0.942575in}}%
\pgfpathlineto{\pgfqpoint{2.176161in}{0.955921in}}%
\pgfpathlineto{\pgfqpoint{2.179285in}{0.969093in}}%
\pgfpathlineto{\pgfqpoint{2.182621in}{0.982302in}}%
\pgfpathlineto{\pgfqpoint{2.186081in}{0.995169in}}%
\pgfpathlineto{\pgfqpoint{2.189706in}{1.007846in}}%
\pgfpathlineto{\pgfqpoint{2.193503in}{1.020323in}}%
\pgfpathlineto{\pgfqpoint{2.197477in}{1.032590in}}%
\pgfpathlineto{\pgfqpoint{2.201632in}{1.044635in}}%
\pgfpathlineto{\pgfqpoint{2.205976in}{1.056444in}}%
\pgfpathlineto{\pgfqpoint{2.210517in}{1.068006in}}%
\pgfpathlineto{\pgfqpoint{2.215264in}{1.079305in}}%
\pgfpathlineto{\pgfqpoint{2.220227in}{1.090324in}}%
\pgfpathlineto{\pgfqpoint{2.225418in}{1.101044in}}%
\pgfpathlineto{\pgfqpoint{2.230851in}{1.111442in}}%
\pgfpathlineto{\pgfqpoint{2.236540in}{1.121487in}}%
\pgfpathlineto{\pgfqpoint{2.242631in}{1.131348in}}%
\pgfpathlineto{\pgfqpoint{2.248984in}{1.140705in}}%
\pgfusepath{stroke}%
\end{pgfscope}%
\begin{pgfscope}%
\pgfpathmoveto{\pgfqpoint{2.326000in}{0.100000in}}%
\pgfpathcurveto{\pgfqpoint{2.615602in}{0.100000in}}{\pgfqpoint{2.893381in}{0.157530in}}{\pgfqpoint{3.098161in}{0.259920in}}%
\pgfpathcurveto{\pgfqpoint{3.302940in}{0.362309in}}{\pgfqpoint{3.418000in}{0.501199in}}{\pgfqpoint{3.418000in}{0.646000in}}%
\pgfpathcurveto{\pgfqpoint{3.418000in}{0.790801in}}{\pgfqpoint{3.302940in}{0.929691in}}{\pgfqpoint{3.098161in}{1.032080in}}%
\pgfpathcurveto{\pgfqpoint{2.893381in}{1.134470in}}{\pgfqpoint{2.615602in}{1.192000in}}{\pgfqpoint{2.326000in}{1.192000in}}%
\pgfpathcurveto{\pgfqpoint{2.036398in}{1.192000in}}{\pgfqpoint{1.758619in}{1.134470in}}{\pgfqpoint{1.553839in}{1.032080in}}%
\pgfpathcurveto{\pgfqpoint{1.349060in}{0.929691in}}{\pgfqpoint{1.234000in}{0.790801in}}{\pgfqpoint{1.234000in}{0.646000in}}%
\pgfpathcurveto{\pgfqpoint{1.234000in}{0.501199in}}{\pgfqpoint{1.349060in}{0.362309in}}{\pgfqpoint{1.553839in}{0.259920in}}%
\pgfpathcurveto{\pgfqpoint{1.758619in}{0.157530in}}{\pgfqpoint{2.036398in}{0.100000in}}{\pgfqpoint{2.326000in}{0.100000in}}%
\pgfpathlineto{\pgfqpoint{2.326000in}{0.100000in}}%
\pgfpathclose%
\pgfusepath{clip}%
\pgfsetrectcap%
\pgfsetroundjoin%
\pgfsetlinewidth{0.451687pt}%
\definecolor{currentstroke}{rgb}{0.690196,0.690196,0.690196}%
\pgfsetstrokecolor{currentstroke}%
\pgfsetstrokeopacity{0.250000}%
\pgfsetdash{}{0pt}%
\pgfpathmoveto{\pgfqpoint{2.326000in}{0.151295in}}%
\pgfpathlineto{\pgfqpoint{2.326000in}{0.160652in}}%
\pgfpathlineto{\pgfqpoint{2.326000in}{0.170513in}}%
\pgfpathlineto{\pgfqpoint{2.326000in}{0.180558in}}%
\pgfpathlineto{\pgfqpoint{2.326000in}{0.190956in}}%
\pgfpathlineto{\pgfqpoint{2.326000in}{0.201676in}}%
\pgfpathlineto{\pgfqpoint{2.326000in}{0.212695in}}%
\pgfpathlineto{\pgfqpoint{2.326000in}{0.223994in}}%
\pgfpathlineto{\pgfqpoint{2.326000in}{0.235556in}}%
\pgfpathlineto{\pgfqpoint{2.326000in}{0.247365in}}%
\pgfpathlineto{\pgfqpoint{2.326000in}{0.259410in}}%
\pgfpathlineto{\pgfqpoint{2.326000in}{0.271677in}}%
\pgfpathlineto{\pgfqpoint{2.326000in}{0.284154in}}%
\pgfpathlineto{\pgfqpoint{2.326000in}{0.296831in}}%
\pgfpathlineto{\pgfqpoint{2.326000in}{0.309698in}}%
\pgfpathlineto{\pgfqpoint{2.326000in}{0.322907in}}%
\pgfpathlineto{\pgfqpoint{2.326000in}{0.336079in}}%
\pgfpathlineto{\pgfqpoint{2.326000in}{0.349425in}}%
\pgfpathlineto{\pgfqpoint{2.326000in}{0.362931in}}%
\pgfpathlineto{\pgfqpoint{2.326000in}{0.376589in}}%
\pgfpathlineto{\pgfqpoint{2.326000in}{0.390387in}}%
\pgfpathlineto{\pgfqpoint{2.326000in}{0.404317in}}%
\pgfpathlineto{\pgfqpoint{2.326000in}{0.418370in}}%
\pgfpathlineto{\pgfqpoint{2.326000in}{0.432537in}}%
\pgfpathlineto{\pgfqpoint{2.326000in}{0.446809in}}%
\pgfpathlineto{\pgfqpoint{2.326000in}{0.461179in}}%
\pgfpathlineto{\pgfqpoint{2.326000in}{0.475639in}}%
\pgfpathlineto{\pgfqpoint{2.326000in}{0.490181in}}%
\pgfpathlineto{\pgfqpoint{2.326000in}{0.504798in}}%
\pgfpathlineto{\pgfqpoint{2.326000in}{0.519482in}}%
\pgfpathlineto{\pgfqpoint{2.326000in}{0.534226in}}%
\pgfpathlineto{\pgfqpoint{2.326000in}{0.549023in}}%
\pgfpathlineto{\pgfqpoint{2.326000in}{0.564029in}}%
\pgfpathlineto{\pgfqpoint{2.326000in}{0.578836in}}%
\pgfpathlineto{\pgfqpoint{2.326000in}{0.593701in}}%
\pgfpathlineto{\pgfqpoint{2.326000in}{0.608611in}}%
\pgfpathlineto{\pgfqpoint{2.326000in}{0.623554in}}%
\pgfpathlineto{\pgfqpoint{2.326000in}{0.638516in}}%
\pgfpathlineto{\pgfqpoint{2.326000in}{0.653484in}}%
\pgfpathlineto{\pgfqpoint{2.326000in}{0.668446in}}%
\pgfpathlineto{\pgfqpoint{2.326000in}{0.683389in}}%
\pgfpathlineto{\pgfqpoint{2.326000in}{0.698299in}}%
\pgfpathlineto{\pgfqpoint{2.326000in}{0.713164in}}%
\pgfpathlineto{\pgfqpoint{2.326000in}{0.727971in}}%
\pgfpathlineto{\pgfqpoint{2.326000in}{0.742977in}}%
\pgfpathlineto{\pgfqpoint{2.326000in}{0.757774in}}%
\pgfpathlineto{\pgfqpoint{2.326000in}{0.772518in}}%
\pgfpathlineto{\pgfqpoint{2.326000in}{0.787202in}}%
\pgfpathlineto{\pgfqpoint{2.326000in}{0.801819in}}%
\pgfpathlineto{\pgfqpoint{2.326000in}{0.816361in}}%
\pgfpathlineto{\pgfqpoint{2.326000in}{0.830821in}}%
\pgfpathlineto{\pgfqpoint{2.326000in}{0.845191in}}%
\pgfpathlineto{\pgfqpoint{2.326000in}{0.859463in}}%
\pgfpathlineto{\pgfqpoint{2.326000in}{0.873630in}}%
\pgfpathlineto{\pgfqpoint{2.326000in}{0.887683in}}%
\pgfpathlineto{\pgfqpoint{2.326000in}{0.901613in}}%
\pgfpathlineto{\pgfqpoint{2.326000in}{0.915411in}}%
\pgfpathlineto{\pgfqpoint{2.326000in}{0.929069in}}%
\pgfpathlineto{\pgfqpoint{2.326000in}{0.942575in}}%
\pgfpathlineto{\pgfqpoint{2.326000in}{0.955921in}}%
\pgfpathlineto{\pgfqpoint{2.326000in}{0.969093in}}%
\pgfpathlineto{\pgfqpoint{2.326000in}{0.982302in}}%
\pgfpathlineto{\pgfqpoint{2.326000in}{0.995169in}}%
\pgfpathlineto{\pgfqpoint{2.326000in}{1.007846in}}%
\pgfpathlineto{\pgfqpoint{2.326000in}{1.020323in}}%
\pgfpathlineto{\pgfqpoint{2.326000in}{1.032590in}}%
\pgfpathlineto{\pgfqpoint{2.326000in}{1.044635in}}%
\pgfpathlineto{\pgfqpoint{2.326000in}{1.056444in}}%
\pgfpathlineto{\pgfqpoint{2.326000in}{1.068006in}}%
\pgfpathlineto{\pgfqpoint{2.326000in}{1.079305in}}%
\pgfpathlineto{\pgfqpoint{2.326000in}{1.090324in}}%
\pgfpathlineto{\pgfqpoint{2.326000in}{1.101044in}}%
\pgfpathlineto{\pgfqpoint{2.326000in}{1.111442in}}%
\pgfpathlineto{\pgfqpoint{2.326000in}{1.121487in}}%
\pgfpathlineto{\pgfqpoint{2.326000in}{1.131348in}}%
\pgfpathlineto{\pgfqpoint{2.326000in}{1.140705in}}%
\pgfusepath{stroke}%
\end{pgfscope}%
\begin{pgfscope}%
\pgfpathmoveto{\pgfqpoint{2.326000in}{0.100000in}}%
\pgfpathcurveto{\pgfqpoint{2.615602in}{0.100000in}}{\pgfqpoint{2.893381in}{0.157530in}}{\pgfqpoint{3.098161in}{0.259920in}}%
\pgfpathcurveto{\pgfqpoint{3.302940in}{0.362309in}}{\pgfqpoint{3.418000in}{0.501199in}}{\pgfqpoint{3.418000in}{0.646000in}}%
\pgfpathcurveto{\pgfqpoint{3.418000in}{0.790801in}}{\pgfqpoint{3.302940in}{0.929691in}}{\pgfqpoint{3.098161in}{1.032080in}}%
\pgfpathcurveto{\pgfqpoint{2.893381in}{1.134470in}}{\pgfqpoint{2.615602in}{1.192000in}}{\pgfqpoint{2.326000in}{1.192000in}}%
\pgfpathcurveto{\pgfqpoint{2.036398in}{1.192000in}}{\pgfqpoint{1.758619in}{1.134470in}}{\pgfqpoint{1.553839in}{1.032080in}}%
\pgfpathcurveto{\pgfqpoint{1.349060in}{0.929691in}}{\pgfqpoint{1.234000in}{0.790801in}}{\pgfqpoint{1.234000in}{0.646000in}}%
\pgfpathcurveto{\pgfqpoint{1.234000in}{0.501199in}}{\pgfqpoint{1.349060in}{0.362309in}}{\pgfqpoint{1.553839in}{0.259920in}}%
\pgfpathcurveto{\pgfqpoint{1.758619in}{0.157530in}}{\pgfqpoint{2.036398in}{0.100000in}}{\pgfqpoint{2.326000in}{0.100000in}}%
\pgfpathlineto{\pgfqpoint{2.326000in}{0.100000in}}%
\pgfpathclose%
\pgfusepath{clip}%
\pgfsetrectcap%
\pgfsetroundjoin%
\pgfsetlinewidth{0.451687pt}%
\definecolor{currentstroke}{rgb}{0.690196,0.690196,0.690196}%
\pgfsetstrokecolor{currentstroke}%
\pgfsetstrokeopacity{0.250000}%
\pgfsetdash{}{0pt}%
\pgfpathmoveto{\pgfqpoint{2.403016in}{0.151295in}}%
\pgfpathlineto{\pgfqpoint{2.409369in}{0.160652in}}%
\pgfpathlineto{\pgfqpoint{2.415460in}{0.170513in}}%
\pgfpathlineto{\pgfqpoint{2.421149in}{0.180558in}}%
\pgfpathlineto{\pgfqpoint{2.426582in}{0.190956in}}%
\pgfpathlineto{\pgfqpoint{2.431773in}{0.201676in}}%
\pgfpathlineto{\pgfqpoint{2.436736in}{0.212695in}}%
\pgfpathlineto{\pgfqpoint{2.441483in}{0.223994in}}%
\pgfpathlineto{\pgfqpoint{2.446024in}{0.235556in}}%
\pgfpathlineto{\pgfqpoint{2.450368in}{0.247365in}}%
\pgfpathlineto{\pgfqpoint{2.454523in}{0.259410in}}%
\pgfpathlineto{\pgfqpoint{2.458497in}{0.271677in}}%
\pgfpathlineto{\pgfqpoint{2.462294in}{0.284154in}}%
\pgfpathlineto{\pgfqpoint{2.465919in}{0.296831in}}%
\pgfpathlineto{\pgfqpoint{2.469379in}{0.309698in}}%
\pgfpathlineto{\pgfqpoint{2.472715in}{0.322907in}}%
\pgfpathlineto{\pgfqpoint{2.475839in}{0.336079in}}%
\pgfpathlineto{\pgfqpoint{2.478810in}{0.349425in}}%
\pgfpathlineto{\pgfqpoint{2.481631in}{0.362931in}}%
\pgfpathlineto{\pgfqpoint{2.484301in}{0.376589in}}%
\pgfpathlineto{\pgfqpoint{2.486824in}{0.390387in}}%
\pgfpathlineto{\pgfqpoint{2.489199in}{0.404317in}}%
\pgfpathlineto{\pgfqpoint{2.491429in}{0.418370in}}%
\pgfpathlineto{\pgfqpoint{2.493514in}{0.432537in}}%
\pgfpathlineto{\pgfqpoint{2.495456in}{0.446809in}}%
\pgfpathlineto{\pgfqpoint{2.497256in}{0.461179in}}%
\pgfpathlineto{\pgfqpoint{2.498914in}{0.475639in}}%
\pgfpathlineto{\pgfqpoint{2.500431in}{0.490181in}}%
\pgfpathlineto{\pgfqpoint{2.501809in}{0.504798in}}%
\pgfpathlineto{\pgfqpoint{2.503046in}{0.519482in}}%
\pgfpathlineto{\pgfqpoint{2.504146in}{0.534226in}}%
\pgfpathlineto{\pgfqpoint{2.505106in}{0.549023in}}%
\pgfpathlineto{\pgfqpoint{2.505937in}{0.564029in}}%
\pgfpathlineto{\pgfqpoint{2.506618in}{0.578836in}}%
\pgfpathlineto{\pgfqpoint{2.507163in}{0.593701in}}%
\pgfpathlineto{\pgfqpoint{2.507573in}{0.608611in}}%
\pgfpathlineto{\pgfqpoint{2.507846in}{0.623554in}}%
\pgfpathlineto{\pgfqpoint{2.507983in}{0.638516in}}%
\pgfpathlineto{\pgfqpoint{2.507983in}{0.653484in}}%
\pgfpathlineto{\pgfqpoint{2.507846in}{0.668446in}}%
\pgfpathlineto{\pgfqpoint{2.507573in}{0.683389in}}%
\pgfpathlineto{\pgfqpoint{2.507163in}{0.698299in}}%
\pgfpathlineto{\pgfqpoint{2.506618in}{0.713164in}}%
\pgfpathlineto{\pgfqpoint{2.505937in}{0.727971in}}%
\pgfpathlineto{\pgfqpoint{2.505106in}{0.742977in}}%
\pgfpathlineto{\pgfqpoint{2.504146in}{0.757774in}}%
\pgfpathlineto{\pgfqpoint{2.503046in}{0.772518in}}%
\pgfpathlineto{\pgfqpoint{2.501809in}{0.787202in}}%
\pgfpathlineto{\pgfqpoint{2.500431in}{0.801819in}}%
\pgfpathlineto{\pgfqpoint{2.498914in}{0.816361in}}%
\pgfpathlineto{\pgfqpoint{2.497256in}{0.830821in}}%
\pgfpathlineto{\pgfqpoint{2.495456in}{0.845191in}}%
\pgfpathlineto{\pgfqpoint{2.493514in}{0.859463in}}%
\pgfpathlineto{\pgfqpoint{2.491429in}{0.873630in}}%
\pgfpathlineto{\pgfqpoint{2.489199in}{0.887683in}}%
\pgfpathlineto{\pgfqpoint{2.486824in}{0.901613in}}%
\pgfpathlineto{\pgfqpoint{2.484301in}{0.915411in}}%
\pgfpathlineto{\pgfqpoint{2.481631in}{0.929069in}}%
\pgfpathlineto{\pgfqpoint{2.478810in}{0.942575in}}%
\pgfpathlineto{\pgfqpoint{2.475839in}{0.955921in}}%
\pgfpathlineto{\pgfqpoint{2.472715in}{0.969093in}}%
\pgfpathlineto{\pgfqpoint{2.469379in}{0.982302in}}%
\pgfpathlineto{\pgfqpoint{2.465919in}{0.995169in}}%
\pgfpathlineto{\pgfqpoint{2.462294in}{1.007846in}}%
\pgfpathlineto{\pgfqpoint{2.458497in}{1.020323in}}%
\pgfpathlineto{\pgfqpoint{2.454523in}{1.032590in}}%
\pgfpathlineto{\pgfqpoint{2.450368in}{1.044635in}}%
\pgfpathlineto{\pgfqpoint{2.446024in}{1.056444in}}%
\pgfpathlineto{\pgfqpoint{2.441483in}{1.068006in}}%
\pgfpathlineto{\pgfqpoint{2.436736in}{1.079305in}}%
\pgfpathlineto{\pgfqpoint{2.431773in}{1.090324in}}%
\pgfpathlineto{\pgfqpoint{2.426582in}{1.101044in}}%
\pgfpathlineto{\pgfqpoint{2.421149in}{1.111442in}}%
\pgfpathlineto{\pgfqpoint{2.415460in}{1.121487in}}%
\pgfpathlineto{\pgfqpoint{2.409369in}{1.131348in}}%
\pgfpathlineto{\pgfqpoint{2.403016in}{1.140705in}}%
\pgfusepath{stroke}%
\end{pgfscope}%
\begin{pgfscope}%
\pgfpathmoveto{\pgfqpoint{2.326000in}{0.100000in}}%
\pgfpathcurveto{\pgfqpoint{2.615602in}{0.100000in}}{\pgfqpoint{2.893381in}{0.157530in}}{\pgfqpoint{3.098161in}{0.259920in}}%
\pgfpathcurveto{\pgfqpoint{3.302940in}{0.362309in}}{\pgfqpoint{3.418000in}{0.501199in}}{\pgfqpoint{3.418000in}{0.646000in}}%
\pgfpathcurveto{\pgfqpoint{3.418000in}{0.790801in}}{\pgfqpoint{3.302940in}{0.929691in}}{\pgfqpoint{3.098161in}{1.032080in}}%
\pgfpathcurveto{\pgfqpoint{2.893381in}{1.134470in}}{\pgfqpoint{2.615602in}{1.192000in}}{\pgfqpoint{2.326000in}{1.192000in}}%
\pgfpathcurveto{\pgfqpoint{2.036398in}{1.192000in}}{\pgfqpoint{1.758619in}{1.134470in}}{\pgfqpoint{1.553839in}{1.032080in}}%
\pgfpathcurveto{\pgfqpoint{1.349060in}{0.929691in}}{\pgfqpoint{1.234000in}{0.790801in}}{\pgfqpoint{1.234000in}{0.646000in}}%
\pgfpathcurveto{\pgfqpoint{1.234000in}{0.501199in}}{\pgfqpoint{1.349060in}{0.362309in}}{\pgfqpoint{1.553839in}{0.259920in}}%
\pgfpathcurveto{\pgfqpoint{1.758619in}{0.157530in}}{\pgfqpoint{2.036398in}{0.100000in}}{\pgfqpoint{2.326000in}{0.100000in}}%
\pgfpathlineto{\pgfqpoint{2.326000in}{0.100000in}}%
\pgfpathclose%
\pgfusepath{clip}%
\pgfsetrectcap%
\pgfsetroundjoin%
\pgfsetlinewidth{0.451687pt}%
\definecolor{currentstroke}{rgb}{0.690196,0.690196,0.690196}%
\pgfsetstrokecolor{currentstroke}%
\pgfsetstrokeopacity{0.250000}%
\pgfsetdash{}{0pt}%
\pgfpathmoveto{\pgfqpoint{2.480031in}{0.151295in}}%
\pgfpathlineto{\pgfqpoint{2.492737in}{0.160652in}}%
\pgfpathlineto{\pgfqpoint{2.504920in}{0.170513in}}%
\pgfpathlineto{\pgfqpoint{2.516298in}{0.180558in}}%
\pgfpathlineto{\pgfqpoint{2.527164in}{0.190956in}}%
\pgfpathlineto{\pgfqpoint{2.537547in}{0.201676in}}%
\pgfpathlineto{\pgfqpoint{2.547473in}{0.212695in}}%
\pgfpathlineto{\pgfqpoint{2.556966in}{0.223994in}}%
\pgfpathlineto{\pgfqpoint{2.566048in}{0.235556in}}%
\pgfpathlineto{\pgfqpoint{2.574736in}{0.247365in}}%
\pgfpathlineto{\pgfqpoint{2.583047in}{0.259410in}}%
\pgfpathlineto{\pgfqpoint{2.590993in}{0.271677in}}%
\pgfpathlineto{\pgfqpoint{2.598587in}{0.284154in}}%
\pgfpathlineto{\pgfqpoint{2.605839in}{0.296831in}}%
\pgfpathlineto{\pgfqpoint{2.612757in}{0.309698in}}%
\pgfpathlineto{\pgfqpoint{2.619430in}{0.322907in}}%
\pgfpathlineto{\pgfqpoint{2.625678in}{0.336079in}}%
\pgfpathlineto{\pgfqpoint{2.631621in}{0.349425in}}%
\pgfpathlineto{\pgfqpoint{2.637261in}{0.362931in}}%
\pgfpathlineto{\pgfqpoint{2.642602in}{0.376589in}}%
\pgfpathlineto{\pgfqpoint{2.647647in}{0.390387in}}%
\pgfpathlineto{\pgfqpoint{2.652398in}{0.404317in}}%
\pgfpathlineto{\pgfqpoint{2.656858in}{0.418370in}}%
\pgfpathlineto{\pgfqpoint{2.661029in}{0.432537in}}%
\pgfpathlineto{\pgfqpoint{2.664913in}{0.446809in}}%
\pgfpathlineto{\pgfqpoint{2.668512in}{0.461179in}}%
\pgfpathlineto{\pgfqpoint{2.671828in}{0.475639in}}%
\pgfpathlineto{\pgfqpoint{2.674863in}{0.490181in}}%
\pgfpathlineto{\pgfqpoint{2.677617in}{0.504798in}}%
\pgfpathlineto{\pgfqpoint{2.680093in}{0.519482in}}%
\pgfpathlineto{\pgfqpoint{2.682291in}{0.534226in}}%
\pgfpathlineto{\pgfqpoint{2.684212in}{0.549023in}}%
\pgfpathlineto{\pgfqpoint{2.685875in}{0.564029in}}%
\pgfpathlineto{\pgfqpoint{2.687236in}{0.578836in}}%
\pgfpathlineto{\pgfqpoint{2.688326in}{0.593701in}}%
\pgfpathlineto{\pgfqpoint{2.689146in}{0.608611in}}%
\pgfpathlineto{\pgfqpoint{2.689692in}{0.623554in}}%
\pgfpathlineto{\pgfqpoint{2.689966in}{0.638516in}}%
\pgfpathlineto{\pgfqpoint{2.689966in}{0.653484in}}%
\pgfpathlineto{\pgfqpoint{2.689692in}{0.668446in}}%
\pgfpathlineto{\pgfqpoint{2.689146in}{0.683389in}}%
\pgfpathlineto{\pgfqpoint{2.688326in}{0.698299in}}%
\pgfpathlineto{\pgfqpoint{2.687236in}{0.713164in}}%
\pgfpathlineto{\pgfqpoint{2.685875in}{0.727971in}}%
\pgfpathlineto{\pgfqpoint{2.684212in}{0.742977in}}%
\pgfpathlineto{\pgfqpoint{2.682291in}{0.757774in}}%
\pgfpathlineto{\pgfqpoint{2.680093in}{0.772518in}}%
\pgfpathlineto{\pgfqpoint{2.677617in}{0.787202in}}%
\pgfpathlineto{\pgfqpoint{2.674863in}{0.801819in}}%
\pgfpathlineto{\pgfqpoint{2.671828in}{0.816361in}}%
\pgfpathlineto{\pgfqpoint{2.668512in}{0.830821in}}%
\pgfpathlineto{\pgfqpoint{2.664913in}{0.845191in}}%
\pgfpathlineto{\pgfqpoint{2.661029in}{0.859463in}}%
\pgfpathlineto{\pgfqpoint{2.656858in}{0.873630in}}%
\pgfpathlineto{\pgfqpoint{2.652398in}{0.887683in}}%
\pgfpathlineto{\pgfqpoint{2.647647in}{0.901613in}}%
\pgfpathlineto{\pgfqpoint{2.642602in}{0.915411in}}%
\pgfpathlineto{\pgfqpoint{2.637261in}{0.929069in}}%
\pgfpathlineto{\pgfqpoint{2.631621in}{0.942575in}}%
\pgfpathlineto{\pgfqpoint{2.625678in}{0.955921in}}%
\pgfpathlineto{\pgfqpoint{2.619430in}{0.969093in}}%
\pgfpathlineto{\pgfqpoint{2.612757in}{0.982302in}}%
\pgfpathlineto{\pgfqpoint{2.605839in}{0.995169in}}%
\pgfpathlineto{\pgfqpoint{2.598587in}{1.007846in}}%
\pgfpathlineto{\pgfqpoint{2.590993in}{1.020323in}}%
\pgfpathlineto{\pgfqpoint{2.583047in}{1.032590in}}%
\pgfpathlineto{\pgfqpoint{2.574736in}{1.044635in}}%
\pgfpathlineto{\pgfqpoint{2.566048in}{1.056444in}}%
\pgfpathlineto{\pgfqpoint{2.556966in}{1.068006in}}%
\pgfpathlineto{\pgfqpoint{2.547473in}{1.079305in}}%
\pgfpathlineto{\pgfqpoint{2.537547in}{1.090324in}}%
\pgfpathlineto{\pgfqpoint{2.527164in}{1.101044in}}%
\pgfpathlineto{\pgfqpoint{2.516298in}{1.111442in}}%
\pgfpathlineto{\pgfqpoint{2.504920in}{1.121487in}}%
\pgfpathlineto{\pgfqpoint{2.492737in}{1.131348in}}%
\pgfpathlineto{\pgfqpoint{2.480031in}{1.140705in}}%
\pgfusepath{stroke}%
\end{pgfscope}%
\begin{pgfscope}%
\pgfpathmoveto{\pgfqpoint{2.326000in}{0.100000in}}%
\pgfpathcurveto{\pgfqpoint{2.615602in}{0.100000in}}{\pgfqpoint{2.893381in}{0.157530in}}{\pgfqpoint{3.098161in}{0.259920in}}%
\pgfpathcurveto{\pgfqpoint{3.302940in}{0.362309in}}{\pgfqpoint{3.418000in}{0.501199in}}{\pgfqpoint{3.418000in}{0.646000in}}%
\pgfpathcurveto{\pgfqpoint{3.418000in}{0.790801in}}{\pgfqpoint{3.302940in}{0.929691in}}{\pgfqpoint{3.098161in}{1.032080in}}%
\pgfpathcurveto{\pgfqpoint{2.893381in}{1.134470in}}{\pgfqpoint{2.615602in}{1.192000in}}{\pgfqpoint{2.326000in}{1.192000in}}%
\pgfpathcurveto{\pgfqpoint{2.036398in}{1.192000in}}{\pgfqpoint{1.758619in}{1.134470in}}{\pgfqpoint{1.553839in}{1.032080in}}%
\pgfpathcurveto{\pgfqpoint{1.349060in}{0.929691in}}{\pgfqpoint{1.234000in}{0.790801in}}{\pgfqpoint{1.234000in}{0.646000in}}%
\pgfpathcurveto{\pgfqpoint{1.234000in}{0.501199in}}{\pgfqpoint{1.349060in}{0.362309in}}{\pgfqpoint{1.553839in}{0.259920in}}%
\pgfpathcurveto{\pgfqpoint{1.758619in}{0.157530in}}{\pgfqpoint{2.036398in}{0.100000in}}{\pgfqpoint{2.326000in}{0.100000in}}%
\pgfpathlineto{\pgfqpoint{2.326000in}{0.100000in}}%
\pgfpathclose%
\pgfusepath{clip}%
\pgfsetrectcap%
\pgfsetroundjoin%
\pgfsetlinewidth{0.451687pt}%
\definecolor{currentstroke}{rgb}{0.690196,0.690196,0.690196}%
\pgfsetstrokecolor{currentstroke}%
\pgfsetstrokeopacity{0.250000}%
\pgfsetdash{}{0pt}%
\pgfpathmoveto{\pgfqpoint{2.557047in}{0.151295in}}%
\pgfpathlineto{\pgfqpoint{2.576106in}{0.160652in}}%
\pgfpathlineto{\pgfqpoint{2.594380in}{0.170513in}}%
\pgfpathlineto{\pgfqpoint{2.611447in}{0.180558in}}%
\pgfpathlineto{\pgfqpoint{2.627746in}{0.190956in}}%
\pgfpathlineto{\pgfqpoint{2.643320in}{0.201676in}}%
\pgfpathlineto{\pgfqpoint{2.658209in}{0.212695in}}%
\pgfpathlineto{\pgfqpoint{2.672449in}{0.223994in}}%
\pgfpathlineto{\pgfqpoint{2.686071in}{0.235556in}}%
\pgfpathlineto{\pgfqpoint{2.699104in}{0.247365in}}%
\pgfpathlineto{\pgfqpoint{2.711570in}{0.259410in}}%
\pgfpathlineto{\pgfqpoint{2.723490in}{0.271677in}}%
\pgfpathlineto{\pgfqpoint{2.734881in}{0.284154in}}%
\pgfpathlineto{\pgfqpoint{2.745758in}{0.296831in}}%
\pgfpathlineto{\pgfqpoint{2.756136in}{0.309698in}}%
\pgfpathlineto{\pgfqpoint{2.766144in}{0.322907in}}%
\pgfpathlineto{\pgfqpoint{2.775517in}{0.336079in}}%
\pgfpathlineto{\pgfqpoint{2.784431in}{0.349425in}}%
\pgfpathlineto{\pgfqpoint{2.792892in}{0.362931in}}%
\pgfpathlineto{\pgfqpoint{2.800904in}{0.376589in}}%
\pgfpathlineto{\pgfqpoint{2.808471in}{0.390387in}}%
\pgfpathlineto{\pgfqpoint{2.815597in}{0.404317in}}%
\pgfpathlineto{\pgfqpoint{2.822287in}{0.418370in}}%
\pgfpathlineto{\pgfqpoint{2.828543in}{0.432537in}}%
\pgfpathlineto{\pgfqpoint{2.834369in}{0.446809in}}%
\pgfpathlineto{\pgfqpoint{2.839768in}{0.461179in}}%
\pgfpathlineto{\pgfqpoint{2.844742in}{0.475639in}}%
\pgfpathlineto{\pgfqpoint{2.849294in}{0.490181in}}%
\pgfpathlineto{\pgfqpoint{2.853426in}{0.504798in}}%
\pgfpathlineto{\pgfqpoint{2.857139in}{0.519482in}}%
\pgfpathlineto{\pgfqpoint{2.860437in}{0.534226in}}%
\pgfpathlineto{\pgfqpoint{2.863319in}{0.549023in}}%
\pgfpathlineto{\pgfqpoint{2.865812in}{0.564029in}}%
\pgfpathlineto{\pgfqpoint{2.867853in}{0.578836in}}%
\pgfpathlineto{\pgfqpoint{2.869490in}{0.593701in}}%
\pgfpathlineto{\pgfqpoint{2.870718in}{0.608611in}}%
\pgfpathlineto{\pgfqpoint{2.871538in}{0.623554in}}%
\pgfpathlineto{\pgfqpoint{2.871949in}{0.638516in}}%
\pgfpathlineto{\pgfqpoint{2.871949in}{0.653484in}}%
\pgfpathlineto{\pgfqpoint{2.871538in}{0.668446in}}%
\pgfpathlineto{\pgfqpoint{2.870718in}{0.683389in}}%
\pgfpathlineto{\pgfqpoint{2.869490in}{0.698299in}}%
\pgfpathlineto{\pgfqpoint{2.867853in}{0.713164in}}%
\pgfpathlineto{\pgfqpoint{2.865812in}{0.727971in}}%
\pgfpathlineto{\pgfqpoint{2.863319in}{0.742977in}}%
\pgfpathlineto{\pgfqpoint{2.860437in}{0.757774in}}%
\pgfpathlineto{\pgfqpoint{2.857139in}{0.772518in}}%
\pgfpathlineto{\pgfqpoint{2.853426in}{0.787202in}}%
\pgfpathlineto{\pgfqpoint{2.849294in}{0.801819in}}%
\pgfpathlineto{\pgfqpoint{2.844742in}{0.816361in}}%
\pgfpathlineto{\pgfqpoint{2.839768in}{0.830821in}}%
\pgfpathlineto{\pgfqpoint{2.834369in}{0.845191in}}%
\pgfpathlineto{\pgfqpoint{2.828543in}{0.859463in}}%
\pgfpathlineto{\pgfqpoint{2.822287in}{0.873630in}}%
\pgfpathlineto{\pgfqpoint{2.815597in}{0.887683in}}%
\pgfpathlineto{\pgfqpoint{2.808471in}{0.901613in}}%
\pgfpathlineto{\pgfqpoint{2.800904in}{0.915411in}}%
\pgfpathlineto{\pgfqpoint{2.792892in}{0.929069in}}%
\pgfpathlineto{\pgfqpoint{2.784431in}{0.942575in}}%
\pgfpathlineto{\pgfqpoint{2.775517in}{0.955921in}}%
\pgfpathlineto{\pgfqpoint{2.766144in}{0.969093in}}%
\pgfpathlineto{\pgfqpoint{2.756136in}{0.982302in}}%
\pgfpathlineto{\pgfqpoint{2.745758in}{0.995169in}}%
\pgfpathlineto{\pgfqpoint{2.734881in}{1.007846in}}%
\pgfpathlineto{\pgfqpoint{2.723490in}{1.020323in}}%
\pgfpathlineto{\pgfqpoint{2.711570in}{1.032590in}}%
\pgfpathlineto{\pgfqpoint{2.699104in}{1.044635in}}%
\pgfpathlineto{\pgfqpoint{2.686071in}{1.056444in}}%
\pgfpathlineto{\pgfqpoint{2.672449in}{1.068006in}}%
\pgfpathlineto{\pgfqpoint{2.658209in}{1.079305in}}%
\pgfpathlineto{\pgfqpoint{2.643320in}{1.090324in}}%
\pgfpathlineto{\pgfqpoint{2.627746in}{1.101044in}}%
\pgfpathlineto{\pgfqpoint{2.611447in}{1.111442in}}%
\pgfpathlineto{\pgfqpoint{2.594380in}{1.121487in}}%
\pgfpathlineto{\pgfqpoint{2.576106in}{1.131348in}}%
\pgfpathlineto{\pgfqpoint{2.557047in}{1.140705in}}%
\pgfusepath{stroke}%
\end{pgfscope}%
\begin{pgfscope}%
\pgfpathmoveto{\pgfqpoint{2.326000in}{0.100000in}}%
\pgfpathcurveto{\pgfqpoint{2.615602in}{0.100000in}}{\pgfqpoint{2.893381in}{0.157530in}}{\pgfqpoint{3.098161in}{0.259920in}}%
\pgfpathcurveto{\pgfqpoint{3.302940in}{0.362309in}}{\pgfqpoint{3.418000in}{0.501199in}}{\pgfqpoint{3.418000in}{0.646000in}}%
\pgfpathcurveto{\pgfqpoint{3.418000in}{0.790801in}}{\pgfqpoint{3.302940in}{0.929691in}}{\pgfqpoint{3.098161in}{1.032080in}}%
\pgfpathcurveto{\pgfqpoint{2.893381in}{1.134470in}}{\pgfqpoint{2.615602in}{1.192000in}}{\pgfqpoint{2.326000in}{1.192000in}}%
\pgfpathcurveto{\pgfqpoint{2.036398in}{1.192000in}}{\pgfqpoint{1.758619in}{1.134470in}}{\pgfqpoint{1.553839in}{1.032080in}}%
\pgfpathcurveto{\pgfqpoint{1.349060in}{0.929691in}}{\pgfqpoint{1.234000in}{0.790801in}}{\pgfqpoint{1.234000in}{0.646000in}}%
\pgfpathcurveto{\pgfqpoint{1.234000in}{0.501199in}}{\pgfqpoint{1.349060in}{0.362309in}}{\pgfqpoint{1.553839in}{0.259920in}}%
\pgfpathcurveto{\pgfqpoint{1.758619in}{0.157530in}}{\pgfqpoint{2.036398in}{0.100000in}}{\pgfqpoint{2.326000in}{0.100000in}}%
\pgfpathlineto{\pgfqpoint{2.326000in}{0.100000in}}%
\pgfpathclose%
\pgfusepath{clip}%
\pgfsetrectcap%
\pgfsetroundjoin%
\pgfsetlinewidth{0.451687pt}%
\definecolor{currentstroke}{rgb}{0.690196,0.690196,0.690196}%
\pgfsetstrokecolor{currentstroke}%
\pgfsetstrokeopacity{0.250000}%
\pgfsetdash{}{0pt}%
\pgfpathmoveto{\pgfqpoint{2.634063in}{0.151295in}}%
\pgfpathlineto{\pgfqpoint{2.659475in}{0.160652in}}%
\pgfpathlineto{\pgfqpoint{2.683840in}{0.170513in}}%
\pgfpathlineto{\pgfqpoint{2.706596in}{0.180558in}}%
\pgfpathlineto{\pgfqpoint{2.728328in}{0.190956in}}%
\pgfpathlineto{\pgfqpoint{2.749093in}{0.201676in}}%
\pgfpathlineto{\pgfqpoint{2.768945in}{0.212695in}}%
\pgfpathlineto{\pgfqpoint{2.787932in}{0.223994in}}%
\pgfpathlineto{\pgfqpoint{2.806095in}{0.235556in}}%
\pgfpathlineto{\pgfqpoint{2.823472in}{0.247365in}}%
\pgfpathlineto{\pgfqpoint{2.840093in}{0.259410in}}%
\pgfpathlineto{\pgfqpoint{2.855986in}{0.271677in}}%
\pgfpathlineto{\pgfqpoint{2.871174in}{0.284154in}}%
\pgfpathlineto{\pgfqpoint{2.885678in}{0.296831in}}%
\pgfpathlineto{\pgfqpoint{2.899515in}{0.309698in}}%
\pgfpathlineto{\pgfqpoint{2.912859in}{0.322907in}}%
\pgfpathlineto{\pgfqpoint{2.925356in}{0.336079in}}%
\pgfpathlineto{\pgfqpoint{2.937241in}{0.349425in}}%
\pgfpathlineto{\pgfqpoint{2.948523in}{0.362931in}}%
\pgfpathlineto{\pgfqpoint{2.959205in}{0.376589in}}%
\pgfpathlineto{\pgfqpoint{2.969295in}{0.390387in}}%
\pgfpathlineto{\pgfqpoint{2.978796in}{0.404317in}}%
\pgfpathlineto{\pgfqpoint{2.987716in}{0.418370in}}%
\pgfpathlineto{\pgfqpoint{2.996057in}{0.432537in}}%
\pgfpathlineto{\pgfqpoint{3.003825in}{0.446809in}}%
\pgfpathlineto{\pgfqpoint{3.011024in}{0.461179in}}%
\pgfpathlineto{\pgfqpoint{3.017656in}{0.475639in}}%
\pgfpathlineto{\pgfqpoint{3.023725in}{0.490181in}}%
\pgfpathlineto{\pgfqpoint{3.029234in}{0.504798in}}%
\pgfpathlineto{\pgfqpoint{3.034186in}{0.519482in}}%
\pgfpathlineto{\pgfqpoint{3.038582in}{0.534226in}}%
\pgfpathlineto{\pgfqpoint{3.042425in}{0.549023in}}%
\pgfpathlineto{\pgfqpoint{3.045749in}{0.564029in}}%
\pgfpathlineto{\pgfqpoint{3.048471in}{0.578836in}}%
\pgfpathlineto{\pgfqpoint{3.050653in}{0.593701in}}%
\pgfpathlineto{\pgfqpoint{3.052291in}{0.608611in}}%
\pgfpathlineto{\pgfqpoint{3.053385in}{0.623554in}}%
\pgfpathlineto{\pgfqpoint{3.053932in}{0.638516in}}%
\pgfpathlineto{\pgfqpoint{3.053932in}{0.653484in}}%
\pgfpathlineto{\pgfqpoint{3.053385in}{0.668446in}}%
\pgfpathlineto{\pgfqpoint{3.052291in}{0.683389in}}%
\pgfpathlineto{\pgfqpoint{3.050653in}{0.698299in}}%
\pgfpathlineto{\pgfqpoint{3.048471in}{0.713164in}}%
\pgfpathlineto{\pgfqpoint{3.045749in}{0.727971in}}%
\pgfpathlineto{\pgfqpoint{3.042425in}{0.742977in}}%
\pgfpathlineto{\pgfqpoint{3.038582in}{0.757774in}}%
\pgfpathlineto{\pgfqpoint{3.034186in}{0.772518in}}%
\pgfpathlineto{\pgfqpoint{3.029234in}{0.787202in}}%
\pgfpathlineto{\pgfqpoint{3.023725in}{0.801819in}}%
\pgfpathlineto{\pgfqpoint{3.017656in}{0.816361in}}%
\pgfpathlineto{\pgfqpoint{3.011024in}{0.830821in}}%
\pgfpathlineto{\pgfqpoint{3.003825in}{0.845191in}}%
\pgfpathlineto{\pgfqpoint{2.996057in}{0.859463in}}%
\pgfpathlineto{\pgfqpoint{2.987716in}{0.873630in}}%
\pgfpathlineto{\pgfqpoint{2.978796in}{0.887683in}}%
\pgfpathlineto{\pgfqpoint{2.969295in}{0.901613in}}%
\pgfpathlineto{\pgfqpoint{2.959205in}{0.915411in}}%
\pgfpathlineto{\pgfqpoint{2.948523in}{0.929069in}}%
\pgfpathlineto{\pgfqpoint{2.937241in}{0.942575in}}%
\pgfpathlineto{\pgfqpoint{2.925356in}{0.955921in}}%
\pgfpathlineto{\pgfqpoint{2.912859in}{0.969093in}}%
\pgfpathlineto{\pgfqpoint{2.899515in}{0.982302in}}%
\pgfpathlineto{\pgfqpoint{2.885678in}{0.995169in}}%
\pgfpathlineto{\pgfqpoint{2.871174in}{1.007846in}}%
\pgfpathlineto{\pgfqpoint{2.855986in}{1.020323in}}%
\pgfpathlineto{\pgfqpoint{2.840093in}{1.032590in}}%
\pgfpathlineto{\pgfqpoint{2.823472in}{1.044635in}}%
\pgfpathlineto{\pgfqpoint{2.806095in}{1.056444in}}%
\pgfpathlineto{\pgfqpoint{2.787932in}{1.068006in}}%
\pgfpathlineto{\pgfqpoint{2.768945in}{1.079305in}}%
\pgfpathlineto{\pgfqpoint{2.749093in}{1.090324in}}%
\pgfpathlineto{\pgfqpoint{2.728328in}{1.101044in}}%
\pgfpathlineto{\pgfqpoint{2.706596in}{1.111442in}}%
\pgfpathlineto{\pgfqpoint{2.683840in}{1.121487in}}%
\pgfpathlineto{\pgfqpoint{2.659475in}{1.131348in}}%
\pgfpathlineto{\pgfqpoint{2.634063in}{1.140705in}}%
\pgfusepath{stroke}%
\end{pgfscope}%
\begin{pgfscope}%
\pgfpathmoveto{\pgfqpoint{2.326000in}{0.100000in}}%
\pgfpathcurveto{\pgfqpoint{2.615602in}{0.100000in}}{\pgfqpoint{2.893381in}{0.157530in}}{\pgfqpoint{3.098161in}{0.259920in}}%
\pgfpathcurveto{\pgfqpoint{3.302940in}{0.362309in}}{\pgfqpoint{3.418000in}{0.501199in}}{\pgfqpoint{3.418000in}{0.646000in}}%
\pgfpathcurveto{\pgfqpoint{3.418000in}{0.790801in}}{\pgfqpoint{3.302940in}{0.929691in}}{\pgfqpoint{3.098161in}{1.032080in}}%
\pgfpathcurveto{\pgfqpoint{2.893381in}{1.134470in}}{\pgfqpoint{2.615602in}{1.192000in}}{\pgfqpoint{2.326000in}{1.192000in}}%
\pgfpathcurveto{\pgfqpoint{2.036398in}{1.192000in}}{\pgfqpoint{1.758619in}{1.134470in}}{\pgfqpoint{1.553839in}{1.032080in}}%
\pgfpathcurveto{\pgfqpoint{1.349060in}{0.929691in}}{\pgfqpoint{1.234000in}{0.790801in}}{\pgfqpoint{1.234000in}{0.646000in}}%
\pgfpathcurveto{\pgfqpoint{1.234000in}{0.501199in}}{\pgfqpoint{1.349060in}{0.362309in}}{\pgfqpoint{1.553839in}{0.259920in}}%
\pgfpathcurveto{\pgfqpoint{1.758619in}{0.157530in}}{\pgfqpoint{2.036398in}{0.100000in}}{\pgfqpoint{2.326000in}{0.100000in}}%
\pgfpathlineto{\pgfqpoint{2.326000in}{0.100000in}}%
\pgfpathclose%
\pgfusepath{clip}%
\pgfsetrectcap%
\pgfsetroundjoin%
\pgfsetlinewidth{0.451687pt}%
\definecolor{currentstroke}{rgb}{0.690196,0.690196,0.690196}%
\pgfsetstrokecolor{currentstroke}%
\pgfsetstrokeopacity{0.250000}%
\pgfsetdash{}{0pt}%
\pgfpathmoveto{\pgfqpoint{2.711079in}{0.151295in}}%
\pgfpathlineto{\pgfqpoint{2.742843in}{0.160652in}}%
\pgfpathlineto{\pgfqpoint{2.773300in}{0.170513in}}%
\pgfpathlineto{\pgfqpoint{2.801745in}{0.180558in}}%
\pgfpathlineto{\pgfqpoint{2.828910in}{0.190956in}}%
\pgfpathlineto{\pgfqpoint{2.854866in}{0.201676in}}%
\pgfpathlineto{\pgfqpoint{2.879681in}{0.212695in}}%
\pgfpathlineto{\pgfqpoint{2.903415in}{0.223994in}}%
\pgfpathlineto{\pgfqpoint{2.926119in}{0.235556in}}%
\pgfpathlineto{\pgfqpoint{2.947840in}{0.247365in}}%
\pgfpathlineto{\pgfqpoint{2.968616in}{0.259410in}}%
\pgfpathlineto{\pgfqpoint{2.988483in}{0.271677in}}%
\pgfpathlineto{\pgfqpoint{3.007468in}{0.284154in}}%
\pgfpathlineto{\pgfqpoint{3.025597in}{0.296831in}}%
\pgfpathlineto{\pgfqpoint{3.042894in}{0.309698in}}%
\pgfpathlineto{\pgfqpoint{3.059574in}{0.322907in}}%
\pgfpathlineto{\pgfqpoint{3.075194in}{0.336079in}}%
\pgfpathlineto{\pgfqpoint{3.090052in}{0.349425in}}%
\pgfpathlineto{\pgfqpoint{3.104153in}{0.362931in}}%
\pgfpathlineto{\pgfqpoint{3.117506in}{0.376589in}}%
\pgfpathlineto{\pgfqpoint{3.130118in}{0.390387in}}%
\pgfpathlineto{\pgfqpoint{3.141996in}{0.404317in}}%
\pgfpathlineto{\pgfqpoint{3.153145in}{0.418370in}}%
\pgfpathlineto{\pgfqpoint{3.163572in}{0.432537in}}%
\pgfpathlineto{\pgfqpoint{3.173282in}{0.446809in}}%
\pgfpathlineto{\pgfqpoint{3.182280in}{0.461179in}}%
\pgfpathlineto{\pgfqpoint{3.190570in}{0.475639in}}%
\pgfpathlineto{\pgfqpoint{3.198156in}{0.490181in}}%
\pgfpathlineto{\pgfqpoint{3.205043in}{0.504798in}}%
\pgfpathlineto{\pgfqpoint{3.211232in}{0.519482in}}%
\pgfpathlineto{\pgfqpoint{3.216728in}{0.534226in}}%
\pgfpathlineto{\pgfqpoint{3.221531in}{0.549023in}}%
\pgfpathlineto{\pgfqpoint{3.225686in}{0.564029in}}%
\pgfpathlineto{\pgfqpoint{3.229089in}{0.578836in}}%
\pgfpathlineto{\pgfqpoint{3.231816in}{0.593701in}}%
\pgfpathlineto{\pgfqpoint{3.233864in}{0.608611in}}%
\pgfpathlineto{\pgfqpoint{3.235231in}{0.623554in}}%
\pgfpathlineto{\pgfqpoint{3.235915in}{0.638516in}}%
\pgfpathlineto{\pgfqpoint{3.235915in}{0.653484in}}%
\pgfpathlineto{\pgfqpoint{3.235231in}{0.668446in}}%
\pgfpathlineto{\pgfqpoint{3.233864in}{0.683389in}}%
\pgfpathlineto{\pgfqpoint{3.231816in}{0.698299in}}%
\pgfpathlineto{\pgfqpoint{3.229089in}{0.713164in}}%
\pgfpathlineto{\pgfqpoint{3.225686in}{0.727971in}}%
\pgfpathlineto{\pgfqpoint{3.221531in}{0.742977in}}%
\pgfpathlineto{\pgfqpoint{3.216728in}{0.757774in}}%
\pgfpathlineto{\pgfqpoint{3.211232in}{0.772518in}}%
\pgfpathlineto{\pgfqpoint{3.205043in}{0.787202in}}%
\pgfpathlineto{\pgfqpoint{3.198156in}{0.801819in}}%
\pgfpathlineto{\pgfqpoint{3.190570in}{0.816361in}}%
\pgfpathlineto{\pgfqpoint{3.182280in}{0.830821in}}%
\pgfpathlineto{\pgfqpoint{3.173282in}{0.845191in}}%
\pgfpathlineto{\pgfqpoint{3.163572in}{0.859463in}}%
\pgfpathlineto{\pgfqpoint{3.153145in}{0.873630in}}%
\pgfpathlineto{\pgfqpoint{3.141996in}{0.887683in}}%
\pgfpathlineto{\pgfqpoint{3.130118in}{0.901613in}}%
\pgfpathlineto{\pgfqpoint{3.117506in}{0.915411in}}%
\pgfpathlineto{\pgfqpoint{3.104153in}{0.929069in}}%
\pgfpathlineto{\pgfqpoint{3.090052in}{0.942575in}}%
\pgfpathlineto{\pgfqpoint{3.075194in}{0.955921in}}%
\pgfpathlineto{\pgfqpoint{3.059574in}{0.969093in}}%
\pgfpathlineto{\pgfqpoint{3.042894in}{0.982302in}}%
\pgfpathlineto{\pgfqpoint{3.025597in}{0.995169in}}%
\pgfpathlineto{\pgfqpoint{3.007468in}{1.007846in}}%
\pgfpathlineto{\pgfqpoint{2.988483in}{1.020323in}}%
\pgfpathlineto{\pgfqpoint{2.968616in}{1.032590in}}%
\pgfpathlineto{\pgfqpoint{2.947840in}{1.044635in}}%
\pgfpathlineto{\pgfqpoint{2.926119in}{1.056444in}}%
\pgfpathlineto{\pgfqpoint{2.903415in}{1.068006in}}%
\pgfpathlineto{\pgfqpoint{2.879681in}{1.079305in}}%
\pgfpathlineto{\pgfqpoint{2.854866in}{1.090324in}}%
\pgfpathlineto{\pgfqpoint{2.828910in}{1.101044in}}%
\pgfpathlineto{\pgfqpoint{2.801745in}{1.111442in}}%
\pgfpathlineto{\pgfqpoint{2.773300in}{1.121487in}}%
\pgfpathlineto{\pgfqpoint{2.742843in}{1.131348in}}%
\pgfpathlineto{\pgfqpoint{2.711079in}{1.140705in}}%
\pgfusepath{stroke}%
\end{pgfscope}%
\begin{pgfscope}%
\pgfpathmoveto{\pgfqpoint{2.326000in}{0.100000in}}%
\pgfpathcurveto{\pgfqpoint{2.615602in}{0.100000in}}{\pgfqpoint{2.893381in}{0.157530in}}{\pgfqpoint{3.098161in}{0.259920in}}%
\pgfpathcurveto{\pgfqpoint{3.302940in}{0.362309in}}{\pgfqpoint{3.418000in}{0.501199in}}{\pgfqpoint{3.418000in}{0.646000in}}%
\pgfpathcurveto{\pgfqpoint{3.418000in}{0.790801in}}{\pgfqpoint{3.302940in}{0.929691in}}{\pgfqpoint{3.098161in}{1.032080in}}%
\pgfpathcurveto{\pgfqpoint{2.893381in}{1.134470in}}{\pgfqpoint{2.615602in}{1.192000in}}{\pgfqpoint{2.326000in}{1.192000in}}%
\pgfpathcurveto{\pgfqpoint{2.036398in}{1.192000in}}{\pgfqpoint{1.758619in}{1.134470in}}{\pgfqpoint{1.553839in}{1.032080in}}%
\pgfpathcurveto{\pgfqpoint{1.349060in}{0.929691in}}{\pgfqpoint{1.234000in}{0.790801in}}{\pgfqpoint{1.234000in}{0.646000in}}%
\pgfpathcurveto{\pgfqpoint{1.234000in}{0.501199in}}{\pgfqpoint{1.349060in}{0.362309in}}{\pgfqpoint{1.553839in}{0.259920in}}%
\pgfpathcurveto{\pgfqpoint{1.758619in}{0.157530in}}{\pgfqpoint{2.036398in}{0.100000in}}{\pgfqpoint{2.326000in}{0.100000in}}%
\pgfpathlineto{\pgfqpoint{2.326000in}{0.100000in}}%
\pgfpathclose%
\pgfusepath{clip}%
\pgfsetrectcap%
\pgfsetroundjoin%
\pgfsetlinewidth{0.451687pt}%
\definecolor{currentstroke}{rgb}{0.690196,0.690196,0.690196}%
\pgfsetstrokecolor{currentstroke}%
\pgfsetstrokeopacity{0.250000}%
\pgfsetdash{}{0pt}%
\pgfpathmoveto{\pgfqpoint{1.863906in}{0.151295in}}%
\pgfpathlineto{\pgfqpoint{1.876228in}{0.151295in}}%
\pgfpathlineto{\pgfqpoint{1.888551in}{0.151295in}}%
\pgfpathlineto{\pgfqpoint{1.900873in}{0.151295in}}%
\pgfpathlineto{\pgfqpoint{1.913196in}{0.151295in}}%
\pgfpathlineto{\pgfqpoint{1.925518in}{0.151295in}}%
\pgfpathlineto{\pgfqpoint{1.937841in}{0.151295in}}%
\pgfpathlineto{\pgfqpoint{1.950163in}{0.151295in}}%
\pgfpathlineto{\pgfqpoint{1.962486in}{0.151295in}}%
\pgfpathlineto{\pgfqpoint{1.974808in}{0.151295in}}%
\pgfpathlineto{\pgfqpoint{1.987131in}{0.151295in}}%
\pgfpathlineto{\pgfqpoint{1.999453in}{0.151295in}}%
\pgfpathlineto{\pgfqpoint{2.011776in}{0.151295in}}%
\pgfpathlineto{\pgfqpoint{2.024098in}{0.151295in}}%
\pgfpathlineto{\pgfqpoint{2.036421in}{0.151295in}}%
\pgfpathlineto{\pgfqpoint{2.048743in}{0.151295in}}%
\pgfpathlineto{\pgfqpoint{2.061066in}{0.151295in}}%
\pgfpathlineto{\pgfqpoint{2.073388in}{0.151295in}}%
\pgfpathlineto{\pgfqpoint{2.085711in}{0.151295in}}%
\pgfpathlineto{\pgfqpoint{2.098033in}{0.151295in}}%
\pgfpathlineto{\pgfqpoint{2.110356in}{0.151295in}}%
\pgfpathlineto{\pgfqpoint{2.122678in}{0.151295in}}%
\pgfpathlineto{\pgfqpoint{2.135001in}{0.151295in}}%
\pgfpathlineto{\pgfqpoint{2.147323in}{0.151295in}}%
\pgfpathlineto{\pgfqpoint{2.159646in}{0.151295in}}%
\pgfpathlineto{\pgfqpoint{2.171969in}{0.151295in}}%
\pgfpathlineto{\pgfqpoint{2.184291in}{0.151295in}}%
\pgfpathlineto{\pgfqpoint{2.196614in}{0.151295in}}%
\pgfpathlineto{\pgfqpoint{2.208936in}{0.151295in}}%
\pgfpathlineto{\pgfqpoint{2.221259in}{0.151295in}}%
\pgfpathlineto{\pgfqpoint{2.233581in}{0.151295in}}%
\pgfpathlineto{\pgfqpoint{2.245904in}{0.151295in}}%
\pgfpathlineto{\pgfqpoint{2.258226in}{0.151295in}}%
\pgfpathlineto{\pgfqpoint{2.270549in}{0.151295in}}%
\pgfpathlineto{\pgfqpoint{2.282871in}{0.151295in}}%
\pgfpathlineto{\pgfqpoint{2.295194in}{0.151295in}}%
\pgfpathlineto{\pgfqpoint{2.307516in}{0.151295in}}%
\pgfpathlineto{\pgfqpoint{2.319839in}{0.151295in}}%
\pgfpathlineto{\pgfqpoint{2.332161in}{0.151295in}}%
\pgfpathlineto{\pgfqpoint{2.344484in}{0.151295in}}%
\pgfpathlineto{\pgfqpoint{2.356806in}{0.151295in}}%
\pgfpathlineto{\pgfqpoint{2.369129in}{0.151295in}}%
\pgfpathlineto{\pgfqpoint{2.381451in}{0.151295in}}%
\pgfpathlineto{\pgfqpoint{2.393774in}{0.151295in}}%
\pgfpathlineto{\pgfqpoint{2.406096in}{0.151295in}}%
\pgfpathlineto{\pgfqpoint{2.418419in}{0.151295in}}%
\pgfpathlineto{\pgfqpoint{2.430741in}{0.151295in}}%
\pgfpathlineto{\pgfqpoint{2.443064in}{0.151295in}}%
\pgfpathlineto{\pgfqpoint{2.455386in}{0.151295in}}%
\pgfpathlineto{\pgfqpoint{2.467709in}{0.151295in}}%
\pgfpathlineto{\pgfqpoint{2.480031in}{0.151295in}}%
\pgfpathlineto{\pgfqpoint{2.492354in}{0.151295in}}%
\pgfpathlineto{\pgfqpoint{2.504677in}{0.151295in}}%
\pgfpathlineto{\pgfqpoint{2.516999in}{0.151295in}}%
\pgfpathlineto{\pgfqpoint{2.529322in}{0.151295in}}%
\pgfpathlineto{\pgfqpoint{2.541644in}{0.151295in}}%
\pgfpathlineto{\pgfqpoint{2.553967in}{0.151295in}}%
\pgfpathlineto{\pgfqpoint{2.566289in}{0.151295in}}%
\pgfpathlineto{\pgfqpoint{2.578612in}{0.151295in}}%
\pgfpathlineto{\pgfqpoint{2.590934in}{0.151295in}}%
\pgfpathlineto{\pgfqpoint{2.603257in}{0.151295in}}%
\pgfpathlineto{\pgfqpoint{2.615579in}{0.151295in}}%
\pgfpathlineto{\pgfqpoint{2.627902in}{0.151295in}}%
\pgfpathlineto{\pgfqpoint{2.640224in}{0.151295in}}%
\pgfpathlineto{\pgfqpoint{2.652547in}{0.151295in}}%
\pgfpathlineto{\pgfqpoint{2.664869in}{0.151295in}}%
\pgfpathlineto{\pgfqpoint{2.677192in}{0.151295in}}%
\pgfpathlineto{\pgfqpoint{2.689514in}{0.151295in}}%
\pgfpathlineto{\pgfqpoint{2.701837in}{0.151295in}}%
\pgfpathlineto{\pgfqpoint{2.714159in}{0.151295in}}%
\pgfpathlineto{\pgfqpoint{2.726482in}{0.151295in}}%
\pgfpathlineto{\pgfqpoint{2.738804in}{0.151295in}}%
\pgfpathlineto{\pgfqpoint{2.751127in}{0.151295in}}%
\pgfpathlineto{\pgfqpoint{2.763449in}{0.151295in}}%
\pgfpathlineto{\pgfqpoint{2.775772in}{0.151295in}}%
\pgfpathlineto{\pgfqpoint{2.788094in}{0.151295in}}%
\pgfusepath{stroke}%
\end{pgfscope}%
\begin{pgfscope}%
\pgfpathmoveto{\pgfqpoint{2.326000in}{0.100000in}}%
\pgfpathcurveto{\pgfqpoint{2.615602in}{0.100000in}}{\pgfqpoint{2.893381in}{0.157530in}}{\pgfqpoint{3.098161in}{0.259920in}}%
\pgfpathcurveto{\pgfqpoint{3.302940in}{0.362309in}}{\pgfqpoint{3.418000in}{0.501199in}}{\pgfqpoint{3.418000in}{0.646000in}}%
\pgfpathcurveto{\pgfqpoint{3.418000in}{0.790801in}}{\pgfqpoint{3.302940in}{0.929691in}}{\pgfqpoint{3.098161in}{1.032080in}}%
\pgfpathcurveto{\pgfqpoint{2.893381in}{1.134470in}}{\pgfqpoint{2.615602in}{1.192000in}}{\pgfqpoint{2.326000in}{1.192000in}}%
\pgfpathcurveto{\pgfqpoint{2.036398in}{1.192000in}}{\pgfqpoint{1.758619in}{1.134470in}}{\pgfqpoint{1.553839in}{1.032080in}}%
\pgfpathcurveto{\pgfqpoint{1.349060in}{0.929691in}}{\pgfqpoint{1.234000in}{0.790801in}}{\pgfqpoint{1.234000in}{0.646000in}}%
\pgfpathcurveto{\pgfqpoint{1.234000in}{0.501199in}}{\pgfqpoint{1.349060in}{0.362309in}}{\pgfqpoint{1.553839in}{0.259920in}}%
\pgfpathcurveto{\pgfqpoint{1.758619in}{0.157530in}}{\pgfqpoint{2.036398in}{0.100000in}}{\pgfqpoint{2.326000in}{0.100000in}}%
\pgfpathlineto{\pgfqpoint{2.326000in}{0.100000in}}%
\pgfpathclose%
\pgfusepath{clip}%
\pgfsetrectcap%
\pgfsetroundjoin%
\pgfsetlinewidth{0.451687pt}%
\definecolor{currentstroke}{rgb}{0.690196,0.690196,0.690196}%
\pgfsetstrokecolor{currentstroke}%
\pgfsetstrokeopacity{0.250000}%
\pgfsetdash{}{0pt}%
\pgfpathmoveto{\pgfqpoint{1.619329in}{0.229743in}}%
\pgfpathlineto{\pgfqpoint{1.638174in}{0.229743in}}%
\pgfpathlineto{\pgfqpoint{1.657018in}{0.229743in}}%
\pgfpathlineto{\pgfqpoint{1.675863in}{0.229743in}}%
\pgfpathlineto{\pgfqpoint{1.694707in}{0.229743in}}%
\pgfpathlineto{\pgfqpoint{1.713552in}{0.229743in}}%
\pgfpathlineto{\pgfqpoint{1.732396in}{0.229743in}}%
\pgfpathlineto{\pgfqpoint{1.751241in}{0.229743in}}%
\pgfpathlineto{\pgfqpoint{1.770086in}{0.229743in}}%
\pgfpathlineto{\pgfqpoint{1.788930in}{0.229743in}}%
\pgfpathlineto{\pgfqpoint{1.807775in}{0.229743in}}%
\pgfpathlineto{\pgfqpoint{1.826619in}{0.229743in}}%
\pgfpathlineto{\pgfqpoint{1.845464in}{0.229743in}}%
\pgfpathlineto{\pgfqpoint{1.864308in}{0.229743in}}%
\pgfpathlineto{\pgfqpoint{1.883153in}{0.229743in}}%
\pgfpathlineto{\pgfqpoint{1.901997in}{0.229743in}}%
\pgfpathlineto{\pgfqpoint{1.920842in}{0.229743in}}%
\pgfpathlineto{\pgfqpoint{1.939687in}{0.229743in}}%
\pgfpathlineto{\pgfqpoint{1.958531in}{0.229743in}}%
\pgfpathlineto{\pgfqpoint{1.977376in}{0.229743in}}%
\pgfpathlineto{\pgfqpoint{1.996220in}{0.229743in}}%
\pgfpathlineto{\pgfqpoint{2.015065in}{0.229743in}}%
\pgfpathlineto{\pgfqpoint{2.033909in}{0.229743in}}%
\pgfpathlineto{\pgfqpoint{2.052754in}{0.229743in}}%
\pgfpathlineto{\pgfqpoint{2.071598in}{0.229743in}}%
\pgfpathlineto{\pgfqpoint{2.090443in}{0.229743in}}%
\pgfpathlineto{\pgfqpoint{2.109288in}{0.229743in}}%
\pgfpathlineto{\pgfqpoint{2.128132in}{0.229743in}}%
\pgfpathlineto{\pgfqpoint{2.146977in}{0.229743in}}%
\pgfpathlineto{\pgfqpoint{2.165821in}{0.229743in}}%
\pgfpathlineto{\pgfqpoint{2.184666in}{0.229743in}}%
\pgfpathlineto{\pgfqpoint{2.203510in}{0.229743in}}%
\pgfpathlineto{\pgfqpoint{2.222355in}{0.229743in}}%
\pgfpathlineto{\pgfqpoint{2.241199in}{0.229743in}}%
\pgfpathlineto{\pgfqpoint{2.260044in}{0.229743in}}%
\pgfpathlineto{\pgfqpoint{2.278889in}{0.229743in}}%
\pgfpathlineto{\pgfqpoint{2.297733in}{0.229743in}}%
\pgfpathlineto{\pgfqpoint{2.316578in}{0.229743in}}%
\pgfpathlineto{\pgfqpoint{2.335422in}{0.229743in}}%
\pgfpathlineto{\pgfqpoint{2.354267in}{0.229743in}}%
\pgfpathlineto{\pgfqpoint{2.373111in}{0.229743in}}%
\pgfpathlineto{\pgfqpoint{2.391956in}{0.229743in}}%
\pgfpathlineto{\pgfqpoint{2.410801in}{0.229743in}}%
\pgfpathlineto{\pgfqpoint{2.429645in}{0.229743in}}%
\pgfpathlineto{\pgfqpoint{2.448490in}{0.229743in}}%
\pgfpathlineto{\pgfqpoint{2.467334in}{0.229743in}}%
\pgfpathlineto{\pgfqpoint{2.486179in}{0.229743in}}%
\pgfpathlineto{\pgfqpoint{2.505023in}{0.229743in}}%
\pgfpathlineto{\pgfqpoint{2.523868in}{0.229743in}}%
\pgfpathlineto{\pgfqpoint{2.542712in}{0.229743in}}%
\pgfpathlineto{\pgfqpoint{2.561557in}{0.229743in}}%
\pgfpathlineto{\pgfqpoint{2.580402in}{0.229743in}}%
\pgfpathlineto{\pgfqpoint{2.599246in}{0.229743in}}%
\pgfpathlineto{\pgfqpoint{2.618091in}{0.229743in}}%
\pgfpathlineto{\pgfqpoint{2.636935in}{0.229743in}}%
\pgfpathlineto{\pgfqpoint{2.655780in}{0.229743in}}%
\pgfpathlineto{\pgfqpoint{2.674624in}{0.229743in}}%
\pgfpathlineto{\pgfqpoint{2.693469in}{0.229743in}}%
\pgfpathlineto{\pgfqpoint{2.712313in}{0.229743in}}%
\pgfpathlineto{\pgfqpoint{2.731158in}{0.229743in}}%
\pgfpathlineto{\pgfqpoint{2.750003in}{0.229743in}}%
\pgfpathlineto{\pgfqpoint{2.768847in}{0.229743in}}%
\pgfpathlineto{\pgfqpoint{2.787692in}{0.229743in}}%
\pgfpathlineto{\pgfqpoint{2.806536in}{0.229743in}}%
\pgfpathlineto{\pgfqpoint{2.825381in}{0.229743in}}%
\pgfpathlineto{\pgfqpoint{2.844225in}{0.229743in}}%
\pgfpathlineto{\pgfqpoint{2.863070in}{0.229743in}}%
\pgfpathlineto{\pgfqpoint{2.881914in}{0.229743in}}%
\pgfpathlineto{\pgfqpoint{2.900759in}{0.229743in}}%
\pgfpathlineto{\pgfqpoint{2.919604in}{0.229743in}}%
\pgfpathlineto{\pgfqpoint{2.938448in}{0.229743in}}%
\pgfpathlineto{\pgfqpoint{2.957293in}{0.229743in}}%
\pgfpathlineto{\pgfqpoint{2.976137in}{0.229743in}}%
\pgfpathlineto{\pgfqpoint{2.994982in}{0.229743in}}%
\pgfpathlineto{\pgfqpoint{3.013826in}{0.229743in}}%
\pgfpathlineto{\pgfqpoint{3.032671in}{0.229743in}}%
\pgfusepath{stroke}%
\end{pgfscope}%
\begin{pgfscope}%
\pgfpathmoveto{\pgfqpoint{2.326000in}{0.100000in}}%
\pgfpathcurveto{\pgfqpoint{2.615602in}{0.100000in}}{\pgfqpoint{2.893381in}{0.157530in}}{\pgfqpoint{3.098161in}{0.259920in}}%
\pgfpathcurveto{\pgfqpoint{3.302940in}{0.362309in}}{\pgfqpoint{3.418000in}{0.501199in}}{\pgfqpoint{3.418000in}{0.646000in}}%
\pgfpathcurveto{\pgfqpoint{3.418000in}{0.790801in}}{\pgfqpoint{3.302940in}{0.929691in}}{\pgfqpoint{3.098161in}{1.032080in}}%
\pgfpathcurveto{\pgfqpoint{2.893381in}{1.134470in}}{\pgfqpoint{2.615602in}{1.192000in}}{\pgfqpoint{2.326000in}{1.192000in}}%
\pgfpathcurveto{\pgfqpoint{2.036398in}{1.192000in}}{\pgfqpoint{1.758619in}{1.134470in}}{\pgfqpoint{1.553839in}{1.032080in}}%
\pgfpathcurveto{\pgfqpoint{1.349060in}{0.929691in}}{\pgfqpoint{1.234000in}{0.790801in}}{\pgfqpoint{1.234000in}{0.646000in}}%
\pgfpathcurveto{\pgfqpoint{1.234000in}{0.501199in}}{\pgfqpoint{1.349060in}{0.362309in}}{\pgfqpoint{1.553839in}{0.259920in}}%
\pgfpathcurveto{\pgfqpoint{1.758619in}{0.157530in}}{\pgfqpoint{2.036398in}{0.100000in}}{\pgfqpoint{2.326000in}{0.100000in}}%
\pgfpathlineto{\pgfqpoint{2.326000in}{0.100000in}}%
\pgfpathclose%
\pgfusepath{clip}%
\pgfsetrectcap%
\pgfsetroundjoin%
\pgfsetlinewidth{0.451687pt}%
\definecolor{currentstroke}{rgb}{0.690196,0.690196,0.690196}%
\pgfsetstrokecolor{currentstroke}%
\pgfsetstrokeopacity{0.250000}%
\pgfsetdash{}{0pt}%
\pgfpathmoveto{\pgfqpoint{1.445711in}{0.322907in}}%
\pgfpathlineto{\pgfqpoint{1.469186in}{0.322907in}}%
\pgfpathlineto{\pgfqpoint{1.492660in}{0.322907in}}%
\pgfpathlineto{\pgfqpoint{1.516134in}{0.322907in}}%
\pgfpathlineto{\pgfqpoint{1.539609in}{0.322907in}}%
\pgfpathlineto{\pgfqpoint{1.563083in}{0.322907in}}%
\pgfpathlineto{\pgfqpoint{1.586558in}{0.322907in}}%
\pgfpathlineto{\pgfqpoint{1.610032in}{0.322907in}}%
\pgfpathlineto{\pgfqpoint{1.633506in}{0.322907in}}%
\pgfpathlineto{\pgfqpoint{1.656981in}{0.322907in}}%
\pgfpathlineto{\pgfqpoint{1.680455in}{0.322907in}}%
\pgfpathlineto{\pgfqpoint{1.703929in}{0.322907in}}%
\pgfpathlineto{\pgfqpoint{1.727404in}{0.322907in}}%
\pgfpathlineto{\pgfqpoint{1.750878in}{0.322907in}}%
\pgfpathlineto{\pgfqpoint{1.774352in}{0.322907in}}%
\pgfpathlineto{\pgfqpoint{1.797827in}{0.322907in}}%
\pgfpathlineto{\pgfqpoint{1.821301in}{0.322907in}}%
\pgfpathlineto{\pgfqpoint{1.844776in}{0.322907in}}%
\pgfpathlineto{\pgfqpoint{1.868250in}{0.322907in}}%
\pgfpathlineto{\pgfqpoint{1.891724in}{0.322907in}}%
\pgfpathlineto{\pgfqpoint{1.915199in}{0.322907in}}%
\pgfpathlineto{\pgfqpoint{1.938673in}{0.322907in}}%
\pgfpathlineto{\pgfqpoint{1.962147in}{0.322907in}}%
\pgfpathlineto{\pgfqpoint{1.985622in}{0.322907in}}%
\pgfpathlineto{\pgfqpoint{2.009096in}{0.322907in}}%
\pgfpathlineto{\pgfqpoint{2.032570in}{0.322907in}}%
\pgfpathlineto{\pgfqpoint{2.056045in}{0.322907in}}%
\pgfpathlineto{\pgfqpoint{2.079519in}{0.322907in}}%
\pgfpathlineto{\pgfqpoint{2.102994in}{0.322907in}}%
\pgfpathlineto{\pgfqpoint{2.126468in}{0.322907in}}%
\pgfpathlineto{\pgfqpoint{2.149942in}{0.322907in}}%
\pgfpathlineto{\pgfqpoint{2.173417in}{0.322907in}}%
\pgfpathlineto{\pgfqpoint{2.196891in}{0.322907in}}%
\pgfpathlineto{\pgfqpoint{2.220365in}{0.322907in}}%
\pgfpathlineto{\pgfqpoint{2.243840in}{0.322907in}}%
\pgfpathlineto{\pgfqpoint{2.267314in}{0.322907in}}%
\pgfpathlineto{\pgfqpoint{2.290788in}{0.322907in}}%
\pgfpathlineto{\pgfqpoint{2.314263in}{0.322907in}}%
\pgfpathlineto{\pgfqpoint{2.337737in}{0.322907in}}%
\pgfpathlineto{\pgfqpoint{2.361212in}{0.322907in}}%
\pgfpathlineto{\pgfqpoint{2.384686in}{0.322907in}}%
\pgfpathlineto{\pgfqpoint{2.408160in}{0.322907in}}%
\pgfpathlineto{\pgfqpoint{2.431635in}{0.322907in}}%
\pgfpathlineto{\pgfqpoint{2.455109in}{0.322907in}}%
\pgfpathlineto{\pgfqpoint{2.478583in}{0.322907in}}%
\pgfpathlineto{\pgfqpoint{2.502058in}{0.322907in}}%
\pgfpathlineto{\pgfqpoint{2.525532in}{0.322907in}}%
\pgfpathlineto{\pgfqpoint{2.549006in}{0.322907in}}%
\pgfpathlineto{\pgfqpoint{2.572481in}{0.322907in}}%
\pgfpathlineto{\pgfqpoint{2.595955in}{0.322907in}}%
\pgfpathlineto{\pgfqpoint{2.619430in}{0.322907in}}%
\pgfpathlineto{\pgfqpoint{2.642904in}{0.322907in}}%
\pgfpathlineto{\pgfqpoint{2.666378in}{0.322907in}}%
\pgfpathlineto{\pgfqpoint{2.689853in}{0.322907in}}%
\pgfpathlineto{\pgfqpoint{2.713327in}{0.322907in}}%
\pgfpathlineto{\pgfqpoint{2.736801in}{0.322907in}}%
\pgfpathlineto{\pgfqpoint{2.760276in}{0.322907in}}%
\pgfpathlineto{\pgfqpoint{2.783750in}{0.322907in}}%
\pgfpathlineto{\pgfqpoint{2.807224in}{0.322907in}}%
\pgfpathlineto{\pgfqpoint{2.830699in}{0.322907in}}%
\pgfpathlineto{\pgfqpoint{2.854173in}{0.322907in}}%
\pgfpathlineto{\pgfqpoint{2.877648in}{0.322907in}}%
\pgfpathlineto{\pgfqpoint{2.901122in}{0.322907in}}%
\pgfpathlineto{\pgfqpoint{2.924596in}{0.322907in}}%
\pgfpathlineto{\pgfqpoint{2.948071in}{0.322907in}}%
\pgfpathlineto{\pgfqpoint{2.971545in}{0.322907in}}%
\pgfpathlineto{\pgfqpoint{2.995019in}{0.322907in}}%
\pgfpathlineto{\pgfqpoint{3.018494in}{0.322907in}}%
\pgfpathlineto{\pgfqpoint{3.041968in}{0.322907in}}%
\pgfpathlineto{\pgfqpoint{3.065442in}{0.322907in}}%
\pgfpathlineto{\pgfqpoint{3.088917in}{0.322907in}}%
\pgfpathlineto{\pgfqpoint{3.112391in}{0.322907in}}%
\pgfpathlineto{\pgfqpoint{3.135866in}{0.322907in}}%
\pgfpathlineto{\pgfqpoint{3.159340in}{0.322907in}}%
\pgfpathlineto{\pgfqpoint{3.182814in}{0.322907in}}%
\pgfpathlineto{\pgfqpoint{3.206289in}{0.322907in}}%
\pgfusepath{stroke}%
\end{pgfscope}%
\begin{pgfscope}%
\pgfpathmoveto{\pgfqpoint{2.326000in}{0.100000in}}%
\pgfpathcurveto{\pgfqpoint{2.615602in}{0.100000in}}{\pgfqpoint{2.893381in}{0.157530in}}{\pgfqpoint{3.098161in}{0.259920in}}%
\pgfpathcurveto{\pgfqpoint{3.302940in}{0.362309in}}{\pgfqpoint{3.418000in}{0.501199in}}{\pgfqpoint{3.418000in}{0.646000in}}%
\pgfpathcurveto{\pgfqpoint{3.418000in}{0.790801in}}{\pgfqpoint{3.302940in}{0.929691in}}{\pgfqpoint{3.098161in}{1.032080in}}%
\pgfpathcurveto{\pgfqpoint{2.893381in}{1.134470in}}{\pgfqpoint{2.615602in}{1.192000in}}{\pgfqpoint{2.326000in}{1.192000in}}%
\pgfpathcurveto{\pgfqpoint{2.036398in}{1.192000in}}{\pgfqpoint{1.758619in}{1.134470in}}{\pgfqpoint{1.553839in}{1.032080in}}%
\pgfpathcurveto{\pgfqpoint{1.349060in}{0.929691in}}{\pgfqpoint{1.234000in}{0.790801in}}{\pgfqpoint{1.234000in}{0.646000in}}%
\pgfpathcurveto{\pgfqpoint{1.234000in}{0.501199in}}{\pgfqpoint{1.349060in}{0.362309in}}{\pgfqpoint{1.553839in}{0.259920in}}%
\pgfpathcurveto{\pgfqpoint{1.758619in}{0.157530in}}{\pgfqpoint{2.036398in}{0.100000in}}{\pgfqpoint{2.326000in}{0.100000in}}%
\pgfpathlineto{\pgfqpoint{2.326000in}{0.100000in}}%
\pgfpathclose%
\pgfusepath{clip}%
\pgfsetrectcap%
\pgfsetroundjoin%
\pgfsetlinewidth{0.451687pt}%
\definecolor{currentstroke}{rgb}{0.690196,0.690196,0.690196}%
\pgfsetstrokecolor{currentstroke}%
\pgfsetstrokeopacity{0.250000}%
\pgfsetdash{}{0pt}%
\pgfpathmoveto{\pgfqpoint{1.327062in}{0.425440in}}%
\pgfpathlineto{\pgfqpoint{1.353701in}{0.425440in}}%
\pgfpathlineto{\pgfqpoint{1.380339in}{0.425440in}}%
\pgfpathlineto{\pgfqpoint{1.406977in}{0.425440in}}%
\pgfpathlineto{\pgfqpoint{1.433616in}{0.425440in}}%
\pgfpathlineto{\pgfqpoint{1.460254in}{0.425440in}}%
\pgfpathlineto{\pgfqpoint{1.486892in}{0.425440in}}%
\pgfpathlineto{\pgfqpoint{1.513531in}{0.425440in}}%
\pgfpathlineto{\pgfqpoint{1.540169in}{0.425440in}}%
\pgfpathlineto{\pgfqpoint{1.566807in}{0.425440in}}%
\pgfpathlineto{\pgfqpoint{1.593446in}{0.425440in}}%
\pgfpathlineto{\pgfqpoint{1.620084in}{0.425440in}}%
\pgfpathlineto{\pgfqpoint{1.646722in}{0.425440in}}%
\pgfpathlineto{\pgfqpoint{1.673361in}{0.425440in}}%
\pgfpathlineto{\pgfqpoint{1.699999in}{0.425440in}}%
\pgfpathlineto{\pgfqpoint{1.726637in}{0.425440in}}%
\pgfpathlineto{\pgfqpoint{1.753276in}{0.425440in}}%
\pgfpathlineto{\pgfqpoint{1.779914in}{0.425440in}}%
\pgfpathlineto{\pgfqpoint{1.806552in}{0.425440in}}%
\pgfpathlineto{\pgfqpoint{1.833191in}{0.425440in}}%
\pgfpathlineto{\pgfqpoint{1.859829in}{0.425440in}}%
\pgfpathlineto{\pgfqpoint{1.886467in}{0.425440in}}%
\pgfpathlineto{\pgfqpoint{1.913106in}{0.425440in}}%
\pgfpathlineto{\pgfqpoint{1.939744in}{0.425440in}}%
\pgfpathlineto{\pgfqpoint{1.966382in}{0.425440in}}%
\pgfpathlineto{\pgfqpoint{1.993021in}{0.425440in}}%
\pgfpathlineto{\pgfqpoint{2.019659in}{0.425440in}}%
\pgfpathlineto{\pgfqpoint{2.046297in}{0.425440in}}%
\pgfpathlineto{\pgfqpoint{2.072936in}{0.425440in}}%
\pgfpathlineto{\pgfqpoint{2.099574in}{0.425440in}}%
\pgfpathlineto{\pgfqpoint{2.126212in}{0.425440in}}%
\pgfpathlineto{\pgfqpoint{2.152851in}{0.425440in}}%
\pgfpathlineto{\pgfqpoint{2.179489in}{0.425440in}}%
\pgfpathlineto{\pgfqpoint{2.206127in}{0.425440in}}%
\pgfpathlineto{\pgfqpoint{2.232766in}{0.425440in}}%
\pgfpathlineto{\pgfqpoint{2.259404in}{0.425440in}}%
\pgfpathlineto{\pgfqpoint{2.286042in}{0.425440in}}%
\pgfpathlineto{\pgfqpoint{2.312681in}{0.425440in}}%
\pgfpathlineto{\pgfqpoint{2.339319in}{0.425440in}}%
\pgfpathlineto{\pgfqpoint{2.365958in}{0.425440in}}%
\pgfpathlineto{\pgfqpoint{2.392596in}{0.425440in}}%
\pgfpathlineto{\pgfqpoint{2.419234in}{0.425440in}}%
\pgfpathlineto{\pgfqpoint{2.445873in}{0.425440in}}%
\pgfpathlineto{\pgfqpoint{2.472511in}{0.425440in}}%
\pgfpathlineto{\pgfqpoint{2.499149in}{0.425440in}}%
\pgfpathlineto{\pgfqpoint{2.525788in}{0.425440in}}%
\pgfpathlineto{\pgfqpoint{2.552426in}{0.425440in}}%
\pgfpathlineto{\pgfqpoint{2.579064in}{0.425440in}}%
\pgfpathlineto{\pgfqpoint{2.605703in}{0.425440in}}%
\pgfpathlineto{\pgfqpoint{2.632341in}{0.425440in}}%
\pgfpathlineto{\pgfqpoint{2.658979in}{0.425440in}}%
\pgfpathlineto{\pgfqpoint{2.685618in}{0.425440in}}%
\pgfpathlineto{\pgfqpoint{2.712256in}{0.425440in}}%
\pgfpathlineto{\pgfqpoint{2.738894in}{0.425440in}}%
\pgfpathlineto{\pgfqpoint{2.765533in}{0.425440in}}%
\pgfpathlineto{\pgfqpoint{2.792171in}{0.425440in}}%
\pgfpathlineto{\pgfqpoint{2.818809in}{0.425440in}}%
\pgfpathlineto{\pgfqpoint{2.845448in}{0.425440in}}%
\pgfpathlineto{\pgfqpoint{2.872086in}{0.425440in}}%
\pgfpathlineto{\pgfqpoint{2.898724in}{0.425440in}}%
\pgfpathlineto{\pgfqpoint{2.925363in}{0.425440in}}%
\pgfpathlineto{\pgfqpoint{2.952001in}{0.425440in}}%
\pgfpathlineto{\pgfqpoint{2.978639in}{0.425440in}}%
\pgfpathlineto{\pgfqpoint{3.005278in}{0.425440in}}%
\pgfpathlineto{\pgfqpoint{3.031916in}{0.425440in}}%
\pgfpathlineto{\pgfqpoint{3.058554in}{0.425440in}}%
\pgfpathlineto{\pgfqpoint{3.085193in}{0.425440in}}%
\pgfpathlineto{\pgfqpoint{3.111831in}{0.425440in}}%
\pgfpathlineto{\pgfqpoint{3.138469in}{0.425440in}}%
\pgfpathlineto{\pgfqpoint{3.165108in}{0.425440in}}%
\pgfpathlineto{\pgfqpoint{3.191746in}{0.425440in}}%
\pgfpathlineto{\pgfqpoint{3.218384in}{0.425440in}}%
\pgfpathlineto{\pgfqpoint{3.245023in}{0.425440in}}%
\pgfpathlineto{\pgfqpoint{3.271661in}{0.425440in}}%
\pgfpathlineto{\pgfqpoint{3.298299in}{0.425440in}}%
\pgfpathlineto{\pgfqpoint{3.324938in}{0.425440in}}%
\pgfusepath{stroke}%
\end{pgfscope}%
\begin{pgfscope}%
\pgfpathmoveto{\pgfqpoint{2.326000in}{0.100000in}}%
\pgfpathcurveto{\pgfqpoint{2.615602in}{0.100000in}}{\pgfqpoint{2.893381in}{0.157530in}}{\pgfqpoint{3.098161in}{0.259920in}}%
\pgfpathcurveto{\pgfqpoint{3.302940in}{0.362309in}}{\pgfqpoint{3.418000in}{0.501199in}}{\pgfqpoint{3.418000in}{0.646000in}}%
\pgfpathcurveto{\pgfqpoint{3.418000in}{0.790801in}}{\pgfqpoint{3.302940in}{0.929691in}}{\pgfqpoint{3.098161in}{1.032080in}}%
\pgfpathcurveto{\pgfqpoint{2.893381in}{1.134470in}}{\pgfqpoint{2.615602in}{1.192000in}}{\pgfqpoint{2.326000in}{1.192000in}}%
\pgfpathcurveto{\pgfqpoint{2.036398in}{1.192000in}}{\pgfqpoint{1.758619in}{1.134470in}}{\pgfqpoint{1.553839in}{1.032080in}}%
\pgfpathcurveto{\pgfqpoint{1.349060in}{0.929691in}}{\pgfqpoint{1.234000in}{0.790801in}}{\pgfqpoint{1.234000in}{0.646000in}}%
\pgfpathcurveto{\pgfqpoint{1.234000in}{0.501199in}}{\pgfqpoint{1.349060in}{0.362309in}}{\pgfqpoint{1.553839in}{0.259920in}}%
\pgfpathcurveto{\pgfqpoint{1.758619in}{0.157530in}}{\pgfqpoint{2.036398in}{0.100000in}}{\pgfqpoint{2.326000in}{0.100000in}}%
\pgfpathlineto{\pgfqpoint{2.326000in}{0.100000in}}%
\pgfpathclose%
\pgfusepath{clip}%
\pgfsetrectcap%
\pgfsetroundjoin%
\pgfsetlinewidth{0.451687pt}%
\definecolor{currentstroke}{rgb}{0.690196,0.690196,0.690196}%
\pgfsetstrokecolor{currentstroke}%
\pgfsetstrokeopacity{0.250000}%
\pgfsetdash{}{0pt}%
\pgfpathmoveto{\pgfqpoint{1.257127in}{0.534226in}}%
\pgfpathlineto{\pgfqpoint{1.285630in}{0.534226in}}%
\pgfpathlineto{\pgfqpoint{1.314133in}{0.534226in}}%
\pgfpathlineto{\pgfqpoint{1.342637in}{0.534226in}}%
\pgfpathlineto{\pgfqpoint{1.371140in}{0.534226in}}%
\pgfpathlineto{\pgfqpoint{1.399643in}{0.534226in}}%
\pgfpathlineto{\pgfqpoint{1.428146in}{0.534226in}}%
\pgfpathlineto{\pgfqpoint{1.456650in}{0.534226in}}%
\pgfpathlineto{\pgfqpoint{1.485153in}{0.534226in}}%
\pgfpathlineto{\pgfqpoint{1.513656in}{0.534226in}}%
\pgfpathlineto{\pgfqpoint{1.542160in}{0.534226in}}%
\pgfpathlineto{\pgfqpoint{1.570663in}{0.534226in}}%
\pgfpathlineto{\pgfqpoint{1.599166in}{0.534226in}}%
\pgfpathlineto{\pgfqpoint{1.627669in}{0.534226in}}%
\pgfpathlineto{\pgfqpoint{1.656173in}{0.534226in}}%
\pgfpathlineto{\pgfqpoint{1.684676in}{0.534226in}}%
\pgfpathlineto{\pgfqpoint{1.713179in}{0.534226in}}%
\pgfpathlineto{\pgfqpoint{1.741683in}{0.534226in}}%
\pgfpathlineto{\pgfqpoint{1.770186in}{0.534226in}}%
\pgfpathlineto{\pgfqpoint{1.798689in}{0.534226in}}%
\pgfpathlineto{\pgfqpoint{1.827192in}{0.534226in}}%
\pgfpathlineto{\pgfqpoint{1.855696in}{0.534226in}}%
\pgfpathlineto{\pgfqpoint{1.884199in}{0.534226in}}%
\pgfpathlineto{\pgfqpoint{1.912702in}{0.534226in}}%
\pgfpathlineto{\pgfqpoint{1.941206in}{0.534226in}}%
\pgfpathlineto{\pgfqpoint{1.969709in}{0.534226in}}%
\pgfpathlineto{\pgfqpoint{1.998212in}{0.534226in}}%
\pgfpathlineto{\pgfqpoint{2.026715in}{0.534226in}}%
\pgfpathlineto{\pgfqpoint{2.055219in}{0.534226in}}%
\pgfpathlineto{\pgfqpoint{2.083722in}{0.534226in}}%
\pgfpathlineto{\pgfqpoint{2.112225in}{0.534226in}}%
\pgfpathlineto{\pgfqpoint{2.140729in}{0.534226in}}%
\pgfpathlineto{\pgfqpoint{2.169232in}{0.534226in}}%
\pgfpathlineto{\pgfqpoint{2.197735in}{0.534226in}}%
\pgfpathlineto{\pgfqpoint{2.226238in}{0.534226in}}%
\pgfpathlineto{\pgfqpoint{2.254742in}{0.534226in}}%
\pgfpathlineto{\pgfqpoint{2.283245in}{0.534226in}}%
\pgfpathlineto{\pgfqpoint{2.311748in}{0.534226in}}%
\pgfpathlineto{\pgfqpoint{2.340252in}{0.534226in}}%
\pgfpathlineto{\pgfqpoint{2.368755in}{0.534226in}}%
\pgfpathlineto{\pgfqpoint{2.397258in}{0.534226in}}%
\pgfpathlineto{\pgfqpoint{2.425762in}{0.534226in}}%
\pgfpathlineto{\pgfqpoint{2.454265in}{0.534226in}}%
\pgfpathlineto{\pgfqpoint{2.482768in}{0.534226in}}%
\pgfpathlineto{\pgfqpoint{2.511271in}{0.534226in}}%
\pgfpathlineto{\pgfqpoint{2.539775in}{0.534226in}}%
\pgfpathlineto{\pgfqpoint{2.568278in}{0.534226in}}%
\pgfpathlineto{\pgfqpoint{2.596781in}{0.534226in}}%
\pgfpathlineto{\pgfqpoint{2.625285in}{0.534226in}}%
\pgfpathlineto{\pgfqpoint{2.653788in}{0.534226in}}%
\pgfpathlineto{\pgfqpoint{2.682291in}{0.534226in}}%
\pgfpathlineto{\pgfqpoint{2.710794in}{0.534226in}}%
\pgfpathlineto{\pgfqpoint{2.739298in}{0.534226in}}%
\pgfpathlineto{\pgfqpoint{2.767801in}{0.534226in}}%
\pgfpathlineto{\pgfqpoint{2.796304in}{0.534226in}}%
\pgfpathlineto{\pgfqpoint{2.824808in}{0.534226in}}%
\pgfpathlineto{\pgfqpoint{2.853311in}{0.534226in}}%
\pgfpathlineto{\pgfqpoint{2.881814in}{0.534226in}}%
\pgfpathlineto{\pgfqpoint{2.910317in}{0.534226in}}%
\pgfpathlineto{\pgfqpoint{2.938821in}{0.534226in}}%
\pgfpathlineto{\pgfqpoint{2.967324in}{0.534226in}}%
\pgfpathlineto{\pgfqpoint{2.995827in}{0.534226in}}%
\pgfpathlineto{\pgfqpoint{3.024331in}{0.534226in}}%
\pgfpathlineto{\pgfqpoint{3.052834in}{0.534226in}}%
\pgfpathlineto{\pgfqpoint{3.081337in}{0.534226in}}%
\pgfpathlineto{\pgfqpoint{3.109840in}{0.534226in}}%
\pgfpathlineto{\pgfqpoint{3.138344in}{0.534226in}}%
\pgfpathlineto{\pgfqpoint{3.166847in}{0.534226in}}%
\pgfpathlineto{\pgfqpoint{3.195350in}{0.534226in}}%
\pgfpathlineto{\pgfqpoint{3.223854in}{0.534226in}}%
\pgfpathlineto{\pgfqpoint{3.252357in}{0.534226in}}%
\pgfpathlineto{\pgfqpoint{3.280860in}{0.534226in}}%
\pgfpathlineto{\pgfqpoint{3.309363in}{0.534226in}}%
\pgfpathlineto{\pgfqpoint{3.337867in}{0.534226in}}%
\pgfpathlineto{\pgfqpoint{3.366370in}{0.534226in}}%
\pgfpathlineto{\pgfqpoint{3.394873in}{0.534226in}}%
\pgfusepath{stroke}%
\end{pgfscope}%
\begin{pgfscope}%
\pgfpathmoveto{\pgfqpoint{2.326000in}{0.100000in}}%
\pgfpathcurveto{\pgfqpoint{2.615602in}{0.100000in}}{\pgfqpoint{2.893381in}{0.157530in}}{\pgfqpoint{3.098161in}{0.259920in}}%
\pgfpathcurveto{\pgfqpoint{3.302940in}{0.362309in}}{\pgfqpoint{3.418000in}{0.501199in}}{\pgfqpoint{3.418000in}{0.646000in}}%
\pgfpathcurveto{\pgfqpoint{3.418000in}{0.790801in}}{\pgfqpoint{3.302940in}{0.929691in}}{\pgfqpoint{3.098161in}{1.032080in}}%
\pgfpathcurveto{\pgfqpoint{2.893381in}{1.134470in}}{\pgfqpoint{2.615602in}{1.192000in}}{\pgfqpoint{2.326000in}{1.192000in}}%
\pgfpathcurveto{\pgfqpoint{2.036398in}{1.192000in}}{\pgfqpoint{1.758619in}{1.134470in}}{\pgfqpoint{1.553839in}{1.032080in}}%
\pgfpathcurveto{\pgfqpoint{1.349060in}{0.929691in}}{\pgfqpoint{1.234000in}{0.790801in}}{\pgfqpoint{1.234000in}{0.646000in}}%
\pgfpathcurveto{\pgfqpoint{1.234000in}{0.501199in}}{\pgfqpoint{1.349060in}{0.362309in}}{\pgfqpoint{1.553839in}{0.259920in}}%
\pgfpathcurveto{\pgfqpoint{1.758619in}{0.157530in}}{\pgfqpoint{2.036398in}{0.100000in}}{\pgfqpoint{2.326000in}{0.100000in}}%
\pgfpathlineto{\pgfqpoint{2.326000in}{0.100000in}}%
\pgfpathclose%
\pgfusepath{clip}%
\pgfsetrectcap%
\pgfsetroundjoin%
\pgfsetlinewidth{0.451687pt}%
\definecolor{currentstroke}{rgb}{0.690196,0.690196,0.690196}%
\pgfsetstrokecolor{currentstroke}%
\pgfsetstrokeopacity{0.250000}%
\pgfsetdash{}{0pt}%
\pgfpathmoveto{\pgfqpoint{1.234000in}{0.646000in}}%
\pgfpathlineto{\pgfqpoint{1.263120in}{0.646000in}}%
\pgfpathlineto{\pgfqpoint{1.292240in}{0.646000in}}%
\pgfpathlineto{\pgfqpoint{1.321360in}{0.646000in}}%
\pgfpathlineto{\pgfqpoint{1.350480in}{0.646000in}}%
\pgfpathlineto{\pgfqpoint{1.379600in}{0.646000in}}%
\pgfpathlineto{\pgfqpoint{1.408720in}{0.646000in}}%
\pgfpathlineto{\pgfqpoint{1.437840in}{0.646000in}}%
\pgfpathlineto{\pgfqpoint{1.466960in}{0.646000in}}%
\pgfpathlineto{\pgfqpoint{1.496080in}{0.646000in}}%
\pgfpathlineto{\pgfqpoint{1.525200in}{0.646000in}}%
\pgfpathlineto{\pgfqpoint{1.554320in}{0.646000in}}%
\pgfpathlineto{\pgfqpoint{1.583440in}{0.646000in}}%
\pgfpathlineto{\pgfqpoint{1.612560in}{0.646000in}}%
\pgfpathlineto{\pgfqpoint{1.641680in}{0.646000in}}%
\pgfpathlineto{\pgfqpoint{1.670800in}{0.646000in}}%
\pgfpathlineto{\pgfqpoint{1.699920in}{0.646000in}}%
\pgfpathlineto{\pgfqpoint{1.729040in}{0.646000in}}%
\pgfpathlineto{\pgfqpoint{1.758160in}{0.646000in}}%
\pgfpathlineto{\pgfqpoint{1.787280in}{0.646000in}}%
\pgfpathlineto{\pgfqpoint{1.816400in}{0.646000in}}%
\pgfpathlineto{\pgfqpoint{1.845520in}{0.646000in}}%
\pgfpathlineto{\pgfqpoint{1.874640in}{0.646000in}}%
\pgfpathlineto{\pgfqpoint{1.903760in}{0.646000in}}%
\pgfpathlineto{\pgfqpoint{1.932880in}{0.646000in}}%
\pgfpathlineto{\pgfqpoint{1.962000in}{0.646000in}}%
\pgfpathlineto{\pgfqpoint{1.991120in}{0.646000in}}%
\pgfpathlineto{\pgfqpoint{2.020240in}{0.646000in}}%
\pgfpathlineto{\pgfqpoint{2.049360in}{0.646000in}}%
\pgfpathlineto{\pgfqpoint{2.078480in}{0.646000in}}%
\pgfpathlineto{\pgfqpoint{2.107600in}{0.646000in}}%
\pgfpathlineto{\pgfqpoint{2.136720in}{0.646000in}}%
\pgfpathlineto{\pgfqpoint{2.165840in}{0.646000in}}%
\pgfpathlineto{\pgfqpoint{2.194960in}{0.646000in}}%
\pgfpathlineto{\pgfqpoint{2.224080in}{0.646000in}}%
\pgfpathlineto{\pgfqpoint{2.253200in}{0.646000in}}%
\pgfpathlineto{\pgfqpoint{2.282320in}{0.646000in}}%
\pgfpathlineto{\pgfqpoint{2.311440in}{0.646000in}}%
\pgfpathlineto{\pgfqpoint{2.340560in}{0.646000in}}%
\pgfpathlineto{\pgfqpoint{2.369680in}{0.646000in}}%
\pgfpathlineto{\pgfqpoint{2.398800in}{0.646000in}}%
\pgfpathlineto{\pgfqpoint{2.427920in}{0.646000in}}%
\pgfpathlineto{\pgfqpoint{2.457040in}{0.646000in}}%
\pgfpathlineto{\pgfqpoint{2.486160in}{0.646000in}}%
\pgfpathlineto{\pgfqpoint{2.515280in}{0.646000in}}%
\pgfpathlineto{\pgfqpoint{2.544400in}{0.646000in}}%
\pgfpathlineto{\pgfqpoint{2.573520in}{0.646000in}}%
\pgfpathlineto{\pgfqpoint{2.602640in}{0.646000in}}%
\pgfpathlineto{\pgfqpoint{2.631760in}{0.646000in}}%
\pgfpathlineto{\pgfqpoint{2.660880in}{0.646000in}}%
\pgfpathlineto{\pgfqpoint{2.690000in}{0.646000in}}%
\pgfpathlineto{\pgfqpoint{2.719120in}{0.646000in}}%
\pgfpathlineto{\pgfqpoint{2.748240in}{0.646000in}}%
\pgfpathlineto{\pgfqpoint{2.777360in}{0.646000in}}%
\pgfpathlineto{\pgfqpoint{2.806480in}{0.646000in}}%
\pgfpathlineto{\pgfqpoint{2.835600in}{0.646000in}}%
\pgfpathlineto{\pgfqpoint{2.864720in}{0.646000in}}%
\pgfpathlineto{\pgfqpoint{2.893840in}{0.646000in}}%
\pgfpathlineto{\pgfqpoint{2.922960in}{0.646000in}}%
\pgfpathlineto{\pgfqpoint{2.952080in}{0.646000in}}%
\pgfpathlineto{\pgfqpoint{2.981200in}{0.646000in}}%
\pgfpathlineto{\pgfqpoint{3.010320in}{0.646000in}}%
\pgfpathlineto{\pgfqpoint{3.039440in}{0.646000in}}%
\pgfpathlineto{\pgfqpoint{3.068560in}{0.646000in}}%
\pgfpathlineto{\pgfqpoint{3.097680in}{0.646000in}}%
\pgfpathlineto{\pgfqpoint{3.126800in}{0.646000in}}%
\pgfpathlineto{\pgfqpoint{3.155920in}{0.646000in}}%
\pgfpathlineto{\pgfqpoint{3.185040in}{0.646000in}}%
\pgfpathlineto{\pgfqpoint{3.214160in}{0.646000in}}%
\pgfpathlineto{\pgfqpoint{3.243280in}{0.646000in}}%
\pgfpathlineto{\pgfqpoint{3.272400in}{0.646000in}}%
\pgfpathlineto{\pgfqpoint{3.301520in}{0.646000in}}%
\pgfpathlineto{\pgfqpoint{3.330640in}{0.646000in}}%
\pgfpathlineto{\pgfqpoint{3.359760in}{0.646000in}}%
\pgfpathlineto{\pgfqpoint{3.388880in}{0.646000in}}%
\pgfpathlineto{\pgfqpoint{3.418000in}{0.646000in}}%
\pgfusepath{stroke}%
\end{pgfscope}%
\begin{pgfscope}%
\pgfpathmoveto{\pgfqpoint{2.326000in}{0.100000in}}%
\pgfpathcurveto{\pgfqpoint{2.615602in}{0.100000in}}{\pgfqpoint{2.893381in}{0.157530in}}{\pgfqpoint{3.098161in}{0.259920in}}%
\pgfpathcurveto{\pgfqpoint{3.302940in}{0.362309in}}{\pgfqpoint{3.418000in}{0.501199in}}{\pgfqpoint{3.418000in}{0.646000in}}%
\pgfpathcurveto{\pgfqpoint{3.418000in}{0.790801in}}{\pgfqpoint{3.302940in}{0.929691in}}{\pgfqpoint{3.098161in}{1.032080in}}%
\pgfpathcurveto{\pgfqpoint{2.893381in}{1.134470in}}{\pgfqpoint{2.615602in}{1.192000in}}{\pgfqpoint{2.326000in}{1.192000in}}%
\pgfpathcurveto{\pgfqpoint{2.036398in}{1.192000in}}{\pgfqpoint{1.758619in}{1.134470in}}{\pgfqpoint{1.553839in}{1.032080in}}%
\pgfpathcurveto{\pgfqpoint{1.349060in}{0.929691in}}{\pgfqpoint{1.234000in}{0.790801in}}{\pgfqpoint{1.234000in}{0.646000in}}%
\pgfpathcurveto{\pgfqpoint{1.234000in}{0.501199in}}{\pgfqpoint{1.349060in}{0.362309in}}{\pgfqpoint{1.553839in}{0.259920in}}%
\pgfpathcurveto{\pgfqpoint{1.758619in}{0.157530in}}{\pgfqpoint{2.036398in}{0.100000in}}{\pgfqpoint{2.326000in}{0.100000in}}%
\pgfpathlineto{\pgfqpoint{2.326000in}{0.100000in}}%
\pgfpathclose%
\pgfusepath{clip}%
\pgfsetrectcap%
\pgfsetroundjoin%
\pgfsetlinewidth{0.451687pt}%
\definecolor{currentstroke}{rgb}{0.690196,0.690196,0.690196}%
\pgfsetstrokecolor{currentstroke}%
\pgfsetstrokeopacity{0.250000}%
\pgfsetdash{}{0pt}%
\pgfpathmoveto{\pgfqpoint{1.257127in}{0.757774in}}%
\pgfpathlineto{\pgfqpoint{1.285630in}{0.757774in}}%
\pgfpathlineto{\pgfqpoint{1.314133in}{0.757774in}}%
\pgfpathlineto{\pgfqpoint{1.342637in}{0.757774in}}%
\pgfpathlineto{\pgfqpoint{1.371140in}{0.757774in}}%
\pgfpathlineto{\pgfqpoint{1.399643in}{0.757774in}}%
\pgfpathlineto{\pgfqpoint{1.428146in}{0.757774in}}%
\pgfpathlineto{\pgfqpoint{1.456650in}{0.757774in}}%
\pgfpathlineto{\pgfqpoint{1.485153in}{0.757774in}}%
\pgfpathlineto{\pgfqpoint{1.513656in}{0.757774in}}%
\pgfpathlineto{\pgfqpoint{1.542160in}{0.757774in}}%
\pgfpathlineto{\pgfqpoint{1.570663in}{0.757774in}}%
\pgfpathlineto{\pgfqpoint{1.599166in}{0.757774in}}%
\pgfpathlineto{\pgfqpoint{1.627669in}{0.757774in}}%
\pgfpathlineto{\pgfqpoint{1.656173in}{0.757774in}}%
\pgfpathlineto{\pgfqpoint{1.684676in}{0.757774in}}%
\pgfpathlineto{\pgfqpoint{1.713179in}{0.757774in}}%
\pgfpathlineto{\pgfqpoint{1.741683in}{0.757774in}}%
\pgfpathlineto{\pgfqpoint{1.770186in}{0.757774in}}%
\pgfpathlineto{\pgfqpoint{1.798689in}{0.757774in}}%
\pgfpathlineto{\pgfqpoint{1.827192in}{0.757774in}}%
\pgfpathlineto{\pgfqpoint{1.855696in}{0.757774in}}%
\pgfpathlineto{\pgfqpoint{1.884199in}{0.757774in}}%
\pgfpathlineto{\pgfqpoint{1.912702in}{0.757774in}}%
\pgfpathlineto{\pgfqpoint{1.941206in}{0.757774in}}%
\pgfpathlineto{\pgfqpoint{1.969709in}{0.757774in}}%
\pgfpathlineto{\pgfqpoint{1.998212in}{0.757774in}}%
\pgfpathlineto{\pgfqpoint{2.026715in}{0.757774in}}%
\pgfpathlineto{\pgfqpoint{2.055219in}{0.757774in}}%
\pgfpathlineto{\pgfqpoint{2.083722in}{0.757774in}}%
\pgfpathlineto{\pgfqpoint{2.112225in}{0.757774in}}%
\pgfpathlineto{\pgfqpoint{2.140729in}{0.757774in}}%
\pgfpathlineto{\pgfqpoint{2.169232in}{0.757774in}}%
\pgfpathlineto{\pgfqpoint{2.197735in}{0.757774in}}%
\pgfpathlineto{\pgfqpoint{2.226238in}{0.757774in}}%
\pgfpathlineto{\pgfqpoint{2.254742in}{0.757774in}}%
\pgfpathlineto{\pgfqpoint{2.283245in}{0.757774in}}%
\pgfpathlineto{\pgfqpoint{2.311748in}{0.757774in}}%
\pgfpathlineto{\pgfqpoint{2.340252in}{0.757774in}}%
\pgfpathlineto{\pgfqpoint{2.368755in}{0.757774in}}%
\pgfpathlineto{\pgfqpoint{2.397258in}{0.757774in}}%
\pgfpathlineto{\pgfqpoint{2.425762in}{0.757774in}}%
\pgfpathlineto{\pgfqpoint{2.454265in}{0.757774in}}%
\pgfpathlineto{\pgfqpoint{2.482768in}{0.757774in}}%
\pgfpathlineto{\pgfqpoint{2.511271in}{0.757774in}}%
\pgfpathlineto{\pgfqpoint{2.539775in}{0.757774in}}%
\pgfpathlineto{\pgfqpoint{2.568278in}{0.757774in}}%
\pgfpathlineto{\pgfqpoint{2.596781in}{0.757774in}}%
\pgfpathlineto{\pgfqpoint{2.625285in}{0.757774in}}%
\pgfpathlineto{\pgfqpoint{2.653788in}{0.757774in}}%
\pgfpathlineto{\pgfqpoint{2.682291in}{0.757774in}}%
\pgfpathlineto{\pgfqpoint{2.710794in}{0.757774in}}%
\pgfpathlineto{\pgfqpoint{2.739298in}{0.757774in}}%
\pgfpathlineto{\pgfqpoint{2.767801in}{0.757774in}}%
\pgfpathlineto{\pgfqpoint{2.796304in}{0.757774in}}%
\pgfpathlineto{\pgfqpoint{2.824808in}{0.757774in}}%
\pgfpathlineto{\pgfqpoint{2.853311in}{0.757774in}}%
\pgfpathlineto{\pgfqpoint{2.881814in}{0.757774in}}%
\pgfpathlineto{\pgfqpoint{2.910317in}{0.757774in}}%
\pgfpathlineto{\pgfqpoint{2.938821in}{0.757774in}}%
\pgfpathlineto{\pgfqpoint{2.967324in}{0.757774in}}%
\pgfpathlineto{\pgfqpoint{2.995827in}{0.757774in}}%
\pgfpathlineto{\pgfqpoint{3.024331in}{0.757774in}}%
\pgfpathlineto{\pgfqpoint{3.052834in}{0.757774in}}%
\pgfpathlineto{\pgfqpoint{3.081337in}{0.757774in}}%
\pgfpathlineto{\pgfqpoint{3.109840in}{0.757774in}}%
\pgfpathlineto{\pgfqpoint{3.138344in}{0.757774in}}%
\pgfpathlineto{\pgfqpoint{3.166847in}{0.757774in}}%
\pgfpathlineto{\pgfqpoint{3.195350in}{0.757774in}}%
\pgfpathlineto{\pgfqpoint{3.223854in}{0.757774in}}%
\pgfpathlineto{\pgfqpoint{3.252357in}{0.757774in}}%
\pgfpathlineto{\pgfqpoint{3.280860in}{0.757774in}}%
\pgfpathlineto{\pgfqpoint{3.309363in}{0.757774in}}%
\pgfpathlineto{\pgfqpoint{3.337867in}{0.757774in}}%
\pgfpathlineto{\pgfqpoint{3.366370in}{0.757774in}}%
\pgfpathlineto{\pgfqpoint{3.394873in}{0.757774in}}%
\pgfusepath{stroke}%
\end{pgfscope}%
\begin{pgfscope}%
\pgfpathmoveto{\pgfqpoint{2.326000in}{0.100000in}}%
\pgfpathcurveto{\pgfqpoint{2.615602in}{0.100000in}}{\pgfqpoint{2.893381in}{0.157530in}}{\pgfqpoint{3.098161in}{0.259920in}}%
\pgfpathcurveto{\pgfqpoint{3.302940in}{0.362309in}}{\pgfqpoint{3.418000in}{0.501199in}}{\pgfqpoint{3.418000in}{0.646000in}}%
\pgfpathcurveto{\pgfqpoint{3.418000in}{0.790801in}}{\pgfqpoint{3.302940in}{0.929691in}}{\pgfqpoint{3.098161in}{1.032080in}}%
\pgfpathcurveto{\pgfqpoint{2.893381in}{1.134470in}}{\pgfqpoint{2.615602in}{1.192000in}}{\pgfqpoint{2.326000in}{1.192000in}}%
\pgfpathcurveto{\pgfqpoint{2.036398in}{1.192000in}}{\pgfqpoint{1.758619in}{1.134470in}}{\pgfqpoint{1.553839in}{1.032080in}}%
\pgfpathcurveto{\pgfqpoint{1.349060in}{0.929691in}}{\pgfqpoint{1.234000in}{0.790801in}}{\pgfqpoint{1.234000in}{0.646000in}}%
\pgfpathcurveto{\pgfqpoint{1.234000in}{0.501199in}}{\pgfqpoint{1.349060in}{0.362309in}}{\pgfqpoint{1.553839in}{0.259920in}}%
\pgfpathcurveto{\pgfqpoint{1.758619in}{0.157530in}}{\pgfqpoint{2.036398in}{0.100000in}}{\pgfqpoint{2.326000in}{0.100000in}}%
\pgfpathlineto{\pgfqpoint{2.326000in}{0.100000in}}%
\pgfpathclose%
\pgfusepath{clip}%
\pgfsetrectcap%
\pgfsetroundjoin%
\pgfsetlinewidth{0.451687pt}%
\definecolor{currentstroke}{rgb}{0.690196,0.690196,0.690196}%
\pgfsetstrokecolor{currentstroke}%
\pgfsetstrokeopacity{0.250000}%
\pgfsetdash{}{0pt}%
\pgfpathmoveto{\pgfqpoint{1.327062in}{0.866560in}}%
\pgfpathlineto{\pgfqpoint{1.353701in}{0.866560in}}%
\pgfpathlineto{\pgfqpoint{1.380339in}{0.866560in}}%
\pgfpathlineto{\pgfqpoint{1.406977in}{0.866560in}}%
\pgfpathlineto{\pgfqpoint{1.433616in}{0.866560in}}%
\pgfpathlineto{\pgfqpoint{1.460254in}{0.866560in}}%
\pgfpathlineto{\pgfqpoint{1.486892in}{0.866560in}}%
\pgfpathlineto{\pgfqpoint{1.513531in}{0.866560in}}%
\pgfpathlineto{\pgfqpoint{1.540169in}{0.866560in}}%
\pgfpathlineto{\pgfqpoint{1.566807in}{0.866560in}}%
\pgfpathlineto{\pgfqpoint{1.593446in}{0.866560in}}%
\pgfpathlineto{\pgfqpoint{1.620084in}{0.866560in}}%
\pgfpathlineto{\pgfqpoint{1.646722in}{0.866560in}}%
\pgfpathlineto{\pgfqpoint{1.673361in}{0.866560in}}%
\pgfpathlineto{\pgfqpoint{1.699999in}{0.866560in}}%
\pgfpathlineto{\pgfqpoint{1.726637in}{0.866560in}}%
\pgfpathlineto{\pgfqpoint{1.753276in}{0.866560in}}%
\pgfpathlineto{\pgfqpoint{1.779914in}{0.866560in}}%
\pgfpathlineto{\pgfqpoint{1.806552in}{0.866560in}}%
\pgfpathlineto{\pgfqpoint{1.833191in}{0.866560in}}%
\pgfpathlineto{\pgfqpoint{1.859829in}{0.866560in}}%
\pgfpathlineto{\pgfqpoint{1.886467in}{0.866560in}}%
\pgfpathlineto{\pgfqpoint{1.913106in}{0.866560in}}%
\pgfpathlineto{\pgfqpoint{1.939744in}{0.866560in}}%
\pgfpathlineto{\pgfqpoint{1.966382in}{0.866560in}}%
\pgfpathlineto{\pgfqpoint{1.993021in}{0.866560in}}%
\pgfpathlineto{\pgfqpoint{2.019659in}{0.866560in}}%
\pgfpathlineto{\pgfqpoint{2.046297in}{0.866560in}}%
\pgfpathlineto{\pgfqpoint{2.072936in}{0.866560in}}%
\pgfpathlineto{\pgfqpoint{2.099574in}{0.866560in}}%
\pgfpathlineto{\pgfqpoint{2.126212in}{0.866560in}}%
\pgfpathlineto{\pgfqpoint{2.152851in}{0.866560in}}%
\pgfpathlineto{\pgfqpoint{2.179489in}{0.866560in}}%
\pgfpathlineto{\pgfqpoint{2.206127in}{0.866560in}}%
\pgfpathlineto{\pgfqpoint{2.232766in}{0.866560in}}%
\pgfpathlineto{\pgfqpoint{2.259404in}{0.866560in}}%
\pgfpathlineto{\pgfqpoint{2.286042in}{0.866560in}}%
\pgfpathlineto{\pgfqpoint{2.312681in}{0.866560in}}%
\pgfpathlineto{\pgfqpoint{2.339319in}{0.866560in}}%
\pgfpathlineto{\pgfqpoint{2.365958in}{0.866560in}}%
\pgfpathlineto{\pgfqpoint{2.392596in}{0.866560in}}%
\pgfpathlineto{\pgfqpoint{2.419234in}{0.866560in}}%
\pgfpathlineto{\pgfqpoint{2.445873in}{0.866560in}}%
\pgfpathlineto{\pgfqpoint{2.472511in}{0.866560in}}%
\pgfpathlineto{\pgfqpoint{2.499149in}{0.866560in}}%
\pgfpathlineto{\pgfqpoint{2.525788in}{0.866560in}}%
\pgfpathlineto{\pgfqpoint{2.552426in}{0.866560in}}%
\pgfpathlineto{\pgfqpoint{2.579064in}{0.866560in}}%
\pgfpathlineto{\pgfqpoint{2.605703in}{0.866560in}}%
\pgfpathlineto{\pgfqpoint{2.632341in}{0.866560in}}%
\pgfpathlineto{\pgfqpoint{2.658979in}{0.866560in}}%
\pgfpathlineto{\pgfqpoint{2.685618in}{0.866560in}}%
\pgfpathlineto{\pgfqpoint{2.712256in}{0.866560in}}%
\pgfpathlineto{\pgfqpoint{2.738894in}{0.866560in}}%
\pgfpathlineto{\pgfqpoint{2.765533in}{0.866560in}}%
\pgfpathlineto{\pgfqpoint{2.792171in}{0.866560in}}%
\pgfpathlineto{\pgfqpoint{2.818809in}{0.866560in}}%
\pgfpathlineto{\pgfqpoint{2.845448in}{0.866560in}}%
\pgfpathlineto{\pgfqpoint{2.872086in}{0.866560in}}%
\pgfpathlineto{\pgfqpoint{2.898724in}{0.866560in}}%
\pgfpathlineto{\pgfqpoint{2.925363in}{0.866560in}}%
\pgfpathlineto{\pgfqpoint{2.952001in}{0.866560in}}%
\pgfpathlineto{\pgfqpoint{2.978639in}{0.866560in}}%
\pgfpathlineto{\pgfqpoint{3.005278in}{0.866560in}}%
\pgfpathlineto{\pgfqpoint{3.031916in}{0.866560in}}%
\pgfpathlineto{\pgfqpoint{3.058554in}{0.866560in}}%
\pgfpathlineto{\pgfqpoint{3.085193in}{0.866560in}}%
\pgfpathlineto{\pgfqpoint{3.111831in}{0.866560in}}%
\pgfpathlineto{\pgfqpoint{3.138469in}{0.866560in}}%
\pgfpathlineto{\pgfqpoint{3.165108in}{0.866560in}}%
\pgfpathlineto{\pgfqpoint{3.191746in}{0.866560in}}%
\pgfpathlineto{\pgfqpoint{3.218384in}{0.866560in}}%
\pgfpathlineto{\pgfqpoint{3.245023in}{0.866560in}}%
\pgfpathlineto{\pgfqpoint{3.271661in}{0.866560in}}%
\pgfpathlineto{\pgfqpoint{3.298299in}{0.866560in}}%
\pgfpathlineto{\pgfqpoint{3.324938in}{0.866560in}}%
\pgfusepath{stroke}%
\end{pgfscope}%
\begin{pgfscope}%
\pgfpathmoveto{\pgfqpoint{2.326000in}{0.100000in}}%
\pgfpathcurveto{\pgfqpoint{2.615602in}{0.100000in}}{\pgfqpoint{2.893381in}{0.157530in}}{\pgfqpoint{3.098161in}{0.259920in}}%
\pgfpathcurveto{\pgfqpoint{3.302940in}{0.362309in}}{\pgfqpoint{3.418000in}{0.501199in}}{\pgfqpoint{3.418000in}{0.646000in}}%
\pgfpathcurveto{\pgfqpoint{3.418000in}{0.790801in}}{\pgfqpoint{3.302940in}{0.929691in}}{\pgfqpoint{3.098161in}{1.032080in}}%
\pgfpathcurveto{\pgfqpoint{2.893381in}{1.134470in}}{\pgfqpoint{2.615602in}{1.192000in}}{\pgfqpoint{2.326000in}{1.192000in}}%
\pgfpathcurveto{\pgfqpoint{2.036398in}{1.192000in}}{\pgfqpoint{1.758619in}{1.134470in}}{\pgfqpoint{1.553839in}{1.032080in}}%
\pgfpathcurveto{\pgfqpoint{1.349060in}{0.929691in}}{\pgfqpoint{1.234000in}{0.790801in}}{\pgfqpoint{1.234000in}{0.646000in}}%
\pgfpathcurveto{\pgfqpoint{1.234000in}{0.501199in}}{\pgfqpoint{1.349060in}{0.362309in}}{\pgfqpoint{1.553839in}{0.259920in}}%
\pgfpathcurveto{\pgfqpoint{1.758619in}{0.157530in}}{\pgfqpoint{2.036398in}{0.100000in}}{\pgfqpoint{2.326000in}{0.100000in}}%
\pgfpathlineto{\pgfqpoint{2.326000in}{0.100000in}}%
\pgfpathclose%
\pgfusepath{clip}%
\pgfsetrectcap%
\pgfsetroundjoin%
\pgfsetlinewidth{0.451687pt}%
\definecolor{currentstroke}{rgb}{0.690196,0.690196,0.690196}%
\pgfsetstrokecolor{currentstroke}%
\pgfsetstrokeopacity{0.250000}%
\pgfsetdash{}{0pt}%
\pgfpathmoveto{\pgfqpoint{1.445711in}{0.969093in}}%
\pgfpathlineto{\pgfqpoint{1.469186in}{0.969093in}}%
\pgfpathlineto{\pgfqpoint{1.492660in}{0.969093in}}%
\pgfpathlineto{\pgfqpoint{1.516134in}{0.969093in}}%
\pgfpathlineto{\pgfqpoint{1.539609in}{0.969093in}}%
\pgfpathlineto{\pgfqpoint{1.563083in}{0.969093in}}%
\pgfpathlineto{\pgfqpoint{1.586558in}{0.969093in}}%
\pgfpathlineto{\pgfqpoint{1.610032in}{0.969093in}}%
\pgfpathlineto{\pgfqpoint{1.633506in}{0.969093in}}%
\pgfpathlineto{\pgfqpoint{1.656981in}{0.969093in}}%
\pgfpathlineto{\pgfqpoint{1.680455in}{0.969093in}}%
\pgfpathlineto{\pgfqpoint{1.703929in}{0.969093in}}%
\pgfpathlineto{\pgfqpoint{1.727404in}{0.969093in}}%
\pgfpathlineto{\pgfqpoint{1.750878in}{0.969093in}}%
\pgfpathlineto{\pgfqpoint{1.774352in}{0.969093in}}%
\pgfpathlineto{\pgfqpoint{1.797827in}{0.969093in}}%
\pgfpathlineto{\pgfqpoint{1.821301in}{0.969093in}}%
\pgfpathlineto{\pgfqpoint{1.844776in}{0.969093in}}%
\pgfpathlineto{\pgfqpoint{1.868250in}{0.969093in}}%
\pgfpathlineto{\pgfqpoint{1.891724in}{0.969093in}}%
\pgfpathlineto{\pgfqpoint{1.915199in}{0.969093in}}%
\pgfpathlineto{\pgfqpoint{1.938673in}{0.969093in}}%
\pgfpathlineto{\pgfqpoint{1.962147in}{0.969093in}}%
\pgfpathlineto{\pgfqpoint{1.985622in}{0.969093in}}%
\pgfpathlineto{\pgfqpoint{2.009096in}{0.969093in}}%
\pgfpathlineto{\pgfqpoint{2.032570in}{0.969093in}}%
\pgfpathlineto{\pgfqpoint{2.056045in}{0.969093in}}%
\pgfpathlineto{\pgfqpoint{2.079519in}{0.969093in}}%
\pgfpathlineto{\pgfqpoint{2.102994in}{0.969093in}}%
\pgfpathlineto{\pgfqpoint{2.126468in}{0.969093in}}%
\pgfpathlineto{\pgfqpoint{2.149942in}{0.969093in}}%
\pgfpathlineto{\pgfqpoint{2.173417in}{0.969093in}}%
\pgfpathlineto{\pgfqpoint{2.196891in}{0.969093in}}%
\pgfpathlineto{\pgfqpoint{2.220365in}{0.969093in}}%
\pgfpathlineto{\pgfqpoint{2.243840in}{0.969093in}}%
\pgfpathlineto{\pgfqpoint{2.267314in}{0.969093in}}%
\pgfpathlineto{\pgfqpoint{2.290788in}{0.969093in}}%
\pgfpathlineto{\pgfqpoint{2.314263in}{0.969093in}}%
\pgfpathlineto{\pgfqpoint{2.337737in}{0.969093in}}%
\pgfpathlineto{\pgfqpoint{2.361212in}{0.969093in}}%
\pgfpathlineto{\pgfqpoint{2.384686in}{0.969093in}}%
\pgfpathlineto{\pgfqpoint{2.408160in}{0.969093in}}%
\pgfpathlineto{\pgfqpoint{2.431635in}{0.969093in}}%
\pgfpathlineto{\pgfqpoint{2.455109in}{0.969093in}}%
\pgfpathlineto{\pgfqpoint{2.478583in}{0.969093in}}%
\pgfpathlineto{\pgfqpoint{2.502058in}{0.969093in}}%
\pgfpathlineto{\pgfqpoint{2.525532in}{0.969093in}}%
\pgfpathlineto{\pgfqpoint{2.549006in}{0.969093in}}%
\pgfpathlineto{\pgfqpoint{2.572481in}{0.969093in}}%
\pgfpathlineto{\pgfqpoint{2.595955in}{0.969093in}}%
\pgfpathlineto{\pgfqpoint{2.619430in}{0.969093in}}%
\pgfpathlineto{\pgfqpoint{2.642904in}{0.969093in}}%
\pgfpathlineto{\pgfqpoint{2.666378in}{0.969093in}}%
\pgfpathlineto{\pgfqpoint{2.689853in}{0.969093in}}%
\pgfpathlineto{\pgfqpoint{2.713327in}{0.969093in}}%
\pgfpathlineto{\pgfqpoint{2.736801in}{0.969093in}}%
\pgfpathlineto{\pgfqpoint{2.760276in}{0.969093in}}%
\pgfpathlineto{\pgfqpoint{2.783750in}{0.969093in}}%
\pgfpathlineto{\pgfqpoint{2.807224in}{0.969093in}}%
\pgfpathlineto{\pgfqpoint{2.830699in}{0.969093in}}%
\pgfpathlineto{\pgfqpoint{2.854173in}{0.969093in}}%
\pgfpathlineto{\pgfqpoint{2.877648in}{0.969093in}}%
\pgfpathlineto{\pgfqpoint{2.901122in}{0.969093in}}%
\pgfpathlineto{\pgfqpoint{2.924596in}{0.969093in}}%
\pgfpathlineto{\pgfqpoint{2.948071in}{0.969093in}}%
\pgfpathlineto{\pgfqpoint{2.971545in}{0.969093in}}%
\pgfpathlineto{\pgfqpoint{2.995019in}{0.969093in}}%
\pgfpathlineto{\pgfqpoint{3.018494in}{0.969093in}}%
\pgfpathlineto{\pgfqpoint{3.041968in}{0.969093in}}%
\pgfpathlineto{\pgfqpoint{3.065442in}{0.969093in}}%
\pgfpathlineto{\pgfqpoint{3.088917in}{0.969093in}}%
\pgfpathlineto{\pgfqpoint{3.112391in}{0.969093in}}%
\pgfpathlineto{\pgfqpoint{3.135866in}{0.969093in}}%
\pgfpathlineto{\pgfqpoint{3.159340in}{0.969093in}}%
\pgfpathlineto{\pgfqpoint{3.182814in}{0.969093in}}%
\pgfpathlineto{\pgfqpoint{3.206289in}{0.969093in}}%
\pgfusepath{stroke}%
\end{pgfscope}%
\begin{pgfscope}%
\pgfpathmoveto{\pgfqpoint{2.326000in}{0.100000in}}%
\pgfpathcurveto{\pgfqpoint{2.615602in}{0.100000in}}{\pgfqpoint{2.893381in}{0.157530in}}{\pgfqpoint{3.098161in}{0.259920in}}%
\pgfpathcurveto{\pgfqpoint{3.302940in}{0.362309in}}{\pgfqpoint{3.418000in}{0.501199in}}{\pgfqpoint{3.418000in}{0.646000in}}%
\pgfpathcurveto{\pgfqpoint{3.418000in}{0.790801in}}{\pgfqpoint{3.302940in}{0.929691in}}{\pgfqpoint{3.098161in}{1.032080in}}%
\pgfpathcurveto{\pgfqpoint{2.893381in}{1.134470in}}{\pgfqpoint{2.615602in}{1.192000in}}{\pgfqpoint{2.326000in}{1.192000in}}%
\pgfpathcurveto{\pgfqpoint{2.036398in}{1.192000in}}{\pgfqpoint{1.758619in}{1.134470in}}{\pgfqpoint{1.553839in}{1.032080in}}%
\pgfpathcurveto{\pgfqpoint{1.349060in}{0.929691in}}{\pgfqpoint{1.234000in}{0.790801in}}{\pgfqpoint{1.234000in}{0.646000in}}%
\pgfpathcurveto{\pgfqpoint{1.234000in}{0.501199in}}{\pgfqpoint{1.349060in}{0.362309in}}{\pgfqpoint{1.553839in}{0.259920in}}%
\pgfpathcurveto{\pgfqpoint{1.758619in}{0.157530in}}{\pgfqpoint{2.036398in}{0.100000in}}{\pgfqpoint{2.326000in}{0.100000in}}%
\pgfpathlineto{\pgfqpoint{2.326000in}{0.100000in}}%
\pgfpathclose%
\pgfusepath{clip}%
\pgfsetrectcap%
\pgfsetroundjoin%
\pgfsetlinewidth{0.451687pt}%
\definecolor{currentstroke}{rgb}{0.690196,0.690196,0.690196}%
\pgfsetstrokecolor{currentstroke}%
\pgfsetstrokeopacity{0.250000}%
\pgfsetdash{}{0pt}%
\pgfpathmoveto{\pgfqpoint{1.619329in}{1.062257in}}%
\pgfpathlineto{\pgfqpoint{1.638174in}{1.062257in}}%
\pgfpathlineto{\pgfqpoint{1.657018in}{1.062257in}}%
\pgfpathlineto{\pgfqpoint{1.675863in}{1.062257in}}%
\pgfpathlineto{\pgfqpoint{1.694707in}{1.062257in}}%
\pgfpathlineto{\pgfqpoint{1.713552in}{1.062257in}}%
\pgfpathlineto{\pgfqpoint{1.732396in}{1.062257in}}%
\pgfpathlineto{\pgfqpoint{1.751241in}{1.062257in}}%
\pgfpathlineto{\pgfqpoint{1.770086in}{1.062257in}}%
\pgfpathlineto{\pgfqpoint{1.788930in}{1.062257in}}%
\pgfpathlineto{\pgfqpoint{1.807775in}{1.062257in}}%
\pgfpathlineto{\pgfqpoint{1.826619in}{1.062257in}}%
\pgfpathlineto{\pgfqpoint{1.845464in}{1.062257in}}%
\pgfpathlineto{\pgfqpoint{1.864308in}{1.062257in}}%
\pgfpathlineto{\pgfqpoint{1.883153in}{1.062257in}}%
\pgfpathlineto{\pgfqpoint{1.901997in}{1.062257in}}%
\pgfpathlineto{\pgfqpoint{1.920842in}{1.062257in}}%
\pgfpathlineto{\pgfqpoint{1.939687in}{1.062257in}}%
\pgfpathlineto{\pgfqpoint{1.958531in}{1.062257in}}%
\pgfpathlineto{\pgfqpoint{1.977376in}{1.062257in}}%
\pgfpathlineto{\pgfqpoint{1.996220in}{1.062257in}}%
\pgfpathlineto{\pgfqpoint{2.015065in}{1.062257in}}%
\pgfpathlineto{\pgfqpoint{2.033909in}{1.062257in}}%
\pgfpathlineto{\pgfqpoint{2.052754in}{1.062257in}}%
\pgfpathlineto{\pgfqpoint{2.071598in}{1.062257in}}%
\pgfpathlineto{\pgfqpoint{2.090443in}{1.062257in}}%
\pgfpathlineto{\pgfqpoint{2.109288in}{1.062257in}}%
\pgfpathlineto{\pgfqpoint{2.128132in}{1.062257in}}%
\pgfpathlineto{\pgfqpoint{2.146977in}{1.062257in}}%
\pgfpathlineto{\pgfqpoint{2.165821in}{1.062257in}}%
\pgfpathlineto{\pgfqpoint{2.184666in}{1.062257in}}%
\pgfpathlineto{\pgfqpoint{2.203510in}{1.062257in}}%
\pgfpathlineto{\pgfqpoint{2.222355in}{1.062257in}}%
\pgfpathlineto{\pgfqpoint{2.241199in}{1.062257in}}%
\pgfpathlineto{\pgfqpoint{2.260044in}{1.062257in}}%
\pgfpathlineto{\pgfqpoint{2.278889in}{1.062257in}}%
\pgfpathlineto{\pgfqpoint{2.297733in}{1.062257in}}%
\pgfpathlineto{\pgfqpoint{2.316578in}{1.062257in}}%
\pgfpathlineto{\pgfqpoint{2.335422in}{1.062257in}}%
\pgfpathlineto{\pgfqpoint{2.354267in}{1.062257in}}%
\pgfpathlineto{\pgfqpoint{2.373111in}{1.062257in}}%
\pgfpathlineto{\pgfqpoint{2.391956in}{1.062257in}}%
\pgfpathlineto{\pgfqpoint{2.410801in}{1.062257in}}%
\pgfpathlineto{\pgfqpoint{2.429645in}{1.062257in}}%
\pgfpathlineto{\pgfqpoint{2.448490in}{1.062257in}}%
\pgfpathlineto{\pgfqpoint{2.467334in}{1.062257in}}%
\pgfpathlineto{\pgfqpoint{2.486179in}{1.062257in}}%
\pgfpathlineto{\pgfqpoint{2.505023in}{1.062257in}}%
\pgfpathlineto{\pgfqpoint{2.523868in}{1.062257in}}%
\pgfpathlineto{\pgfqpoint{2.542712in}{1.062257in}}%
\pgfpathlineto{\pgfqpoint{2.561557in}{1.062257in}}%
\pgfpathlineto{\pgfqpoint{2.580402in}{1.062257in}}%
\pgfpathlineto{\pgfqpoint{2.599246in}{1.062257in}}%
\pgfpathlineto{\pgfqpoint{2.618091in}{1.062257in}}%
\pgfpathlineto{\pgfqpoint{2.636935in}{1.062257in}}%
\pgfpathlineto{\pgfqpoint{2.655780in}{1.062257in}}%
\pgfpathlineto{\pgfqpoint{2.674624in}{1.062257in}}%
\pgfpathlineto{\pgfqpoint{2.693469in}{1.062257in}}%
\pgfpathlineto{\pgfqpoint{2.712313in}{1.062257in}}%
\pgfpathlineto{\pgfqpoint{2.731158in}{1.062257in}}%
\pgfpathlineto{\pgfqpoint{2.750003in}{1.062257in}}%
\pgfpathlineto{\pgfqpoint{2.768847in}{1.062257in}}%
\pgfpathlineto{\pgfqpoint{2.787692in}{1.062257in}}%
\pgfpathlineto{\pgfqpoint{2.806536in}{1.062257in}}%
\pgfpathlineto{\pgfqpoint{2.825381in}{1.062257in}}%
\pgfpathlineto{\pgfqpoint{2.844225in}{1.062257in}}%
\pgfpathlineto{\pgfqpoint{2.863070in}{1.062257in}}%
\pgfpathlineto{\pgfqpoint{2.881914in}{1.062257in}}%
\pgfpathlineto{\pgfqpoint{2.900759in}{1.062257in}}%
\pgfpathlineto{\pgfqpoint{2.919604in}{1.062257in}}%
\pgfpathlineto{\pgfqpoint{2.938448in}{1.062257in}}%
\pgfpathlineto{\pgfqpoint{2.957293in}{1.062257in}}%
\pgfpathlineto{\pgfqpoint{2.976137in}{1.062257in}}%
\pgfpathlineto{\pgfqpoint{2.994982in}{1.062257in}}%
\pgfpathlineto{\pgfqpoint{3.013826in}{1.062257in}}%
\pgfpathlineto{\pgfqpoint{3.032671in}{1.062257in}}%
\pgfusepath{stroke}%
\end{pgfscope}%
\begin{pgfscope}%
\pgfpathmoveto{\pgfqpoint{2.326000in}{0.100000in}}%
\pgfpathcurveto{\pgfqpoint{2.615602in}{0.100000in}}{\pgfqpoint{2.893381in}{0.157530in}}{\pgfqpoint{3.098161in}{0.259920in}}%
\pgfpathcurveto{\pgfqpoint{3.302940in}{0.362309in}}{\pgfqpoint{3.418000in}{0.501199in}}{\pgfqpoint{3.418000in}{0.646000in}}%
\pgfpathcurveto{\pgfqpoint{3.418000in}{0.790801in}}{\pgfqpoint{3.302940in}{0.929691in}}{\pgfqpoint{3.098161in}{1.032080in}}%
\pgfpathcurveto{\pgfqpoint{2.893381in}{1.134470in}}{\pgfqpoint{2.615602in}{1.192000in}}{\pgfqpoint{2.326000in}{1.192000in}}%
\pgfpathcurveto{\pgfqpoint{2.036398in}{1.192000in}}{\pgfqpoint{1.758619in}{1.134470in}}{\pgfqpoint{1.553839in}{1.032080in}}%
\pgfpathcurveto{\pgfqpoint{1.349060in}{0.929691in}}{\pgfqpoint{1.234000in}{0.790801in}}{\pgfqpoint{1.234000in}{0.646000in}}%
\pgfpathcurveto{\pgfqpoint{1.234000in}{0.501199in}}{\pgfqpoint{1.349060in}{0.362309in}}{\pgfqpoint{1.553839in}{0.259920in}}%
\pgfpathcurveto{\pgfqpoint{1.758619in}{0.157530in}}{\pgfqpoint{2.036398in}{0.100000in}}{\pgfqpoint{2.326000in}{0.100000in}}%
\pgfpathlineto{\pgfqpoint{2.326000in}{0.100000in}}%
\pgfpathclose%
\pgfusepath{clip}%
\pgfsetrectcap%
\pgfsetroundjoin%
\pgfsetlinewidth{0.451687pt}%
\definecolor{currentstroke}{rgb}{0.690196,0.690196,0.690196}%
\pgfsetstrokecolor{currentstroke}%
\pgfsetstrokeopacity{0.250000}%
\pgfsetdash{}{0pt}%
\pgfpathmoveto{\pgfqpoint{1.863906in}{1.140705in}}%
\pgfpathlineto{\pgfqpoint{1.876228in}{1.140705in}}%
\pgfpathlineto{\pgfqpoint{1.888551in}{1.140705in}}%
\pgfpathlineto{\pgfqpoint{1.900873in}{1.140705in}}%
\pgfpathlineto{\pgfqpoint{1.913196in}{1.140705in}}%
\pgfpathlineto{\pgfqpoint{1.925518in}{1.140705in}}%
\pgfpathlineto{\pgfqpoint{1.937841in}{1.140705in}}%
\pgfpathlineto{\pgfqpoint{1.950163in}{1.140705in}}%
\pgfpathlineto{\pgfqpoint{1.962486in}{1.140705in}}%
\pgfpathlineto{\pgfqpoint{1.974808in}{1.140705in}}%
\pgfpathlineto{\pgfqpoint{1.987131in}{1.140705in}}%
\pgfpathlineto{\pgfqpoint{1.999453in}{1.140705in}}%
\pgfpathlineto{\pgfqpoint{2.011776in}{1.140705in}}%
\pgfpathlineto{\pgfqpoint{2.024098in}{1.140705in}}%
\pgfpathlineto{\pgfqpoint{2.036421in}{1.140705in}}%
\pgfpathlineto{\pgfqpoint{2.048743in}{1.140705in}}%
\pgfpathlineto{\pgfqpoint{2.061066in}{1.140705in}}%
\pgfpathlineto{\pgfqpoint{2.073388in}{1.140705in}}%
\pgfpathlineto{\pgfqpoint{2.085711in}{1.140705in}}%
\pgfpathlineto{\pgfqpoint{2.098033in}{1.140705in}}%
\pgfpathlineto{\pgfqpoint{2.110356in}{1.140705in}}%
\pgfpathlineto{\pgfqpoint{2.122678in}{1.140705in}}%
\pgfpathlineto{\pgfqpoint{2.135001in}{1.140705in}}%
\pgfpathlineto{\pgfqpoint{2.147323in}{1.140705in}}%
\pgfpathlineto{\pgfqpoint{2.159646in}{1.140705in}}%
\pgfpathlineto{\pgfqpoint{2.171969in}{1.140705in}}%
\pgfpathlineto{\pgfqpoint{2.184291in}{1.140705in}}%
\pgfpathlineto{\pgfqpoint{2.196614in}{1.140705in}}%
\pgfpathlineto{\pgfqpoint{2.208936in}{1.140705in}}%
\pgfpathlineto{\pgfqpoint{2.221259in}{1.140705in}}%
\pgfpathlineto{\pgfqpoint{2.233581in}{1.140705in}}%
\pgfpathlineto{\pgfqpoint{2.245904in}{1.140705in}}%
\pgfpathlineto{\pgfqpoint{2.258226in}{1.140705in}}%
\pgfpathlineto{\pgfqpoint{2.270549in}{1.140705in}}%
\pgfpathlineto{\pgfqpoint{2.282871in}{1.140705in}}%
\pgfpathlineto{\pgfqpoint{2.295194in}{1.140705in}}%
\pgfpathlineto{\pgfqpoint{2.307516in}{1.140705in}}%
\pgfpathlineto{\pgfqpoint{2.319839in}{1.140705in}}%
\pgfpathlineto{\pgfqpoint{2.332161in}{1.140705in}}%
\pgfpathlineto{\pgfqpoint{2.344484in}{1.140705in}}%
\pgfpathlineto{\pgfqpoint{2.356806in}{1.140705in}}%
\pgfpathlineto{\pgfqpoint{2.369129in}{1.140705in}}%
\pgfpathlineto{\pgfqpoint{2.381451in}{1.140705in}}%
\pgfpathlineto{\pgfqpoint{2.393774in}{1.140705in}}%
\pgfpathlineto{\pgfqpoint{2.406096in}{1.140705in}}%
\pgfpathlineto{\pgfqpoint{2.418419in}{1.140705in}}%
\pgfpathlineto{\pgfqpoint{2.430741in}{1.140705in}}%
\pgfpathlineto{\pgfqpoint{2.443064in}{1.140705in}}%
\pgfpathlineto{\pgfqpoint{2.455386in}{1.140705in}}%
\pgfpathlineto{\pgfqpoint{2.467709in}{1.140705in}}%
\pgfpathlineto{\pgfqpoint{2.480031in}{1.140705in}}%
\pgfpathlineto{\pgfqpoint{2.492354in}{1.140705in}}%
\pgfpathlineto{\pgfqpoint{2.504677in}{1.140705in}}%
\pgfpathlineto{\pgfqpoint{2.516999in}{1.140705in}}%
\pgfpathlineto{\pgfqpoint{2.529322in}{1.140705in}}%
\pgfpathlineto{\pgfqpoint{2.541644in}{1.140705in}}%
\pgfpathlineto{\pgfqpoint{2.553967in}{1.140705in}}%
\pgfpathlineto{\pgfqpoint{2.566289in}{1.140705in}}%
\pgfpathlineto{\pgfqpoint{2.578612in}{1.140705in}}%
\pgfpathlineto{\pgfqpoint{2.590934in}{1.140705in}}%
\pgfpathlineto{\pgfqpoint{2.603257in}{1.140705in}}%
\pgfpathlineto{\pgfqpoint{2.615579in}{1.140705in}}%
\pgfpathlineto{\pgfqpoint{2.627902in}{1.140705in}}%
\pgfpathlineto{\pgfqpoint{2.640224in}{1.140705in}}%
\pgfpathlineto{\pgfqpoint{2.652547in}{1.140705in}}%
\pgfpathlineto{\pgfqpoint{2.664869in}{1.140705in}}%
\pgfpathlineto{\pgfqpoint{2.677192in}{1.140705in}}%
\pgfpathlineto{\pgfqpoint{2.689514in}{1.140705in}}%
\pgfpathlineto{\pgfqpoint{2.701837in}{1.140705in}}%
\pgfpathlineto{\pgfqpoint{2.714159in}{1.140705in}}%
\pgfpathlineto{\pgfqpoint{2.726482in}{1.140705in}}%
\pgfpathlineto{\pgfqpoint{2.738804in}{1.140705in}}%
\pgfpathlineto{\pgfqpoint{2.751127in}{1.140705in}}%
\pgfpathlineto{\pgfqpoint{2.763449in}{1.140705in}}%
\pgfpathlineto{\pgfqpoint{2.775772in}{1.140705in}}%
\pgfpathlineto{\pgfqpoint{2.788094in}{1.140705in}}%
\pgfusepath{stroke}%
\end{pgfscope}%
\begin{pgfscope}%
\pgfsetrectcap%
\pgfsetmiterjoin%
\pgfsetlinewidth{0.803000pt}%
\definecolor{currentstroke}{rgb}{0.000000,0.000000,0.000000}%
\pgfsetstrokecolor{currentstroke}%
\pgfsetdash{}{0pt}%
\pgfpathmoveto{\pgfqpoint{2.326000in}{0.100000in}}%
\pgfpathcurveto{\pgfqpoint{2.615602in}{0.100000in}}{\pgfqpoint{2.893381in}{0.157530in}}{\pgfqpoint{3.098161in}{0.259920in}}%
\pgfpathcurveto{\pgfqpoint{3.302940in}{0.362309in}}{\pgfqpoint{3.418000in}{0.501199in}}{\pgfqpoint{3.418000in}{0.646000in}}%
\pgfpathcurveto{\pgfqpoint{3.418000in}{0.790801in}}{\pgfqpoint{3.302940in}{0.929691in}}{\pgfqpoint{3.098161in}{1.032080in}}%
\pgfpathcurveto{\pgfqpoint{2.893381in}{1.134470in}}{\pgfqpoint{2.615602in}{1.192000in}}{\pgfqpoint{2.326000in}{1.192000in}}%
\pgfpathcurveto{\pgfqpoint{2.036398in}{1.192000in}}{\pgfqpoint{1.758619in}{1.134470in}}{\pgfqpoint{1.553839in}{1.032080in}}%
\pgfpathcurveto{\pgfqpoint{1.349060in}{0.929691in}}{\pgfqpoint{1.234000in}{0.790801in}}{\pgfqpoint{1.234000in}{0.646000in}}%
\pgfpathcurveto{\pgfqpoint{1.234000in}{0.501199in}}{\pgfqpoint{1.349060in}{0.362309in}}{\pgfqpoint{1.553839in}{0.259920in}}%
\pgfpathcurveto{\pgfqpoint{1.758619in}{0.157530in}}{\pgfqpoint{2.036398in}{0.100000in}}{\pgfqpoint{2.326000in}{0.100000in}}%
\pgfpathlineto{\pgfqpoint{2.326000in}{0.100000in}}%
\pgfpathclose%
\pgfusepath{stroke}%
\end{pgfscope}%
\begin{pgfscope}%
\definecolor{textcolor}{rgb}{0.000000,0.000000,0.000000}%
\pgfsetstrokecolor{textcolor}%
\pgfsetfillcolor{textcolor}%
\pgftext[x=2.326000in,y=1.233667in,,base]{\color{textcolor}{\rmfamily\fontsize{8.200000}{9.840000}\selectfont\catcode`\^=\active\def^{\ifmmode\sp\else\^{}\fi}\catcode`\%=\active\def%{\%}Four Modes}}%
\end{pgfscope}%
\begin{pgfscope}%
\pgfsetbuttcap%
\pgfsetmiterjoin%
\definecolor{currentfill}{rgb}{1.000000,1.000000,1.000000}%
\pgfsetfillcolor{currentfill}%
\pgfsetlinewidth{0.000000pt}%
\definecolor{currentstroke}{rgb}{0.000000,0.000000,0.000000}%
\pgfsetstrokecolor{currentstroke}%
\pgfsetstrokeopacity{0.000000}%
\pgfsetdash{}{0pt}%
\pgfpathmoveto{\pgfqpoint{4.594000in}{0.100000in}}%
\pgfpathcurveto{\pgfqpoint{4.883602in}{0.100000in}}{\pgfqpoint{5.161381in}{0.157530in}}{\pgfqpoint{5.366161in}{0.259920in}}%
\pgfpathcurveto{\pgfqpoint{5.570940in}{0.362309in}}{\pgfqpoint{5.686000in}{0.501199in}}{\pgfqpoint{5.686000in}{0.646000in}}%
\pgfpathcurveto{\pgfqpoint{5.686000in}{0.790801in}}{\pgfqpoint{5.570940in}{0.929691in}}{\pgfqpoint{5.366161in}{1.032080in}}%
\pgfpathcurveto{\pgfqpoint{5.161381in}{1.134470in}}{\pgfqpoint{4.883602in}{1.192000in}}{\pgfqpoint{4.594000in}{1.192000in}}%
\pgfpathcurveto{\pgfqpoint{4.304398in}{1.192000in}}{\pgfqpoint{4.026619in}{1.134470in}}{\pgfqpoint{3.821839in}{1.032080in}}%
\pgfpathcurveto{\pgfqpoint{3.617060in}{0.929691in}}{\pgfqpoint{3.502000in}{0.790801in}}{\pgfqpoint{3.502000in}{0.646000in}}%
\pgfpathcurveto{\pgfqpoint{3.502000in}{0.501199in}}{\pgfqpoint{3.617060in}{0.362309in}}{\pgfqpoint{3.821839in}{0.259920in}}%
\pgfpathcurveto{\pgfqpoint{4.026619in}{0.157530in}}{\pgfqpoint{4.304398in}{0.100000in}}{\pgfqpoint{4.594000in}{0.100000in}}%
\pgfpathlineto{\pgfqpoint{4.594000in}{0.100000in}}%
\pgfpathclose%
\pgfusepath{fill}%
\end{pgfscope}%
\begin{pgfscope}%
\pgfpathmoveto{\pgfqpoint{4.594000in}{0.100000in}}%
\pgfpathcurveto{\pgfqpoint{4.883602in}{0.100000in}}{\pgfqpoint{5.161381in}{0.157530in}}{\pgfqpoint{5.366161in}{0.259920in}}%
\pgfpathcurveto{\pgfqpoint{5.570940in}{0.362309in}}{\pgfqpoint{5.686000in}{0.501199in}}{\pgfqpoint{5.686000in}{0.646000in}}%
\pgfpathcurveto{\pgfqpoint{5.686000in}{0.790801in}}{\pgfqpoint{5.570940in}{0.929691in}}{\pgfqpoint{5.366161in}{1.032080in}}%
\pgfpathcurveto{\pgfqpoint{5.161381in}{1.134470in}}{\pgfqpoint{4.883602in}{1.192000in}}{\pgfqpoint{4.594000in}{1.192000in}}%
\pgfpathcurveto{\pgfqpoint{4.304398in}{1.192000in}}{\pgfqpoint{4.026619in}{1.134470in}}{\pgfqpoint{3.821839in}{1.032080in}}%
\pgfpathcurveto{\pgfqpoint{3.617060in}{0.929691in}}{\pgfqpoint{3.502000in}{0.790801in}}{\pgfqpoint{3.502000in}{0.646000in}}%
\pgfpathcurveto{\pgfqpoint{3.502000in}{0.501199in}}{\pgfqpoint{3.617060in}{0.362309in}}{\pgfqpoint{3.821839in}{0.259920in}}%
\pgfpathcurveto{\pgfqpoint{4.026619in}{0.157530in}}{\pgfqpoint{4.304398in}{0.100000in}}{\pgfqpoint{4.594000in}{0.100000in}}%
\pgfpathlineto{\pgfqpoint{4.594000in}{0.100000in}}%
\pgfpathclose%
\pgfusepath{clip}%
\pgfsetbuttcap%
\pgfsetroundjoin%
\definecolor{currentfill}{rgb}{0.121569,0.466667,0.705882}%
\pgfsetfillcolor{currentfill}%
\pgfsetfillopacity{0.580000}%
\pgfsetlinewidth{0.000000pt}%
\definecolor{currentstroke}{rgb}{0.121569,0.466667,0.705882}%
\pgfsetstrokecolor{currentstroke}%
\pgfsetstrokeopacity{0.580000}%
\pgfsetdash{}{0pt}%
\pgfsys@defobject{currentmarker}{\pgfqpoint{-0.014232in}{-0.014232in}}{\pgfqpoint{0.014232in}{0.014232in}}{%
\pgfpathmoveto{\pgfqpoint{0.000000in}{-0.014232in}}%
\pgfpathcurveto{\pgfqpoint{0.003774in}{-0.014232in}}{\pgfqpoint{0.007395in}{-0.012732in}}{\pgfqpoint{0.010063in}{-0.010063in}}%
\pgfpathcurveto{\pgfqpoint{0.012732in}{-0.007395in}}{\pgfqpoint{0.014232in}{-0.003774in}}{\pgfqpoint{0.014232in}{0.000000in}}%
\pgfpathcurveto{\pgfqpoint{0.014232in}{0.003774in}}{\pgfqpoint{0.012732in}{0.007395in}}{\pgfqpoint{0.010063in}{0.010063in}}%
\pgfpathcurveto{\pgfqpoint{0.007395in}{0.012732in}}{\pgfqpoint{0.003774in}{0.014232in}}{\pgfqpoint{0.000000in}{0.014232in}}%
\pgfpathcurveto{\pgfqpoint{-0.003774in}{0.014232in}}{\pgfqpoint{-0.007395in}{0.012732in}}{\pgfqpoint{-0.010063in}{0.010063in}}%
\pgfpathcurveto{\pgfqpoint{-0.012732in}{0.007395in}}{\pgfqpoint{-0.014232in}{0.003774in}}{\pgfqpoint{-0.014232in}{0.000000in}}%
\pgfpathcurveto{\pgfqpoint{-0.014232in}{-0.003774in}}{\pgfqpoint{-0.012732in}{-0.007395in}}{\pgfqpoint{-0.010063in}{-0.010063in}}%
\pgfpathcurveto{\pgfqpoint{-0.007395in}{-0.012732in}}{\pgfqpoint{-0.003774in}{-0.014232in}}{\pgfqpoint{0.000000in}{-0.014232in}}%
\pgfpathlineto{\pgfqpoint{0.000000in}{-0.014232in}}%
\pgfpathclose%
\pgfusepath{fill}%
}%
\begin{pgfscope}%
\pgfsys@transformshift{4.594000in}{1.141458in}%
\pgfsys@useobject{currentmarker}{}%
\end{pgfscope}%
\begin{pgfscope}%
\pgfsys@transformshift{4.594000in}{1.141458in}%
\pgfsys@useobject{currentmarker}{}%
\end{pgfscope}%
\begin{pgfscope}%
\pgfsys@transformshift{4.594000in}{1.141458in}%
\pgfsys@useobject{currentmarker}{}%
\end{pgfscope}%
\begin{pgfscope}%
\pgfsys@transformshift{4.594000in}{1.141458in}%
\pgfsys@useobject{currentmarker}{}%
\end{pgfscope}%
\begin{pgfscope}%
\pgfsys@transformshift{4.594000in}{1.141458in}%
\pgfsys@useobject{currentmarker}{}%
\end{pgfscope}%
\begin{pgfscope}%
\pgfsys@transformshift{4.594000in}{1.141458in}%
\pgfsys@useobject{currentmarker}{}%
\end{pgfscope}%
\begin{pgfscope}%
\pgfsys@transformshift{4.594000in}{1.141458in}%
\pgfsys@useobject{currentmarker}{}%
\end{pgfscope}%
\begin{pgfscope}%
\pgfsys@transformshift{4.594000in}{1.141458in}%
\pgfsys@useobject{currentmarker}{}%
\end{pgfscope}%
\begin{pgfscope}%
\pgfsys@transformshift{4.594000in}{1.141458in}%
\pgfsys@useobject{currentmarker}{}%
\end{pgfscope}%
\begin{pgfscope}%
\pgfsys@transformshift{4.594000in}{1.141458in}%
\pgfsys@useobject{currentmarker}{}%
\end{pgfscope}%
\begin{pgfscope}%
\pgfsys@transformshift{5.088522in}{1.085721in}%
\pgfsys@useobject{currentmarker}{}%
\end{pgfscope}%
\begin{pgfscope}%
\pgfsys@transformshift{5.088522in}{1.085721in}%
\pgfsys@useobject{currentmarker}{}%
\end{pgfscope}%
\begin{pgfscope}%
\pgfsys@transformshift{5.088522in}{1.085721in}%
\pgfsys@useobject{currentmarker}{}%
\end{pgfscope}%
\begin{pgfscope}%
\pgfsys@transformshift{5.088522in}{1.085721in}%
\pgfsys@useobject{currentmarker}{}%
\end{pgfscope}%
\begin{pgfscope}%
\pgfsys@transformshift{5.088522in}{1.085721in}%
\pgfsys@useobject{currentmarker}{}%
\end{pgfscope}%
\begin{pgfscope}%
\pgfsys@transformshift{5.088522in}{1.085721in}%
\pgfsys@useobject{currentmarker}{}%
\end{pgfscope}%
\begin{pgfscope}%
\pgfsys@transformshift{5.088522in}{1.085721in}%
\pgfsys@useobject{currentmarker}{}%
\end{pgfscope}%
\begin{pgfscope}%
\pgfsys@transformshift{5.088522in}{1.085721in}%
\pgfsys@useobject{currentmarker}{}%
\end{pgfscope}%
\begin{pgfscope}%
\pgfsys@transformshift{5.088522in}{1.085721in}%
\pgfsys@useobject{currentmarker}{}%
\end{pgfscope}%
\begin{pgfscope}%
\pgfsys@transformshift{5.088522in}{1.085721in}%
\pgfsys@useobject{currentmarker}{}%
\end{pgfscope}%
\begin{pgfscope}%
\pgfsys@transformshift{4.238355in}{1.041301in}%
\pgfsys@useobject{currentmarker}{}%
\end{pgfscope}%
\begin{pgfscope}%
\pgfsys@transformshift{4.238355in}{1.041301in}%
\pgfsys@useobject{currentmarker}{}%
\end{pgfscope}%
\begin{pgfscope}%
\pgfsys@transformshift{4.238355in}{1.041301in}%
\pgfsys@useobject{currentmarker}{}%
\end{pgfscope}%
\begin{pgfscope}%
\pgfsys@transformshift{4.238355in}{1.041301in}%
\pgfsys@useobject{currentmarker}{}%
\end{pgfscope}%
\begin{pgfscope}%
\pgfsys@transformshift{4.238355in}{1.041301in}%
\pgfsys@useobject{currentmarker}{}%
\end{pgfscope}%
\begin{pgfscope}%
\pgfsys@transformshift{4.238355in}{1.041301in}%
\pgfsys@useobject{currentmarker}{}%
\end{pgfscope}%
\begin{pgfscope}%
\pgfsys@transformshift{4.238355in}{1.041301in}%
\pgfsys@useobject{currentmarker}{}%
\end{pgfscope}%
\begin{pgfscope}%
\pgfsys@transformshift{4.238355in}{1.041301in}%
\pgfsys@useobject{currentmarker}{}%
\end{pgfscope}%
\begin{pgfscope}%
\pgfsys@transformshift{4.238355in}{1.041301in}%
\pgfsys@useobject{currentmarker}{}%
\end{pgfscope}%
\begin{pgfscope}%
\pgfsys@transformshift{4.238355in}{1.041301in}%
\pgfsys@useobject{currentmarker}{}%
\end{pgfscope}%
\begin{pgfscope}%
\pgfsys@transformshift{4.835654in}{1.001883in}%
\pgfsys@useobject{currentmarker}{}%
\end{pgfscope}%
\begin{pgfscope}%
\pgfsys@transformshift{4.835654in}{1.001883in}%
\pgfsys@useobject{currentmarker}{}%
\end{pgfscope}%
\begin{pgfscope}%
\pgfsys@transformshift{4.835654in}{1.001883in}%
\pgfsys@useobject{currentmarker}{}%
\end{pgfscope}%
\begin{pgfscope}%
\pgfsys@transformshift{4.835654in}{1.001883in}%
\pgfsys@useobject{currentmarker}{}%
\end{pgfscope}%
\begin{pgfscope}%
\pgfsys@transformshift{4.835654in}{1.001883in}%
\pgfsys@useobject{currentmarker}{}%
\end{pgfscope}%
\begin{pgfscope}%
\pgfsys@transformshift{4.835654in}{1.001883in}%
\pgfsys@useobject{currentmarker}{}%
\end{pgfscope}%
\begin{pgfscope}%
\pgfsys@transformshift{4.835654in}{1.001883in}%
\pgfsys@useobject{currentmarker}{}%
\end{pgfscope}%
\begin{pgfscope}%
\pgfsys@transformshift{4.835654in}{1.001883in}%
\pgfsys@useobject{currentmarker}{}%
\end{pgfscope}%
\begin{pgfscope}%
\pgfsys@transformshift{4.835654in}{1.001883in}%
\pgfsys@useobject{currentmarker}{}%
\end{pgfscope}%
\begin{pgfscope}%
\pgfsys@transformshift{4.835654in}{1.001883in}%
\pgfsys@useobject{currentmarker}{}%
\end{pgfscope}%
\begin{pgfscope}%
\pgfsys@transformshift{3.757608in}{0.965337in}%
\pgfsys@useobject{currentmarker}{}%
\end{pgfscope}%
\begin{pgfscope}%
\pgfsys@transformshift{3.757608in}{0.965337in}%
\pgfsys@useobject{currentmarker}{}%
\end{pgfscope}%
\begin{pgfscope}%
\pgfsys@transformshift{3.757608in}{0.965337in}%
\pgfsys@useobject{currentmarker}{}%
\end{pgfscope}%
\begin{pgfscope}%
\pgfsys@transformshift{3.757608in}{0.965337in}%
\pgfsys@useobject{currentmarker}{}%
\end{pgfscope}%
\begin{pgfscope}%
\pgfsys@transformshift{3.757608in}{0.965337in}%
\pgfsys@useobject{currentmarker}{}%
\end{pgfscope}%
\begin{pgfscope}%
\pgfsys@transformshift{3.757608in}{0.965337in}%
\pgfsys@useobject{currentmarker}{}%
\end{pgfscope}%
\begin{pgfscope}%
\pgfsys@transformshift{3.757608in}{0.965337in}%
\pgfsys@useobject{currentmarker}{}%
\end{pgfscope}%
\begin{pgfscope}%
\pgfsys@transformshift{3.757608in}{0.965337in}%
\pgfsys@useobject{currentmarker}{}%
\end{pgfscope}%
\begin{pgfscope}%
\pgfsys@transformshift{3.757608in}{0.965337in}%
\pgfsys@useobject{currentmarker}{}%
\end{pgfscope}%
\begin{pgfscope}%
\pgfsys@transformshift{3.757608in}{0.965337in}%
\pgfsys@useobject{currentmarker}{}%
\end{pgfscope}%
\begin{pgfscope}%
\pgfsys@transformshift{4.426044in}{0.931081in}%
\pgfsys@useobject{currentmarker}{}%
\end{pgfscope}%
\begin{pgfscope}%
\pgfsys@transformshift{4.426044in}{0.931081in}%
\pgfsys@useobject{currentmarker}{}%
\end{pgfscope}%
\begin{pgfscope}%
\pgfsys@transformshift{4.426044in}{0.931081in}%
\pgfsys@useobject{currentmarker}{}%
\end{pgfscope}%
\begin{pgfscope}%
\pgfsys@transformshift{4.426044in}{0.931081in}%
\pgfsys@useobject{currentmarker}{}%
\end{pgfscope}%
\begin{pgfscope}%
\pgfsys@transformshift{4.426044in}{0.931081in}%
\pgfsys@useobject{currentmarker}{}%
\end{pgfscope}%
\begin{pgfscope}%
\pgfsys@transformshift{4.426044in}{0.931081in}%
\pgfsys@useobject{currentmarker}{}%
\end{pgfscope}%
\begin{pgfscope}%
\pgfsys@transformshift{4.426044in}{0.931081in}%
\pgfsys@useobject{currentmarker}{}%
\end{pgfscope}%
\begin{pgfscope}%
\pgfsys@transformshift{4.426044in}{0.931081in}%
\pgfsys@useobject{currentmarker}{}%
\end{pgfscope}%
\begin{pgfscope}%
\pgfsys@transformshift{4.426044in}{0.931081in}%
\pgfsys@useobject{currentmarker}{}%
\end{pgfscope}%
\begin{pgfscope}%
\pgfsys@transformshift{4.426044in}{0.931081in}%
\pgfsys@useobject{currentmarker}{}%
\end{pgfscope}%
\begin{pgfscope}%
\pgfsys@transformshift{5.159183in}{0.898268in}%
\pgfsys@useobject{currentmarker}{}%
\end{pgfscope}%
\begin{pgfscope}%
\pgfsys@transformshift{5.159183in}{0.898268in}%
\pgfsys@useobject{currentmarker}{}%
\end{pgfscope}%
\begin{pgfscope}%
\pgfsys@transformshift{5.159183in}{0.898268in}%
\pgfsys@useobject{currentmarker}{}%
\end{pgfscope}%
\begin{pgfscope}%
\pgfsys@transformshift{5.159183in}{0.898268in}%
\pgfsys@useobject{currentmarker}{}%
\end{pgfscope}%
\begin{pgfscope}%
\pgfsys@transformshift{5.159183in}{0.898268in}%
\pgfsys@useobject{currentmarker}{}%
\end{pgfscope}%
\begin{pgfscope}%
\pgfsys@transformshift{5.159183in}{0.898268in}%
\pgfsys@useobject{currentmarker}{}%
\end{pgfscope}%
\begin{pgfscope}%
\pgfsys@transformshift{5.159183in}{0.898268in}%
\pgfsys@useobject{currentmarker}{}%
\end{pgfscope}%
\begin{pgfscope}%
\pgfsys@transformshift{5.159183in}{0.898268in}%
\pgfsys@useobject{currentmarker}{}%
\end{pgfscope}%
\begin{pgfscope}%
\pgfsys@transformshift{5.159183in}{0.898268in}%
\pgfsys@useobject{currentmarker}{}%
\end{pgfscope}%
\begin{pgfscope}%
\pgfsys@transformshift{5.159183in}{0.898268in}%
\pgfsys@useobject{currentmarker}{}%
\end{pgfscope}%
\begin{pgfscope}%
\pgfsys@transformshift{3.942217in}{0.866560in}%
\pgfsys@useobject{currentmarker}{}%
\end{pgfscope}%
\begin{pgfscope}%
\pgfsys@transformshift{3.942217in}{0.866560in}%
\pgfsys@useobject{currentmarker}{}%
\end{pgfscope}%
\begin{pgfscope}%
\pgfsys@transformshift{3.942217in}{0.866560in}%
\pgfsys@useobject{currentmarker}{}%
\end{pgfscope}%
\begin{pgfscope}%
\pgfsys@transformshift{3.942217in}{0.866560in}%
\pgfsys@useobject{currentmarker}{}%
\end{pgfscope}%
\begin{pgfscope}%
\pgfsys@transformshift{3.942217in}{0.866560in}%
\pgfsys@useobject{currentmarker}{}%
\end{pgfscope}%
\begin{pgfscope}%
\pgfsys@transformshift{3.942217in}{0.866560in}%
\pgfsys@useobject{currentmarker}{}%
\end{pgfscope}%
\begin{pgfscope}%
\pgfsys@transformshift{3.942217in}{0.866560in}%
\pgfsys@useobject{currentmarker}{}%
\end{pgfscope}%
\begin{pgfscope}%
\pgfsys@transformshift{3.942217in}{0.866560in}%
\pgfsys@useobject{currentmarker}{}%
\end{pgfscope}%
\begin{pgfscope}%
\pgfsys@transformshift{3.942217in}{0.866560in}%
\pgfsys@useobject{currentmarker}{}%
\end{pgfscope}%
\begin{pgfscope}%
\pgfsys@transformshift{3.942217in}{0.866560in}%
\pgfsys@useobject{currentmarker}{}%
\end{pgfscope}%
\begin{pgfscope}%
\pgfsys@transformshift{4.708127in}{0.835712in}%
\pgfsys@useobject{currentmarker}{}%
\end{pgfscope}%
\begin{pgfscope}%
\pgfsys@transformshift{4.708127in}{0.835712in}%
\pgfsys@useobject{currentmarker}{}%
\end{pgfscope}%
\begin{pgfscope}%
\pgfsys@transformshift{4.708127in}{0.835712in}%
\pgfsys@useobject{currentmarker}{}%
\end{pgfscope}%
\begin{pgfscope}%
\pgfsys@transformshift{4.708127in}{0.835712in}%
\pgfsys@useobject{currentmarker}{}%
\end{pgfscope}%
\begin{pgfscope}%
\pgfsys@transformshift{4.708127in}{0.835712in}%
\pgfsys@useobject{currentmarker}{}%
\end{pgfscope}%
\begin{pgfscope}%
\pgfsys@transformshift{4.708127in}{0.835712in}%
\pgfsys@useobject{currentmarker}{}%
\end{pgfscope}%
\begin{pgfscope}%
\pgfsys@transformshift{4.708127in}{0.835712in}%
\pgfsys@useobject{currentmarker}{}%
\end{pgfscope}%
\begin{pgfscope}%
\pgfsys@transformshift{4.708127in}{0.835712in}%
\pgfsys@useobject{currentmarker}{}%
\end{pgfscope}%
\begin{pgfscope}%
\pgfsys@transformshift{4.708127in}{0.835712in}%
\pgfsys@useobject{currentmarker}{}%
\end{pgfscope}%
\begin{pgfscope}%
\pgfsys@transformshift{4.708127in}{0.835712in}%
\pgfsys@useobject{currentmarker}{}%
\end{pgfscope}%
\begin{pgfscope}%
\pgfsys@transformshift{5.508208in}{0.805536in}%
\pgfsys@useobject{currentmarker}{}%
\end{pgfscope}%
\begin{pgfscope}%
\pgfsys@transformshift{5.508208in}{0.805536in}%
\pgfsys@useobject{currentmarker}{}%
\end{pgfscope}%
\begin{pgfscope}%
\pgfsys@transformshift{5.508208in}{0.805536in}%
\pgfsys@useobject{currentmarker}{}%
\end{pgfscope}%
\begin{pgfscope}%
\pgfsys@transformshift{5.508208in}{0.805536in}%
\pgfsys@useobject{currentmarker}{}%
\end{pgfscope}%
\begin{pgfscope}%
\pgfsys@transformshift{5.508208in}{0.805536in}%
\pgfsys@useobject{currentmarker}{}%
\end{pgfscope}%
\begin{pgfscope}%
\pgfsys@transformshift{5.508208in}{0.805536in}%
\pgfsys@useobject{currentmarker}{}%
\end{pgfscope}%
\begin{pgfscope}%
\pgfsys@transformshift{5.508208in}{0.805536in}%
\pgfsys@useobject{currentmarker}{}%
\end{pgfscope}%
\begin{pgfscope}%
\pgfsys@transformshift{5.508208in}{0.805536in}%
\pgfsys@useobject{currentmarker}{}%
\end{pgfscope}%
\begin{pgfscope}%
\pgfsys@transformshift{5.508208in}{0.805536in}%
\pgfsys@useobject{currentmarker}{}%
\end{pgfscope}%
\begin{pgfscope}%
\pgfsys@transformshift{5.508208in}{0.805536in}%
\pgfsys@useobject{currentmarker}{}%
\end{pgfscope}%
\begin{pgfscope}%
\pgfsys@transformshift{4.211444in}{0.775884in}%
\pgfsys@useobject{currentmarker}{}%
\end{pgfscope}%
\begin{pgfscope}%
\pgfsys@transformshift{4.211444in}{0.775884in}%
\pgfsys@useobject{currentmarker}{}%
\end{pgfscope}%
\begin{pgfscope}%
\pgfsys@transformshift{4.211444in}{0.775884in}%
\pgfsys@useobject{currentmarker}{}%
\end{pgfscope}%
\begin{pgfscope}%
\pgfsys@transformshift{4.211444in}{0.775884in}%
\pgfsys@useobject{currentmarker}{}%
\end{pgfscope}%
\begin{pgfscope}%
\pgfsys@transformshift{4.211444in}{0.775884in}%
\pgfsys@useobject{currentmarker}{}%
\end{pgfscope}%
\begin{pgfscope}%
\pgfsys@transformshift{4.211444in}{0.775884in}%
\pgfsys@useobject{currentmarker}{}%
\end{pgfscope}%
\begin{pgfscope}%
\pgfsys@transformshift{4.211444in}{0.775884in}%
\pgfsys@useobject{currentmarker}{}%
\end{pgfscope}%
\begin{pgfscope}%
\pgfsys@transformshift{4.211444in}{0.775884in}%
\pgfsys@useobject{currentmarker}{}%
\end{pgfscope}%
\begin{pgfscope}%
\pgfsys@transformshift{4.211444in}{0.775884in}%
\pgfsys@useobject{currentmarker}{}%
\end{pgfscope}%
\begin{pgfscope}%
\pgfsys@transformshift{4.211444in}{0.775884in}%
\pgfsys@useobject{currentmarker}{}%
\end{pgfscope}%
\begin{pgfscope}%
\pgfsys@transformshift{5.026808in}{0.746632in}%
\pgfsys@useobject{currentmarker}{}%
\end{pgfscope}%
\begin{pgfscope}%
\pgfsys@transformshift{5.026808in}{0.746632in}%
\pgfsys@useobject{currentmarker}{}%
\end{pgfscope}%
\begin{pgfscope}%
\pgfsys@transformshift{5.026808in}{0.746632in}%
\pgfsys@useobject{currentmarker}{}%
\end{pgfscope}%
\begin{pgfscope}%
\pgfsys@transformshift{5.026808in}{0.746632in}%
\pgfsys@useobject{currentmarker}{}%
\end{pgfscope}%
\begin{pgfscope}%
\pgfsys@transformshift{5.026808in}{0.746632in}%
\pgfsys@useobject{currentmarker}{}%
\end{pgfscope}%
\begin{pgfscope}%
\pgfsys@transformshift{5.026808in}{0.746632in}%
\pgfsys@useobject{currentmarker}{}%
\end{pgfscope}%
\begin{pgfscope}%
\pgfsys@transformshift{5.026808in}{0.746632in}%
\pgfsys@useobject{currentmarker}{}%
\end{pgfscope}%
\begin{pgfscope}%
\pgfsys@transformshift{5.026808in}{0.746632in}%
\pgfsys@useobject{currentmarker}{}%
\end{pgfscope}%
\begin{pgfscope}%
\pgfsys@transformshift{5.026808in}{0.746632in}%
\pgfsys@useobject{currentmarker}{}%
\end{pgfscope}%
\begin{pgfscope}%
\pgfsys@transformshift{5.026808in}{0.746632in}%
\pgfsys@useobject{currentmarker}{}%
\end{pgfscope}%
\begin{pgfscope}%
\pgfsys@transformshift{3.692412in}{0.717569in}%
\pgfsys@useobject{currentmarker}{}%
\end{pgfscope}%
\begin{pgfscope}%
\pgfsys@transformshift{3.692412in}{0.717569in}%
\pgfsys@useobject{currentmarker}{}%
\end{pgfscope}%
\begin{pgfscope}%
\pgfsys@transformshift{3.692412in}{0.717569in}%
\pgfsys@useobject{currentmarker}{}%
\end{pgfscope}%
\begin{pgfscope}%
\pgfsys@transformshift{3.692412in}{0.717569in}%
\pgfsys@useobject{currentmarker}{}%
\end{pgfscope}%
\begin{pgfscope}%
\pgfsys@transformshift{3.692412in}{0.717569in}%
\pgfsys@useobject{currentmarker}{}%
\end{pgfscope}%
\begin{pgfscope}%
\pgfsys@transformshift{3.692412in}{0.717569in}%
\pgfsys@useobject{currentmarker}{}%
\end{pgfscope}%
\begin{pgfscope}%
\pgfsys@transformshift{3.692412in}{0.717569in}%
\pgfsys@useobject{currentmarker}{}%
\end{pgfscope}%
\begin{pgfscope}%
\pgfsys@transformshift{3.692412in}{0.717569in}%
\pgfsys@useobject{currentmarker}{}%
\end{pgfscope}%
\begin{pgfscope}%
\pgfsys@transformshift{3.692412in}{0.717569in}%
\pgfsys@useobject{currentmarker}{}%
\end{pgfscope}%
\begin{pgfscope}%
\pgfsys@transformshift{3.692412in}{0.717569in}%
\pgfsys@useobject{currentmarker}{}%
\end{pgfscope}%
\begin{pgfscope}%
\pgfsys@transformshift{4.519012in}{0.688904in}%
\pgfsys@useobject{currentmarker}{}%
\end{pgfscope}%
\begin{pgfscope}%
\pgfsys@transformshift{4.519012in}{0.688904in}%
\pgfsys@useobject{currentmarker}{}%
\end{pgfscope}%
\begin{pgfscope}%
\pgfsys@transformshift{4.519012in}{0.688904in}%
\pgfsys@useobject{currentmarker}{}%
\end{pgfscope}%
\begin{pgfscope}%
\pgfsys@transformshift{4.519012in}{0.688904in}%
\pgfsys@useobject{currentmarker}{}%
\end{pgfscope}%
\begin{pgfscope}%
\pgfsys@transformshift{4.519012in}{0.688904in}%
\pgfsys@useobject{currentmarker}{}%
\end{pgfscope}%
\begin{pgfscope}%
\pgfsys@transformshift{4.519012in}{0.688904in}%
\pgfsys@useobject{currentmarker}{}%
\end{pgfscope}%
\begin{pgfscope}%
\pgfsys@transformshift{4.519012in}{0.688904in}%
\pgfsys@useobject{currentmarker}{}%
\end{pgfscope}%
\begin{pgfscope}%
\pgfsys@transformshift{4.519012in}{0.688904in}%
\pgfsys@useobject{currentmarker}{}%
\end{pgfscope}%
\begin{pgfscope}%
\pgfsys@transformshift{4.519012in}{0.688904in}%
\pgfsys@useobject{currentmarker}{}%
\end{pgfscope}%
\begin{pgfscope}%
\pgfsys@transformshift{4.519012in}{0.688904in}%
\pgfsys@useobject{currentmarker}{}%
\end{pgfscope}%
\begin{pgfscope}%
\pgfsys@transformshift{5.352733in}{0.660295in}%
\pgfsys@useobject{currentmarker}{}%
\end{pgfscope}%
\begin{pgfscope}%
\pgfsys@transformshift{5.352733in}{0.660295in}%
\pgfsys@useobject{currentmarker}{}%
\end{pgfscope}%
\begin{pgfscope}%
\pgfsys@transformshift{5.352733in}{0.660295in}%
\pgfsys@useobject{currentmarker}{}%
\end{pgfscope}%
\begin{pgfscope}%
\pgfsys@transformshift{5.352733in}{0.660295in}%
\pgfsys@useobject{currentmarker}{}%
\end{pgfscope}%
\begin{pgfscope}%
\pgfsys@transformshift{5.352733in}{0.660295in}%
\pgfsys@useobject{currentmarker}{}%
\end{pgfscope}%
\begin{pgfscope}%
\pgfsys@transformshift{5.352733in}{0.660295in}%
\pgfsys@useobject{currentmarker}{}%
\end{pgfscope}%
\begin{pgfscope}%
\pgfsys@transformshift{5.352733in}{0.660295in}%
\pgfsys@useobject{currentmarker}{}%
\end{pgfscope}%
\begin{pgfscope}%
\pgfsys@transformshift{5.352733in}{0.660295in}%
\pgfsys@useobject{currentmarker}{}%
\end{pgfscope}%
\begin{pgfscope}%
\pgfsys@transformshift{5.352733in}{0.660295in}%
\pgfsys@useobject{currentmarker}{}%
\end{pgfscope}%
\begin{pgfscope}%
\pgfsys@transformshift{5.352733in}{0.660295in}%
\pgfsys@useobject{currentmarker}{}%
\end{pgfscope}%
\begin{pgfscope}%
\pgfsys@transformshift{4.003409in}{0.631705in}%
\pgfsys@useobject{currentmarker}{}%
\end{pgfscope}%
\begin{pgfscope}%
\pgfsys@transformshift{4.003409in}{0.631705in}%
\pgfsys@useobject{currentmarker}{}%
\end{pgfscope}%
\begin{pgfscope}%
\pgfsys@transformshift{4.003409in}{0.631705in}%
\pgfsys@useobject{currentmarker}{}%
\end{pgfscope}%
\begin{pgfscope}%
\pgfsys@transformshift{4.003409in}{0.631705in}%
\pgfsys@useobject{currentmarker}{}%
\end{pgfscope}%
\begin{pgfscope}%
\pgfsys@transformshift{4.003409in}{0.631705in}%
\pgfsys@useobject{currentmarker}{}%
\end{pgfscope}%
\begin{pgfscope}%
\pgfsys@transformshift{4.003409in}{0.631705in}%
\pgfsys@useobject{currentmarker}{}%
\end{pgfscope}%
\begin{pgfscope}%
\pgfsys@transformshift{4.003409in}{0.631705in}%
\pgfsys@useobject{currentmarker}{}%
\end{pgfscope}%
\begin{pgfscope}%
\pgfsys@transformshift{4.003409in}{0.631705in}%
\pgfsys@useobject{currentmarker}{}%
\end{pgfscope}%
\begin{pgfscope}%
\pgfsys@transformshift{4.003409in}{0.631705in}%
\pgfsys@useobject{currentmarker}{}%
\end{pgfscope}%
\begin{pgfscope}%
\pgfsys@transformshift{4.003409in}{0.631705in}%
\pgfsys@useobject{currentmarker}{}%
\end{pgfscope}%
\begin{pgfscope}%
\pgfsys@transformshift{4.836668in}{0.603096in}%
\pgfsys@useobject{currentmarker}{}%
\end{pgfscope}%
\begin{pgfscope}%
\pgfsys@transformshift{4.836668in}{0.603096in}%
\pgfsys@useobject{currentmarker}{}%
\end{pgfscope}%
\begin{pgfscope}%
\pgfsys@transformshift{4.836668in}{0.603096in}%
\pgfsys@useobject{currentmarker}{}%
\end{pgfscope}%
\begin{pgfscope}%
\pgfsys@transformshift{4.836668in}{0.603096in}%
\pgfsys@useobject{currentmarker}{}%
\end{pgfscope}%
\begin{pgfscope}%
\pgfsys@transformshift{4.836668in}{0.603096in}%
\pgfsys@useobject{currentmarker}{}%
\end{pgfscope}%
\begin{pgfscope}%
\pgfsys@transformshift{4.836668in}{0.603096in}%
\pgfsys@useobject{currentmarker}{}%
\end{pgfscope}%
\begin{pgfscope}%
\pgfsys@transformshift{4.836668in}{0.603096in}%
\pgfsys@useobject{currentmarker}{}%
\end{pgfscope}%
\begin{pgfscope}%
\pgfsys@transformshift{4.836668in}{0.603096in}%
\pgfsys@useobject{currentmarker}{}%
\end{pgfscope}%
\begin{pgfscope}%
\pgfsys@transformshift{4.836668in}{0.603096in}%
\pgfsys@useobject{currentmarker}{}%
\end{pgfscope}%
\begin{pgfscope}%
\pgfsys@transformshift{4.836668in}{0.603096in}%
\pgfsys@useobject{currentmarker}{}%
\end{pgfscope}%
\begin{pgfscope}%
\pgfsys@transformshift{5.662336in}{0.574431in}%
\pgfsys@useobject{currentmarker}{}%
\end{pgfscope}%
\begin{pgfscope}%
\pgfsys@transformshift{5.662336in}{0.574431in}%
\pgfsys@useobject{currentmarker}{}%
\end{pgfscope}%
\begin{pgfscope}%
\pgfsys@transformshift{5.662336in}{0.574431in}%
\pgfsys@useobject{currentmarker}{}%
\end{pgfscope}%
\begin{pgfscope}%
\pgfsys@transformshift{5.662336in}{0.574431in}%
\pgfsys@useobject{currentmarker}{}%
\end{pgfscope}%
\begin{pgfscope}%
\pgfsys@transformshift{5.662336in}{0.574431in}%
\pgfsys@useobject{currentmarker}{}%
\end{pgfscope}%
\begin{pgfscope}%
\pgfsys@transformshift{5.662336in}{0.574431in}%
\pgfsys@useobject{currentmarker}{}%
\end{pgfscope}%
\begin{pgfscope}%
\pgfsys@transformshift{5.662336in}{0.574431in}%
\pgfsys@useobject{currentmarker}{}%
\end{pgfscope}%
\begin{pgfscope}%
\pgfsys@transformshift{5.662336in}{0.574431in}%
\pgfsys@useobject{currentmarker}{}%
\end{pgfscope}%
\begin{pgfscope}%
\pgfsys@transformshift{5.662336in}{0.574431in}%
\pgfsys@useobject{currentmarker}{}%
\end{pgfscope}%
\begin{pgfscope}%
\pgfsys@transformshift{5.662336in}{0.574431in}%
\pgfsys@useobject{currentmarker}{}%
\end{pgfscope}%
\begin{pgfscope}%
\pgfsys@transformshift{4.326510in}{0.545368in}%
\pgfsys@useobject{currentmarker}{}%
\end{pgfscope}%
\begin{pgfscope}%
\pgfsys@transformshift{4.326510in}{0.545368in}%
\pgfsys@useobject{currentmarker}{}%
\end{pgfscope}%
\begin{pgfscope}%
\pgfsys@transformshift{4.326510in}{0.545368in}%
\pgfsys@useobject{currentmarker}{}%
\end{pgfscope}%
\begin{pgfscope}%
\pgfsys@transformshift{4.326510in}{0.545368in}%
\pgfsys@useobject{currentmarker}{}%
\end{pgfscope}%
\begin{pgfscope}%
\pgfsys@transformshift{4.326510in}{0.545368in}%
\pgfsys@useobject{currentmarker}{}%
\end{pgfscope}%
\begin{pgfscope}%
\pgfsys@transformshift{4.326510in}{0.545368in}%
\pgfsys@useobject{currentmarker}{}%
\end{pgfscope}%
\begin{pgfscope}%
\pgfsys@transformshift{4.326510in}{0.545368in}%
\pgfsys@useobject{currentmarker}{}%
\end{pgfscope}%
\begin{pgfscope}%
\pgfsys@transformshift{4.326510in}{0.545368in}%
\pgfsys@useobject{currentmarker}{}%
\end{pgfscope}%
\begin{pgfscope}%
\pgfsys@transformshift{4.326510in}{0.545368in}%
\pgfsys@useobject{currentmarker}{}%
\end{pgfscope}%
\begin{pgfscope}%
\pgfsys@transformshift{4.326510in}{0.545368in}%
\pgfsys@useobject{currentmarker}{}%
\end{pgfscope}%
\begin{pgfscope}%
\pgfsys@transformshift{5.139927in}{0.516116in}%
\pgfsys@useobject{currentmarker}{}%
\end{pgfscope}%
\begin{pgfscope}%
\pgfsys@transformshift{5.139927in}{0.516116in}%
\pgfsys@useobject{currentmarker}{}%
\end{pgfscope}%
\begin{pgfscope}%
\pgfsys@transformshift{5.139927in}{0.516116in}%
\pgfsys@useobject{currentmarker}{}%
\end{pgfscope}%
\begin{pgfscope}%
\pgfsys@transformshift{5.139927in}{0.516116in}%
\pgfsys@useobject{currentmarker}{}%
\end{pgfscope}%
\begin{pgfscope}%
\pgfsys@transformshift{5.139927in}{0.516116in}%
\pgfsys@useobject{currentmarker}{}%
\end{pgfscope}%
\begin{pgfscope}%
\pgfsys@transformshift{5.139927in}{0.516116in}%
\pgfsys@useobject{currentmarker}{}%
\end{pgfscope}%
\begin{pgfscope}%
\pgfsys@transformshift{5.139927in}{0.516116in}%
\pgfsys@useobject{currentmarker}{}%
\end{pgfscope}%
\begin{pgfscope}%
\pgfsys@transformshift{5.139927in}{0.516116in}%
\pgfsys@useobject{currentmarker}{}%
\end{pgfscope}%
\begin{pgfscope}%
\pgfsys@transformshift{5.139927in}{0.516116in}%
\pgfsys@useobject{currentmarker}{}%
\end{pgfscope}%
\begin{pgfscope}%
\pgfsys@transformshift{5.139927in}{0.516116in}%
\pgfsys@useobject{currentmarker}{}%
\end{pgfscope}%
\begin{pgfscope}%
\pgfsys@transformshift{3.840652in}{0.486464in}%
\pgfsys@useobject{currentmarker}{}%
\end{pgfscope}%
\begin{pgfscope}%
\pgfsys@transformshift{3.840652in}{0.486464in}%
\pgfsys@useobject{currentmarker}{}%
\end{pgfscope}%
\begin{pgfscope}%
\pgfsys@transformshift{3.840652in}{0.486464in}%
\pgfsys@useobject{currentmarker}{}%
\end{pgfscope}%
\begin{pgfscope}%
\pgfsys@transformshift{3.840652in}{0.486464in}%
\pgfsys@useobject{currentmarker}{}%
\end{pgfscope}%
\begin{pgfscope}%
\pgfsys@transformshift{3.840652in}{0.486464in}%
\pgfsys@useobject{currentmarker}{}%
\end{pgfscope}%
\begin{pgfscope}%
\pgfsys@transformshift{3.840652in}{0.486464in}%
\pgfsys@useobject{currentmarker}{}%
\end{pgfscope}%
\begin{pgfscope}%
\pgfsys@transformshift{3.840652in}{0.486464in}%
\pgfsys@useobject{currentmarker}{}%
\end{pgfscope}%
\begin{pgfscope}%
\pgfsys@transformshift{3.840652in}{0.486464in}%
\pgfsys@useobject{currentmarker}{}%
\end{pgfscope}%
\begin{pgfscope}%
\pgfsys@transformshift{3.840652in}{0.486464in}%
\pgfsys@useobject{currentmarker}{}%
\end{pgfscope}%
\begin{pgfscope}%
\pgfsys@transformshift{3.840652in}{0.486464in}%
\pgfsys@useobject{currentmarker}{}%
\end{pgfscope}%
\begin{pgfscope}%
\pgfsys@transformshift{4.637593in}{0.456288in}%
\pgfsys@useobject{currentmarker}{}%
\end{pgfscope}%
\begin{pgfscope}%
\pgfsys@transformshift{4.637593in}{0.456288in}%
\pgfsys@useobject{currentmarker}{}%
\end{pgfscope}%
\begin{pgfscope}%
\pgfsys@transformshift{4.637593in}{0.456288in}%
\pgfsys@useobject{currentmarker}{}%
\end{pgfscope}%
\begin{pgfscope}%
\pgfsys@transformshift{4.637593in}{0.456288in}%
\pgfsys@useobject{currentmarker}{}%
\end{pgfscope}%
\begin{pgfscope}%
\pgfsys@transformshift{4.637593in}{0.456288in}%
\pgfsys@useobject{currentmarker}{}%
\end{pgfscope}%
\begin{pgfscope}%
\pgfsys@transformshift{4.637593in}{0.456288in}%
\pgfsys@useobject{currentmarker}{}%
\end{pgfscope}%
\begin{pgfscope}%
\pgfsys@transformshift{4.637593in}{0.456288in}%
\pgfsys@useobject{currentmarker}{}%
\end{pgfscope}%
\begin{pgfscope}%
\pgfsys@transformshift{4.637593in}{0.456288in}%
\pgfsys@useobject{currentmarker}{}%
\end{pgfscope}%
\begin{pgfscope}%
\pgfsys@transformshift{4.637593in}{0.456288in}%
\pgfsys@useobject{currentmarker}{}%
\end{pgfscope}%
\begin{pgfscope}%
\pgfsys@transformshift{4.637593in}{0.456288in}%
\pgfsys@useobject{currentmarker}{}%
\end{pgfscope}%
\begin{pgfscope}%
\pgfsys@transformshift{5.399648in}{0.425440in}%
\pgfsys@useobject{currentmarker}{}%
\end{pgfscope}%
\begin{pgfscope}%
\pgfsys@transformshift{5.399648in}{0.425440in}%
\pgfsys@useobject{currentmarker}{}%
\end{pgfscope}%
\begin{pgfscope}%
\pgfsys@transformshift{5.399648in}{0.425440in}%
\pgfsys@useobject{currentmarker}{}%
\end{pgfscope}%
\begin{pgfscope}%
\pgfsys@transformshift{5.399648in}{0.425440in}%
\pgfsys@useobject{currentmarker}{}%
\end{pgfscope}%
\begin{pgfscope}%
\pgfsys@transformshift{5.399648in}{0.425440in}%
\pgfsys@useobject{currentmarker}{}%
\end{pgfscope}%
\begin{pgfscope}%
\pgfsys@transformshift{5.399648in}{0.425440in}%
\pgfsys@useobject{currentmarker}{}%
\end{pgfscope}%
\begin{pgfscope}%
\pgfsys@transformshift{5.399648in}{0.425440in}%
\pgfsys@useobject{currentmarker}{}%
\end{pgfscope}%
\begin{pgfscope}%
\pgfsys@transformshift{5.399648in}{0.425440in}%
\pgfsys@useobject{currentmarker}{}%
\end{pgfscope}%
\begin{pgfscope}%
\pgfsys@transformshift{5.399648in}{0.425440in}%
\pgfsys@useobject{currentmarker}{}%
\end{pgfscope}%
\begin{pgfscope}%
\pgfsys@transformshift{5.399648in}{0.425440in}%
\pgfsys@useobject{currentmarker}{}%
\end{pgfscope}%
\begin{pgfscope}%
\pgfsys@transformshift{4.177987in}{0.393732in}%
\pgfsys@useobject{currentmarker}{}%
\end{pgfscope}%
\begin{pgfscope}%
\pgfsys@transformshift{4.177987in}{0.393732in}%
\pgfsys@useobject{currentmarker}{}%
\end{pgfscope}%
\begin{pgfscope}%
\pgfsys@transformshift{4.177987in}{0.393732in}%
\pgfsys@useobject{currentmarker}{}%
\end{pgfscope}%
\begin{pgfscope}%
\pgfsys@transformshift{4.177987in}{0.393732in}%
\pgfsys@useobject{currentmarker}{}%
\end{pgfscope}%
\begin{pgfscope}%
\pgfsys@transformshift{4.177987in}{0.393732in}%
\pgfsys@useobject{currentmarker}{}%
\end{pgfscope}%
\begin{pgfscope}%
\pgfsys@transformshift{4.177987in}{0.393732in}%
\pgfsys@useobject{currentmarker}{}%
\end{pgfscope}%
\begin{pgfscope}%
\pgfsys@transformshift{4.177987in}{0.393732in}%
\pgfsys@useobject{currentmarker}{}%
\end{pgfscope}%
\begin{pgfscope}%
\pgfsys@transformshift{4.177987in}{0.393732in}%
\pgfsys@useobject{currentmarker}{}%
\end{pgfscope}%
\begin{pgfscope}%
\pgfsys@transformshift{4.177987in}{0.393732in}%
\pgfsys@useobject{currentmarker}{}%
\end{pgfscope}%
\begin{pgfscope}%
\pgfsys@transformshift{4.177987in}{0.393732in}%
\pgfsys@useobject{currentmarker}{}%
\end{pgfscope}%
\begin{pgfscope}%
\pgfsys@transformshift{4.905408in}{0.360919in}%
\pgfsys@useobject{currentmarker}{}%
\end{pgfscope}%
\begin{pgfscope}%
\pgfsys@transformshift{4.905408in}{0.360919in}%
\pgfsys@useobject{currentmarker}{}%
\end{pgfscope}%
\begin{pgfscope}%
\pgfsys@transformshift{4.905408in}{0.360919in}%
\pgfsys@useobject{currentmarker}{}%
\end{pgfscope}%
\begin{pgfscope}%
\pgfsys@transformshift{4.905408in}{0.360919in}%
\pgfsys@useobject{currentmarker}{}%
\end{pgfscope}%
\begin{pgfscope}%
\pgfsys@transformshift{4.905408in}{0.360919in}%
\pgfsys@useobject{currentmarker}{}%
\end{pgfscope}%
\begin{pgfscope}%
\pgfsys@transformshift{4.905408in}{0.360919in}%
\pgfsys@useobject{currentmarker}{}%
\end{pgfscope}%
\begin{pgfscope}%
\pgfsys@transformshift{4.905408in}{0.360919in}%
\pgfsys@useobject{currentmarker}{}%
\end{pgfscope}%
\begin{pgfscope}%
\pgfsys@transformshift{4.905408in}{0.360919in}%
\pgfsys@useobject{currentmarker}{}%
\end{pgfscope}%
\begin{pgfscope}%
\pgfsys@transformshift{4.905408in}{0.360919in}%
\pgfsys@useobject{currentmarker}{}%
\end{pgfscope}%
\begin{pgfscope}%
\pgfsys@transformshift{4.905408in}{0.360919in}%
\pgfsys@useobject{currentmarker}{}%
\end{pgfscope}%
\begin{pgfscope}%
\pgfsys@transformshift{3.795317in}{0.326663in}%
\pgfsys@useobject{currentmarker}{}%
\end{pgfscope}%
\begin{pgfscope}%
\pgfsys@transformshift{3.795317in}{0.326663in}%
\pgfsys@useobject{currentmarker}{}%
\end{pgfscope}%
\begin{pgfscope}%
\pgfsys@transformshift{3.795317in}{0.326663in}%
\pgfsys@useobject{currentmarker}{}%
\end{pgfscope}%
\begin{pgfscope}%
\pgfsys@transformshift{3.795317in}{0.326663in}%
\pgfsys@useobject{currentmarker}{}%
\end{pgfscope}%
\begin{pgfscope}%
\pgfsys@transformshift{3.795317in}{0.326663in}%
\pgfsys@useobject{currentmarker}{}%
\end{pgfscope}%
\begin{pgfscope}%
\pgfsys@transformshift{3.795317in}{0.326663in}%
\pgfsys@useobject{currentmarker}{}%
\end{pgfscope}%
\begin{pgfscope}%
\pgfsys@transformshift{3.795317in}{0.326663in}%
\pgfsys@useobject{currentmarker}{}%
\end{pgfscope}%
\begin{pgfscope}%
\pgfsys@transformshift{3.795317in}{0.326663in}%
\pgfsys@useobject{currentmarker}{}%
\end{pgfscope}%
\begin{pgfscope}%
\pgfsys@transformshift{3.795317in}{0.326663in}%
\pgfsys@useobject{currentmarker}{}%
\end{pgfscope}%
\begin{pgfscope}%
\pgfsys@transformshift{3.795317in}{0.326663in}%
\pgfsys@useobject{currentmarker}{}%
\end{pgfscope}%
\begin{pgfscope}%
\pgfsys@transformshift{4.479906in}{0.290117in}%
\pgfsys@useobject{currentmarker}{}%
\end{pgfscope}%
\begin{pgfscope}%
\pgfsys@transformshift{4.479906in}{0.290117in}%
\pgfsys@useobject{currentmarker}{}%
\end{pgfscope}%
\begin{pgfscope}%
\pgfsys@transformshift{4.479906in}{0.290117in}%
\pgfsys@useobject{currentmarker}{}%
\end{pgfscope}%
\begin{pgfscope}%
\pgfsys@transformshift{4.479906in}{0.290117in}%
\pgfsys@useobject{currentmarker}{}%
\end{pgfscope}%
\begin{pgfscope}%
\pgfsys@transformshift{4.479906in}{0.290117in}%
\pgfsys@useobject{currentmarker}{}%
\end{pgfscope}%
\begin{pgfscope}%
\pgfsys@transformshift{4.479906in}{0.290117in}%
\pgfsys@useobject{currentmarker}{}%
\end{pgfscope}%
\begin{pgfscope}%
\pgfsys@transformshift{4.479906in}{0.290117in}%
\pgfsys@useobject{currentmarker}{}%
\end{pgfscope}%
\begin{pgfscope}%
\pgfsys@transformshift{4.479906in}{0.290117in}%
\pgfsys@useobject{currentmarker}{}%
\end{pgfscope}%
\begin{pgfscope}%
\pgfsys@transformshift{4.479906in}{0.290117in}%
\pgfsys@useobject{currentmarker}{}%
\end{pgfscope}%
\begin{pgfscope}%
\pgfsys@transformshift{4.479906in}{0.290117in}%
\pgfsys@useobject{currentmarker}{}%
\end{pgfscope}%
\begin{pgfscope}%
\pgfsys@transformshift{5.065669in}{0.250699in}%
\pgfsys@useobject{currentmarker}{}%
\end{pgfscope}%
\begin{pgfscope}%
\pgfsys@transformshift{5.065669in}{0.250699in}%
\pgfsys@useobject{currentmarker}{}%
\end{pgfscope}%
\begin{pgfscope}%
\pgfsys@transformshift{5.065669in}{0.250699in}%
\pgfsys@useobject{currentmarker}{}%
\end{pgfscope}%
\begin{pgfscope}%
\pgfsys@transformshift{5.065669in}{0.250699in}%
\pgfsys@useobject{currentmarker}{}%
\end{pgfscope}%
\begin{pgfscope}%
\pgfsys@transformshift{5.065669in}{0.250699in}%
\pgfsys@useobject{currentmarker}{}%
\end{pgfscope}%
\begin{pgfscope}%
\pgfsys@transformshift{5.065669in}{0.250699in}%
\pgfsys@useobject{currentmarker}{}%
\end{pgfscope}%
\begin{pgfscope}%
\pgfsys@transformshift{5.065669in}{0.250699in}%
\pgfsys@useobject{currentmarker}{}%
\end{pgfscope}%
\begin{pgfscope}%
\pgfsys@transformshift{5.065669in}{0.250699in}%
\pgfsys@useobject{currentmarker}{}%
\end{pgfscope}%
\begin{pgfscope}%
\pgfsys@transformshift{5.065669in}{0.250699in}%
\pgfsys@useobject{currentmarker}{}%
\end{pgfscope}%
\begin{pgfscope}%
\pgfsys@transformshift{5.065669in}{0.250699in}%
\pgfsys@useobject{currentmarker}{}%
\end{pgfscope}%
\begin{pgfscope}%
\pgfsys@transformshift{4.199187in}{0.206279in}%
\pgfsys@useobject{currentmarker}{}%
\end{pgfscope}%
\begin{pgfscope}%
\pgfsys@transformshift{4.199187in}{0.206279in}%
\pgfsys@useobject{currentmarker}{}%
\end{pgfscope}%
\begin{pgfscope}%
\pgfsys@transformshift{4.199187in}{0.206279in}%
\pgfsys@useobject{currentmarker}{}%
\end{pgfscope}%
\begin{pgfscope}%
\pgfsys@transformshift{4.199187in}{0.206279in}%
\pgfsys@useobject{currentmarker}{}%
\end{pgfscope}%
\begin{pgfscope}%
\pgfsys@transformshift{4.199187in}{0.206279in}%
\pgfsys@useobject{currentmarker}{}%
\end{pgfscope}%
\begin{pgfscope}%
\pgfsys@transformshift{4.199187in}{0.206279in}%
\pgfsys@useobject{currentmarker}{}%
\end{pgfscope}%
\begin{pgfscope}%
\pgfsys@transformshift{4.199187in}{0.206279in}%
\pgfsys@useobject{currentmarker}{}%
\end{pgfscope}%
\begin{pgfscope}%
\pgfsys@transformshift{4.199187in}{0.206279in}%
\pgfsys@useobject{currentmarker}{}%
\end{pgfscope}%
\begin{pgfscope}%
\pgfsys@transformshift{4.199187in}{0.206279in}%
\pgfsys@useobject{currentmarker}{}%
\end{pgfscope}%
\begin{pgfscope}%
\pgfsys@transformshift{4.199187in}{0.206279in}%
\pgfsys@useobject{currentmarker}{}%
\end{pgfscope}%
\begin{pgfscope}%
\pgfsys@transformshift{4.664677in}{0.150542in}%
\pgfsys@useobject{currentmarker}{}%
\end{pgfscope}%
\begin{pgfscope}%
\pgfsys@transformshift{4.664677in}{0.150542in}%
\pgfsys@useobject{currentmarker}{}%
\end{pgfscope}%
\begin{pgfscope}%
\pgfsys@transformshift{4.664677in}{0.150542in}%
\pgfsys@useobject{currentmarker}{}%
\end{pgfscope}%
\begin{pgfscope}%
\pgfsys@transformshift{4.664677in}{0.150542in}%
\pgfsys@useobject{currentmarker}{}%
\end{pgfscope}%
\begin{pgfscope}%
\pgfsys@transformshift{4.664677in}{0.150542in}%
\pgfsys@useobject{currentmarker}{}%
\end{pgfscope}%
\begin{pgfscope}%
\pgfsys@transformshift{4.664677in}{0.150542in}%
\pgfsys@useobject{currentmarker}{}%
\end{pgfscope}%
\begin{pgfscope}%
\pgfsys@transformshift{4.664677in}{0.150542in}%
\pgfsys@useobject{currentmarker}{}%
\end{pgfscope}%
\begin{pgfscope}%
\pgfsys@transformshift{4.664677in}{0.150542in}%
\pgfsys@useobject{currentmarker}{}%
\end{pgfscope}%
\begin{pgfscope}%
\pgfsys@transformshift{4.664677in}{0.150542in}%
\pgfsys@useobject{currentmarker}{}%
\end{pgfscope}%
\begin{pgfscope}%
\pgfsys@transformshift{4.664677in}{0.150542in}%
\pgfsys@useobject{currentmarker}{}%
\end{pgfscope}%
\end{pgfscope}%
\begin{pgfscope}%
\pgfpathmoveto{\pgfqpoint{4.594000in}{0.100000in}}%
\pgfpathcurveto{\pgfqpoint{4.883602in}{0.100000in}}{\pgfqpoint{5.161381in}{0.157530in}}{\pgfqpoint{5.366161in}{0.259920in}}%
\pgfpathcurveto{\pgfqpoint{5.570940in}{0.362309in}}{\pgfqpoint{5.686000in}{0.501199in}}{\pgfqpoint{5.686000in}{0.646000in}}%
\pgfpathcurveto{\pgfqpoint{5.686000in}{0.790801in}}{\pgfqpoint{5.570940in}{0.929691in}}{\pgfqpoint{5.366161in}{1.032080in}}%
\pgfpathcurveto{\pgfqpoint{5.161381in}{1.134470in}}{\pgfqpoint{4.883602in}{1.192000in}}{\pgfqpoint{4.594000in}{1.192000in}}%
\pgfpathcurveto{\pgfqpoint{4.304398in}{1.192000in}}{\pgfqpoint{4.026619in}{1.134470in}}{\pgfqpoint{3.821839in}{1.032080in}}%
\pgfpathcurveto{\pgfqpoint{3.617060in}{0.929691in}}{\pgfqpoint{3.502000in}{0.790801in}}{\pgfqpoint{3.502000in}{0.646000in}}%
\pgfpathcurveto{\pgfqpoint{3.502000in}{0.501199in}}{\pgfqpoint{3.617060in}{0.362309in}}{\pgfqpoint{3.821839in}{0.259920in}}%
\pgfpathcurveto{\pgfqpoint{4.026619in}{0.157530in}}{\pgfqpoint{4.304398in}{0.100000in}}{\pgfqpoint{4.594000in}{0.100000in}}%
\pgfpathlineto{\pgfqpoint{4.594000in}{0.100000in}}%
\pgfpathclose%
\pgfusepath{clip}%
\pgfsetrectcap%
\pgfsetroundjoin%
\pgfsetlinewidth{0.451687pt}%
\definecolor{currentstroke}{rgb}{0.690196,0.690196,0.690196}%
\pgfsetstrokecolor{currentstroke}%
\pgfsetstrokeopacity{0.250000}%
\pgfsetdash{}{0pt}%
\pgfpathmoveto{\pgfqpoint{4.208921in}{0.151295in}}%
\pgfpathlineto{\pgfqpoint{4.177157in}{0.160652in}}%
\pgfpathlineto{\pgfqpoint{4.146700in}{0.170513in}}%
\pgfpathlineto{\pgfqpoint{4.118255in}{0.180558in}}%
\pgfpathlineto{\pgfqpoint{4.091090in}{0.190956in}}%
\pgfpathlineto{\pgfqpoint{4.065134in}{0.201676in}}%
\pgfpathlineto{\pgfqpoint{4.040319in}{0.212695in}}%
\pgfpathlineto{\pgfqpoint{4.016585in}{0.223994in}}%
\pgfpathlineto{\pgfqpoint{3.993881in}{0.235556in}}%
\pgfpathlineto{\pgfqpoint{3.972160in}{0.247365in}}%
\pgfpathlineto{\pgfqpoint{3.951384in}{0.259410in}}%
\pgfpathlineto{\pgfqpoint{3.931517in}{0.271677in}}%
\pgfpathlineto{\pgfqpoint{3.912532in}{0.284154in}}%
\pgfpathlineto{\pgfqpoint{3.894403in}{0.296831in}}%
\pgfpathlineto{\pgfqpoint{3.877106in}{0.309698in}}%
\pgfpathlineto{\pgfqpoint{3.860426in}{0.322907in}}%
\pgfpathlineto{\pgfqpoint{3.844806in}{0.336079in}}%
\pgfpathlineto{\pgfqpoint{3.829948in}{0.349425in}}%
\pgfpathlineto{\pgfqpoint{3.815847in}{0.362931in}}%
\pgfpathlineto{\pgfqpoint{3.802494in}{0.376589in}}%
\pgfpathlineto{\pgfqpoint{3.789882in}{0.390387in}}%
\pgfpathlineto{\pgfqpoint{3.778004in}{0.404317in}}%
\pgfpathlineto{\pgfqpoint{3.766855in}{0.418370in}}%
\pgfpathlineto{\pgfqpoint{3.756428in}{0.432537in}}%
\pgfpathlineto{\pgfqpoint{3.746718in}{0.446809in}}%
\pgfpathlineto{\pgfqpoint{3.737720in}{0.461179in}}%
\pgfpathlineto{\pgfqpoint{3.729430in}{0.475639in}}%
\pgfpathlineto{\pgfqpoint{3.721844in}{0.490181in}}%
\pgfpathlineto{\pgfqpoint{3.714957in}{0.504798in}}%
\pgfpathlineto{\pgfqpoint{3.708768in}{0.519482in}}%
\pgfpathlineto{\pgfqpoint{3.703272in}{0.534226in}}%
\pgfpathlineto{\pgfqpoint{3.698469in}{0.549023in}}%
\pgfpathlineto{\pgfqpoint{3.694314in}{0.564029in}}%
\pgfpathlineto{\pgfqpoint{3.690911in}{0.578836in}}%
\pgfpathlineto{\pgfqpoint{3.688184in}{0.593701in}}%
\pgfpathlineto{\pgfqpoint{3.686136in}{0.608611in}}%
\pgfpathlineto{\pgfqpoint{3.684769in}{0.623554in}}%
\pgfpathlineto{\pgfqpoint{3.684085in}{0.638516in}}%
\pgfpathlineto{\pgfqpoint{3.684085in}{0.653484in}}%
\pgfpathlineto{\pgfqpoint{3.684769in}{0.668446in}}%
\pgfpathlineto{\pgfqpoint{3.686136in}{0.683389in}}%
\pgfpathlineto{\pgfqpoint{3.688184in}{0.698299in}}%
\pgfpathlineto{\pgfqpoint{3.690911in}{0.713164in}}%
\pgfpathlineto{\pgfqpoint{3.694314in}{0.727971in}}%
\pgfpathlineto{\pgfqpoint{3.698469in}{0.742977in}}%
\pgfpathlineto{\pgfqpoint{3.703272in}{0.757774in}}%
\pgfpathlineto{\pgfqpoint{3.708768in}{0.772518in}}%
\pgfpathlineto{\pgfqpoint{3.714957in}{0.787202in}}%
\pgfpathlineto{\pgfqpoint{3.721844in}{0.801819in}}%
\pgfpathlineto{\pgfqpoint{3.729430in}{0.816361in}}%
\pgfpathlineto{\pgfqpoint{3.737720in}{0.830821in}}%
\pgfpathlineto{\pgfqpoint{3.746718in}{0.845191in}}%
\pgfpathlineto{\pgfqpoint{3.756428in}{0.859463in}}%
\pgfpathlineto{\pgfqpoint{3.766855in}{0.873630in}}%
\pgfpathlineto{\pgfqpoint{3.778004in}{0.887683in}}%
\pgfpathlineto{\pgfqpoint{3.789882in}{0.901613in}}%
\pgfpathlineto{\pgfqpoint{3.802494in}{0.915411in}}%
\pgfpathlineto{\pgfqpoint{3.815847in}{0.929069in}}%
\pgfpathlineto{\pgfqpoint{3.829948in}{0.942575in}}%
\pgfpathlineto{\pgfqpoint{3.844806in}{0.955921in}}%
\pgfpathlineto{\pgfqpoint{3.860426in}{0.969093in}}%
\pgfpathlineto{\pgfqpoint{3.877106in}{0.982302in}}%
\pgfpathlineto{\pgfqpoint{3.894403in}{0.995169in}}%
\pgfpathlineto{\pgfqpoint{3.912532in}{1.007846in}}%
\pgfpathlineto{\pgfqpoint{3.931517in}{1.020323in}}%
\pgfpathlineto{\pgfqpoint{3.951384in}{1.032590in}}%
\pgfpathlineto{\pgfqpoint{3.972160in}{1.044635in}}%
\pgfpathlineto{\pgfqpoint{3.993881in}{1.056444in}}%
\pgfpathlineto{\pgfqpoint{4.016585in}{1.068006in}}%
\pgfpathlineto{\pgfqpoint{4.040319in}{1.079305in}}%
\pgfpathlineto{\pgfqpoint{4.065134in}{1.090324in}}%
\pgfpathlineto{\pgfqpoint{4.091090in}{1.101044in}}%
\pgfpathlineto{\pgfqpoint{4.118255in}{1.111442in}}%
\pgfpathlineto{\pgfqpoint{4.146700in}{1.121487in}}%
\pgfpathlineto{\pgfqpoint{4.177157in}{1.131348in}}%
\pgfpathlineto{\pgfqpoint{4.208921in}{1.140705in}}%
\pgfusepath{stroke}%
\end{pgfscope}%
\begin{pgfscope}%
\pgfpathmoveto{\pgfqpoint{4.594000in}{0.100000in}}%
\pgfpathcurveto{\pgfqpoint{4.883602in}{0.100000in}}{\pgfqpoint{5.161381in}{0.157530in}}{\pgfqpoint{5.366161in}{0.259920in}}%
\pgfpathcurveto{\pgfqpoint{5.570940in}{0.362309in}}{\pgfqpoint{5.686000in}{0.501199in}}{\pgfqpoint{5.686000in}{0.646000in}}%
\pgfpathcurveto{\pgfqpoint{5.686000in}{0.790801in}}{\pgfqpoint{5.570940in}{0.929691in}}{\pgfqpoint{5.366161in}{1.032080in}}%
\pgfpathcurveto{\pgfqpoint{5.161381in}{1.134470in}}{\pgfqpoint{4.883602in}{1.192000in}}{\pgfqpoint{4.594000in}{1.192000in}}%
\pgfpathcurveto{\pgfqpoint{4.304398in}{1.192000in}}{\pgfqpoint{4.026619in}{1.134470in}}{\pgfqpoint{3.821839in}{1.032080in}}%
\pgfpathcurveto{\pgfqpoint{3.617060in}{0.929691in}}{\pgfqpoint{3.502000in}{0.790801in}}{\pgfqpoint{3.502000in}{0.646000in}}%
\pgfpathcurveto{\pgfqpoint{3.502000in}{0.501199in}}{\pgfqpoint{3.617060in}{0.362309in}}{\pgfqpoint{3.821839in}{0.259920in}}%
\pgfpathcurveto{\pgfqpoint{4.026619in}{0.157530in}}{\pgfqpoint{4.304398in}{0.100000in}}{\pgfqpoint{4.594000in}{0.100000in}}%
\pgfpathlineto{\pgfqpoint{4.594000in}{0.100000in}}%
\pgfpathclose%
\pgfusepath{clip}%
\pgfsetrectcap%
\pgfsetroundjoin%
\pgfsetlinewidth{0.451687pt}%
\definecolor{currentstroke}{rgb}{0.690196,0.690196,0.690196}%
\pgfsetstrokecolor{currentstroke}%
\pgfsetstrokeopacity{0.250000}%
\pgfsetdash{}{0pt}%
\pgfpathmoveto{\pgfqpoint{4.285937in}{0.151295in}}%
\pgfpathlineto{\pgfqpoint{4.260525in}{0.160652in}}%
\pgfpathlineto{\pgfqpoint{4.236160in}{0.170513in}}%
\pgfpathlineto{\pgfqpoint{4.213404in}{0.180558in}}%
\pgfpathlineto{\pgfqpoint{4.191672in}{0.190956in}}%
\pgfpathlineto{\pgfqpoint{4.170907in}{0.201676in}}%
\pgfpathlineto{\pgfqpoint{4.151055in}{0.212695in}}%
\pgfpathlineto{\pgfqpoint{4.132068in}{0.223994in}}%
\pgfpathlineto{\pgfqpoint{4.113905in}{0.235556in}}%
\pgfpathlineto{\pgfqpoint{4.096528in}{0.247365in}}%
\pgfpathlineto{\pgfqpoint{4.079907in}{0.259410in}}%
\pgfpathlineto{\pgfqpoint{4.064014in}{0.271677in}}%
\pgfpathlineto{\pgfqpoint{4.048826in}{0.284154in}}%
\pgfpathlineto{\pgfqpoint{4.034322in}{0.296831in}}%
\pgfpathlineto{\pgfqpoint{4.020485in}{0.309698in}}%
\pgfpathlineto{\pgfqpoint{4.007141in}{0.322907in}}%
\pgfpathlineto{\pgfqpoint{3.994644in}{0.336079in}}%
\pgfpathlineto{\pgfqpoint{3.982759in}{0.349425in}}%
\pgfpathlineto{\pgfqpoint{3.971477in}{0.362931in}}%
\pgfpathlineto{\pgfqpoint{3.960795in}{0.376589in}}%
\pgfpathlineto{\pgfqpoint{3.950705in}{0.390387in}}%
\pgfpathlineto{\pgfqpoint{3.941204in}{0.404317in}}%
\pgfpathlineto{\pgfqpoint{3.932284in}{0.418370in}}%
\pgfpathlineto{\pgfqpoint{3.923943in}{0.432537in}}%
\pgfpathlineto{\pgfqpoint{3.916175in}{0.446809in}}%
\pgfpathlineto{\pgfqpoint{3.908976in}{0.461179in}}%
\pgfpathlineto{\pgfqpoint{3.902344in}{0.475639in}}%
\pgfpathlineto{\pgfqpoint{3.896275in}{0.490181in}}%
\pgfpathlineto{\pgfqpoint{3.890766in}{0.504798in}}%
\pgfpathlineto{\pgfqpoint{3.885814in}{0.519482in}}%
\pgfpathlineto{\pgfqpoint{3.881418in}{0.534226in}}%
\pgfpathlineto{\pgfqpoint{3.877575in}{0.549023in}}%
\pgfpathlineto{\pgfqpoint{3.874251in}{0.564029in}}%
\pgfpathlineto{\pgfqpoint{3.871529in}{0.578836in}}%
\pgfpathlineto{\pgfqpoint{3.869347in}{0.593701in}}%
\pgfpathlineto{\pgfqpoint{3.867709in}{0.608611in}}%
\pgfpathlineto{\pgfqpoint{3.866615in}{0.623554in}}%
\pgfpathlineto{\pgfqpoint{3.866068in}{0.638516in}}%
\pgfpathlineto{\pgfqpoint{3.866068in}{0.653484in}}%
\pgfpathlineto{\pgfqpoint{3.866615in}{0.668446in}}%
\pgfpathlineto{\pgfqpoint{3.867709in}{0.683389in}}%
\pgfpathlineto{\pgfqpoint{3.869347in}{0.698299in}}%
\pgfpathlineto{\pgfqpoint{3.871529in}{0.713164in}}%
\pgfpathlineto{\pgfqpoint{3.874251in}{0.727971in}}%
\pgfpathlineto{\pgfqpoint{3.877575in}{0.742977in}}%
\pgfpathlineto{\pgfqpoint{3.881418in}{0.757774in}}%
\pgfpathlineto{\pgfqpoint{3.885814in}{0.772518in}}%
\pgfpathlineto{\pgfqpoint{3.890766in}{0.787202in}}%
\pgfpathlineto{\pgfqpoint{3.896275in}{0.801819in}}%
\pgfpathlineto{\pgfqpoint{3.902344in}{0.816361in}}%
\pgfpathlineto{\pgfqpoint{3.908976in}{0.830821in}}%
\pgfpathlineto{\pgfqpoint{3.916175in}{0.845191in}}%
\pgfpathlineto{\pgfqpoint{3.923943in}{0.859463in}}%
\pgfpathlineto{\pgfqpoint{3.932284in}{0.873630in}}%
\pgfpathlineto{\pgfqpoint{3.941204in}{0.887683in}}%
\pgfpathlineto{\pgfqpoint{3.950705in}{0.901613in}}%
\pgfpathlineto{\pgfqpoint{3.960795in}{0.915411in}}%
\pgfpathlineto{\pgfqpoint{3.971477in}{0.929069in}}%
\pgfpathlineto{\pgfqpoint{3.982759in}{0.942575in}}%
\pgfpathlineto{\pgfqpoint{3.994644in}{0.955921in}}%
\pgfpathlineto{\pgfqpoint{4.007141in}{0.969093in}}%
\pgfpathlineto{\pgfqpoint{4.020485in}{0.982302in}}%
\pgfpathlineto{\pgfqpoint{4.034322in}{0.995169in}}%
\pgfpathlineto{\pgfqpoint{4.048826in}{1.007846in}}%
\pgfpathlineto{\pgfqpoint{4.064014in}{1.020323in}}%
\pgfpathlineto{\pgfqpoint{4.079907in}{1.032590in}}%
\pgfpathlineto{\pgfqpoint{4.096528in}{1.044635in}}%
\pgfpathlineto{\pgfqpoint{4.113905in}{1.056444in}}%
\pgfpathlineto{\pgfqpoint{4.132068in}{1.068006in}}%
\pgfpathlineto{\pgfqpoint{4.151055in}{1.079305in}}%
\pgfpathlineto{\pgfqpoint{4.170907in}{1.090324in}}%
\pgfpathlineto{\pgfqpoint{4.191672in}{1.101044in}}%
\pgfpathlineto{\pgfqpoint{4.213404in}{1.111442in}}%
\pgfpathlineto{\pgfqpoint{4.236160in}{1.121487in}}%
\pgfpathlineto{\pgfqpoint{4.260525in}{1.131348in}}%
\pgfpathlineto{\pgfqpoint{4.285937in}{1.140705in}}%
\pgfusepath{stroke}%
\end{pgfscope}%
\begin{pgfscope}%
\pgfpathmoveto{\pgfqpoint{4.594000in}{0.100000in}}%
\pgfpathcurveto{\pgfqpoint{4.883602in}{0.100000in}}{\pgfqpoint{5.161381in}{0.157530in}}{\pgfqpoint{5.366161in}{0.259920in}}%
\pgfpathcurveto{\pgfqpoint{5.570940in}{0.362309in}}{\pgfqpoint{5.686000in}{0.501199in}}{\pgfqpoint{5.686000in}{0.646000in}}%
\pgfpathcurveto{\pgfqpoint{5.686000in}{0.790801in}}{\pgfqpoint{5.570940in}{0.929691in}}{\pgfqpoint{5.366161in}{1.032080in}}%
\pgfpathcurveto{\pgfqpoint{5.161381in}{1.134470in}}{\pgfqpoint{4.883602in}{1.192000in}}{\pgfqpoint{4.594000in}{1.192000in}}%
\pgfpathcurveto{\pgfqpoint{4.304398in}{1.192000in}}{\pgfqpoint{4.026619in}{1.134470in}}{\pgfqpoint{3.821839in}{1.032080in}}%
\pgfpathcurveto{\pgfqpoint{3.617060in}{0.929691in}}{\pgfqpoint{3.502000in}{0.790801in}}{\pgfqpoint{3.502000in}{0.646000in}}%
\pgfpathcurveto{\pgfqpoint{3.502000in}{0.501199in}}{\pgfqpoint{3.617060in}{0.362309in}}{\pgfqpoint{3.821839in}{0.259920in}}%
\pgfpathcurveto{\pgfqpoint{4.026619in}{0.157530in}}{\pgfqpoint{4.304398in}{0.100000in}}{\pgfqpoint{4.594000in}{0.100000in}}%
\pgfpathlineto{\pgfqpoint{4.594000in}{0.100000in}}%
\pgfpathclose%
\pgfusepath{clip}%
\pgfsetrectcap%
\pgfsetroundjoin%
\pgfsetlinewidth{0.451687pt}%
\definecolor{currentstroke}{rgb}{0.690196,0.690196,0.690196}%
\pgfsetstrokecolor{currentstroke}%
\pgfsetstrokeopacity{0.250000}%
\pgfsetdash{}{0pt}%
\pgfpathmoveto{\pgfqpoint{4.362953in}{0.151295in}}%
\pgfpathlineto{\pgfqpoint{4.343894in}{0.160652in}}%
\pgfpathlineto{\pgfqpoint{4.325620in}{0.170513in}}%
\pgfpathlineto{\pgfqpoint{4.308553in}{0.180558in}}%
\pgfpathlineto{\pgfqpoint{4.292254in}{0.190956in}}%
\pgfpathlineto{\pgfqpoint{4.276680in}{0.201676in}}%
\pgfpathlineto{\pgfqpoint{4.261791in}{0.212695in}}%
\pgfpathlineto{\pgfqpoint{4.247551in}{0.223994in}}%
\pgfpathlineto{\pgfqpoint{4.233929in}{0.235556in}}%
\pgfpathlineto{\pgfqpoint{4.220896in}{0.247365in}}%
\pgfpathlineto{\pgfqpoint{4.208430in}{0.259410in}}%
\pgfpathlineto{\pgfqpoint{4.196510in}{0.271677in}}%
\pgfpathlineto{\pgfqpoint{4.185119in}{0.284154in}}%
\pgfpathlineto{\pgfqpoint{4.174242in}{0.296831in}}%
\pgfpathlineto{\pgfqpoint{4.163864in}{0.309698in}}%
\pgfpathlineto{\pgfqpoint{4.153856in}{0.322907in}}%
\pgfpathlineto{\pgfqpoint{4.144483in}{0.336079in}}%
\pgfpathlineto{\pgfqpoint{4.135569in}{0.349425in}}%
\pgfpathlineto{\pgfqpoint{4.127108in}{0.362931in}}%
\pgfpathlineto{\pgfqpoint{4.119096in}{0.376589in}}%
\pgfpathlineto{\pgfqpoint{4.111529in}{0.390387in}}%
\pgfpathlineto{\pgfqpoint{4.104403in}{0.404317in}}%
\pgfpathlineto{\pgfqpoint{4.097713in}{0.418370in}}%
\pgfpathlineto{\pgfqpoint{4.091457in}{0.432537in}}%
\pgfpathlineto{\pgfqpoint{4.085631in}{0.446809in}}%
\pgfpathlineto{\pgfqpoint{4.080232in}{0.461179in}}%
\pgfpathlineto{\pgfqpoint{4.075258in}{0.475639in}}%
\pgfpathlineto{\pgfqpoint{4.070706in}{0.490181in}}%
\pgfpathlineto{\pgfqpoint{4.066574in}{0.504798in}}%
\pgfpathlineto{\pgfqpoint{4.062861in}{0.519482in}}%
\pgfpathlineto{\pgfqpoint{4.059563in}{0.534226in}}%
\pgfpathlineto{\pgfqpoint{4.056681in}{0.549023in}}%
\pgfpathlineto{\pgfqpoint{4.054188in}{0.564029in}}%
\pgfpathlineto{\pgfqpoint{4.052147in}{0.578836in}}%
\pgfpathlineto{\pgfqpoint{4.050510in}{0.593701in}}%
\pgfpathlineto{\pgfqpoint{4.049282in}{0.608611in}}%
\pgfpathlineto{\pgfqpoint{4.048462in}{0.623554in}}%
\pgfpathlineto{\pgfqpoint{4.048051in}{0.638516in}}%
\pgfpathlineto{\pgfqpoint{4.048051in}{0.653484in}}%
\pgfpathlineto{\pgfqpoint{4.048462in}{0.668446in}}%
\pgfpathlineto{\pgfqpoint{4.049282in}{0.683389in}}%
\pgfpathlineto{\pgfqpoint{4.050510in}{0.698299in}}%
\pgfpathlineto{\pgfqpoint{4.052147in}{0.713164in}}%
\pgfpathlineto{\pgfqpoint{4.054188in}{0.727971in}}%
\pgfpathlineto{\pgfqpoint{4.056681in}{0.742977in}}%
\pgfpathlineto{\pgfqpoint{4.059563in}{0.757774in}}%
\pgfpathlineto{\pgfqpoint{4.062861in}{0.772518in}}%
\pgfpathlineto{\pgfqpoint{4.066574in}{0.787202in}}%
\pgfpathlineto{\pgfqpoint{4.070706in}{0.801819in}}%
\pgfpathlineto{\pgfqpoint{4.075258in}{0.816361in}}%
\pgfpathlineto{\pgfqpoint{4.080232in}{0.830821in}}%
\pgfpathlineto{\pgfqpoint{4.085631in}{0.845191in}}%
\pgfpathlineto{\pgfqpoint{4.091457in}{0.859463in}}%
\pgfpathlineto{\pgfqpoint{4.097713in}{0.873630in}}%
\pgfpathlineto{\pgfqpoint{4.104403in}{0.887683in}}%
\pgfpathlineto{\pgfqpoint{4.111529in}{0.901613in}}%
\pgfpathlineto{\pgfqpoint{4.119096in}{0.915411in}}%
\pgfpathlineto{\pgfqpoint{4.127108in}{0.929069in}}%
\pgfpathlineto{\pgfqpoint{4.135569in}{0.942575in}}%
\pgfpathlineto{\pgfqpoint{4.144483in}{0.955921in}}%
\pgfpathlineto{\pgfqpoint{4.153856in}{0.969093in}}%
\pgfpathlineto{\pgfqpoint{4.163864in}{0.982302in}}%
\pgfpathlineto{\pgfqpoint{4.174242in}{0.995169in}}%
\pgfpathlineto{\pgfqpoint{4.185119in}{1.007846in}}%
\pgfpathlineto{\pgfqpoint{4.196510in}{1.020323in}}%
\pgfpathlineto{\pgfqpoint{4.208430in}{1.032590in}}%
\pgfpathlineto{\pgfqpoint{4.220896in}{1.044635in}}%
\pgfpathlineto{\pgfqpoint{4.233929in}{1.056444in}}%
\pgfpathlineto{\pgfqpoint{4.247551in}{1.068006in}}%
\pgfpathlineto{\pgfqpoint{4.261791in}{1.079305in}}%
\pgfpathlineto{\pgfqpoint{4.276680in}{1.090324in}}%
\pgfpathlineto{\pgfqpoint{4.292254in}{1.101044in}}%
\pgfpathlineto{\pgfqpoint{4.308553in}{1.111442in}}%
\pgfpathlineto{\pgfqpoint{4.325620in}{1.121487in}}%
\pgfpathlineto{\pgfqpoint{4.343894in}{1.131348in}}%
\pgfpathlineto{\pgfqpoint{4.362953in}{1.140705in}}%
\pgfusepath{stroke}%
\end{pgfscope}%
\begin{pgfscope}%
\pgfpathmoveto{\pgfqpoint{4.594000in}{0.100000in}}%
\pgfpathcurveto{\pgfqpoint{4.883602in}{0.100000in}}{\pgfqpoint{5.161381in}{0.157530in}}{\pgfqpoint{5.366161in}{0.259920in}}%
\pgfpathcurveto{\pgfqpoint{5.570940in}{0.362309in}}{\pgfqpoint{5.686000in}{0.501199in}}{\pgfqpoint{5.686000in}{0.646000in}}%
\pgfpathcurveto{\pgfqpoint{5.686000in}{0.790801in}}{\pgfqpoint{5.570940in}{0.929691in}}{\pgfqpoint{5.366161in}{1.032080in}}%
\pgfpathcurveto{\pgfqpoint{5.161381in}{1.134470in}}{\pgfqpoint{4.883602in}{1.192000in}}{\pgfqpoint{4.594000in}{1.192000in}}%
\pgfpathcurveto{\pgfqpoint{4.304398in}{1.192000in}}{\pgfqpoint{4.026619in}{1.134470in}}{\pgfqpoint{3.821839in}{1.032080in}}%
\pgfpathcurveto{\pgfqpoint{3.617060in}{0.929691in}}{\pgfqpoint{3.502000in}{0.790801in}}{\pgfqpoint{3.502000in}{0.646000in}}%
\pgfpathcurveto{\pgfqpoint{3.502000in}{0.501199in}}{\pgfqpoint{3.617060in}{0.362309in}}{\pgfqpoint{3.821839in}{0.259920in}}%
\pgfpathcurveto{\pgfqpoint{4.026619in}{0.157530in}}{\pgfqpoint{4.304398in}{0.100000in}}{\pgfqpoint{4.594000in}{0.100000in}}%
\pgfpathlineto{\pgfqpoint{4.594000in}{0.100000in}}%
\pgfpathclose%
\pgfusepath{clip}%
\pgfsetrectcap%
\pgfsetroundjoin%
\pgfsetlinewidth{0.451687pt}%
\definecolor{currentstroke}{rgb}{0.690196,0.690196,0.690196}%
\pgfsetstrokecolor{currentstroke}%
\pgfsetstrokeopacity{0.250000}%
\pgfsetdash{}{0pt}%
\pgfpathmoveto{\pgfqpoint{4.439969in}{0.151295in}}%
\pgfpathlineto{\pgfqpoint{4.427263in}{0.160652in}}%
\pgfpathlineto{\pgfqpoint{4.415080in}{0.170513in}}%
\pgfpathlineto{\pgfqpoint{4.403702in}{0.180558in}}%
\pgfpathlineto{\pgfqpoint{4.392836in}{0.190956in}}%
\pgfpathlineto{\pgfqpoint{4.382453in}{0.201676in}}%
\pgfpathlineto{\pgfqpoint{4.372527in}{0.212695in}}%
\pgfpathlineto{\pgfqpoint{4.363034in}{0.223994in}}%
\pgfpathlineto{\pgfqpoint{4.353952in}{0.235556in}}%
\pgfpathlineto{\pgfqpoint{4.345264in}{0.247365in}}%
\pgfpathlineto{\pgfqpoint{4.336953in}{0.259410in}}%
\pgfpathlineto{\pgfqpoint{4.329007in}{0.271677in}}%
\pgfpathlineto{\pgfqpoint{4.321413in}{0.284154in}}%
\pgfpathlineto{\pgfqpoint{4.314161in}{0.296831in}}%
\pgfpathlineto{\pgfqpoint{4.307243in}{0.309698in}}%
\pgfpathlineto{\pgfqpoint{4.300570in}{0.322907in}}%
\pgfpathlineto{\pgfqpoint{4.294322in}{0.336079in}}%
\pgfpathlineto{\pgfqpoint{4.288379in}{0.349425in}}%
\pgfpathlineto{\pgfqpoint{4.282739in}{0.362931in}}%
\pgfpathlineto{\pgfqpoint{4.277398in}{0.376589in}}%
\pgfpathlineto{\pgfqpoint{4.272353in}{0.390387in}}%
\pgfpathlineto{\pgfqpoint{4.267602in}{0.404317in}}%
\pgfpathlineto{\pgfqpoint{4.263142in}{0.418370in}}%
\pgfpathlineto{\pgfqpoint{4.258971in}{0.432537in}}%
\pgfpathlineto{\pgfqpoint{4.255087in}{0.446809in}}%
\pgfpathlineto{\pgfqpoint{4.251488in}{0.461179in}}%
\pgfpathlineto{\pgfqpoint{4.248172in}{0.475639in}}%
\pgfpathlineto{\pgfqpoint{4.245137in}{0.490181in}}%
\pgfpathlineto{\pgfqpoint{4.242383in}{0.504798in}}%
\pgfpathlineto{\pgfqpoint{4.239907in}{0.519482in}}%
\pgfpathlineto{\pgfqpoint{4.237709in}{0.534226in}}%
\pgfpathlineto{\pgfqpoint{4.235788in}{0.549023in}}%
\pgfpathlineto{\pgfqpoint{4.234125in}{0.564029in}}%
\pgfpathlineto{\pgfqpoint{4.232764in}{0.578836in}}%
\pgfpathlineto{\pgfqpoint{4.231674in}{0.593701in}}%
\pgfpathlineto{\pgfqpoint{4.230854in}{0.608611in}}%
\pgfpathlineto{\pgfqpoint{4.230308in}{0.623554in}}%
\pgfpathlineto{\pgfqpoint{4.230034in}{0.638516in}}%
\pgfpathlineto{\pgfqpoint{4.230034in}{0.653484in}}%
\pgfpathlineto{\pgfqpoint{4.230308in}{0.668446in}}%
\pgfpathlineto{\pgfqpoint{4.230854in}{0.683389in}}%
\pgfpathlineto{\pgfqpoint{4.231674in}{0.698299in}}%
\pgfpathlineto{\pgfqpoint{4.232764in}{0.713164in}}%
\pgfpathlineto{\pgfqpoint{4.234125in}{0.727971in}}%
\pgfpathlineto{\pgfqpoint{4.235788in}{0.742977in}}%
\pgfpathlineto{\pgfqpoint{4.237709in}{0.757774in}}%
\pgfpathlineto{\pgfqpoint{4.239907in}{0.772518in}}%
\pgfpathlineto{\pgfqpoint{4.242383in}{0.787202in}}%
\pgfpathlineto{\pgfqpoint{4.245137in}{0.801819in}}%
\pgfpathlineto{\pgfqpoint{4.248172in}{0.816361in}}%
\pgfpathlineto{\pgfqpoint{4.251488in}{0.830821in}}%
\pgfpathlineto{\pgfqpoint{4.255087in}{0.845191in}}%
\pgfpathlineto{\pgfqpoint{4.258971in}{0.859463in}}%
\pgfpathlineto{\pgfqpoint{4.263142in}{0.873630in}}%
\pgfpathlineto{\pgfqpoint{4.267602in}{0.887683in}}%
\pgfpathlineto{\pgfqpoint{4.272353in}{0.901613in}}%
\pgfpathlineto{\pgfqpoint{4.277398in}{0.915411in}}%
\pgfpathlineto{\pgfqpoint{4.282739in}{0.929069in}}%
\pgfpathlineto{\pgfqpoint{4.288379in}{0.942575in}}%
\pgfpathlineto{\pgfqpoint{4.294322in}{0.955921in}}%
\pgfpathlineto{\pgfqpoint{4.300570in}{0.969093in}}%
\pgfpathlineto{\pgfqpoint{4.307243in}{0.982302in}}%
\pgfpathlineto{\pgfqpoint{4.314161in}{0.995169in}}%
\pgfpathlineto{\pgfqpoint{4.321413in}{1.007846in}}%
\pgfpathlineto{\pgfqpoint{4.329007in}{1.020323in}}%
\pgfpathlineto{\pgfqpoint{4.336953in}{1.032590in}}%
\pgfpathlineto{\pgfqpoint{4.345264in}{1.044635in}}%
\pgfpathlineto{\pgfqpoint{4.353952in}{1.056444in}}%
\pgfpathlineto{\pgfqpoint{4.363034in}{1.068006in}}%
\pgfpathlineto{\pgfqpoint{4.372527in}{1.079305in}}%
\pgfpathlineto{\pgfqpoint{4.382453in}{1.090324in}}%
\pgfpathlineto{\pgfqpoint{4.392836in}{1.101044in}}%
\pgfpathlineto{\pgfqpoint{4.403702in}{1.111442in}}%
\pgfpathlineto{\pgfqpoint{4.415080in}{1.121487in}}%
\pgfpathlineto{\pgfqpoint{4.427263in}{1.131348in}}%
\pgfpathlineto{\pgfqpoint{4.439969in}{1.140705in}}%
\pgfusepath{stroke}%
\end{pgfscope}%
\begin{pgfscope}%
\pgfpathmoveto{\pgfqpoint{4.594000in}{0.100000in}}%
\pgfpathcurveto{\pgfqpoint{4.883602in}{0.100000in}}{\pgfqpoint{5.161381in}{0.157530in}}{\pgfqpoint{5.366161in}{0.259920in}}%
\pgfpathcurveto{\pgfqpoint{5.570940in}{0.362309in}}{\pgfqpoint{5.686000in}{0.501199in}}{\pgfqpoint{5.686000in}{0.646000in}}%
\pgfpathcurveto{\pgfqpoint{5.686000in}{0.790801in}}{\pgfqpoint{5.570940in}{0.929691in}}{\pgfqpoint{5.366161in}{1.032080in}}%
\pgfpathcurveto{\pgfqpoint{5.161381in}{1.134470in}}{\pgfqpoint{4.883602in}{1.192000in}}{\pgfqpoint{4.594000in}{1.192000in}}%
\pgfpathcurveto{\pgfqpoint{4.304398in}{1.192000in}}{\pgfqpoint{4.026619in}{1.134470in}}{\pgfqpoint{3.821839in}{1.032080in}}%
\pgfpathcurveto{\pgfqpoint{3.617060in}{0.929691in}}{\pgfqpoint{3.502000in}{0.790801in}}{\pgfqpoint{3.502000in}{0.646000in}}%
\pgfpathcurveto{\pgfqpoint{3.502000in}{0.501199in}}{\pgfqpoint{3.617060in}{0.362309in}}{\pgfqpoint{3.821839in}{0.259920in}}%
\pgfpathcurveto{\pgfqpoint{4.026619in}{0.157530in}}{\pgfqpoint{4.304398in}{0.100000in}}{\pgfqpoint{4.594000in}{0.100000in}}%
\pgfpathlineto{\pgfqpoint{4.594000in}{0.100000in}}%
\pgfpathclose%
\pgfusepath{clip}%
\pgfsetrectcap%
\pgfsetroundjoin%
\pgfsetlinewidth{0.451687pt}%
\definecolor{currentstroke}{rgb}{0.690196,0.690196,0.690196}%
\pgfsetstrokecolor{currentstroke}%
\pgfsetstrokeopacity{0.250000}%
\pgfsetdash{}{0pt}%
\pgfpathmoveto{\pgfqpoint{4.516984in}{0.151295in}}%
\pgfpathlineto{\pgfqpoint{4.510631in}{0.160652in}}%
\pgfpathlineto{\pgfqpoint{4.504540in}{0.170513in}}%
\pgfpathlineto{\pgfqpoint{4.498851in}{0.180558in}}%
\pgfpathlineto{\pgfqpoint{4.493418in}{0.190956in}}%
\pgfpathlineto{\pgfqpoint{4.488227in}{0.201676in}}%
\pgfpathlineto{\pgfqpoint{4.483264in}{0.212695in}}%
\pgfpathlineto{\pgfqpoint{4.478517in}{0.223994in}}%
\pgfpathlineto{\pgfqpoint{4.473976in}{0.235556in}}%
\pgfpathlineto{\pgfqpoint{4.469632in}{0.247365in}}%
\pgfpathlineto{\pgfqpoint{4.465477in}{0.259410in}}%
\pgfpathlineto{\pgfqpoint{4.461503in}{0.271677in}}%
\pgfpathlineto{\pgfqpoint{4.457706in}{0.284154in}}%
\pgfpathlineto{\pgfqpoint{4.454081in}{0.296831in}}%
\pgfpathlineto{\pgfqpoint{4.450621in}{0.309698in}}%
\pgfpathlineto{\pgfqpoint{4.447285in}{0.322907in}}%
\pgfpathlineto{\pgfqpoint{4.444161in}{0.336079in}}%
\pgfpathlineto{\pgfqpoint{4.441190in}{0.349425in}}%
\pgfpathlineto{\pgfqpoint{4.438369in}{0.362931in}}%
\pgfpathlineto{\pgfqpoint{4.435699in}{0.376589in}}%
\pgfpathlineto{\pgfqpoint{4.433176in}{0.390387in}}%
\pgfpathlineto{\pgfqpoint{4.430801in}{0.404317in}}%
\pgfpathlineto{\pgfqpoint{4.428571in}{0.418370in}}%
\pgfpathlineto{\pgfqpoint{4.426486in}{0.432537in}}%
\pgfpathlineto{\pgfqpoint{4.424544in}{0.446809in}}%
\pgfpathlineto{\pgfqpoint{4.422744in}{0.461179in}}%
\pgfpathlineto{\pgfqpoint{4.421086in}{0.475639in}}%
\pgfpathlineto{\pgfqpoint{4.419569in}{0.490181in}}%
\pgfpathlineto{\pgfqpoint{4.418191in}{0.504798in}}%
\pgfpathlineto{\pgfqpoint{4.416954in}{0.519482in}}%
\pgfpathlineto{\pgfqpoint{4.415854in}{0.534226in}}%
\pgfpathlineto{\pgfqpoint{4.414894in}{0.549023in}}%
\pgfpathlineto{\pgfqpoint{4.414063in}{0.564029in}}%
\pgfpathlineto{\pgfqpoint{4.413382in}{0.578836in}}%
\pgfpathlineto{\pgfqpoint{4.412837in}{0.593701in}}%
\pgfpathlineto{\pgfqpoint{4.412427in}{0.608611in}}%
\pgfpathlineto{\pgfqpoint{4.412154in}{0.623554in}}%
\pgfpathlineto{\pgfqpoint{4.412017in}{0.638516in}}%
\pgfpathlineto{\pgfqpoint{4.412017in}{0.653484in}}%
\pgfpathlineto{\pgfqpoint{4.412154in}{0.668446in}}%
\pgfpathlineto{\pgfqpoint{4.412427in}{0.683389in}}%
\pgfpathlineto{\pgfqpoint{4.412837in}{0.698299in}}%
\pgfpathlineto{\pgfqpoint{4.413382in}{0.713164in}}%
\pgfpathlineto{\pgfqpoint{4.414063in}{0.727971in}}%
\pgfpathlineto{\pgfqpoint{4.414894in}{0.742977in}}%
\pgfpathlineto{\pgfqpoint{4.415854in}{0.757774in}}%
\pgfpathlineto{\pgfqpoint{4.416954in}{0.772518in}}%
\pgfpathlineto{\pgfqpoint{4.418191in}{0.787202in}}%
\pgfpathlineto{\pgfqpoint{4.419569in}{0.801819in}}%
\pgfpathlineto{\pgfqpoint{4.421086in}{0.816361in}}%
\pgfpathlineto{\pgfqpoint{4.422744in}{0.830821in}}%
\pgfpathlineto{\pgfqpoint{4.424544in}{0.845191in}}%
\pgfpathlineto{\pgfqpoint{4.426486in}{0.859463in}}%
\pgfpathlineto{\pgfqpoint{4.428571in}{0.873630in}}%
\pgfpathlineto{\pgfqpoint{4.430801in}{0.887683in}}%
\pgfpathlineto{\pgfqpoint{4.433176in}{0.901613in}}%
\pgfpathlineto{\pgfqpoint{4.435699in}{0.915411in}}%
\pgfpathlineto{\pgfqpoint{4.438369in}{0.929069in}}%
\pgfpathlineto{\pgfqpoint{4.441190in}{0.942575in}}%
\pgfpathlineto{\pgfqpoint{4.444161in}{0.955921in}}%
\pgfpathlineto{\pgfqpoint{4.447285in}{0.969093in}}%
\pgfpathlineto{\pgfqpoint{4.450621in}{0.982302in}}%
\pgfpathlineto{\pgfqpoint{4.454081in}{0.995169in}}%
\pgfpathlineto{\pgfqpoint{4.457706in}{1.007846in}}%
\pgfpathlineto{\pgfqpoint{4.461503in}{1.020323in}}%
\pgfpathlineto{\pgfqpoint{4.465477in}{1.032590in}}%
\pgfpathlineto{\pgfqpoint{4.469632in}{1.044635in}}%
\pgfpathlineto{\pgfqpoint{4.473976in}{1.056444in}}%
\pgfpathlineto{\pgfqpoint{4.478517in}{1.068006in}}%
\pgfpathlineto{\pgfqpoint{4.483264in}{1.079305in}}%
\pgfpathlineto{\pgfqpoint{4.488227in}{1.090324in}}%
\pgfpathlineto{\pgfqpoint{4.493418in}{1.101044in}}%
\pgfpathlineto{\pgfqpoint{4.498851in}{1.111442in}}%
\pgfpathlineto{\pgfqpoint{4.504540in}{1.121487in}}%
\pgfpathlineto{\pgfqpoint{4.510631in}{1.131348in}}%
\pgfpathlineto{\pgfqpoint{4.516984in}{1.140705in}}%
\pgfusepath{stroke}%
\end{pgfscope}%
\begin{pgfscope}%
\pgfpathmoveto{\pgfqpoint{4.594000in}{0.100000in}}%
\pgfpathcurveto{\pgfqpoint{4.883602in}{0.100000in}}{\pgfqpoint{5.161381in}{0.157530in}}{\pgfqpoint{5.366161in}{0.259920in}}%
\pgfpathcurveto{\pgfqpoint{5.570940in}{0.362309in}}{\pgfqpoint{5.686000in}{0.501199in}}{\pgfqpoint{5.686000in}{0.646000in}}%
\pgfpathcurveto{\pgfqpoint{5.686000in}{0.790801in}}{\pgfqpoint{5.570940in}{0.929691in}}{\pgfqpoint{5.366161in}{1.032080in}}%
\pgfpathcurveto{\pgfqpoint{5.161381in}{1.134470in}}{\pgfqpoint{4.883602in}{1.192000in}}{\pgfqpoint{4.594000in}{1.192000in}}%
\pgfpathcurveto{\pgfqpoint{4.304398in}{1.192000in}}{\pgfqpoint{4.026619in}{1.134470in}}{\pgfqpoint{3.821839in}{1.032080in}}%
\pgfpathcurveto{\pgfqpoint{3.617060in}{0.929691in}}{\pgfqpoint{3.502000in}{0.790801in}}{\pgfqpoint{3.502000in}{0.646000in}}%
\pgfpathcurveto{\pgfqpoint{3.502000in}{0.501199in}}{\pgfqpoint{3.617060in}{0.362309in}}{\pgfqpoint{3.821839in}{0.259920in}}%
\pgfpathcurveto{\pgfqpoint{4.026619in}{0.157530in}}{\pgfqpoint{4.304398in}{0.100000in}}{\pgfqpoint{4.594000in}{0.100000in}}%
\pgfpathlineto{\pgfqpoint{4.594000in}{0.100000in}}%
\pgfpathclose%
\pgfusepath{clip}%
\pgfsetrectcap%
\pgfsetroundjoin%
\pgfsetlinewidth{0.451687pt}%
\definecolor{currentstroke}{rgb}{0.690196,0.690196,0.690196}%
\pgfsetstrokecolor{currentstroke}%
\pgfsetstrokeopacity{0.250000}%
\pgfsetdash{}{0pt}%
\pgfpathmoveto{\pgfqpoint{4.594000in}{0.151295in}}%
\pgfpathlineto{\pgfqpoint{4.594000in}{0.160652in}}%
\pgfpathlineto{\pgfqpoint{4.594000in}{0.170513in}}%
\pgfpathlineto{\pgfqpoint{4.594000in}{0.180558in}}%
\pgfpathlineto{\pgfqpoint{4.594000in}{0.190956in}}%
\pgfpathlineto{\pgfqpoint{4.594000in}{0.201676in}}%
\pgfpathlineto{\pgfqpoint{4.594000in}{0.212695in}}%
\pgfpathlineto{\pgfqpoint{4.594000in}{0.223994in}}%
\pgfpathlineto{\pgfqpoint{4.594000in}{0.235556in}}%
\pgfpathlineto{\pgfqpoint{4.594000in}{0.247365in}}%
\pgfpathlineto{\pgfqpoint{4.594000in}{0.259410in}}%
\pgfpathlineto{\pgfqpoint{4.594000in}{0.271677in}}%
\pgfpathlineto{\pgfqpoint{4.594000in}{0.284154in}}%
\pgfpathlineto{\pgfqpoint{4.594000in}{0.296831in}}%
\pgfpathlineto{\pgfqpoint{4.594000in}{0.309698in}}%
\pgfpathlineto{\pgfqpoint{4.594000in}{0.322907in}}%
\pgfpathlineto{\pgfqpoint{4.594000in}{0.336079in}}%
\pgfpathlineto{\pgfqpoint{4.594000in}{0.349425in}}%
\pgfpathlineto{\pgfqpoint{4.594000in}{0.362931in}}%
\pgfpathlineto{\pgfqpoint{4.594000in}{0.376589in}}%
\pgfpathlineto{\pgfqpoint{4.594000in}{0.390387in}}%
\pgfpathlineto{\pgfqpoint{4.594000in}{0.404317in}}%
\pgfpathlineto{\pgfqpoint{4.594000in}{0.418370in}}%
\pgfpathlineto{\pgfqpoint{4.594000in}{0.432537in}}%
\pgfpathlineto{\pgfqpoint{4.594000in}{0.446809in}}%
\pgfpathlineto{\pgfqpoint{4.594000in}{0.461179in}}%
\pgfpathlineto{\pgfqpoint{4.594000in}{0.475639in}}%
\pgfpathlineto{\pgfqpoint{4.594000in}{0.490181in}}%
\pgfpathlineto{\pgfqpoint{4.594000in}{0.504798in}}%
\pgfpathlineto{\pgfqpoint{4.594000in}{0.519482in}}%
\pgfpathlineto{\pgfqpoint{4.594000in}{0.534226in}}%
\pgfpathlineto{\pgfqpoint{4.594000in}{0.549023in}}%
\pgfpathlineto{\pgfqpoint{4.594000in}{0.564029in}}%
\pgfpathlineto{\pgfqpoint{4.594000in}{0.578836in}}%
\pgfpathlineto{\pgfqpoint{4.594000in}{0.593701in}}%
\pgfpathlineto{\pgfqpoint{4.594000in}{0.608611in}}%
\pgfpathlineto{\pgfqpoint{4.594000in}{0.623554in}}%
\pgfpathlineto{\pgfqpoint{4.594000in}{0.638516in}}%
\pgfpathlineto{\pgfqpoint{4.594000in}{0.653484in}}%
\pgfpathlineto{\pgfqpoint{4.594000in}{0.668446in}}%
\pgfpathlineto{\pgfqpoint{4.594000in}{0.683389in}}%
\pgfpathlineto{\pgfqpoint{4.594000in}{0.698299in}}%
\pgfpathlineto{\pgfqpoint{4.594000in}{0.713164in}}%
\pgfpathlineto{\pgfqpoint{4.594000in}{0.727971in}}%
\pgfpathlineto{\pgfqpoint{4.594000in}{0.742977in}}%
\pgfpathlineto{\pgfqpoint{4.594000in}{0.757774in}}%
\pgfpathlineto{\pgfqpoint{4.594000in}{0.772518in}}%
\pgfpathlineto{\pgfqpoint{4.594000in}{0.787202in}}%
\pgfpathlineto{\pgfqpoint{4.594000in}{0.801819in}}%
\pgfpathlineto{\pgfqpoint{4.594000in}{0.816361in}}%
\pgfpathlineto{\pgfqpoint{4.594000in}{0.830821in}}%
\pgfpathlineto{\pgfqpoint{4.594000in}{0.845191in}}%
\pgfpathlineto{\pgfqpoint{4.594000in}{0.859463in}}%
\pgfpathlineto{\pgfqpoint{4.594000in}{0.873630in}}%
\pgfpathlineto{\pgfqpoint{4.594000in}{0.887683in}}%
\pgfpathlineto{\pgfqpoint{4.594000in}{0.901613in}}%
\pgfpathlineto{\pgfqpoint{4.594000in}{0.915411in}}%
\pgfpathlineto{\pgfqpoint{4.594000in}{0.929069in}}%
\pgfpathlineto{\pgfqpoint{4.594000in}{0.942575in}}%
\pgfpathlineto{\pgfqpoint{4.594000in}{0.955921in}}%
\pgfpathlineto{\pgfqpoint{4.594000in}{0.969093in}}%
\pgfpathlineto{\pgfqpoint{4.594000in}{0.982302in}}%
\pgfpathlineto{\pgfqpoint{4.594000in}{0.995169in}}%
\pgfpathlineto{\pgfqpoint{4.594000in}{1.007846in}}%
\pgfpathlineto{\pgfqpoint{4.594000in}{1.020323in}}%
\pgfpathlineto{\pgfqpoint{4.594000in}{1.032590in}}%
\pgfpathlineto{\pgfqpoint{4.594000in}{1.044635in}}%
\pgfpathlineto{\pgfqpoint{4.594000in}{1.056444in}}%
\pgfpathlineto{\pgfqpoint{4.594000in}{1.068006in}}%
\pgfpathlineto{\pgfqpoint{4.594000in}{1.079305in}}%
\pgfpathlineto{\pgfqpoint{4.594000in}{1.090324in}}%
\pgfpathlineto{\pgfqpoint{4.594000in}{1.101044in}}%
\pgfpathlineto{\pgfqpoint{4.594000in}{1.111442in}}%
\pgfpathlineto{\pgfqpoint{4.594000in}{1.121487in}}%
\pgfpathlineto{\pgfqpoint{4.594000in}{1.131348in}}%
\pgfpathlineto{\pgfqpoint{4.594000in}{1.140705in}}%
\pgfusepath{stroke}%
\end{pgfscope}%
\begin{pgfscope}%
\pgfpathmoveto{\pgfqpoint{4.594000in}{0.100000in}}%
\pgfpathcurveto{\pgfqpoint{4.883602in}{0.100000in}}{\pgfqpoint{5.161381in}{0.157530in}}{\pgfqpoint{5.366161in}{0.259920in}}%
\pgfpathcurveto{\pgfqpoint{5.570940in}{0.362309in}}{\pgfqpoint{5.686000in}{0.501199in}}{\pgfqpoint{5.686000in}{0.646000in}}%
\pgfpathcurveto{\pgfqpoint{5.686000in}{0.790801in}}{\pgfqpoint{5.570940in}{0.929691in}}{\pgfqpoint{5.366161in}{1.032080in}}%
\pgfpathcurveto{\pgfqpoint{5.161381in}{1.134470in}}{\pgfqpoint{4.883602in}{1.192000in}}{\pgfqpoint{4.594000in}{1.192000in}}%
\pgfpathcurveto{\pgfqpoint{4.304398in}{1.192000in}}{\pgfqpoint{4.026619in}{1.134470in}}{\pgfqpoint{3.821839in}{1.032080in}}%
\pgfpathcurveto{\pgfqpoint{3.617060in}{0.929691in}}{\pgfqpoint{3.502000in}{0.790801in}}{\pgfqpoint{3.502000in}{0.646000in}}%
\pgfpathcurveto{\pgfqpoint{3.502000in}{0.501199in}}{\pgfqpoint{3.617060in}{0.362309in}}{\pgfqpoint{3.821839in}{0.259920in}}%
\pgfpathcurveto{\pgfqpoint{4.026619in}{0.157530in}}{\pgfqpoint{4.304398in}{0.100000in}}{\pgfqpoint{4.594000in}{0.100000in}}%
\pgfpathlineto{\pgfqpoint{4.594000in}{0.100000in}}%
\pgfpathclose%
\pgfusepath{clip}%
\pgfsetrectcap%
\pgfsetroundjoin%
\pgfsetlinewidth{0.451687pt}%
\definecolor{currentstroke}{rgb}{0.690196,0.690196,0.690196}%
\pgfsetstrokecolor{currentstroke}%
\pgfsetstrokeopacity{0.250000}%
\pgfsetdash{}{0pt}%
\pgfpathmoveto{\pgfqpoint{4.671016in}{0.151295in}}%
\pgfpathlineto{\pgfqpoint{4.677369in}{0.160652in}}%
\pgfpathlineto{\pgfqpoint{4.683460in}{0.170513in}}%
\pgfpathlineto{\pgfqpoint{4.689149in}{0.180558in}}%
\pgfpathlineto{\pgfqpoint{4.694582in}{0.190956in}}%
\pgfpathlineto{\pgfqpoint{4.699773in}{0.201676in}}%
\pgfpathlineto{\pgfqpoint{4.704736in}{0.212695in}}%
\pgfpathlineto{\pgfqpoint{4.709483in}{0.223994in}}%
\pgfpathlineto{\pgfqpoint{4.714024in}{0.235556in}}%
\pgfpathlineto{\pgfqpoint{4.718368in}{0.247365in}}%
\pgfpathlineto{\pgfqpoint{4.722523in}{0.259410in}}%
\pgfpathlineto{\pgfqpoint{4.726497in}{0.271677in}}%
\pgfpathlineto{\pgfqpoint{4.730294in}{0.284154in}}%
\pgfpathlineto{\pgfqpoint{4.733919in}{0.296831in}}%
\pgfpathlineto{\pgfqpoint{4.737379in}{0.309698in}}%
\pgfpathlineto{\pgfqpoint{4.740715in}{0.322907in}}%
\pgfpathlineto{\pgfqpoint{4.743839in}{0.336079in}}%
\pgfpathlineto{\pgfqpoint{4.746810in}{0.349425in}}%
\pgfpathlineto{\pgfqpoint{4.749631in}{0.362931in}}%
\pgfpathlineto{\pgfqpoint{4.752301in}{0.376589in}}%
\pgfpathlineto{\pgfqpoint{4.754824in}{0.390387in}}%
\pgfpathlineto{\pgfqpoint{4.757199in}{0.404317in}}%
\pgfpathlineto{\pgfqpoint{4.759429in}{0.418370in}}%
\pgfpathlineto{\pgfqpoint{4.761514in}{0.432537in}}%
\pgfpathlineto{\pgfqpoint{4.763456in}{0.446809in}}%
\pgfpathlineto{\pgfqpoint{4.765256in}{0.461179in}}%
\pgfpathlineto{\pgfqpoint{4.766914in}{0.475639in}}%
\pgfpathlineto{\pgfqpoint{4.768431in}{0.490181in}}%
\pgfpathlineto{\pgfqpoint{4.769809in}{0.504798in}}%
\pgfpathlineto{\pgfqpoint{4.771046in}{0.519482in}}%
\pgfpathlineto{\pgfqpoint{4.772146in}{0.534226in}}%
\pgfpathlineto{\pgfqpoint{4.773106in}{0.549023in}}%
\pgfpathlineto{\pgfqpoint{4.773937in}{0.564029in}}%
\pgfpathlineto{\pgfqpoint{4.774618in}{0.578836in}}%
\pgfpathlineto{\pgfqpoint{4.775163in}{0.593701in}}%
\pgfpathlineto{\pgfqpoint{4.775573in}{0.608611in}}%
\pgfpathlineto{\pgfqpoint{4.775846in}{0.623554in}}%
\pgfpathlineto{\pgfqpoint{4.775983in}{0.638516in}}%
\pgfpathlineto{\pgfqpoint{4.775983in}{0.653484in}}%
\pgfpathlineto{\pgfqpoint{4.775846in}{0.668446in}}%
\pgfpathlineto{\pgfqpoint{4.775573in}{0.683389in}}%
\pgfpathlineto{\pgfqpoint{4.775163in}{0.698299in}}%
\pgfpathlineto{\pgfqpoint{4.774618in}{0.713164in}}%
\pgfpathlineto{\pgfqpoint{4.773937in}{0.727971in}}%
\pgfpathlineto{\pgfqpoint{4.773106in}{0.742977in}}%
\pgfpathlineto{\pgfqpoint{4.772146in}{0.757774in}}%
\pgfpathlineto{\pgfqpoint{4.771046in}{0.772518in}}%
\pgfpathlineto{\pgfqpoint{4.769809in}{0.787202in}}%
\pgfpathlineto{\pgfqpoint{4.768431in}{0.801819in}}%
\pgfpathlineto{\pgfqpoint{4.766914in}{0.816361in}}%
\pgfpathlineto{\pgfqpoint{4.765256in}{0.830821in}}%
\pgfpathlineto{\pgfqpoint{4.763456in}{0.845191in}}%
\pgfpathlineto{\pgfqpoint{4.761514in}{0.859463in}}%
\pgfpathlineto{\pgfqpoint{4.759429in}{0.873630in}}%
\pgfpathlineto{\pgfqpoint{4.757199in}{0.887683in}}%
\pgfpathlineto{\pgfqpoint{4.754824in}{0.901613in}}%
\pgfpathlineto{\pgfqpoint{4.752301in}{0.915411in}}%
\pgfpathlineto{\pgfqpoint{4.749631in}{0.929069in}}%
\pgfpathlineto{\pgfqpoint{4.746810in}{0.942575in}}%
\pgfpathlineto{\pgfqpoint{4.743839in}{0.955921in}}%
\pgfpathlineto{\pgfqpoint{4.740715in}{0.969093in}}%
\pgfpathlineto{\pgfqpoint{4.737379in}{0.982302in}}%
\pgfpathlineto{\pgfqpoint{4.733919in}{0.995169in}}%
\pgfpathlineto{\pgfqpoint{4.730294in}{1.007846in}}%
\pgfpathlineto{\pgfqpoint{4.726497in}{1.020323in}}%
\pgfpathlineto{\pgfqpoint{4.722523in}{1.032590in}}%
\pgfpathlineto{\pgfqpoint{4.718368in}{1.044635in}}%
\pgfpathlineto{\pgfqpoint{4.714024in}{1.056444in}}%
\pgfpathlineto{\pgfqpoint{4.709483in}{1.068006in}}%
\pgfpathlineto{\pgfqpoint{4.704736in}{1.079305in}}%
\pgfpathlineto{\pgfqpoint{4.699773in}{1.090324in}}%
\pgfpathlineto{\pgfqpoint{4.694582in}{1.101044in}}%
\pgfpathlineto{\pgfqpoint{4.689149in}{1.111442in}}%
\pgfpathlineto{\pgfqpoint{4.683460in}{1.121487in}}%
\pgfpathlineto{\pgfqpoint{4.677369in}{1.131348in}}%
\pgfpathlineto{\pgfqpoint{4.671016in}{1.140705in}}%
\pgfusepath{stroke}%
\end{pgfscope}%
\begin{pgfscope}%
\pgfpathmoveto{\pgfqpoint{4.594000in}{0.100000in}}%
\pgfpathcurveto{\pgfqpoint{4.883602in}{0.100000in}}{\pgfqpoint{5.161381in}{0.157530in}}{\pgfqpoint{5.366161in}{0.259920in}}%
\pgfpathcurveto{\pgfqpoint{5.570940in}{0.362309in}}{\pgfqpoint{5.686000in}{0.501199in}}{\pgfqpoint{5.686000in}{0.646000in}}%
\pgfpathcurveto{\pgfqpoint{5.686000in}{0.790801in}}{\pgfqpoint{5.570940in}{0.929691in}}{\pgfqpoint{5.366161in}{1.032080in}}%
\pgfpathcurveto{\pgfqpoint{5.161381in}{1.134470in}}{\pgfqpoint{4.883602in}{1.192000in}}{\pgfqpoint{4.594000in}{1.192000in}}%
\pgfpathcurveto{\pgfqpoint{4.304398in}{1.192000in}}{\pgfqpoint{4.026619in}{1.134470in}}{\pgfqpoint{3.821839in}{1.032080in}}%
\pgfpathcurveto{\pgfqpoint{3.617060in}{0.929691in}}{\pgfqpoint{3.502000in}{0.790801in}}{\pgfqpoint{3.502000in}{0.646000in}}%
\pgfpathcurveto{\pgfqpoint{3.502000in}{0.501199in}}{\pgfqpoint{3.617060in}{0.362309in}}{\pgfqpoint{3.821839in}{0.259920in}}%
\pgfpathcurveto{\pgfqpoint{4.026619in}{0.157530in}}{\pgfqpoint{4.304398in}{0.100000in}}{\pgfqpoint{4.594000in}{0.100000in}}%
\pgfpathlineto{\pgfqpoint{4.594000in}{0.100000in}}%
\pgfpathclose%
\pgfusepath{clip}%
\pgfsetrectcap%
\pgfsetroundjoin%
\pgfsetlinewidth{0.451687pt}%
\definecolor{currentstroke}{rgb}{0.690196,0.690196,0.690196}%
\pgfsetstrokecolor{currentstroke}%
\pgfsetstrokeopacity{0.250000}%
\pgfsetdash{}{0pt}%
\pgfpathmoveto{\pgfqpoint{4.748031in}{0.151295in}}%
\pgfpathlineto{\pgfqpoint{4.760737in}{0.160652in}}%
\pgfpathlineto{\pgfqpoint{4.772920in}{0.170513in}}%
\pgfpathlineto{\pgfqpoint{4.784298in}{0.180558in}}%
\pgfpathlineto{\pgfqpoint{4.795164in}{0.190956in}}%
\pgfpathlineto{\pgfqpoint{4.805547in}{0.201676in}}%
\pgfpathlineto{\pgfqpoint{4.815473in}{0.212695in}}%
\pgfpathlineto{\pgfqpoint{4.824966in}{0.223994in}}%
\pgfpathlineto{\pgfqpoint{4.834048in}{0.235556in}}%
\pgfpathlineto{\pgfqpoint{4.842736in}{0.247365in}}%
\pgfpathlineto{\pgfqpoint{4.851047in}{0.259410in}}%
\pgfpathlineto{\pgfqpoint{4.858993in}{0.271677in}}%
\pgfpathlineto{\pgfqpoint{4.866587in}{0.284154in}}%
\pgfpathlineto{\pgfqpoint{4.873839in}{0.296831in}}%
\pgfpathlineto{\pgfqpoint{4.880757in}{0.309698in}}%
\pgfpathlineto{\pgfqpoint{4.887430in}{0.322907in}}%
\pgfpathlineto{\pgfqpoint{4.893678in}{0.336079in}}%
\pgfpathlineto{\pgfqpoint{4.899621in}{0.349425in}}%
\pgfpathlineto{\pgfqpoint{4.905261in}{0.362931in}}%
\pgfpathlineto{\pgfqpoint{4.910602in}{0.376589in}}%
\pgfpathlineto{\pgfqpoint{4.915647in}{0.390387in}}%
\pgfpathlineto{\pgfqpoint{4.920398in}{0.404317in}}%
\pgfpathlineto{\pgfqpoint{4.924858in}{0.418370in}}%
\pgfpathlineto{\pgfqpoint{4.929029in}{0.432537in}}%
\pgfpathlineto{\pgfqpoint{4.932913in}{0.446809in}}%
\pgfpathlineto{\pgfqpoint{4.936512in}{0.461179in}}%
\pgfpathlineto{\pgfqpoint{4.939828in}{0.475639in}}%
\pgfpathlineto{\pgfqpoint{4.942863in}{0.490181in}}%
\pgfpathlineto{\pgfqpoint{4.945617in}{0.504798in}}%
\pgfpathlineto{\pgfqpoint{4.948093in}{0.519482in}}%
\pgfpathlineto{\pgfqpoint{4.950291in}{0.534226in}}%
\pgfpathlineto{\pgfqpoint{4.952212in}{0.549023in}}%
\pgfpathlineto{\pgfqpoint{4.953875in}{0.564029in}}%
\pgfpathlineto{\pgfqpoint{4.955236in}{0.578836in}}%
\pgfpathlineto{\pgfqpoint{4.956326in}{0.593701in}}%
\pgfpathlineto{\pgfqpoint{4.957146in}{0.608611in}}%
\pgfpathlineto{\pgfqpoint{4.957692in}{0.623554in}}%
\pgfpathlineto{\pgfqpoint{4.957966in}{0.638516in}}%
\pgfpathlineto{\pgfqpoint{4.957966in}{0.653484in}}%
\pgfpathlineto{\pgfqpoint{4.957692in}{0.668446in}}%
\pgfpathlineto{\pgfqpoint{4.957146in}{0.683389in}}%
\pgfpathlineto{\pgfqpoint{4.956326in}{0.698299in}}%
\pgfpathlineto{\pgfqpoint{4.955236in}{0.713164in}}%
\pgfpathlineto{\pgfqpoint{4.953875in}{0.727971in}}%
\pgfpathlineto{\pgfqpoint{4.952212in}{0.742977in}}%
\pgfpathlineto{\pgfqpoint{4.950291in}{0.757774in}}%
\pgfpathlineto{\pgfqpoint{4.948093in}{0.772518in}}%
\pgfpathlineto{\pgfqpoint{4.945617in}{0.787202in}}%
\pgfpathlineto{\pgfqpoint{4.942863in}{0.801819in}}%
\pgfpathlineto{\pgfqpoint{4.939828in}{0.816361in}}%
\pgfpathlineto{\pgfqpoint{4.936512in}{0.830821in}}%
\pgfpathlineto{\pgfqpoint{4.932913in}{0.845191in}}%
\pgfpathlineto{\pgfqpoint{4.929029in}{0.859463in}}%
\pgfpathlineto{\pgfqpoint{4.924858in}{0.873630in}}%
\pgfpathlineto{\pgfqpoint{4.920398in}{0.887683in}}%
\pgfpathlineto{\pgfqpoint{4.915647in}{0.901613in}}%
\pgfpathlineto{\pgfqpoint{4.910602in}{0.915411in}}%
\pgfpathlineto{\pgfqpoint{4.905261in}{0.929069in}}%
\pgfpathlineto{\pgfqpoint{4.899621in}{0.942575in}}%
\pgfpathlineto{\pgfqpoint{4.893678in}{0.955921in}}%
\pgfpathlineto{\pgfqpoint{4.887430in}{0.969093in}}%
\pgfpathlineto{\pgfqpoint{4.880757in}{0.982302in}}%
\pgfpathlineto{\pgfqpoint{4.873839in}{0.995169in}}%
\pgfpathlineto{\pgfqpoint{4.866587in}{1.007846in}}%
\pgfpathlineto{\pgfqpoint{4.858993in}{1.020323in}}%
\pgfpathlineto{\pgfqpoint{4.851047in}{1.032590in}}%
\pgfpathlineto{\pgfqpoint{4.842736in}{1.044635in}}%
\pgfpathlineto{\pgfqpoint{4.834048in}{1.056444in}}%
\pgfpathlineto{\pgfqpoint{4.824966in}{1.068006in}}%
\pgfpathlineto{\pgfqpoint{4.815473in}{1.079305in}}%
\pgfpathlineto{\pgfqpoint{4.805547in}{1.090324in}}%
\pgfpathlineto{\pgfqpoint{4.795164in}{1.101044in}}%
\pgfpathlineto{\pgfqpoint{4.784298in}{1.111442in}}%
\pgfpathlineto{\pgfqpoint{4.772920in}{1.121487in}}%
\pgfpathlineto{\pgfqpoint{4.760737in}{1.131348in}}%
\pgfpathlineto{\pgfqpoint{4.748031in}{1.140705in}}%
\pgfusepath{stroke}%
\end{pgfscope}%
\begin{pgfscope}%
\pgfpathmoveto{\pgfqpoint{4.594000in}{0.100000in}}%
\pgfpathcurveto{\pgfqpoint{4.883602in}{0.100000in}}{\pgfqpoint{5.161381in}{0.157530in}}{\pgfqpoint{5.366161in}{0.259920in}}%
\pgfpathcurveto{\pgfqpoint{5.570940in}{0.362309in}}{\pgfqpoint{5.686000in}{0.501199in}}{\pgfqpoint{5.686000in}{0.646000in}}%
\pgfpathcurveto{\pgfqpoint{5.686000in}{0.790801in}}{\pgfqpoint{5.570940in}{0.929691in}}{\pgfqpoint{5.366161in}{1.032080in}}%
\pgfpathcurveto{\pgfqpoint{5.161381in}{1.134470in}}{\pgfqpoint{4.883602in}{1.192000in}}{\pgfqpoint{4.594000in}{1.192000in}}%
\pgfpathcurveto{\pgfqpoint{4.304398in}{1.192000in}}{\pgfqpoint{4.026619in}{1.134470in}}{\pgfqpoint{3.821839in}{1.032080in}}%
\pgfpathcurveto{\pgfqpoint{3.617060in}{0.929691in}}{\pgfqpoint{3.502000in}{0.790801in}}{\pgfqpoint{3.502000in}{0.646000in}}%
\pgfpathcurveto{\pgfqpoint{3.502000in}{0.501199in}}{\pgfqpoint{3.617060in}{0.362309in}}{\pgfqpoint{3.821839in}{0.259920in}}%
\pgfpathcurveto{\pgfqpoint{4.026619in}{0.157530in}}{\pgfqpoint{4.304398in}{0.100000in}}{\pgfqpoint{4.594000in}{0.100000in}}%
\pgfpathlineto{\pgfqpoint{4.594000in}{0.100000in}}%
\pgfpathclose%
\pgfusepath{clip}%
\pgfsetrectcap%
\pgfsetroundjoin%
\pgfsetlinewidth{0.451687pt}%
\definecolor{currentstroke}{rgb}{0.690196,0.690196,0.690196}%
\pgfsetstrokecolor{currentstroke}%
\pgfsetstrokeopacity{0.250000}%
\pgfsetdash{}{0pt}%
\pgfpathmoveto{\pgfqpoint{4.825047in}{0.151295in}}%
\pgfpathlineto{\pgfqpoint{4.844106in}{0.160652in}}%
\pgfpathlineto{\pgfqpoint{4.862380in}{0.170513in}}%
\pgfpathlineto{\pgfqpoint{4.879447in}{0.180558in}}%
\pgfpathlineto{\pgfqpoint{4.895746in}{0.190956in}}%
\pgfpathlineto{\pgfqpoint{4.911320in}{0.201676in}}%
\pgfpathlineto{\pgfqpoint{4.926209in}{0.212695in}}%
\pgfpathlineto{\pgfqpoint{4.940449in}{0.223994in}}%
\pgfpathlineto{\pgfqpoint{4.954071in}{0.235556in}}%
\pgfpathlineto{\pgfqpoint{4.967104in}{0.247365in}}%
\pgfpathlineto{\pgfqpoint{4.979570in}{0.259410in}}%
\pgfpathlineto{\pgfqpoint{4.991490in}{0.271677in}}%
\pgfpathlineto{\pgfqpoint{5.002881in}{0.284154in}}%
\pgfpathlineto{\pgfqpoint{5.013758in}{0.296831in}}%
\pgfpathlineto{\pgfqpoint{5.024136in}{0.309698in}}%
\pgfpathlineto{\pgfqpoint{5.034144in}{0.322907in}}%
\pgfpathlineto{\pgfqpoint{5.043517in}{0.336079in}}%
\pgfpathlineto{\pgfqpoint{5.052431in}{0.349425in}}%
\pgfpathlineto{\pgfqpoint{5.060892in}{0.362931in}}%
\pgfpathlineto{\pgfqpoint{5.068904in}{0.376589in}}%
\pgfpathlineto{\pgfqpoint{5.076471in}{0.390387in}}%
\pgfpathlineto{\pgfqpoint{5.083597in}{0.404317in}}%
\pgfpathlineto{\pgfqpoint{5.090287in}{0.418370in}}%
\pgfpathlineto{\pgfqpoint{5.096543in}{0.432537in}}%
\pgfpathlineto{\pgfqpoint{5.102369in}{0.446809in}}%
\pgfpathlineto{\pgfqpoint{5.107768in}{0.461179in}}%
\pgfpathlineto{\pgfqpoint{5.112742in}{0.475639in}}%
\pgfpathlineto{\pgfqpoint{5.117294in}{0.490181in}}%
\pgfpathlineto{\pgfqpoint{5.121426in}{0.504798in}}%
\pgfpathlineto{\pgfqpoint{5.125139in}{0.519482in}}%
\pgfpathlineto{\pgfqpoint{5.128437in}{0.534226in}}%
\pgfpathlineto{\pgfqpoint{5.131319in}{0.549023in}}%
\pgfpathlineto{\pgfqpoint{5.133812in}{0.564029in}}%
\pgfpathlineto{\pgfqpoint{5.135853in}{0.578836in}}%
\pgfpathlineto{\pgfqpoint{5.137490in}{0.593701in}}%
\pgfpathlineto{\pgfqpoint{5.138718in}{0.608611in}}%
\pgfpathlineto{\pgfqpoint{5.139538in}{0.623554in}}%
\pgfpathlineto{\pgfqpoint{5.139949in}{0.638516in}}%
\pgfpathlineto{\pgfqpoint{5.139949in}{0.653484in}}%
\pgfpathlineto{\pgfqpoint{5.139538in}{0.668446in}}%
\pgfpathlineto{\pgfqpoint{5.138718in}{0.683389in}}%
\pgfpathlineto{\pgfqpoint{5.137490in}{0.698299in}}%
\pgfpathlineto{\pgfqpoint{5.135853in}{0.713164in}}%
\pgfpathlineto{\pgfqpoint{5.133812in}{0.727971in}}%
\pgfpathlineto{\pgfqpoint{5.131319in}{0.742977in}}%
\pgfpathlineto{\pgfqpoint{5.128437in}{0.757774in}}%
\pgfpathlineto{\pgfqpoint{5.125139in}{0.772518in}}%
\pgfpathlineto{\pgfqpoint{5.121426in}{0.787202in}}%
\pgfpathlineto{\pgfqpoint{5.117294in}{0.801819in}}%
\pgfpathlineto{\pgfqpoint{5.112742in}{0.816361in}}%
\pgfpathlineto{\pgfqpoint{5.107768in}{0.830821in}}%
\pgfpathlineto{\pgfqpoint{5.102369in}{0.845191in}}%
\pgfpathlineto{\pgfqpoint{5.096543in}{0.859463in}}%
\pgfpathlineto{\pgfqpoint{5.090287in}{0.873630in}}%
\pgfpathlineto{\pgfqpoint{5.083597in}{0.887683in}}%
\pgfpathlineto{\pgfqpoint{5.076471in}{0.901613in}}%
\pgfpathlineto{\pgfqpoint{5.068904in}{0.915411in}}%
\pgfpathlineto{\pgfqpoint{5.060892in}{0.929069in}}%
\pgfpathlineto{\pgfqpoint{5.052431in}{0.942575in}}%
\pgfpathlineto{\pgfqpoint{5.043517in}{0.955921in}}%
\pgfpathlineto{\pgfqpoint{5.034144in}{0.969093in}}%
\pgfpathlineto{\pgfqpoint{5.024136in}{0.982302in}}%
\pgfpathlineto{\pgfqpoint{5.013758in}{0.995169in}}%
\pgfpathlineto{\pgfqpoint{5.002881in}{1.007846in}}%
\pgfpathlineto{\pgfqpoint{4.991490in}{1.020323in}}%
\pgfpathlineto{\pgfqpoint{4.979570in}{1.032590in}}%
\pgfpathlineto{\pgfqpoint{4.967104in}{1.044635in}}%
\pgfpathlineto{\pgfqpoint{4.954071in}{1.056444in}}%
\pgfpathlineto{\pgfqpoint{4.940449in}{1.068006in}}%
\pgfpathlineto{\pgfqpoint{4.926209in}{1.079305in}}%
\pgfpathlineto{\pgfqpoint{4.911320in}{1.090324in}}%
\pgfpathlineto{\pgfqpoint{4.895746in}{1.101044in}}%
\pgfpathlineto{\pgfqpoint{4.879447in}{1.111442in}}%
\pgfpathlineto{\pgfqpoint{4.862380in}{1.121487in}}%
\pgfpathlineto{\pgfqpoint{4.844106in}{1.131348in}}%
\pgfpathlineto{\pgfqpoint{4.825047in}{1.140705in}}%
\pgfusepath{stroke}%
\end{pgfscope}%
\begin{pgfscope}%
\pgfpathmoveto{\pgfqpoint{4.594000in}{0.100000in}}%
\pgfpathcurveto{\pgfqpoint{4.883602in}{0.100000in}}{\pgfqpoint{5.161381in}{0.157530in}}{\pgfqpoint{5.366161in}{0.259920in}}%
\pgfpathcurveto{\pgfqpoint{5.570940in}{0.362309in}}{\pgfqpoint{5.686000in}{0.501199in}}{\pgfqpoint{5.686000in}{0.646000in}}%
\pgfpathcurveto{\pgfqpoint{5.686000in}{0.790801in}}{\pgfqpoint{5.570940in}{0.929691in}}{\pgfqpoint{5.366161in}{1.032080in}}%
\pgfpathcurveto{\pgfqpoint{5.161381in}{1.134470in}}{\pgfqpoint{4.883602in}{1.192000in}}{\pgfqpoint{4.594000in}{1.192000in}}%
\pgfpathcurveto{\pgfqpoint{4.304398in}{1.192000in}}{\pgfqpoint{4.026619in}{1.134470in}}{\pgfqpoint{3.821839in}{1.032080in}}%
\pgfpathcurveto{\pgfqpoint{3.617060in}{0.929691in}}{\pgfqpoint{3.502000in}{0.790801in}}{\pgfqpoint{3.502000in}{0.646000in}}%
\pgfpathcurveto{\pgfqpoint{3.502000in}{0.501199in}}{\pgfqpoint{3.617060in}{0.362309in}}{\pgfqpoint{3.821839in}{0.259920in}}%
\pgfpathcurveto{\pgfqpoint{4.026619in}{0.157530in}}{\pgfqpoint{4.304398in}{0.100000in}}{\pgfqpoint{4.594000in}{0.100000in}}%
\pgfpathlineto{\pgfqpoint{4.594000in}{0.100000in}}%
\pgfpathclose%
\pgfusepath{clip}%
\pgfsetrectcap%
\pgfsetroundjoin%
\pgfsetlinewidth{0.451687pt}%
\definecolor{currentstroke}{rgb}{0.690196,0.690196,0.690196}%
\pgfsetstrokecolor{currentstroke}%
\pgfsetstrokeopacity{0.250000}%
\pgfsetdash{}{0pt}%
\pgfpathmoveto{\pgfqpoint{4.902063in}{0.151295in}}%
\pgfpathlineto{\pgfqpoint{4.927475in}{0.160652in}}%
\pgfpathlineto{\pgfqpoint{4.951840in}{0.170513in}}%
\pgfpathlineto{\pgfqpoint{4.974596in}{0.180558in}}%
\pgfpathlineto{\pgfqpoint{4.996328in}{0.190956in}}%
\pgfpathlineto{\pgfqpoint{5.017093in}{0.201676in}}%
\pgfpathlineto{\pgfqpoint{5.036945in}{0.212695in}}%
\pgfpathlineto{\pgfqpoint{5.055932in}{0.223994in}}%
\pgfpathlineto{\pgfqpoint{5.074095in}{0.235556in}}%
\pgfpathlineto{\pgfqpoint{5.091472in}{0.247365in}}%
\pgfpathlineto{\pgfqpoint{5.108093in}{0.259410in}}%
\pgfpathlineto{\pgfqpoint{5.123986in}{0.271677in}}%
\pgfpathlineto{\pgfqpoint{5.139174in}{0.284154in}}%
\pgfpathlineto{\pgfqpoint{5.153678in}{0.296831in}}%
\pgfpathlineto{\pgfqpoint{5.167515in}{0.309698in}}%
\pgfpathlineto{\pgfqpoint{5.180859in}{0.322907in}}%
\pgfpathlineto{\pgfqpoint{5.193356in}{0.336079in}}%
\pgfpathlineto{\pgfqpoint{5.205241in}{0.349425in}}%
\pgfpathlineto{\pgfqpoint{5.216523in}{0.362931in}}%
\pgfpathlineto{\pgfqpoint{5.227205in}{0.376589in}}%
\pgfpathlineto{\pgfqpoint{5.237295in}{0.390387in}}%
\pgfpathlineto{\pgfqpoint{5.246796in}{0.404317in}}%
\pgfpathlineto{\pgfqpoint{5.255716in}{0.418370in}}%
\pgfpathlineto{\pgfqpoint{5.264057in}{0.432537in}}%
\pgfpathlineto{\pgfqpoint{5.271825in}{0.446809in}}%
\pgfpathlineto{\pgfqpoint{5.279024in}{0.461179in}}%
\pgfpathlineto{\pgfqpoint{5.285656in}{0.475639in}}%
\pgfpathlineto{\pgfqpoint{5.291725in}{0.490181in}}%
\pgfpathlineto{\pgfqpoint{5.297234in}{0.504798in}}%
\pgfpathlineto{\pgfqpoint{5.302186in}{0.519482in}}%
\pgfpathlineto{\pgfqpoint{5.306582in}{0.534226in}}%
\pgfpathlineto{\pgfqpoint{5.310425in}{0.549023in}}%
\pgfpathlineto{\pgfqpoint{5.313749in}{0.564029in}}%
\pgfpathlineto{\pgfqpoint{5.316471in}{0.578836in}}%
\pgfpathlineto{\pgfqpoint{5.318653in}{0.593701in}}%
\pgfpathlineto{\pgfqpoint{5.320291in}{0.608611in}}%
\pgfpathlineto{\pgfqpoint{5.321385in}{0.623554in}}%
\pgfpathlineto{\pgfqpoint{5.321932in}{0.638516in}}%
\pgfpathlineto{\pgfqpoint{5.321932in}{0.653484in}}%
\pgfpathlineto{\pgfqpoint{5.321385in}{0.668446in}}%
\pgfpathlineto{\pgfqpoint{5.320291in}{0.683389in}}%
\pgfpathlineto{\pgfqpoint{5.318653in}{0.698299in}}%
\pgfpathlineto{\pgfqpoint{5.316471in}{0.713164in}}%
\pgfpathlineto{\pgfqpoint{5.313749in}{0.727971in}}%
\pgfpathlineto{\pgfqpoint{5.310425in}{0.742977in}}%
\pgfpathlineto{\pgfqpoint{5.306582in}{0.757774in}}%
\pgfpathlineto{\pgfqpoint{5.302186in}{0.772518in}}%
\pgfpathlineto{\pgfqpoint{5.297234in}{0.787202in}}%
\pgfpathlineto{\pgfqpoint{5.291725in}{0.801819in}}%
\pgfpathlineto{\pgfqpoint{5.285656in}{0.816361in}}%
\pgfpathlineto{\pgfqpoint{5.279024in}{0.830821in}}%
\pgfpathlineto{\pgfqpoint{5.271825in}{0.845191in}}%
\pgfpathlineto{\pgfqpoint{5.264057in}{0.859463in}}%
\pgfpathlineto{\pgfqpoint{5.255716in}{0.873630in}}%
\pgfpathlineto{\pgfqpoint{5.246796in}{0.887683in}}%
\pgfpathlineto{\pgfqpoint{5.237295in}{0.901613in}}%
\pgfpathlineto{\pgfqpoint{5.227205in}{0.915411in}}%
\pgfpathlineto{\pgfqpoint{5.216523in}{0.929069in}}%
\pgfpathlineto{\pgfqpoint{5.205241in}{0.942575in}}%
\pgfpathlineto{\pgfqpoint{5.193356in}{0.955921in}}%
\pgfpathlineto{\pgfqpoint{5.180859in}{0.969093in}}%
\pgfpathlineto{\pgfqpoint{5.167515in}{0.982302in}}%
\pgfpathlineto{\pgfqpoint{5.153678in}{0.995169in}}%
\pgfpathlineto{\pgfqpoint{5.139174in}{1.007846in}}%
\pgfpathlineto{\pgfqpoint{5.123986in}{1.020323in}}%
\pgfpathlineto{\pgfqpoint{5.108093in}{1.032590in}}%
\pgfpathlineto{\pgfqpoint{5.091472in}{1.044635in}}%
\pgfpathlineto{\pgfqpoint{5.074095in}{1.056444in}}%
\pgfpathlineto{\pgfqpoint{5.055932in}{1.068006in}}%
\pgfpathlineto{\pgfqpoint{5.036945in}{1.079305in}}%
\pgfpathlineto{\pgfqpoint{5.017093in}{1.090324in}}%
\pgfpathlineto{\pgfqpoint{4.996328in}{1.101044in}}%
\pgfpathlineto{\pgfqpoint{4.974596in}{1.111442in}}%
\pgfpathlineto{\pgfqpoint{4.951840in}{1.121487in}}%
\pgfpathlineto{\pgfqpoint{4.927475in}{1.131348in}}%
\pgfpathlineto{\pgfqpoint{4.902063in}{1.140705in}}%
\pgfusepath{stroke}%
\end{pgfscope}%
\begin{pgfscope}%
\pgfpathmoveto{\pgfqpoint{4.594000in}{0.100000in}}%
\pgfpathcurveto{\pgfqpoint{4.883602in}{0.100000in}}{\pgfqpoint{5.161381in}{0.157530in}}{\pgfqpoint{5.366161in}{0.259920in}}%
\pgfpathcurveto{\pgfqpoint{5.570940in}{0.362309in}}{\pgfqpoint{5.686000in}{0.501199in}}{\pgfqpoint{5.686000in}{0.646000in}}%
\pgfpathcurveto{\pgfqpoint{5.686000in}{0.790801in}}{\pgfqpoint{5.570940in}{0.929691in}}{\pgfqpoint{5.366161in}{1.032080in}}%
\pgfpathcurveto{\pgfqpoint{5.161381in}{1.134470in}}{\pgfqpoint{4.883602in}{1.192000in}}{\pgfqpoint{4.594000in}{1.192000in}}%
\pgfpathcurveto{\pgfqpoint{4.304398in}{1.192000in}}{\pgfqpoint{4.026619in}{1.134470in}}{\pgfqpoint{3.821839in}{1.032080in}}%
\pgfpathcurveto{\pgfqpoint{3.617060in}{0.929691in}}{\pgfqpoint{3.502000in}{0.790801in}}{\pgfqpoint{3.502000in}{0.646000in}}%
\pgfpathcurveto{\pgfqpoint{3.502000in}{0.501199in}}{\pgfqpoint{3.617060in}{0.362309in}}{\pgfqpoint{3.821839in}{0.259920in}}%
\pgfpathcurveto{\pgfqpoint{4.026619in}{0.157530in}}{\pgfqpoint{4.304398in}{0.100000in}}{\pgfqpoint{4.594000in}{0.100000in}}%
\pgfpathlineto{\pgfqpoint{4.594000in}{0.100000in}}%
\pgfpathclose%
\pgfusepath{clip}%
\pgfsetrectcap%
\pgfsetroundjoin%
\pgfsetlinewidth{0.451687pt}%
\definecolor{currentstroke}{rgb}{0.690196,0.690196,0.690196}%
\pgfsetstrokecolor{currentstroke}%
\pgfsetstrokeopacity{0.250000}%
\pgfsetdash{}{0pt}%
\pgfpathmoveto{\pgfqpoint{4.979079in}{0.151295in}}%
\pgfpathlineto{\pgfqpoint{5.010843in}{0.160652in}}%
\pgfpathlineto{\pgfqpoint{5.041300in}{0.170513in}}%
\pgfpathlineto{\pgfqpoint{5.069745in}{0.180558in}}%
\pgfpathlineto{\pgfqpoint{5.096910in}{0.190956in}}%
\pgfpathlineto{\pgfqpoint{5.122866in}{0.201676in}}%
\pgfpathlineto{\pgfqpoint{5.147681in}{0.212695in}}%
\pgfpathlineto{\pgfqpoint{5.171415in}{0.223994in}}%
\pgfpathlineto{\pgfqpoint{5.194119in}{0.235556in}}%
\pgfpathlineto{\pgfqpoint{5.215840in}{0.247365in}}%
\pgfpathlineto{\pgfqpoint{5.236616in}{0.259410in}}%
\pgfpathlineto{\pgfqpoint{5.256483in}{0.271677in}}%
\pgfpathlineto{\pgfqpoint{5.275468in}{0.284154in}}%
\pgfpathlineto{\pgfqpoint{5.293597in}{0.296831in}}%
\pgfpathlineto{\pgfqpoint{5.310894in}{0.309698in}}%
\pgfpathlineto{\pgfqpoint{5.327574in}{0.322907in}}%
\pgfpathlineto{\pgfqpoint{5.343194in}{0.336079in}}%
\pgfpathlineto{\pgfqpoint{5.358052in}{0.349425in}}%
\pgfpathlineto{\pgfqpoint{5.372153in}{0.362931in}}%
\pgfpathlineto{\pgfqpoint{5.385506in}{0.376589in}}%
\pgfpathlineto{\pgfqpoint{5.398118in}{0.390387in}}%
\pgfpathlineto{\pgfqpoint{5.409996in}{0.404317in}}%
\pgfpathlineto{\pgfqpoint{5.421145in}{0.418370in}}%
\pgfpathlineto{\pgfqpoint{5.431572in}{0.432537in}}%
\pgfpathlineto{\pgfqpoint{5.441282in}{0.446809in}}%
\pgfpathlineto{\pgfqpoint{5.450280in}{0.461179in}}%
\pgfpathlineto{\pgfqpoint{5.458570in}{0.475639in}}%
\pgfpathlineto{\pgfqpoint{5.466156in}{0.490181in}}%
\pgfpathlineto{\pgfqpoint{5.473043in}{0.504798in}}%
\pgfpathlineto{\pgfqpoint{5.479232in}{0.519482in}}%
\pgfpathlineto{\pgfqpoint{5.484728in}{0.534226in}}%
\pgfpathlineto{\pgfqpoint{5.489531in}{0.549023in}}%
\pgfpathlineto{\pgfqpoint{5.493686in}{0.564029in}}%
\pgfpathlineto{\pgfqpoint{5.497089in}{0.578836in}}%
\pgfpathlineto{\pgfqpoint{5.499816in}{0.593701in}}%
\pgfpathlineto{\pgfqpoint{5.501864in}{0.608611in}}%
\pgfpathlineto{\pgfqpoint{5.503231in}{0.623554in}}%
\pgfpathlineto{\pgfqpoint{5.503915in}{0.638516in}}%
\pgfpathlineto{\pgfqpoint{5.503915in}{0.653484in}}%
\pgfpathlineto{\pgfqpoint{5.503231in}{0.668446in}}%
\pgfpathlineto{\pgfqpoint{5.501864in}{0.683389in}}%
\pgfpathlineto{\pgfqpoint{5.499816in}{0.698299in}}%
\pgfpathlineto{\pgfqpoint{5.497089in}{0.713164in}}%
\pgfpathlineto{\pgfqpoint{5.493686in}{0.727971in}}%
\pgfpathlineto{\pgfqpoint{5.489531in}{0.742977in}}%
\pgfpathlineto{\pgfqpoint{5.484728in}{0.757774in}}%
\pgfpathlineto{\pgfqpoint{5.479232in}{0.772518in}}%
\pgfpathlineto{\pgfqpoint{5.473043in}{0.787202in}}%
\pgfpathlineto{\pgfqpoint{5.466156in}{0.801819in}}%
\pgfpathlineto{\pgfqpoint{5.458570in}{0.816361in}}%
\pgfpathlineto{\pgfqpoint{5.450280in}{0.830821in}}%
\pgfpathlineto{\pgfqpoint{5.441282in}{0.845191in}}%
\pgfpathlineto{\pgfqpoint{5.431572in}{0.859463in}}%
\pgfpathlineto{\pgfqpoint{5.421145in}{0.873630in}}%
\pgfpathlineto{\pgfqpoint{5.409996in}{0.887683in}}%
\pgfpathlineto{\pgfqpoint{5.398118in}{0.901613in}}%
\pgfpathlineto{\pgfqpoint{5.385506in}{0.915411in}}%
\pgfpathlineto{\pgfqpoint{5.372153in}{0.929069in}}%
\pgfpathlineto{\pgfqpoint{5.358052in}{0.942575in}}%
\pgfpathlineto{\pgfqpoint{5.343194in}{0.955921in}}%
\pgfpathlineto{\pgfqpoint{5.327574in}{0.969093in}}%
\pgfpathlineto{\pgfqpoint{5.310894in}{0.982302in}}%
\pgfpathlineto{\pgfqpoint{5.293597in}{0.995169in}}%
\pgfpathlineto{\pgfqpoint{5.275468in}{1.007846in}}%
\pgfpathlineto{\pgfqpoint{5.256483in}{1.020323in}}%
\pgfpathlineto{\pgfqpoint{5.236616in}{1.032590in}}%
\pgfpathlineto{\pgfqpoint{5.215840in}{1.044635in}}%
\pgfpathlineto{\pgfqpoint{5.194119in}{1.056444in}}%
\pgfpathlineto{\pgfqpoint{5.171415in}{1.068006in}}%
\pgfpathlineto{\pgfqpoint{5.147681in}{1.079305in}}%
\pgfpathlineto{\pgfqpoint{5.122866in}{1.090324in}}%
\pgfpathlineto{\pgfqpoint{5.096910in}{1.101044in}}%
\pgfpathlineto{\pgfqpoint{5.069745in}{1.111442in}}%
\pgfpathlineto{\pgfqpoint{5.041300in}{1.121487in}}%
\pgfpathlineto{\pgfqpoint{5.010843in}{1.131348in}}%
\pgfpathlineto{\pgfqpoint{4.979079in}{1.140705in}}%
\pgfusepath{stroke}%
\end{pgfscope}%
\begin{pgfscope}%
\pgfpathmoveto{\pgfqpoint{4.594000in}{0.100000in}}%
\pgfpathcurveto{\pgfqpoint{4.883602in}{0.100000in}}{\pgfqpoint{5.161381in}{0.157530in}}{\pgfqpoint{5.366161in}{0.259920in}}%
\pgfpathcurveto{\pgfqpoint{5.570940in}{0.362309in}}{\pgfqpoint{5.686000in}{0.501199in}}{\pgfqpoint{5.686000in}{0.646000in}}%
\pgfpathcurveto{\pgfqpoint{5.686000in}{0.790801in}}{\pgfqpoint{5.570940in}{0.929691in}}{\pgfqpoint{5.366161in}{1.032080in}}%
\pgfpathcurveto{\pgfqpoint{5.161381in}{1.134470in}}{\pgfqpoint{4.883602in}{1.192000in}}{\pgfqpoint{4.594000in}{1.192000in}}%
\pgfpathcurveto{\pgfqpoint{4.304398in}{1.192000in}}{\pgfqpoint{4.026619in}{1.134470in}}{\pgfqpoint{3.821839in}{1.032080in}}%
\pgfpathcurveto{\pgfqpoint{3.617060in}{0.929691in}}{\pgfqpoint{3.502000in}{0.790801in}}{\pgfqpoint{3.502000in}{0.646000in}}%
\pgfpathcurveto{\pgfqpoint{3.502000in}{0.501199in}}{\pgfqpoint{3.617060in}{0.362309in}}{\pgfqpoint{3.821839in}{0.259920in}}%
\pgfpathcurveto{\pgfqpoint{4.026619in}{0.157530in}}{\pgfqpoint{4.304398in}{0.100000in}}{\pgfqpoint{4.594000in}{0.100000in}}%
\pgfpathlineto{\pgfqpoint{4.594000in}{0.100000in}}%
\pgfpathclose%
\pgfusepath{clip}%
\pgfsetrectcap%
\pgfsetroundjoin%
\pgfsetlinewidth{0.451687pt}%
\definecolor{currentstroke}{rgb}{0.690196,0.690196,0.690196}%
\pgfsetstrokecolor{currentstroke}%
\pgfsetstrokeopacity{0.250000}%
\pgfsetdash{}{0pt}%
\pgfpathmoveto{\pgfqpoint{4.131906in}{0.151295in}}%
\pgfpathlineto{\pgfqpoint{4.144228in}{0.151295in}}%
\pgfpathlineto{\pgfqpoint{4.156551in}{0.151295in}}%
\pgfpathlineto{\pgfqpoint{4.168873in}{0.151295in}}%
\pgfpathlineto{\pgfqpoint{4.181196in}{0.151295in}}%
\pgfpathlineto{\pgfqpoint{4.193518in}{0.151295in}}%
\pgfpathlineto{\pgfqpoint{4.205841in}{0.151295in}}%
\pgfpathlineto{\pgfqpoint{4.218163in}{0.151295in}}%
\pgfpathlineto{\pgfqpoint{4.230486in}{0.151295in}}%
\pgfpathlineto{\pgfqpoint{4.242808in}{0.151295in}}%
\pgfpathlineto{\pgfqpoint{4.255131in}{0.151295in}}%
\pgfpathlineto{\pgfqpoint{4.267453in}{0.151295in}}%
\pgfpathlineto{\pgfqpoint{4.279776in}{0.151295in}}%
\pgfpathlineto{\pgfqpoint{4.292098in}{0.151295in}}%
\pgfpathlineto{\pgfqpoint{4.304421in}{0.151295in}}%
\pgfpathlineto{\pgfqpoint{4.316743in}{0.151295in}}%
\pgfpathlineto{\pgfqpoint{4.329066in}{0.151295in}}%
\pgfpathlineto{\pgfqpoint{4.341388in}{0.151295in}}%
\pgfpathlineto{\pgfqpoint{4.353711in}{0.151295in}}%
\pgfpathlineto{\pgfqpoint{4.366033in}{0.151295in}}%
\pgfpathlineto{\pgfqpoint{4.378356in}{0.151295in}}%
\pgfpathlineto{\pgfqpoint{4.390678in}{0.151295in}}%
\pgfpathlineto{\pgfqpoint{4.403001in}{0.151295in}}%
\pgfpathlineto{\pgfqpoint{4.415323in}{0.151295in}}%
\pgfpathlineto{\pgfqpoint{4.427646in}{0.151295in}}%
\pgfpathlineto{\pgfqpoint{4.439969in}{0.151295in}}%
\pgfpathlineto{\pgfqpoint{4.452291in}{0.151295in}}%
\pgfpathlineto{\pgfqpoint{4.464614in}{0.151295in}}%
\pgfpathlineto{\pgfqpoint{4.476936in}{0.151295in}}%
\pgfpathlineto{\pgfqpoint{4.489259in}{0.151295in}}%
\pgfpathlineto{\pgfqpoint{4.501581in}{0.151295in}}%
\pgfpathlineto{\pgfqpoint{4.513904in}{0.151295in}}%
\pgfpathlineto{\pgfqpoint{4.526226in}{0.151295in}}%
\pgfpathlineto{\pgfqpoint{4.538549in}{0.151295in}}%
\pgfpathlineto{\pgfqpoint{4.550871in}{0.151295in}}%
\pgfpathlineto{\pgfqpoint{4.563194in}{0.151295in}}%
\pgfpathlineto{\pgfqpoint{4.575516in}{0.151295in}}%
\pgfpathlineto{\pgfqpoint{4.587839in}{0.151295in}}%
\pgfpathlineto{\pgfqpoint{4.600161in}{0.151295in}}%
\pgfpathlineto{\pgfqpoint{4.612484in}{0.151295in}}%
\pgfpathlineto{\pgfqpoint{4.624806in}{0.151295in}}%
\pgfpathlineto{\pgfqpoint{4.637129in}{0.151295in}}%
\pgfpathlineto{\pgfqpoint{4.649451in}{0.151295in}}%
\pgfpathlineto{\pgfqpoint{4.661774in}{0.151295in}}%
\pgfpathlineto{\pgfqpoint{4.674096in}{0.151295in}}%
\pgfpathlineto{\pgfqpoint{4.686419in}{0.151295in}}%
\pgfpathlineto{\pgfqpoint{4.698741in}{0.151295in}}%
\pgfpathlineto{\pgfqpoint{4.711064in}{0.151295in}}%
\pgfpathlineto{\pgfqpoint{4.723386in}{0.151295in}}%
\pgfpathlineto{\pgfqpoint{4.735709in}{0.151295in}}%
\pgfpathlineto{\pgfqpoint{4.748031in}{0.151295in}}%
\pgfpathlineto{\pgfqpoint{4.760354in}{0.151295in}}%
\pgfpathlineto{\pgfqpoint{4.772677in}{0.151295in}}%
\pgfpathlineto{\pgfqpoint{4.784999in}{0.151295in}}%
\pgfpathlineto{\pgfqpoint{4.797322in}{0.151295in}}%
\pgfpathlineto{\pgfqpoint{4.809644in}{0.151295in}}%
\pgfpathlineto{\pgfqpoint{4.821967in}{0.151295in}}%
\pgfpathlineto{\pgfqpoint{4.834289in}{0.151295in}}%
\pgfpathlineto{\pgfqpoint{4.846612in}{0.151295in}}%
\pgfpathlineto{\pgfqpoint{4.858934in}{0.151295in}}%
\pgfpathlineto{\pgfqpoint{4.871257in}{0.151295in}}%
\pgfpathlineto{\pgfqpoint{4.883579in}{0.151295in}}%
\pgfpathlineto{\pgfqpoint{4.895902in}{0.151295in}}%
\pgfpathlineto{\pgfqpoint{4.908224in}{0.151295in}}%
\pgfpathlineto{\pgfqpoint{4.920547in}{0.151295in}}%
\pgfpathlineto{\pgfqpoint{4.932869in}{0.151295in}}%
\pgfpathlineto{\pgfqpoint{4.945192in}{0.151295in}}%
\pgfpathlineto{\pgfqpoint{4.957514in}{0.151295in}}%
\pgfpathlineto{\pgfqpoint{4.969837in}{0.151295in}}%
\pgfpathlineto{\pgfqpoint{4.982159in}{0.151295in}}%
\pgfpathlineto{\pgfqpoint{4.994482in}{0.151295in}}%
\pgfpathlineto{\pgfqpoint{5.006804in}{0.151295in}}%
\pgfpathlineto{\pgfqpoint{5.019127in}{0.151295in}}%
\pgfpathlineto{\pgfqpoint{5.031449in}{0.151295in}}%
\pgfpathlineto{\pgfqpoint{5.043772in}{0.151295in}}%
\pgfpathlineto{\pgfqpoint{5.056094in}{0.151295in}}%
\pgfusepath{stroke}%
\end{pgfscope}%
\begin{pgfscope}%
\pgfpathmoveto{\pgfqpoint{4.594000in}{0.100000in}}%
\pgfpathcurveto{\pgfqpoint{4.883602in}{0.100000in}}{\pgfqpoint{5.161381in}{0.157530in}}{\pgfqpoint{5.366161in}{0.259920in}}%
\pgfpathcurveto{\pgfqpoint{5.570940in}{0.362309in}}{\pgfqpoint{5.686000in}{0.501199in}}{\pgfqpoint{5.686000in}{0.646000in}}%
\pgfpathcurveto{\pgfqpoint{5.686000in}{0.790801in}}{\pgfqpoint{5.570940in}{0.929691in}}{\pgfqpoint{5.366161in}{1.032080in}}%
\pgfpathcurveto{\pgfqpoint{5.161381in}{1.134470in}}{\pgfqpoint{4.883602in}{1.192000in}}{\pgfqpoint{4.594000in}{1.192000in}}%
\pgfpathcurveto{\pgfqpoint{4.304398in}{1.192000in}}{\pgfqpoint{4.026619in}{1.134470in}}{\pgfqpoint{3.821839in}{1.032080in}}%
\pgfpathcurveto{\pgfqpoint{3.617060in}{0.929691in}}{\pgfqpoint{3.502000in}{0.790801in}}{\pgfqpoint{3.502000in}{0.646000in}}%
\pgfpathcurveto{\pgfqpoint{3.502000in}{0.501199in}}{\pgfqpoint{3.617060in}{0.362309in}}{\pgfqpoint{3.821839in}{0.259920in}}%
\pgfpathcurveto{\pgfqpoint{4.026619in}{0.157530in}}{\pgfqpoint{4.304398in}{0.100000in}}{\pgfqpoint{4.594000in}{0.100000in}}%
\pgfpathlineto{\pgfqpoint{4.594000in}{0.100000in}}%
\pgfpathclose%
\pgfusepath{clip}%
\pgfsetrectcap%
\pgfsetroundjoin%
\pgfsetlinewidth{0.451687pt}%
\definecolor{currentstroke}{rgb}{0.690196,0.690196,0.690196}%
\pgfsetstrokecolor{currentstroke}%
\pgfsetstrokeopacity{0.250000}%
\pgfsetdash{}{0pt}%
\pgfpathmoveto{\pgfqpoint{3.887329in}{0.229743in}}%
\pgfpathlineto{\pgfqpoint{3.906174in}{0.229743in}}%
\pgfpathlineto{\pgfqpoint{3.925018in}{0.229743in}}%
\pgfpathlineto{\pgfqpoint{3.943863in}{0.229743in}}%
\pgfpathlineto{\pgfqpoint{3.962707in}{0.229743in}}%
\pgfpathlineto{\pgfqpoint{3.981552in}{0.229743in}}%
\pgfpathlineto{\pgfqpoint{4.000396in}{0.229743in}}%
\pgfpathlineto{\pgfqpoint{4.019241in}{0.229743in}}%
\pgfpathlineto{\pgfqpoint{4.038086in}{0.229743in}}%
\pgfpathlineto{\pgfqpoint{4.056930in}{0.229743in}}%
\pgfpathlineto{\pgfqpoint{4.075775in}{0.229743in}}%
\pgfpathlineto{\pgfqpoint{4.094619in}{0.229743in}}%
\pgfpathlineto{\pgfqpoint{4.113464in}{0.229743in}}%
\pgfpathlineto{\pgfqpoint{4.132308in}{0.229743in}}%
\pgfpathlineto{\pgfqpoint{4.151153in}{0.229743in}}%
\pgfpathlineto{\pgfqpoint{4.169997in}{0.229743in}}%
\pgfpathlineto{\pgfqpoint{4.188842in}{0.229743in}}%
\pgfpathlineto{\pgfqpoint{4.207687in}{0.229743in}}%
\pgfpathlineto{\pgfqpoint{4.226531in}{0.229743in}}%
\pgfpathlineto{\pgfqpoint{4.245376in}{0.229743in}}%
\pgfpathlineto{\pgfqpoint{4.264220in}{0.229743in}}%
\pgfpathlineto{\pgfqpoint{4.283065in}{0.229743in}}%
\pgfpathlineto{\pgfqpoint{4.301909in}{0.229743in}}%
\pgfpathlineto{\pgfqpoint{4.320754in}{0.229743in}}%
\pgfpathlineto{\pgfqpoint{4.339598in}{0.229743in}}%
\pgfpathlineto{\pgfqpoint{4.358443in}{0.229743in}}%
\pgfpathlineto{\pgfqpoint{4.377288in}{0.229743in}}%
\pgfpathlineto{\pgfqpoint{4.396132in}{0.229743in}}%
\pgfpathlineto{\pgfqpoint{4.414977in}{0.229743in}}%
\pgfpathlineto{\pgfqpoint{4.433821in}{0.229743in}}%
\pgfpathlineto{\pgfqpoint{4.452666in}{0.229743in}}%
\pgfpathlineto{\pgfqpoint{4.471510in}{0.229743in}}%
\pgfpathlineto{\pgfqpoint{4.490355in}{0.229743in}}%
\pgfpathlineto{\pgfqpoint{4.509199in}{0.229743in}}%
\pgfpathlineto{\pgfqpoint{4.528044in}{0.229743in}}%
\pgfpathlineto{\pgfqpoint{4.546889in}{0.229743in}}%
\pgfpathlineto{\pgfqpoint{4.565733in}{0.229743in}}%
\pgfpathlineto{\pgfqpoint{4.584578in}{0.229743in}}%
\pgfpathlineto{\pgfqpoint{4.603422in}{0.229743in}}%
\pgfpathlineto{\pgfqpoint{4.622267in}{0.229743in}}%
\pgfpathlineto{\pgfqpoint{4.641111in}{0.229743in}}%
\pgfpathlineto{\pgfqpoint{4.659956in}{0.229743in}}%
\pgfpathlineto{\pgfqpoint{4.678801in}{0.229743in}}%
\pgfpathlineto{\pgfqpoint{4.697645in}{0.229743in}}%
\pgfpathlineto{\pgfqpoint{4.716490in}{0.229743in}}%
\pgfpathlineto{\pgfqpoint{4.735334in}{0.229743in}}%
\pgfpathlineto{\pgfqpoint{4.754179in}{0.229743in}}%
\pgfpathlineto{\pgfqpoint{4.773023in}{0.229743in}}%
\pgfpathlineto{\pgfqpoint{4.791868in}{0.229743in}}%
\pgfpathlineto{\pgfqpoint{4.810712in}{0.229743in}}%
\pgfpathlineto{\pgfqpoint{4.829557in}{0.229743in}}%
\pgfpathlineto{\pgfqpoint{4.848402in}{0.229743in}}%
\pgfpathlineto{\pgfqpoint{4.867246in}{0.229743in}}%
\pgfpathlineto{\pgfqpoint{4.886091in}{0.229743in}}%
\pgfpathlineto{\pgfqpoint{4.904935in}{0.229743in}}%
\pgfpathlineto{\pgfqpoint{4.923780in}{0.229743in}}%
\pgfpathlineto{\pgfqpoint{4.942624in}{0.229743in}}%
\pgfpathlineto{\pgfqpoint{4.961469in}{0.229743in}}%
\pgfpathlineto{\pgfqpoint{4.980313in}{0.229743in}}%
\pgfpathlineto{\pgfqpoint{4.999158in}{0.229743in}}%
\pgfpathlineto{\pgfqpoint{5.018003in}{0.229743in}}%
\pgfpathlineto{\pgfqpoint{5.036847in}{0.229743in}}%
\pgfpathlineto{\pgfqpoint{5.055692in}{0.229743in}}%
\pgfpathlineto{\pgfqpoint{5.074536in}{0.229743in}}%
\pgfpathlineto{\pgfqpoint{5.093381in}{0.229743in}}%
\pgfpathlineto{\pgfqpoint{5.112225in}{0.229743in}}%
\pgfpathlineto{\pgfqpoint{5.131070in}{0.229743in}}%
\pgfpathlineto{\pgfqpoint{5.149914in}{0.229743in}}%
\pgfpathlineto{\pgfqpoint{5.168759in}{0.229743in}}%
\pgfpathlineto{\pgfqpoint{5.187604in}{0.229743in}}%
\pgfpathlineto{\pgfqpoint{5.206448in}{0.229743in}}%
\pgfpathlineto{\pgfqpoint{5.225293in}{0.229743in}}%
\pgfpathlineto{\pgfqpoint{5.244137in}{0.229743in}}%
\pgfpathlineto{\pgfqpoint{5.262982in}{0.229743in}}%
\pgfpathlineto{\pgfqpoint{5.281826in}{0.229743in}}%
\pgfpathlineto{\pgfqpoint{5.300671in}{0.229743in}}%
\pgfusepath{stroke}%
\end{pgfscope}%
\begin{pgfscope}%
\pgfpathmoveto{\pgfqpoint{4.594000in}{0.100000in}}%
\pgfpathcurveto{\pgfqpoint{4.883602in}{0.100000in}}{\pgfqpoint{5.161381in}{0.157530in}}{\pgfqpoint{5.366161in}{0.259920in}}%
\pgfpathcurveto{\pgfqpoint{5.570940in}{0.362309in}}{\pgfqpoint{5.686000in}{0.501199in}}{\pgfqpoint{5.686000in}{0.646000in}}%
\pgfpathcurveto{\pgfqpoint{5.686000in}{0.790801in}}{\pgfqpoint{5.570940in}{0.929691in}}{\pgfqpoint{5.366161in}{1.032080in}}%
\pgfpathcurveto{\pgfqpoint{5.161381in}{1.134470in}}{\pgfqpoint{4.883602in}{1.192000in}}{\pgfqpoint{4.594000in}{1.192000in}}%
\pgfpathcurveto{\pgfqpoint{4.304398in}{1.192000in}}{\pgfqpoint{4.026619in}{1.134470in}}{\pgfqpoint{3.821839in}{1.032080in}}%
\pgfpathcurveto{\pgfqpoint{3.617060in}{0.929691in}}{\pgfqpoint{3.502000in}{0.790801in}}{\pgfqpoint{3.502000in}{0.646000in}}%
\pgfpathcurveto{\pgfqpoint{3.502000in}{0.501199in}}{\pgfqpoint{3.617060in}{0.362309in}}{\pgfqpoint{3.821839in}{0.259920in}}%
\pgfpathcurveto{\pgfqpoint{4.026619in}{0.157530in}}{\pgfqpoint{4.304398in}{0.100000in}}{\pgfqpoint{4.594000in}{0.100000in}}%
\pgfpathlineto{\pgfqpoint{4.594000in}{0.100000in}}%
\pgfpathclose%
\pgfusepath{clip}%
\pgfsetrectcap%
\pgfsetroundjoin%
\pgfsetlinewidth{0.451687pt}%
\definecolor{currentstroke}{rgb}{0.690196,0.690196,0.690196}%
\pgfsetstrokecolor{currentstroke}%
\pgfsetstrokeopacity{0.250000}%
\pgfsetdash{}{0pt}%
\pgfpathmoveto{\pgfqpoint{3.713711in}{0.322907in}}%
\pgfpathlineto{\pgfqpoint{3.737186in}{0.322907in}}%
\pgfpathlineto{\pgfqpoint{3.760660in}{0.322907in}}%
\pgfpathlineto{\pgfqpoint{3.784134in}{0.322907in}}%
\pgfpathlineto{\pgfqpoint{3.807609in}{0.322907in}}%
\pgfpathlineto{\pgfqpoint{3.831083in}{0.322907in}}%
\pgfpathlineto{\pgfqpoint{3.854558in}{0.322907in}}%
\pgfpathlineto{\pgfqpoint{3.878032in}{0.322907in}}%
\pgfpathlineto{\pgfqpoint{3.901506in}{0.322907in}}%
\pgfpathlineto{\pgfqpoint{3.924981in}{0.322907in}}%
\pgfpathlineto{\pgfqpoint{3.948455in}{0.322907in}}%
\pgfpathlineto{\pgfqpoint{3.971929in}{0.322907in}}%
\pgfpathlineto{\pgfqpoint{3.995404in}{0.322907in}}%
\pgfpathlineto{\pgfqpoint{4.018878in}{0.322907in}}%
\pgfpathlineto{\pgfqpoint{4.042352in}{0.322907in}}%
\pgfpathlineto{\pgfqpoint{4.065827in}{0.322907in}}%
\pgfpathlineto{\pgfqpoint{4.089301in}{0.322907in}}%
\pgfpathlineto{\pgfqpoint{4.112776in}{0.322907in}}%
\pgfpathlineto{\pgfqpoint{4.136250in}{0.322907in}}%
\pgfpathlineto{\pgfqpoint{4.159724in}{0.322907in}}%
\pgfpathlineto{\pgfqpoint{4.183199in}{0.322907in}}%
\pgfpathlineto{\pgfqpoint{4.206673in}{0.322907in}}%
\pgfpathlineto{\pgfqpoint{4.230147in}{0.322907in}}%
\pgfpathlineto{\pgfqpoint{4.253622in}{0.322907in}}%
\pgfpathlineto{\pgfqpoint{4.277096in}{0.322907in}}%
\pgfpathlineto{\pgfqpoint{4.300570in}{0.322907in}}%
\pgfpathlineto{\pgfqpoint{4.324045in}{0.322907in}}%
\pgfpathlineto{\pgfqpoint{4.347519in}{0.322907in}}%
\pgfpathlineto{\pgfqpoint{4.370994in}{0.322907in}}%
\pgfpathlineto{\pgfqpoint{4.394468in}{0.322907in}}%
\pgfpathlineto{\pgfqpoint{4.417942in}{0.322907in}}%
\pgfpathlineto{\pgfqpoint{4.441417in}{0.322907in}}%
\pgfpathlineto{\pgfqpoint{4.464891in}{0.322907in}}%
\pgfpathlineto{\pgfqpoint{4.488365in}{0.322907in}}%
\pgfpathlineto{\pgfqpoint{4.511840in}{0.322907in}}%
\pgfpathlineto{\pgfqpoint{4.535314in}{0.322907in}}%
\pgfpathlineto{\pgfqpoint{4.558788in}{0.322907in}}%
\pgfpathlineto{\pgfqpoint{4.582263in}{0.322907in}}%
\pgfpathlineto{\pgfqpoint{4.605737in}{0.322907in}}%
\pgfpathlineto{\pgfqpoint{4.629212in}{0.322907in}}%
\pgfpathlineto{\pgfqpoint{4.652686in}{0.322907in}}%
\pgfpathlineto{\pgfqpoint{4.676160in}{0.322907in}}%
\pgfpathlineto{\pgfqpoint{4.699635in}{0.322907in}}%
\pgfpathlineto{\pgfqpoint{4.723109in}{0.322907in}}%
\pgfpathlineto{\pgfqpoint{4.746583in}{0.322907in}}%
\pgfpathlineto{\pgfqpoint{4.770058in}{0.322907in}}%
\pgfpathlineto{\pgfqpoint{4.793532in}{0.322907in}}%
\pgfpathlineto{\pgfqpoint{4.817006in}{0.322907in}}%
\pgfpathlineto{\pgfqpoint{4.840481in}{0.322907in}}%
\pgfpathlineto{\pgfqpoint{4.863955in}{0.322907in}}%
\pgfpathlineto{\pgfqpoint{4.887430in}{0.322907in}}%
\pgfpathlineto{\pgfqpoint{4.910904in}{0.322907in}}%
\pgfpathlineto{\pgfqpoint{4.934378in}{0.322907in}}%
\pgfpathlineto{\pgfqpoint{4.957853in}{0.322907in}}%
\pgfpathlineto{\pgfqpoint{4.981327in}{0.322907in}}%
\pgfpathlineto{\pgfqpoint{5.004801in}{0.322907in}}%
\pgfpathlineto{\pgfqpoint{5.028276in}{0.322907in}}%
\pgfpathlineto{\pgfqpoint{5.051750in}{0.322907in}}%
\pgfpathlineto{\pgfqpoint{5.075224in}{0.322907in}}%
\pgfpathlineto{\pgfqpoint{5.098699in}{0.322907in}}%
\pgfpathlineto{\pgfqpoint{5.122173in}{0.322907in}}%
\pgfpathlineto{\pgfqpoint{5.145648in}{0.322907in}}%
\pgfpathlineto{\pgfqpoint{5.169122in}{0.322907in}}%
\pgfpathlineto{\pgfqpoint{5.192596in}{0.322907in}}%
\pgfpathlineto{\pgfqpoint{5.216071in}{0.322907in}}%
\pgfpathlineto{\pgfqpoint{5.239545in}{0.322907in}}%
\pgfpathlineto{\pgfqpoint{5.263019in}{0.322907in}}%
\pgfpathlineto{\pgfqpoint{5.286494in}{0.322907in}}%
\pgfpathlineto{\pgfqpoint{5.309968in}{0.322907in}}%
\pgfpathlineto{\pgfqpoint{5.333442in}{0.322907in}}%
\pgfpathlineto{\pgfqpoint{5.356917in}{0.322907in}}%
\pgfpathlineto{\pgfqpoint{5.380391in}{0.322907in}}%
\pgfpathlineto{\pgfqpoint{5.403866in}{0.322907in}}%
\pgfpathlineto{\pgfqpoint{5.427340in}{0.322907in}}%
\pgfpathlineto{\pgfqpoint{5.450814in}{0.322907in}}%
\pgfpathlineto{\pgfqpoint{5.474289in}{0.322907in}}%
\pgfusepath{stroke}%
\end{pgfscope}%
\begin{pgfscope}%
\pgfpathmoveto{\pgfqpoint{4.594000in}{0.100000in}}%
\pgfpathcurveto{\pgfqpoint{4.883602in}{0.100000in}}{\pgfqpoint{5.161381in}{0.157530in}}{\pgfqpoint{5.366161in}{0.259920in}}%
\pgfpathcurveto{\pgfqpoint{5.570940in}{0.362309in}}{\pgfqpoint{5.686000in}{0.501199in}}{\pgfqpoint{5.686000in}{0.646000in}}%
\pgfpathcurveto{\pgfqpoint{5.686000in}{0.790801in}}{\pgfqpoint{5.570940in}{0.929691in}}{\pgfqpoint{5.366161in}{1.032080in}}%
\pgfpathcurveto{\pgfqpoint{5.161381in}{1.134470in}}{\pgfqpoint{4.883602in}{1.192000in}}{\pgfqpoint{4.594000in}{1.192000in}}%
\pgfpathcurveto{\pgfqpoint{4.304398in}{1.192000in}}{\pgfqpoint{4.026619in}{1.134470in}}{\pgfqpoint{3.821839in}{1.032080in}}%
\pgfpathcurveto{\pgfqpoint{3.617060in}{0.929691in}}{\pgfqpoint{3.502000in}{0.790801in}}{\pgfqpoint{3.502000in}{0.646000in}}%
\pgfpathcurveto{\pgfqpoint{3.502000in}{0.501199in}}{\pgfqpoint{3.617060in}{0.362309in}}{\pgfqpoint{3.821839in}{0.259920in}}%
\pgfpathcurveto{\pgfqpoint{4.026619in}{0.157530in}}{\pgfqpoint{4.304398in}{0.100000in}}{\pgfqpoint{4.594000in}{0.100000in}}%
\pgfpathlineto{\pgfqpoint{4.594000in}{0.100000in}}%
\pgfpathclose%
\pgfusepath{clip}%
\pgfsetrectcap%
\pgfsetroundjoin%
\pgfsetlinewidth{0.451687pt}%
\definecolor{currentstroke}{rgb}{0.690196,0.690196,0.690196}%
\pgfsetstrokecolor{currentstroke}%
\pgfsetstrokeopacity{0.250000}%
\pgfsetdash{}{0pt}%
\pgfpathmoveto{\pgfqpoint{3.595062in}{0.425440in}}%
\pgfpathlineto{\pgfqpoint{3.621701in}{0.425440in}}%
\pgfpathlineto{\pgfqpoint{3.648339in}{0.425440in}}%
\pgfpathlineto{\pgfqpoint{3.674977in}{0.425440in}}%
\pgfpathlineto{\pgfqpoint{3.701616in}{0.425440in}}%
\pgfpathlineto{\pgfqpoint{3.728254in}{0.425440in}}%
\pgfpathlineto{\pgfqpoint{3.754892in}{0.425440in}}%
\pgfpathlineto{\pgfqpoint{3.781531in}{0.425440in}}%
\pgfpathlineto{\pgfqpoint{3.808169in}{0.425440in}}%
\pgfpathlineto{\pgfqpoint{3.834807in}{0.425440in}}%
\pgfpathlineto{\pgfqpoint{3.861446in}{0.425440in}}%
\pgfpathlineto{\pgfqpoint{3.888084in}{0.425440in}}%
\pgfpathlineto{\pgfqpoint{3.914722in}{0.425440in}}%
\pgfpathlineto{\pgfqpoint{3.941361in}{0.425440in}}%
\pgfpathlineto{\pgfqpoint{3.967999in}{0.425440in}}%
\pgfpathlineto{\pgfqpoint{3.994637in}{0.425440in}}%
\pgfpathlineto{\pgfqpoint{4.021276in}{0.425440in}}%
\pgfpathlineto{\pgfqpoint{4.047914in}{0.425440in}}%
\pgfpathlineto{\pgfqpoint{4.074552in}{0.425440in}}%
\pgfpathlineto{\pgfqpoint{4.101191in}{0.425440in}}%
\pgfpathlineto{\pgfqpoint{4.127829in}{0.425440in}}%
\pgfpathlineto{\pgfqpoint{4.154467in}{0.425440in}}%
\pgfpathlineto{\pgfqpoint{4.181106in}{0.425440in}}%
\pgfpathlineto{\pgfqpoint{4.207744in}{0.425440in}}%
\pgfpathlineto{\pgfqpoint{4.234382in}{0.425440in}}%
\pgfpathlineto{\pgfqpoint{4.261021in}{0.425440in}}%
\pgfpathlineto{\pgfqpoint{4.287659in}{0.425440in}}%
\pgfpathlineto{\pgfqpoint{4.314297in}{0.425440in}}%
\pgfpathlineto{\pgfqpoint{4.340936in}{0.425440in}}%
\pgfpathlineto{\pgfqpoint{4.367574in}{0.425440in}}%
\pgfpathlineto{\pgfqpoint{4.394212in}{0.425440in}}%
\pgfpathlineto{\pgfqpoint{4.420851in}{0.425440in}}%
\pgfpathlineto{\pgfqpoint{4.447489in}{0.425440in}}%
\pgfpathlineto{\pgfqpoint{4.474127in}{0.425440in}}%
\pgfpathlineto{\pgfqpoint{4.500766in}{0.425440in}}%
\pgfpathlineto{\pgfqpoint{4.527404in}{0.425440in}}%
\pgfpathlineto{\pgfqpoint{4.554042in}{0.425440in}}%
\pgfpathlineto{\pgfqpoint{4.580681in}{0.425440in}}%
\pgfpathlineto{\pgfqpoint{4.607319in}{0.425440in}}%
\pgfpathlineto{\pgfqpoint{4.633958in}{0.425440in}}%
\pgfpathlineto{\pgfqpoint{4.660596in}{0.425440in}}%
\pgfpathlineto{\pgfqpoint{4.687234in}{0.425440in}}%
\pgfpathlineto{\pgfqpoint{4.713873in}{0.425440in}}%
\pgfpathlineto{\pgfqpoint{4.740511in}{0.425440in}}%
\pgfpathlineto{\pgfqpoint{4.767149in}{0.425440in}}%
\pgfpathlineto{\pgfqpoint{4.793788in}{0.425440in}}%
\pgfpathlineto{\pgfqpoint{4.820426in}{0.425440in}}%
\pgfpathlineto{\pgfqpoint{4.847064in}{0.425440in}}%
\pgfpathlineto{\pgfqpoint{4.873703in}{0.425440in}}%
\pgfpathlineto{\pgfqpoint{4.900341in}{0.425440in}}%
\pgfpathlineto{\pgfqpoint{4.926979in}{0.425440in}}%
\pgfpathlineto{\pgfqpoint{4.953618in}{0.425440in}}%
\pgfpathlineto{\pgfqpoint{4.980256in}{0.425440in}}%
\pgfpathlineto{\pgfqpoint{5.006894in}{0.425440in}}%
\pgfpathlineto{\pgfqpoint{5.033533in}{0.425440in}}%
\pgfpathlineto{\pgfqpoint{5.060171in}{0.425440in}}%
\pgfpathlineto{\pgfqpoint{5.086809in}{0.425440in}}%
\pgfpathlineto{\pgfqpoint{5.113448in}{0.425440in}}%
\pgfpathlineto{\pgfqpoint{5.140086in}{0.425440in}}%
\pgfpathlineto{\pgfqpoint{5.166724in}{0.425440in}}%
\pgfpathlineto{\pgfqpoint{5.193363in}{0.425440in}}%
\pgfpathlineto{\pgfqpoint{5.220001in}{0.425440in}}%
\pgfpathlineto{\pgfqpoint{5.246639in}{0.425440in}}%
\pgfpathlineto{\pgfqpoint{5.273278in}{0.425440in}}%
\pgfpathlineto{\pgfqpoint{5.299916in}{0.425440in}}%
\pgfpathlineto{\pgfqpoint{5.326554in}{0.425440in}}%
\pgfpathlineto{\pgfqpoint{5.353193in}{0.425440in}}%
\pgfpathlineto{\pgfqpoint{5.379831in}{0.425440in}}%
\pgfpathlineto{\pgfqpoint{5.406469in}{0.425440in}}%
\pgfpathlineto{\pgfqpoint{5.433108in}{0.425440in}}%
\pgfpathlineto{\pgfqpoint{5.459746in}{0.425440in}}%
\pgfpathlineto{\pgfqpoint{5.486384in}{0.425440in}}%
\pgfpathlineto{\pgfqpoint{5.513023in}{0.425440in}}%
\pgfpathlineto{\pgfqpoint{5.539661in}{0.425440in}}%
\pgfpathlineto{\pgfqpoint{5.566299in}{0.425440in}}%
\pgfpathlineto{\pgfqpoint{5.592938in}{0.425440in}}%
\pgfusepath{stroke}%
\end{pgfscope}%
\begin{pgfscope}%
\pgfpathmoveto{\pgfqpoint{4.594000in}{0.100000in}}%
\pgfpathcurveto{\pgfqpoint{4.883602in}{0.100000in}}{\pgfqpoint{5.161381in}{0.157530in}}{\pgfqpoint{5.366161in}{0.259920in}}%
\pgfpathcurveto{\pgfqpoint{5.570940in}{0.362309in}}{\pgfqpoint{5.686000in}{0.501199in}}{\pgfqpoint{5.686000in}{0.646000in}}%
\pgfpathcurveto{\pgfqpoint{5.686000in}{0.790801in}}{\pgfqpoint{5.570940in}{0.929691in}}{\pgfqpoint{5.366161in}{1.032080in}}%
\pgfpathcurveto{\pgfqpoint{5.161381in}{1.134470in}}{\pgfqpoint{4.883602in}{1.192000in}}{\pgfqpoint{4.594000in}{1.192000in}}%
\pgfpathcurveto{\pgfqpoint{4.304398in}{1.192000in}}{\pgfqpoint{4.026619in}{1.134470in}}{\pgfqpoint{3.821839in}{1.032080in}}%
\pgfpathcurveto{\pgfqpoint{3.617060in}{0.929691in}}{\pgfqpoint{3.502000in}{0.790801in}}{\pgfqpoint{3.502000in}{0.646000in}}%
\pgfpathcurveto{\pgfqpoint{3.502000in}{0.501199in}}{\pgfqpoint{3.617060in}{0.362309in}}{\pgfqpoint{3.821839in}{0.259920in}}%
\pgfpathcurveto{\pgfqpoint{4.026619in}{0.157530in}}{\pgfqpoint{4.304398in}{0.100000in}}{\pgfqpoint{4.594000in}{0.100000in}}%
\pgfpathlineto{\pgfqpoint{4.594000in}{0.100000in}}%
\pgfpathclose%
\pgfusepath{clip}%
\pgfsetrectcap%
\pgfsetroundjoin%
\pgfsetlinewidth{0.451687pt}%
\definecolor{currentstroke}{rgb}{0.690196,0.690196,0.690196}%
\pgfsetstrokecolor{currentstroke}%
\pgfsetstrokeopacity{0.250000}%
\pgfsetdash{}{0pt}%
\pgfpathmoveto{\pgfqpoint{3.525127in}{0.534226in}}%
\pgfpathlineto{\pgfqpoint{3.553630in}{0.534226in}}%
\pgfpathlineto{\pgfqpoint{3.582133in}{0.534226in}}%
\pgfpathlineto{\pgfqpoint{3.610637in}{0.534226in}}%
\pgfpathlineto{\pgfqpoint{3.639140in}{0.534226in}}%
\pgfpathlineto{\pgfqpoint{3.667643in}{0.534226in}}%
\pgfpathlineto{\pgfqpoint{3.696146in}{0.534226in}}%
\pgfpathlineto{\pgfqpoint{3.724650in}{0.534226in}}%
\pgfpathlineto{\pgfqpoint{3.753153in}{0.534226in}}%
\pgfpathlineto{\pgfqpoint{3.781656in}{0.534226in}}%
\pgfpathlineto{\pgfqpoint{3.810160in}{0.534226in}}%
\pgfpathlineto{\pgfqpoint{3.838663in}{0.534226in}}%
\pgfpathlineto{\pgfqpoint{3.867166in}{0.534226in}}%
\pgfpathlineto{\pgfqpoint{3.895669in}{0.534226in}}%
\pgfpathlineto{\pgfqpoint{3.924173in}{0.534226in}}%
\pgfpathlineto{\pgfqpoint{3.952676in}{0.534226in}}%
\pgfpathlineto{\pgfqpoint{3.981179in}{0.534226in}}%
\pgfpathlineto{\pgfqpoint{4.009683in}{0.534226in}}%
\pgfpathlineto{\pgfqpoint{4.038186in}{0.534226in}}%
\pgfpathlineto{\pgfqpoint{4.066689in}{0.534226in}}%
\pgfpathlineto{\pgfqpoint{4.095192in}{0.534226in}}%
\pgfpathlineto{\pgfqpoint{4.123696in}{0.534226in}}%
\pgfpathlineto{\pgfqpoint{4.152199in}{0.534226in}}%
\pgfpathlineto{\pgfqpoint{4.180702in}{0.534226in}}%
\pgfpathlineto{\pgfqpoint{4.209206in}{0.534226in}}%
\pgfpathlineto{\pgfqpoint{4.237709in}{0.534226in}}%
\pgfpathlineto{\pgfqpoint{4.266212in}{0.534226in}}%
\pgfpathlineto{\pgfqpoint{4.294715in}{0.534226in}}%
\pgfpathlineto{\pgfqpoint{4.323219in}{0.534226in}}%
\pgfpathlineto{\pgfqpoint{4.351722in}{0.534226in}}%
\pgfpathlineto{\pgfqpoint{4.380225in}{0.534226in}}%
\pgfpathlineto{\pgfqpoint{4.408729in}{0.534226in}}%
\pgfpathlineto{\pgfqpoint{4.437232in}{0.534226in}}%
\pgfpathlineto{\pgfqpoint{4.465735in}{0.534226in}}%
\pgfpathlineto{\pgfqpoint{4.494238in}{0.534226in}}%
\pgfpathlineto{\pgfqpoint{4.522742in}{0.534226in}}%
\pgfpathlineto{\pgfqpoint{4.551245in}{0.534226in}}%
\pgfpathlineto{\pgfqpoint{4.579748in}{0.534226in}}%
\pgfpathlineto{\pgfqpoint{4.608252in}{0.534226in}}%
\pgfpathlineto{\pgfqpoint{4.636755in}{0.534226in}}%
\pgfpathlineto{\pgfqpoint{4.665258in}{0.534226in}}%
\pgfpathlineto{\pgfqpoint{4.693762in}{0.534226in}}%
\pgfpathlineto{\pgfqpoint{4.722265in}{0.534226in}}%
\pgfpathlineto{\pgfqpoint{4.750768in}{0.534226in}}%
\pgfpathlineto{\pgfqpoint{4.779271in}{0.534226in}}%
\pgfpathlineto{\pgfqpoint{4.807775in}{0.534226in}}%
\pgfpathlineto{\pgfqpoint{4.836278in}{0.534226in}}%
\pgfpathlineto{\pgfqpoint{4.864781in}{0.534226in}}%
\pgfpathlineto{\pgfqpoint{4.893285in}{0.534226in}}%
\pgfpathlineto{\pgfqpoint{4.921788in}{0.534226in}}%
\pgfpathlineto{\pgfqpoint{4.950291in}{0.534226in}}%
\pgfpathlineto{\pgfqpoint{4.978794in}{0.534226in}}%
\pgfpathlineto{\pgfqpoint{5.007298in}{0.534226in}}%
\pgfpathlineto{\pgfqpoint{5.035801in}{0.534226in}}%
\pgfpathlineto{\pgfqpoint{5.064304in}{0.534226in}}%
\pgfpathlineto{\pgfqpoint{5.092808in}{0.534226in}}%
\pgfpathlineto{\pgfqpoint{5.121311in}{0.534226in}}%
\pgfpathlineto{\pgfqpoint{5.149814in}{0.534226in}}%
\pgfpathlineto{\pgfqpoint{5.178317in}{0.534226in}}%
\pgfpathlineto{\pgfqpoint{5.206821in}{0.534226in}}%
\pgfpathlineto{\pgfqpoint{5.235324in}{0.534226in}}%
\pgfpathlineto{\pgfqpoint{5.263827in}{0.534226in}}%
\pgfpathlineto{\pgfqpoint{5.292331in}{0.534226in}}%
\pgfpathlineto{\pgfqpoint{5.320834in}{0.534226in}}%
\pgfpathlineto{\pgfqpoint{5.349337in}{0.534226in}}%
\pgfpathlineto{\pgfqpoint{5.377840in}{0.534226in}}%
\pgfpathlineto{\pgfqpoint{5.406344in}{0.534226in}}%
\pgfpathlineto{\pgfqpoint{5.434847in}{0.534226in}}%
\pgfpathlineto{\pgfqpoint{5.463350in}{0.534226in}}%
\pgfpathlineto{\pgfqpoint{5.491854in}{0.534226in}}%
\pgfpathlineto{\pgfqpoint{5.520357in}{0.534226in}}%
\pgfpathlineto{\pgfqpoint{5.548860in}{0.534226in}}%
\pgfpathlineto{\pgfqpoint{5.577363in}{0.534226in}}%
\pgfpathlineto{\pgfqpoint{5.605867in}{0.534226in}}%
\pgfpathlineto{\pgfqpoint{5.634370in}{0.534226in}}%
\pgfpathlineto{\pgfqpoint{5.662873in}{0.534226in}}%
\pgfusepath{stroke}%
\end{pgfscope}%
\begin{pgfscope}%
\pgfpathmoveto{\pgfqpoint{4.594000in}{0.100000in}}%
\pgfpathcurveto{\pgfqpoint{4.883602in}{0.100000in}}{\pgfqpoint{5.161381in}{0.157530in}}{\pgfqpoint{5.366161in}{0.259920in}}%
\pgfpathcurveto{\pgfqpoint{5.570940in}{0.362309in}}{\pgfqpoint{5.686000in}{0.501199in}}{\pgfqpoint{5.686000in}{0.646000in}}%
\pgfpathcurveto{\pgfqpoint{5.686000in}{0.790801in}}{\pgfqpoint{5.570940in}{0.929691in}}{\pgfqpoint{5.366161in}{1.032080in}}%
\pgfpathcurveto{\pgfqpoint{5.161381in}{1.134470in}}{\pgfqpoint{4.883602in}{1.192000in}}{\pgfqpoint{4.594000in}{1.192000in}}%
\pgfpathcurveto{\pgfqpoint{4.304398in}{1.192000in}}{\pgfqpoint{4.026619in}{1.134470in}}{\pgfqpoint{3.821839in}{1.032080in}}%
\pgfpathcurveto{\pgfqpoint{3.617060in}{0.929691in}}{\pgfqpoint{3.502000in}{0.790801in}}{\pgfqpoint{3.502000in}{0.646000in}}%
\pgfpathcurveto{\pgfqpoint{3.502000in}{0.501199in}}{\pgfqpoint{3.617060in}{0.362309in}}{\pgfqpoint{3.821839in}{0.259920in}}%
\pgfpathcurveto{\pgfqpoint{4.026619in}{0.157530in}}{\pgfqpoint{4.304398in}{0.100000in}}{\pgfqpoint{4.594000in}{0.100000in}}%
\pgfpathlineto{\pgfqpoint{4.594000in}{0.100000in}}%
\pgfpathclose%
\pgfusepath{clip}%
\pgfsetrectcap%
\pgfsetroundjoin%
\pgfsetlinewidth{0.451687pt}%
\definecolor{currentstroke}{rgb}{0.690196,0.690196,0.690196}%
\pgfsetstrokecolor{currentstroke}%
\pgfsetstrokeopacity{0.250000}%
\pgfsetdash{}{0pt}%
\pgfpathmoveto{\pgfqpoint{3.502000in}{0.646000in}}%
\pgfpathlineto{\pgfqpoint{3.531120in}{0.646000in}}%
\pgfpathlineto{\pgfqpoint{3.560240in}{0.646000in}}%
\pgfpathlineto{\pgfqpoint{3.589360in}{0.646000in}}%
\pgfpathlineto{\pgfqpoint{3.618480in}{0.646000in}}%
\pgfpathlineto{\pgfqpoint{3.647600in}{0.646000in}}%
\pgfpathlineto{\pgfqpoint{3.676720in}{0.646000in}}%
\pgfpathlineto{\pgfqpoint{3.705840in}{0.646000in}}%
\pgfpathlineto{\pgfqpoint{3.734960in}{0.646000in}}%
\pgfpathlineto{\pgfqpoint{3.764080in}{0.646000in}}%
\pgfpathlineto{\pgfqpoint{3.793200in}{0.646000in}}%
\pgfpathlineto{\pgfqpoint{3.822320in}{0.646000in}}%
\pgfpathlineto{\pgfqpoint{3.851440in}{0.646000in}}%
\pgfpathlineto{\pgfqpoint{3.880560in}{0.646000in}}%
\pgfpathlineto{\pgfqpoint{3.909680in}{0.646000in}}%
\pgfpathlineto{\pgfqpoint{3.938800in}{0.646000in}}%
\pgfpathlineto{\pgfqpoint{3.967920in}{0.646000in}}%
\pgfpathlineto{\pgfqpoint{3.997040in}{0.646000in}}%
\pgfpathlineto{\pgfqpoint{4.026160in}{0.646000in}}%
\pgfpathlineto{\pgfqpoint{4.055280in}{0.646000in}}%
\pgfpathlineto{\pgfqpoint{4.084400in}{0.646000in}}%
\pgfpathlineto{\pgfqpoint{4.113520in}{0.646000in}}%
\pgfpathlineto{\pgfqpoint{4.142640in}{0.646000in}}%
\pgfpathlineto{\pgfqpoint{4.171760in}{0.646000in}}%
\pgfpathlineto{\pgfqpoint{4.200880in}{0.646000in}}%
\pgfpathlineto{\pgfqpoint{4.230000in}{0.646000in}}%
\pgfpathlineto{\pgfqpoint{4.259120in}{0.646000in}}%
\pgfpathlineto{\pgfqpoint{4.288240in}{0.646000in}}%
\pgfpathlineto{\pgfqpoint{4.317360in}{0.646000in}}%
\pgfpathlineto{\pgfqpoint{4.346480in}{0.646000in}}%
\pgfpathlineto{\pgfqpoint{4.375600in}{0.646000in}}%
\pgfpathlineto{\pgfqpoint{4.404720in}{0.646000in}}%
\pgfpathlineto{\pgfqpoint{4.433840in}{0.646000in}}%
\pgfpathlineto{\pgfqpoint{4.462960in}{0.646000in}}%
\pgfpathlineto{\pgfqpoint{4.492080in}{0.646000in}}%
\pgfpathlineto{\pgfqpoint{4.521200in}{0.646000in}}%
\pgfpathlineto{\pgfqpoint{4.550320in}{0.646000in}}%
\pgfpathlineto{\pgfqpoint{4.579440in}{0.646000in}}%
\pgfpathlineto{\pgfqpoint{4.608560in}{0.646000in}}%
\pgfpathlineto{\pgfqpoint{4.637680in}{0.646000in}}%
\pgfpathlineto{\pgfqpoint{4.666800in}{0.646000in}}%
\pgfpathlineto{\pgfqpoint{4.695920in}{0.646000in}}%
\pgfpathlineto{\pgfqpoint{4.725040in}{0.646000in}}%
\pgfpathlineto{\pgfqpoint{4.754160in}{0.646000in}}%
\pgfpathlineto{\pgfqpoint{4.783280in}{0.646000in}}%
\pgfpathlineto{\pgfqpoint{4.812400in}{0.646000in}}%
\pgfpathlineto{\pgfqpoint{4.841520in}{0.646000in}}%
\pgfpathlineto{\pgfqpoint{4.870640in}{0.646000in}}%
\pgfpathlineto{\pgfqpoint{4.899760in}{0.646000in}}%
\pgfpathlineto{\pgfqpoint{4.928880in}{0.646000in}}%
\pgfpathlineto{\pgfqpoint{4.958000in}{0.646000in}}%
\pgfpathlineto{\pgfqpoint{4.987120in}{0.646000in}}%
\pgfpathlineto{\pgfqpoint{5.016240in}{0.646000in}}%
\pgfpathlineto{\pgfqpoint{5.045360in}{0.646000in}}%
\pgfpathlineto{\pgfqpoint{5.074480in}{0.646000in}}%
\pgfpathlineto{\pgfqpoint{5.103600in}{0.646000in}}%
\pgfpathlineto{\pgfqpoint{5.132720in}{0.646000in}}%
\pgfpathlineto{\pgfqpoint{5.161840in}{0.646000in}}%
\pgfpathlineto{\pgfqpoint{5.190960in}{0.646000in}}%
\pgfpathlineto{\pgfqpoint{5.220080in}{0.646000in}}%
\pgfpathlineto{\pgfqpoint{5.249200in}{0.646000in}}%
\pgfpathlineto{\pgfqpoint{5.278320in}{0.646000in}}%
\pgfpathlineto{\pgfqpoint{5.307440in}{0.646000in}}%
\pgfpathlineto{\pgfqpoint{5.336560in}{0.646000in}}%
\pgfpathlineto{\pgfqpoint{5.365680in}{0.646000in}}%
\pgfpathlineto{\pgfqpoint{5.394800in}{0.646000in}}%
\pgfpathlineto{\pgfqpoint{5.423920in}{0.646000in}}%
\pgfpathlineto{\pgfqpoint{5.453040in}{0.646000in}}%
\pgfpathlineto{\pgfqpoint{5.482160in}{0.646000in}}%
\pgfpathlineto{\pgfqpoint{5.511280in}{0.646000in}}%
\pgfpathlineto{\pgfqpoint{5.540400in}{0.646000in}}%
\pgfpathlineto{\pgfqpoint{5.569520in}{0.646000in}}%
\pgfpathlineto{\pgfqpoint{5.598640in}{0.646000in}}%
\pgfpathlineto{\pgfqpoint{5.627760in}{0.646000in}}%
\pgfpathlineto{\pgfqpoint{5.656880in}{0.646000in}}%
\pgfpathlineto{\pgfqpoint{5.686000in}{0.646000in}}%
\pgfusepath{stroke}%
\end{pgfscope}%
\begin{pgfscope}%
\pgfpathmoveto{\pgfqpoint{4.594000in}{0.100000in}}%
\pgfpathcurveto{\pgfqpoint{4.883602in}{0.100000in}}{\pgfqpoint{5.161381in}{0.157530in}}{\pgfqpoint{5.366161in}{0.259920in}}%
\pgfpathcurveto{\pgfqpoint{5.570940in}{0.362309in}}{\pgfqpoint{5.686000in}{0.501199in}}{\pgfqpoint{5.686000in}{0.646000in}}%
\pgfpathcurveto{\pgfqpoint{5.686000in}{0.790801in}}{\pgfqpoint{5.570940in}{0.929691in}}{\pgfqpoint{5.366161in}{1.032080in}}%
\pgfpathcurveto{\pgfqpoint{5.161381in}{1.134470in}}{\pgfqpoint{4.883602in}{1.192000in}}{\pgfqpoint{4.594000in}{1.192000in}}%
\pgfpathcurveto{\pgfqpoint{4.304398in}{1.192000in}}{\pgfqpoint{4.026619in}{1.134470in}}{\pgfqpoint{3.821839in}{1.032080in}}%
\pgfpathcurveto{\pgfqpoint{3.617060in}{0.929691in}}{\pgfqpoint{3.502000in}{0.790801in}}{\pgfqpoint{3.502000in}{0.646000in}}%
\pgfpathcurveto{\pgfqpoint{3.502000in}{0.501199in}}{\pgfqpoint{3.617060in}{0.362309in}}{\pgfqpoint{3.821839in}{0.259920in}}%
\pgfpathcurveto{\pgfqpoint{4.026619in}{0.157530in}}{\pgfqpoint{4.304398in}{0.100000in}}{\pgfqpoint{4.594000in}{0.100000in}}%
\pgfpathlineto{\pgfqpoint{4.594000in}{0.100000in}}%
\pgfpathclose%
\pgfusepath{clip}%
\pgfsetrectcap%
\pgfsetroundjoin%
\pgfsetlinewidth{0.451687pt}%
\definecolor{currentstroke}{rgb}{0.690196,0.690196,0.690196}%
\pgfsetstrokecolor{currentstroke}%
\pgfsetstrokeopacity{0.250000}%
\pgfsetdash{}{0pt}%
\pgfpathmoveto{\pgfqpoint{3.525127in}{0.757774in}}%
\pgfpathlineto{\pgfqpoint{3.553630in}{0.757774in}}%
\pgfpathlineto{\pgfqpoint{3.582133in}{0.757774in}}%
\pgfpathlineto{\pgfqpoint{3.610637in}{0.757774in}}%
\pgfpathlineto{\pgfqpoint{3.639140in}{0.757774in}}%
\pgfpathlineto{\pgfqpoint{3.667643in}{0.757774in}}%
\pgfpathlineto{\pgfqpoint{3.696146in}{0.757774in}}%
\pgfpathlineto{\pgfqpoint{3.724650in}{0.757774in}}%
\pgfpathlineto{\pgfqpoint{3.753153in}{0.757774in}}%
\pgfpathlineto{\pgfqpoint{3.781656in}{0.757774in}}%
\pgfpathlineto{\pgfqpoint{3.810160in}{0.757774in}}%
\pgfpathlineto{\pgfqpoint{3.838663in}{0.757774in}}%
\pgfpathlineto{\pgfqpoint{3.867166in}{0.757774in}}%
\pgfpathlineto{\pgfqpoint{3.895669in}{0.757774in}}%
\pgfpathlineto{\pgfqpoint{3.924173in}{0.757774in}}%
\pgfpathlineto{\pgfqpoint{3.952676in}{0.757774in}}%
\pgfpathlineto{\pgfqpoint{3.981179in}{0.757774in}}%
\pgfpathlineto{\pgfqpoint{4.009683in}{0.757774in}}%
\pgfpathlineto{\pgfqpoint{4.038186in}{0.757774in}}%
\pgfpathlineto{\pgfqpoint{4.066689in}{0.757774in}}%
\pgfpathlineto{\pgfqpoint{4.095192in}{0.757774in}}%
\pgfpathlineto{\pgfqpoint{4.123696in}{0.757774in}}%
\pgfpathlineto{\pgfqpoint{4.152199in}{0.757774in}}%
\pgfpathlineto{\pgfqpoint{4.180702in}{0.757774in}}%
\pgfpathlineto{\pgfqpoint{4.209206in}{0.757774in}}%
\pgfpathlineto{\pgfqpoint{4.237709in}{0.757774in}}%
\pgfpathlineto{\pgfqpoint{4.266212in}{0.757774in}}%
\pgfpathlineto{\pgfqpoint{4.294715in}{0.757774in}}%
\pgfpathlineto{\pgfqpoint{4.323219in}{0.757774in}}%
\pgfpathlineto{\pgfqpoint{4.351722in}{0.757774in}}%
\pgfpathlineto{\pgfqpoint{4.380225in}{0.757774in}}%
\pgfpathlineto{\pgfqpoint{4.408729in}{0.757774in}}%
\pgfpathlineto{\pgfqpoint{4.437232in}{0.757774in}}%
\pgfpathlineto{\pgfqpoint{4.465735in}{0.757774in}}%
\pgfpathlineto{\pgfqpoint{4.494238in}{0.757774in}}%
\pgfpathlineto{\pgfqpoint{4.522742in}{0.757774in}}%
\pgfpathlineto{\pgfqpoint{4.551245in}{0.757774in}}%
\pgfpathlineto{\pgfqpoint{4.579748in}{0.757774in}}%
\pgfpathlineto{\pgfqpoint{4.608252in}{0.757774in}}%
\pgfpathlineto{\pgfqpoint{4.636755in}{0.757774in}}%
\pgfpathlineto{\pgfqpoint{4.665258in}{0.757774in}}%
\pgfpathlineto{\pgfqpoint{4.693762in}{0.757774in}}%
\pgfpathlineto{\pgfqpoint{4.722265in}{0.757774in}}%
\pgfpathlineto{\pgfqpoint{4.750768in}{0.757774in}}%
\pgfpathlineto{\pgfqpoint{4.779271in}{0.757774in}}%
\pgfpathlineto{\pgfqpoint{4.807775in}{0.757774in}}%
\pgfpathlineto{\pgfqpoint{4.836278in}{0.757774in}}%
\pgfpathlineto{\pgfqpoint{4.864781in}{0.757774in}}%
\pgfpathlineto{\pgfqpoint{4.893285in}{0.757774in}}%
\pgfpathlineto{\pgfqpoint{4.921788in}{0.757774in}}%
\pgfpathlineto{\pgfqpoint{4.950291in}{0.757774in}}%
\pgfpathlineto{\pgfqpoint{4.978794in}{0.757774in}}%
\pgfpathlineto{\pgfqpoint{5.007298in}{0.757774in}}%
\pgfpathlineto{\pgfqpoint{5.035801in}{0.757774in}}%
\pgfpathlineto{\pgfqpoint{5.064304in}{0.757774in}}%
\pgfpathlineto{\pgfqpoint{5.092808in}{0.757774in}}%
\pgfpathlineto{\pgfqpoint{5.121311in}{0.757774in}}%
\pgfpathlineto{\pgfqpoint{5.149814in}{0.757774in}}%
\pgfpathlineto{\pgfqpoint{5.178317in}{0.757774in}}%
\pgfpathlineto{\pgfqpoint{5.206821in}{0.757774in}}%
\pgfpathlineto{\pgfqpoint{5.235324in}{0.757774in}}%
\pgfpathlineto{\pgfqpoint{5.263827in}{0.757774in}}%
\pgfpathlineto{\pgfqpoint{5.292331in}{0.757774in}}%
\pgfpathlineto{\pgfqpoint{5.320834in}{0.757774in}}%
\pgfpathlineto{\pgfqpoint{5.349337in}{0.757774in}}%
\pgfpathlineto{\pgfqpoint{5.377840in}{0.757774in}}%
\pgfpathlineto{\pgfqpoint{5.406344in}{0.757774in}}%
\pgfpathlineto{\pgfqpoint{5.434847in}{0.757774in}}%
\pgfpathlineto{\pgfqpoint{5.463350in}{0.757774in}}%
\pgfpathlineto{\pgfqpoint{5.491854in}{0.757774in}}%
\pgfpathlineto{\pgfqpoint{5.520357in}{0.757774in}}%
\pgfpathlineto{\pgfqpoint{5.548860in}{0.757774in}}%
\pgfpathlineto{\pgfqpoint{5.577363in}{0.757774in}}%
\pgfpathlineto{\pgfqpoint{5.605867in}{0.757774in}}%
\pgfpathlineto{\pgfqpoint{5.634370in}{0.757774in}}%
\pgfpathlineto{\pgfqpoint{5.662873in}{0.757774in}}%
\pgfusepath{stroke}%
\end{pgfscope}%
\begin{pgfscope}%
\pgfpathmoveto{\pgfqpoint{4.594000in}{0.100000in}}%
\pgfpathcurveto{\pgfqpoint{4.883602in}{0.100000in}}{\pgfqpoint{5.161381in}{0.157530in}}{\pgfqpoint{5.366161in}{0.259920in}}%
\pgfpathcurveto{\pgfqpoint{5.570940in}{0.362309in}}{\pgfqpoint{5.686000in}{0.501199in}}{\pgfqpoint{5.686000in}{0.646000in}}%
\pgfpathcurveto{\pgfqpoint{5.686000in}{0.790801in}}{\pgfqpoint{5.570940in}{0.929691in}}{\pgfqpoint{5.366161in}{1.032080in}}%
\pgfpathcurveto{\pgfqpoint{5.161381in}{1.134470in}}{\pgfqpoint{4.883602in}{1.192000in}}{\pgfqpoint{4.594000in}{1.192000in}}%
\pgfpathcurveto{\pgfqpoint{4.304398in}{1.192000in}}{\pgfqpoint{4.026619in}{1.134470in}}{\pgfqpoint{3.821839in}{1.032080in}}%
\pgfpathcurveto{\pgfqpoint{3.617060in}{0.929691in}}{\pgfqpoint{3.502000in}{0.790801in}}{\pgfqpoint{3.502000in}{0.646000in}}%
\pgfpathcurveto{\pgfqpoint{3.502000in}{0.501199in}}{\pgfqpoint{3.617060in}{0.362309in}}{\pgfqpoint{3.821839in}{0.259920in}}%
\pgfpathcurveto{\pgfqpoint{4.026619in}{0.157530in}}{\pgfqpoint{4.304398in}{0.100000in}}{\pgfqpoint{4.594000in}{0.100000in}}%
\pgfpathlineto{\pgfqpoint{4.594000in}{0.100000in}}%
\pgfpathclose%
\pgfusepath{clip}%
\pgfsetrectcap%
\pgfsetroundjoin%
\pgfsetlinewidth{0.451687pt}%
\definecolor{currentstroke}{rgb}{0.690196,0.690196,0.690196}%
\pgfsetstrokecolor{currentstroke}%
\pgfsetstrokeopacity{0.250000}%
\pgfsetdash{}{0pt}%
\pgfpathmoveto{\pgfqpoint{3.595062in}{0.866560in}}%
\pgfpathlineto{\pgfqpoint{3.621701in}{0.866560in}}%
\pgfpathlineto{\pgfqpoint{3.648339in}{0.866560in}}%
\pgfpathlineto{\pgfqpoint{3.674977in}{0.866560in}}%
\pgfpathlineto{\pgfqpoint{3.701616in}{0.866560in}}%
\pgfpathlineto{\pgfqpoint{3.728254in}{0.866560in}}%
\pgfpathlineto{\pgfqpoint{3.754892in}{0.866560in}}%
\pgfpathlineto{\pgfqpoint{3.781531in}{0.866560in}}%
\pgfpathlineto{\pgfqpoint{3.808169in}{0.866560in}}%
\pgfpathlineto{\pgfqpoint{3.834807in}{0.866560in}}%
\pgfpathlineto{\pgfqpoint{3.861446in}{0.866560in}}%
\pgfpathlineto{\pgfqpoint{3.888084in}{0.866560in}}%
\pgfpathlineto{\pgfqpoint{3.914722in}{0.866560in}}%
\pgfpathlineto{\pgfqpoint{3.941361in}{0.866560in}}%
\pgfpathlineto{\pgfqpoint{3.967999in}{0.866560in}}%
\pgfpathlineto{\pgfqpoint{3.994637in}{0.866560in}}%
\pgfpathlineto{\pgfqpoint{4.021276in}{0.866560in}}%
\pgfpathlineto{\pgfqpoint{4.047914in}{0.866560in}}%
\pgfpathlineto{\pgfqpoint{4.074552in}{0.866560in}}%
\pgfpathlineto{\pgfqpoint{4.101191in}{0.866560in}}%
\pgfpathlineto{\pgfqpoint{4.127829in}{0.866560in}}%
\pgfpathlineto{\pgfqpoint{4.154467in}{0.866560in}}%
\pgfpathlineto{\pgfqpoint{4.181106in}{0.866560in}}%
\pgfpathlineto{\pgfqpoint{4.207744in}{0.866560in}}%
\pgfpathlineto{\pgfqpoint{4.234382in}{0.866560in}}%
\pgfpathlineto{\pgfqpoint{4.261021in}{0.866560in}}%
\pgfpathlineto{\pgfqpoint{4.287659in}{0.866560in}}%
\pgfpathlineto{\pgfqpoint{4.314297in}{0.866560in}}%
\pgfpathlineto{\pgfqpoint{4.340936in}{0.866560in}}%
\pgfpathlineto{\pgfqpoint{4.367574in}{0.866560in}}%
\pgfpathlineto{\pgfqpoint{4.394212in}{0.866560in}}%
\pgfpathlineto{\pgfqpoint{4.420851in}{0.866560in}}%
\pgfpathlineto{\pgfqpoint{4.447489in}{0.866560in}}%
\pgfpathlineto{\pgfqpoint{4.474127in}{0.866560in}}%
\pgfpathlineto{\pgfqpoint{4.500766in}{0.866560in}}%
\pgfpathlineto{\pgfqpoint{4.527404in}{0.866560in}}%
\pgfpathlineto{\pgfqpoint{4.554042in}{0.866560in}}%
\pgfpathlineto{\pgfqpoint{4.580681in}{0.866560in}}%
\pgfpathlineto{\pgfqpoint{4.607319in}{0.866560in}}%
\pgfpathlineto{\pgfqpoint{4.633958in}{0.866560in}}%
\pgfpathlineto{\pgfqpoint{4.660596in}{0.866560in}}%
\pgfpathlineto{\pgfqpoint{4.687234in}{0.866560in}}%
\pgfpathlineto{\pgfqpoint{4.713873in}{0.866560in}}%
\pgfpathlineto{\pgfqpoint{4.740511in}{0.866560in}}%
\pgfpathlineto{\pgfqpoint{4.767149in}{0.866560in}}%
\pgfpathlineto{\pgfqpoint{4.793788in}{0.866560in}}%
\pgfpathlineto{\pgfqpoint{4.820426in}{0.866560in}}%
\pgfpathlineto{\pgfqpoint{4.847064in}{0.866560in}}%
\pgfpathlineto{\pgfqpoint{4.873703in}{0.866560in}}%
\pgfpathlineto{\pgfqpoint{4.900341in}{0.866560in}}%
\pgfpathlineto{\pgfqpoint{4.926979in}{0.866560in}}%
\pgfpathlineto{\pgfqpoint{4.953618in}{0.866560in}}%
\pgfpathlineto{\pgfqpoint{4.980256in}{0.866560in}}%
\pgfpathlineto{\pgfqpoint{5.006894in}{0.866560in}}%
\pgfpathlineto{\pgfqpoint{5.033533in}{0.866560in}}%
\pgfpathlineto{\pgfqpoint{5.060171in}{0.866560in}}%
\pgfpathlineto{\pgfqpoint{5.086809in}{0.866560in}}%
\pgfpathlineto{\pgfqpoint{5.113448in}{0.866560in}}%
\pgfpathlineto{\pgfqpoint{5.140086in}{0.866560in}}%
\pgfpathlineto{\pgfqpoint{5.166724in}{0.866560in}}%
\pgfpathlineto{\pgfqpoint{5.193363in}{0.866560in}}%
\pgfpathlineto{\pgfqpoint{5.220001in}{0.866560in}}%
\pgfpathlineto{\pgfqpoint{5.246639in}{0.866560in}}%
\pgfpathlineto{\pgfqpoint{5.273278in}{0.866560in}}%
\pgfpathlineto{\pgfqpoint{5.299916in}{0.866560in}}%
\pgfpathlineto{\pgfqpoint{5.326554in}{0.866560in}}%
\pgfpathlineto{\pgfqpoint{5.353193in}{0.866560in}}%
\pgfpathlineto{\pgfqpoint{5.379831in}{0.866560in}}%
\pgfpathlineto{\pgfqpoint{5.406469in}{0.866560in}}%
\pgfpathlineto{\pgfqpoint{5.433108in}{0.866560in}}%
\pgfpathlineto{\pgfqpoint{5.459746in}{0.866560in}}%
\pgfpathlineto{\pgfqpoint{5.486384in}{0.866560in}}%
\pgfpathlineto{\pgfqpoint{5.513023in}{0.866560in}}%
\pgfpathlineto{\pgfqpoint{5.539661in}{0.866560in}}%
\pgfpathlineto{\pgfqpoint{5.566299in}{0.866560in}}%
\pgfpathlineto{\pgfqpoint{5.592938in}{0.866560in}}%
\pgfusepath{stroke}%
\end{pgfscope}%
\begin{pgfscope}%
\pgfpathmoveto{\pgfqpoint{4.594000in}{0.100000in}}%
\pgfpathcurveto{\pgfqpoint{4.883602in}{0.100000in}}{\pgfqpoint{5.161381in}{0.157530in}}{\pgfqpoint{5.366161in}{0.259920in}}%
\pgfpathcurveto{\pgfqpoint{5.570940in}{0.362309in}}{\pgfqpoint{5.686000in}{0.501199in}}{\pgfqpoint{5.686000in}{0.646000in}}%
\pgfpathcurveto{\pgfqpoint{5.686000in}{0.790801in}}{\pgfqpoint{5.570940in}{0.929691in}}{\pgfqpoint{5.366161in}{1.032080in}}%
\pgfpathcurveto{\pgfqpoint{5.161381in}{1.134470in}}{\pgfqpoint{4.883602in}{1.192000in}}{\pgfqpoint{4.594000in}{1.192000in}}%
\pgfpathcurveto{\pgfqpoint{4.304398in}{1.192000in}}{\pgfqpoint{4.026619in}{1.134470in}}{\pgfqpoint{3.821839in}{1.032080in}}%
\pgfpathcurveto{\pgfqpoint{3.617060in}{0.929691in}}{\pgfqpoint{3.502000in}{0.790801in}}{\pgfqpoint{3.502000in}{0.646000in}}%
\pgfpathcurveto{\pgfqpoint{3.502000in}{0.501199in}}{\pgfqpoint{3.617060in}{0.362309in}}{\pgfqpoint{3.821839in}{0.259920in}}%
\pgfpathcurveto{\pgfqpoint{4.026619in}{0.157530in}}{\pgfqpoint{4.304398in}{0.100000in}}{\pgfqpoint{4.594000in}{0.100000in}}%
\pgfpathlineto{\pgfqpoint{4.594000in}{0.100000in}}%
\pgfpathclose%
\pgfusepath{clip}%
\pgfsetrectcap%
\pgfsetroundjoin%
\pgfsetlinewidth{0.451687pt}%
\definecolor{currentstroke}{rgb}{0.690196,0.690196,0.690196}%
\pgfsetstrokecolor{currentstroke}%
\pgfsetstrokeopacity{0.250000}%
\pgfsetdash{}{0pt}%
\pgfpathmoveto{\pgfqpoint{3.713711in}{0.969093in}}%
\pgfpathlineto{\pgfqpoint{3.737186in}{0.969093in}}%
\pgfpathlineto{\pgfqpoint{3.760660in}{0.969093in}}%
\pgfpathlineto{\pgfqpoint{3.784134in}{0.969093in}}%
\pgfpathlineto{\pgfqpoint{3.807609in}{0.969093in}}%
\pgfpathlineto{\pgfqpoint{3.831083in}{0.969093in}}%
\pgfpathlineto{\pgfqpoint{3.854558in}{0.969093in}}%
\pgfpathlineto{\pgfqpoint{3.878032in}{0.969093in}}%
\pgfpathlineto{\pgfqpoint{3.901506in}{0.969093in}}%
\pgfpathlineto{\pgfqpoint{3.924981in}{0.969093in}}%
\pgfpathlineto{\pgfqpoint{3.948455in}{0.969093in}}%
\pgfpathlineto{\pgfqpoint{3.971929in}{0.969093in}}%
\pgfpathlineto{\pgfqpoint{3.995404in}{0.969093in}}%
\pgfpathlineto{\pgfqpoint{4.018878in}{0.969093in}}%
\pgfpathlineto{\pgfqpoint{4.042352in}{0.969093in}}%
\pgfpathlineto{\pgfqpoint{4.065827in}{0.969093in}}%
\pgfpathlineto{\pgfqpoint{4.089301in}{0.969093in}}%
\pgfpathlineto{\pgfqpoint{4.112776in}{0.969093in}}%
\pgfpathlineto{\pgfqpoint{4.136250in}{0.969093in}}%
\pgfpathlineto{\pgfqpoint{4.159724in}{0.969093in}}%
\pgfpathlineto{\pgfqpoint{4.183199in}{0.969093in}}%
\pgfpathlineto{\pgfqpoint{4.206673in}{0.969093in}}%
\pgfpathlineto{\pgfqpoint{4.230147in}{0.969093in}}%
\pgfpathlineto{\pgfqpoint{4.253622in}{0.969093in}}%
\pgfpathlineto{\pgfqpoint{4.277096in}{0.969093in}}%
\pgfpathlineto{\pgfqpoint{4.300570in}{0.969093in}}%
\pgfpathlineto{\pgfqpoint{4.324045in}{0.969093in}}%
\pgfpathlineto{\pgfqpoint{4.347519in}{0.969093in}}%
\pgfpathlineto{\pgfqpoint{4.370994in}{0.969093in}}%
\pgfpathlineto{\pgfqpoint{4.394468in}{0.969093in}}%
\pgfpathlineto{\pgfqpoint{4.417942in}{0.969093in}}%
\pgfpathlineto{\pgfqpoint{4.441417in}{0.969093in}}%
\pgfpathlineto{\pgfqpoint{4.464891in}{0.969093in}}%
\pgfpathlineto{\pgfqpoint{4.488365in}{0.969093in}}%
\pgfpathlineto{\pgfqpoint{4.511840in}{0.969093in}}%
\pgfpathlineto{\pgfqpoint{4.535314in}{0.969093in}}%
\pgfpathlineto{\pgfqpoint{4.558788in}{0.969093in}}%
\pgfpathlineto{\pgfqpoint{4.582263in}{0.969093in}}%
\pgfpathlineto{\pgfqpoint{4.605737in}{0.969093in}}%
\pgfpathlineto{\pgfqpoint{4.629212in}{0.969093in}}%
\pgfpathlineto{\pgfqpoint{4.652686in}{0.969093in}}%
\pgfpathlineto{\pgfqpoint{4.676160in}{0.969093in}}%
\pgfpathlineto{\pgfqpoint{4.699635in}{0.969093in}}%
\pgfpathlineto{\pgfqpoint{4.723109in}{0.969093in}}%
\pgfpathlineto{\pgfqpoint{4.746583in}{0.969093in}}%
\pgfpathlineto{\pgfqpoint{4.770058in}{0.969093in}}%
\pgfpathlineto{\pgfqpoint{4.793532in}{0.969093in}}%
\pgfpathlineto{\pgfqpoint{4.817006in}{0.969093in}}%
\pgfpathlineto{\pgfqpoint{4.840481in}{0.969093in}}%
\pgfpathlineto{\pgfqpoint{4.863955in}{0.969093in}}%
\pgfpathlineto{\pgfqpoint{4.887430in}{0.969093in}}%
\pgfpathlineto{\pgfqpoint{4.910904in}{0.969093in}}%
\pgfpathlineto{\pgfqpoint{4.934378in}{0.969093in}}%
\pgfpathlineto{\pgfqpoint{4.957853in}{0.969093in}}%
\pgfpathlineto{\pgfqpoint{4.981327in}{0.969093in}}%
\pgfpathlineto{\pgfqpoint{5.004801in}{0.969093in}}%
\pgfpathlineto{\pgfqpoint{5.028276in}{0.969093in}}%
\pgfpathlineto{\pgfqpoint{5.051750in}{0.969093in}}%
\pgfpathlineto{\pgfqpoint{5.075224in}{0.969093in}}%
\pgfpathlineto{\pgfqpoint{5.098699in}{0.969093in}}%
\pgfpathlineto{\pgfqpoint{5.122173in}{0.969093in}}%
\pgfpathlineto{\pgfqpoint{5.145648in}{0.969093in}}%
\pgfpathlineto{\pgfqpoint{5.169122in}{0.969093in}}%
\pgfpathlineto{\pgfqpoint{5.192596in}{0.969093in}}%
\pgfpathlineto{\pgfqpoint{5.216071in}{0.969093in}}%
\pgfpathlineto{\pgfqpoint{5.239545in}{0.969093in}}%
\pgfpathlineto{\pgfqpoint{5.263019in}{0.969093in}}%
\pgfpathlineto{\pgfqpoint{5.286494in}{0.969093in}}%
\pgfpathlineto{\pgfqpoint{5.309968in}{0.969093in}}%
\pgfpathlineto{\pgfqpoint{5.333442in}{0.969093in}}%
\pgfpathlineto{\pgfqpoint{5.356917in}{0.969093in}}%
\pgfpathlineto{\pgfqpoint{5.380391in}{0.969093in}}%
\pgfpathlineto{\pgfqpoint{5.403866in}{0.969093in}}%
\pgfpathlineto{\pgfqpoint{5.427340in}{0.969093in}}%
\pgfpathlineto{\pgfqpoint{5.450814in}{0.969093in}}%
\pgfpathlineto{\pgfqpoint{5.474289in}{0.969093in}}%
\pgfusepath{stroke}%
\end{pgfscope}%
\begin{pgfscope}%
\pgfpathmoveto{\pgfqpoint{4.594000in}{0.100000in}}%
\pgfpathcurveto{\pgfqpoint{4.883602in}{0.100000in}}{\pgfqpoint{5.161381in}{0.157530in}}{\pgfqpoint{5.366161in}{0.259920in}}%
\pgfpathcurveto{\pgfqpoint{5.570940in}{0.362309in}}{\pgfqpoint{5.686000in}{0.501199in}}{\pgfqpoint{5.686000in}{0.646000in}}%
\pgfpathcurveto{\pgfqpoint{5.686000in}{0.790801in}}{\pgfqpoint{5.570940in}{0.929691in}}{\pgfqpoint{5.366161in}{1.032080in}}%
\pgfpathcurveto{\pgfqpoint{5.161381in}{1.134470in}}{\pgfqpoint{4.883602in}{1.192000in}}{\pgfqpoint{4.594000in}{1.192000in}}%
\pgfpathcurveto{\pgfqpoint{4.304398in}{1.192000in}}{\pgfqpoint{4.026619in}{1.134470in}}{\pgfqpoint{3.821839in}{1.032080in}}%
\pgfpathcurveto{\pgfqpoint{3.617060in}{0.929691in}}{\pgfqpoint{3.502000in}{0.790801in}}{\pgfqpoint{3.502000in}{0.646000in}}%
\pgfpathcurveto{\pgfqpoint{3.502000in}{0.501199in}}{\pgfqpoint{3.617060in}{0.362309in}}{\pgfqpoint{3.821839in}{0.259920in}}%
\pgfpathcurveto{\pgfqpoint{4.026619in}{0.157530in}}{\pgfqpoint{4.304398in}{0.100000in}}{\pgfqpoint{4.594000in}{0.100000in}}%
\pgfpathlineto{\pgfqpoint{4.594000in}{0.100000in}}%
\pgfpathclose%
\pgfusepath{clip}%
\pgfsetrectcap%
\pgfsetroundjoin%
\pgfsetlinewidth{0.451687pt}%
\definecolor{currentstroke}{rgb}{0.690196,0.690196,0.690196}%
\pgfsetstrokecolor{currentstroke}%
\pgfsetstrokeopacity{0.250000}%
\pgfsetdash{}{0pt}%
\pgfpathmoveto{\pgfqpoint{3.887329in}{1.062257in}}%
\pgfpathlineto{\pgfqpoint{3.906174in}{1.062257in}}%
\pgfpathlineto{\pgfqpoint{3.925018in}{1.062257in}}%
\pgfpathlineto{\pgfqpoint{3.943863in}{1.062257in}}%
\pgfpathlineto{\pgfqpoint{3.962707in}{1.062257in}}%
\pgfpathlineto{\pgfqpoint{3.981552in}{1.062257in}}%
\pgfpathlineto{\pgfqpoint{4.000396in}{1.062257in}}%
\pgfpathlineto{\pgfqpoint{4.019241in}{1.062257in}}%
\pgfpathlineto{\pgfqpoint{4.038086in}{1.062257in}}%
\pgfpathlineto{\pgfqpoint{4.056930in}{1.062257in}}%
\pgfpathlineto{\pgfqpoint{4.075775in}{1.062257in}}%
\pgfpathlineto{\pgfqpoint{4.094619in}{1.062257in}}%
\pgfpathlineto{\pgfqpoint{4.113464in}{1.062257in}}%
\pgfpathlineto{\pgfqpoint{4.132308in}{1.062257in}}%
\pgfpathlineto{\pgfqpoint{4.151153in}{1.062257in}}%
\pgfpathlineto{\pgfqpoint{4.169997in}{1.062257in}}%
\pgfpathlineto{\pgfqpoint{4.188842in}{1.062257in}}%
\pgfpathlineto{\pgfqpoint{4.207687in}{1.062257in}}%
\pgfpathlineto{\pgfqpoint{4.226531in}{1.062257in}}%
\pgfpathlineto{\pgfqpoint{4.245376in}{1.062257in}}%
\pgfpathlineto{\pgfqpoint{4.264220in}{1.062257in}}%
\pgfpathlineto{\pgfqpoint{4.283065in}{1.062257in}}%
\pgfpathlineto{\pgfqpoint{4.301909in}{1.062257in}}%
\pgfpathlineto{\pgfqpoint{4.320754in}{1.062257in}}%
\pgfpathlineto{\pgfqpoint{4.339598in}{1.062257in}}%
\pgfpathlineto{\pgfqpoint{4.358443in}{1.062257in}}%
\pgfpathlineto{\pgfqpoint{4.377288in}{1.062257in}}%
\pgfpathlineto{\pgfqpoint{4.396132in}{1.062257in}}%
\pgfpathlineto{\pgfqpoint{4.414977in}{1.062257in}}%
\pgfpathlineto{\pgfqpoint{4.433821in}{1.062257in}}%
\pgfpathlineto{\pgfqpoint{4.452666in}{1.062257in}}%
\pgfpathlineto{\pgfqpoint{4.471510in}{1.062257in}}%
\pgfpathlineto{\pgfqpoint{4.490355in}{1.062257in}}%
\pgfpathlineto{\pgfqpoint{4.509199in}{1.062257in}}%
\pgfpathlineto{\pgfqpoint{4.528044in}{1.062257in}}%
\pgfpathlineto{\pgfqpoint{4.546889in}{1.062257in}}%
\pgfpathlineto{\pgfqpoint{4.565733in}{1.062257in}}%
\pgfpathlineto{\pgfqpoint{4.584578in}{1.062257in}}%
\pgfpathlineto{\pgfqpoint{4.603422in}{1.062257in}}%
\pgfpathlineto{\pgfqpoint{4.622267in}{1.062257in}}%
\pgfpathlineto{\pgfqpoint{4.641111in}{1.062257in}}%
\pgfpathlineto{\pgfqpoint{4.659956in}{1.062257in}}%
\pgfpathlineto{\pgfqpoint{4.678801in}{1.062257in}}%
\pgfpathlineto{\pgfqpoint{4.697645in}{1.062257in}}%
\pgfpathlineto{\pgfqpoint{4.716490in}{1.062257in}}%
\pgfpathlineto{\pgfqpoint{4.735334in}{1.062257in}}%
\pgfpathlineto{\pgfqpoint{4.754179in}{1.062257in}}%
\pgfpathlineto{\pgfqpoint{4.773023in}{1.062257in}}%
\pgfpathlineto{\pgfqpoint{4.791868in}{1.062257in}}%
\pgfpathlineto{\pgfqpoint{4.810712in}{1.062257in}}%
\pgfpathlineto{\pgfqpoint{4.829557in}{1.062257in}}%
\pgfpathlineto{\pgfqpoint{4.848402in}{1.062257in}}%
\pgfpathlineto{\pgfqpoint{4.867246in}{1.062257in}}%
\pgfpathlineto{\pgfqpoint{4.886091in}{1.062257in}}%
\pgfpathlineto{\pgfqpoint{4.904935in}{1.062257in}}%
\pgfpathlineto{\pgfqpoint{4.923780in}{1.062257in}}%
\pgfpathlineto{\pgfqpoint{4.942624in}{1.062257in}}%
\pgfpathlineto{\pgfqpoint{4.961469in}{1.062257in}}%
\pgfpathlineto{\pgfqpoint{4.980313in}{1.062257in}}%
\pgfpathlineto{\pgfqpoint{4.999158in}{1.062257in}}%
\pgfpathlineto{\pgfqpoint{5.018003in}{1.062257in}}%
\pgfpathlineto{\pgfqpoint{5.036847in}{1.062257in}}%
\pgfpathlineto{\pgfqpoint{5.055692in}{1.062257in}}%
\pgfpathlineto{\pgfqpoint{5.074536in}{1.062257in}}%
\pgfpathlineto{\pgfqpoint{5.093381in}{1.062257in}}%
\pgfpathlineto{\pgfqpoint{5.112225in}{1.062257in}}%
\pgfpathlineto{\pgfqpoint{5.131070in}{1.062257in}}%
\pgfpathlineto{\pgfqpoint{5.149914in}{1.062257in}}%
\pgfpathlineto{\pgfqpoint{5.168759in}{1.062257in}}%
\pgfpathlineto{\pgfqpoint{5.187604in}{1.062257in}}%
\pgfpathlineto{\pgfqpoint{5.206448in}{1.062257in}}%
\pgfpathlineto{\pgfqpoint{5.225293in}{1.062257in}}%
\pgfpathlineto{\pgfqpoint{5.244137in}{1.062257in}}%
\pgfpathlineto{\pgfqpoint{5.262982in}{1.062257in}}%
\pgfpathlineto{\pgfqpoint{5.281826in}{1.062257in}}%
\pgfpathlineto{\pgfqpoint{5.300671in}{1.062257in}}%
\pgfusepath{stroke}%
\end{pgfscope}%
\begin{pgfscope}%
\pgfpathmoveto{\pgfqpoint{4.594000in}{0.100000in}}%
\pgfpathcurveto{\pgfqpoint{4.883602in}{0.100000in}}{\pgfqpoint{5.161381in}{0.157530in}}{\pgfqpoint{5.366161in}{0.259920in}}%
\pgfpathcurveto{\pgfqpoint{5.570940in}{0.362309in}}{\pgfqpoint{5.686000in}{0.501199in}}{\pgfqpoint{5.686000in}{0.646000in}}%
\pgfpathcurveto{\pgfqpoint{5.686000in}{0.790801in}}{\pgfqpoint{5.570940in}{0.929691in}}{\pgfqpoint{5.366161in}{1.032080in}}%
\pgfpathcurveto{\pgfqpoint{5.161381in}{1.134470in}}{\pgfqpoint{4.883602in}{1.192000in}}{\pgfqpoint{4.594000in}{1.192000in}}%
\pgfpathcurveto{\pgfqpoint{4.304398in}{1.192000in}}{\pgfqpoint{4.026619in}{1.134470in}}{\pgfqpoint{3.821839in}{1.032080in}}%
\pgfpathcurveto{\pgfqpoint{3.617060in}{0.929691in}}{\pgfqpoint{3.502000in}{0.790801in}}{\pgfqpoint{3.502000in}{0.646000in}}%
\pgfpathcurveto{\pgfqpoint{3.502000in}{0.501199in}}{\pgfqpoint{3.617060in}{0.362309in}}{\pgfqpoint{3.821839in}{0.259920in}}%
\pgfpathcurveto{\pgfqpoint{4.026619in}{0.157530in}}{\pgfqpoint{4.304398in}{0.100000in}}{\pgfqpoint{4.594000in}{0.100000in}}%
\pgfpathlineto{\pgfqpoint{4.594000in}{0.100000in}}%
\pgfpathclose%
\pgfusepath{clip}%
\pgfsetrectcap%
\pgfsetroundjoin%
\pgfsetlinewidth{0.451687pt}%
\definecolor{currentstroke}{rgb}{0.690196,0.690196,0.690196}%
\pgfsetstrokecolor{currentstroke}%
\pgfsetstrokeopacity{0.250000}%
\pgfsetdash{}{0pt}%
\pgfpathmoveto{\pgfqpoint{4.131906in}{1.140705in}}%
\pgfpathlineto{\pgfqpoint{4.144228in}{1.140705in}}%
\pgfpathlineto{\pgfqpoint{4.156551in}{1.140705in}}%
\pgfpathlineto{\pgfqpoint{4.168873in}{1.140705in}}%
\pgfpathlineto{\pgfqpoint{4.181196in}{1.140705in}}%
\pgfpathlineto{\pgfqpoint{4.193518in}{1.140705in}}%
\pgfpathlineto{\pgfqpoint{4.205841in}{1.140705in}}%
\pgfpathlineto{\pgfqpoint{4.218163in}{1.140705in}}%
\pgfpathlineto{\pgfqpoint{4.230486in}{1.140705in}}%
\pgfpathlineto{\pgfqpoint{4.242808in}{1.140705in}}%
\pgfpathlineto{\pgfqpoint{4.255131in}{1.140705in}}%
\pgfpathlineto{\pgfqpoint{4.267453in}{1.140705in}}%
\pgfpathlineto{\pgfqpoint{4.279776in}{1.140705in}}%
\pgfpathlineto{\pgfqpoint{4.292098in}{1.140705in}}%
\pgfpathlineto{\pgfqpoint{4.304421in}{1.140705in}}%
\pgfpathlineto{\pgfqpoint{4.316743in}{1.140705in}}%
\pgfpathlineto{\pgfqpoint{4.329066in}{1.140705in}}%
\pgfpathlineto{\pgfqpoint{4.341388in}{1.140705in}}%
\pgfpathlineto{\pgfqpoint{4.353711in}{1.140705in}}%
\pgfpathlineto{\pgfqpoint{4.366033in}{1.140705in}}%
\pgfpathlineto{\pgfqpoint{4.378356in}{1.140705in}}%
\pgfpathlineto{\pgfqpoint{4.390678in}{1.140705in}}%
\pgfpathlineto{\pgfqpoint{4.403001in}{1.140705in}}%
\pgfpathlineto{\pgfqpoint{4.415323in}{1.140705in}}%
\pgfpathlineto{\pgfqpoint{4.427646in}{1.140705in}}%
\pgfpathlineto{\pgfqpoint{4.439969in}{1.140705in}}%
\pgfpathlineto{\pgfqpoint{4.452291in}{1.140705in}}%
\pgfpathlineto{\pgfqpoint{4.464614in}{1.140705in}}%
\pgfpathlineto{\pgfqpoint{4.476936in}{1.140705in}}%
\pgfpathlineto{\pgfqpoint{4.489259in}{1.140705in}}%
\pgfpathlineto{\pgfqpoint{4.501581in}{1.140705in}}%
\pgfpathlineto{\pgfqpoint{4.513904in}{1.140705in}}%
\pgfpathlineto{\pgfqpoint{4.526226in}{1.140705in}}%
\pgfpathlineto{\pgfqpoint{4.538549in}{1.140705in}}%
\pgfpathlineto{\pgfqpoint{4.550871in}{1.140705in}}%
\pgfpathlineto{\pgfqpoint{4.563194in}{1.140705in}}%
\pgfpathlineto{\pgfqpoint{4.575516in}{1.140705in}}%
\pgfpathlineto{\pgfqpoint{4.587839in}{1.140705in}}%
\pgfpathlineto{\pgfqpoint{4.600161in}{1.140705in}}%
\pgfpathlineto{\pgfqpoint{4.612484in}{1.140705in}}%
\pgfpathlineto{\pgfqpoint{4.624806in}{1.140705in}}%
\pgfpathlineto{\pgfqpoint{4.637129in}{1.140705in}}%
\pgfpathlineto{\pgfqpoint{4.649451in}{1.140705in}}%
\pgfpathlineto{\pgfqpoint{4.661774in}{1.140705in}}%
\pgfpathlineto{\pgfqpoint{4.674096in}{1.140705in}}%
\pgfpathlineto{\pgfqpoint{4.686419in}{1.140705in}}%
\pgfpathlineto{\pgfqpoint{4.698741in}{1.140705in}}%
\pgfpathlineto{\pgfqpoint{4.711064in}{1.140705in}}%
\pgfpathlineto{\pgfqpoint{4.723386in}{1.140705in}}%
\pgfpathlineto{\pgfqpoint{4.735709in}{1.140705in}}%
\pgfpathlineto{\pgfqpoint{4.748031in}{1.140705in}}%
\pgfpathlineto{\pgfqpoint{4.760354in}{1.140705in}}%
\pgfpathlineto{\pgfqpoint{4.772677in}{1.140705in}}%
\pgfpathlineto{\pgfqpoint{4.784999in}{1.140705in}}%
\pgfpathlineto{\pgfqpoint{4.797322in}{1.140705in}}%
\pgfpathlineto{\pgfqpoint{4.809644in}{1.140705in}}%
\pgfpathlineto{\pgfqpoint{4.821967in}{1.140705in}}%
\pgfpathlineto{\pgfqpoint{4.834289in}{1.140705in}}%
\pgfpathlineto{\pgfqpoint{4.846612in}{1.140705in}}%
\pgfpathlineto{\pgfqpoint{4.858934in}{1.140705in}}%
\pgfpathlineto{\pgfqpoint{4.871257in}{1.140705in}}%
\pgfpathlineto{\pgfqpoint{4.883579in}{1.140705in}}%
\pgfpathlineto{\pgfqpoint{4.895902in}{1.140705in}}%
\pgfpathlineto{\pgfqpoint{4.908224in}{1.140705in}}%
\pgfpathlineto{\pgfqpoint{4.920547in}{1.140705in}}%
\pgfpathlineto{\pgfqpoint{4.932869in}{1.140705in}}%
\pgfpathlineto{\pgfqpoint{4.945192in}{1.140705in}}%
\pgfpathlineto{\pgfqpoint{4.957514in}{1.140705in}}%
\pgfpathlineto{\pgfqpoint{4.969837in}{1.140705in}}%
\pgfpathlineto{\pgfqpoint{4.982159in}{1.140705in}}%
\pgfpathlineto{\pgfqpoint{4.994482in}{1.140705in}}%
\pgfpathlineto{\pgfqpoint{5.006804in}{1.140705in}}%
\pgfpathlineto{\pgfqpoint{5.019127in}{1.140705in}}%
\pgfpathlineto{\pgfqpoint{5.031449in}{1.140705in}}%
\pgfpathlineto{\pgfqpoint{5.043772in}{1.140705in}}%
\pgfpathlineto{\pgfqpoint{5.056094in}{1.140705in}}%
\pgfusepath{stroke}%
\end{pgfscope}%
\begin{pgfscope}%
\pgfsetrectcap%
\pgfsetmiterjoin%
\pgfsetlinewidth{0.803000pt}%
\definecolor{currentstroke}{rgb}{0.000000,0.000000,0.000000}%
\pgfsetstrokecolor{currentstroke}%
\pgfsetdash{}{0pt}%
\pgfpathmoveto{\pgfqpoint{4.594000in}{0.100000in}}%
\pgfpathcurveto{\pgfqpoint{4.883602in}{0.100000in}}{\pgfqpoint{5.161381in}{0.157530in}}{\pgfqpoint{5.366161in}{0.259920in}}%
\pgfpathcurveto{\pgfqpoint{5.570940in}{0.362309in}}{\pgfqpoint{5.686000in}{0.501199in}}{\pgfqpoint{5.686000in}{0.646000in}}%
\pgfpathcurveto{\pgfqpoint{5.686000in}{0.790801in}}{\pgfqpoint{5.570940in}{0.929691in}}{\pgfqpoint{5.366161in}{1.032080in}}%
\pgfpathcurveto{\pgfqpoint{5.161381in}{1.134470in}}{\pgfqpoint{4.883602in}{1.192000in}}{\pgfqpoint{4.594000in}{1.192000in}}%
\pgfpathcurveto{\pgfqpoint{4.304398in}{1.192000in}}{\pgfqpoint{4.026619in}{1.134470in}}{\pgfqpoint{3.821839in}{1.032080in}}%
\pgfpathcurveto{\pgfqpoint{3.617060in}{0.929691in}}{\pgfqpoint{3.502000in}{0.790801in}}{\pgfqpoint{3.502000in}{0.646000in}}%
\pgfpathcurveto{\pgfqpoint{3.502000in}{0.501199in}}{\pgfqpoint{3.617060in}{0.362309in}}{\pgfqpoint{3.821839in}{0.259920in}}%
\pgfpathcurveto{\pgfqpoint{4.026619in}{0.157530in}}{\pgfqpoint{4.304398in}{0.100000in}}{\pgfqpoint{4.594000in}{0.100000in}}%
\pgfpathlineto{\pgfqpoint{4.594000in}{0.100000in}}%
\pgfpathclose%
\pgfusepath{stroke}%
\end{pgfscope}%
\begin{pgfscope}%
\definecolor{textcolor}{rgb}{0.000000,0.000000,0.000000}%
\pgfsetstrokecolor{textcolor}%
\pgfsetfillcolor{textcolor}%
\pgftext[x=4.594000in,y=1.233667in,,base]{\color{textcolor}{\rmfamily\fontsize{8.200000}{9.840000}\selectfont\catcode`\^=\active\def^{\ifmmode\sp\else\^{}\fi}\catcode`\%=\active\def%{\%}Duplicates}}%
\end{pgfscope}%
\end{pgfpicture}%
\makeatother%
\endgroup%

\caption{Mollweide projections of the five synthetic $S^2$ configurations. Antipodal, girdle, and four-mode structures can have small resultant length. However, they remain highly nonuniform at relevant scales.}
\label{fig:clouds}
\end{figure*}

The synthetic values use the constructions in Appendix~\ref{app:add_experiments}. We compute each quantity with the routine in Section~\ref{sec:gess}. Only the heat Gram matrix is sphere-specific. The heat kernel is
\begin{equation*}
  k_t(x,y)=\sum_{\ell=0}^{\infty}(2\ell+1)e^{-\ell(\ell+1)t}P_\ell(x^\top y),
\end{equation*}
where $P_\ell$ is the Legendre polynomial. For each weighted cloud, we first compute
\begin{equation*}
  B_\ell=(2\ell+1)\sum_{i,j}w_iw_jP_\ell(x_i^\top x_j)
\end{equation*}
once and evaluate
\begin{equation*}
  \widehat A_w(t)=\sum_{\ell\ge0}e^{-2\ell(\ell+1)t}B_\ell.
\end{equation*}
As a similarity-sensitive baseline, we use Leinster-Cobbold-style diversity with geodesic Gaussian similarity:
\begin{equation*}
  D_{\rm LC}^{(\alpha)}=\left[\sum_{i,j}w_iw_j\exp\{-(d_g(x_i,x_j)/\alpha)^2\}\right]^{-1}.
\end{equation*}
This baseline matches the $e^{-1}$ convention at $d_g=\alpha$, but it is not derived from the heat semigroup.
We also report ordinary $\ell_2$ ESS. The real-data table includes entropy/perplexity ESS, $\ess_H=\exp\{-\sum_iw_i\log w_i\}$. This weight-only baseline appears in ESS surveys \citep{martino_2017_EffectiveSampleSize}. The heat-kernel MMD curve equals $\widehat A_w(t)-1$, as noted in Section~\ref{sec:related-work}. Figure~\ref{fig:asymptotics} uses this unlogged curve to assess both asymptotic approximations.


\begin{figure*}[t]
\centering
%\includegraphics[width=0.96\textwidth]{hkep_fig_asymptotics.pdf}
\providecommand{\mathdefault}[1]{#1}
\providecommand{\mathregular}[1]{#1}
%% Creator: Matplotlib, PGF backend
%%
%% To include the figure in your LaTeX document, write
%%   \input{<filename>.pgf}
%%
%% Make sure the required packages are loaded in your preamble
%%   \usepackage{pgf}
%%
%% Also ensure that all the required font packages are loaded; for instance,
%% the lmodern package is sometimes necessary when using math font.
%%   \usepackage{lmodern}
%%
%% Figures using additional raster images can only be included by \input if
%% they are in the same directory as the main LaTeX file. For loading figures
%% from other directories you can use the `import` package
%%   \usepackage{import}
%%
%% and then include the figures with
%%   \import{<path to file>}{<filename>.pgf}
%%
%% Matplotlib used the following preamble
%%   \def\mathdefault#1{#1}
%%   \everymath=\expandafter{\the\everymath\displaystyle}
%%   \IfFileExists{scrextend.sty}{
%%     \usepackage[fontsize=8.000000pt]{scrextend}
%%   }{
%%     \renewcommand{\normalsize}{\fontsize{8.000000}{9.600000}\selectfont}
%%     \normalsize
%%   }
%%   \usepackage{amsmath,amssymb}
%%   \makeatletter\@ifpackageloaded{underscore}{}{\usepackage[strings]{underscore}}\makeatother
%%
\begingroup%
\makeatletter%
\begin{pgfpicture}%
\pgfpathrectangle{\pgfpointorigin}{\pgfqpoint{7.023821in}{3.017306in}}%
\pgfusepath{use as bounding box, clip}%
\begin{pgfscope}%
\pgfsetbuttcap%
\pgfsetmiterjoin%
\definecolor{currentfill}{rgb}{1.000000,1.000000,1.000000}%
\pgfsetfillcolor{currentfill}%
\pgfsetlinewidth{0.000000pt}%
\definecolor{currentstroke}{rgb}{1.000000,1.000000,1.000000}%
\pgfsetstrokecolor{currentstroke}%
\pgfsetdash{}{0pt}%
\pgfpathmoveto{\pgfqpoint{0.000000in}{0.000000in}}%
\pgfpathlineto{\pgfqpoint{7.023821in}{0.000000in}}%
\pgfpathlineto{\pgfqpoint{7.023821in}{3.017306in}}%
\pgfpathlineto{\pgfqpoint{0.000000in}{3.017306in}}%
\pgfpathlineto{\pgfqpoint{0.000000in}{0.000000in}}%
\pgfpathclose%
\pgfusepath{fill}%
\end{pgfscope}%
\begin{pgfscope}%
\pgfsetbuttcap%
\pgfsetmiterjoin%
\definecolor{currentfill}{rgb}{1.000000,1.000000,1.000000}%
\pgfsetfillcolor{currentfill}%
\pgfsetlinewidth{0.000000pt}%
\definecolor{currentstroke}{rgb}{0.000000,0.000000,0.000000}%
\pgfsetstrokecolor{currentstroke}%
\pgfsetstrokeopacity{0.000000}%
\pgfsetdash{}{0pt}%
\pgfpathmoveto{\pgfqpoint{0.553821in}{0.829139in}}%
\pgfpathlineto{\pgfqpoint{3.299510in}{0.829139in}}%
\pgfpathlineto{\pgfqpoint{3.299510in}{2.750639in}}%
\pgfpathlineto{\pgfqpoint{0.553821in}{2.750639in}}%
\pgfpathlineto{\pgfqpoint{0.553821in}{0.829139in}}%
\pgfpathclose%
\pgfusepath{fill}%
\end{pgfscope}%
\begin{pgfscope}%
\pgfpathrectangle{\pgfqpoint{0.553821in}{0.829139in}}{\pgfqpoint{2.745690in}{1.921500in}}%
\pgfusepath{clip}%
\pgfsetrectcap%
\pgfsetroundjoin%
\pgfsetlinewidth{0.803000pt}%
\definecolor{currentstroke}{rgb}{0.690196,0.690196,0.690196}%
\pgfsetstrokecolor{currentstroke}%
\pgfsetstrokeopacity{0.300000}%
\pgfsetdash{}{0pt}%
\pgfpathmoveto{\pgfqpoint{0.821002in}{0.829139in}}%
\pgfpathlineto{\pgfqpoint{0.821002in}{2.750639in}}%
\pgfusepath{stroke}%
\end{pgfscope}%
\begin{pgfscope}%
\pgfsetbuttcap%
\pgfsetroundjoin%
\definecolor{currentfill}{rgb}{0.000000,0.000000,0.000000}%
\pgfsetfillcolor{currentfill}%
\pgfsetlinewidth{0.803000pt}%
\definecolor{currentstroke}{rgb}{0.000000,0.000000,0.000000}%
\pgfsetstrokecolor{currentstroke}%
\pgfsetdash{}{0pt}%
\pgfsys@defobject{currentmarker}{\pgfqpoint{0.000000in}{-0.048611in}}{\pgfqpoint{0.000000in}{0.000000in}}{%
\pgfpathmoveto{\pgfqpoint{0.000000in}{0.000000in}}%
\pgfpathlineto{\pgfqpoint{0.000000in}{-0.048611in}}%
\pgfusepath{stroke,fill}%
}%
\begin{pgfscope}%
\pgfsys@transformshift{0.821002in}{0.829139in}%
\pgfsys@useobject{currentmarker}{}%
\end{pgfscope}%
\end{pgfscope}%
\begin{pgfscope}%
\definecolor{textcolor}{rgb}{0.000000,0.000000,0.000000}%
\pgfsetstrokecolor{textcolor}%
\pgfsetfillcolor{textcolor}%
\pgftext[x=0.821002in,y=0.731917in,,top]{\color{textcolor}{\rmfamily\fontsize{7.000000}{8.400000}\selectfont\catcode`\^=\active\def^{\ifmmode\sp\else\^{}\fi}\catcode`\%=\active\def%{\%}$\mathdefault{10^{-3}}$}}%
\end{pgfscope}%
\begin{pgfscope}%
\pgfpathrectangle{\pgfqpoint{0.553821in}{0.829139in}}{\pgfqpoint{2.745690in}{1.921500in}}%
\pgfusepath{clip}%
\pgfsetrectcap%
\pgfsetroundjoin%
\pgfsetlinewidth{0.803000pt}%
\definecolor{currentstroke}{rgb}{0.690196,0.690196,0.690196}%
\pgfsetstrokecolor{currentstroke}%
\pgfsetstrokeopacity{0.300000}%
\pgfsetdash{}{0pt}%
\pgfpathmoveto{\pgfqpoint{2.290176in}{0.829139in}}%
\pgfpathlineto{\pgfqpoint{2.290176in}{2.750639in}}%
\pgfusepath{stroke}%
\end{pgfscope}%
\begin{pgfscope}%
\pgfsetbuttcap%
\pgfsetroundjoin%
\definecolor{currentfill}{rgb}{0.000000,0.000000,0.000000}%
\pgfsetfillcolor{currentfill}%
\pgfsetlinewidth{0.803000pt}%
\definecolor{currentstroke}{rgb}{0.000000,0.000000,0.000000}%
\pgfsetstrokecolor{currentstroke}%
\pgfsetdash{}{0pt}%
\pgfsys@defobject{currentmarker}{\pgfqpoint{0.000000in}{-0.048611in}}{\pgfqpoint{0.000000in}{0.000000in}}{%
\pgfpathmoveto{\pgfqpoint{0.000000in}{0.000000in}}%
\pgfpathlineto{\pgfqpoint{0.000000in}{-0.048611in}}%
\pgfusepath{stroke,fill}%
}%
\begin{pgfscope}%
\pgfsys@transformshift{2.290176in}{0.829139in}%
\pgfsys@useobject{currentmarker}{}%
\end{pgfscope}%
\end{pgfscope}%
\begin{pgfscope}%
\definecolor{textcolor}{rgb}{0.000000,0.000000,0.000000}%
\pgfsetstrokecolor{textcolor}%
\pgfsetfillcolor{textcolor}%
\pgftext[x=2.290176in,y=0.731917in,,top]{\color{textcolor}{\rmfamily\fontsize{7.000000}{8.400000}\selectfont\catcode`\^=\active\def^{\ifmmode\sp\else\^{}\fi}\catcode`\%=\active\def%{\%}$\mathdefault{10^{-2}}$}}%
\end{pgfscope}%
\begin{pgfscope}%
\pgfpathrectangle{\pgfqpoint{0.553821in}{0.829139in}}{\pgfqpoint{2.745690in}{1.921500in}}%
\pgfusepath{clip}%
\pgfsetrectcap%
\pgfsetroundjoin%
\pgfsetlinewidth{0.803000pt}%
\definecolor{currentstroke}{rgb}{0.690196,0.690196,0.690196}%
\pgfsetstrokecolor{currentstroke}%
\pgfsetstrokeopacity{0.300000}%
\pgfsetdash{}{0pt}%
\pgfpathmoveto{\pgfqpoint{0.593425in}{0.829139in}}%
\pgfpathlineto{\pgfqpoint{0.593425in}{2.750639in}}%
\pgfusepath{stroke}%
\end{pgfscope}%
\begin{pgfscope}%
\pgfsetbuttcap%
\pgfsetroundjoin%
\definecolor{currentfill}{rgb}{0.000000,0.000000,0.000000}%
\pgfsetfillcolor{currentfill}%
\pgfsetlinewidth{0.602250pt}%
\definecolor{currentstroke}{rgb}{0.000000,0.000000,0.000000}%
\pgfsetstrokecolor{currentstroke}%
\pgfsetdash{}{0pt}%
\pgfsys@defobject{currentmarker}{\pgfqpoint{0.000000in}{-0.027778in}}{\pgfqpoint{0.000000in}{0.000000in}}{%
\pgfpathmoveto{\pgfqpoint{0.000000in}{0.000000in}}%
\pgfpathlineto{\pgfqpoint{0.000000in}{-0.027778in}}%
\pgfusepath{stroke,fill}%
}%
\begin{pgfscope}%
\pgfsys@transformshift{0.593425in}{0.829139in}%
\pgfsys@useobject{currentmarker}{}%
\end{pgfscope}%
\end{pgfscope}%
\begin{pgfscope}%
\pgfpathrectangle{\pgfqpoint{0.553821in}{0.829139in}}{\pgfqpoint{2.745690in}{1.921500in}}%
\pgfusepath{clip}%
\pgfsetrectcap%
\pgfsetroundjoin%
\pgfsetlinewidth{0.803000pt}%
\definecolor{currentstroke}{rgb}{0.690196,0.690196,0.690196}%
\pgfsetstrokecolor{currentstroke}%
\pgfsetstrokeopacity{0.300000}%
\pgfsetdash{}{0pt}%
\pgfpathmoveto{\pgfqpoint{0.678625in}{0.829139in}}%
\pgfpathlineto{\pgfqpoint{0.678625in}{2.750639in}}%
\pgfusepath{stroke}%
\end{pgfscope}%
\begin{pgfscope}%
\pgfsetbuttcap%
\pgfsetroundjoin%
\definecolor{currentfill}{rgb}{0.000000,0.000000,0.000000}%
\pgfsetfillcolor{currentfill}%
\pgfsetlinewidth{0.602250pt}%
\definecolor{currentstroke}{rgb}{0.000000,0.000000,0.000000}%
\pgfsetstrokecolor{currentstroke}%
\pgfsetdash{}{0pt}%
\pgfsys@defobject{currentmarker}{\pgfqpoint{0.000000in}{-0.027778in}}{\pgfqpoint{0.000000in}{0.000000in}}{%
\pgfpathmoveto{\pgfqpoint{0.000000in}{0.000000in}}%
\pgfpathlineto{\pgfqpoint{0.000000in}{-0.027778in}}%
\pgfusepath{stroke,fill}%
}%
\begin{pgfscope}%
\pgfsys@transformshift{0.678625in}{0.829139in}%
\pgfsys@useobject{currentmarker}{}%
\end{pgfscope}%
\end{pgfscope}%
\begin{pgfscope}%
\pgfpathrectangle{\pgfqpoint{0.553821in}{0.829139in}}{\pgfqpoint{2.745690in}{1.921500in}}%
\pgfusepath{clip}%
\pgfsetrectcap%
\pgfsetroundjoin%
\pgfsetlinewidth{0.803000pt}%
\definecolor{currentstroke}{rgb}{0.690196,0.690196,0.690196}%
\pgfsetstrokecolor{currentstroke}%
\pgfsetstrokeopacity{0.300000}%
\pgfsetdash{}{0pt}%
\pgfpathmoveto{\pgfqpoint{0.753777in}{0.829139in}}%
\pgfpathlineto{\pgfqpoint{0.753777in}{2.750639in}}%
\pgfusepath{stroke}%
\end{pgfscope}%
\begin{pgfscope}%
\pgfsetbuttcap%
\pgfsetroundjoin%
\definecolor{currentfill}{rgb}{0.000000,0.000000,0.000000}%
\pgfsetfillcolor{currentfill}%
\pgfsetlinewidth{0.602250pt}%
\definecolor{currentstroke}{rgb}{0.000000,0.000000,0.000000}%
\pgfsetstrokecolor{currentstroke}%
\pgfsetdash{}{0pt}%
\pgfsys@defobject{currentmarker}{\pgfqpoint{0.000000in}{-0.027778in}}{\pgfqpoint{0.000000in}{0.000000in}}{%
\pgfpathmoveto{\pgfqpoint{0.000000in}{0.000000in}}%
\pgfpathlineto{\pgfqpoint{0.000000in}{-0.027778in}}%
\pgfusepath{stroke,fill}%
}%
\begin{pgfscope}%
\pgfsys@transformshift{0.753777in}{0.829139in}%
\pgfsys@useobject{currentmarker}{}%
\end{pgfscope}%
\end{pgfscope}%
\begin{pgfscope}%
\pgfpathrectangle{\pgfqpoint{0.553821in}{0.829139in}}{\pgfqpoint{2.745690in}{1.921500in}}%
\pgfusepath{clip}%
\pgfsetrectcap%
\pgfsetroundjoin%
\pgfsetlinewidth{0.803000pt}%
\definecolor{currentstroke}{rgb}{0.690196,0.690196,0.690196}%
\pgfsetstrokecolor{currentstroke}%
\pgfsetstrokeopacity{0.300000}%
\pgfsetdash{}{0pt}%
\pgfpathmoveto{\pgfqpoint{1.263268in}{0.829139in}}%
\pgfpathlineto{\pgfqpoint{1.263268in}{2.750639in}}%
\pgfusepath{stroke}%
\end{pgfscope}%
\begin{pgfscope}%
\pgfsetbuttcap%
\pgfsetroundjoin%
\definecolor{currentfill}{rgb}{0.000000,0.000000,0.000000}%
\pgfsetfillcolor{currentfill}%
\pgfsetlinewidth{0.602250pt}%
\definecolor{currentstroke}{rgb}{0.000000,0.000000,0.000000}%
\pgfsetstrokecolor{currentstroke}%
\pgfsetdash{}{0pt}%
\pgfsys@defobject{currentmarker}{\pgfqpoint{0.000000in}{-0.027778in}}{\pgfqpoint{0.000000in}{0.000000in}}{%
\pgfpathmoveto{\pgfqpoint{0.000000in}{0.000000in}}%
\pgfpathlineto{\pgfqpoint{0.000000in}{-0.027778in}}%
\pgfusepath{stroke,fill}%
}%
\begin{pgfscope}%
\pgfsys@transformshift{1.263268in}{0.829139in}%
\pgfsys@useobject{currentmarker}{}%
\end{pgfscope}%
\end{pgfscope}%
\begin{pgfscope}%
\pgfpathrectangle{\pgfqpoint{0.553821in}{0.829139in}}{\pgfqpoint{2.745690in}{1.921500in}}%
\pgfusepath{clip}%
\pgfsetrectcap%
\pgfsetroundjoin%
\pgfsetlinewidth{0.803000pt}%
\definecolor{currentstroke}{rgb}{0.690196,0.690196,0.690196}%
\pgfsetstrokecolor{currentstroke}%
\pgfsetstrokeopacity{0.300000}%
\pgfsetdash{}{0pt}%
\pgfpathmoveto{\pgfqpoint{1.521976in}{0.829139in}}%
\pgfpathlineto{\pgfqpoint{1.521976in}{2.750639in}}%
\pgfusepath{stroke}%
\end{pgfscope}%
\begin{pgfscope}%
\pgfsetbuttcap%
\pgfsetroundjoin%
\definecolor{currentfill}{rgb}{0.000000,0.000000,0.000000}%
\pgfsetfillcolor{currentfill}%
\pgfsetlinewidth{0.602250pt}%
\definecolor{currentstroke}{rgb}{0.000000,0.000000,0.000000}%
\pgfsetstrokecolor{currentstroke}%
\pgfsetdash{}{0pt}%
\pgfsys@defobject{currentmarker}{\pgfqpoint{0.000000in}{-0.027778in}}{\pgfqpoint{0.000000in}{0.000000in}}{%
\pgfpathmoveto{\pgfqpoint{0.000000in}{0.000000in}}%
\pgfpathlineto{\pgfqpoint{0.000000in}{-0.027778in}}%
\pgfusepath{stroke,fill}%
}%
\begin{pgfscope}%
\pgfsys@transformshift{1.521976in}{0.829139in}%
\pgfsys@useobject{currentmarker}{}%
\end{pgfscope}%
\end{pgfscope}%
\begin{pgfscope}%
\pgfpathrectangle{\pgfqpoint{0.553821in}{0.829139in}}{\pgfqpoint{2.745690in}{1.921500in}}%
\pgfusepath{clip}%
\pgfsetrectcap%
\pgfsetroundjoin%
\pgfsetlinewidth{0.803000pt}%
\definecolor{currentstroke}{rgb}{0.690196,0.690196,0.690196}%
\pgfsetstrokecolor{currentstroke}%
\pgfsetstrokeopacity{0.300000}%
\pgfsetdash{}{0pt}%
\pgfpathmoveto{\pgfqpoint{1.705533in}{0.829139in}}%
\pgfpathlineto{\pgfqpoint{1.705533in}{2.750639in}}%
\pgfusepath{stroke}%
\end{pgfscope}%
\begin{pgfscope}%
\pgfsetbuttcap%
\pgfsetroundjoin%
\definecolor{currentfill}{rgb}{0.000000,0.000000,0.000000}%
\pgfsetfillcolor{currentfill}%
\pgfsetlinewidth{0.602250pt}%
\definecolor{currentstroke}{rgb}{0.000000,0.000000,0.000000}%
\pgfsetstrokecolor{currentstroke}%
\pgfsetdash{}{0pt}%
\pgfsys@defobject{currentmarker}{\pgfqpoint{0.000000in}{-0.027778in}}{\pgfqpoint{0.000000in}{0.000000in}}{%
\pgfpathmoveto{\pgfqpoint{0.000000in}{0.000000in}}%
\pgfpathlineto{\pgfqpoint{0.000000in}{-0.027778in}}%
\pgfusepath{stroke,fill}%
}%
\begin{pgfscope}%
\pgfsys@transformshift{1.705533in}{0.829139in}%
\pgfsys@useobject{currentmarker}{}%
\end{pgfscope}%
\end{pgfscope}%
\begin{pgfscope}%
\pgfpathrectangle{\pgfqpoint{0.553821in}{0.829139in}}{\pgfqpoint{2.745690in}{1.921500in}}%
\pgfusepath{clip}%
\pgfsetrectcap%
\pgfsetroundjoin%
\pgfsetlinewidth{0.803000pt}%
\definecolor{currentstroke}{rgb}{0.690196,0.690196,0.690196}%
\pgfsetstrokecolor{currentstroke}%
\pgfsetstrokeopacity{0.300000}%
\pgfsetdash{}{0pt}%
\pgfpathmoveto{\pgfqpoint{1.847911in}{0.829139in}}%
\pgfpathlineto{\pgfqpoint{1.847911in}{2.750639in}}%
\pgfusepath{stroke}%
\end{pgfscope}%
\begin{pgfscope}%
\pgfsetbuttcap%
\pgfsetroundjoin%
\definecolor{currentfill}{rgb}{0.000000,0.000000,0.000000}%
\pgfsetfillcolor{currentfill}%
\pgfsetlinewidth{0.602250pt}%
\definecolor{currentstroke}{rgb}{0.000000,0.000000,0.000000}%
\pgfsetstrokecolor{currentstroke}%
\pgfsetdash{}{0pt}%
\pgfsys@defobject{currentmarker}{\pgfqpoint{0.000000in}{-0.027778in}}{\pgfqpoint{0.000000in}{0.000000in}}{%
\pgfpathmoveto{\pgfqpoint{0.000000in}{0.000000in}}%
\pgfpathlineto{\pgfqpoint{0.000000in}{-0.027778in}}%
\pgfusepath{stroke,fill}%
}%
\begin{pgfscope}%
\pgfsys@transformshift{1.847911in}{0.829139in}%
\pgfsys@useobject{currentmarker}{}%
\end{pgfscope}%
\end{pgfscope}%
\begin{pgfscope}%
\pgfpathrectangle{\pgfqpoint{0.553821in}{0.829139in}}{\pgfqpoint{2.745690in}{1.921500in}}%
\pgfusepath{clip}%
\pgfsetrectcap%
\pgfsetroundjoin%
\pgfsetlinewidth{0.803000pt}%
\definecolor{currentstroke}{rgb}{0.690196,0.690196,0.690196}%
\pgfsetstrokecolor{currentstroke}%
\pgfsetstrokeopacity{0.300000}%
\pgfsetdash{}{0pt}%
\pgfpathmoveto{\pgfqpoint{1.964241in}{0.829139in}}%
\pgfpathlineto{\pgfqpoint{1.964241in}{2.750639in}}%
\pgfusepath{stroke}%
\end{pgfscope}%
\begin{pgfscope}%
\pgfsetbuttcap%
\pgfsetroundjoin%
\definecolor{currentfill}{rgb}{0.000000,0.000000,0.000000}%
\pgfsetfillcolor{currentfill}%
\pgfsetlinewidth{0.602250pt}%
\definecolor{currentstroke}{rgb}{0.000000,0.000000,0.000000}%
\pgfsetstrokecolor{currentstroke}%
\pgfsetdash{}{0pt}%
\pgfsys@defobject{currentmarker}{\pgfqpoint{0.000000in}{-0.027778in}}{\pgfqpoint{0.000000in}{0.000000in}}{%
\pgfpathmoveto{\pgfqpoint{0.000000in}{0.000000in}}%
\pgfpathlineto{\pgfqpoint{0.000000in}{-0.027778in}}%
\pgfusepath{stroke,fill}%
}%
\begin{pgfscope}%
\pgfsys@transformshift{1.964241in}{0.829139in}%
\pgfsys@useobject{currentmarker}{}%
\end{pgfscope}%
\end{pgfscope}%
\begin{pgfscope}%
\pgfpathrectangle{\pgfqpoint{0.553821in}{0.829139in}}{\pgfqpoint{2.745690in}{1.921500in}}%
\pgfusepath{clip}%
\pgfsetrectcap%
\pgfsetroundjoin%
\pgfsetlinewidth{0.803000pt}%
\definecolor{currentstroke}{rgb}{0.690196,0.690196,0.690196}%
\pgfsetstrokecolor{currentstroke}%
\pgfsetstrokeopacity{0.300000}%
\pgfsetdash{}{0pt}%
\pgfpathmoveto{\pgfqpoint{2.062598in}{0.829139in}}%
\pgfpathlineto{\pgfqpoint{2.062598in}{2.750639in}}%
\pgfusepath{stroke}%
\end{pgfscope}%
\begin{pgfscope}%
\pgfsetbuttcap%
\pgfsetroundjoin%
\definecolor{currentfill}{rgb}{0.000000,0.000000,0.000000}%
\pgfsetfillcolor{currentfill}%
\pgfsetlinewidth{0.602250pt}%
\definecolor{currentstroke}{rgb}{0.000000,0.000000,0.000000}%
\pgfsetstrokecolor{currentstroke}%
\pgfsetdash{}{0pt}%
\pgfsys@defobject{currentmarker}{\pgfqpoint{0.000000in}{-0.027778in}}{\pgfqpoint{0.000000in}{0.000000in}}{%
\pgfpathmoveto{\pgfqpoint{0.000000in}{0.000000in}}%
\pgfpathlineto{\pgfqpoint{0.000000in}{-0.027778in}}%
\pgfusepath{stroke,fill}%
}%
\begin{pgfscope}%
\pgfsys@transformshift{2.062598in}{0.829139in}%
\pgfsys@useobject{currentmarker}{}%
\end{pgfscope}%
\end{pgfscope}%
\begin{pgfscope}%
\pgfpathrectangle{\pgfqpoint{0.553821in}{0.829139in}}{\pgfqpoint{2.745690in}{1.921500in}}%
\pgfusepath{clip}%
\pgfsetrectcap%
\pgfsetroundjoin%
\pgfsetlinewidth{0.803000pt}%
\definecolor{currentstroke}{rgb}{0.690196,0.690196,0.690196}%
\pgfsetstrokecolor{currentstroke}%
\pgfsetstrokeopacity{0.300000}%
\pgfsetdash{}{0pt}%
\pgfpathmoveto{\pgfqpoint{2.147798in}{0.829139in}}%
\pgfpathlineto{\pgfqpoint{2.147798in}{2.750639in}}%
\pgfusepath{stroke}%
\end{pgfscope}%
\begin{pgfscope}%
\pgfsetbuttcap%
\pgfsetroundjoin%
\definecolor{currentfill}{rgb}{0.000000,0.000000,0.000000}%
\pgfsetfillcolor{currentfill}%
\pgfsetlinewidth{0.602250pt}%
\definecolor{currentstroke}{rgb}{0.000000,0.000000,0.000000}%
\pgfsetstrokecolor{currentstroke}%
\pgfsetdash{}{0pt}%
\pgfsys@defobject{currentmarker}{\pgfqpoint{0.000000in}{-0.027778in}}{\pgfqpoint{0.000000in}{0.000000in}}{%
\pgfpathmoveto{\pgfqpoint{0.000000in}{0.000000in}}%
\pgfpathlineto{\pgfqpoint{0.000000in}{-0.027778in}}%
\pgfusepath{stroke,fill}%
}%
\begin{pgfscope}%
\pgfsys@transformshift{2.147798in}{0.829139in}%
\pgfsys@useobject{currentmarker}{}%
\end{pgfscope}%
\end{pgfscope}%
\begin{pgfscope}%
\pgfpathrectangle{\pgfqpoint{0.553821in}{0.829139in}}{\pgfqpoint{2.745690in}{1.921500in}}%
\pgfusepath{clip}%
\pgfsetrectcap%
\pgfsetroundjoin%
\pgfsetlinewidth{0.803000pt}%
\definecolor{currentstroke}{rgb}{0.690196,0.690196,0.690196}%
\pgfsetstrokecolor{currentstroke}%
\pgfsetstrokeopacity{0.300000}%
\pgfsetdash{}{0pt}%
\pgfpathmoveto{\pgfqpoint{2.222950in}{0.829139in}}%
\pgfpathlineto{\pgfqpoint{2.222950in}{2.750639in}}%
\pgfusepath{stroke}%
\end{pgfscope}%
\begin{pgfscope}%
\pgfsetbuttcap%
\pgfsetroundjoin%
\definecolor{currentfill}{rgb}{0.000000,0.000000,0.000000}%
\pgfsetfillcolor{currentfill}%
\pgfsetlinewidth{0.602250pt}%
\definecolor{currentstroke}{rgb}{0.000000,0.000000,0.000000}%
\pgfsetstrokecolor{currentstroke}%
\pgfsetdash{}{0pt}%
\pgfsys@defobject{currentmarker}{\pgfqpoint{0.000000in}{-0.027778in}}{\pgfqpoint{0.000000in}{0.000000in}}{%
\pgfpathmoveto{\pgfqpoint{0.000000in}{0.000000in}}%
\pgfpathlineto{\pgfqpoint{0.000000in}{-0.027778in}}%
\pgfusepath{stroke,fill}%
}%
\begin{pgfscope}%
\pgfsys@transformshift{2.222950in}{0.829139in}%
\pgfsys@useobject{currentmarker}{}%
\end{pgfscope}%
\end{pgfscope}%
\begin{pgfscope}%
\pgfpathrectangle{\pgfqpoint{0.553821in}{0.829139in}}{\pgfqpoint{2.745690in}{1.921500in}}%
\pgfusepath{clip}%
\pgfsetrectcap%
\pgfsetroundjoin%
\pgfsetlinewidth{0.803000pt}%
\definecolor{currentstroke}{rgb}{0.690196,0.690196,0.690196}%
\pgfsetstrokecolor{currentstroke}%
\pgfsetstrokeopacity{0.300000}%
\pgfsetdash{}{0pt}%
\pgfpathmoveto{\pgfqpoint{2.732441in}{0.829139in}}%
\pgfpathlineto{\pgfqpoint{2.732441in}{2.750639in}}%
\pgfusepath{stroke}%
\end{pgfscope}%
\begin{pgfscope}%
\pgfsetbuttcap%
\pgfsetroundjoin%
\definecolor{currentfill}{rgb}{0.000000,0.000000,0.000000}%
\pgfsetfillcolor{currentfill}%
\pgfsetlinewidth{0.602250pt}%
\definecolor{currentstroke}{rgb}{0.000000,0.000000,0.000000}%
\pgfsetstrokecolor{currentstroke}%
\pgfsetdash{}{0pt}%
\pgfsys@defobject{currentmarker}{\pgfqpoint{0.000000in}{-0.027778in}}{\pgfqpoint{0.000000in}{0.000000in}}{%
\pgfpathmoveto{\pgfqpoint{0.000000in}{0.000000in}}%
\pgfpathlineto{\pgfqpoint{0.000000in}{-0.027778in}}%
\pgfusepath{stroke,fill}%
}%
\begin{pgfscope}%
\pgfsys@transformshift{2.732441in}{0.829139in}%
\pgfsys@useobject{currentmarker}{}%
\end{pgfscope}%
\end{pgfscope}%
\begin{pgfscope}%
\pgfpathrectangle{\pgfqpoint{0.553821in}{0.829139in}}{\pgfqpoint{2.745690in}{1.921500in}}%
\pgfusepath{clip}%
\pgfsetrectcap%
\pgfsetroundjoin%
\pgfsetlinewidth{0.803000pt}%
\definecolor{currentstroke}{rgb}{0.690196,0.690196,0.690196}%
\pgfsetstrokecolor{currentstroke}%
\pgfsetstrokeopacity{0.300000}%
\pgfsetdash{}{0pt}%
\pgfpathmoveto{\pgfqpoint{2.991150in}{0.829139in}}%
\pgfpathlineto{\pgfqpoint{2.991150in}{2.750639in}}%
\pgfusepath{stroke}%
\end{pgfscope}%
\begin{pgfscope}%
\pgfsetbuttcap%
\pgfsetroundjoin%
\definecolor{currentfill}{rgb}{0.000000,0.000000,0.000000}%
\pgfsetfillcolor{currentfill}%
\pgfsetlinewidth{0.602250pt}%
\definecolor{currentstroke}{rgb}{0.000000,0.000000,0.000000}%
\pgfsetstrokecolor{currentstroke}%
\pgfsetdash{}{0pt}%
\pgfsys@defobject{currentmarker}{\pgfqpoint{0.000000in}{-0.027778in}}{\pgfqpoint{0.000000in}{0.000000in}}{%
\pgfpathmoveto{\pgfqpoint{0.000000in}{0.000000in}}%
\pgfpathlineto{\pgfqpoint{0.000000in}{-0.027778in}}%
\pgfusepath{stroke,fill}%
}%
\begin{pgfscope}%
\pgfsys@transformshift{2.991150in}{0.829139in}%
\pgfsys@useobject{currentmarker}{}%
\end{pgfscope}%
\end{pgfscope}%
\begin{pgfscope}%
\pgfpathrectangle{\pgfqpoint{0.553821in}{0.829139in}}{\pgfqpoint{2.745690in}{1.921500in}}%
\pgfusepath{clip}%
\pgfsetrectcap%
\pgfsetroundjoin%
\pgfsetlinewidth{0.803000pt}%
\definecolor{currentstroke}{rgb}{0.690196,0.690196,0.690196}%
\pgfsetstrokecolor{currentstroke}%
\pgfsetstrokeopacity{0.300000}%
\pgfsetdash{}{0pt}%
\pgfpathmoveto{\pgfqpoint{3.174706in}{0.829139in}}%
\pgfpathlineto{\pgfqpoint{3.174706in}{2.750639in}}%
\pgfusepath{stroke}%
\end{pgfscope}%
\begin{pgfscope}%
\pgfsetbuttcap%
\pgfsetroundjoin%
\definecolor{currentfill}{rgb}{0.000000,0.000000,0.000000}%
\pgfsetfillcolor{currentfill}%
\pgfsetlinewidth{0.602250pt}%
\definecolor{currentstroke}{rgb}{0.000000,0.000000,0.000000}%
\pgfsetstrokecolor{currentstroke}%
\pgfsetdash{}{0pt}%
\pgfsys@defobject{currentmarker}{\pgfqpoint{0.000000in}{-0.027778in}}{\pgfqpoint{0.000000in}{0.000000in}}{%
\pgfpathmoveto{\pgfqpoint{0.000000in}{0.000000in}}%
\pgfpathlineto{\pgfqpoint{0.000000in}{-0.027778in}}%
\pgfusepath{stroke,fill}%
}%
\begin{pgfscope}%
\pgfsys@transformshift{3.174706in}{0.829139in}%
\pgfsys@useobject{currentmarker}{}%
\end{pgfscope}%
\end{pgfscope}%
\begin{pgfscope}%
\definecolor{textcolor}{rgb}{0.000000,0.000000,0.000000}%
\pgfsetstrokecolor{textcolor}%
\pgfsetfillcolor{textcolor}%
\pgftext[x=1.926666in,y=0.581206in,,top]{\color{textcolor}{\rmfamily\fontsize{8.000000}{9.600000}\selectfont\catcode`\^=\active\def^{\ifmmode\sp\else\^{}\fi}\catcode`\%=\active\def%{\%}Diffusion Time $t$}}%
\end{pgfscope}%
\begin{pgfscope}%
\pgfpathrectangle{\pgfqpoint{0.553821in}{0.829139in}}{\pgfqpoint{2.745690in}{1.921500in}}%
\pgfusepath{clip}%
\pgfsetrectcap%
\pgfsetroundjoin%
\pgfsetlinewidth{0.803000pt}%
\definecolor{currentstroke}{rgb}{0.690196,0.690196,0.690196}%
\pgfsetstrokecolor{currentstroke}%
\pgfsetstrokeopacity{0.300000}%
\pgfsetdash{}{0pt}%
\pgfpathmoveto{\pgfqpoint{0.553821in}{2.545163in}}%
\pgfpathlineto{\pgfqpoint{3.299510in}{2.545163in}}%
\pgfusepath{stroke}%
\end{pgfscope}%
\begin{pgfscope}%
\pgfsetbuttcap%
\pgfsetroundjoin%
\definecolor{currentfill}{rgb}{0.000000,0.000000,0.000000}%
\pgfsetfillcolor{currentfill}%
\pgfsetlinewidth{0.803000pt}%
\definecolor{currentstroke}{rgb}{0.000000,0.000000,0.000000}%
\pgfsetstrokecolor{currentstroke}%
\pgfsetdash{}{0pt}%
\pgfsys@defobject{currentmarker}{\pgfqpoint{-0.048611in}{0.000000in}}{\pgfqpoint{-0.000000in}{0.000000in}}{%
\pgfpathmoveto{\pgfqpoint{-0.000000in}{0.000000in}}%
\pgfpathlineto{\pgfqpoint{-0.048611in}{0.000000in}}%
\pgfusepath{stroke,fill}%
}%
\begin{pgfscope}%
\pgfsys@transformshift{0.553821in}{2.545163in}%
\pgfsys@useobject{currentmarker}{}%
\end{pgfscope}%
\end{pgfscope}%
\begin{pgfscope}%
\definecolor{textcolor}{rgb}{0.000000,0.000000,0.000000}%
\pgfsetstrokecolor{textcolor}%
\pgfsetfillcolor{textcolor}%
\pgftext[x=0.291667in, y=2.507809in, left, base]{\color{textcolor}{\rmfamily\fontsize{7.000000}{8.400000}\selectfont\catcode`\^=\active\def^{\ifmmode\sp\else\^{}\fi}\catcode`\%=\active\def%{\%}$\mathdefault{10^{2}}$}}%
\end{pgfscope}%
\begin{pgfscope}%
\pgfpathrectangle{\pgfqpoint{0.553821in}{0.829139in}}{\pgfqpoint{2.745690in}{1.921500in}}%
\pgfusepath{clip}%
\pgfsetrectcap%
\pgfsetroundjoin%
\pgfsetlinewidth{0.803000pt}%
\definecolor{currentstroke}{rgb}{0.690196,0.690196,0.690196}%
\pgfsetstrokecolor{currentstroke}%
\pgfsetstrokeopacity{0.300000}%
\pgfsetdash{}{0pt}%
\pgfpathmoveto{\pgfqpoint{0.553821in}{1.332847in}}%
\pgfpathlineto{\pgfqpoint{3.299510in}{1.332847in}}%
\pgfusepath{stroke}%
\end{pgfscope}%
\begin{pgfscope}%
\pgfsetbuttcap%
\pgfsetroundjoin%
\definecolor{currentfill}{rgb}{0.000000,0.000000,0.000000}%
\pgfsetfillcolor{currentfill}%
\pgfsetlinewidth{0.602250pt}%
\definecolor{currentstroke}{rgb}{0.000000,0.000000,0.000000}%
\pgfsetstrokecolor{currentstroke}%
\pgfsetdash{}{0pt}%
\pgfsys@defobject{currentmarker}{\pgfqpoint{-0.027778in}{0.000000in}}{\pgfqpoint{-0.000000in}{0.000000in}}{%
\pgfpathmoveto{\pgfqpoint{-0.000000in}{0.000000in}}%
\pgfpathlineto{\pgfqpoint{-0.027778in}{0.000000in}}%
\pgfusepath{stroke,fill}%
}%
\begin{pgfscope}%
\pgfsys@transformshift{0.553821in}{1.332847in}%
\pgfsys@useobject{currentmarker}{}%
\end{pgfscope}%
\end{pgfscope}%
\begin{pgfscope}%
\pgfpathrectangle{\pgfqpoint{0.553821in}{0.829139in}}{\pgfqpoint{2.745690in}{1.921500in}}%
\pgfusepath{clip}%
\pgfsetrectcap%
\pgfsetroundjoin%
\pgfsetlinewidth{0.803000pt}%
\definecolor{currentstroke}{rgb}{0.690196,0.690196,0.690196}%
\pgfsetstrokecolor{currentstroke}%
\pgfsetstrokeopacity{0.300000}%
\pgfsetdash{}{0pt}%
\pgfpathmoveto{\pgfqpoint{0.553821in}{1.638265in}}%
\pgfpathlineto{\pgfqpoint{3.299510in}{1.638265in}}%
\pgfusepath{stroke}%
\end{pgfscope}%
\begin{pgfscope}%
\pgfsetbuttcap%
\pgfsetroundjoin%
\definecolor{currentfill}{rgb}{0.000000,0.000000,0.000000}%
\pgfsetfillcolor{currentfill}%
\pgfsetlinewidth{0.602250pt}%
\definecolor{currentstroke}{rgb}{0.000000,0.000000,0.000000}%
\pgfsetstrokecolor{currentstroke}%
\pgfsetdash{}{0pt}%
\pgfsys@defobject{currentmarker}{\pgfqpoint{-0.027778in}{0.000000in}}{\pgfqpoint{-0.000000in}{0.000000in}}{%
\pgfpathmoveto{\pgfqpoint{-0.000000in}{0.000000in}}%
\pgfpathlineto{\pgfqpoint{-0.027778in}{0.000000in}}%
\pgfusepath{stroke,fill}%
}%
\begin{pgfscope}%
\pgfsys@transformshift{0.553821in}{1.638265in}%
\pgfsys@useobject{currentmarker}{}%
\end{pgfscope}%
\end{pgfscope}%
\begin{pgfscope}%
\pgfpathrectangle{\pgfqpoint{0.553821in}{0.829139in}}{\pgfqpoint{2.745690in}{1.921500in}}%
\pgfusepath{clip}%
\pgfsetrectcap%
\pgfsetroundjoin%
\pgfsetlinewidth{0.803000pt}%
\definecolor{currentstroke}{rgb}{0.690196,0.690196,0.690196}%
\pgfsetstrokecolor{currentstroke}%
\pgfsetstrokeopacity{0.300000}%
\pgfsetdash{}{0pt}%
\pgfpathmoveto{\pgfqpoint{0.553821in}{1.854963in}}%
\pgfpathlineto{\pgfqpoint{3.299510in}{1.854963in}}%
\pgfusepath{stroke}%
\end{pgfscope}%
\begin{pgfscope}%
\pgfsetbuttcap%
\pgfsetroundjoin%
\definecolor{currentfill}{rgb}{0.000000,0.000000,0.000000}%
\pgfsetfillcolor{currentfill}%
\pgfsetlinewidth{0.602250pt}%
\definecolor{currentstroke}{rgb}{0.000000,0.000000,0.000000}%
\pgfsetstrokecolor{currentstroke}%
\pgfsetdash{}{0pt}%
\pgfsys@defobject{currentmarker}{\pgfqpoint{-0.027778in}{0.000000in}}{\pgfqpoint{-0.000000in}{0.000000in}}{%
\pgfpathmoveto{\pgfqpoint{-0.000000in}{0.000000in}}%
\pgfpathlineto{\pgfqpoint{-0.027778in}{0.000000in}}%
\pgfusepath{stroke,fill}%
}%
\begin{pgfscope}%
\pgfsys@transformshift{0.553821in}{1.854963in}%
\pgfsys@useobject{currentmarker}{}%
\end{pgfscope}%
\end{pgfscope}%
\begin{pgfscope}%
\pgfpathrectangle{\pgfqpoint{0.553821in}{0.829139in}}{\pgfqpoint{2.745690in}{1.921500in}}%
\pgfusepath{clip}%
\pgfsetrectcap%
\pgfsetroundjoin%
\pgfsetlinewidth{0.803000pt}%
\definecolor{currentstroke}{rgb}{0.690196,0.690196,0.690196}%
\pgfsetstrokecolor{currentstroke}%
\pgfsetstrokeopacity{0.300000}%
\pgfsetdash{}{0pt}%
\pgfpathmoveto{\pgfqpoint{0.553821in}{2.023047in}}%
\pgfpathlineto{\pgfqpoint{3.299510in}{2.023047in}}%
\pgfusepath{stroke}%
\end{pgfscope}%
\begin{pgfscope}%
\pgfsetbuttcap%
\pgfsetroundjoin%
\definecolor{currentfill}{rgb}{0.000000,0.000000,0.000000}%
\pgfsetfillcolor{currentfill}%
\pgfsetlinewidth{0.602250pt}%
\definecolor{currentstroke}{rgb}{0.000000,0.000000,0.000000}%
\pgfsetstrokecolor{currentstroke}%
\pgfsetdash{}{0pt}%
\pgfsys@defobject{currentmarker}{\pgfqpoint{-0.027778in}{0.000000in}}{\pgfqpoint{-0.000000in}{0.000000in}}{%
\pgfpathmoveto{\pgfqpoint{-0.000000in}{0.000000in}}%
\pgfpathlineto{\pgfqpoint{-0.027778in}{0.000000in}}%
\pgfusepath{stroke,fill}%
}%
\begin{pgfscope}%
\pgfsys@transformshift{0.553821in}{2.023047in}%
\pgfsys@useobject{currentmarker}{}%
\end{pgfscope}%
\end{pgfscope}%
\begin{pgfscope}%
\pgfpathrectangle{\pgfqpoint{0.553821in}{0.829139in}}{\pgfqpoint{2.745690in}{1.921500in}}%
\pgfusepath{clip}%
\pgfsetrectcap%
\pgfsetroundjoin%
\pgfsetlinewidth{0.803000pt}%
\definecolor{currentstroke}{rgb}{0.690196,0.690196,0.690196}%
\pgfsetstrokecolor{currentstroke}%
\pgfsetstrokeopacity{0.300000}%
\pgfsetdash{}{0pt}%
\pgfpathmoveto{\pgfqpoint{0.553821in}{2.160381in}}%
\pgfpathlineto{\pgfqpoint{3.299510in}{2.160381in}}%
\pgfusepath{stroke}%
\end{pgfscope}%
\begin{pgfscope}%
\pgfsetbuttcap%
\pgfsetroundjoin%
\definecolor{currentfill}{rgb}{0.000000,0.000000,0.000000}%
\pgfsetfillcolor{currentfill}%
\pgfsetlinewidth{0.602250pt}%
\definecolor{currentstroke}{rgb}{0.000000,0.000000,0.000000}%
\pgfsetstrokecolor{currentstroke}%
\pgfsetdash{}{0pt}%
\pgfsys@defobject{currentmarker}{\pgfqpoint{-0.027778in}{0.000000in}}{\pgfqpoint{-0.000000in}{0.000000in}}{%
\pgfpathmoveto{\pgfqpoint{-0.000000in}{0.000000in}}%
\pgfpathlineto{\pgfqpoint{-0.027778in}{0.000000in}}%
\pgfusepath{stroke,fill}%
}%
\begin{pgfscope}%
\pgfsys@transformshift{0.553821in}{2.160381in}%
\pgfsys@useobject{currentmarker}{}%
\end{pgfscope}%
\end{pgfscope}%
\begin{pgfscope}%
\pgfpathrectangle{\pgfqpoint{0.553821in}{0.829139in}}{\pgfqpoint{2.745690in}{1.921500in}}%
\pgfusepath{clip}%
\pgfsetrectcap%
\pgfsetroundjoin%
\pgfsetlinewidth{0.803000pt}%
\definecolor{currentstroke}{rgb}{0.690196,0.690196,0.690196}%
\pgfsetstrokecolor{currentstroke}%
\pgfsetstrokeopacity{0.300000}%
\pgfsetdash{}{0pt}%
\pgfpathmoveto{\pgfqpoint{0.553821in}{2.276496in}}%
\pgfpathlineto{\pgfqpoint{3.299510in}{2.276496in}}%
\pgfusepath{stroke}%
\end{pgfscope}%
\begin{pgfscope}%
\pgfsetbuttcap%
\pgfsetroundjoin%
\definecolor{currentfill}{rgb}{0.000000,0.000000,0.000000}%
\pgfsetfillcolor{currentfill}%
\pgfsetlinewidth{0.602250pt}%
\definecolor{currentstroke}{rgb}{0.000000,0.000000,0.000000}%
\pgfsetstrokecolor{currentstroke}%
\pgfsetdash{}{0pt}%
\pgfsys@defobject{currentmarker}{\pgfqpoint{-0.027778in}{0.000000in}}{\pgfqpoint{-0.000000in}{0.000000in}}{%
\pgfpathmoveto{\pgfqpoint{-0.000000in}{0.000000in}}%
\pgfpathlineto{\pgfqpoint{-0.027778in}{0.000000in}}%
\pgfusepath{stroke,fill}%
}%
\begin{pgfscope}%
\pgfsys@transformshift{0.553821in}{2.276496in}%
\pgfsys@useobject{currentmarker}{}%
\end{pgfscope}%
\end{pgfscope}%
\begin{pgfscope}%
\pgfpathrectangle{\pgfqpoint{0.553821in}{0.829139in}}{\pgfqpoint{2.745690in}{1.921500in}}%
\pgfusepath{clip}%
\pgfsetrectcap%
\pgfsetroundjoin%
\pgfsetlinewidth{0.803000pt}%
\definecolor{currentstroke}{rgb}{0.690196,0.690196,0.690196}%
\pgfsetstrokecolor{currentstroke}%
\pgfsetstrokeopacity{0.300000}%
\pgfsetdash{}{0pt}%
\pgfpathmoveto{\pgfqpoint{0.553821in}{2.377079in}}%
\pgfpathlineto{\pgfqpoint{3.299510in}{2.377079in}}%
\pgfusepath{stroke}%
\end{pgfscope}%
\begin{pgfscope}%
\pgfsetbuttcap%
\pgfsetroundjoin%
\definecolor{currentfill}{rgb}{0.000000,0.000000,0.000000}%
\pgfsetfillcolor{currentfill}%
\pgfsetlinewidth{0.602250pt}%
\definecolor{currentstroke}{rgb}{0.000000,0.000000,0.000000}%
\pgfsetstrokecolor{currentstroke}%
\pgfsetdash{}{0pt}%
\pgfsys@defobject{currentmarker}{\pgfqpoint{-0.027778in}{0.000000in}}{\pgfqpoint{-0.000000in}{0.000000in}}{%
\pgfpathmoveto{\pgfqpoint{-0.000000in}{0.000000in}}%
\pgfpathlineto{\pgfqpoint{-0.027778in}{0.000000in}}%
\pgfusepath{stroke,fill}%
}%
\begin{pgfscope}%
\pgfsys@transformshift{0.553821in}{2.377079in}%
\pgfsys@useobject{currentmarker}{}%
\end{pgfscope}%
\end{pgfscope}%
\begin{pgfscope}%
\pgfpathrectangle{\pgfqpoint{0.553821in}{0.829139in}}{\pgfqpoint{2.745690in}{1.921500in}}%
\pgfusepath{clip}%
\pgfsetrectcap%
\pgfsetroundjoin%
\pgfsetlinewidth{0.803000pt}%
\definecolor{currentstroke}{rgb}{0.690196,0.690196,0.690196}%
\pgfsetstrokecolor{currentstroke}%
\pgfsetstrokeopacity{0.300000}%
\pgfsetdash{}{0pt}%
\pgfpathmoveto{\pgfqpoint{0.553821in}{2.465800in}}%
\pgfpathlineto{\pgfqpoint{3.299510in}{2.465800in}}%
\pgfusepath{stroke}%
\end{pgfscope}%
\begin{pgfscope}%
\pgfsetbuttcap%
\pgfsetroundjoin%
\definecolor{currentfill}{rgb}{0.000000,0.000000,0.000000}%
\pgfsetfillcolor{currentfill}%
\pgfsetlinewidth{0.602250pt}%
\definecolor{currentstroke}{rgb}{0.000000,0.000000,0.000000}%
\pgfsetstrokecolor{currentstroke}%
\pgfsetdash{}{0pt}%
\pgfsys@defobject{currentmarker}{\pgfqpoint{-0.027778in}{0.000000in}}{\pgfqpoint{-0.000000in}{0.000000in}}{%
\pgfpathmoveto{\pgfqpoint{-0.000000in}{0.000000in}}%
\pgfpathlineto{\pgfqpoint{-0.027778in}{0.000000in}}%
\pgfusepath{stroke,fill}%
}%
\begin{pgfscope}%
\pgfsys@transformshift{0.553821in}{2.465800in}%
\pgfsys@useobject{currentmarker}{}%
\end{pgfscope}%
\end{pgfscope}%
\begin{pgfscope}%
\definecolor{textcolor}{rgb}{0.000000,0.000000,0.000000}%
\pgfsetstrokecolor{textcolor}%
\pgfsetfillcolor{textcolor}%
\pgftext[x=0.236111in,y=1.789889in,,bottom,rotate=90.000000]{\color{textcolor}{\rmfamily\fontsize{8.000000}{9.600000}\selectfont\catcode`\^=\active\def^{\ifmmode\sp\else\^{}\fi}\catcode`\%=\active\def%{\%}Heat Power $\widehat A_w(t)$}}%
\end{pgfscope}%
\begin{pgfscope}%
\pgfpathrectangle{\pgfqpoint{0.553821in}{0.829139in}}{\pgfqpoint{2.745690in}{1.921500in}}%
\pgfusepath{clip}%
\pgfsetrectcap%
\pgfsetroundjoin%
\pgfsetlinewidth{1.706375pt}%
\definecolor{currentstroke}{rgb}{0.121569,0.466667,0.705882}%
\pgfsetstrokecolor{currentstroke}%
\pgfsetdash{}{0pt}%
\pgfpathmoveto{\pgfqpoint{0.678625in}{2.663298in}}%
\pgfpathlineto{\pgfqpoint{0.706671in}{2.656689in}}%
\pgfpathlineto{\pgfqpoint{0.734716in}{2.649868in}}%
\pgfpathlineto{\pgfqpoint{0.762762in}{2.642829in}}%
\pgfpathlineto{\pgfqpoint{0.790808in}{2.635564in}}%
\pgfpathlineto{\pgfqpoint{0.818854in}{2.628067in}}%
\pgfpathlineto{\pgfqpoint{0.846900in}{2.620332in}}%
\pgfpathlineto{\pgfqpoint{0.874946in}{2.612352in}}%
\pgfpathlineto{\pgfqpoint{0.902992in}{2.604121in}}%
\pgfpathlineto{\pgfqpoint{0.931038in}{2.595633in}}%
\pgfpathlineto{\pgfqpoint{0.959083in}{2.586881in}}%
\pgfpathlineto{\pgfqpoint{0.987129in}{2.577859in}}%
\pgfpathlineto{\pgfqpoint{1.015175in}{2.568562in}}%
\pgfpathlineto{\pgfqpoint{1.043221in}{2.558984in}}%
\pgfpathlineto{\pgfqpoint{1.071267in}{2.549119in}}%
\pgfpathlineto{\pgfqpoint{1.099313in}{2.538962in}}%
\pgfpathlineto{\pgfqpoint{1.127359in}{2.528507in}}%
\pgfpathlineto{\pgfqpoint{1.155404in}{2.517750in}}%
\pgfpathlineto{\pgfqpoint{1.183450in}{2.506685in}}%
\pgfpathlineto{\pgfqpoint{1.211496in}{2.495307in}}%
\pgfpathlineto{\pgfqpoint{1.239542in}{2.483612in}}%
\pgfpathlineto{\pgfqpoint{1.267588in}{2.471596in}}%
\pgfpathlineto{\pgfqpoint{1.295634in}{2.459255in}}%
\pgfpathlineto{\pgfqpoint{1.323680in}{2.446584in}}%
\pgfpathlineto{\pgfqpoint{1.351725in}{2.433580in}}%
\pgfpathlineto{\pgfqpoint{1.379771in}{2.420240in}}%
\pgfpathlineto{\pgfqpoint{1.407817in}{2.406560in}}%
\pgfpathlineto{\pgfqpoint{1.435863in}{2.392538in}}%
\pgfpathlineto{\pgfqpoint{1.463909in}{2.378171in}}%
\pgfpathlineto{\pgfqpoint{1.491955in}{2.363458in}}%
\pgfpathlineto{\pgfqpoint{1.520001in}{2.348396in}}%
\pgfpathlineto{\pgfqpoint{1.548046in}{2.332983in}}%
\pgfpathlineto{\pgfqpoint{1.576092in}{2.317219in}}%
\pgfpathlineto{\pgfqpoint{1.604138in}{2.301103in}}%
\pgfpathlineto{\pgfqpoint{1.632184in}{2.284633in}}%
\pgfpathlineto{\pgfqpoint{1.660230in}{2.267811in}}%
\pgfpathlineto{\pgfqpoint{1.688276in}{2.250636in}}%
\pgfpathlineto{\pgfqpoint{1.716322in}{2.233108in}}%
\pgfpathlineto{\pgfqpoint{1.744367in}{2.215229in}}%
\pgfpathlineto{\pgfqpoint{1.772413in}{2.197000in}}%
\pgfpathlineto{\pgfqpoint{1.800459in}{2.178422in}}%
\pgfpathlineto{\pgfqpoint{1.828505in}{2.159497in}}%
\pgfpathlineto{\pgfqpoint{1.856551in}{2.140228in}}%
\pgfpathlineto{\pgfqpoint{1.884597in}{2.120616in}}%
\pgfpathlineto{\pgfqpoint{1.912643in}{2.100666in}}%
\pgfpathlineto{\pgfqpoint{1.940688in}{2.080379in}}%
\pgfpathlineto{\pgfqpoint{1.968734in}{2.059760in}}%
\pgfpathlineto{\pgfqpoint{1.996780in}{2.038812in}}%
\pgfpathlineto{\pgfqpoint{2.024826in}{2.017540in}}%
\pgfpathlineto{\pgfqpoint{2.052872in}{1.995946in}}%
\pgfpathlineto{\pgfqpoint{2.080918in}{1.974038in}}%
\pgfpathlineto{\pgfqpoint{2.108964in}{1.951818in}}%
\pgfpathlineto{\pgfqpoint{2.137009in}{1.929291in}}%
\pgfpathlineto{\pgfqpoint{2.165055in}{1.906464in}}%
\pgfpathlineto{\pgfqpoint{2.193101in}{1.883342in}}%
\pgfpathlineto{\pgfqpoint{2.221147in}{1.859930in}}%
\pgfpathlineto{\pgfqpoint{2.249193in}{1.836233in}}%
\pgfpathlineto{\pgfqpoint{2.277239in}{1.812259in}}%
\pgfpathlineto{\pgfqpoint{2.305285in}{1.788012in}}%
\pgfpathlineto{\pgfqpoint{2.333330in}{1.763499in}}%
\pgfpathlineto{\pgfqpoint{2.361376in}{1.738726in}}%
\pgfpathlineto{\pgfqpoint{2.389422in}{1.713700in}}%
\pgfpathlineto{\pgfqpoint{2.417468in}{1.688427in}}%
\pgfpathlineto{\pgfqpoint{2.445514in}{1.662913in}}%
\pgfpathlineto{\pgfqpoint{2.473560in}{1.637166in}}%
\pgfpathlineto{\pgfqpoint{2.501606in}{1.611191in}}%
\pgfpathlineto{\pgfqpoint{2.529652in}{1.584995in}}%
\pgfpathlineto{\pgfqpoint{2.557697in}{1.558586in}}%
\pgfpathlineto{\pgfqpoint{2.585743in}{1.531968in}}%
\pgfpathlineto{\pgfqpoint{2.613789in}{1.505151in}}%
\pgfpathlineto{\pgfqpoint{2.641835in}{1.478139in}}%
\pgfpathlineto{\pgfqpoint{2.669881in}{1.450939in}}%
\pgfpathlineto{\pgfqpoint{2.697927in}{1.423559in}}%
\pgfpathlineto{\pgfqpoint{2.725973in}{1.396004in}}%
\pgfpathlineto{\pgfqpoint{2.754018in}{1.368282in}}%
\pgfpathlineto{\pgfqpoint{2.782064in}{1.340398in}}%
\pgfpathlineto{\pgfqpoint{2.810110in}{1.312360in}}%
\pgfpathlineto{\pgfqpoint{2.838156in}{1.284174in}}%
\pgfpathlineto{\pgfqpoint{2.866202in}{1.255845in}}%
\pgfpathlineto{\pgfqpoint{2.894248in}{1.227381in}}%
\pgfpathlineto{\pgfqpoint{2.922294in}{1.198787in}}%
\pgfpathlineto{\pgfqpoint{2.950339in}{1.170071in}}%
\pgfpathlineto{\pgfqpoint{2.978385in}{1.141237in}}%
\pgfpathlineto{\pgfqpoint{3.006431in}{1.112292in}}%
\pgfpathlineto{\pgfqpoint{3.034477in}{1.083243in}}%
\pgfpathlineto{\pgfqpoint{3.062523in}{1.054094in}}%
\pgfpathlineto{\pgfqpoint{3.090569in}{1.024853in}}%
\pgfpathlineto{\pgfqpoint{3.118615in}{0.995525in}}%
\pgfpathlineto{\pgfqpoint{3.146660in}{0.966115in}}%
\pgfpathlineto{\pgfqpoint{3.174706in}{0.936630in}}%
\pgfusepath{stroke}%
\end{pgfscope}%
\begin{pgfscope}%
\pgfpathrectangle{\pgfqpoint{0.553821in}{0.829139in}}{\pgfqpoint{2.745690in}{1.921500in}}%
\pgfusepath{clip}%
\pgfsetbuttcap%
\pgfsetroundjoin%
\pgfsetlinewidth{1.505625pt}%
\definecolor{currentstroke}{rgb}{1.000000,0.498039,0.054902}%
\pgfsetstrokecolor{currentstroke}%
\pgfsetdash{{5.550000pt}{2.400000pt}}{0.000000pt}%
\pgfpathmoveto{\pgfqpoint{0.678625in}{2.662896in}}%
\pgfpathlineto{\pgfqpoint{0.706671in}{2.656269in}}%
\pgfpathlineto{\pgfqpoint{0.734716in}{2.649429in}}%
\pgfpathlineto{\pgfqpoint{0.762762in}{2.642370in}}%
\pgfpathlineto{\pgfqpoint{0.790808in}{2.635085in}}%
\pgfpathlineto{\pgfqpoint{0.818854in}{2.627567in}}%
\pgfpathlineto{\pgfqpoint{0.846900in}{2.619809in}}%
\pgfpathlineto{\pgfqpoint{0.874946in}{2.611806in}}%
\pgfpathlineto{\pgfqpoint{0.902992in}{2.603550in}}%
\pgfpathlineto{\pgfqpoint{0.931038in}{2.595036in}}%
\pgfpathlineto{\pgfqpoint{0.959083in}{2.586257in}}%
\pgfpathlineto{\pgfqpoint{0.987129in}{2.577208in}}%
\pgfpathlineto{\pgfqpoint{1.015175in}{2.567881in}}%
\pgfpathlineto{\pgfqpoint{1.043221in}{2.558273in}}%
\pgfpathlineto{\pgfqpoint{1.071267in}{2.548376in}}%
\pgfpathlineto{\pgfqpoint{1.099313in}{2.538185in}}%
\pgfpathlineto{\pgfqpoint{1.127359in}{2.527696in}}%
\pgfpathlineto{\pgfqpoint{1.155404in}{2.516902in}}%
\pgfpathlineto{\pgfqpoint{1.183450in}{2.505798in}}%
\pgfpathlineto{\pgfqpoint{1.211496in}{2.494381in}}%
\pgfpathlineto{\pgfqpoint{1.239542in}{2.482644in}}%
\pgfpathlineto{\pgfqpoint{1.267588in}{2.470585in}}%
\pgfpathlineto{\pgfqpoint{1.295634in}{2.458198in}}%
\pgfpathlineto{\pgfqpoint{1.323680in}{2.445480in}}%
\pgfpathlineto{\pgfqpoint{1.351725in}{2.432426in}}%
\pgfpathlineto{\pgfqpoint{1.379771in}{2.419034in}}%
\pgfpathlineto{\pgfqpoint{1.407817in}{2.405300in}}%
\pgfpathlineto{\pgfqpoint{1.435863in}{2.391221in}}%
\pgfpathlineto{\pgfqpoint{1.463909in}{2.376795in}}%
\pgfpathlineto{\pgfqpoint{1.491955in}{2.362020in}}%
\pgfpathlineto{\pgfqpoint{1.520001in}{2.346893in}}%
\pgfpathlineto{\pgfqpoint{1.548046in}{2.331413in}}%
\pgfpathlineto{\pgfqpoint{1.576092in}{2.315578in}}%
\pgfpathlineto{\pgfqpoint{1.604138in}{2.299388in}}%
\pgfpathlineto{\pgfqpoint{1.632184in}{2.282842in}}%
\pgfpathlineto{\pgfqpoint{1.660230in}{2.265939in}}%
\pgfpathlineto{\pgfqpoint{1.688276in}{2.248680in}}%
\pgfpathlineto{\pgfqpoint{1.716322in}{2.231064in}}%
\pgfpathlineto{\pgfqpoint{1.744367in}{2.213093in}}%
\pgfpathlineto{\pgfqpoint{1.772413in}{2.194768in}}%
\pgfpathlineto{\pgfqpoint{1.800459in}{2.176090in}}%
\pgfpathlineto{\pgfqpoint{1.828505in}{2.157060in}}%
\pgfpathlineto{\pgfqpoint{1.856551in}{2.137681in}}%
\pgfpathlineto{\pgfqpoint{1.884597in}{2.117955in}}%
\pgfpathlineto{\pgfqpoint{1.912643in}{2.097885in}}%
\pgfpathlineto{\pgfqpoint{1.940688in}{2.077473in}}%
\pgfpathlineto{\pgfqpoint{1.968734in}{2.056724in}}%
\pgfpathlineto{\pgfqpoint{1.996780in}{2.035639in}}%
\pgfpathlineto{\pgfqpoint{2.024826in}{2.014224in}}%
\pgfpathlineto{\pgfqpoint{2.052872in}{1.992482in}}%
\pgfpathlineto{\pgfqpoint{2.080918in}{1.970417in}}%
\pgfpathlineto{\pgfqpoint{2.108964in}{1.948034in}}%
\pgfpathlineto{\pgfqpoint{2.137009in}{1.925338in}}%
\pgfpathlineto{\pgfqpoint{2.165055in}{1.902333in}}%
\pgfpathlineto{\pgfqpoint{2.193101in}{1.879025in}}%
\pgfpathlineto{\pgfqpoint{2.221147in}{1.855419in}}%
\pgfpathlineto{\pgfqpoint{2.249193in}{1.831519in}}%
\pgfpathlineto{\pgfqpoint{2.277239in}{1.807333in}}%
\pgfpathlineto{\pgfqpoint{2.305285in}{1.782864in}}%
\pgfpathlineto{\pgfqpoint{2.333330in}{1.758120in}}%
\pgfpathlineto{\pgfqpoint{2.361376in}{1.733105in}}%
\pgfpathlineto{\pgfqpoint{2.389422in}{1.707826in}}%
\pgfpathlineto{\pgfqpoint{2.417468in}{1.682289in}}%
\pgfpathlineto{\pgfqpoint{2.445514in}{1.656499in}}%
\pgfpathlineto{\pgfqpoint{2.473560in}{1.630463in}}%
\pgfpathlineto{\pgfqpoint{2.501606in}{1.604187in}}%
\pgfpathlineto{\pgfqpoint{2.529652in}{1.577676in}}%
\pgfpathlineto{\pgfqpoint{2.557697in}{1.550937in}}%
\pgfpathlineto{\pgfqpoint{2.585743in}{1.523976in}}%
\pgfpathlineto{\pgfqpoint{2.613789in}{1.496799in}}%
\pgfpathlineto{\pgfqpoint{2.641835in}{1.469411in}}%
\pgfpathlineto{\pgfqpoint{2.669881in}{1.441819in}}%
\pgfpathlineto{\pgfqpoint{2.697927in}{1.414028in}}%
\pgfpathlineto{\pgfqpoint{2.725973in}{1.386045in}}%
\pgfpathlineto{\pgfqpoint{2.754018in}{1.357874in}}%
\pgfpathlineto{\pgfqpoint{2.782064in}{1.329522in}}%
\pgfpathlineto{\pgfqpoint{2.810110in}{1.300995in}}%
\pgfpathlineto{\pgfqpoint{2.838156in}{1.272296in}}%
\pgfpathlineto{\pgfqpoint{2.866202in}{1.243433in}}%
\pgfpathlineto{\pgfqpoint{2.894248in}{1.214410in}}%
\pgfpathlineto{\pgfqpoint{2.922294in}{1.185233in}}%
\pgfpathlineto{\pgfqpoint{2.950339in}{1.155906in}}%
\pgfpathlineto{\pgfqpoint{2.978385in}{1.126434in}}%
\pgfpathlineto{\pgfqpoint{3.006431in}{1.096823in}}%
\pgfpathlineto{\pgfqpoint{3.034477in}{1.067077in}}%
\pgfpathlineto{\pgfqpoint{3.062523in}{1.037200in}}%
\pgfpathlineto{\pgfqpoint{3.090569in}{1.007198in}}%
\pgfpathlineto{\pgfqpoint{3.118615in}{0.977074in}}%
\pgfpathlineto{\pgfqpoint{3.146660in}{0.946834in}}%
\pgfpathlineto{\pgfqpoint{3.174706in}{0.916480in}}%
\pgfusepath{stroke}%
\end{pgfscope}%
\begin{pgfscope}%
\pgfsetrectcap%
\pgfsetmiterjoin%
\pgfsetlinewidth{0.803000pt}%
\definecolor{currentstroke}{rgb}{0.000000,0.000000,0.000000}%
\pgfsetstrokecolor{currentstroke}%
\pgfsetdash{}{0pt}%
\pgfpathmoveto{\pgfqpoint{0.553821in}{0.829139in}}%
\pgfpathlineto{\pgfqpoint{0.553821in}{2.750639in}}%
\pgfusepath{stroke}%
\end{pgfscope}%
\begin{pgfscope}%
\pgfsetrectcap%
\pgfsetmiterjoin%
\pgfsetlinewidth{0.803000pt}%
\definecolor{currentstroke}{rgb}{0.000000,0.000000,0.000000}%
\pgfsetstrokecolor{currentstroke}%
\pgfsetdash{}{0pt}%
\pgfpathmoveto{\pgfqpoint{3.299510in}{0.829139in}}%
\pgfpathlineto{\pgfqpoint{3.299510in}{2.750639in}}%
\pgfusepath{stroke}%
\end{pgfscope}%
\begin{pgfscope}%
\pgfsetrectcap%
\pgfsetmiterjoin%
\pgfsetlinewidth{0.803000pt}%
\definecolor{currentstroke}{rgb}{0.000000,0.000000,0.000000}%
\pgfsetstrokecolor{currentstroke}%
\pgfsetdash{}{0pt}%
\pgfpathmoveto{\pgfqpoint{0.553821in}{0.829139in}}%
\pgfpathlineto{\pgfqpoint{3.299510in}{0.829139in}}%
\pgfusepath{stroke}%
\end{pgfscope}%
\begin{pgfscope}%
\pgfsetrectcap%
\pgfsetmiterjoin%
\pgfsetlinewidth{0.803000pt}%
\definecolor{currentstroke}{rgb}{0.000000,0.000000,0.000000}%
\pgfsetstrokecolor{currentstroke}%
\pgfsetdash{}{0pt}%
\pgfpathmoveto{\pgfqpoint{0.553821in}{2.750639in}}%
\pgfpathlineto{\pgfqpoint{3.299510in}{2.750639in}}%
\pgfusepath{stroke}%
\end{pgfscope}%
\begin{pgfscope}%
\definecolor{textcolor}{rgb}{0.000000,0.000000,0.000000}%
\pgfsetstrokecolor{textcolor}%
\pgfsetfillcolor{textcolor}%
\pgftext[x=1.926666in,y=2.833972in,,base]{\color{textcolor}{\rmfamily\fontsize{8.000000}{9.600000}\selectfont\catcode`\^=\active\def^{\ifmmode\sp\else\^{}\fi}\catcode`\%=\active\def%{\%}(a) Fine-Scale Overlap}}%
\end{pgfscope}%
\begin{pgfscope}%
\pgfsetbuttcap%
\pgfsetmiterjoin%
\definecolor{currentfill}{rgb}{1.000000,1.000000,1.000000}%
\pgfsetfillcolor{currentfill}%
\pgfsetlinewidth{0.000000pt}%
\definecolor{currentstroke}{rgb}{0.000000,0.000000,0.000000}%
\pgfsetstrokecolor{currentstroke}%
\pgfsetstrokeopacity{0.000000}%
\pgfsetdash{}{0pt}%
\pgfpathmoveto{\pgfqpoint{4.178131in}{0.829139in}}%
\pgfpathlineto{\pgfqpoint{6.923821in}{0.829139in}}%
\pgfpathlineto{\pgfqpoint{6.923821in}{2.750639in}}%
\pgfpathlineto{\pgfqpoint{4.178131in}{2.750639in}}%
\pgfpathlineto{\pgfqpoint{4.178131in}{0.829139in}}%
\pgfpathclose%
\pgfusepath{fill}%
\end{pgfscope}%
\begin{pgfscope}%
\pgfpathrectangle{\pgfqpoint{4.178131in}{0.829139in}}{\pgfqpoint{2.745690in}{1.921500in}}%
\pgfusepath{clip}%
\pgfsetrectcap%
\pgfsetroundjoin%
\pgfsetlinewidth{0.803000pt}%
\definecolor{currentstroke}{rgb}{0.690196,0.690196,0.690196}%
\pgfsetstrokecolor{currentstroke}%
\pgfsetstrokeopacity{0.300000}%
\pgfsetdash{}{0pt}%
\pgfpathmoveto{\pgfqpoint{4.302269in}{0.829139in}}%
\pgfpathlineto{\pgfqpoint{4.302269in}{2.750639in}}%
\pgfusepath{stroke}%
\end{pgfscope}%
\begin{pgfscope}%
\pgfsetbuttcap%
\pgfsetroundjoin%
\definecolor{currentfill}{rgb}{0.000000,0.000000,0.000000}%
\pgfsetfillcolor{currentfill}%
\pgfsetlinewidth{0.803000pt}%
\definecolor{currentstroke}{rgb}{0.000000,0.000000,0.000000}%
\pgfsetstrokecolor{currentstroke}%
\pgfsetdash{}{0pt}%
\pgfsys@defobject{currentmarker}{\pgfqpoint{0.000000in}{-0.048611in}}{\pgfqpoint{0.000000in}{0.000000in}}{%
\pgfpathmoveto{\pgfqpoint{0.000000in}{0.000000in}}%
\pgfpathlineto{\pgfqpoint{0.000000in}{-0.048611in}}%
\pgfusepath{stroke,fill}%
}%
\begin{pgfscope}%
\pgfsys@transformshift{4.302269in}{0.829139in}%
\pgfsys@useobject{currentmarker}{}%
\end{pgfscope}%
\end{pgfscope}%
\begin{pgfscope}%
\definecolor{textcolor}{rgb}{0.000000,0.000000,0.000000}%
\pgfsetstrokecolor{textcolor}%
\pgfsetfillcolor{textcolor}%
\pgftext[x=4.302269in,y=0.731917in,,top]{\color{textcolor}{\rmfamily\fontsize{7.000000}{8.400000}\selectfont\catcode`\^=\active\def^{\ifmmode\sp\else\^{}\fi}\catcode`\%=\active\def%{\%}0.0}}%
\end{pgfscope}%
\begin{pgfscope}%
\pgfpathrectangle{\pgfqpoint{4.178131in}{0.829139in}}{\pgfqpoint{2.745690in}{1.921500in}}%
\pgfusepath{clip}%
\pgfsetrectcap%
\pgfsetroundjoin%
\pgfsetlinewidth{0.803000pt}%
\definecolor{currentstroke}{rgb}{0.690196,0.690196,0.690196}%
\pgfsetstrokecolor{currentstroke}%
\pgfsetstrokeopacity{0.300000}%
\pgfsetdash{}{0pt}%
\pgfpathmoveto{\pgfqpoint{4.718394in}{0.829139in}}%
\pgfpathlineto{\pgfqpoint{4.718394in}{2.750639in}}%
\pgfusepath{stroke}%
\end{pgfscope}%
\begin{pgfscope}%
\pgfsetbuttcap%
\pgfsetroundjoin%
\definecolor{currentfill}{rgb}{0.000000,0.000000,0.000000}%
\pgfsetfillcolor{currentfill}%
\pgfsetlinewidth{0.803000pt}%
\definecolor{currentstroke}{rgb}{0.000000,0.000000,0.000000}%
\pgfsetstrokecolor{currentstroke}%
\pgfsetdash{}{0pt}%
\pgfsys@defobject{currentmarker}{\pgfqpoint{0.000000in}{-0.048611in}}{\pgfqpoint{0.000000in}{0.000000in}}{%
\pgfpathmoveto{\pgfqpoint{0.000000in}{0.000000in}}%
\pgfpathlineto{\pgfqpoint{0.000000in}{-0.048611in}}%
\pgfusepath{stroke,fill}%
}%
\begin{pgfscope}%
\pgfsys@transformshift{4.718394in}{0.829139in}%
\pgfsys@useobject{currentmarker}{}%
\end{pgfscope}%
\end{pgfscope}%
\begin{pgfscope}%
\definecolor{textcolor}{rgb}{0.000000,0.000000,0.000000}%
\pgfsetstrokecolor{textcolor}%
\pgfsetfillcolor{textcolor}%
\pgftext[x=4.718394in,y=0.731917in,,top]{\color{textcolor}{\rmfamily\fontsize{7.000000}{8.400000}\selectfont\catcode`\^=\active\def^{\ifmmode\sp\else\^{}\fi}\catcode`\%=\active\def%{\%}0.5}}%
\end{pgfscope}%
\begin{pgfscope}%
\pgfpathrectangle{\pgfqpoint{4.178131in}{0.829139in}}{\pgfqpoint{2.745690in}{1.921500in}}%
\pgfusepath{clip}%
\pgfsetrectcap%
\pgfsetroundjoin%
\pgfsetlinewidth{0.803000pt}%
\definecolor{currentstroke}{rgb}{0.690196,0.690196,0.690196}%
\pgfsetstrokecolor{currentstroke}%
\pgfsetstrokeopacity{0.300000}%
\pgfsetdash{}{0pt}%
\pgfpathmoveto{\pgfqpoint{5.134518in}{0.829139in}}%
\pgfpathlineto{\pgfqpoint{5.134518in}{2.750639in}}%
\pgfusepath{stroke}%
\end{pgfscope}%
\begin{pgfscope}%
\pgfsetbuttcap%
\pgfsetroundjoin%
\definecolor{currentfill}{rgb}{0.000000,0.000000,0.000000}%
\pgfsetfillcolor{currentfill}%
\pgfsetlinewidth{0.803000pt}%
\definecolor{currentstroke}{rgb}{0.000000,0.000000,0.000000}%
\pgfsetstrokecolor{currentstroke}%
\pgfsetdash{}{0pt}%
\pgfsys@defobject{currentmarker}{\pgfqpoint{0.000000in}{-0.048611in}}{\pgfqpoint{0.000000in}{0.000000in}}{%
\pgfpathmoveto{\pgfqpoint{0.000000in}{0.000000in}}%
\pgfpathlineto{\pgfqpoint{0.000000in}{-0.048611in}}%
\pgfusepath{stroke,fill}%
}%
\begin{pgfscope}%
\pgfsys@transformshift{5.134518in}{0.829139in}%
\pgfsys@useobject{currentmarker}{}%
\end{pgfscope}%
\end{pgfscope}%
\begin{pgfscope}%
\definecolor{textcolor}{rgb}{0.000000,0.000000,0.000000}%
\pgfsetstrokecolor{textcolor}%
\pgfsetfillcolor{textcolor}%
\pgftext[x=5.134518in,y=0.731917in,,top]{\color{textcolor}{\rmfamily\fontsize{7.000000}{8.400000}\selectfont\catcode`\^=\active\def^{\ifmmode\sp\else\^{}\fi}\catcode`\%=\active\def%{\%}1.0}}%
\end{pgfscope}%
\begin{pgfscope}%
\pgfpathrectangle{\pgfqpoint{4.178131in}{0.829139in}}{\pgfqpoint{2.745690in}{1.921500in}}%
\pgfusepath{clip}%
\pgfsetrectcap%
\pgfsetroundjoin%
\pgfsetlinewidth{0.803000pt}%
\definecolor{currentstroke}{rgb}{0.690196,0.690196,0.690196}%
\pgfsetstrokecolor{currentstroke}%
\pgfsetstrokeopacity{0.300000}%
\pgfsetdash{}{0pt}%
\pgfpathmoveto{\pgfqpoint{5.550643in}{0.829139in}}%
\pgfpathlineto{\pgfqpoint{5.550643in}{2.750639in}}%
\pgfusepath{stroke}%
\end{pgfscope}%
\begin{pgfscope}%
\pgfsetbuttcap%
\pgfsetroundjoin%
\definecolor{currentfill}{rgb}{0.000000,0.000000,0.000000}%
\pgfsetfillcolor{currentfill}%
\pgfsetlinewidth{0.803000pt}%
\definecolor{currentstroke}{rgb}{0.000000,0.000000,0.000000}%
\pgfsetstrokecolor{currentstroke}%
\pgfsetdash{}{0pt}%
\pgfsys@defobject{currentmarker}{\pgfqpoint{0.000000in}{-0.048611in}}{\pgfqpoint{0.000000in}{0.000000in}}{%
\pgfpathmoveto{\pgfqpoint{0.000000in}{0.000000in}}%
\pgfpathlineto{\pgfqpoint{0.000000in}{-0.048611in}}%
\pgfusepath{stroke,fill}%
}%
\begin{pgfscope}%
\pgfsys@transformshift{5.550643in}{0.829139in}%
\pgfsys@useobject{currentmarker}{}%
\end{pgfscope}%
\end{pgfscope}%
\begin{pgfscope}%
\definecolor{textcolor}{rgb}{0.000000,0.000000,0.000000}%
\pgfsetstrokecolor{textcolor}%
\pgfsetfillcolor{textcolor}%
\pgftext[x=5.550643in,y=0.731917in,,top]{\color{textcolor}{\rmfamily\fontsize{7.000000}{8.400000}\selectfont\catcode`\^=\active\def^{\ifmmode\sp\else\^{}\fi}\catcode`\%=\active\def%{\%}1.5}}%
\end{pgfscope}%
\begin{pgfscope}%
\pgfpathrectangle{\pgfqpoint{4.178131in}{0.829139in}}{\pgfqpoint{2.745690in}{1.921500in}}%
\pgfusepath{clip}%
\pgfsetrectcap%
\pgfsetroundjoin%
\pgfsetlinewidth{0.803000pt}%
\definecolor{currentstroke}{rgb}{0.690196,0.690196,0.690196}%
\pgfsetstrokecolor{currentstroke}%
\pgfsetstrokeopacity{0.300000}%
\pgfsetdash{}{0pt}%
\pgfpathmoveto{\pgfqpoint{5.966768in}{0.829139in}}%
\pgfpathlineto{\pgfqpoint{5.966768in}{2.750639in}}%
\pgfusepath{stroke}%
\end{pgfscope}%
\begin{pgfscope}%
\pgfsetbuttcap%
\pgfsetroundjoin%
\definecolor{currentfill}{rgb}{0.000000,0.000000,0.000000}%
\pgfsetfillcolor{currentfill}%
\pgfsetlinewidth{0.803000pt}%
\definecolor{currentstroke}{rgb}{0.000000,0.000000,0.000000}%
\pgfsetstrokecolor{currentstroke}%
\pgfsetdash{}{0pt}%
\pgfsys@defobject{currentmarker}{\pgfqpoint{0.000000in}{-0.048611in}}{\pgfqpoint{0.000000in}{0.000000in}}{%
\pgfpathmoveto{\pgfqpoint{0.000000in}{0.000000in}}%
\pgfpathlineto{\pgfqpoint{0.000000in}{-0.048611in}}%
\pgfusepath{stroke,fill}%
}%
\begin{pgfscope}%
\pgfsys@transformshift{5.966768in}{0.829139in}%
\pgfsys@useobject{currentmarker}{}%
\end{pgfscope}%
\end{pgfscope}%
\begin{pgfscope}%
\definecolor{textcolor}{rgb}{0.000000,0.000000,0.000000}%
\pgfsetstrokecolor{textcolor}%
\pgfsetfillcolor{textcolor}%
\pgftext[x=5.966768in,y=0.731917in,,top]{\color{textcolor}{\rmfamily\fontsize{7.000000}{8.400000}\selectfont\catcode`\^=\active\def^{\ifmmode\sp\else\^{}\fi}\catcode`\%=\active\def%{\%}2.0}}%
\end{pgfscope}%
\begin{pgfscope}%
\pgfpathrectangle{\pgfqpoint{4.178131in}{0.829139in}}{\pgfqpoint{2.745690in}{1.921500in}}%
\pgfusepath{clip}%
\pgfsetrectcap%
\pgfsetroundjoin%
\pgfsetlinewidth{0.803000pt}%
\definecolor{currentstroke}{rgb}{0.690196,0.690196,0.690196}%
\pgfsetstrokecolor{currentstroke}%
\pgfsetstrokeopacity{0.300000}%
\pgfsetdash{}{0pt}%
\pgfpathmoveto{\pgfqpoint{6.382892in}{0.829139in}}%
\pgfpathlineto{\pgfqpoint{6.382892in}{2.750639in}}%
\pgfusepath{stroke}%
\end{pgfscope}%
\begin{pgfscope}%
\pgfsetbuttcap%
\pgfsetroundjoin%
\definecolor{currentfill}{rgb}{0.000000,0.000000,0.000000}%
\pgfsetfillcolor{currentfill}%
\pgfsetlinewidth{0.803000pt}%
\definecolor{currentstroke}{rgb}{0.000000,0.000000,0.000000}%
\pgfsetstrokecolor{currentstroke}%
\pgfsetdash{}{0pt}%
\pgfsys@defobject{currentmarker}{\pgfqpoint{0.000000in}{-0.048611in}}{\pgfqpoint{0.000000in}{0.000000in}}{%
\pgfpathmoveto{\pgfqpoint{0.000000in}{0.000000in}}%
\pgfpathlineto{\pgfqpoint{0.000000in}{-0.048611in}}%
\pgfusepath{stroke,fill}%
}%
\begin{pgfscope}%
\pgfsys@transformshift{6.382892in}{0.829139in}%
\pgfsys@useobject{currentmarker}{}%
\end{pgfscope}%
\end{pgfscope}%
\begin{pgfscope}%
\definecolor{textcolor}{rgb}{0.000000,0.000000,0.000000}%
\pgfsetstrokecolor{textcolor}%
\pgfsetfillcolor{textcolor}%
\pgftext[x=6.382892in,y=0.731917in,,top]{\color{textcolor}{\rmfamily\fontsize{7.000000}{8.400000}\selectfont\catcode`\^=\active\def^{\ifmmode\sp\else\^{}\fi}\catcode`\%=\active\def%{\%}2.5}}%
\end{pgfscope}%
\begin{pgfscope}%
\pgfpathrectangle{\pgfqpoint{4.178131in}{0.829139in}}{\pgfqpoint{2.745690in}{1.921500in}}%
\pgfusepath{clip}%
\pgfsetrectcap%
\pgfsetroundjoin%
\pgfsetlinewidth{0.803000pt}%
\definecolor{currentstroke}{rgb}{0.690196,0.690196,0.690196}%
\pgfsetstrokecolor{currentstroke}%
\pgfsetstrokeopacity{0.300000}%
\pgfsetdash{}{0pt}%
\pgfpathmoveto{\pgfqpoint{6.799017in}{0.829139in}}%
\pgfpathlineto{\pgfqpoint{6.799017in}{2.750639in}}%
\pgfusepath{stroke}%
\end{pgfscope}%
\begin{pgfscope}%
\pgfsetbuttcap%
\pgfsetroundjoin%
\definecolor{currentfill}{rgb}{0.000000,0.000000,0.000000}%
\pgfsetfillcolor{currentfill}%
\pgfsetlinewidth{0.803000pt}%
\definecolor{currentstroke}{rgb}{0.000000,0.000000,0.000000}%
\pgfsetstrokecolor{currentstroke}%
\pgfsetdash{}{0pt}%
\pgfsys@defobject{currentmarker}{\pgfqpoint{0.000000in}{-0.048611in}}{\pgfqpoint{0.000000in}{0.000000in}}{%
\pgfpathmoveto{\pgfqpoint{0.000000in}{0.000000in}}%
\pgfpathlineto{\pgfqpoint{0.000000in}{-0.048611in}}%
\pgfusepath{stroke,fill}%
}%
\begin{pgfscope}%
\pgfsys@transformshift{6.799017in}{0.829139in}%
\pgfsys@useobject{currentmarker}{}%
\end{pgfscope}%
\end{pgfscope}%
\begin{pgfscope}%
\definecolor{textcolor}{rgb}{0.000000,0.000000,0.000000}%
\pgfsetstrokecolor{textcolor}%
\pgfsetfillcolor{textcolor}%
\pgftext[x=6.799017in,y=0.731917in,,top]{\color{textcolor}{\rmfamily\fontsize{7.000000}{8.400000}\selectfont\catcode`\^=\active\def^{\ifmmode\sp\else\^{}\fi}\catcode`\%=\active\def%{\%}3.0}}%
\end{pgfscope}%
\begin{pgfscope}%
\definecolor{textcolor}{rgb}{0.000000,0.000000,0.000000}%
\pgfsetstrokecolor{textcolor}%
\pgfsetfillcolor{textcolor}%
\pgftext[x=5.550976in,y=0.589942in,,top]{\color{textcolor}{\rmfamily\fontsize{8.000000}{9.600000}\selectfont\catcode`\^=\active\def^{\ifmmode\sp\else\^{}\fi}\catcode`\%=\active\def%{\%}Diffusion Time $t$}}%
\end{pgfscope}%
\begin{pgfscope}%
\pgfpathrectangle{\pgfqpoint{4.178131in}{0.829139in}}{\pgfqpoint{2.745690in}{1.921500in}}%
\pgfusepath{clip}%
\pgfsetrectcap%
\pgfsetroundjoin%
\pgfsetlinewidth{0.803000pt}%
\definecolor{currentstroke}{rgb}{0.690196,0.690196,0.690196}%
\pgfsetstrokecolor{currentstroke}%
\pgfsetstrokeopacity{0.300000}%
\pgfsetdash{}{0pt}%
\pgfpathmoveto{\pgfqpoint{4.178131in}{0.849890in}}%
\pgfpathlineto{\pgfqpoint{6.923821in}{0.849890in}}%
\pgfusepath{stroke}%
\end{pgfscope}%
\begin{pgfscope}%
\pgfsetbuttcap%
\pgfsetroundjoin%
\definecolor{currentfill}{rgb}{0.000000,0.000000,0.000000}%
\pgfsetfillcolor{currentfill}%
\pgfsetlinewidth{0.803000pt}%
\definecolor{currentstroke}{rgb}{0.000000,0.000000,0.000000}%
\pgfsetstrokecolor{currentstroke}%
\pgfsetdash{}{0pt}%
\pgfsys@defobject{currentmarker}{\pgfqpoint{-0.048611in}{0.000000in}}{\pgfqpoint{-0.000000in}{0.000000in}}{%
\pgfpathmoveto{\pgfqpoint{-0.000000in}{0.000000in}}%
\pgfpathlineto{\pgfqpoint{-0.048611in}{0.000000in}}%
\pgfusepath{stroke,fill}%
}%
\begin{pgfscope}%
\pgfsys@transformshift{4.178131in}{0.849890in}%
\pgfsys@useobject{currentmarker}{}%
\end{pgfscope}%
\end{pgfscope}%
\begin{pgfscope}%
\definecolor{textcolor}{rgb}{0.000000,0.000000,0.000000}%
\pgfsetstrokecolor{textcolor}%
\pgfsetfillcolor{textcolor}%
\pgftext[x=3.840745in, y=0.811764in, left, base]{\color{textcolor}{\rmfamily\fontsize{7.000000}{8.400000}\selectfont\catcode`\^=\active\def^{\ifmmode\sp\else\^{}\fi}\catcode`\%=\active\def%{\%}$\mathdefault{10^{-5}}$}}%
\end{pgfscope}%
\begin{pgfscope}%
\pgfpathrectangle{\pgfqpoint{4.178131in}{0.829139in}}{\pgfqpoint{2.745690in}{1.921500in}}%
\pgfusepath{clip}%
\pgfsetrectcap%
\pgfsetroundjoin%
\pgfsetlinewidth{0.803000pt}%
\definecolor{currentstroke}{rgb}{0.690196,0.690196,0.690196}%
\pgfsetstrokecolor{currentstroke}%
\pgfsetstrokeopacity{0.300000}%
\pgfsetdash{}{0pt}%
\pgfpathmoveto{\pgfqpoint{4.178131in}{1.106587in}}%
\pgfpathlineto{\pgfqpoint{6.923821in}{1.106587in}}%
\pgfusepath{stroke}%
\end{pgfscope}%
\begin{pgfscope}%
\pgfsetbuttcap%
\pgfsetroundjoin%
\definecolor{currentfill}{rgb}{0.000000,0.000000,0.000000}%
\pgfsetfillcolor{currentfill}%
\pgfsetlinewidth{0.803000pt}%
\definecolor{currentstroke}{rgb}{0.000000,0.000000,0.000000}%
\pgfsetstrokecolor{currentstroke}%
\pgfsetdash{}{0pt}%
\pgfsys@defobject{currentmarker}{\pgfqpoint{-0.048611in}{0.000000in}}{\pgfqpoint{-0.000000in}{0.000000in}}{%
\pgfpathmoveto{\pgfqpoint{-0.000000in}{0.000000in}}%
\pgfpathlineto{\pgfqpoint{-0.048611in}{0.000000in}}%
\pgfusepath{stroke,fill}%
}%
\begin{pgfscope}%
\pgfsys@transformshift{4.178131in}{1.106587in}%
\pgfsys@useobject{currentmarker}{}%
\end{pgfscope}%
\end{pgfscope}%
\begin{pgfscope}%
\definecolor{textcolor}{rgb}{0.000000,0.000000,0.000000}%
\pgfsetstrokecolor{textcolor}%
\pgfsetfillcolor{textcolor}%
\pgftext[x=3.840745in, y=1.068462in, left, base]{\color{textcolor}{\rmfamily\fontsize{7.000000}{8.400000}\selectfont\catcode`\^=\active\def^{\ifmmode\sp\else\^{}\fi}\catcode`\%=\active\def%{\%}$\mathdefault{10^{-4}}$}}%
\end{pgfscope}%
\begin{pgfscope}%
\pgfpathrectangle{\pgfqpoint{4.178131in}{0.829139in}}{\pgfqpoint{2.745690in}{1.921500in}}%
\pgfusepath{clip}%
\pgfsetrectcap%
\pgfsetroundjoin%
\pgfsetlinewidth{0.803000pt}%
\definecolor{currentstroke}{rgb}{0.690196,0.690196,0.690196}%
\pgfsetstrokecolor{currentstroke}%
\pgfsetstrokeopacity{0.300000}%
\pgfsetdash{}{0pt}%
\pgfpathmoveto{\pgfqpoint{4.178131in}{1.363284in}}%
\pgfpathlineto{\pgfqpoint{6.923821in}{1.363284in}}%
\pgfusepath{stroke}%
\end{pgfscope}%
\begin{pgfscope}%
\pgfsetbuttcap%
\pgfsetroundjoin%
\definecolor{currentfill}{rgb}{0.000000,0.000000,0.000000}%
\pgfsetfillcolor{currentfill}%
\pgfsetlinewidth{0.803000pt}%
\definecolor{currentstroke}{rgb}{0.000000,0.000000,0.000000}%
\pgfsetstrokecolor{currentstroke}%
\pgfsetdash{}{0pt}%
\pgfsys@defobject{currentmarker}{\pgfqpoint{-0.048611in}{0.000000in}}{\pgfqpoint{-0.000000in}{0.000000in}}{%
\pgfpathmoveto{\pgfqpoint{-0.000000in}{0.000000in}}%
\pgfpathlineto{\pgfqpoint{-0.048611in}{0.000000in}}%
\pgfusepath{stroke,fill}%
}%
\begin{pgfscope}%
\pgfsys@transformshift{4.178131in}{1.363284in}%
\pgfsys@useobject{currentmarker}{}%
\end{pgfscope}%
\end{pgfscope}%
\begin{pgfscope}%
\definecolor{textcolor}{rgb}{0.000000,0.000000,0.000000}%
\pgfsetstrokecolor{textcolor}%
\pgfsetfillcolor{textcolor}%
\pgftext[x=3.840745in, y=1.325159in, left, base]{\color{textcolor}{\rmfamily\fontsize{7.000000}{8.400000}\selectfont\catcode`\^=\active\def^{\ifmmode\sp\else\^{}\fi}\catcode`\%=\active\def%{\%}$\mathdefault{10^{-3}}$}}%
\end{pgfscope}%
\begin{pgfscope}%
\pgfpathrectangle{\pgfqpoint{4.178131in}{0.829139in}}{\pgfqpoint{2.745690in}{1.921500in}}%
\pgfusepath{clip}%
\pgfsetrectcap%
\pgfsetroundjoin%
\pgfsetlinewidth{0.803000pt}%
\definecolor{currentstroke}{rgb}{0.690196,0.690196,0.690196}%
\pgfsetstrokecolor{currentstroke}%
\pgfsetstrokeopacity{0.300000}%
\pgfsetdash{}{0pt}%
\pgfpathmoveto{\pgfqpoint{4.178131in}{1.619982in}}%
\pgfpathlineto{\pgfqpoint{6.923821in}{1.619982in}}%
\pgfusepath{stroke}%
\end{pgfscope}%
\begin{pgfscope}%
\pgfsetbuttcap%
\pgfsetroundjoin%
\definecolor{currentfill}{rgb}{0.000000,0.000000,0.000000}%
\pgfsetfillcolor{currentfill}%
\pgfsetlinewidth{0.803000pt}%
\definecolor{currentstroke}{rgb}{0.000000,0.000000,0.000000}%
\pgfsetstrokecolor{currentstroke}%
\pgfsetdash{}{0pt}%
\pgfsys@defobject{currentmarker}{\pgfqpoint{-0.048611in}{0.000000in}}{\pgfqpoint{-0.000000in}{0.000000in}}{%
\pgfpathmoveto{\pgfqpoint{-0.000000in}{0.000000in}}%
\pgfpathlineto{\pgfqpoint{-0.048611in}{0.000000in}}%
\pgfusepath{stroke,fill}%
}%
\begin{pgfscope}%
\pgfsys@transformshift{4.178131in}{1.619982in}%
\pgfsys@useobject{currentmarker}{}%
\end{pgfscope}%
\end{pgfscope}%
\begin{pgfscope}%
\definecolor{textcolor}{rgb}{0.000000,0.000000,0.000000}%
\pgfsetstrokecolor{textcolor}%
\pgfsetfillcolor{textcolor}%
\pgftext[x=3.840745in, y=1.581856in, left, base]{\color{textcolor}{\rmfamily\fontsize{7.000000}{8.400000}\selectfont\catcode`\^=\active\def^{\ifmmode\sp\else\^{}\fi}\catcode`\%=\active\def%{\%}$\mathdefault{10^{-2}}$}}%
\end{pgfscope}%
\begin{pgfscope}%
\pgfpathrectangle{\pgfqpoint{4.178131in}{0.829139in}}{\pgfqpoint{2.745690in}{1.921500in}}%
\pgfusepath{clip}%
\pgfsetrectcap%
\pgfsetroundjoin%
\pgfsetlinewidth{0.803000pt}%
\definecolor{currentstroke}{rgb}{0.690196,0.690196,0.690196}%
\pgfsetstrokecolor{currentstroke}%
\pgfsetstrokeopacity{0.300000}%
\pgfsetdash{}{0pt}%
\pgfpathmoveto{\pgfqpoint{4.178131in}{1.876679in}}%
\pgfpathlineto{\pgfqpoint{6.923821in}{1.876679in}}%
\pgfusepath{stroke}%
\end{pgfscope}%
\begin{pgfscope}%
\pgfsetbuttcap%
\pgfsetroundjoin%
\definecolor{currentfill}{rgb}{0.000000,0.000000,0.000000}%
\pgfsetfillcolor{currentfill}%
\pgfsetlinewidth{0.803000pt}%
\definecolor{currentstroke}{rgb}{0.000000,0.000000,0.000000}%
\pgfsetstrokecolor{currentstroke}%
\pgfsetdash{}{0pt}%
\pgfsys@defobject{currentmarker}{\pgfqpoint{-0.048611in}{0.000000in}}{\pgfqpoint{-0.000000in}{0.000000in}}{%
\pgfpathmoveto{\pgfqpoint{-0.000000in}{0.000000in}}%
\pgfpathlineto{\pgfqpoint{-0.048611in}{0.000000in}}%
\pgfusepath{stroke,fill}%
}%
\begin{pgfscope}%
\pgfsys@transformshift{4.178131in}{1.876679in}%
\pgfsys@useobject{currentmarker}{}%
\end{pgfscope}%
\end{pgfscope}%
\begin{pgfscope}%
\definecolor{textcolor}{rgb}{0.000000,0.000000,0.000000}%
\pgfsetstrokecolor{textcolor}%
\pgfsetfillcolor{textcolor}%
\pgftext[x=3.840745in, y=1.838554in, left, base]{\color{textcolor}{\rmfamily\fontsize{7.000000}{8.400000}\selectfont\catcode`\^=\active\def^{\ifmmode\sp\else\^{}\fi}\catcode`\%=\active\def%{\%}$\mathdefault{10^{-1}}$}}%
\end{pgfscope}%
\begin{pgfscope}%
\pgfpathrectangle{\pgfqpoint{4.178131in}{0.829139in}}{\pgfqpoint{2.745690in}{1.921500in}}%
\pgfusepath{clip}%
\pgfsetrectcap%
\pgfsetroundjoin%
\pgfsetlinewidth{0.803000pt}%
\definecolor{currentstroke}{rgb}{0.690196,0.690196,0.690196}%
\pgfsetstrokecolor{currentstroke}%
\pgfsetstrokeopacity{0.300000}%
\pgfsetdash{}{0pt}%
\pgfpathmoveto{\pgfqpoint{4.178131in}{2.133376in}}%
\pgfpathlineto{\pgfqpoint{6.923821in}{2.133376in}}%
\pgfusepath{stroke}%
\end{pgfscope}%
\begin{pgfscope}%
\pgfsetbuttcap%
\pgfsetroundjoin%
\definecolor{currentfill}{rgb}{0.000000,0.000000,0.000000}%
\pgfsetfillcolor{currentfill}%
\pgfsetlinewidth{0.803000pt}%
\definecolor{currentstroke}{rgb}{0.000000,0.000000,0.000000}%
\pgfsetstrokecolor{currentstroke}%
\pgfsetdash{}{0pt}%
\pgfsys@defobject{currentmarker}{\pgfqpoint{-0.048611in}{0.000000in}}{\pgfqpoint{-0.000000in}{0.000000in}}{%
\pgfpathmoveto{\pgfqpoint{-0.000000in}{0.000000in}}%
\pgfpathlineto{\pgfqpoint{-0.048611in}{0.000000in}}%
\pgfusepath{stroke,fill}%
}%
\begin{pgfscope}%
\pgfsys@transformshift{4.178131in}{2.133376in}%
\pgfsys@useobject{currentmarker}{}%
\end{pgfscope}%
\end{pgfscope}%
\begin{pgfscope}%
\definecolor{textcolor}{rgb}{0.000000,0.000000,0.000000}%
\pgfsetstrokecolor{textcolor}%
\pgfsetfillcolor{textcolor}%
\pgftext[x=3.915977in, y=2.096023in, left, base]{\color{textcolor}{\rmfamily\fontsize{7.000000}{8.400000}\selectfont\catcode`\^=\active\def^{\ifmmode\sp\else\^{}\fi}\catcode`\%=\active\def%{\%}$\mathdefault{10^{0}}$}}%
\end{pgfscope}%
\begin{pgfscope}%
\pgfpathrectangle{\pgfqpoint{4.178131in}{0.829139in}}{\pgfqpoint{2.745690in}{1.921500in}}%
\pgfusepath{clip}%
\pgfsetrectcap%
\pgfsetroundjoin%
\pgfsetlinewidth{0.803000pt}%
\definecolor{currentstroke}{rgb}{0.690196,0.690196,0.690196}%
\pgfsetstrokecolor{currentstroke}%
\pgfsetstrokeopacity{0.300000}%
\pgfsetdash{}{0pt}%
\pgfpathmoveto{\pgfqpoint{4.178131in}{2.390074in}}%
\pgfpathlineto{\pgfqpoint{6.923821in}{2.390074in}}%
\pgfusepath{stroke}%
\end{pgfscope}%
\begin{pgfscope}%
\pgfsetbuttcap%
\pgfsetroundjoin%
\definecolor{currentfill}{rgb}{0.000000,0.000000,0.000000}%
\pgfsetfillcolor{currentfill}%
\pgfsetlinewidth{0.803000pt}%
\definecolor{currentstroke}{rgb}{0.000000,0.000000,0.000000}%
\pgfsetstrokecolor{currentstroke}%
\pgfsetdash{}{0pt}%
\pgfsys@defobject{currentmarker}{\pgfqpoint{-0.048611in}{0.000000in}}{\pgfqpoint{-0.000000in}{0.000000in}}{%
\pgfpathmoveto{\pgfqpoint{-0.000000in}{0.000000in}}%
\pgfpathlineto{\pgfqpoint{-0.048611in}{0.000000in}}%
\pgfusepath{stroke,fill}%
}%
\begin{pgfscope}%
\pgfsys@transformshift{4.178131in}{2.390074in}%
\pgfsys@useobject{currentmarker}{}%
\end{pgfscope}%
\end{pgfscope}%
\begin{pgfscope}%
\definecolor{textcolor}{rgb}{0.000000,0.000000,0.000000}%
\pgfsetstrokecolor{textcolor}%
\pgfsetfillcolor{textcolor}%
\pgftext[x=3.915977in, y=2.352720in, left, base]{\color{textcolor}{\rmfamily\fontsize{7.000000}{8.400000}\selectfont\catcode`\^=\active\def^{\ifmmode\sp\else\^{}\fi}\catcode`\%=\active\def%{\%}$\mathdefault{10^{1}}$}}%
\end{pgfscope}%
\begin{pgfscope}%
\pgfpathrectangle{\pgfqpoint{4.178131in}{0.829139in}}{\pgfqpoint{2.745690in}{1.921500in}}%
\pgfusepath{clip}%
\pgfsetrectcap%
\pgfsetroundjoin%
\pgfsetlinewidth{0.803000pt}%
\definecolor{currentstroke}{rgb}{0.690196,0.690196,0.690196}%
\pgfsetstrokecolor{currentstroke}%
\pgfsetstrokeopacity{0.300000}%
\pgfsetdash{}{0pt}%
\pgfpathmoveto{\pgfqpoint{4.178131in}{2.646771in}}%
\pgfpathlineto{\pgfqpoint{6.923821in}{2.646771in}}%
\pgfusepath{stroke}%
\end{pgfscope}%
\begin{pgfscope}%
\pgfsetbuttcap%
\pgfsetroundjoin%
\definecolor{currentfill}{rgb}{0.000000,0.000000,0.000000}%
\pgfsetfillcolor{currentfill}%
\pgfsetlinewidth{0.803000pt}%
\definecolor{currentstroke}{rgb}{0.000000,0.000000,0.000000}%
\pgfsetstrokecolor{currentstroke}%
\pgfsetdash{}{0pt}%
\pgfsys@defobject{currentmarker}{\pgfqpoint{-0.048611in}{0.000000in}}{\pgfqpoint{-0.000000in}{0.000000in}}{%
\pgfpathmoveto{\pgfqpoint{-0.000000in}{0.000000in}}%
\pgfpathlineto{\pgfqpoint{-0.048611in}{0.000000in}}%
\pgfusepath{stroke,fill}%
}%
\begin{pgfscope}%
\pgfsys@transformshift{4.178131in}{2.646771in}%
\pgfsys@useobject{currentmarker}{}%
\end{pgfscope}%
\end{pgfscope}%
\begin{pgfscope}%
\definecolor{textcolor}{rgb}{0.000000,0.000000,0.000000}%
\pgfsetstrokecolor{textcolor}%
\pgfsetfillcolor{textcolor}%
\pgftext[x=3.915977in, y=2.609418in, left, base]{\color{textcolor}{\rmfamily\fontsize{7.000000}{8.400000}\selectfont\catcode`\^=\active\def^{\ifmmode\sp\else\^{}\fi}\catcode`\%=\active\def%{\%}$\mathdefault{10^{2}}$}}%
\end{pgfscope}%
\begin{pgfscope}%
\pgfpathrectangle{\pgfqpoint{4.178131in}{0.829139in}}{\pgfqpoint{2.745690in}{1.921500in}}%
\pgfusepath{clip}%
\pgfsetrectcap%
\pgfsetroundjoin%
\pgfsetlinewidth{0.803000pt}%
\definecolor{currentstroke}{rgb}{0.690196,0.690196,0.690196}%
\pgfsetstrokecolor{currentstroke}%
\pgfsetstrokeopacity{0.300000}%
\pgfsetdash{}{0pt}%
\pgfpathmoveto{\pgfqpoint{4.178131in}{0.838144in}}%
\pgfpathlineto{\pgfqpoint{6.923821in}{0.838144in}}%
\pgfusepath{stroke}%
\end{pgfscope}%
\begin{pgfscope}%
\pgfsetbuttcap%
\pgfsetroundjoin%
\definecolor{currentfill}{rgb}{0.000000,0.000000,0.000000}%
\pgfsetfillcolor{currentfill}%
\pgfsetlinewidth{0.602250pt}%
\definecolor{currentstroke}{rgb}{0.000000,0.000000,0.000000}%
\pgfsetstrokecolor{currentstroke}%
\pgfsetdash{}{0pt}%
\pgfsys@defobject{currentmarker}{\pgfqpoint{-0.027778in}{0.000000in}}{\pgfqpoint{-0.000000in}{0.000000in}}{%
\pgfpathmoveto{\pgfqpoint{-0.000000in}{0.000000in}}%
\pgfpathlineto{\pgfqpoint{-0.027778in}{0.000000in}}%
\pgfusepath{stroke,fill}%
}%
\begin{pgfscope}%
\pgfsys@transformshift{4.178131in}{0.838144in}%
\pgfsys@useobject{currentmarker}{}%
\end{pgfscope}%
\end{pgfscope}%
\begin{pgfscope}%
\pgfpathrectangle{\pgfqpoint{4.178131in}{0.829139in}}{\pgfqpoint{2.745690in}{1.921500in}}%
\pgfusepath{clip}%
\pgfsetrectcap%
\pgfsetroundjoin%
\pgfsetlinewidth{0.803000pt}%
\definecolor{currentstroke}{rgb}{0.690196,0.690196,0.690196}%
\pgfsetstrokecolor{currentstroke}%
\pgfsetstrokeopacity{0.300000}%
\pgfsetdash{}{0pt}%
\pgfpathmoveto{\pgfqpoint{4.178131in}{0.927163in}}%
\pgfpathlineto{\pgfqpoint{6.923821in}{0.927163in}}%
\pgfusepath{stroke}%
\end{pgfscope}%
\begin{pgfscope}%
\pgfsetbuttcap%
\pgfsetroundjoin%
\definecolor{currentfill}{rgb}{0.000000,0.000000,0.000000}%
\pgfsetfillcolor{currentfill}%
\pgfsetlinewidth{0.602250pt}%
\definecolor{currentstroke}{rgb}{0.000000,0.000000,0.000000}%
\pgfsetstrokecolor{currentstroke}%
\pgfsetdash{}{0pt}%
\pgfsys@defobject{currentmarker}{\pgfqpoint{-0.027778in}{0.000000in}}{\pgfqpoint{-0.000000in}{0.000000in}}{%
\pgfpathmoveto{\pgfqpoint{-0.000000in}{0.000000in}}%
\pgfpathlineto{\pgfqpoint{-0.027778in}{0.000000in}}%
\pgfusepath{stroke,fill}%
}%
\begin{pgfscope}%
\pgfsys@transformshift{4.178131in}{0.927163in}%
\pgfsys@useobject{currentmarker}{}%
\end{pgfscope}%
\end{pgfscope}%
\begin{pgfscope}%
\pgfpathrectangle{\pgfqpoint{4.178131in}{0.829139in}}{\pgfqpoint{2.745690in}{1.921500in}}%
\pgfusepath{clip}%
\pgfsetrectcap%
\pgfsetroundjoin%
\pgfsetlinewidth{0.803000pt}%
\definecolor{currentstroke}{rgb}{0.690196,0.690196,0.690196}%
\pgfsetstrokecolor{currentstroke}%
\pgfsetstrokeopacity{0.300000}%
\pgfsetdash{}{0pt}%
\pgfpathmoveto{\pgfqpoint{4.178131in}{0.972365in}}%
\pgfpathlineto{\pgfqpoint{6.923821in}{0.972365in}}%
\pgfusepath{stroke}%
\end{pgfscope}%
\begin{pgfscope}%
\pgfsetbuttcap%
\pgfsetroundjoin%
\definecolor{currentfill}{rgb}{0.000000,0.000000,0.000000}%
\pgfsetfillcolor{currentfill}%
\pgfsetlinewidth{0.602250pt}%
\definecolor{currentstroke}{rgb}{0.000000,0.000000,0.000000}%
\pgfsetstrokecolor{currentstroke}%
\pgfsetdash{}{0pt}%
\pgfsys@defobject{currentmarker}{\pgfqpoint{-0.027778in}{0.000000in}}{\pgfqpoint{-0.000000in}{0.000000in}}{%
\pgfpathmoveto{\pgfqpoint{-0.000000in}{0.000000in}}%
\pgfpathlineto{\pgfqpoint{-0.027778in}{0.000000in}}%
\pgfusepath{stroke,fill}%
}%
\begin{pgfscope}%
\pgfsys@transformshift{4.178131in}{0.972365in}%
\pgfsys@useobject{currentmarker}{}%
\end{pgfscope}%
\end{pgfscope}%
\begin{pgfscope}%
\pgfpathrectangle{\pgfqpoint{4.178131in}{0.829139in}}{\pgfqpoint{2.745690in}{1.921500in}}%
\pgfusepath{clip}%
\pgfsetrectcap%
\pgfsetroundjoin%
\pgfsetlinewidth{0.803000pt}%
\definecolor{currentstroke}{rgb}{0.690196,0.690196,0.690196}%
\pgfsetstrokecolor{currentstroke}%
\pgfsetstrokeopacity{0.300000}%
\pgfsetdash{}{0pt}%
\pgfpathmoveto{\pgfqpoint{4.178131in}{1.004437in}}%
\pgfpathlineto{\pgfqpoint{6.923821in}{1.004437in}}%
\pgfusepath{stroke}%
\end{pgfscope}%
\begin{pgfscope}%
\pgfsetbuttcap%
\pgfsetroundjoin%
\definecolor{currentfill}{rgb}{0.000000,0.000000,0.000000}%
\pgfsetfillcolor{currentfill}%
\pgfsetlinewidth{0.602250pt}%
\definecolor{currentstroke}{rgb}{0.000000,0.000000,0.000000}%
\pgfsetstrokecolor{currentstroke}%
\pgfsetdash{}{0pt}%
\pgfsys@defobject{currentmarker}{\pgfqpoint{-0.027778in}{0.000000in}}{\pgfqpoint{-0.000000in}{0.000000in}}{%
\pgfpathmoveto{\pgfqpoint{-0.000000in}{0.000000in}}%
\pgfpathlineto{\pgfqpoint{-0.027778in}{0.000000in}}%
\pgfusepath{stroke,fill}%
}%
\begin{pgfscope}%
\pgfsys@transformshift{4.178131in}{1.004437in}%
\pgfsys@useobject{currentmarker}{}%
\end{pgfscope}%
\end{pgfscope}%
\begin{pgfscope}%
\pgfpathrectangle{\pgfqpoint{4.178131in}{0.829139in}}{\pgfqpoint{2.745690in}{1.921500in}}%
\pgfusepath{clip}%
\pgfsetrectcap%
\pgfsetroundjoin%
\pgfsetlinewidth{0.803000pt}%
\definecolor{currentstroke}{rgb}{0.690196,0.690196,0.690196}%
\pgfsetstrokecolor{currentstroke}%
\pgfsetstrokeopacity{0.300000}%
\pgfsetdash{}{0pt}%
\pgfpathmoveto{\pgfqpoint{4.178131in}{1.029313in}}%
\pgfpathlineto{\pgfqpoint{6.923821in}{1.029313in}}%
\pgfusepath{stroke}%
\end{pgfscope}%
\begin{pgfscope}%
\pgfsetbuttcap%
\pgfsetroundjoin%
\definecolor{currentfill}{rgb}{0.000000,0.000000,0.000000}%
\pgfsetfillcolor{currentfill}%
\pgfsetlinewidth{0.602250pt}%
\definecolor{currentstroke}{rgb}{0.000000,0.000000,0.000000}%
\pgfsetstrokecolor{currentstroke}%
\pgfsetdash{}{0pt}%
\pgfsys@defobject{currentmarker}{\pgfqpoint{-0.027778in}{0.000000in}}{\pgfqpoint{-0.000000in}{0.000000in}}{%
\pgfpathmoveto{\pgfqpoint{-0.000000in}{0.000000in}}%
\pgfpathlineto{\pgfqpoint{-0.027778in}{0.000000in}}%
\pgfusepath{stroke,fill}%
}%
\begin{pgfscope}%
\pgfsys@transformshift{4.178131in}{1.029313in}%
\pgfsys@useobject{currentmarker}{}%
\end{pgfscope}%
\end{pgfscope}%
\begin{pgfscope}%
\pgfpathrectangle{\pgfqpoint{4.178131in}{0.829139in}}{\pgfqpoint{2.745690in}{1.921500in}}%
\pgfusepath{clip}%
\pgfsetrectcap%
\pgfsetroundjoin%
\pgfsetlinewidth{0.803000pt}%
\definecolor{currentstroke}{rgb}{0.690196,0.690196,0.690196}%
\pgfsetstrokecolor{currentstroke}%
\pgfsetstrokeopacity{0.300000}%
\pgfsetdash{}{0pt}%
\pgfpathmoveto{\pgfqpoint{4.178131in}{1.049639in}}%
\pgfpathlineto{\pgfqpoint{6.923821in}{1.049639in}}%
\pgfusepath{stroke}%
\end{pgfscope}%
\begin{pgfscope}%
\pgfsetbuttcap%
\pgfsetroundjoin%
\definecolor{currentfill}{rgb}{0.000000,0.000000,0.000000}%
\pgfsetfillcolor{currentfill}%
\pgfsetlinewidth{0.602250pt}%
\definecolor{currentstroke}{rgb}{0.000000,0.000000,0.000000}%
\pgfsetstrokecolor{currentstroke}%
\pgfsetdash{}{0pt}%
\pgfsys@defobject{currentmarker}{\pgfqpoint{-0.027778in}{0.000000in}}{\pgfqpoint{-0.000000in}{0.000000in}}{%
\pgfpathmoveto{\pgfqpoint{-0.000000in}{0.000000in}}%
\pgfpathlineto{\pgfqpoint{-0.027778in}{0.000000in}}%
\pgfusepath{stroke,fill}%
}%
\begin{pgfscope}%
\pgfsys@transformshift{4.178131in}{1.049639in}%
\pgfsys@useobject{currentmarker}{}%
\end{pgfscope}%
\end{pgfscope}%
\begin{pgfscope}%
\pgfpathrectangle{\pgfqpoint{4.178131in}{0.829139in}}{\pgfqpoint{2.745690in}{1.921500in}}%
\pgfusepath{clip}%
\pgfsetrectcap%
\pgfsetroundjoin%
\pgfsetlinewidth{0.803000pt}%
\definecolor{currentstroke}{rgb}{0.690196,0.690196,0.690196}%
\pgfsetstrokecolor{currentstroke}%
\pgfsetstrokeopacity{0.300000}%
\pgfsetdash{}{0pt}%
\pgfpathmoveto{\pgfqpoint{4.178131in}{1.066824in}}%
\pgfpathlineto{\pgfqpoint{6.923821in}{1.066824in}}%
\pgfusepath{stroke}%
\end{pgfscope}%
\begin{pgfscope}%
\pgfsetbuttcap%
\pgfsetroundjoin%
\definecolor{currentfill}{rgb}{0.000000,0.000000,0.000000}%
\pgfsetfillcolor{currentfill}%
\pgfsetlinewidth{0.602250pt}%
\definecolor{currentstroke}{rgb}{0.000000,0.000000,0.000000}%
\pgfsetstrokecolor{currentstroke}%
\pgfsetdash{}{0pt}%
\pgfsys@defobject{currentmarker}{\pgfqpoint{-0.027778in}{0.000000in}}{\pgfqpoint{-0.000000in}{0.000000in}}{%
\pgfpathmoveto{\pgfqpoint{-0.000000in}{0.000000in}}%
\pgfpathlineto{\pgfqpoint{-0.027778in}{0.000000in}}%
\pgfusepath{stroke,fill}%
}%
\begin{pgfscope}%
\pgfsys@transformshift{4.178131in}{1.066824in}%
\pgfsys@useobject{currentmarker}{}%
\end{pgfscope}%
\end{pgfscope}%
\begin{pgfscope}%
\pgfpathrectangle{\pgfqpoint{4.178131in}{0.829139in}}{\pgfqpoint{2.745690in}{1.921500in}}%
\pgfusepath{clip}%
\pgfsetrectcap%
\pgfsetroundjoin%
\pgfsetlinewidth{0.803000pt}%
\definecolor{currentstroke}{rgb}{0.690196,0.690196,0.690196}%
\pgfsetstrokecolor{currentstroke}%
\pgfsetstrokeopacity{0.300000}%
\pgfsetdash{}{0pt}%
\pgfpathmoveto{\pgfqpoint{4.178131in}{1.081711in}}%
\pgfpathlineto{\pgfqpoint{6.923821in}{1.081711in}}%
\pgfusepath{stroke}%
\end{pgfscope}%
\begin{pgfscope}%
\pgfsetbuttcap%
\pgfsetroundjoin%
\definecolor{currentfill}{rgb}{0.000000,0.000000,0.000000}%
\pgfsetfillcolor{currentfill}%
\pgfsetlinewidth{0.602250pt}%
\definecolor{currentstroke}{rgb}{0.000000,0.000000,0.000000}%
\pgfsetstrokecolor{currentstroke}%
\pgfsetdash{}{0pt}%
\pgfsys@defobject{currentmarker}{\pgfqpoint{-0.027778in}{0.000000in}}{\pgfqpoint{-0.000000in}{0.000000in}}{%
\pgfpathmoveto{\pgfqpoint{-0.000000in}{0.000000in}}%
\pgfpathlineto{\pgfqpoint{-0.027778in}{0.000000in}}%
\pgfusepath{stroke,fill}%
}%
\begin{pgfscope}%
\pgfsys@transformshift{4.178131in}{1.081711in}%
\pgfsys@useobject{currentmarker}{}%
\end{pgfscope}%
\end{pgfscope}%
\begin{pgfscope}%
\pgfpathrectangle{\pgfqpoint{4.178131in}{0.829139in}}{\pgfqpoint{2.745690in}{1.921500in}}%
\pgfusepath{clip}%
\pgfsetrectcap%
\pgfsetroundjoin%
\pgfsetlinewidth{0.803000pt}%
\definecolor{currentstroke}{rgb}{0.690196,0.690196,0.690196}%
\pgfsetstrokecolor{currentstroke}%
\pgfsetstrokeopacity{0.300000}%
\pgfsetdash{}{0pt}%
\pgfpathmoveto{\pgfqpoint{4.178131in}{1.094841in}}%
\pgfpathlineto{\pgfqpoint{6.923821in}{1.094841in}}%
\pgfusepath{stroke}%
\end{pgfscope}%
\begin{pgfscope}%
\pgfsetbuttcap%
\pgfsetroundjoin%
\definecolor{currentfill}{rgb}{0.000000,0.000000,0.000000}%
\pgfsetfillcolor{currentfill}%
\pgfsetlinewidth{0.602250pt}%
\definecolor{currentstroke}{rgb}{0.000000,0.000000,0.000000}%
\pgfsetstrokecolor{currentstroke}%
\pgfsetdash{}{0pt}%
\pgfsys@defobject{currentmarker}{\pgfqpoint{-0.027778in}{0.000000in}}{\pgfqpoint{-0.000000in}{0.000000in}}{%
\pgfpathmoveto{\pgfqpoint{-0.000000in}{0.000000in}}%
\pgfpathlineto{\pgfqpoint{-0.027778in}{0.000000in}}%
\pgfusepath{stroke,fill}%
}%
\begin{pgfscope}%
\pgfsys@transformshift{4.178131in}{1.094841in}%
\pgfsys@useobject{currentmarker}{}%
\end{pgfscope}%
\end{pgfscope}%
\begin{pgfscope}%
\pgfpathrectangle{\pgfqpoint{4.178131in}{0.829139in}}{\pgfqpoint{2.745690in}{1.921500in}}%
\pgfusepath{clip}%
\pgfsetrectcap%
\pgfsetroundjoin%
\pgfsetlinewidth{0.803000pt}%
\definecolor{currentstroke}{rgb}{0.690196,0.690196,0.690196}%
\pgfsetstrokecolor{currentstroke}%
\pgfsetstrokeopacity{0.300000}%
\pgfsetdash{}{0pt}%
\pgfpathmoveto{\pgfqpoint{4.178131in}{1.183861in}}%
\pgfpathlineto{\pgfqpoint{6.923821in}{1.183861in}}%
\pgfusepath{stroke}%
\end{pgfscope}%
\begin{pgfscope}%
\pgfsetbuttcap%
\pgfsetroundjoin%
\definecolor{currentfill}{rgb}{0.000000,0.000000,0.000000}%
\pgfsetfillcolor{currentfill}%
\pgfsetlinewidth{0.602250pt}%
\definecolor{currentstroke}{rgb}{0.000000,0.000000,0.000000}%
\pgfsetstrokecolor{currentstroke}%
\pgfsetdash{}{0pt}%
\pgfsys@defobject{currentmarker}{\pgfqpoint{-0.027778in}{0.000000in}}{\pgfqpoint{-0.000000in}{0.000000in}}{%
\pgfpathmoveto{\pgfqpoint{-0.000000in}{0.000000in}}%
\pgfpathlineto{\pgfqpoint{-0.027778in}{0.000000in}}%
\pgfusepath{stroke,fill}%
}%
\begin{pgfscope}%
\pgfsys@transformshift{4.178131in}{1.183861in}%
\pgfsys@useobject{currentmarker}{}%
\end{pgfscope}%
\end{pgfscope}%
\begin{pgfscope}%
\pgfpathrectangle{\pgfqpoint{4.178131in}{0.829139in}}{\pgfqpoint{2.745690in}{1.921500in}}%
\pgfusepath{clip}%
\pgfsetrectcap%
\pgfsetroundjoin%
\pgfsetlinewidth{0.803000pt}%
\definecolor{currentstroke}{rgb}{0.690196,0.690196,0.690196}%
\pgfsetstrokecolor{currentstroke}%
\pgfsetstrokeopacity{0.300000}%
\pgfsetdash{}{0pt}%
\pgfpathmoveto{\pgfqpoint{4.178131in}{1.229063in}}%
\pgfpathlineto{\pgfqpoint{6.923821in}{1.229063in}}%
\pgfusepath{stroke}%
\end{pgfscope}%
\begin{pgfscope}%
\pgfsetbuttcap%
\pgfsetroundjoin%
\definecolor{currentfill}{rgb}{0.000000,0.000000,0.000000}%
\pgfsetfillcolor{currentfill}%
\pgfsetlinewidth{0.602250pt}%
\definecolor{currentstroke}{rgb}{0.000000,0.000000,0.000000}%
\pgfsetstrokecolor{currentstroke}%
\pgfsetdash{}{0pt}%
\pgfsys@defobject{currentmarker}{\pgfqpoint{-0.027778in}{0.000000in}}{\pgfqpoint{-0.000000in}{0.000000in}}{%
\pgfpathmoveto{\pgfqpoint{-0.000000in}{0.000000in}}%
\pgfpathlineto{\pgfqpoint{-0.027778in}{0.000000in}}%
\pgfusepath{stroke,fill}%
}%
\begin{pgfscope}%
\pgfsys@transformshift{4.178131in}{1.229063in}%
\pgfsys@useobject{currentmarker}{}%
\end{pgfscope}%
\end{pgfscope}%
\begin{pgfscope}%
\pgfpathrectangle{\pgfqpoint{4.178131in}{0.829139in}}{\pgfqpoint{2.745690in}{1.921500in}}%
\pgfusepath{clip}%
\pgfsetrectcap%
\pgfsetroundjoin%
\pgfsetlinewidth{0.803000pt}%
\definecolor{currentstroke}{rgb}{0.690196,0.690196,0.690196}%
\pgfsetstrokecolor{currentstroke}%
\pgfsetstrokeopacity{0.300000}%
\pgfsetdash{}{0pt}%
\pgfpathmoveto{\pgfqpoint{4.178131in}{1.261134in}}%
\pgfpathlineto{\pgfqpoint{6.923821in}{1.261134in}}%
\pgfusepath{stroke}%
\end{pgfscope}%
\begin{pgfscope}%
\pgfsetbuttcap%
\pgfsetroundjoin%
\definecolor{currentfill}{rgb}{0.000000,0.000000,0.000000}%
\pgfsetfillcolor{currentfill}%
\pgfsetlinewidth{0.602250pt}%
\definecolor{currentstroke}{rgb}{0.000000,0.000000,0.000000}%
\pgfsetstrokecolor{currentstroke}%
\pgfsetdash{}{0pt}%
\pgfsys@defobject{currentmarker}{\pgfqpoint{-0.027778in}{0.000000in}}{\pgfqpoint{-0.000000in}{0.000000in}}{%
\pgfpathmoveto{\pgfqpoint{-0.000000in}{0.000000in}}%
\pgfpathlineto{\pgfqpoint{-0.027778in}{0.000000in}}%
\pgfusepath{stroke,fill}%
}%
\begin{pgfscope}%
\pgfsys@transformshift{4.178131in}{1.261134in}%
\pgfsys@useobject{currentmarker}{}%
\end{pgfscope}%
\end{pgfscope}%
\begin{pgfscope}%
\pgfpathrectangle{\pgfqpoint{4.178131in}{0.829139in}}{\pgfqpoint{2.745690in}{1.921500in}}%
\pgfusepath{clip}%
\pgfsetrectcap%
\pgfsetroundjoin%
\pgfsetlinewidth{0.803000pt}%
\definecolor{currentstroke}{rgb}{0.690196,0.690196,0.690196}%
\pgfsetstrokecolor{currentstroke}%
\pgfsetstrokeopacity{0.300000}%
\pgfsetdash{}{0pt}%
\pgfpathmoveto{\pgfqpoint{4.178131in}{1.286011in}}%
\pgfpathlineto{\pgfqpoint{6.923821in}{1.286011in}}%
\pgfusepath{stroke}%
\end{pgfscope}%
\begin{pgfscope}%
\pgfsetbuttcap%
\pgfsetroundjoin%
\definecolor{currentfill}{rgb}{0.000000,0.000000,0.000000}%
\pgfsetfillcolor{currentfill}%
\pgfsetlinewidth{0.602250pt}%
\definecolor{currentstroke}{rgb}{0.000000,0.000000,0.000000}%
\pgfsetstrokecolor{currentstroke}%
\pgfsetdash{}{0pt}%
\pgfsys@defobject{currentmarker}{\pgfqpoint{-0.027778in}{0.000000in}}{\pgfqpoint{-0.000000in}{0.000000in}}{%
\pgfpathmoveto{\pgfqpoint{-0.000000in}{0.000000in}}%
\pgfpathlineto{\pgfqpoint{-0.027778in}{0.000000in}}%
\pgfusepath{stroke,fill}%
}%
\begin{pgfscope}%
\pgfsys@transformshift{4.178131in}{1.286011in}%
\pgfsys@useobject{currentmarker}{}%
\end{pgfscope}%
\end{pgfscope}%
\begin{pgfscope}%
\pgfpathrectangle{\pgfqpoint{4.178131in}{0.829139in}}{\pgfqpoint{2.745690in}{1.921500in}}%
\pgfusepath{clip}%
\pgfsetrectcap%
\pgfsetroundjoin%
\pgfsetlinewidth{0.803000pt}%
\definecolor{currentstroke}{rgb}{0.690196,0.690196,0.690196}%
\pgfsetstrokecolor{currentstroke}%
\pgfsetstrokeopacity{0.300000}%
\pgfsetdash{}{0pt}%
\pgfpathmoveto{\pgfqpoint{4.178131in}{1.306336in}}%
\pgfpathlineto{\pgfqpoint{6.923821in}{1.306336in}}%
\pgfusepath{stroke}%
\end{pgfscope}%
\begin{pgfscope}%
\pgfsetbuttcap%
\pgfsetroundjoin%
\definecolor{currentfill}{rgb}{0.000000,0.000000,0.000000}%
\pgfsetfillcolor{currentfill}%
\pgfsetlinewidth{0.602250pt}%
\definecolor{currentstroke}{rgb}{0.000000,0.000000,0.000000}%
\pgfsetstrokecolor{currentstroke}%
\pgfsetdash{}{0pt}%
\pgfsys@defobject{currentmarker}{\pgfqpoint{-0.027778in}{0.000000in}}{\pgfqpoint{-0.000000in}{0.000000in}}{%
\pgfpathmoveto{\pgfqpoint{-0.000000in}{0.000000in}}%
\pgfpathlineto{\pgfqpoint{-0.027778in}{0.000000in}}%
\pgfusepath{stroke,fill}%
}%
\begin{pgfscope}%
\pgfsys@transformshift{4.178131in}{1.306336in}%
\pgfsys@useobject{currentmarker}{}%
\end{pgfscope}%
\end{pgfscope}%
\begin{pgfscope}%
\pgfpathrectangle{\pgfqpoint{4.178131in}{0.829139in}}{\pgfqpoint{2.745690in}{1.921500in}}%
\pgfusepath{clip}%
\pgfsetrectcap%
\pgfsetroundjoin%
\pgfsetlinewidth{0.803000pt}%
\definecolor{currentstroke}{rgb}{0.690196,0.690196,0.690196}%
\pgfsetstrokecolor{currentstroke}%
\pgfsetstrokeopacity{0.300000}%
\pgfsetdash{}{0pt}%
\pgfpathmoveto{\pgfqpoint{4.178131in}{1.323521in}}%
\pgfpathlineto{\pgfqpoint{6.923821in}{1.323521in}}%
\pgfusepath{stroke}%
\end{pgfscope}%
\begin{pgfscope}%
\pgfsetbuttcap%
\pgfsetroundjoin%
\definecolor{currentfill}{rgb}{0.000000,0.000000,0.000000}%
\pgfsetfillcolor{currentfill}%
\pgfsetlinewidth{0.602250pt}%
\definecolor{currentstroke}{rgb}{0.000000,0.000000,0.000000}%
\pgfsetstrokecolor{currentstroke}%
\pgfsetdash{}{0pt}%
\pgfsys@defobject{currentmarker}{\pgfqpoint{-0.027778in}{0.000000in}}{\pgfqpoint{-0.000000in}{0.000000in}}{%
\pgfpathmoveto{\pgfqpoint{-0.000000in}{0.000000in}}%
\pgfpathlineto{\pgfqpoint{-0.027778in}{0.000000in}}%
\pgfusepath{stroke,fill}%
}%
\begin{pgfscope}%
\pgfsys@transformshift{4.178131in}{1.323521in}%
\pgfsys@useobject{currentmarker}{}%
\end{pgfscope}%
\end{pgfscope}%
\begin{pgfscope}%
\pgfpathrectangle{\pgfqpoint{4.178131in}{0.829139in}}{\pgfqpoint{2.745690in}{1.921500in}}%
\pgfusepath{clip}%
\pgfsetrectcap%
\pgfsetroundjoin%
\pgfsetlinewidth{0.803000pt}%
\definecolor{currentstroke}{rgb}{0.690196,0.690196,0.690196}%
\pgfsetstrokecolor{currentstroke}%
\pgfsetstrokeopacity{0.300000}%
\pgfsetdash{}{0pt}%
\pgfpathmoveto{\pgfqpoint{4.178131in}{1.338408in}}%
\pgfpathlineto{\pgfqpoint{6.923821in}{1.338408in}}%
\pgfusepath{stroke}%
\end{pgfscope}%
\begin{pgfscope}%
\pgfsetbuttcap%
\pgfsetroundjoin%
\definecolor{currentfill}{rgb}{0.000000,0.000000,0.000000}%
\pgfsetfillcolor{currentfill}%
\pgfsetlinewidth{0.602250pt}%
\definecolor{currentstroke}{rgb}{0.000000,0.000000,0.000000}%
\pgfsetstrokecolor{currentstroke}%
\pgfsetdash{}{0pt}%
\pgfsys@defobject{currentmarker}{\pgfqpoint{-0.027778in}{0.000000in}}{\pgfqpoint{-0.000000in}{0.000000in}}{%
\pgfpathmoveto{\pgfqpoint{-0.000000in}{0.000000in}}%
\pgfpathlineto{\pgfqpoint{-0.027778in}{0.000000in}}%
\pgfusepath{stroke,fill}%
}%
\begin{pgfscope}%
\pgfsys@transformshift{4.178131in}{1.338408in}%
\pgfsys@useobject{currentmarker}{}%
\end{pgfscope}%
\end{pgfscope}%
\begin{pgfscope}%
\pgfpathrectangle{\pgfqpoint{4.178131in}{0.829139in}}{\pgfqpoint{2.745690in}{1.921500in}}%
\pgfusepath{clip}%
\pgfsetrectcap%
\pgfsetroundjoin%
\pgfsetlinewidth{0.803000pt}%
\definecolor{currentstroke}{rgb}{0.690196,0.690196,0.690196}%
\pgfsetstrokecolor{currentstroke}%
\pgfsetstrokeopacity{0.300000}%
\pgfsetdash{}{0pt}%
\pgfpathmoveto{\pgfqpoint{4.178131in}{1.351539in}}%
\pgfpathlineto{\pgfqpoint{6.923821in}{1.351539in}}%
\pgfusepath{stroke}%
\end{pgfscope}%
\begin{pgfscope}%
\pgfsetbuttcap%
\pgfsetroundjoin%
\definecolor{currentfill}{rgb}{0.000000,0.000000,0.000000}%
\pgfsetfillcolor{currentfill}%
\pgfsetlinewidth{0.602250pt}%
\definecolor{currentstroke}{rgb}{0.000000,0.000000,0.000000}%
\pgfsetstrokecolor{currentstroke}%
\pgfsetdash{}{0pt}%
\pgfsys@defobject{currentmarker}{\pgfqpoint{-0.027778in}{0.000000in}}{\pgfqpoint{-0.000000in}{0.000000in}}{%
\pgfpathmoveto{\pgfqpoint{-0.000000in}{0.000000in}}%
\pgfpathlineto{\pgfqpoint{-0.027778in}{0.000000in}}%
\pgfusepath{stroke,fill}%
}%
\begin{pgfscope}%
\pgfsys@transformshift{4.178131in}{1.351539in}%
\pgfsys@useobject{currentmarker}{}%
\end{pgfscope}%
\end{pgfscope}%
\begin{pgfscope}%
\pgfpathrectangle{\pgfqpoint{4.178131in}{0.829139in}}{\pgfqpoint{2.745690in}{1.921500in}}%
\pgfusepath{clip}%
\pgfsetrectcap%
\pgfsetroundjoin%
\pgfsetlinewidth{0.803000pt}%
\definecolor{currentstroke}{rgb}{0.690196,0.690196,0.690196}%
\pgfsetstrokecolor{currentstroke}%
\pgfsetstrokeopacity{0.300000}%
\pgfsetdash{}{0pt}%
\pgfpathmoveto{\pgfqpoint{4.178131in}{1.440558in}}%
\pgfpathlineto{\pgfqpoint{6.923821in}{1.440558in}}%
\pgfusepath{stroke}%
\end{pgfscope}%
\begin{pgfscope}%
\pgfsetbuttcap%
\pgfsetroundjoin%
\definecolor{currentfill}{rgb}{0.000000,0.000000,0.000000}%
\pgfsetfillcolor{currentfill}%
\pgfsetlinewidth{0.602250pt}%
\definecolor{currentstroke}{rgb}{0.000000,0.000000,0.000000}%
\pgfsetstrokecolor{currentstroke}%
\pgfsetdash{}{0pt}%
\pgfsys@defobject{currentmarker}{\pgfqpoint{-0.027778in}{0.000000in}}{\pgfqpoint{-0.000000in}{0.000000in}}{%
\pgfpathmoveto{\pgfqpoint{-0.000000in}{0.000000in}}%
\pgfpathlineto{\pgfqpoint{-0.027778in}{0.000000in}}%
\pgfusepath{stroke,fill}%
}%
\begin{pgfscope}%
\pgfsys@transformshift{4.178131in}{1.440558in}%
\pgfsys@useobject{currentmarker}{}%
\end{pgfscope}%
\end{pgfscope}%
\begin{pgfscope}%
\pgfpathrectangle{\pgfqpoint{4.178131in}{0.829139in}}{\pgfqpoint{2.745690in}{1.921500in}}%
\pgfusepath{clip}%
\pgfsetrectcap%
\pgfsetroundjoin%
\pgfsetlinewidth{0.803000pt}%
\definecolor{currentstroke}{rgb}{0.690196,0.690196,0.690196}%
\pgfsetstrokecolor{currentstroke}%
\pgfsetstrokeopacity{0.300000}%
\pgfsetdash{}{0pt}%
\pgfpathmoveto{\pgfqpoint{4.178131in}{1.485760in}}%
\pgfpathlineto{\pgfqpoint{6.923821in}{1.485760in}}%
\pgfusepath{stroke}%
\end{pgfscope}%
\begin{pgfscope}%
\pgfsetbuttcap%
\pgfsetroundjoin%
\definecolor{currentfill}{rgb}{0.000000,0.000000,0.000000}%
\pgfsetfillcolor{currentfill}%
\pgfsetlinewidth{0.602250pt}%
\definecolor{currentstroke}{rgb}{0.000000,0.000000,0.000000}%
\pgfsetstrokecolor{currentstroke}%
\pgfsetdash{}{0pt}%
\pgfsys@defobject{currentmarker}{\pgfqpoint{-0.027778in}{0.000000in}}{\pgfqpoint{-0.000000in}{0.000000in}}{%
\pgfpathmoveto{\pgfqpoint{-0.000000in}{0.000000in}}%
\pgfpathlineto{\pgfqpoint{-0.027778in}{0.000000in}}%
\pgfusepath{stroke,fill}%
}%
\begin{pgfscope}%
\pgfsys@transformshift{4.178131in}{1.485760in}%
\pgfsys@useobject{currentmarker}{}%
\end{pgfscope}%
\end{pgfscope}%
\begin{pgfscope}%
\pgfpathrectangle{\pgfqpoint{4.178131in}{0.829139in}}{\pgfqpoint{2.745690in}{1.921500in}}%
\pgfusepath{clip}%
\pgfsetrectcap%
\pgfsetroundjoin%
\pgfsetlinewidth{0.803000pt}%
\definecolor{currentstroke}{rgb}{0.690196,0.690196,0.690196}%
\pgfsetstrokecolor{currentstroke}%
\pgfsetstrokeopacity{0.300000}%
\pgfsetdash{}{0pt}%
\pgfpathmoveto{\pgfqpoint{4.178131in}{1.517832in}}%
\pgfpathlineto{\pgfqpoint{6.923821in}{1.517832in}}%
\pgfusepath{stroke}%
\end{pgfscope}%
\begin{pgfscope}%
\pgfsetbuttcap%
\pgfsetroundjoin%
\definecolor{currentfill}{rgb}{0.000000,0.000000,0.000000}%
\pgfsetfillcolor{currentfill}%
\pgfsetlinewidth{0.602250pt}%
\definecolor{currentstroke}{rgb}{0.000000,0.000000,0.000000}%
\pgfsetstrokecolor{currentstroke}%
\pgfsetdash{}{0pt}%
\pgfsys@defobject{currentmarker}{\pgfqpoint{-0.027778in}{0.000000in}}{\pgfqpoint{-0.000000in}{0.000000in}}{%
\pgfpathmoveto{\pgfqpoint{-0.000000in}{0.000000in}}%
\pgfpathlineto{\pgfqpoint{-0.027778in}{0.000000in}}%
\pgfusepath{stroke,fill}%
}%
\begin{pgfscope}%
\pgfsys@transformshift{4.178131in}{1.517832in}%
\pgfsys@useobject{currentmarker}{}%
\end{pgfscope}%
\end{pgfscope}%
\begin{pgfscope}%
\pgfpathrectangle{\pgfqpoint{4.178131in}{0.829139in}}{\pgfqpoint{2.745690in}{1.921500in}}%
\pgfusepath{clip}%
\pgfsetrectcap%
\pgfsetroundjoin%
\pgfsetlinewidth{0.803000pt}%
\definecolor{currentstroke}{rgb}{0.690196,0.690196,0.690196}%
\pgfsetstrokecolor{currentstroke}%
\pgfsetstrokeopacity{0.300000}%
\pgfsetdash{}{0pt}%
\pgfpathmoveto{\pgfqpoint{4.178131in}{1.542708in}}%
\pgfpathlineto{\pgfqpoint{6.923821in}{1.542708in}}%
\pgfusepath{stroke}%
\end{pgfscope}%
\begin{pgfscope}%
\pgfsetbuttcap%
\pgfsetroundjoin%
\definecolor{currentfill}{rgb}{0.000000,0.000000,0.000000}%
\pgfsetfillcolor{currentfill}%
\pgfsetlinewidth{0.602250pt}%
\definecolor{currentstroke}{rgb}{0.000000,0.000000,0.000000}%
\pgfsetstrokecolor{currentstroke}%
\pgfsetdash{}{0pt}%
\pgfsys@defobject{currentmarker}{\pgfqpoint{-0.027778in}{0.000000in}}{\pgfqpoint{-0.000000in}{0.000000in}}{%
\pgfpathmoveto{\pgfqpoint{-0.000000in}{0.000000in}}%
\pgfpathlineto{\pgfqpoint{-0.027778in}{0.000000in}}%
\pgfusepath{stroke,fill}%
}%
\begin{pgfscope}%
\pgfsys@transformshift{4.178131in}{1.542708in}%
\pgfsys@useobject{currentmarker}{}%
\end{pgfscope}%
\end{pgfscope}%
\begin{pgfscope}%
\pgfpathrectangle{\pgfqpoint{4.178131in}{0.829139in}}{\pgfqpoint{2.745690in}{1.921500in}}%
\pgfusepath{clip}%
\pgfsetrectcap%
\pgfsetroundjoin%
\pgfsetlinewidth{0.803000pt}%
\definecolor{currentstroke}{rgb}{0.690196,0.690196,0.690196}%
\pgfsetstrokecolor{currentstroke}%
\pgfsetstrokeopacity{0.300000}%
\pgfsetdash{}{0pt}%
\pgfpathmoveto{\pgfqpoint{4.178131in}{1.563034in}}%
\pgfpathlineto{\pgfqpoint{6.923821in}{1.563034in}}%
\pgfusepath{stroke}%
\end{pgfscope}%
\begin{pgfscope}%
\pgfsetbuttcap%
\pgfsetroundjoin%
\definecolor{currentfill}{rgb}{0.000000,0.000000,0.000000}%
\pgfsetfillcolor{currentfill}%
\pgfsetlinewidth{0.602250pt}%
\definecolor{currentstroke}{rgb}{0.000000,0.000000,0.000000}%
\pgfsetstrokecolor{currentstroke}%
\pgfsetdash{}{0pt}%
\pgfsys@defobject{currentmarker}{\pgfqpoint{-0.027778in}{0.000000in}}{\pgfqpoint{-0.000000in}{0.000000in}}{%
\pgfpathmoveto{\pgfqpoint{-0.000000in}{0.000000in}}%
\pgfpathlineto{\pgfqpoint{-0.027778in}{0.000000in}}%
\pgfusepath{stroke,fill}%
}%
\begin{pgfscope}%
\pgfsys@transformshift{4.178131in}{1.563034in}%
\pgfsys@useobject{currentmarker}{}%
\end{pgfscope}%
\end{pgfscope}%
\begin{pgfscope}%
\pgfpathrectangle{\pgfqpoint{4.178131in}{0.829139in}}{\pgfqpoint{2.745690in}{1.921500in}}%
\pgfusepath{clip}%
\pgfsetrectcap%
\pgfsetroundjoin%
\pgfsetlinewidth{0.803000pt}%
\definecolor{currentstroke}{rgb}{0.690196,0.690196,0.690196}%
\pgfsetstrokecolor{currentstroke}%
\pgfsetstrokeopacity{0.300000}%
\pgfsetdash{}{0pt}%
\pgfpathmoveto{\pgfqpoint{4.178131in}{1.580219in}}%
\pgfpathlineto{\pgfqpoint{6.923821in}{1.580219in}}%
\pgfusepath{stroke}%
\end{pgfscope}%
\begin{pgfscope}%
\pgfsetbuttcap%
\pgfsetroundjoin%
\definecolor{currentfill}{rgb}{0.000000,0.000000,0.000000}%
\pgfsetfillcolor{currentfill}%
\pgfsetlinewidth{0.602250pt}%
\definecolor{currentstroke}{rgb}{0.000000,0.000000,0.000000}%
\pgfsetstrokecolor{currentstroke}%
\pgfsetdash{}{0pt}%
\pgfsys@defobject{currentmarker}{\pgfqpoint{-0.027778in}{0.000000in}}{\pgfqpoint{-0.000000in}{0.000000in}}{%
\pgfpathmoveto{\pgfqpoint{-0.000000in}{0.000000in}}%
\pgfpathlineto{\pgfqpoint{-0.027778in}{0.000000in}}%
\pgfusepath{stroke,fill}%
}%
\begin{pgfscope}%
\pgfsys@transformshift{4.178131in}{1.580219in}%
\pgfsys@useobject{currentmarker}{}%
\end{pgfscope}%
\end{pgfscope}%
\begin{pgfscope}%
\pgfpathrectangle{\pgfqpoint{4.178131in}{0.829139in}}{\pgfqpoint{2.745690in}{1.921500in}}%
\pgfusepath{clip}%
\pgfsetrectcap%
\pgfsetroundjoin%
\pgfsetlinewidth{0.803000pt}%
\definecolor{currentstroke}{rgb}{0.690196,0.690196,0.690196}%
\pgfsetstrokecolor{currentstroke}%
\pgfsetstrokeopacity{0.300000}%
\pgfsetdash{}{0pt}%
\pgfpathmoveto{\pgfqpoint{4.178131in}{1.595105in}}%
\pgfpathlineto{\pgfqpoint{6.923821in}{1.595105in}}%
\pgfusepath{stroke}%
\end{pgfscope}%
\begin{pgfscope}%
\pgfsetbuttcap%
\pgfsetroundjoin%
\definecolor{currentfill}{rgb}{0.000000,0.000000,0.000000}%
\pgfsetfillcolor{currentfill}%
\pgfsetlinewidth{0.602250pt}%
\definecolor{currentstroke}{rgb}{0.000000,0.000000,0.000000}%
\pgfsetstrokecolor{currentstroke}%
\pgfsetdash{}{0pt}%
\pgfsys@defobject{currentmarker}{\pgfqpoint{-0.027778in}{0.000000in}}{\pgfqpoint{-0.000000in}{0.000000in}}{%
\pgfpathmoveto{\pgfqpoint{-0.000000in}{0.000000in}}%
\pgfpathlineto{\pgfqpoint{-0.027778in}{0.000000in}}%
\pgfusepath{stroke,fill}%
}%
\begin{pgfscope}%
\pgfsys@transformshift{4.178131in}{1.595105in}%
\pgfsys@useobject{currentmarker}{}%
\end{pgfscope}%
\end{pgfscope}%
\begin{pgfscope}%
\pgfpathrectangle{\pgfqpoint{4.178131in}{0.829139in}}{\pgfqpoint{2.745690in}{1.921500in}}%
\pgfusepath{clip}%
\pgfsetrectcap%
\pgfsetroundjoin%
\pgfsetlinewidth{0.803000pt}%
\definecolor{currentstroke}{rgb}{0.690196,0.690196,0.690196}%
\pgfsetstrokecolor{currentstroke}%
\pgfsetstrokeopacity{0.300000}%
\pgfsetdash{}{0pt}%
\pgfpathmoveto{\pgfqpoint{4.178131in}{1.608236in}}%
\pgfpathlineto{\pgfqpoint{6.923821in}{1.608236in}}%
\pgfusepath{stroke}%
\end{pgfscope}%
\begin{pgfscope}%
\pgfsetbuttcap%
\pgfsetroundjoin%
\definecolor{currentfill}{rgb}{0.000000,0.000000,0.000000}%
\pgfsetfillcolor{currentfill}%
\pgfsetlinewidth{0.602250pt}%
\definecolor{currentstroke}{rgb}{0.000000,0.000000,0.000000}%
\pgfsetstrokecolor{currentstroke}%
\pgfsetdash{}{0pt}%
\pgfsys@defobject{currentmarker}{\pgfqpoint{-0.027778in}{0.000000in}}{\pgfqpoint{-0.000000in}{0.000000in}}{%
\pgfpathmoveto{\pgfqpoint{-0.000000in}{0.000000in}}%
\pgfpathlineto{\pgfqpoint{-0.027778in}{0.000000in}}%
\pgfusepath{stroke,fill}%
}%
\begin{pgfscope}%
\pgfsys@transformshift{4.178131in}{1.608236in}%
\pgfsys@useobject{currentmarker}{}%
\end{pgfscope}%
\end{pgfscope}%
\begin{pgfscope}%
\pgfpathrectangle{\pgfqpoint{4.178131in}{0.829139in}}{\pgfqpoint{2.745690in}{1.921500in}}%
\pgfusepath{clip}%
\pgfsetrectcap%
\pgfsetroundjoin%
\pgfsetlinewidth{0.803000pt}%
\definecolor{currentstroke}{rgb}{0.690196,0.690196,0.690196}%
\pgfsetstrokecolor{currentstroke}%
\pgfsetstrokeopacity{0.300000}%
\pgfsetdash{}{0pt}%
\pgfpathmoveto{\pgfqpoint{4.178131in}{1.697255in}}%
\pgfpathlineto{\pgfqpoint{6.923821in}{1.697255in}}%
\pgfusepath{stroke}%
\end{pgfscope}%
\begin{pgfscope}%
\pgfsetbuttcap%
\pgfsetroundjoin%
\definecolor{currentfill}{rgb}{0.000000,0.000000,0.000000}%
\pgfsetfillcolor{currentfill}%
\pgfsetlinewidth{0.602250pt}%
\definecolor{currentstroke}{rgb}{0.000000,0.000000,0.000000}%
\pgfsetstrokecolor{currentstroke}%
\pgfsetdash{}{0pt}%
\pgfsys@defobject{currentmarker}{\pgfqpoint{-0.027778in}{0.000000in}}{\pgfqpoint{-0.000000in}{0.000000in}}{%
\pgfpathmoveto{\pgfqpoint{-0.000000in}{0.000000in}}%
\pgfpathlineto{\pgfqpoint{-0.027778in}{0.000000in}}%
\pgfusepath{stroke,fill}%
}%
\begin{pgfscope}%
\pgfsys@transformshift{4.178131in}{1.697255in}%
\pgfsys@useobject{currentmarker}{}%
\end{pgfscope}%
\end{pgfscope}%
\begin{pgfscope}%
\pgfpathrectangle{\pgfqpoint{4.178131in}{0.829139in}}{\pgfqpoint{2.745690in}{1.921500in}}%
\pgfusepath{clip}%
\pgfsetrectcap%
\pgfsetroundjoin%
\pgfsetlinewidth{0.803000pt}%
\definecolor{currentstroke}{rgb}{0.690196,0.690196,0.690196}%
\pgfsetstrokecolor{currentstroke}%
\pgfsetstrokeopacity{0.300000}%
\pgfsetdash{}{0pt}%
\pgfpathmoveto{\pgfqpoint{4.178131in}{1.742457in}}%
\pgfpathlineto{\pgfqpoint{6.923821in}{1.742457in}}%
\pgfusepath{stroke}%
\end{pgfscope}%
\begin{pgfscope}%
\pgfsetbuttcap%
\pgfsetroundjoin%
\definecolor{currentfill}{rgb}{0.000000,0.000000,0.000000}%
\pgfsetfillcolor{currentfill}%
\pgfsetlinewidth{0.602250pt}%
\definecolor{currentstroke}{rgb}{0.000000,0.000000,0.000000}%
\pgfsetstrokecolor{currentstroke}%
\pgfsetdash{}{0pt}%
\pgfsys@defobject{currentmarker}{\pgfqpoint{-0.027778in}{0.000000in}}{\pgfqpoint{-0.000000in}{0.000000in}}{%
\pgfpathmoveto{\pgfqpoint{-0.000000in}{0.000000in}}%
\pgfpathlineto{\pgfqpoint{-0.027778in}{0.000000in}}%
\pgfusepath{stroke,fill}%
}%
\begin{pgfscope}%
\pgfsys@transformshift{4.178131in}{1.742457in}%
\pgfsys@useobject{currentmarker}{}%
\end{pgfscope}%
\end{pgfscope}%
\begin{pgfscope}%
\pgfpathrectangle{\pgfqpoint{4.178131in}{0.829139in}}{\pgfqpoint{2.745690in}{1.921500in}}%
\pgfusepath{clip}%
\pgfsetrectcap%
\pgfsetroundjoin%
\pgfsetlinewidth{0.803000pt}%
\definecolor{currentstroke}{rgb}{0.690196,0.690196,0.690196}%
\pgfsetstrokecolor{currentstroke}%
\pgfsetstrokeopacity{0.300000}%
\pgfsetdash{}{0pt}%
\pgfpathmoveto{\pgfqpoint{4.178131in}{1.774529in}}%
\pgfpathlineto{\pgfqpoint{6.923821in}{1.774529in}}%
\pgfusepath{stroke}%
\end{pgfscope}%
\begin{pgfscope}%
\pgfsetbuttcap%
\pgfsetroundjoin%
\definecolor{currentfill}{rgb}{0.000000,0.000000,0.000000}%
\pgfsetfillcolor{currentfill}%
\pgfsetlinewidth{0.602250pt}%
\definecolor{currentstroke}{rgb}{0.000000,0.000000,0.000000}%
\pgfsetstrokecolor{currentstroke}%
\pgfsetdash{}{0pt}%
\pgfsys@defobject{currentmarker}{\pgfqpoint{-0.027778in}{0.000000in}}{\pgfqpoint{-0.000000in}{0.000000in}}{%
\pgfpathmoveto{\pgfqpoint{-0.000000in}{0.000000in}}%
\pgfpathlineto{\pgfqpoint{-0.027778in}{0.000000in}}%
\pgfusepath{stroke,fill}%
}%
\begin{pgfscope}%
\pgfsys@transformshift{4.178131in}{1.774529in}%
\pgfsys@useobject{currentmarker}{}%
\end{pgfscope}%
\end{pgfscope}%
\begin{pgfscope}%
\pgfpathrectangle{\pgfqpoint{4.178131in}{0.829139in}}{\pgfqpoint{2.745690in}{1.921500in}}%
\pgfusepath{clip}%
\pgfsetrectcap%
\pgfsetroundjoin%
\pgfsetlinewidth{0.803000pt}%
\definecolor{currentstroke}{rgb}{0.690196,0.690196,0.690196}%
\pgfsetstrokecolor{currentstroke}%
\pgfsetstrokeopacity{0.300000}%
\pgfsetdash{}{0pt}%
\pgfpathmoveto{\pgfqpoint{4.178131in}{1.799405in}}%
\pgfpathlineto{\pgfqpoint{6.923821in}{1.799405in}}%
\pgfusepath{stroke}%
\end{pgfscope}%
\begin{pgfscope}%
\pgfsetbuttcap%
\pgfsetroundjoin%
\definecolor{currentfill}{rgb}{0.000000,0.000000,0.000000}%
\pgfsetfillcolor{currentfill}%
\pgfsetlinewidth{0.602250pt}%
\definecolor{currentstroke}{rgb}{0.000000,0.000000,0.000000}%
\pgfsetstrokecolor{currentstroke}%
\pgfsetdash{}{0pt}%
\pgfsys@defobject{currentmarker}{\pgfqpoint{-0.027778in}{0.000000in}}{\pgfqpoint{-0.000000in}{0.000000in}}{%
\pgfpathmoveto{\pgfqpoint{-0.000000in}{0.000000in}}%
\pgfpathlineto{\pgfqpoint{-0.027778in}{0.000000in}}%
\pgfusepath{stroke,fill}%
}%
\begin{pgfscope}%
\pgfsys@transformshift{4.178131in}{1.799405in}%
\pgfsys@useobject{currentmarker}{}%
\end{pgfscope}%
\end{pgfscope}%
\begin{pgfscope}%
\pgfpathrectangle{\pgfqpoint{4.178131in}{0.829139in}}{\pgfqpoint{2.745690in}{1.921500in}}%
\pgfusepath{clip}%
\pgfsetrectcap%
\pgfsetroundjoin%
\pgfsetlinewidth{0.803000pt}%
\definecolor{currentstroke}{rgb}{0.690196,0.690196,0.690196}%
\pgfsetstrokecolor{currentstroke}%
\pgfsetstrokeopacity{0.300000}%
\pgfsetdash{}{0pt}%
\pgfpathmoveto{\pgfqpoint{4.178131in}{1.819731in}}%
\pgfpathlineto{\pgfqpoint{6.923821in}{1.819731in}}%
\pgfusepath{stroke}%
\end{pgfscope}%
\begin{pgfscope}%
\pgfsetbuttcap%
\pgfsetroundjoin%
\definecolor{currentfill}{rgb}{0.000000,0.000000,0.000000}%
\pgfsetfillcolor{currentfill}%
\pgfsetlinewidth{0.602250pt}%
\definecolor{currentstroke}{rgb}{0.000000,0.000000,0.000000}%
\pgfsetstrokecolor{currentstroke}%
\pgfsetdash{}{0pt}%
\pgfsys@defobject{currentmarker}{\pgfqpoint{-0.027778in}{0.000000in}}{\pgfqpoint{-0.000000in}{0.000000in}}{%
\pgfpathmoveto{\pgfqpoint{-0.000000in}{0.000000in}}%
\pgfpathlineto{\pgfqpoint{-0.027778in}{0.000000in}}%
\pgfusepath{stroke,fill}%
}%
\begin{pgfscope}%
\pgfsys@transformshift{4.178131in}{1.819731in}%
\pgfsys@useobject{currentmarker}{}%
\end{pgfscope}%
\end{pgfscope}%
\begin{pgfscope}%
\pgfpathrectangle{\pgfqpoint{4.178131in}{0.829139in}}{\pgfqpoint{2.745690in}{1.921500in}}%
\pgfusepath{clip}%
\pgfsetrectcap%
\pgfsetroundjoin%
\pgfsetlinewidth{0.803000pt}%
\definecolor{currentstroke}{rgb}{0.690196,0.690196,0.690196}%
\pgfsetstrokecolor{currentstroke}%
\pgfsetstrokeopacity{0.300000}%
\pgfsetdash{}{0pt}%
\pgfpathmoveto{\pgfqpoint{4.178131in}{1.836916in}}%
\pgfpathlineto{\pgfqpoint{6.923821in}{1.836916in}}%
\pgfusepath{stroke}%
\end{pgfscope}%
\begin{pgfscope}%
\pgfsetbuttcap%
\pgfsetroundjoin%
\definecolor{currentfill}{rgb}{0.000000,0.000000,0.000000}%
\pgfsetfillcolor{currentfill}%
\pgfsetlinewidth{0.602250pt}%
\definecolor{currentstroke}{rgb}{0.000000,0.000000,0.000000}%
\pgfsetstrokecolor{currentstroke}%
\pgfsetdash{}{0pt}%
\pgfsys@defobject{currentmarker}{\pgfqpoint{-0.027778in}{0.000000in}}{\pgfqpoint{-0.000000in}{0.000000in}}{%
\pgfpathmoveto{\pgfqpoint{-0.000000in}{0.000000in}}%
\pgfpathlineto{\pgfqpoint{-0.027778in}{0.000000in}}%
\pgfusepath{stroke,fill}%
}%
\begin{pgfscope}%
\pgfsys@transformshift{4.178131in}{1.836916in}%
\pgfsys@useobject{currentmarker}{}%
\end{pgfscope}%
\end{pgfscope}%
\begin{pgfscope}%
\pgfpathrectangle{\pgfqpoint{4.178131in}{0.829139in}}{\pgfqpoint{2.745690in}{1.921500in}}%
\pgfusepath{clip}%
\pgfsetrectcap%
\pgfsetroundjoin%
\pgfsetlinewidth{0.803000pt}%
\definecolor{currentstroke}{rgb}{0.690196,0.690196,0.690196}%
\pgfsetstrokecolor{currentstroke}%
\pgfsetstrokeopacity{0.300000}%
\pgfsetdash{}{0pt}%
\pgfpathmoveto{\pgfqpoint{4.178131in}{1.851802in}}%
\pgfpathlineto{\pgfqpoint{6.923821in}{1.851802in}}%
\pgfusepath{stroke}%
\end{pgfscope}%
\begin{pgfscope}%
\pgfsetbuttcap%
\pgfsetroundjoin%
\definecolor{currentfill}{rgb}{0.000000,0.000000,0.000000}%
\pgfsetfillcolor{currentfill}%
\pgfsetlinewidth{0.602250pt}%
\definecolor{currentstroke}{rgb}{0.000000,0.000000,0.000000}%
\pgfsetstrokecolor{currentstroke}%
\pgfsetdash{}{0pt}%
\pgfsys@defobject{currentmarker}{\pgfqpoint{-0.027778in}{0.000000in}}{\pgfqpoint{-0.000000in}{0.000000in}}{%
\pgfpathmoveto{\pgfqpoint{-0.000000in}{0.000000in}}%
\pgfpathlineto{\pgfqpoint{-0.027778in}{0.000000in}}%
\pgfusepath{stroke,fill}%
}%
\begin{pgfscope}%
\pgfsys@transformshift{4.178131in}{1.851802in}%
\pgfsys@useobject{currentmarker}{}%
\end{pgfscope}%
\end{pgfscope}%
\begin{pgfscope}%
\pgfpathrectangle{\pgfqpoint{4.178131in}{0.829139in}}{\pgfqpoint{2.745690in}{1.921500in}}%
\pgfusepath{clip}%
\pgfsetrectcap%
\pgfsetroundjoin%
\pgfsetlinewidth{0.803000pt}%
\definecolor{currentstroke}{rgb}{0.690196,0.690196,0.690196}%
\pgfsetstrokecolor{currentstroke}%
\pgfsetstrokeopacity{0.300000}%
\pgfsetdash{}{0pt}%
\pgfpathmoveto{\pgfqpoint{4.178131in}{1.864933in}}%
\pgfpathlineto{\pgfqpoint{6.923821in}{1.864933in}}%
\pgfusepath{stroke}%
\end{pgfscope}%
\begin{pgfscope}%
\pgfsetbuttcap%
\pgfsetroundjoin%
\definecolor{currentfill}{rgb}{0.000000,0.000000,0.000000}%
\pgfsetfillcolor{currentfill}%
\pgfsetlinewidth{0.602250pt}%
\definecolor{currentstroke}{rgb}{0.000000,0.000000,0.000000}%
\pgfsetstrokecolor{currentstroke}%
\pgfsetdash{}{0pt}%
\pgfsys@defobject{currentmarker}{\pgfqpoint{-0.027778in}{0.000000in}}{\pgfqpoint{-0.000000in}{0.000000in}}{%
\pgfpathmoveto{\pgfqpoint{-0.000000in}{0.000000in}}%
\pgfpathlineto{\pgfqpoint{-0.027778in}{0.000000in}}%
\pgfusepath{stroke,fill}%
}%
\begin{pgfscope}%
\pgfsys@transformshift{4.178131in}{1.864933in}%
\pgfsys@useobject{currentmarker}{}%
\end{pgfscope}%
\end{pgfscope}%
\begin{pgfscope}%
\pgfpathrectangle{\pgfqpoint{4.178131in}{0.829139in}}{\pgfqpoint{2.745690in}{1.921500in}}%
\pgfusepath{clip}%
\pgfsetrectcap%
\pgfsetroundjoin%
\pgfsetlinewidth{0.803000pt}%
\definecolor{currentstroke}{rgb}{0.690196,0.690196,0.690196}%
\pgfsetstrokecolor{currentstroke}%
\pgfsetstrokeopacity{0.300000}%
\pgfsetdash{}{0pt}%
\pgfpathmoveto{\pgfqpoint{4.178131in}{1.953953in}}%
\pgfpathlineto{\pgfqpoint{6.923821in}{1.953953in}}%
\pgfusepath{stroke}%
\end{pgfscope}%
\begin{pgfscope}%
\pgfsetbuttcap%
\pgfsetroundjoin%
\definecolor{currentfill}{rgb}{0.000000,0.000000,0.000000}%
\pgfsetfillcolor{currentfill}%
\pgfsetlinewidth{0.602250pt}%
\definecolor{currentstroke}{rgb}{0.000000,0.000000,0.000000}%
\pgfsetstrokecolor{currentstroke}%
\pgfsetdash{}{0pt}%
\pgfsys@defobject{currentmarker}{\pgfqpoint{-0.027778in}{0.000000in}}{\pgfqpoint{-0.000000in}{0.000000in}}{%
\pgfpathmoveto{\pgfqpoint{-0.000000in}{0.000000in}}%
\pgfpathlineto{\pgfqpoint{-0.027778in}{0.000000in}}%
\pgfusepath{stroke,fill}%
}%
\begin{pgfscope}%
\pgfsys@transformshift{4.178131in}{1.953953in}%
\pgfsys@useobject{currentmarker}{}%
\end{pgfscope}%
\end{pgfscope}%
\begin{pgfscope}%
\pgfpathrectangle{\pgfqpoint{4.178131in}{0.829139in}}{\pgfqpoint{2.745690in}{1.921500in}}%
\pgfusepath{clip}%
\pgfsetrectcap%
\pgfsetroundjoin%
\pgfsetlinewidth{0.803000pt}%
\definecolor{currentstroke}{rgb}{0.690196,0.690196,0.690196}%
\pgfsetstrokecolor{currentstroke}%
\pgfsetstrokeopacity{0.300000}%
\pgfsetdash{}{0pt}%
\pgfpathmoveto{\pgfqpoint{4.178131in}{1.999155in}}%
\pgfpathlineto{\pgfqpoint{6.923821in}{1.999155in}}%
\pgfusepath{stroke}%
\end{pgfscope}%
\begin{pgfscope}%
\pgfsetbuttcap%
\pgfsetroundjoin%
\definecolor{currentfill}{rgb}{0.000000,0.000000,0.000000}%
\pgfsetfillcolor{currentfill}%
\pgfsetlinewidth{0.602250pt}%
\definecolor{currentstroke}{rgb}{0.000000,0.000000,0.000000}%
\pgfsetstrokecolor{currentstroke}%
\pgfsetdash{}{0pt}%
\pgfsys@defobject{currentmarker}{\pgfqpoint{-0.027778in}{0.000000in}}{\pgfqpoint{-0.000000in}{0.000000in}}{%
\pgfpathmoveto{\pgfqpoint{-0.000000in}{0.000000in}}%
\pgfpathlineto{\pgfqpoint{-0.027778in}{0.000000in}}%
\pgfusepath{stroke,fill}%
}%
\begin{pgfscope}%
\pgfsys@transformshift{4.178131in}{1.999155in}%
\pgfsys@useobject{currentmarker}{}%
\end{pgfscope}%
\end{pgfscope}%
\begin{pgfscope}%
\pgfpathrectangle{\pgfqpoint{4.178131in}{0.829139in}}{\pgfqpoint{2.745690in}{1.921500in}}%
\pgfusepath{clip}%
\pgfsetrectcap%
\pgfsetroundjoin%
\pgfsetlinewidth{0.803000pt}%
\definecolor{currentstroke}{rgb}{0.690196,0.690196,0.690196}%
\pgfsetstrokecolor{currentstroke}%
\pgfsetstrokeopacity{0.300000}%
\pgfsetdash{}{0pt}%
\pgfpathmoveto{\pgfqpoint{4.178131in}{2.031226in}}%
\pgfpathlineto{\pgfqpoint{6.923821in}{2.031226in}}%
\pgfusepath{stroke}%
\end{pgfscope}%
\begin{pgfscope}%
\pgfsetbuttcap%
\pgfsetroundjoin%
\definecolor{currentfill}{rgb}{0.000000,0.000000,0.000000}%
\pgfsetfillcolor{currentfill}%
\pgfsetlinewidth{0.602250pt}%
\definecolor{currentstroke}{rgb}{0.000000,0.000000,0.000000}%
\pgfsetstrokecolor{currentstroke}%
\pgfsetdash{}{0pt}%
\pgfsys@defobject{currentmarker}{\pgfqpoint{-0.027778in}{0.000000in}}{\pgfqpoint{-0.000000in}{0.000000in}}{%
\pgfpathmoveto{\pgfqpoint{-0.000000in}{0.000000in}}%
\pgfpathlineto{\pgfqpoint{-0.027778in}{0.000000in}}%
\pgfusepath{stroke,fill}%
}%
\begin{pgfscope}%
\pgfsys@transformshift{4.178131in}{2.031226in}%
\pgfsys@useobject{currentmarker}{}%
\end{pgfscope}%
\end{pgfscope}%
\begin{pgfscope}%
\pgfpathrectangle{\pgfqpoint{4.178131in}{0.829139in}}{\pgfqpoint{2.745690in}{1.921500in}}%
\pgfusepath{clip}%
\pgfsetrectcap%
\pgfsetroundjoin%
\pgfsetlinewidth{0.803000pt}%
\definecolor{currentstroke}{rgb}{0.690196,0.690196,0.690196}%
\pgfsetstrokecolor{currentstroke}%
\pgfsetstrokeopacity{0.300000}%
\pgfsetdash{}{0pt}%
\pgfpathmoveto{\pgfqpoint{4.178131in}{2.056103in}}%
\pgfpathlineto{\pgfqpoint{6.923821in}{2.056103in}}%
\pgfusepath{stroke}%
\end{pgfscope}%
\begin{pgfscope}%
\pgfsetbuttcap%
\pgfsetroundjoin%
\definecolor{currentfill}{rgb}{0.000000,0.000000,0.000000}%
\pgfsetfillcolor{currentfill}%
\pgfsetlinewidth{0.602250pt}%
\definecolor{currentstroke}{rgb}{0.000000,0.000000,0.000000}%
\pgfsetstrokecolor{currentstroke}%
\pgfsetdash{}{0pt}%
\pgfsys@defobject{currentmarker}{\pgfqpoint{-0.027778in}{0.000000in}}{\pgfqpoint{-0.000000in}{0.000000in}}{%
\pgfpathmoveto{\pgfqpoint{-0.000000in}{0.000000in}}%
\pgfpathlineto{\pgfqpoint{-0.027778in}{0.000000in}}%
\pgfusepath{stroke,fill}%
}%
\begin{pgfscope}%
\pgfsys@transformshift{4.178131in}{2.056103in}%
\pgfsys@useobject{currentmarker}{}%
\end{pgfscope}%
\end{pgfscope}%
\begin{pgfscope}%
\pgfpathrectangle{\pgfqpoint{4.178131in}{0.829139in}}{\pgfqpoint{2.745690in}{1.921500in}}%
\pgfusepath{clip}%
\pgfsetrectcap%
\pgfsetroundjoin%
\pgfsetlinewidth{0.803000pt}%
\definecolor{currentstroke}{rgb}{0.690196,0.690196,0.690196}%
\pgfsetstrokecolor{currentstroke}%
\pgfsetstrokeopacity{0.300000}%
\pgfsetdash{}{0pt}%
\pgfpathmoveto{\pgfqpoint{4.178131in}{2.076428in}}%
\pgfpathlineto{\pgfqpoint{6.923821in}{2.076428in}}%
\pgfusepath{stroke}%
\end{pgfscope}%
\begin{pgfscope}%
\pgfsetbuttcap%
\pgfsetroundjoin%
\definecolor{currentfill}{rgb}{0.000000,0.000000,0.000000}%
\pgfsetfillcolor{currentfill}%
\pgfsetlinewidth{0.602250pt}%
\definecolor{currentstroke}{rgb}{0.000000,0.000000,0.000000}%
\pgfsetstrokecolor{currentstroke}%
\pgfsetdash{}{0pt}%
\pgfsys@defobject{currentmarker}{\pgfqpoint{-0.027778in}{0.000000in}}{\pgfqpoint{-0.000000in}{0.000000in}}{%
\pgfpathmoveto{\pgfqpoint{-0.000000in}{0.000000in}}%
\pgfpathlineto{\pgfqpoint{-0.027778in}{0.000000in}}%
\pgfusepath{stroke,fill}%
}%
\begin{pgfscope}%
\pgfsys@transformshift{4.178131in}{2.076428in}%
\pgfsys@useobject{currentmarker}{}%
\end{pgfscope}%
\end{pgfscope}%
\begin{pgfscope}%
\pgfpathrectangle{\pgfqpoint{4.178131in}{0.829139in}}{\pgfqpoint{2.745690in}{1.921500in}}%
\pgfusepath{clip}%
\pgfsetrectcap%
\pgfsetroundjoin%
\pgfsetlinewidth{0.803000pt}%
\definecolor{currentstroke}{rgb}{0.690196,0.690196,0.690196}%
\pgfsetstrokecolor{currentstroke}%
\pgfsetstrokeopacity{0.300000}%
\pgfsetdash{}{0pt}%
\pgfpathmoveto{\pgfqpoint{4.178131in}{2.093613in}}%
\pgfpathlineto{\pgfqpoint{6.923821in}{2.093613in}}%
\pgfusepath{stroke}%
\end{pgfscope}%
\begin{pgfscope}%
\pgfsetbuttcap%
\pgfsetroundjoin%
\definecolor{currentfill}{rgb}{0.000000,0.000000,0.000000}%
\pgfsetfillcolor{currentfill}%
\pgfsetlinewidth{0.602250pt}%
\definecolor{currentstroke}{rgb}{0.000000,0.000000,0.000000}%
\pgfsetstrokecolor{currentstroke}%
\pgfsetdash{}{0pt}%
\pgfsys@defobject{currentmarker}{\pgfqpoint{-0.027778in}{0.000000in}}{\pgfqpoint{-0.000000in}{0.000000in}}{%
\pgfpathmoveto{\pgfqpoint{-0.000000in}{0.000000in}}%
\pgfpathlineto{\pgfqpoint{-0.027778in}{0.000000in}}%
\pgfusepath{stroke,fill}%
}%
\begin{pgfscope}%
\pgfsys@transformshift{4.178131in}{2.093613in}%
\pgfsys@useobject{currentmarker}{}%
\end{pgfscope}%
\end{pgfscope}%
\begin{pgfscope}%
\pgfpathrectangle{\pgfqpoint{4.178131in}{0.829139in}}{\pgfqpoint{2.745690in}{1.921500in}}%
\pgfusepath{clip}%
\pgfsetrectcap%
\pgfsetroundjoin%
\pgfsetlinewidth{0.803000pt}%
\definecolor{currentstroke}{rgb}{0.690196,0.690196,0.690196}%
\pgfsetstrokecolor{currentstroke}%
\pgfsetstrokeopacity{0.300000}%
\pgfsetdash{}{0pt}%
\pgfpathmoveto{\pgfqpoint{4.178131in}{2.108500in}}%
\pgfpathlineto{\pgfqpoint{6.923821in}{2.108500in}}%
\pgfusepath{stroke}%
\end{pgfscope}%
\begin{pgfscope}%
\pgfsetbuttcap%
\pgfsetroundjoin%
\definecolor{currentfill}{rgb}{0.000000,0.000000,0.000000}%
\pgfsetfillcolor{currentfill}%
\pgfsetlinewidth{0.602250pt}%
\definecolor{currentstroke}{rgb}{0.000000,0.000000,0.000000}%
\pgfsetstrokecolor{currentstroke}%
\pgfsetdash{}{0pt}%
\pgfsys@defobject{currentmarker}{\pgfqpoint{-0.027778in}{0.000000in}}{\pgfqpoint{-0.000000in}{0.000000in}}{%
\pgfpathmoveto{\pgfqpoint{-0.000000in}{0.000000in}}%
\pgfpathlineto{\pgfqpoint{-0.027778in}{0.000000in}}%
\pgfusepath{stroke,fill}%
}%
\begin{pgfscope}%
\pgfsys@transformshift{4.178131in}{2.108500in}%
\pgfsys@useobject{currentmarker}{}%
\end{pgfscope}%
\end{pgfscope}%
\begin{pgfscope}%
\pgfpathrectangle{\pgfqpoint{4.178131in}{0.829139in}}{\pgfqpoint{2.745690in}{1.921500in}}%
\pgfusepath{clip}%
\pgfsetrectcap%
\pgfsetroundjoin%
\pgfsetlinewidth{0.803000pt}%
\definecolor{currentstroke}{rgb}{0.690196,0.690196,0.690196}%
\pgfsetstrokecolor{currentstroke}%
\pgfsetstrokeopacity{0.300000}%
\pgfsetdash{}{0pt}%
\pgfpathmoveto{\pgfqpoint{4.178131in}{2.121631in}}%
\pgfpathlineto{\pgfqpoint{6.923821in}{2.121631in}}%
\pgfusepath{stroke}%
\end{pgfscope}%
\begin{pgfscope}%
\pgfsetbuttcap%
\pgfsetroundjoin%
\definecolor{currentfill}{rgb}{0.000000,0.000000,0.000000}%
\pgfsetfillcolor{currentfill}%
\pgfsetlinewidth{0.602250pt}%
\definecolor{currentstroke}{rgb}{0.000000,0.000000,0.000000}%
\pgfsetstrokecolor{currentstroke}%
\pgfsetdash{}{0pt}%
\pgfsys@defobject{currentmarker}{\pgfqpoint{-0.027778in}{0.000000in}}{\pgfqpoint{-0.000000in}{0.000000in}}{%
\pgfpathmoveto{\pgfqpoint{-0.000000in}{0.000000in}}%
\pgfpathlineto{\pgfqpoint{-0.027778in}{0.000000in}}%
\pgfusepath{stroke,fill}%
}%
\begin{pgfscope}%
\pgfsys@transformshift{4.178131in}{2.121631in}%
\pgfsys@useobject{currentmarker}{}%
\end{pgfscope}%
\end{pgfscope}%
\begin{pgfscope}%
\pgfpathrectangle{\pgfqpoint{4.178131in}{0.829139in}}{\pgfqpoint{2.745690in}{1.921500in}}%
\pgfusepath{clip}%
\pgfsetrectcap%
\pgfsetroundjoin%
\pgfsetlinewidth{0.803000pt}%
\definecolor{currentstroke}{rgb}{0.690196,0.690196,0.690196}%
\pgfsetstrokecolor{currentstroke}%
\pgfsetstrokeopacity{0.300000}%
\pgfsetdash{}{0pt}%
\pgfpathmoveto{\pgfqpoint{4.178131in}{2.210650in}}%
\pgfpathlineto{\pgfqpoint{6.923821in}{2.210650in}}%
\pgfusepath{stroke}%
\end{pgfscope}%
\begin{pgfscope}%
\pgfsetbuttcap%
\pgfsetroundjoin%
\definecolor{currentfill}{rgb}{0.000000,0.000000,0.000000}%
\pgfsetfillcolor{currentfill}%
\pgfsetlinewidth{0.602250pt}%
\definecolor{currentstroke}{rgb}{0.000000,0.000000,0.000000}%
\pgfsetstrokecolor{currentstroke}%
\pgfsetdash{}{0pt}%
\pgfsys@defobject{currentmarker}{\pgfqpoint{-0.027778in}{0.000000in}}{\pgfqpoint{-0.000000in}{0.000000in}}{%
\pgfpathmoveto{\pgfqpoint{-0.000000in}{0.000000in}}%
\pgfpathlineto{\pgfqpoint{-0.027778in}{0.000000in}}%
\pgfusepath{stroke,fill}%
}%
\begin{pgfscope}%
\pgfsys@transformshift{4.178131in}{2.210650in}%
\pgfsys@useobject{currentmarker}{}%
\end{pgfscope}%
\end{pgfscope}%
\begin{pgfscope}%
\pgfpathrectangle{\pgfqpoint{4.178131in}{0.829139in}}{\pgfqpoint{2.745690in}{1.921500in}}%
\pgfusepath{clip}%
\pgfsetrectcap%
\pgfsetroundjoin%
\pgfsetlinewidth{0.803000pt}%
\definecolor{currentstroke}{rgb}{0.690196,0.690196,0.690196}%
\pgfsetstrokecolor{currentstroke}%
\pgfsetstrokeopacity{0.300000}%
\pgfsetdash{}{0pt}%
\pgfpathmoveto{\pgfqpoint{4.178131in}{2.255852in}}%
\pgfpathlineto{\pgfqpoint{6.923821in}{2.255852in}}%
\pgfusepath{stroke}%
\end{pgfscope}%
\begin{pgfscope}%
\pgfsetbuttcap%
\pgfsetroundjoin%
\definecolor{currentfill}{rgb}{0.000000,0.000000,0.000000}%
\pgfsetfillcolor{currentfill}%
\pgfsetlinewidth{0.602250pt}%
\definecolor{currentstroke}{rgb}{0.000000,0.000000,0.000000}%
\pgfsetstrokecolor{currentstroke}%
\pgfsetdash{}{0pt}%
\pgfsys@defobject{currentmarker}{\pgfqpoint{-0.027778in}{0.000000in}}{\pgfqpoint{-0.000000in}{0.000000in}}{%
\pgfpathmoveto{\pgfqpoint{-0.000000in}{0.000000in}}%
\pgfpathlineto{\pgfqpoint{-0.027778in}{0.000000in}}%
\pgfusepath{stroke,fill}%
}%
\begin{pgfscope}%
\pgfsys@transformshift{4.178131in}{2.255852in}%
\pgfsys@useobject{currentmarker}{}%
\end{pgfscope}%
\end{pgfscope}%
\begin{pgfscope}%
\pgfpathrectangle{\pgfqpoint{4.178131in}{0.829139in}}{\pgfqpoint{2.745690in}{1.921500in}}%
\pgfusepath{clip}%
\pgfsetrectcap%
\pgfsetroundjoin%
\pgfsetlinewidth{0.803000pt}%
\definecolor{currentstroke}{rgb}{0.690196,0.690196,0.690196}%
\pgfsetstrokecolor{currentstroke}%
\pgfsetstrokeopacity{0.300000}%
\pgfsetdash{}{0pt}%
\pgfpathmoveto{\pgfqpoint{4.178131in}{2.287924in}}%
\pgfpathlineto{\pgfqpoint{6.923821in}{2.287924in}}%
\pgfusepath{stroke}%
\end{pgfscope}%
\begin{pgfscope}%
\pgfsetbuttcap%
\pgfsetroundjoin%
\definecolor{currentfill}{rgb}{0.000000,0.000000,0.000000}%
\pgfsetfillcolor{currentfill}%
\pgfsetlinewidth{0.602250pt}%
\definecolor{currentstroke}{rgb}{0.000000,0.000000,0.000000}%
\pgfsetstrokecolor{currentstroke}%
\pgfsetdash{}{0pt}%
\pgfsys@defobject{currentmarker}{\pgfqpoint{-0.027778in}{0.000000in}}{\pgfqpoint{-0.000000in}{0.000000in}}{%
\pgfpathmoveto{\pgfqpoint{-0.000000in}{0.000000in}}%
\pgfpathlineto{\pgfqpoint{-0.027778in}{0.000000in}}%
\pgfusepath{stroke,fill}%
}%
\begin{pgfscope}%
\pgfsys@transformshift{4.178131in}{2.287924in}%
\pgfsys@useobject{currentmarker}{}%
\end{pgfscope}%
\end{pgfscope}%
\begin{pgfscope}%
\pgfpathrectangle{\pgfqpoint{4.178131in}{0.829139in}}{\pgfqpoint{2.745690in}{1.921500in}}%
\pgfusepath{clip}%
\pgfsetrectcap%
\pgfsetroundjoin%
\pgfsetlinewidth{0.803000pt}%
\definecolor{currentstroke}{rgb}{0.690196,0.690196,0.690196}%
\pgfsetstrokecolor{currentstroke}%
\pgfsetstrokeopacity{0.300000}%
\pgfsetdash{}{0pt}%
\pgfpathmoveto{\pgfqpoint{4.178131in}{2.312800in}}%
\pgfpathlineto{\pgfqpoint{6.923821in}{2.312800in}}%
\pgfusepath{stroke}%
\end{pgfscope}%
\begin{pgfscope}%
\pgfsetbuttcap%
\pgfsetroundjoin%
\definecolor{currentfill}{rgb}{0.000000,0.000000,0.000000}%
\pgfsetfillcolor{currentfill}%
\pgfsetlinewidth{0.602250pt}%
\definecolor{currentstroke}{rgb}{0.000000,0.000000,0.000000}%
\pgfsetstrokecolor{currentstroke}%
\pgfsetdash{}{0pt}%
\pgfsys@defobject{currentmarker}{\pgfqpoint{-0.027778in}{0.000000in}}{\pgfqpoint{-0.000000in}{0.000000in}}{%
\pgfpathmoveto{\pgfqpoint{-0.000000in}{0.000000in}}%
\pgfpathlineto{\pgfqpoint{-0.027778in}{0.000000in}}%
\pgfusepath{stroke,fill}%
}%
\begin{pgfscope}%
\pgfsys@transformshift{4.178131in}{2.312800in}%
\pgfsys@useobject{currentmarker}{}%
\end{pgfscope}%
\end{pgfscope}%
\begin{pgfscope}%
\pgfpathrectangle{\pgfqpoint{4.178131in}{0.829139in}}{\pgfqpoint{2.745690in}{1.921500in}}%
\pgfusepath{clip}%
\pgfsetrectcap%
\pgfsetroundjoin%
\pgfsetlinewidth{0.803000pt}%
\definecolor{currentstroke}{rgb}{0.690196,0.690196,0.690196}%
\pgfsetstrokecolor{currentstroke}%
\pgfsetstrokeopacity{0.300000}%
\pgfsetdash{}{0pt}%
\pgfpathmoveto{\pgfqpoint{4.178131in}{2.333126in}}%
\pgfpathlineto{\pgfqpoint{6.923821in}{2.333126in}}%
\pgfusepath{stroke}%
\end{pgfscope}%
\begin{pgfscope}%
\pgfsetbuttcap%
\pgfsetroundjoin%
\definecolor{currentfill}{rgb}{0.000000,0.000000,0.000000}%
\pgfsetfillcolor{currentfill}%
\pgfsetlinewidth{0.602250pt}%
\definecolor{currentstroke}{rgb}{0.000000,0.000000,0.000000}%
\pgfsetstrokecolor{currentstroke}%
\pgfsetdash{}{0pt}%
\pgfsys@defobject{currentmarker}{\pgfqpoint{-0.027778in}{0.000000in}}{\pgfqpoint{-0.000000in}{0.000000in}}{%
\pgfpathmoveto{\pgfqpoint{-0.000000in}{0.000000in}}%
\pgfpathlineto{\pgfqpoint{-0.027778in}{0.000000in}}%
\pgfusepath{stroke,fill}%
}%
\begin{pgfscope}%
\pgfsys@transformshift{4.178131in}{2.333126in}%
\pgfsys@useobject{currentmarker}{}%
\end{pgfscope}%
\end{pgfscope}%
\begin{pgfscope}%
\pgfpathrectangle{\pgfqpoint{4.178131in}{0.829139in}}{\pgfqpoint{2.745690in}{1.921500in}}%
\pgfusepath{clip}%
\pgfsetrectcap%
\pgfsetroundjoin%
\pgfsetlinewidth{0.803000pt}%
\definecolor{currentstroke}{rgb}{0.690196,0.690196,0.690196}%
\pgfsetstrokecolor{currentstroke}%
\pgfsetstrokeopacity{0.300000}%
\pgfsetdash{}{0pt}%
\pgfpathmoveto{\pgfqpoint{4.178131in}{2.350311in}}%
\pgfpathlineto{\pgfqpoint{6.923821in}{2.350311in}}%
\pgfusepath{stroke}%
\end{pgfscope}%
\begin{pgfscope}%
\pgfsetbuttcap%
\pgfsetroundjoin%
\definecolor{currentfill}{rgb}{0.000000,0.000000,0.000000}%
\pgfsetfillcolor{currentfill}%
\pgfsetlinewidth{0.602250pt}%
\definecolor{currentstroke}{rgb}{0.000000,0.000000,0.000000}%
\pgfsetstrokecolor{currentstroke}%
\pgfsetdash{}{0pt}%
\pgfsys@defobject{currentmarker}{\pgfqpoint{-0.027778in}{0.000000in}}{\pgfqpoint{-0.000000in}{0.000000in}}{%
\pgfpathmoveto{\pgfqpoint{-0.000000in}{0.000000in}}%
\pgfpathlineto{\pgfqpoint{-0.027778in}{0.000000in}}%
\pgfusepath{stroke,fill}%
}%
\begin{pgfscope}%
\pgfsys@transformshift{4.178131in}{2.350311in}%
\pgfsys@useobject{currentmarker}{}%
\end{pgfscope}%
\end{pgfscope}%
\begin{pgfscope}%
\pgfpathrectangle{\pgfqpoint{4.178131in}{0.829139in}}{\pgfqpoint{2.745690in}{1.921500in}}%
\pgfusepath{clip}%
\pgfsetrectcap%
\pgfsetroundjoin%
\pgfsetlinewidth{0.803000pt}%
\definecolor{currentstroke}{rgb}{0.690196,0.690196,0.690196}%
\pgfsetstrokecolor{currentstroke}%
\pgfsetstrokeopacity{0.300000}%
\pgfsetdash{}{0pt}%
\pgfpathmoveto{\pgfqpoint{4.178131in}{2.365197in}}%
\pgfpathlineto{\pgfqpoint{6.923821in}{2.365197in}}%
\pgfusepath{stroke}%
\end{pgfscope}%
\begin{pgfscope}%
\pgfsetbuttcap%
\pgfsetroundjoin%
\definecolor{currentfill}{rgb}{0.000000,0.000000,0.000000}%
\pgfsetfillcolor{currentfill}%
\pgfsetlinewidth{0.602250pt}%
\definecolor{currentstroke}{rgb}{0.000000,0.000000,0.000000}%
\pgfsetstrokecolor{currentstroke}%
\pgfsetdash{}{0pt}%
\pgfsys@defobject{currentmarker}{\pgfqpoint{-0.027778in}{0.000000in}}{\pgfqpoint{-0.000000in}{0.000000in}}{%
\pgfpathmoveto{\pgfqpoint{-0.000000in}{0.000000in}}%
\pgfpathlineto{\pgfqpoint{-0.027778in}{0.000000in}}%
\pgfusepath{stroke,fill}%
}%
\begin{pgfscope}%
\pgfsys@transformshift{4.178131in}{2.365197in}%
\pgfsys@useobject{currentmarker}{}%
\end{pgfscope}%
\end{pgfscope}%
\begin{pgfscope}%
\pgfpathrectangle{\pgfqpoint{4.178131in}{0.829139in}}{\pgfqpoint{2.745690in}{1.921500in}}%
\pgfusepath{clip}%
\pgfsetrectcap%
\pgfsetroundjoin%
\pgfsetlinewidth{0.803000pt}%
\definecolor{currentstroke}{rgb}{0.690196,0.690196,0.690196}%
\pgfsetstrokecolor{currentstroke}%
\pgfsetstrokeopacity{0.300000}%
\pgfsetdash{}{0pt}%
\pgfpathmoveto{\pgfqpoint{4.178131in}{2.378328in}}%
\pgfpathlineto{\pgfqpoint{6.923821in}{2.378328in}}%
\pgfusepath{stroke}%
\end{pgfscope}%
\begin{pgfscope}%
\pgfsetbuttcap%
\pgfsetroundjoin%
\definecolor{currentfill}{rgb}{0.000000,0.000000,0.000000}%
\pgfsetfillcolor{currentfill}%
\pgfsetlinewidth{0.602250pt}%
\definecolor{currentstroke}{rgb}{0.000000,0.000000,0.000000}%
\pgfsetstrokecolor{currentstroke}%
\pgfsetdash{}{0pt}%
\pgfsys@defobject{currentmarker}{\pgfqpoint{-0.027778in}{0.000000in}}{\pgfqpoint{-0.000000in}{0.000000in}}{%
\pgfpathmoveto{\pgfqpoint{-0.000000in}{0.000000in}}%
\pgfpathlineto{\pgfqpoint{-0.027778in}{0.000000in}}%
\pgfusepath{stroke,fill}%
}%
\begin{pgfscope}%
\pgfsys@transformshift{4.178131in}{2.378328in}%
\pgfsys@useobject{currentmarker}{}%
\end{pgfscope}%
\end{pgfscope}%
\begin{pgfscope}%
\pgfpathrectangle{\pgfqpoint{4.178131in}{0.829139in}}{\pgfqpoint{2.745690in}{1.921500in}}%
\pgfusepath{clip}%
\pgfsetrectcap%
\pgfsetroundjoin%
\pgfsetlinewidth{0.803000pt}%
\definecolor{currentstroke}{rgb}{0.690196,0.690196,0.690196}%
\pgfsetstrokecolor{currentstroke}%
\pgfsetstrokeopacity{0.300000}%
\pgfsetdash{}{0pt}%
\pgfpathmoveto{\pgfqpoint{4.178131in}{2.467347in}}%
\pgfpathlineto{\pgfqpoint{6.923821in}{2.467347in}}%
\pgfusepath{stroke}%
\end{pgfscope}%
\begin{pgfscope}%
\pgfsetbuttcap%
\pgfsetroundjoin%
\definecolor{currentfill}{rgb}{0.000000,0.000000,0.000000}%
\pgfsetfillcolor{currentfill}%
\pgfsetlinewidth{0.602250pt}%
\definecolor{currentstroke}{rgb}{0.000000,0.000000,0.000000}%
\pgfsetstrokecolor{currentstroke}%
\pgfsetdash{}{0pt}%
\pgfsys@defobject{currentmarker}{\pgfqpoint{-0.027778in}{0.000000in}}{\pgfqpoint{-0.000000in}{0.000000in}}{%
\pgfpathmoveto{\pgfqpoint{-0.000000in}{0.000000in}}%
\pgfpathlineto{\pgfqpoint{-0.027778in}{0.000000in}}%
\pgfusepath{stroke,fill}%
}%
\begin{pgfscope}%
\pgfsys@transformshift{4.178131in}{2.467347in}%
\pgfsys@useobject{currentmarker}{}%
\end{pgfscope}%
\end{pgfscope}%
\begin{pgfscope}%
\pgfpathrectangle{\pgfqpoint{4.178131in}{0.829139in}}{\pgfqpoint{2.745690in}{1.921500in}}%
\pgfusepath{clip}%
\pgfsetrectcap%
\pgfsetroundjoin%
\pgfsetlinewidth{0.803000pt}%
\definecolor{currentstroke}{rgb}{0.690196,0.690196,0.690196}%
\pgfsetstrokecolor{currentstroke}%
\pgfsetstrokeopacity{0.300000}%
\pgfsetdash{}{0pt}%
\pgfpathmoveto{\pgfqpoint{4.178131in}{2.512549in}}%
\pgfpathlineto{\pgfqpoint{6.923821in}{2.512549in}}%
\pgfusepath{stroke}%
\end{pgfscope}%
\begin{pgfscope}%
\pgfsetbuttcap%
\pgfsetroundjoin%
\definecolor{currentfill}{rgb}{0.000000,0.000000,0.000000}%
\pgfsetfillcolor{currentfill}%
\pgfsetlinewidth{0.602250pt}%
\definecolor{currentstroke}{rgb}{0.000000,0.000000,0.000000}%
\pgfsetstrokecolor{currentstroke}%
\pgfsetdash{}{0pt}%
\pgfsys@defobject{currentmarker}{\pgfqpoint{-0.027778in}{0.000000in}}{\pgfqpoint{-0.000000in}{0.000000in}}{%
\pgfpathmoveto{\pgfqpoint{-0.000000in}{0.000000in}}%
\pgfpathlineto{\pgfqpoint{-0.027778in}{0.000000in}}%
\pgfusepath{stroke,fill}%
}%
\begin{pgfscope}%
\pgfsys@transformshift{4.178131in}{2.512549in}%
\pgfsys@useobject{currentmarker}{}%
\end{pgfscope}%
\end{pgfscope}%
\begin{pgfscope}%
\pgfpathrectangle{\pgfqpoint{4.178131in}{0.829139in}}{\pgfqpoint{2.745690in}{1.921500in}}%
\pgfusepath{clip}%
\pgfsetrectcap%
\pgfsetroundjoin%
\pgfsetlinewidth{0.803000pt}%
\definecolor{currentstroke}{rgb}{0.690196,0.690196,0.690196}%
\pgfsetstrokecolor{currentstroke}%
\pgfsetstrokeopacity{0.300000}%
\pgfsetdash{}{0pt}%
\pgfpathmoveto{\pgfqpoint{4.178131in}{2.544621in}}%
\pgfpathlineto{\pgfqpoint{6.923821in}{2.544621in}}%
\pgfusepath{stroke}%
\end{pgfscope}%
\begin{pgfscope}%
\pgfsetbuttcap%
\pgfsetroundjoin%
\definecolor{currentfill}{rgb}{0.000000,0.000000,0.000000}%
\pgfsetfillcolor{currentfill}%
\pgfsetlinewidth{0.602250pt}%
\definecolor{currentstroke}{rgb}{0.000000,0.000000,0.000000}%
\pgfsetstrokecolor{currentstroke}%
\pgfsetdash{}{0pt}%
\pgfsys@defobject{currentmarker}{\pgfqpoint{-0.027778in}{0.000000in}}{\pgfqpoint{-0.000000in}{0.000000in}}{%
\pgfpathmoveto{\pgfqpoint{-0.000000in}{0.000000in}}%
\pgfpathlineto{\pgfqpoint{-0.027778in}{0.000000in}}%
\pgfusepath{stroke,fill}%
}%
\begin{pgfscope}%
\pgfsys@transformshift{4.178131in}{2.544621in}%
\pgfsys@useobject{currentmarker}{}%
\end{pgfscope}%
\end{pgfscope}%
\begin{pgfscope}%
\pgfpathrectangle{\pgfqpoint{4.178131in}{0.829139in}}{\pgfqpoint{2.745690in}{1.921500in}}%
\pgfusepath{clip}%
\pgfsetrectcap%
\pgfsetroundjoin%
\pgfsetlinewidth{0.803000pt}%
\definecolor{currentstroke}{rgb}{0.690196,0.690196,0.690196}%
\pgfsetstrokecolor{currentstroke}%
\pgfsetstrokeopacity{0.300000}%
\pgfsetdash{}{0pt}%
\pgfpathmoveto{\pgfqpoint{4.178131in}{2.569497in}}%
\pgfpathlineto{\pgfqpoint{6.923821in}{2.569497in}}%
\pgfusepath{stroke}%
\end{pgfscope}%
\begin{pgfscope}%
\pgfsetbuttcap%
\pgfsetroundjoin%
\definecolor{currentfill}{rgb}{0.000000,0.000000,0.000000}%
\pgfsetfillcolor{currentfill}%
\pgfsetlinewidth{0.602250pt}%
\definecolor{currentstroke}{rgb}{0.000000,0.000000,0.000000}%
\pgfsetstrokecolor{currentstroke}%
\pgfsetdash{}{0pt}%
\pgfsys@defobject{currentmarker}{\pgfqpoint{-0.027778in}{0.000000in}}{\pgfqpoint{-0.000000in}{0.000000in}}{%
\pgfpathmoveto{\pgfqpoint{-0.000000in}{0.000000in}}%
\pgfpathlineto{\pgfqpoint{-0.027778in}{0.000000in}}%
\pgfusepath{stroke,fill}%
}%
\begin{pgfscope}%
\pgfsys@transformshift{4.178131in}{2.569497in}%
\pgfsys@useobject{currentmarker}{}%
\end{pgfscope}%
\end{pgfscope}%
\begin{pgfscope}%
\pgfpathrectangle{\pgfqpoint{4.178131in}{0.829139in}}{\pgfqpoint{2.745690in}{1.921500in}}%
\pgfusepath{clip}%
\pgfsetrectcap%
\pgfsetroundjoin%
\pgfsetlinewidth{0.803000pt}%
\definecolor{currentstroke}{rgb}{0.690196,0.690196,0.690196}%
\pgfsetstrokecolor{currentstroke}%
\pgfsetstrokeopacity{0.300000}%
\pgfsetdash{}{0pt}%
\pgfpathmoveto{\pgfqpoint{4.178131in}{2.589823in}}%
\pgfpathlineto{\pgfqpoint{6.923821in}{2.589823in}}%
\pgfusepath{stroke}%
\end{pgfscope}%
\begin{pgfscope}%
\pgfsetbuttcap%
\pgfsetroundjoin%
\definecolor{currentfill}{rgb}{0.000000,0.000000,0.000000}%
\pgfsetfillcolor{currentfill}%
\pgfsetlinewidth{0.602250pt}%
\definecolor{currentstroke}{rgb}{0.000000,0.000000,0.000000}%
\pgfsetstrokecolor{currentstroke}%
\pgfsetdash{}{0pt}%
\pgfsys@defobject{currentmarker}{\pgfqpoint{-0.027778in}{0.000000in}}{\pgfqpoint{-0.000000in}{0.000000in}}{%
\pgfpathmoveto{\pgfqpoint{-0.000000in}{0.000000in}}%
\pgfpathlineto{\pgfqpoint{-0.027778in}{0.000000in}}%
\pgfusepath{stroke,fill}%
}%
\begin{pgfscope}%
\pgfsys@transformshift{4.178131in}{2.589823in}%
\pgfsys@useobject{currentmarker}{}%
\end{pgfscope}%
\end{pgfscope}%
\begin{pgfscope}%
\pgfpathrectangle{\pgfqpoint{4.178131in}{0.829139in}}{\pgfqpoint{2.745690in}{1.921500in}}%
\pgfusepath{clip}%
\pgfsetrectcap%
\pgfsetroundjoin%
\pgfsetlinewidth{0.803000pt}%
\definecolor{currentstroke}{rgb}{0.690196,0.690196,0.690196}%
\pgfsetstrokecolor{currentstroke}%
\pgfsetstrokeopacity{0.300000}%
\pgfsetdash{}{0pt}%
\pgfpathmoveto{\pgfqpoint{4.178131in}{2.607008in}}%
\pgfpathlineto{\pgfqpoint{6.923821in}{2.607008in}}%
\pgfusepath{stroke}%
\end{pgfscope}%
\begin{pgfscope}%
\pgfsetbuttcap%
\pgfsetroundjoin%
\definecolor{currentfill}{rgb}{0.000000,0.000000,0.000000}%
\pgfsetfillcolor{currentfill}%
\pgfsetlinewidth{0.602250pt}%
\definecolor{currentstroke}{rgb}{0.000000,0.000000,0.000000}%
\pgfsetstrokecolor{currentstroke}%
\pgfsetdash{}{0pt}%
\pgfsys@defobject{currentmarker}{\pgfqpoint{-0.027778in}{0.000000in}}{\pgfqpoint{-0.000000in}{0.000000in}}{%
\pgfpathmoveto{\pgfqpoint{-0.000000in}{0.000000in}}%
\pgfpathlineto{\pgfqpoint{-0.027778in}{0.000000in}}%
\pgfusepath{stroke,fill}%
}%
\begin{pgfscope}%
\pgfsys@transformshift{4.178131in}{2.607008in}%
\pgfsys@useobject{currentmarker}{}%
\end{pgfscope}%
\end{pgfscope}%
\begin{pgfscope}%
\pgfpathrectangle{\pgfqpoint{4.178131in}{0.829139in}}{\pgfqpoint{2.745690in}{1.921500in}}%
\pgfusepath{clip}%
\pgfsetrectcap%
\pgfsetroundjoin%
\pgfsetlinewidth{0.803000pt}%
\definecolor{currentstroke}{rgb}{0.690196,0.690196,0.690196}%
\pgfsetstrokecolor{currentstroke}%
\pgfsetstrokeopacity{0.300000}%
\pgfsetdash{}{0pt}%
\pgfpathmoveto{\pgfqpoint{4.178131in}{2.621894in}}%
\pgfpathlineto{\pgfqpoint{6.923821in}{2.621894in}}%
\pgfusepath{stroke}%
\end{pgfscope}%
\begin{pgfscope}%
\pgfsetbuttcap%
\pgfsetroundjoin%
\definecolor{currentfill}{rgb}{0.000000,0.000000,0.000000}%
\pgfsetfillcolor{currentfill}%
\pgfsetlinewidth{0.602250pt}%
\definecolor{currentstroke}{rgb}{0.000000,0.000000,0.000000}%
\pgfsetstrokecolor{currentstroke}%
\pgfsetdash{}{0pt}%
\pgfsys@defobject{currentmarker}{\pgfqpoint{-0.027778in}{0.000000in}}{\pgfqpoint{-0.000000in}{0.000000in}}{%
\pgfpathmoveto{\pgfqpoint{-0.000000in}{0.000000in}}%
\pgfpathlineto{\pgfqpoint{-0.027778in}{0.000000in}}%
\pgfusepath{stroke,fill}%
}%
\begin{pgfscope}%
\pgfsys@transformshift{4.178131in}{2.621894in}%
\pgfsys@useobject{currentmarker}{}%
\end{pgfscope}%
\end{pgfscope}%
\begin{pgfscope}%
\pgfpathrectangle{\pgfqpoint{4.178131in}{0.829139in}}{\pgfqpoint{2.745690in}{1.921500in}}%
\pgfusepath{clip}%
\pgfsetrectcap%
\pgfsetroundjoin%
\pgfsetlinewidth{0.803000pt}%
\definecolor{currentstroke}{rgb}{0.690196,0.690196,0.690196}%
\pgfsetstrokecolor{currentstroke}%
\pgfsetstrokeopacity{0.300000}%
\pgfsetdash{}{0pt}%
\pgfpathmoveto{\pgfqpoint{4.178131in}{2.635025in}}%
\pgfpathlineto{\pgfqpoint{6.923821in}{2.635025in}}%
\pgfusepath{stroke}%
\end{pgfscope}%
\begin{pgfscope}%
\pgfsetbuttcap%
\pgfsetroundjoin%
\definecolor{currentfill}{rgb}{0.000000,0.000000,0.000000}%
\pgfsetfillcolor{currentfill}%
\pgfsetlinewidth{0.602250pt}%
\definecolor{currentstroke}{rgb}{0.000000,0.000000,0.000000}%
\pgfsetstrokecolor{currentstroke}%
\pgfsetdash{}{0pt}%
\pgfsys@defobject{currentmarker}{\pgfqpoint{-0.027778in}{0.000000in}}{\pgfqpoint{-0.000000in}{0.000000in}}{%
\pgfpathmoveto{\pgfqpoint{-0.000000in}{0.000000in}}%
\pgfpathlineto{\pgfqpoint{-0.027778in}{0.000000in}}%
\pgfusepath{stroke,fill}%
}%
\begin{pgfscope}%
\pgfsys@transformshift{4.178131in}{2.635025in}%
\pgfsys@useobject{currentmarker}{}%
\end{pgfscope}%
\end{pgfscope}%
\begin{pgfscope}%
\pgfpathrectangle{\pgfqpoint{4.178131in}{0.829139in}}{\pgfqpoint{2.745690in}{1.921500in}}%
\pgfusepath{clip}%
\pgfsetrectcap%
\pgfsetroundjoin%
\pgfsetlinewidth{0.803000pt}%
\definecolor{currentstroke}{rgb}{0.690196,0.690196,0.690196}%
\pgfsetstrokecolor{currentstroke}%
\pgfsetstrokeopacity{0.300000}%
\pgfsetdash{}{0pt}%
\pgfpathmoveto{\pgfqpoint{4.178131in}{2.724045in}}%
\pgfpathlineto{\pgfqpoint{6.923821in}{2.724045in}}%
\pgfusepath{stroke}%
\end{pgfscope}%
\begin{pgfscope}%
\pgfsetbuttcap%
\pgfsetroundjoin%
\definecolor{currentfill}{rgb}{0.000000,0.000000,0.000000}%
\pgfsetfillcolor{currentfill}%
\pgfsetlinewidth{0.602250pt}%
\definecolor{currentstroke}{rgb}{0.000000,0.000000,0.000000}%
\pgfsetstrokecolor{currentstroke}%
\pgfsetdash{}{0pt}%
\pgfsys@defobject{currentmarker}{\pgfqpoint{-0.027778in}{0.000000in}}{\pgfqpoint{-0.000000in}{0.000000in}}{%
\pgfpathmoveto{\pgfqpoint{-0.000000in}{0.000000in}}%
\pgfpathlineto{\pgfqpoint{-0.027778in}{0.000000in}}%
\pgfusepath{stroke,fill}%
}%
\begin{pgfscope}%
\pgfsys@transformshift{4.178131in}{2.724045in}%
\pgfsys@useobject{currentmarker}{}%
\end{pgfscope}%
\end{pgfscope}%
\begin{pgfscope}%
\definecolor{textcolor}{rgb}{0.000000,0.000000,0.000000}%
\pgfsetstrokecolor{textcolor}%
\pgfsetfillcolor{textcolor}%
\pgftext[x=3.785189in,y=1.789889in,,bottom,rotate=90.000000]{\color{textcolor}{\rmfamily\fontsize{8.000000}{9.600000}\selectfont\catcode`\^=\active\def^{\ifmmode\sp\else\^{}\fi}\catcode`\%=\active\def%{\%}Unlogged Power $\widehat A_w(t)-1$}}%
\end{pgfscope}%
\begin{pgfscope}%
\pgfpathrectangle{\pgfqpoint{4.178131in}{0.829139in}}{\pgfqpoint{2.745690in}{1.921500in}}%
\pgfusepath{clip}%
\pgfsetrectcap%
\pgfsetroundjoin%
\pgfsetlinewidth{1.706375pt}%
\definecolor{currentstroke}{rgb}{0.121569,0.466667,0.705882}%
\pgfsetstrokecolor{currentstroke}%
\pgfsetdash{}{0pt}%
\pgfpathmoveto{\pgfqpoint{4.302935in}{2.663298in}}%
\pgfpathlineto{\pgfqpoint{4.304091in}{2.631099in}}%
\pgfpathlineto{\pgfqpoint{4.305562in}{2.600559in}}%
\pgfpathlineto{\pgfqpoint{4.307253in}{2.573426in}}%
\pgfpathlineto{\pgfqpoint{4.308970in}{2.551238in}}%
\pgfpathlineto{\pgfqpoint{4.310761in}{2.531926in}}%
\pgfpathlineto{\pgfqpoint{4.313030in}{2.511340in}}%
\pgfpathlineto{\pgfqpoint{4.315121in}{2.495131in}}%
\pgfpathlineto{\pgfqpoint{4.317619in}{2.478316in}}%
\pgfpathlineto{\pgfqpoint{4.320603in}{2.460936in}}%
\pgfpathlineto{\pgfqpoint{4.324166in}{2.443028in}}%
\pgfpathlineto{\pgfqpoint{4.326919in}{2.430811in}}%
\pgfpathlineto{\pgfqpoint{4.330017in}{2.418379in}}%
\pgfpathlineto{\pgfqpoint{4.333505in}{2.405737in}}%
\pgfpathlineto{\pgfqpoint{4.337432in}{2.392888in}}%
\pgfpathlineto{\pgfqpoint{4.341852in}{2.379830in}}%
\pgfpathlineto{\pgfqpoint{4.346828in}{2.366563in}}%
\pgfpathlineto{\pgfqpoint{4.352429in}{2.353080in}}%
\pgfpathlineto{\pgfqpoint{4.358734in}{2.339373in}}%
\pgfpathlineto{\pgfqpoint{4.365832in}{2.325431in}}%
\pgfpathlineto{\pgfqpoint{4.373822in}{2.311237in}}%
\pgfpathlineto{\pgfqpoint{4.378186in}{2.304040in}}%
\pgfpathlineto{\pgfqpoint{4.382817in}{2.296772in}}%
\pgfpathlineto{\pgfqpoint{4.387730in}{2.289429in}}%
\pgfpathlineto{\pgfqpoint{4.392942in}{2.282008in}}%
\pgfpathlineto{\pgfqpoint{4.398473in}{2.274505in}}%
\pgfpathlineto{\pgfqpoint{4.404340in}{2.266916in}}%
\pgfpathlineto{\pgfqpoint{4.410566in}{2.259234in}}%
\pgfpathlineto{\pgfqpoint{4.417171in}{2.251454in}}%
\pgfpathlineto{\pgfqpoint{4.424179in}{2.243571in}}%
\pgfpathlineto{\pgfqpoint{4.431615in}{2.235577in}}%
\pgfpathlineto{\pgfqpoint{4.439504in}{2.227464in}}%
\pgfpathlineto{\pgfqpoint{4.447874in}{2.219224in}}%
\pgfpathlineto{\pgfqpoint{4.456755in}{2.210848in}}%
\pgfpathlineto{\pgfqpoint{4.466178in}{2.202325in}}%
\pgfpathlineto{\pgfqpoint{4.476175in}{2.193644in}}%
\pgfpathlineto{\pgfqpoint{4.486782in}{2.184791in}}%
\pgfpathlineto{\pgfqpoint{4.498036in}{2.175753in}}%
\pgfpathlineto{\pgfqpoint{4.509976in}{2.166513in}}%
\pgfpathlineto{\pgfqpoint{4.522645in}{2.157054in}}%
\pgfpathlineto{\pgfqpoint{4.536086in}{2.147356in}}%
\pgfpathlineto{\pgfqpoint{4.550348in}{2.137397in}}%
\pgfpathlineto{\pgfqpoint{4.565479in}{2.127152in}}%
\pgfpathlineto{\pgfqpoint{4.581533in}{2.116593in}}%
\pgfpathlineto{\pgfqpoint{4.598566in}{2.105690in}}%
\pgfpathlineto{\pgfqpoint{4.616638in}{2.094409in}}%
\pgfpathlineto{\pgfqpoint{4.635812in}{2.082712in}}%
\pgfpathlineto{\pgfqpoint{4.656155in}{2.070557in}}%
\pgfpathlineto{\pgfqpoint{4.677740in}{2.057898in}}%
\pgfpathlineto{\pgfqpoint{4.724939in}{2.030868in}}%
\pgfpathlineto{\pgfqpoint{4.778071in}{2.001177in}}%
\pgfpathlineto{\pgfqpoint{4.837882in}{1.968324in}}%
\pgfpathlineto{\pgfqpoint{4.905211in}{1.931751in}}%
\pgfpathlineto{\pgfqpoint{5.022403in}{1.868588in}}%
\pgfpathlineto{\pgfqpoint{5.270492in}{1.735498in}}%
\pgfpathlineto{\pgfqpoint{6.799017in}{0.916480in}}%
\pgfpathlineto{\pgfqpoint{6.799017in}{0.916480in}}%
\pgfusepath{stroke}%
\end{pgfscope}%
\begin{pgfscope}%
\pgfpathrectangle{\pgfqpoint{4.178131in}{0.829139in}}{\pgfqpoint{2.745690in}{1.921500in}}%
\pgfusepath{clip}%
\pgfsetbuttcap%
\pgfsetroundjoin%
\pgfsetlinewidth{1.505625pt}%
\definecolor{currentstroke}{rgb}{1.000000,0.498039,0.054902}%
\pgfsetstrokecolor{currentstroke}%
\pgfsetdash{{5.550000pt}{2.400000pt}}{0.000000pt}%
\pgfpathmoveto{\pgfqpoint{4.302935in}{2.253910in}}%
\pgfpathlineto{\pgfqpoint{6.799017in}{0.916480in}}%
\pgfpathlineto{\pgfqpoint{6.799017in}{0.916480in}}%
\pgfusepath{stroke}%
\end{pgfscope}%
\begin{pgfscope}%
\pgfpathrectangle{\pgfqpoint{4.178131in}{0.829139in}}{\pgfqpoint{2.745690in}{1.921500in}}%
\pgfusepath{clip}%
\pgfsetbuttcap%
\pgfsetroundjoin%
\pgfsetlinewidth{1.706375pt}%
\definecolor{currentstroke}{rgb}{0.172549,0.627451,0.172549}%
\pgfsetstrokecolor{currentstroke}%
\pgfsetdash{{1.700000pt}{2.805000pt}}{0.000000pt}%
\pgfpathmoveto{\pgfqpoint{4.302935in}{2.360843in}}%
\pgfpathlineto{\pgfqpoint{4.317619in}{2.343535in}}%
\pgfpathlineto{\pgfqpoint{4.331709in}{2.327433in}}%
\pgfpathlineto{\pgfqpoint{4.346828in}{2.310720in}}%
\pgfpathlineto{\pgfqpoint{4.362178in}{2.294351in}}%
\pgfpathlineto{\pgfqpoint{4.378186in}{2.277925in}}%
\pgfpathlineto{\pgfqpoint{4.392942in}{2.263359in}}%
\pgfpathlineto{\pgfqpoint{4.404340in}{2.252479in}}%
\pgfpathlineto{\pgfqpoint{4.417171in}{2.240606in}}%
\pgfpathlineto{\pgfqpoint{4.431615in}{2.227701in}}%
\pgfpathlineto{\pgfqpoint{4.447874in}{2.213730in}}%
\pgfpathlineto{\pgfqpoint{4.466178in}{2.198665in}}%
\pgfpathlineto{\pgfqpoint{4.486782in}{2.182480in}}%
\pgfpathlineto{\pgfqpoint{4.509976in}{2.165141in}}%
\pgfpathlineto{\pgfqpoint{4.522645in}{2.156024in}}%
\pgfpathlineto{\pgfqpoint{4.536086in}{2.146598in}}%
\pgfpathlineto{\pgfqpoint{4.550348in}{2.136849in}}%
\pgfpathlineto{\pgfqpoint{4.565479in}{2.126765in}}%
\pgfpathlineto{\pgfqpoint{4.581533in}{2.116326in}}%
\pgfpathlineto{\pgfqpoint{4.598566in}{2.105511in}}%
\pgfpathlineto{\pgfqpoint{4.616638in}{2.094291in}}%
\pgfpathlineto{\pgfqpoint{4.635812in}{2.082637in}}%
\pgfpathlineto{\pgfqpoint{4.656155in}{2.070510in}}%
\pgfpathlineto{\pgfqpoint{4.700641in}{2.044670in}}%
\pgfpathlineto{\pgfqpoint{4.750719in}{2.016380in}}%
\pgfpathlineto{\pgfqpoint{4.807091in}{1.985178in}}%
\pgfpathlineto{\pgfqpoint{4.870550in}{1.950538in}}%
\pgfpathlineto{\pgfqpoint{4.981005in}{1.890857in}}%
\pgfpathlineto{\pgfqpoint{5.162372in}{1.793460in}}%
\pgfpathlineto{\pgfqpoint{5.857071in}{1.421186in}}%
\pgfpathlineto{\pgfqpoint{6.799017in}{0.916480in}}%
\pgfpathlineto{\pgfqpoint{6.799017in}{0.916480in}}%
\pgfusepath{stroke}%
\end{pgfscope}%
\begin{pgfscope}%
\pgfsetrectcap%
\pgfsetmiterjoin%
\pgfsetlinewidth{0.803000pt}%
\definecolor{currentstroke}{rgb}{0.000000,0.000000,0.000000}%
\pgfsetstrokecolor{currentstroke}%
\pgfsetdash{}{0pt}%
\pgfpathmoveto{\pgfqpoint{4.178131in}{0.829139in}}%
\pgfpathlineto{\pgfqpoint{4.178131in}{2.750639in}}%
\pgfusepath{stroke}%
\end{pgfscope}%
\begin{pgfscope}%
\pgfsetrectcap%
\pgfsetmiterjoin%
\pgfsetlinewidth{0.803000pt}%
\definecolor{currentstroke}{rgb}{0.000000,0.000000,0.000000}%
\pgfsetstrokecolor{currentstroke}%
\pgfsetdash{}{0pt}%
\pgfpathmoveto{\pgfqpoint{6.923821in}{0.829139in}}%
\pgfpathlineto{\pgfqpoint{6.923821in}{2.750639in}}%
\pgfusepath{stroke}%
\end{pgfscope}%
\begin{pgfscope}%
\pgfsetrectcap%
\pgfsetmiterjoin%
\pgfsetlinewidth{0.803000pt}%
\definecolor{currentstroke}{rgb}{0.000000,0.000000,0.000000}%
\pgfsetstrokecolor{currentstroke}%
\pgfsetdash{}{0pt}%
\pgfpathmoveto{\pgfqpoint{4.178131in}{0.829139in}}%
\pgfpathlineto{\pgfqpoint{6.923821in}{0.829139in}}%
\pgfusepath{stroke}%
\end{pgfscope}%
\begin{pgfscope}%
\pgfsetrectcap%
\pgfsetmiterjoin%
\pgfsetlinewidth{0.803000pt}%
\definecolor{currentstroke}{rgb}{0.000000,0.000000,0.000000}%
\pgfsetstrokecolor{currentstroke}%
\pgfsetdash{}{0pt}%
\pgfpathmoveto{\pgfqpoint{4.178131in}{2.750639in}}%
\pgfpathlineto{\pgfqpoint{6.923821in}{2.750639in}}%
\pgfusepath{stroke}%
\end{pgfscope}%
\begin{pgfscope}%
\definecolor{textcolor}{rgb}{0.000000,0.000000,0.000000}%
\pgfsetstrokecolor{textcolor}%
\pgfsetfillcolor{textcolor}%
\pgftext[x=5.550976in,y=2.833972in,,base]{\color{textcolor}{\rmfamily\fontsize{8.000000}{9.600000}\selectfont\catcode`\^=\active\def^{\ifmmode\sp\else\^{}\fi}\catcode`\%=\active\def%{\%}(b) Low-Frequency Decay}}%
\end{pgfscope}%
\begin{pgfscope}%
\pgfsetrectcap%
\pgfsetroundjoin%
\pgfsetlinewidth{1.706375pt}%
\definecolor{currentstroke}{rgb}{0.121569,0.466667,0.705882}%
\pgfsetstrokecolor{currentstroke}%
\pgfsetdash{}{0pt}%
\pgfpathmoveto{\pgfqpoint{1.282119in}{0.197222in}}%
\pgfpathlineto{\pgfqpoint{1.369619in}{0.197222in}}%
\pgfpathlineto{\pgfqpoint{1.457119in}{0.197222in}}%
\pgfusepath{stroke}%
\end{pgfscope}%
\begin{pgfscope}%
\definecolor{textcolor}{rgb}{0.000000,0.000000,0.000000}%
\pgfsetstrokecolor{textcolor}%
\pgfsetfillcolor{textcolor}%
\pgftext[x=1.534897in,y=0.163194in,left,base]{\color{textcolor}{\rmfamily\fontsize{7.000000}{8.400000}\selectfont\catcode`\^=\active\def^{\ifmmode\sp\else\^{}\fi}\catcode`\%=\active\def%{\%}$\widehat A_w(t)$}}%
\end{pgfscope}%
\begin{pgfscope}%
\pgfsetbuttcap%
\pgfsetroundjoin%
\pgfsetlinewidth{1.505625pt}%
\definecolor{currentstroke}{rgb}{1.000000,0.498039,0.054902}%
\pgfsetstrokecolor{currentstroke}%
\pgfsetdash{{5.550000pt}{2.400000pt}}{0.000000pt}%
\pgfpathmoveto{\pgfqpoint{1.942734in}{0.197222in}}%
\pgfpathlineto{\pgfqpoint{2.030234in}{0.197222in}}%
\pgfpathlineto{\pgfqpoint{2.117734in}{0.197222in}}%
\pgfusepath{stroke}%
\end{pgfscope}%
\begin{pgfscope}%
\definecolor{textcolor}{rgb}{0.000000,0.000000,0.000000}%
\pgfsetstrokecolor{textcolor}%
\pgfsetfillcolor{textcolor}%
\pgftext[x=2.195511in,y=0.163194in,left,base]{\color{textcolor}{\rmfamily\fontsize{7.000000}{8.400000}\selectfont\catcode`\^=\active\def^{\ifmmode\sp\else\^{}\fi}\catcode`\%=\active\def%{\%}Small-Time Guide}}%
\end{pgfscope}%
\begin{pgfscope}%
\pgfsetrectcap%
\pgfsetroundjoin%
\pgfsetlinewidth{1.706375pt}%
\definecolor{currentstroke}{rgb}{0.121569,0.466667,0.705882}%
\pgfsetstrokecolor{currentstroke}%
\pgfsetdash{}{0pt}%
\pgfpathmoveto{\pgfqpoint{3.198660in}{0.197222in}}%
\pgfpathlineto{\pgfqpoint{3.286160in}{0.197222in}}%
\pgfpathlineto{\pgfqpoint{3.373660in}{0.197222in}}%
\pgfusepath{stroke}%
\end{pgfscope}%
\begin{pgfscope}%
\definecolor{textcolor}{rgb}{0.000000,0.000000,0.000000}%
\pgfsetstrokecolor{textcolor}%
\pgfsetfillcolor{textcolor}%
\pgftext[x=3.451438in,y=0.163194in,left,base]{\color{textcolor}{\rmfamily\fontsize{7.000000}{8.400000}\selectfont\catcode`\^=\active\def^{\ifmmode\sp\else\^{}\fi}\catcode`\%=\active\def%{\%}$\widehat A_w(t)-1$}}%
\end{pgfscope}%
\begin{pgfscope}%
\pgfsetbuttcap%
\pgfsetroundjoin%
\pgfsetlinewidth{1.505625pt}%
\definecolor{currentstroke}{rgb}{1.000000,0.498039,0.054902}%
\pgfsetstrokecolor{currentstroke}%
\pgfsetdash{{5.550000pt}{2.400000pt}}{0.000000pt}%
\pgfpathmoveto{\pgfqpoint{4.052026in}{0.197222in}}%
\pgfpathlineto{\pgfqpoint{4.139526in}{0.197222in}}%
\pgfpathlineto{\pgfqpoint{4.227026in}{0.197222in}}%
\pgfusepath{stroke}%
\end{pgfscope}%
\begin{pgfscope}%
\definecolor{textcolor}{rgb}{0.000000,0.000000,0.000000}%
\pgfsetstrokecolor{textcolor}%
\pgfsetfillcolor{textcolor}%
\pgftext[x=4.304804in,y=0.163194in,left,base]{\color{textcolor}{\rmfamily\fontsize{7.000000}{8.400000}\selectfont\catcode`\^=\active\def^{\ifmmode\sp\else\^{}\fi}\catcode`\%=\active\def%{\%}Degree 1}}%
\end{pgfscope}%
\begin{pgfscope}%
\pgfsetbuttcap%
\pgfsetroundjoin%
\pgfsetlinewidth{1.706375pt}%
\definecolor{currentstroke}{rgb}{0.172549,0.627451,0.172549}%
\pgfsetstrokecolor{currentstroke}%
\pgfsetdash{{1.700000pt}{2.805000pt}}{0.000000pt}%
\pgfpathmoveto{\pgfqpoint{4.844794in}{0.197222in}}%
\pgfpathlineto{\pgfqpoint{4.932294in}{0.197222in}}%
\pgfpathlineto{\pgfqpoint{5.019794in}{0.197222in}}%
\pgfusepath{stroke}%
\end{pgfscope}%
\begin{pgfscope}%
\definecolor{textcolor}{rgb}{0.000000,0.000000,0.000000}%
\pgfsetstrokecolor{textcolor}%
\pgfsetfillcolor{textcolor}%
\pgftext[x=5.097572in,y=0.163194in,left,base]{\color{textcolor}{\rmfamily\fontsize{7.000000}{8.400000}\selectfont\catcode`\^=\active\def^{\ifmmode\sp\else\^{}\fi}\catcode`\%=\active\def%{\%}Degrees 1+2}}%
\end{pgfscope}%
\end{pgfpicture}%
\makeatother%
\endgroup%

\caption{Two-regime validation for a unimodal vMF cloud on $S^2$. Panel (a) compares $\widehat A_w(t)$ to the leading small-time pairwise heat-overlap approximation. Panel (b) compares $\widehat A_w(t)-1$ to the low-frequency approximation. The unlogged power is used because it is the cross-term-free object in the spherical expansion.}
\label{fig:asymptotics}
\end{figure*}



\paragraph{Synthetic configurations.}

Each synthetic example targets a different limitation. The unimodal cloud is a concentrated reference. The antipodal and girdle examples have nearly zero first moments but remain nonuniform. The four-mode example tests sensitivity to modal separation. The duplicate example tests atom-splitting invariance, which ordinary ESS lacks. We compare five equal-weight configurations with $n=300$. These are a vMF cloud, an antipodal mixture, a girdle, four separated modes, and duplicates at 30 quasi-uniform locations. Table~\ref{tab:synthetic} reports selected scales. Ordinary ESS is 300 in every row, while $R$ misses symmetric nonuniformity. The energy $B_2$ detects axial structure in the antipodal and girdle examples. In contrast, gESS gives a scale-dependent effective number. The duplicate example has 300 labels at only 30 locations. At $15^\circ$, its gESS is about 30, consistent with merge invariance and limited overlap between locations.

\begin{table}[t]
\centering
\caption{Synthetic examples on $S^2$. Here $R=\|\sum_iw_ix_i\|$ and $B_2$ is the degree-2/Bingham energy.
The baseline is $D_{\rm LC}^{(15)}$, and $\alpha=\sqrt{8t}$ is the angular scale in degrees.}
\label{tab:synthetic}
\small
\resizebox{\columnwidth}{!}{%
\begin{tabular}{lrrrrrr}
\toprule
Example & $R$ & $B_2$ & ESS & $\gess(15^\circ)$ & $D_{\rm LC}^{(15)}$ & $\gess(45^\circ)$ \\
\midrule
Unimodal vMF & 0.993 & 4.79 & 300 & 1.4 & 1.4 & 1.0 \\
Antipodal & 0.006 & 4.78 & 300 & 2.9 & 2.9 & 2.1 \\
Girdle & 0.029 & 1.24 & 300 & 13.3 & 13.4 & 4.4 \\
Four modes & 0.007 & 0.00 & 300 & 5.5 & 5.5 & 4.1 \\
Duplicates & 0.005 & 0.00 & 300 & 29.7 & 29.7 & 6.8 \\
\bottomrule
\end{tabular}%
}
\end{table}

\paragraph{Real-data embedding illustration.}

For a data example, we use the handwritten digits dataset\footnote{\url{https://doi.org/10.24432/C50P49}} from
\textsf{scikit-learn} \citep{pedregosa_2011_ScikitlearnMachineLearning}. The $8\times8$ images are standardized and projected onto three principal components. We normalize these projections onto $S^2$. The example compares reweighting the same embeddings with restricting them to a subset. Table~\ref{tab:digits} shows that nonuniform reliability-style weights reduce ordinary and geometric ESS. The 0-vs-1 subset also differs from the full cloud in directional structure.

\begin{table}[t]
\centering
\caption{Normalized digit embeddings after PCA to $\R^3$ and projection onto $S^2$. The weighted row uses deterministic weights proportional to $\exp(2.5x_1)$.}
\label{tab:digits}
\small
\resizebox{\columnwidth}{!}{%
\begin{tabular}{lrrrrrrr}
\toprule
Data & $n$ & $R$ & $B_2$ & ESS & $\ess_H$ & $\gess(15^\circ)$ & $D_{\rm LC}^{(15)}$ \\
\midrule
Digits all & 300 & 0.152 & 0.07 & 300 & 300 & 36.7 & 36.8 \\
Digits weighted & 300 & 0.649 & 0.53 & 123 & 156 & 15.4 & 15.5 \\
Digits 0/1 & 60 & 0.540 & 0.81 & 60 & 60 & 9.8 & 9.9 \\
\bottomrule
\end{tabular}%
}
\end{table}


%----%----%----%----%----%----%----%----%----%----%----%----%----%
\section{DISCUSSION}

Heat-kernel entropy profiles summarize weighted empirical measures on compact manifolds. For order-two R\'enyi entropy, the pairwise formula also gives a geometry-aware ESS. Its scale limits do not require a spherical parametrization. The fine-scale gESS limit is the ESS of coalesced support locations. At large times, the expansion identifies structures that persist under diffusion.

The theory assumes a compact, connected, boundaryless manifold with normalized Riemannian volume. Theorem~\ref{thm:pairwise_monotone} proves monotonicity for $\widehat U_2(t)$. It does not prove gESS monotonicity on nonhomogeneous manifolds because the diagonal normalization varies with location. Exact computation also requires heat-kernel evaluations. Harmonic expansions are natural on spheres and compact Lie groups. General manifolds may instead use spectral truncation, graph Laplacians \citep{chung_1997_SpectralGraphTheory}, diffusion maps \citep{coifman_2006_DiffusionMaps}, or numerical heat solvers.

The statistical results require a positive lower scale. Small scales reveal fine structure but have weaker stability because heat-kernel suprema and derivatives grow as $t\downarrow0$. We therefore state consistency on $[t_0,T]$ with $t_0>0$. The experiments describe profile behavior without calibrated inference or confidence bands. Further work could develop full-profile inference and allow heavy-tailed importance weights. Extensions to noncompact manifolds would also require a suitable reference measure. Another direction is joint use with maximum-entropy hierarchies for model-based uncertainty decomposition.

%----%----%----%----%----%----%----%----%----%----%----%----%----%

%\subsubsection*{Acknowledgements}
%All acknowledgments go at the end of the paper, including thanks to reviewers who gave useful comments, to colleagues who contributed to the ideas, and to funding agencies and corporate sponsors that provided financial support.  To preserve the anonymity, please include acknowledgments \emph{only} in the camera-ready papers. The acknowledgements do not count against the 9-page page limit in the camera-ready.




% If you use BibTeX in apalike style, activate the following line:
\bibliographystyle{apalike}
\bibliography{references}



%%%%%%%%%%%%%%%%%%%%%%%%%%%%%%%%%%%%%%%%%%%%%%%%%%%%%%%%%%%%




%\clearpage
\appendix
%\thispagestyle{empty}

% Supplementary material: To improve readability, you must use a single-column format for the supplementary material.
%\onecolumn
%\aistatstitle{Supplementary M}


\section{PROOFS}\label{app:proofs}

\subsection{Proof of Theorem~\ref{thm:pairwise_monotone}}
By bilinearity and the definition of $\widehat p_t$,
\begin{equation*}
\int \widehat p_t(y)^2\dd\nu(y)
=\sum_{i,j}w_iw_j\int k_t(y,x_i)k_t(y,x_j)\dd\nu(y).
\end{equation*}
The heat kernel is symmetric and satisfies the semigroup identity. Hence, the integral is $k_{2t}(x_i,x_j)$, proving \eqref{eq:pairwise_identity}.

For monotonicity, $\widehat p_t$ solves $\partial_t\widehat p_t=\Delta\widehat p_t$. Under our convention, the spectrum of $-\Delta$ is nonnegative. Since $\M$ is compact without boundary,
\begin{align*}
  \frac{d}{dt}\int \widehat p_t^2\dd\nu
  &=2\int \widehat p_t\Delta\widehat p_t\dd\nu
   =-2\int\|\nabla\widehat p_t\|^2\dd\nu\le0.
\end{align*}
Therefore, $\widehat D_2(t)=\log\|\widehat p_t\|_2^2$ is nonincreasing, while $\widehat U_2(t)$ is nondecreasing. Compactness and connectedness imply that heat flow converges in $L^2(\nu)$ to the constant density $1$. Thus, $\|\widehat p_t\|_2^2\to1$.

\subsection{Proof of Theorem~\ref{thm:gess}}
The semigroup identity gives
\begin{align*}
  k_{2t}(x,y)&=\int k_t(x,z)k_t(y,z)\dd\nu(z)\\
  &=\ip{k_t(x,\cdot)}{k_t(y,\cdot)}_{L^2(\nu)}.
\end{align*}
The heat kernel is strictly positive on a connected compact manifold for $t>0$. Hence, $s_t(x,y)>0$. The Cauchy-Schwarz inequality gives
\begin{equation*}
  k_{2t}(x,y)^2\le k_{2t}(x,x)k_{2t}(y,y),
\end{equation*}
which proves $s_t(x,y)\le1$ and $s_t(x,x)=1$.

Since $s_t(x_i,x_j)\le1$ and $\sum_{i,j}w_iw_j=1$,
\begin{equation*}
  \sum_{i,j}w_iw_js_t(x_i,x_j)\le1,
\end{equation*}
so $\gess_w(t)\ge1$. Since $s_t(x_i,x_i)=1$ and off-diagonal terms are nonnegative,
\begin{equation*}
  \sum_{i,j}w_iw_js_t(x_i,x_j)\ge\sum_iw_i^2,
\end{equation*}
so $\gess_w(t)\le1/\sum_iw_i^2$.

Merging duplicate atoms leaves the double sum unchanged because their terms have the same similarity value. For distinct support points $z_a$ with masses $\alpha_a$, the denominator is
\begin{equation*}
  \sum_{a,b}\alpha_a\alpha_bs_t(z_a,z_b).
\end{equation*}
As $t\downarrow0$, $s_t(z_a,z_a)=1$, while $s_t(z_a,z_b)\to0$ for $a\ne b$. Gaussian upper and diagonal lower bounds make this convergence uniform away from the diagonal \citep{grigoryan_2009_HeatKernelAnalysis}. The denominator therefore converges to $\sum_a\alpha_a^2$. As $t\to\infty$, the spectral expansion and gap give uniform convergence of $k_{2t}(x,y)$ to $1$. Hence, $s_t(x,y)\to1$, and the denominator converges to $1$.


\subsection{Proof of Theorem~\ref{thm:small_time} and Proposition~\ref{prop:merging}}
Let $K_t$ be the heat kernel with respect to Riemannian volume. Since $\nu=\vol/V$, the normalized-volume kernel is $k_t=VK_t$. The diagonal Minakshisundaram-Pleijel expansion gives
\citep{rosenberg_1997_LaplacianRiemannianManifold, grigoryan_2009_HeatKernelAnalysis}
\begin{equation*}
  K_s(x,x)=(4\pi s)^{-m/2}\left[1+\frac{s}{6}\Scal(x)+O(s^2)\right].
\end{equation*}
With $s=2t$,
\begin{equation*}
  k_{2t}(z_a,z_a)=V(8\pi t)^{-m/2}\left[1+\frac{t}{3}\Scal(z_a)+O(t^2)\right].
\end{equation*}
The finite, distinct support points $z_a$ have a positive minimum distance. Gaussian heat-kernel bounds make the off-diagonal contribution $V(8\pi t)^{-m/2}O(e^{-c/t})$ for some $c>0$ \citep{grigoryan_2009_HeatKernelAnalysis}. Summing the diagonal terms gives \eqref{eq:small_time_expansion}. Taylor expansion of $\log(S_\alpha+\epsilon)$ and $(S_\alpha+\epsilon)^{-1}$ gives the remaining statements. Here, $\epsilon=(t/3)C_\alpha+O(t^2)+O(e^{-c/t})$.

For Proposition~\ref{prop:merging}, use the near-diagonal heat-kernel parametrix uniformly in a normal neighborhood:
\begin{equation*}
  K_s(x,y)=(4\pi s)^{-m/2}\exp\left\{-\frac{d_g(x,y)^2}{4s}\right\} \{u_0(x,y)+O(s)\}.
\end{equation*}
Here, $u_0(x,y)$ is the square root of the inverse Jacobian determinant of the exponential map. It satisfies $u_0(x,y)=1+O(d_g(x,y)^2)$. Combine this relation with the diagonal expansion in the denominator of $s_t(x,y)$, and set $s=2t$. This gives \eqref{eq:local_overlap}. Substitution into $a^2+(1-a)^2+2a(1-a)s_t(x,y)$ proves \eqref{eq:two_atom_gess}.

\subsection{Proof of Theorem~\ref{thm:spectral}}
The heat kernel has the spectral representation
\begin{equation*}
  k_t(x,y)=\sum_{r\ge0}e^{-\lambda_rt}\varphi_r(x)\varphi_r(y),
\end{equation*}
where $\varphi_0\equiv1$. Substituting this into \eqref{eq:pt_def} gives
\begin{equation*}
  \widehat p_t(y)=1+\sum_{r\ge1}e^{-\lambda_rt}\left(\sum_iw_i\varphi_r(x_i)\right)\varphi_r(y),
\end{equation*}
which is \eqref{eq:pt_spectral}. Orthogonality implies
\begin{equation*}
  \int\widehat p_t^2\dd\nu=1+\sum_{r\ge1}e^{-2\lambda_rt}\widehat a_r^2,
\end{equation*}
which proves \eqref{eq:D2_spectral} and \eqref{eq:Aminus_spectral}. Applying $\log(1+z)=z+O(z^2)$ and separating the lowest nonzero eigenspace gives the large-time expansion.

\subsection{Proof of Corollary~\ref{cor:sphere}}
\label{app:sphere}
For $S^{p-1}$ with normalized surface measure, the coordinate functions satisfy
\begin{equation*}
  \int x_ax_b\dd\nu(x)=\frac{\delta_{ab}}{p}.
\end{equation*}
Thus, $\sqrt p\,x_a$ for $a=1,\ldots,p$ form an orthonormal basis of the degree-1 eigenspace. Its energy is
\begin{equation*}
  \sum_{a=1}^p\left(\sum_iw_i\sqrt p\,x_{ia}\right)^2=p\left\|\sum_iw_ix_i\right\|^2.
\end{equation*}

For degree 2, let $A$ and $B$ be symmetric traceless matrices. The fourth-moment identity on the sphere gives
\begin{equation*}
  \int (x^\top A x)(x^\top Bx)\dd\nu(x)
  =\frac{2\tr(AB)}{p(p+2)}.
\end{equation*}
Hence, $A\mapsto \sqrt{p(p+2)/2}\,x^\top A x$ is an isometry into the degree-2 harmonic subspace. Its domain is the Frobenius space of symmetric traceless matrices. The map is onto because both spaces have dimension $p(p+1)/2-1$. Since $Q_w-I_p/p$ is the traceless part of $Q_w$, the degree-2 energy is
\begin{equation*}
  \frac{p(p+2)}{2}\left\|Q_w-\frac{I_p}{p}\right\|_F^2.
\end{equation*}
This proves \eqref{eq:B1} and \eqref{eq:B2}.


\subsection{Proof of Theorem~\ref{thm:consistency}}
\label{app:weighted_rate}
We first prove deterministic stability. For fixed $t\in[t_0,T]$, define
\begin{equation*}
  f_{P,t}(x)=\int k_{2t}(x,y)\dd P(y).
\end{equation*}
Compactness and $t\ge t_0>0$ uniformly bound $k_{2t}$ and its first spatial derivatives on $\M^2\times[t_0,T]$. Thus, $f_{P,t}$ is uniformly Lipschitz with constant at most $L=L(t_0,T,\M)$. For probability measures $P,Q$,
\begin{align*}
  A_P(t)-A_Q(t)
  &=\iint k_{2t}(x,y)\dd(P-Q)(x)\dd P(y) \\
  &\quad+\iint k_{2t}(x,y)\dd Q(x)\dd(P-Q)(y).
\end{align*}
Kantorovich-Rubinstein duality bounds each term by $L W_1(P,Q)$. This proves \eqref{eq:W1_stability} with $C=2L$, uniformly over $[t_0,T]$. Since $A_P(t)=\|p_t\|_2^2\ge1$, the logarithm is one-Lipschitz on the relevant range.

For the random weighted rate, set $g_t(x,y)=k_{2t}(x,y)$, write
$S_n=\sum_iw_i^2$, and let $F_n(t)=\widehat A_w(t)$. Smoothness on the
compact set $\M^2\times[t_0,T]$ gives $0<g_t\le B$ and a uniform
time-Lipschitz constant $L_t$. If
$d_P(t)=\E\{g_t(X,X)\}$, independence gives
\begin{equation*}
  \E F_n(t)=(1-S_n)A_P(t)+S_nd_P(t).
\end{equation*}
Hence $\sup_t|\E F_n(t)-A_P(t)|\le C S_n$.

Replacing $X_i$ while holding the other observations fixed changes
$F_n(t)$ by at most $4Bw_i$. McDiarmid's inequality therefore yields,
uniformly at every fixed $t$,
\begin{equation}
  \Pr\{|F_n(t)-\E F_n(t)|>u\}
  \le 2\exp\{-c u^2/(B^2S_n)\}.
  \label{eq:mcdiarmid_weighted}
\end{equation}
Choose a time grid of mesh $u/(4L_t)$ and cardinality at most $C/u$.
Both $F_n(t)$ and $\E F_n(t)$ are uniformly Lipschitz in $t$. Taking
$u=K\sqrt{S_n\log(1/S_n)}$ in \eqref{eq:mcdiarmid_weighted}, applying a
union bound on the grid, and then interpolating proves
\begin{equation*}
  \sup_{t\in[t_0,T]}|F_n(t)-\E F_n(t)|
  =O_p\{\sqrt{S_n\log(1/S_n)}\}.
\end{equation*}
Since $S_n=o\{\sqrt{S_n\log(1/S_n)}\}$, the bias bound proves
\eqref{eq:weighted_rate}.

Now suppose $P=\nu$ and fix $t>0$. Then $A_\nu(t)=1$ and
$\int g_t(x,y)\dd\nu(y)=1$, so $g_t(x,y)-1$ is canonical. The
off-diagonal part of $F_n(t)-1$ has variance at most
$C\sum_{i\ne j}w_i^2w_j^2\le CS_n^2$. The diagonal part has mean
$O(S_n)$ and variance at most $C\sum_iw_i^4\le CS_n^2$. Thus
$F_n(t)-1=O_p(S_n)$, as asserted in the main text.

For self-normalized importance sampling, write $Z_n=n^{-1}\sum_i\rho(X_i)$ and
\begin{equation*}
  N_n(t)=n^{-2}\sum_{i,j}\rho(X_i)\rho(X_j)g_t(X_i,X_j).
\end{equation*}
The off-diagonal expectation of $N_n(t)$ is $(1-n^{-1})A_P(t)$, and
the bounded diagonal contributes $O(n^{-1})$ uniformly in $t$. Replacing
one observation changes $N_n(t)$ by at most $C/n$. The same grid and
bounded-differences argument therefore gives
\begin{equation*}
  \sup_t|N_n(t)-A_P(t)|=O_p\{\sqrt{\log n/n}\}.
\end{equation*}
Boundedness of $\rho$ also gives $Z_n=1+O_p(n^{-1/2})$. On $Z_n>0$, $\widehat A_w(t)=N_n(t)/Z_n^2$. Uniform boundedness of $A_P(t)$ gives the claimed rate for this ratio. Finally, $Q\{\rho(X)=0\}<1$, so $\Pr(Z_n=0)=Q\{\rho(X)=0\}^n\to0$. The convention on this event is asymptotically irrelevant.


\section{ADDITIONAL EXPERIMENTS AND REPRODUCIBILITY}
\label{app:add_experiments}

This appendix gives the details needed to reproduce the figures and tables. It also reports results beyond the main-text summaries. Replications assess profile stability, while a truncation study measures numerical sensitivity. A self-normalized importance-sampling example compares ordinary ESS with gESS. Complete data profiles extend Table~\ref{tab:digits}.

\paragraph{Synthetic configurations on $S^2$.}
All five examples in Table~\ref{tab:synthetic} use equal weights and $n=300$ labeled atoms. The unimodal sample follows a vMF distribution on $S^2$. Its mean direction is $e_3=(0,0,1)^\top$, and its concentration is $\kappa=140$. The antipodal example combines 150 observations from this distribution with 150 from its antipodal component. The second component has mean direction $-e_3$ and concentration $\kappa=140$. For the girdle, draw $\theta\sim\mathrm{Unif}(0,2\pi)$ and $z\sim N(0,0.035^2)$. Clip $z$ to $[-0.12,0.12]$, then set
\[
  (\sqrt{1-z^2}\cos\theta,\sqrt{1-z^2}\sin\theta,z)^\top.
\]
The four-mode example uses four tetrahedral directions:
\[
  \frac{(1,1,1)}{\sqrt3},\quad
  \frac{(1,-1,-1)}{\sqrt3},\quad
  \frac{(-1,1,-1)}{\sqrt3},\quad
  \frac{(-1,-1,1)}{\sqrt3}.
\]
Each center contributes 75 vMF observations with $\kappa=140$. The duplicate example constructs 30 quasi-uniform Fibonacci-sphere locations. Repeating each location ten times gives 300 labels at 30 distinct support points. This construction tests exact atom-splitting invariance.

\paragraph{Angular scale, heat-kernel computation, and truncation.}
On $S^2$, we express diffusion time through angular resolution:
\[
  \alpha=\sqrt{8t}.
\]
We convert radians to degrees by multiplying by $180/\pi$. Under the small-time law $\exp\{-d_g(x,y)^2/(8t)\}$, distance $d_g(x,y)=\alpha$ gives leading overlap $e^{-1}$. We evaluate the heat kernel through the spectral series
\[
  k_t(x,y)=\sum_{\ell=0}^{\infty}(2\ell+1)e^{-\ell(\ell+1)t}
  P_\ell(x^\top y).
\]
The main figures and tables truncate the series at $L=180$. The replication, importance-sampling, and digit-profile figures use $L=120$, $140$, and $160$. For $k_{2t}$, the diagonal truncation tail is bounded by
\[
  \sum_{\ell>L}(2\ell+1)e^{-2t\ell(\ell+1)}.
\]
Finer diffusion scales require $L$ to increase on the order of $t^{-1/2}$. Given the heat-kernel values, evaluating $K$ scales costs $O(n^2K)$.

\paragraph{Replicated synthetic profiles.}

Figure~\ref{fig:app_replicates} uses 30 independent variants of the synthetic constructions, with $n=180$ per replication. The vMF concentration is $\kappa=100$ to make variation visible. The girdle standard deviation is $0.045$, with clipping to $[-0.15,0.15]$ before the same embedding. The duplicate example repeats each of 20 Fibonacci-sphere locations nine times.
Curves report medians and shaded interquartile ranges.


\begin{figure*}[t]
\centering
%\includegraphics[width=0.96\textwidth]{hkep_app_replicates.png}
\providecommand{\mathdefault}[1]{#1}
\providecommand{\mathregular}[1]{#1}
%% Creator: Matplotlib, PGF backend
%%
%% To include the figure in your LaTeX document, write
%%   \input{<filename>.pgf}
%%
%% Make sure the required packages are loaded in your preamble
%%   \usepackage{pgf}
%%
%% Also ensure that all the required font packages are loaded; for instance,
%% the lmodern package is sometimes necessary when using math font.
%%   \usepackage{lmodern}
%%
%% Figures using additional raster images can only be included by \input if
%% they are in the same directory as the main LaTeX file. For loading figures
%% from other directories you can use the `import` package
%%   \usepackage{import}
%%
%% and then include the figures with
%%   \import{<path to file>}{<filename>.pgf}
%%
%% Matplotlib used the following preamble
%%   \def\mathdefault#1{#1}
%%   \everymath=\expandafter{\the\everymath\displaystyle}
%%   \IfFileExists{scrextend.sty}{
%%     \usepackage[fontsize=8.000000pt]{scrextend}
%%   }{
%%     \renewcommand{\normalsize}{\fontsize{8.000000}{9.600000}\selectfont}
%%     \normalsize
%%   }
%%   \usepackage{amsmath,amssymb}
%%   \makeatletter\@ifpackageloaded{underscore}{}{\usepackage[strings]{underscore}}\makeatother
%%
\begingroup%
\makeatletter%
\begin{pgfpicture}%
\pgfpathrectangle{\pgfpointorigin}{\pgfqpoint{6.932269in}{3.090056in}}%
\pgfusepath{use as bounding box, clip}%
\begin{pgfscope}%
\pgfsetbuttcap%
\pgfsetmiterjoin%
\definecolor{currentfill}{rgb}{1.000000,1.000000,1.000000}%
\pgfsetfillcolor{currentfill}%
\pgfsetlinewidth{0.000000pt}%
\definecolor{currentstroke}{rgb}{1.000000,1.000000,1.000000}%
\pgfsetstrokecolor{currentstroke}%
\pgfsetdash{}{0pt}%
\pgfpathmoveto{\pgfqpoint{0.000000in}{-0.000000in}}%
\pgfpathlineto{\pgfqpoint{6.932269in}{-0.000000in}}%
\pgfpathlineto{\pgfqpoint{6.932269in}{3.090056in}}%
\pgfpathlineto{\pgfqpoint{0.000000in}{3.090056in}}%
\pgfpathlineto{\pgfqpoint{0.000000in}{-0.000000in}}%
\pgfpathclose%
\pgfusepath{fill}%
\end{pgfscope}%
\begin{pgfscope}%
\pgfsetbuttcap%
\pgfsetmiterjoin%
\definecolor{currentfill}{rgb}{1.000000,1.000000,1.000000}%
\pgfsetfillcolor{currentfill}%
\pgfsetlinewidth{0.000000pt}%
\definecolor{currentstroke}{rgb}{0.000000,0.000000,0.000000}%
\pgfsetstrokecolor{currentstroke}%
\pgfsetstrokeopacity{0.000000}%
\pgfsetdash{}{0pt}%
\pgfpathmoveto{\pgfqpoint{0.462269in}{0.807389in}}%
\pgfpathlineto{\pgfqpoint{3.231834in}{0.807389in}}%
\pgfpathlineto{\pgfqpoint{3.231834in}{2.823389in}}%
\pgfpathlineto{\pgfqpoint{0.462269in}{2.823389in}}%
\pgfpathlineto{\pgfqpoint{0.462269in}{0.807389in}}%
\pgfpathclose%
\pgfusepath{fill}%
\end{pgfscope}%
\begin{pgfscope}%
\pgfpathrectangle{\pgfqpoint{0.462269in}{0.807389in}}{\pgfqpoint{2.769565in}{2.016000in}}%
\pgfusepath{clip}%
\pgfsetbuttcap%
\pgfsetroundjoin%
\definecolor{currentfill}{rgb}{0.121569,0.466667,0.705882}%
\pgfsetfillcolor{currentfill}%
\pgfsetfillopacity{0.150000}%
\pgfsetlinewidth{0.000000pt}%
\definecolor{currentstroke}{rgb}{0.121569,0.466667,0.705882}%
\pgfsetstrokecolor{currentstroke}%
\pgfsetstrokeopacity{0.150000}%
\pgfsetdash{}{0pt}%
\pgfsys@defobject{currentmarker}{\pgfqpoint{0.588158in}{0.899025in}}{\pgfqpoint{3.105945in}{1.113591in}}{%
\pgfpathmoveto{\pgfqpoint{0.588158in}{1.113591in}}%
\pgfpathlineto{\pgfqpoint{0.588158in}{1.090634in}}%
\pgfpathlineto{\pgfqpoint{0.634784in}{1.071101in}}%
\pgfpathlineto{\pgfqpoint{0.681410in}{1.053353in}}%
\pgfpathlineto{\pgfqpoint{0.728036in}{1.037399in}}%
\pgfpathlineto{\pgfqpoint{0.774661in}{1.023055in}}%
\pgfpathlineto{\pgfqpoint{0.821287in}{1.010039in}}%
\pgfpathlineto{\pgfqpoint{0.867913in}{0.998367in}}%
\pgfpathlineto{\pgfqpoint{0.914538in}{0.987904in}}%
\pgfpathlineto{\pgfqpoint{0.961164in}{0.978527in}}%
\pgfpathlineto{\pgfqpoint{1.007790in}{0.970126in}}%
\pgfpathlineto{\pgfqpoint{1.054415in}{0.962602in}}%
\pgfpathlineto{\pgfqpoint{1.101041in}{0.955877in}}%
\pgfpathlineto{\pgfqpoint{1.147667in}{0.949947in}}%
\pgfpathlineto{\pgfqpoint{1.194292in}{0.944628in}}%
\pgfpathlineto{\pgfqpoint{1.240918in}{0.939838in}}%
\pgfpathlineto{\pgfqpoint{1.287544in}{0.935506in}}%
\pgfpathlineto{\pgfqpoint{1.334169in}{0.931628in}}%
\pgfpathlineto{\pgfqpoint{1.380795in}{0.928159in}}%
\pgfpathlineto{\pgfqpoint{1.427421in}{0.925082in}}%
\pgfpathlineto{\pgfqpoint{1.474046in}{0.922316in}}%
\pgfpathlineto{\pgfqpoint{1.520672in}{0.919827in}}%
\pgfpathlineto{\pgfqpoint{1.567298in}{0.917599in}}%
\pgfpathlineto{\pgfqpoint{1.613923in}{0.915603in}}%
\pgfpathlineto{\pgfqpoint{1.660549in}{0.913816in}}%
\pgfpathlineto{\pgfqpoint{1.707175in}{0.912215in}}%
\pgfpathlineto{\pgfqpoint{1.753800in}{0.910782in}}%
\pgfpathlineto{\pgfqpoint{1.800426in}{0.909499in}}%
\pgfpathlineto{\pgfqpoint{1.847052in}{0.908347in}}%
\pgfpathlineto{\pgfqpoint{1.893677in}{0.907314in}}%
\pgfpathlineto{\pgfqpoint{1.940303in}{0.906390in}}%
\pgfpathlineto{\pgfqpoint{1.986929in}{0.905564in}}%
\pgfpathlineto{\pgfqpoint{2.033554in}{0.904825in}}%
\pgfpathlineto{\pgfqpoint{2.080180in}{0.904164in}}%
\pgfpathlineto{\pgfqpoint{2.126806in}{0.903573in}}%
\pgfpathlineto{\pgfqpoint{2.173432in}{0.903044in}}%
\pgfpathlineto{\pgfqpoint{2.220057in}{0.902571in}}%
\pgfpathlineto{\pgfqpoint{2.266683in}{0.902147in}}%
\pgfpathlineto{\pgfqpoint{2.313309in}{0.901768in}}%
\pgfpathlineto{\pgfqpoint{2.359934in}{0.901429in}}%
\pgfpathlineto{\pgfqpoint{2.406560in}{0.901126in}}%
\pgfpathlineto{\pgfqpoint{2.453186in}{0.900854in}}%
\pgfpathlineto{\pgfqpoint{2.499811in}{0.900611in}}%
\pgfpathlineto{\pgfqpoint{2.546437in}{0.900394in}}%
\pgfpathlineto{\pgfqpoint{2.593063in}{0.900199in}}%
\pgfpathlineto{\pgfqpoint{2.639688in}{0.900025in}}%
\pgfpathlineto{\pgfqpoint{2.686314in}{0.899869in}}%
\pgfpathlineto{\pgfqpoint{2.732940in}{0.899729in}}%
\pgfpathlineto{\pgfqpoint{2.779565in}{0.899604in}}%
\pgfpathlineto{\pgfqpoint{2.826191in}{0.899492in}}%
\pgfpathlineto{\pgfqpoint{2.872817in}{0.899392in}}%
\pgfpathlineto{\pgfqpoint{2.919442in}{0.899302in}}%
\pgfpathlineto{\pgfqpoint{2.966068in}{0.899222in}}%
\pgfpathlineto{\pgfqpoint{3.012694in}{0.899149in}}%
\pgfpathlineto{\pgfqpoint{3.059319in}{0.899084in}}%
\pgfpathlineto{\pgfqpoint{3.105945in}{0.899025in}}%
\pgfpathlineto{\pgfqpoint{3.105945in}{0.899079in}}%
\pgfpathlineto{\pgfqpoint{3.105945in}{0.899079in}}%
\pgfpathlineto{\pgfqpoint{3.059319in}{0.899146in}}%
\pgfpathlineto{\pgfqpoint{3.012694in}{0.899221in}}%
\pgfpathlineto{\pgfqpoint{2.966068in}{0.899304in}}%
\pgfpathlineto{\pgfqpoint{2.919442in}{0.899397in}}%
\pgfpathlineto{\pgfqpoint{2.872817in}{0.899500in}}%
\pgfpathlineto{\pgfqpoint{2.826191in}{0.899616in}}%
\pgfpathlineto{\pgfqpoint{2.779565in}{0.899744in}}%
\pgfpathlineto{\pgfqpoint{2.732940in}{0.899888in}}%
\pgfpathlineto{\pgfqpoint{2.686314in}{0.900048in}}%
\pgfpathlineto{\pgfqpoint{2.639688in}{0.900227in}}%
\pgfpathlineto{\pgfqpoint{2.593063in}{0.900427in}}%
\pgfpathlineto{\pgfqpoint{2.546437in}{0.900650in}}%
\pgfpathlineto{\pgfqpoint{2.499811in}{0.900899in}}%
\pgfpathlineto{\pgfqpoint{2.453186in}{0.901178in}}%
\pgfpathlineto{\pgfqpoint{2.406560in}{0.901489in}}%
\pgfpathlineto{\pgfqpoint{2.359934in}{0.901837in}}%
\pgfpathlineto{\pgfqpoint{2.313309in}{0.902225in}}%
\pgfpathlineto{\pgfqpoint{2.266683in}{0.902658in}}%
\pgfpathlineto{\pgfqpoint{2.220057in}{0.903143in}}%
\pgfpathlineto{\pgfqpoint{2.173432in}{0.903684in}}%
\pgfpathlineto{\pgfqpoint{2.126806in}{0.904289in}}%
\pgfpathlineto{\pgfqpoint{2.080180in}{0.904964in}}%
\pgfpathlineto{\pgfqpoint{2.033554in}{0.905718in}}%
\pgfpathlineto{\pgfqpoint{1.986929in}{0.906560in}}%
\pgfpathlineto{\pgfqpoint{1.940303in}{0.907501in}}%
\pgfpathlineto{\pgfqpoint{1.893677in}{0.908552in}}%
\pgfpathlineto{\pgfqpoint{1.847052in}{0.909725in}}%
\pgfpathlineto{\pgfqpoint{1.800426in}{0.911035in}}%
\pgfpathlineto{\pgfqpoint{1.753800in}{0.912497in}}%
\pgfpathlineto{\pgfqpoint{1.707175in}{0.914129in}}%
\pgfpathlineto{\pgfqpoint{1.660549in}{0.915951in}}%
\pgfpathlineto{\pgfqpoint{1.613923in}{0.917985in}}%
\pgfpathlineto{\pgfqpoint{1.567298in}{0.920254in}}%
\pgfpathlineto{\pgfqpoint{1.520672in}{0.922785in}}%
\pgfpathlineto{\pgfqpoint{1.474046in}{0.925608in}}%
\pgfpathlineto{\pgfqpoint{1.427421in}{0.928757in}}%
\pgfpathlineto{\pgfqpoint{1.380795in}{0.932268in}}%
\pgfpathlineto{\pgfqpoint{1.334169in}{0.936181in}}%
\pgfpathlineto{\pgfqpoint{1.287544in}{0.940543in}}%
\pgfpathlineto{\pgfqpoint{1.240918in}{0.945451in}}%
\pgfpathlineto{\pgfqpoint{1.194292in}{0.950927in}}%
\pgfpathlineto{\pgfqpoint{1.147667in}{0.957033in}}%
\pgfpathlineto{\pgfqpoint{1.101041in}{0.963841in}}%
\pgfpathlineto{\pgfqpoint{1.054415in}{0.971428in}}%
\pgfpathlineto{\pgfqpoint{1.007790in}{0.979893in}}%
\pgfpathlineto{\pgfqpoint{0.961164in}{0.989360in}}%
\pgfpathlineto{\pgfqpoint{0.914538in}{0.999906in}}%
\pgfpathlineto{\pgfqpoint{0.867913in}{1.011459in}}%
\pgfpathlineto{\pgfqpoint{0.821287in}{1.024299in}}%
\pgfpathlineto{\pgfqpoint{0.774661in}{1.038564in}}%
\pgfpathlineto{\pgfqpoint{0.728036in}{1.054403in}}%
\pgfpathlineto{\pgfqpoint{0.681410in}{1.071985in}}%
\pgfpathlineto{\pgfqpoint{0.634784in}{1.091632in}}%
\pgfpathlineto{\pgfqpoint{0.588158in}{1.113591in}}%
\pgfpathlineto{\pgfqpoint{0.588158in}{1.113591in}}%
\pgfpathclose%
\pgfusepath{fill}%
}%
\begin{pgfscope}%
\pgfsys@transformshift{0.000000in}{0.000000in}%
\pgfsys@useobject{currentmarker}{}%
\end{pgfscope}%
\end{pgfscope}%
\begin{pgfscope}%
\pgfpathrectangle{\pgfqpoint{0.462269in}{0.807389in}}{\pgfqpoint{2.769565in}{2.016000in}}%
\pgfusepath{clip}%
\pgfsetbuttcap%
\pgfsetroundjoin%
\definecolor{currentfill}{rgb}{1.000000,0.498039,0.054902}%
\pgfsetfillcolor{currentfill}%
\pgfsetfillopacity{0.150000}%
\pgfsetlinewidth{0.000000pt}%
\definecolor{currentstroke}{rgb}{1.000000,0.498039,0.054902}%
\pgfsetstrokecolor{currentstroke}%
\pgfsetstrokeopacity{0.150000}%
\pgfsetdash{}{0pt}%
\pgfsys@defobject{currentmarker}{\pgfqpoint{0.588158in}{0.924369in}}{\pgfqpoint{3.105945in}{1.374210in}}{%
\pgfpathmoveto{\pgfqpoint{0.588158in}{1.374210in}}%
\pgfpathlineto{\pgfqpoint{0.588158in}{1.323716in}}%
\pgfpathlineto{\pgfqpoint{0.634784in}{1.285163in}}%
\pgfpathlineto{\pgfqpoint{0.681410in}{1.250272in}}%
\pgfpathlineto{\pgfqpoint{0.728036in}{1.218990in}}%
\pgfpathlineto{\pgfqpoint{0.774661in}{1.190806in}}%
\pgfpathlineto{\pgfqpoint{0.821287in}{1.165377in}}%
\pgfpathlineto{\pgfqpoint{0.867913in}{1.142387in}}%
\pgfpathlineto{\pgfqpoint{0.914538in}{1.122082in}}%
\pgfpathlineto{\pgfqpoint{0.961164in}{1.103463in}}%
\pgfpathlineto{\pgfqpoint{1.007790in}{1.086817in}}%
\pgfpathlineto{\pgfqpoint{1.054415in}{1.071825in}}%
\pgfpathlineto{\pgfqpoint{1.101041in}{1.058324in}}%
\pgfpathlineto{\pgfqpoint{1.147667in}{1.046195in}}%
\pgfpathlineto{\pgfqpoint{1.194292in}{1.035288in}}%
\pgfpathlineto{\pgfqpoint{1.240918in}{1.025407in}}%
\pgfpathlineto{\pgfqpoint{1.287544in}{1.016568in}}%
\pgfpathlineto{\pgfqpoint{1.334169in}{1.008625in}}%
\pgfpathlineto{\pgfqpoint{1.380795in}{1.001490in}}%
\pgfpathlineto{\pgfqpoint{1.427421in}{0.995081in}}%
\pgfpathlineto{\pgfqpoint{1.474046in}{0.989326in}}%
\pgfpathlineto{\pgfqpoint{1.520672in}{0.984160in}}%
\pgfpathlineto{\pgfqpoint{1.567298in}{0.979524in}}%
\pgfpathlineto{\pgfqpoint{1.613923in}{0.975363in}}%
\pgfpathlineto{\pgfqpoint{1.660549in}{0.971630in}}%
\pgfpathlineto{\pgfqpoint{1.707175in}{0.968261in}}%
\pgfpathlineto{\pgfqpoint{1.753800in}{0.965242in}}%
\pgfpathlineto{\pgfqpoint{1.800426in}{0.962528in}}%
\pgfpathlineto{\pgfqpoint{1.847052in}{0.960089in}}%
\pgfpathlineto{\pgfqpoint{1.893677in}{0.957907in}}%
\pgfpathlineto{\pgfqpoint{1.940303in}{0.955955in}}%
\pgfpathlineto{\pgfqpoint{1.986929in}{0.954210in}}%
\pgfpathlineto{\pgfqpoint{2.033554in}{0.952648in}}%
\pgfpathlineto{\pgfqpoint{2.080180in}{0.951252in}}%
\pgfpathlineto{\pgfqpoint{2.126806in}{0.950015in}}%
\pgfpathlineto{\pgfqpoint{2.173432in}{0.948907in}}%
\pgfpathlineto{\pgfqpoint{2.220057in}{0.947913in}}%
\pgfpathlineto{\pgfqpoint{2.266683in}{0.947022in}}%
\pgfpathlineto{\pgfqpoint{2.313309in}{0.946223in}}%
\pgfpathlineto{\pgfqpoint{2.359934in}{0.945508in}}%
\pgfpathlineto{\pgfqpoint{2.406560in}{0.944867in}}%
\pgfpathlineto{\pgfqpoint{2.453186in}{0.944292in}}%
\pgfpathlineto{\pgfqpoint{2.499811in}{0.943776in}}%
\pgfpathlineto{\pgfqpoint{2.546437in}{0.943309in}}%
\pgfpathlineto{\pgfqpoint{2.593063in}{0.942878in}}%
\pgfpathlineto{\pgfqpoint{2.639688in}{0.942463in}}%
\pgfpathlineto{\pgfqpoint{2.686314in}{0.942029in}}%
\pgfpathlineto{\pgfqpoint{2.732940in}{0.941519in}}%
\pgfpathlineto{\pgfqpoint{2.779565in}{0.940850in}}%
\pgfpathlineto{\pgfqpoint{2.826191in}{0.939915in}}%
\pgfpathlineto{\pgfqpoint{2.872817in}{0.938594in}}%
\pgfpathlineto{\pgfqpoint{2.919442in}{0.936776in}}%
\pgfpathlineto{\pgfqpoint{2.966068in}{0.934399in}}%
\pgfpathlineto{\pgfqpoint{3.012694in}{0.931469in}}%
\pgfpathlineto{\pgfqpoint{3.059319in}{0.928074in}}%
\pgfpathlineto{\pgfqpoint{3.105945in}{0.924369in}}%
\pgfpathlineto{\pgfqpoint{3.105945in}{0.924386in}}%
\pgfpathlineto{\pgfqpoint{3.105945in}{0.924386in}}%
\pgfpathlineto{\pgfqpoint{3.059319in}{0.928102in}}%
\pgfpathlineto{\pgfqpoint{3.012694in}{0.931511in}}%
\pgfpathlineto{\pgfqpoint{2.966068in}{0.934461in}}%
\pgfpathlineto{\pgfqpoint{2.919442in}{0.936862in}}%
\pgfpathlineto{\pgfqpoint{2.872817in}{0.938706in}}%
\pgfpathlineto{\pgfqpoint{2.826191in}{0.940057in}}%
\pgfpathlineto{\pgfqpoint{2.779565in}{0.941022in}}%
\pgfpathlineto{\pgfqpoint{2.732940in}{0.941722in}}%
\pgfpathlineto{\pgfqpoint{2.686314in}{0.942264in}}%
\pgfpathlineto{\pgfqpoint{2.639688in}{0.942732in}}%
\pgfpathlineto{\pgfqpoint{2.593063in}{0.943183in}}%
\pgfpathlineto{\pgfqpoint{2.546437in}{0.943654in}}%
\pgfpathlineto{\pgfqpoint{2.499811in}{0.944165in}}%
\pgfpathlineto{\pgfqpoint{2.453186in}{0.944732in}}%
\pgfpathlineto{\pgfqpoint{2.406560in}{0.945363in}}%
\pgfpathlineto{\pgfqpoint{2.359934in}{0.946067in}}%
\pgfpathlineto{\pgfqpoint{2.313309in}{0.946853in}}%
\pgfpathlineto{\pgfqpoint{2.266683in}{0.947732in}}%
\pgfpathlineto{\pgfqpoint{2.220057in}{0.948714in}}%
\pgfpathlineto{\pgfqpoint{2.173432in}{0.949810in}}%
\pgfpathlineto{\pgfqpoint{2.126806in}{0.951035in}}%
\pgfpathlineto{\pgfqpoint{2.080180in}{0.952403in}}%
\pgfpathlineto{\pgfqpoint{2.033554in}{0.953932in}}%
\pgfpathlineto{\pgfqpoint{1.986929in}{0.955641in}}%
\pgfpathlineto{\pgfqpoint{1.940303in}{0.957550in}}%
\pgfpathlineto{\pgfqpoint{1.893677in}{0.959683in}}%
\pgfpathlineto{\pgfqpoint{1.847052in}{0.962066in}}%
\pgfpathlineto{\pgfqpoint{1.800426in}{0.964727in}}%
\pgfpathlineto{\pgfqpoint{1.753800in}{0.967700in}}%
\pgfpathlineto{\pgfqpoint{1.707175in}{0.971021in}}%
\pgfpathlineto{\pgfqpoint{1.660549in}{0.974730in}}%
\pgfpathlineto{\pgfqpoint{1.613923in}{0.978872in}}%
\pgfpathlineto{\pgfqpoint{1.567298in}{0.983498in}}%
\pgfpathlineto{\pgfqpoint{1.520672in}{0.988663in}}%
\pgfpathlineto{\pgfqpoint{1.474046in}{0.994431in}}%
\pgfpathlineto{\pgfqpoint{1.427421in}{1.000870in}}%
\pgfpathlineto{\pgfqpoint{1.380795in}{1.008058in}}%
\pgfpathlineto{\pgfqpoint{1.334169in}{1.016081in}}%
\pgfpathlineto{\pgfqpoint{1.287544in}{1.025035in}}%
\pgfpathlineto{\pgfqpoint{1.240918in}{1.035024in}}%
\pgfpathlineto{\pgfqpoint{1.194292in}{1.046165in}}%
\pgfpathlineto{\pgfqpoint{1.147667in}{1.058588in}}%
\pgfpathlineto{\pgfqpoint{1.101041in}{1.072434in}}%
\pgfpathlineto{\pgfqpoint{1.054415in}{1.087858in}}%
\pgfpathlineto{\pgfqpoint{1.007790in}{1.105030in}}%
\pgfpathlineto{\pgfqpoint{0.961164in}{1.124135in}}%
\pgfpathlineto{\pgfqpoint{0.914538in}{1.145374in}}%
\pgfpathlineto{\pgfqpoint{0.867913in}{1.168964in}}%
\pgfpathlineto{\pgfqpoint{0.821287in}{1.195138in}}%
\pgfpathlineto{\pgfqpoint{0.774661in}{1.224149in}}%
\pgfpathlineto{\pgfqpoint{0.728036in}{1.256264in}}%
\pgfpathlineto{\pgfqpoint{0.681410in}{1.291772in}}%
\pgfpathlineto{\pgfqpoint{0.634784in}{1.330979in}}%
\pgfpathlineto{\pgfqpoint{0.588158in}{1.374210in}}%
\pgfpathlineto{\pgfqpoint{0.588158in}{1.374210in}}%
\pgfpathclose%
\pgfusepath{fill}%
}%
\begin{pgfscope}%
\pgfsys@transformshift{0.000000in}{0.000000in}%
\pgfsys@useobject{currentmarker}{}%
\end{pgfscope}%
\end{pgfscope}%
\begin{pgfscope}%
\pgfpathrectangle{\pgfqpoint{0.462269in}{0.807389in}}{\pgfqpoint{2.769565in}{2.016000in}}%
\pgfusepath{clip}%
\pgfsetbuttcap%
\pgfsetroundjoin%
\definecolor{currentfill}{rgb}{0.172549,0.627451,0.172549}%
\pgfsetfillcolor{currentfill}%
\pgfsetfillopacity{0.150000}%
\pgfsetlinewidth{0.000000pt}%
\definecolor{currentstroke}{rgb}{0.172549,0.627451,0.172549}%
\pgfsetstrokecolor{currentstroke}%
\pgfsetstrokeopacity{0.150000}%
\pgfsetdash{}{0pt}%
\pgfsys@defobject{currentmarker}{\pgfqpoint{0.588158in}{0.926336in}}{\pgfqpoint{3.105945in}{2.731753in}}{%
\pgfpathmoveto{\pgfqpoint{0.588158in}{2.731753in}}%
\pgfpathlineto{\pgfqpoint{0.588158in}{2.631510in}}%
\pgfpathlineto{\pgfqpoint{0.634784in}{2.520468in}}%
\pgfpathlineto{\pgfqpoint{0.681410in}{2.416407in}}%
\pgfpathlineto{\pgfqpoint{0.728036in}{2.320040in}}%
\pgfpathlineto{\pgfqpoint{0.774661in}{2.230277in}}%
\pgfpathlineto{\pgfqpoint{0.821287in}{2.146474in}}%
\pgfpathlineto{\pgfqpoint{0.867913in}{2.068292in}}%
\pgfpathlineto{\pgfqpoint{0.914538in}{1.995375in}}%
\pgfpathlineto{\pgfqpoint{0.961164in}{1.927365in}}%
\pgfpathlineto{\pgfqpoint{1.007790in}{1.865113in}}%
\pgfpathlineto{\pgfqpoint{1.054415in}{1.805776in}}%
\pgfpathlineto{\pgfqpoint{1.101041in}{1.750586in}}%
\pgfpathlineto{\pgfqpoint{1.147667in}{1.699501in}}%
\pgfpathlineto{\pgfqpoint{1.194292in}{1.651706in}}%
\pgfpathlineto{\pgfqpoint{1.240918in}{1.606878in}}%
\pgfpathlineto{\pgfqpoint{1.287544in}{1.563704in}}%
\pgfpathlineto{\pgfqpoint{1.334169in}{1.523431in}}%
\pgfpathlineto{\pgfqpoint{1.380795in}{1.486212in}}%
\pgfpathlineto{\pgfqpoint{1.427421in}{1.451193in}}%
\pgfpathlineto{\pgfqpoint{1.474046in}{1.418231in}}%
\pgfpathlineto{\pgfqpoint{1.520672in}{1.387195in}}%
\pgfpathlineto{\pgfqpoint{1.567298in}{1.357966in}}%
\pgfpathlineto{\pgfqpoint{1.613923in}{1.330404in}}%
\pgfpathlineto{\pgfqpoint{1.660549in}{1.304412in}}%
\pgfpathlineto{\pgfqpoint{1.707175in}{1.279907in}}%
\pgfpathlineto{\pgfqpoint{1.753800in}{1.256966in}}%
\pgfpathlineto{\pgfqpoint{1.800426in}{1.235351in}}%
\pgfpathlineto{\pgfqpoint{1.847052in}{1.214902in}}%
\pgfpathlineto{\pgfqpoint{1.893677in}{1.195596in}}%
\pgfpathlineto{\pgfqpoint{1.940303in}{1.177358in}}%
\pgfpathlineto{\pgfqpoint{1.986929in}{1.160118in}}%
\pgfpathlineto{\pgfqpoint{2.033554in}{1.143908in}}%
\pgfpathlineto{\pgfqpoint{2.080180in}{1.128573in}}%
\pgfpathlineto{\pgfqpoint{2.126806in}{1.114012in}}%
\pgfpathlineto{\pgfqpoint{2.173432in}{1.100154in}}%
\pgfpathlineto{\pgfqpoint{2.220057in}{1.086994in}}%
\pgfpathlineto{\pgfqpoint{2.266683in}{1.074450in}}%
\pgfpathlineto{\pgfqpoint{2.313309in}{1.062492in}}%
\pgfpathlineto{\pgfqpoint{2.359934in}{1.051099in}}%
\pgfpathlineto{\pgfqpoint{2.406560in}{1.040297in}}%
\pgfpathlineto{\pgfqpoint{2.453186in}{1.030114in}}%
\pgfpathlineto{\pgfqpoint{2.499811in}{1.020389in}}%
\pgfpathlineto{\pgfqpoint{2.546437in}{1.011114in}}%
\pgfpathlineto{\pgfqpoint{2.593063in}{1.002237in}}%
\pgfpathlineto{\pgfqpoint{2.639688in}{0.993761in}}%
\pgfpathlineto{\pgfqpoint{2.686314in}{0.985647in}}%
\pgfpathlineto{\pgfqpoint{2.732940in}{0.977855in}}%
\pgfpathlineto{\pgfqpoint{2.779565in}{0.970372in}}%
\pgfpathlineto{\pgfqpoint{2.826191in}{0.963175in}}%
\pgfpathlineto{\pgfqpoint{2.872817in}{0.956253in}}%
\pgfpathlineto{\pgfqpoint{2.919442in}{0.949607in}}%
\pgfpathlineto{\pgfqpoint{2.966068in}{0.943257in}}%
\pgfpathlineto{\pgfqpoint{3.012694in}{0.937233in}}%
\pgfpathlineto{\pgfqpoint{3.059319in}{0.931578in}}%
\pgfpathlineto{\pgfqpoint{3.105945in}{0.926336in}}%
\pgfpathlineto{\pgfqpoint{3.105945in}{0.926668in}}%
\pgfpathlineto{\pgfqpoint{3.105945in}{0.926668in}}%
\pgfpathlineto{\pgfqpoint{3.059319in}{0.931995in}}%
\pgfpathlineto{\pgfqpoint{3.012694in}{0.937746in}}%
\pgfpathlineto{\pgfqpoint{2.966068in}{0.943878in}}%
\pgfpathlineto{\pgfqpoint{2.919442in}{0.950355in}}%
\pgfpathlineto{\pgfqpoint{2.872817in}{0.957139in}}%
\pgfpathlineto{\pgfqpoint{2.826191in}{0.964210in}}%
\pgfpathlineto{\pgfqpoint{2.779565in}{0.971564in}}%
\pgfpathlineto{\pgfqpoint{2.732940in}{0.979213in}}%
\pgfpathlineto{\pgfqpoint{2.686314in}{0.987179in}}%
\pgfpathlineto{\pgfqpoint{2.639688in}{0.995507in}}%
\pgfpathlineto{\pgfqpoint{2.593063in}{1.004226in}}%
\pgfpathlineto{\pgfqpoint{2.546437in}{1.013368in}}%
\pgfpathlineto{\pgfqpoint{2.499811in}{1.022931in}}%
\pgfpathlineto{\pgfqpoint{2.453186in}{1.032972in}}%
\pgfpathlineto{\pgfqpoint{2.406560in}{1.043523in}}%
\pgfpathlineto{\pgfqpoint{2.359934in}{1.054660in}}%
\pgfpathlineto{\pgfqpoint{2.313309in}{1.066423in}}%
\pgfpathlineto{\pgfqpoint{2.266683in}{1.078831in}}%
\pgfpathlineto{\pgfqpoint{2.220057in}{1.091884in}}%
\pgfpathlineto{\pgfqpoint{2.173432in}{1.105608in}}%
\pgfpathlineto{\pgfqpoint{2.126806in}{1.120075in}}%
\pgfpathlineto{\pgfqpoint{2.080180in}{1.135278in}}%
\pgfpathlineto{\pgfqpoint{2.033554in}{1.151337in}}%
\pgfpathlineto{\pgfqpoint{1.986929in}{1.168348in}}%
\pgfpathlineto{\pgfqpoint{1.940303in}{1.186395in}}%
\pgfpathlineto{\pgfqpoint{1.893677in}{1.205425in}}%
\pgfpathlineto{\pgfqpoint{1.847052in}{1.225427in}}%
\pgfpathlineto{\pgfqpoint{1.800426in}{1.246571in}}%
\pgfpathlineto{\pgfqpoint{1.753800in}{1.268971in}}%
\pgfpathlineto{\pgfqpoint{1.707175in}{1.292823in}}%
\pgfpathlineto{\pgfqpoint{1.660549in}{1.318262in}}%
\pgfpathlineto{\pgfqpoint{1.613923in}{1.345540in}}%
\pgfpathlineto{\pgfqpoint{1.567298in}{1.374499in}}%
\pgfpathlineto{\pgfqpoint{1.520672in}{1.404783in}}%
\pgfpathlineto{\pgfqpoint{1.474046in}{1.437417in}}%
\pgfpathlineto{\pgfqpoint{1.427421in}{1.471940in}}%
\pgfpathlineto{\pgfqpoint{1.380795in}{1.508784in}}%
\pgfpathlineto{\pgfqpoint{1.334169in}{1.548123in}}%
\pgfpathlineto{\pgfqpoint{1.287544in}{1.590561in}}%
\pgfpathlineto{\pgfqpoint{1.240918in}{1.636383in}}%
\pgfpathlineto{\pgfqpoint{1.194292in}{1.685910in}}%
\pgfpathlineto{\pgfqpoint{1.147667in}{1.739067in}}%
\pgfpathlineto{\pgfqpoint{1.101041in}{1.795649in}}%
\pgfpathlineto{\pgfqpoint{1.054415in}{1.854557in}}%
\pgfpathlineto{\pgfqpoint{1.007790in}{1.917575in}}%
\pgfpathlineto{\pgfqpoint{0.961164in}{1.984999in}}%
\pgfpathlineto{\pgfqpoint{0.914538in}{2.057136in}}%
\pgfpathlineto{\pgfqpoint{0.867913in}{2.134310in}}%
\pgfpathlineto{\pgfqpoint{0.821287in}{2.216851in}}%
\pgfpathlineto{\pgfqpoint{0.774661in}{2.305993in}}%
\pgfpathlineto{\pgfqpoint{0.728036in}{2.401573in}}%
\pgfpathlineto{\pgfqpoint{0.681410in}{2.503763in}}%
\pgfpathlineto{\pgfqpoint{0.634784in}{2.613836in}}%
\pgfpathlineto{\pgfqpoint{0.588158in}{2.731753in}}%
\pgfpathlineto{\pgfqpoint{0.588158in}{2.731753in}}%
\pgfpathclose%
\pgfusepath{fill}%
}%
\begin{pgfscope}%
\pgfsys@transformshift{0.000000in}{0.000000in}%
\pgfsys@useobject{currentmarker}{}%
\end{pgfscope}%
\end{pgfscope}%
\begin{pgfscope}%
\pgfpathrectangle{\pgfqpoint{0.462269in}{0.807389in}}{\pgfqpoint{2.769565in}{2.016000in}}%
\pgfusepath{clip}%
\pgfsetbuttcap%
\pgfsetroundjoin%
\definecolor{currentfill}{rgb}{0.839216,0.152941,0.156863}%
\pgfsetfillcolor{currentfill}%
\pgfsetfillopacity{0.150000}%
\pgfsetlinewidth{0.000000pt}%
\definecolor{currentstroke}{rgb}{0.839216,0.152941,0.156863}%
\pgfsetstrokecolor{currentstroke}%
\pgfsetstrokeopacity{0.150000}%
\pgfsetdash{}{0pt}%
\pgfsys@defobject{currentmarker}{\pgfqpoint{0.588158in}{0.927597in}}{\pgfqpoint{3.105945in}{1.797802in}}{%
\pgfpathmoveto{\pgfqpoint{0.588158in}{1.797802in}}%
\pgfpathlineto{\pgfqpoint{0.588158in}{1.733297in}}%
\pgfpathlineto{\pgfqpoint{0.634784in}{1.664784in}}%
\pgfpathlineto{\pgfqpoint{0.681410in}{1.602894in}}%
\pgfpathlineto{\pgfqpoint{0.728036in}{1.546549in}}%
\pgfpathlineto{\pgfqpoint{0.774661in}{1.495280in}}%
\pgfpathlineto{\pgfqpoint{0.821287in}{1.447952in}}%
\pgfpathlineto{\pgfqpoint{0.867913in}{1.405108in}}%
\pgfpathlineto{\pgfqpoint{0.914538in}{1.366083in}}%
\pgfpathlineto{\pgfqpoint{0.961164in}{1.331305in}}%
\pgfpathlineto{\pgfqpoint{1.007790in}{1.299943in}}%
\pgfpathlineto{\pgfqpoint{1.054415in}{1.271650in}}%
\pgfpathlineto{\pgfqpoint{1.101041in}{1.246149in}}%
\pgfpathlineto{\pgfqpoint{1.147667in}{1.223181in}}%
\pgfpathlineto{\pgfqpoint{1.194292in}{1.202505in}}%
\pgfpathlineto{\pgfqpoint{1.240918in}{1.183875in}}%
\pgfpathlineto{\pgfqpoint{1.287544in}{1.167064in}}%
\pgfpathlineto{\pgfqpoint{1.334169in}{1.151964in}}%
\pgfpathlineto{\pgfqpoint{1.380795in}{1.138406in}}%
\pgfpathlineto{\pgfqpoint{1.427421in}{1.126267in}}%
\pgfpathlineto{\pgfqpoint{1.474046in}{1.115453in}}%
\pgfpathlineto{\pgfqpoint{1.520672in}{1.105736in}}%
\pgfpathlineto{\pgfqpoint{1.567298in}{1.097008in}}%
\pgfpathlineto{\pgfqpoint{1.613923in}{1.089170in}}%
\pgfpathlineto{\pgfqpoint{1.660549in}{1.082133in}}%
\pgfpathlineto{\pgfqpoint{1.707175in}{1.075816in}}%
\pgfpathlineto{\pgfqpoint{1.753800in}{1.070147in}}%
\pgfpathlineto{\pgfqpoint{1.800426in}{1.065055in}}%
\pgfpathlineto{\pgfqpoint{1.847052in}{1.060482in}}%
\pgfpathlineto{\pgfqpoint{1.893677in}{1.056364in}}%
\pgfpathlineto{\pgfqpoint{1.940303in}{1.052672in}}%
\pgfpathlineto{\pgfqpoint{1.986929in}{1.049365in}}%
\pgfpathlineto{\pgfqpoint{2.033554in}{1.046403in}}%
\pgfpathlineto{\pgfqpoint{2.080180in}{1.043748in}}%
\pgfpathlineto{\pgfqpoint{2.126806in}{1.041364in}}%
\pgfpathlineto{\pgfqpoint{2.173432in}{1.039214in}}%
\pgfpathlineto{\pgfqpoint{2.220057in}{1.037250in}}%
\pgfpathlineto{\pgfqpoint{2.266683in}{1.035413in}}%
\pgfpathlineto{\pgfqpoint{2.313309in}{1.033615in}}%
\pgfpathlineto{\pgfqpoint{2.359934in}{1.031731in}}%
\pgfpathlineto{\pgfqpoint{2.406560in}{1.029594in}}%
\pgfpathlineto{\pgfqpoint{2.453186in}{1.026993in}}%
\pgfpathlineto{\pgfqpoint{2.499811in}{1.023691in}}%
\pgfpathlineto{\pgfqpoint{2.546437in}{1.019456in}}%
\pgfpathlineto{\pgfqpoint{2.593063in}{1.014117in}}%
\pgfpathlineto{\pgfqpoint{2.639688in}{1.007630in}}%
\pgfpathlineto{\pgfqpoint{2.686314in}{1.000064in}}%
\pgfpathlineto{\pgfqpoint{2.732940in}{0.991624in}}%
\pgfpathlineto{\pgfqpoint{2.779565in}{0.982628in}}%
\pgfpathlineto{\pgfqpoint{2.826191in}{0.973429in}}%
\pgfpathlineto{\pgfqpoint{2.872817in}{0.964367in}}%
\pgfpathlineto{\pgfqpoint{2.919442in}{0.955703in}}%
\pgfpathlineto{\pgfqpoint{2.966068in}{0.947623in}}%
\pgfpathlineto{\pgfqpoint{3.012694in}{0.940227in}}%
\pgfpathlineto{\pgfqpoint{3.059319in}{0.933554in}}%
\pgfpathlineto{\pgfqpoint{3.105945in}{0.927597in}}%
\pgfpathlineto{\pgfqpoint{3.105945in}{0.927604in}}%
\pgfpathlineto{\pgfqpoint{3.105945in}{0.927604in}}%
\pgfpathlineto{\pgfqpoint{3.059319in}{0.933563in}}%
\pgfpathlineto{\pgfqpoint{3.012694in}{0.940238in}}%
\pgfpathlineto{\pgfqpoint{2.966068in}{0.947635in}}%
\pgfpathlineto{\pgfqpoint{2.919442in}{0.955714in}}%
\pgfpathlineto{\pgfqpoint{2.872817in}{0.964384in}}%
\pgfpathlineto{\pgfqpoint{2.826191in}{0.973465in}}%
\pgfpathlineto{\pgfqpoint{2.779565in}{0.982688in}}%
\pgfpathlineto{\pgfqpoint{2.732940in}{0.991732in}}%
\pgfpathlineto{\pgfqpoint{2.686314in}{1.000247in}}%
\pgfpathlineto{\pgfqpoint{2.639688in}{1.007916in}}%
\pgfpathlineto{\pgfqpoint{2.593063in}{1.014510in}}%
\pgfpathlineto{\pgfqpoint{2.546437in}{1.019955in}}%
\pgfpathlineto{\pgfqpoint{2.499811in}{1.024314in}}%
\pgfpathlineto{\pgfqpoint{2.453186in}{1.027755in}}%
\pgfpathlineto{\pgfqpoint{2.406560in}{1.030503in}}%
\pgfpathlineto{\pgfqpoint{2.359934in}{1.032792in}}%
\pgfpathlineto{\pgfqpoint{2.313309in}{1.034834in}}%
\pgfpathlineto{\pgfqpoint{2.266683in}{1.036800in}}%
\pgfpathlineto{\pgfqpoint{2.220057in}{1.038816in}}%
\pgfpathlineto{\pgfqpoint{2.173432in}{1.040973in}}%
\pgfpathlineto{\pgfqpoint{2.126806in}{1.043335in}}%
\pgfpathlineto{\pgfqpoint{2.080180in}{1.045951in}}%
\pgfpathlineto{\pgfqpoint{2.033554in}{1.048864in}}%
\pgfpathlineto{\pgfqpoint{1.986929in}{1.052111in}}%
\pgfpathlineto{\pgfqpoint{1.940303in}{1.055734in}}%
\pgfpathlineto{\pgfqpoint{1.893677in}{1.059776in}}%
\pgfpathlineto{\pgfqpoint{1.847052in}{1.064284in}}%
\pgfpathlineto{\pgfqpoint{1.800426in}{1.069312in}}%
\pgfpathlineto{\pgfqpoint{1.753800in}{1.074919in}}%
\pgfpathlineto{\pgfqpoint{1.707175in}{1.081169in}}%
\pgfpathlineto{\pgfqpoint{1.660549in}{1.088134in}}%
\pgfpathlineto{\pgfqpoint{1.613923in}{1.095895in}}%
\pgfpathlineto{\pgfqpoint{1.567298in}{1.104538in}}%
\pgfpathlineto{\pgfqpoint{1.520672in}{1.114184in}}%
\pgfpathlineto{\pgfqpoint{1.474046in}{1.124938in}}%
\pgfpathlineto{\pgfqpoint{1.427421in}{1.136910in}}%
\pgfpathlineto{\pgfqpoint{1.380795in}{1.150231in}}%
\pgfpathlineto{\pgfqpoint{1.334169in}{1.165185in}}%
\pgfpathlineto{\pgfqpoint{1.287544in}{1.181861in}}%
\pgfpathlineto{\pgfqpoint{1.240918in}{1.200424in}}%
\pgfpathlineto{\pgfqpoint{1.194292in}{1.221076in}}%
\pgfpathlineto{\pgfqpoint{1.147667in}{1.244037in}}%
\pgfpathlineto{\pgfqpoint{1.101041in}{1.269368in}}%
\pgfpathlineto{\pgfqpoint{1.054415in}{1.296753in}}%
\pgfpathlineto{\pgfqpoint{1.007790in}{1.327080in}}%
\pgfpathlineto{\pgfqpoint{0.961164in}{1.360579in}}%
\pgfpathlineto{\pgfqpoint{0.914538in}{1.397591in}}%
\pgfpathlineto{\pgfqpoint{0.867913in}{1.439214in}}%
\pgfpathlineto{\pgfqpoint{0.821287in}{1.485466in}}%
\pgfpathlineto{\pgfqpoint{0.774661in}{1.536535in}}%
\pgfpathlineto{\pgfqpoint{0.728036in}{1.592850in}}%
\pgfpathlineto{\pgfqpoint{0.681410in}{1.654857in}}%
\pgfpathlineto{\pgfqpoint{0.634784in}{1.723018in}}%
\pgfpathlineto{\pgfqpoint{0.588158in}{1.797802in}}%
\pgfpathlineto{\pgfqpoint{0.588158in}{1.797802in}}%
\pgfpathclose%
\pgfusepath{fill}%
}%
\begin{pgfscope}%
\pgfsys@transformshift{0.000000in}{0.000000in}%
\pgfsys@useobject{currentmarker}{}%
\end{pgfscope}%
\end{pgfscope}%
\begin{pgfscope}%
\pgfpathrectangle{\pgfqpoint{0.462269in}{0.807389in}}{\pgfqpoint{2.769565in}{2.016000in}}%
\pgfusepath{clip}%
\pgfsetbuttcap%
\pgfsetroundjoin%
\definecolor{currentfill}{rgb}{0.580392,0.403922,0.741176}%
\pgfsetfillcolor{currentfill}%
\pgfsetfillopacity{0.150000}%
\pgfsetlinewidth{0.000000pt}%
\definecolor{currentstroke}{rgb}{0.580392,0.403922,0.741176}%
\pgfsetstrokecolor{currentstroke}%
\pgfsetstrokeopacity{0.150000}%
\pgfsetdash{}{0pt}%
\pgfsys@defobject{currentmarker}{\pgfqpoint{0.588158in}{0.927628in}}{\pgfqpoint{3.105945in}{1.677321in}}{%
\pgfpathmoveto{\pgfqpoint{0.588158in}{1.677321in}}%
\pgfpathlineto{\pgfqpoint{0.588158in}{1.677321in}}%
\pgfpathlineto{\pgfqpoint{0.634784in}{1.677321in}}%
\pgfpathlineto{\pgfqpoint{0.681410in}{1.677321in}}%
\pgfpathlineto{\pgfqpoint{0.728036in}{1.677321in}}%
\pgfpathlineto{\pgfqpoint{0.774661in}{1.677321in}}%
\pgfpathlineto{\pgfqpoint{0.821287in}{1.677321in}}%
\pgfpathlineto{\pgfqpoint{0.867913in}{1.677321in}}%
\pgfpathlineto{\pgfqpoint{0.914538in}{1.677321in}}%
\pgfpathlineto{\pgfqpoint{0.961164in}{1.677321in}}%
\pgfpathlineto{\pgfqpoint{1.007790in}{1.677321in}}%
\pgfpathlineto{\pgfqpoint{1.054415in}{1.677321in}}%
\pgfpathlineto{\pgfqpoint{1.101041in}{1.677321in}}%
\pgfpathlineto{\pgfqpoint{1.147667in}{1.677321in}}%
\pgfpathlineto{\pgfqpoint{1.194292in}{1.677321in}}%
\pgfpathlineto{\pgfqpoint{1.240918in}{1.677321in}}%
\pgfpathlineto{\pgfqpoint{1.287544in}{1.677319in}}%
\pgfpathlineto{\pgfqpoint{1.334169in}{1.677313in}}%
\pgfpathlineto{\pgfqpoint{1.380795in}{1.677293in}}%
\pgfpathlineto{\pgfqpoint{1.427421in}{1.677233in}}%
\pgfpathlineto{\pgfqpoint{1.474046in}{1.677072in}}%
\pgfpathlineto{\pgfqpoint{1.520672in}{1.676685in}}%
\pgfpathlineto{\pgfqpoint{1.567298in}{1.675839in}}%
\pgfpathlineto{\pgfqpoint{1.613923in}{1.674140in}}%
\pgfpathlineto{\pgfqpoint{1.660549in}{1.670990in}}%
\pgfpathlineto{\pgfqpoint{1.707175in}{1.665555in}}%
\pgfpathlineto{\pgfqpoint{1.753800in}{1.656793in}}%
\pgfpathlineto{\pgfqpoint{1.800426in}{1.643552in}}%
\pgfpathlineto{\pgfqpoint{1.847052in}{1.624756in}}%
\pgfpathlineto{\pgfqpoint{1.893677in}{1.599651in}}%
\pgfpathlineto{\pgfqpoint{1.940303in}{1.568046in}}%
\pgfpathlineto{\pgfqpoint{1.986929in}{1.530459in}}%
\pgfpathlineto{\pgfqpoint{2.033554in}{1.488081in}}%
\pgfpathlineto{\pgfqpoint{2.080180in}{1.442569in}}%
\pgfpathlineto{\pgfqpoint{2.126806in}{1.395736in}}%
\pgfpathlineto{\pgfqpoint{2.173432in}{1.349253in}}%
\pgfpathlineto{\pgfqpoint{2.220057in}{1.304448in}}%
\pgfpathlineto{\pgfqpoint{2.266683in}{1.262231in}}%
\pgfpathlineto{\pgfqpoint{2.313309in}{1.223120in}}%
\pgfpathlineto{\pgfqpoint{2.359934in}{1.187323in}}%
\pgfpathlineto{\pgfqpoint{2.406560in}{1.154827in}}%
\pgfpathlineto{\pgfqpoint{2.453186in}{1.125491in}}%
\pgfpathlineto{\pgfqpoint{2.499811in}{1.099100in}}%
\pgfpathlineto{\pgfqpoint{2.546437in}{1.075411in}}%
\pgfpathlineto{\pgfqpoint{2.593063in}{1.054179in}}%
\pgfpathlineto{\pgfqpoint{2.639688in}{1.035167in}}%
\pgfpathlineto{\pgfqpoint{2.686314in}{1.018154in}}%
\pgfpathlineto{\pgfqpoint{2.732940in}{1.002940in}}%
\pgfpathlineto{\pgfqpoint{2.779565in}{0.989342in}}%
\pgfpathlineto{\pgfqpoint{2.826191in}{0.977196in}}%
\pgfpathlineto{\pgfqpoint{2.872817in}{0.966356in}}%
\pgfpathlineto{\pgfqpoint{2.919442in}{0.956689in}}%
\pgfpathlineto{\pgfqpoint{2.966068in}{0.948080in}}%
\pgfpathlineto{\pgfqpoint{3.012694in}{0.940425in}}%
\pgfpathlineto{\pgfqpoint{3.059319in}{0.933634in}}%
\pgfpathlineto{\pgfqpoint{3.105945in}{0.927628in}}%
\pgfpathlineto{\pgfqpoint{3.105945in}{0.927628in}}%
\pgfpathlineto{\pgfqpoint{3.105945in}{0.927628in}}%
\pgfpathlineto{\pgfqpoint{3.059319in}{0.933634in}}%
\pgfpathlineto{\pgfqpoint{3.012694in}{0.940425in}}%
\pgfpathlineto{\pgfqpoint{2.966068in}{0.948080in}}%
\pgfpathlineto{\pgfqpoint{2.919442in}{0.956689in}}%
\pgfpathlineto{\pgfqpoint{2.872817in}{0.966356in}}%
\pgfpathlineto{\pgfqpoint{2.826191in}{0.977196in}}%
\pgfpathlineto{\pgfqpoint{2.779565in}{0.989342in}}%
\pgfpathlineto{\pgfqpoint{2.732940in}{1.002940in}}%
\pgfpathlineto{\pgfqpoint{2.686314in}{1.018154in}}%
\pgfpathlineto{\pgfqpoint{2.639688in}{1.035167in}}%
\pgfpathlineto{\pgfqpoint{2.593063in}{1.054179in}}%
\pgfpathlineto{\pgfqpoint{2.546437in}{1.075411in}}%
\pgfpathlineto{\pgfqpoint{2.499811in}{1.099100in}}%
\pgfpathlineto{\pgfqpoint{2.453186in}{1.125491in}}%
\pgfpathlineto{\pgfqpoint{2.406560in}{1.154827in}}%
\pgfpathlineto{\pgfqpoint{2.359934in}{1.187323in}}%
\pgfpathlineto{\pgfqpoint{2.313309in}{1.223120in}}%
\pgfpathlineto{\pgfqpoint{2.266683in}{1.262231in}}%
\pgfpathlineto{\pgfqpoint{2.220057in}{1.304448in}}%
\pgfpathlineto{\pgfqpoint{2.173432in}{1.349253in}}%
\pgfpathlineto{\pgfqpoint{2.126806in}{1.395736in}}%
\pgfpathlineto{\pgfqpoint{2.080180in}{1.442569in}}%
\pgfpathlineto{\pgfqpoint{2.033554in}{1.488081in}}%
\pgfpathlineto{\pgfqpoint{1.986929in}{1.530459in}}%
\pgfpathlineto{\pgfqpoint{1.940303in}{1.568046in}}%
\pgfpathlineto{\pgfqpoint{1.893677in}{1.599651in}}%
\pgfpathlineto{\pgfqpoint{1.847052in}{1.624756in}}%
\pgfpathlineto{\pgfqpoint{1.800426in}{1.643552in}}%
\pgfpathlineto{\pgfqpoint{1.753800in}{1.656793in}}%
\pgfpathlineto{\pgfqpoint{1.707175in}{1.665555in}}%
\pgfpathlineto{\pgfqpoint{1.660549in}{1.670990in}}%
\pgfpathlineto{\pgfqpoint{1.613923in}{1.674140in}}%
\pgfpathlineto{\pgfqpoint{1.567298in}{1.675839in}}%
\pgfpathlineto{\pgfqpoint{1.520672in}{1.676685in}}%
\pgfpathlineto{\pgfqpoint{1.474046in}{1.677072in}}%
\pgfpathlineto{\pgfqpoint{1.427421in}{1.677233in}}%
\pgfpathlineto{\pgfqpoint{1.380795in}{1.677293in}}%
\pgfpathlineto{\pgfqpoint{1.334169in}{1.677313in}}%
\pgfpathlineto{\pgfqpoint{1.287544in}{1.677319in}}%
\pgfpathlineto{\pgfqpoint{1.240918in}{1.677321in}}%
\pgfpathlineto{\pgfqpoint{1.194292in}{1.677321in}}%
\pgfpathlineto{\pgfqpoint{1.147667in}{1.677321in}}%
\pgfpathlineto{\pgfqpoint{1.101041in}{1.677321in}}%
\pgfpathlineto{\pgfqpoint{1.054415in}{1.677321in}}%
\pgfpathlineto{\pgfqpoint{1.007790in}{1.677321in}}%
\pgfpathlineto{\pgfqpoint{0.961164in}{1.677321in}}%
\pgfpathlineto{\pgfqpoint{0.914538in}{1.677321in}}%
\pgfpathlineto{\pgfqpoint{0.867913in}{1.677321in}}%
\pgfpathlineto{\pgfqpoint{0.821287in}{1.677321in}}%
\pgfpathlineto{\pgfqpoint{0.774661in}{1.677321in}}%
\pgfpathlineto{\pgfqpoint{0.728036in}{1.677321in}}%
\pgfpathlineto{\pgfqpoint{0.681410in}{1.677321in}}%
\pgfpathlineto{\pgfqpoint{0.634784in}{1.677321in}}%
\pgfpathlineto{\pgfqpoint{0.588158in}{1.677321in}}%
\pgfpathlineto{\pgfqpoint{0.588158in}{1.677321in}}%
\pgfpathclose%
\pgfusepath{fill}%
}%
\begin{pgfscope}%
\pgfsys@transformshift{0.000000in}{0.000000in}%
\pgfsys@useobject{currentmarker}{}%
\end{pgfscope}%
\end{pgfscope}%
\begin{pgfscope}%
\pgfpathrectangle{\pgfqpoint{0.462269in}{0.807389in}}{\pgfqpoint{2.769565in}{2.016000in}}%
\pgfusepath{clip}%
\pgfsetrectcap%
\pgfsetroundjoin%
\pgfsetlinewidth{0.803000pt}%
\definecolor{currentstroke}{rgb}{0.690196,0.690196,0.690196}%
\pgfsetstrokecolor{currentstroke}%
\pgfsetstrokeopacity{0.300000}%
\pgfsetdash{}{0pt}%
\pgfpathmoveto{\pgfqpoint{1.170719in}{0.807389in}}%
\pgfpathlineto{\pgfqpoint{1.170719in}{2.823389in}}%
\pgfusepath{stroke}%
\end{pgfscope}%
\begin{pgfscope}%
\pgfsetbuttcap%
\pgfsetroundjoin%
\definecolor{currentfill}{rgb}{0.000000,0.000000,0.000000}%
\pgfsetfillcolor{currentfill}%
\pgfsetlinewidth{0.803000pt}%
\definecolor{currentstroke}{rgb}{0.000000,0.000000,0.000000}%
\pgfsetstrokecolor{currentstroke}%
\pgfsetdash{}{0pt}%
\pgfsys@defobject{currentmarker}{\pgfqpoint{0.000000in}{-0.048611in}}{\pgfqpoint{0.000000in}{0.000000in}}{%
\pgfpathmoveto{\pgfqpoint{0.000000in}{0.000000in}}%
\pgfpathlineto{\pgfqpoint{0.000000in}{-0.048611in}}%
\pgfusepath{stroke,fill}%
}%
\begin{pgfscope}%
\pgfsys@transformshift{1.170719in}{0.807389in}%
\pgfsys@useobject{currentmarker}{}%
\end{pgfscope}%
\end{pgfscope}%
\begin{pgfscope}%
\definecolor{textcolor}{rgb}{0.000000,0.000000,0.000000}%
\pgfsetstrokecolor{textcolor}%
\pgfsetfillcolor{textcolor}%
\pgftext[x=1.170719in,y=0.710167in,,top]{\color{textcolor}{\rmfamily\fontsize{7.000000}{8.400000}\selectfont\catcode`\^=\active\def^{\ifmmode\sp\else\^{}\fi}\catcode`\%=\active\def%{\%}$\mathdefault{10^{1}}$}}%
\end{pgfscope}%
\begin{pgfscope}%
\pgfpathrectangle{\pgfqpoint{0.462269in}{0.807389in}}{\pgfqpoint{2.769565in}{2.016000in}}%
\pgfusepath{clip}%
\pgfsetrectcap%
\pgfsetroundjoin%
\pgfsetlinewidth{0.803000pt}%
\definecolor{currentstroke}{rgb}{0.690196,0.690196,0.690196}%
\pgfsetstrokecolor{currentstroke}%
\pgfsetstrokeopacity{0.300000}%
\pgfsetdash{}{0pt}%
\pgfpathmoveto{\pgfqpoint{3.105945in}{0.807389in}}%
\pgfpathlineto{\pgfqpoint{3.105945in}{2.823389in}}%
\pgfusepath{stroke}%
\end{pgfscope}%
\begin{pgfscope}%
\pgfsetbuttcap%
\pgfsetroundjoin%
\definecolor{currentfill}{rgb}{0.000000,0.000000,0.000000}%
\pgfsetfillcolor{currentfill}%
\pgfsetlinewidth{0.803000pt}%
\definecolor{currentstroke}{rgb}{0.000000,0.000000,0.000000}%
\pgfsetstrokecolor{currentstroke}%
\pgfsetdash{}{0pt}%
\pgfsys@defobject{currentmarker}{\pgfqpoint{0.000000in}{-0.048611in}}{\pgfqpoint{0.000000in}{0.000000in}}{%
\pgfpathmoveto{\pgfqpoint{0.000000in}{0.000000in}}%
\pgfpathlineto{\pgfqpoint{0.000000in}{-0.048611in}}%
\pgfusepath{stroke,fill}%
}%
\begin{pgfscope}%
\pgfsys@transformshift{3.105945in}{0.807389in}%
\pgfsys@useobject{currentmarker}{}%
\end{pgfscope}%
\end{pgfscope}%
\begin{pgfscope}%
\definecolor{textcolor}{rgb}{0.000000,0.000000,0.000000}%
\pgfsetstrokecolor{textcolor}%
\pgfsetfillcolor{textcolor}%
\pgftext[x=3.105945in,y=0.710167in,,top]{\color{textcolor}{\rmfamily\fontsize{7.000000}{8.400000}\selectfont\catcode`\^=\active\def^{\ifmmode\sp\else\^{}\fi}\catcode`\%=\active\def%{\%}$\mathdefault{10^{2}}$}}%
\end{pgfscope}%
\begin{pgfscope}%
\pgfsetbuttcap%
\pgfsetroundjoin%
\definecolor{currentfill}{rgb}{0.000000,0.000000,0.000000}%
\pgfsetfillcolor{currentfill}%
\pgfsetlinewidth{0.602250pt}%
\definecolor{currentstroke}{rgb}{0.000000,0.000000,0.000000}%
\pgfsetstrokecolor{currentstroke}%
\pgfsetdash{}{0pt}%
\pgfsys@defobject{currentmarker}{\pgfqpoint{0.000000in}{-0.027778in}}{\pgfqpoint{0.000000in}{0.000000in}}{%
\pgfpathmoveto{\pgfqpoint{0.000000in}{0.000000in}}%
\pgfpathlineto{\pgfqpoint{0.000000in}{-0.027778in}}%
\pgfusepath{stroke,fill}%
}%
\begin{pgfscope}%
\pgfsys@transformshift{0.588158in}{0.807389in}%
\pgfsys@useobject{currentmarker}{}%
\end{pgfscope}%
\end{pgfscope}%
\begin{pgfscope}%
\pgfsetbuttcap%
\pgfsetroundjoin%
\definecolor{currentfill}{rgb}{0.000000,0.000000,0.000000}%
\pgfsetfillcolor{currentfill}%
\pgfsetlinewidth{0.602250pt}%
\definecolor{currentstroke}{rgb}{0.000000,0.000000,0.000000}%
\pgfsetstrokecolor{currentstroke}%
\pgfsetdash{}{0pt}%
\pgfsys@defobject{currentmarker}{\pgfqpoint{0.000000in}{-0.027778in}}{\pgfqpoint{0.000000in}{0.000000in}}{%
\pgfpathmoveto{\pgfqpoint{0.000000in}{0.000000in}}%
\pgfpathlineto{\pgfqpoint{0.000000in}{-0.027778in}}%
\pgfusepath{stroke,fill}%
}%
\begin{pgfscope}%
\pgfsys@transformshift{0.741392in}{0.807389in}%
\pgfsys@useobject{currentmarker}{}%
\end{pgfscope}%
\end{pgfscope}%
\begin{pgfscope}%
\pgfsetbuttcap%
\pgfsetroundjoin%
\definecolor{currentfill}{rgb}{0.000000,0.000000,0.000000}%
\pgfsetfillcolor{currentfill}%
\pgfsetlinewidth{0.602250pt}%
\definecolor{currentstroke}{rgb}{0.000000,0.000000,0.000000}%
\pgfsetstrokecolor{currentstroke}%
\pgfsetdash{}{0pt}%
\pgfsys@defobject{currentmarker}{\pgfqpoint{0.000000in}{-0.027778in}}{\pgfqpoint{0.000000in}{0.000000in}}{%
\pgfpathmoveto{\pgfqpoint{0.000000in}{0.000000in}}%
\pgfpathlineto{\pgfqpoint{0.000000in}{-0.027778in}}%
\pgfusepath{stroke,fill}%
}%
\begin{pgfscope}%
\pgfsys@transformshift{0.870949in}{0.807389in}%
\pgfsys@useobject{currentmarker}{}%
\end{pgfscope}%
\end{pgfscope}%
\begin{pgfscope}%
\pgfsetbuttcap%
\pgfsetroundjoin%
\definecolor{currentfill}{rgb}{0.000000,0.000000,0.000000}%
\pgfsetfillcolor{currentfill}%
\pgfsetlinewidth{0.602250pt}%
\definecolor{currentstroke}{rgb}{0.000000,0.000000,0.000000}%
\pgfsetstrokecolor{currentstroke}%
\pgfsetdash{}{0pt}%
\pgfsys@defobject{currentmarker}{\pgfqpoint{0.000000in}{-0.027778in}}{\pgfqpoint{0.000000in}{0.000000in}}{%
\pgfpathmoveto{\pgfqpoint{0.000000in}{0.000000in}}%
\pgfpathlineto{\pgfqpoint{0.000000in}{-0.027778in}}%
\pgfusepath{stroke,fill}%
}%
\begin{pgfscope}%
\pgfsys@transformshift{0.983177in}{0.807389in}%
\pgfsys@useobject{currentmarker}{}%
\end{pgfscope}%
\end{pgfscope}%
\begin{pgfscope}%
\pgfsetbuttcap%
\pgfsetroundjoin%
\definecolor{currentfill}{rgb}{0.000000,0.000000,0.000000}%
\pgfsetfillcolor{currentfill}%
\pgfsetlinewidth{0.602250pt}%
\definecolor{currentstroke}{rgb}{0.000000,0.000000,0.000000}%
\pgfsetstrokecolor{currentstroke}%
\pgfsetdash{}{0pt}%
\pgfsys@defobject{currentmarker}{\pgfqpoint{0.000000in}{-0.027778in}}{\pgfqpoint{0.000000in}{0.000000in}}{%
\pgfpathmoveto{\pgfqpoint{0.000000in}{0.000000in}}%
\pgfpathlineto{\pgfqpoint{0.000000in}{-0.027778in}}%
\pgfusepath{stroke,fill}%
}%
\begin{pgfscope}%
\pgfsys@transformshift{1.082168in}{0.807389in}%
\pgfsys@useobject{currentmarker}{}%
\end{pgfscope}%
\end{pgfscope}%
\begin{pgfscope}%
\pgfsetbuttcap%
\pgfsetroundjoin%
\definecolor{currentfill}{rgb}{0.000000,0.000000,0.000000}%
\pgfsetfillcolor{currentfill}%
\pgfsetlinewidth{0.602250pt}%
\definecolor{currentstroke}{rgb}{0.000000,0.000000,0.000000}%
\pgfsetstrokecolor{currentstroke}%
\pgfsetdash{}{0pt}%
\pgfsys@defobject{currentmarker}{\pgfqpoint{0.000000in}{-0.027778in}}{\pgfqpoint{0.000000in}{0.000000in}}{%
\pgfpathmoveto{\pgfqpoint{0.000000in}{0.000000in}}%
\pgfpathlineto{\pgfqpoint{0.000000in}{-0.027778in}}%
\pgfusepath{stroke,fill}%
}%
\begin{pgfscope}%
\pgfsys@transformshift{1.753280in}{0.807389in}%
\pgfsys@useobject{currentmarker}{}%
\end{pgfscope}%
\end{pgfscope}%
\begin{pgfscope}%
\pgfsetbuttcap%
\pgfsetroundjoin%
\definecolor{currentfill}{rgb}{0.000000,0.000000,0.000000}%
\pgfsetfillcolor{currentfill}%
\pgfsetlinewidth{0.602250pt}%
\definecolor{currentstroke}{rgb}{0.000000,0.000000,0.000000}%
\pgfsetstrokecolor{currentstroke}%
\pgfsetdash{}{0pt}%
\pgfsys@defobject{currentmarker}{\pgfqpoint{0.000000in}{-0.027778in}}{\pgfqpoint{0.000000in}{0.000000in}}{%
\pgfpathmoveto{\pgfqpoint{0.000000in}{0.000000in}}%
\pgfpathlineto{\pgfqpoint{0.000000in}{-0.027778in}}%
\pgfusepath{stroke,fill}%
}%
\begin{pgfscope}%
\pgfsys@transformshift{2.094057in}{0.807389in}%
\pgfsys@useobject{currentmarker}{}%
\end{pgfscope}%
\end{pgfscope}%
\begin{pgfscope}%
\pgfsetbuttcap%
\pgfsetroundjoin%
\definecolor{currentfill}{rgb}{0.000000,0.000000,0.000000}%
\pgfsetfillcolor{currentfill}%
\pgfsetlinewidth{0.602250pt}%
\definecolor{currentstroke}{rgb}{0.000000,0.000000,0.000000}%
\pgfsetstrokecolor{currentstroke}%
\pgfsetdash{}{0pt}%
\pgfsys@defobject{currentmarker}{\pgfqpoint{0.000000in}{-0.027778in}}{\pgfqpoint{0.000000in}{0.000000in}}{%
\pgfpathmoveto{\pgfqpoint{0.000000in}{0.000000in}}%
\pgfpathlineto{\pgfqpoint{0.000000in}{-0.027778in}}%
\pgfusepath{stroke,fill}%
}%
\begin{pgfscope}%
\pgfsys@transformshift{2.335841in}{0.807389in}%
\pgfsys@useobject{currentmarker}{}%
\end{pgfscope}%
\end{pgfscope}%
\begin{pgfscope}%
\pgfsetbuttcap%
\pgfsetroundjoin%
\definecolor{currentfill}{rgb}{0.000000,0.000000,0.000000}%
\pgfsetfillcolor{currentfill}%
\pgfsetlinewidth{0.602250pt}%
\definecolor{currentstroke}{rgb}{0.000000,0.000000,0.000000}%
\pgfsetstrokecolor{currentstroke}%
\pgfsetdash{}{0pt}%
\pgfsys@defobject{currentmarker}{\pgfqpoint{0.000000in}{-0.027778in}}{\pgfqpoint{0.000000in}{0.000000in}}{%
\pgfpathmoveto{\pgfqpoint{0.000000in}{0.000000in}}%
\pgfpathlineto{\pgfqpoint{0.000000in}{-0.027778in}}%
\pgfusepath{stroke,fill}%
}%
\begin{pgfscope}%
\pgfsys@transformshift{2.523384in}{0.807389in}%
\pgfsys@useobject{currentmarker}{}%
\end{pgfscope}%
\end{pgfscope}%
\begin{pgfscope}%
\pgfsetbuttcap%
\pgfsetroundjoin%
\definecolor{currentfill}{rgb}{0.000000,0.000000,0.000000}%
\pgfsetfillcolor{currentfill}%
\pgfsetlinewidth{0.602250pt}%
\definecolor{currentstroke}{rgb}{0.000000,0.000000,0.000000}%
\pgfsetstrokecolor{currentstroke}%
\pgfsetdash{}{0pt}%
\pgfsys@defobject{currentmarker}{\pgfqpoint{0.000000in}{-0.027778in}}{\pgfqpoint{0.000000in}{0.000000in}}{%
\pgfpathmoveto{\pgfqpoint{0.000000in}{0.000000in}}%
\pgfpathlineto{\pgfqpoint{0.000000in}{-0.027778in}}%
\pgfusepath{stroke,fill}%
}%
\begin{pgfscope}%
\pgfsys@transformshift{2.676618in}{0.807389in}%
\pgfsys@useobject{currentmarker}{}%
\end{pgfscope}%
\end{pgfscope}%
\begin{pgfscope}%
\pgfsetbuttcap%
\pgfsetroundjoin%
\definecolor{currentfill}{rgb}{0.000000,0.000000,0.000000}%
\pgfsetfillcolor{currentfill}%
\pgfsetlinewidth{0.602250pt}%
\definecolor{currentstroke}{rgb}{0.000000,0.000000,0.000000}%
\pgfsetstrokecolor{currentstroke}%
\pgfsetdash{}{0pt}%
\pgfsys@defobject{currentmarker}{\pgfqpoint{0.000000in}{-0.027778in}}{\pgfqpoint{0.000000in}{0.000000in}}{%
\pgfpathmoveto{\pgfqpoint{0.000000in}{0.000000in}}%
\pgfpathlineto{\pgfqpoint{0.000000in}{-0.027778in}}%
\pgfusepath{stroke,fill}%
}%
\begin{pgfscope}%
\pgfsys@transformshift{2.806175in}{0.807389in}%
\pgfsys@useobject{currentmarker}{}%
\end{pgfscope}%
\end{pgfscope}%
\begin{pgfscope}%
\pgfsetbuttcap%
\pgfsetroundjoin%
\definecolor{currentfill}{rgb}{0.000000,0.000000,0.000000}%
\pgfsetfillcolor{currentfill}%
\pgfsetlinewidth{0.602250pt}%
\definecolor{currentstroke}{rgb}{0.000000,0.000000,0.000000}%
\pgfsetstrokecolor{currentstroke}%
\pgfsetdash{}{0pt}%
\pgfsys@defobject{currentmarker}{\pgfqpoint{0.000000in}{-0.027778in}}{\pgfqpoint{0.000000in}{0.000000in}}{%
\pgfpathmoveto{\pgfqpoint{0.000000in}{0.000000in}}%
\pgfpathlineto{\pgfqpoint{0.000000in}{-0.027778in}}%
\pgfusepath{stroke,fill}%
}%
\begin{pgfscope}%
\pgfsys@transformshift{2.918402in}{0.807389in}%
\pgfsys@useobject{currentmarker}{}%
\end{pgfscope}%
\end{pgfscope}%
\begin{pgfscope}%
\pgfsetbuttcap%
\pgfsetroundjoin%
\definecolor{currentfill}{rgb}{0.000000,0.000000,0.000000}%
\pgfsetfillcolor{currentfill}%
\pgfsetlinewidth{0.602250pt}%
\definecolor{currentstroke}{rgb}{0.000000,0.000000,0.000000}%
\pgfsetstrokecolor{currentstroke}%
\pgfsetdash{}{0pt}%
\pgfsys@defobject{currentmarker}{\pgfqpoint{0.000000in}{-0.027778in}}{\pgfqpoint{0.000000in}{0.000000in}}{%
\pgfpathmoveto{\pgfqpoint{0.000000in}{0.000000in}}%
\pgfpathlineto{\pgfqpoint{0.000000in}{-0.027778in}}%
\pgfusepath{stroke,fill}%
}%
\begin{pgfscope}%
\pgfsys@transformshift{3.017394in}{0.807389in}%
\pgfsys@useobject{currentmarker}{}%
\end{pgfscope}%
\end{pgfscope}%
\begin{pgfscope}%
\definecolor{textcolor}{rgb}{0.000000,0.000000,0.000000}%
\pgfsetstrokecolor{textcolor}%
\pgfsetfillcolor{textcolor}%
\pgftext[x=1.847052in,y=0.561000in,,top]{\color{textcolor}{\rmfamily\fontsize{8.000000}{9.600000}\selectfont\catcode`\^=\active\def^{\ifmmode\sp\else\^{}\fi}\catcode`\%=\active\def%{\%}Angular Scale $\alpha$ (Degrees)}}%
\end{pgfscope}%
\begin{pgfscope}%
\pgfpathrectangle{\pgfqpoint{0.462269in}{0.807389in}}{\pgfqpoint{2.769565in}{2.016000in}}%
\pgfusepath{clip}%
\pgfsetrectcap%
\pgfsetroundjoin%
\pgfsetlinewidth{0.803000pt}%
\definecolor{currentstroke}{rgb}{0.690196,0.690196,0.690196}%
\pgfsetstrokecolor{currentstroke}%
\pgfsetstrokeopacity{0.300000}%
\pgfsetdash{}{0pt}%
\pgfpathmoveto{\pgfqpoint{0.462269in}{0.857690in}}%
\pgfpathlineto{\pgfqpoint{3.231834in}{0.857690in}}%
\pgfusepath{stroke}%
\end{pgfscope}%
\begin{pgfscope}%
\pgfsetbuttcap%
\pgfsetroundjoin%
\definecolor{currentfill}{rgb}{0.000000,0.000000,0.000000}%
\pgfsetfillcolor{currentfill}%
\pgfsetlinewidth{0.803000pt}%
\definecolor{currentstroke}{rgb}{0.000000,0.000000,0.000000}%
\pgfsetstrokecolor{currentstroke}%
\pgfsetdash{}{0pt}%
\pgfsys@defobject{currentmarker}{\pgfqpoint{-0.048611in}{0.000000in}}{\pgfqpoint{-0.000000in}{0.000000in}}{%
\pgfpathmoveto{\pgfqpoint{-0.000000in}{0.000000in}}%
\pgfpathlineto{\pgfqpoint{-0.048611in}{0.000000in}}%
\pgfusepath{stroke,fill}%
}%
\begin{pgfscope}%
\pgfsys@transformshift{0.462269in}{0.857690in}%
\pgfsys@useobject{currentmarker}{}%
\end{pgfscope}%
\end{pgfscope}%
\begin{pgfscope}%
\definecolor{textcolor}{rgb}{0.000000,0.000000,0.000000}%
\pgfsetstrokecolor{textcolor}%
\pgfsetfillcolor{textcolor}%
\pgftext[x=0.309684in, y=0.823932in, left, base]{\color{textcolor}{\rmfamily\fontsize{7.000000}{8.400000}\selectfont\catcode`\^=\active\def^{\ifmmode\sp\else\^{}\fi}\catcode`\%=\active\def%{\%}0}}%
\end{pgfscope}%
\begin{pgfscope}%
\pgfpathrectangle{\pgfqpoint{0.462269in}{0.807389in}}{\pgfqpoint{2.769565in}{2.016000in}}%
\pgfusepath{clip}%
\pgfsetrectcap%
\pgfsetroundjoin%
\pgfsetlinewidth{0.803000pt}%
\definecolor{currentstroke}{rgb}{0.690196,0.690196,0.690196}%
\pgfsetstrokecolor{currentstroke}%
\pgfsetstrokeopacity{0.300000}%
\pgfsetdash{}{0pt}%
\pgfpathmoveto{\pgfqpoint{0.462269in}{1.267506in}}%
\pgfpathlineto{\pgfqpoint{3.231834in}{1.267506in}}%
\pgfusepath{stroke}%
\end{pgfscope}%
\begin{pgfscope}%
\pgfsetbuttcap%
\pgfsetroundjoin%
\definecolor{currentfill}{rgb}{0.000000,0.000000,0.000000}%
\pgfsetfillcolor{currentfill}%
\pgfsetlinewidth{0.803000pt}%
\definecolor{currentstroke}{rgb}{0.000000,0.000000,0.000000}%
\pgfsetstrokecolor{currentstroke}%
\pgfsetdash{}{0pt}%
\pgfsys@defobject{currentmarker}{\pgfqpoint{-0.048611in}{0.000000in}}{\pgfqpoint{-0.000000in}{0.000000in}}{%
\pgfpathmoveto{\pgfqpoint{-0.000000in}{0.000000in}}%
\pgfpathlineto{\pgfqpoint{-0.048611in}{0.000000in}}%
\pgfusepath{stroke,fill}%
}%
\begin{pgfscope}%
\pgfsys@transformshift{0.462269in}{1.267506in}%
\pgfsys@useobject{currentmarker}{}%
\end{pgfscope}%
\end{pgfscope}%
\begin{pgfscope}%
\definecolor{textcolor}{rgb}{0.000000,0.000000,0.000000}%
\pgfsetstrokecolor{textcolor}%
\pgfsetfillcolor{textcolor}%
\pgftext[x=0.254321in, y=1.233748in, left, base]{\color{textcolor}{\rmfamily\fontsize{7.000000}{8.400000}\selectfont\catcode`\^=\active\def^{\ifmmode\sp\else\^{}\fi}\catcode`\%=\active\def%{\%}10}}%
\end{pgfscope}%
\begin{pgfscope}%
\pgfpathrectangle{\pgfqpoint{0.462269in}{0.807389in}}{\pgfqpoint{2.769565in}{2.016000in}}%
\pgfusepath{clip}%
\pgfsetrectcap%
\pgfsetroundjoin%
\pgfsetlinewidth{0.803000pt}%
\definecolor{currentstroke}{rgb}{0.690196,0.690196,0.690196}%
\pgfsetstrokecolor{currentstroke}%
\pgfsetstrokeopacity{0.300000}%
\pgfsetdash{}{0pt}%
\pgfpathmoveto{\pgfqpoint{0.462269in}{1.677321in}}%
\pgfpathlineto{\pgfqpoint{3.231834in}{1.677321in}}%
\pgfusepath{stroke}%
\end{pgfscope}%
\begin{pgfscope}%
\pgfsetbuttcap%
\pgfsetroundjoin%
\definecolor{currentfill}{rgb}{0.000000,0.000000,0.000000}%
\pgfsetfillcolor{currentfill}%
\pgfsetlinewidth{0.803000pt}%
\definecolor{currentstroke}{rgb}{0.000000,0.000000,0.000000}%
\pgfsetstrokecolor{currentstroke}%
\pgfsetdash{}{0pt}%
\pgfsys@defobject{currentmarker}{\pgfqpoint{-0.048611in}{0.000000in}}{\pgfqpoint{-0.000000in}{0.000000in}}{%
\pgfpathmoveto{\pgfqpoint{-0.000000in}{0.000000in}}%
\pgfpathlineto{\pgfqpoint{-0.048611in}{0.000000in}}%
\pgfusepath{stroke,fill}%
}%
\begin{pgfscope}%
\pgfsys@transformshift{0.462269in}{1.677321in}%
\pgfsys@useobject{currentmarker}{}%
\end{pgfscope}%
\end{pgfscope}%
\begin{pgfscope}%
\definecolor{textcolor}{rgb}{0.000000,0.000000,0.000000}%
\pgfsetstrokecolor{textcolor}%
\pgfsetfillcolor{textcolor}%
\pgftext[x=0.254321in, y=1.643563in, left, base]{\color{textcolor}{\rmfamily\fontsize{7.000000}{8.400000}\selectfont\catcode`\^=\active\def^{\ifmmode\sp\else\^{}\fi}\catcode`\%=\active\def%{\%}20}}%
\end{pgfscope}%
\begin{pgfscope}%
\pgfpathrectangle{\pgfqpoint{0.462269in}{0.807389in}}{\pgfqpoint{2.769565in}{2.016000in}}%
\pgfusepath{clip}%
\pgfsetrectcap%
\pgfsetroundjoin%
\pgfsetlinewidth{0.803000pt}%
\definecolor{currentstroke}{rgb}{0.690196,0.690196,0.690196}%
\pgfsetstrokecolor{currentstroke}%
\pgfsetstrokeopacity{0.300000}%
\pgfsetdash{}{0pt}%
\pgfpathmoveto{\pgfqpoint{0.462269in}{2.087137in}}%
\pgfpathlineto{\pgfqpoint{3.231834in}{2.087137in}}%
\pgfusepath{stroke}%
\end{pgfscope}%
\begin{pgfscope}%
\pgfsetbuttcap%
\pgfsetroundjoin%
\definecolor{currentfill}{rgb}{0.000000,0.000000,0.000000}%
\pgfsetfillcolor{currentfill}%
\pgfsetlinewidth{0.803000pt}%
\definecolor{currentstroke}{rgb}{0.000000,0.000000,0.000000}%
\pgfsetstrokecolor{currentstroke}%
\pgfsetdash{}{0pt}%
\pgfsys@defobject{currentmarker}{\pgfqpoint{-0.048611in}{0.000000in}}{\pgfqpoint{-0.000000in}{0.000000in}}{%
\pgfpathmoveto{\pgfqpoint{-0.000000in}{0.000000in}}%
\pgfpathlineto{\pgfqpoint{-0.048611in}{0.000000in}}%
\pgfusepath{stroke,fill}%
}%
\begin{pgfscope}%
\pgfsys@transformshift{0.462269in}{2.087137in}%
\pgfsys@useobject{currentmarker}{}%
\end{pgfscope}%
\end{pgfscope}%
\begin{pgfscope}%
\definecolor{textcolor}{rgb}{0.000000,0.000000,0.000000}%
\pgfsetstrokecolor{textcolor}%
\pgfsetfillcolor{textcolor}%
\pgftext[x=0.254321in, y=2.053379in, left, base]{\color{textcolor}{\rmfamily\fontsize{7.000000}{8.400000}\selectfont\catcode`\^=\active\def^{\ifmmode\sp\else\^{}\fi}\catcode`\%=\active\def%{\%}30}}%
\end{pgfscope}%
\begin{pgfscope}%
\pgfpathrectangle{\pgfqpoint{0.462269in}{0.807389in}}{\pgfqpoint{2.769565in}{2.016000in}}%
\pgfusepath{clip}%
\pgfsetrectcap%
\pgfsetroundjoin%
\pgfsetlinewidth{0.803000pt}%
\definecolor{currentstroke}{rgb}{0.690196,0.690196,0.690196}%
\pgfsetstrokecolor{currentstroke}%
\pgfsetstrokeopacity{0.300000}%
\pgfsetdash{}{0pt}%
\pgfpathmoveto{\pgfqpoint{0.462269in}{2.496952in}}%
\pgfpathlineto{\pgfqpoint{3.231834in}{2.496952in}}%
\pgfusepath{stroke}%
\end{pgfscope}%
\begin{pgfscope}%
\pgfsetbuttcap%
\pgfsetroundjoin%
\definecolor{currentfill}{rgb}{0.000000,0.000000,0.000000}%
\pgfsetfillcolor{currentfill}%
\pgfsetlinewidth{0.803000pt}%
\definecolor{currentstroke}{rgb}{0.000000,0.000000,0.000000}%
\pgfsetstrokecolor{currentstroke}%
\pgfsetdash{}{0pt}%
\pgfsys@defobject{currentmarker}{\pgfqpoint{-0.048611in}{0.000000in}}{\pgfqpoint{-0.000000in}{0.000000in}}{%
\pgfpathmoveto{\pgfqpoint{-0.000000in}{0.000000in}}%
\pgfpathlineto{\pgfqpoint{-0.048611in}{0.000000in}}%
\pgfusepath{stroke,fill}%
}%
\begin{pgfscope}%
\pgfsys@transformshift{0.462269in}{2.496952in}%
\pgfsys@useobject{currentmarker}{}%
\end{pgfscope}%
\end{pgfscope}%
\begin{pgfscope}%
\definecolor{textcolor}{rgb}{0.000000,0.000000,0.000000}%
\pgfsetstrokecolor{textcolor}%
\pgfsetfillcolor{textcolor}%
\pgftext[x=0.254321in, y=2.463194in, left, base]{\color{textcolor}{\rmfamily\fontsize{7.000000}{8.400000}\selectfont\catcode`\^=\active\def^{\ifmmode\sp\else\^{}\fi}\catcode`\%=\active\def%{\%}40}}%
\end{pgfscope}%
\begin{pgfscope}%
\definecolor{textcolor}{rgb}{0.000000,0.000000,0.000000}%
\pgfsetstrokecolor{textcolor}%
\pgfsetfillcolor{textcolor}%
\pgftext[x=0.198766in,y=1.815389in,,bottom,rotate=90.000000]{\color{textcolor}{\rmfamily\fontsize{8.000000}{9.600000}\selectfont\catcode`\^=\active\def^{\ifmmode\sp\else\^{}\fi}\catcode`\%=\active\def%{\%}Geometric ESS}}%
\end{pgfscope}%
\begin{pgfscope}%
\pgfpathrectangle{\pgfqpoint{0.462269in}{0.807389in}}{\pgfqpoint{2.769565in}{2.016000in}}%
\pgfusepath{clip}%
\pgfsetrectcap%
\pgfsetroundjoin%
\pgfsetlinewidth{1.806750pt}%
\definecolor{currentstroke}{rgb}{0.121569,0.466667,0.705882}%
\pgfsetstrokecolor{currentstroke}%
\pgfsetdash{}{0pt}%
\pgfpathmoveto{\pgfqpoint{0.588158in}{1.102343in}}%
\pgfpathlineto{\pgfqpoint{0.634784in}{1.081381in}}%
\pgfpathlineto{\pgfqpoint{0.681410in}{1.062531in}}%
\pgfpathlineto{\pgfqpoint{0.728036in}{1.045587in}}%
\pgfpathlineto{\pgfqpoint{0.774661in}{1.030363in}}%
\pgfpathlineto{\pgfqpoint{0.821287in}{1.016690in}}%
\pgfpathlineto{\pgfqpoint{0.867913in}{1.004438in}}%
\pgfpathlineto{\pgfqpoint{0.914538in}{0.993468in}}%
\pgfpathlineto{\pgfqpoint{0.961164in}{0.983624in}}%
\pgfpathlineto{\pgfqpoint{1.007790in}{0.974792in}}%
\pgfpathlineto{\pgfqpoint{1.054415in}{0.966870in}}%
\pgfpathlineto{\pgfqpoint{1.101041in}{0.959707in}}%
\pgfpathlineto{\pgfqpoint{1.147667in}{0.953311in}}%
\pgfpathlineto{\pgfqpoint{1.194292in}{0.947623in}}%
\pgfpathlineto{\pgfqpoint{1.240918in}{0.942519in}}%
\pgfpathlineto{\pgfqpoint{1.287544in}{0.937950in}}%
\pgfpathlineto{\pgfqpoint{1.334169in}{0.933864in}}%
\pgfpathlineto{\pgfqpoint{1.380795in}{0.930199in}}%
\pgfpathlineto{\pgfqpoint{1.427421in}{0.926911in}}%
\pgfpathlineto{\pgfqpoint{1.474046in}{0.923961in}}%
\pgfpathlineto{\pgfqpoint{1.520672in}{0.921316in}}%
\pgfpathlineto{\pgfqpoint{1.567298in}{0.918943in}}%
\pgfpathlineto{\pgfqpoint{1.613923in}{0.916786in}}%
\pgfpathlineto{\pgfqpoint{1.660549in}{0.914858in}}%
\pgfpathlineto{\pgfqpoint{1.707175in}{0.913138in}}%
\pgfpathlineto{\pgfqpoint{1.753800in}{0.911606in}}%
\pgfpathlineto{\pgfqpoint{1.800426in}{0.910235in}}%
\pgfpathlineto{\pgfqpoint{1.847052in}{0.909008in}}%
\pgfpathlineto{\pgfqpoint{1.893677in}{0.907909in}}%
\pgfpathlineto{\pgfqpoint{1.940303in}{0.906925in}}%
\pgfpathlineto{\pgfqpoint{1.986929in}{0.906044in}}%
\pgfpathlineto{\pgfqpoint{2.033554in}{0.905256in}}%
\pgfpathlineto{\pgfqpoint{2.080180in}{0.904550in}}%
\pgfpathlineto{\pgfqpoint{2.126806in}{0.903919in}}%
\pgfpathlineto{\pgfqpoint{2.173432in}{0.903353in}}%
\pgfpathlineto{\pgfqpoint{2.220057in}{0.902848in}}%
\pgfpathlineto{\pgfqpoint{2.266683in}{0.902395in}}%
\pgfpathlineto{\pgfqpoint{2.313309in}{0.901989in}}%
\pgfpathlineto{\pgfqpoint{2.359934in}{0.901627in}}%
\pgfpathlineto{\pgfqpoint{2.406560in}{0.901302in}}%
\pgfpathlineto{\pgfqpoint{2.453186in}{0.901011in}}%
\pgfpathlineto{\pgfqpoint{2.499811in}{0.900751in}}%
\pgfpathlineto{\pgfqpoint{2.546437in}{0.900518in}}%
\pgfpathlineto{\pgfqpoint{2.593063in}{0.900310in}}%
\pgfpathlineto{\pgfqpoint{2.639688in}{0.900123in}}%
\pgfpathlineto{\pgfqpoint{2.686314in}{0.899956in}}%
\pgfpathlineto{\pgfqpoint{2.732940in}{0.899807in}}%
\pgfpathlineto{\pgfqpoint{2.779565in}{0.899673in}}%
\pgfpathlineto{\pgfqpoint{2.826191in}{0.899553in}}%
\pgfpathlineto{\pgfqpoint{2.872817in}{0.899445in}}%
\pgfpathlineto{\pgfqpoint{2.919442in}{0.899349in}}%
\pgfpathlineto{\pgfqpoint{2.966068in}{0.899262in}}%
\pgfpathlineto{\pgfqpoint{3.012694in}{0.899184in}}%
\pgfpathlineto{\pgfqpoint{3.059319in}{0.899114in}}%
\pgfpathlineto{\pgfqpoint{3.105945in}{0.899051in}}%
\pgfusepath{stroke}%
\end{pgfscope}%
\begin{pgfscope}%
\pgfpathrectangle{\pgfqpoint{0.462269in}{0.807389in}}{\pgfqpoint{2.769565in}{2.016000in}}%
\pgfusepath{clip}%
\pgfsetrectcap%
\pgfsetroundjoin%
\pgfsetlinewidth{1.806750pt}%
\definecolor{currentstroke}{rgb}{1.000000,0.498039,0.054902}%
\pgfsetstrokecolor{currentstroke}%
\pgfsetdash{}{0pt}%
\pgfpathmoveto{\pgfqpoint{0.588158in}{1.347942in}}%
\pgfpathlineto{\pgfqpoint{0.634784in}{1.307172in}}%
\pgfpathlineto{\pgfqpoint{0.681410in}{1.270362in}}%
\pgfpathlineto{\pgfqpoint{0.728036in}{1.236995in}}%
\pgfpathlineto{\pgfqpoint{0.774661in}{1.206790in}}%
\pgfpathlineto{\pgfqpoint{0.821287in}{1.179485in}}%
\pgfpathlineto{\pgfqpoint{0.867913in}{1.154831in}}%
\pgfpathlineto{\pgfqpoint{0.914538in}{1.132870in}}%
\pgfpathlineto{\pgfqpoint{0.961164in}{1.113099in}}%
\pgfpathlineto{\pgfqpoint{1.007790in}{1.095280in}}%
\pgfpathlineto{\pgfqpoint{1.054415in}{1.079234in}}%
\pgfpathlineto{\pgfqpoint{1.101041in}{1.064796in}}%
\pgfpathlineto{\pgfqpoint{1.147667in}{1.051814in}}%
\pgfpathlineto{\pgfqpoint{1.194292in}{1.040150in}}%
\pgfpathlineto{\pgfqpoint{1.240918in}{1.029674in}}%
\pgfpathlineto{\pgfqpoint{1.287544in}{1.020272in}}%
\pgfpathlineto{\pgfqpoint{1.334169in}{1.011837in}}%
\pgfpathlineto{\pgfqpoint{1.380795in}{1.004272in}}%
\pgfpathlineto{\pgfqpoint{1.427421in}{0.997485in}}%
\pgfpathlineto{\pgfqpoint{1.474046in}{0.991372in}}%
\pgfpathlineto{\pgfqpoint{1.520672in}{0.985899in}}%
\pgfpathlineto{\pgfqpoint{1.567298in}{0.981000in}}%
\pgfpathlineto{\pgfqpoint{1.613923in}{0.976615in}}%
\pgfpathlineto{\pgfqpoint{1.660549in}{0.972724in}}%
\pgfpathlineto{\pgfqpoint{1.707175in}{0.969255in}}%
\pgfpathlineto{\pgfqpoint{1.753800in}{0.966144in}}%
\pgfpathlineto{\pgfqpoint{1.800426in}{0.963356in}}%
\pgfpathlineto{\pgfqpoint{1.847052in}{0.960856in}}%
\pgfpathlineto{\pgfqpoint{1.893677in}{0.958616in}}%
\pgfpathlineto{\pgfqpoint{1.940303in}{0.956608in}}%
\pgfpathlineto{\pgfqpoint{1.986929in}{0.954809in}}%
\pgfpathlineto{\pgfqpoint{2.033554in}{0.953197in}}%
\pgfpathlineto{\pgfqpoint{2.080180in}{0.951752in}}%
\pgfpathlineto{\pgfqpoint{2.126806in}{0.950458in}}%
\pgfpathlineto{\pgfqpoint{2.173432in}{0.949299in}}%
\pgfpathlineto{\pgfqpoint{2.220057in}{0.948261in}}%
\pgfpathlineto{\pgfqpoint{2.266683in}{0.947331in}}%
\pgfpathlineto{\pgfqpoint{2.313309in}{0.946498in}}%
\pgfpathlineto{\pgfqpoint{2.359934in}{0.945751in}}%
\pgfpathlineto{\pgfqpoint{2.406560in}{0.945083in}}%
\pgfpathlineto{\pgfqpoint{2.453186in}{0.944484in}}%
\pgfpathlineto{\pgfqpoint{2.499811in}{0.943946in}}%
\pgfpathlineto{\pgfqpoint{2.546437in}{0.943460in}}%
\pgfpathlineto{\pgfqpoint{2.593063in}{0.943012in}}%
\pgfpathlineto{\pgfqpoint{2.639688in}{0.942582in}}%
\pgfpathlineto{\pgfqpoint{2.686314in}{0.942133in}}%
\pgfpathlineto{\pgfqpoint{2.732940in}{0.941609in}}%
\pgfpathlineto{\pgfqpoint{2.779565in}{0.940927in}}%
\pgfpathlineto{\pgfqpoint{2.826191in}{0.939978in}}%
\pgfpathlineto{\pgfqpoint{2.872817in}{0.938643in}}%
\pgfpathlineto{\pgfqpoint{2.919442in}{0.936815in}}%
\pgfpathlineto{\pgfqpoint{2.966068in}{0.934428in}}%
\pgfpathlineto{\pgfqpoint{3.012694in}{0.931489in}}%
\pgfpathlineto{\pgfqpoint{3.059319in}{0.928087in}}%
\pgfpathlineto{\pgfqpoint{3.105945in}{0.924375in}}%
\pgfusepath{stroke}%
\end{pgfscope}%
\begin{pgfscope}%
\pgfpathrectangle{\pgfqpoint{0.462269in}{0.807389in}}{\pgfqpoint{2.769565in}{2.016000in}}%
\pgfusepath{clip}%
\pgfsetrectcap%
\pgfsetroundjoin%
\pgfsetlinewidth{1.806750pt}%
\definecolor{currentstroke}{rgb}{0.172549,0.627451,0.172549}%
\pgfsetstrokecolor{currentstroke}%
\pgfsetdash{}{0pt}%
\pgfpathmoveto{\pgfqpoint{0.588158in}{2.664650in}}%
\pgfpathlineto{\pgfqpoint{0.634784in}{2.555612in}}%
\pgfpathlineto{\pgfqpoint{0.681410in}{2.453033in}}%
\pgfpathlineto{\pgfqpoint{0.728036in}{2.356708in}}%
\pgfpathlineto{\pgfqpoint{0.774661in}{2.266386in}}%
\pgfpathlineto{\pgfqpoint{0.821287in}{2.181927in}}%
\pgfpathlineto{\pgfqpoint{0.867913in}{2.102830in}}%
\pgfpathlineto{\pgfqpoint{0.914538in}{2.028494in}}%
\pgfpathlineto{\pgfqpoint{0.961164in}{1.958772in}}%
\pgfpathlineto{\pgfqpoint{1.007790in}{1.893149in}}%
\pgfpathlineto{\pgfqpoint{1.054415in}{1.831425in}}%
\pgfpathlineto{\pgfqpoint{1.101041in}{1.773801in}}%
\pgfpathlineto{\pgfqpoint{1.147667in}{1.720942in}}%
\pgfpathlineto{\pgfqpoint{1.194292in}{1.671240in}}%
\pgfpathlineto{\pgfqpoint{1.240918in}{1.624739in}}%
\pgfpathlineto{\pgfqpoint{1.287544in}{1.581185in}}%
\pgfpathlineto{\pgfqpoint{1.334169in}{1.540347in}}%
\pgfpathlineto{\pgfqpoint{1.380795in}{1.501916in}}%
\pgfpathlineto{\pgfqpoint{1.427421in}{1.465755in}}%
\pgfpathlineto{\pgfqpoint{1.474046in}{1.431924in}}%
\pgfpathlineto{\pgfqpoint{1.520672in}{1.400086in}}%
\pgfpathlineto{\pgfqpoint{1.567298in}{1.369753in}}%
\pgfpathlineto{\pgfqpoint{1.613923in}{1.341107in}}%
\pgfpathlineto{\pgfqpoint{1.660549in}{1.314180in}}%
\pgfpathlineto{\pgfqpoint{1.707175in}{1.288851in}}%
\pgfpathlineto{\pgfqpoint{1.753800in}{1.265007in}}%
\pgfpathlineto{\pgfqpoint{1.800426in}{1.242909in}}%
\pgfpathlineto{\pgfqpoint{1.847052in}{1.222032in}}%
\pgfpathlineto{\pgfqpoint{1.893677in}{1.202273in}}%
\pgfpathlineto{\pgfqpoint{1.940303in}{1.183486in}}%
\pgfpathlineto{\pgfqpoint{1.986929in}{1.165564in}}%
\pgfpathlineto{\pgfqpoint{2.033554in}{1.148841in}}%
\pgfpathlineto{\pgfqpoint{2.080180in}{1.133094in}}%
\pgfpathlineto{\pgfqpoint{2.126806in}{1.118083in}}%
\pgfpathlineto{\pgfqpoint{2.173432in}{1.103744in}}%
\pgfpathlineto{\pgfqpoint{2.220057in}{1.090076in}}%
\pgfpathlineto{\pgfqpoint{2.266683in}{1.077137in}}%
\pgfpathlineto{\pgfqpoint{2.313309in}{1.064881in}}%
\pgfpathlineto{\pgfqpoint{2.359934in}{1.053263in}}%
\pgfpathlineto{\pgfqpoint{2.406560in}{1.042194in}}%
\pgfpathlineto{\pgfqpoint{2.453186in}{1.031680in}}%
\pgfpathlineto{\pgfqpoint{2.499811in}{1.021710in}}%
\pgfpathlineto{\pgfqpoint{2.546437in}{1.012258in}}%
\pgfpathlineto{\pgfqpoint{2.593063in}{1.003253in}}%
\pgfpathlineto{\pgfqpoint{2.639688in}{0.994661in}}%
\pgfpathlineto{\pgfqpoint{2.686314in}{0.986444in}}%
\pgfpathlineto{\pgfqpoint{2.732940in}{0.978567in}}%
\pgfpathlineto{\pgfqpoint{2.779565in}{0.971001in}}%
\pgfpathlineto{\pgfqpoint{2.826191in}{0.963721in}}%
\pgfpathlineto{\pgfqpoint{2.872817in}{0.956719in}}%
\pgfpathlineto{\pgfqpoint{2.919442in}{0.949999in}}%
\pgfpathlineto{\pgfqpoint{2.966068in}{0.943581in}}%
\pgfpathlineto{\pgfqpoint{3.012694in}{0.937499in}}%
\pgfpathlineto{\pgfqpoint{3.059319in}{0.931802in}}%
\pgfpathlineto{\pgfqpoint{3.105945in}{0.926520in}}%
\pgfusepath{stroke}%
\end{pgfscope}%
\begin{pgfscope}%
\pgfpathrectangle{\pgfqpoint{0.462269in}{0.807389in}}{\pgfqpoint{2.769565in}{2.016000in}}%
\pgfusepath{clip}%
\pgfsetrectcap%
\pgfsetroundjoin%
\pgfsetlinewidth{1.806750pt}%
\definecolor{currentstroke}{rgb}{0.839216,0.152941,0.156863}%
\pgfsetstrokecolor{currentstroke}%
\pgfsetdash{}{0pt}%
\pgfpathmoveto{\pgfqpoint{0.588158in}{1.756980in}}%
\pgfpathlineto{\pgfqpoint{0.634784in}{1.686458in}}%
\pgfpathlineto{\pgfqpoint{0.681410in}{1.622096in}}%
\pgfpathlineto{\pgfqpoint{0.728036in}{1.563488in}}%
\pgfpathlineto{\pgfqpoint{0.774661in}{1.510223in}}%
\pgfpathlineto{\pgfqpoint{0.821287in}{1.461900in}}%
\pgfpathlineto{\pgfqpoint{0.867913in}{1.418323in}}%
\pgfpathlineto{\pgfqpoint{0.914538in}{1.379242in}}%
\pgfpathlineto{\pgfqpoint{0.961164in}{1.343495in}}%
\pgfpathlineto{\pgfqpoint{1.007790in}{1.311162in}}%
\pgfpathlineto{\pgfqpoint{1.054415in}{1.281956in}}%
\pgfpathlineto{\pgfqpoint{1.101041in}{1.255592in}}%
\pgfpathlineto{\pgfqpoint{1.147667in}{1.231810in}}%
\pgfpathlineto{\pgfqpoint{1.194292in}{1.210370in}}%
\pgfpathlineto{\pgfqpoint{1.240918in}{1.191053in}}%
\pgfpathlineto{\pgfqpoint{1.287544in}{1.173659in}}%
\pgfpathlineto{\pgfqpoint{1.334169in}{1.158038in}}%
\pgfpathlineto{\pgfqpoint{1.380795in}{1.144151in}}%
\pgfpathlineto{\pgfqpoint{1.427421in}{1.131642in}}%
\pgfpathlineto{\pgfqpoint{1.474046in}{1.120379in}}%
\pgfpathlineto{\pgfqpoint{1.520672in}{1.110242in}}%
\pgfpathlineto{\pgfqpoint{1.567298in}{1.101121in}}%
\pgfpathlineto{\pgfqpoint{1.613923in}{1.092916in}}%
\pgfpathlineto{\pgfqpoint{1.660549in}{1.085539in}}%
\pgfpathlineto{\pgfqpoint{1.707175in}{1.078907in}}%
\pgfpathlineto{\pgfqpoint{1.753800in}{1.072948in}}%
\pgfpathlineto{\pgfqpoint{1.800426in}{1.067594in}}%
\pgfpathlineto{\pgfqpoint{1.847052in}{1.062770in}}%
\pgfpathlineto{\pgfqpoint{1.893677in}{1.058429in}}%
\pgfpathlineto{\pgfqpoint{1.940303in}{1.054537in}}%
\pgfpathlineto{\pgfqpoint{1.986929in}{1.051047in}}%
\pgfpathlineto{\pgfqpoint{2.033554in}{1.047918in}}%
\pgfpathlineto{\pgfqpoint{2.080180in}{1.045111in}}%
\pgfpathlineto{\pgfqpoint{2.126806in}{1.042588in}}%
\pgfpathlineto{\pgfqpoint{2.173432in}{1.040310in}}%
\pgfpathlineto{\pgfqpoint{2.220057in}{1.038230in}}%
\pgfpathlineto{\pgfqpoint{2.266683in}{1.036283in}}%
\pgfpathlineto{\pgfqpoint{2.313309in}{1.034382in}}%
\pgfpathlineto{\pgfqpoint{2.359934in}{1.032404in}}%
\pgfpathlineto{\pgfqpoint{2.406560in}{1.030174in}}%
\pgfpathlineto{\pgfqpoint{2.453186in}{1.027483in}}%
\pgfpathlineto{\pgfqpoint{2.499811in}{1.024092in}}%
\pgfpathlineto{\pgfqpoint{2.546437in}{1.019777in}}%
\pgfpathlineto{\pgfqpoint{2.593063in}{1.014368in}}%
\pgfpathlineto{\pgfqpoint{2.639688in}{1.007806in}}%
\pgfpathlineto{\pgfqpoint{2.686314in}{1.000174in}}%
\pgfpathlineto{\pgfqpoint{2.732940in}{0.991682in}}%
\pgfpathlineto{\pgfqpoint{2.779565in}{0.982655in}}%
\pgfpathlineto{\pgfqpoint{2.826191in}{0.973441in}}%
\pgfpathlineto{\pgfqpoint{2.872817in}{0.964372in}}%
\pgfpathlineto{\pgfqpoint{2.919442in}{0.955709in}}%
\pgfpathlineto{\pgfqpoint{2.966068in}{0.947629in}}%
\pgfpathlineto{\pgfqpoint{3.012694in}{0.940234in}}%
\pgfpathlineto{\pgfqpoint{3.059319in}{0.933560in}}%
\pgfpathlineto{\pgfqpoint{3.105945in}{0.927602in}}%
\pgfusepath{stroke}%
\end{pgfscope}%
\begin{pgfscope}%
\pgfpathrectangle{\pgfqpoint{0.462269in}{0.807389in}}{\pgfqpoint{2.769565in}{2.016000in}}%
\pgfusepath{clip}%
\pgfsetrectcap%
\pgfsetroundjoin%
\pgfsetlinewidth{1.806750pt}%
\definecolor{currentstroke}{rgb}{0.580392,0.403922,0.741176}%
\pgfsetstrokecolor{currentstroke}%
\pgfsetdash{}{0pt}%
\pgfpathmoveto{\pgfqpoint{0.588158in}{1.677321in}}%
\pgfpathlineto{\pgfqpoint{0.634784in}{1.677321in}}%
\pgfpathlineto{\pgfqpoint{0.681410in}{1.677321in}}%
\pgfpathlineto{\pgfqpoint{0.728036in}{1.677321in}}%
\pgfpathlineto{\pgfqpoint{0.774661in}{1.677321in}}%
\pgfpathlineto{\pgfqpoint{0.821287in}{1.677321in}}%
\pgfpathlineto{\pgfqpoint{0.867913in}{1.677321in}}%
\pgfpathlineto{\pgfqpoint{0.914538in}{1.677321in}}%
\pgfpathlineto{\pgfqpoint{0.961164in}{1.677321in}}%
\pgfpathlineto{\pgfqpoint{1.007790in}{1.677321in}}%
\pgfpathlineto{\pgfqpoint{1.054415in}{1.677321in}}%
\pgfpathlineto{\pgfqpoint{1.101041in}{1.677321in}}%
\pgfpathlineto{\pgfqpoint{1.147667in}{1.677321in}}%
\pgfpathlineto{\pgfqpoint{1.194292in}{1.677321in}}%
\pgfpathlineto{\pgfqpoint{1.240918in}{1.677321in}}%
\pgfpathlineto{\pgfqpoint{1.287544in}{1.677319in}}%
\pgfpathlineto{\pgfqpoint{1.334169in}{1.677313in}}%
\pgfpathlineto{\pgfqpoint{1.380795in}{1.677293in}}%
\pgfpathlineto{\pgfqpoint{1.427421in}{1.677233in}}%
\pgfpathlineto{\pgfqpoint{1.474046in}{1.677072in}}%
\pgfpathlineto{\pgfqpoint{1.520672in}{1.676685in}}%
\pgfpathlineto{\pgfqpoint{1.567298in}{1.675839in}}%
\pgfpathlineto{\pgfqpoint{1.613923in}{1.674140in}}%
\pgfpathlineto{\pgfqpoint{1.660549in}{1.670990in}}%
\pgfpathlineto{\pgfqpoint{1.707175in}{1.665555in}}%
\pgfpathlineto{\pgfqpoint{1.753800in}{1.656793in}}%
\pgfpathlineto{\pgfqpoint{1.800426in}{1.643552in}}%
\pgfpathlineto{\pgfqpoint{1.847052in}{1.624756in}}%
\pgfpathlineto{\pgfqpoint{1.893677in}{1.599651in}}%
\pgfpathlineto{\pgfqpoint{1.940303in}{1.568046in}}%
\pgfpathlineto{\pgfqpoint{1.986929in}{1.530459in}}%
\pgfpathlineto{\pgfqpoint{2.033554in}{1.488081in}}%
\pgfpathlineto{\pgfqpoint{2.080180in}{1.442569in}}%
\pgfpathlineto{\pgfqpoint{2.126806in}{1.395736in}}%
\pgfpathlineto{\pgfqpoint{2.173432in}{1.349253in}}%
\pgfpathlineto{\pgfqpoint{2.220057in}{1.304448in}}%
\pgfpathlineto{\pgfqpoint{2.266683in}{1.262231in}}%
\pgfpathlineto{\pgfqpoint{2.313309in}{1.223120in}}%
\pgfpathlineto{\pgfqpoint{2.359934in}{1.187323in}}%
\pgfpathlineto{\pgfqpoint{2.406560in}{1.154827in}}%
\pgfpathlineto{\pgfqpoint{2.453186in}{1.125491in}}%
\pgfpathlineto{\pgfqpoint{2.499811in}{1.099100in}}%
\pgfpathlineto{\pgfqpoint{2.546437in}{1.075411in}}%
\pgfpathlineto{\pgfqpoint{2.593063in}{1.054179in}}%
\pgfpathlineto{\pgfqpoint{2.639688in}{1.035167in}}%
\pgfpathlineto{\pgfqpoint{2.686314in}{1.018154in}}%
\pgfpathlineto{\pgfqpoint{2.732940in}{1.002940in}}%
\pgfpathlineto{\pgfqpoint{2.779565in}{0.989342in}}%
\pgfpathlineto{\pgfqpoint{2.826191in}{0.977196in}}%
\pgfpathlineto{\pgfqpoint{2.872817in}{0.966356in}}%
\pgfpathlineto{\pgfqpoint{2.919442in}{0.956689in}}%
\pgfpathlineto{\pgfqpoint{2.966068in}{0.948080in}}%
\pgfpathlineto{\pgfqpoint{3.012694in}{0.940425in}}%
\pgfpathlineto{\pgfqpoint{3.059319in}{0.933634in}}%
\pgfpathlineto{\pgfqpoint{3.105945in}{0.927628in}}%
\pgfusepath{stroke}%
\end{pgfscope}%
\begin{pgfscope}%
\pgfsetrectcap%
\pgfsetmiterjoin%
\pgfsetlinewidth{0.803000pt}%
\definecolor{currentstroke}{rgb}{0.000000,0.000000,0.000000}%
\pgfsetstrokecolor{currentstroke}%
\pgfsetdash{}{0pt}%
\pgfpathmoveto{\pgfqpoint{0.462269in}{0.807389in}}%
\pgfpathlineto{\pgfqpoint{0.462269in}{2.823389in}}%
\pgfusepath{stroke}%
\end{pgfscope}%
\begin{pgfscope}%
\pgfsetrectcap%
\pgfsetmiterjoin%
\pgfsetlinewidth{0.803000pt}%
\definecolor{currentstroke}{rgb}{0.000000,0.000000,0.000000}%
\pgfsetstrokecolor{currentstroke}%
\pgfsetdash{}{0pt}%
\pgfpathmoveto{\pgfqpoint{3.231834in}{0.807389in}}%
\pgfpathlineto{\pgfqpoint{3.231834in}{2.823389in}}%
\pgfusepath{stroke}%
\end{pgfscope}%
\begin{pgfscope}%
\pgfsetrectcap%
\pgfsetmiterjoin%
\pgfsetlinewidth{0.803000pt}%
\definecolor{currentstroke}{rgb}{0.000000,0.000000,0.000000}%
\pgfsetstrokecolor{currentstroke}%
\pgfsetdash{}{0pt}%
\pgfpathmoveto{\pgfqpoint{0.462269in}{0.807389in}}%
\pgfpathlineto{\pgfqpoint{3.231834in}{0.807389in}}%
\pgfusepath{stroke}%
\end{pgfscope}%
\begin{pgfscope}%
\pgfsetrectcap%
\pgfsetmiterjoin%
\pgfsetlinewidth{0.803000pt}%
\definecolor{currentstroke}{rgb}{0.000000,0.000000,0.000000}%
\pgfsetstrokecolor{currentstroke}%
\pgfsetdash{}{0pt}%
\pgfpathmoveto{\pgfqpoint{0.462269in}{2.823389in}}%
\pgfpathlineto{\pgfqpoint{3.231834in}{2.823389in}}%
\pgfusepath{stroke}%
\end{pgfscope}%
\begin{pgfscope}%
\definecolor{textcolor}{rgb}{0.000000,0.000000,0.000000}%
\pgfsetstrokecolor{textcolor}%
\pgfsetfillcolor{textcolor}%
\pgftext[x=1.847052in,y=2.906722in,,base]{\color{textcolor}{\rmfamily\fontsize{8.000000}{9.600000}\selectfont\catcode`\^=\active\def^{\ifmmode\sp\else\^{}\fi}\catcode`\%=\active\def%{\%}(a) Replicated gESS Profiles}}%
\end{pgfscope}%
\begin{pgfscope}%
\pgfsetbuttcap%
\pgfsetmiterjoin%
\definecolor{currentfill}{rgb}{1.000000,1.000000,1.000000}%
\pgfsetfillcolor{currentfill}%
\pgfsetlinewidth{0.000000pt}%
\definecolor{currentstroke}{rgb}{0.000000,0.000000,0.000000}%
\pgfsetstrokecolor{currentstroke}%
\pgfsetstrokeopacity{0.000000}%
\pgfsetdash{}{0pt}%
\pgfpathmoveto{\pgfqpoint{4.062704in}{0.807389in}}%
\pgfpathlineto{\pgfqpoint{6.832269in}{0.807389in}}%
\pgfpathlineto{\pgfqpoint{6.832269in}{2.823389in}}%
\pgfpathlineto{\pgfqpoint{4.062704in}{2.823389in}}%
\pgfpathlineto{\pgfqpoint{4.062704in}{0.807389in}}%
\pgfpathclose%
\pgfusepath{fill}%
\end{pgfscope}%
\begin{pgfscope}%
\pgfpathrectangle{\pgfqpoint{4.062704in}{0.807389in}}{\pgfqpoint{2.769565in}{2.016000in}}%
\pgfusepath{clip}%
\pgfsetbuttcap%
\pgfsetroundjoin%
\definecolor{currentfill}{rgb}{0.121569,0.466667,0.705882}%
\pgfsetfillcolor{currentfill}%
\pgfsetfillopacity{0.150000}%
\pgfsetlinewidth{0.000000pt}%
\definecolor{currentstroke}{rgb}{0.121569,0.466667,0.705882}%
\pgfsetstrokecolor{currentstroke}%
\pgfsetstrokeopacity{0.150000}%
\pgfsetdash{}{0pt}%
\pgfsys@defobject{currentmarker}{\pgfqpoint{4.188593in}{0.899025in}}{\pgfqpoint{6.706380in}{1.975438in}}{%
\pgfpathmoveto{\pgfqpoint{4.188593in}{0.901000in}}%
\pgfpathlineto{\pgfqpoint{4.188593in}{0.899025in}}%
\pgfpathlineto{\pgfqpoint{4.235219in}{0.899498in}}%
\pgfpathlineto{\pgfqpoint{4.281845in}{0.899997in}}%
\pgfpathlineto{\pgfqpoint{4.328470in}{0.900546in}}%
\pgfpathlineto{\pgfqpoint{4.375096in}{0.901151in}}%
\pgfpathlineto{\pgfqpoint{4.421722in}{0.901800in}}%
\pgfpathlineto{\pgfqpoint{4.468347in}{0.902521in}}%
\pgfpathlineto{\pgfqpoint{4.514973in}{0.903323in}}%
\pgfpathlineto{\pgfqpoint{4.561599in}{0.904217in}}%
\pgfpathlineto{\pgfqpoint{4.608224in}{0.905213in}}%
\pgfpathlineto{\pgfqpoint{4.654850in}{0.906323in}}%
\pgfpathlineto{\pgfqpoint{4.701476in}{0.907567in}}%
\pgfpathlineto{\pgfqpoint{4.748101in}{0.908984in}}%
\pgfpathlineto{\pgfqpoint{4.794727in}{0.910563in}}%
\pgfpathlineto{\pgfqpoint{4.841353in}{0.912314in}}%
\pgfpathlineto{\pgfqpoint{4.887978in}{0.914248in}}%
\pgfpathlineto{\pgfqpoint{4.934604in}{0.916408in}}%
\pgfpathlineto{\pgfqpoint{4.981230in}{0.918820in}}%
\pgfpathlineto{\pgfqpoint{5.027855in}{0.921529in}}%
\pgfpathlineto{\pgfqpoint{5.074481in}{0.924546in}}%
\pgfpathlineto{\pgfqpoint{5.121107in}{0.927901in}}%
\pgfpathlineto{\pgfqpoint{5.167732in}{0.931645in}}%
\pgfpathlineto{\pgfqpoint{5.214358in}{0.935820in}}%
\pgfpathlineto{\pgfqpoint{5.260984in}{0.940477in}}%
\pgfpathlineto{\pgfqpoint{5.307610in}{0.945669in}}%
\pgfpathlineto{\pgfqpoint{5.354235in}{0.951458in}}%
\pgfpathlineto{\pgfqpoint{5.400861in}{0.957911in}}%
\pgfpathlineto{\pgfqpoint{5.447487in}{0.965095in}}%
\pgfpathlineto{\pgfqpoint{5.494112in}{0.973097in}}%
\pgfpathlineto{\pgfqpoint{5.540738in}{0.982011in}}%
\pgfpathlineto{\pgfqpoint{5.587364in}{0.991937in}}%
\pgfpathlineto{\pgfqpoint{5.633989in}{1.002984in}}%
\pgfpathlineto{\pgfqpoint{5.680615in}{1.015273in}}%
\pgfpathlineto{\pgfqpoint{5.727241in}{1.028936in}}%
\pgfpathlineto{\pgfqpoint{5.773866in}{1.044118in}}%
\pgfpathlineto{\pgfqpoint{5.820492in}{1.060975in}}%
\pgfpathlineto{\pgfqpoint{5.867118in}{1.079679in}}%
\pgfpathlineto{\pgfqpoint{5.913743in}{1.100414in}}%
\pgfpathlineto{\pgfqpoint{5.960369in}{1.123379in}}%
\pgfpathlineto{\pgfqpoint{6.006995in}{1.148784in}}%
\pgfpathlineto{\pgfqpoint{6.053620in}{1.176854in}}%
\pgfpathlineto{\pgfqpoint{6.100246in}{1.207829in}}%
\pgfpathlineto{\pgfqpoint{6.146872in}{1.241954in}}%
\pgfpathlineto{\pgfqpoint{6.193497in}{1.279485in}}%
\pgfpathlineto{\pgfqpoint{6.240123in}{1.320679in}}%
\pgfpathlineto{\pgfqpoint{6.286749in}{1.365791in}}%
\pgfpathlineto{\pgfqpoint{6.333374in}{1.415065in}}%
\pgfpathlineto{\pgfqpoint{6.380000in}{1.468727in}}%
\pgfpathlineto{\pgfqpoint{6.426626in}{1.526969in}}%
\pgfpathlineto{\pgfqpoint{6.473251in}{1.589934in}}%
\pgfpathlineto{\pgfqpoint{6.519877in}{1.657696in}}%
\pgfpathlineto{\pgfqpoint{6.566503in}{1.730232in}}%
\pgfpathlineto{\pgfqpoint{6.613128in}{1.807389in}}%
\pgfpathlineto{\pgfqpoint{6.659754in}{1.888841in}}%
\pgfpathlineto{\pgfqpoint{6.706380in}{1.974026in}}%
\pgfpathlineto{\pgfqpoint{6.706380in}{1.975438in}}%
\pgfpathlineto{\pgfqpoint{6.706380in}{1.975438in}}%
\pgfpathlineto{\pgfqpoint{6.659754in}{1.890357in}}%
\pgfpathlineto{\pgfqpoint{6.613128in}{1.809000in}}%
\pgfpathlineto{\pgfqpoint{6.566503in}{1.731929in}}%
\pgfpathlineto{\pgfqpoint{6.519877in}{1.659470in}}%
\pgfpathlineto{\pgfqpoint{6.473251in}{1.591779in}}%
\pgfpathlineto{\pgfqpoint{6.426626in}{1.528878in}}%
\pgfpathlineto{\pgfqpoint{6.380000in}{1.470694in}}%
\pgfpathlineto{\pgfqpoint{6.333374in}{1.417085in}}%
\pgfpathlineto{\pgfqpoint{6.286749in}{1.367857in}}%
\pgfpathlineto{\pgfqpoint{6.240123in}{1.322788in}}%
\pgfpathlineto{\pgfqpoint{6.193497in}{1.281631in}}%
\pgfpathlineto{\pgfqpoint{6.146872in}{1.244133in}}%
\pgfpathlineto{\pgfqpoint{6.100246in}{1.210037in}}%
\pgfpathlineto{\pgfqpoint{6.053620in}{1.179088in}}%
\pgfpathlineto{\pgfqpoint{6.006995in}{1.151039in}}%
\pgfpathlineto{\pgfqpoint{5.960369in}{1.125654in}}%
\pgfpathlineto{\pgfqpoint{5.913743in}{1.102708in}}%
\pgfpathlineto{\pgfqpoint{5.867118in}{1.081988in}}%
\pgfpathlineto{\pgfqpoint{5.820492in}{1.063296in}}%
\pgfpathlineto{\pgfqpoint{5.773866in}{1.046449in}}%
\pgfpathlineto{\pgfqpoint{5.727241in}{1.031275in}}%
\pgfpathlineto{\pgfqpoint{5.680615in}{1.017617in}}%
\pgfpathlineto{\pgfqpoint{5.633989in}{1.005332in}}%
\pgfpathlineto{\pgfqpoint{5.587364in}{0.994287in}}%
\pgfpathlineto{\pgfqpoint{5.540738in}{0.984361in}}%
\pgfpathlineto{\pgfqpoint{5.494112in}{0.975445in}}%
\pgfpathlineto{\pgfqpoint{5.447487in}{0.967438in}}%
\pgfpathlineto{\pgfqpoint{5.400861in}{0.960250in}}%
\pgfpathlineto{\pgfqpoint{5.354235in}{0.953799in}}%
\pgfpathlineto{\pgfqpoint{5.307610in}{0.948010in}}%
\pgfpathlineto{\pgfqpoint{5.260984in}{0.942817in}}%
\pgfpathlineto{\pgfqpoint{5.214358in}{0.938158in}}%
\pgfpathlineto{\pgfqpoint{5.167732in}{0.933978in}}%
\pgfpathlineto{\pgfqpoint{5.121107in}{0.930230in}}%
\pgfpathlineto{\pgfqpoint{5.074481in}{0.926867in}}%
\pgfpathlineto{\pgfqpoint{5.027855in}{0.923849in}}%
\pgfpathlineto{\pgfqpoint{4.981230in}{0.921142in}}%
\pgfpathlineto{\pgfqpoint{4.934604in}{0.918712in}}%
\pgfpathlineto{\pgfqpoint{4.887978in}{0.916531in}}%
\pgfpathlineto{\pgfqpoint{4.841353in}{0.914591in}}%
\pgfpathlineto{\pgfqpoint{4.794727in}{0.912851in}}%
\pgfpathlineto{\pgfqpoint{4.748101in}{0.911288in}}%
\pgfpathlineto{\pgfqpoint{4.701476in}{0.909885in}}%
\pgfpathlineto{\pgfqpoint{4.654850in}{0.908623in}}%
\pgfpathlineto{\pgfqpoint{4.608224in}{0.907491in}}%
\pgfpathlineto{\pgfqpoint{4.561599in}{0.906478in}}%
\pgfpathlineto{\pgfqpoint{4.514973in}{0.905566in}}%
\pgfpathlineto{\pgfqpoint{4.468347in}{0.904711in}}%
\pgfpathlineto{\pgfqpoint{4.421722in}{0.903935in}}%
\pgfpathlineto{\pgfqpoint{4.375096in}{0.903230in}}%
\pgfpathlineto{\pgfqpoint{4.328470in}{0.902586in}}%
\pgfpathlineto{\pgfqpoint{4.281845in}{0.901998in}}%
\pgfpathlineto{\pgfqpoint{4.235219in}{0.901471in}}%
\pgfpathlineto{\pgfqpoint{4.188593in}{0.901000in}}%
\pgfpathlineto{\pgfqpoint{4.188593in}{0.901000in}}%
\pgfpathclose%
\pgfusepath{fill}%
}%
\begin{pgfscope}%
\pgfsys@transformshift{0.000000in}{0.000000in}%
\pgfsys@useobject{currentmarker}{}%
\end{pgfscope}%
\end{pgfscope}%
\begin{pgfscope}%
\pgfpathrectangle{\pgfqpoint{4.062704in}{0.807389in}}{\pgfqpoint{2.769565in}{2.016000in}}%
\pgfusepath{clip}%
\pgfsetbuttcap%
\pgfsetroundjoin%
\definecolor{currentfill}{rgb}{1.000000,0.498039,0.054902}%
\pgfsetfillcolor{currentfill}%
\pgfsetfillopacity{0.150000}%
\pgfsetlinewidth{0.000000pt}%
\definecolor{currentstroke}{rgb}{1.000000,0.498039,0.054902}%
\pgfsetstrokecolor{currentstroke}%
\pgfsetstrokeopacity{0.150000}%
\pgfsetdash{}{0pt}%
\pgfsys@defobject{currentmarker}{\pgfqpoint{4.188593in}{0.919076in}}{\pgfqpoint{6.706380in}{2.645869in}}{%
\pgfpathmoveto{\pgfqpoint{4.188593in}{0.923420in}}%
\pgfpathlineto{\pgfqpoint{4.188593in}{0.919076in}}%
\pgfpathlineto{\pgfqpoint{4.235219in}{0.920072in}}%
\pgfpathlineto{\pgfqpoint{4.281845in}{0.921142in}}%
\pgfpathlineto{\pgfqpoint{4.328470in}{0.922332in}}%
\pgfpathlineto{\pgfqpoint{4.375096in}{0.923635in}}%
\pgfpathlineto{\pgfqpoint{4.421722in}{0.925061in}}%
\pgfpathlineto{\pgfqpoint{4.468347in}{0.926615in}}%
\pgfpathlineto{\pgfqpoint{4.514973in}{0.928401in}}%
\pgfpathlineto{\pgfqpoint{4.561599in}{0.930303in}}%
\pgfpathlineto{\pgfqpoint{4.608224in}{0.932431in}}%
\pgfpathlineto{\pgfqpoint{4.654850in}{0.934784in}}%
\pgfpathlineto{\pgfqpoint{4.701476in}{0.937387in}}%
\pgfpathlineto{\pgfqpoint{4.748101in}{0.940280in}}%
\pgfpathlineto{\pgfqpoint{4.794727in}{0.943491in}}%
\pgfpathlineto{\pgfqpoint{4.841353in}{0.947029in}}%
\pgfpathlineto{\pgfqpoint{4.887978in}{0.950981in}}%
\pgfpathlineto{\pgfqpoint{4.934604in}{0.955378in}}%
\pgfpathlineto{\pgfqpoint{4.981230in}{0.960271in}}%
\pgfpathlineto{\pgfqpoint{5.027855in}{0.965718in}}%
\pgfpathlineto{\pgfqpoint{5.074481in}{0.971785in}}%
\pgfpathlineto{\pgfqpoint{5.121107in}{0.978545in}}%
\pgfpathlineto{\pgfqpoint{5.167732in}{0.986075in}}%
\pgfpathlineto{\pgfqpoint{5.214358in}{0.994467in}}%
\pgfpathlineto{\pgfqpoint{5.260984in}{1.003817in}}%
\pgfpathlineto{\pgfqpoint{5.307610in}{1.014211in}}%
\pgfpathlineto{\pgfqpoint{5.354235in}{1.025796in}}%
\pgfpathlineto{\pgfqpoint{5.400861in}{1.038693in}}%
\pgfpathlineto{\pgfqpoint{5.447487in}{1.053049in}}%
\pgfpathlineto{\pgfqpoint{5.494112in}{1.069045in}}%
\pgfpathlineto{\pgfqpoint{5.540738in}{1.086865in}}%
\pgfpathlineto{\pgfqpoint{5.587364in}{1.106707in}}%
\pgfpathlineto{\pgfqpoint{5.633989in}{1.128791in}}%
\pgfpathlineto{\pgfqpoint{5.680615in}{1.153358in}}%
\pgfpathlineto{\pgfqpoint{5.727241in}{1.180711in}}%
\pgfpathlineto{\pgfqpoint{5.773866in}{1.211102in}}%
\pgfpathlineto{\pgfqpoint{5.820492in}{1.244839in}}%
\pgfpathlineto{\pgfqpoint{5.867118in}{1.282258in}}%
\pgfpathlineto{\pgfqpoint{5.913743in}{1.323734in}}%
\pgfpathlineto{\pgfqpoint{5.960369in}{1.369666in}}%
\pgfpathlineto{\pgfqpoint{6.006995in}{1.420475in}}%
\pgfpathlineto{\pgfqpoint{6.053620in}{1.476610in}}%
\pgfpathlineto{\pgfqpoint{6.100246in}{1.538537in}}%
\pgfpathlineto{\pgfqpoint{6.146872in}{1.606723in}}%
\pgfpathlineto{\pgfqpoint{6.193497in}{1.681586in}}%
\pgfpathlineto{\pgfqpoint{6.240123in}{1.763452in}}%
\pgfpathlineto{\pgfqpoint{6.286749in}{1.852378in}}%
\pgfpathlineto{\pgfqpoint{6.333374in}{1.947961in}}%
\pgfpathlineto{\pgfqpoint{6.380000in}{2.049064in}}%
\pgfpathlineto{\pgfqpoint{6.426626in}{2.153565in}}%
\pgfpathlineto{\pgfqpoint{6.473251in}{2.258248in}}%
\pgfpathlineto{\pgfqpoint{6.519877in}{2.358964in}}%
\pgfpathlineto{\pgfqpoint{6.566503in}{2.451232in}}%
\pgfpathlineto{\pgfqpoint{6.613128in}{2.531123in}}%
\pgfpathlineto{\pgfqpoint{6.659754in}{2.596086in}}%
\pgfpathlineto{\pgfqpoint{6.706380in}{2.645411in}}%
\pgfpathlineto{\pgfqpoint{6.706380in}{2.645869in}}%
\pgfpathlineto{\pgfqpoint{6.706380in}{2.645869in}}%
\pgfpathlineto{\pgfqpoint{6.659754in}{2.596759in}}%
\pgfpathlineto{\pgfqpoint{6.613128in}{2.532081in}}%
\pgfpathlineto{\pgfqpoint{6.566503in}{2.452510in}}%
\pgfpathlineto{\pgfqpoint{6.519877in}{2.360559in}}%
\pgfpathlineto{\pgfqpoint{6.473251in}{2.260155in}}%
\pgfpathlineto{\pgfqpoint{6.426626in}{2.155756in}}%
\pgfpathlineto{\pgfqpoint{6.380000in}{2.051482in}}%
\pgfpathlineto{\pgfqpoint{6.333374in}{1.950551in}}%
\pgfpathlineto{\pgfqpoint{6.286749in}{1.855094in}}%
\pgfpathlineto{\pgfqpoint{6.240123in}{1.766259in}}%
\pgfpathlineto{\pgfqpoint{6.193497in}{1.684461in}}%
\pgfpathlineto{\pgfqpoint{6.146872in}{1.609651in}}%
\pgfpathlineto{\pgfqpoint{6.100246in}{1.541521in}}%
\pgfpathlineto{\pgfqpoint{6.053620in}{1.479644in}}%
\pgfpathlineto{\pgfqpoint{6.006995in}{1.423555in}}%
\pgfpathlineto{\pgfqpoint{5.960369in}{1.372790in}}%
\pgfpathlineto{\pgfqpoint{5.913743in}{1.326901in}}%
\pgfpathlineto{\pgfqpoint{5.867118in}{1.285464in}}%
\pgfpathlineto{\pgfqpoint{5.820492in}{1.248084in}}%
\pgfpathlineto{\pgfqpoint{5.773866in}{1.214392in}}%
\pgfpathlineto{\pgfqpoint{5.727241in}{1.184047in}}%
\pgfpathlineto{\pgfqpoint{5.680615in}{1.156734in}}%
\pgfpathlineto{\pgfqpoint{5.633989in}{1.132167in}}%
\pgfpathlineto{\pgfqpoint{5.587364in}{1.110083in}}%
\pgfpathlineto{\pgfqpoint{5.540738in}{1.090239in}}%
\pgfpathlineto{\pgfqpoint{5.494112in}{1.072414in}}%
\pgfpathlineto{\pgfqpoint{5.447487in}{1.056409in}}%
\pgfpathlineto{\pgfqpoint{5.400861in}{1.042043in}}%
\pgfpathlineto{\pgfqpoint{5.354235in}{1.029152in}}%
\pgfpathlineto{\pgfqpoint{5.307610in}{1.017586in}}%
\pgfpathlineto{\pgfqpoint{5.260984in}{1.007213in}}%
\pgfpathlineto{\pgfqpoint{5.214358in}{0.997910in}}%
\pgfpathlineto{\pgfqpoint{5.167732in}{0.989568in}}%
\pgfpathlineto{\pgfqpoint{5.121107in}{0.982089in}}%
\pgfpathlineto{\pgfqpoint{5.074481in}{0.975384in}}%
\pgfpathlineto{\pgfqpoint{5.027855in}{0.969373in}}%
\pgfpathlineto{\pgfqpoint{4.981230in}{0.963983in}}%
\pgfpathlineto{\pgfqpoint{4.934604in}{0.959151in}}%
\pgfpathlineto{\pgfqpoint{4.887978in}{0.954818in}}%
\pgfpathlineto{\pgfqpoint{4.841353in}{0.950930in}}%
\pgfpathlineto{\pgfqpoint{4.794727in}{0.947442in}}%
\pgfpathlineto{\pgfqpoint{4.748101in}{0.944310in}}%
\pgfpathlineto{\pgfqpoint{4.701476in}{0.941494in}}%
\pgfpathlineto{\pgfqpoint{4.654850in}{0.938962in}}%
\pgfpathlineto{\pgfqpoint{4.608224in}{0.936679in}}%
\pgfpathlineto{\pgfqpoint{4.561599in}{0.934619in}}%
\pgfpathlineto{\pgfqpoint{4.514973in}{0.932754in}}%
\pgfpathlineto{\pgfqpoint{4.468347in}{0.931061in}}%
\pgfpathlineto{\pgfqpoint{4.421722in}{0.929517in}}%
\pgfpathlineto{\pgfqpoint{4.375096in}{0.928104in}}%
\pgfpathlineto{\pgfqpoint{4.328470in}{0.926803in}}%
\pgfpathlineto{\pgfqpoint{4.281845in}{0.925599in}}%
\pgfpathlineto{\pgfqpoint{4.235219in}{0.924475in}}%
\pgfpathlineto{\pgfqpoint{4.188593in}{0.923420in}}%
\pgfpathlineto{\pgfqpoint{4.188593in}{0.923420in}}%
\pgfpathclose%
\pgfusepath{fill}%
}%
\begin{pgfscope}%
\pgfsys@transformshift{0.000000in}{0.000000in}%
\pgfsys@useobject{currentmarker}{}%
\end{pgfscope}%
\end{pgfscope}%
\begin{pgfscope}%
\pgfpathrectangle{\pgfqpoint{4.062704in}{0.807389in}}{\pgfqpoint{2.769565in}{2.016000in}}%
\pgfusepath{clip}%
\pgfsetbuttcap%
\pgfsetroundjoin%
\definecolor{currentfill}{rgb}{0.172549,0.627451,0.172549}%
\pgfsetfillcolor{currentfill}%
\pgfsetfillopacity{0.150000}%
\pgfsetlinewidth{0.000000pt}%
\definecolor{currentstroke}{rgb}{0.172549,0.627451,0.172549}%
\pgfsetstrokecolor{currentstroke}%
\pgfsetstrokeopacity{0.150000}%
\pgfsetdash{}{0pt}%
\pgfsys@defobject{currentmarker}{\pgfqpoint{4.188593in}{1.031579in}}{\pgfqpoint{6.706380in}{2.706316in}}{%
\pgfpathmoveto{\pgfqpoint{4.188593in}{1.040202in}}%
\pgfpathlineto{\pgfqpoint{4.188593in}{1.031579in}}%
\pgfpathlineto{\pgfqpoint{4.235219in}{1.038799in}}%
\pgfpathlineto{\pgfqpoint{4.281845in}{1.046363in}}%
\pgfpathlineto{\pgfqpoint{4.328470in}{1.054425in}}%
\pgfpathlineto{\pgfqpoint{4.375096in}{1.062960in}}%
\pgfpathlineto{\pgfqpoint{4.421722in}{1.071975in}}%
\pgfpathlineto{\pgfqpoint{4.468347in}{1.081513in}}%
\pgfpathlineto{\pgfqpoint{4.514973in}{1.091617in}}%
\pgfpathlineto{\pgfqpoint{4.561599in}{1.102329in}}%
\pgfpathlineto{\pgfqpoint{4.608224in}{1.113971in}}%
\pgfpathlineto{\pgfqpoint{4.654850in}{1.126030in}}%
\pgfpathlineto{\pgfqpoint{4.701476in}{1.138892in}}%
\pgfpathlineto{\pgfqpoint{4.748101in}{1.152706in}}%
\pgfpathlineto{\pgfqpoint{4.794727in}{1.167379in}}%
\pgfpathlineto{\pgfqpoint{4.841353in}{1.182931in}}%
\pgfpathlineto{\pgfqpoint{4.887978in}{1.198913in}}%
\pgfpathlineto{\pgfqpoint{4.934604in}{1.215931in}}%
\pgfpathlineto{\pgfqpoint{4.981230in}{1.234265in}}%
\pgfpathlineto{\pgfqpoint{5.027855in}{1.253652in}}%
\pgfpathlineto{\pgfqpoint{5.074481in}{1.274150in}}%
\pgfpathlineto{\pgfqpoint{5.121107in}{1.295817in}}%
\pgfpathlineto{\pgfqpoint{5.167732in}{1.318717in}}%
\pgfpathlineto{\pgfqpoint{5.214358in}{1.342892in}}%
\pgfpathlineto{\pgfqpoint{5.260984in}{1.368407in}}%
\pgfpathlineto{\pgfqpoint{5.307610in}{1.395345in}}%
\pgfpathlineto{\pgfqpoint{5.354235in}{1.424004in}}%
\pgfpathlineto{\pgfqpoint{5.400861in}{1.454303in}}%
\pgfpathlineto{\pgfqpoint{5.447487in}{1.486193in}}%
\pgfpathlineto{\pgfqpoint{5.494112in}{1.519817in}}%
\pgfpathlineto{\pgfqpoint{5.540738in}{1.555239in}}%
\pgfpathlineto{\pgfqpoint{5.587364in}{1.592509in}}%
\pgfpathlineto{\pgfqpoint{5.633989in}{1.631930in}}%
\pgfpathlineto{\pgfqpoint{5.680615in}{1.673355in}}%
\pgfpathlineto{\pgfqpoint{5.727241in}{1.716670in}}%
\pgfpathlineto{\pgfqpoint{5.773866in}{1.761785in}}%
\pgfpathlineto{\pgfqpoint{5.820492in}{1.808805in}}%
\pgfpathlineto{\pgfqpoint{5.867118in}{1.857507in}}%
\pgfpathlineto{\pgfqpoint{5.913743in}{1.907814in}}%
\pgfpathlineto{\pgfqpoint{5.960369in}{1.959654in}}%
\pgfpathlineto{\pgfqpoint{6.006995in}{2.013228in}}%
\pgfpathlineto{\pgfqpoint{6.053620in}{2.068850in}}%
\pgfpathlineto{\pgfqpoint{6.100246in}{2.125515in}}%
\pgfpathlineto{\pgfqpoint{6.146872in}{2.183045in}}%
\pgfpathlineto{\pgfqpoint{6.193497in}{2.240831in}}%
\pgfpathlineto{\pgfqpoint{6.240123in}{2.298655in}}%
\pgfpathlineto{\pgfqpoint{6.286749in}{2.355794in}}%
\pgfpathlineto{\pgfqpoint{6.333374in}{2.411314in}}%
\pgfpathlineto{\pgfqpoint{6.380000in}{2.464432in}}%
\pgfpathlineto{\pgfqpoint{6.426626in}{2.514121in}}%
\pgfpathlineto{\pgfqpoint{6.473251in}{2.559316in}}%
\pgfpathlineto{\pgfqpoint{6.519877in}{2.599072in}}%
\pgfpathlineto{\pgfqpoint{6.566503in}{2.632787in}}%
\pgfpathlineto{\pgfqpoint{6.613128in}{2.660207in}}%
\pgfpathlineto{\pgfqpoint{6.659754in}{2.681571in}}%
\pgfpathlineto{\pgfqpoint{6.706380in}{2.697531in}}%
\pgfpathlineto{\pgfqpoint{6.706380in}{2.706316in}}%
\pgfpathlineto{\pgfqpoint{6.706380in}{2.706316in}}%
\pgfpathlineto{\pgfqpoint{6.659754in}{2.691740in}}%
\pgfpathlineto{\pgfqpoint{6.613128in}{2.671690in}}%
\pgfpathlineto{\pgfqpoint{6.566503in}{2.645527in}}%
\pgfpathlineto{\pgfqpoint{6.519877in}{2.613062in}}%
\pgfpathlineto{\pgfqpoint{6.473251in}{2.574415in}}%
\pgfpathlineto{\pgfqpoint{6.426626in}{2.530161in}}%
\pgfpathlineto{\pgfqpoint{6.380000in}{2.481211in}}%
\pgfpathlineto{\pgfqpoint{6.333374in}{2.428620in}}%
\pgfpathlineto{\pgfqpoint{6.286749in}{2.373469in}}%
\pgfpathlineto{\pgfqpoint{6.240123in}{2.316878in}}%
\pgfpathlineto{\pgfqpoint{6.193497in}{2.259573in}}%
\pgfpathlineto{\pgfqpoint{6.146872in}{2.202201in}}%
\pgfpathlineto{\pgfqpoint{6.100246in}{2.144989in}}%
\pgfpathlineto{\pgfqpoint{6.053620in}{2.088572in}}%
\pgfpathlineto{\pgfqpoint{6.006995in}{2.033266in}}%
\pgfpathlineto{\pgfqpoint{5.960369in}{1.979549in}}%
\pgfpathlineto{\pgfqpoint{5.913743in}{1.927559in}}%
\pgfpathlineto{\pgfqpoint{5.867118in}{1.877286in}}%
\pgfpathlineto{\pgfqpoint{5.820492in}{1.828633in}}%
\pgfpathlineto{\pgfqpoint{5.773866in}{1.781646in}}%
\pgfpathlineto{\pgfqpoint{5.727241in}{1.736484in}}%
\pgfpathlineto{\pgfqpoint{5.680615in}{1.693016in}}%
\pgfpathlineto{\pgfqpoint{5.633989in}{1.651473in}}%
\pgfpathlineto{\pgfqpoint{5.587364in}{1.611927in}}%
\pgfpathlineto{\pgfqpoint{5.540738in}{1.574356in}}%
\pgfpathlineto{\pgfqpoint{5.494112in}{1.538459in}}%
\pgfpathlineto{\pgfqpoint{5.447487in}{1.504084in}}%
\pgfpathlineto{\pgfqpoint{5.400861in}{1.471395in}}%
\pgfpathlineto{\pgfqpoint{5.354235in}{1.440390in}}%
\pgfpathlineto{\pgfqpoint{5.307610in}{1.411140in}}%
\pgfpathlineto{\pgfqpoint{5.260984in}{1.383581in}}%
\pgfpathlineto{\pgfqpoint{5.214358in}{1.357745in}}%
\pgfpathlineto{\pgfqpoint{5.167732in}{1.333250in}}%
\pgfpathlineto{\pgfqpoint{5.121107in}{1.309662in}}%
\pgfpathlineto{\pgfqpoint{5.074481in}{1.287675in}}%
\pgfpathlineto{\pgfqpoint{5.027855in}{1.266750in}}%
\pgfpathlineto{\pgfqpoint{4.981230in}{1.247024in}}%
\pgfpathlineto{\pgfqpoint{4.934604in}{1.228428in}}%
\pgfpathlineto{\pgfqpoint{4.887978in}{1.211084in}}%
\pgfpathlineto{\pgfqpoint{4.841353in}{1.194902in}}%
\pgfpathlineto{\pgfqpoint{4.794727in}{1.179802in}}%
\pgfpathlineto{\pgfqpoint{4.748101in}{1.165571in}}%
\pgfpathlineto{\pgfqpoint{4.701476in}{1.152009in}}%
\pgfpathlineto{\pgfqpoint{4.654850in}{1.138741in}}%
\pgfpathlineto{\pgfqpoint{4.608224in}{1.126208in}}%
\pgfpathlineto{\pgfqpoint{4.561599in}{1.114363in}}%
\pgfpathlineto{\pgfqpoint{4.514973in}{1.103160in}}%
\pgfpathlineto{\pgfqpoint{4.468347in}{1.092558in}}%
\pgfpathlineto{\pgfqpoint{4.421722in}{1.082513in}}%
\pgfpathlineto{\pgfqpoint{4.375096in}{1.073108in}}%
\pgfpathlineto{\pgfqpoint{4.328470in}{1.064207in}}%
\pgfpathlineto{\pgfqpoint{4.281845in}{1.055744in}}%
\pgfpathlineto{\pgfqpoint{4.235219in}{1.047773in}}%
\pgfpathlineto{\pgfqpoint{4.188593in}{1.040202in}}%
\pgfpathlineto{\pgfqpoint{4.188593in}{1.040202in}}%
\pgfpathclose%
\pgfusepath{fill}%
}%
\begin{pgfscope}%
\pgfsys@transformshift{0.000000in}{0.000000in}%
\pgfsys@useobject{currentmarker}{}%
\end{pgfscope}%
\end{pgfscope}%
\begin{pgfscope}%
\pgfpathrectangle{\pgfqpoint{4.062704in}{0.807389in}}{\pgfqpoint{2.769565in}{2.016000in}}%
\pgfusepath{clip}%
\pgfsetbuttcap%
\pgfsetroundjoin%
\definecolor{currentfill}{rgb}{0.839216,0.152941,0.156863}%
\pgfsetfillcolor{currentfill}%
\pgfsetfillopacity{0.150000}%
\pgfsetlinewidth{0.000000pt}%
\definecolor{currentstroke}{rgb}{0.839216,0.152941,0.156863}%
\pgfsetstrokecolor{currentstroke}%
\pgfsetstrokeopacity{0.150000}%
\pgfsetdash{}{0pt}%
\pgfsys@defobject{currentmarker}{\pgfqpoint{4.188593in}{0.954310in}}{\pgfqpoint{6.706380in}{2.731128in}}{%
\pgfpathmoveto{\pgfqpoint{4.188593in}{0.959859in}}%
\pgfpathlineto{\pgfqpoint{4.188593in}{0.954310in}}%
\pgfpathlineto{\pgfqpoint{4.235219in}{0.956558in}}%
\pgfpathlineto{\pgfqpoint{4.281845in}{0.959007in}}%
\pgfpathlineto{\pgfqpoint{4.328470in}{0.961629in}}%
\pgfpathlineto{\pgfqpoint{4.375096in}{0.964445in}}%
\pgfpathlineto{\pgfqpoint{4.421722in}{0.967375in}}%
\pgfpathlineto{\pgfqpoint{4.468347in}{0.970566in}}%
\pgfpathlineto{\pgfqpoint{4.514973in}{0.974004in}}%
\pgfpathlineto{\pgfqpoint{4.561599in}{0.977875in}}%
\pgfpathlineto{\pgfqpoint{4.608224in}{0.982143in}}%
\pgfpathlineto{\pgfqpoint{4.654850in}{0.986853in}}%
\pgfpathlineto{\pgfqpoint{4.701476in}{0.992060in}}%
\pgfpathlineto{\pgfqpoint{4.748101in}{0.997828in}}%
\pgfpathlineto{\pgfqpoint{4.794727in}{1.004226in}}%
\pgfpathlineto{\pgfqpoint{4.841353in}{1.011319in}}%
\pgfpathlineto{\pgfqpoint{4.887978in}{1.019178in}}%
\pgfpathlineto{\pgfqpoint{4.934604in}{1.027924in}}%
\pgfpathlineto{\pgfqpoint{4.981230in}{1.037664in}}%
\pgfpathlineto{\pgfqpoint{5.027855in}{1.048533in}}%
\pgfpathlineto{\pgfqpoint{5.074481in}{1.060701in}}%
\pgfpathlineto{\pgfqpoint{5.121107in}{1.074250in}}%
\pgfpathlineto{\pgfqpoint{5.167732in}{1.089341in}}%
\pgfpathlineto{\pgfqpoint{5.214358in}{1.106153in}}%
\pgfpathlineto{\pgfqpoint{5.260984in}{1.124882in}}%
\pgfpathlineto{\pgfqpoint{5.307610in}{1.145748in}}%
\pgfpathlineto{\pgfqpoint{5.354235in}{1.168994in}}%
\pgfpathlineto{\pgfqpoint{5.400861in}{1.194878in}}%
\pgfpathlineto{\pgfqpoint{5.447487in}{1.223701in}}%
\pgfpathlineto{\pgfqpoint{5.494112in}{1.255767in}}%
\pgfpathlineto{\pgfqpoint{5.540738in}{1.291467in}}%
\pgfpathlineto{\pgfqpoint{5.587364in}{1.331207in}}%
\pgfpathlineto{\pgfqpoint{5.633989in}{1.375427in}}%
\pgfpathlineto{\pgfqpoint{5.680615in}{1.424603in}}%
\pgfpathlineto{\pgfqpoint{5.727241in}{1.479250in}}%
\pgfpathlineto{\pgfqpoint{5.773866in}{1.539905in}}%
\pgfpathlineto{\pgfqpoint{5.820492in}{1.607097in}}%
\pgfpathlineto{\pgfqpoint{5.867118in}{1.681284in}}%
\pgfpathlineto{\pgfqpoint{5.913743in}{1.762748in}}%
\pgfpathlineto{\pgfqpoint{5.960369in}{1.851436in}}%
\pgfpathlineto{\pgfqpoint{6.006995in}{1.946748in}}%
\pgfpathlineto{\pgfqpoint{6.053620in}{2.047311in}}%
\pgfpathlineto{\pgfqpoint{6.100246in}{2.150810in}}%
\pgfpathlineto{\pgfqpoint{6.146872in}{2.253949in}}%
\pgfpathlineto{\pgfqpoint{6.193497in}{2.352765in}}%
\pgfpathlineto{\pgfqpoint{6.240123in}{2.443357in}}%
\pgfpathlineto{\pgfqpoint{6.286749in}{2.522176in}}%
\pgfpathlineto{\pgfqpoint{6.333374in}{2.586891in}}%
\pgfpathlineto{\pgfqpoint{6.380000in}{2.636879in}}%
\pgfpathlineto{\pgfqpoint{6.426626in}{2.673069in}}%
\pgfpathlineto{\pgfqpoint{6.473251in}{2.697638in}}%
\pgfpathlineto{\pgfqpoint{6.519877in}{2.713150in}}%
\pgfpathlineto{\pgfqpoint{6.566503in}{2.722281in}}%
\pgfpathlineto{\pgfqpoint{6.613128in}{2.727251in}}%
\pgfpathlineto{\pgfqpoint{6.659754in}{2.729779in}}%
\pgfpathlineto{\pgfqpoint{6.706380in}{2.730947in}}%
\pgfpathlineto{\pgfqpoint{6.706380in}{2.731128in}}%
\pgfpathlineto{\pgfqpoint{6.706380in}{2.731128in}}%
\pgfpathlineto{\pgfqpoint{6.659754in}{2.729985in}}%
\pgfpathlineto{\pgfqpoint{6.613128in}{2.727505in}}%
\pgfpathlineto{\pgfqpoint{6.566503in}{2.722522in}}%
\pgfpathlineto{\pgfqpoint{6.519877in}{2.713351in}}%
\pgfpathlineto{\pgfqpoint{6.473251in}{2.697920in}}%
\pgfpathlineto{\pgfqpoint{6.426626in}{2.673623in}}%
\pgfpathlineto{\pgfqpoint{6.380000in}{2.637721in}}%
\pgfpathlineto{\pgfqpoint{6.333374in}{2.588260in}}%
\pgfpathlineto{\pgfqpoint{6.286749in}{2.524292in}}%
\pgfpathlineto{\pgfqpoint{6.240123in}{2.446338in}}%
\pgfpathlineto{\pgfqpoint{6.193497in}{2.356466in}}%
\pgfpathlineto{\pgfqpoint{6.146872in}{2.258190in}}%
\pgfpathlineto{\pgfqpoint{6.100246in}{2.155582in}}%
\pgfpathlineto{\pgfqpoint{6.053620in}{2.052575in}}%
\pgfpathlineto{\pgfqpoint{6.006995in}{1.952396in}}%
\pgfpathlineto{\pgfqpoint{5.960369in}{1.857366in}}%
\pgfpathlineto{\pgfqpoint{5.913743in}{1.768875in}}%
\pgfpathlineto{\pgfqpoint{5.867118in}{1.687543in}}%
\pgfpathlineto{\pgfqpoint{5.820492in}{1.613444in}}%
\pgfpathlineto{\pgfqpoint{5.773866in}{1.546309in}}%
\pgfpathlineto{\pgfqpoint{5.727241in}{1.485690in}}%
\pgfpathlineto{\pgfqpoint{5.680615in}{1.431065in}}%
\pgfpathlineto{\pgfqpoint{5.633989in}{1.381901in}}%
\pgfpathlineto{\pgfqpoint{5.587364in}{1.337686in}}%
\pgfpathlineto{\pgfqpoint{5.540738in}{1.297945in}}%
\pgfpathlineto{\pgfqpoint{5.494112in}{1.262237in}}%
\pgfpathlineto{\pgfqpoint{5.447487in}{1.230165in}}%
\pgfpathlineto{\pgfqpoint{5.400861in}{1.201365in}}%
\pgfpathlineto{\pgfqpoint{5.354235in}{1.175507in}}%
\pgfpathlineto{\pgfqpoint{5.307610in}{1.152294in}}%
\pgfpathlineto{\pgfqpoint{5.260984in}{1.131457in}}%
\pgfpathlineto{\pgfqpoint{5.214358in}{1.112752in}}%
\pgfpathlineto{\pgfqpoint{5.167732in}{1.095960in}}%
\pgfpathlineto{\pgfqpoint{5.121107in}{1.080900in}}%
\pgfpathlineto{\pgfqpoint{5.074481in}{1.067388in}}%
\pgfpathlineto{\pgfqpoint{5.027855in}{1.055252in}}%
\pgfpathlineto{\pgfqpoint{4.981230in}{1.044348in}}%
\pgfpathlineto{\pgfqpoint{4.934604in}{1.034615in}}%
\pgfpathlineto{\pgfqpoint{4.887978in}{1.025883in}}%
\pgfpathlineto{\pgfqpoint{4.841353in}{1.018033in}}%
\pgfpathlineto{\pgfqpoint{4.794727in}{1.010971in}}%
\pgfpathlineto{\pgfqpoint{4.748101in}{1.004609in}}%
\pgfpathlineto{\pgfqpoint{4.701476in}{0.998818in}}%
\pgfpathlineto{\pgfqpoint{4.654850in}{0.993393in}}%
\pgfpathlineto{\pgfqpoint{4.608224in}{0.988473in}}%
\pgfpathlineto{\pgfqpoint{4.561599in}{0.983987in}}%
\pgfpathlineto{\pgfqpoint{4.514973in}{0.979893in}}%
\pgfpathlineto{\pgfqpoint{4.468347in}{0.976272in}}%
\pgfpathlineto{\pgfqpoint{4.421722in}{0.972992in}}%
\pgfpathlineto{\pgfqpoint{4.375096in}{0.969975in}}%
\pgfpathlineto{\pgfqpoint{4.328470in}{0.967184in}}%
\pgfpathlineto{\pgfqpoint{4.281845in}{0.964587in}}%
\pgfpathlineto{\pgfqpoint{4.235219in}{0.962155in}}%
\pgfpathlineto{\pgfqpoint{4.188593in}{0.959859in}}%
\pgfpathlineto{\pgfqpoint{4.188593in}{0.959859in}}%
\pgfpathclose%
\pgfusepath{fill}%
}%
\begin{pgfscope}%
\pgfsys@transformshift{0.000000in}{0.000000in}%
\pgfsys@useobject{currentmarker}{}%
\end{pgfscope}%
\end{pgfscope}%
\begin{pgfscope}%
\pgfpathrectangle{\pgfqpoint{4.062704in}{0.807389in}}{\pgfqpoint{2.769565in}{2.016000in}}%
\pgfusepath{clip}%
\pgfsetbuttcap%
\pgfsetroundjoin%
\definecolor{currentfill}{rgb}{0.580392,0.403922,0.741176}%
\pgfsetfillcolor{currentfill}%
\pgfsetfillopacity{0.150000}%
\pgfsetlinewidth{0.000000pt}%
\definecolor{currentstroke}{rgb}{0.580392,0.403922,0.741176}%
\pgfsetstrokecolor{currentstroke}%
\pgfsetstrokeopacity{0.150000}%
\pgfsetdash{}{0pt}%
\pgfsys@defobject{currentmarker}{\pgfqpoint{4.188593in}{0.949495in}}{\pgfqpoint{6.706380in}{2.731753in}}{%
\pgfpathmoveto{\pgfqpoint{4.188593in}{0.949495in}}%
\pgfpathlineto{\pgfqpoint{4.188593in}{0.949495in}}%
\pgfpathlineto{\pgfqpoint{4.235219in}{0.957763in}}%
\pgfpathlineto{\pgfqpoint{4.281845in}{0.966999in}}%
\pgfpathlineto{\pgfqpoint{4.328470in}{0.977318in}}%
\pgfpathlineto{\pgfqpoint{4.375096in}{0.988845in}}%
\pgfpathlineto{\pgfqpoint{4.421722in}{1.001722in}}%
\pgfpathlineto{\pgfqpoint{4.468347in}{1.016106in}}%
\pgfpathlineto{\pgfqpoint{4.514973in}{1.032174in}}%
\pgfpathlineto{\pgfqpoint{4.561599in}{1.050121in}}%
\pgfpathlineto{\pgfqpoint{4.608224in}{1.070168in}}%
\pgfpathlineto{\pgfqpoint{4.654850in}{1.092559in}}%
\pgfpathlineto{\pgfqpoint{4.701476in}{1.117566in}}%
\pgfpathlineto{\pgfqpoint{4.748101in}{1.145494in}}%
\pgfpathlineto{\pgfqpoint{4.794727in}{1.176683in}}%
\pgfpathlineto{\pgfqpoint{4.841353in}{1.211510in}}%
\pgfpathlineto{\pgfqpoint{4.887978in}{1.250398in}}%
\pgfpathlineto{\pgfqpoint{4.934604in}{1.293814in}}%
\pgfpathlineto{\pgfqpoint{4.981230in}{1.342275in}}%
\pgfpathlineto{\pgfqpoint{5.027855in}{1.396347in}}%
\pgfpathlineto{\pgfqpoint{5.074481in}{1.456624in}}%
\pgfpathlineto{\pgfqpoint{5.121107in}{1.523705in}}%
\pgfpathlineto{\pgfqpoint{5.167732in}{1.598120in}}%
\pgfpathlineto{\pgfqpoint{5.214358in}{1.680222in}}%
\pgfpathlineto{\pgfqpoint{5.260984in}{1.770024in}}%
\pgfpathlineto{\pgfqpoint{5.307610in}{1.866980in}}%
\pgfpathlineto{\pgfqpoint{5.354235in}{1.969773in}}%
\pgfpathlineto{\pgfqpoint{5.400861in}{2.076142in}}%
\pgfpathlineto{\pgfqpoint{5.447487in}{2.182883in}}%
\pgfpathlineto{\pgfqpoint{5.494112in}{2.286100in}}%
\pgfpathlineto{\pgfqpoint{5.540738in}{2.381733in}}%
\pgfpathlineto{\pgfqpoint{5.587364in}{2.466259in}}%
\pgfpathlineto{\pgfqpoint{5.633989in}{2.537335in}}%
\pgfpathlineto{\pgfqpoint{5.680615in}{2.594150in}}%
\pgfpathlineto{\pgfqpoint{5.727241in}{2.637369in}}%
\pgfpathlineto{\pgfqpoint{5.773866in}{2.668743in}}%
\pgfpathlineto{\pgfqpoint{5.820492in}{2.690571in}}%
\pgfpathlineto{\pgfqpoint{5.867118in}{2.705208in}}%
\pgfpathlineto{\pgfqpoint{5.913743in}{2.714733in}}%
\pgfpathlineto{\pgfqpoint{5.960369in}{2.720801in}}%
\pgfpathlineto{\pgfqpoint{6.006995in}{2.724623in}}%
\pgfpathlineto{\pgfqpoint{6.053620in}{2.727028in}}%
\pgfpathlineto{\pgfqpoint{6.100246in}{2.728559in}}%
\pgfpathlineto{\pgfqpoint{6.146872in}{2.729550in}}%
\pgfpathlineto{\pgfqpoint{6.193497in}{2.730209in}}%
\pgfpathlineto{\pgfqpoint{6.240123in}{2.730657in}}%
\pgfpathlineto{\pgfqpoint{6.286749in}{2.730969in}}%
\pgfpathlineto{\pgfqpoint{6.333374in}{2.731190in}}%
\pgfpathlineto{\pgfqpoint{6.380000in}{2.731348in}}%
\pgfpathlineto{\pgfqpoint{6.426626in}{2.731462in}}%
\pgfpathlineto{\pgfqpoint{6.473251in}{2.731545in}}%
\pgfpathlineto{\pgfqpoint{6.519877in}{2.731608in}}%
\pgfpathlineto{\pgfqpoint{6.566503in}{2.731655in}}%
\pgfpathlineto{\pgfqpoint{6.613128in}{2.731693in}}%
\pgfpathlineto{\pgfqpoint{6.659754in}{2.731725in}}%
\pgfpathlineto{\pgfqpoint{6.706380in}{2.731753in}}%
\pgfpathlineto{\pgfqpoint{6.706380in}{2.731753in}}%
\pgfpathlineto{\pgfqpoint{6.706380in}{2.731753in}}%
\pgfpathlineto{\pgfqpoint{6.659754in}{2.731725in}}%
\pgfpathlineto{\pgfqpoint{6.613128in}{2.731693in}}%
\pgfpathlineto{\pgfqpoint{6.566503in}{2.731655in}}%
\pgfpathlineto{\pgfqpoint{6.519877in}{2.731608in}}%
\pgfpathlineto{\pgfqpoint{6.473251in}{2.731545in}}%
\pgfpathlineto{\pgfqpoint{6.426626in}{2.731462in}}%
\pgfpathlineto{\pgfqpoint{6.380000in}{2.731348in}}%
\pgfpathlineto{\pgfqpoint{6.333374in}{2.731190in}}%
\pgfpathlineto{\pgfqpoint{6.286749in}{2.730969in}}%
\pgfpathlineto{\pgfqpoint{6.240123in}{2.730657in}}%
\pgfpathlineto{\pgfqpoint{6.193497in}{2.730209in}}%
\pgfpathlineto{\pgfqpoint{6.146872in}{2.729550in}}%
\pgfpathlineto{\pgfqpoint{6.100246in}{2.728559in}}%
\pgfpathlineto{\pgfqpoint{6.053620in}{2.727028in}}%
\pgfpathlineto{\pgfqpoint{6.006995in}{2.724623in}}%
\pgfpathlineto{\pgfqpoint{5.960369in}{2.720801in}}%
\pgfpathlineto{\pgfqpoint{5.913743in}{2.714733in}}%
\pgfpathlineto{\pgfqpoint{5.867118in}{2.705208in}}%
\pgfpathlineto{\pgfqpoint{5.820492in}{2.690571in}}%
\pgfpathlineto{\pgfqpoint{5.773866in}{2.668743in}}%
\pgfpathlineto{\pgfqpoint{5.727241in}{2.637369in}}%
\pgfpathlineto{\pgfqpoint{5.680615in}{2.594150in}}%
\pgfpathlineto{\pgfqpoint{5.633989in}{2.537335in}}%
\pgfpathlineto{\pgfqpoint{5.587364in}{2.466259in}}%
\pgfpathlineto{\pgfqpoint{5.540738in}{2.381733in}}%
\pgfpathlineto{\pgfqpoint{5.494112in}{2.286100in}}%
\pgfpathlineto{\pgfqpoint{5.447487in}{2.182883in}}%
\pgfpathlineto{\pgfqpoint{5.400861in}{2.076142in}}%
\pgfpathlineto{\pgfqpoint{5.354235in}{1.969773in}}%
\pgfpathlineto{\pgfqpoint{5.307610in}{1.866980in}}%
\pgfpathlineto{\pgfqpoint{5.260984in}{1.770024in}}%
\pgfpathlineto{\pgfqpoint{5.214358in}{1.680222in}}%
\pgfpathlineto{\pgfqpoint{5.167732in}{1.598120in}}%
\pgfpathlineto{\pgfqpoint{5.121107in}{1.523705in}}%
\pgfpathlineto{\pgfqpoint{5.074481in}{1.456624in}}%
\pgfpathlineto{\pgfqpoint{5.027855in}{1.396347in}}%
\pgfpathlineto{\pgfqpoint{4.981230in}{1.342275in}}%
\pgfpathlineto{\pgfqpoint{4.934604in}{1.293814in}}%
\pgfpathlineto{\pgfqpoint{4.887978in}{1.250398in}}%
\pgfpathlineto{\pgfqpoint{4.841353in}{1.211510in}}%
\pgfpathlineto{\pgfqpoint{4.794727in}{1.176683in}}%
\pgfpathlineto{\pgfqpoint{4.748101in}{1.145494in}}%
\pgfpathlineto{\pgfqpoint{4.701476in}{1.117566in}}%
\pgfpathlineto{\pgfqpoint{4.654850in}{1.092559in}}%
\pgfpathlineto{\pgfqpoint{4.608224in}{1.070168in}}%
\pgfpathlineto{\pgfqpoint{4.561599in}{1.050121in}}%
\pgfpathlineto{\pgfqpoint{4.514973in}{1.032174in}}%
\pgfpathlineto{\pgfqpoint{4.468347in}{1.016106in}}%
\pgfpathlineto{\pgfqpoint{4.421722in}{1.001722in}}%
\pgfpathlineto{\pgfqpoint{4.375096in}{0.988845in}}%
\pgfpathlineto{\pgfqpoint{4.328470in}{0.977318in}}%
\pgfpathlineto{\pgfqpoint{4.281845in}{0.966999in}}%
\pgfpathlineto{\pgfqpoint{4.235219in}{0.957763in}}%
\pgfpathlineto{\pgfqpoint{4.188593in}{0.949495in}}%
\pgfpathlineto{\pgfqpoint{4.188593in}{0.949495in}}%
\pgfpathclose%
\pgfusepath{fill}%
}%
\begin{pgfscope}%
\pgfsys@transformshift{0.000000in}{0.000000in}%
\pgfsys@useobject{currentmarker}{}%
\end{pgfscope}%
\end{pgfscope}%
\begin{pgfscope}%
\pgfpathrectangle{\pgfqpoint{4.062704in}{0.807389in}}{\pgfqpoint{2.769565in}{2.016000in}}%
\pgfusepath{clip}%
\pgfsetrectcap%
\pgfsetroundjoin%
\pgfsetlinewidth{0.803000pt}%
\definecolor{currentstroke}{rgb}{0.690196,0.690196,0.690196}%
\pgfsetstrokecolor{currentstroke}%
\pgfsetstrokeopacity{0.300000}%
\pgfsetdash{}{0pt}%
\pgfpathmoveto{\pgfqpoint{4.771154in}{0.807389in}}%
\pgfpathlineto{\pgfqpoint{4.771154in}{2.823389in}}%
\pgfusepath{stroke}%
\end{pgfscope}%
\begin{pgfscope}%
\pgfsetbuttcap%
\pgfsetroundjoin%
\definecolor{currentfill}{rgb}{0.000000,0.000000,0.000000}%
\pgfsetfillcolor{currentfill}%
\pgfsetlinewidth{0.803000pt}%
\definecolor{currentstroke}{rgb}{0.000000,0.000000,0.000000}%
\pgfsetstrokecolor{currentstroke}%
\pgfsetdash{}{0pt}%
\pgfsys@defobject{currentmarker}{\pgfqpoint{0.000000in}{-0.048611in}}{\pgfqpoint{0.000000in}{0.000000in}}{%
\pgfpathmoveto{\pgfqpoint{0.000000in}{0.000000in}}%
\pgfpathlineto{\pgfqpoint{0.000000in}{-0.048611in}}%
\pgfusepath{stroke,fill}%
}%
\begin{pgfscope}%
\pgfsys@transformshift{4.771154in}{0.807389in}%
\pgfsys@useobject{currentmarker}{}%
\end{pgfscope}%
\end{pgfscope}%
\begin{pgfscope}%
\definecolor{textcolor}{rgb}{0.000000,0.000000,0.000000}%
\pgfsetstrokecolor{textcolor}%
\pgfsetfillcolor{textcolor}%
\pgftext[x=4.771154in,y=0.710167in,,top]{\color{textcolor}{\rmfamily\fontsize{7.000000}{8.400000}\selectfont\catcode`\^=\active\def^{\ifmmode\sp\else\^{}\fi}\catcode`\%=\active\def%{\%}$\mathdefault{10^{1}}$}}%
\end{pgfscope}%
\begin{pgfscope}%
\pgfpathrectangle{\pgfqpoint{4.062704in}{0.807389in}}{\pgfqpoint{2.769565in}{2.016000in}}%
\pgfusepath{clip}%
\pgfsetrectcap%
\pgfsetroundjoin%
\pgfsetlinewidth{0.803000pt}%
\definecolor{currentstroke}{rgb}{0.690196,0.690196,0.690196}%
\pgfsetstrokecolor{currentstroke}%
\pgfsetstrokeopacity{0.300000}%
\pgfsetdash{}{0pt}%
\pgfpathmoveto{\pgfqpoint{6.706380in}{0.807389in}}%
\pgfpathlineto{\pgfqpoint{6.706380in}{2.823389in}}%
\pgfusepath{stroke}%
\end{pgfscope}%
\begin{pgfscope}%
\pgfsetbuttcap%
\pgfsetroundjoin%
\definecolor{currentfill}{rgb}{0.000000,0.000000,0.000000}%
\pgfsetfillcolor{currentfill}%
\pgfsetlinewidth{0.803000pt}%
\definecolor{currentstroke}{rgb}{0.000000,0.000000,0.000000}%
\pgfsetstrokecolor{currentstroke}%
\pgfsetdash{}{0pt}%
\pgfsys@defobject{currentmarker}{\pgfqpoint{0.000000in}{-0.048611in}}{\pgfqpoint{0.000000in}{0.000000in}}{%
\pgfpathmoveto{\pgfqpoint{0.000000in}{0.000000in}}%
\pgfpathlineto{\pgfqpoint{0.000000in}{-0.048611in}}%
\pgfusepath{stroke,fill}%
}%
\begin{pgfscope}%
\pgfsys@transformshift{6.706380in}{0.807389in}%
\pgfsys@useobject{currentmarker}{}%
\end{pgfscope}%
\end{pgfscope}%
\begin{pgfscope}%
\definecolor{textcolor}{rgb}{0.000000,0.000000,0.000000}%
\pgfsetstrokecolor{textcolor}%
\pgfsetfillcolor{textcolor}%
\pgftext[x=6.706380in,y=0.710167in,,top]{\color{textcolor}{\rmfamily\fontsize{7.000000}{8.400000}\selectfont\catcode`\^=\active\def^{\ifmmode\sp\else\^{}\fi}\catcode`\%=\active\def%{\%}$\mathdefault{10^{2}}$}}%
\end{pgfscope}%
\begin{pgfscope}%
\pgfsetbuttcap%
\pgfsetroundjoin%
\definecolor{currentfill}{rgb}{0.000000,0.000000,0.000000}%
\pgfsetfillcolor{currentfill}%
\pgfsetlinewidth{0.602250pt}%
\definecolor{currentstroke}{rgb}{0.000000,0.000000,0.000000}%
\pgfsetstrokecolor{currentstroke}%
\pgfsetdash{}{0pt}%
\pgfsys@defobject{currentmarker}{\pgfqpoint{0.000000in}{-0.027778in}}{\pgfqpoint{0.000000in}{0.000000in}}{%
\pgfpathmoveto{\pgfqpoint{0.000000in}{0.000000in}}%
\pgfpathlineto{\pgfqpoint{0.000000in}{-0.027778in}}%
\pgfusepath{stroke,fill}%
}%
\begin{pgfscope}%
\pgfsys@transformshift{4.188593in}{0.807389in}%
\pgfsys@useobject{currentmarker}{}%
\end{pgfscope}%
\end{pgfscope}%
\begin{pgfscope}%
\pgfsetbuttcap%
\pgfsetroundjoin%
\definecolor{currentfill}{rgb}{0.000000,0.000000,0.000000}%
\pgfsetfillcolor{currentfill}%
\pgfsetlinewidth{0.602250pt}%
\definecolor{currentstroke}{rgb}{0.000000,0.000000,0.000000}%
\pgfsetstrokecolor{currentstroke}%
\pgfsetdash{}{0pt}%
\pgfsys@defobject{currentmarker}{\pgfqpoint{0.000000in}{-0.027778in}}{\pgfqpoint{0.000000in}{0.000000in}}{%
\pgfpathmoveto{\pgfqpoint{0.000000in}{0.000000in}}%
\pgfpathlineto{\pgfqpoint{0.000000in}{-0.027778in}}%
\pgfusepath{stroke,fill}%
}%
\begin{pgfscope}%
\pgfsys@transformshift{4.341827in}{0.807389in}%
\pgfsys@useobject{currentmarker}{}%
\end{pgfscope}%
\end{pgfscope}%
\begin{pgfscope}%
\pgfsetbuttcap%
\pgfsetroundjoin%
\definecolor{currentfill}{rgb}{0.000000,0.000000,0.000000}%
\pgfsetfillcolor{currentfill}%
\pgfsetlinewidth{0.602250pt}%
\definecolor{currentstroke}{rgb}{0.000000,0.000000,0.000000}%
\pgfsetstrokecolor{currentstroke}%
\pgfsetdash{}{0pt}%
\pgfsys@defobject{currentmarker}{\pgfqpoint{0.000000in}{-0.027778in}}{\pgfqpoint{0.000000in}{0.000000in}}{%
\pgfpathmoveto{\pgfqpoint{0.000000in}{0.000000in}}%
\pgfpathlineto{\pgfqpoint{0.000000in}{-0.027778in}}%
\pgfusepath{stroke,fill}%
}%
\begin{pgfscope}%
\pgfsys@transformshift{4.471384in}{0.807389in}%
\pgfsys@useobject{currentmarker}{}%
\end{pgfscope}%
\end{pgfscope}%
\begin{pgfscope}%
\pgfsetbuttcap%
\pgfsetroundjoin%
\definecolor{currentfill}{rgb}{0.000000,0.000000,0.000000}%
\pgfsetfillcolor{currentfill}%
\pgfsetlinewidth{0.602250pt}%
\definecolor{currentstroke}{rgb}{0.000000,0.000000,0.000000}%
\pgfsetstrokecolor{currentstroke}%
\pgfsetdash{}{0pt}%
\pgfsys@defobject{currentmarker}{\pgfqpoint{0.000000in}{-0.027778in}}{\pgfqpoint{0.000000in}{0.000000in}}{%
\pgfpathmoveto{\pgfqpoint{0.000000in}{0.000000in}}%
\pgfpathlineto{\pgfqpoint{0.000000in}{-0.027778in}}%
\pgfusepath{stroke,fill}%
}%
\begin{pgfscope}%
\pgfsys@transformshift{4.583611in}{0.807389in}%
\pgfsys@useobject{currentmarker}{}%
\end{pgfscope}%
\end{pgfscope}%
\begin{pgfscope}%
\pgfsetbuttcap%
\pgfsetroundjoin%
\definecolor{currentfill}{rgb}{0.000000,0.000000,0.000000}%
\pgfsetfillcolor{currentfill}%
\pgfsetlinewidth{0.602250pt}%
\definecolor{currentstroke}{rgb}{0.000000,0.000000,0.000000}%
\pgfsetstrokecolor{currentstroke}%
\pgfsetdash{}{0pt}%
\pgfsys@defobject{currentmarker}{\pgfqpoint{0.000000in}{-0.027778in}}{\pgfqpoint{0.000000in}{0.000000in}}{%
\pgfpathmoveto{\pgfqpoint{0.000000in}{0.000000in}}%
\pgfpathlineto{\pgfqpoint{0.000000in}{-0.027778in}}%
\pgfusepath{stroke,fill}%
}%
\begin{pgfscope}%
\pgfsys@transformshift{4.682603in}{0.807389in}%
\pgfsys@useobject{currentmarker}{}%
\end{pgfscope}%
\end{pgfscope}%
\begin{pgfscope}%
\pgfsetbuttcap%
\pgfsetroundjoin%
\definecolor{currentfill}{rgb}{0.000000,0.000000,0.000000}%
\pgfsetfillcolor{currentfill}%
\pgfsetlinewidth{0.602250pt}%
\definecolor{currentstroke}{rgb}{0.000000,0.000000,0.000000}%
\pgfsetstrokecolor{currentstroke}%
\pgfsetdash{}{0pt}%
\pgfsys@defobject{currentmarker}{\pgfqpoint{0.000000in}{-0.027778in}}{\pgfqpoint{0.000000in}{0.000000in}}{%
\pgfpathmoveto{\pgfqpoint{0.000000in}{0.000000in}}%
\pgfpathlineto{\pgfqpoint{0.000000in}{-0.027778in}}%
\pgfusepath{stroke,fill}%
}%
\begin{pgfscope}%
\pgfsys@transformshift{5.353715in}{0.807389in}%
\pgfsys@useobject{currentmarker}{}%
\end{pgfscope}%
\end{pgfscope}%
\begin{pgfscope}%
\pgfsetbuttcap%
\pgfsetroundjoin%
\definecolor{currentfill}{rgb}{0.000000,0.000000,0.000000}%
\pgfsetfillcolor{currentfill}%
\pgfsetlinewidth{0.602250pt}%
\definecolor{currentstroke}{rgb}{0.000000,0.000000,0.000000}%
\pgfsetstrokecolor{currentstroke}%
\pgfsetdash{}{0pt}%
\pgfsys@defobject{currentmarker}{\pgfqpoint{0.000000in}{-0.027778in}}{\pgfqpoint{0.000000in}{0.000000in}}{%
\pgfpathmoveto{\pgfqpoint{0.000000in}{0.000000in}}%
\pgfpathlineto{\pgfqpoint{0.000000in}{-0.027778in}}%
\pgfusepath{stroke,fill}%
}%
\begin{pgfscope}%
\pgfsys@transformshift{5.694492in}{0.807389in}%
\pgfsys@useobject{currentmarker}{}%
\end{pgfscope}%
\end{pgfscope}%
\begin{pgfscope}%
\pgfsetbuttcap%
\pgfsetroundjoin%
\definecolor{currentfill}{rgb}{0.000000,0.000000,0.000000}%
\pgfsetfillcolor{currentfill}%
\pgfsetlinewidth{0.602250pt}%
\definecolor{currentstroke}{rgb}{0.000000,0.000000,0.000000}%
\pgfsetstrokecolor{currentstroke}%
\pgfsetdash{}{0pt}%
\pgfsys@defobject{currentmarker}{\pgfqpoint{0.000000in}{-0.027778in}}{\pgfqpoint{0.000000in}{0.000000in}}{%
\pgfpathmoveto{\pgfqpoint{0.000000in}{0.000000in}}%
\pgfpathlineto{\pgfqpoint{0.000000in}{-0.027778in}}%
\pgfusepath{stroke,fill}%
}%
\begin{pgfscope}%
\pgfsys@transformshift{5.936276in}{0.807389in}%
\pgfsys@useobject{currentmarker}{}%
\end{pgfscope}%
\end{pgfscope}%
\begin{pgfscope}%
\pgfsetbuttcap%
\pgfsetroundjoin%
\definecolor{currentfill}{rgb}{0.000000,0.000000,0.000000}%
\pgfsetfillcolor{currentfill}%
\pgfsetlinewidth{0.602250pt}%
\definecolor{currentstroke}{rgb}{0.000000,0.000000,0.000000}%
\pgfsetstrokecolor{currentstroke}%
\pgfsetdash{}{0pt}%
\pgfsys@defobject{currentmarker}{\pgfqpoint{0.000000in}{-0.027778in}}{\pgfqpoint{0.000000in}{0.000000in}}{%
\pgfpathmoveto{\pgfqpoint{0.000000in}{0.000000in}}%
\pgfpathlineto{\pgfqpoint{0.000000in}{-0.027778in}}%
\pgfusepath{stroke,fill}%
}%
\begin{pgfscope}%
\pgfsys@transformshift{6.123819in}{0.807389in}%
\pgfsys@useobject{currentmarker}{}%
\end{pgfscope}%
\end{pgfscope}%
\begin{pgfscope}%
\pgfsetbuttcap%
\pgfsetroundjoin%
\definecolor{currentfill}{rgb}{0.000000,0.000000,0.000000}%
\pgfsetfillcolor{currentfill}%
\pgfsetlinewidth{0.602250pt}%
\definecolor{currentstroke}{rgb}{0.000000,0.000000,0.000000}%
\pgfsetstrokecolor{currentstroke}%
\pgfsetdash{}{0pt}%
\pgfsys@defobject{currentmarker}{\pgfqpoint{0.000000in}{-0.027778in}}{\pgfqpoint{0.000000in}{0.000000in}}{%
\pgfpathmoveto{\pgfqpoint{0.000000in}{0.000000in}}%
\pgfpathlineto{\pgfqpoint{0.000000in}{-0.027778in}}%
\pgfusepath{stroke,fill}%
}%
\begin{pgfscope}%
\pgfsys@transformshift{6.277052in}{0.807389in}%
\pgfsys@useobject{currentmarker}{}%
\end{pgfscope}%
\end{pgfscope}%
\begin{pgfscope}%
\pgfsetbuttcap%
\pgfsetroundjoin%
\definecolor{currentfill}{rgb}{0.000000,0.000000,0.000000}%
\pgfsetfillcolor{currentfill}%
\pgfsetlinewidth{0.602250pt}%
\definecolor{currentstroke}{rgb}{0.000000,0.000000,0.000000}%
\pgfsetstrokecolor{currentstroke}%
\pgfsetdash{}{0pt}%
\pgfsys@defobject{currentmarker}{\pgfqpoint{0.000000in}{-0.027778in}}{\pgfqpoint{0.000000in}{0.000000in}}{%
\pgfpathmoveto{\pgfqpoint{0.000000in}{0.000000in}}%
\pgfpathlineto{\pgfqpoint{0.000000in}{-0.027778in}}%
\pgfusepath{stroke,fill}%
}%
\begin{pgfscope}%
\pgfsys@transformshift{6.406610in}{0.807389in}%
\pgfsys@useobject{currentmarker}{}%
\end{pgfscope}%
\end{pgfscope}%
\begin{pgfscope}%
\pgfsetbuttcap%
\pgfsetroundjoin%
\definecolor{currentfill}{rgb}{0.000000,0.000000,0.000000}%
\pgfsetfillcolor{currentfill}%
\pgfsetlinewidth{0.602250pt}%
\definecolor{currentstroke}{rgb}{0.000000,0.000000,0.000000}%
\pgfsetstrokecolor{currentstroke}%
\pgfsetdash{}{0pt}%
\pgfsys@defobject{currentmarker}{\pgfqpoint{0.000000in}{-0.027778in}}{\pgfqpoint{0.000000in}{0.000000in}}{%
\pgfpathmoveto{\pgfqpoint{0.000000in}{0.000000in}}%
\pgfpathlineto{\pgfqpoint{0.000000in}{-0.027778in}}%
\pgfusepath{stroke,fill}%
}%
\begin{pgfscope}%
\pgfsys@transformshift{6.518837in}{0.807389in}%
\pgfsys@useobject{currentmarker}{}%
\end{pgfscope}%
\end{pgfscope}%
\begin{pgfscope}%
\pgfsetbuttcap%
\pgfsetroundjoin%
\definecolor{currentfill}{rgb}{0.000000,0.000000,0.000000}%
\pgfsetfillcolor{currentfill}%
\pgfsetlinewidth{0.602250pt}%
\definecolor{currentstroke}{rgb}{0.000000,0.000000,0.000000}%
\pgfsetstrokecolor{currentstroke}%
\pgfsetdash{}{0pt}%
\pgfsys@defobject{currentmarker}{\pgfqpoint{0.000000in}{-0.027778in}}{\pgfqpoint{0.000000in}{0.000000in}}{%
\pgfpathmoveto{\pgfqpoint{0.000000in}{0.000000in}}%
\pgfpathlineto{\pgfqpoint{0.000000in}{-0.027778in}}%
\pgfusepath{stroke,fill}%
}%
\begin{pgfscope}%
\pgfsys@transformshift{6.617829in}{0.807389in}%
\pgfsys@useobject{currentmarker}{}%
\end{pgfscope}%
\end{pgfscope}%
\begin{pgfscope}%
\definecolor{textcolor}{rgb}{0.000000,0.000000,0.000000}%
\pgfsetstrokecolor{textcolor}%
\pgfsetfillcolor{textcolor}%
\pgftext[x=5.447487in,y=0.561000in,,top]{\color{textcolor}{\rmfamily\fontsize{8.000000}{9.600000}\selectfont\catcode`\^=\active\def^{\ifmmode\sp\else\^{}\fi}\catcode`\%=\active\def%{\%}Angular Scale $\alpha$ (Degrees)}}%
\end{pgfscope}%
\begin{pgfscope}%
\pgfpathrectangle{\pgfqpoint{4.062704in}{0.807389in}}{\pgfqpoint{2.769565in}{2.016000in}}%
\pgfusepath{clip}%
\pgfsetrectcap%
\pgfsetroundjoin%
\pgfsetlinewidth{0.803000pt}%
\definecolor{currentstroke}{rgb}{0.690196,0.690196,0.690196}%
\pgfsetstrokecolor{currentstroke}%
\pgfsetstrokeopacity{0.300000}%
\pgfsetdash{}{0pt}%
\pgfpathmoveto{\pgfqpoint{4.062704in}{0.878986in}}%
\pgfpathlineto{\pgfqpoint{6.832269in}{0.878986in}}%
\pgfusepath{stroke}%
\end{pgfscope}%
\begin{pgfscope}%
\pgfsetbuttcap%
\pgfsetroundjoin%
\definecolor{currentfill}{rgb}{0.000000,0.000000,0.000000}%
\pgfsetfillcolor{currentfill}%
\pgfsetlinewidth{0.803000pt}%
\definecolor{currentstroke}{rgb}{0.000000,0.000000,0.000000}%
\pgfsetstrokecolor{currentstroke}%
\pgfsetdash{}{0pt}%
\pgfsys@defobject{currentmarker}{\pgfqpoint{-0.048611in}{0.000000in}}{\pgfqpoint{-0.000000in}{0.000000in}}{%
\pgfpathmoveto{\pgfqpoint{-0.000000in}{0.000000in}}%
\pgfpathlineto{\pgfqpoint{-0.048611in}{0.000000in}}%
\pgfusepath{stroke,fill}%
}%
\begin{pgfscope}%
\pgfsys@transformshift{4.062704in}{0.878986in}%
\pgfsys@useobject{currentmarker}{}%
\end{pgfscope}%
\end{pgfscope}%
\begin{pgfscope}%
\definecolor{textcolor}{rgb}{0.000000,0.000000,0.000000}%
\pgfsetstrokecolor{textcolor}%
\pgfsetfillcolor{textcolor}%
\pgftext[x=3.823313in, y=0.845229in, left, base]{\color{textcolor}{\rmfamily\fontsize{7.000000}{8.400000}\selectfont\catcode`\^=\active\def^{\ifmmode\sp\else\^{}\fi}\catcode`\%=\active\def%{\%}0.0}}%
\end{pgfscope}%
\begin{pgfscope}%
\pgfpathrectangle{\pgfqpoint{4.062704in}{0.807389in}}{\pgfqpoint{2.769565in}{2.016000in}}%
\pgfusepath{clip}%
\pgfsetrectcap%
\pgfsetroundjoin%
\pgfsetlinewidth{0.803000pt}%
\definecolor{currentstroke}{rgb}{0.690196,0.690196,0.690196}%
\pgfsetstrokecolor{currentstroke}%
\pgfsetstrokeopacity{0.300000}%
\pgfsetdash{}{0pt}%
\pgfpathmoveto{\pgfqpoint{4.062704in}{1.249567in}}%
\pgfpathlineto{\pgfqpoint{6.832269in}{1.249567in}}%
\pgfusepath{stroke}%
\end{pgfscope}%
\begin{pgfscope}%
\pgfsetbuttcap%
\pgfsetroundjoin%
\definecolor{currentfill}{rgb}{0.000000,0.000000,0.000000}%
\pgfsetfillcolor{currentfill}%
\pgfsetlinewidth{0.803000pt}%
\definecolor{currentstroke}{rgb}{0.000000,0.000000,0.000000}%
\pgfsetstrokecolor{currentstroke}%
\pgfsetdash{}{0pt}%
\pgfsys@defobject{currentmarker}{\pgfqpoint{-0.048611in}{0.000000in}}{\pgfqpoint{-0.000000in}{0.000000in}}{%
\pgfpathmoveto{\pgfqpoint{-0.000000in}{0.000000in}}%
\pgfpathlineto{\pgfqpoint{-0.048611in}{0.000000in}}%
\pgfusepath{stroke,fill}%
}%
\begin{pgfscope}%
\pgfsys@transformshift{4.062704in}{1.249567in}%
\pgfsys@useobject{currentmarker}{}%
\end{pgfscope}%
\end{pgfscope}%
\begin{pgfscope}%
\definecolor{textcolor}{rgb}{0.000000,0.000000,0.000000}%
\pgfsetstrokecolor{textcolor}%
\pgfsetfillcolor{textcolor}%
\pgftext[x=3.823313in, y=1.215810in, left, base]{\color{textcolor}{\rmfamily\fontsize{7.000000}{8.400000}\selectfont\catcode`\^=\active\def^{\ifmmode\sp\else\^{}\fi}\catcode`\%=\active\def%{\%}0.2}}%
\end{pgfscope}%
\begin{pgfscope}%
\pgfpathrectangle{\pgfqpoint{4.062704in}{0.807389in}}{\pgfqpoint{2.769565in}{2.016000in}}%
\pgfusepath{clip}%
\pgfsetrectcap%
\pgfsetroundjoin%
\pgfsetlinewidth{0.803000pt}%
\definecolor{currentstroke}{rgb}{0.690196,0.690196,0.690196}%
\pgfsetstrokecolor{currentstroke}%
\pgfsetstrokeopacity{0.300000}%
\pgfsetdash{}{0pt}%
\pgfpathmoveto{\pgfqpoint{4.062704in}{1.620149in}}%
\pgfpathlineto{\pgfqpoint{6.832269in}{1.620149in}}%
\pgfusepath{stroke}%
\end{pgfscope}%
\begin{pgfscope}%
\pgfsetbuttcap%
\pgfsetroundjoin%
\definecolor{currentfill}{rgb}{0.000000,0.000000,0.000000}%
\pgfsetfillcolor{currentfill}%
\pgfsetlinewidth{0.803000pt}%
\definecolor{currentstroke}{rgb}{0.000000,0.000000,0.000000}%
\pgfsetstrokecolor{currentstroke}%
\pgfsetdash{}{0pt}%
\pgfsys@defobject{currentmarker}{\pgfqpoint{-0.048611in}{0.000000in}}{\pgfqpoint{-0.000000in}{0.000000in}}{%
\pgfpathmoveto{\pgfqpoint{-0.000000in}{0.000000in}}%
\pgfpathlineto{\pgfqpoint{-0.048611in}{0.000000in}}%
\pgfusepath{stroke,fill}%
}%
\begin{pgfscope}%
\pgfsys@transformshift{4.062704in}{1.620149in}%
\pgfsys@useobject{currentmarker}{}%
\end{pgfscope}%
\end{pgfscope}%
\begin{pgfscope}%
\definecolor{textcolor}{rgb}{0.000000,0.000000,0.000000}%
\pgfsetstrokecolor{textcolor}%
\pgfsetfillcolor{textcolor}%
\pgftext[x=3.823313in, y=1.586391in, left, base]{\color{textcolor}{\rmfamily\fontsize{7.000000}{8.400000}\selectfont\catcode`\^=\active\def^{\ifmmode\sp\else\^{}\fi}\catcode`\%=\active\def%{\%}0.4}}%
\end{pgfscope}%
\begin{pgfscope}%
\pgfpathrectangle{\pgfqpoint{4.062704in}{0.807389in}}{\pgfqpoint{2.769565in}{2.016000in}}%
\pgfusepath{clip}%
\pgfsetrectcap%
\pgfsetroundjoin%
\pgfsetlinewidth{0.803000pt}%
\definecolor{currentstroke}{rgb}{0.690196,0.690196,0.690196}%
\pgfsetstrokecolor{currentstroke}%
\pgfsetstrokeopacity{0.300000}%
\pgfsetdash{}{0pt}%
\pgfpathmoveto{\pgfqpoint{4.062704in}{1.990730in}}%
\pgfpathlineto{\pgfqpoint{6.832269in}{1.990730in}}%
\pgfusepath{stroke}%
\end{pgfscope}%
\begin{pgfscope}%
\pgfsetbuttcap%
\pgfsetroundjoin%
\definecolor{currentfill}{rgb}{0.000000,0.000000,0.000000}%
\pgfsetfillcolor{currentfill}%
\pgfsetlinewidth{0.803000pt}%
\definecolor{currentstroke}{rgb}{0.000000,0.000000,0.000000}%
\pgfsetstrokecolor{currentstroke}%
\pgfsetdash{}{0pt}%
\pgfsys@defobject{currentmarker}{\pgfqpoint{-0.048611in}{0.000000in}}{\pgfqpoint{-0.000000in}{0.000000in}}{%
\pgfpathmoveto{\pgfqpoint{-0.000000in}{0.000000in}}%
\pgfpathlineto{\pgfqpoint{-0.048611in}{0.000000in}}%
\pgfusepath{stroke,fill}%
}%
\begin{pgfscope}%
\pgfsys@transformshift{4.062704in}{1.990730in}%
\pgfsys@useobject{currentmarker}{}%
\end{pgfscope}%
\end{pgfscope}%
\begin{pgfscope}%
\definecolor{textcolor}{rgb}{0.000000,0.000000,0.000000}%
\pgfsetstrokecolor{textcolor}%
\pgfsetfillcolor{textcolor}%
\pgftext[x=3.823313in, y=1.956972in, left, base]{\color{textcolor}{\rmfamily\fontsize{7.000000}{8.400000}\selectfont\catcode`\^=\active\def^{\ifmmode\sp\else\^{}\fi}\catcode`\%=\active\def%{\%}0.6}}%
\end{pgfscope}%
\begin{pgfscope}%
\pgfpathrectangle{\pgfqpoint{4.062704in}{0.807389in}}{\pgfqpoint{2.769565in}{2.016000in}}%
\pgfusepath{clip}%
\pgfsetrectcap%
\pgfsetroundjoin%
\pgfsetlinewidth{0.803000pt}%
\definecolor{currentstroke}{rgb}{0.690196,0.690196,0.690196}%
\pgfsetstrokecolor{currentstroke}%
\pgfsetstrokeopacity{0.300000}%
\pgfsetdash{}{0pt}%
\pgfpathmoveto{\pgfqpoint{4.062704in}{2.361311in}}%
\pgfpathlineto{\pgfqpoint{6.832269in}{2.361311in}}%
\pgfusepath{stroke}%
\end{pgfscope}%
\begin{pgfscope}%
\pgfsetbuttcap%
\pgfsetroundjoin%
\definecolor{currentfill}{rgb}{0.000000,0.000000,0.000000}%
\pgfsetfillcolor{currentfill}%
\pgfsetlinewidth{0.803000pt}%
\definecolor{currentstroke}{rgb}{0.000000,0.000000,0.000000}%
\pgfsetstrokecolor{currentstroke}%
\pgfsetdash{}{0pt}%
\pgfsys@defobject{currentmarker}{\pgfqpoint{-0.048611in}{0.000000in}}{\pgfqpoint{-0.000000in}{0.000000in}}{%
\pgfpathmoveto{\pgfqpoint{-0.000000in}{0.000000in}}%
\pgfpathlineto{\pgfqpoint{-0.048611in}{0.000000in}}%
\pgfusepath{stroke,fill}%
}%
\begin{pgfscope}%
\pgfsys@transformshift{4.062704in}{2.361311in}%
\pgfsys@useobject{currentmarker}{}%
\end{pgfscope}%
\end{pgfscope}%
\begin{pgfscope}%
\definecolor{textcolor}{rgb}{0.000000,0.000000,0.000000}%
\pgfsetstrokecolor{textcolor}%
\pgfsetfillcolor{textcolor}%
\pgftext[x=3.823313in, y=2.327553in, left, base]{\color{textcolor}{\rmfamily\fontsize{7.000000}{8.400000}\selectfont\catcode`\^=\active\def^{\ifmmode\sp\else\^{}\fi}\catcode`\%=\active\def%{\%}0.8}}%
\end{pgfscope}%
\begin{pgfscope}%
\pgfpathrectangle{\pgfqpoint{4.062704in}{0.807389in}}{\pgfqpoint{2.769565in}{2.016000in}}%
\pgfusepath{clip}%
\pgfsetrectcap%
\pgfsetroundjoin%
\pgfsetlinewidth{0.803000pt}%
\definecolor{currentstroke}{rgb}{0.690196,0.690196,0.690196}%
\pgfsetstrokecolor{currentstroke}%
\pgfsetstrokeopacity{0.300000}%
\pgfsetdash{}{0pt}%
\pgfpathmoveto{\pgfqpoint{4.062704in}{2.731892in}}%
\pgfpathlineto{\pgfqpoint{6.832269in}{2.731892in}}%
\pgfusepath{stroke}%
\end{pgfscope}%
\begin{pgfscope}%
\pgfsetbuttcap%
\pgfsetroundjoin%
\definecolor{currentfill}{rgb}{0.000000,0.000000,0.000000}%
\pgfsetfillcolor{currentfill}%
\pgfsetlinewidth{0.803000pt}%
\definecolor{currentstroke}{rgb}{0.000000,0.000000,0.000000}%
\pgfsetstrokecolor{currentstroke}%
\pgfsetdash{}{0pt}%
\pgfsys@defobject{currentmarker}{\pgfqpoint{-0.048611in}{0.000000in}}{\pgfqpoint{-0.000000in}{0.000000in}}{%
\pgfpathmoveto{\pgfqpoint{-0.000000in}{0.000000in}}%
\pgfpathlineto{\pgfqpoint{-0.048611in}{0.000000in}}%
\pgfusepath{stroke,fill}%
}%
\begin{pgfscope}%
\pgfsys@transformshift{4.062704in}{2.731892in}%
\pgfsys@useobject{currentmarker}{}%
\end{pgfscope}%
\end{pgfscope}%
\begin{pgfscope}%
\definecolor{textcolor}{rgb}{0.000000,0.000000,0.000000}%
\pgfsetstrokecolor{textcolor}%
\pgfsetfillcolor{textcolor}%
\pgftext[x=3.823313in, y=2.698134in, left, base]{\color{textcolor}{\rmfamily\fontsize{7.000000}{8.400000}\selectfont\catcode`\^=\active\def^{\ifmmode\sp\else\^{}\fi}\catcode`\%=\active\def%{\%}1.0}}%
\end{pgfscope}%
\begin{pgfscope}%
\definecolor{textcolor}{rgb}{0.000000,0.000000,0.000000}%
\pgfsetstrokecolor{textcolor}%
\pgfsetfillcolor{textcolor}%
\pgftext[x=3.767757in,y=1.815389in,,bottom,rotate=90.000000]{\color{textcolor}{\rmfamily\fontsize{8.000000}{9.600000}\selectfont\catcode`\^=\active\def^{\ifmmode\sp\else\^{}\fi}\catcode`\%=\active\def%{\%}Effective Occupied Volume $\widehat U_2(t)$}}%
\end{pgfscope}%
\begin{pgfscope}%
\pgfpathrectangle{\pgfqpoint{4.062704in}{0.807389in}}{\pgfqpoint{2.769565in}{2.016000in}}%
\pgfusepath{clip}%
\pgfsetrectcap%
\pgfsetroundjoin%
\pgfsetlinewidth{1.806750pt}%
\definecolor{currentstroke}{rgb}{0.121569,0.466667,0.705882}%
\pgfsetstrokecolor{currentstroke}%
\pgfsetdash{}{0pt}%
\pgfpathmoveto{\pgfqpoint{4.188593in}{0.900033in}}%
\pgfpathlineto{\pgfqpoint{4.235219in}{0.900486in}}%
\pgfpathlineto{\pgfqpoint{4.281845in}{0.900982in}}%
\pgfpathlineto{\pgfqpoint{4.328470in}{0.901528in}}%
\pgfpathlineto{\pgfqpoint{4.375096in}{0.902130in}}%
\pgfpathlineto{\pgfqpoint{4.421722in}{0.902796in}}%
\pgfpathlineto{\pgfqpoint{4.468347in}{0.903536in}}%
\pgfpathlineto{\pgfqpoint{4.514973in}{0.904363in}}%
\pgfpathlineto{\pgfqpoint{4.561599in}{0.905281in}}%
\pgfpathlineto{\pgfqpoint{4.608224in}{0.906301in}}%
\pgfpathlineto{\pgfqpoint{4.654850in}{0.907436in}}%
\pgfpathlineto{\pgfqpoint{4.701476in}{0.908682in}}%
\pgfpathlineto{\pgfqpoint{4.748101in}{0.910078in}}%
\pgfpathlineto{\pgfqpoint{4.794727in}{0.911651in}}%
\pgfpathlineto{\pgfqpoint{4.841353in}{0.913401in}}%
\pgfpathlineto{\pgfqpoint{4.887978in}{0.915356in}}%
\pgfpathlineto{\pgfqpoint{4.934604in}{0.917540in}}%
\pgfpathlineto{\pgfqpoint{4.981230in}{0.919973in}}%
\pgfpathlineto{\pgfqpoint{5.027855in}{0.922684in}}%
\pgfpathlineto{\pgfqpoint{5.074481in}{0.925705in}}%
\pgfpathlineto{\pgfqpoint{5.121107in}{0.929073in}}%
\pgfpathlineto{\pgfqpoint{5.167732in}{0.932826in}}%
\pgfpathlineto{\pgfqpoint{5.214358in}{0.936981in}}%
\pgfpathlineto{\pgfqpoint{5.260984in}{0.941619in}}%
\pgfpathlineto{\pgfqpoint{5.307610in}{0.946797in}}%
\pgfpathlineto{\pgfqpoint{5.354235in}{0.952583in}}%
\pgfpathlineto{\pgfqpoint{5.400861in}{0.959032in}}%
\pgfpathlineto{\pgfqpoint{5.447487in}{0.966219in}}%
\pgfpathlineto{\pgfqpoint{5.494112in}{0.974225in}}%
\pgfpathlineto{\pgfqpoint{5.540738in}{0.983142in}}%
\pgfpathlineto{\pgfqpoint{5.587364in}{0.993069in}}%
\pgfpathlineto{\pgfqpoint{5.633989in}{1.004116in}}%
\pgfpathlineto{\pgfqpoint{5.680615in}{1.016405in}}%
\pgfpathlineto{\pgfqpoint{5.727241in}{1.030066in}}%
\pgfpathlineto{\pgfqpoint{5.773866in}{1.045245in}}%
\pgfpathlineto{\pgfqpoint{5.820492in}{1.062099in}}%
\pgfpathlineto{\pgfqpoint{5.867118in}{1.080797in}}%
\pgfpathlineto{\pgfqpoint{5.913743in}{1.101526in}}%
\pgfpathlineto{\pgfqpoint{5.960369in}{1.124482in}}%
\pgfpathlineto{\pgfqpoint{6.006995in}{1.149878in}}%
\pgfpathlineto{\pgfqpoint{6.053620in}{1.177939in}}%
\pgfpathlineto{\pgfqpoint{6.100246in}{1.208903in}}%
\pgfpathlineto{\pgfqpoint{6.146872in}{1.243015in}}%
\pgfpathlineto{\pgfqpoint{6.193497in}{1.280531in}}%
\pgfpathlineto{\pgfqpoint{6.240123in}{1.321707in}}%
\pgfpathlineto{\pgfqpoint{6.286749in}{1.366799in}}%
\pgfpathlineto{\pgfqpoint{6.333374in}{1.416051in}}%
\pgfpathlineto{\pgfqpoint{6.380000in}{1.469688in}}%
\pgfpathlineto{\pgfqpoint{6.426626in}{1.527902in}}%
\pgfpathlineto{\pgfqpoint{6.473251in}{1.590836in}}%
\pgfpathlineto{\pgfqpoint{6.519877in}{1.658564in}}%
\pgfpathlineto{\pgfqpoint{6.566503in}{1.731062in}}%
\pgfpathlineto{\pgfqpoint{6.613128in}{1.808177in}}%
\pgfpathlineto{\pgfqpoint{6.659754in}{1.889582in}}%
\pgfpathlineto{\pgfqpoint{6.706380in}{1.974717in}}%
\pgfusepath{stroke}%
\end{pgfscope}%
\begin{pgfscope}%
\pgfpathrectangle{\pgfqpoint{4.062704in}{0.807389in}}{\pgfqpoint{2.769565in}{2.016000in}}%
\pgfusepath{clip}%
\pgfsetrectcap%
\pgfsetroundjoin%
\pgfsetlinewidth{1.806750pt}%
\definecolor{currentstroke}{rgb}{1.000000,0.498039,0.054902}%
\pgfsetstrokecolor{currentstroke}%
\pgfsetdash{}{0pt}%
\pgfpathmoveto{\pgfqpoint{4.188593in}{0.921160in}}%
\pgfpathlineto{\pgfqpoint{4.235219in}{0.922187in}}%
\pgfpathlineto{\pgfqpoint{4.281845in}{0.923300in}}%
\pgfpathlineto{\pgfqpoint{4.328470in}{0.924492in}}%
\pgfpathlineto{\pgfqpoint{4.375096in}{0.925778in}}%
\pgfpathlineto{\pgfqpoint{4.421722in}{0.927173in}}%
\pgfpathlineto{\pgfqpoint{4.468347in}{0.928696in}}%
\pgfpathlineto{\pgfqpoint{4.514973in}{0.930417in}}%
\pgfpathlineto{\pgfqpoint{4.561599in}{0.932314in}}%
\pgfpathlineto{\pgfqpoint{4.608224in}{0.934405in}}%
\pgfpathlineto{\pgfqpoint{4.654850in}{0.936714in}}%
\pgfpathlineto{\pgfqpoint{4.701476in}{0.939271in}}%
\pgfpathlineto{\pgfqpoint{4.748101in}{0.942107in}}%
\pgfpathlineto{\pgfqpoint{4.794727in}{0.945257in}}%
\pgfpathlineto{\pgfqpoint{4.841353in}{0.948760in}}%
\pgfpathlineto{\pgfqpoint{4.887978in}{0.952660in}}%
\pgfpathlineto{\pgfqpoint{4.934604in}{0.957003in}}%
\pgfpathlineto{\pgfqpoint{4.981230in}{0.961844in}}%
\pgfpathlineto{\pgfqpoint{5.027855in}{0.967236in}}%
\pgfpathlineto{\pgfqpoint{5.074481in}{0.973228in}}%
\pgfpathlineto{\pgfqpoint{5.121107in}{0.979913in}}%
\pgfpathlineto{\pgfqpoint{5.167732in}{0.987373in}}%
\pgfpathlineto{\pgfqpoint{5.214358in}{0.995695in}}%
\pgfpathlineto{\pgfqpoint{5.260984in}{1.005016in}}%
\pgfpathlineto{\pgfqpoint{5.307610in}{1.015426in}}%
\pgfpathlineto{\pgfqpoint{5.354235in}{1.027028in}}%
\pgfpathlineto{\pgfqpoint{5.400861in}{1.039954in}}%
\pgfpathlineto{\pgfqpoint{5.447487in}{1.054353in}}%
\pgfpathlineto{\pgfqpoint{5.494112in}{1.070390in}}%
\pgfpathlineto{\pgfqpoint{5.540738in}{1.088246in}}%
\pgfpathlineto{\pgfqpoint{5.587364in}{1.108120in}}%
\pgfpathlineto{\pgfqpoint{5.633989in}{1.130233in}}%
\pgfpathlineto{\pgfqpoint{5.680615in}{1.154826in}}%
\pgfpathlineto{\pgfqpoint{5.727241in}{1.182162in}}%
\pgfpathlineto{\pgfqpoint{5.773866in}{1.212531in}}%
\pgfpathlineto{\pgfqpoint{5.820492in}{1.246247in}}%
\pgfpathlineto{\pgfqpoint{5.867118in}{1.283652in}}%
\pgfpathlineto{\pgfqpoint{5.913743in}{1.325113in}}%
\pgfpathlineto{\pgfqpoint{5.960369in}{1.371027in}}%
\pgfpathlineto{\pgfqpoint{6.006995in}{1.421819in}}%
\pgfpathlineto{\pgfqpoint{6.053620in}{1.477935in}}%
\pgfpathlineto{\pgfqpoint{6.100246in}{1.539841in}}%
\pgfpathlineto{\pgfqpoint{6.146872in}{1.608004in}}%
\pgfpathlineto{\pgfqpoint{6.193497in}{1.682850in}}%
\pgfpathlineto{\pgfqpoint{6.240123in}{1.764691in}}%
\pgfpathlineto{\pgfqpoint{6.286749in}{1.853581in}}%
\pgfpathlineto{\pgfqpoint{6.333374in}{1.949111in}}%
\pgfpathlineto{\pgfqpoint{6.380000in}{2.050138in}}%
\pgfpathlineto{\pgfqpoint{6.426626in}{2.154538in}}%
\pgfpathlineto{\pgfqpoint{6.473251in}{2.259094in}}%
\pgfpathlineto{\pgfqpoint{6.519877in}{2.359679in}}%
\pgfpathlineto{\pgfqpoint{6.566503in}{2.451824in}}%
\pgfpathlineto{\pgfqpoint{6.613128in}{2.531583in}}%
\pgfpathlineto{\pgfqpoint{6.659754in}{2.596388in}}%
\pgfpathlineto{\pgfqpoint{6.706380in}{2.645582in}}%
\pgfusepath{stroke}%
\end{pgfscope}%
\begin{pgfscope}%
\pgfpathrectangle{\pgfqpoint{4.062704in}{0.807389in}}{\pgfqpoint{2.769565in}{2.016000in}}%
\pgfusepath{clip}%
\pgfsetrectcap%
\pgfsetroundjoin%
\pgfsetlinewidth{1.806750pt}%
\definecolor{currentstroke}{rgb}{0.172549,0.627451,0.172549}%
\pgfsetstrokecolor{currentstroke}%
\pgfsetdash{}{0pt}%
\pgfpathmoveto{\pgfqpoint{4.188593in}{1.034430in}}%
\pgfpathlineto{\pgfqpoint{4.235219in}{1.042177in}}%
\pgfpathlineto{\pgfqpoint{4.281845in}{1.050296in}}%
\pgfpathlineto{\pgfqpoint{4.328470in}{1.058824in}}%
\pgfpathlineto{\pgfqpoint{4.375096in}{1.067800in}}%
\pgfpathlineto{\pgfqpoint{4.421722in}{1.077284in}}%
\pgfpathlineto{\pgfqpoint{4.468347in}{1.087291in}}%
\pgfpathlineto{\pgfqpoint{4.514973in}{1.097807in}}%
\pgfpathlineto{\pgfqpoint{4.561599in}{1.108887in}}%
\pgfpathlineto{\pgfqpoint{4.608224in}{1.120511in}}%
\pgfpathlineto{\pgfqpoint{4.654850in}{1.132714in}}%
\pgfpathlineto{\pgfqpoint{4.701476in}{1.145649in}}%
\pgfpathlineto{\pgfqpoint{4.748101in}{1.159678in}}%
\pgfpathlineto{\pgfqpoint{4.794727in}{1.174474in}}%
\pgfpathlineto{\pgfqpoint{4.841353in}{1.190178in}}%
\pgfpathlineto{\pgfqpoint{4.887978in}{1.206835in}}%
\pgfpathlineto{\pgfqpoint{4.934604in}{1.224493in}}%
\pgfpathlineto{\pgfqpoint{4.981230in}{1.243142in}}%
\pgfpathlineto{\pgfqpoint{5.027855in}{1.262845in}}%
\pgfpathlineto{\pgfqpoint{5.074481in}{1.283803in}}%
\pgfpathlineto{\pgfqpoint{5.121107in}{1.305964in}}%
\pgfpathlineto{\pgfqpoint{5.167732in}{1.329078in}}%
\pgfpathlineto{\pgfqpoint{5.214358in}{1.353395in}}%
\pgfpathlineto{\pgfqpoint{5.260984in}{1.379109in}}%
\pgfpathlineto{\pgfqpoint{5.307610in}{1.406283in}}%
\pgfpathlineto{\pgfqpoint{5.354235in}{1.434980in}}%
\pgfpathlineto{\pgfqpoint{5.400861in}{1.465816in}}%
\pgfpathlineto{\pgfqpoint{5.447487in}{1.498313in}}%
\pgfpathlineto{\pgfqpoint{5.494112in}{1.532481in}}%
\pgfpathlineto{\pgfqpoint{5.540738in}{1.568202in}}%
\pgfpathlineto{\pgfqpoint{5.587364in}{1.605356in}}%
\pgfpathlineto{\pgfqpoint{5.633989in}{1.644908in}}%
\pgfpathlineto{\pgfqpoint{5.680615in}{1.686611in}}%
\pgfpathlineto{\pgfqpoint{5.727241in}{1.729974in}}%
\pgfpathlineto{\pgfqpoint{5.773866in}{1.774856in}}%
\pgfpathlineto{\pgfqpoint{5.820492in}{1.821303in}}%
\pgfpathlineto{\pgfqpoint{5.867118in}{1.869639in}}%
\pgfpathlineto{\pgfqpoint{5.913743in}{1.919814in}}%
\pgfpathlineto{\pgfqpoint{5.960369in}{1.971748in}}%
\pgfpathlineto{\pgfqpoint{6.006995in}{2.025015in}}%
\pgfpathlineto{\pgfqpoint{6.053620in}{2.079659in}}%
\pgfpathlineto{\pgfqpoint{6.100246in}{2.135630in}}%
\pgfpathlineto{\pgfqpoint{6.146872in}{2.192762in}}%
\pgfpathlineto{\pgfqpoint{6.193497in}{2.250410in}}%
\pgfpathlineto{\pgfqpoint{6.240123in}{2.308044in}}%
\pgfpathlineto{\pgfqpoint{6.286749in}{2.364983in}}%
\pgfpathlineto{\pgfqpoint{6.333374in}{2.420389in}}%
\pgfpathlineto{\pgfqpoint{6.380000in}{2.473277in}}%
\pgfpathlineto{\pgfqpoint{6.426626in}{2.522581in}}%
\pgfpathlineto{\pgfqpoint{6.473251in}{2.567251in}}%
\pgfpathlineto{\pgfqpoint{6.519877in}{2.606398in}}%
\pgfpathlineto{\pgfqpoint{6.566503in}{2.639431in}}%
\pgfpathlineto{\pgfqpoint{6.613128in}{2.666163in}}%
\pgfpathlineto{\pgfqpoint{6.659754in}{2.687023in}}%
\pgfpathlineto{\pgfqpoint{6.706380in}{2.702416in}}%
\pgfusepath{stroke}%
\end{pgfscope}%
\begin{pgfscope}%
\pgfpathrectangle{\pgfqpoint{4.062704in}{0.807389in}}{\pgfqpoint{2.769565in}{2.016000in}}%
\pgfusepath{clip}%
\pgfsetrectcap%
\pgfsetroundjoin%
\pgfsetlinewidth{1.806750pt}%
\definecolor{currentstroke}{rgb}{0.839216,0.152941,0.156863}%
\pgfsetstrokecolor{currentstroke}%
\pgfsetdash{}{0pt}%
\pgfpathmoveto{\pgfqpoint{4.188593in}{0.956348in}}%
\pgfpathlineto{\pgfqpoint{4.235219in}{0.958641in}}%
\pgfpathlineto{\pgfqpoint{4.281845in}{0.961069in}}%
\pgfpathlineto{\pgfqpoint{4.328470in}{0.963661in}}%
\pgfpathlineto{\pgfqpoint{4.375096in}{0.966448in}}%
\pgfpathlineto{\pgfqpoint{4.421722in}{0.969463in}}%
\pgfpathlineto{\pgfqpoint{4.468347in}{0.972777in}}%
\pgfpathlineto{\pgfqpoint{4.514973in}{0.976463in}}%
\pgfpathlineto{\pgfqpoint{4.561599in}{0.980420in}}%
\pgfpathlineto{\pgfqpoint{4.608224in}{0.984760in}}%
\pgfpathlineto{\pgfqpoint{4.654850in}{0.989538in}}%
\pgfpathlineto{\pgfqpoint{4.701476in}{0.994808in}}%
\pgfpathlineto{\pgfqpoint{4.748101in}{1.000634in}}%
\pgfpathlineto{\pgfqpoint{4.794727in}{1.007082in}}%
\pgfpathlineto{\pgfqpoint{4.841353in}{1.014231in}}%
\pgfpathlineto{\pgfqpoint{4.887978in}{1.022166in}}%
\pgfpathlineto{\pgfqpoint{4.934604in}{1.030998in}}%
\pgfpathlineto{\pgfqpoint{4.981230in}{1.040911in}}%
\pgfpathlineto{\pgfqpoint{5.027855in}{1.051926in}}%
\pgfpathlineto{\pgfqpoint{5.074481in}{1.064174in}}%
\pgfpathlineto{\pgfqpoint{5.121107in}{1.077797in}}%
\pgfpathlineto{\pgfqpoint{5.167732in}{1.092956in}}%
\pgfpathlineto{\pgfqpoint{5.214358in}{1.109829in}}%
\pgfpathlineto{\pgfqpoint{5.260984in}{1.128613in}}%
\pgfpathlineto{\pgfqpoint{5.307610in}{1.149528in}}%
\pgfpathlineto{\pgfqpoint{5.354235in}{1.172816in}}%
\pgfpathlineto{\pgfqpoint{5.400861in}{1.198746in}}%
\pgfpathlineto{\pgfqpoint{5.447487in}{1.227591in}}%
\pgfpathlineto{\pgfqpoint{5.494112in}{1.259684in}}%
\pgfpathlineto{\pgfqpoint{5.540738in}{1.295413in}}%
\pgfpathlineto{\pgfqpoint{5.587364in}{1.335177in}}%
\pgfpathlineto{\pgfqpoint{5.633989in}{1.379413in}}%
\pgfpathlineto{\pgfqpoint{5.680615in}{1.428601in}}%
\pgfpathlineto{\pgfqpoint{5.727241in}{1.483251in}}%
\pgfpathlineto{\pgfqpoint{5.773866in}{1.543899in}}%
\pgfpathlineto{\pgfqpoint{5.820492in}{1.611068in}}%
\pgfpathlineto{\pgfqpoint{5.867118in}{1.685211in}}%
\pgfpathlineto{\pgfqpoint{5.913743in}{1.766604in}}%
\pgfpathlineto{\pgfqpoint{5.960369in}{1.855197in}}%
\pgfpathlineto{\pgfqpoint{6.006995in}{1.950355in}}%
\pgfpathlineto{\pgfqpoint{6.053620in}{2.050692in}}%
\pgfpathlineto{\pgfqpoint{6.100246in}{2.153883in}}%
\pgfpathlineto{\pgfqpoint{6.146872in}{2.256671in}}%
\pgfpathlineto{\pgfqpoint{6.193497in}{2.355126in}}%
\pgfpathlineto{\pgfqpoint{6.240123in}{2.445196in}}%
\pgfpathlineto{\pgfqpoint{6.286749in}{2.523446in}}%
\pgfpathlineto{\pgfqpoint{6.333374in}{2.587630in}}%
\pgfpathlineto{\pgfqpoint{6.380000in}{2.637260in}}%
\pgfpathlineto{\pgfqpoint{6.426626in}{2.673250in}}%
\pgfpathlineto{\pgfqpoint{6.473251in}{2.697721in}}%
\pgfpathlineto{\pgfqpoint{6.519877in}{2.713268in}}%
\pgfpathlineto{\pgfqpoint{6.566503in}{2.722411in}}%
\pgfpathlineto{\pgfqpoint{6.613128in}{2.727409in}}%
\pgfpathlineto{\pgfqpoint{6.659754in}{2.729923in}}%
\pgfpathlineto{\pgfqpoint{6.706380in}{2.731078in}}%
\pgfusepath{stroke}%
\end{pgfscope}%
\begin{pgfscope}%
\pgfpathrectangle{\pgfqpoint{4.062704in}{0.807389in}}{\pgfqpoint{2.769565in}{2.016000in}}%
\pgfusepath{clip}%
\pgfsetrectcap%
\pgfsetroundjoin%
\pgfsetlinewidth{1.806750pt}%
\definecolor{currentstroke}{rgb}{0.580392,0.403922,0.741176}%
\pgfsetstrokecolor{currentstroke}%
\pgfsetdash{}{0pt}%
\pgfpathmoveto{\pgfqpoint{4.188593in}{0.949495in}}%
\pgfpathlineto{\pgfqpoint{4.235219in}{0.957763in}}%
\pgfpathlineto{\pgfqpoint{4.281845in}{0.966999in}}%
\pgfpathlineto{\pgfqpoint{4.328470in}{0.977318in}}%
\pgfpathlineto{\pgfqpoint{4.375096in}{0.988845in}}%
\pgfpathlineto{\pgfqpoint{4.421722in}{1.001722in}}%
\pgfpathlineto{\pgfqpoint{4.468347in}{1.016106in}}%
\pgfpathlineto{\pgfqpoint{4.514973in}{1.032174in}}%
\pgfpathlineto{\pgfqpoint{4.561599in}{1.050121in}}%
\pgfpathlineto{\pgfqpoint{4.608224in}{1.070168in}}%
\pgfpathlineto{\pgfqpoint{4.654850in}{1.092559in}}%
\pgfpathlineto{\pgfqpoint{4.701476in}{1.117566in}}%
\pgfpathlineto{\pgfqpoint{4.748101in}{1.145494in}}%
\pgfpathlineto{\pgfqpoint{4.794727in}{1.176683in}}%
\pgfpathlineto{\pgfqpoint{4.841353in}{1.211510in}}%
\pgfpathlineto{\pgfqpoint{4.887978in}{1.250398in}}%
\pgfpathlineto{\pgfqpoint{4.934604in}{1.293814in}}%
\pgfpathlineto{\pgfqpoint{4.981230in}{1.342275in}}%
\pgfpathlineto{\pgfqpoint{5.027855in}{1.396347in}}%
\pgfpathlineto{\pgfqpoint{5.074481in}{1.456624in}}%
\pgfpathlineto{\pgfqpoint{5.121107in}{1.523705in}}%
\pgfpathlineto{\pgfqpoint{5.167732in}{1.598120in}}%
\pgfpathlineto{\pgfqpoint{5.214358in}{1.680222in}}%
\pgfpathlineto{\pgfqpoint{5.260984in}{1.770024in}}%
\pgfpathlineto{\pgfqpoint{5.307610in}{1.866980in}}%
\pgfpathlineto{\pgfqpoint{5.354235in}{1.969773in}}%
\pgfpathlineto{\pgfqpoint{5.400861in}{2.076142in}}%
\pgfpathlineto{\pgfqpoint{5.447487in}{2.182883in}}%
\pgfpathlineto{\pgfqpoint{5.494112in}{2.286100in}}%
\pgfpathlineto{\pgfqpoint{5.540738in}{2.381733in}}%
\pgfpathlineto{\pgfqpoint{5.587364in}{2.466259in}}%
\pgfpathlineto{\pgfqpoint{5.633989in}{2.537335in}}%
\pgfpathlineto{\pgfqpoint{5.680615in}{2.594150in}}%
\pgfpathlineto{\pgfqpoint{5.727241in}{2.637369in}}%
\pgfpathlineto{\pgfqpoint{5.773866in}{2.668743in}}%
\pgfpathlineto{\pgfqpoint{5.820492in}{2.690571in}}%
\pgfpathlineto{\pgfqpoint{5.867118in}{2.705208in}}%
\pgfpathlineto{\pgfqpoint{5.913743in}{2.714733in}}%
\pgfpathlineto{\pgfqpoint{5.960369in}{2.720801in}}%
\pgfpathlineto{\pgfqpoint{6.006995in}{2.724623in}}%
\pgfpathlineto{\pgfqpoint{6.053620in}{2.727028in}}%
\pgfpathlineto{\pgfqpoint{6.100246in}{2.728559in}}%
\pgfpathlineto{\pgfqpoint{6.146872in}{2.729550in}}%
\pgfpathlineto{\pgfqpoint{6.193497in}{2.730209in}}%
\pgfpathlineto{\pgfqpoint{6.240123in}{2.730657in}}%
\pgfpathlineto{\pgfqpoint{6.286749in}{2.730969in}}%
\pgfpathlineto{\pgfqpoint{6.333374in}{2.731190in}}%
\pgfpathlineto{\pgfqpoint{6.380000in}{2.731348in}}%
\pgfpathlineto{\pgfqpoint{6.426626in}{2.731462in}}%
\pgfpathlineto{\pgfqpoint{6.473251in}{2.731545in}}%
\pgfpathlineto{\pgfqpoint{6.519877in}{2.731608in}}%
\pgfpathlineto{\pgfqpoint{6.566503in}{2.731655in}}%
\pgfpathlineto{\pgfqpoint{6.613128in}{2.731693in}}%
\pgfpathlineto{\pgfqpoint{6.659754in}{2.731725in}}%
\pgfpathlineto{\pgfqpoint{6.706380in}{2.731753in}}%
\pgfusepath{stroke}%
\end{pgfscope}%
\begin{pgfscope}%
\pgfsetrectcap%
\pgfsetmiterjoin%
\pgfsetlinewidth{0.803000pt}%
\definecolor{currentstroke}{rgb}{0.000000,0.000000,0.000000}%
\pgfsetstrokecolor{currentstroke}%
\pgfsetdash{}{0pt}%
\pgfpathmoveto{\pgfqpoint{4.062704in}{0.807389in}}%
\pgfpathlineto{\pgfqpoint{4.062704in}{2.823389in}}%
\pgfusepath{stroke}%
\end{pgfscope}%
\begin{pgfscope}%
\pgfsetrectcap%
\pgfsetmiterjoin%
\pgfsetlinewidth{0.803000pt}%
\definecolor{currentstroke}{rgb}{0.000000,0.000000,0.000000}%
\pgfsetstrokecolor{currentstroke}%
\pgfsetdash{}{0pt}%
\pgfpathmoveto{\pgfqpoint{6.832269in}{0.807389in}}%
\pgfpathlineto{\pgfqpoint{6.832269in}{2.823389in}}%
\pgfusepath{stroke}%
\end{pgfscope}%
\begin{pgfscope}%
\pgfsetrectcap%
\pgfsetmiterjoin%
\pgfsetlinewidth{0.803000pt}%
\definecolor{currentstroke}{rgb}{0.000000,0.000000,0.000000}%
\pgfsetstrokecolor{currentstroke}%
\pgfsetdash{}{0pt}%
\pgfpathmoveto{\pgfqpoint{4.062704in}{0.807389in}}%
\pgfpathlineto{\pgfqpoint{6.832269in}{0.807389in}}%
\pgfusepath{stroke}%
\end{pgfscope}%
\begin{pgfscope}%
\pgfsetrectcap%
\pgfsetmiterjoin%
\pgfsetlinewidth{0.803000pt}%
\definecolor{currentstroke}{rgb}{0.000000,0.000000,0.000000}%
\pgfsetstrokecolor{currentstroke}%
\pgfsetdash{}{0pt}%
\pgfpathmoveto{\pgfqpoint{4.062704in}{2.823389in}}%
\pgfpathlineto{\pgfqpoint{6.832269in}{2.823389in}}%
\pgfusepath{stroke}%
\end{pgfscope}%
\begin{pgfscope}%
\definecolor{textcolor}{rgb}{0.000000,0.000000,0.000000}%
\pgfsetstrokecolor{textcolor}%
\pgfsetfillcolor{textcolor}%
\pgftext[x=5.447487in,y=2.906722in,,base]{\color{textcolor}{\rmfamily\fontsize{8.000000}{9.600000}\selectfont\catcode`\^=\active\def^{\ifmmode\sp\else\^{}\fi}\catcode`\%=\active\def%{\%}(b) Replicated Volume Profiles}}%
\end{pgfscope}%
\begin{pgfscope}%
\pgfsetrectcap%
\pgfsetroundjoin%
\pgfsetlinewidth{1.806750pt}%
\definecolor{currentstroke}{rgb}{0.121569,0.466667,0.705882}%
\pgfsetstrokecolor{currentstroke}%
\pgfsetdash{}{0pt}%
\pgfpathmoveto{\pgfqpoint{1.366181in}{0.191821in}}%
\pgfpathlineto{\pgfqpoint{1.439098in}{0.191821in}}%
\pgfpathlineto{\pgfqpoint{1.512014in}{0.191821in}}%
\pgfusepath{stroke}%
\end{pgfscope}%
\begin{pgfscope}%
\definecolor{textcolor}{rgb}{0.000000,0.000000,0.000000}%
\pgfsetstrokecolor{textcolor}%
\pgfsetfillcolor{textcolor}%
\pgftext[x=1.589792in,y=0.157793in,left,base]{\color{textcolor}{\rmfamily\fontsize{7.000000}{8.400000}\selectfont\catcode`\^=\active\def^{\ifmmode\sp\else\^{}\fi}\catcode`\%=\active\def%{\%}Unimodal vMF}}%
\end{pgfscope}%
\begin{pgfscope}%
\pgfsetrectcap%
\pgfsetroundjoin%
\pgfsetlinewidth{1.806750pt}%
\definecolor{currentstroke}{rgb}{1.000000,0.498039,0.054902}%
\pgfsetstrokecolor{currentstroke}%
\pgfsetdash{}{0pt}%
\pgfpathmoveto{\pgfqpoint{2.417168in}{0.191821in}}%
\pgfpathlineto{\pgfqpoint{2.490085in}{0.191821in}}%
\pgfpathlineto{\pgfqpoint{2.563002in}{0.191821in}}%
\pgfusepath{stroke}%
\end{pgfscope}%
\begin{pgfscope}%
\definecolor{textcolor}{rgb}{0.000000,0.000000,0.000000}%
\pgfsetstrokecolor{textcolor}%
\pgfsetfillcolor{textcolor}%
\pgftext[x=2.640780in,y=0.157793in,left,base]{\color{textcolor}{\rmfamily\fontsize{7.000000}{8.400000}\selectfont\catcode`\^=\active\def^{\ifmmode\sp\else\^{}\fi}\catcode`\%=\active\def%{\%}Antipodal}}%
\end{pgfscope}%
\begin{pgfscope}%
\pgfsetrectcap%
\pgfsetroundjoin%
\pgfsetlinewidth{1.806750pt}%
\definecolor{currentstroke}{rgb}{0.172549,0.627451,0.172549}%
\pgfsetstrokecolor{currentstroke}%
\pgfsetdash{}{0pt}%
\pgfpathmoveto{\pgfqpoint{3.214296in}{0.191821in}}%
\pgfpathlineto{\pgfqpoint{3.287213in}{0.191821in}}%
\pgfpathlineto{\pgfqpoint{3.360130in}{0.191821in}}%
\pgfusepath{stroke}%
\end{pgfscope}%
\begin{pgfscope}%
\definecolor{textcolor}{rgb}{0.000000,0.000000,0.000000}%
\pgfsetstrokecolor{textcolor}%
\pgfsetfillcolor{textcolor}%
\pgftext[x=3.437907in,y=0.157793in,left,base]{\color{textcolor}{\rmfamily\fontsize{7.000000}{8.400000}\selectfont\catcode`\^=\active\def^{\ifmmode\sp\else\^{}\fi}\catcode`\%=\active\def%{\%}Girdle}}%
\end{pgfscope}%
\begin{pgfscope}%
\pgfsetrectcap%
\pgfsetroundjoin%
\pgfsetlinewidth{1.806750pt}%
\definecolor{currentstroke}{rgb}{0.839216,0.152941,0.156863}%
\pgfsetstrokecolor{currentstroke}%
\pgfsetdash{}{0pt}%
\pgfpathmoveto{\pgfqpoint{3.828360in}{0.191821in}}%
\pgfpathlineto{\pgfqpoint{3.901277in}{0.191821in}}%
\pgfpathlineto{\pgfqpoint{3.974194in}{0.191821in}}%
\pgfusepath{stroke}%
\end{pgfscope}%
\begin{pgfscope}%
\definecolor{textcolor}{rgb}{0.000000,0.000000,0.000000}%
\pgfsetstrokecolor{textcolor}%
\pgfsetfillcolor{textcolor}%
\pgftext[x=4.051971in,y=0.157793in,left,base]{\color{textcolor}{\rmfamily\fontsize{7.000000}{8.400000}\selectfont\catcode`\^=\active\def^{\ifmmode\sp\else\^{}\fi}\catcode`\%=\active\def%{\%}Four modes}}%
\end{pgfscope}%
\begin{pgfscope}%
\pgfsetrectcap%
\pgfsetroundjoin%
\pgfsetlinewidth{1.806750pt}%
\definecolor{currentstroke}{rgb}{0.580392,0.403922,0.741176}%
\pgfsetstrokecolor{currentstroke}%
\pgfsetdash{}{0pt}%
\pgfpathmoveto{\pgfqpoint{4.704019in}{0.191821in}}%
\pgfpathlineto{\pgfqpoint{4.776935in}{0.191821in}}%
\pgfpathlineto{\pgfqpoint{4.849852in}{0.191821in}}%
\pgfusepath{stroke}%
\end{pgfscope}%
\begin{pgfscope}%
\definecolor{textcolor}{rgb}{0.000000,0.000000,0.000000}%
\pgfsetstrokecolor{textcolor}%
\pgfsetfillcolor{textcolor}%
\pgftext[x=4.927630in,y=0.157793in,left,base]{\color{textcolor}{\rmfamily\fontsize{7.000000}{8.400000}\selectfont\catcode`\^=\active\def^{\ifmmode\sp\else\^{}\fi}\catcode`\%=\active\def%{\%}Duplicates}}%
\end{pgfscope}%
\end{pgfpicture}%
\makeatother%
\endgroup%

\caption{Replicated synthetic profiles on $S^2$. Each curve is a median over 30 replications, with an interquartile band.
The left panel shows gESS. The right panel shows effective occupied volume. Profile shapes are stable across replications.}
\label{fig:app_replicates}
\end{figure*}

\paragraph{Spectral truncation sensitivity.}
Figure~\ref{fig:app_truncation} varies $L_{\max}$ in the $S^2$ expansion. The reference uses $L_{\max}=360$. The figure reports relative error in $\widehat A_w(t)$ for a four-mode cloud. Fine angular scales require more terms, consistent with $L_{\max}$ of order $t^{-1/2}$. The main-paper cutoff is adequate for the reported scales. The finest scales are also statistically and numerically less stable.

\begin{figure*}[t]
\centering
%\includegraphics[width=0.98\linewidth]{hkep_app_truncation.png}
\providecommand{\mathdefault}[1]{#1}
\providecommand{\mathregular}[1]{#1}
%% Creator: Matplotlib, PGF backend
%%
%% To include the figure in your LaTeX document, write
%%   \input{<filename>.pgf}
%%
%% Make sure the required packages are loaded in your preamble
%%   \usepackage{pgf}
%%
%% Also ensure that all the required font packages are loaded; for instance,
%% the lmodern package is sometimes necessary when using math font.
%%   \usepackage{lmodern}
%%
%% Figures using additional raster images can only be included by \input if
%% they are in the same directory as the main LaTeX file. For loading figures
%% from other directories you can use the `import` package
%%   \usepackage{import}
%%
%% and then include the figures with
%%   \import{<path to file>}{<filename>.pgf}
%%
%% Matplotlib used the following preamble
%%   \def\mathdefault#1{#1}
%%   \everymath=\expandafter{\the\everymath\displaystyle}
%%   \IfFileExists{scrextend.sty}{
%%     \usepackage[fontsize=8.000000pt]{scrextend}
%%   }{
%%     \renewcommand{\normalsize}{\fontsize{8.000000}{9.600000}\selectfont}
%%     \normalsize
%%   }
%%   \usepackage{amsmath,amssymb}
%%   \makeatletter\@ifpackageloaded{underscore}{}{\usepackage[strings]{underscore}}\makeatother
%%
\begingroup%
\makeatletter%
\begin{pgfpicture}%
\pgfpathrectangle{\pgfpointorigin}{\pgfqpoint{5.342374in}{2.819916in}}%
\pgfusepath{use as bounding box, clip}%
\begin{pgfscope}%
\pgfsetbuttcap%
\pgfsetmiterjoin%
\definecolor{currentfill}{rgb}{1.000000,1.000000,1.000000}%
\pgfsetfillcolor{currentfill}%
\pgfsetlinewidth{0.000000pt}%
\definecolor{currentstroke}{rgb}{1.000000,1.000000,1.000000}%
\pgfsetstrokecolor{currentstroke}%
\pgfsetdash{}{0pt}%
\pgfpathmoveto{\pgfqpoint{0.000000in}{0.000000in}}%
\pgfpathlineto{\pgfqpoint{5.342374in}{0.000000in}}%
\pgfpathlineto{\pgfqpoint{5.342374in}{2.819916in}}%
\pgfpathlineto{\pgfqpoint{0.000000in}{2.819916in}}%
\pgfpathlineto{\pgfqpoint{0.000000in}{0.000000in}}%
\pgfpathclose%
\pgfusepath{fill}%
\end{pgfscope}%
\begin{pgfscope}%
\pgfsetbuttcap%
\pgfsetmiterjoin%
\definecolor{currentfill}{rgb}{1.000000,1.000000,1.000000}%
\pgfsetfillcolor{currentfill}%
\pgfsetlinewidth{0.000000pt}%
\definecolor{currentstroke}{rgb}{0.000000,0.000000,0.000000}%
\pgfsetstrokecolor{currentstroke}%
\pgfsetstrokeopacity{0.000000}%
\pgfsetdash{}{0pt}%
\pgfpathmoveto{\pgfqpoint{0.676315in}{0.668050in}}%
\pgfpathlineto{\pgfqpoint{5.242374in}{0.668050in}}%
\pgfpathlineto{\pgfqpoint{5.242374in}{2.536721in}}%
\pgfpathlineto{\pgfqpoint{0.676315in}{2.536721in}}%
\pgfpathlineto{\pgfqpoint{0.676315in}{0.668050in}}%
\pgfpathclose%
\pgfusepath{fill}%
\end{pgfscope}%
\begin{pgfscope}%
\pgfpathrectangle{\pgfqpoint{0.676315in}{0.668050in}}{\pgfqpoint{4.566060in}{1.868671in}}%
\pgfusepath{clip}%
\pgfsetrectcap%
\pgfsetroundjoin%
\pgfsetlinewidth{0.803000pt}%
\definecolor{currentstroke}{rgb}{0.690196,0.690196,0.690196}%
\pgfsetstrokecolor{currentstroke}%
\pgfsetstrokeopacity{0.300000}%
\pgfsetdash{}{0pt}%
\pgfpathmoveto{\pgfqpoint{0.823704in}{0.668050in}}%
\pgfpathlineto{\pgfqpoint{0.823704in}{2.536721in}}%
\pgfusepath{stroke}%
\end{pgfscope}%
\begin{pgfscope}%
\pgfsetbuttcap%
\pgfsetroundjoin%
\definecolor{currentfill}{rgb}{0.000000,0.000000,0.000000}%
\pgfsetfillcolor{currentfill}%
\pgfsetlinewidth{0.803000pt}%
\definecolor{currentstroke}{rgb}{0.000000,0.000000,0.000000}%
\pgfsetstrokecolor{currentstroke}%
\pgfsetdash{}{0pt}%
\pgfsys@defobject{currentmarker}{\pgfqpoint{0.000000in}{-0.048611in}}{\pgfqpoint{0.000000in}{0.000000in}}{%
\pgfpathmoveto{\pgfqpoint{0.000000in}{0.000000in}}%
\pgfpathlineto{\pgfqpoint{0.000000in}{-0.048611in}}%
\pgfusepath{stroke,fill}%
}%
\begin{pgfscope}%
\pgfsys@transformshift{0.823704in}{0.668050in}%
\pgfsys@useobject{currentmarker}{}%
\end{pgfscope}%
\end{pgfscope}%
\begin{pgfscope}%
\definecolor{textcolor}{rgb}{0.000000,0.000000,0.000000}%
\pgfsetstrokecolor{textcolor}%
\pgfsetfillcolor{textcolor}%
\pgftext[x=0.823704in,y=0.570828in,,top]{\color{textcolor}{\rmfamily\fontsize{7.000000}{8.400000}\selectfont\catcode`\^=\active\def^{\ifmmode\sp\else\^{}\fi}\catcode`\%=\active\def%{\%}0}}%
\end{pgfscope}%
\begin{pgfscope}%
\pgfpathrectangle{\pgfqpoint{0.676315in}{0.668050in}}{\pgfqpoint{4.566060in}{1.868671in}}%
\pgfusepath{clip}%
\pgfsetrectcap%
\pgfsetroundjoin%
\pgfsetlinewidth{0.803000pt}%
\definecolor{currentstroke}{rgb}{0.690196,0.690196,0.690196}%
\pgfsetstrokecolor{currentstroke}%
\pgfsetstrokeopacity{0.300000}%
\pgfsetdash{}{0pt}%
\pgfpathmoveto{\pgfqpoint{1.575690in}{0.668050in}}%
\pgfpathlineto{\pgfqpoint{1.575690in}{2.536721in}}%
\pgfusepath{stroke}%
\end{pgfscope}%
\begin{pgfscope}%
\pgfsetbuttcap%
\pgfsetroundjoin%
\definecolor{currentfill}{rgb}{0.000000,0.000000,0.000000}%
\pgfsetfillcolor{currentfill}%
\pgfsetlinewidth{0.803000pt}%
\definecolor{currentstroke}{rgb}{0.000000,0.000000,0.000000}%
\pgfsetstrokecolor{currentstroke}%
\pgfsetdash{}{0pt}%
\pgfsys@defobject{currentmarker}{\pgfqpoint{0.000000in}{-0.048611in}}{\pgfqpoint{0.000000in}{0.000000in}}{%
\pgfpathmoveto{\pgfqpoint{0.000000in}{0.000000in}}%
\pgfpathlineto{\pgfqpoint{0.000000in}{-0.048611in}}%
\pgfusepath{stroke,fill}%
}%
\begin{pgfscope}%
\pgfsys@transformshift{1.575690in}{0.668050in}%
\pgfsys@useobject{currentmarker}{}%
\end{pgfscope}%
\end{pgfscope}%
\begin{pgfscope}%
\definecolor{textcolor}{rgb}{0.000000,0.000000,0.000000}%
\pgfsetstrokecolor{textcolor}%
\pgfsetfillcolor{textcolor}%
\pgftext[x=1.575690in,y=0.570828in,,top]{\color{textcolor}{\rmfamily\fontsize{7.000000}{8.400000}\selectfont\catcode`\^=\active\def^{\ifmmode\sp\else\^{}\fi}\catcode`\%=\active\def%{\%}50}}%
\end{pgfscope}%
\begin{pgfscope}%
\pgfpathrectangle{\pgfqpoint{0.676315in}{0.668050in}}{\pgfqpoint{4.566060in}{1.868671in}}%
\pgfusepath{clip}%
\pgfsetrectcap%
\pgfsetroundjoin%
\pgfsetlinewidth{0.803000pt}%
\definecolor{currentstroke}{rgb}{0.690196,0.690196,0.690196}%
\pgfsetstrokecolor{currentstroke}%
\pgfsetstrokeopacity{0.300000}%
\pgfsetdash{}{0pt}%
\pgfpathmoveto{\pgfqpoint{2.327676in}{0.668050in}}%
\pgfpathlineto{\pgfqpoint{2.327676in}{2.536721in}}%
\pgfusepath{stroke}%
\end{pgfscope}%
\begin{pgfscope}%
\pgfsetbuttcap%
\pgfsetroundjoin%
\definecolor{currentfill}{rgb}{0.000000,0.000000,0.000000}%
\pgfsetfillcolor{currentfill}%
\pgfsetlinewidth{0.803000pt}%
\definecolor{currentstroke}{rgb}{0.000000,0.000000,0.000000}%
\pgfsetstrokecolor{currentstroke}%
\pgfsetdash{}{0pt}%
\pgfsys@defobject{currentmarker}{\pgfqpoint{0.000000in}{-0.048611in}}{\pgfqpoint{0.000000in}{0.000000in}}{%
\pgfpathmoveto{\pgfqpoint{0.000000in}{0.000000in}}%
\pgfpathlineto{\pgfqpoint{0.000000in}{-0.048611in}}%
\pgfusepath{stroke,fill}%
}%
\begin{pgfscope}%
\pgfsys@transformshift{2.327676in}{0.668050in}%
\pgfsys@useobject{currentmarker}{}%
\end{pgfscope}%
\end{pgfscope}%
\begin{pgfscope}%
\definecolor{textcolor}{rgb}{0.000000,0.000000,0.000000}%
\pgfsetstrokecolor{textcolor}%
\pgfsetfillcolor{textcolor}%
\pgftext[x=2.327676in,y=0.570828in,,top]{\color{textcolor}{\rmfamily\fontsize{7.000000}{8.400000}\selectfont\catcode`\^=\active\def^{\ifmmode\sp\else\^{}\fi}\catcode`\%=\active\def%{\%}100}}%
\end{pgfscope}%
\begin{pgfscope}%
\pgfpathrectangle{\pgfqpoint{0.676315in}{0.668050in}}{\pgfqpoint{4.566060in}{1.868671in}}%
\pgfusepath{clip}%
\pgfsetrectcap%
\pgfsetroundjoin%
\pgfsetlinewidth{0.803000pt}%
\definecolor{currentstroke}{rgb}{0.690196,0.690196,0.690196}%
\pgfsetstrokecolor{currentstroke}%
\pgfsetstrokeopacity{0.300000}%
\pgfsetdash{}{0pt}%
\pgfpathmoveto{\pgfqpoint{3.079662in}{0.668050in}}%
\pgfpathlineto{\pgfqpoint{3.079662in}{2.536721in}}%
\pgfusepath{stroke}%
\end{pgfscope}%
\begin{pgfscope}%
\pgfsetbuttcap%
\pgfsetroundjoin%
\definecolor{currentfill}{rgb}{0.000000,0.000000,0.000000}%
\pgfsetfillcolor{currentfill}%
\pgfsetlinewidth{0.803000pt}%
\definecolor{currentstroke}{rgb}{0.000000,0.000000,0.000000}%
\pgfsetstrokecolor{currentstroke}%
\pgfsetdash{}{0pt}%
\pgfsys@defobject{currentmarker}{\pgfqpoint{0.000000in}{-0.048611in}}{\pgfqpoint{0.000000in}{0.000000in}}{%
\pgfpathmoveto{\pgfqpoint{0.000000in}{0.000000in}}%
\pgfpathlineto{\pgfqpoint{0.000000in}{-0.048611in}}%
\pgfusepath{stroke,fill}%
}%
\begin{pgfscope}%
\pgfsys@transformshift{3.079662in}{0.668050in}%
\pgfsys@useobject{currentmarker}{}%
\end{pgfscope}%
\end{pgfscope}%
\begin{pgfscope}%
\definecolor{textcolor}{rgb}{0.000000,0.000000,0.000000}%
\pgfsetstrokecolor{textcolor}%
\pgfsetfillcolor{textcolor}%
\pgftext[x=3.079662in,y=0.570828in,,top]{\color{textcolor}{\rmfamily\fontsize{7.000000}{8.400000}\selectfont\catcode`\^=\active\def^{\ifmmode\sp\else\^{}\fi}\catcode`\%=\active\def%{\%}150}}%
\end{pgfscope}%
\begin{pgfscope}%
\pgfpathrectangle{\pgfqpoint{0.676315in}{0.668050in}}{\pgfqpoint{4.566060in}{1.868671in}}%
\pgfusepath{clip}%
\pgfsetrectcap%
\pgfsetroundjoin%
\pgfsetlinewidth{0.803000pt}%
\definecolor{currentstroke}{rgb}{0.690196,0.690196,0.690196}%
\pgfsetstrokecolor{currentstroke}%
\pgfsetstrokeopacity{0.300000}%
\pgfsetdash{}{0pt}%
\pgfpathmoveto{\pgfqpoint{3.831648in}{0.668050in}}%
\pgfpathlineto{\pgfqpoint{3.831648in}{2.536721in}}%
\pgfusepath{stroke}%
\end{pgfscope}%
\begin{pgfscope}%
\pgfsetbuttcap%
\pgfsetroundjoin%
\definecolor{currentfill}{rgb}{0.000000,0.000000,0.000000}%
\pgfsetfillcolor{currentfill}%
\pgfsetlinewidth{0.803000pt}%
\definecolor{currentstroke}{rgb}{0.000000,0.000000,0.000000}%
\pgfsetstrokecolor{currentstroke}%
\pgfsetdash{}{0pt}%
\pgfsys@defobject{currentmarker}{\pgfqpoint{0.000000in}{-0.048611in}}{\pgfqpoint{0.000000in}{0.000000in}}{%
\pgfpathmoveto{\pgfqpoint{0.000000in}{0.000000in}}%
\pgfpathlineto{\pgfqpoint{0.000000in}{-0.048611in}}%
\pgfusepath{stroke,fill}%
}%
\begin{pgfscope}%
\pgfsys@transformshift{3.831648in}{0.668050in}%
\pgfsys@useobject{currentmarker}{}%
\end{pgfscope}%
\end{pgfscope}%
\begin{pgfscope}%
\definecolor{textcolor}{rgb}{0.000000,0.000000,0.000000}%
\pgfsetstrokecolor{textcolor}%
\pgfsetfillcolor{textcolor}%
\pgftext[x=3.831648in,y=0.570828in,,top]{\color{textcolor}{\rmfamily\fontsize{7.000000}{8.400000}\selectfont\catcode`\^=\active\def^{\ifmmode\sp\else\^{}\fi}\catcode`\%=\active\def%{\%}200}}%
\end{pgfscope}%
\begin{pgfscope}%
\pgfpathrectangle{\pgfqpoint{0.676315in}{0.668050in}}{\pgfqpoint{4.566060in}{1.868671in}}%
\pgfusepath{clip}%
\pgfsetrectcap%
\pgfsetroundjoin%
\pgfsetlinewidth{0.803000pt}%
\definecolor{currentstroke}{rgb}{0.690196,0.690196,0.690196}%
\pgfsetstrokecolor{currentstroke}%
\pgfsetstrokeopacity{0.300000}%
\pgfsetdash{}{0pt}%
\pgfpathmoveto{\pgfqpoint{4.583635in}{0.668050in}}%
\pgfpathlineto{\pgfqpoint{4.583635in}{2.536721in}}%
\pgfusepath{stroke}%
\end{pgfscope}%
\begin{pgfscope}%
\pgfsetbuttcap%
\pgfsetroundjoin%
\definecolor{currentfill}{rgb}{0.000000,0.000000,0.000000}%
\pgfsetfillcolor{currentfill}%
\pgfsetlinewidth{0.803000pt}%
\definecolor{currentstroke}{rgb}{0.000000,0.000000,0.000000}%
\pgfsetstrokecolor{currentstroke}%
\pgfsetdash{}{0pt}%
\pgfsys@defobject{currentmarker}{\pgfqpoint{0.000000in}{-0.048611in}}{\pgfqpoint{0.000000in}{0.000000in}}{%
\pgfpathmoveto{\pgfqpoint{0.000000in}{0.000000in}}%
\pgfpathlineto{\pgfqpoint{0.000000in}{-0.048611in}}%
\pgfusepath{stroke,fill}%
}%
\begin{pgfscope}%
\pgfsys@transformshift{4.583635in}{0.668050in}%
\pgfsys@useobject{currentmarker}{}%
\end{pgfscope}%
\end{pgfscope}%
\begin{pgfscope}%
\definecolor{textcolor}{rgb}{0.000000,0.000000,0.000000}%
\pgfsetstrokecolor{textcolor}%
\pgfsetfillcolor{textcolor}%
\pgftext[x=4.583635in,y=0.570828in,,top]{\color{textcolor}{\rmfamily\fontsize{7.000000}{8.400000}\selectfont\catcode`\^=\active\def^{\ifmmode\sp\else\^{}\fi}\catcode`\%=\active\def%{\%}250}}%
\end{pgfscope}%
\begin{pgfscope}%
\definecolor{textcolor}{rgb}{0.000000,0.000000,0.000000}%
\pgfsetstrokecolor{textcolor}%
\pgfsetfillcolor{textcolor}%
\pgftext[x=2.959345in,y=0.428853in,,top]{\color{textcolor}{\rmfamily\fontsize{8.000000}{9.600000}\selectfont\catcode`\^=\active\def^{\ifmmode\sp\else\^{}\fi}\catcode`\%=\active\def%{\%}Harmonic Truncation $L_{\max}$}}%
\end{pgfscope}%
\begin{pgfscope}%
\pgfpathrectangle{\pgfqpoint{0.676315in}{0.668050in}}{\pgfqpoint{4.566060in}{1.868671in}}%
\pgfusepath{clip}%
\pgfsetrectcap%
\pgfsetroundjoin%
\pgfsetlinewidth{0.803000pt}%
\definecolor{currentstroke}{rgb}{0.690196,0.690196,0.690196}%
\pgfsetstrokecolor{currentstroke}%
\pgfsetstrokeopacity{0.300000}%
\pgfsetdash{}{0pt}%
\pgfpathmoveto{\pgfqpoint{0.676315in}{0.814490in}}%
\pgfpathlineto{\pgfqpoint{5.242374in}{0.814490in}}%
\pgfusepath{stroke}%
\end{pgfscope}%
\begin{pgfscope}%
\pgfsetbuttcap%
\pgfsetroundjoin%
\definecolor{currentfill}{rgb}{0.000000,0.000000,0.000000}%
\pgfsetfillcolor{currentfill}%
\pgfsetlinewidth{0.803000pt}%
\definecolor{currentstroke}{rgb}{0.000000,0.000000,0.000000}%
\pgfsetstrokecolor{currentstroke}%
\pgfsetdash{}{0pt}%
\pgfsys@defobject{currentmarker}{\pgfqpoint{-0.048611in}{0.000000in}}{\pgfqpoint{-0.000000in}{0.000000in}}{%
\pgfpathmoveto{\pgfqpoint{-0.000000in}{0.000000in}}%
\pgfpathlineto{\pgfqpoint{-0.048611in}{0.000000in}}%
\pgfusepath{stroke,fill}%
}%
\begin{pgfscope}%
\pgfsys@transformshift{0.676315in}{0.814490in}%
\pgfsys@useobject{currentmarker}{}%
\end{pgfscope}%
\end{pgfscope}%
\begin{pgfscope}%
\definecolor{textcolor}{rgb}{0.000000,0.000000,0.000000}%
\pgfsetstrokecolor{textcolor}%
\pgfsetfillcolor{textcolor}%
\pgftext[x=0.291667in, y=0.776364in, left, base]{\color{textcolor}{\rmfamily\fontsize{7.000000}{8.400000}\selectfont\catcode`\^=\active\def^{\ifmmode\sp\else\^{}\fi}\catcode`\%=\active\def%{\%}$\mathdefault{10^{-15}}$}}%
\end{pgfscope}%
\begin{pgfscope}%
\pgfpathrectangle{\pgfqpoint{0.676315in}{0.668050in}}{\pgfqpoint{4.566060in}{1.868671in}}%
\pgfusepath{clip}%
\pgfsetrectcap%
\pgfsetroundjoin%
\pgfsetlinewidth{0.803000pt}%
\definecolor{currentstroke}{rgb}{0.690196,0.690196,0.690196}%
\pgfsetstrokecolor{currentstroke}%
\pgfsetstrokeopacity{0.300000}%
\pgfsetdash{}{0pt}%
\pgfpathmoveto{\pgfqpoint{0.676315in}{1.039603in}}%
\pgfpathlineto{\pgfqpoint{5.242374in}{1.039603in}}%
\pgfusepath{stroke}%
\end{pgfscope}%
\begin{pgfscope}%
\pgfsetbuttcap%
\pgfsetroundjoin%
\definecolor{currentfill}{rgb}{0.000000,0.000000,0.000000}%
\pgfsetfillcolor{currentfill}%
\pgfsetlinewidth{0.803000pt}%
\definecolor{currentstroke}{rgb}{0.000000,0.000000,0.000000}%
\pgfsetstrokecolor{currentstroke}%
\pgfsetdash{}{0pt}%
\pgfsys@defobject{currentmarker}{\pgfqpoint{-0.048611in}{0.000000in}}{\pgfqpoint{-0.000000in}{0.000000in}}{%
\pgfpathmoveto{\pgfqpoint{-0.000000in}{0.000000in}}%
\pgfpathlineto{\pgfqpoint{-0.048611in}{0.000000in}}%
\pgfusepath{stroke,fill}%
}%
\begin{pgfscope}%
\pgfsys@transformshift{0.676315in}{1.039603in}%
\pgfsys@useobject{currentmarker}{}%
\end{pgfscope}%
\end{pgfscope}%
\begin{pgfscope}%
\definecolor{textcolor}{rgb}{0.000000,0.000000,0.000000}%
\pgfsetstrokecolor{textcolor}%
\pgfsetfillcolor{textcolor}%
\pgftext[x=0.291667in, y=1.001478in, left, base]{\color{textcolor}{\rmfamily\fontsize{7.000000}{8.400000}\selectfont\catcode`\^=\active\def^{\ifmmode\sp\else\^{}\fi}\catcode`\%=\active\def%{\%}$\mathdefault{10^{-13}}$}}%
\end{pgfscope}%
\begin{pgfscope}%
\pgfpathrectangle{\pgfqpoint{0.676315in}{0.668050in}}{\pgfqpoint{4.566060in}{1.868671in}}%
\pgfusepath{clip}%
\pgfsetrectcap%
\pgfsetroundjoin%
\pgfsetlinewidth{0.803000pt}%
\definecolor{currentstroke}{rgb}{0.690196,0.690196,0.690196}%
\pgfsetstrokecolor{currentstroke}%
\pgfsetstrokeopacity{0.300000}%
\pgfsetdash{}{0pt}%
\pgfpathmoveto{\pgfqpoint{0.676315in}{1.264716in}}%
\pgfpathlineto{\pgfqpoint{5.242374in}{1.264716in}}%
\pgfusepath{stroke}%
\end{pgfscope}%
\begin{pgfscope}%
\pgfsetbuttcap%
\pgfsetroundjoin%
\definecolor{currentfill}{rgb}{0.000000,0.000000,0.000000}%
\pgfsetfillcolor{currentfill}%
\pgfsetlinewidth{0.803000pt}%
\definecolor{currentstroke}{rgb}{0.000000,0.000000,0.000000}%
\pgfsetstrokecolor{currentstroke}%
\pgfsetdash{}{0pt}%
\pgfsys@defobject{currentmarker}{\pgfqpoint{-0.048611in}{0.000000in}}{\pgfqpoint{-0.000000in}{0.000000in}}{%
\pgfpathmoveto{\pgfqpoint{-0.000000in}{0.000000in}}%
\pgfpathlineto{\pgfqpoint{-0.048611in}{0.000000in}}%
\pgfusepath{stroke,fill}%
}%
\begin{pgfscope}%
\pgfsys@transformshift{0.676315in}{1.264716in}%
\pgfsys@useobject{currentmarker}{}%
\end{pgfscope}%
\end{pgfscope}%
\begin{pgfscope}%
\definecolor{textcolor}{rgb}{0.000000,0.000000,0.000000}%
\pgfsetstrokecolor{textcolor}%
\pgfsetfillcolor{textcolor}%
\pgftext[x=0.291667in, y=1.226591in, left, base]{\color{textcolor}{\rmfamily\fontsize{7.000000}{8.400000}\selectfont\catcode`\^=\active\def^{\ifmmode\sp\else\^{}\fi}\catcode`\%=\active\def%{\%}$\mathdefault{10^{-11}}$}}%
\end{pgfscope}%
\begin{pgfscope}%
\pgfpathrectangle{\pgfqpoint{0.676315in}{0.668050in}}{\pgfqpoint{4.566060in}{1.868671in}}%
\pgfusepath{clip}%
\pgfsetrectcap%
\pgfsetroundjoin%
\pgfsetlinewidth{0.803000pt}%
\definecolor{currentstroke}{rgb}{0.690196,0.690196,0.690196}%
\pgfsetstrokecolor{currentstroke}%
\pgfsetstrokeopacity{0.300000}%
\pgfsetdash{}{0pt}%
\pgfpathmoveto{\pgfqpoint{0.676315in}{1.489829in}}%
\pgfpathlineto{\pgfqpoint{5.242374in}{1.489829in}}%
\pgfusepath{stroke}%
\end{pgfscope}%
\begin{pgfscope}%
\pgfsetbuttcap%
\pgfsetroundjoin%
\definecolor{currentfill}{rgb}{0.000000,0.000000,0.000000}%
\pgfsetfillcolor{currentfill}%
\pgfsetlinewidth{0.803000pt}%
\definecolor{currentstroke}{rgb}{0.000000,0.000000,0.000000}%
\pgfsetstrokecolor{currentstroke}%
\pgfsetdash{}{0pt}%
\pgfsys@defobject{currentmarker}{\pgfqpoint{-0.048611in}{0.000000in}}{\pgfqpoint{-0.000000in}{0.000000in}}{%
\pgfpathmoveto{\pgfqpoint{-0.000000in}{0.000000in}}%
\pgfpathlineto{\pgfqpoint{-0.048611in}{0.000000in}}%
\pgfusepath{stroke,fill}%
}%
\begin{pgfscope}%
\pgfsys@transformshift{0.676315in}{1.489829in}%
\pgfsys@useobject{currentmarker}{}%
\end{pgfscope}%
\end{pgfscope}%
\begin{pgfscope}%
\definecolor{textcolor}{rgb}{0.000000,0.000000,0.000000}%
\pgfsetstrokecolor{textcolor}%
\pgfsetfillcolor{textcolor}%
\pgftext[x=0.338928in, y=1.451704in, left, base]{\color{textcolor}{\rmfamily\fontsize{7.000000}{8.400000}\selectfont\catcode`\^=\active\def^{\ifmmode\sp\else\^{}\fi}\catcode`\%=\active\def%{\%}$\mathdefault{10^{-9}}$}}%
\end{pgfscope}%
\begin{pgfscope}%
\pgfpathrectangle{\pgfqpoint{0.676315in}{0.668050in}}{\pgfqpoint{4.566060in}{1.868671in}}%
\pgfusepath{clip}%
\pgfsetrectcap%
\pgfsetroundjoin%
\pgfsetlinewidth{0.803000pt}%
\definecolor{currentstroke}{rgb}{0.690196,0.690196,0.690196}%
\pgfsetstrokecolor{currentstroke}%
\pgfsetstrokeopacity{0.300000}%
\pgfsetdash{}{0pt}%
\pgfpathmoveto{\pgfqpoint{0.676315in}{1.714942in}}%
\pgfpathlineto{\pgfqpoint{5.242374in}{1.714942in}}%
\pgfusepath{stroke}%
\end{pgfscope}%
\begin{pgfscope}%
\pgfsetbuttcap%
\pgfsetroundjoin%
\definecolor{currentfill}{rgb}{0.000000,0.000000,0.000000}%
\pgfsetfillcolor{currentfill}%
\pgfsetlinewidth{0.803000pt}%
\definecolor{currentstroke}{rgb}{0.000000,0.000000,0.000000}%
\pgfsetstrokecolor{currentstroke}%
\pgfsetdash{}{0pt}%
\pgfsys@defobject{currentmarker}{\pgfqpoint{-0.048611in}{0.000000in}}{\pgfqpoint{-0.000000in}{0.000000in}}{%
\pgfpathmoveto{\pgfqpoint{-0.000000in}{0.000000in}}%
\pgfpathlineto{\pgfqpoint{-0.048611in}{0.000000in}}%
\pgfusepath{stroke,fill}%
}%
\begin{pgfscope}%
\pgfsys@transformshift{0.676315in}{1.714942in}%
\pgfsys@useobject{currentmarker}{}%
\end{pgfscope}%
\end{pgfscope}%
\begin{pgfscope}%
\definecolor{textcolor}{rgb}{0.000000,0.000000,0.000000}%
\pgfsetstrokecolor{textcolor}%
\pgfsetfillcolor{textcolor}%
\pgftext[x=0.338928in, y=1.676817in, left, base]{\color{textcolor}{\rmfamily\fontsize{7.000000}{8.400000}\selectfont\catcode`\^=\active\def^{\ifmmode\sp\else\^{}\fi}\catcode`\%=\active\def%{\%}$\mathdefault{10^{-7}}$}}%
\end{pgfscope}%
\begin{pgfscope}%
\pgfpathrectangle{\pgfqpoint{0.676315in}{0.668050in}}{\pgfqpoint{4.566060in}{1.868671in}}%
\pgfusepath{clip}%
\pgfsetrectcap%
\pgfsetroundjoin%
\pgfsetlinewidth{0.803000pt}%
\definecolor{currentstroke}{rgb}{0.690196,0.690196,0.690196}%
\pgfsetstrokecolor{currentstroke}%
\pgfsetstrokeopacity{0.300000}%
\pgfsetdash{}{0pt}%
\pgfpathmoveto{\pgfqpoint{0.676315in}{1.940056in}}%
\pgfpathlineto{\pgfqpoint{5.242374in}{1.940056in}}%
\pgfusepath{stroke}%
\end{pgfscope}%
\begin{pgfscope}%
\pgfsetbuttcap%
\pgfsetroundjoin%
\definecolor{currentfill}{rgb}{0.000000,0.000000,0.000000}%
\pgfsetfillcolor{currentfill}%
\pgfsetlinewidth{0.803000pt}%
\definecolor{currentstroke}{rgb}{0.000000,0.000000,0.000000}%
\pgfsetstrokecolor{currentstroke}%
\pgfsetdash{}{0pt}%
\pgfsys@defobject{currentmarker}{\pgfqpoint{-0.048611in}{0.000000in}}{\pgfqpoint{-0.000000in}{0.000000in}}{%
\pgfpathmoveto{\pgfqpoint{-0.000000in}{0.000000in}}%
\pgfpathlineto{\pgfqpoint{-0.048611in}{0.000000in}}%
\pgfusepath{stroke,fill}%
}%
\begin{pgfscope}%
\pgfsys@transformshift{0.676315in}{1.940056in}%
\pgfsys@useobject{currentmarker}{}%
\end{pgfscope}%
\end{pgfscope}%
\begin{pgfscope}%
\definecolor{textcolor}{rgb}{0.000000,0.000000,0.000000}%
\pgfsetstrokecolor{textcolor}%
\pgfsetfillcolor{textcolor}%
\pgftext[x=0.338928in, y=1.901930in, left, base]{\color{textcolor}{\rmfamily\fontsize{7.000000}{8.400000}\selectfont\catcode`\^=\active\def^{\ifmmode\sp\else\^{}\fi}\catcode`\%=\active\def%{\%}$\mathdefault{10^{-5}}$}}%
\end{pgfscope}%
\begin{pgfscope}%
\pgfpathrectangle{\pgfqpoint{0.676315in}{0.668050in}}{\pgfqpoint{4.566060in}{1.868671in}}%
\pgfusepath{clip}%
\pgfsetrectcap%
\pgfsetroundjoin%
\pgfsetlinewidth{0.803000pt}%
\definecolor{currentstroke}{rgb}{0.690196,0.690196,0.690196}%
\pgfsetstrokecolor{currentstroke}%
\pgfsetstrokeopacity{0.300000}%
\pgfsetdash{}{0pt}%
\pgfpathmoveto{\pgfqpoint{0.676315in}{2.165169in}}%
\pgfpathlineto{\pgfqpoint{5.242374in}{2.165169in}}%
\pgfusepath{stroke}%
\end{pgfscope}%
\begin{pgfscope}%
\pgfsetbuttcap%
\pgfsetroundjoin%
\definecolor{currentfill}{rgb}{0.000000,0.000000,0.000000}%
\pgfsetfillcolor{currentfill}%
\pgfsetlinewidth{0.803000pt}%
\definecolor{currentstroke}{rgb}{0.000000,0.000000,0.000000}%
\pgfsetstrokecolor{currentstroke}%
\pgfsetdash{}{0pt}%
\pgfsys@defobject{currentmarker}{\pgfqpoint{-0.048611in}{0.000000in}}{\pgfqpoint{-0.000000in}{0.000000in}}{%
\pgfpathmoveto{\pgfqpoint{-0.000000in}{0.000000in}}%
\pgfpathlineto{\pgfqpoint{-0.048611in}{0.000000in}}%
\pgfusepath{stroke,fill}%
}%
\begin{pgfscope}%
\pgfsys@transformshift{0.676315in}{2.165169in}%
\pgfsys@useobject{currentmarker}{}%
\end{pgfscope}%
\end{pgfscope}%
\begin{pgfscope}%
\definecolor{textcolor}{rgb}{0.000000,0.000000,0.000000}%
\pgfsetstrokecolor{textcolor}%
\pgfsetfillcolor{textcolor}%
\pgftext[x=0.338928in, y=2.127043in, left, base]{\color{textcolor}{\rmfamily\fontsize{7.000000}{8.400000}\selectfont\catcode`\^=\active\def^{\ifmmode\sp\else\^{}\fi}\catcode`\%=\active\def%{\%}$\mathdefault{10^{-3}}$}}%
\end{pgfscope}%
\begin{pgfscope}%
\pgfpathrectangle{\pgfqpoint{0.676315in}{0.668050in}}{\pgfqpoint{4.566060in}{1.868671in}}%
\pgfusepath{clip}%
\pgfsetrectcap%
\pgfsetroundjoin%
\pgfsetlinewidth{0.803000pt}%
\definecolor{currentstroke}{rgb}{0.690196,0.690196,0.690196}%
\pgfsetstrokecolor{currentstroke}%
\pgfsetstrokeopacity{0.300000}%
\pgfsetdash{}{0pt}%
\pgfpathmoveto{\pgfqpoint{0.676315in}{2.390282in}}%
\pgfpathlineto{\pgfqpoint{5.242374in}{2.390282in}}%
\pgfusepath{stroke}%
\end{pgfscope}%
\begin{pgfscope}%
\pgfsetbuttcap%
\pgfsetroundjoin%
\definecolor{currentfill}{rgb}{0.000000,0.000000,0.000000}%
\pgfsetfillcolor{currentfill}%
\pgfsetlinewidth{0.803000pt}%
\definecolor{currentstroke}{rgb}{0.000000,0.000000,0.000000}%
\pgfsetstrokecolor{currentstroke}%
\pgfsetdash{}{0pt}%
\pgfsys@defobject{currentmarker}{\pgfqpoint{-0.048611in}{0.000000in}}{\pgfqpoint{-0.000000in}{0.000000in}}{%
\pgfpathmoveto{\pgfqpoint{-0.000000in}{0.000000in}}%
\pgfpathlineto{\pgfqpoint{-0.048611in}{0.000000in}}%
\pgfusepath{stroke,fill}%
}%
\begin{pgfscope}%
\pgfsys@transformshift{0.676315in}{2.390282in}%
\pgfsys@useobject{currentmarker}{}%
\end{pgfscope}%
\end{pgfscope}%
\begin{pgfscope}%
\definecolor{textcolor}{rgb}{0.000000,0.000000,0.000000}%
\pgfsetstrokecolor{textcolor}%
\pgfsetfillcolor{textcolor}%
\pgftext[x=0.338928in, y=2.352157in, left, base]{\color{textcolor}{\rmfamily\fontsize{7.000000}{8.400000}\selectfont\catcode`\^=\active\def^{\ifmmode\sp\else\^{}\fi}\catcode`\%=\active\def%{\%}$\mathdefault{10^{-1}}$}}%
\end{pgfscope}%
\begin{pgfscope}%
\definecolor{textcolor}{rgb}{0.000000,0.000000,0.000000}%
\pgfsetstrokecolor{textcolor}%
\pgfsetfillcolor{textcolor}%
\pgftext[x=0.236111in,y=1.602386in,,bottom,rotate=90.000000]{\color{textcolor}{\rmfamily\fontsize{8.000000}{9.600000}\selectfont\catcode`\^=\active\def^{\ifmmode\sp\else\^{}\fi}\catcode`\%=\active\def%{\%}Relative Error in $\widehat A_w(t)$}}%
\end{pgfscope}%
\begin{pgfscope}%
\pgfpathrectangle{\pgfqpoint{0.676315in}{0.668050in}}{\pgfqpoint{4.566060in}{1.868671in}}%
\pgfusepath{clip}%
\pgfsetrectcap%
\pgfsetroundjoin%
\pgfsetlinewidth{1.505625pt}%
\definecolor{currentstroke}{rgb}{0.121569,0.466667,0.705882}%
\pgfsetstrokecolor{currentstroke}%
\pgfsetdash{}{0pt}%
\pgfpathmoveto{\pgfqpoint{0.883863in}{2.487442in}}%
\pgfpathlineto{\pgfqpoint{0.913942in}{2.477135in}}%
\pgfpathlineto{\pgfqpoint{0.944022in}{2.463505in}}%
\pgfpathlineto{\pgfqpoint{0.974101in}{2.436520in}}%
\pgfpathlineto{\pgfqpoint{1.004181in}{2.423493in}}%
\pgfpathlineto{\pgfqpoint{1.049300in}{2.390848in}}%
\pgfpathlineto{\pgfqpoint{1.124498in}{2.353280in}}%
\pgfpathlineto{\pgfqpoint{1.274896in}{2.302625in}}%
\pgfpathlineto{\pgfqpoint{1.425293in}{2.224882in}}%
\pgfpathlineto{\pgfqpoint{1.726087in}{2.033945in}}%
\pgfpathlineto{\pgfqpoint{2.026882in}{1.791628in}}%
\pgfpathlineto{\pgfqpoint{2.327676in}{1.444579in}}%
\pgfpathlineto{\pgfqpoint{2.929265in}{0.701933in}}%
\pgfpathlineto{\pgfqpoint{3.831648in}{0.701933in}}%
\pgfpathlineto{\pgfqpoint{5.034826in}{0.701933in}}%
\pgfusepath{stroke}%
\end{pgfscope}%
\begin{pgfscope}%
\pgfpathrectangle{\pgfqpoint{0.676315in}{0.668050in}}{\pgfqpoint{4.566060in}{1.868671in}}%
\pgfusepath{clip}%
\pgfsetbuttcap%
\pgfsetroundjoin%
\definecolor{currentfill}{rgb}{0.121569,0.466667,0.705882}%
\pgfsetfillcolor{currentfill}%
\pgfsetlinewidth{1.003750pt}%
\definecolor{currentstroke}{rgb}{0.121569,0.466667,0.705882}%
\pgfsetstrokecolor{currentstroke}%
\pgfsetdash{}{0pt}%
\pgfsys@defobject{currentmarker}{\pgfqpoint{-0.015278in}{-0.015278in}}{\pgfqpoint{0.015278in}{0.015278in}}{%
\pgfpathmoveto{\pgfqpoint{0.000000in}{-0.015278in}}%
\pgfpathcurveto{\pgfqpoint{0.004052in}{-0.015278in}}{\pgfqpoint{0.007938in}{-0.013668in}}{\pgfqpoint{0.010803in}{-0.010803in}}%
\pgfpathcurveto{\pgfqpoint{0.013668in}{-0.007938in}}{\pgfqpoint{0.015278in}{-0.004052in}}{\pgfqpoint{0.015278in}{0.000000in}}%
\pgfpathcurveto{\pgfqpoint{0.015278in}{0.004052in}}{\pgfqpoint{0.013668in}{0.007938in}}{\pgfqpoint{0.010803in}{0.010803in}}%
\pgfpathcurveto{\pgfqpoint{0.007938in}{0.013668in}}{\pgfqpoint{0.004052in}{0.015278in}}{\pgfqpoint{0.000000in}{0.015278in}}%
\pgfpathcurveto{\pgfqpoint{-0.004052in}{0.015278in}}{\pgfqpoint{-0.007938in}{0.013668in}}{\pgfqpoint{-0.010803in}{0.010803in}}%
\pgfpathcurveto{\pgfqpoint{-0.013668in}{0.007938in}}{\pgfqpoint{-0.015278in}{0.004052in}}{\pgfqpoint{-0.015278in}{0.000000in}}%
\pgfpathcurveto{\pgfqpoint{-0.015278in}{-0.004052in}}{\pgfqpoint{-0.013668in}{-0.007938in}}{\pgfqpoint{-0.010803in}{-0.010803in}}%
\pgfpathcurveto{\pgfqpoint{-0.007938in}{-0.013668in}}{\pgfqpoint{-0.004052in}{-0.015278in}}{\pgfqpoint{0.000000in}{-0.015278in}}%
\pgfpathlineto{\pgfqpoint{0.000000in}{-0.015278in}}%
\pgfpathclose%
\pgfusepath{stroke,fill}%
}%
\begin{pgfscope}%
\pgfsys@transformshift{0.883863in}{2.487442in}%
\pgfsys@useobject{currentmarker}{}%
\end{pgfscope}%
\begin{pgfscope}%
\pgfsys@transformshift{0.913942in}{2.477135in}%
\pgfsys@useobject{currentmarker}{}%
\end{pgfscope}%
\begin{pgfscope}%
\pgfsys@transformshift{0.944022in}{2.463505in}%
\pgfsys@useobject{currentmarker}{}%
\end{pgfscope}%
\begin{pgfscope}%
\pgfsys@transformshift{0.974101in}{2.436520in}%
\pgfsys@useobject{currentmarker}{}%
\end{pgfscope}%
\begin{pgfscope}%
\pgfsys@transformshift{1.004181in}{2.423493in}%
\pgfsys@useobject{currentmarker}{}%
\end{pgfscope}%
\begin{pgfscope}%
\pgfsys@transformshift{1.049300in}{2.390848in}%
\pgfsys@useobject{currentmarker}{}%
\end{pgfscope}%
\begin{pgfscope}%
\pgfsys@transformshift{1.124498in}{2.353280in}%
\pgfsys@useobject{currentmarker}{}%
\end{pgfscope}%
\begin{pgfscope}%
\pgfsys@transformshift{1.274896in}{2.302625in}%
\pgfsys@useobject{currentmarker}{}%
\end{pgfscope}%
\begin{pgfscope}%
\pgfsys@transformshift{1.425293in}{2.224882in}%
\pgfsys@useobject{currentmarker}{}%
\end{pgfscope}%
\begin{pgfscope}%
\pgfsys@transformshift{1.726087in}{2.033945in}%
\pgfsys@useobject{currentmarker}{}%
\end{pgfscope}%
\begin{pgfscope}%
\pgfsys@transformshift{2.026882in}{1.791628in}%
\pgfsys@useobject{currentmarker}{}%
\end{pgfscope}%
\begin{pgfscope}%
\pgfsys@transformshift{2.327676in}{1.444579in}%
\pgfsys@useobject{currentmarker}{}%
\end{pgfscope}%
\begin{pgfscope}%
\pgfsys@transformshift{2.929265in}{0.701933in}%
\pgfsys@useobject{currentmarker}{}%
\end{pgfscope}%
\begin{pgfscope}%
\pgfsys@transformshift{3.831648in}{0.701933in}%
\pgfsys@useobject{currentmarker}{}%
\end{pgfscope}%
\begin{pgfscope}%
\pgfsys@transformshift{5.034826in}{0.701933in}%
\pgfsys@useobject{currentmarker}{}%
\end{pgfscope}%
\end{pgfscope}%
\begin{pgfscope}%
\pgfpathrectangle{\pgfqpoint{0.676315in}{0.668050in}}{\pgfqpoint{4.566060in}{1.868671in}}%
\pgfusepath{clip}%
\pgfsetrectcap%
\pgfsetroundjoin%
\pgfsetlinewidth{1.505625pt}%
\definecolor{currentstroke}{rgb}{1.000000,0.498039,0.054902}%
\pgfsetstrokecolor{currentstroke}%
\pgfsetdash{}{0pt}%
\pgfpathmoveto{\pgfqpoint{0.883863in}{2.478350in}}%
\pgfpathlineto{\pgfqpoint{0.913942in}{2.462157in}}%
\pgfpathlineto{\pgfqpoint{0.944022in}{2.439397in}}%
\pgfpathlineto{\pgfqpoint{0.974101in}{2.393402in}}%
\pgfpathlineto{\pgfqpoint{1.004181in}{2.369439in}}%
\pgfpathlineto{\pgfqpoint{1.049300in}{2.300338in}}%
\pgfpathlineto{\pgfqpoint{1.124498in}{2.184174in}}%
\pgfpathlineto{\pgfqpoint{1.274896in}{1.996125in}}%
\pgfpathlineto{\pgfqpoint{1.425293in}{1.707153in}}%
\pgfpathlineto{\pgfqpoint{1.726087in}{0.943590in}}%
\pgfpathlineto{\pgfqpoint{2.026882in}{0.701933in}}%
\pgfpathlineto{\pgfqpoint{2.327676in}{0.701933in}}%
\pgfpathlineto{\pgfqpoint{2.929265in}{0.701933in}}%
\pgfpathlineto{\pgfqpoint{3.831648in}{0.701933in}}%
\pgfpathlineto{\pgfqpoint{5.034826in}{0.701933in}}%
\pgfusepath{stroke}%
\end{pgfscope}%
\begin{pgfscope}%
\pgfpathrectangle{\pgfqpoint{0.676315in}{0.668050in}}{\pgfqpoint{4.566060in}{1.868671in}}%
\pgfusepath{clip}%
\pgfsetbuttcap%
\pgfsetroundjoin%
\definecolor{currentfill}{rgb}{1.000000,0.498039,0.054902}%
\pgfsetfillcolor{currentfill}%
\pgfsetlinewidth{1.003750pt}%
\definecolor{currentstroke}{rgb}{1.000000,0.498039,0.054902}%
\pgfsetstrokecolor{currentstroke}%
\pgfsetdash{}{0pt}%
\pgfsys@defobject{currentmarker}{\pgfqpoint{-0.015278in}{-0.015278in}}{\pgfqpoint{0.015278in}{0.015278in}}{%
\pgfpathmoveto{\pgfqpoint{0.000000in}{-0.015278in}}%
\pgfpathcurveto{\pgfqpoint{0.004052in}{-0.015278in}}{\pgfqpoint{0.007938in}{-0.013668in}}{\pgfqpoint{0.010803in}{-0.010803in}}%
\pgfpathcurveto{\pgfqpoint{0.013668in}{-0.007938in}}{\pgfqpoint{0.015278in}{-0.004052in}}{\pgfqpoint{0.015278in}{0.000000in}}%
\pgfpathcurveto{\pgfqpoint{0.015278in}{0.004052in}}{\pgfqpoint{0.013668in}{0.007938in}}{\pgfqpoint{0.010803in}{0.010803in}}%
\pgfpathcurveto{\pgfqpoint{0.007938in}{0.013668in}}{\pgfqpoint{0.004052in}{0.015278in}}{\pgfqpoint{0.000000in}{0.015278in}}%
\pgfpathcurveto{\pgfqpoint{-0.004052in}{0.015278in}}{\pgfqpoint{-0.007938in}{0.013668in}}{\pgfqpoint{-0.010803in}{0.010803in}}%
\pgfpathcurveto{\pgfqpoint{-0.013668in}{0.007938in}}{\pgfqpoint{-0.015278in}{0.004052in}}{\pgfqpoint{-0.015278in}{0.000000in}}%
\pgfpathcurveto{\pgfqpoint{-0.015278in}{-0.004052in}}{\pgfqpoint{-0.013668in}{-0.007938in}}{\pgfqpoint{-0.010803in}{-0.010803in}}%
\pgfpathcurveto{\pgfqpoint{-0.007938in}{-0.013668in}}{\pgfqpoint{-0.004052in}{-0.015278in}}{\pgfqpoint{0.000000in}{-0.015278in}}%
\pgfpathlineto{\pgfqpoint{0.000000in}{-0.015278in}}%
\pgfpathclose%
\pgfusepath{stroke,fill}%
}%
\begin{pgfscope}%
\pgfsys@transformshift{0.883863in}{2.478350in}%
\pgfsys@useobject{currentmarker}{}%
\end{pgfscope}%
\begin{pgfscope}%
\pgfsys@transformshift{0.913942in}{2.462157in}%
\pgfsys@useobject{currentmarker}{}%
\end{pgfscope}%
\begin{pgfscope}%
\pgfsys@transformshift{0.944022in}{2.439397in}%
\pgfsys@useobject{currentmarker}{}%
\end{pgfscope}%
\begin{pgfscope}%
\pgfsys@transformshift{0.974101in}{2.393402in}%
\pgfsys@useobject{currentmarker}{}%
\end{pgfscope}%
\begin{pgfscope}%
\pgfsys@transformshift{1.004181in}{2.369439in}%
\pgfsys@useobject{currentmarker}{}%
\end{pgfscope}%
\begin{pgfscope}%
\pgfsys@transformshift{1.049300in}{2.300338in}%
\pgfsys@useobject{currentmarker}{}%
\end{pgfscope}%
\begin{pgfscope}%
\pgfsys@transformshift{1.124498in}{2.184174in}%
\pgfsys@useobject{currentmarker}{}%
\end{pgfscope}%
\begin{pgfscope}%
\pgfsys@transformshift{1.274896in}{1.996125in}%
\pgfsys@useobject{currentmarker}{}%
\end{pgfscope}%
\begin{pgfscope}%
\pgfsys@transformshift{1.425293in}{1.707153in}%
\pgfsys@useobject{currentmarker}{}%
\end{pgfscope}%
\begin{pgfscope}%
\pgfsys@transformshift{1.726087in}{0.943590in}%
\pgfsys@useobject{currentmarker}{}%
\end{pgfscope}%
\begin{pgfscope}%
\pgfsys@transformshift{2.026882in}{0.701933in}%
\pgfsys@useobject{currentmarker}{}%
\end{pgfscope}%
\begin{pgfscope}%
\pgfsys@transformshift{2.327676in}{0.701933in}%
\pgfsys@useobject{currentmarker}{}%
\end{pgfscope}%
\begin{pgfscope}%
\pgfsys@transformshift{2.929265in}{0.701933in}%
\pgfsys@useobject{currentmarker}{}%
\end{pgfscope}%
\begin{pgfscope}%
\pgfsys@transformshift{3.831648in}{0.701933in}%
\pgfsys@useobject{currentmarker}{}%
\end{pgfscope}%
\begin{pgfscope}%
\pgfsys@transformshift{5.034826in}{0.701933in}%
\pgfsys@useobject{currentmarker}{}%
\end{pgfscope}%
\end{pgfscope}%
\begin{pgfscope}%
\pgfpathrectangle{\pgfqpoint{0.676315in}{0.668050in}}{\pgfqpoint{4.566060in}{1.868671in}}%
\pgfusepath{clip}%
\pgfsetrectcap%
\pgfsetroundjoin%
\pgfsetlinewidth{1.505625pt}%
\definecolor{currentstroke}{rgb}{0.172549,0.627451,0.172549}%
\pgfsetstrokecolor{currentstroke}%
\pgfsetdash{}{0pt}%
\pgfpathmoveto{\pgfqpoint{0.883863in}{2.464179in}}%
\pgfpathlineto{\pgfqpoint{0.913942in}{2.439782in}}%
\pgfpathlineto{\pgfqpoint{0.944022in}{2.402930in}}%
\pgfpathlineto{\pgfqpoint{0.974101in}{2.331518in}}%
\pgfpathlineto{\pgfqpoint{1.004181in}{2.292941in}}%
\pgfpathlineto{\pgfqpoint{1.049300in}{2.179560in}}%
\pgfpathlineto{\pgfqpoint{1.124498in}{1.960641in}}%
\pgfpathlineto{\pgfqpoint{1.274896in}{1.532917in}}%
\pgfpathlineto{\pgfqpoint{1.425293in}{0.903302in}}%
\pgfpathlineto{\pgfqpoint{1.726087in}{0.701933in}}%
\pgfpathlineto{\pgfqpoint{2.026882in}{0.701933in}}%
\pgfpathlineto{\pgfqpoint{2.327676in}{0.701933in}}%
\pgfpathlineto{\pgfqpoint{2.929265in}{0.701933in}}%
\pgfpathlineto{\pgfqpoint{3.831648in}{0.701933in}}%
\pgfpathlineto{\pgfqpoint{5.034826in}{0.701933in}}%
\pgfusepath{stroke}%
\end{pgfscope}%
\begin{pgfscope}%
\pgfpathrectangle{\pgfqpoint{0.676315in}{0.668050in}}{\pgfqpoint{4.566060in}{1.868671in}}%
\pgfusepath{clip}%
\pgfsetbuttcap%
\pgfsetroundjoin%
\definecolor{currentfill}{rgb}{0.172549,0.627451,0.172549}%
\pgfsetfillcolor{currentfill}%
\pgfsetlinewidth{1.003750pt}%
\definecolor{currentstroke}{rgb}{0.172549,0.627451,0.172549}%
\pgfsetstrokecolor{currentstroke}%
\pgfsetdash{}{0pt}%
\pgfsys@defobject{currentmarker}{\pgfqpoint{-0.015278in}{-0.015278in}}{\pgfqpoint{0.015278in}{0.015278in}}{%
\pgfpathmoveto{\pgfqpoint{0.000000in}{-0.015278in}}%
\pgfpathcurveto{\pgfqpoint{0.004052in}{-0.015278in}}{\pgfqpoint{0.007938in}{-0.013668in}}{\pgfqpoint{0.010803in}{-0.010803in}}%
\pgfpathcurveto{\pgfqpoint{0.013668in}{-0.007938in}}{\pgfqpoint{0.015278in}{-0.004052in}}{\pgfqpoint{0.015278in}{0.000000in}}%
\pgfpathcurveto{\pgfqpoint{0.015278in}{0.004052in}}{\pgfqpoint{0.013668in}{0.007938in}}{\pgfqpoint{0.010803in}{0.010803in}}%
\pgfpathcurveto{\pgfqpoint{0.007938in}{0.013668in}}{\pgfqpoint{0.004052in}{0.015278in}}{\pgfqpoint{0.000000in}{0.015278in}}%
\pgfpathcurveto{\pgfqpoint{-0.004052in}{0.015278in}}{\pgfqpoint{-0.007938in}{0.013668in}}{\pgfqpoint{-0.010803in}{0.010803in}}%
\pgfpathcurveto{\pgfqpoint{-0.013668in}{0.007938in}}{\pgfqpoint{-0.015278in}{0.004052in}}{\pgfqpoint{-0.015278in}{0.000000in}}%
\pgfpathcurveto{\pgfqpoint{-0.015278in}{-0.004052in}}{\pgfqpoint{-0.013668in}{-0.007938in}}{\pgfqpoint{-0.010803in}{-0.010803in}}%
\pgfpathcurveto{\pgfqpoint{-0.007938in}{-0.013668in}}{\pgfqpoint{-0.004052in}{-0.015278in}}{\pgfqpoint{0.000000in}{-0.015278in}}%
\pgfpathlineto{\pgfqpoint{0.000000in}{-0.015278in}}%
\pgfpathclose%
\pgfusepath{stroke,fill}%
}%
\begin{pgfscope}%
\pgfsys@transformshift{0.883863in}{2.464179in}%
\pgfsys@useobject{currentmarker}{}%
\end{pgfscope}%
\begin{pgfscope}%
\pgfsys@transformshift{0.913942in}{2.439782in}%
\pgfsys@useobject{currentmarker}{}%
\end{pgfscope}%
\begin{pgfscope}%
\pgfsys@transformshift{0.944022in}{2.402930in}%
\pgfsys@useobject{currentmarker}{}%
\end{pgfscope}%
\begin{pgfscope}%
\pgfsys@transformshift{0.974101in}{2.331518in}%
\pgfsys@useobject{currentmarker}{}%
\end{pgfscope}%
\begin{pgfscope}%
\pgfsys@transformshift{1.004181in}{2.292941in}%
\pgfsys@useobject{currentmarker}{}%
\end{pgfscope}%
\begin{pgfscope}%
\pgfsys@transformshift{1.049300in}{2.179560in}%
\pgfsys@useobject{currentmarker}{}%
\end{pgfscope}%
\begin{pgfscope}%
\pgfsys@transformshift{1.124498in}{1.960641in}%
\pgfsys@useobject{currentmarker}{}%
\end{pgfscope}%
\begin{pgfscope}%
\pgfsys@transformshift{1.274896in}{1.532917in}%
\pgfsys@useobject{currentmarker}{}%
\end{pgfscope}%
\begin{pgfscope}%
\pgfsys@transformshift{1.425293in}{0.903302in}%
\pgfsys@useobject{currentmarker}{}%
\end{pgfscope}%
\begin{pgfscope}%
\pgfsys@transformshift{1.726087in}{0.701933in}%
\pgfsys@useobject{currentmarker}{}%
\end{pgfscope}%
\begin{pgfscope}%
\pgfsys@transformshift{2.026882in}{0.701933in}%
\pgfsys@useobject{currentmarker}{}%
\end{pgfscope}%
\begin{pgfscope}%
\pgfsys@transformshift{2.327676in}{0.701933in}%
\pgfsys@useobject{currentmarker}{}%
\end{pgfscope}%
\begin{pgfscope}%
\pgfsys@transformshift{2.929265in}{0.701933in}%
\pgfsys@useobject{currentmarker}{}%
\end{pgfscope}%
\begin{pgfscope}%
\pgfsys@transformshift{3.831648in}{0.701933in}%
\pgfsys@useobject{currentmarker}{}%
\end{pgfscope}%
\begin{pgfscope}%
\pgfsys@transformshift{5.034826in}{0.701933in}%
\pgfsys@useobject{currentmarker}{}%
\end{pgfscope}%
\end{pgfscope}%
\begin{pgfscope}%
\pgfpathrectangle{\pgfqpoint{0.676315in}{0.668050in}}{\pgfqpoint{4.566060in}{1.868671in}}%
\pgfusepath{clip}%
\pgfsetrectcap%
\pgfsetroundjoin%
\pgfsetlinewidth{1.505625pt}%
\definecolor{currentstroke}{rgb}{0.839216,0.152941,0.156863}%
\pgfsetstrokecolor{currentstroke}%
\pgfsetdash{}{0pt}%
\pgfpathmoveto{\pgfqpoint{0.883863in}{2.381783in}}%
\pgfpathlineto{\pgfqpoint{0.913942in}{2.321513in}}%
\pgfpathlineto{\pgfqpoint{0.944022in}{2.197721in}}%
\pgfpathlineto{\pgfqpoint{0.974101in}{2.006653in}}%
\pgfpathlineto{\pgfqpoint{1.004181in}{1.871043in}}%
\pgfpathlineto{\pgfqpoint{1.049300in}{1.529141in}}%
\pgfpathlineto{\pgfqpoint{1.124498in}{0.818129in}}%
\pgfpathlineto{\pgfqpoint{1.274896in}{0.750363in}}%
\pgfpathlineto{\pgfqpoint{1.425293in}{0.701933in}}%
\pgfpathlineto{\pgfqpoint{1.726087in}{0.701933in}}%
\pgfpathlineto{\pgfqpoint{2.026882in}{0.701933in}}%
\pgfpathlineto{\pgfqpoint{2.327676in}{0.701933in}}%
\pgfpathlineto{\pgfqpoint{2.929265in}{0.701933in}}%
\pgfpathlineto{\pgfqpoint{3.831648in}{0.701933in}}%
\pgfpathlineto{\pgfqpoint{5.034826in}{0.701933in}}%
\pgfusepath{stroke}%
\end{pgfscope}%
\begin{pgfscope}%
\pgfpathrectangle{\pgfqpoint{0.676315in}{0.668050in}}{\pgfqpoint{4.566060in}{1.868671in}}%
\pgfusepath{clip}%
\pgfsetbuttcap%
\pgfsetroundjoin%
\definecolor{currentfill}{rgb}{0.839216,0.152941,0.156863}%
\pgfsetfillcolor{currentfill}%
\pgfsetlinewidth{1.003750pt}%
\definecolor{currentstroke}{rgb}{0.839216,0.152941,0.156863}%
\pgfsetstrokecolor{currentstroke}%
\pgfsetdash{}{0pt}%
\pgfsys@defobject{currentmarker}{\pgfqpoint{-0.015278in}{-0.015278in}}{\pgfqpoint{0.015278in}{0.015278in}}{%
\pgfpathmoveto{\pgfqpoint{0.000000in}{-0.015278in}}%
\pgfpathcurveto{\pgfqpoint{0.004052in}{-0.015278in}}{\pgfqpoint{0.007938in}{-0.013668in}}{\pgfqpoint{0.010803in}{-0.010803in}}%
\pgfpathcurveto{\pgfqpoint{0.013668in}{-0.007938in}}{\pgfqpoint{0.015278in}{-0.004052in}}{\pgfqpoint{0.015278in}{0.000000in}}%
\pgfpathcurveto{\pgfqpoint{0.015278in}{0.004052in}}{\pgfqpoint{0.013668in}{0.007938in}}{\pgfqpoint{0.010803in}{0.010803in}}%
\pgfpathcurveto{\pgfqpoint{0.007938in}{0.013668in}}{\pgfqpoint{0.004052in}{0.015278in}}{\pgfqpoint{0.000000in}{0.015278in}}%
\pgfpathcurveto{\pgfqpoint{-0.004052in}{0.015278in}}{\pgfqpoint{-0.007938in}{0.013668in}}{\pgfqpoint{-0.010803in}{0.010803in}}%
\pgfpathcurveto{\pgfqpoint{-0.013668in}{0.007938in}}{\pgfqpoint{-0.015278in}{0.004052in}}{\pgfqpoint{-0.015278in}{0.000000in}}%
\pgfpathcurveto{\pgfqpoint{-0.015278in}{-0.004052in}}{\pgfqpoint{-0.013668in}{-0.007938in}}{\pgfqpoint{-0.010803in}{-0.010803in}}%
\pgfpathcurveto{\pgfqpoint{-0.007938in}{-0.013668in}}{\pgfqpoint{-0.004052in}{-0.015278in}}{\pgfqpoint{0.000000in}{-0.015278in}}%
\pgfpathlineto{\pgfqpoint{0.000000in}{-0.015278in}}%
\pgfpathclose%
\pgfusepath{stroke,fill}%
}%
\begin{pgfscope}%
\pgfsys@transformshift{0.883863in}{2.381783in}%
\pgfsys@useobject{currentmarker}{}%
\end{pgfscope}%
\begin{pgfscope}%
\pgfsys@transformshift{0.913942in}{2.321513in}%
\pgfsys@useobject{currentmarker}{}%
\end{pgfscope}%
\begin{pgfscope}%
\pgfsys@transformshift{0.944022in}{2.197721in}%
\pgfsys@useobject{currentmarker}{}%
\end{pgfscope}%
\begin{pgfscope}%
\pgfsys@transformshift{0.974101in}{2.006653in}%
\pgfsys@useobject{currentmarker}{}%
\end{pgfscope}%
\begin{pgfscope}%
\pgfsys@transformshift{1.004181in}{1.871043in}%
\pgfsys@useobject{currentmarker}{}%
\end{pgfscope}%
\begin{pgfscope}%
\pgfsys@transformshift{1.049300in}{1.529141in}%
\pgfsys@useobject{currentmarker}{}%
\end{pgfscope}%
\begin{pgfscope}%
\pgfsys@transformshift{1.124498in}{0.818129in}%
\pgfsys@useobject{currentmarker}{}%
\end{pgfscope}%
\begin{pgfscope}%
\pgfsys@transformshift{1.274896in}{0.750363in}%
\pgfsys@useobject{currentmarker}{}%
\end{pgfscope}%
\begin{pgfscope}%
\pgfsys@transformshift{1.425293in}{0.701933in}%
\pgfsys@useobject{currentmarker}{}%
\end{pgfscope}%
\begin{pgfscope}%
\pgfsys@transformshift{1.726087in}{0.701933in}%
\pgfsys@useobject{currentmarker}{}%
\end{pgfscope}%
\begin{pgfscope}%
\pgfsys@transformshift{2.026882in}{0.701933in}%
\pgfsys@useobject{currentmarker}{}%
\end{pgfscope}%
\begin{pgfscope}%
\pgfsys@transformshift{2.327676in}{0.701933in}%
\pgfsys@useobject{currentmarker}{}%
\end{pgfscope}%
\begin{pgfscope}%
\pgfsys@transformshift{2.929265in}{0.701933in}%
\pgfsys@useobject{currentmarker}{}%
\end{pgfscope}%
\begin{pgfscope}%
\pgfsys@transformshift{3.831648in}{0.701933in}%
\pgfsys@useobject{currentmarker}{}%
\end{pgfscope}%
\begin{pgfscope}%
\pgfsys@transformshift{5.034826in}{0.701933in}%
\pgfsys@useobject{currentmarker}{}%
\end{pgfscope}%
\end{pgfscope}%
\begin{pgfscope}%
\pgfpathrectangle{\pgfqpoint{0.676315in}{0.668050in}}{\pgfqpoint{4.566060in}{1.868671in}}%
\pgfusepath{clip}%
\pgfsetrectcap%
\pgfsetroundjoin%
\pgfsetlinewidth{1.505625pt}%
\definecolor{currentstroke}{rgb}{0.580392,0.403922,0.741176}%
\pgfsetstrokecolor{currentstroke}%
\pgfsetdash{}{0pt}%
\pgfpathmoveto{\pgfqpoint{0.883863in}{1.998034in}}%
\pgfpathlineto{\pgfqpoint{0.913942in}{1.804746in}}%
\pgfpathlineto{\pgfqpoint{0.944022in}{1.329405in}}%
\pgfpathlineto{\pgfqpoint{0.974101in}{0.768495in}}%
\pgfpathlineto{\pgfqpoint{1.004181in}{0.734612in}}%
\pgfpathlineto{\pgfqpoint{1.049300in}{0.701933in}}%
\pgfpathlineto{\pgfqpoint{1.124498in}{0.701933in}}%
\pgfpathlineto{\pgfqpoint{1.274896in}{0.701933in}}%
\pgfpathlineto{\pgfqpoint{1.425293in}{0.701933in}}%
\pgfpathlineto{\pgfqpoint{1.726087in}{0.701933in}}%
\pgfpathlineto{\pgfqpoint{2.026882in}{0.701933in}}%
\pgfpathlineto{\pgfqpoint{2.327676in}{0.701933in}}%
\pgfpathlineto{\pgfqpoint{2.929265in}{0.701933in}}%
\pgfpathlineto{\pgfqpoint{3.831648in}{0.701933in}}%
\pgfpathlineto{\pgfqpoint{5.034826in}{0.701933in}}%
\pgfusepath{stroke}%
\end{pgfscope}%
\begin{pgfscope}%
\pgfpathrectangle{\pgfqpoint{0.676315in}{0.668050in}}{\pgfqpoint{4.566060in}{1.868671in}}%
\pgfusepath{clip}%
\pgfsetbuttcap%
\pgfsetroundjoin%
\definecolor{currentfill}{rgb}{0.580392,0.403922,0.741176}%
\pgfsetfillcolor{currentfill}%
\pgfsetlinewidth{1.003750pt}%
\definecolor{currentstroke}{rgb}{0.580392,0.403922,0.741176}%
\pgfsetstrokecolor{currentstroke}%
\pgfsetdash{}{0pt}%
\pgfsys@defobject{currentmarker}{\pgfqpoint{-0.015278in}{-0.015278in}}{\pgfqpoint{0.015278in}{0.015278in}}{%
\pgfpathmoveto{\pgfqpoint{0.000000in}{-0.015278in}}%
\pgfpathcurveto{\pgfqpoint{0.004052in}{-0.015278in}}{\pgfqpoint{0.007938in}{-0.013668in}}{\pgfqpoint{0.010803in}{-0.010803in}}%
\pgfpathcurveto{\pgfqpoint{0.013668in}{-0.007938in}}{\pgfqpoint{0.015278in}{-0.004052in}}{\pgfqpoint{0.015278in}{0.000000in}}%
\pgfpathcurveto{\pgfqpoint{0.015278in}{0.004052in}}{\pgfqpoint{0.013668in}{0.007938in}}{\pgfqpoint{0.010803in}{0.010803in}}%
\pgfpathcurveto{\pgfqpoint{0.007938in}{0.013668in}}{\pgfqpoint{0.004052in}{0.015278in}}{\pgfqpoint{0.000000in}{0.015278in}}%
\pgfpathcurveto{\pgfqpoint{-0.004052in}{0.015278in}}{\pgfqpoint{-0.007938in}{0.013668in}}{\pgfqpoint{-0.010803in}{0.010803in}}%
\pgfpathcurveto{\pgfqpoint{-0.013668in}{0.007938in}}{\pgfqpoint{-0.015278in}{0.004052in}}{\pgfqpoint{-0.015278in}{0.000000in}}%
\pgfpathcurveto{\pgfqpoint{-0.015278in}{-0.004052in}}{\pgfqpoint{-0.013668in}{-0.007938in}}{\pgfqpoint{-0.010803in}{-0.010803in}}%
\pgfpathcurveto{\pgfqpoint{-0.007938in}{-0.013668in}}{\pgfqpoint{-0.004052in}{-0.015278in}}{\pgfqpoint{0.000000in}{-0.015278in}}%
\pgfpathlineto{\pgfqpoint{0.000000in}{-0.015278in}}%
\pgfpathclose%
\pgfusepath{stroke,fill}%
}%
\begin{pgfscope}%
\pgfsys@transformshift{0.883863in}{1.998034in}%
\pgfsys@useobject{currentmarker}{}%
\end{pgfscope}%
\begin{pgfscope}%
\pgfsys@transformshift{0.913942in}{1.804746in}%
\pgfsys@useobject{currentmarker}{}%
\end{pgfscope}%
\begin{pgfscope}%
\pgfsys@transformshift{0.944022in}{1.329405in}%
\pgfsys@useobject{currentmarker}{}%
\end{pgfscope}%
\begin{pgfscope}%
\pgfsys@transformshift{0.974101in}{0.768495in}%
\pgfsys@useobject{currentmarker}{}%
\end{pgfscope}%
\begin{pgfscope}%
\pgfsys@transformshift{1.004181in}{0.734612in}%
\pgfsys@useobject{currentmarker}{}%
\end{pgfscope}%
\begin{pgfscope}%
\pgfsys@transformshift{1.049300in}{0.701933in}%
\pgfsys@useobject{currentmarker}{}%
\end{pgfscope}%
\begin{pgfscope}%
\pgfsys@transformshift{1.124498in}{0.701933in}%
\pgfsys@useobject{currentmarker}{}%
\end{pgfscope}%
\begin{pgfscope}%
\pgfsys@transformshift{1.274896in}{0.701933in}%
\pgfsys@useobject{currentmarker}{}%
\end{pgfscope}%
\begin{pgfscope}%
\pgfsys@transformshift{1.425293in}{0.701933in}%
\pgfsys@useobject{currentmarker}{}%
\end{pgfscope}%
\begin{pgfscope}%
\pgfsys@transformshift{1.726087in}{0.701933in}%
\pgfsys@useobject{currentmarker}{}%
\end{pgfscope}%
\begin{pgfscope}%
\pgfsys@transformshift{2.026882in}{0.701933in}%
\pgfsys@useobject{currentmarker}{}%
\end{pgfscope}%
\begin{pgfscope}%
\pgfsys@transformshift{2.327676in}{0.701933in}%
\pgfsys@useobject{currentmarker}{}%
\end{pgfscope}%
\begin{pgfscope}%
\pgfsys@transformshift{2.929265in}{0.701933in}%
\pgfsys@useobject{currentmarker}{}%
\end{pgfscope}%
\begin{pgfscope}%
\pgfsys@transformshift{3.831648in}{0.701933in}%
\pgfsys@useobject{currentmarker}{}%
\end{pgfscope}%
\begin{pgfscope}%
\pgfsys@transformshift{5.034826in}{0.701933in}%
\pgfsys@useobject{currentmarker}{}%
\end{pgfscope}%
\end{pgfscope}%
\begin{pgfscope}%
\pgfsetrectcap%
\pgfsetmiterjoin%
\pgfsetlinewidth{0.803000pt}%
\definecolor{currentstroke}{rgb}{0.000000,0.000000,0.000000}%
\pgfsetstrokecolor{currentstroke}%
\pgfsetdash{}{0pt}%
\pgfpathmoveto{\pgfqpoint{0.676315in}{0.668050in}}%
\pgfpathlineto{\pgfqpoint{0.676315in}{2.536721in}}%
\pgfusepath{stroke}%
\end{pgfscope}%
\begin{pgfscope}%
\pgfsetrectcap%
\pgfsetmiterjoin%
\pgfsetlinewidth{0.803000pt}%
\definecolor{currentstroke}{rgb}{0.000000,0.000000,0.000000}%
\pgfsetstrokecolor{currentstroke}%
\pgfsetdash{}{0pt}%
\pgfpathmoveto{\pgfqpoint{5.242374in}{0.668050in}}%
\pgfpathlineto{\pgfqpoint{5.242374in}{2.536721in}}%
\pgfusepath{stroke}%
\end{pgfscope}%
\begin{pgfscope}%
\pgfsetrectcap%
\pgfsetmiterjoin%
\pgfsetlinewidth{0.803000pt}%
\definecolor{currentstroke}{rgb}{0.000000,0.000000,0.000000}%
\pgfsetstrokecolor{currentstroke}%
\pgfsetdash{}{0pt}%
\pgfpathmoveto{\pgfqpoint{0.676315in}{0.668050in}}%
\pgfpathlineto{\pgfqpoint{5.242374in}{0.668050in}}%
\pgfusepath{stroke}%
\end{pgfscope}%
\begin{pgfscope}%
\pgfsetrectcap%
\pgfsetmiterjoin%
\pgfsetlinewidth{0.803000pt}%
\definecolor{currentstroke}{rgb}{0.000000,0.000000,0.000000}%
\pgfsetstrokecolor{currentstroke}%
\pgfsetdash{}{0pt}%
\pgfpathmoveto{\pgfqpoint{0.676315in}{2.536721in}}%
\pgfpathlineto{\pgfqpoint{5.242374in}{2.536721in}}%
\pgfusepath{stroke}%
\end{pgfscope}%
\begin{pgfscope}%
\definecolor{textcolor}{rgb}{0.000000,0.000000,0.000000}%
\pgfsetstrokecolor{textcolor}%
\pgfsetfillcolor{textcolor}%
\pgftext[x=2.959345in,y=2.620055in,,base]{\color{textcolor}{\rmfamily\fontsize{8.000000}{9.600000}\selectfont\catcode`\^=\active\def^{\ifmmode\sp\else\^{}\fi}\catcode`\%=\active\def%{\%}Truncation Sensitivity of the $S^2$ Heat Series}}%
\end{pgfscope}%
\begin{pgfscope}%
\pgfsetrectcap%
\pgfsetroundjoin%
\pgfsetlinewidth{1.505625pt}%
\definecolor{currentstroke}{rgb}{0.121569,0.466667,0.705882}%
\pgfsetstrokecolor{currentstroke}%
\pgfsetdash{}{0pt}%
\pgfpathmoveto{\pgfqpoint{1.254978in}{0.191821in}}%
\pgfpathlineto{\pgfqpoint{1.332756in}{0.191821in}}%
\pgfpathlineto{\pgfqpoint{1.410534in}{0.191821in}}%
\pgfusepath{stroke}%
\end{pgfscope}%
\begin{pgfscope}%
\pgfsetbuttcap%
\pgfsetroundjoin%
\definecolor{currentfill}{rgb}{0.121569,0.466667,0.705882}%
\pgfsetfillcolor{currentfill}%
\pgfsetlinewidth{1.003750pt}%
\definecolor{currentstroke}{rgb}{0.121569,0.466667,0.705882}%
\pgfsetstrokecolor{currentstroke}%
\pgfsetdash{}{0pt}%
\pgfsys@defobject{currentmarker}{\pgfqpoint{-0.015278in}{-0.015278in}}{\pgfqpoint{0.015278in}{0.015278in}}{%
\pgfpathmoveto{\pgfqpoint{0.000000in}{-0.015278in}}%
\pgfpathcurveto{\pgfqpoint{0.004052in}{-0.015278in}}{\pgfqpoint{0.007938in}{-0.013668in}}{\pgfqpoint{0.010803in}{-0.010803in}}%
\pgfpathcurveto{\pgfqpoint{0.013668in}{-0.007938in}}{\pgfqpoint{0.015278in}{-0.004052in}}{\pgfqpoint{0.015278in}{0.000000in}}%
\pgfpathcurveto{\pgfqpoint{0.015278in}{0.004052in}}{\pgfqpoint{0.013668in}{0.007938in}}{\pgfqpoint{0.010803in}{0.010803in}}%
\pgfpathcurveto{\pgfqpoint{0.007938in}{0.013668in}}{\pgfqpoint{0.004052in}{0.015278in}}{\pgfqpoint{0.000000in}{0.015278in}}%
\pgfpathcurveto{\pgfqpoint{-0.004052in}{0.015278in}}{\pgfqpoint{-0.007938in}{0.013668in}}{\pgfqpoint{-0.010803in}{0.010803in}}%
\pgfpathcurveto{\pgfqpoint{-0.013668in}{0.007938in}}{\pgfqpoint{-0.015278in}{0.004052in}}{\pgfqpoint{-0.015278in}{0.000000in}}%
\pgfpathcurveto{\pgfqpoint{-0.015278in}{-0.004052in}}{\pgfqpoint{-0.013668in}{-0.007938in}}{\pgfqpoint{-0.010803in}{-0.010803in}}%
\pgfpathcurveto{\pgfqpoint{-0.007938in}{-0.013668in}}{\pgfqpoint{-0.004052in}{-0.015278in}}{\pgfqpoint{0.000000in}{-0.015278in}}%
\pgfpathlineto{\pgfqpoint{0.000000in}{-0.015278in}}%
\pgfpathclose%
\pgfusepath{stroke,fill}%
}%
\begin{pgfscope}%
\pgfsys@transformshift{1.332756in}{0.191821in}%
\pgfsys@useobject{currentmarker}{}%
\end{pgfscope}%
\end{pgfscope}%
\begin{pgfscope}%
\definecolor{textcolor}{rgb}{0.000000,0.000000,0.000000}%
\pgfsetstrokecolor{textcolor}%
\pgfsetfillcolor{textcolor}%
\pgftext[x=1.488311in,y=0.157793in,left,base]{\color{textcolor}{\rmfamily\fontsize{7.000000}{8.400000}\selectfont\catcode`\^=\active\def^{\ifmmode\sp\else\^{}\fi}\catcode`\%=\active\def%{\%}$\alpha=5^\circ$}}%
\end{pgfscope}%
\begin{pgfscope}%
\pgfsetrectcap%
\pgfsetroundjoin%
\pgfsetlinewidth{1.505625pt}%
\definecolor{currentstroke}{rgb}{1.000000,0.498039,0.054902}%
\pgfsetstrokecolor{currentstroke}%
\pgfsetdash{}{0pt}%
\pgfpathmoveto{\pgfqpoint{1.909934in}{0.191821in}}%
\pgfpathlineto{\pgfqpoint{1.987712in}{0.191821in}}%
\pgfpathlineto{\pgfqpoint{2.065490in}{0.191821in}}%
\pgfusepath{stroke}%
\end{pgfscope}%
\begin{pgfscope}%
\pgfsetbuttcap%
\pgfsetroundjoin%
\definecolor{currentfill}{rgb}{1.000000,0.498039,0.054902}%
\pgfsetfillcolor{currentfill}%
\pgfsetlinewidth{1.003750pt}%
\definecolor{currentstroke}{rgb}{1.000000,0.498039,0.054902}%
\pgfsetstrokecolor{currentstroke}%
\pgfsetdash{}{0pt}%
\pgfsys@defobject{currentmarker}{\pgfqpoint{-0.015278in}{-0.015278in}}{\pgfqpoint{0.015278in}{0.015278in}}{%
\pgfpathmoveto{\pgfqpoint{0.000000in}{-0.015278in}}%
\pgfpathcurveto{\pgfqpoint{0.004052in}{-0.015278in}}{\pgfqpoint{0.007938in}{-0.013668in}}{\pgfqpoint{0.010803in}{-0.010803in}}%
\pgfpathcurveto{\pgfqpoint{0.013668in}{-0.007938in}}{\pgfqpoint{0.015278in}{-0.004052in}}{\pgfqpoint{0.015278in}{0.000000in}}%
\pgfpathcurveto{\pgfqpoint{0.015278in}{0.004052in}}{\pgfqpoint{0.013668in}{0.007938in}}{\pgfqpoint{0.010803in}{0.010803in}}%
\pgfpathcurveto{\pgfqpoint{0.007938in}{0.013668in}}{\pgfqpoint{0.004052in}{0.015278in}}{\pgfqpoint{0.000000in}{0.015278in}}%
\pgfpathcurveto{\pgfqpoint{-0.004052in}{0.015278in}}{\pgfqpoint{-0.007938in}{0.013668in}}{\pgfqpoint{-0.010803in}{0.010803in}}%
\pgfpathcurveto{\pgfqpoint{-0.013668in}{0.007938in}}{\pgfqpoint{-0.015278in}{0.004052in}}{\pgfqpoint{-0.015278in}{0.000000in}}%
\pgfpathcurveto{\pgfqpoint{-0.015278in}{-0.004052in}}{\pgfqpoint{-0.013668in}{-0.007938in}}{\pgfqpoint{-0.010803in}{-0.010803in}}%
\pgfpathcurveto{\pgfqpoint{-0.007938in}{-0.013668in}}{\pgfqpoint{-0.004052in}{-0.015278in}}{\pgfqpoint{0.000000in}{-0.015278in}}%
\pgfpathlineto{\pgfqpoint{0.000000in}{-0.015278in}}%
\pgfpathclose%
\pgfusepath{stroke,fill}%
}%
\begin{pgfscope}%
\pgfsys@transformshift{1.987712in}{0.191821in}%
\pgfsys@useobject{currentmarker}{}%
\end{pgfscope}%
\end{pgfscope}%
\begin{pgfscope}%
\definecolor{textcolor}{rgb}{0.000000,0.000000,0.000000}%
\pgfsetstrokecolor{textcolor}%
\pgfsetfillcolor{textcolor}%
\pgftext[x=2.143268in,y=0.157793in,left,base]{\color{textcolor}{\rmfamily\fontsize{7.000000}{8.400000}\selectfont\catcode`\^=\active\def^{\ifmmode\sp\else\^{}\fi}\catcode`\%=\active\def%{\%}$\alpha=10^\circ$}}%
\end{pgfscope}%
\begin{pgfscope}%
\pgfsetrectcap%
\pgfsetroundjoin%
\pgfsetlinewidth{1.505625pt}%
\definecolor{currentstroke}{rgb}{0.172549,0.627451,0.172549}%
\pgfsetstrokecolor{currentstroke}%
\pgfsetdash{}{0pt}%
\pgfpathmoveto{\pgfqpoint{2.620253in}{0.191821in}}%
\pgfpathlineto{\pgfqpoint{2.698031in}{0.191821in}}%
\pgfpathlineto{\pgfqpoint{2.775809in}{0.191821in}}%
\pgfusepath{stroke}%
\end{pgfscope}%
\begin{pgfscope}%
\pgfsetbuttcap%
\pgfsetroundjoin%
\definecolor{currentfill}{rgb}{0.172549,0.627451,0.172549}%
\pgfsetfillcolor{currentfill}%
\pgfsetlinewidth{1.003750pt}%
\definecolor{currentstroke}{rgb}{0.172549,0.627451,0.172549}%
\pgfsetstrokecolor{currentstroke}%
\pgfsetdash{}{0pt}%
\pgfsys@defobject{currentmarker}{\pgfqpoint{-0.015278in}{-0.015278in}}{\pgfqpoint{0.015278in}{0.015278in}}{%
\pgfpathmoveto{\pgfqpoint{0.000000in}{-0.015278in}}%
\pgfpathcurveto{\pgfqpoint{0.004052in}{-0.015278in}}{\pgfqpoint{0.007938in}{-0.013668in}}{\pgfqpoint{0.010803in}{-0.010803in}}%
\pgfpathcurveto{\pgfqpoint{0.013668in}{-0.007938in}}{\pgfqpoint{0.015278in}{-0.004052in}}{\pgfqpoint{0.015278in}{0.000000in}}%
\pgfpathcurveto{\pgfqpoint{0.015278in}{0.004052in}}{\pgfqpoint{0.013668in}{0.007938in}}{\pgfqpoint{0.010803in}{0.010803in}}%
\pgfpathcurveto{\pgfqpoint{0.007938in}{0.013668in}}{\pgfqpoint{0.004052in}{0.015278in}}{\pgfqpoint{0.000000in}{0.015278in}}%
\pgfpathcurveto{\pgfqpoint{-0.004052in}{0.015278in}}{\pgfqpoint{-0.007938in}{0.013668in}}{\pgfqpoint{-0.010803in}{0.010803in}}%
\pgfpathcurveto{\pgfqpoint{-0.013668in}{0.007938in}}{\pgfqpoint{-0.015278in}{0.004052in}}{\pgfqpoint{-0.015278in}{0.000000in}}%
\pgfpathcurveto{\pgfqpoint{-0.015278in}{-0.004052in}}{\pgfqpoint{-0.013668in}{-0.007938in}}{\pgfqpoint{-0.010803in}{-0.010803in}}%
\pgfpathcurveto{\pgfqpoint{-0.007938in}{-0.013668in}}{\pgfqpoint{-0.004052in}{-0.015278in}}{\pgfqpoint{0.000000in}{-0.015278in}}%
\pgfpathlineto{\pgfqpoint{0.000000in}{-0.015278in}}%
\pgfpathclose%
\pgfusepath{stroke,fill}%
}%
\begin{pgfscope}%
\pgfsys@transformshift{2.698031in}{0.191821in}%
\pgfsys@useobject{currentmarker}{}%
\end{pgfscope}%
\end{pgfscope}%
\begin{pgfscope}%
\definecolor{textcolor}{rgb}{0.000000,0.000000,0.000000}%
\pgfsetstrokecolor{textcolor}%
\pgfsetfillcolor{textcolor}%
\pgftext[x=2.853587in,y=0.157793in,left,base]{\color{textcolor}{\rmfamily\fontsize{7.000000}{8.400000}\selectfont\catcode`\^=\active\def^{\ifmmode\sp\else\^{}\fi}\catcode`\%=\active\def%{\%}$\alpha=15^\circ$}}%
\end{pgfscope}%
\begin{pgfscope}%
\pgfsetrectcap%
\pgfsetroundjoin%
\pgfsetlinewidth{1.505625pt}%
\definecolor{currentstroke}{rgb}{0.839216,0.152941,0.156863}%
\pgfsetstrokecolor{currentstroke}%
\pgfsetdash{}{0pt}%
\pgfpathmoveto{\pgfqpoint{3.330573in}{0.191821in}}%
\pgfpathlineto{\pgfqpoint{3.408350in}{0.191821in}}%
\pgfpathlineto{\pgfqpoint{3.486128in}{0.191821in}}%
\pgfusepath{stroke}%
\end{pgfscope}%
\begin{pgfscope}%
\pgfsetbuttcap%
\pgfsetroundjoin%
\definecolor{currentfill}{rgb}{0.839216,0.152941,0.156863}%
\pgfsetfillcolor{currentfill}%
\pgfsetlinewidth{1.003750pt}%
\definecolor{currentstroke}{rgb}{0.839216,0.152941,0.156863}%
\pgfsetstrokecolor{currentstroke}%
\pgfsetdash{}{0pt}%
\pgfsys@defobject{currentmarker}{\pgfqpoint{-0.015278in}{-0.015278in}}{\pgfqpoint{0.015278in}{0.015278in}}{%
\pgfpathmoveto{\pgfqpoint{0.000000in}{-0.015278in}}%
\pgfpathcurveto{\pgfqpoint{0.004052in}{-0.015278in}}{\pgfqpoint{0.007938in}{-0.013668in}}{\pgfqpoint{0.010803in}{-0.010803in}}%
\pgfpathcurveto{\pgfqpoint{0.013668in}{-0.007938in}}{\pgfqpoint{0.015278in}{-0.004052in}}{\pgfqpoint{0.015278in}{0.000000in}}%
\pgfpathcurveto{\pgfqpoint{0.015278in}{0.004052in}}{\pgfqpoint{0.013668in}{0.007938in}}{\pgfqpoint{0.010803in}{0.010803in}}%
\pgfpathcurveto{\pgfqpoint{0.007938in}{0.013668in}}{\pgfqpoint{0.004052in}{0.015278in}}{\pgfqpoint{0.000000in}{0.015278in}}%
\pgfpathcurveto{\pgfqpoint{-0.004052in}{0.015278in}}{\pgfqpoint{-0.007938in}{0.013668in}}{\pgfqpoint{-0.010803in}{0.010803in}}%
\pgfpathcurveto{\pgfqpoint{-0.013668in}{0.007938in}}{\pgfqpoint{-0.015278in}{0.004052in}}{\pgfqpoint{-0.015278in}{0.000000in}}%
\pgfpathcurveto{\pgfqpoint{-0.015278in}{-0.004052in}}{\pgfqpoint{-0.013668in}{-0.007938in}}{\pgfqpoint{-0.010803in}{-0.010803in}}%
\pgfpathcurveto{\pgfqpoint{-0.007938in}{-0.013668in}}{\pgfqpoint{-0.004052in}{-0.015278in}}{\pgfqpoint{0.000000in}{-0.015278in}}%
\pgfpathlineto{\pgfqpoint{0.000000in}{-0.015278in}}%
\pgfpathclose%
\pgfusepath{stroke,fill}%
}%
\begin{pgfscope}%
\pgfsys@transformshift{3.408350in}{0.191821in}%
\pgfsys@useobject{currentmarker}{}%
\end{pgfscope}%
\end{pgfscope}%
\begin{pgfscope}%
\definecolor{textcolor}{rgb}{0.000000,0.000000,0.000000}%
\pgfsetstrokecolor{textcolor}%
\pgfsetfillcolor{textcolor}%
\pgftext[x=3.563906in,y=0.157793in,left,base]{\color{textcolor}{\rmfamily\fontsize{7.000000}{8.400000}\selectfont\catcode`\^=\active\def^{\ifmmode\sp\else\^{}\fi}\catcode`\%=\active\def%{\%}$\alpha=30^\circ$}}%
\end{pgfscope}%
\begin{pgfscope}%
\pgfsetrectcap%
\pgfsetroundjoin%
\pgfsetlinewidth{1.505625pt}%
\definecolor{currentstroke}{rgb}{0.580392,0.403922,0.741176}%
\pgfsetstrokecolor{currentstroke}%
\pgfsetdash{}{0pt}%
\pgfpathmoveto{\pgfqpoint{4.040892in}{0.191821in}}%
\pgfpathlineto{\pgfqpoint{4.118670in}{0.191821in}}%
\pgfpathlineto{\pgfqpoint{4.196447in}{0.191821in}}%
\pgfusepath{stroke}%
\end{pgfscope}%
\begin{pgfscope}%
\pgfsetbuttcap%
\pgfsetroundjoin%
\definecolor{currentfill}{rgb}{0.580392,0.403922,0.741176}%
\pgfsetfillcolor{currentfill}%
\pgfsetlinewidth{1.003750pt}%
\definecolor{currentstroke}{rgb}{0.580392,0.403922,0.741176}%
\pgfsetstrokecolor{currentstroke}%
\pgfsetdash{}{0pt}%
\pgfsys@defobject{currentmarker}{\pgfqpoint{-0.015278in}{-0.015278in}}{\pgfqpoint{0.015278in}{0.015278in}}{%
\pgfpathmoveto{\pgfqpoint{0.000000in}{-0.015278in}}%
\pgfpathcurveto{\pgfqpoint{0.004052in}{-0.015278in}}{\pgfqpoint{0.007938in}{-0.013668in}}{\pgfqpoint{0.010803in}{-0.010803in}}%
\pgfpathcurveto{\pgfqpoint{0.013668in}{-0.007938in}}{\pgfqpoint{0.015278in}{-0.004052in}}{\pgfqpoint{0.015278in}{0.000000in}}%
\pgfpathcurveto{\pgfqpoint{0.015278in}{0.004052in}}{\pgfqpoint{0.013668in}{0.007938in}}{\pgfqpoint{0.010803in}{0.010803in}}%
\pgfpathcurveto{\pgfqpoint{0.007938in}{0.013668in}}{\pgfqpoint{0.004052in}{0.015278in}}{\pgfqpoint{0.000000in}{0.015278in}}%
\pgfpathcurveto{\pgfqpoint{-0.004052in}{0.015278in}}{\pgfqpoint{-0.007938in}{0.013668in}}{\pgfqpoint{-0.010803in}{0.010803in}}%
\pgfpathcurveto{\pgfqpoint{-0.013668in}{0.007938in}}{\pgfqpoint{-0.015278in}{0.004052in}}{\pgfqpoint{-0.015278in}{0.000000in}}%
\pgfpathcurveto{\pgfqpoint{-0.015278in}{-0.004052in}}{\pgfqpoint{-0.013668in}{-0.007938in}}{\pgfqpoint{-0.010803in}{-0.010803in}}%
\pgfpathcurveto{\pgfqpoint{-0.007938in}{-0.013668in}}{\pgfqpoint{-0.004052in}{-0.015278in}}{\pgfqpoint{0.000000in}{-0.015278in}}%
\pgfpathlineto{\pgfqpoint{0.000000in}{-0.015278in}}%
\pgfpathclose%
\pgfusepath{stroke,fill}%
}%
\begin{pgfscope}%
\pgfsys@transformshift{4.118670in}{0.191821in}%
\pgfsys@useobject{currentmarker}{}%
\end{pgfscope}%
\end{pgfscope}%
\begin{pgfscope}%
\definecolor{textcolor}{rgb}{0.000000,0.000000,0.000000}%
\pgfsetstrokecolor{textcolor}%
\pgfsetfillcolor{textcolor}%
\pgftext[x=4.274225in,y=0.157793in,left,base]{\color{textcolor}{\rmfamily\fontsize{7.000000}{8.400000}\selectfont\catcode`\^=\active\def^{\ifmmode\sp\else\^{}\fi}\catcode`\%=\active\def%{\%}$\alpha=60^\circ$}}%
\end{pgfscope}%
\end{pgfpicture}%
\makeatother%
\endgroup%

\caption{Sensitivity of the $S^2$ heat profile to harmonic truncation. The plot compares $\widehat A_w(t)$ with a longer-series reference.
Smaller angular scales need more harmonics because heat flow has not yet damped the high frequencies.}
\label{fig:app_truncation}
\end{figure*}

\paragraph{Self-normalized importance-sampling diagnostic.}
Figure~\ref{fig:app_snis} considers self-normalized importance sampling with a uniform proposal on $S^2$. The target is an equal antipodal mixture of two vMF components. Their means are $e_3$ and $-e_3$, with $\kappa=18$. This concentration gives nontrivial importance weights and a clearly bimodal target. The density ratio is bounded because the proposal is uniform and $S^2$ is compact. Thus, Theorem~\ref{thm:consistency} applies. Ordinary ESS does not reveal whether high-weight particles occupy one region, two regions, or a diffuse set. The gESS curve adds this geometry and approaches one as the antipodal components blur together.

\begin{figure*}[t]
\centering
%\includegraphics[width=0.98\linewidth]{hkep_app_snis.png}
\providecommand{\mathdefault}[1]{#1}
\providecommand{\mathregular}[1]{#1}
%% Creator: Matplotlib, PGF backend
%%
%% To include the figure in your LaTeX document, write
%%   \input{<filename>.pgf}
%%
%% Make sure the required packages are loaded in your preamble
%%   \usepackage{pgf}
%%
%% Also ensure that all the required font packages are loaded; for instance,
%% the lmodern package is sometimes necessary when using math font.
%%   \usepackage{lmodern}
%%
%% Figures using additional raster images can only be included by \input if
%% they are in the same directory as the main LaTeX file. For loading figures
%% from other directories you can use the `import` package
%%   \usepackage{import}
%%
%% and then include the figures with
%%   \import{<path to file>}{<filename>.pgf}
%%
%% Matplotlib used the following preamble
%%   \def\mathdefault#1{#1}
%%   \everymath=\expandafter{\the\everymath\displaystyle}
%%   \IfFileExists{scrextend.sty}{
%%     \usepackage[fontsize=8.000000pt]{scrextend}
%%   }{
%%     \renewcommand{\normalsize}{\fontsize{8.000000}{9.600000}\selectfont}
%%     \normalsize
%%   }
%%   \usepackage{amsmath,amssymb}
%%   \makeatletter\@ifpackageloaded{underscore}{}{\usepackage[strings]{underscore}}\makeatother
%%
\begingroup%
\makeatletter%
\begin{pgfpicture}%
\pgfpathrectangle{\pgfpointorigin}{\pgfqpoint{5.431158in}{3.399883in}}%
\pgfusepath{use as bounding box, clip}%
\begin{pgfscope}%
\pgfsetbuttcap%
\pgfsetmiterjoin%
\definecolor{currentfill}{rgb}{1.000000,1.000000,1.000000}%
\pgfsetfillcolor{currentfill}%
\pgfsetlinewidth{0.000000pt}%
\definecolor{currentstroke}{rgb}{1.000000,1.000000,1.000000}%
\pgfsetstrokecolor{currentstroke}%
\pgfsetdash{}{0pt}%
\pgfpathmoveto{\pgfqpoint{0.000000in}{0.000000in}}%
\pgfpathlineto{\pgfqpoint{5.431158in}{0.000000in}}%
\pgfpathlineto{\pgfqpoint{5.431158in}{3.399883in}}%
\pgfpathlineto{\pgfqpoint{0.000000in}{3.399883in}}%
\pgfpathlineto{\pgfqpoint{0.000000in}{0.000000in}}%
\pgfpathclose%
\pgfusepath{fill}%
\end{pgfscope}%
\begin{pgfscope}%
\pgfsetbuttcap%
\pgfsetmiterjoin%
\definecolor{currentfill}{rgb}{1.000000,1.000000,1.000000}%
\pgfsetfillcolor{currentfill}%
\pgfsetlinewidth{0.000000pt}%
\definecolor{currentstroke}{rgb}{0.000000,0.000000,0.000000}%
\pgfsetstrokecolor{currentstroke}%
\pgfsetstrokeopacity{0.000000}%
\pgfsetdash{}{0pt}%
\pgfpathmoveto{\pgfqpoint{0.462269in}{0.457500in}}%
\pgfpathlineto{\pgfqpoint{5.331158in}{0.457500in}}%
\pgfpathlineto{\pgfqpoint{5.331158in}{3.139389in}}%
\pgfpathlineto{\pgfqpoint{0.462269in}{3.139389in}}%
\pgfpathlineto{\pgfqpoint{0.462269in}{0.457500in}}%
\pgfpathclose%
\pgfusepath{fill}%
\end{pgfscope}%
\begin{pgfscope}%
\pgfpathrectangle{\pgfqpoint{0.462269in}{0.457500in}}{\pgfqpoint{4.868889in}{2.681889in}}%
\pgfusepath{clip}%
\pgfsetbuttcap%
\pgfsetroundjoin%
\definecolor{currentfill}{rgb}{0.121569,0.466667,0.705882}%
\pgfsetfillcolor{currentfill}%
\pgfsetfillopacity{0.200000}%
\pgfsetlinewidth{0.000000pt}%
\definecolor{currentstroke}{rgb}{0.000000,0.000000,0.000000}%
\pgfsetstrokecolor{currentstroke}%
\pgfsetdash{}{0pt}%
\pgfpathmoveto{\pgfqpoint{0.683582in}{2.101731in}}%
\pgfpathlineto{\pgfqpoint{0.683582in}{1.907344in}}%
\pgfpathlineto{\pgfqpoint{0.758604in}{1.849076in}}%
\pgfpathlineto{\pgfqpoint{0.833625in}{1.777103in}}%
\pgfpathlineto{\pgfqpoint{0.908646in}{1.726294in}}%
\pgfpathlineto{\pgfqpoint{0.983668in}{1.660214in}}%
\pgfpathlineto{\pgfqpoint{1.058689in}{1.597069in}}%
\pgfpathlineto{\pgfqpoint{1.133711in}{1.532985in}}%
\pgfpathlineto{\pgfqpoint{1.208732in}{1.471867in}}%
\pgfpathlineto{\pgfqpoint{1.283754in}{1.413162in}}%
\pgfpathlineto{\pgfqpoint{1.358775in}{1.357495in}}%
\pgfpathlineto{\pgfqpoint{1.433796in}{1.303764in}}%
\pgfpathlineto{\pgfqpoint{1.508818in}{1.251271in}}%
\pgfpathlineto{\pgfqpoint{1.583839in}{1.201105in}}%
\pgfpathlineto{\pgfqpoint{1.658861in}{1.153515in}}%
\pgfpathlineto{\pgfqpoint{1.733882in}{1.108773in}}%
\pgfpathlineto{\pgfqpoint{1.808903in}{1.066913in}}%
\pgfpathlineto{\pgfqpoint{1.883925in}{1.027922in}}%
\pgfpathlineto{\pgfqpoint{1.958946in}{0.991747in}}%
\pgfpathlineto{\pgfqpoint{2.033968in}{0.958300in}}%
\pgfpathlineto{\pgfqpoint{2.108989in}{0.927467in}}%
\pgfpathlineto{\pgfqpoint{2.184010in}{0.899117in}}%
\pgfpathlineto{\pgfqpoint{2.259032in}{0.873115in}}%
\pgfpathlineto{\pgfqpoint{2.334053in}{0.849536in}}%
\pgfpathlineto{\pgfqpoint{2.409075in}{0.827943in}}%
\pgfpathlineto{\pgfqpoint{2.484096in}{0.808199in}}%
\pgfpathlineto{\pgfqpoint{2.559117in}{0.790293in}}%
\pgfpathlineto{\pgfqpoint{2.634139in}{0.774075in}}%
\pgfpathlineto{\pgfqpoint{2.709160in}{0.759118in}}%
\pgfpathlineto{\pgfqpoint{2.784182in}{0.745013in}}%
\pgfpathlineto{\pgfqpoint{2.859203in}{0.731923in}}%
\pgfpathlineto{\pgfqpoint{2.934224in}{0.718950in}}%
\pgfpathlineto{\pgfqpoint{3.009246in}{0.708960in}}%
\pgfpathlineto{\pgfqpoint{3.084267in}{0.699903in}}%
\pgfpathlineto{\pgfqpoint{3.159289in}{0.691671in}}%
\pgfpathlineto{\pgfqpoint{3.234310in}{0.683941in}}%
\pgfpathlineto{\pgfqpoint{3.309331in}{0.676928in}}%
\pgfpathlineto{\pgfqpoint{3.384353in}{0.670592in}}%
\pgfpathlineto{\pgfqpoint{3.459374in}{0.664937in}}%
\pgfpathlineto{\pgfqpoint{3.534396in}{0.659884in}}%
\pgfpathlineto{\pgfqpoint{3.609417in}{0.655326in}}%
\pgfpathlineto{\pgfqpoint{3.684438in}{0.651212in}}%
\pgfpathlineto{\pgfqpoint{3.759460in}{0.647554in}}%
\pgfpathlineto{\pgfqpoint{3.834481in}{0.644309in}}%
\pgfpathlineto{\pgfqpoint{3.909503in}{0.640988in}}%
\pgfpathlineto{\pgfqpoint{3.984524in}{0.637813in}}%
\pgfpathlineto{\pgfqpoint{4.059545in}{0.635213in}}%
\pgfpathlineto{\pgfqpoint{4.134567in}{0.632687in}}%
\pgfpathlineto{\pgfqpoint{4.209588in}{0.630133in}}%
\pgfpathlineto{\pgfqpoint{4.284610in}{0.627550in}}%
\pgfpathlineto{\pgfqpoint{4.359631in}{0.624694in}}%
\pgfpathlineto{\pgfqpoint{4.434652in}{0.621522in}}%
\pgfpathlineto{\pgfqpoint{4.509674in}{0.617932in}}%
\pgfpathlineto{\pgfqpoint{4.584695in}{0.613867in}}%
\pgfpathlineto{\pgfqpoint{4.659717in}{0.609274in}}%
\pgfpathlineto{\pgfqpoint{4.734738in}{0.604280in}}%
\pgfpathlineto{\pgfqpoint{4.809759in}{0.599069in}}%
\pgfpathlineto{\pgfqpoint{4.884781in}{0.593811in}}%
\pgfpathlineto{\pgfqpoint{4.959802in}{0.588682in}}%
\pgfpathlineto{\pgfqpoint{5.034824in}{0.583837in}}%
\pgfpathlineto{\pgfqpoint{5.109845in}{0.579404in}}%
\pgfpathlineto{\pgfqpoint{5.109845in}{0.580153in}}%
\pgfpathlineto{\pgfqpoint{5.109845in}{0.580153in}}%
\pgfpathlineto{\pgfqpoint{5.034824in}{0.584811in}}%
\pgfpathlineto{\pgfqpoint{4.959802in}{0.589926in}}%
\pgfpathlineto{\pgfqpoint{4.884781in}{0.595334in}}%
\pgfpathlineto{\pgfqpoint{4.809759in}{0.600888in}}%
\pgfpathlineto{\pgfqpoint{4.734738in}{0.606480in}}%
\pgfpathlineto{\pgfqpoint{4.659717in}{0.611819in}}%
\pgfpathlineto{\pgfqpoint{4.584695in}{0.616747in}}%
\pgfpathlineto{\pgfqpoint{4.509674in}{0.621196in}}%
\pgfpathlineto{\pgfqpoint{4.434652in}{0.625152in}}%
\pgfpathlineto{\pgfqpoint{4.359631in}{0.628669in}}%
\pgfpathlineto{\pgfqpoint{4.284610in}{0.631850in}}%
\pgfpathlineto{\pgfqpoint{4.209588in}{0.634827in}}%
\pgfpathlineto{\pgfqpoint{4.134567in}{0.637733in}}%
\pgfpathlineto{\pgfqpoint{4.059545in}{0.640693in}}%
\pgfpathlineto{\pgfqpoint{3.984524in}{0.643801in}}%
\pgfpathlineto{\pgfqpoint{3.909503in}{0.647131in}}%
\pgfpathlineto{\pgfqpoint{3.834481in}{0.650787in}}%
\pgfpathlineto{\pgfqpoint{3.759460in}{0.654976in}}%
\pgfpathlineto{\pgfqpoint{3.684438in}{0.659486in}}%
\pgfpathlineto{\pgfqpoint{3.609417in}{0.664492in}}%
\pgfpathlineto{\pgfqpoint{3.534396in}{0.670054in}}%
\pgfpathlineto{\pgfqpoint{3.459374in}{0.676235in}}%
\pgfpathlineto{\pgfqpoint{3.384353in}{0.683102in}}%
\pgfpathlineto{\pgfqpoint{3.309331in}{0.690728in}}%
\pgfpathlineto{\pgfqpoint{3.234310in}{0.699193in}}%
\pgfpathlineto{\pgfqpoint{3.159289in}{0.708585in}}%
\pgfpathlineto{\pgfqpoint{3.084267in}{0.718997in}}%
\pgfpathlineto{\pgfqpoint{3.009246in}{0.730534in}}%
\pgfpathlineto{\pgfqpoint{2.934224in}{0.743275in}}%
\pgfpathlineto{\pgfqpoint{2.859203in}{0.757357in}}%
\pgfpathlineto{\pgfqpoint{2.784182in}{0.772903in}}%
\pgfpathlineto{\pgfqpoint{2.709160in}{0.790044in}}%
\pgfpathlineto{\pgfqpoint{2.634139in}{0.808918in}}%
\pgfpathlineto{\pgfqpoint{2.559117in}{0.829751in}}%
\pgfpathlineto{\pgfqpoint{2.484096in}{0.852854in}}%
\pgfpathlineto{\pgfqpoint{2.409075in}{0.878121in}}%
\pgfpathlineto{\pgfqpoint{2.334053in}{0.904643in}}%
\pgfpathlineto{\pgfqpoint{2.259032in}{0.934030in}}%
\pgfpathlineto{\pgfqpoint{2.184010in}{0.966130in}}%
\pgfpathlineto{\pgfqpoint{2.108989in}{1.000817in}}%
\pgfpathlineto{\pgfqpoint{2.033968in}{1.038161in}}%
\pgfpathlineto{\pgfqpoint{1.958946in}{1.078212in}}%
\pgfpathlineto{\pgfqpoint{1.883925in}{1.120362in}}%
\pgfpathlineto{\pgfqpoint{1.808903in}{1.164610in}}%
\pgfpathlineto{\pgfqpoint{1.733882in}{1.210688in}}%
\pgfpathlineto{\pgfqpoint{1.658861in}{1.261681in}}%
\pgfpathlineto{\pgfqpoint{1.583839in}{1.316443in}}%
\pgfpathlineto{\pgfqpoint{1.508818in}{1.373447in}}%
\pgfpathlineto{\pgfqpoint{1.433796in}{1.432501in}}%
\pgfpathlineto{\pgfqpoint{1.358775in}{1.493248in}}%
\pgfpathlineto{\pgfqpoint{1.283754in}{1.555735in}}%
\pgfpathlineto{\pgfqpoint{1.208732in}{1.629306in}}%
\pgfpathlineto{\pgfqpoint{1.133711in}{1.704374in}}%
\pgfpathlineto{\pgfqpoint{1.058689in}{1.775718in}}%
\pgfpathlineto{\pgfqpoint{0.983668in}{1.843875in}}%
\pgfpathlineto{\pgfqpoint{0.908646in}{1.911625in}}%
\pgfpathlineto{\pgfqpoint{0.833625in}{1.974240in}}%
\pgfpathlineto{\pgfqpoint{0.758604in}{2.036590in}}%
\pgfpathlineto{\pgfqpoint{0.683582in}{2.101731in}}%
\pgfpathlineto{\pgfqpoint{0.683582in}{2.101731in}}%
\pgfpathclose%
\pgfusepath{fill}%
\end{pgfscope}%
\begin{pgfscope}%
\pgfpathrectangle{\pgfqpoint{0.462269in}{0.457500in}}{\pgfqpoint{4.868889in}{2.681889in}}%
\pgfusepath{clip}%
\pgfsetrectcap%
\pgfsetroundjoin%
\pgfsetlinewidth{0.803000pt}%
\definecolor{currentstroke}{rgb}{0.690196,0.690196,0.690196}%
\pgfsetstrokecolor{currentstroke}%
\pgfsetstrokeopacity{0.300000}%
\pgfsetdash{}{0pt}%
\pgfpathmoveto{\pgfqpoint{1.648969in}{0.457500in}}%
\pgfpathlineto{\pgfqpoint{1.648969in}{3.139389in}}%
\pgfusepath{stroke}%
\end{pgfscope}%
\begin{pgfscope}%
\pgfsetbuttcap%
\pgfsetroundjoin%
\definecolor{currentfill}{rgb}{0.000000,0.000000,0.000000}%
\pgfsetfillcolor{currentfill}%
\pgfsetlinewidth{0.803000pt}%
\definecolor{currentstroke}{rgb}{0.000000,0.000000,0.000000}%
\pgfsetstrokecolor{currentstroke}%
\pgfsetdash{}{0pt}%
\pgfsys@defobject{currentmarker}{\pgfqpoint{0.000000in}{-0.048611in}}{\pgfqpoint{0.000000in}{0.000000in}}{%
\pgfpathmoveto{\pgfqpoint{0.000000in}{0.000000in}}%
\pgfpathlineto{\pgfqpoint{0.000000in}{-0.048611in}}%
\pgfusepath{stroke,fill}%
}%
\begin{pgfscope}%
\pgfsys@transformshift{1.648969in}{0.457500in}%
\pgfsys@useobject{currentmarker}{}%
\end{pgfscope}%
\end{pgfscope}%
\begin{pgfscope}%
\definecolor{textcolor}{rgb}{0.000000,0.000000,0.000000}%
\pgfsetstrokecolor{textcolor}%
\pgfsetfillcolor{textcolor}%
\pgftext[x=1.648969in,y=0.360278in,,top]{\color{textcolor}{\rmfamily\fontsize{7.000000}{8.400000}\selectfont\catcode`\^=\active\def^{\ifmmode\sp\else\^{}\fi}\catcode`\%=\active\def%{\%}$\mathdefault{10^{1}}$}}%
\end{pgfscope}%
\begin{pgfscope}%
\pgfpathrectangle{\pgfqpoint{0.462269in}{0.457500in}}{\pgfqpoint{4.868889in}{2.681889in}}%
\pgfusepath{clip}%
\pgfsetrectcap%
\pgfsetroundjoin%
\pgfsetlinewidth{0.803000pt}%
\definecolor{currentstroke}{rgb}{0.690196,0.690196,0.690196}%
\pgfsetstrokecolor{currentstroke}%
\pgfsetstrokeopacity{0.300000}%
\pgfsetdash{}{0pt}%
\pgfpathmoveto{\pgfqpoint{4.855915in}{0.457500in}}%
\pgfpathlineto{\pgfqpoint{4.855915in}{3.139389in}}%
\pgfusepath{stroke}%
\end{pgfscope}%
\begin{pgfscope}%
\pgfsetbuttcap%
\pgfsetroundjoin%
\definecolor{currentfill}{rgb}{0.000000,0.000000,0.000000}%
\pgfsetfillcolor{currentfill}%
\pgfsetlinewidth{0.803000pt}%
\definecolor{currentstroke}{rgb}{0.000000,0.000000,0.000000}%
\pgfsetstrokecolor{currentstroke}%
\pgfsetdash{}{0pt}%
\pgfsys@defobject{currentmarker}{\pgfqpoint{0.000000in}{-0.048611in}}{\pgfqpoint{0.000000in}{0.000000in}}{%
\pgfpathmoveto{\pgfqpoint{0.000000in}{0.000000in}}%
\pgfpathlineto{\pgfqpoint{0.000000in}{-0.048611in}}%
\pgfusepath{stroke,fill}%
}%
\begin{pgfscope}%
\pgfsys@transformshift{4.855915in}{0.457500in}%
\pgfsys@useobject{currentmarker}{}%
\end{pgfscope}%
\end{pgfscope}%
\begin{pgfscope}%
\definecolor{textcolor}{rgb}{0.000000,0.000000,0.000000}%
\pgfsetstrokecolor{textcolor}%
\pgfsetfillcolor{textcolor}%
\pgftext[x=4.855915in,y=0.360278in,,top]{\color{textcolor}{\rmfamily\fontsize{7.000000}{8.400000}\selectfont\catcode`\^=\active\def^{\ifmmode\sp\else\^{}\fi}\catcode`\%=\active\def%{\%}$\mathdefault{10^{2}}$}}%
\end{pgfscope}%
\begin{pgfscope}%
\pgfsetbuttcap%
\pgfsetroundjoin%
\definecolor{currentfill}{rgb}{0.000000,0.000000,0.000000}%
\pgfsetfillcolor{currentfill}%
\pgfsetlinewidth{0.602250pt}%
\definecolor{currentstroke}{rgb}{0.000000,0.000000,0.000000}%
\pgfsetstrokecolor{currentstroke}%
\pgfsetdash{}{0pt}%
\pgfsys@defobject{currentmarker}{\pgfqpoint{0.000000in}{-0.027778in}}{\pgfqpoint{0.000000in}{0.000000in}}{%
\pgfpathmoveto{\pgfqpoint{0.000000in}{0.000000in}}%
\pgfpathlineto{\pgfqpoint{0.000000in}{-0.027778in}}%
\pgfusepath{stroke,fill}%
}%
\begin{pgfscope}%
\pgfsys@transformshift{0.683582in}{0.457500in}%
\pgfsys@useobject{currentmarker}{}%
\end{pgfscope}%
\end{pgfscope}%
\begin{pgfscope}%
\pgfsetbuttcap%
\pgfsetroundjoin%
\definecolor{currentfill}{rgb}{0.000000,0.000000,0.000000}%
\pgfsetfillcolor{currentfill}%
\pgfsetlinewidth{0.602250pt}%
\definecolor{currentstroke}{rgb}{0.000000,0.000000,0.000000}%
\pgfsetstrokecolor{currentstroke}%
\pgfsetdash{}{0pt}%
\pgfsys@defobject{currentmarker}{\pgfqpoint{0.000000in}{-0.027778in}}{\pgfqpoint{0.000000in}{0.000000in}}{%
\pgfpathmoveto{\pgfqpoint{0.000000in}{0.000000in}}%
\pgfpathlineto{\pgfqpoint{0.000000in}{-0.027778in}}%
\pgfusepath{stroke,fill}%
}%
\begin{pgfscope}%
\pgfsys@transformshift{0.937512in}{0.457500in}%
\pgfsys@useobject{currentmarker}{}%
\end{pgfscope}%
\end{pgfscope}%
\begin{pgfscope}%
\pgfsetbuttcap%
\pgfsetroundjoin%
\definecolor{currentfill}{rgb}{0.000000,0.000000,0.000000}%
\pgfsetfillcolor{currentfill}%
\pgfsetlinewidth{0.602250pt}%
\definecolor{currentstroke}{rgb}{0.000000,0.000000,0.000000}%
\pgfsetstrokecolor{currentstroke}%
\pgfsetdash{}{0pt}%
\pgfsys@defobject{currentmarker}{\pgfqpoint{0.000000in}{-0.027778in}}{\pgfqpoint{0.000000in}{0.000000in}}{%
\pgfpathmoveto{\pgfqpoint{0.000000in}{0.000000in}}%
\pgfpathlineto{\pgfqpoint{0.000000in}{-0.027778in}}%
\pgfusepath{stroke,fill}%
}%
\begin{pgfscope}%
\pgfsys@transformshift{1.152207in}{0.457500in}%
\pgfsys@useobject{currentmarker}{}%
\end{pgfscope}%
\end{pgfscope}%
\begin{pgfscope}%
\pgfsetbuttcap%
\pgfsetroundjoin%
\definecolor{currentfill}{rgb}{0.000000,0.000000,0.000000}%
\pgfsetfillcolor{currentfill}%
\pgfsetlinewidth{0.602250pt}%
\definecolor{currentstroke}{rgb}{0.000000,0.000000,0.000000}%
\pgfsetstrokecolor{currentstroke}%
\pgfsetdash{}{0pt}%
\pgfsys@defobject{currentmarker}{\pgfqpoint{0.000000in}{-0.027778in}}{\pgfqpoint{0.000000in}{0.000000in}}{%
\pgfpathmoveto{\pgfqpoint{0.000000in}{0.000000in}}%
\pgfpathlineto{\pgfqpoint{0.000000in}{-0.027778in}}%
\pgfusepath{stroke,fill}%
}%
\begin{pgfscope}%
\pgfsys@transformshift{1.338184in}{0.457500in}%
\pgfsys@useobject{currentmarker}{}%
\end{pgfscope}%
\end{pgfscope}%
\begin{pgfscope}%
\pgfsetbuttcap%
\pgfsetroundjoin%
\definecolor{currentfill}{rgb}{0.000000,0.000000,0.000000}%
\pgfsetfillcolor{currentfill}%
\pgfsetlinewidth{0.602250pt}%
\definecolor{currentstroke}{rgb}{0.000000,0.000000,0.000000}%
\pgfsetstrokecolor{currentstroke}%
\pgfsetdash{}{0pt}%
\pgfsys@defobject{currentmarker}{\pgfqpoint{0.000000in}{-0.027778in}}{\pgfqpoint{0.000000in}{0.000000in}}{%
\pgfpathmoveto{\pgfqpoint{0.000000in}{0.000000in}}%
\pgfpathlineto{\pgfqpoint{0.000000in}{-0.027778in}}%
\pgfusepath{stroke,fill}%
}%
\begin{pgfscope}%
\pgfsys@transformshift{1.502227in}{0.457500in}%
\pgfsys@useobject{currentmarker}{}%
\end{pgfscope}%
\end{pgfscope}%
\begin{pgfscope}%
\pgfsetbuttcap%
\pgfsetroundjoin%
\definecolor{currentfill}{rgb}{0.000000,0.000000,0.000000}%
\pgfsetfillcolor{currentfill}%
\pgfsetlinewidth{0.602250pt}%
\definecolor{currentstroke}{rgb}{0.000000,0.000000,0.000000}%
\pgfsetstrokecolor{currentstroke}%
\pgfsetdash{}{0pt}%
\pgfsys@defobject{currentmarker}{\pgfqpoint{0.000000in}{-0.027778in}}{\pgfqpoint{0.000000in}{0.000000in}}{%
\pgfpathmoveto{\pgfqpoint{0.000000in}{0.000000in}}%
\pgfpathlineto{\pgfqpoint{0.000000in}{-0.027778in}}%
\pgfusepath{stroke,fill}%
}%
\begin{pgfscope}%
\pgfsys@transformshift{2.614356in}{0.457500in}%
\pgfsys@useobject{currentmarker}{}%
\end{pgfscope}%
\end{pgfscope}%
\begin{pgfscope}%
\pgfsetbuttcap%
\pgfsetroundjoin%
\definecolor{currentfill}{rgb}{0.000000,0.000000,0.000000}%
\pgfsetfillcolor{currentfill}%
\pgfsetlinewidth{0.602250pt}%
\definecolor{currentstroke}{rgb}{0.000000,0.000000,0.000000}%
\pgfsetstrokecolor{currentstroke}%
\pgfsetdash{}{0pt}%
\pgfsys@defobject{currentmarker}{\pgfqpoint{0.000000in}{-0.027778in}}{\pgfqpoint{0.000000in}{0.000000in}}{%
\pgfpathmoveto{\pgfqpoint{0.000000in}{0.000000in}}%
\pgfpathlineto{\pgfqpoint{0.000000in}{-0.027778in}}%
\pgfusepath{stroke,fill}%
}%
\begin{pgfscope}%
\pgfsys@transformshift{3.179071in}{0.457500in}%
\pgfsys@useobject{currentmarker}{}%
\end{pgfscope}%
\end{pgfscope}%
\begin{pgfscope}%
\pgfsetbuttcap%
\pgfsetroundjoin%
\definecolor{currentfill}{rgb}{0.000000,0.000000,0.000000}%
\pgfsetfillcolor{currentfill}%
\pgfsetlinewidth{0.602250pt}%
\definecolor{currentstroke}{rgb}{0.000000,0.000000,0.000000}%
\pgfsetstrokecolor{currentstroke}%
\pgfsetdash{}{0pt}%
\pgfsys@defobject{currentmarker}{\pgfqpoint{0.000000in}{-0.027778in}}{\pgfqpoint{0.000000in}{0.000000in}}{%
\pgfpathmoveto{\pgfqpoint{0.000000in}{0.000000in}}%
\pgfpathlineto{\pgfqpoint{0.000000in}{-0.027778in}}%
\pgfusepath{stroke,fill}%
}%
\begin{pgfscope}%
\pgfsys@transformshift{3.579743in}{0.457500in}%
\pgfsys@useobject{currentmarker}{}%
\end{pgfscope}%
\end{pgfscope}%
\begin{pgfscope}%
\pgfsetbuttcap%
\pgfsetroundjoin%
\definecolor{currentfill}{rgb}{0.000000,0.000000,0.000000}%
\pgfsetfillcolor{currentfill}%
\pgfsetlinewidth{0.602250pt}%
\definecolor{currentstroke}{rgb}{0.000000,0.000000,0.000000}%
\pgfsetstrokecolor{currentstroke}%
\pgfsetdash{}{0pt}%
\pgfsys@defobject{currentmarker}{\pgfqpoint{0.000000in}{-0.027778in}}{\pgfqpoint{0.000000in}{0.000000in}}{%
\pgfpathmoveto{\pgfqpoint{0.000000in}{0.000000in}}%
\pgfpathlineto{\pgfqpoint{0.000000in}{-0.027778in}}%
\pgfusepath{stroke,fill}%
}%
\begin{pgfscope}%
\pgfsys@transformshift{3.890528in}{0.457500in}%
\pgfsys@useobject{currentmarker}{}%
\end{pgfscope}%
\end{pgfscope}%
\begin{pgfscope}%
\pgfsetbuttcap%
\pgfsetroundjoin%
\definecolor{currentfill}{rgb}{0.000000,0.000000,0.000000}%
\pgfsetfillcolor{currentfill}%
\pgfsetlinewidth{0.602250pt}%
\definecolor{currentstroke}{rgb}{0.000000,0.000000,0.000000}%
\pgfsetstrokecolor{currentstroke}%
\pgfsetdash{}{0pt}%
\pgfsys@defobject{currentmarker}{\pgfqpoint{0.000000in}{-0.027778in}}{\pgfqpoint{0.000000in}{0.000000in}}{%
\pgfpathmoveto{\pgfqpoint{0.000000in}{0.000000in}}%
\pgfpathlineto{\pgfqpoint{0.000000in}{-0.027778in}}%
\pgfusepath{stroke,fill}%
}%
\begin{pgfscope}%
\pgfsys@transformshift{4.144458in}{0.457500in}%
\pgfsys@useobject{currentmarker}{}%
\end{pgfscope}%
\end{pgfscope}%
\begin{pgfscope}%
\pgfsetbuttcap%
\pgfsetroundjoin%
\definecolor{currentfill}{rgb}{0.000000,0.000000,0.000000}%
\pgfsetfillcolor{currentfill}%
\pgfsetlinewidth{0.602250pt}%
\definecolor{currentstroke}{rgb}{0.000000,0.000000,0.000000}%
\pgfsetstrokecolor{currentstroke}%
\pgfsetdash{}{0pt}%
\pgfsys@defobject{currentmarker}{\pgfqpoint{0.000000in}{-0.027778in}}{\pgfqpoint{0.000000in}{0.000000in}}{%
\pgfpathmoveto{\pgfqpoint{0.000000in}{0.000000in}}%
\pgfpathlineto{\pgfqpoint{0.000000in}{-0.027778in}}%
\pgfusepath{stroke,fill}%
}%
\begin{pgfscope}%
\pgfsys@transformshift{4.359153in}{0.457500in}%
\pgfsys@useobject{currentmarker}{}%
\end{pgfscope}%
\end{pgfscope}%
\begin{pgfscope}%
\pgfsetbuttcap%
\pgfsetroundjoin%
\definecolor{currentfill}{rgb}{0.000000,0.000000,0.000000}%
\pgfsetfillcolor{currentfill}%
\pgfsetlinewidth{0.602250pt}%
\definecolor{currentstroke}{rgb}{0.000000,0.000000,0.000000}%
\pgfsetstrokecolor{currentstroke}%
\pgfsetdash{}{0pt}%
\pgfsys@defobject{currentmarker}{\pgfqpoint{0.000000in}{-0.027778in}}{\pgfqpoint{0.000000in}{0.000000in}}{%
\pgfpathmoveto{\pgfqpoint{0.000000in}{0.000000in}}%
\pgfpathlineto{\pgfqpoint{0.000000in}{-0.027778in}}%
\pgfusepath{stroke,fill}%
}%
\begin{pgfscope}%
\pgfsys@transformshift{4.545130in}{0.457500in}%
\pgfsys@useobject{currentmarker}{}%
\end{pgfscope}%
\end{pgfscope}%
\begin{pgfscope}%
\pgfsetbuttcap%
\pgfsetroundjoin%
\definecolor{currentfill}{rgb}{0.000000,0.000000,0.000000}%
\pgfsetfillcolor{currentfill}%
\pgfsetlinewidth{0.602250pt}%
\definecolor{currentstroke}{rgb}{0.000000,0.000000,0.000000}%
\pgfsetstrokecolor{currentstroke}%
\pgfsetdash{}{0pt}%
\pgfsys@defobject{currentmarker}{\pgfqpoint{0.000000in}{-0.027778in}}{\pgfqpoint{0.000000in}{0.000000in}}{%
\pgfpathmoveto{\pgfqpoint{0.000000in}{0.000000in}}%
\pgfpathlineto{\pgfqpoint{0.000000in}{-0.027778in}}%
\pgfusepath{stroke,fill}%
}%
\begin{pgfscope}%
\pgfsys@transformshift{4.709173in}{0.457500in}%
\pgfsys@useobject{currentmarker}{}%
\end{pgfscope}%
\end{pgfscope}%
\begin{pgfscope}%
\definecolor{textcolor}{rgb}{0.000000,0.000000,0.000000}%
\pgfsetstrokecolor{textcolor}%
\pgfsetfillcolor{textcolor}%
\pgftext[x=2.896714in,y=0.211111in,,top]{\color{textcolor}{\rmfamily\fontsize{8.000000}{9.600000}\selectfont\catcode`\^=\active\def^{\ifmmode\sp\else\^{}\fi}\catcode`\%=\active\def%{\%}Angular Scale $\alpha$ (Degrees)}}%
\end{pgfscope}%
\begin{pgfscope}%
\pgfpathrectangle{\pgfqpoint{0.462269in}{0.457500in}}{\pgfqpoint{4.868889in}{2.681889in}}%
\pgfusepath{clip}%
\pgfsetrectcap%
\pgfsetroundjoin%
\pgfsetlinewidth{0.803000pt}%
\definecolor{currentstroke}{rgb}{0.690196,0.690196,0.690196}%
\pgfsetstrokecolor{currentstroke}%
\pgfsetstrokeopacity{0.300000}%
\pgfsetdash{}{0pt}%
\pgfpathmoveto{\pgfqpoint{0.462269in}{0.505289in}}%
\pgfpathlineto{\pgfqpoint{5.331158in}{0.505289in}}%
\pgfusepath{stroke}%
\end{pgfscope}%
\begin{pgfscope}%
\pgfsetbuttcap%
\pgfsetroundjoin%
\definecolor{currentfill}{rgb}{0.000000,0.000000,0.000000}%
\pgfsetfillcolor{currentfill}%
\pgfsetlinewidth{0.803000pt}%
\definecolor{currentstroke}{rgb}{0.000000,0.000000,0.000000}%
\pgfsetstrokecolor{currentstroke}%
\pgfsetdash{}{0pt}%
\pgfsys@defobject{currentmarker}{\pgfqpoint{-0.048611in}{0.000000in}}{\pgfqpoint{-0.000000in}{0.000000in}}{%
\pgfpathmoveto{\pgfqpoint{-0.000000in}{0.000000in}}%
\pgfpathlineto{\pgfqpoint{-0.048611in}{0.000000in}}%
\pgfusepath{stroke,fill}%
}%
\begin{pgfscope}%
\pgfsys@transformshift{0.462269in}{0.505289in}%
\pgfsys@useobject{currentmarker}{}%
\end{pgfscope}%
\end{pgfscope}%
\begin{pgfscope}%
\definecolor{textcolor}{rgb}{0.000000,0.000000,0.000000}%
\pgfsetstrokecolor{textcolor}%
\pgfsetfillcolor{textcolor}%
\pgftext[x=0.309684in, y=0.471532in, left, base]{\color{textcolor}{\rmfamily\fontsize{7.000000}{8.400000}\selectfont\catcode`\^=\active\def^{\ifmmode\sp\else\^{}\fi}\catcode`\%=\active\def%{\%}0}}%
\end{pgfscope}%
\begin{pgfscope}%
\pgfpathrectangle{\pgfqpoint{0.462269in}{0.457500in}}{\pgfqpoint{4.868889in}{2.681889in}}%
\pgfusepath{clip}%
\pgfsetrectcap%
\pgfsetroundjoin%
\pgfsetlinewidth{0.803000pt}%
\definecolor{currentstroke}{rgb}{0.690196,0.690196,0.690196}%
\pgfsetstrokecolor{currentstroke}%
\pgfsetstrokeopacity{0.300000}%
\pgfsetdash{}{0pt}%
\pgfpathmoveto{\pgfqpoint{0.462269in}{1.066683in}}%
\pgfpathlineto{\pgfqpoint{5.331158in}{1.066683in}}%
\pgfusepath{stroke}%
\end{pgfscope}%
\begin{pgfscope}%
\pgfsetbuttcap%
\pgfsetroundjoin%
\definecolor{currentfill}{rgb}{0.000000,0.000000,0.000000}%
\pgfsetfillcolor{currentfill}%
\pgfsetlinewidth{0.803000pt}%
\definecolor{currentstroke}{rgb}{0.000000,0.000000,0.000000}%
\pgfsetstrokecolor{currentstroke}%
\pgfsetdash{}{0pt}%
\pgfsys@defobject{currentmarker}{\pgfqpoint{-0.048611in}{0.000000in}}{\pgfqpoint{-0.000000in}{0.000000in}}{%
\pgfpathmoveto{\pgfqpoint{-0.000000in}{0.000000in}}%
\pgfpathlineto{\pgfqpoint{-0.048611in}{0.000000in}}%
\pgfusepath{stroke,fill}%
}%
\begin{pgfscope}%
\pgfsys@transformshift{0.462269in}{1.066683in}%
\pgfsys@useobject{currentmarker}{}%
\end{pgfscope}%
\end{pgfscope}%
\begin{pgfscope}%
\definecolor{textcolor}{rgb}{0.000000,0.000000,0.000000}%
\pgfsetstrokecolor{textcolor}%
\pgfsetfillcolor{textcolor}%
\pgftext[x=0.254321in, y=1.032926in, left, base]{\color{textcolor}{\rmfamily\fontsize{7.000000}{8.400000}\selectfont\catcode`\^=\active\def^{\ifmmode\sp\else\^{}\fi}\catcode`\%=\active\def%{\%}10}}%
\end{pgfscope}%
\begin{pgfscope}%
\pgfpathrectangle{\pgfqpoint{0.462269in}{0.457500in}}{\pgfqpoint{4.868889in}{2.681889in}}%
\pgfusepath{clip}%
\pgfsetrectcap%
\pgfsetroundjoin%
\pgfsetlinewidth{0.803000pt}%
\definecolor{currentstroke}{rgb}{0.690196,0.690196,0.690196}%
\pgfsetstrokecolor{currentstroke}%
\pgfsetstrokeopacity{0.300000}%
\pgfsetdash{}{0pt}%
\pgfpathmoveto{\pgfqpoint{0.462269in}{1.628077in}}%
\pgfpathlineto{\pgfqpoint{5.331158in}{1.628077in}}%
\pgfusepath{stroke}%
\end{pgfscope}%
\begin{pgfscope}%
\pgfsetbuttcap%
\pgfsetroundjoin%
\definecolor{currentfill}{rgb}{0.000000,0.000000,0.000000}%
\pgfsetfillcolor{currentfill}%
\pgfsetlinewidth{0.803000pt}%
\definecolor{currentstroke}{rgb}{0.000000,0.000000,0.000000}%
\pgfsetstrokecolor{currentstroke}%
\pgfsetdash{}{0pt}%
\pgfsys@defobject{currentmarker}{\pgfqpoint{-0.048611in}{0.000000in}}{\pgfqpoint{-0.000000in}{0.000000in}}{%
\pgfpathmoveto{\pgfqpoint{-0.000000in}{0.000000in}}%
\pgfpathlineto{\pgfqpoint{-0.048611in}{0.000000in}}%
\pgfusepath{stroke,fill}%
}%
\begin{pgfscope}%
\pgfsys@transformshift{0.462269in}{1.628077in}%
\pgfsys@useobject{currentmarker}{}%
\end{pgfscope}%
\end{pgfscope}%
\begin{pgfscope}%
\definecolor{textcolor}{rgb}{0.000000,0.000000,0.000000}%
\pgfsetstrokecolor{textcolor}%
\pgfsetfillcolor{textcolor}%
\pgftext[x=0.254321in, y=1.594320in, left, base]{\color{textcolor}{\rmfamily\fontsize{7.000000}{8.400000}\selectfont\catcode`\^=\active\def^{\ifmmode\sp\else\^{}\fi}\catcode`\%=\active\def%{\%}20}}%
\end{pgfscope}%
\begin{pgfscope}%
\pgfpathrectangle{\pgfqpoint{0.462269in}{0.457500in}}{\pgfqpoint{4.868889in}{2.681889in}}%
\pgfusepath{clip}%
\pgfsetrectcap%
\pgfsetroundjoin%
\pgfsetlinewidth{0.803000pt}%
\definecolor{currentstroke}{rgb}{0.690196,0.690196,0.690196}%
\pgfsetstrokecolor{currentstroke}%
\pgfsetstrokeopacity{0.300000}%
\pgfsetdash{}{0pt}%
\pgfpathmoveto{\pgfqpoint{0.462269in}{2.189472in}}%
\pgfpathlineto{\pgfqpoint{5.331158in}{2.189472in}}%
\pgfusepath{stroke}%
\end{pgfscope}%
\begin{pgfscope}%
\pgfsetbuttcap%
\pgfsetroundjoin%
\definecolor{currentfill}{rgb}{0.000000,0.000000,0.000000}%
\pgfsetfillcolor{currentfill}%
\pgfsetlinewidth{0.803000pt}%
\definecolor{currentstroke}{rgb}{0.000000,0.000000,0.000000}%
\pgfsetstrokecolor{currentstroke}%
\pgfsetdash{}{0pt}%
\pgfsys@defobject{currentmarker}{\pgfqpoint{-0.048611in}{0.000000in}}{\pgfqpoint{-0.000000in}{0.000000in}}{%
\pgfpathmoveto{\pgfqpoint{-0.000000in}{0.000000in}}%
\pgfpathlineto{\pgfqpoint{-0.048611in}{0.000000in}}%
\pgfusepath{stroke,fill}%
}%
\begin{pgfscope}%
\pgfsys@transformshift{0.462269in}{2.189472in}%
\pgfsys@useobject{currentmarker}{}%
\end{pgfscope}%
\end{pgfscope}%
\begin{pgfscope}%
\definecolor{textcolor}{rgb}{0.000000,0.000000,0.000000}%
\pgfsetstrokecolor{textcolor}%
\pgfsetfillcolor{textcolor}%
\pgftext[x=0.254321in, y=2.155714in, left, base]{\color{textcolor}{\rmfamily\fontsize{7.000000}{8.400000}\selectfont\catcode`\^=\active\def^{\ifmmode\sp\else\^{}\fi}\catcode`\%=\active\def%{\%}30}}%
\end{pgfscope}%
\begin{pgfscope}%
\pgfpathrectangle{\pgfqpoint{0.462269in}{0.457500in}}{\pgfqpoint{4.868889in}{2.681889in}}%
\pgfusepath{clip}%
\pgfsetrectcap%
\pgfsetroundjoin%
\pgfsetlinewidth{0.803000pt}%
\definecolor{currentstroke}{rgb}{0.690196,0.690196,0.690196}%
\pgfsetstrokecolor{currentstroke}%
\pgfsetstrokeopacity{0.300000}%
\pgfsetdash{}{0pt}%
\pgfpathmoveto{\pgfqpoint{0.462269in}{2.750866in}}%
\pgfpathlineto{\pgfqpoint{5.331158in}{2.750866in}}%
\pgfusepath{stroke}%
\end{pgfscope}%
\begin{pgfscope}%
\pgfsetbuttcap%
\pgfsetroundjoin%
\definecolor{currentfill}{rgb}{0.000000,0.000000,0.000000}%
\pgfsetfillcolor{currentfill}%
\pgfsetlinewidth{0.803000pt}%
\definecolor{currentstroke}{rgb}{0.000000,0.000000,0.000000}%
\pgfsetstrokecolor{currentstroke}%
\pgfsetdash{}{0pt}%
\pgfsys@defobject{currentmarker}{\pgfqpoint{-0.048611in}{0.000000in}}{\pgfqpoint{-0.000000in}{0.000000in}}{%
\pgfpathmoveto{\pgfqpoint{-0.000000in}{0.000000in}}%
\pgfpathlineto{\pgfqpoint{-0.048611in}{0.000000in}}%
\pgfusepath{stroke,fill}%
}%
\begin{pgfscope}%
\pgfsys@transformshift{0.462269in}{2.750866in}%
\pgfsys@useobject{currentmarker}{}%
\end{pgfscope}%
\end{pgfscope}%
\begin{pgfscope}%
\definecolor{textcolor}{rgb}{0.000000,0.000000,0.000000}%
\pgfsetstrokecolor{textcolor}%
\pgfsetfillcolor{textcolor}%
\pgftext[x=0.254321in, y=2.717108in, left, base]{\color{textcolor}{\rmfamily\fontsize{7.000000}{8.400000}\selectfont\catcode`\^=\active\def^{\ifmmode\sp\else\^{}\fi}\catcode`\%=\active\def%{\%}40}}%
\end{pgfscope}%
\begin{pgfscope}%
\definecolor{textcolor}{rgb}{0.000000,0.000000,0.000000}%
\pgfsetstrokecolor{textcolor}%
\pgfsetfillcolor{textcolor}%
\pgftext[x=0.198766in,y=1.798444in,,bottom,rotate=90.000000]{\color{textcolor}{\rmfamily\fontsize{8.000000}{9.600000}\selectfont\catcode`\^=\active\def^{\ifmmode\sp\else\^{}\fi}\catcode`\%=\active\def%{\%}Effective Number}}%
\end{pgfscope}%
\begin{pgfscope}%
\pgfpathrectangle{\pgfqpoint{0.462269in}{0.457500in}}{\pgfqpoint{4.868889in}{2.681889in}}%
\pgfusepath{clip}%
\pgfsetrectcap%
\pgfsetroundjoin%
\pgfsetlinewidth{2.007500pt}%
\definecolor{currentstroke}{rgb}{0.121569,0.466667,0.705882}%
\pgfsetstrokecolor{currentstroke}%
\pgfsetdash{}{0pt}%
\pgfpathmoveto{\pgfqpoint{0.683582in}{2.015789in}}%
\pgfpathlineto{\pgfqpoint{0.758604in}{1.951216in}}%
\pgfpathlineto{\pgfqpoint{0.833625in}{1.892504in}}%
\pgfpathlineto{\pgfqpoint{0.908646in}{1.819530in}}%
\pgfpathlineto{\pgfqpoint{0.983668in}{1.743327in}}%
\pgfpathlineto{\pgfqpoint{1.058689in}{1.673211in}}%
\pgfpathlineto{\pgfqpoint{1.133711in}{1.612457in}}%
\pgfpathlineto{\pgfqpoint{1.208732in}{1.557837in}}%
\pgfpathlineto{\pgfqpoint{1.283754in}{1.498185in}}%
\pgfpathlineto{\pgfqpoint{1.358775in}{1.433108in}}%
\pgfpathlineto{\pgfqpoint{1.433796in}{1.372000in}}%
\pgfpathlineto{\pgfqpoint{1.508818in}{1.317106in}}%
\pgfpathlineto{\pgfqpoint{1.583839in}{1.264229in}}%
\pgfpathlineto{\pgfqpoint{1.658861in}{1.213545in}}%
\pgfpathlineto{\pgfqpoint{1.733882in}{1.165205in}}%
\pgfpathlineto{\pgfqpoint{1.808903in}{1.118892in}}%
\pgfpathlineto{\pgfqpoint{1.883925in}{1.074387in}}%
\pgfpathlineto{\pgfqpoint{1.958946in}{1.034217in}}%
\pgfpathlineto{\pgfqpoint{2.033968in}{0.997056in}}%
\pgfpathlineto{\pgfqpoint{2.108989in}{0.962755in}}%
\pgfpathlineto{\pgfqpoint{2.184010in}{0.932051in}}%
\pgfpathlineto{\pgfqpoint{2.259032in}{0.903067in}}%
\pgfpathlineto{\pgfqpoint{2.334053in}{0.876267in}}%
\pgfpathlineto{\pgfqpoint{2.409075in}{0.851613in}}%
\pgfpathlineto{\pgfqpoint{2.484096in}{0.829634in}}%
\pgfpathlineto{\pgfqpoint{2.559117in}{0.809576in}}%
\pgfpathlineto{\pgfqpoint{2.634139in}{0.791197in}}%
\pgfpathlineto{\pgfqpoint{2.709160in}{0.774378in}}%
\pgfpathlineto{\pgfqpoint{2.784182in}{0.759047in}}%
\pgfpathlineto{\pgfqpoint{2.859203in}{0.744906in}}%
\pgfpathlineto{\pgfqpoint{2.934224in}{0.732031in}}%
\pgfpathlineto{\pgfqpoint{3.009246in}{0.720335in}}%
\pgfpathlineto{\pgfqpoint{3.084267in}{0.709722in}}%
\pgfpathlineto{\pgfqpoint{3.159289in}{0.700101in}}%
\pgfpathlineto{\pgfqpoint{3.234310in}{0.691386in}}%
\pgfpathlineto{\pgfqpoint{3.309331in}{0.683498in}}%
\pgfpathlineto{\pgfqpoint{3.384353in}{0.676318in}}%
\pgfpathlineto{\pgfqpoint{3.459374in}{0.669878in}}%
\pgfpathlineto{\pgfqpoint{3.534396in}{0.664243in}}%
\pgfpathlineto{\pgfqpoint{3.609417in}{0.659240in}}%
\pgfpathlineto{\pgfqpoint{3.684438in}{0.654722in}}%
\pgfpathlineto{\pgfqpoint{3.759460in}{0.650639in}}%
\pgfpathlineto{\pgfqpoint{3.834481in}{0.646944in}}%
\pgfpathlineto{\pgfqpoint{3.909503in}{0.643587in}}%
\pgfpathlineto{\pgfqpoint{3.984524in}{0.640414in}}%
\pgfpathlineto{\pgfqpoint{4.059545in}{0.637418in}}%
\pgfpathlineto{\pgfqpoint{4.134567in}{0.634671in}}%
\pgfpathlineto{\pgfqpoint{4.209588in}{0.632100in}}%
\pgfpathlineto{\pgfqpoint{4.284610in}{0.629411in}}%
\pgfpathlineto{\pgfqpoint{4.359631in}{0.626494in}}%
\pgfpathlineto{\pgfqpoint{4.434652in}{0.623128in}}%
\pgfpathlineto{\pgfqpoint{4.509674in}{0.619282in}}%
\pgfpathlineto{\pgfqpoint{4.584695in}{0.615158in}}%
\pgfpathlineto{\pgfqpoint{4.659717in}{0.610538in}}%
\pgfpathlineto{\pgfqpoint{4.734738in}{0.605443in}}%
\pgfpathlineto{\pgfqpoint{4.809759in}{0.600108in}}%
\pgfpathlineto{\pgfqpoint{4.884781in}{0.594725in}}%
\pgfpathlineto{\pgfqpoint{4.959802in}{0.589450in}}%
\pgfpathlineto{\pgfqpoint{5.034824in}{0.584458in}}%
\pgfpathlineto{\pgfqpoint{5.109845in}{0.579888in}}%
\pgfusepath{stroke}%
\end{pgfscope}%
\begin{pgfscope}%
\pgfpathrectangle{\pgfqpoint{0.462269in}{0.457500in}}{\pgfqpoint{4.868889in}{2.681889in}}%
\pgfusepath{clip}%
\pgfsetbuttcap%
\pgfsetroundjoin%
\pgfsetlinewidth{1.505625pt}%
\definecolor{currentstroke}{rgb}{0.000000,0.000000,0.000000}%
\pgfsetstrokecolor{currentstroke}%
\pgfsetdash{{5.550000pt}{2.400000pt}}{0.000000pt}%
\pgfpathmoveto{\pgfqpoint{0.462269in}{3.017485in}}%
\pgfpathlineto{\pgfqpoint{5.331158in}{3.017485in}}%
\pgfusepath{stroke}%
\end{pgfscope}%
\begin{pgfscope}%
\pgfsetrectcap%
\pgfsetmiterjoin%
\pgfsetlinewidth{0.803000pt}%
\definecolor{currentstroke}{rgb}{0.000000,0.000000,0.000000}%
\pgfsetstrokecolor{currentstroke}%
\pgfsetdash{}{0pt}%
\pgfpathmoveto{\pgfqpoint{0.462269in}{0.457500in}}%
\pgfpathlineto{\pgfqpoint{0.462269in}{3.139389in}}%
\pgfusepath{stroke}%
\end{pgfscope}%
\begin{pgfscope}%
\pgfsetrectcap%
\pgfsetmiterjoin%
\pgfsetlinewidth{0.803000pt}%
\definecolor{currentstroke}{rgb}{0.000000,0.000000,0.000000}%
\pgfsetstrokecolor{currentstroke}%
\pgfsetdash{}{0pt}%
\pgfpathmoveto{\pgfqpoint{5.331158in}{0.457500in}}%
\pgfpathlineto{\pgfqpoint{5.331158in}{3.139389in}}%
\pgfusepath{stroke}%
\end{pgfscope}%
\begin{pgfscope}%
\pgfsetrectcap%
\pgfsetmiterjoin%
\pgfsetlinewidth{0.803000pt}%
\definecolor{currentstroke}{rgb}{0.000000,0.000000,0.000000}%
\pgfsetstrokecolor{currentstroke}%
\pgfsetdash{}{0pt}%
\pgfpathmoveto{\pgfqpoint{0.462269in}{0.457500in}}%
\pgfpathlineto{\pgfqpoint{5.331158in}{0.457500in}}%
\pgfusepath{stroke}%
\end{pgfscope}%
\begin{pgfscope}%
\pgfsetrectcap%
\pgfsetmiterjoin%
\pgfsetlinewidth{0.803000pt}%
\definecolor{currentstroke}{rgb}{0.000000,0.000000,0.000000}%
\pgfsetstrokecolor{currentstroke}%
\pgfsetdash{}{0pt}%
\pgfpathmoveto{\pgfqpoint{0.462269in}{3.139389in}}%
\pgfpathlineto{\pgfqpoint{5.331158in}{3.139389in}}%
\pgfusepath{stroke}%
\end{pgfscope}%
\begin{pgfscope}%
\definecolor{textcolor}{rgb}{0.000000,0.000000,0.000000}%
\pgfsetstrokecolor{textcolor}%
\pgfsetfillcolor{textcolor}%
\pgftext[x=2.896714in,y=3.222722in,,base]{\color{textcolor}{\rmfamily\fontsize{8.000000}{9.600000}\selectfont\catcode`\^=\active\def^{\ifmmode\sp\else\^{}\fi}\catcode`\%=\active\def%{\%}Self-Normalized IS for an Antipodal Target}}%
\end{pgfscope}%
\begin{pgfscope}%
\pgfsetrectcap%
\pgfsetroundjoin%
\pgfsetlinewidth{2.007500pt}%
\definecolor{currentstroke}{rgb}{0.121569,0.466667,0.705882}%
\pgfsetstrokecolor{currentstroke}%
\pgfsetdash{}{0pt}%
\pgfpathmoveto{\pgfqpoint{0.562269in}{0.772932in}}%
\pgfpathlineto{\pgfqpoint{0.673380in}{0.772932in}}%
\pgfpathlineto{\pgfqpoint{0.784491in}{0.772932in}}%
\pgfusepath{stroke}%
\end{pgfscope}%
\begin{pgfscope}%
\definecolor{textcolor}{rgb}{0.000000,0.000000,0.000000}%
\pgfsetstrokecolor{textcolor}%
\pgfsetfillcolor{textcolor}%
\pgftext[x=0.873380in,y=0.734043in,left,base]{\color{textcolor}{\rmfamily\fontsize{8.000000}{9.600000}\selectfont\catcode`\^=\active\def^{\ifmmode\sp\else\^{}\fi}\catcode`\%=\active\def%{\%}gESS profile}}%
\end{pgfscope}%
\begin{pgfscope}%
\pgfsetbuttcap%
\pgfsetroundjoin%
\pgfsetlinewidth{1.505625pt}%
\definecolor{currentstroke}{rgb}{0.000000,0.000000,0.000000}%
\pgfsetstrokecolor{currentstroke}%
\pgfsetdash{{5.550000pt}{2.400000pt}}{0.000000pt}%
\pgfpathmoveto{\pgfqpoint{0.562269in}{0.617994in}}%
\pgfpathlineto{\pgfqpoint{0.673380in}{0.617994in}}%
\pgfpathlineto{\pgfqpoint{0.784491in}{0.617994in}}%
\pgfusepath{stroke}%
\end{pgfscope}%
\begin{pgfscope}%
\definecolor{textcolor}{rgb}{0.000000,0.000000,0.000000}%
\pgfsetstrokecolor{textcolor}%
\pgfsetfillcolor{textcolor}%
\pgftext[x=0.873380in,y=0.579105in,left,base]{\color{textcolor}{\rmfamily\fontsize{8.000000}{9.600000}\selectfont\catcode`\^=\active\def^{\ifmmode\sp\else\^{}\fi}\catcode`\%=\active\def%{\%}median ordinary ESS}}%
\end{pgfscope}%
\end{pgfpicture}%
\makeatother%
\endgroup%

\caption{Self-normalized importance sampling for an antipodal two-mode target on $S^2$. The solid curve is median gESS over 40 trials, with an interquartile band.
The dashed line is median ordinary ESS. Only gESS shows how particle distinguishability changes with angular scale.}
\label{fig:app_snis}
\end{figure*}

\paragraph{Real-data normalized embedding example.}
Table~\ref{tab:digits} uses the scikit-learn handwritten digits dataset. We standardize the $8\times8$ images featurewise and fit three principal components to the full dataset. The projected points are normalized to unit length in $\mathbb R^3$. We retain the first 30 observations of each digit. ``Digits all'' assigns equal weights to these 300 embeddings. ``Digits weighted'' uses the same points with weights proportional to $\exp(2.5x_1)$. Here, $x_1$ is the first normalized coordinate. ``Digits 0/1'' uses 60 equal-weight embeddings from digits 0 and 1. This example studies reweighting a normalized representation. It does not assume that the original images are spherical data.

Figure~\ref{fig:app_digits_profiles} gives the full profiles for Table~\ref{tab:digits}. The table reports $15^\circ$, while these curves cover a range of scales. The reliability-weighted embedding has lower gESS throughout much of this range. The two-class subset also differs in profile shape from the ten-class cloud. A single weight-only ESS cannot show either difference.

\begin{figure*}[t]
\centering
%\includegraphics[width=0.96\textwidth]{hkep_app_digits_profiles.png}
\providecommand{\mathdefault}[1]{#1}
\providecommand{\mathregular}[1]{#1}
%% Creator: Matplotlib, PGF backend
%%
%% To include the figure in your LaTeX document, write
%%   \input{<filename>.pgf}
%%
%% Make sure the required packages are loaded in your preamble
%%   \usepackage{pgf}
%%
%% Also ensure that all the required font packages are loaded; for instance,
%% the lmodern package is sometimes necessary when using math font.
%%   \usepackage{lmodern}
%%
%% Figures using additional raster images can only be included by \input if
%% they are in the same directory as the main LaTeX file. For loading figures
%% from other directories you can use the `import` package
%%   \usepackage{import}
%%
%% and then include the figures with
%%   \import{<path to file>}{<filename>.pgf}
%%
%% Matplotlib used the following preamble
%%   \def\mathdefault#1{#1}
%%   \everymath=\expandafter{\the\everymath\displaystyle}
%%   \IfFileExists{scrextend.sty}{
%%     \usepackage[fontsize=8.000000pt]{scrextend}
%%   }{
%%     \renewcommand{\normalsize}{\fontsize{8.000000}{9.600000}\selectfont}
%%     \normalsize
%%   }
%%   \usepackage{amsmath,amssymb}
%%   \makeatletter\@ifpackageloaded{underscore}{}{\usepackage[strings]{underscore}}\makeatother
%%
\begingroup%
\makeatletter%
\begin{pgfpicture}%
\pgfpathrectangle{\pgfpointorigin}{\pgfqpoint{6.987632in}{3.046056in}}%
\pgfusepath{use as bounding box, clip}%
\begin{pgfscope}%
\pgfsetbuttcap%
\pgfsetmiterjoin%
\definecolor{currentfill}{rgb}{1.000000,1.000000,1.000000}%
\pgfsetfillcolor{currentfill}%
\pgfsetlinewidth{0.000000pt}%
\definecolor{currentstroke}{rgb}{1.000000,1.000000,1.000000}%
\pgfsetstrokecolor{currentstroke}%
\pgfsetdash{}{0pt}%
\pgfpathmoveto{\pgfqpoint{0.000000in}{0.000000in}}%
\pgfpathlineto{\pgfqpoint{6.987632in}{0.000000in}}%
\pgfpathlineto{\pgfqpoint{6.987632in}{3.046056in}}%
\pgfpathlineto{\pgfqpoint{0.000000in}{3.046056in}}%
\pgfpathlineto{\pgfqpoint{0.000000in}{0.000000in}}%
\pgfpathclose%
\pgfusepath{fill}%
\end{pgfscope}%
\begin{pgfscope}%
\pgfsetbuttcap%
\pgfsetmiterjoin%
\definecolor{currentfill}{rgb}{1.000000,1.000000,1.000000}%
\pgfsetfillcolor{currentfill}%
\pgfsetlinewidth{0.000000pt}%
\definecolor{currentstroke}{rgb}{0.000000,0.000000,0.000000}%
\pgfsetstrokecolor{currentstroke}%
\pgfsetstrokeopacity{0.000000}%
\pgfsetdash{}{0pt}%
\pgfpathmoveto{\pgfqpoint{0.517632in}{0.795389in}}%
\pgfpathlineto{\pgfqpoint{3.287197in}{0.795389in}}%
\pgfpathlineto{\pgfqpoint{3.287197in}{2.779389in}}%
\pgfpathlineto{\pgfqpoint{0.517632in}{2.779389in}}%
\pgfpathlineto{\pgfqpoint{0.517632in}{0.795389in}}%
\pgfpathclose%
\pgfusepath{fill}%
\end{pgfscope}%
\begin{pgfscope}%
\pgfpathrectangle{\pgfqpoint{0.517632in}{0.795389in}}{\pgfqpoint{2.769565in}{1.984000in}}%
\pgfusepath{clip}%
\pgfsetrectcap%
\pgfsetroundjoin%
\pgfsetlinewidth{0.803000pt}%
\definecolor{currentstroke}{rgb}{0.690196,0.690196,0.690196}%
\pgfsetstrokecolor{currentstroke}%
\pgfsetstrokeopacity{0.300000}%
\pgfsetdash{}{0pt}%
\pgfpathmoveto{\pgfqpoint{1.167258in}{0.795389in}}%
\pgfpathlineto{\pgfqpoint{1.167258in}{2.779389in}}%
\pgfusepath{stroke}%
\end{pgfscope}%
\begin{pgfscope}%
\pgfsetbuttcap%
\pgfsetroundjoin%
\definecolor{currentfill}{rgb}{0.000000,0.000000,0.000000}%
\pgfsetfillcolor{currentfill}%
\pgfsetlinewidth{0.803000pt}%
\definecolor{currentstroke}{rgb}{0.000000,0.000000,0.000000}%
\pgfsetstrokecolor{currentstroke}%
\pgfsetdash{}{0pt}%
\pgfsys@defobject{currentmarker}{\pgfqpoint{0.000000in}{-0.048611in}}{\pgfqpoint{0.000000in}{0.000000in}}{%
\pgfpathmoveto{\pgfqpoint{0.000000in}{0.000000in}}%
\pgfpathlineto{\pgfqpoint{0.000000in}{-0.048611in}}%
\pgfusepath{stroke,fill}%
}%
\begin{pgfscope}%
\pgfsys@transformshift{1.167258in}{0.795389in}%
\pgfsys@useobject{currentmarker}{}%
\end{pgfscope}%
\end{pgfscope}%
\begin{pgfscope}%
\definecolor{textcolor}{rgb}{0.000000,0.000000,0.000000}%
\pgfsetstrokecolor{textcolor}%
\pgfsetfillcolor{textcolor}%
\pgftext[x=1.167258in,y=0.698167in,,top]{\color{textcolor}{\rmfamily\fontsize{7.000000}{8.400000}\selectfont\catcode`\^=\active\def^{\ifmmode\sp\else\^{}\fi}\catcode`\%=\active\def%{\%}$\mathdefault{10^{1}}$}}%
\end{pgfscope}%
\begin{pgfscope}%
\pgfpathrectangle{\pgfqpoint{0.517632in}{0.795389in}}{\pgfqpoint{2.769565in}{1.984000in}}%
\pgfusepath{clip}%
\pgfsetrectcap%
\pgfsetroundjoin%
\pgfsetlinewidth{0.803000pt}%
\definecolor{currentstroke}{rgb}{0.690196,0.690196,0.690196}%
\pgfsetstrokecolor{currentstroke}%
\pgfsetstrokeopacity{0.300000}%
\pgfsetdash{}{0pt}%
\pgfpathmoveto{\pgfqpoint{2.907072in}{0.795389in}}%
\pgfpathlineto{\pgfqpoint{2.907072in}{2.779389in}}%
\pgfusepath{stroke}%
\end{pgfscope}%
\begin{pgfscope}%
\pgfsetbuttcap%
\pgfsetroundjoin%
\definecolor{currentfill}{rgb}{0.000000,0.000000,0.000000}%
\pgfsetfillcolor{currentfill}%
\pgfsetlinewidth{0.803000pt}%
\definecolor{currentstroke}{rgb}{0.000000,0.000000,0.000000}%
\pgfsetstrokecolor{currentstroke}%
\pgfsetdash{}{0pt}%
\pgfsys@defobject{currentmarker}{\pgfqpoint{0.000000in}{-0.048611in}}{\pgfqpoint{0.000000in}{0.000000in}}{%
\pgfpathmoveto{\pgfqpoint{0.000000in}{0.000000in}}%
\pgfpathlineto{\pgfqpoint{0.000000in}{-0.048611in}}%
\pgfusepath{stroke,fill}%
}%
\begin{pgfscope}%
\pgfsys@transformshift{2.907072in}{0.795389in}%
\pgfsys@useobject{currentmarker}{}%
\end{pgfscope}%
\end{pgfscope}%
\begin{pgfscope}%
\definecolor{textcolor}{rgb}{0.000000,0.000000,0.000000}%
\pgfsetstrokecolor{textcolor}%
\pgfsetfillcolor{textcolor}%
\pgftext[x=2.907072in,y=0.698167in,,top]{\color{textcolor}{\rmfamily\fontsize{7.000000}{8.400000}\selectfont\catcode`\^=\active\def^{\ifmmode\sp\else\^{}\fi}\catcode`\%=\active\def%{\%}$\mathdefault{10^{2}}$}}%
\end{pgfscope}%
\begin{pgfscope}%
\pgfsetbuttcap%
\pgfsetroundjoin%
\definecolor{currentfill}{rgb}{0.000000,0.000000,0.000000}%
\pgfsetfillcolor{currentfill}%
\pgfsetlinewidth{0.602250pt}%
\definecolor{currentstroke}{rgb}{0.000000,0.000000,0.000000}%
\pgfsetstrokecolor{currentstroke}%
\pgfsetdash{}{0pt}%
\pgfsys@defobject{currentmarker}{\pgfqpoint{0.000000in}{-0.027778in}}{\pgfqpoint{0.000000in}{0.000000in}}{%
\pgfpathmoveto{\pgfqpoint{0.000000in}{0.000000in}}%
\pgfpathlineto{\pgfqpoint{0.000000in}{-0.027778in}}%
\pgfusepath{stroke,fill}%
}%
\begin{pgfscope}%
\pgfsys@transformshift{0.643521in}{0.795389in}%
\pgfsys@useobject{currentmarker}{}%
\end{pgfscope}%
\end{pgfscope}%
\begin{pgfscope}%
\pgfsetbuttcap%
\pgfsetroundjoin%
\definecolor{currentfill}{rgb}{0.000000,0.000000,0.000000}%
\pgfsetfillcolor{currentfill}%
\pgfsetlinewidth{0.602250pt}%
\definecolor{currentstroke}{rgb}{0.000000,0.000000,0.000000}%
\pgfsetstrokecolor{currentstroke}%
\pgfsetdash{}{0pt}%
\pgfsys@defobject{currentmarker}{\pgfqpoint{0.000000in}{-0.027778in}}{\pgfqpoint{0.000000in}{0.000000in}}{%
\pgfpathmoveto{\pgfqpoint{0.000000in}{0.000000in}}%
\pgfpathlineto{\pgfqpoint{0.000000in}{-0.027778in}}%
\pgfusepath{stroke,fill}%
}%
\begin{pgfscope}%
\pgfsys@transformshift{0.781282in}{0.795389in}%
\pgfsys@useobject{currentmarker}{}%
\end{pgfscope}%
\end{pgfscope}%
\begin{pgfscope}%
\pgfsetbuttcap%
\pgfsetroundjoin%
\definecolor{currentfill}{rgb}{0.000000,0.000000,0.000000}%
\pgfsetfillcolor{currentfill}%
\pgfsetlinewidth{0.602250pt}%
\definecolor{currentstroke}{rgb}{0.000000,0.000000,0.000000}%
\pgfsetstrokecolor{currentstroke}%
\pgfsetdash{}{0pt}%
\pgfsys@defobject{currentmarker}{\pgfqpoint{0.000000in}{-0.027778in}}{\pgfqpoint{0.000000in}{0.000000in}}{%
\pgfpathmoveto{\pgfqpoint{0.000000in}{0.000000in}}%
\pgfpathlineto{\pgfqpoint{0.000000in}{-0.027778in}}%
\pgfusepath{stroke,fill}%
}%
\begin{pgfscope}%
\pgfsys@transformshift{0.897757in}{0.795389in}%
\pgfsys@useobject{currentmarker}{}%
\end{pgfscope}%
\end{pgfscope}%
\begin{pgfscope}%
\pgfsetbuttcap%
\pgfsetroundjoin%
\definecolor{currentfill}{rgb}{0.000000,0.000000,0.000000}%
\pgfsetfillcolor{currentfill}%
\pgfsetlinewidth{0.602250pt}%
\definecolor{currentstroke}{rgb}{0.000000,0.000000,0.000000}%
\pgfsetstrokecolor{currentstroke}%
\pgfsetdash{}{0pt}%
\pgfsys@defobject{currentmarker}{\pgfqpoint{0.000000in}{-0.027778in}}{\pgfqpoint{0.000000in}{0.000000in}}{%
\pgfpathmoveto{\pgfqpoint{0.000000in}{0.000000in}}%
\pgfpathlineto{\pgfqpoint{0.000000in}{-0.027778in}}%
\pgfusepath{stroke,fill}%
}%
\begin{pgfscope}%
\pgfsys@transformshift{0.998652in}{0.795389in}%
\pgfsys@useobject{currentmarker}{}%
\end{pgfscope}%
\end{pgfscope}%
\begin{pgfscope}%
\pgfsetbuttcap%
\pgfsetroundjoin%
\definecolor{currentfill}{rgb}{0.000000,0.000000,0.000000}%
\pgfsetfillcolor{currentfill}%
\pgfsetlinewidth{0.602250pt}%
\definecolor{currentstroke}{rgb}{0.000000,0.000000,0.000000}%
\pgfsetstrokecolor{currentstroke}%
\pgfsetdash{}{0pt}%
\pgfsys@defobject{currentmarker}{\pgfqpoint{0.000000in}{-0.027778in}}{\pgfqpoint{0.000000in}{0.000000in}}{%
\pgfpathmoveto{\pgfqpoint{0.000000in}{0.000000in}}%
\pgfpathlineto{\pgfqpoint{0.000000in}{-0.027778in}}%
\pgfusepath{stroke,fill}%
}%
\begin{pgfscope}%
\pgfsys@transformshift{1.087648in}{0.795389in}%
\pgfsys@useobject{currentmarker}{}%
\end{pgfscope}%
\end{pgfscope}%
\begin{pgfscope}%
\pgfsetbuttcap%
\pgfsetroundjoin%
\definecolor{currentfill}{rgb}{0.000000,0.000000,0.000000}%
\pgfsetfillcolor{currentfill}%
\pgfsetlinewidth{0.602250pt}%
\definecolor{currentstroke}{rgb}{0.000000,0.000000,0.000000}%
\pgfsetstrokecolor{currentstroke}%
\pgfsetdash{}{0pt}%
\pgfsys@defobject{currentmarker}{\pgfqpoint{0.000000in}{-0.027778in}}{\pgfqpoint{0.000000in}{0.000000in}}{%
\pgfpathmoveto{\pgfqpoint{0.000000in}{0.000000in}}%
\pgfpathlineto{\pgfqpoint{0.000000in}{-0.027778in}}%
\pgfusepath{stroke,fill}%
}%
\begin{pgfscope}%
\pgfsys@transformshift{1.690994in}{0.795389in}%
\pgfsys@useobject{currentmarker}{}%
\end{pgfscope}%
\end{pgfscope}%
\begin{pgfscope}%
\pgfsetbuttcap%
\pgfsetroundjoin%
\definecolor{currentfill}{rgb}{0.000000,0.000000,0.000000}%
\pgfsetfillcolor{currentfill}%
\pgfsetlinewidth{0.602250pt}%
\definecolor{currentstroke}{rgb}{0.000000,0.000000,0.000000}%
\pgfsetstrokecolor{currentstroke}%
\pgfsetdash{}{0pt}%
\pgfsys@defobject{currentmarker}{\pgfqpoint{0.000000in}{-0.027778in}}{\pgfqpoint{0.000000in}{0.000000in}}{%
\pgfpathmoveto{\pgfqpoint{0.000000in}{0.000000in}}%
\pgfpathlineto{\pgfqpoint{0.000000in}{-0.027778in}}%
\pgfusepath{stroke,fill}%
}%
\begin{pgfscope}%
\pgfsys@transformshift{1.997360in}{0.795389in}%
\pgfsys@useobject{currentmarker}{}%
\end{pgfscope}%
\end{pgfscope}%
\begin{pgfscope}%
\pgfsetbuttcap%
\pgfsetroundjoin%
\definecolor{currentfill}{rgb}{0.000000,0.000000,0.000000}%
\pgfsetfillcolor{currentfill}%
\pgfsetlinewidth{0.602250pt}%
\definecolor{currentstroke}{rgb}{0.000000,0.000000,0.000000}%
\pgfsetstrokecolor{currentstroke}%
\pgfsetdash{}{0pt}%
\pgfsys@defobject{currentmarker}{\pgfqpoint{0.000000in}{-0.027778in}}{\pgfqpoint{0.000000in}{0.000000in}}{%
\pgfpathmoveto{\pgfqpoint{0.000000in}{0.000000in}}%
\pgfpathlineto{\pgfqpoint{0.000000in}{-0.027778in}}%
\pgfusepath{stroke,fill}%
}%
\begin{pgfscope}%
\pgfsys@transformshift{2.214730in}{0.795389in}%
\pgfsys@useobject{currentmarker}{}%
\end{pgfscope}%
\end{pgfscope}%
\begin{pgfscope}%
\pgfsetbuttcap%
\pgfsetroundjoin%
\definecolor{currentfill}{rgb}{0.000000,0.000000,0.000000}%
\pgfsetfillcolor{currentfill}%
\pgfsetlinewidth{0.602250pt}%
\definecolor{currentstroke}{rgb}{0.000000,0.000000,0.000000}%
\pgfsetstrokecolor{currentstroke}%
\pgfsetdash{}{0pt}%
\pgfsys@defobject{currentmarker}{\pgfqpoint{0.000000in}{-0.027778in}}{\pgfqpoint{0.000000in}{0.000000in}}{%
\pgfpathmoveto{\pgfqpoint{0.000000in}{0.000000in}}%
\pgfpathlineto{\pgfqpoint{0.000000in}{-0.027778in}}%
\pgfusepath{stroke,fill}%
}%
\begin{pgfscope}%
\pgfsys@transformshift{2.383336in}{0.795389in}%
\pgfsys@useobject{currentmarker}{}%
\end{pgfscope}%
\end{pgfscope}%
\begin{pgfscope}%
\pgfsetbuttcap%
\pgfsetroundjoin%
\definecolor{currentfill}{rgb}{0.000000,0.000000,0.000000}%
\pgfsetfillcolor{currentfill}%
\pgfsetlinewidth{0.602250pt}%
\definecolor{currentstroke}{rgb}{0.000000,0.000000,0.000000}%
\pgfsetstrokecolor{currentstroke}%
\pgfsetdash{}{0pt}%
\pgfsys@defobject{currentmarker}{\pgfqpoint{0.000000in}{-0.027778in}}{\pgfqpoint{0.000000in}{0.000000in}}{%
\pgfpathmoveto{\pgfqpoint{0.000000in}{0.000000in}}%
\pgfpathlineto{\pgfqpoint{0.000000in}{-0.027778in}}%
\pgfusepath{stroke,fill}%
}%
\begin{pgfscope}%
\pgfsys@transformshift{2.521097in}{0.795389in}%
\pgfsys@useobject{currentmarker}{}%
\end{pgfscope}%
\end{pgfscope}%
\begin{pgfscope}%
\pgfsetbuttcap%
\pgfsetroundjoin%
\definecolor{currentfill}{rgb}{0.000000,0.000000,0.000000}%
\pgfsetfillcolor{currentfill}%
\pgfsetlinewidth{0.602250pt}%
\definecolor{currentstroke}{rgb}{0.000000,0.000000,0.000000}%
\pgfsetstrokecolor{currentstroke}%
\pgfsetdash{}{0pt}%
\pgfsys@defobject{currentmarker}{\pgfqpoint{0.000000in}{-0.027778in}}{\pgfqpoint{0.000000in}{0.000000in}}{%
\pgfpathmoveto{\pgfqpoint{0.000000in}{0.000000in}}%
\pgfpathlineto{\pgfqpoint{0.000000in}{-0.027778in}}%
\pgfusepath{stroke,fill}%
}%
\begin{pgfscope}%
\pgfsys@transformshift{2.637572in}{0.795389in}%
\pgfsys@useobject{currentmarker}{}%
\end{pgfscope}%
\end{pgfscope}%
\begin{pgfscope}%
\pgfsetbuttcap%
\pgfsetroundjoin%
\definecolor{currentfill}{rgb}{0.000000,0.000000,0.000000}%
\pgfsetfillcolor{currentfill}%
\pgfsetlinewidth{0.602250pt}%
\definecolor{currentstroke}{rgb}{0.000000,0.000000,0.000000}%
\pgfsetstrokecolor{currentstroke}%
\pgfsetdash{}{0pt}%
\pgfsys@defobject{currentmarker}{\pgfqpoint{0.000000in}{-0.027778in}}{\pgfqpoint{0.000000in}{0.000000in}}{%
\pgfpathmoveto{\pgfqpoint{0.000000in}{0.000000in}}%
\pgfpathlineto{\pgfqpoint{0.000000in}{-0.027778in}}%
\pgfusepath{stroke,fill}%
}%
\begin{pgfscope}%
\pgfsys@transformshift{2.738467in}{0.795389in}%
\pgfsys@useobject{currentmarker}{}%
\end{pgfscope}%
\end{pgfscope}%
\begin{pgfscope}%
\pgfsetbuttcap%
\pgfsetroundjoin%
\definecolor{currentfill}{rgb}{0.000000,0.000000,0.000000}%
\pgfsetfillcolor{currentfill}%
\pgfsetlinewidth{0.602250pt}%
\definecolor{currentstroke}{rgb}{0.000000,0.000000,0.000000}%
\pgfsetstrokecolor{currentstroke}%
\pgfsetdash{}{0pt}%
\pgfsys@defobject{currentmarker}{\pgfqpoint{0.000000in}{-0.027778in}}{\pgfqpoint{0.000000in}{0.000000in}}{%
\pgfpathmoveto{\pgfqpoint{0.000000in}{0.000000in}}%
\pgfpathlineto{\pgfqpoint{0.000000in}{-0.027778in}}%
\pgfusepath{stroke,fill}%
}%
\begin{pgfscope}%
\pgfsys@transformshift{2.827463in}{0.795389in}%
\pgfsys@useobject{currentmarker}{}%
\end{pgfscope}%
\end{pgfscope}%
\begin{pgfscope}%
\definecolor{textcolor}{rgb}{0.000000,0.000000,0.000000}%
\pgfsetstrokecolor{textcolor}%
\pgfsetfillcolor{textcolor}%
\pgftext[x=1.902415in,y=0.549000in,,top]{\color{textcolor}{\rmfamily\fontsize{8.000000}{9.600000}\selectfont\catcode`\^=\active\def^{\ifmmode\sp\else\^{}\fi}\catcode`\%=\active\def%{\%}Angular Scale $\alpha$ (Degrees)}}%
\end{pgfscope}%
\begin{pgfscope}%
\pgfpathrectangle{\pgfqpoint{0.517632in}{0.795389in}}{\pgfqpoint{2.769565in}{1.984000in}}%
\pgfusepath{clip}%
\pgfsetrectcap%
\pgfsetroundjoin%
\pgfsetlinewidth{0.803000pt}%
\definecolor{currentstroke}{rgb}{0.690196,0.690196,0.690196}%
\pgfsetstrokecolor{currentstroke}%
\pgfsetstrokeopacity{0.300000}%
\pgfsetdash{}{0pt}%
\pgfpathmoveto{\pgfqpoint{0.517632in}{0.872760in}}%
\pgfpathlineto{\pgfqpoint{3.287197in}{0.872760in}}%
\pgfusepath{stroke}%
\end{pgfscope}%
\begin{pgfscope}%
\pgfsetbuttcap%
\pgfsetroundjoin%
\definecolor{currentfill}{rgb}{0.000000,0.000000,0.000000}%
\pgfsetfillcolor{currentfill}%
\pgfsetlinewidth{0.803000pt}%
\definecolor{currentstroke}{rgb}{0.000000,0.000000,0.000000}%
\pgfsetstrokecolor{currentstroke}%
\pgfsetdash{}{0pt}%
\pgfsys@defobject{currentmarker}{\pgfqpoint{-0.048611in}{0.000000in}}{\pgfqpoint{-0.000000in}{0.000000in}}{%
\pgfpathmoveto{\pgfqpoint{-0.000000in}{0.000000in}}%
\pgfpathlineto{\pgfqpoint{-0.048611in}{0.000000in}}%
\pgfusepath{stroke,fill}%
}%
\begin{pgfscope}%
\pgfsys@transformshift{0.517632in}{0.872760in}%
\pgfsys@useobject{currentmarker}{}%
\end{pgfscope}%
\end{pgfscope}%
\begin{pgfscope}%
\definecolor{textcolor}{rgb}{0.000000,0.000000,0.000000}%
\pgfsetstrokecolor{textcolor}%
\pgfsetfillcolor{textcolor}%
\pgftext[x=0.365047in, y=0.839002in, left, base]{\color{textcolor}{\rmfamily\fontsize{7.000000}{8.400000}\selectfont\catcode`\^=\active\def^{\ifmmode\sp\else\^{}\fi}\catcode`\%=\active\def%{\%}0}}%
\end{pgfscope}%
\begin{pgfscope}%
\pgfpathrectangle{\pgfqpoint{0.517632in}{0.795389in}}{\pgfqpoint{2.769565in}{1.984000in}}%
\pgfusepath{clip}%
\pgfsetrectcap%
\pgfsetroundjoin%
\pgfsetlinewidth{0.803000pt}%
\definecolor{currentstroke}{rgb}{0.690196,0.690196,0.690196}%
\pgfsetstrokecolor{currentstroke}%
\pgfsetstrokeopacity{0.300000}%
\pgfsetdash{}{0pt}%
\pgfpathmoveto{\pgfqpoint{0.517632in}{1.109316in}}%
\pgfpathlineto{\pgfqpoint{3.287197in}{1.109316in}}%
\pgfusepath{stroke}%
\end{pgfscope}%
\begin{pgfscope}%
\pgfsetbuttcap%
\pgfsetroundjoin%
\definecolor{currentfill}{rgb}{0.000000,0.000000,0.000000}%
\pgfsetfillcolor{currentfill}%
\pgfsetlinewidth{0.803000pt}%
\definecolor{currentstroke}{rgb}{0.000000,0.000000,0.000000}%
\pgfsetstrokecolor{currentstroke}%
\pgfsetdash{}{0pt}%
\pgfsys@defobject{currentmarker}{\pgfqpoint{-0.048611in}{0.000000in}}{\pgfqpoint{-0.000000in}{0.000000in}}{%
\pgfpathmoveto{\pgfqpoint{-0.000000in}{0.000000in}}%
\pgfpathlineto{\pgfqpoint{-0.048611in}{0.000000in}}%
\pgfusepath{stroke,fill}%
}%
\begin{pgfscope}%
\pgfsys@transformshift{0.517632in}{1.109316in}%
\pgfsys@useobject{currentmarker}{}%
\end{pgfscope}%
\end{pgfscope}%
\begin{pgfscope}%
\definecolor{textcolor}{rgb}{0.000000,0.000000,0.000000}%
\pgfsetstrokecolor{textcolor}%
\pgfsetfillcolor{textcolor}%
\pgftext[x=0.309684in, y=1.075558in, left, base]{\color{textcolor}{\rmfamily\fontsize{7.000000}{8.400000}\selectfont\catcode`\^=\active\def^{\ifmmode\sp\else\^{}\fi}\catcode`\%=\active\def%{\%}20}}%
\end{pgfscope}%
\begin{pgfscope}%
\pgfpathrectangle{\pgfqpoint{0.517632in}{0.795389in}}{\pgfqpoint{2.769565in}{1.984000in}}%
\pgfusepath{clip}%
\pgfsetrectcap%
\pgfsetroundjoin%
\pgfsetlinewidth{0.803000pt}%
\definecolor{currentstroke}{rgb}{0.690196,0.690196,0.690196}%
\pgfsetstrokecolor{currentstroke}%
\pgfsetstrokeopacity{0.300000}%
\pgfsetdash{}{0pt}%
\pgfpathmoveto{\pgfqpoint{0.517632in}{1.345872in}}%
\pgfpathlineto{\pgfqpoint{3.287197in}{1.345872in}}%
\pgfusepath{stroke}%
\end{pgfscope}%
\begin{pgfscope}%
\pgfsetbuttcap%
\pgfsetroundjoin%
\definecolor{currentfill}{rgb}{0.000000,0.000000,0.000000}%
\pgfsetfillcolor{currentfill}%
\pgfsetlinewidth{0.803000pt}%
\definecolor{currentstroke}{rgb}{0.000000,0.000000,0.000000}%
\pgfsetstrokecolor{currentstroke}%
\pgfsetdash{}{0pt}%
\pgfsys@defobject{currentmarker}{\pgfqpoint{-0.048611in}{0.000000in}}{\pgfqpoint{-0.000000in}{0.000000in}}{%
\pgfpathmoveto{\pgfqpoint{-0.000000in}{0.000000in}}%
\pgfpathlineto{\pgfqpoint{-0.048611in}{0.000000in}}%
\pgfusepath{stroke,fill}%
}%
\begin{pgfscope}%
\pgfsys@transformshift{0.517632in}{1.345872in}%
\pgfsys@useobject{currentmarker}{}%
\end{pgfscope}%
\end{pgfscope}%
\begin{pgfscope}%
\definecolor{textcolor}{rgb}{0.000000,0.000000,0.000000}%
\pgfsetstrokecolor{textcolor}%
\pgfsetfillcolor{textcolor}%
\pgftext[x=0.309684in, y=1.312114in, left, base]{\color{textcolor}{\rmfamily\fontsize{7.000000}{8.400000}\selectfont\catcode`\^=\active\def^{\ifmmode\sp\else\^{}\fi}\catcode`\%=\active\def%{\%}40}}%
\end{pgfscope}%
\begin{pgfscope}%
\pgfpathrectangle{\pgfqpoint{0.517632in}{0.795389in}}{\pgfqpoint{2.769565in}{1.984000in}}%
\pgfusepath{clip}%
\pgfsetrectcap%
\pgfsetroundjoin%
\pgfsetlinewidth{0.803000pt}%
\definecolor{currentstroke}{rgb}{0.690196,0.690196,0.690196}%
\pgfsetstrokecolor{currentstroke}%
\pgfsetstrokeopacity{0.300000}%
\pgfsetdash{}{0pt}%
\pgfpathmoveto{\pgfqpoint{0.517632in}{1.582427in}}%
\pgfpathlineto{\pgfqpoint{3.287197in}{1.582427in}}%
\pgfusepath{stroke}%
\end{pgfscope}%
\begin{pgfscope}%
\pgfsetbuttcap%
\pgfsetroundjoin%
\definecolor{currentfill}{rgb}{0.000000,0.000000,0.000000}%
\pgfsetfillcolor{currentfill}%
\pgfsetlinewidth{0.803000pt}%
\definecolor{currentstroke}{rgb}{0.000000,0.000000,0.000000}%
\pgfsetstrokecolor{currentstroke}%
\pgfsetdash{}{0pt}%
\pgfsys@defobject{currentmarker}{\pgfqpoint{-0.048611in}{0.000000in}}{\pgfqpoint{-0.000000in}{0.000000in}}{%
\pgfpathmoveto{\pgfqpoint{-0.000000in}{0.000000in}}%
\pgfpathlineto{\pgfqpoint{-0.048611in}{0.000000in}}%
\pgfusepath{stroke,fill}%
}%
\begin{pgfscope}%
\pgfsys@transformshift{0.517632in}{1.582427in}%
\pgfsys@useobject{currentmarker}{}%
\end{pgfscope}%
\end{pgfscope}%
\begin{pgfscope}%
\definecolor{textcolor}{rgb}{0.000000,0.000000,0.000000}%
\pgfsetstrokecolor{textcolor}%
\pgfsetfillcolor{textcolor}%
\pgftext[x=0.309684in, y=1.548670in, left, base]{\color{textcolor}{\rmfamily\fontsize{7.000000}{8.400000}\selectfont\catcode`\^=\active\def^{\ifmmode\sp\else\^{}\fi}\catcode`\%=\active\def%{\%}60}}%
\end{pgfscope}%
\begin{pgfscope}%
\pgfpathrectangle{\pgfqpoint{0.517632in}{0.795389in}}{\pgfqpoint{2.769565in}{1.984000in}}%
\pgfusepath{clip}%
\pgfsetrectcap%
\pgfsetroundjoin%
\pgfsetlinewidth{0.803000pt}%
\definecolor{currentstroke}{rgb}{0.690196,0.690196,0.690196}%
\pgfsetstrokecolor{currentstroke}%
\pgfsetstrokeopacity{0.300000}%
\pgfsetdash{}{0pt}%
\pgfpathmoveto{\pgfqpoint{0.517632in}{1.818983in}}%
\pgfpathlineto{\pgfqpoint{3.287197in}{1.818983in}}%
\pgfusepath{stroke}%
\end{pgfscope}%
\begin{pgfscope}%
\pgfsetbuttcap%
\pgfsetroundjoin%
\definecolor{currentfill}{rgb}{0.000000,0.000000,0.000000}%
\pgfsetfillcolor{currentfill}%
\pgfsetlinewidth{0.803000pt}%
\definecolor{currentstroke}{rgb}{0.000000,0.000000,0.000000}%
\pgfsetstrokecolor{currentstroke}%
\pgfsetdash{}{0pt}%
\pgfsys@defobject{currentmarker}{\pgfqpoint{-0.048611in}{0.000000in}}{\pgfqpoint{-0.000000in}{0.000000in}}{%
\pgfpathmoveto{\pgfqpoint{-0.000000in}{0.000000in}}%
\pgfpathlineto{\pgfqpoint{-0.048611in}{0.000000in}}%
\pgfusepath{stroke,fill}%
}%
\begin{pgfscope}%
\pgfsys@transformshift{0.517632in}{1.818983in}%
\pgfsys@useobject{currentmarker}{}%
\end{pgfscope}%
\end{pgfscope}%
\begin{pgfscope}%
\definecolor{textcolor}{rgb}{0.000000,0.000000,0.000000}%
\pgfsetstrokecolor{textcolor}%
\pgfsetfillcolor{textcolor}%
\pgftext[x=0.309684in, y=1.785225in, left, base]{\color{textcolor}{\rmfamily\fontsize{7.000000}{8.400000}\selectfont\catcode`\^=\active\def^{\ifmmode\sp\else\^{}\fi}\catcode`\%=\active\def%{\%}80}}%
\end{pgfscope}%
\begin{pgfscope}%
\pgfpathrectangle{\pgfqpoint{0.517632in}{0.795389in}}{\pgfqpoint{2.769565in}{1.984000in}}%
\pgfusepath{clip}%
\pgfsetrectcap%
\pgfsetroundjoin%
\pgfsetlinewidth{0.803000pt}%
\definecolor{currentstroke}{rgb}{0.690196,0.690196,0.690196}%
\pgfsetstrokecolor{currentstroke}%
\pgfsetstrokeopacity{0.300000}%
\pgfsetdash{}{0pt}%
\pgfpathmoveto{\pgfqpoint{0.517632in}{2.055539in}}%
\pgfpathlineto{\pgfqpoint{3.287197in}{2.055539in}}%
\pgfusepath{stroke}%
\end{pgfscope}%
\begin{pgfscope}%
\pgfsetbuttcap%
\pgfsetroundjoin%
\definecolor{currentfill}{rgb}{0.000000,0.000000,0.000000}%
\pgfsetfillcolor{currentfill}%
\pgfsetlinewidth{0.803000pt}%
\definecolor{currentstroke}{rgb}{0.000000,0.000000,0.000000}%
\pgfsetstrokecolor{currentstroke}%
\pgfsetdash{}{0pt}%
\pgfsys@defobject{currentmarker}{\pgfqpoint{-0.048611in}{0.000000in}}{\pgfqpoint{-0.000000in}{0.000000in}}{%
\pgfpathmoveto{\pgfqpoint{-0.000000in}{0.000000in}}%
\pgfpathlineto{\pgfqpoint{-0.048611in}{0.000000in}}%
\pgfusepath{stroke,fill}%
}%
\begin{pgfscope}%
\pgfsys@transformshift{0.517632in}{2.055539in}%
\pgfsys@useobject{currentmarker}{}%
\end{pgfscope}%
\end{pgfscope}%
\begin{pgfscope}%
\definecolor{textcolor}{rgb}{0.000000,0.000000,0.000000}%
\pgfsetstrokecolor{textcolor}%
\pgfsetfillcolor{textcolor}%
\pgftext[x=0.254321in, y=2.021781in, left, base]{\color{textcolor}{\rmfamily\fontsize{7.000000}{8.400000}\selectfont\catcode`\^=\active\def^{\ifmmode\sp\else\^{}\fi}\catcode`\%=\active\def%{\%}100}}%
\end{pgfscope}%
\begin{pgfscope}%
\pgfpathrectangle{\pgfqpoint{0.517632in}{0.795389in}}{\pgfqpoint{2.769565in}{1.984000in}}%
\pgfusepath{clip}%
\pgfsetrectcap%
\pgfsetroundjoin%
\pgfsetlinewidth{0.803000pt}%
\definecolor{currentstroke}{rgb}{0.690196,0.690196,0.690196}%
\pgfsetstrokecolor{currentstroke}%
\pgfsetstrokeopacity{0.300000}%
\pgfsetdash{}{0pt}%
\pgfpathmoveto{\pgfqpoint{0.517632in}{2.292094in}}%
\pgfpathlineto{\pgfqpoint{3.287197in}{2.292094in}}%
\pgfusepath{stroke}%
\end{pgfscope}%
\begin{pgfscope}%
\pgfsetbuttcap%
\pgfsetroundjoin%
\definecolor{currentfill}{rgb}{0.000000,0.000000,0.000000}%
\pgfsetfillcolor{currentfill}%
\pgfsetlinewidth{0.803000pt}%
\definecolor{currentstroke}{rgb}{0.000000,0.000000,0.000000}%
\pgfsetstrokecolor{currentstroke}%
\pgfsetdash{}{0pt}%
\pgfsys@defobject{currentmarker}{\pgfqpoint{-0.048611in}{0.000000in}}{\pgfqpoint{-0.000000in}{0.000000in}}{%
\pgfpathmoveto{\pgfqpoint{-0.000000in}{0.000000in}}%
\pgfpathlineto{\pgfqpoint{-0.048611in}{0.000000in}}%
\pgfusepath{stroke,fill}%
}%
\begin{pgfscope}%
\pgfsys@transformshift{0.517632in}{2.292094in}%
\pgfsys@useobject{currentmarker}{}%
\end{pgfscope}%
\end{pgfscope}%
\begin{pgfscope}%
\definecolor{textcolor}{rgb}{0.000000,0.000000,0.000000}%
\pgfsetstrokecolor{textcolor}%
\pgfsetfillcolor{textcolor}%
\pgftext[x=0.254321in, y=2.258337in, left, base]{\color{textcolor}{\rmfamily\fontsize{7.000000}{8.400000}\selectfont\catcode`\^=\active\def^{\ifmmode\sp\else\^{}\fi}\catcode`\%=\active\def%{\%}120}}%
\end{pgfscope}%
\begin{pgfscope}%
\pgfpathrectangle{\pgfqpoint{0.517632in}{0.795389in}}{\pgfqpoint{2.769565in}{1.984000in}}%
\pgfusepath{clip}%
\pgfsetrectcap%
\pgfsetroundjoin%
\pgfsetlinewidth{0.803000pt}%
\definecolor{currentstroke}{rgb}{0.690196,0.690196,0.690196}%
\pgfsetstrokecolor{currentstroke}%
\pgfsetstrokeopacity{0.300000}%
\pgfsetdash{}{0pt}%
\pgfpathmoveto{\pgfqpoint{0.517632in}{2.528650in}}%
\pgfpathlineto{\pgfqpoint{3.287197in}{2.528650in}}%
\pgfusepath{stroke}%
\end{pgfscope}%
\begin{pgfscope}%
\pgfsetbuttcap%
\pgfsetroundjoin%
\definecolor{currentfill}{rgb}{0.000000,0.000000,0.000000}%
\pgfsetfillcolor{currentfill}%
\pgfsetlinewidth{0.803000pt}%
\definecolor{currentstroke}{rgb}{0.000000,0.000000,0.000000}%
\pgfsetstrokecolor{currentstroke}%
\pgfsetdash{}{0pt}%
\pgfsys@defobject{currentmarker}{\pgfqpoint{-0.048611in}{0.000000in}}{\pgfqpoint{-0.000000in}{0.000000in}}{%
\pgfpathmoveto{\pgfqpoint{-0.000000in}{0.000000in}}%
\pgfpathlineto{\pgfqpoint{-0.048611in}{0.000000in}}%
\pgfusepath{stroke,fill}%
}%
\begin{pgfscope}%
\pgfsys@transformshift{0.517632in}{2.528650in}%
\pgfsys@useobject{currentmarker}{}%
\end{pgfscope}%
\end{pgfscope}%
\begin{pgfscope}%
\definecolor{textcolor}{rgb}{0.000000,0.000000,0.000000}%
\pgfsetstrokecolor{textcolor}%
\pgfsetfillcolor{textcolor}%
\pgftext[x=0.254321in, y=2.494892in, left, base]{\color{textcolor}{\rmfamily\fontsize{7.000000}{8.400000}\selectfont\catcode`\^=\active\def^{\ifmmode\sp\else\^{}\fi}\catcode`\%=\active\def%{\%}140}}%
\end{pgfscope}%
\begin{pgfscope}%
\pgfpathrectangle{\pgfqpoint{0.517632in}{0.795389in}}{\pgfqpoint{2.769565in}{1.984000in}}%
\pgfusepath{clip}%
\pgfsetrectcap%
\pgfsetroundjoin%
\pgfsetlinewidth{0.803000pt}%
\definecolor{currentstroke}{rgb}{0.690196,0.690196,0.690196}%
\pgfsetstrokecolor{currentstroke}%
\pgfsetstrokeopacity{0.300000}%
\pgfsetdash{}{0pt}%
\pgfpathmoveto{\pgfqpoint{0.517632in}{2.765206in}}%
\pgfpathlineto{\pgfqpoint{3.287197in}{2.765206in}}%
\pgfusepath{stroke}%
\end{pgfscope}%
\begin{pgfscope}%
\pgfsetbuttcap%
\pgfsetroundjoin%
\definecolor{currentfill}{rgb}{0.000000,0.000000,0.000000}%
\pgfsetfillcolor{currentfill}%
\pgfsetlinewidth{0.803000pt}%
\definecolor{currentstroke}{rgb}{0.000000,0.000000,0.000000}%
\pgfsetstrokecolor{currentstroke}%
\pgfsetdash{}{0pt}%
\pgfsys@defobject{currentmarker}{\pgfqpoint{-0.048611in}{0.000000in}}{\pgfqpoint{-0.000000in}{0.000000in}}{%
\pgfpathmoveto{\pgfqpoint{-0.000000in}{0.000000in}}%
\pgfpathlineto{\pgfqpoint{-0.048611in}{0.000000in}}%
\pgfusepath{stroke,fill}%
}%
\begin{pgfscope}%
\pgfsys@transformshift{0.517632in}{2.765206in}%
\pgfsys@useobject{currentmarker}{}%
\end{pgfscope}%
\end{pgfscope}%
\begin{pgfscope}%
\definecolor{textcolor}{rgb}{0.000000,0.000000,0.000000}%
\pgfsetstrokecolor{textcolor}%
\pgfsetfillcolor{textcolor}%
\pgftext[x=0.254321in, y=2.731448in, left, base]{\color{textcolor}{\rmfamily\fontsize{7.000000}{8.400000}\selectfont\catcode`\^=\active\def^{\ifmmode\sp\else\^{}\fi}\catcode`\%=\active\def%{\%}160}}%
\end{pgfscope}%
\begin{pgfscope}%
\definecolor{textcolor}{rgb}{0.000000,0.000000,0.000000}%
\pgfsetstrokecolor{textcolor}%
\pgfsetfillcolor{textcolor}%
\pgftext[x=0.198766in,y=1.787389in,,bottom,rotate=90.000000]{\color{textcolor}{\rmfamily\fontsize{8.000000}{9.600000}\selectfont\catcode`\^=\active\def^{\ifmmode\sp\else\^{}\fi}\catcode`\%=\active\def%{\%}Geometric ESS}}%
\end{pgfscope}%
\begin{pgfscope}%
\pgfpathrectangle{\pgfqpoint{0.517632in}{0.795389in}}{\pgfqpoint{2.769565in}{1.984000in}}%
\pgfusepath{clip}%
\pgfsetrectcap%
\pgfsetroundjoin%
\pgfsetlinewidth{1.806750pt}%
\definecolor{currentstroke}{rgb}{0.121569,0.466667,0.705882}%
\pgfsetstrokecolor{currentstroke}%
\pgfsetdash{}{0pt}%
\pgfpathmoveto{\pgfqpoint{0.643521in}{2.689207in}}%
\pgfpathlineto{\pgfqpoint{0.677546in}{2.608806in}}%
\pgfpathlineto{\pgfqpoint{0.711570in}{2.529299in}}%
\pgfpathlineto{\pgfqpoint{0.745594in}{2.450949in}}%
\pgfpathlineto{\pgfqpoint{0.779618in}{2.374000in}}%
\pgfpathlineto{\pgfqpoint{0.813642in}{2.298676in}}%
\pgfpathlineto{\pgfqpoint{0.847666in}{2.225185in}}%
\pgfpathlineto{\pgfqpoint{0.881690in}{2.153714in}}%
\pgfpathlineto{\pgfqpoint{0.915715in}{2.084428in}}%
\pgfpathlineto{\pgfqpoint{0.949739in}{2.017469in}}%
\pgfpathlineto{\pgfqpoint{0.983763in}{1.952958in}}%
\pgfpathlineto{\pgfqpoint{1.017787in}{1.890990in}}%
\pgfpathlineto{\pgfqpoint{1.051811in}{1.831635in}}%
\pgfpathlineto{\pgfqpoint{1.085835in}{1.774939in}}%
\pgfpathlineto{\pgfqpoint{1.119859in}{1.720924in}}%
\pgfpathlineto{\pgfqpoint{1.153884in}{1.669584in}}%
\pgfpathlineto{\pgfqpoint{1.187908in}{1.620896in}}%
\pgfpathlineto{\pgfqpoint{1.221932in}{1.574812in}}%
\pgfpathlineto{\pgfqpoint{1.255956in}{1.531268in}}%
\pgfpathlineto{\pgfqpoint{1.289980in}{1.490181in}}%
\pgfpathlineto{\pgfqpoint{1.324004in}{1.451459in}}%
\pgfpathlineto{\pgfqpoint{1.358028in}{1.414999in}}%
\pgfpathlineto{\pgfqpoint{1.392053in}{1.380691in}}%
\pgfpathlineto{\pgfqpoint{1.426077in}{1.348423in}}%
\pgfpathlineto{\pgfqpoint{1.460101in}{1.318080in}}%
\pgfpathlineto{\pgfqpoint{1.494125in}{1.289550in}}%
\pgfpathlineto{\pgfqpoint{1.528149in}{1.262725in}}%
\pgfpathlineto{\pgfqpoint{1.562173in}{1.237499in}}%
\pgfpathlineto{\pgfqpoint{1.596197in}{1.213772in}}%
\pgfpathlineto{\pgfqpoint{1.630222in}{1.191452in}}%
\pgfpathlineto{\pgfqpoint{1.664246in}{1.170451in}}%
\pgfpathlineto{\pgfqpoint{1.698270in}{1.150687in}}%
\pgfpathlineto{\pgfqpoint{1.732294in}{1.132087in}}%
\pgfpathlineto{\pgfqpoint{1.766318in}{1.114582in}}%
\pgfpathlineto{\pgfqpoint{1.800342in}{1.098109in}}%
\pgfpathlineto{\pgfqpoint{1.834366in}{1.082612in}}%
\pgfpathlineto{\pgfqpoint{1.868391in}{1.068039in}}%
\pgfpathlineto{\pgfqpoint{1.902415in}{1.054343in}}%
\pgfpathlineto{\pgfqpoint{1.936439in}{1.041482in}}%
\pgfpathlineto{\pgfqpoint{1.970463in}{1.029415in}}%
\pgfpathlineto{\pgfqpoint{2.004487in}{1.018108in}}%
\pgfpathlineto{\pgfqpoint{2.038511in}{1.007525in}}%
\pgfpathlineto{\pgfqpoint{2.072535in}{0.997635in}}%
\pgfpathlineto{\pgfqpoint{2.106560in}{0.988404in}}%
\pgfpathlineto{\pgfqpoint{2.140584in}{0.979804in}}%
\pgfpathlineto{\pgfqpoint{2.174608in}{0.971802in}}%
\pgfpathlineto{\pgfqpoint{2.208632in}{0.964369in}}%
\pgfpathlineto{\pgfqpoint{2.242656in}{0.957473in}}%
\pgfpathlineto{\pgfqpoint{2.276680in}{0.951085in}}%
\pgfpathlineto{\pgfqpoint{2.310704in}{0.945174in}}%
\pgfpathlineto{\pgfqpoint{2.344729in}{0.939711in}}%
\pgfpathlineto{\pgfqpoint{2.378753in}{0.934666in}}%
\pgfpathlineto{\pgfqpoint{2.412777in}{0.930012in}}%
\pgfpathlineto{\pgfqpoint{2.446801in}{0.925722in}}%
\pgfpathlineto{\pgfqpoint{2.480825in}{0.921770in}}%
\pgfpathlineto{\pgfqpoint{2.514849in}{0.918132in}}%
\pgfpathlineto{\pgfqpoint{2.548873in}{0.914785in}}%
\pgfpathlineto{\pgfqpoint{2.582898in}{0.911708in}}%
\pgfpathlineto{\pgfqpoint{2.616922in}{0.908880in}}%
\pgfpathlineto{\pgfqpoint{2.650946in}{0.906284in}}%
\pgfpathlineto{\pgfqpoint{2.684970in}{0.903902in}}%
\pgfpathlineto{\pgfqpoint{2.718994in}{0.901718in}}%
\pgfpathlineto{\pgfqpoint{2.753018in}{0.899719in}}%
\pgfpathlineto{\pgfqpoint{2.787042in}{0.897890in}}%
\pgfpathlineto{\pgfqpoint{2.821067in}{0.896219in}}%
\pgfpathlineto{\pgfqpoint{2.855091in}{0.894696in}}%
\pgfpathlineto{\pgfqpoint{2.889115in}{0.893310in}}%
\pgfpathlineto{\pgfqpoint{2.923139in}{0.892053in}}%
\pgfpathlineto{\pgfqpoint{2.957163in}{0.890917in}}%
\pgfpathlineto{\pgfqpoint{2.991187in}{0.889895in}}%
\pgfpathlineto{\pgfqpoint{3.025211in}{0.888983in}}%
\pgfpathlineto{\pgfqpoint{3.059236in}{0.888176in}}%
\pgfpathlineto{\pgfqpoint{3.093260in}{0.887469in}}%
\pgfpathlineto{\pgfqpoint{3.127284in}{0.886859in}}%
\pgfpathlineto{\pgfqpoint{3.161308in}{0.886341in}}%
\pgfusepath{stroke}%
\end{pgfscope}%
\begin{pgfscope}%
\pgfpathrectangle{\pgfqpoint{0.517632in}{0.795389in}}{\pgfqpoint{2.769565in}{1.984000in}}%
\pgfusepath{clip}%
\pgfsetrectcap%
\pgfsetroundjoin%
\pgfsetlinewidth{1.806750pt}%
\definecolor{currentstroke}{rgb}{1.000000,0.498039,0.054902}%
\pgfsetstrokecolor{currentstroke}%
\pgfsetdash{}{0pt}%
\pgfpathmoveto{\pgfqpoint{0.643521in}{1.605651in}}%
\pgfpathlineto{\pgfqpoint{0.677546in}{1.574830in}}%
\pgfpathlineto{\pgfqpoint{0.711570in}{1.544140in}}%
\pgfpathlineto{\pgfqpoint{0.745594in}{1.513697in}}%
\pgfpathlineto{\pgfqpoint{0.779618in}{1.483613in}}%
\pgfpathlineto{\pgfqpoint{0.813642in}{1.454001in}}%
\pgfpathlineto{\pgfqpoint{0.847666in}{1.424967in}}%
\pgfpathlineto{\pgfqpoint{0.881690in}{1.396611in}}%
\pgfpathlineto{\pgfqpoint{0.915715in}{1.369022in}}%
\pgfpathlineto{\pgfqpoint{0.949739in}{1.342278in}}%
\pgfpathlineto{\pgfqpoint{0.983763in}{1.316446in}}%
\pgfpathlineto{\pgfqpoint{1.017787in}{1.291577in}}%
\pgfpathlineto{\pgfqpoint{1.051811in}{1.267712in}}%
\pgfpathlineto{\pgfqpoint{1.085835in}{1.244877in}}%
\pgfpathlineto{\pgfqpoint{1.119859in}{1.223090in}}%
\pgfpathlineto{\pgfqpoint{1.153884in}{1.202355in}}%
\pgfpathlineto{\pgfqpoint{1.187908in}{1.182669in}}%
\pgfpathlineto{\pgfqpoint{1.221932in}{1.164018in}}%
\pgfpathlineto{\pgfqpoint{1.255956in}{1.146385in}}%
\pgfpathlineto{\pgfqpoint{1.289980in}{1.129742in}}%
\pgfpathlineto{\pgfqpoint{1.324004in}{1.114057in}}%
\pgfpathlineto{\pgfqpoint{1.358028in}{1.099295in}}%
\pgfpathlineto{\pgfqpoint{1.392053in}{1.085415in}}%
\pgfpathlineto{\pgfqpoint{1.426077in}{1.072374in}}%
\pgfpathlineto{\pgfqpoint{1.460101in}{1.060127in}}%
\pgfpathlineto{\pgfqpoint{1.494125in}{1.048628in}}%
\pgfpathlineto{\pgfqpoint{1.528149in}{1.037831in}}%
\pgfpathlineto{\pgfqpoint{1.562173in}{1.027690in}}%
\pgfpathlineto{\pgfqpoint{1.596197in}{1.018159in}}%
\pgfpathlineto{\pgfqpoint{1.630222in}{1.009198in}}%
\pgfpathlineto{\pgfqpoint{1.664246in}{1.000765in}}%
\pgfpathlineto{\pgfqpoint{1.698270in}{0.992824in}}%
\pgfpathlineto{\pgfqpoint{1.732294in}{0.985341in}}%
\pgfpathlineto{\pgfqpoint{1.766318in}{0.978288in}}%
\pgfpathlineto{\pgfqpoint{1.800342in}{0.971638in}}%
\pgfpathlineto{\pgfqpoint{1.834366in}{0.965368in}}%
\pgfpathlineto{\pgfqpoint{1.868391in}{0.959460in}}%
\pgfpathlineto{\pgfqpoint{1.902415in}{0.953897in}}%
\pgfpathlineto{\pgfqpoint{1.936439in}{0.948664in}}%
\pgfpathlineto{\pgfqpoint{1.970463in}{0.943749in}}%
\pgfpathlineto{\pgfqpoint{2.004487in}{0.939139in}}%
\pgfpathlineto{\pgfqpoint{2.038511in}{0.934824in}}%
\pgfpathlineto{\pgfqpoint{2.072535in}{0.930792in}}%
\pgfpathlineto{\pgfqpoint{2.106560in}{0.927033in}}%
\pgfpathlineto{\pgfqpoint{2.140584in}{0.923534in}}%
\pgfpathlineto{\pgfqpoint{2.174608in}{0.920285in}}%
\pgfpathlineto{\pgfqpoint{2.208632in}{0.917273in}}%
\pgfpathlineto{\pgfqpoint{2.242656in}{0.914485in}}%
\pgfpathlineto{\pgfqpoint{2.276680in}{0.911909in}}%
\pgfpathlineto{\pgfqpoint{2.310704in}{0.909533in}}%
\pgfpathlineto{\pgfqpoint{2.344729in}{0.907342in}}%
\pgfpathlineto{\pgfqpoint{2.378753in}{0.905325in}}%
\pgfpathlineto{\pgfqpoint{2.412777in}{0.903470in}}%
\pgfpathlineto{\pgfqpoint{2.446801in}{0.901764in}}%
\pgfpathlineto{\pgfqpoint{2.480825in}{0.900196in}}%
\pgfpathlineto{\pgfqpoint{2.514849in}{0.898755in}}%
\pgfpathlineto{\pgfqpoint{2.548873in}{0.897431in}}%
\pgfpathlineto{\pgfqpoint{2.582898in}{0.896215in}}%
\pgfpathlineto{\pgfqpoint{2.616922in}{0.895097in}}%
\pgfpathlineto{\pgfqpoint{2.650946in}{0.894069in}}%
\pgfpathlineto{\pgfqpoint{2.684970in}{0.893123in}}%
\pgfpathlineto{\pgfqpoint{2.718994in}{0.892252in}}%
\pgfpathlineto{\pgfqpoint{2.753018in}{0.891450in}}%
\pgfpathlineto{\pgfqpoint{2.787042in}{0.890710in}}%
\pgfpathlineto{\pgfqpoint{2.821067in}{0.890027in}}%
\pgfpathlineto{\pgfqpoint{2.855091in}{0.889396in}}%
\pgfpathlineto{\pgfqpoint{2.889115in}{0.888812in}}%
\pgfpathlineto{\pgfqpoint{2.923139in}{0.888272in}}%
\pgfpathlineto{\pgfqpoint{2.957163in}{0.887773in}}%
\pgfpathlineto{\pgfqpoint{2.991187in}{0.887312in}}%
\pgfpathlineto{\pgfqpoint{3.025211in}{0.886889in}}%
\pgfpathlineto{\pgfqpoint{3.059236in}{0.886502in}}%
\pgfpathlineto{\pgfqpoint{3.093260in}{0.886153in}}%
\pgfpathlineto{\pgfqpoint{3.127284in}{0.885842in}}%
\pgfpathlineto{\pgfqpoint{3.161308in}{0.885571in}}%
\pgfusepath{stroke}%
\end{pgfscope}%
\begin{pgfscope}%
\pgfpathrectangle{\pgfqpoint{0.517632in}{0.795389in}}{\pgfqpoint{2.769565in}{1.984000in}}%
\pgfusepath{clip}%
\pgfsetrectcap%
\pgfsetroundjoin%
\pgfsetlinewidth{1.806750pt}%
\definecolor{currentstroke}{rgb}{0.172549,0.627451,0.172549}%
\pgfsetstrokecolor{currentstroke}%
\pgfsetdash{}{0pt}%
\pgfpathmoveto{\pgfqpoint{0.643521in}{1.299938in}}%
\pgfpathlineto{\pgfqpoint{0.677546in}{1.284165in}}%
\pgfpathlineto{\pgfqpoint{0.711570in}{1.268224in}}%
\pgfpathlineto{\pgfqpoint{0.745594in}{1.252187in}}%
\pgfpathlineto{\pgfqpoint{0.779618in}{1.236134in}}%
\pgfpathlineto{\pgfqpoint{0.813642in}{1.220142in}}%
\pgfpathlineto{\pgfqpoint{0.847666in}{1.204286in}}%
\pgfpathlineto{\pgfqpoint{0.881690in}{1.188639in}}%
\pgfpathlineto{\pgfqpoint{0.915715in}{1.173263in}}%
\pgfpathlineto{\pgfqpoint{0.949739in}{1.158215in}}%
\pgfpathlineto{\pgfqpoint{0.983763in}{1.143546in}}%
\pgfpathlineto{\pgfqpoint{1.017787in}{1.129298in}}%
\pgfpathlineto{\pgfqpoint{1.051811in}{1.115507in}}%
\pgfpathlineto{\pgfqpoint{1.085835in}{1.102203in}}%
\pgfpathlineto{\pgfqpoint{1.119859in}{1.089411in}}%
\pgfpathlineto{\pgfqpoint{1.153884in}{1.077153in}}%
\pgfpathlineto{\pgfqpoint{1.187908in}{1.065442in}}%
\pgfpathlineto{\pgfqpoint{1.221932in}{1.054291in}}%
\pgfpathlineto{\pgfqpoint{1.255956in}{1.043703in}}%
\pgfpathlineto{\pgfqpoint{1.289980in}{1.033679in}}%
\pgfpathlineto{\pgfqpoint{1.324004in}{1.024215in}}%
\pgfpathlineto{\pgfqpoint{1.358028in}{1.015301in}}%
\pgfpathlineto{\pgfqpoint{1.392053in}{1.006923in}}%
\pgfpathlineto{\pgfqpoint{1.426077in}{0.999066in}}%
\pgfpathlineto{\pgfqpoint{1.460101in}{0.991708in}}%
\pgfpathlineto{\pgfqpoint{1.494125in}{0.984829in}}%
\pgfpathlineto{\pgfqpoint{1.528149in}{0.978405in}}%
\pgfpathlineto{\pgfqpoint{1.562173in}{0.972412in}}%
\pgfpathlineto{\pgfqpoint{1.596197in}{0.966826in}}%
\pgfpathlineto{\pgfqpoint{1.630222in}{0.961622in}}%
\pgfpathlineto{\pgfqpoint{1.664246in}{0.956777in}}%
\pgfpathlineto{\pgfqpoint{1.698270in}{0.952268in}}%
\pgfpathlineto{\pgfqpoint{1.732294in}{0.948073in}}%
\pgfpathlineto{\pgfqpoint{1.766318in}{0.944172in}}%
\pgfpathlineto{\pgfqpoint{1.800342in}{0.940544in}}%
\pgfpathlineto{\pgfqpoint{1.834366in}{0.937169in}}%
\pgfpathlineto{\pgfqpoint{1.868391in}{0.934030in}}%
\pgfpathlineto{\pgfqpoint{1.902415in}{0.931110in}}%
\pgfpathlineto{\pgfqpoint{1.936439in}{0.928392in}}%
\pgfpathlineto{\pgfqpoint{1.970463in}{0.925860in}}%
\pgfpathlineto{\pgfqpoint{2.004487in}{0.923499in}}%
\pgfpathlineto{\pgfqpoint{2.038511in}{0.921296in}}%
\pgfpathlineto{\pgfqpoint{2.072535in}{0.919235in}}%
\pgfpathlineto{\pgfqpoint{2.106560in}{0.917304in}}%
\pgfpathlineto{\pgfqpoint{2.140584in}{0.915491in}}%
\pgfpathlineto{\pgfqpoint{2.174608in}{0.913783in}}%
\pgfpathlineto{\pgfqpoint{2.208632in}{0.912168in}}%
\pgfpathlineto{\pgfqpoint{2.242656in}{0.910636in}}%
\pgfpathlineto{\pgfqpoint{2.276680in}{0.909176in}}%
\pgfpathlineto{\pgfqpoint{2.310704in}{0.907779in}}%
\pgfpathlineto{\pgfqpoint{2.344729in}{0.906435in}}%
\pgfpathlineto{\pgfqpoint{2.378753in}{0.905139in}}%
\pgfpathlineto{\pgfqpoint{2.412777in}{0.903883in}}%
\pgfpathlineto{\pgfqpoint{2.446801in}{0.902663in}}%
\pgfpathlineto{\pgfqpoint{2.480825in}{0.901476in}}%
\pgfpathlineto{\pgfqpoint{2.514849in}{0.900319in}}%
\pgfpathlineto{\pgfqpoint{2.548873in}{0.899193in}}%
\pgfpathlineto{\pgfqpoint{2.582898in}{0.898097in}}%
\pgfpathlineto{\pgfqpoint{2.616922in}{0.897035in}}%
\pgfpathlineto{\pgfqpoint{2.650946in}{0.896007in}}%
\pgfpathlineto{\pgfqpoint{2.684970in}{0.895017in}}%
\pgfpathlineto{\pgfqpoint{2.718994in}{0.894066in}}%
\pgfpathlineto{\pgfqpoint{2.753018in}{0.893157in}}%
\pgfpathlineto{\pgfqpoint{2.787042in}{0.892292in}}%
\pgfpathlineto{\pgfqpoint{2.821067in}{0.891471in}}%
\pgfpathlineto{\pgfqpoint{2.855091in}{0.890695in}}%
\pgfpathlineto{\pgfqpoint{2.889115in}{0.889964in}}%
\pgfpathlineto{\pgfqpoint{2.923139in}{0.889278in}}%
\pgfpathlineto{\pgfqpoint{2.957163in}{0.888638in}}%
\pgfpathlineto{\pgfqpoint{2.991187in}{0.888045in}}%
\pgfpathlineto{\pgfqpoint{3.025211in}{0.887498in}}%
\pgfpathlineto{\pgfqpoint{3.059236in}{0.887000in}}%
\pgfpathlineto{\pgfqpoint{3.093260in}{0.886552in}}%
\pgfpathlineto{\pgfqpoint{3.127284in}{0.886155in}}%
\pgfpathlineto{\pgfqpoint{3.161308in}{0.885811in}}%
\pgfusepath{stroke}%
\end{pgfscope}%
\begin{pgfscope}%
\pgfsetrectcap%
\pgfsetmiterjoin%
\pgfsetlinewidth{0.803000pt}%
\definecolor{currentstroke}{rgb}{0.000000,0.000000,0.000000}%
\pgfsetstrokecolor{currentstroke}%
\pgfsetdash{}{0pt}%
\pgfpathmoveto{\pgfqpoint{0.517632in}{0.795389in}}%
\pgfpathlineto{\pgfqpoint{0.517632in}{2.779389in}}%
\pgfusepath{stroke}%
\end{pgfscope}%
\begin{pgfscope}%
\pgfsetrectcap%
\pgfsetmiterjoin%
\pgfsetlinewidth{0.803000pt}%
\definecolor{currentstroke}{rgb}{0.000000,0.000000,0.000000}%
\pgfsetstrokecolor{currentstroke}%
\pgfsetdash{}{0pt}%
\pgfpathmoveto{\pgfqpoint{3.287197in}{0.795389in}}%
\pgfpathlineto{\pgfqpoint{3.287197in}{2.779389in}}%
\pgfusepath{stroke}%
\end{pgfscope}%
\begin{pgfscope}%
\pgfsetrectcap%
\pgfsetmiterjoin%
\pgfsetlinewidth{0.803000pt}%
\definecolor{currentstroke}{rgb}{0.000000,0.000000,0.000000}%
\pgfsetstrokecolor{currentstroke}%
\pgfsetdash{}{0pt}%
\pgfpathmoveto{\pgfqpoint{0.517632in}{0.795389in}}%
\pgfpathlineto{\pgfqpoint{3.287197in}{0.795389in}}%
\pgfusepath{stroke}%
\end{pgfscope}%
\begin{pgfscope}%
\pgfsetrectcap%
\pgfsetmiterjoin%
\pgfsetlinewidth{0.803000pt}%
\definecolor{currentstroke}{rgb}{0.000000,0.000000,0.000000}%
\pgfsetstrokecolor{currentstroke}%
\pgfsetdash{}{0pt}%
\pgfpathmoveto{\pgfqpoint{0.517632in}{2.779389in}}%
\pgfpathlineto{\pgfqpoint{3.287197in}{2.779389in}}%
\pgfusepath{stroke}%
\end{pgfscope}%
\begin{pgfscope}%
\definecolor{textcolor}{rgb}{0.000000,0.000000,0.000000}%
\pgfsetstrokecolor{textcolor}%
\pgfsetfillcolor{textcolor}%
\pgftext[x=1.902415in,y=2.862722in,,base]{\color{textcolor}{\rmfamily\fontsize{8.000000}{9.600000}\selectfont\catcode`\^=\active\def^{\ifmmode\sp\else\^{}\fi}\catcode`\%=\active\def%{\%}(a) Normalized-Embedding gESS}}%
\end{pgfscope}%
\begin{pgfscope}%
\pgfsetbuttcap%
\pgfsetmiterjoin%
\definecolor{currentfill}{rgb}{1.000000,1.000000,1.000000}%
\pgfsetfillcolor{currentfill}%
\pgfsetlinewidth{0.000000pt}%
\definecolor{currentstroke}{rgb}{0.000000,0.000000,0.000000}%
\pgfsetstrokecolor{currentstroke}%
\pgfsetstrokeopacity{0.000000}%
\pgfsetdash{}{0pt}%
\pgfpathmoveto{\pgfqpoint{4.118067in}{0.795389in}}%
\pgfpathlineto{\pgfqpoint{6.887632in}{0.795389in}}%
\pgfpathlineto{\pgfqpoint{6.887632in}{2.779389in}}%
\pgfpathlineto{\pgfqpoint{4.118067in}{2.779389in}}%
\pgfpathlineto{\pgfqpoint{4.118067in}{0.795389in}}%
\pgfpathclose%
\pgfusepath{fill}%
\end{pgfscope}%
\begin{pgfscope}%
\pgfpathrectangle{\pgfqpoint{4.118067in}{0.795389in}}{\pgfqpoint{2.769565in}{1.984000in}}%
\pgfusepath{clip}%
\pgfsetrectcap%
\pgfsetroundjoin%
\pgfsetlinewidth{0.803000pt}%
\definecolor{currentstroke}{rgb}{0.690196,0.690196,0.690196}%
\pgfsetstrokecolor{currentstroke}%
\pgfsetstrokeopacity{0.300000}%
\pgfsetdash{}{0pt}%
\pgfpathmoveto{\pgfqpoint{4.767693in}{0.795389in}}%
\pgfpathlineto{\pgfqpoint{4.767693in}{2.779389in}}%
\pgfusepath{stroke}%
\end{pgfscope}%
\begin{pgfscope}%
\pgfsetbuttcap%
\pgfsetroundjoin%
\definecolor{currentfill}{rgb}{0.000000,0.000000,0.000000}%
\pgfsetfillcolor{currentfill}%
\pgfsetlinewidth{0.803000pt}%
\definecolor{currentstroke}{rgb}{0.000000,0.000000,0.000000}%
\pgfsetstrokecolor{currentstroke}%
\pgfsetdash{}{0pt}%
\pgfsys@defobject{currentmarker}{\pgfqpoint{0.000000in}{-0.048611in}}{\pgfqpoint{0.000000in}{0.000000in}}{%
\pgfpathmoveto{\pgfqpoint{0.000000in}{0.000000in}}%
\pgfpathlineto{\pgfqpoint{0.000000in}{-0.048611in}}%
\pgfusepath{stroke,fill}%
}%
\begin{pgfscope}%
\pgfsys@transformshift{4.767693in}{0.795389in}%
\pgfsys@useobject{currentmarker}{}%
\end{pgfscope}%
\end{pgfscope}%
\begin{pgfscope}%
\definecolor{textcolor}{rgb}{0.000000,0.000000,0.000000}%
\pgfsetstrokecolor{textcolor}%
\pgfsetfillcolor{textcolor}%
\pgftext[x=4.767693in,y=0.698167in,,top]{\color{textcolor}{\rmfamily\fontsize{7.000000}{8.400000}\selectfont\catcode`\^=\active\def^{\ifmmode\sp\else\^{}\fi}\catcode`\%=\active\def%{\%}$\mathdefault{10^{1}}$}}%
\end{pgfscope}%
\begin{pgfscope}%
\pgfpathrectangle{\pgfqpoint{4.118067in}{0.795389in}}{\pgfqpoint{2.769565in}{1.984000in}}%
\pgfusepath{clip}%
\pgfsetrectcap%
\pgfsetroundjoin%
\pgfsetlinewidth{0.803000pt}%
\definecolor{currentstroke}{rgb}{0.690196,0.690196,0.690196}%
\pgfsetstrokecolor{currentstroke}%
\pgfsetstrokeopacity{0.300000}%
\pgfsetdash{}{0pt}%
\pgfpathmoveto{\pgfqpoint{6.507507in}{0.795389in}}%
\pgfpathlineto{\pgfqpoint{6.507507in}{2.779389in}}%
\pgfusepath{stroke}%
\end{pgfscope}%
\begin{pgfscope}%
\pgfsetbuttcap%
\pgfsetroundjoin%
\definecolor{currentfill}{rgb}{0.000000,0.000000,0.000000}%
\pgfsetfillcolor{currentfill}%
\pgfsetlinewidth{0.803000pt}%
\definecolor{currentstroke}{rgb}{0.000000,0.000000,0.000000}%
\pgfsetstrokecolor{currentstroke}%
\pgfsetdash{}{0pt}%
\pgfsys@defobject{currentmarker}{\pgfqpoint{0.000000in}{-0.048611in}}{\pgfqpoint{0.000000in}{0.000000in}}{%
\pgfpathmoveto{\pgfqpoint{0.000000in}{0.000000in}}%
\pgfpathlineto{\pgfqpoint{0.000000in}{-0.048611in}}%
\pgfusepath{stroke,fill}%
}%
\begin{pgfscope}%
\pgfsys@transformshift{6.507507in}{0.795389in}%
\pgfsys@useobject{currentmarker}{}%
\end{pgfscope}%
\end{pgfscope}%
\begin{pgfscope}%
\definecolor{textcolor}{rgb}{0.000000,0.000000,0.000000}%
\pgfsetstrokecolor{textcolor}%
\pgfsetfillcolor{textcolor}%
\pgftext[x=6.507507in,y=0.698167in,,top]{\color{textcolor}{\rmfamily\fontsize{7.000000}{8.400000}\selectfont\catcode`\^=\active\def^{\ifmmode\sp\else\^{}\fi}\catcode`\%=\active\def%{\%}$\mathdefault{10^{2}}$}}%
\end{pgfscope}%
\begin{pgfscope}%
\pgfsetbuttcap%
\pgfsetroundjoin%
\definecolor{currentfill}{rgb}{0.000000,0.000000,0.000000}%
\pgfsetfillcolor{currentfill}%
\pgfsetlinewidth{0.602250pt}%
\definecolor{currentstroke}{rgb}{0.000000,0.000000,0.000000}%
\pgfsetstrokecolor{currentstroke}%
\pgfsetdash{}{0pt}%
\pgfsys@defobject{currentmarker}{\pgfqpoint{0.000000in}{-0.027778in}}{\pgfqpoint{0.000000in}{0.000000in}}{%
\pgfpathmoveto{\pgfqpoint{0.000000in}{0.000000in}}%
\pgfpathlineto{\pgfqpoint{0.000000in}{-0.027778in}}%
\pgfusepath{stroke,fill}%
}%
\begin{pgfscope}%
\pgfsys@transformshift{4.243956in}{0.795389in}%
\pgfsys@useobject{currentmarker}{}%
\end{pgfscope}%
\end{pgfscope}%
\begin{pgfscope}%
\pgfsetbuttcap%
\pgfsetroundjoin%
\definecolor{currentfill}{rgb}{0.000000,0.000000,0.000000}%
\pgfsetfillcolor{currentfill}%
\pgfsetlinewidth{0.602250pt}%
\definecolor{currentstroke}{rgb}{0.000000,0.000000,0.000000}%
\pgfsetstrokecolor{currentstroke}%
\pgfsetdash{}{0pt}%
\pgfsys@defobject{currentmarker}{\pgfqpoint{0.000000in}{-0.027778in}}{\pgfqpoint{0.000000in}{0.000000in}}{%
\pgfpathmoveto{\pgfqpoint{0.000000in}{0.000000in}}%
\pgfpathlineto{\pgfqpoint{0.000000in}{-0.027778in}}%
\pgfusepath{stroke,fill}%
}%
\begin{pgfscope}%
\pgfsys@transformshift{4.381717in}{0.795389in}%
\pgfsys@useobject{currentmarker}{}%
\end{pgfscope}%
\end{pgfscope}%
\begin{pgfscope}%
\pgfsetbuttcap%
\pgfsetroundjoin%
\definecolor{currentfill}{rgb}{0.000000,0.000000,0.000000}%
\pgfsetfillcolor{currentfill}%
\pgfsetlinewidth{0.602250pt}%
\definecolor{currentstroke}{rgb}{0.000000,0.000000,0.000000}%
\pgfsetstrokecolor{currentstroke}%
\pgfsetdash{}{0pt}%
\pgfsys@defobject{currentmarker}{\pgfqpoint{0.000000in}{-0.027778in}}{\pgfqpoint{0.000000in}{0.000000in}}{%
\pgfpathmoveto{\pgfqpoint{0.000000in}{0.000000in}}%
\pgfpathlineto{\pgfqpoint{0.000000in}{-0.027778in}}%
\pgfusepath{stroke,fill}%
}%
\begin{pgfscope}%
\pgfsys@transformshift{4.498192in}{0.795389in}%
\pgfsys@useobject{currentmarker}{}%
\end{pgfscope}%
\end{pgfscope}%
\begin{pgfscope}%
\pgfsetbuttcap%
\pgfsetroundjoin%
\definecolor{currentfill}{rgb}{0.000000,0.000000,0.000000}%
\pgfsetfillcolor{currentfill}%
\pgfsetlinewidth{0.602250pt}%
\definecolor{currentstroke}{rgb}{0.000000,0.000000,0.000000}%
\pgfsetstrokecolor{currentstroke}%
\pgfsetdash{}{0pt}%
\pgfsys@defobject{currentmarker}{\pgfqpoint{0.000000in}{-0.027778in}}{\pgfqpoint{0.000000in}{0.000000in}}{%
\pgfpathmoveto{\pgfqpoint{0.000000in}{0.000000in}}%
\pgfpathlineto{\pgfqpoint{0.000000in}{-0.027778in}}%
\pgfusepath{stroke,fill}%
}%
\begin{pgfscope}%
\pgfsys@transformshift{4.599087in}{0.795389in}%
\pgfsys@useobject{currentmarker}{}%
\end{pgfscope}%
\end{pgfscope}%
\begin{pgfscope}%
\pgfsetbuttcap%
\pgfsetroundjoin%
\definecolor{currentfill}{rgb}{0.000000,0.000000,0.000000}%
\pgfsetfillcolor{currentfill}%
\pgfsetlinewidth{0.602250pt}%
\definecolor{currentstroke}{rgb}{0.000000,0.000000,0.000000}%
\pgfsetstrokecolor{currentstroke}%
\pgfsetdash{}{0pt}%
\pgfsys@defobject{currentmarker}{\pgfqpoint{0.000000in}{-0.027778in}}{\pgfqpoint{0.000000in}{0.000000in}}{%
\pgfpathmoveto{\pgfqpoint{0.000000in}{0.000000in}}%
\pgfpathlineto{\pgfqpoint{0.000000in}{-0.027778in}}%
\pgfusepath{stroke,fill}%
}%
\begin{pgfscope}%
\pgfsys@transformshift{4.688083in}{0.795389in}%
\pgfsys@useobject{currentmarker}{}%
\end{pgfscope}%
\end{pgfscope}%
\begin{pgfscope}%
\pgfsetbuttcap%
\pgfsetroundjoin%
\definecolor{currentfill}{rgb}{0.000000,0.000000,0.000000}%
\pgfsetfillcolor{currentfill}%
\pgfsetlinewidth{0.602250pt}%
\definecolor{currentstroke}{rgb}{0.000000,0.000000,0.000000}%
\pgfsetstrokecolor{currentstroke}%
\pgfsetdash{}{0pt}%
\pgfsys@defobject{currentmarker}{\pgfqpoint{0.000000in}{-0.027778in}}{\pgfqpoint{0.000000in}{0.000000in}}{%
\pgfpathmoveto{\pgfqpoint{0.000000in}{0.000000in}}%
\pgfpathlineto{\pgfqpoint{0.000000in}{-0.027778in}}%
\pgfusepath{stroke,fill}%
}%
\begin{pgfscope}%
\pgfsys@transformshift{5.291429in}{0.795389in}%
\pgfsys@useobject{currentmarker}{}%
\end{pgfscope}%
\end{pgfscope}%
\begin{pgfscope}%
\pgfsetbuttcap%
\pgfsetroundjoin%
\definecolor{currentfill}{rgb}{0.000000,0.000000,0.000000}%
\pgfsetfillcolor{currentfill}%
\pgfsetlinewidth{0.602250pt}%
\definecolor{currentstroke}{rgb}{0.000000,0.000000,0.000000}%
\pgfsetstrokecolor{currentstroke}%
\pgfsetdash{}{0pt}%
\pgfsys@defobject{currentmarker}{\pgfqpoint{0.000000in}{-0.027778in}}{\pgfqpoint{0.000000in}{0.000000in}}{%
\pgfpathmoveto{\pgfqpoint{0.000000in}{0.000000in}}%
\pgfpathlineto{\pgfqpoint{0.000000in}{-0.027778in}}%
\pgfusepath{stroke,fill}%
}%
\begin{pgfscope}%
\pgfsys@transformshift{5.597795in}{0.795389in}%
\pgfsys@useobject{currentmarker}{}%
\end{pgfscope}%
\end{pgfscope}%
\begin{pgfscope}%
\pgfsetbuttcap%
\pgfsetroundjoin%
\definecolor{currentfill}{rgb}{0.000000,0.000000,0.000000}%
\pgfsetfillcolor{currentfill}%
\pgfsetlinewidth{0.602250pt}%
\definecolor{currentstroke}{rgb}{0.000000,0.000000,0.000000}%
\pgfsetstrokecolor{currentstroke}%
\pgfsetdash{}{0pt}%
\pgfsys@defobject{currentmarker}{\pgfqpoint{0.000000in}{-0.027778in}}{\pgfqpoint{0.000000in}{0.000000in}}{%
\pgfpathmoveto{\pgfqpoint{0.000000in}{0.000000in}}%
\pgfpathlineto{\pgfqpoint{0.000000in}{-0.027778in}}%
\pgfusepath{stroke,fill}%
}%
\begin{pgfscope}%
\pgfsys@transformshift{5.815165in}{0.795389in}%
\pgfsys@useobject{currentmarker}{}%
\end{pgfscope}%
\end{pgfscope}%
\begin{pgfscope}%
\pgfsetbuttcap%
\pgfsetroundjoin%
\definecolor{currentfill}{rgb}{0.000000,0.000000,0.000000}%
\pgfsetfillcolor{currentfill}%
\pgfsetlinewidth{0.602250pt}%
\definecolor{currentstroke}{rgb}{0.000000,0.000000,0.000000}%
\pgfsetstrokecolor{currentstroke}%
\pgfsetdash{}{0pt}%
\pgfsys@defobject{currentmarker}{\pgfqpoint{0.000000in}{-0.027778in}}{\pgfqpoint{0.000000in}{0.000000in}}{%
\pgfpathmoveto{\pgfqpoint{0.000000in}{0.000000in}}%
\pgfpathlineto{\pgfqpoint{0.000000in}{-0.027778in}}%
\pgfusepath{stroke,fill}%
}%
\begin{pgfscope}%
\pgfsys@transformshift{5.983771in}{0.795389in}%
\pgfsys@useobject{currentmarker}{}%
\end{pgfscope}%
\end{pgfscope}%
\begin{pgfscope}%
\pgfsetbuttcap%
\pgfsetroundjoin%
\definecolor{currentfill}{rgb}{0.000000,0.000000,0.000000}%
\pgfsetfillcolor{currentfill}%
\pgfsetlinewidth{0.602250pt}%
\definecolor{currentstroke}{rgb}{0.000000,0.000000,0.000000}%
\pgfsetstrokecolor{currentstroke}%
\pgfsetdash{}{0pt}%
\pgfsys@defobject{currentmarker}{\pgfqpoint{0.000000in}{-0.027778in}}{\pgfqpoint{0.000000in}{0.000000in}}{%
\pgfpathmoveto{\pgfqpoint{0.000000in}{0.000000in}}%
\pgfpathlineto{\pgfqpoint{0.000000in}{-0.027778in}}%
\pgfusepath{stroke,fill}%
}%
\begin{pgfscope}%
\pgfsys@transformshift{6.121531in}{0.795389in}%
\pgfsys@useobject{currentmarker}{}%
\end{pgfscope}%
\end{pgfscope}%
\begin{pgfscope}%
\pgfsetbuttcap%
\pgfsetroundjoin%
\definecolor{currentfill}{rgb}{0.000000,0.000000,0.000000}%
\pgfsetfillcolor{currentfill}%
\pgfsetlinewidth{0.602250pt}%
\definecolor{currentstroke}{rgb}{0.000000,0.000000,0.000000}%
\pgfsetstrokecolor{currentstroke}%
\pgfsetdash{}{0pt}%
\pgfsys@defobject{currentmarker}{\pgfqpoint{0.000000in}{-0.027778in}}{\pgfqpoint{0.000000in}{0.000000in}}{%
\pgfpathmoveto{\pgfqpoint{0.000000in}{0.000000in}}%
\pgfpathlineto{\pgfqpoint{0.000000in}{-0.027778in}}%
\pgfusepath{stroke,fill}%
}%
\begin{pgfscope}%
\pgfsys@transformshift{6.238006in}{0.795389in}%
\pgfsys@useobject{currentmarker}{}%
\end{pgfscope}%
\end{pgfscope}%
\begin{pgfscope}%
\pgfsetbuttcap%
\pgfsetroundjoin%
\definecolor{currentfill}{rgb}{0.000000,0.000000,0.000000}%
\pgfsetfillcolor{currentfill}%
\pgfsetlinewidth{0.602250pt}%
\definecolor{currentstroke}{rgb}{0.000000,0.000000,0.000000}%
\pgfsetstrokecolor{currentstroke}%
\pgfsetdash{}{0pt}%
\pgfsys@defobject{currentmarker}{\pgfqpoint{0.000000in}{-0.027778in}}{\pgfqpoint{0.000000in}{0.000000in}}{%
\pgfpathmoveto{\pgfqpoint{0.000000in}{0.000000in}}%
\pgfpathlineto{\pgfqpoint{0.000000in}{-0.027778in}}%
\pgfusepath{stroke,fill}%
}%
\begin{pgfscope}%
\pgfsys@transformshift{6.338902in}{0.795389in}%
\pgfsys@useobject{currentmarker}{}%
\end{pgfscope}%
\end{pgfscope}%
\begin{pgfscope}%
\pgfsetbuttcap%
\pgfsetroundjoin%
\definecolor{currentfill}{rgb}{0.000000,0.000000,0.000000}%
\pgfsetfillcolor{currentfill}%
\pgfsetlinewidth{0.602250pt}%
\definecolor{currentstroke}{rgb}{0.000000,0.000000,0.000000}%
\pgfsetstrokecolor{currentstroke}%
\pgfsetdash{}{0pt}%
\pgfsys@defobject{currentmarker}{\pgfqpoint{0.000000in}{-0.027778in}}{\pgfqpoint{0.000000in}{0.000000in}}{%
\pgfpathmoveto{\pgfqpoint{0.000000in}{0.000000in}}%
\pgfpathlineto{\pgfqpoint{0.000000in}{-0.027778in}}%
\pgfusepath{stroke,fill}%
}%
\begin{pgfscope}%
\pgfsys@transformshift{6.427898in}{0.795389in}%
\pgfsys@useobject{currentmarker}{}%
\end{pgfscope}%
\end{pgfscope}%
\begin{pgfscope}%
\definecolor{textcolor}{rgb}{0.000000,0.000000,0.000000}%
\pgfsetstrokecolor{textcolor}%
\pgfsetfillcolor{textcolor}%
\pgftext[x=5.502849in,y=0.549000in,,top]{\color{textcolor}{\rmfamily\fontsize{8.000000}{9.600000}\selectfont\catcode`\^=\active\def^{\ifmmode\sp\else\^{}\fi}\catcode`\%=\active\def%{\%}Angular Scale $\alpha$ (Degrees)}}%
\end{pgfscope}%
\begin{pgfscope}%
\pgfpathrectangle{\pgfqpoint{4.118067in}{0.795389in}}{\pgfqpoint{2.769565in}{1.984000in}}%
\pgfusepath{clip}%
\pgfsetrectcap%
\pgfsetroundjoin%
\pgfsetlinewidth{0.803000pt}%
\definecolor{currentstroke}{rgb}{0.690196,0.690196,0.690196}%
\pgfsetstrokecolor{currentstroke}%
\pgfsetstrokeopacity{0.300000}%
\pgfsetdash{}{0pt}%
\pgfpathmoveto{\pgfqpoint{4.118067in}{1.140789in}}%
\pgfpathlineto{\pgfqpoint{6.887632in}{1.140789in}}%
\pgfusepath{stroke}%
\end{pgfscope}%
\begin{pgfscope}%
\pgfsetbuttcap%
\pgfsetroundjoin%
\definecolor{currentfill}{rgb}{0.000000,0.000000,0.000000}%
\pgfsetfillcolor{currentfill}%
\pgfsetlinewidth{0.803000pt}%
\definecolor{currentstroke}{rgb}{0.000000,0.000000,0.000000}%
\pgfsetstrokecolor{currentstroke}%
\pgfsetdash{}{0pt}%
\pgfsys@defobject{currentmarker}{\pgfqpoint{-0.048611in}{0.000000in}}{\pgfqpoint{-0.000000in}{0.000000in}}{%
\pgfpathmoveto{\pgfqpoint{-0.000000in}{0.000000in}}%
\pgfpathlineto{\pgfqpoint{-0.048611in}{0.000000in}}%
\pgfusepath{stroke,fill}%
}%
\begin{pgfscope}%
\pgfsys@transformshift{4.118067in}{1.140789in}%
\pgfsys@useobject{currentmarker}{}%
\end{pgfscope}%
\end{pgfscope}%
\begin{pgfscope}%
\definecolor{textcolor}{rgb}{0.000000,0.000000,0.000000}%
\pgfsetstrokecolor{textcolor}%
\pgfsetfillcolor{textcolor}%
\pgftext[x=3.878676in, y=1.107031in, left, base]{\color{textcolor}{\rmfamily\fontsize{7.000000}{8.400000}\selectfont\catcode`\^=\active\def^{\ifmmode\sp\else\^{}\fi}\catcode`\%=\active\def%{\%}0.2}}%
\end{pgfscope}%
\begin{pgfscope}%
\pgfpathrectangle{\pgfqpoint{4.118067in}{0.795389in}}{\pgfqpoint{2.769565in}{1.984000in}}%
\pgfusepath{clip}%
\pgfsetrectcap%
\pgfsetroundjoin%
\pgfsetlinewidth{0.803000pt}%
\definecolor{currentstroke}{rgb}{0.690196,0.690196,0.690196}%
\pgfsetstrokecolor{currentstroke}%
\pgfsetstrokeopacity{0.300000}%
\pgfsetdash{}{0pt}%
\pgfpathmoveto{\pgfqpoint{4.118067in}{1.529594in}}%
\pgfpathlineto{\pgfqpoint{6.887632in}{1.529594in}}%
\pgfusepath{stroke}%
\end{pgfscope}%
\begin{pgfscope}%
\pgfsetbuttcap%
\pgfsetroundjoin%
\definecolor{currentfill}{rgb}{0.000000,0.000000,0.000000}%
\pgfsetfillcolor{currentfill}%
\pgfsetlinewidth{0.803000pt}%
\definecolor{currentstroke}{rgb}{0.000000,0.000000,0.000000}%
\pgfsetstrokecolor{currentstroke}%
\pgfsetdash{}{0pt}%
\pgfsys@defobject{currentmarker}{\pgfqpoint{-0.048611in}{0.000000in}}{\pgfqpoint{-0.000000in}{0.000000in}}{%
\pgfpathmoveto{\pgfqpoint{-0.000000in}{0.000000in}}%
\pgfpathlineto{\pgfqpoint{-0.048611in}{0.000000in}}%
\pgfusepath{stroke,fill}%
}%
\begin{pgfscope}%
\pgfsys@transformshift{4.118067in}{1.529594in}%
\pgfsys@useobject{currentmarker}{}%
\end{pgfscope}%
\end{pgfscope}%
\begin{pgfscope}%
\definecolor{textcolor}{rgb}{0.000000,0.000000,0.000000}%
\pgfsetstrokecolor{textcolor}%
\pgfsetfillcolor{textcolor}%
\pgftext[x=3.878676in, y=1.495837in, left, base]{\color{textcolor}{\rmfamily\fontsize{7.000000}{8.400000}\selectfont\catcode`\^=\active\def^{\ifmmode\sp\else\^{}\fi}\catcode`\%=\active\def%{\%}0.4}}%
\end{pgfscope}%
\begin{pgfscope}%
\pgfpathrectangle{\pgfqpoint{4.118067in}{0.795389in}}{\pgfqpoint{2.769565in}{1.984000in}}%
\pgfusepath{clip}%
\pgfsetrectcap%
\pgfsetroundjoin%
\pgfsetlinewidth{0.803000pt}%
\definecolor{currentstroke}{rgb}{0.690196,0.690196,0.690196}%
\pgfsetstrokecolor{currentstroke}%
\pgfsetstrokeopacity{0.300000}%
\pgfsetdash{}{0pt}%
\pgfpathmoveto{\pgfqpoint{4.118067in}{1.918400in}}%
\pgfpathlineto{\pgfqpoint{6.887632in}{1.918400in}}%
\pgfusepath{stroke}%
\end{pgfscope}%
\begin{pgfscope}%
\pgfsetbuttcap%
\pgfsetroundjoin%
\definecolor{currentfill}{rgb}{0.000000,0.000000,0.000000}%
\pgfsetfillcolor{currentfill}%
\pgfsetlinewidth{0.803000pt}%
\definecolor{currentstroke}{rgb}{0.000000,0.000000,0.000000}%
\pgfsetstrokecolor{currentstroke}%
\pgfsetdash{}{0pt}%
\pgfsys@defobject{currentmarker}{\pgfqpoint{-0.048611in}{0.000000in}}{\pgfqpoint{-0.000000in}{0.000000in}}{%
\pgfpathmoveto{\pgfqpoint{-0.000000in}{0.000000in}}%
\pgfpathlineto{\pgfqpoint{-0.048611in}{0.000000in}}%
\pgfusepath{stroke,fill}%
}%
\begin{pgfscope}%
\pgfsys@transformshift{4.118067in}{1.918400in}%
\pgfsys@useobject{currentmarker}{}%
\end{pgfscope}%
\end{pgfscope}%
\begin{pgfscope}%
\definecolor{textcolor}{rgb}{0.000000,0.000000,0.000000}%
\pgfsetstrokecolor{textcolor}%
\pgfsetfillcolor{textcolor}%
\pgftext[x=3.878676in, y=1.884643in, left, base]{\color{textcolor}{\rmfamily\fontsize{7.000000}{8.400000}\selectfont\catcode`\^=\active\def^{\ifmmode\sp\else\^{}\fi}\catcode`\%=\active\def%{\%}0.6}}%
\end{pgfscope}%
\begin{pgfscope}%
\pgfpathrectangle{\pgfqpoint{4.118067in}{0.795389in}}{\pgfqpoint{2.769565in}{1.984000in}}%
\pgfusepath{clip}%
\pgfsetrectcap%
\pgfsetroundjoin%
\pgfsetlinewidth{0.803000pt}%
\definecolor{currentstroke}{rgb}{0.690196,0.690196,0.690196}%
\pgfsetstrokecolor{currentstroke}%
\pgfsetstrokeopacity{0.300000}%
\pgfsetdash{}{0pt}%
\pgfpathmoveto{\pgfqpoint{4.118067in}{2.307206in}}%
\pgfpathlineto{\pgfqpoint{6.887632in}{2.307206in}}%
\pgfusepath{stroke}%
\end{pgfscope}%
\begin{pgfscope}%
\pgfsetbuttcap%
\pgfsetroundjoin%
\definecolor{currentfill}{rgb}{0.000000,0.000000,0.000000}%
\pgfsetfillcolor{currentfill}%
\pgfsetlinewidth{0.803000pt}%
\definecolor{currentstroke}{rgb}{0.000000,0.000000,0.000000}%
\pgfsetstrokecolor{currentstroke}%
\pgfsetdash{}{0pt}%
\pgfsys@defobject{currentmarker}{\pgfqpoint{-0.048611in}{0.000000in}}{\pgfqpoint{-0.000000in}{0.000000in}}{%
\pgfpathmoveto{\pgfqpoint{-0.000000in}{0.000000in}}%
\pgfpathlineto{\pgfqpoint{-0.048611in}{0.000000in}}%
\pgfusepath{stroke,fill}%
}%
\begin{pgfscope}%
\pgfsys@transformshift{4.118067in}{2.307206in}%
\pgfsys@useobject{currentmarker}{}%
\end{pgfscope}%
\end{pgfscope}%
\begin{pgfscope}%
\definecolor{textcolor}{rgb}{0.000000,0.000000,0.000000}%
\pgfsetstrokecolor{textcolor}%
\pgfsetfillcolor{textcolor}%
\pgftext[x=3.878676in, y=2.273449in, left, base]{\color{textcolor}{\rmfamily\fontsize{7.000000}{8.400000}\selectfont\catcode`\^=\active\def^{\ifmmode\sp\else\^{}\fi}\catcode`\%=\active\def%{\%}0.8}}%
\end{pgfscope}%
\begin{pgfscope}%
\pgfpathrectangle{\pgfqpoint{4.118067in}{0.795389in}}{\pgfqpoint{2.769565in}{1.984000in}}%
\pgfusepath{clip}%
\pgfsetrectcap%
\pgfsetroundjoin%
\pgfsetlinewidth{0.803000pt}%
\definecolor{currentstroke}{rgb}{0.690196,0.690196,0.690196}%
\pgfsetstrokecolor{currentstroke}%
\pgfsetstrokeopacity{0.300000}%
\pgfsetdash{}{0pt}%
\pgfpathmoveto{\pgfqpoint{4.118067in}{2.696012in}}%
\pgfpathlineto{\pgfqpoint{6.887632in}{2.696012in}}%
\pgfusepath{stroke}%
\end{pgfscope}%
\begin{pgfscope}%
\pgfsetbuttcap%
\pgfsetroundjoin%
\definecolor{currentfill}{rgb}{0.000000,0.000000,0.000000}%
\pgfsetfillcolor{currentfill}%
\pgfsetlinewidth{0.803000pt}%
\definecolor{currentstroke}{rgb}{0.000000,0.000000,0.000000}%
\pgfsetstrokecolor{currentstroke}%
\pgfsetdash{}{0pt}%
\pgfsys@defobject{currentmarker}{\pgfqpoint{-0.048611in}{0.000000in}}{\pgfqpoint{-0.000000in}{0.000000in}}{%
\pgfpathmoveto{\pgfqpoint{-0.000000in}{0.000000in}}%
\pgfpathlineto{\pgfqpoint{-0.048611in}{0.000000in}}%
\pgfusepath{stroke,fill}%
}%
\begin{pgfscope}%
\pgfsys@transformshift{4.118067in}{2.696012in}%
\pgfsys@useobject{currentmarker}{}%
\end{pgfscope}%
\end{pgfscope}%
\begin{pgfscope}%
\definecolor{textcolor}{rgb}{0.000000,0.000000,0.000000}%
\pgfsetstrokecolor{textcolor}%
\pgfsetfillcolor{textcolor}%
\pgftext[x=3.878676in, y=2.662255in, left, base]{\color{textcolor}{\rmfamily\fontsize{7.000000}{8.400000}\selectfont\catcode`\^=\active\def^{\ifmmode\sp\else\^{}\fi}\catcode`\%=\active\def%{\%}1.0}}%
\end{pgfscope}%
\begin{pgfscope}%
\definecolor{textcolor}{rgb}{0.000000,0.000000,0.000000}%
\pgfsetstrokecolor{textcolor}%
\pgfsetfillcolor{textcolor}%
\pgftext[x=3.823120in,y=1.787389in,,bottom,rotate=90.000000]{\color{textcolor}{\rmfamily\fontsize{8.000000}{9.600000}\selectfont\catcode`\^=\active\def^{\ifmmode\sp\else\^{}\fi}\catcode`\%=\active\def%{\%}Effective Occupied Volume $\widehat U_2(t)$}}%
\end{pgfscope}%
\begin{pgfscope}%
\pgfpathrectangle{\pgfqpoint{4.118067in}{0.795389in}}{\pgfqpoint{2.769565in}{1.984000in}}%
\pgfusepath{clip}%
\pgfsetrectcap%
\pgfsetroundjoin%
\pgfsetlinewidth{1.806750pt}%
\definecolor{currentstroke}{rgb}{0.121569,0.466667,0.705882}%
\pgfsetstrokecolor{currentstroke}%
\pgfsetdash{}{0pt}%
\pgfpathmoveto{\pgfqpoint{4.243956in}{1.320026in}}%
\pgfpathlineto{\pgfqpoint{4.277980in}{1.346010in}}%
\pgfpathlineto{\pgfqpoint{4.312004in}{1.372181in}}%
\pgfpathlineto{\pgfqpoint{4.346029in}{1.398484in}}%
\pgfpathlineto{\pgfqpoint{4.380053in}{1.424864in}}%
\pgfpathlineto{\pgfqpoint{4.414077in}{1.451273in}}%
\pgfpathlineto{\pgfqpoint{4.448101in}{1.477668in}}%
\pgfpathlineto{\pgfqpoint{4.482125in}{1.504015in}}%
\pgfpathlineto{\pgfqpoint{4.516149in}{1.530288in}}%
\pgfpathlineto{\pgfqpoint{4.550173in}{1.556473in}}%
\pgfpathlineto{\pgfqpoint{4.584198in}{1.582566in}}%
\pgfpathlineto{\pgfqpoint{4.618222in}{1.608574in}}%
\pgfpathlineto{\pgfqpoint{4.652246in}{1.634519in}}%
\pgfpathlineto{\pgfqpoint{4.686270in}{1.660429in}}%
\pgfpathlineto{\pgfqpoint{4.720294in}{1.686343in}}%
\pgfpathlineto{\pgfqpoint{4.754318in}{1.712307in}}%
\pgfpathlineto{\pgfqpoint{4.788342in}{1.738372in}}%
\pgfpathlineto{\pgfqpoint{4.822367in}{1.764587in}}%
\pgfpathlineto{\pgfqpoint{4.856391in}{1.791001in}}%
\pgfpathlineto{\pgfqpoint{4.890415in}{1.817657in}}%
\pgfpathlineto{\pgfqpoint{4.924439in}{1.844591in}}%
\pgfpathlineto{\pgfqpoint{4.958463in}{1.871827in}}%
\pgfpathlineto{\pgfqpoint{4.992487in}{1.899374in}}%
\pgfpathlineto{\pgfqpoint{5.026511in}{1.927231in}}%
\pgfpathlineto{\pgfqpoint{5.060536in}{1.955379in}}%
\pgfpathlineto{\pgfqpoint{5.094560in}{1.983784in}}%
\pgfpathlineto{\pgfqpoint{5.128584in}{2.012399in}}%
\pgfpathlineto{\pgfqpoint{5.162608in}{2.041159in}}%
\pgfpathlineto{\pgfqpoint{5.196632in}{2.069988in}}%
\pgfpathlineto{\pgfqpoint{5.230656in}{2.098794in}}%
\pgfpathlineto{\pgfqpoint{5.264680in}{2.127476in}}%
\pgfpathlineto{\pgfqpoint{5.298705in}{2.155924in}}%
\pgfpathlineto{\pgfqpoint{5.332729in}{2.184017in}}%
\pgfpathlineto{\pgfqpoint{5.366753in}{2.211632in}}%
\pgfpathlineto{\pgfqpoint{5.400777in}{2.238646in}}%
\pgfpathlineto{\pgfqpoint{5.434801in}{2.264936in}}%
\pgfpathlineto{\pgfqpoint{5.468825in}{2.290389in}}%
\pgfpathlineto{\pgfqpoint{5.502849in}{2.314902in}}%
\pgfpathlineto{\pgfqpoint{5.536874in}{2.338388in}}%
\pgfpathlineto{\pgfqpoint{5.570898in}{2.360782in}}%
\pgfpathlineto{\pgfqpoint{5.604922in}{2.382041in}}%
\pgfpathlineto{\pgfqpoint{5.638946in}{2.402145in}}%
\pgfpathlineto{\pgfqpoint{5.672970in}{2.421099in}}%
\pgfpathlineto{\pgfqpoint{5.706994in}{2.438932in}}%
\pgfpathlineto{\pgfqpoint{5.741018in}{2.455688in}}%
\pgfpathlineto{\pgfqpoint{5.775043in}{2.471431in}}%
\pgfpathlineto{\pgfqpoint{5.809067in}{2.486231in}}%
\pgfpathlineto{\pgfqpoint{5.843091in}{2.500167in}}%
\pgfpathlineto{\pgfqpoint{5.877115in}{2.513315in}}%
\pgfpathlineto{\pgfqpoint{5.911139in}{2.525750in}}%
\pgfpathlineto{\pgfqpoint{5.945163in}{2.537539in}}%
\pgfpathlineto{\pgfqpoint{5.979187in}{2.548739in}}%
\pgfpathlineto{\pgfqpoint{6.013212in}{2.559399in}}%
\pgfpathlineto{\pgfqpoint{6.047236in}{2.569557in}}%
\pgfpathlineto{\pgfqpoint{6.081260in}{2.579241in}}%
\pgfpathlineto{\pgfqpoint{6.115284in}{2.588474in}}%
\pgfpathlineto{\pgfqpoint{6.149308in}{2.597270in}}%
\pgfpathlineto{\pgfqpoint{6.183332in}{2.605641in}}%
\pgfpathlineto{\pgfqpoint{6.217356in}{2.613598in}}%
\pgfpathlineto{\pgfqpoint{6.251381in}{2.621147in}}%
\pgfpathlineto{\pgfqpoint{6.285405in}{2.628297in}}%
\pgfpathlineto{\pgfqpoint{6.319429in}{2.635054in}}%
\pgfpathlineto{\pgfqpoint{6.353453in}{2.641426in}}%
\pgfpathlineto{\pgfqpoint{6.387477in}{2.647419in}}%
\pgfpathlineto{\pgfqpoint{6.421501in}{2.653038in}}%
\pgfpathlineto{\pgfqpoint{6.455525in}{2.658289in}}%
\pgfpathlineto{\pgfqpoint{6.489550in}{2.663174in}}%
\pgfpathlineto{\pgfqpoint{6.523574in}{2.667695in}}%
\pgfpathlineto{\pgfqpoint{6.557598in}{2.671854in}}%
\pgfpathlineto{\pgfqpoint{6.591622in}{2.675649in}}%
\pgfpathlineto{\pgfqpoint{6.625646in}{2.679081in}}%
\pgfpathlineto{\pgfqpoint{6.659670in}{2.682150in}}%
\pgfpathlineto{\pgfqpoint{6.693694in}{2.684856in}}%
\pgfpathlineto{\pgfqpoint{6.727719in}{2.687205in}}%
\pgfpathlineto{\pgfqpoint{6.761743in}{2.689207in}}%
\pgfusepath{stroke}%
\end{pgfscope}%
\begin{pgfscope}%
\pgfpathrectangle{\pgfqpoint{4.118067in}{0.795389in}}{\pgfqpoint{2.769565in}{1.984000in}}%
\pgfusepath{clip}%
\pgfsetrectcap%
\pgfsetroundjoin%
\pgfsetlinewidth{1.806750pt}%
\definecolor{currentstroke}{rgb}{1.000000,0.498039,0.054902}%
\pgfsetstrokecolor{currentstroke}%
\pgfsetdash{}{0pt}%
\pgfpathmoveto{\pgfqpoint{4.243956in}{0.981174in}}%
\pgfpathlineto{\pgfqpoint{4.277980in}{0.992212in}}%
\pgfpathlineto{\pgfqpoint{4.312004in}{1.003343in}}%
\pgfpathlineto{\pgfqpoint{4.346029in}{1.014541in}}%
\pgfpathlineto{\pgfqpoint{4.380053in}{1.025777in}}%
\pgfpathlineto{\pgfqpoint{4.414077in}{1.037032in}}%
\pgfpathlineto{\pgfqpoint{4.448101in}{1.048286in}}%
\pgfpathlineto{\pgfqpoint{4.482125in}{1.059529in}}%
\pgfpathlineto{\pgfqpoint{4.516149in}{1.070753in}}%
\pgfpathlineto{\pgfqpoint{4.550173in}{1.081955in}}%
\pgfpathlineto{\pgfqpoint{4.584198in}{1.093140in}}%
\pgfpathlineto{\pgfqpoint{4.618222in}{1.104315in}}%
\pgfpathlineto{\pgfqpoint{4.652246in}{1.115491in}}%
\pgfpathlineto{\pgfqpoint{4.686270in}{1.126685in}}%
\pgfpathlineto{\pgfqpoint{4.720294in}{1.137916in}}%
\pgfpathlineto{\pgfqpoint{4.754318in}{1.149207in}}%
\pgfpathlineto{\pgfqpoint{4.788342in}{1.160585in}}%
\pgfpathlineto{\pgfqpoint{4.822367in}{1.172079in}}%
\pgfpathlineto{\pgfqpoint{4.856391in}{1.183718in}}%
\pgfpathlineto{\pgfqpoint{4.890415in}{1.195536in}}%
\pgfpathlineto{\pgfqpoint{4.924439in}{1.207562in}}%
\pgfpathlineto{\pgfqpoint{4.958463in}{1.219828in}}%
\pgfpathlineto{\pgfqpoint{4.992487in}{1.232360in}}%
\pgfpathlineto{\pgfqpoint{5.026511in}{1.245181in}}%
\pgfpathlineto{\pgfqpoint{5.060536in}{1.258308in}}%
\pgfpathlineto{\pgfqpoint{5.094560in}{1.271751in}}%
\pgfpathlineto{\pgfqpoint{5.128584in}{1.285513in}}%
\pgfpathlineto{\pgfqpoint{5.162608in}{1.299584in}}%
\pgfpathlineto{\pgfqpoint{5.196632in}{1.313947in}}%
\pgfpathlineto{\pgfqpoint{5.230656in}{1.328576in}}%
\pgfpathlineto{\pgfqpoint{5.264680in}{1.343434in}}%
\pgfpathlineto{\pgfqpoint{5.298705in}{1.358480in}}%
\pgfpathlineto{\pgfqpoint{5.332729in}{1.373669in}}%
\pgfpathlineto{\pgfqpoint{5.366753in}{1.388955in}}%
\pgfpathlineto{\pgfqpoint{5.400777in}{1.404295in}}%
\pgfpathlineto{\pgfqpoint{5.434801in}{1.419654in}}%
\pgfpathlineto{\pgfqpoint{5.468825in}{1.435007in}}%
\pgfpathlineto{\pgfqpoint{5.502849in}{1.450344in}}%
\pgfpathlineto{\pgfqpoint{5.536874in}{1.465672in}}%
\pgfpathlineto{\pgfqpoint{5.570898in}{1.481014in}}%
\pgfpathlineto{\pgfqpoint{5.604922in}{1.496415in}}%
\pgfpathlineto{\pgfqpoint{5.638946in}{1.511936in}}%
\pgfpathlineto{\pgfqpoint{5.672970in}{1.527658in}}%
\pgfpathlineto{\pgfqpoint{5.706994in}{1.543677in}}%
\pgfpathlineto{\pgfqpoint{5.741018in}{1.560101in}}%
\pgfpathlineto{\pgfqpoint{5.775043in}{1.577050in}}%
\pgfpathlineto{\pgfqpoint{5.809067in}{1.594653in}}%
\pgfpathlineto{\pgfqpoint{5.843091in}{1.613043in}}%
\pgfpathlineto{\pgfqpoint{5.877115in}{1.632354in}}%
\pgfpathlineto{\pgfqpoint{5.911139in}{1.652722in}}%
\pgfpathlineto{\pgfqpoint{5.945163in}{1.674280in}}%
\pgfpathlineto{\pgfqpoint{5.979187in}{1.697157in}}%
\pgfpathlineto{\pgfqpoint{6.013212in}{1.721475in}}%
\pgfpathlineto{\pgfqpoint{6.047236in}{1.747348in}}%
\pgfpathlineto{\pgfqpoint{6.081260in}{1.774881in}}%
\pgfpathlineto{\pgfqpoint{6.115284in}{1.804165in}}%
\pgfpathlineto{\pgfqpoint{6.149308in}{1.835277in}}%
\pgfpathlineto{\pgfqpoint{6.183332in}{1.868277in}}%
\pgfpathlineto{\pgfqpoint{6.217356in}{1.903202in}}%
\pgfpathlineto{\pgfqpoint{6.251381in}{1.940063in}}%
\pgfpathlineto{\pgfqpoint{6.285405in}{1.978842in}}%
\pgfpathlineto{\pgfqpoint{6.319429in}{2.019481in}}%
\pgfpathlineto{\pgfqpoint{6.353453in}{2.061877in}}%
\pgfpathlineto{\pgfqpoint{6.387477in}{2.105874in}}%
\pgfpathlineto{\pgfqpoint{6.421501in}{2.151250in}}%
\pgfpathlineto{\pgfqpoint{6.455525in}{2.197713in}}%
\pgfpathlineto{\pgfqpoint{6.489550in}{2.244888in}}%
\pgfpathlineto{\pgfqpoint{6.523574in}{2.292315in}}%
\pgfpathlineto{\pgfqpoint{6.557598in}{2.339446in}}%
\pgfpathlineto{\pgfqpoint{6.591622in}{2.385660in}}%
\pgfpathlineto{\pgfqpoint{6.625646in}{2.430277in}}%
\pgfpathlineto{\pgfqpoint{6.659670in}{2.472594in}}%
\pgfpathlineto{\pgfqpoint{6.693694in}{2.511930in}}%
\pgfpathlineto{\pgfqpoint{6.727719in}{2.547680in}}%
\pgfpathlineto{\pgfqpoint{6.761743in}{2.579369in}}%
\pgfusepath{stroke}%
\end{pgfscope}%
\begin{pgfscope}%
\pgfpathrectangle{\pgfqpoint{4.118067in}{0.795389in}}{\pgfqpoint{2.769565in}{1.984000in}}%
\pgfusepath{clip}%
\pgfsetrectcap%
\pgfsetroundjoin%
\pgfsetlinewidth{1.806750pt}%
\definecolor{currentstroke}{rgb}{0.172549,0.627451,0.172549}%
\pgfsetstrokecolor{currentstroke}%
\pgfsetdash{}{0pt}%
\pgfpathmoveto{\pgfqpoint{4.243956in}{0.885571in}}%
\pgfpathlineto{\pgfqpoint{4.277980in}{0.892754in}}%
\pgfpathlineto{\pgfqpoint{4.312004in}{0.900042in}}%
\pgfpathlineto{\pgfqpoint{4.346029in}{0.907414in}}%
\pgfpathlineto{\pgfqpoint{4.380053in}{0.914853in}}%
\pgfpathlineto{\pgfqpoint{4.414077in}{0.922344in}}%
\pgfpathlineto{\pgfqpoint{4.448101in}{0.929873in}}%
\pgfpathlineto{\pgfqpoint{4.482125in}{0.937431in}}%
\pgfpathlineto{\pgfqpoint{4.516149in}{0.945008in}}%
\pgfpathlineto{\pgfqpoint{4.550173in}{0.952598in}}%
\pgfpathlineto{\pgfqpoint{4.584198in}{0.960195in}}%
\pgfpathlineto{\pgfqpoint{4.618222in}{0.967797in}}%
\pgfpathlineto{\pgfqpoint{4.652246in}{0.975403in}}%
\pgfpathlineto{\pgfqpoint{4.686270in}{0.983019in}}%
\pgfpathlineto{\pgfqpoint{4.720294in}{0.990651in}}%
\pgfpathlineto{\pgfqpoint{4.754318in}{0.998315in}}%
\pgfpathlineto{\pgfqpoint{4.788342in}{1.006027in}}%
\pgfpathlineto{\pgfqpoint{4.822367in}{1.013813in}}%
\pgfpathlineto{\pgfqpoint{4.856391in}{1.021703in}}%
\pgfpathlineto{\pgfqpoint{4.890415in}{1.029730in}}%
\pgfpathlineto{\pgfqpoint{4.924439in}{1.037935in}}%
\pgfpathlineto{\pgfqpoint{4.958463in}{1.046360in}}%
\pgfpathlineto{\pgfqpoint{4.992487in}{1.055051in}}%
\pgfpathlineto{\pgfqpoint{5.026511in}{1.064054in}}%
\pgfpathlineto{\pgfqpoint{5.060536in}{1.073419in}}%
\pgfpathlineto{\pgfqpoint{5.094560in}{1.083197in}}%
\pgfpathlineto{\pgfqpoint{5.128584in}{1.093441in}}%
\pgfpathlineto{\pgfqpoint{5.162608in}{1.104204in}}%
\pgfpathlineto{\pgfqpoint{5.196632in}{1.115543in}}%
\pgfpathlineto{\pgfqpoint{5.230656in}{1.127516in}}%
\pgfpathlineto{\pgfqpoint{5.264680in}{1.140185in}}%
\pgfpathlineto{\pgfqpoint{5.298705in}{1.153614in}}%
\pgfpathlineto{\pgfqpoint{5.332729in}{1.167872in}}%
\pgfpathlineto{\pgfqpoint{5.366753in}{1.183029in}}%
\pgfpathlineto{\pgfqpoint{5.400777in}{1.199161in}}%
\pgfpathlineto{\pgfqpoint{5.434801in}{1.216348in}}%
\pgfpathlineto{\pgfqpoint{5.468825in}{1.234670in}}%
\pgfpathlineto{\pgfqpoint{5.502849in}{1.254213in}}%
\pgfpathlineto{\pgfqpoint{5.536874in}{1.275061in}}%
\pgfpathlineto{\pgfqpoint{5.570898in}{1.297301in}}%
\pgfpathlineto{\pgfqpoint{5.604922in}{1.321017in}}%
\pgfpathlineto{\pgfqpoint{5.638946in}{1.346286in}}%
\pgfpathlineto{\pgfqpoint{5.672970in}{1.373183in}}%
\pgfpathlineto{\pgfqpoint{5.706994in}{1.401767in}}%
\pgfpathlineto{\pgfqpoint{5.741018in}{1.432086in}}%
\pgfpathlineto{\pgfqpoint{5.775043in}{1.464167in}}%
\pgfpathlineto{\pgfqpoint{5.809067in}{1.498015in}}%
\pgfpathlineto{\pgfqpoint{5.843091in}{1.533606in}}%
\pgfpathlineto{\pgfqpoint{5.877115in}{1.570883in}}%
\pgfpathlineto{\pgfqpoint{5.911139in}{1.609755in}}%
\pgfpathlineto{\pgfqpoint{5.945163in}{1.650098in}}%
\pgfpathlineto{\pgfqpoint{5.979187in}{1.691750in}}%
\pgfpathlineto{\pgfqpoint{6.013212in}{1.734525in}}%
\pgfpathlineto{\pgfqpoint{6.047236in}{1.778215in}}%
\pgfpathlineto{\pgfqpoint{6.081260in}{1.822602in}}%
\pgfpathlineto{\pgfqpoint{6.115284in}{1.867471in}}%
\pgfpathlineto{\pgfqpoint{6.149308in}{1.912622in}}%
\pgfpathlineto{\pgfqpoint{6.183332in}{1.957878in}}%
\pgfpathlineto{\pgfqpoint{6.217356in}{2.003092in}}%
\pgfpathlineto{\pgfqpoint{6.251381in}{2.048151in}}%
\pgfpathlineto{\pgfqpoint{6.285405in}{2.092965in}}%
\pgfpathlineto{\pgfqpoint{6.319429in}{2.137460in}}%
\pgfpathlineto{\pgfqpoint{6.353453in}{2.181568in}}%
\pgfpathlineto{\pgfqpoint{6.387477in}{2.225206in}}%
\pgfpathlineto{\pgfqpoint{6.421501in}{2.268261in}}%
\pgfpathlineto{\pgfqpoint{6.455525in}{2.310586in}}%
\pgfpathlineto{\pgfqpoint{6.489550in}{2.351981in}}%
\pgfpathlineto{\pgfqpoint{6.523574in}{2.392199in}}%
\pgfpathlineto{\pgfqpoint{6.557598in}{2.430946in}}%
\pgfpathlineto{\pgfqpoint{6.591622in}{2.467889in}}%
\pgfpathlineto{\pgfqpoint{6.625646in}{2.502673in}}%
\pgfpathlineto{\pgfqpoint{6.659670in}{2.534944in}}%
\pgfpathlineto{\pgfqpoint{6.693694in}{2.564369in}}%
\pgfpathlineto{\pgfqpoint{6.727719in}{2.590673in}}%
\pgfpathlineto{\pgfqpoint{6.761743in}{2.613666in}}%
\pgfusepath{stroke}%
\end{pgfscope}%
\begin{pgfscope}%
\pgfsetrectcap%
\pgfsetmiterjoin%
\pgfsetlinewidth{0.803000pt}%
\definecolor{currentstroke}{rgb}{0.000000,0.000000,0.000000}%
\pgfsetstrokecolor{currentstroke}%
\pgfsetdash{}{0pt}%
\pgfpathmoveto{\pgfqpoint{4.118067in}{0.795389in}}%
\pgfpathlineto{\pgfqpoint{4.118067in}{2.779389in}}%
\pgfusepath{stroke}%
\end{pgfscope}%
\begin{pgfscope}%
\pgfsetrectcap%
\pgfsetmiterjoin%
\pgfsetlinewidth{0.803000pt}%
\definecolor{currentstroke}{rgb}{0.000000,0.000000,0.000000}%
\pgfsetstrokecolor{currentstroke}%
\pgfsetdash{}{0pt}%
\pgfpathmoveto{\pgfqpoint{6.887632in}{0.795389in}}%
\pgfpathlineto{\pgfqpoint{6.887632in}{2.779389in}}%
\pgfusepath{stroke}%
\end{pgfscope}%
\begin{pgfscope}%
\pgfsetrectcap%
\pgfsetmiterjoin%
\pgfsetlinewidth{0.803000pt}%
\definecolor{currentstroke}{rgb}{0.000000,0.000000,0.000000}%
\pgfsetstrokecolor{currentstroke}%
\pgfsetdash{}{0pt}%
\pgfpathmoveto{\pgfqpoint{4.118067in}{0.795389in}}%
\pgfpathlineto{\pgfqpoint{6.887632in}{0.795389in}}%
\pgfusepath{stroke}%
\end{pgfscope}%
\begin{pgfscope}%
\pgfsetrectcap%
\pgfsetmiterjoin%
\pgfsetlinewidth{0.803000pt}%
\definecolor{currentstroke}{rgb}{0.000000,0.000000,0.000000}%
\pgfsetstrokecolor{currentstroke}%
\pgfsetdash{}{0pt}%
\pgfpathmoveto{\pgfqpoint{4.118067in}{2.779389in}}%
\pgfpathlineto{\pgfqpoint{6.887632in}{2.779389in}}%
\pgfusepath{stroke}%
\end{pgfscope}%
\begin{pgfscope}%
\definecolor{textcolor}{rgb}{0.000000,0.000000,0.000000}%
\pgfsetstrokecolor{textcolor}%
\pgfsetfillcolor{textcolor}%
\pgftext[x=5.502849in,y=2.862722in,,base]{\color{textcolor}{\rmfamily\fontsize{8.000000}{9.600000}\selectfont\catcode`\^=\active\def^{\ifmmode\sp\else\^{}\fi}\catcode`\%=\active\def%{\%}(b) Normalized-Embedding Volume}}%
\end{pgfscope}%
\begin{pgfscope}%
\pgfsetrectcap%
\pgfsetroundjoin%
\pgfsetlinewidth{1.806750pt}%
\definecolor{currentstroke}{rgb}{0.121569,0.466667,0.705882}%
\pgfsetstrokecolor{currentstroke}%
\pgfsetdash{}{0pt}%
\pgfpathmoveto{\pgfqpoint{2.155988in}{0.197222in}}%
\pgfpathlineto{\pgfqpoint{2.233766in}{0.197222in}}%
\pgfpathlineto{\pgfqpoint{2.311544in}{0.197222in}}%
\pgfusepath{stroke}%
\end{pgfscope}%
\begin{pgfscope}%
\definecolor{textcolor}{rgb}{0.000000,0.000000,0.000000}%
\pgfsetstrokecolor{textcolor}%
\pgfsetfillcolor{textcolor}%
\pgftext[x=2.389322in,y=0.163194in,left,base]{\color{textcolor}{\rmfamily\fontsize{7.000000}{8.400000}\selectfont\catcode`\^=\active\def^{\ifmmode\sp\else\^{}\fi}\catcode`\%=\active\def%{\%}Digits all}}%
\end{pgfscope}%
\begin{pgfscope}%
\pgfsetrectcap%
\pgfsetroundjoin%
\pgfsetlinewidth{1.806750pt}%
\definecolor{currentstroke}{rgb}{1.000000,0.498039,0.054902}%
\pgfsetstrokecolor{currentstroke}%
\pgfsetdash{}{0pt}%
\pgfpathmoveto{\pgfqpoint{2.941214in}{0.197222in}}%
\pgfpathlineto{\pgfqpoint{3.018992in}{0.197222in}}%
\pgfpathlineto{\pgfqpoint{3.096770in}{0.197222in}}%
\pgfusepath{stroke}%
\end{pgfscope}%
\begin{pgfscope}%
\definecolor{textcolor}{rgb}{0.000000,0.000000,0.000000}%
\pgfsetstrokecolor{textcolor}%
\pgfsetfillcolor{textcolor}%
\pgftext[x=3.174548in,y=0.163194in,left,base]{\color{textcolor}{\rmfamily\fontsize{7.000000}{8.400000}\selectfont\catcode`\^=\active\def^{\ifmmode\sp\else\^{}\fi}\catcode`\%=\active\def%{\%}Digits weighted}}%
\end{pgfscope}%
\begin{pgfscope}%
\pgfsetrectcap%
\pgfsetroundjoin%
\pgfsetlinewidth{1.806750pt}%
\definecolor{currentstroke}{rgb}{0.172549,0.627451,0.172549}%
\pgfsetstrokecolor{currentstroke}%
\pgfsetdash{}{0pt}%
\pgfpathmoveto{\pgfqpoint{4.033155in}{0.197222in}}%
\pgfpathlineto{\pgfqpoint{4.110932in}{0.197222in}}%
\pgfpathlineto{\pgfqpoint{4.188710in}{0.197222in}}%
\pgfusepath{stroke}%
\end{pgfscope}%
\begin{pgfscope}%
\definecolor{textcolor}{rgb}{0.000000,0.000000,0.000000}%
\pgfsetstrokecolor{textcolor}%
\pgfsetfillcolor{textcolor}%
\pgftext[x=4.266488in,y=0.163194in,left,base]{\color{textcolor}{\rmfamily\fontsize{7.000000}{8.400000}\selectfont\catcode`\^=\active\def^{\ifmmode\sp\else\^{}\fi}\catcode`\%=\active\def%{\%}Digits 0/1}}%
\end{pgfscope}%
\end{pgfpicture}%
\makeatother%
\endgroup%

\caption{Full profiles for the normalized embeddings in Table~\ref{tab:digits}. The left panel shows gESS, and the right shows effective occupied volume.
The weighted embedding is lower on both measures across many angular scales.}
\label{fig:app_digits_profiles}
\end{figure*}
\end{document}
